\xchapter{Sources}

THIRD PERIOD

FROM CONSTANTINE THE GREAT TO GREGORY THE GREAT.

a. d. 311–590.

SOURCES.

I. Christian Sources: (a) The Acts Of Councils; in the Collectiones conciliorum of Hardouin, Par. 1715 sqq. 12 vols. fol.; Mansi, Flor. et Ven. 1759 sqq. 31 vols. fol.; Fuchs: Bibliothek der Kirchenversammlungen des 4ten und 5ten Jahrh. Leipz. 1780 sqq.; and Bruns: Biblioth. eccl. vol. i. Canones Apost. et Conc. saec. iv.–vii. Berol. 1839.

(b) The Imperial Laws and Decrees referring to the church, in the Codex Theodosianus, collected a.d. 438, the Codex Justinianeus, collected in 529, and the Cod. repetitae praelectionis of 534.

(c) The Official Letters of popes (in the Bullarium Romanum), patriarchs, and bishops.

(d) The writings of all the Church Fathers from the beginning of the 4th century to the end of the 6th. Especially of Eusebius, Athanasius, Basil, the two Gregories, the two Cyrils, Chrysostom, and Theodoret, of the Greek church; and Ambrose, Augustine, Jerome, and Leo the Great, of the Latin. Comp. the Benedictine Editions of the several Fathers; the Maxima Bibliotheca veterum Patrum, Lugd. 1677 sqq. (in all 27 vols. fol.), vols. iii.–xi.; Gallandi: Biblioth. vet. Patrum, etc. Ven. 1765 sqq. (14 vols. fol.), vols. iv.–xii.

(e) Contemporary Church Historians, (1) of the Greek church: Eusebius of Caesarea († about 340): the ninth and tenth books of his H. E. down to 324, and his biography of Constantine the Great, see § 2 infra; Socrates Scholasticus of Constantinople: Histor. ecclesiast. libri vii, a.d. 306–439; Hermias Sozomen of Constantinople: H. eccl. l. ix, a.d. 323–423; Theodoret, bishop of Cyros in Mesopotamia: H. eccl. l. v, a.d. 325–429; the Arian Philostorgius: H. eccl. l. xii, a.d. 318–425, extant only in extracts in Photius Cod. 40; Theodorus Lector, of Constantinople, epitomizer of Socrates, Sozomen, and Theodoret, continuing the latter down to 518, preserved in fragments by Nicephorus Callistus; Evagrius of Antioch: H. eccl. l. vi, a.d. 431–594; Nicephorus Callistus (or Niceph. Callisti), about 1330, author of a church history in 23 books, to a.d. 911 (ed. Fronto Ducaeus, Par. 1630). The historical works of these Greek writers, excepting the last, are also published together under the title: Historiae ecclesiasticae Scriptores, etc., Graec. et Lat., with notes by H. Valesius (and G. Reading), Par. 1659–1673; and Cantabr. 1720, 3 vols. fol. (2) Of the Latin church historians few are important: Rufinus, presb. of Aquileia (†410), translated Eusebius and continued him in two more books to 395; Sulpicius Severus, presb. in Gaul: Hist. Sacra, l. ii, from the creation to a.d. 400; Paulus Orosius, presbyter in Spain: Historiarum libri vii. written about 416, extending from the creation to his own time; Cassiodorus, about 550: Hist. tripartite, l. xii. a mere extract from the works of the Greek church historians, but, with the work of Rufinus, the chief source of historical knowledge through the whole middle age; and Jerome († 419): De viris illustrious, or Catalogus scriptorum eccles., written about 392, continued under the same title by Gennadius, about 495, and by Isidor of Seville, about 630.

(f) For chronology, the Greek Πασχάλιον, or Chronicon Paschale (wrongly called Alexandrinum), primarily a table of the passovers from the beginning of the world to a.d. 354 under Constantius, with later additions down to 628. (Ed. Car. du Fresne Dom. du Cange. Par. 1688, and L. Dindorf, Bonn. 1832, 2 vols.) The Chronicle of Eusebius and Jerome (Χρονικὰ συγγράμματα, παντοδαπὴ ἱστορία), containing an outline of universal history down to 325, mainly after the chronography of Julius Africanus, and an extract from the universal chronicle in tabular form down to 379, long extant only in the free Latin translation and continuation of Jerome (ed. Jos. Scaliger. Lugd. Batav. 1606 and later), since 1792 known also in an Armenian translation (ed. J. Bapt. Aucher. Ven. 1818, and Aug. Mai, Script. vet. nov. coll. 1833. Tom. viii). In continuation of the Latin chronicle of Jerome, the chronicle of Prosper of Aquitania down to 455; that of the spanish bishop Idatius, to 469; and that of Marcellinus Comes, to 534. Comp. Chronica medii aevi post Euseb. atque Hieron., etc. ed. Roesler, Tüb. 1798.

II. Heathen Sources: Ammianus Marcellinus (officer under Julian, honest and impartial): Rerum gestarum libri xiv-xxxi, a.d. 353–378 (the first 13 books are lost), ed. Jac. Gronov. Lugd. Batav. 1693 fol., and J. A. Ernesti, Lips. 1778 and 1835. Eunapius (philosopher and historian; bitter against the Christian emperors): Χρονικὴ ἱστορία, a.d. 268–405, extant only in fragments, ed. Bekker and Niebuhr, Bonn. 1829. Zosimus (court officer under Theodosius II., likewise biassed): Ἱστορία νέα, l. vi, a.d. 284–410, ed. Cellarius 1679, Reitemeier 1784, and Imm. Bekker, Bonn. 1837. Also the writings of Julian the Apostate (against Christianity), Libanius and Symmachus (philosophically tolerant), \&c. Comp. the literature at § 2 and 4.

\xchapter{Later Literature}

LATER LITERATURE.

Besides the contemporary histories named above under 1 (e) among the sources, we should mention particularly Baronius (R.C. of the a.d.Ultramontane school, † 1607): Annales Eccles. vol. iii.–viii. (a heavy and unreadable chronicle, but valuable for reference to original documents). Tillemont (R.C. leaning to Jansenism, † 1698): Mémoires, etc., vol. vi.–xvi. (mostly biographical, minute, and conscientious). Gibbon († 1794): Decline and Fall of the Roman Empire, from ch. xvii. onward (unsurpassed in the skilful use of sources and artistic composition, but skeptical and destitute of sympathy with the genius of Christianity). Schröckh (moderate Lutheran, † 1808): Christl. Kirchengesch. Theil v.–xviii. (A simple and diffuse, but thorough and trustworthy narrative). Neander (Evangel. † 1850): Allg. Gesch. der Chr. Rel. und Kirche. Hamb. vol. iv.–vi., 2d ed. 1846 sqq. Engl. transl. by Torrey, vol. ii. (Profound and genial in the genetic development of Christian doctrine and life, but defective in the political and aesthetic sections, and prolix and careless in style and arrangement). Gieseler (Protest. † 1854): Kirchengesch. Bonn. i. 2. 2d ed. 1845. Engl. transl. by Davidson, and revised by H. B. Smith, N. York, vol. i. and ii. (Critical and reliable in the notes, but meagre, dry, and cold in the text).

Isaac Taylor (Independent): Ancient Christianity, and the Doctrines of the Oxf. Tracts for the Times. Lond. 4th ed. 1844. 2 vols. (Anti-Puseyite). Böhringer (G. Ref.): Kirchengeschichte in Biographieen, vol. i. parts 3 and 4. Zür. 1845 sq. (from Ambrose to Gregory the Great). Carwithen And Lyall: History of the Christian Church from the 4th to the 12th Cent. in the Encycl. Metrop. 1849; published separately in Lond. and Glasg. 1856. J. C. Robertson (Angl.): Hist. of the Christ. Church to the Pontificate of Gregory the Great. Lond. 1854 (pp. 166–516). H. H. Milman (Angl.): History of Christianity from the Birth of Christ to the abolition of Paganism in the Roman Empire. Lond. 1840 (New York, 1844), Book III. and IV. Milman: Hist. of Latin Christianity; including that of the Popes to the Pontificate of Nicholas V. Lond. 1854 sqq. 6 vols., republished in New York, 1860, in 8 vols. (vol. i. a resumé of the first six centuries to Gregory I., the remaining vols. devoted to the middle ages). K. R. Hagenbach (G. Ref.):Die Christl. Kirche vom 4ten his 6ten Jahrh. Leipz. 1855 (2d vol. of his popular “Vorlesungen über die ältere Kirchengesch.”). Albert de Broglie (R.C.): L’église et l’empire romain au IVme siècle. Par. 1855–’66. 6 vols. Ferd. Christ. Baur: Die Christl. Kirche vom Anfang des vierten bis zum Ende des sechsten Jahrhunderts in den Hauptmomenten ihrer Entwicklung. Tüb. 1859 (critical and philosophical). Wm. Bright: A History of the Church from the Edict of Milan, a.d. 313, to the Council of Chalcedon, a.d. 451. Oxf. and Lond. 1860. Arthur P. Stanley: Lectures on the History of the Eastern Church. Lond. 1861 (pp. 512), republished in New York from the 2d Lond. ed. 1862 (a series of graphic pictures of prominent characters and events in the history of the Greek and Russian church, but no complete history).

\xchapter{Introduction and General View}

§ 1. Introduction and General View.

From the Christianity of the Apostles and Martyrs we proceed to the Christianity of the Patriarchs and Emperors.

The third period of the history of the Church, which forms the subject of this volume, extends from the emperor Constantine to the pope Gregory I.; from the beginning of the fourth century to the close of the sixth. During this period Christianity still moves, as in the first three centuries, upon the geographical scene of the Graeco-Roman empire and the ancient classical culture, the countries around the Mediterranean Sea. But its field and its operation are materially enlarged, and even touch the barbarians on the limit of the empire. Above all, its relation to the temporal power, and its social and political position and import, undergo an entire and permanent change. We have here to do with the church of the Graeco-Roman empire, and with the beginning of Christianity among the Germanic barbarians. Let us glance first at the general character and leading events of this important period.

The reign of Constantine the Great marks the transition of the Christian religion from under persecution by the secular government to union with the same; the beginning of the state-church system. The Graeco-Roman heathenism, the most cultivated and powerful form of idolatry, which history knows, surrenders, after three hundred years’ struggle, to Christianity, and dies of incurable consumption, with the confession: Galilean, thou hast conquered! The ruler of the civilized world lays his crown at the feet of the crucified Jesus of Nazareth. The successor of Nero, Domitian, and Diocletian appears in the imperial purple at the council of Nice as protector of the church, and takes his golden throne at the nod of bishops, who still bear the scars of persecution. The despised sect, which, like its Founder in the days of His humiliation, had not where to lay its head, is raised to sovereign authority in the state, enters into the prerogatives of the pagan priesthood, grows rich and powerful, builds countless churches out of the stones of idol temples to the honor of Christ and his martyrs, employs the wisdom of Greece and Rome to vindicate the foolishness of the cross, exerts a molding power upon civil legislation, rules the national life, and leads off the history of the world. But at the same time the church, embracing the mass of the population of the empire, from the Caesar to the meanest slave, and living amidst all its institutions, received into her bosom vast deposits of foreign material from the world and from heathenism, exposing herself to new dangers and imposing upon herself new and heavy labors.

The union of church and state extends its influence, now healthful, now baneful, into every department of our history.

The Christian life of the Nicene and post-Nicene age reveals a mass of worldliness within the church; an entire abatement of chiliasm with its longing after the return of Christ and his glorious reign, and in its stead an easy repose in the present order of things; with a sublime enthusiasm, on the other hand, for the renunciation of self and the world, particularly in the hermitage and the cloister, and with some of the noblest heroes of Christian holiness.

Monasticism, in pursuance of the ascetic tendencies of the previous period, and in opposition to the prevailing secularization of Christianity, sought to save the virgin purity of the church and the glory of martyrdom by retreat from the world into the wilderness; and it carried the ascetic principle to the summit of moral heroism, though not rarely to the borders of fanaticism and brutish stupefaction. It spread with incredible rapidity and irresistible fascination from Egypt over the whole church, east and west, and received the sanction of the greatest church teachers, of an Athanasius, a Basil, a Chrysostom, an Augustine, a Jerome, as the surest and shortest way to heaven.

It soon became a powerful rival of the priesthood, and formed a third order, between the priesthood and the laity. The more extraordinary and eccentric the religion of the anchorets and monks, the more they were venerated among the people. The whole conception of the Christian life from the fourth to the sixteenth century is pervaded with the ascetic and monastic spirit, and pays the highest admiration to the voluntary celibacy, poverty, absolute obedience, and excessive self-punishments of the pillar-saints and the martyrs of the desert; while in the same degree the modest virtues of every-day household and social life are looked upon as an inferior degree of morality.

In this point the old Catholic ethical ideas essentially differ from those of evangelical Protestantism and modern civilization. But, to understand and appreciate them, we must consider them in connection with the corrupt social condition of the rapidly decaying empire of Rome. The Christian spirit in that age, in just its most earnest and vigorous forms, felt compelled to assume in some measure an anti-social, seclusive character, and to prepare itself in the school of privation and solitude for the work of transforming the world and founding a new Christian order of society upon the ruins of the ancient heathenism.

In the development of doctrine the Nicene and post-Nicene age is second in productiveness and importance only to those of the apostles and of the reformation. It is the classical period for the objective fundamental dogmas, which constitute the ecumenical or old Catholic confession of faith. The Greek church produced the symbolical definition of the orthodox view of the holy Trinity and the person of Christ, while the Latin church made considerable advance with the anthropological and soteriological doctrines of sin and grace. The fourth and fifth centuries produced the greatest church fathers, Athanasius and Chrysostom in the East, Jerome and Augustine in the West. All learning and science now came into the service of the church, and all classes of society, from the emperor to the artisan, took the liveliest, even a passionate interest, in the theological controversies. Now, too, for the first time, could ecumenical councils be held, in which the church of the whole Roman empire was represented, and fixed its articles of faith in an authoritative way.

Now also, however, the lines of orthodoxy were more and more strictly drawn; freedom of inquiry was restricted; and all as departure from the state-church system was met not only, as formerly, with spiritual weapons, but also with civil punishments. So early as the fourth century the dominant party, the orthodox as well as the heterodox, with help of the imperial authority practised deposition, confiscation, and banishment upon its opponents. It was but one step thence to the penalties of torture and death, which were ordained in the middle age, and even so lately as the middle of the seventeenth century, by state-church authority, both Protestant and Roman Catholic, and continue in many countries to this day, against religious dissenters of every kind as enemies to the prevailing order of things. Absolute freedom of religion and of worship is in fact logically impossible on the state-church system. It requires the separation of the spiritual and temporal powers. Yet, from the very beginning of political persecution, loud voices rise against it and in behalf of ecclesiastico-religious toleration; though the plea always comes from the oppressed party, which, as soon as it gains the power, is generally found, in lamentable inconsistency, imitating the violence of its former oppressors. The protest springs rather from the sense of personal injury, than from horror of the principle of persecution, or from any clear apprehension of the nature of the gospel and its significant words: “Put up thy sword into the sheath;” “My kingdom is not of this world.”

The organization of the church adapts itself to the political and geographical divisions of the empire. The powers of the hierarchy are enlarged, the bishops become leading officers of the state and acquire a controlling influence in civil and political affairs, though more or less at the expense of their spiritual dignity and independence, especially at the Byzantine court. The episcopal system passes on into the metropolitan and patriarchal. In the fifth century the patriarchs of Rome, Constantinople, Antioch, Alexandria, and Jerusalem stand at the head of Christendom. Among these Rome and Constantinople are the most powerful rivals, and the Roman patriarch already puts forth a claim to universal spiritual supremacy, which subsequently culminates in the mediaeval papacy, though limited to the West and resisted by the constant protest of the Greek church and of all non-Catholic sects. In addition to provincial synods we have now also general synods, but called by the emperors and more or less affected, though not controlled, by political influence.

From the time of Constantine church discipline declines; the whole Roman world having become nominally Christian, and the host of hypocritical professors multiplying beyond all control. Yet the firmness of Ambrose with the emperor Theodosius shows, that noble instances of discipline are not altogether wanting.

Worship appears greatly enriched and adorned; for art now comes into the service of the church. A Christian architecture, a Christian sculpture, a Christian painting, music, and poetry arise, favoring at once devotion and solemnity, and all sorts of superstition and empty display. The introduction of religious images succeeds only after long and violent opposition. The element of priesthood and of mystery is developed, but in connection with a superstitious reliance upon a certain magical operation of outward rites. Church festivals are multiplied and celebrated with great pomp; and not exclusively in honor of Christ, but in connection with an extravagant veneration of martyrs and saints, which borders on idolatry, and often reminds us of the heathen hero-worship not yet uprooted from the general mind. The multiplication and accumulation of religious ceremonies impressed the senses and the imagination, but prejudiced simplicity, spirituality, and fervor in the worship of God. Hence also the beginnings of reaction against ceremonialism and formalism.

Notwithstanding the complete and sudden change of the social and political circumstances of the church, which meets us on the threshold of this period, we have still before us the natural, necessary continuation of the pre-Constantine church in its light and shade, and the gradual transition of the old Graeco-Roman Catholicism into the Germano-Roman Catholicism of the middle age.

Our attention will now for the first time be turned in earnest, not only to Christianity in the Roman empire, but also to Christianity among the Germanic barbarians, who from East and North threaten the empire and the entire civilization of classic antiquity. The church prolonged, indeed, the existence of the Roman empire, gave it a new splendor and elevation, new strength and unity, as well as comfort in misfortune; but could not prevent its final dissolution, first in the West (a.d. 476), afterwards (1453) in the East. But she herself survived the storms of the great migration, brought the pagan invaders under the influence of Christianity, taught the barbarians the arts of peace, planted a higher civilization upon the ruins of the ancient world, and thus gave new proof of the indestructible, all-subduing energy of her life.

In a minute history of the fourth, fifth, and sixth centuries we should mark the following subdivisions:

1. The Constantinian and Athanasian, or the Nicene and Trinitarian age, from 311 to the second general council in 381, distinguished by the conversion of Constantine, the alliance of the empire with the church, and the great Arian and semi-Arian controversy concerning the Divinity of Christ and the Holy Spirit.

2. The post-Nicene, or Christological and Augustinian age, extending to the fourth general council in 451, and including the Nestorian and Eutychian disputes on the person of Christ, and the Pelagian controversy on sin and grace.

3. The age of Leo the Great (440–461), or the rise of the papal supremacy in the West, amidst the barbarian devastations which made an end to the western Roman empire in 476.

4. The Justinian age (527–565), which exhibits the Byzantine state-church despotism at the height of its power, and at the beginning of its decline.

5. The Gregorian age (590–604) forms the transition from the ancient Graeco-Roman to the mediaeval Romano-Germanic Christianity, and will be more properly included in the church history of the middle ages.

\xchapter{Downfall of Heathenism and Victory of Christianity in the Roman Empire}

DOWNFALL OF HEATHENISM AND VICTORY OF CHRISTIANITY IN THE ROMAN EMPIRE.

GENERAL LITERATURE.

J. G. Hoffmann: Ruina Superstitionis Paganae. Vitemb. 1738. Tzschirner: Der Fall des Heidenthums. Leipz. 1829. A. Beugnot: Histoire de la destruction du paganisme en occident. Par. 1835. 2 vols. Et. Chastel (of Geneva): Histoire de la destruction du paganisme dans l’empire d’orient. Par. 1850. E. v. Lasaulx: Der Untergang des Hellenismus u. die Einziehung seiner Tempelgüter durch die christl. Kaiser. Münch. 1854. F. Lübker: Der Fall des Heidenthums. Schwerin, 1856. Ch. Merivale: Conversion of The Roman Empire. New York, 1865.

\xsection{Constantine The Great. a.d. 306-337}

§ 2. Constantine The Great. a.d. 306–337.

1. Contemporary Sources: Lactantius († 330): De mortibus persecutorum, cap. 18 sqq. Eusebius: Hist. Eccl. l. Ix. et x.; also his panegyric and very partial Vita Constantini, in 4 books (Εἰς τόν βίον τοῦ μακαρίου Κωνσταντίνου τοῦ βασιλέως) and his Panegyricus or De laudibus Constantini; in the editions of the hist. works of Euseb. by Valesius, Par. 1659–1673, Amstel. 1695, Cantabr. 1720; Zimmermann, Frcf. 1822; Heinichen, Lips. 1827–30; Burton, Oxon. 1838. Comp. the imperial documents in the Codex Theodos.l. xvi. also the Letters and Treatises of Athanasius († 373), and on the heathen side the Panegyric of Nazarius at Rome (321) and the Caesars of Julian († 363).

2. Later sources: Socrates: Hist. Eccl. l. i. Sozomenus: H. E. l. i et ii. Zosimus (a heathen historian and court-officer, comes et advocatus fisci, under Theodosius II.): ιστορία νέα, l. ii. ed. Bekker, Bonn. 1837. Eusebius and Zosimus present the extremes of partiality for and against Constantine. A just estimate of his character must be formed from the facts admitted by both, and from the effect of his secular and ecclesiastical policy.

3. Modern authorities. Mosheim: De reb. Christ. ante Const. M. etc., last section (p. 958 sqq. In Murdock’s Engl. transl., vol. ii. p. 454–481). Nath. Lardner, in the second part of his great work on the Credibility of the Gospel History, see Works ed. by Kippis, Lond. 1838, vol. iv. p. 3–55. Abbé de Voisin: Dissertation critique sur la vision de Constantin. Par. 1774. Gibbon: l.c. chs. xiv. and xvii.–xxi. Fr. Gusta: Vita di Constantino il Grande. Foligno, 1786. Manso: Das Leben Constantins des Gr. Bresl. 1817. Hug (R.C.): Denkschrift zur Ehrenrettung Constant. Frieb. 1829. Heinichen: Excurs. in Eus. Vitam Const. 1830. Arendt (R.C.): Const. u. sein Verb. zum Christenthum. Tüb. (Quartalschrift) 1834. Milman: Hist. of Christianity, etc., 1840, book iii. ch. 1–4. Jacob Burckhardt: Die Zeit Const. des Gr. Bas. 1853. Albert de Broglie: L’église et l’empire romain au IVme siècle. Par. 1856 (vols. i. and ii.). A. P. Stanley: Lectures on the Hist. of the Eastern Church, 1862, Lect. vi. p. 281 sqq. (Am. Ed.). Theod. Keim: Der Uebertritt Constantins des Gr. zum Christenthum. Zürich, 1862 (an apology for Constantine’s character against Burckhardt’s view).

The last great imperial persecution of the Christians under Diocletian and Galerius, which was aimed at the entire uprooting of the new religion, ended with the edict of toleration of 311 and the tragical ruin of the persecutors.\footnote{Comp. vol. i. § 57. Galerius died soon after of a disgusting and terrible disease (morbus pedicularis), described with great minuteness by Eusebius, H. E. viii. 16, and Lactantius, De mort. persec. c. 33.“His body,” says Gibbon, ch. xiv. “swelled by an intemperate course of life to an unwieldy corpulence, was covered with ulcers and devoured by innumerable swarms of those insects which have given their name to a most loathsome disease.” Diocletian had withdrawn from the throne in 305, and in 313 put an end to his embittered life by suicide. In his retirement he found more pleasure in raising cabbage than he had found in ruling the empire; a confession we may readily believe. (President Lincoln of the United States, during the dark days of the civil war in Dec. 1862, declared that he would gladly exchange his position with any common soldier in the tented field.) Maximin, who kept up the persecution in the East, even after the toleration edict, as long as he could, died likewise a violent death by poison, in 313. In this tragical end of their last three imperial persecutors the Christians saw a palpable judgment of God.} The edict of toleration was an involuntary and irresistible concession of the incurable impotence of heathenism and the indestructible power of Christianity. It left but a step to the downfall of the one and the supremacy of the other in the empire of the Caesars.

This great epoch is marked by the reign of Constantine I.\footnote{His full name in Latin is Caius Flavius Valerius Aurelius Claudius Constantinus Magnus.} He understood the signs of the times and acted accordingly. He was the man for the times, as the times were prepared for him by that Providence which controls both and fits them for each other. He placed himself at the head of true progress, while his nephew, Julian the Apostate, opposed it and was left behind. He was the chief instrument for raising the church from the low estate of oppression and persecution to well deserved honor and power. For this service a thankful posterity has given him the surname of the Great, to which he was entitled, though not by his moral character, yet doubtless by his military and administrative ability, his judicious policy, his appreciation and protection of Christianity, and the far-reaching consequences of his reign. His greatness was not indeed of the first, but of the second order, and is to be measured more by what he did than by what he was. To the Greek church, which honors him even as a canonized saint, he has the same significance as Charlemagne to the Latin.

Constantine, the first Christian Caesar, the founder of Constantinople and the Byzantine empire, and one of the most gifted, energetic, and successful of the Roman emperors, was the first representative of the imposing idea of a Christian theocracy, or of that system of policy which assumes all subjects to be Christians, connects civil and religious rights, and regards church and state as the two arms of one and the same divine government on earth. This idea was more fully developed by his successors, it animated the whole middle age, and is yet working under various forms in these latest times; though it has never been fully realized, whether in the Byzantine, the German, or the Russian empire, the Roman church-state, the Calvinistic republic of Geneva, or the early Puritanic colonies of New England. At the same time, however, Constantine stands also as the type of an undiscriminating and harmful conjunction of Christianity with politics, of the holy symbol of peace with the horrors of war, of the spiritual interests of the kingdom of heaven with the earthly interests of the state.

In judging of this remarkable man and his reign, we must by all means keep to the great historical principle, that all representative characters act, consciously or unconsciously, as the free and responsible organs of the spirit of their age, which moulds them first before they can mould it in turn, and that the spirit of the age itself, whether good or bad or mixed, is but an instrument in the hands of divine Providence, which rules and overrules all the actions and motives of men.

Through a history of three centuries Christianity had already inwardly overcome the world, and thus rendered such an outward revolution, as has attached itself to the name of this prince, both possible and unavoidable. It were extremely superficial to refer so thorough and momentous a change to the personal motives of an individual, be they motives of policy, of piety, or of superstition. But unquestionably every age produces and shapes its own organs, as its own purposes require. So in the case of Constantine. He was distinguished by that genuine political wisdom, which, putting itself at the head of the age, clearly saw that idolatry had outlived itself in the Roman empire, and that Christianity alone could breathe new vigor into it and furnish its moral support. Especially on the point of the external Catholic unity his monarchical politics accorded with the hierarchical episcopacy of the church. Hence from the year 313 he placed himself in close connection with the bishops, made peace and harmony his first object in the Donatist and Arian controversies and applied the predicate “catholic” to the church in all official documents. And as his predecessors were supreme pontiffs of the heathen religion of the empire, so he desired to be looked upon as a sort of bishop, as universal bishop of the external affairs of the church.\footnote{\textgreek{Ἐπίσκοπος τῶν ἐκτος} [\textgreek{πραγμάτων}], viz.: \textgreek{τῆς ἐκκλησίας}, in distinction from the proper bishops, the \textgreek{ἐπίσκοποι τῶν εἴσω τῆς ἐκκλησίας}. Vid. Eus.: Vit Const. iv. 24. Comp. § 24.} All this by no means from mere self-interest, but for the good of the empire, which, now shaken to its foundations and threatened by barbarians on every side, could only by some new bond of unity be consolidated and upheld until at least the seeds of Christianity and civilization should be planted among the barbarians themselves, the representatives of the future. His personal policy thus coincided with the interests of the state. Christianity appeared to him, as it proved in fact, the only efficient power for a political reformation of the empire, from which the ancient spirit of Rome was fast departing, while internal, civil, and religious dissensions and the outward pressure of the barbarians threatened a gradual dissolution of society.

But with the political he united also a religious motive, not clear and deep, indeed, yet honest, and strongly infused with the superstitious disposition to judge of a religion by its outward success and to ascribe a magical virtue to signs and ceremonies. His whole family was swayed by religious sentiment, which manifested itself in very different forms, in the devout pilgrimages of Helena, the fanatical Arianism of Constantia, and Constantius, and the fanatical paganism of Julian. Constantine adopted Christianity first as a superstition, and put it by the side of his heathen superstition, till finally in his conviction the Christian vanquished the pagan, though without itself developing into a pure and enlightened faith.\footnote{A similar view is substantially expressed by the great historian Niebuhr, Vorträge über Röm. Geschichte, 1848. iii. 302. Mosheim, in his work on the First Three Centuries, p. 965 sqq. (Murdock’s Transl. ii. 460 sqq.) labors to prove at length that Constantinewas no hypocrite, but sincerely believed, during the greater part of his life, that the Christian religion was the only true religion. Burckhardt, the most recent biographer of Constantine, represents him as a great politician of decided genius, but destitute of moral principle and religious interest. So also Dr. Baur.}

At first Constantine, like his father, in the spirit of the Neo-Platonic syncretism of dying heathendom, reverenced all the gods as mysterious powers; especially Apollo, the god of the sun, to whom in the year 308 he presented munificent gifts. Nay, so late as the year 321 he enjoined regular consultation of the soothsayers\footnote{The \emph{haruspices}, or interpreters of sacrifices, who foretold future events from the entrails of victims.} in public misfortunes, according to ancient heathen usage; even later, he placed his new residence, Byzantium, under the protection of the God of the Martyrs and the heathen goddess of Fortune;\footnote{According to Eusebius (Vit. Const. l. iii. c. 48) he dedicated Constantinople to “the God of the martyrs,” but, according to Zosimus (Hist. ii. c. 31), to two female deities, probably Mary and Fortuna. Subsequently the city stood under the special protection of the Virgin Mary.} and down to the end of his life he retained the title and the dignity of a Pontifex Maximus, or high-priest of the heathen hierarchy.\footnote{His successors also did the same, down to Gratian, 375, who renounced the title, then become quite empty.} His coins bore on the one side the letters of the name of Christ, on the other the figure of the Sun-god, and the inscription “Sol invictus.” Of course these inconsistencies may be referred also to policy and accommodation to the toleration edict of 313. Nor is it difficult to adduce parallels of persons who, in passing from Judaism to Christianity, or from Romanism to Protestantism, have so wavered between their old and their new position that they might be claimed by both. With his every victory, over his pagan rivals, Galerius, Maxentius, and Licinius, his personal leaning to Christianity and his confidence in the magic power of the sign of the cross increased; yet he did not formally renounce heathenism, and did not receive baptism until, in 337, he was laid upon the bed of death.

He had an imposing and winning person, and was compared by flatterers with Apollo. He was tall, broad-shouldered, handsome, and of a remarkably vigorous and healthy constitution, but given to excessive vanity in his dress and outward demeanor, always wearing an oriental diadem, a helmet studded with jewels, and a purple mantle of silk richly embroidered with pearls and flowers worked in gold,\footnote{Euseb. Laud. Const. c. 5.} His mind was not highly cultivated, but naturally clear, strong, and shrewd, and seldom thrown off its guard. He is said to have combined a cynical contempt of mankind with an inordinate love of praise. He possessed a good knowledge of human nature and administrative energy and tact.

His moral character was not without noble traits, among which a chastity rare for the time,\footnote{All Christian accounts speak of his continence, but Julian insinuates the contrary, and charges him with the old Roman vice of voracious gluttony (Caes. 329, 335).} and a liberality and beneficence bordering on wastefulness were prominent. Many of his laws and regulations breathed the spirit of Christian justice and humanity, promoted the elevation of the female sex, improved the condition of slaves and of unfortunates, and gave free play to the efficiency of the church throughout the whole empire. Altogether he was one of the best, the most fortunate, and the most influential of the Roman emperors, Christian and pagan.

Yet he had great faults. He was far from being so pure and so venerable as Eusebius, blinded by his favor to the church, depicts him, in his bombastic and almost dishonestly eulogistic biography, with the evident intention of setting him up as a model for all future Christian princes. It must, with all regret, be conceded, that his progress in the knowledge of Christianity was not a progress in the practice of its virtues. His love of display and his prodigality, his suspiciousness and his despotism, increased with his power.

The very brightest period of his reign is stained with gross crimes, which even the spirit of the age and the policy of an absolute monarch cannot excuse. After having reached, upon the bloody path of war, the goal of his ambition, the sole possession of the empire, yea, in the very year in which he summoned the great council of Nicaea, he ordered the execution of his conquered rival and brother-in-law, Licinius, in breach of a solemn promise of mercy (324).\footnote{Eusebius justifies this procedure towards an enemy of the Christians by the laws of war. But what becomes of the breach of a solemn pledge? The murder of Crispus and Fausta he passes over in prudent silence, in violation of the highest duty of the historian to relate the truth and the whole truth.} Not satisfied with this, he caused soon afterwards, from political suspicion, the death of the young Licinius, his nephew, a boy of hardly eleven years. But the worst of all is the murder of his eldest son, Crispus, in 326, who had incurred suspicion of political conspiracy, and of adulterous and incestuous purposes towards his step-mother Fausta, but is generally regarded as innocent. This domestic and political tragedy emerged from a vortex of mutual suspicion and rivalry, and calls to mind the conduct of Philip II. towards Don Carlos, of Peter the Great towards his son Alexis, and of Soliman the Great towards his son Mustapha. Later authors assert, though gratuitously, that the emperor, like David, bitterly repented of this sin. He has been frequently charged besides, though it would seem altogether unjustly, with the death of his second wife Fausta (326?), who, after twenty years, of happy wedlock, is said to have been convicted of slandering her stepson Crispus, and of adultery with a slave or one of the imperial guards, and then to have been suffocated in the vapor of an over-heated bath. But the accounts of the cause and manner of her death are so late and discordant as to make Constantine’s part in it at least very doubtful.\footnote{Zosimus, certainly in heathen prejudice and slanderous extravagance, ascribes to Constantine under the instigation of his mother Helena, who was furious at the loss of her favorite grandson, the death of two women, the innocent Fausta and an adulteress, the supposed mother of his three successors; Philostorgius, on the contrary, declares Fausta guilty (H. E. ii. 4; only fragmentary). Then again, older witnesses indirectly contradict this whole view; two orations, namely, of the next following reign, which imply, that Fausta survived the death of her son, the younger Constantine, who outlived his father by three years. Comp. Julian. Orat. i., and Monod. in Const. Jun. c. 4, ad Calcem Eutrop., cited by Gibbon, ch. xviii., notes 25 and 26. Evagrius denies both the murder of Crispus and of Fausta, though only on account of the silence of Eusebius, whose extreme partiality for his imperial friend seriously impairs the value of his narrative. Gibbon and still more decidedly Niebuhr (Vorträge über Röm. Geschichte, iii. 302) are inclined to acquit Constantine of all guilt in the death of Fausta. The latest biographer, Burckhardt (l.c. p. 375) charges him with it rather hastily, without even mentioning the critical difficulties in the way. So also Stanley (l.c. p. 300).}

At all events Christianity did not produce in Constantine a thorough moral transformation. He was concerned more to advance the outward social position of the Christian religion, than to further its inward mission. He was praised and censured in turn by the Christians and Pagans, the Orthodox and the Arians, as they successively experienced his favor or dislike. He bears some resemblance to Peter the Great both in his public acts and his private character, by combining great virtues and merits with monstrous crimes, and he probably died with the same consolation as Peter, whose last words were: “I trust that in respect of the good I have striven to do my people (the church), God will pardon my sins.” It is quite characteristic of his piety that he turned the sacred nails of the Saviour’s cross which Helena brought from Jerusalem, the one into the bit of his war-horse, the other into an ornament of his helmet. Not a decided, pure, and consistent character, he stands on the line of transition between two ages and two religions; and his life bears plain marks of both. When at last on his death bed he submitted to baptism, with the remark, “Now let us cast away all duplicity,” he honestly admitted the conflict of two antagonistic principles which swayed his private character and public life.\footnote{The heathen historians extol the earlier part of his reign, and depreciate the later. Thus Eutropius, x. 6: “In primo imperii tempore optimis principibus, ultimo mediis comparandus.” With this judgment Gibbon agrees (ch. xviii.), presenting in Constantinean inverted Augustus: “In the life of Augustus we behold the tyrant of the republic, converted, almost by imperceptible degrees, into the father of his country and of human kind. In that of Constantine, we may contemplate a hero, who had so long inspired his subjects with love, and his enemies with terror, degenerating into a cruel and dissolute monarch, corrupted by his fortune, or raised by conquest above the necessity of dissimulation.” But this theory of progressive degeneracy, adopted also by F. C. Schlosser in his Weltgeschichte, by Stanley, l.c. p. 297, and many others, is as untenable as the opposite view of a progressive improvement, held by Eusebius, Mosheim, and other ecclesiastical historians. For, on the one hand, the earlier life of Constantinehas such features of cruelty as the surrender of the conquered barbarian kings to the wild beasts in the ampitheatre at Treves in 310 or 311, for which he was lauded by a heathen orator; the ungenerous conduct toward Herculius, his father-in-law; the murder of the infant son of Maxentius; and the triumphal exhibition of the head of Maxentius on his entrance into Rome in 312. On the other hand his most humane laws, such as the abolition of the gladiatorial shows and of licentious and cruel rites, date from his later reign.}

From these general remarks we turn to the leading features of Constantine’s life and reign, so far as they bear upon the history of the church. We shall consider in order his youth and training, the vision of the Cross, the edict of toleration, his legislation in favor of Christianity, his baptism and death.

Constantine, son of the co-emperor Constantius Chlorus, who reigned over Gaul, Spain, and Britain till his death in 306, was born probably in the year 272, either in Britain or at Naissus (now called Nissa), a town of Dardania, in Illyricum.\footnote{According to Baronius (Ann. 306, n. 16) and others he was born in Britain, because an ancient panegyric of 307 says that Constantine ennobled Britain by his birth (tu Britannias nobiles oriendo fecisti); but this may be understood of his royal as well as of his natural birth, since he was there proclaimed Caesar by the soldiers. The other opinion rests also on ancient testimonies, and is held by Pagi, Tillemont, and most of the recent historians.} His mother was Helena, daughter of an innkeeper,\footnote{Ambrose(De obitu Theodos.) calls her \emph{stabulariam}, when Constantius made her acquaintance.} the first wife of Constantius, afterwards divorced, when Constantius, for political reasons, married a daughter of Maximian.\footnote{This is the more probable view, and rests on good authority. Zosimus and even the Paschal Chronicle call Helena the concubine of Constantius, and Constantine illegitimate. But in this case it would be difficult to understand that he was so well treated at the court of Diocletian and elected Caesar without opposition, since Constantius had three sons and three daughters by a legal wife, Theodora. It is possible, however, that Helena was first a concubine and afterwards legally married. Constantine, when emperor, took good care of her position and bestowed upon her the title of Augusta and empress with appropriate honors.} She is described by Christian writers as a discreet and devout woman, and has been honored with a place in the catalogue of saints. Her name is identified with the discovery of the cross and the pious superstitions of the holy places. She lived to a very advanced age and died in the year 326 or 327, in or near the city of Rome. Rising by her beauty and good fortune from obscurity to the splendor of the court, then meeting the fate of Josephine, but restored to imperial dignity by her son, and ending as a saint of the Catholic church: Helena would form an interesting subject for a historical novel illustrating the leading events of the Nicene age and the triumph of Christianity in the Roman empire.

Constantine first distinguished himself in the service of Diocletian in the Egyptian and Persian wars; went afterwards to Gaul and Britain, and in the Praetorium at York was proclaimed emperor by his dying father and by the Roman troops. His father before him held a favorable opinion of the Christians as peaceable and honorable citizens, and protected them in the West during the Diocletian persecution in the East. This respectful tolerant regard descended to Constantine, and the good effects of it, compared with the evil results of the opposite course of his antagonist Galerius, could but encourage him to pursue it. He reasoned, as Eusebius reports from his own mouth, in the following manner: “My father revered the Christian God and uniformly prospered, while the emperors who worshipped the heathen gods, died a miserable death; therefore, that I may enjoy a happy life and reign, I will imitate the example of my father and join myself to the cause of the Christians, who are growing daily, while the heathen are diminishing.” This low utilitarian consideration weighed heavily in the mind of an ambitious captain, who looked forward to the highest seat of power within the gift of his age. Whether his mother, whom he always revered, and who made a pilgrimage to Jerusalem in her eightieth year (a.d. 325), planted the germ of the Christian faith in her son, as Theodoret supposes, or herself became a Christian through his influence, as Eusebius asserts, must remain undecided. According to the heathen Zosimus, whose statement is unquestionably false and malicious, an Egyptian, who came out of Spain (probably the bishop Hosius of Cordova, a native of Egypt, is intended), persuaded him, after the murder of Crispus (which did not occur before 326), that by converting to Christianity he might obtain forgiveness of his sins.

The first public evidence of a positive leaning towards the Christian religion he gave in his contest with the pagan Maxentius, who had usurped the government of Italy and Africa, and is universally represented as a cruel, dissolute tyrant, hated by heathens and Christians alike,\footnote{Even Zosimus gives the most unfavorable account of him.} called by the Roman people to their aid, Constantine marched from Gaul across the Alps with an army of ninety-eight thousand soldiers of every nationality, and defeated Maxentius in three battles; the last in October, 312, at the Milvian bridge, near Rome, where Maxentius found a disgraceful death in the waters of the Tiber.

Here belongs the familiar story of the miraculous cross. The precise day and place cannot be fixed, but the event must have occurred shortly before the final victory over Maxentius in the neighborhood of Rome. As this vision is one of the most noted miracles in church history, and has a representative significance, it deserves a closer examination. It marks for us on the one hand the victory of Christianity over paganism in the Roman empire, and on the other the ominous admixture of foreign, political, and military interests with it.\footnote{“It was,” says Milman (Hist. of Christianity, p. 288, N. York ed.), “the first advance to the military Christianity of the Middle Ages; a modification of the pure religion of the Gospel, if directly opposed to its genuine principles, still apparently indispensable to the social progress of man; through which the Roman empire and the barbarous nations, which were blended together in the vast European and Christian system, must necessarily have passed before they could arrive at a higher civilization and a purer Christianity.”} We need not be surprised that in the Nicene age so great a revolution and transition should have been clothed with a supernatural character.

The occurrence is variously described and is not without serious difficulties. Lactantius, the earliest witness, some three years after the battle, speaks only of a dream by night, in which the emperor was directed (it is not stated by whom, whether by Christ, or by an angel) to stamp on the shields of his soldiers “the heavenly sign of God,” that is, the cross with the name of Christ, and thus to go forth against his enemy.\footnote{De mortibus persecutorum, c. 44 (ed. Lips. II. 278 sq.): “Commonitus est in quiete Constantinus, ut coeleste signum Dei notaret in scutis, atque ita proelium committeret. Fecit ut jussus est, et transverse X litera, summo capite circumflexo Christum in scutis notat [i.e., he ordered the name of Christ or the two first letters X and P to be put on the shields of his soldiers]. Quo signo armatus exercitus capit ferrum.”—This work is indeed by Burckhardt and others denied to Lactantius, but was at all events composed soon after the event, about 314 or 315, while Constantine was as yet on good terms with Licinius, to whom the author, c. 46, ascribes a similar vision of an angel, who is said to have taught him a form of prayer on his expedition against the heathen tyrant Maximin.} Eusebius, on the contrary, gives a more minute account on the authority of a subsequent private communication of the aged Constantine himself under oath—not, however, till the year 338, a year after the death of the emperor, his only witness, and twenty-six years after the event.\footnote{In his Vita Constant. i. 27-30, composed about 338, a work more panegyrical than historical, and abounding in vague declamation and circumlocution. But in his Church History, written before 326, though he has good occasion (l. ix. c. 8, 9), Eusebius says nothing of the occurrence, whether through oversight or ignorance, or of purpose, it is hard to decide. In any case the silence casts suspicion on the details of his subsequent story, and has been urged against it not only by Gibbon, but also by Lardner and others.} On his march from Gaul to Italy (the spot and date are not specified), the emperor, whilst earnestly praying to the true God for light and help at this critical time, saw, together with his army,\footnote{This is probably a mistake or an exaggeration. For if a whole army consisting of many thousand soldiers of every nation had seen the vision of the cross, Eusebius might have cited a number of living witnesses, and Constantine might have dispensed with a solemn oath. But on the other hand the two heathen witnesses (see below) extend the vision likewise to the soldiers.} in clear daylight towards evening, a shining cross in the heavens above the sun) with the inscription: “By this conquer,”\footnote{\textgreek{τούτῳ} [\textgreek{τῷ σημείῳ}] \textgreek{νίκα; } Hac, or Hoc [sc. signo] vince, or vinces. Eusebius leaves the impression that the inscription was in Greek. But Nicephorus and Zonaras say that it was in Latin.} and in the following night Christ himself appeared to him while he slept, and directed him to have a standard prepared in the form of this sign of the cross, and with that to proceed against Maxentius and all other enemies. This account of Eusebius, or rather of Constantine himself, adds to the night dream of Lactantius the preceding vision of the day, and the direction concerning the standard, while Lactantius speaks of the inscription of the initial letters of Christ’s name on the shields of the soldiers. According to Rufinus,\footnote{Hist. Eccl. ix, 9. Comp. the similar account of Sozomenus, H. E. i. 3.} a later historian, who elsewhere depends entirely on Eusebius and can therefore not be regarded as a proper witness in the case, the sign of the cross appeared to Constantine in a dream (which agrees with the account of Lactantius), and upon his awaking in terror, an angel (not Christ) exclaimed to him: “Hoc vince.” Lactantius, Eusebius, and Rufinus are the only Christian writers of the fourth century, who mention the apparition. But we have besides one or two heathen testimonies, which, though vague and obscure, still serve to strengthen the evidence in favor of some actual occurrence. The contemporaneous orator Nazarius, in a panegyric upon the emperor, pronounced March 1, 321, apparently at Rome, speaks of an army of divine warriors and a divine assistance which Constantine received in the engagement with Maxentius, but he converts it to the service of heathenism by recurring to old prodigies, such as the appearance of Castor and Pollux.\footnote{Nazar. Paneg. in Const. c. 14: “In ore denique est omnium Galliarum [this would seem to indicate a pretty general rumor of some supernatural assistance], exercitus visos, qui se divinitus missos prae se ferebant,” etc. Comp. Baronius, Annal. ad ann. 312, n. 11. This historian adduces also (n. 14) another and still older pagan testimony from an anonymous panegyrical orator, who, in 313, speaks of a certain undefined omen which filled the soldiers of Constantinewith misgivings and fears, while it emboldened him to the combat. Baronius and J. H. Newman (in his “Essay on Miracles”) plausibly suppose this \emph{omen} to have been the cross.}

This famous tradition may be explained either as a real miracle implying a personal appearance of Christ,\footnote{This is the view of the older historians, Protestant as well as Catholic. Among more modern writers on the subject it has hardly any advocates of note, except Döllinger (R.C.), J. H.Newman (in his “Essay on Miracles,” published in 1842, before his transition to Romanism, and prefixed to the first volume of his translation of Fleury), and Guericke (Lutheran). Comp. also De Broglie, i. 219 and 442.} or as a pious fraud,\footnote{So more or less distinctly Hoornebeck (of Leyden), Thomasius, Arnold, Lardner, Gibbon, and Waddington. The last writer (Hist. of the Church, vol. i. 171) disposes of it too summarily by the remark that “this flattering fable may very safely be consigned to contempt and oblivion.” Burckhardt, the most recent biographer of Constantine, is of the same opinion. He considers the story as a joint fabrication of Eusebius and the emperor, and of no historical value whatever (Die Zeit Constantins des Gr. 1853, pp. 394 and 395). Lardner saddles the lie exclusively upon the emperor (although he admits him otherwise to have been a sincere Christian), and tries to prove that Eusebius himself hardly believed it.} or as a natural phenomenon in the clouds and an optical illusion,\footnote{This is substantially the theory of J. A. Fabricius (in a special dissertation), Schröckh (vol. v. 83), Manso, Heinichen (in the first Excursus to his ed. of Euseb), Gieseler, Neander, Milman, Robertson, and Stanley. Gieseler (vol. i. § 56, note 29) mentions similar cross-like clouds which appeared in Germany, Dec. 1517 and 1552, and were mistaken by contemporary Lutherans for supernatural signs. Stanley (Lectures on the Eastern Church, p. 288) refers to the natural phenomenon known by the name of “parhelion,” which in an afternoon sky not unfrequently assumes almost the form of the cross. He also brings in, as a new illustration, the Aurora Borealis which appeared in November, 1848, and was variously interpreted, in France as forming the letters L. N., in view of the approaching election of Louis Napoleon, in Rome as the blood of the murdered Rossi crying for vengeance from heaven against his assassins. Mosheim, after a lengthy discussion of the subject in his large work on the ante-Nicene age, comes to no definite conclusion, but favors the hypothesis of a mere dream or a psychological illusion. Neander and Robertson connect with the supposition of a natural phenomenon in the skies a dream of Constantine which reflected the optical vision of the day. Keim, the latest writer on the subject, l.c. p. 89, admits the dream, but denies the cross in the clouds. So Mosheim.} or finally as a prophetic dream.

The propriety of a miracle, parallel to the signs in heaven which preceded the destruction of Jerusalem, might be justified by the significance of the victory as marking a great epoch in history, namely, the downfall of paganism and the establishment of Christianity in the empire. But even if we waive the purely critical objections to the Eusebian narrative, the assumed connection, in this case, of the gentle Prince of peace with the god of battle, and the subserviency of the sacred symbol of redemption to military ambition, is repugnant to the genius of the gospel and to sound Christian feeling, unless we stretch the theory of divine accommodation to the spirit of the age and the passions and interests of individuals beyond the ordinary limits. We should suppose, moreover, that Christ, if he had really appeared to Constantine either in person (according to Eusebius) or through angels (as Rufinus and Sozomen modify it), would have exhorted him to repent and be baptized rather than to construct a military ensign for a bloody battle.\footnote{Dr. Murdock (notes to his translation of Mosheim) raises the additional objection, which has some force from his Puritan standpoint: “If the miracle of the luminous cross was a reality, has not God himself sanctioned the use of the cross as the appointed symbol of our religion? so that there is no superstition in the use of it, but the Catholics are correct and the Protestants in an error on this subject?”} In no case can we ascribe to this occurrence, with Eusebius, Theodoret, and older writers, the character of a sudden and genuine conversion, as to Paul’s vision of Christ on the way to Damascus;\footnote{Theodoret says that Constantinewas called not of men or by men (\textgreek{οὐκ ἀπ  ἀνθρώπου, οὐδὲ δι  ἀνθρώπου, }Gal. i. 1), but from heaven, as the divine apostle Paul was (\textgreek{οὐρανόθεν κατὰ τὸν θεῖον ἀπόστολον}). Hist. Eccl. l. i. c. 2.} for, on the one hand, Constantine was never hostile to Christianity, but most probably friendly to it from his early youth, according to the example of his father; and, on the other, he put off his baptism quite five and twenty years, almost to the hour of his death.

The opposite hypothesis of a mere military stratagem or intentional fraud is still more objectionable, and would compel us either to impute to the first Christian emperor at a venerable age the double crime of falsehood and perjury, or, if Eusebius invented the story, to deny to the “father of church history” all claim to credibility and common respectability. Besides it should be remembered that the older testimony of Lactantius, or whoever was the author of the work on the Deaths of Persecutors, is quite independent of that of Eusebius, and derives additional force from the vague heathen rumors of the time. Finally the Hoc vince which has passed into proverbial significance as a most appropriate motto of the invincible religion of the cross, is too good to be traced to sheer falsehood. Some actual fact, therefore, must be supposed to underlie the tradition, and the question only is this, whether it was an external visible phenomenon or an internal experience.

The hypothesis of a natural formation of the clouds, which Constantine by an optical illusion mistook for a supernatural sign of the cross, besides smacking of the exploded rationalistic explanation of the New Testament miracles, and deriving an important event from a mere accident, leaves the figure of Christ and the Greek or Latin inscription: By this sign thou shalt conquer! altogether unexplained.

We are shut up therefore to the theory of a dream or vision, and an experience within the mind of Constantine. This is supported by the oldest testimony of Lactantius, as well as by the report of Rufinus and Sozomen, and we do not hesitate to regard the Eusebian cross in the skies as originally a part of the dream,\footnote{So Sozomenus, H. E. lib. i. cap. 3, expressly represents it: \textgreek{ὅναρ εἶδε τὸ τοῦ σταυροῦ σημεῖον σελαγίζον}etc. Afterwards he gives, it is true, the fuller report of Eusebius in his own words. Comp. Rufin. ix. 9; Euseb. Vit. Const. i. 29; Lact. De mort. persec. 44, and the allusions of the heathen panegyrists.} which only subsequently assumed the character of an outward objective apparition either in the imagination of Constantine, or by a mistake of the memory of the historian, but in either case without intentional fraud. That the vision was traced to supernatural origin, especially after the happy success, is quite natural and in perfect keeping with the prevailing ideas of the age.\footnote{Licinius before the battle with Maximin had a vision of an angel who taught him a prayer for victory (Lactant. De mort. persec. c. 46). Julianthe Apostate was even more superstitious in this respect than his Christian uncle, and fully addicted to the whole train of omens, presages, prodigies, spectres, dreams, visions, auguries, and oracles (Comp. below, § 4). On his expedition against the Persians he was supposed by Libanius to have been surrounded by a whole army of gods, which, however, in the view of Gregory of Nazianzen, was a host of demons. See Ullmann, Gregory of Naz., p. 100.} Tertullian and other ante-Nicene and Nicene fathers attributed many conversions to nocturnal dreams and visions. Constantine and his friends referred the most important facts of his life, as the knowledge of the approach of hostile armies, the discovery of the holy sepulchre, the founding of Constantinople, to divine revelation through visions and dreams. Nor are we disposed in the least to deny the connection of the vision of the cross with the agency of divine Providence, which controlled this remarkable turning point of history. We may go farther and admit a special providence, or what the old divines call a providentia specialissima; but this does not necessarily imply a violation of the order of nature or an actual miracle in the shape of an objective personal appearance of the Saviour. We may refer to a somewhat similar, though far less important, vision in the life of the pious English Colonel James Gardiner.\footnote{According to the account of his friend, Dr. Philip Doddridge, who learned the facts from Gardiner, as Eusebius from Constantine. When engaged in serious meditation on a Sabbath night in July, 1719, Gardiner “suddenly thought he saw an unusual blaze of light fall on the book while he was reading, which he at first imagined might have happened by some accident in the candle. But lifting up his eyes, he apprehended, to his extreme amazement, that there was before him, as it were suspended in the air, a visible representation of the Lord Jesus Christ upon the cross, surrounded with a glory; and was impressed as if a voice, or something equivalent to a voice, had come to him, to this effect: ’O sinner, did I suffer this for thee, and are these the returns?’ ” After this event he changed from a dissolute worldling to an earnest and godly man. But the whole apparition was probably, after all, merely an inward one. For the report adds as to the voice: “Whether this were an audible voice, or only \emph{a strong impression on} \emph{his mind}, equally striking, he did not seem confident, though he judged it to be the former. He thought he was awake. But everybody knows how easy it is towards midnight to fall into a doze over a dull or even a good book. It is very probable then that this apparition resolves itself into a significant dream which marked an epoch in his life. No reflecting person will on that account doubt the seriousness of Gardiner’s conversion, which was amply proved by his whole subsequent life, even far more than Constantine’s was.} The Bible itself sanctions the general theory of providential or prophetic dreams and nocturnal visions through which divine revelations and admonitions are communicated to men.\footnote{Numbers xii. 6: “I the Lord will make myself known in a vision, and will speak in a dream.” Job xxxiii. 15, 16: “In a dream, in a vision of the night, when deep sleep falleth upon men, in slumberings upon the bed, then he openeth the ears of men and sealeth their instruction.” For actual facts see Gen. xxxi. 10, 24; xxxvii. 5; 1 Kings iii. 5; Dan. ii. 4, 36; vii. 1; Matt. i. 20; ii. 12, 13, 19, 22; Acts x. 17; xxii. 17, 18.}

The facts, therefore, may have been these. Before the battle Constantine, leaning already towards Christianity as probably the best and most hopeful of the various religions, seriously sought in prayer, as he related to Eusebius, the assistance of the God of the Christians, while his heathen antagonist Maxentius, according to Zosimus,\footnote{Histor. ii. 16.} was consulting the sibylline books and offering sacrifice to the idols. Filled with mingled fears and hopes about the issue of the conflict, he fell asleep and saw in a dream the sign of the cross of Christ with a significant inscription and promise of victory. Being already familiar with the general use of this sign among the numerous Christians of the empire, many of whom no doubt were in his own army, he constructed the labarum,\footnote{\textgreek{Λάβωρον}, also \textgreek{λάβουρον}; derived not from \emph{labor}, nor from \textgreek{λάφυρον}, i.e. \emph{praeda,} nor from \textgreek{λαβεῖν}, but probably from a barbarian root, otherwise unknown, and introduced into the Roman terminology, long before Constantine, by the Celtic or Germanic recruits. Comp. Du Cange, Glossar., and Suicer, Thesaur. s. h. v. The labarum, as described by Eusebius, who saw it himself (Vita Const. i. 30), consisted of a long spear overlaid with gold, and a crosspiece of wood, from which hung a square flag of purple cloth embroidered and covered with precious stones. On the of top of the shaft was a crown composed of gold and precious stones, and containing the monogram of Christ (see next note), and just under this crown was a likeness the emperor and his sons in gold. The emperor told Eusebius (I. ii. c. 7) some incredible things about this labarum, e.g. that none of its bearers was ever hurt by the darts of the enemy.} or rather he changed the heathen labarum into a standard of the Christian cross with the Greek monogram of Christ,\footnote{X and P, the first two letters of the name of Christ, so written upon one another as to make the form of the cross: P with x (Rho with Chi on the lower part) or Pwith\textbf{—}(Rho with a dash on the lower part to make a cross)\textbf{,} or αPω(i.e. Christos—Alpha and Omega, the beginning and the end with a chi on the stem to make the cross), and similar forms, of which Münter (Sinnbilder der alten Christen, p. 36 sqq.) has collected from ancient coins, vessels, and tombstones more than twenty. The monogram, as well as the sign of the cross, was in use among the Christians Iong before Constantine, probably as early as the Antonines and Hadrian. Yea, the standards and trophies of victory generally had the appearance of a cross, as Minucius Felix, Tertullian, Justin, and other apologists of the second century told the heathens. According to Killen (Ancient Church, p. 317, note), who quotes Aringhus, Roma subterranea, ii. p. 567, as his authority, the famous monogram (of course in a different sense) is found even before Christ on coins of the Ptolemies. The only thing new, therefore, was the \emph{union} of this symbol, in its \emph{Christian} sense and application, with the Roman \emph{military} \emph{standard.}} which he had also put upon the shields of the soldiers. To this cross-standard, which now took the place of the Roman eagles, he attributed the decisive victory over the heathen Maxentius.

Accordingly, after his triumphal entrance into Rome, he had his statue erected upon the forum with the labarum in his right hand, and the inscription beneath: “By this saving sign, the true token of bravery, I have delivered your city from the yoke of the tyrant.”\footnote{Eus., H. E. ix. 9: \textgreek{Τούτῳ τῷ σωτηριώδει} (salutari, not singulari, as Rufinus has it) \textgreek{σημείῳ, τῶ ἀληθινῷ ἐλέγχῳ τῶς ἀνδρίας , τήν πόλιν ὑμῶν ἀπὸ ζυγοῦ τοῦ τυράννου διασωθεῖσαν ἐλευθέρωσα, κ. τ. λ.} Gibbon, however thinks it more probable, that at least the labarum and the inscription date only from the second or third visit of Constantineto Rome.} Three years afterwards the senate erected to him a triumphal arch of marble, which to this day, within sight of the sublime ruins of the pagan Colosseum, indicates at once the decay of ancient art, and the downfall of heathenism; as the neighboring arch of Titus commemorates the downfall of Judaism and the destruction of the temple. The inscription on this arch of Constantine, however, ascribes his victory over the hated tyrant, not only to his master mind, but indefinitely also to the impulse of Deity;\footnote{“Instinctu Divinitatis et mentis magnitudine.” \emph{Divinitas} may be taken as an ambiguous word like Providence, “which veils Constantine’s passage from Paganism to Christianity.”} by which a Christian would naturally understand the true God, while a heathen, like the orator Nazarius, in his eulogy on Constantine, might take it for the celestial guardian power of the “urbs aeterna.”

At all events the victory of Constantine over Maxentius was a military and political victory of Christianity over heathenism; the intellectual and moral victory having been already accomplished by the literature and life of the church in the preceding period. The emblem of ignominy and oppression\footnote{Cicero says, pro Raberio, c. 5: “Nomen ipsum \emph{crucis} absit non modo a corpore civium Romanorum, sed etiam a cogitatione, oculis, auribus.” With other ancient heathens, however, the Egyptians, the Buddhists, and even the aborigines of Mexico, the cross seems to have been in use as a religious symbol. Socrates relates (H. E. v. 17) that at the destruction of the temple of Serapis, among the hieroglyphic inscriptions forms of crosses were found, which pagans and Christians alike referred to their respective religions. Some of the heathen converts conversant with hieroglyphic characters interpreted the form of the cross to mean \emph{the Life to come}. According to Prescott (Conquest of Mexico, iii. 338-340) the Spaniards found the cross among the objects of worship in the idol temples of Anahnac.} became thenceforward the badge of honor and dominion, and was invested in the emperor’s view, according to the spirit of the church of his day, with a magic virtue.\footnote{Even church teachers long before Constantine, Justin, Tertullian, Minucius Felix, in downright opposition to this pagan antipathy, had found the sign of the cross everywhere on the face of nature and of human life; in the military banners and trophies of victory, in the ship with swelling sails and extended oars, in the plow in the flying bird, in man swimming or praying, in the features of the face and the form of the body with outstretched arms. Hence the daily use of the of the cross by the early Christians. Comp. vol. ii. § 77 (p. 269 sqq.).} It now took the place of the eagle and other field-badges, under which the heathen Romans had conquered the world. It was stamped on the imperial coin, and on the standards, helmets, and shields of the soldiers. Above all military representations of the cross the original imperial labarum shone in the richest decorations of gold and gems; was intrusted to the truest and bravest fifty of the body guard; filled the Christians with the spirit of victory, and spread fear and terror among their enemies; until, under the weak successors of Theodosius II., it fell out of use, and was lodged as a venerable relic in the imperial palace at Constantinople.

After this victory at Rome (which occurred October 27, 312), Constantine, in conjunction with his eastern colleague, Licinius, published in January, 313, from Milan, an edict of religious toleration, which goes a step beyond the edict of the still anti-Christian Galerius in 311, and grants, in the spirit of religious eclecticism, full freedom to all existing forms of worship, with special reference to the Christian.\footnote{This in the second edict of toleration, not the third, as was formerly supposed. An edict of 312 does not exist and rests on a mistake. See vol. ii. § 25, p. 72.} The edict of 313 not only recognized Christianity within existing limits, but allowed every subject of the Roman empire to choose whatever religion he preferred.\footnote{“Haec ordinanda esse credidimus ... ut daremus et Christianis et omnibus liberam potestatem sequendi religionem, quamquisque voluisset ... ut nulli omnino facultatem obnegandam putaremus, qui vel observationi Christianorum, vel ei religioni mentem suam dederet, quam ipse sibi aptissimam esse sentiret ... ut, amotis omnibus ominino conditionibus [by which are meant, no doubt, the restrictions of toleration in the edict of 311], nunc libere ac simpliciter unusquisque eorum qui eandem observandae religioni Christianorum gerunt voluntatem, citra ullam inquietudinem et molestiam sui id ipsum observare contendant.” Lact., De mort, persec. c. 48 (ii. p. 282, ed. Fritzsche). Eusebius gives the edict in a stiff and obscure Greek translation, with some variations, H. E. x. 5. Comp. Niceph. H. E. vii. 41. Also a special essay on the edicts of toleration, by Theod. Keim in the Tübinger Theolog. Jahrbücher for 1852, and Mason, persecution of Diocletian, pp. 299 and 326.} At the same time the church buildings and property confiscated in the Diocletian persecution were ordered to be restored, and private property-owners to be indemnified from the imperial treasury.

In this notable edict, however, we should look in vain for the modern Protestant and Anglo-American theory of religious liberty as one of the universal and inalienable rights of man. Sundry voices, it is true, in the Christian church itself, at that time, as before and after, declared against all compulsion in religion.\footnote{Compare the remarkable passages of Tertullian, cited in vol. ii. § 13, p. 35. Lactantius likewise, in the beginning of the fourth century, says, Instit. div. l. v. c. 19 (i. p. 267 sq. ed. Lips.): “Non est opus vi et injuria, quia religio cogi non potest; verbis potius, quam verberibus res agenda est, ut sit voluntas .... Defendenda religio est, non occidendo, sed moriendo; non saevitia, sed patientia; non scelere, sed fide .... Nam si sanguine, si tormentis, si malo religionem defendere velis, jam non defendetur illa, sed polluetur atque violabitur. Nihil est enim tam voluntarium, quam religio, in qua si animus sacrificantis aversus est, jam sublata, jam nulla est.” Comp. c. 20.} But the spirit of the Roman empire was too absolutistic to abandon the prerogative of a supervision of public worship. The Constantinian toleration was a temporary measure of state policy, which, as indeed the edict expressly states the motive, promised the greatest security to the public peace and the protection of all divine and heavenly powers, for emperor and empire. It was, as the result teaches, but the necessary transition step to a new order of things. It opened the door to the elevation of Christianity, and specifically of Catholic hierarchical Christianity, with its exclusiveness towards heretical and schismatic sects, to be the religion of the state. For, once put on equal footing with heathenism, it must soon, in spite of numerical minority, bear away the victory from a religion which had already inwardly outlived itself.

From this time Constantine decidedly favored the church, though without persecuting or forbidding the pagan religions. He always mentions the Christian church with reverence in his imperial edicts, and uniformly applies to it, as we have already observed, the predicate of catholic. For only as a catholic, thoroughly organized, firmly compacted, and conservative institution did it meet his rigid monarchical interest, and afford the splendid state and court dress he wished for his empire. So early as the year 313 we find the bishop Hosius of Cordova among his counsellors, and heathen writers ascribe to the bishop even a magical influence over the emperor. Lactantius, also, and Eusebius of Caesarea belonged to his confidential circle. He exempted the Christian clergy from military and municipal duty (March, 313); abolished various customs and ordinances offensive to the Christians (315); facilitated the emancipation of Christian slaves (before 316); legalized bequests to catholic churches (321); enjoined the civil observance of Sunday, though not as dies Domini, but as dies Solis, in conformity to his worship of Apollo, and in company with an ordinance for the regular consulting of the haruspex (321); contributed liberally to the building of churches and the support of the clergy; erased the heathen symbols of Jupiter and Apollo, Mars and Hercules from the imperial coins (323); and gave his sons a Christian education.

This mighty example was followed, as might be expected, by a general transition of those subjects, who were more influenced in their conduct by outward circumstances, than by inward conviction and principle. The story, that in one year (324) twelve thousand men, with women and children in proportion, were baptized in Rome, and that the emperor had promised to each convert a white garment and twenty pieces of gold, is at least in accordance with the spirit of that reign, though the fact itself, in all probability, is greatly exaggerated.\footnote{For the Acta St. Silvestri and the H. Eccl. of Nicephorus Callist. vii. 34 (in Baronius, ad ann. 324) are of course not reliable authority on this point.}

Constantine came out with still greater decision, when, by his victory over his Eastern colleague and brother-in-law, Licinius, he became sole head of the whole Roman empire. To strengthen his position, Licinius had gradually placed himself at the head of the heathen party, still very numerous, and had vexed the Christians first with wanton ridicule\footnote{He commanded the Christians, for example, to hold their large assemblies in open fields instead of in the churches, because the fresh air was more wholesome for them than the close atmosphere in a building!} then with exclusion from civil and military office, with banishment, and in some instances perhaps even with bloody persecution. This gave the political strife for the monarchy between himself and Constantine the character also of a war of religions; and the defeat of Licinius in the battle of Adrianople in July, 324, and at Chalcedon in September, was a new triumph of the standard of the cross over the sacrifices of the gods; save that Constantine dishonored himself and his cause by the execution of Licinius and his son.

The emperor now issued a general exhortation to his subjects to embrace the Christian religion, still leaving them, however, to their own free conviction. In the year 325, as patron of the church, he summoned the council of Nice, and himself attended it; banished the Arians, though he afterwards recalled them; and, in his monarchical spirit of uniformity, showed great zeal for the settlement of all theological disputes, while he was blind to their deep significance. He first introduced the practice of subscription to the articles of a written creed and of the infliction of civil punishments for non-conformity. In the years 325–329, in connection with his mother, Helena, he erected magnificent churches on the sacred spots in Jerusalem.

As heathenism had still the preponderance in Rome, where it was hallowed by its great traditions, Constantine, by divine command as he supposed,\footnote{“Jubente Deo,” says he in one of his laws. Cod. Theodos. l. xiii. tit. v. leg. 7. Later writers ascribe the founding of Constantinople to a nocturnal vision of the emperor, and an injunction of the Virgin Mary\emph{,} who was revered as patroness, one might almost suppose as goddess, of the city.} in the year 330, transferred the seat of his government to Byzantium, and thus fixed the policy, already initiated by Domitian, of orientalizing and dividing the empire. In the selection of the unrivalled locality he showed more taste and genius than the founders of Madrid, Vienna, Berlin, St. Petersburg, or Washington. With incredible rapidity, and by all the means within reach of an absolute monarch, he turned this nobly situated town, connecting two seas and two continents, into a splendid residence and a new Christian Rome, “for which now,” as Gregory of Nazianzen expresses it, “sea and land emulate each other, to load it with their treasures, and crown it queen of cities.”\footnote{The Turks still call it emphatically \emph{the city.} For Stambul is a corruption of Istambul, which means: \textgreek{εἰς τὴν πόλιν}.} Here, instead of idol temples and altars, churches and crucifixes rose; though among them the statues of patron deities from all over Greece, mutilated by all sorts of tasteless adaptations, were also gathered in the new metropolis.\footnote{The most offensive of these is the colossal bronze statue of Apollo, pretended to be the work of Phidias, which Constantine set up in the middle of the Forum on a pillar of porphyry, a hundred and twenty feet high, and which, at least according to later interpretations, served to represent the emperor himself with the attributes of Christ and the god of the sun! So says the author of Antiquit. Constant. in Banduri, and J. v. Hammer: Constantinopolis u. der Bosphorus, i. 162 (cited in Milman’s notes to Gibbon). Nothing now remains of the pillar but a mutilated piece.} The main hall in the palace was adorned with representations of the crucifixion and other biblical scenes. The gladiatorial shows, so popular in Rome, were forbidden here, though theatres, amphitheatres, and hippodromes kept their place. It could nowhere be mistaken, that the new imperial residence was as to all outward appearance a Christian city. The smoke of heathen sacrifices never rose from the seven hills of New Rome except during the short reign of Julian the Apostate. It became the residence of a bishop who not only claimed the authority of the apostolic see of neighboring Ephesus, but soon outshone the patriarchate of Alexandria and rivalled for centuries the papal power in ancient Rome.

The emperor diligently attended divine worship, and is portrayed upon medals in the posture of prayer. He kept the Easter vigils with great devotion. He would stand during the longest sermons of his bishops, who always surrounded him, and unfortunately flattered him only too much. And he even himself composed and delivered discourses to his court, in the Latin language, from which they were translated into Greek by interpreters appointed for the purpose.\footnote{Euseb. V. C. iv. 29-33. Burckhardt, l.c. p. 400, gives little credit to this whole account of Eusebius, and thus intimates the charge of deliberate falsehood.} General invitations were issued, and the citizens flocked in great crowds to the palace to hear the imperial preacher, who would in vain try to prevent their loud applause by pointing to heaven as the source of his wisdom. He dwelt mainly on the truth of Christianity, the folly of idolatry, the unity and providence of God, the coming of Christ, and the judgment. At times he would severely rebuke the avarice and rapacity of his courtiers, who would loudly applaud him with their mouths, and belie his exhortation by their works.\footnote{Euseb. Vit. Const. iv. 29 ad finem.} One of these productions is still extant,\footnote{Const. \emph{Oratio ad Sanctorum coetum}, was preserved in Greek translation by Eusebius as an appendix to his biography of the emperor.} in which he recommends Christianity in a characteristic strain, and in proof of its divine origin cites especially the fulfilment of prophecy, including the Sibylline books and the Fourth Eclogue of Virgil, with the contrast between his own happy and brilliant reign and the tragical fate of his persecuting predecessors and colleagues.

Nevertheless he continued in his later years true upon the whole to the toleration principles of the edict of 313, protected the pagan priests and temples in their privileges, and wisely abstained from all violent measures against heathenism, in the persuasion that it would in time die out. He retained many heathens at court and in public office, although he loved to promote Christians to honorable positions. In several cases, however, he prohibited idolatry, where it sanctioned scandalous immorality, as in the obscene worship of Venus in Phenicia; or in places which were specially sacred to the Christians, as the sepulchre of Christ and the grove of Mamre; and he caused a number of deserted temples and images to be destroyed or turned into Christian churches. Eusebius relates several such instances with evident approbation, and praises also his later edicts against various heretics and schismatics, but without mentioning the Arians. In his later years he seems, indeed, to have issued a general prohibition of idolatrous sacrifice; Eusebius speaks of it, and his sons in 341 refer to an edict to that effect; but the repetition of it by his successors proves, that, if issued, it was not carried into general execution under his reign.

With this shrewd, cautious, and moderate policy of Constantine, which contrasts well with the violent fanaticism of his sons, accords the postponement of his own baptism to his last sickness.\footnote{The pretended baptism of Constantine by the Roman bishop Sylvester in 324, and his bestowment of lands on the pope in connection with it, is a mediaeval fiction, still unblushingly defended indeed by Baronius (ad ann. 324, No. 43-49), but long since given up by other Roman Catholic historians, such as Noris, Tillemont, and Valesius. It is sufficiently refuted by the contemporary testimony of Eusebius alone (Vit. Const. iv. 61, 62), who places the baptism of Constantineat the end of his life, and minutely describes it; and Socrates, Sozomen, Ambrose, and Jerome coincide with him.} For this he had the further motives of a superstitious desire, which he himself expresses, to be baptized in the Jordan, whose waters had been sanctified by the Saviour’s baptism, and no doubt also a fear, that he might by relapse forfeit the sacramental remission of sins. He wished to secure all the benefit of baptism as a complete expiation of past sins, with as little risk as possible, and thus to make the best of both worlds. Deathbed baptisms then were to half Christians of that age what deathbed conversions and deathbed communions are now. Yet he presumed to preach the gospel, he called himself the bishop of bishops, he convened the first general council, and made Christianity the religion of the empire, long before his baptism! Strange as this inconsistency appears to us, what shall we think of the court bishops who, from false prudence, relaxed in his favor the otherwise strict discipline of the church, and admitted him, at least tacitly, to the enjoyment of nearly all the privileges of believers, before he had taken upon himself even a single obligation of a catechumen!

When, after a life of almost uninterrupted health, he felt the approach of death, he was received into the number of catechumens by laying on of hands, and then formally admitted by baptism into the full communion of the church in the year 337, the sixty-fifth year of his age, by the Arian (or properly Semi-Arian) bishop Eusebius of Nicomedia, whom he had shortly before recalled from exile together with Arius.\footnote{Hence Jerome says, Constantine was baptized into Arianism. And Dr. Newman, the ex-Tractarian, remarks, that in conferring his benefaction on the church he burdened it with the bequest of an heresy, which outlived his age by many centuries, and still exists in its effects in the divisions of the East (The Arians of the 4th Century, 1854, p. 138). But Eusebius (not the church historian) was probably the nearest bishop, and acted here not as a party leader. Constantine, too, in spite of the influence which the Arians had over him in his later years, considered himself constantly a true adherent of the Nicene faith, and he is reported by Theodoret (H. E. I. 32) to have ordered the recall of Athanasius from exile on his deathbed, in spite of the opposition of the Arian Eusebius. He was in these matters frequently misled by misrepresentations, and cared more for peace than for truth. The deeper significance of the dogmatic controversy was entirely beyond his sphere. Gibbon is right in this matter: “The credulous monarch, unskilled in the stratagems of theological warfare, might be deceived by the modest and specious professions of the heretics, whose sentiments he never perfectly understood; and while he protected Arius, and persecuted Athanasius, he still considered the council of Nice as the bulwark of the Christian faith, and the peculiar glory of his own reign.” Ch. xxi.} His dying testimony then was, as to form, in favor of heretical rather than orthodox Christianity, but merely from accident, not from intention. He meant the Christian as against the heathen religion, and whatever of Arianism may have polluted his baptism, was for the Greek church fully wiped out by the orthodox canonization. After the solemn ceremony he promised to live thenceforth worthily of a disciple of Jesus; refused to wear again the imperial mantle of cunningly woven silk richly ornamented with gold; retained the white baptismal robe; and died a few days after, on Pentecost, May 22, 337, trusting in the mercy of God, and leaving a long, a fortunate, and a brilliant reign, such as none but Augustus, of all his predecessors, had enjoyed. “So passed away the first Christian Emperor, the first Defender of the Faith, the first Imperial patron of the Papal see, and of the whole Eastern Church, the first founder of the Holy Places, Pagan and Christian, orthodox and heretical, liberal and fanatical, not to be imitated or admired, but much to be remembered, and deeply to be studied.”\footnote{Stanley, l.c. p. 320.}

His remains were removed in a golden coffin by a procession of distinguished civilians and the whole army, from Nicomedia to Constantinople, and deposited, with the highest Christian honors, in the church of the Apostles,\footnote{This church became the burial place of the Byzantine emperors, till in the fourth crusade the coffins were rifled and the bodies cast out. Mahomet II. destroyed the church and built in its place the magnificent mosque which bears his name. See von Hammer, i. 390.} while the Roman senate, after its ancient custom, proudly ignoring the great religious revolution of the age, enrolled him among the gods of the heathen Olympus. Soon after his death, Eusebius set him above the greatest princes of all times; from the fifth century he began to be recognized in the East as a saint; and the Greek and Russian church to this day celebrates his memory under the extravagant title of “Isapostolos,” the “Equal of the apostles.”\footnote{Comp the Acta Sact. ad 21 Maii, p. 13 sq. Niebuhr justly remarks: “When certain oriental writers call Constantine“ equal to the Apostles,’ they do not know what they are saying; and to speak of him as a ’saint’ is a profanation of the word.”} The Latin church, on the contrary, with truer tact, has never placed him among the saints, but has been content with naming him “the Great,” in just and grateful remembrance of his services to the cause of Christianity and civilization.

\xsection{The Sons of Constantine. a.d. 337-361}

§ 3. The Sons of Constantine. a.d. 337–361.

For the literature see § 2 and § 4.

With the death of Constantine the monarchy also came, for the present, to an end. The empire was divided among his three sons, Constantine II., Constans, and Constantius. Their accession was not in Christian style, but after the manner of genuine Turkish, oriental despotism; it trod upon the corpses of the numerous kindred of their father, excepting two nephews, Gallus and Julian, who were saved only by sickness and youth from the fury of the soldiers. Three years later followed a war of the brothers for the sole supremacy. Constantine II. was slain by Constans (340), who was in turn murdered by a barbarian field officer and rival, Magnentius (350). After the defeat and the suicide of Magnentius, Constantius, who had hitherto reigned in the East, became sole emperor, and maintained himself through many storms until his natural death (353–361).

The sons of Constantine did their Christian education little honor, and departed from their father’s wise policy of toleration. Constantius, a temperate and chaste, but jealous, vain, and weak prince, entirely under the control of eunuchs, women, and bishops, entered upon a violent suppression of the heathen religion, pillaged and destroyed many temples, gave the booty to the church, or to his eunuchs, flatterers, and worthless favorites, and prohibited, under penalty of death, all sacrifices and worship of images in Rome, Alexandria, and Athens, though the prohibition could not be carried out. Hosts now came over to Christianity, though, of course, for the most part with the lips only, not with the heart. But this emperor proceeded with the same intolerance against the adherents of the Nicene orthodoxy, and punished them with confiscation and banishment. His brothers supported Athanasius, but he himself was a fanatical Arian. In fact, he meddled in all the affairs of the church, which was convulsed during his reign with doctrinal controversy. He summoned a multitude of councils, in Gaul, in Italy, in Illyricum, and in Asia; aspired to the renown of a theologian; and was fond of being called bishop of bishops, though, like his father, he postponed baptism till shortly before his death.

Those were there, it is true, who justified this violent suppression of idolatry, by reference to the extermination of the Canaanites under Joshua.\footnote{So Julius Firmicus Maternus, author of a tract De errore profanarum religionum, written about 348 and dedicated to the emperors Constantius and Constans.} But intelligent church teachers, like Athanasius, Hosius, and Hilary, gave their voice for toleration, though even they mean particularly toleration for orthodoxy, for the sake of which they themselves had been deposed and banished by the Arian power. Athanasius says, for example: “Satan, because there is no truth in him, breaks in with axe and sword. But the Saviour is gentle, and forces no one, to whom he comes, but knocks and speaks to the soul: Open to me, my sister?\footnote{Song of Sol. v. 2.} If we open to him, he enters; but if we will not, he departs. For the truth is not preached by sword and dungeon, by the might of an army, but by persuasion and exhortation. How can there be persuasion where fear of the emperor is uppermost? How exhortation, where the contradicter has to expect banishment and death?” With equal truth Hilary confronts the emperor with the wrong of his course, in the words: “With the gold of the state thou burdenest the sanctuary of God, and what is torn from the temples, or gained by confiscation, or extorted by punishment, thou obtrudest upon God.”

By the laws of history the forced Christianity of Constantius must provoke a reaction of heathenism. And such reaction in fact ensued, though only for a brief period immediately after this emperor’s death.

\xsection{Julian the Apostate, and the Reaction of Paganism. a.d. 361-363}

§ 4. Julian the Apostate, and the Reaction of Paganism. a.d. 361–363.

SOURCES.

These agree in all the principal facts, even to unimportant details, but differ entirely in spirit and in judgment; Julian himself exhibiting the vanity of self-praise, Libanius and Zosimus the extreme of passionate admiration, Gregory and Cyril the opposite extreme of hatred and abhorrence, Ammianus Marcellinus a mixture of praise and censure.

1. Heathen sources: Juliani imperatoris Opera, quae supersunt omnia, ed. by Petavius, Par. 1583; and more completely by Ezech. Spanhemius, Lips. 1696, 2 vols. fol. in one (Spanheim gives the Greek original with a good Latin version, and the Ten Books of Cyril of Alex. against Julian). We have from Julian: Misopogon (Misopwvgon, the Beard-hater, a defence of himself against the accusations of the Antiochians); Caesares (two satires on his predecessors); eight Orationes; sixty-five Epistolae (the latter separately and most completely edited, with shorter fragments, by Heyler, Mog. 1828); and Fragments of his three or seven Books κατὰ Χριστιανῶνin the Reply of Cyril. Libanius: Ἐπιτάφιος ἐπ  Ἰουλιανῷ, in Lib. Opp. ed. Reiske, Altenb. 1791–97. 4 vols. Mamertinus: Gratiarum actio Juliano. The relevant passages in the heathen historians Ammianus Marcellinus (I.c. lib. xxi-xxv. 3), Zosimus and Eunapius.

2. Christian Sources (all in Greek): the early church historians, Socrates (l. iii.), Sozomen (I. v. and vi.), Theodoret (I. iii.). Gregory Naz.: Orationes invectivae in Jul. duae, written some six months after the death of Julian (Opp. tom. i.). Cyril of Alex.: Contra impium Jul. libri x. (in the Opp. Cyr., ed. J. Aubert, Par. 1638, tom. vi., and in Spanheim’s ed. of the works of Julian).

LITERATURE.

Tillemont: Memoires, etc., vol. vii. p. 322–423 (Venice ed.), and Histoire des empereurs Rom. Par. 1690 sqq., vol. iv. 483–576. Abbé De la Bleterie: Vie de l’empereur Julien. Amst. 1735. 2 vols. The same in English, Lond. 1746. W. Warburton: Julian. Lond. 3d ed. 1763. Nath. Lardner: Works, ed. Dr. Kippis, vol. vii. p. 581 sqq. Gibbon: l.c. ch. xxii.–xxiv., particularly xxiii. Neander: Julian u. sein Zeitalter. Leipz. 1812 (his first historical production), and Allg. K. G., iii. (2d ed. 1846), p. 76–148. English ed. Torrey, ii. 37–67. Jondot (R.C.): Histoire de l’empereur Julien. 1817, 2 vols. C. H. Van Herwerden: De Juliano imper. religionis Christ. hoste, eodemque vindice. Lugd. Bat. 1827. G. F. Wiggers: Jul. der Abtrünnige. Leipz. 1837 (in Illgen’s Zeitschr. f. Hist. Theol.). H. Schulze: De philos. et moribus Jul. Strals. 1839. D. Fr. Strauss (author of the mythological “Leben Jesu”): Der Romantiker auf dem Thron der Caesaren, oder Julian der Abtr. Manh. 1847 (containing a clear survey of the various opinions concerning Julian from Libanius and Gregory to Gibbon, Schlosser, Neander, and Ullmann, but hiding a political aim against King Frederick William IV. of Prussia). J. E. Auer (R.C.): Kaiser Jul. der Abtr. im Kampf mit den Kirchenvaetern seiner Zeit. Wien, 1855. W. Mangold: Jul. der Abtr. Stuttg. 1862. C. Semisch: Jul. der Abtr. Bresl. 1862. F. Lübker: Julians Kampf u. Ende. Hamb. 1864.

Notwithstanding this great conversion of the government and of public sentiment, the pagan religion still had many adherents, and retained an important influence through habit and superstition over the rude peasantry, and through literature and learned schools of philosophy and rhetoric at Alexandria, Athens, \&c., over the educated classes. And now, under the lead of one of the most talented, energetic, and notable Roman emperors, it once more made a systematic and vigorous effort to recover its ascendency in the Roman empire. But in the entire failure of this effort heathenism itself gave the strongest proof that it had outlived itself forever. It now became evident during the brief, but interesting and instructive episode of Julian’s reign, that the policy of Constantine was entirely judicious and consistent with the course of history itself, and that Christianity really carried all the moral vigor of the present and all the hopes of the future. At the same time this temporary persecution was a just punishment and wholesome discipline for a secularized church and clergy.\footnote{So Gregory of Naz. regarded it, and Tillemont justly remarks, Mem. vii. 322: “Le grand nombre de pechez dont beaucoup de Chrétiens estoient coupables, fut cause que Dieu donna a ce prince la puissance imperials pour les punir; et sa malice fut comme une verge entre les mains de Dieu pour les corriger.”}

Julian, surnamed the Apostate (Apostata), a nephew of Constantine the Great and cousin of Constantius, was born in the year 331, and was therefore only six years old when his uncle died. The general slaughter of his kindred, not excepting his father, at the change of the throne, could beget neither love for Constantius nor respect for his court Christianity. He afterwards ascribed his escape to the special favor of the old gods. He was systematically spoiled by false education and made the enemy of that very religion which pedantic teachers attempted to force upon his free and independent mind, and which they so poorly recommended by their lives. We have a striking parallel in more recent history in the case of Frederick the Great of Prussia. Julian was jealously watched by the emperor, and kept in rural retirement almost like a prisoner. With his step-brother Gallus, he received a nominally Christian training under the direction of the Arian bishop Eusebius of Nicomedia and several eunuchs; he was baptized; even educated for the clerical order, and ordained a Lector.\footnote{Jul. ad Athen. p. 271; Socr. iii. 1; Sozom. v. 2; Theod. iii. 2.} He prayed, fasted, celebrated the memory of the martyrs, paid the usual reverence to the bishops, besought the blessing of hermits, and read the Scriptures in the church of Nicomedia. Even his plays must wear the hue of devotion. But this despotic and mechanical force-work of a repulsively austere and fiercely polemic type of Christianity roused the intelligent, wakeful, and vigorous spirit of Julian to rebellion, and drove him over towards the heathen side. The Arian pseudo-Christianity of Constantius produced the heathen anti-Christianity of Julian; and the latter was a well-deserved punishment of the former. With enthusiasm and with untiring diligence the young prince studied Homer, Plato, Aristotle, and the Neo-Platonists. The partial prohibition of such reading gave it double zest. He secretly obtained the lectures of the celebrated rhetorician Libanius, afterwards his eulogist, whose productions, however, represent the degeneracy of the heathen literature in that day, covering emptiness with a pompous and tawdry style, attractive only to a vitiated taste. He became acquainted by degrees with the most eminent representatives of heathenism, particularly the Neo-Platonic philosophers, rhetoricians, and priests, like Libanius, Aedesius, Maximus, and Chrysanthius. These confirmed him in his superstitions by sophistries and sorceries of every kind. He gradually became the secret head of the heathen party. Through the favor and mediation of the empress Eusebia he visited for some months the schools of Athens (a.d. 355), where he was initiated in the Eleusinian mysteries, and thus completed his transition to the Grecian idolatry.

This heathenism, however, was not a simple, spontaneous growth; it was all an artificial and morbid production. It was the heathenism of the Neo-Platonic, pantheistic eclecticism, a strange mixture of philosophy, poesy, and superstition, and, in Julian at least, in great part an imitation or caricature of Christianity. It sought to spiritualize and revive the old mythology by uniting with it oriental theosophemes and a few Christian ideas; taught a higher, abstract unity above the multiplicity of the national gods, genii, heroes, and natural powers; believed in immediate communications and revelations of the gods through dreams, visions, oracles, entrails of sacrifices, prodigies; and stood in league with all kinds of magical and theurgic arts.\footnote{Comp. vol. i. § 61.} Julian himself, with all his philosophical intelligence, credited the most insipid legends of the gods, or gave them a deeper, mystic meaning by the most arbitrary allegorical interpretation. He was in intimate personal intercourse with Jupiter, Minerva, Apollo, Hercules, who paid their nocturnal visits to his heated fancy, and assured him of their special protection. And he practised the art of divination as a master.\footnote{Libanius says of him, Epit. p. 582: ... \textgreek{μαντέων τε τοῖς αρίστοις χρώμενος, αὐτός τε ὤν οὐδαμῶν ἐν τῇ τέχνῃ δεύτερος}. Ammanius Marcellinus calls him, xxv. 4, praesagiorum sciscitationi nimiae deditus, superstitiosus magis quam sacrorum legitimus observator. Comp. Sozom. v. 2.} Among the various divinities he worshipped with peculiar devotion the great king Helios, or the god of the sun, whose servant he called himself, and whose ethereal light attracted him even in tender childhood with magic force. He regarded him as the centre of the universe, from which light, life, and salvation proceed upon all creatures.\footnote{Comp. his fourth Oratio, which is devoted to the praise of Helios.} In this view of a supreme divinity he made an approach to the Christian monotheism, but substituted an airy myth and pantheistic fancy for the only true and living God and the personal historical Christ.

His moral character corresponds with the preposterous nature of this system. With all his brilliant talents and stoical virtues, he wanted the genuine simplicity and naturalness, which are the foundation of all true greatness of mind and character. As his worship of Helios was a shadowy reflection of the Christian monotheism, and so far an involuntary tribute to the religion he opposed, so in his artificial and ostentatious asceticism we can only see a caricature of the ecclesiastical monasticism of the age which he so deeply despised for its humility and spirituality. He was full of affectation, vanity, sophistry, loquacity, and a master in the art of dissimulation. Everything he said or wrote was studied and calculated for effect. Instead of discerning the spirit of the age and putting himself at the head of the current of true progress, he identified himself with a party of no vigor nor promise, and thus fell into a false and untenable position, at variance with the mission of a ruler. Great minds, indeed, are always more or less at war with their age, as we may see in the reformers, in the apostles, nay, in Christ himself. But their antagonism proceeds from a clear knowledge of the real wants and a sincere devotion to the best interests of the age; it is all progressive and reformatory, and at last carries the deeper spirit of the age with itself, and raises it to a higher level. The antagonism of Julian, starting with a radical misconception of the tendency of history and animated by selfish ambition, was one of retrogression and reaction, and in addition, was devoted to a bad cause. He had all the faults, and therefore deserved the tragic fate, of a fanatical reactionist.

His apostasy from Christianity, to which he was probably never at heart committed, Julian himself dates as early as his twentieth year, a.d. 351. But while Constantius lived, he concealed his pagan sympathies with consummate hypocrisy, publicly observed Christian ceremonies, while secretly sacrificing to Jupiter and Helios, kept the feast of Epiphany in the church at Vienne so late as January, 361, and praised the emperor in the most extravagant style, though he thoroughly hated him, and after his death all the more bitterly mocked him.\footnote{Comp. Jul. Orat. i. in Constantii laudes; Epist. ad Athenienses, p. 270; Caesares, p. 335 sq. Even heathen authors concede his dissimulation, as Ammianus Marc. xxi. 2, comp. xxii. 5, and Libanius, who excuses him with the plea of regard to his security, Opp. p. 528, ed. Reiske.} For ten years he kept the mask. After December, 355, the student of books astonished the world with brilliant military and executive powers as Caesar in Gaul, which was at that time heavily threatened by the German barbarians; he won the enthusiastic love of the soldiers, and received from them the dignity of Augustus. Then he raised the standard of rebellion against his suspicious and envious imperial cousin and brother-in-law, and in 361 openly declared himself a friend of the gods. By the sudden death of Constantius in the same year he became sole head of the Roman empire, and in December, as the only remaining heir of the house of Constantine,\footnote{His older brother, Gallus, for some time emperor at Antioch, had already been justly deposed by Constantius in 854, and beheaded, for his entire incapacity and his merciless cruelty.} made his entry into Constantinople amidst universal applause and rejoicing over escape from civil war.

He immediately gave himself, with the utmost zeal, to the duties of his high station, unweariedly active as prince, general, judge, orator, high-priest, correspondent, and author. He sought to unite the fame of an Alexander, a Marcus Aurelius, a Plato, and a Diogenes in himself. His only recreation was a change of labor. He would use at once his hand in writing, his ear in hearing, and his voice in speaking. He considered his whole time due to his empire and the culture of his own mind. The eighteen short months of his reign Dec. 361-June 363) comprehend the plans of a life-long administration and most of his literary works. He practised the strictest economy in the public affairs, banished all useless luxury from his court, and dismissed with one decree whole hosts of barbers, cup-bearers, cooks, masters of ceremonies, and other superfluous officers, with whom the palace swarmed, but surrounded himself instead with equally useless pagan mystics, sophists, jugglers, theurgists, soothsayers, babblers, and scoffers, who now streamed from all quarters to the court. In striking contrast with his predecessors, he maintained the simplicity of a philosopher and an ascetic in his manner of life, and gratified his pride and vanity with contempt of the pomp and pleasures of the imperial purple. He lived chiefly on vegetable diet, abstaining now from this food, now from that, according to the taste of the god or goddess to whom the day was consecrated. He wore common clothing, usually slept on the floor, let his beard and nails grow, and, like the strict anachorets of Egypt, neglected the laws of decency and cleanliness.\footnote{In the Misopogon (from \textgreek{μισέω} and \textgreek{πώγων}, the beard-hater, i.e. hater of bearded philosophers), his witty apology to the refined Antiochians for his philosophical beard, p. 338 sq., he boasts of this cynic coarseness, and describes, with great complacence, his long nails, his ink-stained hands, his rough, uncombed beard, inhabited (horribile dictu) by certain \textgreek{θηρία.} It should not be forgotten, however, that contemporary writers give him the credit of a strict chastity, which raises him far above most heathen princes, and which furnishes another proof to the involuntary influence of Christian asceticism upon his life. Libanius asserts in his panegyric, that Julian, before his brief married life, and after the death of his wife, a sister of Constantius, never knew a woman; and Namertinus calls his lectulus, “Vestalium toris purior.” Add to this the testimony of the honest Ammianus Marcellinus, and the silence of Christian antagonists. Comp. Gibbon, c. xxii. note 50; and Carwithen and Lyall: Hist. of the Chr. Ch., etc. p. 54. On the other hand, the Christians accused him of all sorts of secret crimes; for instance, the butchering of boys and girls (Gregor. Orat. iii. p. 91, and Theodor. iii. 26, 27), which was probably an unfounded inference from his fanatical zeal for bloody sacrifices and divinations.} This cynic eccentricity and vain ostentation certainly spoiled his reputation for simplicity and self-denial, and made him ridiculous. It evinced, also, not so much the boldness and wisdom of a reformer, as the pedantry and folly of a reactionist. In military and executive talent and personal bravery he was not inferior to Constantine; while in mind and literary culture he far excelled him, as well as in energy and moral self-control; and, doubtless to his own credit, he closed his public career at the age at which his uncle’s began; but he entirely lacked the clear, sound common sense of his great predecessor, and that practical statesmanship, which discerns the wants of the age, and acts according to them. He had more uncommon sense than common sense, and the latter is often even more important than the former, and indispensable to a good practical statesman. But his greatest fault as a ruler was his utterly false position towards the paramount question of his time: that of religion. This was the cause of that complete failure which made his reign as trackless as a meteor.

The ruling passion of Julian, and the soul of his short but most active, remarkable, and in its negative results instructive reign, was fanatical love of the pagan religion and bitter hatred of the Christian, at a time when the former had already forever given up to the latter the reins of government in the world. He considered it the great mission of his life to restore the worship of the gods, and to reduce the religion of Jesus first to a contemptible sect, and at last, if possible, to utter extinction from the earth. To this he believed himself called by the gods themselves, and in this faith he was confirmed by theurgic arts, visions, and dreams. To this end all the means, which talent, zeal, and power could command, were applied; and the failure must be attributed solely to the intrinsic folly and impracticability of the end itself.

I. To look, first, at the positive side of his plan, the restoration and reformation of heathenism:

He reinstated, in its ancient splendor, the worship of the gods at the public expense; called forth hosts of priests from concealment; conferred upon them all their former privileges, and showed them every honor; enjoined upon the soldiers and civil officers attendance at the forsaken temples and altars; forgot no god or goddess, though himself specially devoted to the worship of Apollo, or the sun; and notwithstanding his parsimony in other respects, caused the rarest birds and whole herds of bulls and lambs to be sacrificed, until the continuance of the species became a subject of concern.\footnote{Ammianus Marc. xxv. 4 ... innumeras sine parsimonia pecudes mactans ut aestemaretur, si revertisset de Parthis, boves jam defuturos.} He removed the cross and the monogram of Christ from the coins and standards, and replaced the former pagan symbols. He surrounded the statues and portraits of the emperors with the signs of idolatry, that every one might be compelled to bow before the gods, who would pay the emperors due respect. He advocated images of the gods on the same grounds on which afterwards the Christian iconolaters defended the images of the saints. If you love the emperor, if you love your father, says he, you like to see his portrait; so the friend of the gods loves to look upon their images, by which he is pervaded with reverence for the invisible gods, who are looking down upon him.

Julian led the way himself with a complete example. He discovered on every occasion the utmost zeal for the heathen religion, and performed, with the most scrupulous devotion, the offices of a pontifex maximus, which had been altogether neglected, although not formally abolished, under his two predecessors. Every morning and evening he sacrificed to the rising and setting sun, or the supreme light-god; every night, to the moon and the stars; every day, to some other divinity. Says Libanius, his heathen admirer: “He received the rising sun with blood, and attended him again with blood at his setting.” As he could not go abroad so often as he would, he turned his palace into a temple and erected altars in his garden, which was kept purer than most chapels. “Wherever there was a temple,” says the same writer, “whether in the city or on the hill or the mountain top, no matter how rough, or difficult of access, he ran to it.” He prostrated himself devoutly before the altars and the images, not allowing the most violent storm to prevent him. Several times in a day, surrounded by priests and dancing women, he sacrificed a hundred bulls, himself furnishing the wood and kindling the flames. He used the knife himself, and as haruspex searched with his own hand the secrets of the future in the reeking entrails.

But his zeal found no echo, and only made him ridiculous in the eyes of cultivated heathens themselves. He complains repeatedly of the indifference of his party, and accuses one of his priests of a secret league with Christian bishops. The spectators at his sacrifices came not from devotion, but from curiosity, and grieved the devout emperor by their rounds of applause, as if he were simply a theatrical actor of religion. Often there were no spectators at all. When he endeavored to restore the oracle of Apollo Daphneus in the famous cypress grove at Antioch, and arranged for a magnificent procession, with libation, dances, and incense, he found in the temple one solitary old priest, and this priest ominously offered in sacrifice—a goose.\footnote{Misopog. p. 362 sq., where Julian himself relates this ludicrous scene, and vents his anger at the Antiochians for squandering the rich incomes of the temple upon Christianity and worldly pleasures. Dr. Baur, l.c. p. 17, justly remarks on Julian’s zeal for idolatry: “Seine ganze persönliche Erscheinung, der Mangel an innerer Haltung in seinem Benehmen gegen Heiden und Christen, die stete Unruhe und schwärmerische Aufregung, in welcher er sich befand, wenn er von Tempel zu Tempel eilte, auf allen Altären opferte und nichts unversucht liess, um den heidnischen Cultus, dessen höchstes Vorbild er selbst als Pontifex maximum sein wollte, in seinem vollen Glanz und Gepränge, mit alten seinen Ceremonien und Mysterien wieder herzustellen, macht einen Eindruck, der es kaum verkennen lässt, wie wenig er sich selbst das Unnatürliche und Erfolglose eines solchen Strebens verbergen konnte.”}

At the same time, however, Julian sought to renovate and transform heathenism by incorporating with it the morals of Christianity; vainly thinking thus to bring it back to its original purity. In this he himself unwittingly and unwillingly bore witness to the poverty of the heathen religion, and paid the highest tribute to the Christian; and the Christians for this reason not inaptly called him an “ape of Christianity.”

In the first place, he proposed to improve the irreclaimable priesthood after the model of the Christian clergy. The priests, as true mediators between the gods and men, should be constantly in the temples, should occupy themselves with holy things, should study no immoral or skeptical books of the school of Epicurus and Pyrrho, but the works of Homer, Pythagoras, Plato, Chrysippus, and Zeno; they should visit no taverns nor theatres, should pursue no dishonorable trade, should give alms, practise hospitality, live in strict chastity and temperance, wear simple clothing, but in their official functions always appear in the costliest garments and most imposing dignity. He borrowed almost every feature of the then prevalent idea of the Christian priesthood, and applied it to the polytheistic religion.\footnote{Julian’s views on the heathen priests are laid down especially in his 49th Epistle to Ursacius, the highpriest of Gaul, p. 429, and in the fragment of an oration, p. 300 sqq., ed. Spanh. Ullmann, in his work on Gregory of Nazianzen, p. 527 sqq., draws an interesting parallel between Gregory’s and Julian’s ideal of a priest.} Then, he borrowed from the constitution and worship of the church a hierarchical system of orders, and a sort of penitential discipline, with excommunication, absolution, and restoration, besides a fixed ritual embracing didactic and musical elements. Mitred priests in purple were to edify the people regularly with sermons; that is, with allegorical expositions and practical applications of tasteless and immoral mythological stories! Every temple was to have a well arranged choir, and the congregation its responses. And finally, Julian established in different provinces monasteries, nunneries, and hospitals for the sick, for orphans, and for foreigners without distinction of religion, appropriated to them considerable sums from the public treasury, and at the same time, though fruitlessly, invited voluntary contributions. He made the noteworthy concession, that the heathens did not help even their own brethren in faith; while the Jews never begged, and “the godless Galileans,” as he malignantly styled the Christians, supplied not only their own, but even the heathen poor, and thus aided the worst of causes by a good practice.

But of course all these attempts to regenerate heathenism by foreign elements were utterly futile. They were like galvanizing a decaying corpse, or grafting fresh scions on a dead trunk, sowing good seed on a rock, or pouring new wine into old bottles, bursting the bottles and wasting the wine.

II. The negative side of Julian’s plan was the suppression and final extinction of Christianity.

In this he proceeded with extraordinary sagacity. He abstained from bloody persecution, because he would not forego the credit of philosophical toleration, nor give the church the glory of a new martyrdom. A history of three centuries also had proved that violent measures were fruitless. According to Libanius it was a principle with him, that fire and sword cannot change a man’s faith, and that persecution only begets hypocrites and martyrs. Finally, he doubtless perceived that the Christians were too numerous to be assailed by a general persecution without danger of a bloody civil war. Hence he oppressed the church “gently,”\footnote{\textgreek{Ἐπιεικῶς ἐβιά ζετο}, as Gregory Nazianzen, Orat. iv., expresses it.} under show of equity and universal toleration. He persecuted not so much the Christians as Christianity, by endeavoring to draw off its confessors. He thought to gain the result of persecution without incurring the personal reproach and the public danger of persecution itself. His disappointments, however, increased his bitterness, and had he returned victorious from the Persian war, he would probably have resorted to open violence. In fact, Gregory Nazianzen and Sozomen, and some heathen writers also, tell of local persecutions in the provinces, particularly at Anthusa and Alexandria, with which the emperor is, at least indirectly, to be charged. His officials acted in those cases, not under public orders indeed, but according to the secret wish of Julian, who ignored their illegal proceedings as long as he could, and then discovered his real views by lenient censure and substantial acquittal of the offending magistrates.

He first, therefore, employed against the Christians of all parties and sects the policy of toleration, in hope of their destroying each other by internal controversies. He permitted the orthodox bishops and all other clergy, who had been banished under Constantius, to return to their dioceses, and left Arians, Apollinarians, Novatians, Macedonians, Donatists, and so on, to themselves. He affected compassion for the “poor, blind, deluded Galileans, who forsook the most glorious privilege of man, the worship of the immortal gods, and instead of them worshipped dead men and dead men’s bones.” He once even suffered himself to be insulted by a blind bishop, Maris of Chalcedon, who, when reminded by him, that the Galilean God could not restore his eyesight, answered: “I thank my God for my blindness, which spares me the painful sight of such an impious Apostate as thou.” He afterwards, however, caused the bishop to be severely punished.\footnote{Socrates: H. E. iii. 12.} So in Antioch, also, he bore with philosophic equanimity the ridicule of the Christian populace, but avenged himself on the inhabitants of the city by unsparing satire in the Misopogon. His whole bearing towards the Christians was instinct with bitter hatred and accompanied with sarcastic mockery.\footnote{Gibbon well says, ch. xxiii.: “He affected to pity the unhappy Christians, but his pity was degraded by contempt, his contempt was embittered by hatred; and the sentiments of Julian were expressed in a style of sarcastic wit, which inflicts a deep and deadly wound whenever it issues from the mouth of a sovereign.”} This betrays itself even in the contemptuous term, Galileans, which he constantly applies to them after the fashion of the Jews, and which he probably also commanded to be given them by others.\footnote{Perhaps there lay at the bottom of this also a secret fear of the name of Christ, as Warburton (p. 35) suggests; since the Neo-Platonists believed in the mysterious virtue of names.} He considered them a sect of fanatics contemptible to men and hateful to the gods, and as atheists in open war with all that was sacred and divine in the world.\footnote{\textgreek{Ἀσεβεῖς, δυσσεβεῖς, ἄθεοι}. Their religion he calls a \textgreek{μωρία} or \textgreek{ἀπόνοια}. Comp. Ep. 7 (ap. Heyler, p. 190).} He sometimes had representatives of different parties dispute in his presence, and then exclaimed: “No wild beasts are so fierce and irreconcilable as the Galilean sectarians.” When he found that toleration was rather profitable than hurtful to the church, and tended to soften the vehemence of doctrinal controversies, he proceeded, for example, to banish Athanasius, who was particularly offensive to him, from Alexandria, and even from Egypt, calling this greatest man of his age an insignificant manikin,\footnote{\textgreek{ Ἄθρωπίσκος εὐτελής}.} and reviling him with vulgar language, because through his influence many prominent heathens, especially heathen women, passed over to Christianity. His toleration, therefore, was neither that of genuine humanity, nor that of religious indifferentism, but a hypocritical mask for a fanatical love of heathenism and a bitter hatred of Christianity.

This appears in his open partiality and injustice against the Christians. His liberal patronage of heathenism was in itself an injury to Christianity. Nothing gave him greater joy than an apostasy, and he held out the temptation of splendid reward; thus himself employing the impure means of proselyting, for which he reproached the Christians. Once he even advocated conversion by violent measures. While he called heathens to all the higher offices, and, in case of their palpable disobedience, inflicted very mild punishment, if any at all, the Christians came to be everywhere disregarded, and their complaints dismissed from the tribunal with a mocking reference to their Master’s precept, to give their enemy their cloak also with their coat, and turn the other cheek to his blows.\footnote{Matt. v. 89, 40.} They were removed from military and civil office, deprived of all their former privileges, oppressed with taxes, and compelled to restore without indemnity the temple property, with all their own improvements on it, and to contribute to the support of the public idolatry. Upon occasion of a controversy between the Arians and the orthodox at Edessa, Julian confiscated the church property and distributed it among his soldiers, under the sarcastic pretence of facilitating the Christians’ entrance into the kingdom of heaven, from which, according to the doctrine of their religion (comp. Matt. xix. 23, 24), riches might exclude them.

Equally unjust and tyrannical was the law, which placed all the state schools under the direction of heathens, and prohibited the Christians teaching the sciences and the arts.\footnote{Gregory of Naz., Orat. iv., censures the emperor bitterly for forbidding the Christians what was the common property of all rational men, as if it were the exclusive possession of the Greeks. Even the heathen Ammianus Marcellinus, xxii. 10, condemns this measure: “Illud autem erat inclemens, obruendum perenni silentio, quod arcebat docere magistros rhetoricos et grammaticos, ritus Christiani cultores.” Gibbon is equally decided. Directly, Julian forbade the Christians only to teach, but indirectly also to learn, the classical literature; as they were of course unwilling to go to heathen schools.} Julian would thus deny Christian youth the advantages of education, and compel them either to sink in ignorance and barbarism, or to imbibe with the study of the classics in the heathen schools the principles of idolatry. In his view the Hellenic writings, especially the works of the poets, were not only literary, but also religious documents to which the heathens had an exclusive claim, and he regarded Christianity irreconcilable with genuine human culture. The Galileans, says he in ridicule, should content themselves with expounding Matthew and Luke in their churches, instead of profaning the glorious Greek authors. For it is preposterous and ungrateful, that they should study the writings of the classics, and yet despise the gods, whom the authors revered; since the gods were in fact the authors and guides of the minds of a Homer, a Hesiod, a Demosthenes, a Thucydides, an Isocrates, and a Lysias, and these writers consecrated their works to Mercury or the muses.\footnote{Epist. 42.} Hence he hated especially the learned church teachers, Basil, Gregory of Nazianzen, Apollinaris of Laodicea, who applied the classical culture to the refutation of heathenism and the defence of Christianity. To evade his interdict, the two Apollinaris produced with all haste Christian imitations of Homer, Pindar, Euripides, and Menander, which were considered by Sozomen equal to the originals, but soon passed into oblivion. Gregory also wrote the tragedy of “The Suffering Christ,” and several hymns, which still exist. Thus these fathers bore witness to the indispensableness of classical literature for a higher Christian education, and the church has ever since maintained the same view.\footnote{Dr. Baur (l.c. p. 42) unjustly charges the fathers with the contradiction of making use of the classics as necessary means of education, and yet of condemning heathenism as a work of Satan. But this was only the one side, which has its element of truth, especially as applied to the heathen \emph{religion}; while on the other side they acknowledged, with Justin M., Clement and Origen, the working of the divine Logos in the Hellenic philosophy and poetry preparing the way for Christianity. The indiscriminate condemnation of classical literature dates from a later period, from Gregory I.}

Julian further sought to promote his cause by literary assaults upon the Christian religion; himself writing, shortly before his death, and in the midst of his preparations for the Persian campaign, a bitter work against it, of which we shall speak more fully in a subsequent section.\footnote{See below, § 9.}

3. To the same hostile design against Christianity is to be referred the favor of Julian to its old hereditary enemy, Judaism.

The emperor, in an official document affected reverence for that ancient popular religion, and sympathy with its adherents, praised their firmness under misfortune, and condemned their oppressors. He exempted the Jews from burdensome taxation, and encouraged them even to return to the holy land and to rebuild the temple on Moriah in its original splendor. He appropriated considerable sums to this object from the public treasury, intrusted his accomplished minister Alypius with the supervision of the building, and promised, if he should return victorious from the Persian war, to honor with his own presence the solemnities of reconsecration and the restoration of the Mosaic sacrificial worship.\footnote{Jul. Epist. 25, which is addressed to the Jews, and is mentioned also by Sozomen, v. 22.}

His real purpose in this undertaking was certainly not to advance the Jewish religion; for in his work against the Christians he speaks with great contempt of the Old Testament, and ranks Moses and Solomon far below the pagan lawgivers and philosophers. His object in the rebuilding of the temple was rather, in the first place, to enhance the splendor of his reign, and thus gratify his personal vanity; and then most probably to put to shame the prophecy of Jesus respecting the destruction of the temple (which, however, was actually fulfilled three hundred years before once for all), to deprive the Christians of their most popular argument against the Jews, and to break the power of the new religion in Jerusalem.\footnote{Gibbon, ch. xxiii.: “The restoration of the Jewish temple was secretly connected with the ruin of the Christian church.”}

The Jews now poured from east and west into the holy city of their fathers, which from the time of Hadrian they had been forbidden to visit, and entered with fanatical zeal upon the great national religious work, in hope of the speedy irruption of the Messianic reign and the fulfilment of all the prophecies. Women, we are told, brought their costly ornaments, turned them into silver shovels and spades, and carried even the earth and stones of the holy spot in their silken aprons. But the united power of heathen emperor and Jewish nation was insufficient to restore a work which had been overthrown by the judgment of God. Repeated attempts at the building were utterly frustrated, as even a contemporary heathen historian of conceded credibility relates, by fiery eruptions from on subterranean vaults;\footnote{Julian himself seems to admit the failure of the work, but, more prudently, is silent as to the cause, in a fragment of an epistle or oration, p. 295, ed. Spanh., according to the usual interpretation of this passage. He here asks: \textgreek{Τί περὶ τοῦ νεὼ φύσουσι, τοῦ παρ  αὐτοῖς, τρίτον ἀνατραπέντος , ἐγειρομένου δὲ οὐδὲ νῦν}:: “What will they [i.e., the Jewish prophets] say of their own temple, which has been \emph{three times} destroyed, and is not even now restored?” “This I have said (he continues) with no wish to reproach them, for I myself, at so late a day, had intended to rebuild it for the honor of him who was worshipped there.” He probably saw in the event a sign of the divine displeasure with the religion of the Jews, or an accidental misfortune, but intended, after his return from the Persian war, to attempt the work anew. It is by no means certain, however, that the threefold destruction of the temple here spoken of refers to Julian’s own reign. He may have meant, and probably did mean, the destruction by the Assyrians and the destruction by the Romans; and as to the third destruction, it may be a mere exaggeration, or may refer to the profanation of the temple by Antiochus, or to his own reign. (Comp. Warburton and Lardner on this point.) The impartial Ammianus Marcellinus, himself a professed pagan, a friend of Julian and his companion in arms, tells us more particularly, lib. xxiii. 1, that Julian, being desirous of perpetuating the memory of his reign by some great work, resolved to rebuild at vast expense the magnificent temple at Jerusalem, and committed the conduct of this enterprise to Alypius at Antioch, and then continues: “Quum itaque rei fortiter instaret Alypius, juvaretque provinciae rector, \emph{metuendi globi flammarum} prope fundamenta crebris assultibus erumpentes fecere locum exustis aliquoties operantibus inaccessum; hocque modo clemento destinatius repellente, cessavit inceptum.” (“Alypius, therefore, set himself vigorously to the work, and was assisted by the governor of the province, when fearful balls of fire broke out near the foundations, and continued their attacks until they made the place inaccessible to the workmen, after repeated scorchings; and thus, the fierce element obstinately repelling them, he gave up his attempt.”) Michaelis, Lardner (who, however, is disposed to doubt the whole story), Gibbon, Guizot, Milman (note on Gibbon), Gieseler, and others, endeavor to explain this as a natural phenomenon, resulting from the bituminous nature of the soil and the subterranean vaults and reservoirs of the temple hill, of which Josephus and Tacitus speak. When Herod, in building the temple, wished to penetrate into the tomb of David, to obtain its treasures, fire likewise broke out and consumed the workmen, according to Joseph. Antiqu. Jud. xvi. 7, § 1. But when Titus undermined the temple, a.d.70, when Hadrian built there the Aelia Capitolina, in 135, and when Omar built a Turkish mosque in 644, no such destructive phenomena occurred as far as we know. We must therefore believe, that Providence itself, by these natural causes, prevented the rebuilding of the national sanctuary of the Jews.} and, perhaps, as Christian writers add, by a violent whirlwind, lightning, earthquake, and miraculous signs, especially a luminous cross, in the heavens,\footnote{Gregory Nazianzen, Socrates, Sozomen, Theodoret, Philostorgius, Rufinus, Ambrose, Chrysostom; all of whom regard the event as supernatural, although they differ somewhat in detail. Theodoret speaks first of a violent whirlwind, which scattered about vast quantities of lime, sand, and other building materials, and was followed by a storm of thunder and lightning; Socrates mentions fire from heaven, which melted the workmen’s tools, spades, axes, and saws; both add an earthquake, which threw up the stones of the old foundations, filled up the excavation, and, as Rufinus has it, threw down the neighboring buildings. At length a calm succeeded the commotion, and according to Gregory a luminous cross surrounded by a circle appeared in the sky, nay, crosses were impressed upon the bodies of the persons present, which were shining by night (Rufinus), and would not wash out (Socrates). Of these writers however, Gregory alone is strictly a contemporary witness, relating the event in the year of its occurrence, 363, and that with the assurance that even the heathens did not call it in question. (Orat. iv. p. 110-113). Next to him come Ambrose, and Chrysostom, who speaks of this event several times. The Greek and Roman church historians, and Warburton, Mosheim, Schröckh, Neander, Guericke, Kurtz, Newman, Robertson, and others, of the Protestant, vindicate the miraculous, or at least providential, character of the remarkable event. Comp. also J. H. Newman (since gone over to Romanism): “Essay on the Miracles recorded in ecclesiastical history,” prefixed to the Oxford Tractarian translation of Fleury’s Eccles. Hist. from 381-400 (Oxford, 1842) I. p. clxxv.–clxxxv. Warburton and Newman defend even the crosses, and refer to similar cases, for instance one in England in 1610, where marks of a cross of a phosphoric nature and resembling meteoric phenomena appeared in connection with lightning and produced by electricity. In Julian’s case they assumed that the immediate cause which set all these various physical agents in motion, as in the case of the destruction of Sodom, was supernatural.} so that the workmen either perished in the flames, or fled from the devoted spot in terror and despair. Thus, instead of depriving the Christians of a support of their faith, Julian only furnished them a new argument in the ruins of this fruitless labor.

The providential frustration of this project is a symbol of the whole reign of Julian, which soon afterward sank into an early grave. As Caesar he had conquered the barbarian enemies of the Roman empire in the West; and now he proposed, as ruler of the world, to humble its enemies in the East, and by the conquest of Persia to win the renown of a second Alexander. He proudly rejected all proposals of peace; crossed the Tigris at the head of an army of sixty-five thousand men, after wintering in Antioch, and after solemn consultation of the oracle; took several fortified towns in Mesopotamia; exposed himself to every hardship and peril of war; restored at the same time, wherever he could, the worship of the heathen gods; but brought the army into a most critical position, and, in an unimportant nocturnal skirmish, received from a hostile arrow a mortal wound. He died soon after, on the 27th of June, 363, in the thirty-second year of his life; according to heathen testimony, in the proud repose and dignity of a Stoic philosopher, conversing of the glory of the soul (the immortality of which, however, he considered at best an uncertain opinion);\footnote{Ammianus, l. xxv. 3. He was himself in the campaign, and served in the body guard of the emperor; thus having the best opportunity for observation.} but according to later and somewhat doubtful Christian accounts, with the hopeless exclamation: “Galilean, thou hast conquered!”\footnote{Sozomen, vi. 2; Theodoret, iii. 25 (\textgreek{Νενίκηκας Γαλιλαῖε} ); then, somewhat differing, Philostorgius, vii. 15. Gregory Nazianzen, on the contrary, who elsewhere presents Julian in the worst light, knows nothing of this exclamation, to which one may apply the Italian maxim: “Se non è vero, è ben trovato.” The above-named historians mention also other incidents of the death, not very credible; e.g. that he threw toward heaven a handful of blood from his wound; that he blasphemed the heathen gods; that Christ appeared to him, \&c. Sozomen quotes also the groundless assertion of Libanius, that the mortal wound was inflicted not by a Persian, but by a Christian, and was not ashamed to add, that he can hardly be blamed who had done this ” noble deed for God and his religion” (\textgreek{διὰ θεὸν καὶ θρησκείαν ἣν ἐπῄνεσεν})! This is, so far as I know, the first instance, within the Christian church, of the vindication of tyrannicide \emph{ad majorem Dei gloriam}.} The parting address to his friends, which Ammianus puts into his mouth, is altogether characteristic. It reminds one of the last hours of Socrates, without the natural simplicity of the original, and with a strong admixture of self-complacence and theatrical affectation. His body was taken, at his own direction, to Tarsus, the birthplace of the apostle Paul, whom he hated more than any other apostle, and a monument was erected to him there, with a simple inscription, which calls him a good ruler and a brave warrior, but says nothing of his religion.

So died, in the prime of life, a prince, who darkened his brilliant military, executive, and literary talents, and a rare energy, by fanatical zeal for a false religion and opposition to the true; perverted them to a useless and wicked end; and earned, instead of immortal honor, the shame of an unsuccessful Apostate. Had he lived longer, he would probably have plunged the empire into the sad distraction of a religious civil war. The Christians were generally expecting a bloody persecution in case of his successful return from the Persian war. We need, therefore, the less wonder that they abhorred his memory. At Antioch they celebrated his death by festal dancings in the churches and theatres.\footnote{Theodor. H. E. iii. 27.} Even the celebrated divine and orator, Gregory Nazianzen, compared him to Pharaoh, Ahab, and Nebuchadnezzar.\footnote{The Christian poet, Prudentius, forms an exception, in his well known just estimate of Julian(Apotheos. 450 sqq.), which Gibbon also cites: \par ——“Ductor fortissimus armis; \par Conditor et legum celeberrimus; ore manuque \par Consultor patriae; sed non consultor habendae \par Religionis; amans tercentûm millia Divûm. \par Perfidus ille Deo, sed non et perfidus orbi.”} It has been reserved for the more impartial historiography of modern times to do justice to his nobler qualities, and to endeavor to excuse, or at least to account for his utterly false position toward Christianity, by his perverted education, the despotism of his predecessor, and the imperfections of the church in his day.

With Julian himself fell also his artificial, galvanized heathenism, “like the baseless fabric of a vision, leaving no wreck behind,” save the great doctrine, that it is impossible to swim against the stream of history or to stop the progress of Christianity. The heathen philosophers and soothsayers, who had basked in his favor, fell back into obscurity. In the dispersion of their dream they found no comfort from their superstition. Libanius charges the guilt upon his own gods, who suffered Constantius to reign twenty years, and Julian hardly twenty months. But the Christians could learn from it, what Gregory Nazianzen had said in the beginning of this reign, that the church had far more to fear from enemies within, than from without.

\xsection{From Jovian to Theodosius. a.d. 363-392}

§ 5. From Jovian to Theodosius. a.d. 363–392.

I. The heathen sources here, besides Ammianus Marcellinus (who unfortunately breaks off at the death of Valens), Zosimus and Eunapius (who are very partial), are: Libanius: Ὑπὲρ τῶν ἱερῶν, or Oratio pro templis (first complete ed. by L. de Sinner, in Novus Patrum Grace. saec. iv. delectus, Par. 1842). Symmachus: Epist. x. 61 (ed. Pareus, Frcf. 1642). On the Christian side: Ambrose: Epist. xvii. and xviii. ad Valentinian. II. Prudentius: Adv. Symmachum. Augustin: De civitate Dei, l. v. c. 24–26 (on the emperors from Jovinian to Theodosius, especially the latter, whom he greatly glorifies). Socr.: l. iii. c. 22 sqq. Sozom.: l. vi. c. 3 sqq. Theodor.: l. iv. c. 1 sqq. Cod. Theodos.: l. ix.–xvi.

II. De la Bleterie: Histoire de l’empereur Jovien. Amsterd. 1740, 2 vols. Gibbon: chap. xxv–xxviii. Schröckh: vii. p. 213 sqq. Stuffken: De Theodosii M. in rem christianam meritis. Lugd. Batav. 1828

From this time heathenism approached, with slow but steady step, its inevitable dissolution, until it found an inglorious grave amid the storms of the great migration and the ruins of the empire of the Caesars, and in its death proclaimed the victory of Christianity. Emperors, bishops, and monks committed indeed manifold injustice in destroying temples and confiscating property; but that injustice was nothing compared with the bloody persecution of Christianity for three hundred years. The heathenism of ancient Greece and Rome died of internal decay, which no human power could prevent.

After Julian, the succession of Christian emperors continued unbroken. On the day of his death, which was also the extinction of the Constantinian family, the general Jovian, a Christian (363–364), was chosen emperor by the army. He concluded with the Persians a disadvantageous but necessary peace, replaced the cross in the labarum, and restored to the church her privileges, but, beyond this, declared universal toleration in the spirit of Constantine. Under the circumstances, this was plainly the wisest policy. Like Constantine, also, he abstained from all interference with the internal affairs of the church, though for himself holding the Nicene faith and warmly favorable to Athanasius. He died in the thirty-third year of his age, after a brief reign of eight months. Augustin says, God took him away sooner than Julian, that no emperor might become a Christian for the sake of Constantine’s good fortune, but only for the sake of eternal life.

His successor, Valentinian I. (died 375), though generally inclined to despotic measures, declared likewise for the policy of religious freedom,\footnote{Cod. Theodos. l. ix. tit. 16, I. 9 (of the year 371): Testes sunt leges a me in exordio imperii mei datae, quibus \emph{unicuique}, \emph{quod animo imbibisset, colendi libera facultas} tributa est. This is confirmed by Ammian. Marc. l. xxx. c. 9.} and, though personally an adherent of the Nicene orthodoxy, kept aloof from the doctrinal controversies; while his brother and co-emperor, Valens, who reigned in the East till 378, favored the Arians and persecuted the Catholics. Both, however, prohibited bloody sacrifices\footnote{Libanius, l.c. (ed. Reiske, ii. 163): \textgreek{τὸ θύειν ἱερεῖα—ἐκωλύθη παρὰ τοῖν ἀδελφοιν, ἀλλ  οὐ τὸ λιανωτόν}. No such law, however, has come down to us.} and divination. Maximin, the representative of Valentinian at Rome, proceeded with savage cruelty against all who were found guilty of the crime of magic, especially the Roman aristocracy. Soothsayers were burnt alive, while their meaner accomplices were beaten to death by straps loaded with lead. In almost every case recorded the magical arts can be traced to pagan religious usages.

Under this reign heathenism was for the first time officially designated as paganismus, that is, peasant-religion; because it had almost entirely died out in the cities, and maintained only a decrepit and obscure existence in retired villages.\footnote{The word \emph{pagani} (from pagus), properly villagers, peasantry, then equivalent to rude, simple, ignorant, \textgreek{ἰδιώτης, ἄφρων}, first occurs in the religious sense in a law of Valentinian, of 368 (Cod. Theodos. l. xvi. tit 2, I. 18), and came into general use under Theodosius, instead of the earlier terms: \emph{gentes, gentiles, nationes, Graeci, cultores simulacrorum}, etc. The English \emph{heathen} and \emph{heathenism} (from \emph{heath),} and the German \emph{Heiden} and \emph{Heidenthum} (from \emph{Heide}), have a similar meaning, and are probably imitations of the Latin \emph{paganismus} in its later usage.} What an inversion of the state of things in the second century, when Celsus contemptuously called Christianity a religion of mechanics and slaves! Of course large exceptions must in both cases be made. Especially in Rome, many of the oldest and most respectable families for a long time still adhered to the heathen traditions, and the city appears to have preserved until the latter part of the fourth century a hundred and fifty-two temples and a hundred and eighty-three smaller chapels and altars of patron deities.\footnote{According to the Descriptiones Urbis of Publicus Victor and Sextus Rufus Festus, which cannot have been composed before, nor long after, the reign of Valentinian. Comp. Beugnot, l.c. i. 266, and Robertson, l.c. p. 260.} But advocates of the old religion—a Themistius, a Libanius, and a Symmachus—limited themselves to the claim of toleration, and thus, in their oppressed condition, became, as formerly the Christians were, and as the persecuted sects in the Catholic church and the Protestant state churches since have been, advocates of religious freedom.

The same toleration continued under Gratian, son and successor of Valentinian (375–383). After a time, however; under the influence of Ambrose, bishop of Milan, this emperor went a step further. He laid aside the title and dignity of Pontifex Maximus, confiscated the temple property, abolished most of the privileges of the priests and vestal virgins, and withdrew, at least in part, the appropriation from the public treasury for their support.\footnote{Cod. Theos. xii. 1, 75; xvi. 10, 20. Symmach. Ep. x. 61. Ambrose, Ep. xvii.} By this step heathenism became, like Christianity before Constantine and now in the American republic, dependent on the voluntary system, while, unlike Christianity, it had no spirit of self-sacrifice, no energy of self-preservation. The withdrawal of the public support cut its lifestring, and left it still to exist for a time by vis inertiae alone. Gratian also, in spite of the protest of the heathen party, removed in 382 the statue and the altar of Victoria, the goddess of victory, in the senate building at Rome, where once the senators used to take their oath, scatter incense, and offer sacrifice; though he was obliged still to tolerate there the elsewhere forbidden sacrifices and the public support of some heathen festivities. Inspired by Ambrose with great zeal for the Catholic faith, he refused freedom to heretics, and prohibited the public assemblies of the Eunomians, Photinians, and Manichaeans.

His brother, Valentinian II. (383–392), rejected the renewed petition of the Romans for the restoration of the altar of Victoria (384). The eloquent and truly venerable prefect Symmachus, who, as princeps senatus and first Pontifex in Rome, was now the spokesman of the heathen party, prayed the emperor in a dignified and elegant address, but in the tone of apologetic diffidence, to make a distinction between his private religion and the religio urbis, to respect the authority of antiquity and the rights of the venerable city, which had attained the dominion of the world under the worship of the gods. But Ambrose of Milan represented to the emperor, in the firm tone of episcopal dignity and conscious success, that the granting of the petition would be a sanctioning of heathenism and a renunciation of his Christian convictions; denied, that the greatness of Rome was due to idolatry, to which indeed her subjugated enemies were likewise addicted; and contrasted the power of Christianity, which had greatly increased under persecution and had produced whole hosts of consecrated virgins and ascetics, with the weakness of heathenism, which, with all its privileges, could hardly maintain the number of its seven vestals, and could show no works of benevolence and mercy for the oppressed. The same petition was renewed in 389 to Theodosius, but again through the influence of Ambrose rejected. The last national sanctuary of the Romans had hopelessly fallen. The triumph, which the heathen party gained under the usurper Eugenius (392–394), lasted but a couple of years; and after his defeat by Theodosius, six hundred of the most distinguished patrician families, the Annii, Probi, Anicii, Olybii, Paulini, Bassi, Gracchi, \&c., are said by Prudentius to have gone over at once to the Christian religion.

\xsection{Theodosius the Great and his Successors. a.d. 392-550}

§ 6. Theodosius the Great and his Successors. a.d. 392–550.

J. R. Stuffken: Diss. de Theod. M. in rem. Christ. meritis. Leyden, 1828. M. Fléchier: Histoire de Theodose le Grand. Par. 1860.

The final suppression of heathenism is usually, though not quite justly, ascribed to the emperor Theodosius I., who, on this account, as well as for his victories over the Goths, his wise legislation, and other services to the empire, bears the distinction of the Great, and deserves, for his personal virtues, to be counted among the best emperors of Rome.\footnote{Gibbon gives a very favorable estimate of his character, and justly charges the heathen Zosimus with gross prejudice against Theodosius. Schlosser and Milman also extol him.} A native of Spain, son of a very worthy general of the same name, he was called by Gratian to be co-emperor in the East in a time of great danger from the threatening barbarians (379), and after the death of Valentinian, he rose to the head of the empire (392–395). He labored for the unity, of the state and the supremacy of the Catholic religion. He was a decided adherent of the Nicene orthodoxy, procured it the victory at the second ecumenical council (381), gave it all the privileges of the state religion, and issued a series of rigid laws against all heretics and schismatics. In his treatment of heathenism, for a time he only enforced the existing prohibition of sacrifice for purposes of magic and divination (385), but gradually extended it to the whole sacrificial worship. In the year 391 he prohibited, under heavy fine, the visiting of a heathen temple for a religious purpose; in the following year, even the private performance of libations and other pagan rites. The practice of idolatry was therefore henceforth a political offence, as Constantius had already, though prematurely, declared it to be, and was subjected to the severest penalties.\footnote{Cod. Theos. xvi. 10, 12.}

Yet Theodosius by no means pressed the execution of these laws in places where the heathen party retained considerable strength; he did not exclude heathens from public office, and allowed them at least full liberty of thought and speech. His countryman, the Christian poet Prudentius, states with approbation, that in the distribution of the secular offices, he looked not at religion, but at merit and talent, and raised the heathen Symmachus to the dignity of consul.\footnote{Prudent. in Symrnachum (written A-D. 403), l. i. v. 617 sqq.: \par “Denique pro meritis terrestribus aequa rependens \par Munera sacricolis summos impertit honores \par Dux bonus, et certare sinit cum laud e suorum, \par Nec pago implicitos [i.e. paganos, heathen] per debita culmina mundi \par Ire viros prohibet: quoniam coelestia nunquam \par Terrenis solitum per iter gradientibus obstant. \par Ipse magistratum tibi consulis, ipse tribunal \par Contulit.”} The emperor likewise appointed the heathen rhetorician, Themistius, prefect of Constantinople, and even intrusted him with the education of his son Arcadius. He acknowledged personal friendship toward Libanius, who addressed to him his celebrated plea for the temples in 384 or 390; though it is doubtful whether he actually delivered it in the imperial presence. In short this emperor stood in such favor with the heathens, that after his death he was enrolled by the Senate, according to ancient custom, among the gods.\footnote{Claudian, who at this period roused pagan poetry from its long sleep and derived his inspiration from the glory of Theodosius and his family, represents his death as an ascension to the gods. De tertio consulatu Honorii, v. 162 sqq.}

Theodosius issued no law for the destruction of temples. He only continued Gratian’s policy of confiscating the temple property and withdrawing entirely the public contribution to the support of idolatry. But in many places, especially in the East, the fanaticism of the monks and the Christian populace broke out in a rage for destruction, which Libanius bitterly laments. He calls these iconoclastic monks “men in black clothes, as voracious as elephants, and insatiably thirsty, but concealing their sensuality under an artificial paleness.” The belief of the Christians, that the heathen gods were living beings, demons,\footnote{Ambrose, Resp. ad Symmachum: “Dii enim gentium daemonia, ut Scriptura docet.” Comp. Ps. xcvi. 5, Septuag.: \textgreek{Πάντες οἱ θεοὶ τῶν ἐθνῶν δαιμόνια}. On this principle especially St. Martin of Tours proceeded in his zeal against the idol temples of Gaul. He asserted that the devil himself frequently assumed the visible form of Jupiter and Mercury, of Minerva and Venus, to protect their sinking sanctuaries. See Sulpit. Severna: Vita B. Martini, c. 4 and 6.} and dwelt in the temples, was the leading influence here, and overshadowed all artistic and archaeological considerations. In Alexandria, a chief seat of the Neo-Platonic mysticism, there arose, at the instigation of the violent and unspiritual bishop Theophilus,\footnote{Gibbon styles him, unfortunately not without reason, “a bold, bad man, whose hands were alternately polluted with gold and with blood.”} a bloody conflict between heathens and Christians, in which the colossal statue and the magnificent temple of Serapis, next to the temple of Jupiter Capitolinus in Rome the proudest monument of heathen architecture,\footnote{See an extended description of the Serapeion in Gibbon, and especially in Milman: Hist. of Christianity, \&c., book iii. c. 8 (p. 377 sqq. N. York ed.).} was destroyed, without verifying the current expectation that upon its destruction the heavens would fall (391). The power of superstition once broken by this decisive blow, the other temples in Egypt soon met a similar fate; though the eloquent ruins of the works of the Pharaohs, the Ptolemies, and the Roman emperors in the valley of the Nile still stand and cast their twilight into the mysterious darkness of antiquity. Marcellus, bishop of Apamea in Syria, accompanied by an armed band of soldiers and gladiators, proceeded with the same zeal against the monuments and vital centres of heathen worship in his diocese, but was burnt alive for it by the enraged heathens, who went unpunished for the murder. In Gaul, St. Martin of Tours, between the years 375 and 400, destroyed a multitude of temples and images, and built churches and cloisters in their stead.

But we also hear important protests from the church against this pious vandalism. Says Chrysostom at Antioch in the beginning of this reign, in his beautiful tract on the martyr Babylas: “Christians are not to destroy error by force and violence, but should work the salvation of men by persuasion, instruction, and love.” In the same spirit says Augustin, though not quite consistently: “Let us first obliterate the idols in the hearts of the heathen, and once they become Christians they will either themselves invite us to the execution of so good a work [the destruction of the idols], or anticipate us in it. Now we must pray for them, and not exasperate them.” Yet he commended the severe laws of the emperors against idolatry.

In the west the work of destruction was not systematically carried on, and the many ruined temples of Greece and Italy at this day prove that even then reason and taste sometimes prevailed over the rude caprice of fanaticism, and that the maxim, It is easier to tear down than to build up, has its exceptions.

With the death of Theodosius the empire again fell into two parts, which were never afterward reunited. The weak sons and successors of this prince, Arcadius in the east (395–408) and Honorius in the west (395–423), and likewise Theodosius II., or the younger (son of Arcadius, 408–450), and Valentinian III. (423–455), repeated and in some cases added to the laws of the previous reign against the heathen. In the year 408, Honorius even issued an edict excluding heathens from civil and military office;\footnote{Cod. Theodos. xvi. 5, 42: “Eos qui Catholicae sectae sunt inimici, intra palatium militare prohibemus. Nullus nobis sit aliqua ratione conjunctus, qui a nobis fide et religione discordat.” According to the somewhat doubtful but usually admitted testimony of Zosimus, l. v. c. 46, this edict was revoked, in consequence of the threatened resignation of a pagan general, Generid, whom Honorius could not dispense with. But Theodosius issued similar laws in the east from 410 to 439. See Gibbon, Milman, Schröckh, and Neander, l.c. The latter erroneously places the edict of Honorius in the year 416, instead of 408.} and in 423 appeared another edict, which questioned the existence of heathens.\footnote{Theodos. II. in Cod. Theodos. xvi. 10, 22: “Paganos, qui supersunt, \emph{quamquam jam nullos} \emph{esse credamus,} promulgatarum legum jamdudum praescripta compescant.” But between 321 and 426 appeared no less than eight laws against apostasy to heathenism; showing that many nominal Christians changed their religion according to circumstances.} But in the first place, such laws, in the then critical condition of the empire amidst the confusion of the great migration, especially in the West, could be but imperfectly enforced; and in the next place, the frequent repetition of them itself proves that heathenism still had its votaries. This fact is witnessed also by various heathen writers. Zosimus wrote his “New History,” down to the year 410, under the reign and at the court of the younger Theodosius (appearing in the high office of comes and advocatus fisci, as he styles himself), in bitter prejudice against the Christian emperors. In many places the Christians, in their work of demolishing the idols, were murdered by the infuriated pagans.

Meantime, however, there was cruelty also on the Christian side. One of the last instances of it was the terrible tragedy of Hypatia. This lady, a teacher of the Neo-Platonic philosophy in Alexandria, distinguished for her beauty, her intelligence, her learning, and her virtue, and esteemed both by Christians and by heathens, was seized in the open street by the Christian populace and fanatical monks, perhaps not without the connivance of the violent bishop Cyril, thrust out from her carriage, dragged to the cathedral, completely stripped, barbarously murdered with shells before the altar, and then torn to pieces and burnt, a.d. 415.\footnote{Socrat. vii. 15 (who considers Cyril guilty); the letters of Synesius, a pupil of Hypatia; and Philostorg. viii. 9. Comp. also Schröckh, vii. 45 sqq. and Wernsdorf: De Hypatia, philosopha Alex. diss. iv. Viteb. 1748. The “Hypatia” of Charles Kingsley is a historical didactic romance, with a polemical aim against the Puseyite overvaluation of patristic Christianity.} Socrates, who relates this, adds: “It brought great censure both on Cyril and on the Alexandrian church.”

\xsection{The Downfall of Heathenism}

§ 7. The Downfall of Heathenism.

The final dissolution of heathenism in the eastern empire may be dated from the middle of the fifth century. In the year 435 Theodosius II. commanded the temples to be destroyed or turned into churches. There still appear some heathens in civil office and at court so late as the beginning of the reign of Justinian I. (527–567). But this despotic emperor prohibited heathenism as a form of worship in the empire on pain of death, and in 529 abolished the last intellectual seminary of it, the philosophical school of Athens, which had stood nine hundred years. At that time just seven philosophers were teaching in that school,\footnote{Damascius of Syria, Simplicius of Cilicia (the most celebrated), Eulalius of Phrygia, Priscianus of Lydia, Isidore of Gaza, Hermias, and Diogenes. They had the courage to prefer exile to the renunciation of their convictions, and found with King Chosroes of Persia a welcome reception, but afterwards returned into the Roman empire under promise of toleration. Comp. Schröckh, xvi. p. 74 sqq.} the shades of the ancient seven sages of Greece,—a striking play of history, like the name of the last west-Roman emperor, Romulus Augustus, or, in contemptuous diminutive, Augustulus, combining the names of the founder of the city and the founder of the empire.

In the West, heathenism maintained itself until near the middle of the sixth century, and even later, partly as a private religious conviction among many cultivated and aristocratic families in Rome, partly even in the full form of worship in the remote provinces and on the mountains of Sicily, Sardinia,\footnote{On these remains of heathenism in the West comp. the citations of Gieseler, i. §79, not. 22 and 23 (i. 2. p. 38-40. Engl. ed. of N. York, i. p. 219 sq.).} and Corsica, and partly in heathen customs and popular usages like the gladiatorial shows still extant in Rome in 404, and the wanton Lupercalia, a sort of heathen carnival, the feast of Lupercus, the god of herds, still celebrated with all its excesses in February, 495. But, in general, it may be said that the Graeco-Roman heathenism, as a system of worship, was buried under the ruins of the western empire, which sunk under the storms of the great migration. It is remarkable that the northern barbarians labored with the same zeal in the destruction of idolatry as in the destruction of the empire, and really promoted the victory of the Christian religion. The Gothic king Alaric, on entering Rome, expressly ordered that the churches of the apostles Peter and Paul should be spared, as inviolable sanctuaries; and he showed a humanity, which Augustin justly attributes to the influence of Christianity (even perverted Arian Christianity) on these barbarous people. The Christian name, he says, which the heathen blaspheme, has effected not the destruction, but the salvation of the city.\footnote{Aug.: De Civit. Dei, l. i. c. 1-6.} Odoacer, who put an end to the western Roman empire in 476, was incited to his expedition into Italy by St. Severin, and, though himself an Arian, showed great regard to the catholic bishops. The same is true of his conqueror and successor, Theodoric the Ostrogoth, who was recognized by the east-Roman emperor Anastasius as king of Italy (a.d. 500), and was likewise an Arian. Thus between the barbarians and the Romans, as between the Romans and the Greeks and in a measure also the Jews, the conquered gave laws to the conquerors. Christianity triumphed over both.

This is the end of Graeco-Roman heathenism, with its wisdom, and beauty. It fell a victim to a slow but steady process of incurable consumption. Its downfall is a sublime tragedy which, with all our abhorrence of idolatry, we cannot witness without a certain sadness. At the first appearance of Christianity it comprised all the wisdom, literature, art, and political power of the civilized world, and led all into the field against the weaponless religion of the crucified Nazarene. After a conflict of four or five centuries it lay prostrate in the dust without hope of resurrection. With the outward protection of the state, it lost all power, and had not even the courage of martyrdom; while the Christian church showed countless hosts of confessors and blood-witnesses, and Judaism lives to-day in spite of all persecution. The expectation, that Christianity would fall about the year 398, after an existence of three hundred and sixty-five years,\footnote{Augustin mentions this story, De Civit. Dei, xviii. 53. Gieseler (vol. i. § 79, not. 17) derives it from a heathen perversion of the Christian (heretical) expectation of the second coming of Christ and the end of the world; referring to Philastr. haer. 106: “Alia est haeresis de anno annunciato ambigens, quod ait propheta Esaias: \emph{Annuntiare annum Dei acceptabilem et diem retributionis.} Putant ergo quidam, quod ex quo venit Dominus usque ad consummationem saeculi non plus nec minus fieri annorum numerum, nisi CCCLXV usque ad Christi Domini iterum de coelo divinam praesentiam.”} turned out in the fulfilment to relate to heathenism itself. The last glimmer of life in the old religion was its pitiable prayer for toleration and its lamentation over the ruin of the empire. Its best elements took refuge in the church and became converted, or at least took Christian names. Now the gods were dethroned, oracles and prodigies ceased, sibylline books were burned, temples were destroyed, or transformed into churches, or still stand as memorials of the victory of Christianity.\footnote{Comp. August.: Epist. 232, where he thus eloquently addresses the heathen: Videtis simulacrorum templa partim sine reparatione collapsa, partim diruta, partim clausa, partim in usus alienos commutata; ipsaque simulacra vel confringi, vel incendi, vel includi, vel destrui; atque ipsas huius saeculi potestates quae aliquando pro simulacris populum Christianum persequebantur, victas et domitas, non a repugnantibus sed a morientibus Christianis, et contra eadem simulacra, pro quibus Christianos occidebant, impetus suos legesque vertisse et imperii nobilissimi eminentissimum culmen ad sepulcrum piscatoris Petri submisso diademate supplicare.”}

But although ancient Greece and Rome have fallen forever, the spirit of Graeco-Roman paganism is not extinct. It still lives in the natural heart of man, which at this day as much as ever needs regeneration by the spirit of God. It lives also in many idolatrous and superstitious usages of the Greek and Roman churches, against which the pure spirit of Christianity has instinctively protested from the beginning, and will protest, till all remains of gross and refined idolatry shall be outwardly as well as inwardly overcome, and baptized and sanctified not only with water, but also with the spirit and fire of the gospel.

Finally the better genius of ancient Greece and Rome still lives in the immortal productions of their poets, philosophers, historians, and orators,—yet no longer an enemy, but a friend and servant of Christ. What is truly great, and noble, and beautiful can never perish. The classic literature had prepared the way for the gospel, in the sphere of natural culture, and was to be turned thenceforth into a weapon for its defence. It passed, like the Old Testament, as a rightful inheritance, into the possession of the Christian church, which saved those precious works of genius through the ravages of the migration of nations and the darkness of the middle ages, and used them as material in the rearing of the temple of modern civilization. The word of the great apostle of the Gentiles was here fulfilled: “All things are yours.” The ancient classics, delivered from the demoniacal possession of idolatry, have come into the service of the only true and living God, once “unknown” to them, but now everywhere revealed, and are thus enabled to fulfil their true mission as the preparatory tutors of youth for Christian learning and culture. This is the noblest, the most worthy, and most complete victory of Christianity, transforming the enemy into friend and ally.

\xchapter{The Literary Triumph of Christianity over Greek and Roman Heathenism}

THE LITERARY TRIUMPH OF CHRISTIANITY OVER GREEK AND ROMAN HEATHENISM.

\xsection{Heathen Polemics. New Objections}

§ 8. Heathen Polemics. New Objections.

I. Comp. The sources at §§ 4 and 5, especially the writings of Julian The Apostate Κατά Χριστιανῶν, and Libanius, ὑπὲρ τῶν ἱερῶν. Also Pseudo-lucian: Philopatris (of the age of Julian or later, comprised in the works of Lucian). Proclus (412–487): xviii ἐπιχειρήματα κατά χριστιανῶν(preserved in the counter work of Joh. Philoponus: De aeternitate mundi, ed. Venet. 1535). In part also the historical works of Eunapius and Zosimus.

II. Marqu. d’Argens: defense du paganisme par l’emper. Julien en grec et en franc. (collected from fragments in Cyril), avec des dissertat. Berl. 1764, sec. ed. Augmentée, 1767. This singular work gave occasion to two against it by G. Fr. Meier, Halle, 1764, And W. Crichton, Halle, 1765, in which the arguments of Julian were refuted anew. Nath. Lardner, in his learned collection of ancient heathen testimonies for the credibility of the Gospel History, treats also largely of Julian. See his collected works, ed. by Dr. Kippis, Lond. 1838, vol. vii. p. 581–652. Schröckh: vi. 354–385. Neander: iii. 77 sqq. (Engl. transl. of Torrey ii. 84–93).

The internal conflict between heathenism and Christianity presents the same spectacle of dissolution on the one hand and conscious power on the other. And here the Nicene age reaped the fruit of the earlier apologists, who ably and fearlessly defended the truth of the true religion and refuted the errors of idolatry in the midst of persecution.\footnote{Comp. vol. i. §§ 60-66.} The literary opposition to Christianity had already virtually exhausted itself, and was now thrown by the great change of circumstances into apology for heathenism; while what was then apology on the Christian side now became triumphant polemics. The last enemy was the Neo-Platonic philosophy, as taught particularly in the schools of Alexandria and Athens even down to the fifth century. This philosophy, however, as we have before remarked,\footnote{Comp. § 4 (p. 42), and vol. i. § 61.} was no longer the product of pure, fresh heathenism, but an artificial syncretism of elements heathen and Christian, Oriental and Hellenic, speculative and theurgic, evincing only the growing weakness of the old religion and the irresistible power of the new.

Besides the old oft-refuted objections, sundry new ones came forward after the time of Constantine, in some cases the very opposite of the earlier ones, touching not so much the Christianity of the Bible as more or less the state-church system of the Nicene and post-Nicene age, and testifying the intrusion of heathen elements into the church. Formerly simplicity and purity of morals were the great ornament of the Christians over against the prevailing corruption; now it could be justly observed that, as the whole world had crowded into the church, it had let in also all the vices of the world. Against those vices, indeed, the genuine virtues of Christianity proved themselves as vigorous as ever. But the heathen either could not or would not look through the outward appearance and discriminate the wheat from the chaff. Again: the Christians of the first three centuries had confessed their faith at the risk of life, maintained it under sufferings and death, and claimed only toleration; now they had to meet reproach from the heathen minority for hypocrisy, selfishness, ambition, intolerance, and the spirit of persecution against heathens, Jews, and heretics. From being suspected as enemies to the emperor and the empire, they now came to be charged in various ways with servile and fawning submission to the Christian rulers. Formerly known as abhorring every kind of idolatry and all pomp in worship, they now appeared in their growing veneration for martyrs and relics to reproduce and even exceed the ancient worship of heroes.

Finally, even the victory of Christianity was branded as a reproach. It was held responsible by the latest heathen historians not only for the frequent public calamities, which had been already charged upon it under Marcus Aurelius and in the time of Tertullian, but also for the decline and fall of the once so mighty Roman empire. But this objection, very popular at the time, is refuted by the simple fact, that the empire in the East, where Christianity earlier and more completely prevailed, outlived by nearly ten centuries the western branch. The dissolution of the west-Roman empire was due rather to its unwieldy extent, the incursion of barbarians, and the decay of morals, which was hastened by the introduction of all the vices of conquered nations, and which had already begun under Augustus, yea, during the glorious period of the republic; for the republic would have lasted much longer if the foundations of public and private virtue had not been undermined.\footnote{Gibbon, too, imputes the fall of the west-Roman empire not, as unjustly charged by Dr. Kurtz (Handbuch der allg. Kirchengesch. i. 2, p. 15, 3d ed.), to Christianity, but almost solely to the pressure of its own weight. Comp. his General Observations on the Fall of the R. Empire in the West, at the close of ch. xxxviii., where he says: “The decline of Rome was the natural and inevitable effect of immoderate greatness. Prosperity ripened the principle of decay; the causes of destruction multiplied with the extent of conquest; and as soon as time or accident had removed the artificial supports, the stupendous fabric yielded to the pressure of its own weight. The story of its ruin is simple and obvious; and instead of inquiring \emph{why} the Roman empire was destroyed, we should rather be surprised that it had subsisted so long.” Gibbon then mentions Christianity also, it is true, or more properly monasticism, which, he thinks, suppressed with its passive virtues the patriotic and martial spirit, and so far contributed to the catastrophe; but adds: “If the decline of the Roman empire was \emph{hastened} [—he says not: \emph{caused—}] by the conversion of Constantine, his victorious religion broke the violence of the fall, and mollified the ferocious temper of the conquerors.” This view is very different from that of Eunapius and Zosimus, with which Kurtz identifies it. Gibbon in general follows more closely Ammianus Marcellinus, whom, with all reason, he holds as a historian far superior to the others.—Lord Byron truthfully expresses the law of decay to which Rome succumbed, in these words from Childe Harold: \par “There is the moral of all human tales; \par ’T is but the same rehearsal of the past: \par First freedom, and then glory—when that fails, \par Wealth, vice, corruption, barbarism at last.”} Taken from a higher point of view, the downfall of Rome was a divine judgment upon the old essentially heathen world, as the destruction of Jerusalem was a judgment upon the Jewish nation for their unbelief. But it was at the same time the inevitable transition to a new creation which Christianity soon began to rear on the ruins of heathendom by the conversion of the barbarian conquerors, and the founding of a higher Christian civilization. This was the best refutation of the last charge of the heathen opponents of the religion of the cross.

\xsection{Julian's Attack upon Christianity}

§ 9. Julian’s Attack upon Christianity.

For Literature comp. § 4 p. 39, 40.

The last direct and systematic attack upon the Christian religion proceeded from the emperor Julian. In his winter evenings at Antioch in 363, to account to the whole world for his apostasy, he wrote a work against the Christians, which survives, at least in fragments, in a refutation of it by Cyril of Alexandria, written about 432. In its three books, perhaps seven (Cyril mentions only three\footnote{In the preface to his refutation, Contra Jul. i. p. 3: \textgreek{Τρία συγγέγραψε βιβλία κατὰ τῶν ἁγίων εὐαγγελίων καὶ κατὰ τῆς εὐαγοῦς τῶν Χριστιανῶν θρησκείας}. But Jerome says, Epist. 83 (tom. iv. p. 655): ” Julianus Augustus \emph{septem libros,} in expeditione Parthica [or rather \emph{before} he left Antioch and started for Persia], adversus Christianos vomuit.”}), it shows no trace of the dispassionate philosophical or historical appreciation of so mighty a phenomenon as Christianity in any case is. Julian had no sense for the fundamental ideas of sin and redemption or the cardinal virtues of humility and love. He stood entirely in the sphere of naturalism, where the natural light of Helios outshines the mild radiance of the King of truth, and the admiration of worldly greatness leaves no room for the recognition of the spiritual glory of self-renunciation. He repeated the arguments of a Celsus and a Porphyry in modified form; expanded them by his larger acquaintance with the Bible, which he had learned according to the letter in his clerical education; and breathed into all the bitter hatred of an Apostate, which agreed ill with his famous toleration and entirely blinded him to all that was good in his opponents. He calls the religion of “the Galilean” an impious human invention and a conglomeration of the worst elements of Judaism and heathenism without the good of either; that is, without the wholesome though somewhat harsh discipline of the former, or the pious belief in the gods, which belongs to the latter. Hence he compares the Christians to leeches, which draw all impure blood and leave the pure. In his view, Jesus, “the dead Jew,” did nothing remarkable during his lifetime, compared with heathen heroes, but to heal lame and blind people and exorcise daemoniacs, which is no very great matter.\footnote{Cyril has omitted the worst passages of Julian respecting Christ, but quotes the following (Contra Jul. l. vi. p. 191, ed. Spanh.), which is very characteristic: “Jesus, who over-persuaded much (\textgreek{ἀναπείσας}) the lowest among you, some few, has now been talked of (\textgreek{ὀνομάζεται}) for three hundred years, though during his life he performed nothing worth mentioning (\textgreek{οὐδὲν ἀκοῆς ἄξιον}), unless it be thought a mighty matter to heal the cripples and blind persons and to exorcise those possessed of demons in the villages of Bethsaida and Bethany (\textgreek{εἰ μή τις εἴεται τοὺς κολλοὺς καὶ τοὺς τυφλοὺς ιάσασθαι, καὶ δαιμονώντας ἐφορκίζειν ἐν Βηθσείδᾳ καὶ ἐν Βηθανίᾳ ταῖς κώμαις τῶν μεγίστων ἔργων εῖναι} )” Dr. Lardner has ingeniously inferred from this passage that, Julian, by conceding to Christ the power of working miracles, and admitting the general truths of the gospel traditions, furnishes an argument for Christianity rather than against it.} He was able to persuade only a few of the ignorant peasantry, not even to gain his own kinsmen.\footnote{Jno. vii. 5.} Neither Matthew, nor Mark, nor Luke, nor Paul called him God. John was the first to venture so far, and procured acceptance for his view by a cunning artifice.\footnote{“Neither Paul,” he says (Cyr. l. x. p. 327), “nor Matthew, nor Luke, nor Mark has dared to call Jesus God. But honest John (\textgreek{ὁ χρηστόσ Ἰωάννης}\emph{),} understanding that a great multitude of men in the cities of Greece and Italy were seized with this distemper; and hearing likewise, as I suppose, that the tombs of Peter and Paul were respected, and frequented, though as yet privately only, however, having heard of it, he then first presumed to advance that doctrine.”} The later Christians perverted his doctrine still more impiously, and have abandoned the Jewish sacrificial worship and ceremonial law, which was given for all time, and was declared irrevocable by Jesus himself.\footnote{Matt. v. 17-19.} A universal religion, with all the peculiarities of different national characters, appeared to him unreasonable and impossible. He endeavored to expose all manner of contradictions and absurdities in the Bible. The Mosaic history of the creation was defective, and not to be compared with the Platonic. Eve was given to Adam for a help, yet she led him astray. Human speech is put into the mouth of the serpent, and the curse is denounced on him, though he leads man on to the knowledge of good and evil, and thus proves himself of great service. Moses represents God as jealous, teaches monotheism, yet polytheism also in calling the angels gods. The moral precepts of the decalogue are found also among the heathen, except the commands, “Thou shalt have no other gods before me,” and, “Remember the Sabbath day.” He prefers Lycurgus and Solon to Moses. As to Samson and David, they were not very remarkable for valor, and exceeded by many Greeks and Egyptians, and all their power was confined within the narrow limits of Judea. The Jews never had any general equal to Alexander or Caesar. Solomon is not to be compared with Theognis, Socrates, and other Greek sages; moreover he is said to have been overcome by women, and therefore does not deserve to be ranked among wise men. Paul was an arch-traitor; calling God now the God of the Jews, now the God of the Gentiles, now both at once; not seldom contradicting the Old Testament, Christ, and himself, and generally accommodating his doctrine to circumstances. The heathen emperor thinks it absurd that Christian baptism should be able to cleanse from gross sins, while it cannot remove a wart, or gout, or any bodily evil. He puts the Bible far below the Hellenic literature, and asserts, that it made men slaves, while the study of the classics educated great heroes and philosophers. The first Christians he styles most contemptible men, and the Christians of his day he charges with ignorance, intolerance, and worshipping dead persons, bones, and the wood of the cross.

With all his sarcastic bitterness against Christianity, Julian undesignedly furnishes some valuable arguments for the historical character of the religion he hated and assailed. The learned and critical Lardner, after a careful analysis of his work against Christianity, thus ably and truthfully sums up Julian’s testimony in favor of it:

“Julian argues against the Jews as well as against the Christians. He has borne a valuable testimony to the history and to the books of the New Testament, as all must acknowledge who have read the extracts just made from his work. He allows that Jesus was born in the reign of Augustus, at the time of the taxing made in Judea by Cyrenius: that the Christian religion had its rise and began to be propagated in the times of the emperors Tiberius and Claudius. He bears witness to the genuineness and authenticity of the four gospels of Matthew, Mark, Luke, and John, and the Acts of the Apostles: and he so quotes them, as to intimate, that these were the only historical books received by Christians as of authority, and the only authentic memoirs of Jesus Christ and his apostles, and the doctrine preached by them. He allows their early date, and even argues for it. He also quotes, or plainly refers to the Acts of the Apostles, to St. Paul’s Epistles to the Romans, the Corinthians, and the Galatians. He does not deny the miracles of Jesus Christ, but allows him to have ’healed the blind, and the lame, and demoniacs,’ and ’to have rebuked the winds, and walked upon the waves of the sea.’ He endeavors indeed to diminish these works; but in vain. The consequence is undeniable: such works are good proofs of a divine mission. He endeavors also to lessen the number of the early believers in Jesus, and yet he acknowledgeth, that there were ’multitudes of such men in Greece and Italy,’ before St. John wrote his gospel. He likewise affects to diminish the quality of the early believers; and yet acknowledgeth, that beside ’menservants, and maidservants,’ Cornelius, a Roman centurion at Caesarea, and Sergius Paulus, proconsul of Cyprus, were converted to the faith of Jesus before the end of the reign of Claudius. And he often speaks with great indignation of Peter and Paul, those two great apostles of Jesus, and successful preachers of his gospel. So that, upon the whole, he has undesignedly borne witness to the truth of many things recorded in the books of the New Testament: he aimed to overthrow the Christian religion, but has confirmed it: his arguments against it are perfectly harmless, and insufficient to unsettle the weakest Christian. He justly excepts to some things introduced into the Christian profession by the late professors of it, in his own time, or sooner; but has not made one objection of moment against the Christian religion, as contained in the genuine and authentic books of the New Testament.”\footnote{Dr. Nathiel Lardner’s Works, ed. by Dr. Kippis in ten vols. Vol. vii. pp. 638 and 639. As against the mythical theory of Strauss and Renan the extract from Lardner has considerable force, as well as his whole work on the credibility of the Gospel History.}

The other works against Christianity are far less important.

The dialogue Philopatris, or The Patriot, is ascribed indeed to the ready scoffer and satirist Lucian (died about 200), and joined to his works; but it is vastly inferior in style and probably belongs to the reign of Julian, or a still later period;\footnote{According to Niebuhr’s view it must have been composed under the emperor Phocas, 968 or 969. Moyle places it in the year 302, Dodwell in the year 261, others in the year 272.} since it combats the church doctrine of the Trinity and of the procession of the Spirit from the Father, though not by argument, but only by ridicule. It is a frivolous derision of the character and doctrines of the Christians in the form of a dialogue between Critias, a professed heathen, and Triephon, an Epicurean, personating a Christian. It represents the Christians as disaffected to the government, dangerous to civil society, and delighting in public calamities. It calls St. Paul a half bald, long-nosed Galilean, who travelled through the air to the third heaven (2 Cor. 12, 1–4).

The last renowned representative of Neo-Platonism, Proclus of Athens (died 487), defended the Platonic doctrine of the eternity of the world, and, without mentioning Christianity, contested the biblical doctrine of the creation and the end of the world in eighteen arguments, which the Christian philosopher, John Philoponus, refuted in the seventh century.

The last heathen historians, Eunapius and Zosimus, of the first half of the fifth century, indirectly assailed Christianity by a one-sided representation of the history of the Roman empire from the time of Constantine, and by tracing its decline to the Christian religion; while, on the contrary, Ammianus Marcellinus (died about 390) presents with honorable impartiality both the dark and the bright sides of the Christian emperors and of the Apostate Julian.\footnote{The more is it to be regretted, that the fisrt thirteen books of his history of the Roman emperors from Nerva to 353 are lost. The remaining eighteen books reach from 353 to 378.}

\xsection{The Heathen Apologetic Literature}

§ 10. The Heathen Apologetic Literature.

After the death of Julian most of the heathen writers, especially the ablest and most estimable, confined themselves to the defence of their religion, and thus became, by reason of their position, advocates of toleration; and, of course, of toleration for the religious syncretism, which in its cooler form degenerates into philosophical indifferentism.

Among these were Themistius, teacher of rhetoric, senator, and prefect of Constantinople, and afterwards preceptor of the young emperor Arcadius; Aurelius Symmachus, rhetorician, senator, and prefect of Rome under Gratian and Valentinian II., the eloquent pleader for the altar of Victoria; and above all, the rhetorician Libanius, friend and admirer of Julian, alternately teaching in Constantinople, Nicomedia, and Antioch. These all belong to the second half of the fourth century, and represent at once the last bloom and the decline of the classic eloquence. They were all more or less devoted to the Neo-Platonic syncretism. They held, that the Deity had implanted in all men a religious nature and want, but had left the particular form of worshiping God to the free will of the several nations and individuals; that all outward constraint, therefore, was contrary to the nature of religion and could only beget hypocrisy. Themistius vindicated this variety of the forms of religion as favorable to religion itself, as many Protestants justify the system of sects. “The rivalry of different religions,” says he in his oration on Jovian, “serves to stimulate zeal for the worship of God. There are different paths, some hard, others easy, some rough, others smooth, leading to the same goal. Leave only one way, and shut up the rest, and you destroy emulation. God would have no such uniformity among men .... The Lord of the universe delights in manifoldness. It is his will, that Syrians, Greeks, Egyptians should worship him, each nation in its own way, and that the Syrians again should divide into small sects, no one of which agrees entirely with another. Why should we thus enforce what is impossible?” In the same style argues Symmachus, who withholds all direct opposition to Christianity and contends only against its exclusive supremacy.

Libanius, in his plea for the temples addressed to Theodosius I. (384 or 390), called to his aid every argument, religious, political, and artistic, in behalf of the heathen sanctuaries, but interspersed bitter remarks against the temple-storming monks. He asserts among other things, that the principles of Christianity itself condemn the use of force in religion, and commend the indulgence of free conviction.

Of course this heathen plea for toleration was but the last desperate defence of a hopeless minority, and an indirect self-condemnation of heathenism for its persecution of the Christian religion in the first three centuries.

\xsection{Christian Apologists and Polemics}

§ 11. Christian Apologists and Polemics.

SOURCES.

I. The Greek Apologists: Eusebius Caes.: Προπαρασκευὴ εὐαγγελική(Preparatio evang.), and Ἀπόδειξις εὐαγγελική(Demonstratio evang.); besides his controversial work against Hierocles; and his Theophany, discovered in 1842 in a Syriac version (ed. Lee, Lond. 1842). Athanasius: Κατὰτῶν Ἑλλήνων(Oratio contra Gentes), and Περὶ τῆς ἐνανθρωπήσεως τοῦ Λόγου(De incarnatione Verbi Dei): two treatises belonging together (Opera, ed. Bened. tom. i. 1 sqq.). Cyril of Alex.: Contra impium Julianum libri X (with extracts from the three books of Julian against Christianity). Theodoret: Graecarum affectionum curatio (Ἑλληνικῶν θεραπευτικὴ παθημάτων), disput. XII.

II. The Latin Apologists: Lactantius: Instit. divin. l. vii (particularly the first three books, de falsa religione, de origine erroris, and de falsa sapientia; the third against the heathen philosophy). Julius Firmicus Maternus: De errore profanarum religionum (not mentioned by the ancients, but edited several times in the sixteenth century, and latterly by F. Münter, Havn. 1826). Ambrose: Ep. 17 and 18 (against Symmachus). Prudentius: In Symmachum (an apologetic poem). Paul. Orosius: Adv. paganos historiarum l. vii (an apologetic universal history, against Eunapius and Zosimus). Augustine: De civitate Dei l. xxii (often separately published). Salvianus: De gubernatione Dei l. viii (the eighth book incomplete).

MODERN LITERATURE.

Comp. in part the apologetic literature at § 63 of vol. i. Also Schrökh: vii., p. 263–355. Neander: iii., 188–195 (Engl. ed. of Torrey, ii., 90–93). Döllinger (R.C.): Hdbuch der K. G., vol. I., part 2, p. 50–91.K. Werner (R.C.): Geschichte der Apolog. und polem. Literatur der christl. Theol. Schaffh. 1861–’65, 4 vols. vol. i.

In the new state of things the defence of Christianity was no longer of so urgent and direct importance as it had been before the time of Constantine. And the theological activity of the church now addressed itself mainly to internal doctrinal controversy. Still the fourth and fifth centuries produced several important apologetic works, which far outshone the corresponding literature of the heathen.

(1) Under Constantine we have Lactantius in Latin, Eusebius and Athanasius in Greek, representing, together with Theodoret, who was a century later, the close of the older apology.

Lactantius prefaces his vindication of Christian truth with a refutation of the heathen superstition and philosophy; and he is more happy in the latter than in the former. He claims freedom for all religions, and represents the transition standpoint of the Constantinian edicts of toleration.

Eusebius, the celebrated historian, collected with diligence and learning in several apologetic works, above all in his “Evangelic Preparation,” the usual arguments against heathenism, and in his “Evangelic Demonstration” the positive evidences of Christianity, laying chief stress upon the prophecies.

With less scholarship, but with far greater speculative compass and acumen, the great Athanasius, in his youthful productions “against the Greeks,” and “on the incarnation of the Logos” (before 325), gave in main outline the argument for the divine origin, the truth, the reasonableness, and the perfection of the Christian religion. These two treatises, particularly the second, are, next to Origen’s doctrinal work De principiis, the first attempt to construct a scientific system of the Christian religion upon certain fundamental ideas of God and world, sin and redemption; and they form the ripe fruit of the positive apology in the Greek church. The Logos, Athanasius teaches, is the image of the living, only true God. Man is the image of the Logos. In communion with him consist the original holiness and blessedness of paradise. Man fell by his own will, and thus came to need redemption. Evil is not a substance of itself, not matter, as the Greeks suppose, nor does it come from the Creator of all things. It is an abuse of freedom on the part of man, and consists in selfishness or self-love, and in the dominion of the sensuous principle over the reason. Sin, as apostasy from God, begets idolatry. Once alienated from God and plunged into finiteness and sensuousness, men deified the powers of nature, or mortal men, or even carnal lusts, as in Aphrodite. The inevitable consequence of sin is death and corruption. The Logos, however, did not forsake men. He gave them the law and the prophets to prepare them for salvation. At last he himself became man, neutralized in human nature the power of sin and death, restored the divine image, uniting us with God and imparting to us his imperishable life. The possibility and legitimacy of the incarnation lie in the original relation of the Logos to the world, which was created and is upheld by him. The incarnation, however, does not suspend the universal reign of the Logos. While he was in man, he was at the same time everywhere active and reposing in the bosom of the Father. The necessity of the incarnation to salvation follows from the fact, that the corruption had entered into human nature itself, and thus must be overcome within that nature. An external redemption, as by preaching God, could profit nothing. “For this reason the Saviour assumed humanity, that man, united with life, might not remain mortal and in death, but imbibing immortality might by the resurrection be immortal. The outward preaching of redemption would have to be continually repeated, and yet death would abide in man.”\footnote{De incarn. c. 44 (Opera ed. Bened. i. p. 86).} The object of the incarnation is, negatively, the annihilation of sin and death; positively, the communication of righteousness and life and the deification of man.\footnote{\textgreek{Ὁ Λόγος ἐνανθρώπησεν, ἲνα ἡμεῖς θεοποιηθῶμεν}.} The miracles of Christ are the proof of his original dominion over nature, and lead men from nature-worship to the worship of God. The death of Jesus was necessary to the blotting out of sin and to the demonstration of his life-power in the resurrection, whereby also the death of believers is now no longer punishment, but a transition to resurrection and glory.—This speculative analysis of the incarnation Athanasius supports by referring to the continuous moral effects of Christianity, which is doing great things every day, calling man from idolatry, magic, and sorceries to the worship of the true God, obliterating sinful and irrational lusts, taming the wild manners of barbarians, inciting to a holy walk, turning the natural fear of death into rejoicing, and lifting the eye of man from earth to heaven, from mortality to resurrection and eternal glory. The benefits of the incarnation are incalculable, like the waves of the sea pursuing one another in constant succession.

(2) Under the sons of Constantine, between the years 343 and 350, Julius Firmicus Maternus, an author otherwise unknown to us,\footnote{It is uncertain whether he was the author of a mathematical and astrological work written some years earlier and published at Basel in 1551, which treats of the influence of the stars upon men, but conjures its readers not to divulge these Egyptian and Babylonian mysteries, as astrology was forbidden at the time. If he were the author, he must have not only wholly changed his religion, but considerably improved his style.} wrote against heathenism with large knowledge of antiquity, but with fanatical zeal, regarding it, now on the principle of Euhemerus, as a deification of mortal men and natural elements, now as a distortion of the biblical history.\footnote{The Egyptian Serapis, for instance, was no other than Joseph, who, being the grand-son of Sara, was named \textgreek{Σαρᾶς ἀπό}.} At the close, quite mistaking the gentle spirit of the New Testament, he urges the sons of Constantine to exterminate heathenism by force, as God commanded the children of Israel to proceed against the Canaanites; and openly counsels them boldly to pillage the temples and to enrich themselves and the church with the stolen goods. This sort of apology fully corresponds with the despotic conduct of Constantius, which induced the reaction of heathenism under Julian.

(3) The attack of Julian upon Christianity brought out no reply on the spot,\footnote{Though Apollinaris wrote a book “Of the Truth” against the emperor and the heathen philosophers, of which Julianis reported to have said sneeringly: \textgreek{Ἀνέγνων, ἔγνων, κατέγνων}:“I have read it, understood it, and condemned it.” To which the Christian bishops rejoined in like tone: \textgreek{Ἀνέγνως, ἀλλ αὐκ ἔγνως , εἰ γάρ ἔγνως οὐκ ἄν κατέγνως}: “You have read, but not understood, for, had you understood you would not have condemned.” So says Sozomen: v. 18. Comp. Schröckh: vi. 355.} but subsequently several refutations, the chief one by Cyril of Alexandria († 444), in ten books “against the impious Julian,” still extant and belonging among his most valuable works. About the same time Theodoret wrote an apologetic and polemic work: “The Healing of the Heathen Affections,” in twelve treatises, in which he endeavors to refute the errors of the false religion by comparison of the prophecies and miracles of the Bible with the heathen oracles, of the apostles with the heroes and lawgivers of antiquity, of the Christian morality with the immorality of the heathen world.

\xsection{Augustine's City of God. Salvianus}

§ 12. Augustine’s City of God. Salvianus.

(4) Among the Latin apologists we must mention Augustine, Orosius, and Salvianus, of the fifth century. They struck a different path from the Greeks, and devoted themselves chiefly to the objection of the heathens, that the overthrow of idolatry and the ascendency of Christianity were chargeable with the misfortunes and the decline of the Roman empire. This objection had already been touched by Tertullian, but now, since the repeated incursions of the barbarians, and especially the capture and sacking of the city of Rome under the Gothic king Alaric in 410, it recurred with peculiar force. By way of historical refutation the Spanish presbyter Orosius, at the suggestion of Augustine, wrote an outline of universal history in the year 417.

Augustine himself answered the charge in his immortal work “On the city of God,” that is, the church of Christ, in twenty-two books, upon which he labored twelve years, from 413 to 426, amidst the storms of the great migration and towards the close of his life. He was not wanting in appreciation of the old Roman virtues, and he attributes to these the former greatness of the empire, and to the decline of them he imputes her growing weakness. But he rose at the same time far above the superficial view, which estimates persons and things by the scale of earthly profit and loss, and of temporary success. “The City of God” is the most powerful, comprehensive, profound, and fertile production in refutation of heathenism and vindication of Christianity, which the ancient church has bequeathed to us, and forms a worthy close to her literary contest with Graeco-Roman paganism.\footnote{Milman says (l.c. book iii. ch. 10) The \emph{City of God} was unquestionably the noblest work, both in its original design and in the fulness of its elaborate execution, which the genius of man had as yet contributed to the support of Christianity.”} It is a grand funeral discourse upon the departing universal empire of heathenism, and a lofty salutation to the approaching universal order of Christianity. While even Jerome deplored in the destruction of the city the downfall of the empire as the omen of the approaching doom of the world,\footnote{Proleg. in Ezek.: In una urbe totus orbis interiit. Epist. 60: Quid salvum est, si Roma perit!} the African father saw in it only a passing revolution preparing the way for new conquests of Christianity. Standing at that remarkable turning-point of history, he considers the origin, progress, and end of the perishable kingdom of this world, and the imperishable kingdom of God, from the fall of man to the final judgment, where at last they fully and forever separate into hell and heaven. The antagonism of the two cities has its root in the highest regions of the spirit world, the distinction of good and evil angels; its historical evolution commences with Cain and Abel, then proceeds in the progress of paganism and Judaism to the birth of Christ, and continues after that great epoch to his return in glory. Upon the whole his philosophy of history is dualistic, and does not rise to the unity and comprehensiveness of the divine plan to which all the kingdoms of this world and even Satan himself are made subservient. He hands the one city over to God, the other to the demons. Yet he softens the rigor of the contrast by the express acknowledgment of shades in the one, and rays of light in the other. In the present order of the world the two cities touch and influence each other at innumerable points; and as not all Jews were citizens of the heavenly Jerusalem, so there were on the other hand true children of God scattered among the heathen like Melchisedek and Job, who were united to the city of God not by a visible, but by an invisible celestial tie. In this sublime contrast Augustine weaves up the whole material of his Scriptural and antiquarian knowledge, his speculation, and his Christian experience, but interweaves also many arbitrary allegorical conceits and empty subtleties. The first ten books he directs against heathenism, showing up the gradual decline of the Roman power as the necessary result of idolatry and of a process of moral dissolution, which commenced with the introduction of foreign vices after the destruction of Carthage; and he represents the calamities and approaching doom of the empire as a mighty preaching of repentance to the heathen, and at the same time as a wholesome trial of the Christians, and as the birth-throes of a new creation. In the last twelve books of this tragedy of history he places in contrast the picture of the supernatural state of God, founded upon a rock, coming forth renovated and strengthened from all the storms and revolutions of time, breathing into wasting humanity an imperishable divine life, and entering at last, after the completion of this earthly work, into the sabbath of eternity, where believers shall rest and see, see and love, love and praise, without end.\footnote{“Ibi vacabimus, ” reads the conclusion, l. xxii. c. 30, “et videbimus; videbimus, et amabimus; amabimus, et laudabimus. Ecce quod erit in fine sine fine. Nam quia alius noster est finis, nisi pervenire ad regnum, cuius nullus est finis.” Tillemont and Schröckh give an extended analysis of the \emph{Civitas Dei.} So also more recently Dr. Baur in his work on the Christian church from the fourth to the sixth century, pp. 43-52. Gibbon, on the other hand, whose great history treats in some sense, though in totally different form and in opposite spirit, the same theme, only touches this work incidentally, notwithstanding his general minuteness. He says in a contemptuous tone, that his knowledge of Augustineis limited to the “Confessions,” and the “City of God.” Of course Augustine’s philosophy of history is almost as flatly opposed to the deism of the English historian, as to the heathen views of his contemporaries Ammianus, Eunapius, and Zosimus.}

Less important, but still noteworthy and peculiar, is the apologetic work of the Gallic presbyter, Salvianus, on providence and the government of the world.\footnote{Of this book: “De gubernatione Dei, et de justo Dei praesentique judicio,” Isaac Taylor has made very large use in his interesting work on “Ancient Christianity” (vol. ii. p. 34 sqq.), to refute the idealized Puseyite view of the Nicene and post-Nicene age. But he ascribes too great importance to it, and forgets that it is an unbalanced picture of the shady side of the church at that time. It is true as far as it goes, and yet leaves a false impression. There are books which by a partial and one-sided representation make even the truth lie.} It was composed about the middle of the fifth century (440–455) in answer at once to the charge that Christianity occasioned all the misfortunes of the times, and to the doubts concerning divine providence, which were spreading among Christians themselves. The blame of the divine judgments he places, however, not upon the heathens, but upon the Christianity of the day, and, in forcible and lively, but turgid and extravagant style, draws an extremely unfavorable picture of the moral condition of the Christians, especially in Gaul, Spain, Italy, and Africa. His apology for Christianity, or rather for the Christian faith in the divine government of the world, was also a polemic against the degenerate Christians. It was certainly unsuited to convert heathens, but well fitted to awaken the church to more dangerous enemies within, and stimulate her to that moral self-reform, which puts the crown upon victory over outward foes. “The church,” says this Jeremiah of his time, “which ought everywhere to propitiate God, what does she, but provoke him to anger?\footnote{“Ipsa Dei ecclesia quae in omnibus esse debet placatrix Dei, quid est aliud quam exacerbatrix Dei? aut, praeter paucissimos quosdam, qui mala fugiunt, quid est aliud pene omnis coetus Christianorum, quam sentina vitiorum?” (P. 91.)} How many may one meet, even in the church, who are not still drunkards, or debauchees, or adulterers, or fornicators, or robbers, or murderers, or the like, or all these at once, without end? It is even a sort of holiness among Christian people, to be less vicious.” From the public worship of God, he continues, and almost during it, they pass to deeds of shame. Scarce a rich man, but would commit murder and fornication. We have lost the whole power of Christianity, and offend God the more, that we sin as Christians. We are worse than the barbarians and heathen. If the Saxon is wild, the Frank faithless, the Goth inhuman, the Alanian drunken, the Hun licentious, they are by reason of their ignorance far less punishable than we, who, knowing the commandments of God, commit all these crimes. He compares the Christians especially of Rome with the Arian Goths and Vandals, to the disparagement of the Romans, who add to the gross sins of nature the refined vices of civilization, passion for theatres, debauchery, and unnatural lewdness. Therefore has the just God given them into the hands of the barbarians and exposed them to the ravages of the migrating hordes.

This horrible picture of the Christendom of the fifth century is undoubtedly in many respects an exaggeration of ascetic and monastic zeal. Yet it is in general not untrue; it presents the dark side of the picture, and enables us to understand more fully on moral and psychological grounds the final dissolution of the western empire of Rome.

\xchapter{Alliance of Church and State and Its Influence on Public Morals and Religion}

ALLIANCE OF CHURCH AND STATE AND ITS INFLUENCE ON PUBLIC MORALS AND RELIGION.

SOURCES.

The church laws of the Christian emperors from Constantine to Justinian, collected in the Codex Theodosianus of the year 438 (edited, with a learned commentary, by Jac. Gothofredus, Lyons, 1668, in six vols. fol.; afterwards by J. D. Ritter, Lips. 1736, in seven vols.; and more recently, with newly discovered books and fragments, by G. Haenel, Bonn, 1842), and in the Codex Justinianeus of 534 (in the numerous editions of the Corpus juris civilis Romani). Also Eusebius: Vita Constant., and H. Eccl. l. x. On the other hand, the lamentations of the church fathers, especially Gregory Naz., Chrysostom, and Augustine (in their sermons), over the secularized Christianity of their time.

LITERATURE.

C. G. de Rhoer: Dissertationes de effectu religionis Christianae in jurisprudentiam Romanam. Groning. 1776. Martini: Die Einführung der christl. Religion als Staatsreligion im röm. Reiche durch Constantin. Münch. 1813. H. O. de Meysenburg: De Christ. religionis vi et effectu in jus civile. Gött. 1828. C. Riffel (R.C.): Gesch. Darstellung des Verhältnisses zwischen Kirche u. Staat. Mainz. 1838, vol. i. Troplong: De l’influence du Christianisme sur le droit civil des Romains. Par. 1843. P. E. Lind: Christendommens inflydelse paa den sociale forfatning. Kjobenh. 1852. B. C. Cooper: The Free Church of Ancient Christendom and its Subjugation by Constantine. Lond. 1851(?)

Comp. also Gibbon, chap. xx. Schröckh, several sections from vol. v. onward. Neander, iii. 273–303. Milman, Anc. Christ. Book iv. ch. 1.

\xsection{The New Position of the Church in the Empire}

§ 13. The New Position of the Church in the Empire.

The previous chapter has shown us how Christianity gradually supplanted the Graeco-Roman heathenism and became the established religion in the empire of the Caesars. Since that time the church and the state, though frequently jarring, have remained united in Europe, either on the hierarchical basis, with the temporal power under the tutelage of the spiritual, or on the caesaro-papal, with the spiritual power merged in the temporal; while in the United States of America, since the end of the eighteenth century, the two powers have stood peacefully but independently side by side. The church could now act upon the state; but so could the state act upon the church; and this mutual influence became a source of both profit and loss, blessing and curse, on either side.

The martyrs and confessors of the first three centuries, in their expectation of the impending end of the world and their desire for the speedy return of the Lord, had never once thought of such a thing as the great and sudden change, which meets us at the beginning of this period in the relation of the Roman state to the Christian church. Tertullian had even held the Christian profession to be irreconcilable with the office of a Roman emperor.\footnote{Apologeticus, c. 21 “Sed et Caesares credidissent, si aut Caesares non essent saeculo necessarii, aut si et Christiani potuissent esse Caesares.”} Nevertheless, clergy and people very soon and very easily accommodated themselves to the new order of things, and recognized in it a reproduction of the theocratic constitution of the people of God under the ancient covenant. Save that the dissenting sects, who derived no benefit from this union, but were rather subject to persecution from the state and from the established Catholicism, the Donatists for an especial instance, protested against the intermeddling of the temporal power with religious concerns.\footnote{Thus the bishop Donatus of Carthage in 347 rejected the imperial commissioners, Paulus and Macarius, with the exclamation: “Quid est imperatori cum ecclesia?” See Optatus Milev.: De schismate Donat. l. iii. c. 3. The Donatists, however, were the first to invoke the imperial intervention in their controversies, and would doubtless have spoken very differently, had the decision turned in their favor.} The heathen, who now came over in a mass, had all along been accustomed to a union of politics with religion, of the imperial with the sacerdotal dignity. They could not imagine a state without some cultus, whatever might be its name. And as heathenism had outlived itself in the empire, and Judaism with its national exclusiveness and its stationary character was totally disqualified, Christianity must take the throne.

The change was as natural and inevitable as it was great. When Constantine planted the standard of the cross upon the forsaken temples of the gods, he but followed the irresistible current of history itself. Christianity had already, without a stroke of sword or of intrigue, achieved over the false religion the internal victory of spirit over matter, of truth over falsehood, of faith over superstition, of the worship of God over idolatry, of morality over corruption. Under a three hundred years’ oppression, it had preserved its irrepressible moral vigor, and abundantly earned its new social position. It could not possibly continue a despised sect, a homeless child of the wilderness, but, like its divine founder on the third day after his crucifixion, it must rise again, take the reins of the world into its hands, and, as an all-transforming principle, take state, science, and art to itself, to breathe into them a higher life and consecrate them to the service of God. The church, of course, continues to the end a servant, as Christ himself came not to be ministered unto, but to minister; and she must at all times suffer persecution, outwardly or inwardly, from the ungodly world. Yet is she also the bride of the Son of God, therefore of royal blood; and she is to make her purifying and sanctifying influence felt upon all orders of natural life and all forms of human society. And from this influence the state, of course, is not excepted. Union with the state is no more necessarily a profanation of holy things than union with science and art, which, in fact, themselves proceed from God, and must subserve his glory.

On the other hand, the state, as a necessary and divine institution for the protection of person and property, for the administration of law and justice, and for the promotion of earthly weal, could not possibly persist forever in her hostility to Christianity, but must at least allow it a legal existence and free play; and if she would attain a higher development and better answer her moral ends than she could in union with idolatry, she must surrender herself to its influence. The kingdom of the Father, to which the state belongs, is not essentially incompatible with the church, the kingdom of the Son; rather does “the Father draw to the Son,” and the Son leads back to the Father, till God become “all in all.” Henceforth should kings again be nursing fathers, and queens nursing mothers to the church,\footnote{Is. xlix. 23.} and the prophecy begin to be fulfilled: “The kingdoms of this world are become the kingdoms of our Lord and of his Christ, and he shall reign forever and ever.”\footnote{Rev. xi. 15.}

The American separation of church and state, even if regarded as the best settlement of the true relation of the two, is not in the least inconsistent with this view. It is not a return to the pre-Constantinian basis, with its spirit of persecution, but rests upon the mutual reverential recognition and support of the two powers, and must be regarded as the continued result of that mighty revolution of the fourth century.

But the elevation of Christianity as the religion of the state presents also an opposite aspect to our contemplation. It involved great risk of degeneracy to the church. The Roman state, with its laws, institutions, and usages, was still deeply rooted in heathenism, and could not be transformed by a magical stroke. The christianizing of the state amounted therefore in great measure to a paganizing and secularizing of the church. The world overcame the church, as much as the church overcame the world, and the temporal gain of Christianity was in many respects cancelled by spiritual loss. The mass of the Roman empire was baptized only with water, not with the Spirit and fire of the gospel, and it smuggled heathen manners and practices into the sanctuary under a new name. The very combination of the cross with the military ensign by Constantine was a most doubtful omen, portending an unhappy mixture of the temporal and the spiritual powers, the kingdom which is of the earth, and that which is from heaven. The settlement of the boundary between the two powers, which, with all their unity, remain as essentially distinct as body and soul, law and gospel, was itself a prolific source of errors and vehement strifes about jurisdiction, which stretch through all the middle age, and still repeat themselves in these latest times, save where the amicable American separation has thus far forestalled collision.

Amidst all the bad consequences of the union of church and state, however, we must not forget that the deeper spirit of the gospel has ever reacted against the evils and abuses of it, whether under an imperial pope or a papal emperor, and has preserved its divine power for the salvation of men under every form of constitution. Though standing and working in the world, and in many ways linked with it, yet is Christianity not of the world, but stands above it.

Nor must we think the degeneracy of the church began with her union with the state.\footnote{This view is now very prevalent in America. It was not formerly so. Jonathan Edwards, in his “History of Redemption,” a practical and edifying survey of church history as an unfolding of the plan of redemption, even saw in the accession of Constantine a type of the future appearing of Christ in the clouds for the redemption of his people, and attributed to it the most beneficent results; to wit: ”(1) The Christian church was thereby wholly delivered from persecution .... (2) God now appeared to execute terrible judgments on their enemies .... (3) Heathenism now was in a great measure abolished throughout the Roman empire .... (4) The Christian church was brought into a state of great peace and prosperity.” ... “This revolution,” he further says, p. 312, “was the greatest that had occurred since the flood. Satan, the prince of darkness, that king and god of the heathen world, was cast out. The roaring lion was conquered by the Lamb of God in the strongest dominion he ever had. This was a remarkable accomplishment of Jerem. x. 11: ’The gods that have not made the heaven and the earth, even they shall perish from the earth and from the heavens.’ ” This work, still much read in America and England, was written, to be sure, Iong before the separation of church and state in New England, viz., in 1739 (first printed in Edinburgh in 1774, twenty-six years after the author’s death). But the great difference of the judgment of this renowned Puritan divine from the prevailing American opinion of the present day is an interesting proof that our view of history is very much determined by the ecclesiastical circumstances in which we live, and at the same time that the whole question of church and state is not at all essential in Christian theology and ethics. In America all confessions, even the Roman Catholics, are satisfied with the separation, while in Europe with few exceptions it is the reverse.} Corruption and apostasy cannot attach to any one fact or personage, be he Constantine or Gregory I. or Gregory VII. They are rooted in the natural heart of man. They revealed themselves, at least in the germ, even in the apostolic age, and are by no means avoided, as the condition of America proves, by the separation of the two powers. We have among ourselves almost all the errors and abuses of the old world, not collected indeed in any one communion, but distributed among our various denominations and sects. The history of the church presents from the beginning a twofold development of good and of evil, an incessant antagonism of light and darkness, truth and falsehood, the mystery of godliness and the mystery of iniquity, Christianity and Antichrist. According to the Lord’s parables of the net and of the tares among the wheat, we cannot expect a complete separation before the final judgment, though in a relative sense the history of the church is a progressive judgment of the church, as the history of the world is a judgment of the world.

\xsection{Rights and Privileges of the Church. Secular Advantages}

§ 14. Rights and Privileges of the Church. Secular Advantages.

The conversion of Constantine and the gradual establishment of Christianity as the religion of the state had first of all the important effect of giving the church not only the usual rights of a legal corporation, which she possesses also in America, and here without distinction of confessions, but at the same time the peculiar privileges, which the heathen worship and priesthood had heretofore enjoyed. These rights and privileges she gradually secured either by tacit concession or through special laws of the Christian emperors as laid down in the collections of the Theodosian and Justinian Codes.\footnote{Comp. § 18.} These were limited, however, as we must here at the outset observe, exclusively to the catholic or orthodox church.\footnote{So early as 326 Constantine promulgated the law (Cod. Theodos. lib. xvi. tit. 5, l. 1): “Privilegia, quae contemplatione religionis indulta sunt, \emph{catholicae tantum legis observatoribus} prodesse oportet. Haereticos autem atque schismaticos non tantum ab his privilegiis alienos esse volumus, sed etiam diversis muneribus constringi et subjici.” Yet he was lenient towards the Novatians, adding in the same year respecting them (C. Theodos. xvi. 5, 2): “Novatianos non adeo comperimus praedamnatos, ut iis quae petiverunt, crederemus minime largienda. Itaque ecclesiae suae domos, et loca sepulcris apta sine inquietudine eos firmiter possidere praecipimus.” Comp. the 8th canon of the Council of Nice, which likewise deals with them indulgently.} The heretical and schismatic sects without distinction, excepting the Arians during their brief ascendency under Arian emperors, were now worse off than they had been before, and were forbidden the free exercise of their worship even under Constantine upon pain of fines and confiscation, and from the time of Theodosius and Justinian upon pain of death. Equal patronage of all Christian parties was totally foreign to the despotic uniformity system of the Byzantine emperors and the ecclesiastical exclusiveness and absolutism of the popes. Nor can it be at all consistently carried out upon the state-church basis; for every concession to dissenters loosens the bond between the church and the state.

The immunities and privileges, which were conferred upon the catholic church in the Roman empire from the time of Constantine by imperial legislation, may be specified as follows:

1. The exemption of the clergy from most public burdens.

Among these were obligatory public services,\footnote{The munera publica, or \textgreek{λειτουργίαι}, attaching in part to the person as a subject of the empire, in part to the possession of property (munera patrimoniorum).} such as military duty, low manual labor, the bearing of costly dignities, and in a measure taxes for the real estate of the church. The exemption,\footnote{Immunitas, \textgreek{ἀλειτουργησία}.} which had been enjoyed, indeed, not by the heathen priests alone, but at least partially by physicians also and rhetoricians, and the Jewish rulers of synagogues, was first granted by Constantine in the year 313 to the catholic clergy in Africa, and afterwards, in 319, extended throughout the empire. But this led many to press into the clerical office without inward call, to the prejudice of the state; and in 320 the emperor made a law prohibiting the wealthy\footnote{The decuriones and curiales.} from entering the ministry, and limiting the increase of the clergy, on the singular ground, that “the rich should bear the burdens of the world, the poor be supported by the property of the church.” Valentinian I. issued a similar law in 364. Under Valentinian II. and Theodosius I. the rich were admitted to the spiritual office on condition of assigning their property to others, who should fulfill the demands of the state in their stead. But these arbitrary laws were certainly not strictly observed.

Constantine also exempted the church from the land tax, but afterwards revoked this immunity; and his successors likewise were not uniform in this matter. Ambrose, though one of the strongest advocates of the rights of the church, accedes to the fact and the justice of the assessment of church lands;\footnote{“Si tributum petit Imperator,” says he in the Orat. de basilicas non tradendis haereticis, “non negamus; agri ecclesiae solvunt tributum, solvimus quae sunt Caesaris Caesari, et qum sunt Dei Deo; tributum Caesaris est; non negatur.” Baronius (ad ann. 387) endeavors to prove that this tribute was meant by Ambrose merely as an act of love, not of duty!} but the hierarchy afterwards claimed for the church a divine right of exemption from all taxation.

2. The enrichment and endowment of the church.

Here again Constantine led the way. He not only restored (in 313) the buildings and estates, which had been confiscated in the Diocletian persecution, but granted the church also the right to receive legacies (321), and himself made liberal contributions in money and grain to the support of the clergy and the building of churches in Africa,\footnote{So early as 314 he caused to be paid to the bishop Caecilian of Carthage 3,000 \emph{folles} (\textgreek{τρισχιλίους φόλεις}£18,000) from the public treasury of the province for the catholic churches in Africa, Numidia, and Mauritania, promising further gifts for similar purposes. Euseb: H. E. x. 6, and Vit. Const. iv. 28.} in the Holy Land, in Nicomedia, Antioch, and Constantinople. Though this, be it remembered, can be no great merit in an absolute monarch, who is lord of the public treasury as he is of his private purse, and can afford to be generous at the expense of his subjects. He and his successors likewise gave to the church the heathen temples and their estates and the public property of heretics; but these more frequently were confiscated to the civil treasury or squandered on favorites. Wealthy subjects, some from pure piety, others from motives of interest, conveyed their property to the church, often to the prejudice of the just claims of their kindred. Bishops and monks not rarely used unworthy influences with widows and dying persons; though Augustine positively rejected every legacy, which deprived a son of his rights. Valentinian I. found it necessary to oppose the legacy-hunting of the clergy, particularly in Rome, with a law of the year 370,\footnote{In an edict to Damasus, bishop of Rome. Cod. Theod. xvi. 2, 20: “Ecclesiastici ... viduaram ac pupillarum domos non adeant,” etc.} and Jerome acknowledges there was good reason for it.\footnote{Epist. 34 (al. 2) ad Nepotianum, where he says of this law: “Nec de lege conqueror, sed doleo, cur meruerimus hanc legem;” and of the clergy of his time: “Ignominia omnium sacerdotum est, propriis studere divitiis,” etc.} The wealth of the church was converted mostly into real estate, or at least secured by it. And the church soon came to own the tenth part of all the landed property. This land, to be sure, had long been worthless or neglected, but under favorable conditions rose in value with uncommon rapidity. At the time of Chrysostom, towards the close of the fourth century, the church of Antioch was strong enough to maintain entirely or in part three thousand widows and consecrated virgins besides many poor, sick, and strangers.\footnote{Chrys. Hom. 66 in Matt. (vii. p. 658).} The metropolitan churches of Rome and Alexandria were the most wealthy. The various churches of Rome in the sixth century, besides enormous treasures in money and gold and silver vases, owned many houses and lands not only in Italy and Sicily, but even in Syria, Asia Minor, and Egypt.\footnote{Comp. the Epistles of Gregory the Great at the end of our period.} And when John, who bears the honorable distinction of the Almsgiver for his unlimited liberality to the poor, became patriarch of Alexandria (606), he found in the church treasury eight thousand pounds of gold, and himself received ten thousand, though be retained hardly an ordinary blanket for himself, and is said on one occasion to have fed seven thousand five hundred poor at once.\footnote{See the Vita S. Joannis Eleemosynarii (the next to the last catholic patriarch of Alexandria) in the Acta Sanct. Bolland. ad 23 Jan.}

The control of the ecclesiastical revenues vested in the bishops. The bishops distributed the funds according, to the prevailing custom into three or four parts: for themselves, for their clergy, for the current expenses of worship, and for the poor. They frequently exposed themselves to the suspicion of avarice and nepotism. The best of them, like Chrysostom and Augustine, were averse to this concernment with earthly property, since it often conflicted with their higher duties; and they preferred the poverty of earlier times, because the present abundant revenues diminished private beneficence.

And most certainly this opulence had two sides. It was a source both of profit and of loss to the church. According to the spirit of its proprietors and its controllers, it might be used for the furtherance of the kingdom of God, the building of churches, the support of the needy, and the founding of charitable institutions for the poor, the sick, for widows and orphans, for destitute strangers and aged persons,\footnote{The \textgreek{πτωχοτροφεῖα, νοσοκομεῖα, ὀρφανοτροφεῖα, γηροκομεῖα} and \textgreek{ξενῶνες} or \textgreek{ξενοδοχεῖα}, as they were called; which all sprang from the church. Especially favored was the \emph{Basilias} for sick and strangers in Caesarea, named after its founder, the bishop Basil the Great. Basil. Ep. 94. Gregor. Naz. Orat. 27 and 30.} or perverted to the fostering of indolence and luxury, and thus promote moral corruption and decay. This was felt by serious minds even in the palmy days of the external power of the hierarchy. Dante, believing Constantine to be the author of the pope’s temporal sovereignty, on the ground of the fictitious donation to Sylvester, bitterly exclaimed:

\xsection{Support of the Clergy}

§ 15. Support of the Clergy.

3. The better support of the clergy was another advantage connected with the new position of Christianity in the empire.

Hitherto the clergy had been entirely dependent on the voluntary contributions of the Christians, and the Christians were for the most part poor. Now they received a fixed income from the church funds and from imperial and municipal treasuries. To this was added the contribution of first-fruits and tithes, which, though not as yet legally enforced, arose as a voluntary custom at a very early period, and probably in churches of Jewish origin existed from the first, after the example of the Jewish law.\footnote{Lev. xxvii. 30-33; Nu. xviii. 20-24; Deut. xiv. 22 sqq. 2 Chron. xxxi. 4 sqq.} Where these means of support were not sufficient, the clergy turned to agriculture or some other occupation; and so late as the fifth century many synods recommended this means of subsistence, although the Apostolical Canons prohibited the engagement of the clergy in secular callings under penalty of deposition.\footnote{. Constit. Apost. lib. viii. cap. 47, can. 6 (p. 239, ed. Ueltzen): \textgreek{Ἐπίσκοπος ἢ πρεσβύτερος ἢ διάκονος κοσμικὰς φροντίδας μὴ ἀναλαμβανέτο· εἰ δὲ μὴ, καθαιρείσθω.}}

This improvement, also, in the external condition of the clergy was often attended with a proportional degeneracy in their moral character. It raised them above oppressive and distracting cares for livelihood, made them independent, and permitted them to devote their whole strength to the duties of their office; but it also favored ease and luxury, allured a host of unworthy persons into the service of the church, and checked the exercise of free giving among the people. The better bishops, like Athanasius, the two Gregories, Basil, Chrysosotom, Theodoret, Ambrose, Augustine, lived in ascetic simplicity, and used their revenues for the public good; while others indulged their vanity, their love of magnificence, and their voluptuousness. The heathen historian Ammianus gives the country clergy in general the credit of simplicity, temperance, and virtue, while he represents the Roman hierarchy, greatly enriched by the gifts of matrons, as extreme in the luxury of their dress and their more than royal banquets;\footnote{Lib. xxvii. c. 3.} and St. Jerome agrees with him.\footnote{Hieron. Ep. 34 (al. 2) et passim.} The distinguished heathen prefect, Praetextatus, said to Pope Damasus, that for the price of the bishopric of Rome he himself might become a Christian at once. The bishops of Constantinople, according to the account of Gregory Nazianzen,\footnote{Orat. 32.} who himself held that see for a short time, were not behind their Roman colleagues in this extravagance, and vied with the most honorable functionaries of the state in pomp and sumptuous diet. The cathedrals of Constantinople and Carthage had hundreds of priests, deacons, deaconesses, subdeacons, prelectors, singers, and janitors.\footnote{The cathedral of Constantinople fell under censure for the excessive number of its clergy and subordinate officers, so that Justinian reduced it to five hundred and twenty-five, of which probably more than half were useless. Comp. Iust. Novell. ciii.}

It is worthy of notice, that, as we have already intimated, the two greatest church fathers gave the preference in principle to the voluntary system in the support of the church and the ministry, which prevailed before the Nicene era, and which has been restored in modern times in the United States of America. Chrysostom no doubt perceived that under existing circumstances the wants of the church could not well be otherwise supplied, but he was decidedly averse to the accumulation of treasure by the church, and said to his hearers in Antioch: “The treasure of the church should be with you all, and it is only your hardness of heart that requires her to hold earthly property and to deal in houses and lands. Ye are unfruitful in good works, and so the ministers of God must meddle in a thousand matters foreign to their office. In the days of the apostles people might likewise have given them houses and lands; why did they prefer to sell the houses and lands and give the proceeds? Because this was without doubt the better way. Your fathers would have preferred that you should give alms of your incomes, but they feared that your avarice might leave the poor to hunger; hence the present order of things.”\footnote{Homil. 85 in Matt. (vii. 808 sq.). Hom. 21 in 1 Cor. 7 (x. 190). Comp. also De sacerdot. l. iii. c. 16.} Augustine desired that his people in Hippo should take back the church property and support the clergy and the poor by free gifts.\footnote{Possidius, in Vita Aug. c. 23: “Alloquebatur plebem Dei, malle se ex collationibus plebes Dei vivere quam illarum possessionum curam vel gubernationem pati, et paratum se esse illis cedere, ut eo modo omnes Dei servi et ministri viverent.”}

\xsection{Episcopal Jurisdiction and Intercession}

§ 16. Episcopal Jurisdiction and Intercession.

4. We proceed to the legal validity, of the episcopal jurisdiction, which likewise dates from the time of Constantine.

After the manner of the Jewish synagogues, and according to the exhortation of St. Paul,\footnote{1 Cor. vi. 1-6.} the Christians were accustomed from the beginning to settle their controversies before the church, rather than carry them before heathen tribunals; but down to the time of Constantine the validity, of the bishop’s decision depended on the voluntary, submission of both parties. Now this decision was invested with the force of law, and in spiritual matters no appeal could be taken from it to the civil court. Constantine himself, so early as 314, rejected such an appeal in the Donatist controversy with the significant declaration: “The judgment of the priests must be regarded as the judgment of Christ himself.”\footnote{“Sacerdotum judicium ita debet haberi, ut si ipse Dominus residens judicet. Optatus Milev.: De schism. Donat. f. 184.} Even a sentence of excommunication was final; and Justinian allowed appeal only to the metropolitan, not to the civil tribunal. Several councils, that of Chalcedon, for example, in 451, went so far as to threaten clergy, who should avoid the episcopal tribunal or appeal from it to the civil, with deposition. Sometimes the bishops called in the help of the state, where the offender contemned the censure of the church. Justinian I. extended the episcopal jurisdiction also to the monasteries. Heraclius subsequently (628) referred even criminal causes among the clergy to the bishops, thus dismissing the clergy thenceforth entirely from the secular courts; though of course holding them liable for the physical penalty, when convicted of capital crime,\footnote{Even Constantine, however, before the council of Nice, had declared, that should he himself detect a bishop in the act of adultery, he would rather throw over him his imperial mantle than bring scandal on the church by punishing a clergyman.} as the ecclesiastical jurisdiction ended with deposition and excommunication. Another privilege, granted by Theodosius to the clergy, was, that they should not be compelled by torture to bear testimony before the civil tribunal.

This elevation of the power and influence of the bishops was a salutary check upon the jurisdiction of the state, and on the whole conduced to the interests of justice and humanity; though it also nourished hierarchical arrogance and entangled the bishops, to the prejudice of their higher functions, in all manner of secular suits, in which they were frequently called into consultation. Chrysostom complains that “the arbitrator undergoes incalculable vexations, much labor, and more difficulties than the public judge. It is hard to discover the right, but harder not to violate it when discovered. Not labor and difficulty alone are connected with office, but also no little danger.”\footnote{De sacerd. l. iii. c. 18, at the beginning.} Augustine, too, who could make better use of his time, felt this part of his official duty a burden, which nevertheless he bore for love to the church.\footnote{In Psalm. xxv. (vol. iv. 115) and Epist. 213, where he complains that before and after noon he was beset and distracted by the members of his church with temporal concerns, though they had promised to leave him undisturbed five days in the week, to finish some theological labors. Comp. Neander, iii. 291 sq. (ed. Torrey, ii. 139 sq.).} Others handed over these matters to a subordinate ecclesiastic, or even, like Silvanus, bishop of Troas, to a layman.\footnote{Socrat. l. vii. c. 37.}

5. Another advantage resulting from the alliance of the church with the empire was the episcopal right of intercession.

The privilege of interceding with the secular power for criminals, prisoners, and unfortunates of every kind had belonged to the heathen priests, and especially to the vestals, and now passed to the Christian ministry, above all to the bishops, and thenceforth became an essential function of their office. A church in Gaul about the year 460 opposed the ordination of a monk to the bishopric, because, being unaccustomed to intercourse with secular magistrates, though he might intercede with the Heavenly Judge for their souls, he could not with the earthly for their bodies. The bishops were regarded particularly as the guardians of widows and orphans, and the control of their property was intrusted to them. Justinian in 529 assigned to them also a supervision of the prisons, which they were to visit on Wednesdays and Fridays, the days of Christ’s passion.

The exercise of this right of intercession, one may well suppose, often obstructed the course of justice; but it also, in innumerable cases, especially in times of cruel, arbitrary despotism, protected the interests of innocence, humanity, and mercy. Sometimes, by the powerful pleadings of bishops with governors and emperors, whole provinces were rescued from oppressive taxation and from the revenge of conquerors. Thus Flavian of Antioch in 387 averted the wrath of Theodosius on occasion of a rebellion, journeying under the double burden of age and sickness even to Constantinople to the emperor himself, and with complete success, as an ambassador of their common Lord, reminding him of the words: “If ye forgive men their trespasses, your heavenly Father will also forgive you.”\footnote{Matt. vi. 14.}

6. With the right of intercession was closely connected the right of asylum in churches.

In former times many of the heathen temples and altars, with some exceptions, were held inviolable as places of refuge; and the Christian churches now inherited also this prerogative. The usage, with some precautions against abuse, was made law by Theodosius II. in 431, and the ill treatment of an unarmed fugitive in any part of the church edifice, or even upon the consecrated ground, was threatened with the penalty of death.\footnote{Cod. Theodos. ix. 45, 1-4. Comp. Socrat. vii. 33.}

Thus slaves found sure refuge from the rage of their masters, debtors from the persecution of inexorable creditors, women and virgins from the approaches of profligates, the conquered from the sword of their enemies, in the holy places, until the bishop by his powerful mediation could procure justice or mercy. The beneficence of this law, which had its root not in superstition alone, but in the nobler sympathies of the people, comes most impressively to view amidst the ragings of the great migration and of the frequent intestine wars.\footnote{“The rash violence of despotism,” says even Gibbon, “was suspended by the mild interposition of the church; and the lives or fortunes of the most eminent subjects might be protected by the mediation of the bishop.”}

\xsection{Legal Sanction of Sunday}

§ 17. Legal Sanction of Sunday.

7. The civil sanction of the observance of Sunday and other festivals of the church.

The state, indeed, should not and cannot enforce this observance upon any one, but may undoubtedly and should prohibit the public disturbance and profanation of the Christian Sabbath, and protect the Christians in their right and duty of its proper observance. Constantine in 321 forbade the sitting of courts and all secular labor in towns on “the venerable day of the sun,” as he expresses himself, perhaps with reference at once to the sun-god, Apollo, and to Christ, the true Sun of righteousness; to his pagan and his Christian subjects. But he distinctly permitted the culture of farms and vineyards in the country, because frequently this could be attended to on no other day so well;\footnote{This exception is entirely unnoticed by many church histories, but stands in the same law of 321 in the Cod. Justin. lib. iii. tit. 12, de feriis, l. 3: “Omnes judices, urbanaeque plebes, et cunctarum artium officia venerabili die Solis quiescant. Ruri tamen positi agrorum culturae libere licenterque inserviant: quoniam frequenter evenit, ut non aptius alio die frumenta sulcis, aut vineae scrobibus mandentur, ne occasione momenti pereat commoditas coelesti provisione concessa.” Such work was formerly permitted, too, on the pagan feast days. Comp. Virgil. Georg. i. v. 268 sqq. Cato, De re rust. c. 2.} though one would suppose that the hard-working peasantry were the very ones who most needed the day of rest. Soon afterward, in June, 321, he allowed the manumission of slaves on Sunday;\footnote{Cod. Theodos. lib. ii. tit. 8. l. 1: “Emancipandi et manumittendi die festo cuncti licentiam habeant, et super his rebus actus non prohibeantur.”} as this, being an act of benevolence, was different from ordinary business, and might be altogether appropriate to the day of resurrection and redemption. According to Eusebius, Constantine also prohibited all military exercises on Sunday, and at the same time enjoined the observance of Friday in memory of the death of Christ.\footnote{Eus. Vit. Const. iv. 18-20. Comp. Sozom. i. 8. In our times military parades and theatrical exhibitions in Paris, Vienna, Berlin, and other European cities are so frequent on no other day as on the Lord’s day! In France, political elections are usually held on the Sabbath!}

Nay, he went so far, in well-meaning but mistaken zeal, as to require of his soldiers, even the pagan ones, the positive observance of Sunday, by pronouncing at a signal the following prayer, which they mechanically learned: “Thee alone we acknowledge as God; thee we confess as king; to thee we call as our helper; from thee we have received victories; through thee we have conquered enemies. Thee we thank for good received; from thee we hope for good to come. Thee we all most humbly beseech to keep our Constantine and his God-fearing sons through long life healthy and victorious.”\footnote{Eus. Vit. Const. l. iv. c. 20. The formulary was prescribed in the Latin language, as Eusebius says in c. 19. He is speaking of the whole army (comp. c. 18), and it may presumed that many of the soldiers were heathen.} Though this formula was held in a deistical generalness, yet the legal injunction of it lay clearly beyond the province of the civil power, trespassed on the rights of conscience, and unavoidably encouraged hypocrisy and empty formalism.

Later emperors declared the profanation of Sunday to be sacrilege, and prohibited also the collecting of taxes and private debts (368 and 386), and even theatrical and circus performances, on Sunday and the high festivals (386 and 425).\footnote{The second law against opening theatres on Sundays and festivals (a.d.425) in the Cod. Theodos. l. xv. tit. 7, I. 5, says expressly: “Omni theatrorum atque circensium voluptate per universas urbes ... denegata, totae Christianorum ac fidelium mentes Dei cultibus occupentur.”} But this interdiction of public amusements, on which a council of Carthage (399 or 401) with reason insisted, was probably never rigidly enforced, and was repeatedly supplanted by the opposite practice, which gradually prevailed all over Europe.\footnote{As Chrysostom, at the end of the fourth century and the beginning of the fifth, often complains that the theatre is better attended than the church; so down to this day the same is true in almost all the large cities on the continent of Europe. Only in England and the United States, under the influence of Calvinism and Puritanism, are the theatres closed on Sunday.}

\xsection{Influence of Christianity on Civil Legislation. The Justinian Code}

§ 18. Influence of Christianity on Civil Legislation. The Justinian Code.

Comp. on this subject particularly the works cited at § 13, sub ii, by Rhoer, Meysenburg, and Troplong; also Gibbon, chap. xliv (an admirable summary of the Roman law), Milman: Lat. Christianity, vol. I. B. iii. chap. 5, and in part the works of Schmidt and Chastel on the influence of Christianity upon society in the Roman empire, quoted in vol. i. § 86.

While in this way the state secured to the church the well-deserved rights of a legal corporation, the church exerted in turn a most beneficent influence on the state, liberating it by degrees from the power of heathen laws and customs, from the spirit of egotism, revenge, and retaliation, and extending its care beyond mere material prosperity to the higher moral interests of society. In the previous period we observed the contrast between Christian morality and heathen corruption in the Roman empire.\footnote{Vol. i §§ 86-93.} We are now to see how the principles of Christian morality gained public recognition, and began at least in some degree to rule the civil and political life.

As early as the second century, under the better heathen emperors, and evidently under the indirect, struggling, yet irresistible influence of the Christian spirit, legislation took a reformatory, humane turn, which was carried by the Christian emperors as far as it could be carried on the basis of the ancient Graeco-Roman civilization. Now, above all, the principle of justice and equity, humanity and love, began to assert itself in the state. For Christianity, with its doctrines of man’s likeness to God, of the infinite value of personality, of the original unity of the human race, and of the common redemption through Christ, first brought the universal rights of man to bear in opposition to the exclusive national spirit, the heartless selfishness, and the political absolutism of the old world, which harshly separated nations and classes, and respected man only as a citizen, while at the same time it denied the right of citizenship to the great mass of slaves, foreigners, and barbarians.\footnote{Comp. Lactantius: Inst. divin. l. v. c. 15.}

Christ himself began his reformation with the lowest orders of the people, with fishermen and taxgatherers, with the poor, the lame, the blind, with demoniacs and sufferers of every kind, and raised them first to the sense of their dignity and their high destiny. So now the church wrought in the state and through the state for the elevation of the oppressed and the needy, and of those classes which under the reign of heathenism were not reckoned at all in the body politic, but were heartlessly trodden under foot. The reformatory motion was thwarted, it is true, to a considerable extent, by popular custom, which is stronger than law, and by the structure of society in the Roman empire, which was still essentially heathen and doomed to dissolution. But reform was at last set in motion, and could not be turned back even by the overthrow of the empire; it propagated itself among the German tribes. And although even in Christian states the old social maladies are ever breaking forth from corrupt human nature, sometimes with the violence of revolution, Christianity is ever coming in to restrain, to purify, to heal, and to console, curbing the wild passions of tyrants and of populace, vindicating the persecuted, mitigating the horrors of war, and repressing incalculable vice in public and in private life among Christian people. The most cursory comparison of Christendom with the most civilized heathen and Mohammedan countries affords ample testimony of this.

Here again the reign of Constantine is a turning point. Though an oriental despot, and but imperfectly possessed with the earnestness of Christian morality, he nevertheless enacted many laws, which distinctly breathe the spirit of Christian justice and humanity: the abolition of the punishment of crucifixion, the prohibition of gladiatorial games and cruel rites, the discouragement of infanticide, and the encouragement of the emancipation of slaves. Eusebius says he improved most of the old laws or replaced them by new ones.\footnote{Vit. Const. l. iv. c. 26, where the most important laws of Constantine are recapitulated. Even the heathen Libanius (Basil. ii. p. 146) records that under Constantine and his sons legislation was much more favorable to the lower classes: though he accounts for this only by the personal clemency of the emperors.} Henceforward we feel beneath the toga of the Roman lawgiver the warmth of a Christian heart. We perceive the influence of the evangelical preaching and exhortations of the father of monasticism out of the Egyptian desert to the rulers of the world, Constantine and his sons: that they should show justice and mercy to the poor, and remember the judgment to come.

Even Julian, with all his hatred of the Christians, could not entirely renounce the influence of his education and of the reigning spirit of the age, but had to borrow from the church many of his measures for the reformation of heathenism. He recognized especially the duty of benevolence toward all men, charity to the poor, and clemency to prisoners; though this was contrary to the heathen sentiment, and though he proved himself anything but benevolent toward the Christians. But then the total failure of his philanthropic plans and measures shows that the true love for man can thrive only in Christian soil. And it is remarkable, that, with all this involuntary concession to Christianity, Julian himself passed not a single law in line with the progress of natural rights and equity.\footnote{Troplong, p. 127. C. Schmidt, 378.}

His successors trod in the footsteps of Constantine, and to the end of the West Roman empire kept the civil legislation under the influence of the Christian spirit, though thus often occasioning conflicts with the still lingering heathen element, and sometimes temporary apostasy and reaction. We observe also, in remarkable contradiction, that while the laws were milder in some respects, they were in others even more severe and bloody than ever before: a paradox to be explained no doubt in part by the despotic character of the Byzantine government, and in part by the disorders of the time.\footnote{Comp. de Rhoer, p. 59 sqq. The origin of this increased severity of penal laws is, at all events, not to be sought in the church; for in the fourth and fifth centuries she was still rather averse to the death penalty. Comp. Ambros. Ep. 25 and 26 (al. 51 and 52), and Augustine, Ep. 153 ad Macedonium.}

It now became necessary to collect the imperial ordinances\footnote{Constitutiones or Leges. If answers to questions, they were called Rescripta; if spontaneous decrees, Edicta.} in a codex or corpus juris. Of the first two attempts of this kind, made in the middle of the fourth century, only some fragments remain.\footnote{The Codex Gregorianus and Codex Hermogenianus; so called from the compilers, two private lawyers. They contained the rescripts and edicts of the heathen emperors from Hadrian to Constantine, and would facilitate a comparison of the heathen legislation with the Christian.} But we have the Codex Theodosianus, which Theodosius II. caused to be made by several jurists between the years 429 and 438. It contains the laws of the Christian emperors from Constantine down, adulterated with many heathen elements; and it was sanctioned by Valentinian III. for the western empire. A hundred years later, in the flourishing period of the Byzantine state-church despotism, Justinian I., who, by the way, cannot be acquitted of the reproach of capricious and fickle law-making, committed to a number of lawyers, under the direction of the renowned Tribonianus,\footnote{Tribonianus, a native of Side in Paphlagonia, was an advocate and a poet, and rose by his talents, and the favor of Justinian, to be quaestor, consul, and at last magister officiorum. Gibbon compares him, both for his comprehensive learning and administrative ability and for his enormous avarice and venality, with Lord Bacon. But in one point these statesmen were very different: while Bacon was a decided Christian in his convictions, Tribonianus was accused of pagan proclivities and of atheism. In a popular tumult in Constantinople the emperor was obliged to dismiss him, but found him indispensable and soon restored him.} the great task of making a complete revised and digested collection of the Roman law from the time of Hadrian to his own reign; and thus arose, in the short period of seven years (527–534), through the combination of the best talent and the best facilities, the celebrated Codex Justinianeus, which thenceforth became the universal law of the Roman empire, the sole text book in the academies at Rome, Constantinople, and Berytus, and the basis of the legal relations of the greater part of Christian Europe to this day.\footnote{The complete \emph{Codex Justinianeus}, which has long outlasted the conquests of that emperor (as Napoleon’s Code has outlasted his), comprises properly three separate works: (1) The \emph{Institutiones}, an elementary text book of jurisprudence, of the year 533. (2) The \emph{Digesta} or \emph{Pandectae} (\textgreek{πάνδεκται}, complete repository), an abstract of the spirit of the whole Roman jurisprudence, according to the decisions of the most distinguished jurists of the earlier times, composed in 530-533. (3) The Codex, first prepared in 528 and 529, but in 534 reconstructed, enlarged, and improved, and hence called \emph{Codex repetitae praelectionis}; containing 4,648 ordinances in 765 titles, in chronological order. To these is added (4) a later Appendix: \emph{Novellae constitutiones} (v\textgreek{εαραὶ διατάξεις}), or simply \emph{Novellae} (a barbarism); that is, 168 decrees of Justinian, subsequently collected from the 1st January, 535, to his death in 565, mostly in Greek, or in both Greek and Latin. Excepting some of the novels of Justinian, the codex was composed in the Latin language, which Justinian and Tribonianus understood; but afterward, as this tongue died out in the East, it was translated into Greek, and sanctioned in this form by the emperor Phocas in 600. The emperor Basil the Macedonian in 876 caused a Greek abstract (\textgreek{πρόχειρον τῶν νόμων}\emph{)} to be prepared, which, under the name of the \emph{Basilicae}, gradually supplanted the book of Justinian in the Byzantine empire. The Pandects have narrowly escaped destruction. Most of the editions and manuscripts of the west (not all, as Gibbon says) are taken from the Codex Florentinus, which was transcribed in the beginning of the seventh century at Constantinople, and afterward carried by the vissitudes of war and trade to Amalfi, to Pisa, and in 1411 to Florence.}

This body of Roman law\footnote{Called \emph{Corpus juris} \emph{Romani}or \emph{C. juris civilis}, in distinction from \emph{Corpus juris canonici}, the Roman Catholic church law, which is based chiefly on the canons of the ancient councils, as the civil law is upon the rescripts and edicts of the emperors.} is an important source of our knowledge of the Christian life in its relations to the state and its influence upon it. It is, to be sure, in great part the legacy of pagan Rome, which was constitutionally endowed with legislative and administrative genius, and thereby as it were predestined to universal empire. But it received essential modification through the orientalizing change in the character of the empire from the time of Constantine, through the infusion of various Germanic elements, through the influence of the law of Moses, and, in its best points, through the spirit of Christianity. The church it fully recognizes as a legitimate institution and of divine authority, and several of its laws were enacted at the direct instance of bishops. So the “Common Law,” the unwritten traditional law of England and America, though descending from the Anglo-Saxon times, therefore from heathen Germandom, has ripened under the influence of Christianity and the church, and betrays this influence even far more plainly than the Roman code, especially in all that regards the individual and personal rights and liberties of man.

\xsection{Elevation of Woman and the Family}

§ 19. Elevation of Woman and the Family.

The benign effect of Christianity on legislation in the Graeco-Roman empire is especially noticeable in the following points:

1. In the treatment of women. From the beginning, Christianity labored, primarily in the silent way of fact, for the elevation of the female sex from the degraded, slavish position, which it occupied in the heathen world;\footnote{On this subject, and on the heathen family life, comp. vol. i. § 91.} and even in this period it produced such illustrious models of female virtue as Nonna, Anthusa, and Monica, who commanded the highest respect of the heathens themselves. The Christian emperors pursued this work, though the Roman legislation stops considerably short of the later Germanic in regard to the rights of woman. Constantine in 321 granted women the same right as men to control their property, except in the sale of their landed estates. At the same time, from regard to their modesty, he prohibited the summoning them in person before the public tribunal. Theodosius I. in 390 was the first to allow the mother a certain right of guardianship, which had formerly been intrusted exclusively to men. Theodosius II. in 439 interdicted, but unfortunately with little success, the scandalous trade of the lenones, who lived by the prostitution of women, and paid a considerable license tax to the state.\footnote{Cod. Theod. lib. xv. tit. 8: de lenonibus.} Woman received protection in various ways against the beastly passion of man. The rape of consecrated virgins and widows was punishable, from the time of Constantine, with death.\footnote{C. Theod. ix. 24: de raptu virginum et viduarum (probably nuns and deaconesses).}

2. In the marriage laws, Constantine gave marriage its due freedom by abolishing the old Roman penalties against celibacy and childlessness.\footnote{C. Theod. viii. 16, 1. Comp. Euseb. Vit. Const. iv. 26.} On the other hand, marriage now came to be restricted under heavy penalties by the introduction of the Old Testament prohibitions of marriage within certain degrees of consanguinity, which subsequently were arbitrarily extended even to the relation of cousin down to the third remove.\footnote{C. Theod. iii. 12: de incestis nuptiis.} Justinian forbade also marriage between godparent and godchild, on the ground of spiritual kinship. But better than all, the dignity and sanctity of marriage were now protected by restrictions upon the boundless liberty of divorce which had obtained from the time of Augustus, and had vastly hastened the decay of public morals. Still, the strict view of the fathers, who, following the word of Christ, recognized adultery alone as a sufficient ground of divorce, could not be carried out in the state.\footnote{C. Theod. iii. 16: de repudiis. Hence Jeromesays in view of this, Ep. 30 (al. 84) ad Oceanum: “Aliae sunt leges Caesarum, aliae Christi; aliud Papinianus [the most celebrated Roman jurist, died a.d.212], aliud Paulus noster praecipit.”} The legislation of the emperors in this matter wavered between the licentiousness of Rome and the doctrine of the church. So late as the fifth century we hear a Christian author complain that men exchange wives as they would garments, and that the bridal chamber is exposed to sale like a shoe on the market! Justinian attempted to bring the public laws up to the wish of the church, but found himself compelled to relax them; and his successor allowed divorce even on the ground of mutual consent.\footnote{Gibbon: “The dignity of marriage was restored by the Christians .... The Christian princes were the first who specified the just causes of a private divorce; their institutions, from Constantineto Justinian, appear to fluctuate between the custom of the empire and the wishes of the church, and the author of the Novels too frequently reforms the jurisprudence of the Code and the Pandects .... The successor of Justinian yielded to the prayers of his unhappy subjects, and restored the liberty of divorce by mutual consent.”}

Concubinage was forbidden from the time of Constantine, and adultery punished as one of the grossest crimes.\footnote{In a law of 326 it is called “facinus atrocissimum, scelus immane.” Cod. Theod. l. ix. tit. 7, 1. 1 sq. And the definition of adultery, too, was now made broader. According to the old Roman law, the idea of adultery on the part of the man was limited to illicit intercourse with the \emph{married} lady of a \emph{free} \emph{citizen}, and was thought punishable not so much for its own sake, as for its encroachment on the rights of another husband. Hence Jerome says, l.c., of the heathen: “Apud illos viris impudicitiae frena laxantur, et solo stupro et adulterio condemnato passim per lupanaria et ancillulas libido permittitur; quasi culpam dignitas faciat, non voluntas. Apud nos quod non licet feminis, aeque non licet viris, et eadem servitus pari conditione censetur.” Yet the law, even under the emperors, still excepted carnal intercourse with a female slave from adultery. Thus the state here also stopped short of the church, and does to this day in countries where the institution of slavery exists.} Yet here also pagan habit ever and anon reacted in practice, and even the law seems to have long tolerated the wild marriage which rested only on mutual agreement, and was entered into without convenant, dowry, or ecclesiastical sanction.\footnote{Even a council at Toledo in 398 conceded so far on this point as to decree, can. 17: “Si quis habens uxorem fidelis concubinam habeat, non communicet. Ceterum is, qui non habet uxorem et pro uxore concubinam habeat, a communione non repellatur, tantum ut unius mulieris aut uxoris aut concubinae, ut ei placuerit, sit conjunctione contentus. Alias vero vivens abjiciatur donec desinat et per poenitentiam, revertatur.”} Solemnization by the church was not required by the state as the condition of a legitimate marriage till the eighth century. Second marriage, also, and mixed marriages with heretics and heathens, continued to be allowed, notwithstanding the disapproval of the stricter church teachers; only marriage with Jews was prohibited, on account of their fanatical hatred of the Christians.\footnote{Cod. Theod. iii. 7, 2; C. Justin. i. 9, 6. A proposal of marriage to a nun was even punished with death (ix. 25, 2).}

3. The power of fathers over their children, which according to the old Roman law extended even to their freedom and life, had been restricted by Alexander Severus under the influence of the monarchical spirit, which is unfavorable to private jurisdiction, and was still further limited under Constantine. This emperor declared the killing of a child by its father, which the Pompeian law left unpunished, to be one of the greatest crimes.\footnote{a.d.318; Valentinian did the same in 374. Cod. Theod. ix. tit. 14 and 15. Comp. the Pandects, lib. xlviii. tit. 8, l ix.} But the cruel and unnatural practice of exposing children and selling them into slavery continued for a long time, especially among the laboring and agricultural classes. Even the indirect measures of Valentinian and Theodosius I. could not eradicate the evil. Theodosius in 391 commanded that children which had been sold as slaves by their father from poverty, should be free, and that without indemnity to the purchasers; and Justinian in 529 gave all exposed children without exception their freedom.\footnote{Cod. Theod. iii. 3, 1; Cod. Just. iv. 43, 1; viii. 52, 3. Gibbon says: “The Roman empire was stained with the blood of infants, till such murders were included, by Valentinian and his colleagues, in the letter and spirit of the Cornelian law. The lessons of jurisprudence and Christianity had been inefficient to eradicate this inhuman practice, till their gentle influence was fortified by the terrors of capital punishment.”}

\xsection{Social Reforms. The Institution of Slavery}

§ 20. Social Reforms. The Institution of Slavery.

4. The institution of slavery\footnote{Comp. vol. i. § 89, and the author’s “Hist. of the Apost. Church,” § 113.} remained throughout the empire, and is recognized in the laws of Justinian as altogether legitimate.\footnote{Instit. lib. i. tit. 5-8; Digest. l. i. tit. 5 and 6, etc.} The Justinian code rests on the broad distinction of the human race into freemen and slaves. It declares, indeed, the natural equality of men, and so far rises above the theory of Aristotle, who regards certain races and classes of men as irrevocably doomed, by their physical and intellectual inferiority, to perpetual servitude; but it destroys the practical value of this concession by insisting as sternly as ever on the inferior legal and social condition of the slave, by degrading his marriage to the disgrace of concubinage, by refusing him all legal remedy in case of adultery, by depriving him of all power over his children, by making him an article of merchandise like irrational beasts of burden, whose transfer from vender to buyer was a legal transaction as valid and frequent as the sale of any other property. The purchase and sale of slaves for from ten to seventy pieces of gold, according to their age, strength, and training, was a daily occurrence.\footnote{The legal price, which, however, was generally under the market price, was thus established under Justinian (Cod. l. vi. tit. xliii. l. 3): Ten pieces of gold for an ordinary male or female slave under ten years; twenty, for slaves over ten; thirty, for such as understood a trade; fifty, for notaries and scribes; sixty, for physicians, and midwives. Eunuchs ranged to seventy pieces.} The number was not limited; many a master owning even two or three thousand slaves.

The barbarian codes do not essentially differ in this respect from the Roman. They, too, recognize slavery as an ordinary condition of mankind and the slave as a marketable commodity. All captives in war became slaves, and thousands of human lives were thus saved from indiscriminate massacre and extermination. The victory of Stilicho over Rhadagaisus threw 200,000 Goths and other Germans into the market, and lowered the price of a slave from twenty-five pieces of gold to one. The capture and sale of men was part of the piratical system along all the shores of Europe. Anglo-Saxons were freely sold in Rome at the time of Gregory the Great. The barbarian codes prohibited as severely as the Justinian code the debasing alliance of the freeman with the slave, but they seem to excel the latter in acknowledging the legality and religious sanctity of marriages between slaves; that of the Lombards on the authority of the Scripture sentence: “Whom God has joined together, let no man put asunder.”

The legal wall of partition, which separated the slaves from free citizens and excluded them from the universal rights of man, was indeed undermined, but by no means broken down, by the ancient church, who taught only the moral and religious equality of men. We find slaveholders even among the bishops and the higher clergy of the empire. Slaves belonged to the papal household at Rome, as we learn incidentally from the acts of a Roman synod held in 501 in consequence of the disputed election of Symmachus, where his opponents insisted upon his slaves being called in as witnesses, while his adherents protested against this extraordinary request, since the civil law excluded the slaves from the right of giving testimony before a court of justice.\footnote{Comp. Hefele: “Conciliengeschichte,” ii. p. 620; and Milman: “Latin Christianity,” vol. i. p. 419 (Am. ed.), who infers from this fact, “that slaves formed the household of the Pope, and that, by law, they were yet liable to torture. This seems clear from the words of Ennodius.”} Among the barbarians, likewise, we read of slaveholding churches, and of special provisions to protect their slaves.\footnote{Comp. Milman, l.c. i. 531.} Constantine issued rigid laws against intermarriage with slaves, all the offspring of which must be slaves; and against fugitive slaves (a.d. 319 and 326), who at that time in great multitudes plundered deserted provinces or joined with hostile barbarians against the empire. But on the other hand he facilitated manumission, permitted it even on Sunday, and gave the clergy the right to emancipate their slaves simply by their own word, without the witnesses and ceremonies required in other cases.\footnote{In two laws of 316 and 321; Corp. Jur. l. i. tit. 13, l. 1 and 2.} By Theodosius and Justinian the liberation of slaves was still further encouraged. The latter emperor abolished the penalty of condemnation to servitude, and by giving to freed persons the rank and rights of citizens, he removed the stain which had formerly attached to that class.\footnote{Cod. Just. vii. 5, 6; Nov. 22, c. 8 (a.d.536), and Nov. 78, praef. 1, 2 (a.d.539).} The spirit of his laws favored the gradual abolition of domestic slavery. In the Byzantine empire in general the differences of rank in society were more equalized, though not so much on Christian principle as in the interest of despotic monarchy. Despotism and extreme democracy meet in predilection for universal equality and uniformity. Neither can suffer any overshadowing greatness, save the majesty of the prince or the will of the people. The one system knows none but slaves; the other, none but masters.

Nor was an entire abolition of slavery at that time at all demanded or desired even by the church. As in the previous period, she still thought it sufficient to insist on the kind Christian treatment of slaves, enjoining upon them obedience for the sake of the Lord, comforting them in their low condition with the thought of their higher moral freedom and equality, and by the religious education of the slaves making an inward preparation for the abolition of the institution. All hasty and violent measures met with decided disapproval. The council of Gangra threatens with the ban every one, who under pretext of religion seduces slaves into contempt of their masters; and the council of Chalcedon, in its fourth canon, on pain of excommunication forbids monasteries to harbor slaves without permission of the masters, lest Christianity be guilty of encouraging insubordination. The church fathers, so far as they enter this subject at all, seem to look upon slavery as at once a necessary evil and a divine instrument of discipline; tracing it to the curse on Ham and Canaan.\footnote{Gen. ix. 25: “Cursed be Canaan; a servant of servants shall he be unto his brethren.” But Christ appeared to remove every curse of sin, and every kind of slavery. The service of God is perfect \emph{freedom.}} It is true, they favor emancipation in individual cases, as an act of Christian love on the part of the master, but not as a right on the part of the slave; and the well-known passage: “If then mayest be made free, use it rather,” they understand not as a challenge to slaves to take the first opportunity to gain their freedom, but, on the contrary, as a challenge to remain in their servitude, since they are at all events inwardly free in Christ, and their outward condition is of no account.\footnote{1 Cor. vii. 21. The Greek fathers supply, with \textgreek{μᾶλλον χρῆσαι}, the word \textgreek{δουλείᾳ}(Chrysostom: \textgreek{μᾶλλον δούλευε}); whereas nearly all modem interpreters (except De Wette, Meyer, Ewald, and Alford) follow Calvin and Grotius in supplying \textgreek{ἐλευθερίᾳ}. Chrysostom, however, mentions this construction, and in another place (Serm. iv. in Genes. tom. v. p. 666) seems himself to favor it. The verb \emph{use} connects itself more naturally with \emph{freedom}, which is a boon and a blessing, than with \emph{bondage}, which is a state of privation. Milman, however, goes too far when he asserts (Lat. Christianity, vol. i. 492): “The abrogation of slavery was not contemplated even as a remote possibility. A general enfranchisement seems never to have dawned on the wisest and best of the Christian writers, notwithstanding the greater facility for manumission, and the sanctity, as it were, assigned to the act by Constantine, by placing it under the special superintendence of the clergy.” Compare against this statement the views of Chrysostom and Augustine, in the text.}

Even St. Chrysostom, though of all the church fathers the nearest to the emancipation theory and the most attentive to the question of slavery in general, does not rise materially above this view.\footnote{The views of Chrysostom on slavery are presented in his Homilies on Genesis and on the Epistles of Paul, and are collected by Möhler in his beautiful article on the Abolition of Slavery (Vermischte Schriften, ii. p. 89 sqq.). Möhler says that since the times of the apostle Paul no one has done a more valuable service to slaves then St. Chrysostom. But he overrates his merit.} According to him mankind were originally created perfectly free and equal, without the addition of a slave. But by the fall man lost the power of self-government, and fell into a threefold bondage: the bondage of woman under man, of slave under master, of subject under ruler. These three relations he considers divine punishments and divine means of discipline. Thus slavery, as a divine arrangement occasioned by the fall, is at once relatively justified and in principle condemned. Now since Christ has delivered us from evil and its consequences, slavery, according to Chrysostom, is in principle abolished in the church, yet only in the sense in which sin and death are abolished. Regenerate Christians are not slaves, but perfectly free men in Christ and brethren among themselves. The exclusive authority of the one and subjection of the other give place to mutual service in love. Consistently carried out, this view leads of course to emancipation. Chrysostom, it is true, does not carry it to that point, but he decidedly condemns all luxurious slaveholding, and thinks one or two servants enough for necessary help, while many patricians had hundreds and thousands. He advises the liberation of superfluous slaves, and the education of all, that in case they should be liberated, they may know how to take care of themselves. He is of opinion that the first Christian community at Jerusalem, in connection with community of goods, emancipated all their slaves;\footnote{Homil. xi. in Acta Apost. (Opera omn., tom. ix. p. 93): \textgreek{Οὐδὲ γὰρ τότε τοῦτο ἧν, ἀλλ  ἐλευθέρους ἴσως ἐπέτρεπον γίνεσθαι}. The monk Nilus, a pupil of Chrysostom, went so far as to declare slaveholding inconsistent with true love to Christ, Ep. lib. i. ep. 142 (quoted by Neander in his chapter on monasticism): \textgreek{Οὐ γὰρ οἷμαι οἰκέτην ἔχειν τὸν φιλόχριστον, εἰδότα τὴν χάριν τὴν πάντας ἐλευθερώσασαν}.} and thus he gives his hearers a hint to follow that example. But of an appeal to slaves to break their bonds, this father shows of course no trace; he rather, after apostolic precedent, exhorts them to conscientious and cheerful obedience for Christ’s sake, as earnestly as he inculcates upon masters humanity and love. The same is true of Ambrose, Augustine, and Peter Chrysologus of Ravenna († 458).

St. Augustine, the noblest representative of the Latin church, in his profound work on the “City of God,” excludes slavery from the original idea of man and the final condition of society, and views it as an evil consequent upon sin, yet under divine direction and control. For God, he says, created man reasonable and lord only over the unreasonable, not over man. The burden of servitude was justly laid upon the sinner. Therefore the term servant is not found in the Scriptures till Noah used it as a curse upon his offending son. Thus it was guilt and not nature that deserved that name. The Latin word servus is supposed to be derived from servare [servire rather], or the preservation of the prisoners of war from death, which itself implies the desert of sin. For even in a just war there is sin on one side, and every victory humbles the conquered by divine judgment, either reforming their sins or punishing them. Daniel saw in the sins of the people the real cause of their captivity. Sin, therefore, is the mother of servitude and first cause of man’s subjection to man; yet this does not come to pass except by the judgment of God, with whom there is no injustice, and who knows how to adjust the various punishments to the merits of the offenders .... The apostle exhorts the servants to obey their masters and to serve them ex animo, with good will; to the end that, if they cannot be made free from their masters, they may make their servitude a freedom to themselves by serving them not in deceitful fear, but in faithful love, until iniquity be overpassed, and all man’s principality and power be annulled, and God be all in all.\footnote{De Civit. Dei, lib. xix. cap. 15.}

As might be expected, after the conversion of the emperors, and of rich and noble families, who owned most slaves, cases of emancipation became more frequent.\footnote{For earlier cases, at the close of the previous period, see vol. i. § 89, at the end.} The biographer of St. Samson Xenodochos, a contemporary of Justinian, says of him: “His troop of slaves he would not keep, still less exercise over his fellow servants a lordly authority; he preferred magnanimously to let them go free, and gave them enough for the necessaries of life.”\footnote{Acta Sanct. Boll. Jun. tom. v. p. 267. According to Palladius, Hist. c. 119, St. Melania had, in concert with her husband Pinius, manumitted as many as eight thousand slaves. Yet it is only the ancient Latin translation that has this almost incredible number.} Salvianus, a Gallic presbyter of the fifth century, says that slaves were emancipated daily.\footnote{Ad Eccles. cath. l. iii. § 7 (Galland. tom. x. p. 71): “In usu quidem quotidiano est, ut servi, etsi non optimae, certe non infirmae servitudinis, Romana a dominis libertate donentur; in qua scilicet et proprietatem peculii capiunt et jus testamentarium consequuntur: ita ut et viventes, cui volunt, res suas tradant, et morientes donatione transcribAnt. Nec solum hoc, sed et illa, quae in servitute positi conquisierant, ex dominorum domo tollere non vetantur.” From this passage it appears that many masters, with a view to set their slaves free, allowed them to earn something; which was not allowed by the Roman law.} On the other hand, very much was done in the church to prevent the increase of slavery; especially in the way of redeeming prisoners, to which sometimes the gold and silver vessels of churches were applied. But we have no reliable statistics for comparing even approximately the proportion of the slaves to the free population at the close of the sixth century with the proportion in the former period.

We infer then, that the Christianity of the Nicene and post-Nicene age, though naturally conservative and decidedly opposed to social revolution and violent measures of reform, yet in its inmost instincts and ultimate tendencies favored the universal freedom of man, and, by elevating the slave to spiritual equality with the master, and uniformly treating him as capable of the same virtues, blessings, and rewards, has placed the hateful institution of human bondage in the way of gradual amelioration and final extinction. This result, however, was not reached in Europe till many centuries after our period, nor by the influence of the church alone, but with the help of various economical and political causes, the unprofitableness of slavery, especially in more northern latitudes, the new relations introduced by the barbarian conquests, the habits of the Teutonic tribes settled within the Roman empire, the attachment of the rural slave to the soil, and the change of the slave into the serf, who was as immovable as the soil, and thus, in some degree independent on the caprice and despotism of his master.

5. The poor and unfortunate in general, above all the widows and orphans, prisoners and sick, who were so terribly neglected in heathen times, now drew the attention of the imperial legislators. Constantine in 315 prohibited the branding of criminals on the forehead, “that the human countenance,” as he said, “formed after the image of heavenly beauty, should not be defaced.”\footnote{Cod. Theod. ix. 40, 1 and 2.} He provided against the inhuman maltreatment of prisoners before their trial.\footnote{C. Theod. ix. tit. 3, de custodia reorum. Comp. later similar laws of the year 409 in l. 7, and of 529 in the Cod. Justin. i. 4, 22.} To deprive poor parents of all pretext for selling or exposing their children, he had them furnished with food and clothing, partly at his own expense and partly at that of the state.\footnote{Comp. the two laws De alimentis quae inopes parentes de publico petere debent, in the Cod. Theod. xi. 27, 1 and 2.} He likewise endeavored, particularly by a law of the year 331, to protect the poor against the venality and extortion of judges, advocates, and tax collectors, who drained the people by their exactions.\footnote{Cod. Theod. I. tit. 7, l. 1: Cessent jam nunc rapaces officialium manus, cessent inquam! nam si moniti non cessaverint, gladiis praecidentur.} In the year 334 he ordered that widows, orphans, the sick, and the poor should not be compelled to appear before a tribunal outside their own province. Valentinian, in 365, exempted widows and orphans from the ignoble poll tax.\footnote{The capitatio plebeja. Cod. Theod. xiii. 10, 1 and 4. Other laws in behalf of widows, Cod. Just. iii. 14; ix. 24.} In 364 he intrusted the bishops with the supervision of the poor. Honorius did the same in 409. Justinian, in 529, as we have before remarked, gave the bishops the oversight of the state prisons, which they were to visit on Wednesdays and Fridays, to bring home to the unfortunates the earnestness and comfort of religion. The same emperor issued laws against usury and inhuman severity in creditors, and secured benevolent and religious foundations by strict laws against alienation of their revenues from the original design of the founders. Several emperors and empresses took the church institutions for the poor and sick, for strangers, widows, and orphans, under their special patronage, exempted them from the usual taxes, and enriched or enlarged them from their private funds.\footnote{Cod. Theod. xi. 16, xiii. 1; Cod. Just. i. 3; Nov. 131. Comp. here in general Chastel: The Charity of the Primitive Churches (transl. by Mathe), pp. 281-293.} Yet in those days, as still in ours, the private beneficence of Christian love took the lead, and the state followed at a distance, rather with ratification and patronage than with independent and original activity.\footnote{Comp. Chastel, l.c., p. 293: “It appears, then, as to charitable institutions, the part of the Christian emperors was much less to found themselves, than to recognize, to regulate, to guarantee, sometimes also to enrich with their private gifts, that which the church had founded. Everywhere the initiative had been taken by religious charity. Public charity only followed in the distance, and when it attempted to go ahead originally and alone, it soon found that it had strayed aside, and was constrained to withdraw.”}

\xsection{Abolition of Gladiatorial Shows}

§ 21. Abolition of Gladiatorial Shows.

6. And finally, one of the greatest and most beautiful victories of Christian humanity over heathen barbarism and cruelty was the abolition of gladiatorial contests, against which the apologists in the second century had already raised the most earnest protest.\footnote{Comp. vol. i. § 88.}

These bloody shows, in which human beings, mostly criminals, prisoners of war, and barbarians, by hundreds and thousands killed one another or were killed in fight with wild beasts for the amusement of the spectators, were still in full favor at the beginning of the period before us. The pagan civilization here proves itself impotent. In its eyes the life of a barbarian is of no other use than to serve the cruel amusement of the Roman people, who wish quietly to behold with their own eyes and enjoy at home the martial bloodshedding of their frontiers. Even the humane Symmachus gave an exhibition of this kind during his consulate (391), and was enraged that twenty-nine Saxon prisoners of war escaped this public shame by suicide.\footnote{Symm. l. ii. Ep. 46. Comp. vii. 4.} While the Vestal virgins existed, it was their special prerogative to cheer on the combatants in the amphitheatre to the bloody work, and to give the signal for the deadly stroke.\footnote{Prudentius Adv. Symmach. ii. 1095: \par Virgo—consurgit ad ictus, \par Et quotiens victor ferrum jugulo inserit, illa \par Delicias ait esse suas, pectusque jacentis \par Virgo modesta jubet, converso pollice, rumpi; \par Ni lateat pars ulla animae vitalibus imis, \par Altius impresso dum palpitat ense secutor.}

The contagion of the thirst for blood, which these spectacles generated, is presented to us in a striking example by Augustine in his Confessions.\footnote{Lib. vi. c. 8.} His friend Alypius, afterward bishop of Tagaste, was induced by some friends in 385 to visit the amphitheatre at Rome, and went resolved to lock himself up against all impressions. “When they reached the spot,” says Augustine, “and took their places on the hired seats, everything already foamed with bloodthirsty delight. But Alypius, with closed eyes, forbade his soul to yield to this sin. O had he but stopped also his ears! For when, on the fall of a gladiator in the contest, the wild shout of the whole multitude fell upon him, overcome by curiosity he opened his eyes, though prepared to despise and resist the sight. But he was smitten with a more grievous wound in the soul than the combatant in the body, and fell more lamentably .... For when he saw the blood, he imbibed at once the love of it, turned not away, fastened his eyes upon it, caught the spirit of rage and vengeance before he knew it, and, fascinated with the murderous game, became drunk with bloodthirsty joy .... He looked, shouted applause, burned, and carried with him thence the frenzy, by which he was drawn to go back, not only with those who had taken him there, but before them, and taking others with him.”

Christianity finally succeeded in closing the amphitheatre. Constantine, who in his earlier reign himself did homage to the popular custom in this matter, and exposed a great multitude of conquered barbarians to death in the amphitheatre at Treves, for which he was highly commended by a heathen orator,\footnote{Eumenii Panegyr. c. 12.} issued in 325, the year of the great council of the church at Nice, the first prohibition of the bloody spectacles, “because they cannot be pleasing in a time of public peace.”\footnote{Cod. Theod. xv. tit. 12, l. 1, de gladiatoribus: “Cruenta spectacula in otio civili et domestica quiete non placent; quapropter omnino gladiatores esse prohibemus.” Comp. Euseb. Vita Const. iv. 25.} But this edict, which is directed to the prefects of Phoenicia, had no permanent effect even in the East, except at Constantinople, which was never stained with the blood of gladiators. In Syria and especially in the West, above all in Rome, the deeply rooted institution continued into the fifth century. Honorius (395–423), who at first considered it indestructible, abolished the gladiatorial shows about 404, and did so at the instance of the heroic self-denial of an eastern monk by the name of Telemachus, who journeyed to Rome expressly to protest against this inhuman barbarity, threw himself into the arena, separated the combatants, and then was torn to pieces by the populace, a martyr to humanity.\footnote{So relates Theodoret: Hist. eccl. l. v. c. 26. For there is no law of Honorius extant on the subject. Yet after this time there is no mention of a gladiatorial contest between man and man.} Yet this put a stop only to the bloody combats of men. Unbloody spectacles of every kind, even on the high festivals of the church and amidst the invasions of the barbarians, as we see by the grievous complaints of a Chrysostom, an Augustine, and a Salvian, were as largely and as passionately attended as ever; and even fights with wild animals, in which human life was generally more or less sacrificed, continued,\footnote{In a law of Leo, of the year 469 (in the Cod. Justin. iii. tit. 12, l. 11), besides the scena theatralis and the circense theatrum, also ferarum lacrymosa spectacula are mentioned as existing. Salvian likewise, in the fifth century (De gubern. Dei, l. vi. p. 51), censures the delight of his contemporaries in such bloody combats of man with wild beasts. So late as the end of the seventh century a prohibition from the Trullan council was called for in the East, In the West, Theodoric appears to have exchanged the beast fights for military displays, whence proceeded the later tournaments. Yet these shows have never become entirely extinct, but remain in the bull fights of Southern Europe, especially in Spain.} and, to the scandal of the Christian name, are tolerated in Spain and South America to this day.

\xsection{Evils of the Union of Church and State. Secularization of the Church}

§ 22. Evils of the Union of Church and State. Secularization of the Church.

We turn now to the dark side of the union of the church with the state; to the consideration of the disadvantages which grew out of their altered relation after the time of Constantine, and which continue to show themselves in the condition of the church in Europe to our own time.

These evil results may be summed up under the general designation of the secularization of the church. By taking in the whole population of the Roman empire the church became, indeed, a church of the masses, a church of the people, but at the same time more or less a church of the world. Christianity became a matter of fashion. The number of hypocrites and formal professors rapidly increased;\footnote{Thus Augustine, for example, Tract. in JoAnn. xxv. c. 10, laments that the church filled itself daily with those who sought Jesus not for Jesus, but for earthly profit. Comp. the similar complaint of Eusebius, Vita Const. l. iv. c. 54.} strict discipline, zeal, self-sacrifice, and brotherly love proportionally ebbed away; and many heathen customs and usages, under altered names, crept into the worship of God and the life of the Christian people. The Roman state had grown up under the influence of idolatry, and was not to be magically transformed at a stroke. With the secularizing process, therefore, a paganizing tendency went hand in hand.

Yet the pure spirit of Christianity could by no means be polluted by this. On the contrary it retained even in the darkest days its faithful and steadfast confessors, conquered new provinces from time to time, constantly reacted, both within the established church and outside of it, in the form of monasticism, against the secular and the pagan influences, and, in its very struggle with the prevailing corruption, produced such church fathers as Athanasius, Chrysostom, and Augustine, such exemplary Christian mothers as Nonna, Anthusa, and Monica, and such extraordinary saints of the desert as Anthony, Pachomius, and Benedict. New enemies and dangers called forth new duties and virtues, which could now unfold themselves on a larger stage, and therefore also on a grander scale. Besides, it must not be forgotten, that the tendency to secularization is by no means to be ascribed only to Constantine and the influence of the state, but to the deeper source of the corrupt heart of man, and did reveal itself, in fact, though within a much narrower compass, long before, under the heathen emperors, especially in the intervals of repose, when the earnestness and zeal of Christian life slumbered and gave scope to a worldly spirit.

The difference between the age after Constantine and the age before consists, therefore, not at all in the cessation of true Christianity and the entrance of false, but in the preponderance of the one over the other. The field of the church was now much larger, but with much good soil it included far more that was stony, barren, and overgrown with weeds. The line between church and world, between regenerate and unregenerate, between those who were Christians in name and those who were Christians in heart, was more or less obliterated, and in place of the former hostility between the two parties there came a fusion of them in the same outward communion of baptism and confession. This brought the conflict between light and darkness, truth and falsehood, Christ and antichrist, into the bosom of Christendom itself.

\xsection{Worldliness and Extravagance}

§23. Worldliness and Extravagance.

The secularization of the church appeared most strikingly in the prevalence of mammon worship and luxury compared with the poverty and simplicity of the primitive Christians. The aristocracy of the later empire had a morbid passion for outward display and the sensual enjoyments of wealth, without the taste, the politeness, or the culture of true civilization. The gentlemen measured their fortune by the number of their marble palaces, baths, slaves, and gilded carriages; the ladies indulged in raiment of silk and gold ornamented with secular or religious figures, and in heavy golden necklaces, bracelets, and rings, and went to church in the same flaunting dress as to the theatre.\footnote{Ammianus Marcellinus gives the most graphic account of the extravagant and tasteless luxury of the Roman aristocracy in the fourth century; which Gibbon has admirably translated and explained in his 31st chapter.} Chrysostom addresses a patrician of Antioch: “You count so and so many acres of land, ten or twenty palaces, as many baths, a thousand or two thousand slaves, carriages plated with silver and gold.”\footnote{Homil. in Matt. 63, § 4 (tom. vii. p. 533), comp. Hom. in 1 Cor. 21, § 6, and many other places in his sermons. Comp. Neander’s Chrysostomus, i. p. 10 sqq. and Is. Taylor’s Anc. Christianity, vol. ii., supplement, p. xxx. sqq.} Gregory Nazianzen, who presided for a time in the second ecumenical council of Constantinople in 381, gives us the following picture, evidently rhetorically colored, yet drawn from life, of the luxury of the degenerate civilization of that period: “We repose in splendor on high and sumptuous cushions, upon the most exquisite covers, which one is almost afraid to touch, and are vexed if we but hear the voice of a moaning pauper; our chamber must breathe the odor of flowers, even rare flowers; our table must flow with the most fragrant and costly ointment, so that we become perfectly effeminate. Slaves must stand ready, richly adorned and in order, with waving, maidenlike hair, and faces shorn perfectly smooth, more adorned throughout than is good for lascivious eyes; some, to hold cups both delicately and firmly with the tips of their fingers, others, to fan fresh air upon the head. Our table must bend under the load of dishes, while all the kingdoms of nature, air, water and earth, furnish copious contributions, and there must be almost no room for the artificial products of cook and baker .... The poor man is content with water; but we fill our goblets with wine to drunkenness, nay, immeasurably beyond it. We refuse one wine, another we pronounce excellent when well flavored, over a third we institute philosophical discussions; nay, we count it a pity, if he does not, as a king, add to the domestic wine a foreign also.”\footnote{Orat. xiv. Comp. Ullmann’s monograph on Gregory, p. 6.} Still more unfavorable are the pictures which, a half century later, the Gallic presbyter, Salvianus, draws of the general moral condition of the Christians in the Roman empire.\footnote{Adv. avarit. and De gubern. Dei, passim. Comp. § 12, at the close.}

It is true, these earnest protests against degeneracy themselves, as well as the honor in which monasticism and ascetic contempt of the world were universally held, attest the existence of a better spirit. But the uncontrollable progress of avarice, prodigality, voluptuousness, theatre going, intemperance, lewdness, in short, of all the heathen vices, which Christianity had come to eradicate, still carried the Roman empire and people with rapid strides toward dissolution, and gave it at last into the hands of the rude, but simple and morally vigorous barbarians. When the Christians were awakened by the crashings of the falling empire, and anxiously asked why God permitted it, Salvian, the Jeremiah of his time, answered: “Think of your vileness and your crimes, and see whether you are worthy of the divine protection.”\footnote{De gubern. Dei, l. iv. c. 12, p. 82.} Nothing but the divine judgment of destruction upon this nominally Christian, but essentially heathen world, could open the way for the moral regeneration of society. There must be new, fresh nations, if the Christian civilization prepared in the old Roman empire was to take firm root and bear ripe fruit.

\xsection{Byzantine Court Christianity}

§ 24. Byzantine Court Christianity.

The unnatural confusion of Christianity with the world culminated in the imperial court of Constantinople, which, it is true, never violated moral decency so grossly as the court of a Nero or a Domitian, but in vain pomp and prodigality far outdid the courts of the better heathen emperors, and degenerated into complete oriental despotism. The household of Constantius, according to the description of Libanius,\footnote{Lib., Epitaph. Julian.} embraced no less than a thousand barbers, a thousand cup bearers, a thousand cooks, and so many eunuchs, that they could be compared only to the insects of a summer day. This boundless luxury was for a time suppressed by the pagan Julian, who delighted in stoical and cynical severity, and was fond of displaying it; but under his Christian successors the same prodigality returned; especially under Theodosius and his sons. These emperors, who prohibited idolatry upon pain of death, called their laws, edicts, and palaces “divine,” bore themselves as gods upon earth, and, on the rare occasions when they showed themselves to the people, unfurled an incredible magnificence and empty splendor.

“When Arcadius,” to borrow a graphic description from a modern historian, “condescended to reveal to the public the majesty of the sovereign, he was preceded by a vast multitude of attendants, dukes, tribunes, civil and military officers, their horses glittering with golden ornaments, with shields of gold set with precious stones, and golden lances. They proclaimed the coming of the emperor, and commanded the ignoble crowd to clear the streets before him. The emperor stood or reclined on a gorgeous chariot, surrounded by his immediate attendants, distinguished by shields with golden bosses set round with golden eyes, and drawn by white mules with gilded trappings; the chariot was set with precious stones, and golden fans vibrated with the movement, and cooled the air. The multitude contemplated at a distance the snow-white cushions, the silken carpets, with dragons inwoven upon them in rich colors. Those who were fortunate enough to catch a glimpse of the emperor, beheld his ears loaded with golden rings, his arms with golden chains, his diadem set with gems of all hues, his purple robes, which, with the diadem, were reserved for the emperor, in all their sutures embroidered with precious stones. The wondering people, on their return to their homes, could talk of nothing but the splendor of the spectacle: the robes, the mules, the carpets, the size and splendor of the jewels. On his return to the palace, the emperor walked on gold; ships were employed with the express purpose of bringing gold dust from remote provinces, which was strewn by the officious care of a host of attendants, so that the emperor rarely set his foot on the bare pavement.”\footnote{Milman: Hist. of Ancient Christianity, p. 440 (Am. ed.). Comp. the sketch of the court of Arcadius, which Montfaucon, in a treatise in the last volume of his Opera Chrys., and Müller: De genio, moribus, et luxu aevi Theodosiani, Copenh. 1798, have drawn, chiefly from the works of Chrysostom.}

The Christianity of the Byzantine court lived in the atmosphere of intrigue, dissimulation, and flattery. Even the court divines and bishops could hardly escape the contamination, though their high office, with its sacred functions, was certainly a protecting wall around them. One of these bishops congratulated Constantine, at the celebration of the third decennium of his reign (the tricennalia), that he had been appointed by God ruler over all in this world, and would reign with the Son of God in the other! This blasphemous flattery was too much even for the vain emperor, and he exhorted the bishop rather to pray God that he might be worthy to be one of his servants in this world and the next.\footnote{Euseb. Vit. Const. iv. 48.} Even the church historian and bishop Eusebius, who elsewhere knew well enough how to value the higher blessings, and lamented the indescribable hypocrisy of the sham Christianity around the emperor,\footnote{V. Const. iv. 54.} suffered himself to be so far blinded by the splendor of the imperial favor, as to see in a banquet, which Constantine gave in his palace to the bishops at the close of the council of Nice, in honor of his twenty years’ reign (the vicennalia), an emblem of the glorious reign of Christ upon the earth!\footnote{V. Const. iii. 15, where Eusebius, at the close of this imperio-episcopal banquet, “which transcended all description,” says: \textgreek{Χριστοῦ βασιλείας ἔδοξεν ἄν τις φαντασιοῦσθαι εἰκόνα, ὄναρ τ  εῖναι ἀλλ  οὐχ ὕπερ τὸ γινόμενον}.}

And these were bishops, of whom many still bore in their body the marks of the Diocletian persecution. So rapidly had changed the spirit of the age. While, on the other hand, the well-known firmness of Ambrose with Theodosius, and the life of Chrysostom, afford delightful proof that there were not wanting, even in this age, bishops of Christian earnestness and courage to rebuke the sins of crowned heads.

\xsection{Intrusion of Politics into Religion}

§ 25. Intrusion of Politics into Religion.

With the union of the church and the state begins the long and tedious history of their collisions and their mutual struggles for the mastery: the state seeking to subject the church to the empire, the church to subject the state to the hierarchy, and both very often transgressing the limits prescribed to their power in that word of the Lord: “Render unto Caesar the things which are Caesar’s, and unto God the things that are God’s.” From the time of Constantine, therefore, the history of the church and that of the world in Europe are so closely interwoven, that neither can be understood without the other. On the one hand, the political rulers, as the highest members and the patrons of the church, claimed a right to a share in her government, and interfered in various ways in her external and internal affairs, either to her profit or to her prejudice. On the other hand, the bishops and patriarchs, as the highest dignitaries and officers of the state religion, became involved in all sorts of secular matters and in the intrigues of the Byzantine court. This mutual intermixture, on the whole, was of more injury than benefit to the church and to religion, and fettered her free and natural development.

Of a separation of religion and politics, of the spiritual power from the temporal, heathen antiquity knew nothing, because it regarded religion itself only from a natural point of view, and subjected it to the purposes of the all-ruling state, the highest known form of human society. The Egyptian kings, as Plutarch tells us, were at the same time priests, or were received into the priesthood at their election. In Greece the civil magistrate had supervision of the priests and sanctuaries.\footnote{This overseer was called \textgreek{βασιλεύς} of the \textgreek{ἱερεῖς} and \textgreek{ἱερά}.} In Rome, after the time of Numa, this supervision was intrusted to a senator, and afterward united with the imperial office. All the pagan emperors, from Augustus\footnote{Augustus took the dignity of Pontifex Maximus after the death of Lepidus, a.u.742, and thenceforth that office remained inherent in the imperial, though it was usually conferred by a decree of the senate. Formerly the pontifex maximus was elected by the people for life, could take no civil office, must never leave Italy, touch a corpse, or contract a second marriage; and he dwelt in the old king’s house, the regia. Augustus himself exercised the office despotically enough, though with great prudence. He nominated and increased at pleasure the members of the sacerdotal college, chose the vestal virgins, determined the authority of the vaticinia, purged the Sibylline books of apocryphal interpolations, continued the reform of the calendar begun by Caesar, and changed the month Sextius into Augustus in his own honor, as Quintius, the birth-month of Julius Caesar, had before been rebaptized Julius. Comp. Charles Merivale: Hist. of the Romans under the Empire, vol. iii. (Lond. 1851), p, 478 sqq. (This work, which stops where Gibbon begins, has been republished in 7 vols. in New York, 1863.)} to Julian the Apostate, were at the same time supreme pontiffs (Pontifices Maximi), the heads of the state religion, emperor-popes. As such they could not only perform all priestly functions, even to offering sacrifices, when superstition or policy prompted them to do so, but they also stood at the head of the highest sacerdotal college (of fifteen or more Pontifices), which in turn regulated and superintended the three lower classes of priests (the Epulones, Quindecemviri, and Augures), the temples and altars, the sacrifices, divinations, feasts, and ceremonies, the exposition of the Sibylline books, the calendar, in short, all public worship, and in part even the affairs of marriage and inheritance.

Now it may easily be supposed that the Christian emperors, who, down to Gratian (about 380), even retained the name and the insignia of the Pontifex Maximus, claimed the same oversight of the Christian religion established in the empire, which their predecessors had had of the heathen; only with this material difference, that they found here a stricter separation between the religious element and the political, the ecclesiastical and the secular, and were obliged to bind themselves to the already existing doctrines, usages, and traditions of the church which claimed divine institution and authority.

\xsection{The Emperor-Papacy and the Hierarchy}

§ 26. The Emperor-Papacy and the Hierarchy.

And this, in point of fact, took place first under Constantine, and developed under his successors, particularly under Justinian, into the system of the Byzantine imperial papacy,\footnote{In England and Scotland the term \emph{Erastianism} is used for this; but is less general, and not properly applicable at all to the Greek church. For the man who furnished the word, Thomas Erastus, a learned and able physician and professor of medicine in Heidelberg (died at Basle in Switzerland, 1583), was an opponent not only of the independence of the church toward the state, but also of the church ban and of the presbyterial constitution and discipline, as advocated by Frederick III., of the Palatinate, and the authors of the Heidelberg Catechism, especially Olevianus, a pupil of Calvin. He was at last excommunicated for his views by the church council in Heidelberg.} or of the supremacy of the state over the church.

Constantine once said to the bishops at a banquet, that he also, as a Christian emperor, was a divinely appointed bishop, a bishop over the external affairs of the church, while the internal affairs belonged to the bishops proper.\footnote{His words, which are to be taken neither in jest and pun (as Neander supposes), nor as mere compliment to the bishops, but in earnest, run thus, in Eusebius: Vita Const. l. iv. c. 24: \textgreek{Ὑμεῖς} (the \textgreek{ἐπίσκοποι} addressed) \textgreek{μέν τῶν εἴσω τῆς ἐκκλησίας, ἐγὼ δὲ τῶν ἐκτὸς ὑπὸ θεοῦ καθεσταμένος ἐπίσκοπος ἅν εἴην}. All depends here on the intrepretation of the antithesis \textgreek{τῶμ εἴσω} and \textgreek{τῶν ἐκτὸς τῆς ἐκκλησίας}. (a) The explanation of Stroth and others takes the genitive as masculine, \textgreek{οἱ εἴσω} denoting Christians, and \textgreek{οἱ ἐκτός} heathens; so that Constantine ascribed to himself only a sort of episcopate in \emph{partibus infidelium}. But this contradicts the connection; for Eusebius says immediately after, that he took a certain religious oversight over all his subjects (\textgreek{τοὺς ἀρχομένους ἅπαντας ἐπεσκόπει}, etc.), and calls him also elsewhere a universal bishop ” (i. 44). (b) Gieseler’s interpretation is not much better (I. 2. § 92, not. 20, Amer. ed. vol. i. p. 371): that \textgreek{οἱ ἐκτός} denotes all his subjects, Christian as well as non-Christian, but only in their civil relations, so far as they are outside the church. This entirely blunts the antithesis with \textgreek{οἱ εἴσω}, and puts into the emperor’s mouth a mere commonplace instead of a new idea; for no one doubted his \emph{political} sovereignty. (c) The genitive is rather to be taken as neuter in both cases, and \textgreek{πραγμάτων} to be supplied. This agrees with usage (we find it in Polybius), and gives a sense which agrees with the view of Eusebius and with the whole practice of Constantine. There is, however, of course, another question: What is the proper distinction between \textgreek{τὰ εἴσω} and \textgreek{τὰ ἐκτός} the \emph{interna} and \emph{externa} of the church, or, what is much the same, between the sacerdotal \emph{jus} \emph{in sacra} and the imperial \emph{jus circa sacra.} This Constantine and his age certainly could not themselves exactly define, since the whole relation was at that time as yet new and undeveloped.} In this pregnant word he expressed the new posture of the civil sovereign toward the church in a characteristic though indefinite and equivocal way. He made there a distinction between two divinely authorized episcopates; one secular or imperial, corresponding with the old office of Pontifex Maximus, and extending over the whole Roman empire, therefore ecumenical or universal; the other spiritual or sacerdotal, divided among the different diocesan bishops, and appearing properly in its unity and totality only in a general council.

Accordingly, though not yet even baptized, he acted as the patron and universal temporal bishop of the church;\footnote{Eusebius in fact calls him a divinely appointed universal bishop, \textgreek{οἷά τις κοινὸς ἐπίσκοπος ἐκ θεοῦ δακεσταμένος , συνόδους τῶν τοῦ θεοῦ λειτουργῶν συνεκρότει}. Vit. Const. i. 44. His son Constantius was fond of being called ” bishop of bishops.”} summoned the first ecumenical council for the settlement of the controversy respecting the divinity of Christ; instituted and deposed bishops; and occasionally even delivered sermons to the people; but on the other hand, with genuine tact (though this was in his earlier period, a.d. 314), kept aloof from the Donatist controversy, and referred to the episcopal tribunal as the highest and last resort in purely spiritual matters. In the exercise of his imperial right of supervision he did not follow any clear insight and definite theory so much as an instinctive impulse of control, a sense of politico-religious duty, and the requirements of the time. His word only raised, did not solve, the question of the relation between the imperial and the sacerdotal episcopacy and the extent of their respective jurisdictions in a Christian state.

This question became thenceforth the problem and the strife of history both sacred and secular, ran through the whole mediaeval conflict between emperor and pope, between imperial and hierarchical episcopacy, and recurs in modified form in every Protestant established church.

In general, from this time forth the prevailing view was, that God has divided all power between the priesthood and the kingdom (sacerdotium et imperium), giving internal or spiritual affairs, especially doctrine and worship, to the former, and external or temporal affairs, such as government and discipline, to the latter.\footnote{Justinian states the Byzantine theory thus, in the preface to the 6th Novel: “Maxima quidem in hominibus sunt dona Dei a superna collata clementia \emph{Sacerdotium} et \emph{Imperium}, et illud quidem divinis ministrans, hoc autem humanis praesidens ac diligentiam exhibens, ex uno eodemque principio utraque procedentia, humanam exornant vitam.” But he then ascribes to the Imperium the supervision of the Sacerdotium, and “maximam sollicitudinem circa vera Dei dogmata et circa Sacerdotum honestatem.” Later Greek emperors, on the ground of their anointing, even claimed a priestly character. Leo the Isaurian, for example, wrote to Pope Gregory II. in 730: \textgreek{βασιλεὺς καὶ ἱερεύς εἰμι} (Mansi xii. 976). This, however, was contested even in the East, and the monk Maximus in 655 answered negatively the question put to him: “Ergo non est omnis Christianus imperator etiam sacerdos?” At first the emperor’s throne stood side by side with the bishop’s in the choir; but Ambrose gave the emperor a seat next to the choir. Yet, after the ancient custom, which the Concilium Quinisext., a.d.692, in its 69th canon, expressly confirmed, the emperors might enter the choir of the church, and lay their oblations in person upon the altar—a privilege which was denied to all the laity, and which implied at least a half-priestly character in the emperor. Gibbon’s statement needs correction accordingly (ch. xx.): “The monarch, whose spiritual rank is less honorable than that of the meanest deacon, was seated below the rails of the sanctuary, and confounded with the rest of the faithful multitude.”} But internal and external here vitally interpenetrate and depend on each other, as soul and body, and frequent reciprocal encroachments and collisions are inevitable upon state-church ground. This becomes manifest in the period before us in many ways, especially in the East, where the Byzantine despotism had freer play, than in the distant West.

The emperors after Constantine (as the popes after them) summoned the general councils, bore the necessary expenses, presided in the councils through commissions, gave to the decisions in doctrine and discipline the force of law for the whole Roman empire, and maintained them by their authority. The emperors nominated or confirmed the most influential metropolitans and patriarchs. They took part in all theological disputes, and thereby inflamed the passion of parties. They protected orthodoxy and punished heresy with the arm of power. Often, however, they took the heretical side, and banished orthodox bishops from their sees. Thus Arianism, Nestorianism, Eutychianism, and Monophysitism successively found favor and protection at court. Even empresses meddled in the internal and external concerns of the church. Justina endeavored with all her might to introduce Arianism in Milan, but met a successful opponent in bishop Ambrose. Eudoxia procured the deposition and banishment of the noble Chrysostom. Theodora, raised from the stage to the throne, ruled the emperor Justinian, and sought by every kind of intrigue to promote the victory of the Monophysite heresy. It is true, the doctrinal decisions proceeded properly from the councils, and could not have maintained themselves long without that sanction. But Basiliscus, Zeno, Justinian I., Heraclius, Constans II., and other emperors issued many purely ecclesiastical edicts and rescripts without consulting the councils, or through the councils by their own influence upon them. Justinian opens his celebrated codex with the imperial creed on the trinity and the imperial anathema against Nestorius, Eutyches, Apollinaris, on the basis certainly of the apostolic church and of the four ecumenical councils, but in the consciousness of absolute legislative and executive authority even over the faith and conscience of all his subjects.

The voice of the catholic church in this period conceded to the Christian emperors in general, with the duty of protecting and supporting the church, the right of supervision over its external affairs, but claimed for the clergy, particularly for the bishops, the right to govern her within, to fix her doctrine, to direct her worship. The new state of things was regarded as a restoration of the Mosaic and Davidic theocracy on Christian soil, and judged accordingly. But in respect to the extent and application of the emperor’s power in the church, opinion was generally determined, consciously or unconsciously, by some special religious interest. Hence we find that catholics and heretics, Athanasians and Arians, justified or condemned the interference of the emperor in the development of doctrine, the appointment and deposition of bishops, and the patronage and persecution of parties, according as they themselves were affected by them. The same Donatists who first appealed to the imperial protection, when the decision went against them denounced all intermeddling of the state with the church. There were bishops who justified even the most arbitrary excesses of the Byzantine despotism in religion by reference to Melchizedek and the pious kings of Israel, and yielded them selves willing tools of the court. But there were never wanting also fearless defenders of the rights of the church against the civil power. Maximus the Confessor declared before his judges in Constantinople, that Melchizedek was a type of Christ alone, not of the emperor.

In general the hierarchy formed a powerful and wholesome check on the imperial papacy, and preserved the freedom and independence of the church toward the temporal power. That age had only the alternative of imperial or episcopal despotism; and of these the latter was the less hurtful and the more profitable, because it represented the higher intellectual and moral interests. Without the hierarchy, the church in the Roman empire and among the barbarians would have been the football of civil and military despots. It was, therefore, of the utmost importance, that the church, at the time of her marriage with the state, had already grown so large and strong as to withstand all material alteration by imperial caprice, and all effort to degrade her into a tool. The Apostolic Constitutions place the bishops even above all kings and magistrates.\footnote{Lib. ii. c. 11, where the bishop is reminded of his exalted position, \textgreek{ὡς θεοὶ τύπον ἔχων ἐν ἀνθρώποις τῷ πάντων ἄρχειν ἀνθρώπων, ἱερέων, βασιλέων, ἀρχόντων}\emph{,} etc. Comp. c. 33 and 34.} Chrysostom says that the first ministers of the state enjoyed no such honor as the ministers of the church. And in general the ministers of the church deserved their honor. Though there were prelates enough who abused their power to sordid ends, still there were men like Athanasius, Basil, Ambrose, Chrysostom, Augustine, Leo, the purest and most venerable characters, which meet us in the fourth and fifth centuries, far surpassing the contemporary emperors. It was the universal opinion that the doctrines and institutions of the church, resting on divine revelation, are above all human power and will. The people looked, in blind faith and superstition, to the clergy as their guides in all matters of conscience, and even the emperors had to pay the bishops, as the fathers of the churches, the greatest reverence, kiss their hands, beg their blessing, and submit to their admonition and discipline. In most cases the emperors were mere tools of parties in the church. Arbitrary laws which were imposed upon the church from without rarely survived their makers, and were condemned by history. For there is a divine authority above all thrones, and kings, and bishops, and a power of truth above all the machinations of falsehood and intrigue.

The Western church, as a whole, preserved her independence far more than the Eastern; partly through the great firmness of the Roman character, partly through the favor of political circumstances, and of remoteness from the influence and the intrigues of the Byzantine court. Here the hierarchical principle developed itself from the time of Leo the Great even to the absolute papacy, which, however, after it fulfilled its mission for the world among the barbarian nations of the middle ages, degenerated into an insufferable tyranny over conscience, and thus exposed itself to destruction. In the Catholic system the freedom and independence of the church involve the supremacy of an exclusive priesthood and papacy; in the Protestant, they can be realized only on the broader basis of the universal priesthood, in the self-government of the Christian people; though this is, as yet, in all Protestant established churches more or less restricted by the power of the state.

\xsection{Restriction of Religious Freedom, and Beginnings of Persecution of Heretics}

§ 27. Restriction of Religious Freedom, and Beginnings of Persecution of Heretics.

Sam. Eliot: History of Liberty. Boston, 1858, 4 vols. Early Christians, vols. i. and ii. The most important facts are scattered through the sections of the larger church histories on the heresies, the doctrinal controversies, and church discipline.

An inevitable consequence of the union of church and state was restriction of religious freedom in faith and worship, and the civil punishment of departure from the doctrine and discipline of the established church.

The church, dominant and recognized by the state, gained indeed external freedom and authority, but in a measure at the expense of inward liberty and self-control. She came, as we have seen in the previous section, under the patronage and supervision of the head of the Christian state, especially in the Byzantine empire. In the first three centuries, the church, with all her external lowliness and oppression, enjoyed the greater liberty within, in the development of her doctrines and institutions, by reason of her entire separation from the state.

But the freedom of error and division was now still more restricted. In the ante-Nicene age, heresy and schism were as much hated and abhorred indeed, as afterward, yet were met only in a moral way, by word and writing, and were punished with excommunication from the rights of the church. Justin Martyr, Tertullian, and even Lactantius were the first advocates of the principle of freedom of conscience, and maintained, against the heathen, that religion was essentially a matter of free will, and could be promoted only by instruction and persuasion not by outward force.\footnote{Just. Mart. Apol. i. 2, 4, 12; Tertull. Apolog. c. 24, 28; Ad Scapul.c. 2; Lactant. Instit. v. 19, 20; Epit. c. 54. Comp. vol. i. § 51.} All they say against the persecution of Christians by the heathen applies in full to the persecution of heretics by the church. After the Nicene age all departures from the reigning state-church faith were not only abhorred and excommunicated as religious errors, but were treated also as crimes against the Christian state, and hence were punished with civil penalties; at first with deposition, banishment, confiscation, and, after Theodosius, even with death.

This persecution of heretics was a natural consequence of the union of religious and civil duties and rights, the confusion of the civil and the ecclesiastical, the judicial and the moral, which came to pass since Constantine. It proceeded from the state and from the emperors, who in this respect showed themselves the successors of the Pontifices Maximi, with their relation to the church reversed. The church, indeed, steadfastly adhered to the principle that, as such, she should employ only spiritual penalties, excommunication in extreme cases; as in fact Christ and the apostles expressly spurned and prohibited all carnal weapons, and would rather suffer and die than use violence. But, involved in the idea of Jewish theocracy and of a state church, she practically confounded in various ways the position of the law and that of the gospel, and in theory approved the application of forcible measures to heretics, and not rarely encouraged and urged the state to it; thus making herself at least indirectly responsible for the persecution. This is especially, true of the Roman church in the times of her greatest power, in the middle age and down to the end of the sixteenth century; and by this course that church has made herself almost more offensive in the eyes of the world and of modern civilization than by her peculiar doctrines and usages. The Protestant reformation dispelled the dream that Christianity was identical with an outward organization, or the papacy, and gave a mighty shock thereby to the principle of ecclesiastical exclusiveness. Yet, properly speaking, it was not till the eighteenth century that a radical revolution of views was accomplished in regard to religious toleration; and the progress of toleration and free worship has gone hand in hand with the gradual loosening of the state-church basis and with the clearer separation of civil and religious rights and of the temporal and spiritual power.

In the, beginning of his reign, Constantine proclaimed full freedom of religion (312), and in the main continued tolerably true to it; at all events he used no violent measures, as his successors did. This toleration, however, was not a matter of fixed principle with him, but merely of temporary policy; a necessary consequence of the incipient separation of the Roman throne from idolatry, and the natural transition from the sole supremacy of the heathen religion to the same supremacy of the Christian. Intolerance directed itself first against heathenism; but as the false religion gradually died out of itself, and at any rate had no moral energy for martyrdom, there resulted no such bloody persecutions of idolatry under the Christian emperors, as there had been of Christianity under their heathen predecessors. Instead of Christianity, the intolerance of the civil power now took up Christian heretics, whom it recognized as such. Constantine even in his day limited the freedom and the privileges which he conferred, to the catholic, that is, the prevailing orthodox hierarchical church, and soon after the Council of Nice, by an edict of the year 326, expressly excluded heretics and schismatics from these privileges.\footnote{Cod. Theod. xvi. 5, 1: Privilegia, quae contemplatione religionis indulta sunt, catholicae tantum legis observatoribus prodesse opportet. Haereticos autem atque schismaticos non tantum ab his privilegiis alienos esse volumus, sed etiam diversis muneribus constringi et subjici.} Accordingly he banished the leaders of Arianism and ordered their writings to be burned, but afterward, wavering in his views of orthodoxy and heterodoxy, and persuaded over by some bishops and his sister, he recalled Arius and banished Athanasius. He himself was baptized shortly before his death by an Arian bishop. His son Constantius was a fanatical persecutor both of idolatry and the Nicene orthodoxy, and endeavored with all his might to establish Arianism alone in the empire. Hence the earnest protest of the orthodox bishops, Hosius, Athanasius, and Hilary, against this despotism and in favor of toleration;\footnote{Comp. § 8, above.} which came, however, we have to remember, from parties who were themselves the sufferers under intolerance, and who did not regard the banishment of the Arians as unjust.

Under Julian the Apostate religious liberty was again proclaimed, but only as the beginning of return to the exclusive establishment of heathenism; the counterpart, therefore, of Constantine’s toleration. After his early death Arianism again prevailed, at least in the East, and showed itself more, intolerant and violent than the catholic orthodoxy.

At last Theodosius the Great, the first emperor who was baptized in the Nicene faith, put an end to the Arian interregnum, proclaimed the exclusive authority of the Nicene creed, and at the same time enacted the first rigid penalties not only against the pagan idolatry, the practice of which was thenceforth a capital crime in the empire, but also against all Christian heresies and sects. The ruling principle of his public life was the unity of the empire and of the orthodox church. Soon after his baptism, in 380, he issued, in connection with his weak coëmperors, Gratian and Valentinian II., to the inhabitants of Constantinople, then the chief seat of Arianism, the following edict: “We, the three emperors, will, that all our subjects steadfastly adhere to the religion which was taught by St. Peter to the Romans, which has been faithfully preserved by tradition, and which is now professed by the pontiff Damasus, of Rome, and Peter, bishop of Alexandria, a man of apostolic holiness. According to the institution of the apostles and the doctrine of the gospel, let us believe in the one Godhead of the Father, the Son, and the Holy Ghost, of equal majesty in the holy Trinity. We order that the adherents of this faith be called Catholic Christians; we brand all the senseless followers of other religions with the infamous name of heretics, and forbid their conventicles assuming the name of churches. Besides the condemnation of divine justice, they must expect the heavy penalties which our authority, guided by heavenly wisdom, shall think proper to inflict.”\footnote{Cod. Theod. xvi, 1, 2. Baronius (Ann.), and even Godefroy call this edict which in this case, to be sure, favored the true doctrine, but involves the absolute despotism of the emperor over faith, an “edictum aureum, pium et salutare.”} In the course of fifteen years this emperor issued at least fifteen penal laws against heretics,\footnote{Comp. Cod. Theod. xvi. tit. v. leg. 6-33, and Godefroy’s Commentary.} by which he gradually deprived them of all right to the exercise of their religion, excluded them from all civil offices, and threatened them with fines, confiscation, banishment, and in some cases, as the Manichaeans, the Audians, and even the Quartodecimanians, with death.

From Theodosius therefore dates the state-church theory of the persecution of heretics, and the embodiment of it in legislation. His primary design, it is true, was rather to terrify and convert, than to punish, the refractory subjects.\footnote{So Sozomen asserts, l. vii. c. 12.}

From the theory, however, to the practice was a single step; and this step his rival and colleague, Maximus, took, when, at the instigation of the unworthy bishop Ithacius, he caused the Spanish bishop, Priscillian, with six respectable adherents of his Manichaean-like sect (two presbyters, two deacons, the poet Latronian, and Euchrocia, a noble matron of Bordeaux), to be tortured and beheaded with the sword at Treves in 385. This was the first shedding of the blood of heretics by a Christian prince for religious opinions. The bishops assembled at Treves, with the exception of Theognistus, approved this act.

But the better feeling of the Christian church shrank from it with horror. The bishops Ambrose of Milan,\footnote{Epist. xxiv. ad Valentin. (tom. ii. p. 891). He would have nothing to do with bishops, “qui aliquos, devios licet a fide, ad necem petebant.”} and Martin of Tours,\footnote{In Sulpic. Sever., Hist. Sacra, ii. 50: “Namque tum Martinus apud Treveros constitutus, non desinebat increpare Ithacium, ut ab accusatione desisteret, Maximum orare, ut sanguine infelicium abstineret: satis superque sufficere, ut episcopali sententia haeretici judicati ecclesiis pellerentur: novum esse et inauditum nefas, ut causam ecclesiae judex saeculi judicaret.” Comp. Sulp. Sev., Dial. iii. c. 11-13, and his Vit. Mart. c. 20.} raised a memorable protest against it, and broke off all communion with Ithacius and the other bishops who had approved the execution. Yet it should not be forgotten that these bishops, at least Ambrose, were committed against the death penalty in general, and in other respects had no indulgence for heathens and heretics.\footnote{Hence Gibbon, ch. xxvii., charges them, not quite groundlessly, with inconsistency: “It is with pleasure that we can observe the human inconsistency of the most illustrious saints and bishops, Ambrose of Milan, and Martin of Tours, who, on this occasion, asserted the cause of toleration. They pitied the unhappy men who had been executed at Treves; they refused to hold communion with their episcopal murderers; and if Martin deviated from that generous resolution, his motives were laudable, and his repentance was exemplary. The bishops of Tours and Milan pronounced, without hesitation, the eternal damnation of heretics; but they were surprised and shocked by the bloody image of their temporal death, and the honest feelings of nature resisted the artificial prejudices of theology.”} The whole thing, too, was irregularly done; on the one hand the bishops appeared as accusers in a criminal cause, and on the other a temporal judge admitted an appeal from the episcopal jurisdiction, and pronounced an opinion in a matter of faith. Subsequently the functions of the temporal and spiritual courts in the trial of heretics were more accurately distinguished.

The execution of the Priscillianists is the only instance of the bloody punishment of heretics in this period, as it is the first in the history of Christianity. But the propriety of violent measures against heresy was thenceforth vindicated even by the best fathers of the church. Chrysostom recommends, indeed, Christian love toward heretics and heathens, and declares against their execution, but approved the prohibition of their assemblies and the confiscation of their churches; and he acted accordingly against the Novatians and the Quartodecimanians, so that many considered his own subsequent misfortunes as condign punishment.\footnote{Hom. xxix. and xlvi. in Matt. Comp. Socrat. H. E. vi. 19. Elsewhere his principle was (in Phocam mart. et c. haer. tom. ii. p. 705): \textgreek{Ἐμοὶ ἔθος ἐστὶ διώκεσθαι καὶ μὴ διώκειν}\emph{;} that is, he himself would rather suffer injury than inflict injury.} Jerome, appealing to Deut. xiii. 6–10, seems to justify even the penalty of death against religious errorists.\footnote{Epist. xxxvii. (al. liii.) ad Riparium Adv. Vigilantium.}

Augustine, who himself belonged nine years to the Manichaean sect, and was wonderfully converted by the grace of God to the Catholic church, without the slightest pressure from without, held at first the truly evangelical view, that heretics and schismatics should not be violently dealt with, but won by instruction and conviction; but after the year 400 he turned and retracted this view, in consequence of his experience with the Donatists, whom he endeavored in vain to convert by disputation and writing, while many submitted to the imperial laws.\footnote{Epist. 93, ad Vincent. § 17: “Mea primitus sententia non erat, nisi neminem ad unitatem Christi esse cogendum, verbo esse agendum, disputatione pugnandum, ratione vincendum, ne fictos catholicos haberemus, quos apertos haereticos noveramus. Sed—he continues § haec opinio mea non contradicentium verbis, sed demonstrantium superabatur exemplis.” Then he adduces his experience with the Donatists. Comp. Retract. ii. 5.} Thenceforth he was led to advocate the persecution of heretics, partly by his doctrine of the Christian state, partly by the seditious excesses of the fanatical Circumcelliones, partly by the hope of a wholesome effect of temporal punishments, and partly by a false interpretation of the Cogite intrare, in the parable of the great supper, Luke xiv. 23.\footnote{The direction: ”\emph{Compel them to come in},” which has often since been abused in defence of coercive measures against heretics, must, of course, be interpreted in harmony with the whole spirit of the gospel, and is only a strong descriptive term in the parable, to signify the fervent zeal in the conversion of the heathen, such as St. Paul manifested without ever resorting to physical coercion.} “It is, indeed, better,” says he, “that men should be brought to serve God by instruction than by fear of punishment or by pain. But because the former means are better, the latter must not therefore be neglected .... Many must often be brought back to their Lord, like wicked servants, by the rod of temporal suffering, before they attain the highest grade of religious development .... The Lord himself orders that the guests be first invited, then compelled, to his great supper.”\footnote{Epist. 185, ad Bonifacium, § 21, § 24.} This father thinks that, if the state be denied the right to punish religious error, neither should she punish any other crime, like murder or adultery, since Paul, in Gal. v. 19, attributes divisions and sects to the same source in the flesh.\footnote{C. Gaudent. Donat. i. § 20. C. Epist. Parmen. i. § 16.} He charges his Donatist opponents with inconsistency in seeming to approve the emperors’ prohibitions of idolatry, but condemning their persecution of Christian heretics. It is to the honor of Augustine’s heart, indeed, that in actual cases he earnestly urged upon the magistrates clemency and humanity, and thus in practice remained true to his noble maxim: “Nothing conquers but truth, the victory of truth is love.”\footnote{“Non vincit nisi veritas, victoria veritatis est caritas.”} But his theory, as Neander justly observes, “contains the germ of the whole system of spiritual despotism, intolerance, and persecution, even to the court of the Inquisition.”\footnote{Kirchengesch. iii. p. 427; Torrey’s ed. ii. p. 217.} The great authority of his name was often afterward made to justify cruelties from which he himself would have shrunk with horror. Soon after him, Leo the Great, the first representative of consistent, exclusive, universal papacy, advocated even the penalty of death for heresy.\footnote{Epist. xv. ad Turribium, where Leo mentions the execution of the Priscillianists with evident approbation: “Etiam mundi principes ita hanc sacrilegam amentiam detestati sunt, ut auctorem ejus cum plerisque discipulis legum publicarum ense prosternerent.”}

Henceforth none but the persecuted parties, from time to time, protested against religious persecution; being made, by their sufferings, if not from principle, at least from policy and self-interest, the advocates of toleration. Thus the Donatist bishop Petilian, in Africa, against whom Augustine wrote, rebukes his Catholic opponents, as formerly his countryman Tertullian had condemned the heathen persecutors of the Christians, for using outward force in matters of conscience; appealing to Christ and the apostles, who never persecuted, but rather suffered and died. “Think you,” says he, “to serve God by killing us with your own hand? Ye err, ye err, if ye, poor mortals, think this; God has not hangmen for priests. Christ teaches us to bear wrong, not to revenge it.” The Donatist bishop Gaudentius says: “God appointed prophets and fishermen, not princes and soldiers, to spread the faith.” Still we cannot forget, that the Donatists were the first who appealed to the imperial tribunal in an ecclesiastical matter, and did not, till after that tribunal had decided against them, turn against the state-church system.

\xchapter{The Rise and Progress of Monasticism}

THE RISE AND PROGRESS OF MONASTICISM.

SOURCES.

1. Greek: Socrates: Hist. Eccles. lib. iv. cap. 23 sqq. Sozomen: H. E. l. i. c. 12–14; iii. 14; vi. 28–34. Palladius (first a monk and disciple of the younger Macarius, then bishop of Helenopolis in Bithynia, ordained by Chrysostom; †431): Historia Lausiaca (Ἱστορία πρὸς Λαῦσον, a court officer under Theodosius II, to whom the work was dedicated), composed about 421, with enthusiastic admiration, from personal acquaintance, of the most celebrated contemporaneous ascetics of Egypt. Theodoret (†457): Historia religiosa, seu ascetica vivendi ratio (φιλόθεος ἱστοπία), biographies of thirty Oriental anchorets and monks, for the most part from personal observation. Nilus the Elder (an anchoret on Mt. Sinai, † about 450): De vita ascetica, De exercitatione monastica, Epistolae 355, and other writings.

2. Latin: Rufinus (†410): Histor. Eremitica, S. Vitae Patrum. Sulpicius Severus (about 400): Dialogi III. (the first dialogue contains a lively and entertaining account of the Egyptian monks, whom he visited; the two others relate to Martin of Tours). Cassianus (†432): Institutiones coenobiales, and Collationes Patrum (spiritual conversations of eastern monks).

Also the ascetic writings of Athanasius (Vita Antonii), Basil, Gregory Nazianzen, Chrysostom, Nilus, Isidore of Pelusium, among the Greek; Ambrose, Augustine, Jerome (his Lives of anchorets, and his letters), Cassiodorus, and Gregory the Great, among the Latin fathers.

LATER LITERATURE.

L. Holstenius (born at Hamburg 1596, a Protest., then a Romanist convert, and librarian of the Vatican): Codex regularum monastic., first Rom. 1661; then, enlarged, Par. and Augsb. in 6 vols. fol. The older Greek Menologia (μηνολόγια), and Menaea (μηναῖα), and the Latin Calendaria and Martyrologia, i.e. church calendars or indices of memorial days (days of the earthly death and heavenly birth) of the saints, with short biographical notices for liturgical use. P. Herbert Rosweyde (Jesuit): Vitae Patrum, sive Historiae Eremiticae, libri x. Antw. 1628. Acta Sanctorum, quotquot toto orbe coluntur, Antw. 1643–1786, 53 vols. fol. (begun by the Jesuit Bollandus, continued by several scholars of his order, called Bollandists, down to the 11th Oct. in the calendar of saints’ days, and resumed in 1845, after long interruption, by Theiner and others). D’achery and Mabillon (Benedictines): Acta Sanctorum ordinis S. Benedicti, Par. 1668–1701, 9 vols. fol. (to 1100). Pet. Helyot (Franciscan): Histoire des ordres monastiques religieux et militaires, Par. 1714–’19, 8 vols. 4to. Alban Butler (R.C.): The Lives of the Fathers, Martyrs, and other principal Saints (arranged according to the Catholic calendar, and completed to the 31st Dec.), first 1745; often since (best ed. Lond. 1812–’13) in 12 vols.; another, Baltimore, 1844, in 4 vols). Gibbon: Chap. xxxvii. (Origin, Progress, and Effects of Monastic Life; very unfavorable, and written in lofty philosophical contempt). Henrion (R.C.): Histoire des ordres religieux, Par. 1835 (deutsch bearbeitet von S. Fehr, Tüb. 1845, 2 vols.). F. v. Biedenfeld: Ursprung u. s. w. saemmtlicher Mönchsorden im Orient u. Occident, Weimar, 1837, 3 vols. Schmidt (R.C.): Die Mönchs-, Nonnen-, u. geistlichen Ritterorden nebst Ordensregeln u. Abbildungen., Augsb. 1838, sqq. H. H. Milman (Anglican): History of Ancient Christianity, 1844, book iii. ch. 11. H. Ruffner (Presbyterian): The Fathers of the Desert, New York, 1850, 2 vols. (full of curious information, in popular form). Count de Montalembert (R.C.): Les Moines d’Occident depuis St. Bénoit jusqu’à St. Bernard, Par. 1860, sqq. (to embrace 6 vols.); transl. into English: The Monks of the West, etc., Edinb. and Lond. 1861, in 2 vols. (vol. i. gives the history of monasticism before St. Benedict, vol. ii. is mainly devoted to St. Benedict; eloquently eulogistic of, and apologetic for, monasticism). Otto Zöckler: Kritische Geschichte der Askese. Frankf. a. M. 1863. Comp. also the relevant sections of Tillemont, Fleury, Schröckh (vols. v. and viii.), Neander, and Gieseler.

\xsection{Origin of Christian Monasticism. Comparison with other forms of Asceticism}

§ 28. Origin of Christian Monasticism. Comparison with other forms of Asceticism.

Hospinian: De origine et progressu monachatus, l. vi., Tig. 1588, and enlarged, Genev. 1669, fol. J. A. Möhler (R.C.): Geschichte des Mönchthums in der Zeit seiner Entstehung u. ersten Ausbildung, 1836 (in his collected works, Regensb. vol. ii. p. 165 sqq.). Isaac Taylor (Independent): Ancient Christianity, Lond. 1844, vol. i. p. 299 sqq. A. Vogel: Ueber das Mönchthum, Berl. 1858 (in the “Deutsche Zeitschrift für christl. Wissenschaft,” etc.). P. Schaff: Ueber den Ursprung und Charakter des Mönchthums (in Dorner’s, etc. “Jahrbücher für deutsche Theol.,” 1861, p. 555 ff.). J. Cropp: Origenes et causae monachatus. Gott. 1863.

In the beginning of the fourth century monasticism appears in the history of the church, and thenceforth occupies a distinguished place. Beginning in Egypt, it spread in an irresistible tide over the East and the West, continued to be the chief repository of the Christian life down to the times of the Reformation, and still remains in the Greek and Roman churches an indispensable institution and the most productive seminary of saints, priests, and missionaries.

With the ascetic tendency in general, monasticism in particular is found by no means only in the Christian church, but in other religions, both before and after Christ, especially in the East. It proceeds from religious seriousness, enthusiasm, and ambition; from a sense of the vanity of the world, and an inclination of noble souls toward solitude, contemplation, and freedom from the bonds of the flesh and the temptations of the world; but it gives this tendency an undue predominance over the social, practical, and world-reforming spirit of religion. Among the Hindoos the ascetic system may be traced back almost to the time of Moses, certainly beyond Alexander the Great, who found it there in full force, and substantially with the same characteristics which it presents at the present day.\footnote{Comp. the occasional notices of the Indian gymnosophists in Strabo (lib. xv. cap. 1, after accounts from the time of Alexander the Great), Arrian (Exped. Alex. l. vii. c. 1-3, and Hist. Ind. c. 11), Plinius (Hist Nat. vii. 2), Diodorus Siculus (lib. ii.), Plutarch (Alex. 64), Porphyry (De abstinent. l. iv.), Lucian (Fugit. 7), Clemens Alex. (Strom. l. i. and iii.), and Augustine(De Civit. Dei, l. xiv. c. 17: “Per opacas Indiae solitudines, quum quidam nudi philosophentur, unde gymnosophistae nominantur; adhibent tamen genitalibus tegmina, quibus per caetera membrorum carent;” and l. xv. 20, where he denies all merit to their celibacy, because it is not “secundum fidem summi boni, qui est Deus”). With these ancient representations agree the narratives of Fon Koueki (about 400, translated by M. A. Rémusat, Par. 1836), Marco Polo (1280), Bernier (1670), Hamilton (1700), Papi, Niebuhr, Orlich, Sonnerat, and others.} Let us consider it a few moments.

The Vedas, portions of which date from the fifteenth century before Christ, the Laws of Menu, which were completed before the rise of Buddhism, that is, six or seven centuries before our era, and the numerous other sacred books of the Indian religion, enjoin by example and precept entire abstraction of thought, seclusion from the world, and a variety of penitential and meritorious acts of self-mortification, by which the devotee assumes a proud superiority over the vulgar herd of mortals, and is absorbed at last into the divine fountain of all being. The ascetic system is essential alike to Brahmanism and Buddhism, the two opposite and yet cognate branches of the Indian religion, which in many respects are similarly related to each other as Judaism is to Christianity, or also as Romanism to Protestantism. Buddhism is a later reformation of Brahmanism; it dates probably from the sixth century before Christ (according to other accounts much earlier), and, although subsequently expelled by the Brahmins from Hindostan, it embraces more followers than any other heathen religion, since it rules in Farther India, nearly all the Indian islands, Japan, Thibet, a great part of China and Central Asia to the borders of Siberia. But the two religions start from opposite principles. Brahmanic asceticism\footnote{The Indian word for it is \emph{tapas,} i.e. the burning out, or the extinction of the individual being and its absorption into the essence of Brahma.} proceeds from a pantheistic view of the world, the Buddhistic from an atheistic and nihilistic, yet very earnest view; the one if; controlled by the idea of the absolute but abstract unity and a feeling of contempt of the world, the other by the idea of the absolute but unreal variety and a feeling of deep grief over the emptiness and nothingness of all existence; the one is predominantly objective, positive, and idealistic, the other more subjective, negative, and realistic; the one aims at an absorption into the universal spirit of Brahm, the other consistently at an absorption into nonentity, if it be true that Buddhism starts from an atheistic rather than a pantheistic or dualistic basis. “Brahmanism”—says a modern writer on the subject\footnote{Ad. Wuttke, in his able and instructive work: Das Geistesleben der Chinesen, Japaner, und Indier(second part of his History of Heathenism), 1853, p. 593.}—“looks back to the beginning, Buddhism to the end; the former loves cosmogony, the latter eschatology. Both reject the existing world; the Brahman despises it, because he contrasts it with the higher being of Brahma, the Buddhist bewails it because of its unrealness; the former sees God in all, the other emptiness in all.” Yet as all extremes meet, the abstract all-entity of Brahmanism and the equally abstract non-entity or vacuity of Buddhism come to the same thing in the end, and may lead to the same ascetic practices. The asceticism of Brahmanism takes more the direction of anchoretism, while that of Buddhism exists generally in the social form of regular convent life.

The Hindoo monks or gymnosophists (naked philosophers), as the Greeks called them, live in woods, caves, on mountains, or rocks, in poverty, celibacy, abstinence, silence: sleeping on straw or the bare ground, crawling on the belly, standing all day on tiptoe, exposed to the pouring rain or scorching sun with four fires kindled around them, presenting a savage and frightful appearance, yet greatly revered by the multitude, especially the women, and performing miracles, not unfrequently completing their austerities by suicide on the stake or in the waves of the Ganges. Thus they are described by the ancients and by modern travellers. The Buddhist monks are less fanatical and extravagant than the Hindoo Yogis and Fakirs. They depend mainly on fasting, prayer, psalmody, intense contemplation, and the use of the whip, to keep their rebellious flesh in subjection. They have a fully developed system of monasticism in connection with their priesthood, and a large number of convents; also nunneries for female devotees. The Buddhist monasticism, especially in Thibet, with its vows of celibacy, poverty, and obedience, its common meals, readings, and various pious exercises, bears such a remarkable resemblance to that of the Roman Catholic church that Roman missionaries thought it could be only explained as a diabolical imitation.\footnote{See the older accounts of Catholic missionaries to Thibet, in Pinkerton’s Collection of Voyages and Travels, vol. vii., and also the recent work of Huc, a French missionary priest of the congregation of St. Lazare: Souvenirs d’un Voyage dans la Tartarie, le Thibet, et la Chine, pendant les années1844-1846. Comp. also on the whole subject the two works of R. \emph{S.} Hardy\emph{:} “Eastern Monachism” and “A Manual of Buddhism in its modern development, translated from Singalese MSS.” Lond. 1850. The striking affinity between Buddhism and Romanism extends, by the way, beyond monkery and convent life to the heirarchical organization, with the Grand Lama for pope, and to the worship, with its ceremonies, feasts, processions, pilgrimages, confessional, a kind of mass, prayers for the dead, extreme unction, \&c. The view is certainly at least plausible, to which the great geographer Carl Ritter (Erdkunde, ii. p. 283-299, 2d ed.) has given the weight of his name, that the Lamaists in Thibet borrowed their religious forms and ceremonies in part from the Nestorian missionaries. But this view is a mere hypothesis, and is rendered improbable by the fact, that Buddhism in Cochin China, Tonquin, and Japan, where no Nestorian missionaries ever were, shows the same striking resemblance to Romanism as the Lamaism of Thibet, Tartary, and North China. Respecting the singular tradition of Prester John, or the Christian priest-king in Eastern Asia, which arose about the eleventh century, and respecting the Nestorian missions, see Ritter, l.c.} But the original always precedes the caricature, and the ascetic system was completed in India long before the introduction of Christianity, even if we should trace this back to St. Bartholomew and St. Thomas.

The Hellenic heathenism was less serious and contemplative, indeed, than the Oriental; yet the Pythagoreans were a kind of monastic society, and the Platonic view of matter and of body not only lies at the bottom of the Gnostic and Manichaean asceticism, but had much to do also with the ethics of Origen and the Alexandrian School.

Judaism, apart from the ancient Nazarites,\footnote{Comp. Num. vi. 1-21.} had its Essenes in Palestine\footnote{Comp. the remarkable description of these Jewish monks by the elder Pliny, Hist. Natur. v. 15: “Gens sola, et in toto orbe praeter caeteros mira, sine ulla femina, omni venere abdicata, sine pecunia, socia palmarum. Ita per seculorum millia (incredibile dictu) gens aeterna est in qua nemo nascitur. Tam foecunda illis aliorum vitae penitentia est.”} and its Therapeutae in Egypt;\footnote{Eusebius, H. E. ii. 17, erroneously takes them for Christians.} though these betray the intrusion of foreign elements into the Mosaic religion, and so find no mention in the New Testament.

Lastly, Mohammedanism, though in mere imitation of Christian and pagan examples, has, as is well known, its dervises and its cloisters.\footnote{H. Ruffner, l.c. vol. i. ch. ii.–ix., gives an extended description of these extra-Christian forms of monasticism, and derives the Christian from them, especially from the Buddhist.}

Now were these earlier phenomena the source, or only analogies, of the Christian monasticism? That a multitude of foreign usages and rites made their way into the church in the age of Constantine, is undeniable. Hence many have held, that monasticism also came from heathenism, and was an apostasy from apostolic Christianity, which Paul had plainly foretold in the Pastoral Epistles.\footnote{So even Calvin, who, in his commentary on 1 Tim. iv. 3, refers Paul’s prophecy of the ascetic apostasy primarily to the Encratites, Gnostics, Montanists, and Manichaeans, but extends it also to the Papists, “quando coelibatum et ciborum abstinentiam severius urgent quam ullum Dei praeceptum.” So, recently, Ruffner, and especially Is. Taylor, who, in his “Ancient Christianity,” vol. i. p. 299 sqq., has a special chapter on The Predicted Ascetic Apostasy. The best modern interpreters, however, are agreed, that the apostle has the heretical Gnostic dualistic asceticism in his eye, which forbade marriage and certain meats as intrinsically impure; whereas the Roman and Greek churches make marriage a sacrament, only subordinate it to celibacy, and limit the prohibition of it to priests and monks. The application of 1 Tim. iv. 1-3 to the Catholic church is, therefore, admissible at most only in a partial and indirect way.} But such a view can hardly be reconciled with the great place of this phenomenon in history; and would, furthermore, involve the entire ancient church, with its greatest and best representatives both east and west, its Athanasius, its Chrysostom, its Jerome, its Augustine, in the predicted apostasy from the faith. And no one will now hold, that these men, who all admired and commended the monastic life, were antichristian errorists, and that the few and almost exclusively negative opponents of that asceticism, as Jovinian, Helvidius, and Vigilantius, were the sole representatives of pure Christianity in the Nicene and next following age.

In this whole matter we must carefully distinguish two forms of asceticism, antagonistic and irreconcilable in spirit and principle, though similar in form: the Gnostic dualistic, and the Catholic. The former of these did certainly come from heathenism; but the latter sprang independently from the Christian spirit of self-denial and longing for moral perfection, and, in spite of all its excrescences, has fulfilled an important mission in the history of the church.

The pagan monachism, the pseudo-Jewish, the heretical Christian, above all the Gnostic and Manichaean, is based on in irreconcilable metaphysical dualism between mind and matter; the Catholic Christian Monachism arises from the moral conflict between the spirit and the flesh. The former is prompted throughout by spiritual pride and selfishness; the latter, by humility and love to God and man. The false asceticism aims at annihilation of the body and pantheistic absorption of the human being in the divine; the Christian strives after the glorification of the body and personal fellowship with the living God in Christ. And the effects of the two are equally different. Though it is also unquestionable, that, notwithstanding this difference of principle, and despite the condemnation of Gnosticism and Manichaeism, the heathen dualism exerted a powerful influence on the Catholic asceticism and its view of the world, particularly upon anchoretism and monasticism in the East, and has been fully overcome only in evangelical Protestantism. The precise degree of this influence, and the exact proportion of Christian and heathen ingredients in the early monachism of the church, were an interesting subject of special investigation.

The germs of the Christian monasticism may be traced as far back as the middle of the second century, and in fact faintly even in the anxious ascetic practices of some of the Jewish Christians in the apostolic age. This asceticism, particularly fasting and celibacy, was commended more or less distinctly by the most eminent ante-Nicene fathers, and was practised, at least partially, by a particular class of Christians (by Origen even to the unnatural extreme of self-emasculation).\footnote{Comp. vol. i. § 94-97.} So early as the Decian persecution, about the year 250, we meet also the first instances of the flight of ascetics or Christian philosophers into the wilderness; though rather in exceptional cases, and by way of escape from personal danger. So long as the church herself was a child of the desert, and stood in abrupt opposition to the persecuting world, the ascetics of both sexes usually lived near the congregations or in the midst of them, often even in the families, seeking there to realize the ideal of Christian perfection. But when, under Constantine, the mass of the population of the empire became nominally Christian, they felt, that in this world-church, especially in such cities as Alexandria, Antioch, and Constantinople, they were not at home, and voluntarily retired into waste and desolate places and mountain clefts, there to work out the salvation of their souls undisturbed.

Thus far monachism is a reaction against the secularizing state-church system and the decay of discipline, and an earnest, well-meant, though mistaken effort to save the virginal purity of the Christian church by transplanting it in the wilderness. The moral corruption of the Roman empire, which had the appearance of Christianity, but was essentially heathen in the whole framework of society, the oppressiveness of taxes\footnote{Lactantius says it was necessary to buy even the liberty of breathing, and according to Zosimus (Hist. ii. 38) the fathers prostituted their daughters to have means to pay their tax.} the extremes of despotism and slavery, of extravagant luxury and hopeless poverty, the repletion of all classes, the decay of all productive energy in science and art, and the threatening incursions of barbarians on the frontiers—all favored the inclination toward solitude in just the most earnest minds.

At the same time, however, monasticism afforded also a compensation for martyrdom, which ceased with the Christianization of the state, and thus gave place to a voluntary martyrdom, a gradual self-destruction, a sort of religious suicide. In the burning deserts and awful caverns of Egypt and Syria, amidst the pains of self-torture, the mortification of natural desires, and relentless battles with hellish monsters, the ascetics now sought to win the crown of heavenly glory, which their predecessors in the times of persecution had more quickly and easily gained by a bloody death.

The native land of the monastic life was Egypt, the land where Oriental and Grecian literature, philosophy, and religion, Christian orthodoxy and Gnostic heresy, met both in friendship and in hostility. Monasticism was favored and promoted here by climate and geographic features, by the oasis-like seclusion of the country, by the bold contrast of barren deserts with the fertile valley of the Nile, by the superstition, the contemplative turn, and the passive endurance of the national character, by the example of the Therapeutae, and by the moral principles of the Alexandrian fathers; especially by Origen’s theory of a higher and lower morality and of the merit of voluntary poverty and celibacy. Aelian says of the Egyptians, that they bear the most exquisite torture without a murmur, and would rather be tormented to death than compromise truth. Such natures, once seized with religious enthusiasm, were eminently qualified for saints of the desert.

\xsection{Development of Monasticism}

§ 29. Development of Monasticism.

In the historical development of the monastic institution we must distinguish four stages. The first three were completed in the fourth century; the remaining one reached maturity in the Latin church of the middle age.

The first stage is an ascetic life as yet not organized nor separated from the church. It comes down from the ante-Nicene age, and has been already noticed. It now took the form, for the most part, of either hermit or coenobite life, but continued in the church itself, especially among the clergy, who might be called half monks.

The second stage is hermit life or anchoretism.\footnote{From \textgreek{ἀναχωρέω}, to retire (from human society), \textgreek{ἀναχωρητής}, \textgreek{ἐρημίτης}(from \textgreek{ἐρημία}, a desert). The word \textgreek{μοναχός}(from \textgreek{μόνος}, alone, and \textgreek{μονάζειν}, to live alone), monachus (whence monk), also points originally to solitary, hermit life, but is commonly synonymous with coenobite or friar.} It arose in the beginning of the fourth century, gave asceticism a fixed and permanent shape, and pushed it to even external separation from the world. It took the prophets Elijah and John the Baptist for its models, and went beyond them. Not content with partial and temporary retirement from common life, which may be united with social intercourse and useful labors, the consistent anchoret secludes himself from all society, even from kindred ascetics, and comes only exceptionally into contact with human affairs, either to receive the visits of admirers of every class, especially of the sick and the needy (which were very frequent in the case of the more celebrated monks), or to appear in the cities on some extraordinary occasion, as a spirit from another world. His clothing is a hair shirt and a wild beast’s skin; his food, bread and salt; his dwelling, a cave; his employment, prayer, affliction of the body, and conflict with satanic powers and wild images of fancy. This mode of life was founded by Paul of Thebes and St. Anthony, and came to perfection in the East. It was too eccentric and unpractical for the West, and hence less frequent there, especially in the rougher climates. To the female sex it was entirely unsuited. There was a class of hermits, the Sarabaites in Egypt, and the Rhemoboths in Syria, who lived in bands of at least two or three together; but their quarrelsomeness, occasional intemperance, and opposition to the clergy, brought them into ill repute.

The third step in the progress of the monastic life brings us to coenobitism or cloister life, monasticism in the ordinary sense of the word.\footnote{\textgreek{Κοινόβιον}, coenobium; from \textgreek{κοινός βίος}, vita communis; then the congregation of monks; sometimes also used for the building. In the same sense \textgreek{μάνδρα}, stable, fold, and \textgreek{μοναστήριον}, claustrum (whence cloister). Also \textgreek{λαύραι}, laurae (literally, streets), that is cells, of which usually a number were built not far apart, so as to form a hamlet. Hence this term is often used in the same sense as monasterium. The singular, \textgreek{λαῦρα}, however, answers to the anchoret life. On this nomenclature of monasticism comp. Du Cange, in the Glossarium mediae et infimae Latinitatis, under the respective words.} It originated likewise in Egypt, from the example of the Essenes and Therapeutae, and was carried by St. Pachomius to the East, and afterward by St. Benedict to the West. Both these ascetics, like the most celebrated order-founders of later days, were originally hermits. Cloister life is a regular organization of the ascetic life on a social basis. It recognizes, at least in a measure, the social element of human nature, and represents it in a narrower sphere secluded from the larger world. As hermit life often led to cloister life, so the cloister life was not only a refuge for the spirit weary of the world, but also in many ways a school for practical life in the church. It formed the transition from isolated to social Christianity. It consists in an association of a number of anchorets of the same sex for mutual advancement in ascetic holiness. The coenobites live, somewhat according to the laws of civilization, under one roof, and under a superintendent or abbot.\footnote{\textgreek{Ἡγούμενος, ἀρχεμανδρίτης , ἀββᾶς}, i.e. father, hence abbot. A female superintendent was called in Syriac \textgreek{ἀμμᾶς}, mother, abbess.} They divide their time between common devotions and manual labor, and devote their surplus provisions to charity; except the mendicant monks, who themselves live by alms. In this modified form monasticism became available to the female sex, to which the solitary desert life was utterly impracticable; and with the cloisters of monks, there appear at once cloisters also of nuns.\footnote{From \emph{nonna}, i.e. casta, chaste, holy. The word is probably of Coptic origin, and occurs as early as in Jerome. The masculine \emph{nonnus}, monk, appears frequently in the middle age. Comp. the examples in Du Cange, s. v.} Between the anchorets and the coenobites no little jealousy reigned; the former charging the latter with ease and conformity to the world; the latter accusing the former of selfishness and misanthropy. The most eminent church teachers generally prefer the cloister life. But the hermits, though their numbers diminished, never became extinct. Many a monk was a hermit first, and then a coenobite; and many a coenobite turned to a hermit.

The same social impulse, finally, which produced monastic congregations, led afterward to monastic orders, unions of a number of cloisters under one rule and a common government. In this fourth and last stage monasticism has done most for the diffusion of Christianity and the advancement of learning,\footnote{Hence Middleton says, not without reason: “By all which I have ever read of the old, and have seen of the modern monks, I take the preference to be clearly due to the last, as having a more regular discipline, more good learning, and less superstition among them than the first.”} has fulfilled its practical mission in the Roman Catholic church, and still wields a mighty influence there. At the same time it became in some sense the cradle of the German reformation. Luther belonged to the order of St. Augustine, and the monastic discipline of Erfurt was to him a preparation for evangelical freedom, as the Mosaic law was to Paul a schoolmaster to lead to Christ. And for this very reason Protestantism is the end of the monastic life.

\xsection{Nature and Aim of Monasticism}

§ 30. Nature and Aim of Monasticism.

Monasticism was from the first distinguished as the contemplative life from the practical.\footnote{\textgreek{Βίος θεωρητικός ,} and \textgreek{βίος πρακτικός}, according to Gregory Nazianzen and others. Throughout the middle age the distinction between the \emph{vita contemplativa} and the \emph{vita activa} was illustrated by the two sisters of Lazarus, Luke x. 38-42.} It passed with the ancient church for the true, the divine, or Christian philosophy,\footnote{\textgreek{Ἡ κατὰ θεὸν} or \textgreek{Χριστὸν φιλοσοφία, ἡ ὑψηλή φιλος}., i.e. in the sense of the ancients, not so much a speculative system, as a mode of life under a particular rule. So in the Pythagoreans, Stoics, Cynics, and Neo-Platonists. Ascetic and philosopher are the same.} an unworldly purely apostolic, angelic life.\footnote{\textgreek{Ἀποστολικὸς βίος , ὁ τῶν ἀγγέλων βίος}, vita angelica; after an unwarranted application of Christ’s word respecting the sexless life of the angels, Matt. xxii. 30, which is not presented here as a model for imitation, but only mentioned as an argument against the Sadducees.} It rests upon an earnest view of life; upon the instinctive struggle after perfect dominion of the spirit over the flesh, reason over sense, the supernatural over the natural, after the highest grade of holiness and an undisturbed communion of the soul with God; but also upon a morbid depreciation of the body, the family, the state, and the divinely established social order of the world. It recognizes the world, indeed, as a creature of God, and the family and property as divine institutions, in opposition to the Gnostic Manichaean asceticism, which ascribes matter as such to an evil principle. But it makes a distinction between two grades of morality: a common and lower grade, democratic, so to speak, which moves in the natural ordinances of God; and a higher, extraordinary, aristocratic grade, which lies beyond them and is attended with special merit. It places the great problem of Christianity not in the transformation, but in the abandonment, of the world. It is an extreme unworldliness, over against the worldliness of the mass of the visible church in union with the state. It demands entire renunciation, not only of sin, but also of property and of marriage, which are lawful in themselves, ordained by God himself, and indispensable to the continuance and welfare of the human race. The poverty of the individual, however, does not exclude the possession of common property; and it is well known, that some monastic orders, especially the Benedictines, have in course of time grown very rich. The coenobite institution requires also absolute obedience to the will of the superior, as the visible representative of Christ. As obedience to orders and sacrifice of self is the first duty of the soldier, and the condition of military success and renown, so also in this spiritual army in its war against the flesh, the world, and the devil, monks are not allowed to have a will of their own. To them may be applied the lines of Tennyson:\footnote{ln his famous battle poem: “The Charge of the Light Brigade at Balaclava,” first ed. 1854.}

Voluntary poverty, voluntary celibacy, and absolute obedience form the three monastic vows, as they are called, and are supposed to constitute a higher virtue and to secure a higher reward in heaven.

But this threefold self-denial is only the negative side of the matter, and a means to an end. It places man beyond the reach of the temptations connected with earthly possessions, married life, and independent will, and facilitates his progress toward heaven. The positive aspect of monasticism is unreserved surrender of the whole man, with all his time and strength, to God; though, as we have said, not within, but without the sphere of society and the order of nature. This devoted life is employed in continual prayer, meditation, fasting, and castigation of the body. Some votaries went so far as to reject all bodily employment, for its interference with devotion. But in general a moderate union of spiritual exercises with scientific studies or with such manual labor as agriculture, basket making, weaving, for their own living and the support of the poor, was held not only lawful but wholesome for monks. It was a proverb, that a laborious monk was beset by only one devil; an idle one, by a legion.

With all the austerities and rigors of asceticism, the monastic life had its spiritual joys and irresistible charms for noble, contemplative, and heaven-aspiring souls, who fled from the turmoil and vain show of the city as a prison, and turned the solitude into a paradise of freedom and sweet communion with God and his saints; while to others the same solitude became a fruitful nursery of idleness, despondency, and the most perilous temptations and ultimate ruin.\footnote{Comp. the truthful remark of Yves de Chartres, of the twelfth century, Ep. 192 (quoted by Montalembert): “Non beatum faciunt hominem secreta sylvarum, cacumina montium, si secum non habet solitudinem mentis, sabbatum cordis, tranquillitatem conscientiae, ascensiones in corde, sine quibus omnem solitudinem comitantur mentis acedia, curiositas, vana gloria, periculosae tentationum procellae.”}

\xsection{Monasticism and the Bible}

§ 31. Monasticism and the Bible.

Monasticism, therefore, claims to be the highest and purest form of Christian piety and virtue, and the surest way to heaven. Then, we should think, it must be preëminently commended in the Bible, and actually exhibited in the life of Christ and the apostles. But just in this biblical support it falls short.

The advocates of it uniformly refer first to the examples of Elijah, Elisha, and John the Baptist;\footnote{So Jerome, Ep. 49 (ed. Ben.), ad Paulinum, where he adduces, besides Elijah and John, Isaiah also and the sons of the prophets, as the fathers of monasticism; and in his Vita Pauli, where, however, he more correctly designates Paul of Thebes and Anthonyas the first hermits, properly so called, in distinction from the prophets. Comp. also Sozomen: H. E., 1. i. c. 12: \textgreek{Ταύτης δὲ τῆς ἀρίστης φιλοσοφίας ἤρξατο, ωὝς τινες λέγουσιν, Ἡλίας ὁ προφήτης καὶ Ἰωάννης ὁ βαπτιστής}. This appeal to the example of Elijah and John the Baptist has become traditional with Catholic writers on the subject. Alban Butler says, under Jan. 15, in the life of Paul of Thebes: “Elias and John the Baptist sanctified the deserts, and Jesus Christ himself was a model of the eremitical state during his forty days’ fast in the wilderness; neither is it to be questioned but the Holy Ghost conducted the saint of this day (Paul of Thebes) into the desert, and was to him an instructor there.”} but these stand upon the legal level of the Old Testament, and are to be looked upon as extraordinary personages of an extraordinary age; and though they may be regarded as types of a partial anchoretism (not of cloister life), still they are nowhere commended to our imitation in this particular, but rather in their influence upon the world.

The next appeal is to a few isolated passages of the New Testament, which do not, indeed, in their literal sense require the renunciation of property and marriage, yet seem to recommend it as a special, exceptional form of piety for those Christians who strive after higher perfection.\footnote{Hence called \emph{consilia evangelica,} in distinction from \emph{mandata} \emph{divina}; after 1 Cor. vii. 25, where Paul does certainly make a similar distinction. The \emph{consilium} and \emph{votum} \emph{paupertatis} is based on Matt. xix. 21; the \emph{votum} \emph{castitatis,} on 1 Cor. vii. 8, 25, 38-40. For the \emph{votum} \emph{obedientiae} no particular text is quoted. The theory appears substantially as early as in Origen, and was in him not merely a personal opinion, but the reflex of a very widely spread practice. Comp. vol. i. § 94 and 95.}

Finally, as respects the spirit of the monastic life, reference is sometimes made even to the poverty of Christ and his apostles, to the silent, contemplative Mary, in contrast with the busy, practical Martha, and to the voluntary community of goods in the first Christian church in Jerusalem.

But this monastic interpretation of primitive Christianity mistakes a few incidental points of outward resemblance for essential identity, measures the spirit of Christianity by some isolated passages, instead of explaining the latter from the former, and is upon the whole a miserable emaciation and caricature. The gospel makes upon all men virtually the same moral demand, and knows no distinction of a religion for the masses and another for the few.

Jesus, the model for all believers, was neither a coenobite, nor an anchoret, nor an ascetic of any kind, but the perfect pattern man for universal imitation. There is not a trace of monkish austerity and ascetic rigor in his life or precepts, but in all his acts and words a wonderful harmony of freedom and purity, of the most comprehensive charity and spotless holiness. He retired to the mountains and into solitude, but only temporarily, and for the purpose of renewing his strength for active work. Amidst the society of his disciples, of both sexes, with kindred and friends, in Cana and Bethany, at the table of publicans and sinners, and in intercourse with all classes of the people, he kept himself unspotted from the world, and transfigured the world into the kingdom of God. His poverty and celibacy have nothing to do with asceticism, but represent, the one the condescension of his redeeming love, the other his ideal uniqueness and his absolutely peculiar relation to the whole church, which alone is fit or worthy to be his bride. No single daughter of Eve could have been an equal partner of the Saviour of mankind, or the representative head of the new creation.

The example of the sister of Lazarus proves only, that the contemplative life may dwell in the same house with the practical, and with the other sex, but justifies no separation from the social ties.

The life of the apostles and primitive Christians in general was anything but a hermit life; else had not the gospel spread so quickly to all the cities of the Roman world. Peter was married, and travelled with his wife as a missionary. Paul assumes one marriage of the clergy as the rule, and notwithstanding his personal and relative preference for celibacy in the then oppressed condition of the church, he is the most zealous advocate of evangelical freedom, in opposition to all legal bondage and anxious asceticism.

Monasticism, therefore, in any case, is not the normal form of Christian piety. It is an abnormal phenomenon, a humanly devised service of God,\footnote{Comp. Col. ii. 16-23.} and not rarely a sad enervation and repulsive distortion of the Christianity of the Bible. And it is to be estimated, therefore, not by the extent of its self-denial, not by its outward acts of self-discipline (which may all be found in heathenism, Judaism, and Mohammedanism as well), but by the Christian spirit of humility and love which animated it. For humility is the groundwork, and love the all-ruling principle, of the Christian life, and the distinctive characteristic of the Christian religion. Without love to God and charity to man, the severest self-punishment and the utmost abandonment of the world are worthless before God.\footnote{Comp. 1 Cor. xiii. 1-3. Comp. p. 168 sq.}

\xsection{Lights and Shades of Monastic Life}

§ 32. Lights and Shades of Monastic Life.

The contrast between pure and normal Bible-Christianity and abnormal Monastic Christianity, will appear more fully if we enter into a close examination of the latter as it actually appeared in the ancient church.

The extraordinary rapidity with which this world-forsaking form of piety spread, bears witness to a high degree of self-denying moral earnestness, which even in its mistakes and vagrancies we must admire. Our age, accustomed and wedded to all possible comforts, but far in advance of the Nicene age in respect to the average morality of the masses, could beget no such ascetic extremes. In our estimate of the diffusion and value of monasticism, the polluting power of the theatre, oppressive taxation, slavery, the multitude of civil wars, and the hopeless condition of the Roman empire, must all come into view. Nor must we, by any means, measure the moral importance of this phenomenon by numbers. Monasticism from the beginning attracted persons of opposite character and from opposite motives. Moral earnestness and religious enthusiasm were accompanied here, as formerly in martyrdom, though even in larger measure than there, with all kinds of sinister motives; indolence, discontent, weariness of life, misanthropy, ambition for spiritual distinction, and every sort of misfortune or accidental circumstance. Palladius, to mention but one illustrious example, tells of Paul the Simple,\footnote{\textgreek{Ἄπλαστος}, lit. \emph{not moulded}; hence \emph{natural, sincere}.} that, from indignation against his wife, whom he detected in an act of infidelity, he hastened, with the current oath of that day, “in the name of Jesus,”\footnote{\textgreek{Μὰ τὸν Ἰησοῦν} \emph{(per Christum,} in Salvian), which now took the place of the pagan oath: \textgreek{μὰ τὸν Δία}, by Jupiter.} into the wilderness; and immediately, though now sixty years old, under the direction of Anthony, he became a very model monk, and attained an astonishing degree of humility, simplicity, and perfect submission of will.

In view of these different motives we need not be surprised that the moral character of the monks varied greatly, and presents opposite extremes. Augustine says he found among the monks and nuns the best and the worst of mankind.

Looking more closely, in the first place, at anchoretism, we meet in its history unquestionably many a heroic character, who attained an incredible mastery over his sensual nature, and, like the Old Testament prophets and John the Baptist, by their mere appearance and their occasional preaching, made an overwhelming impression on his contemporaries, even among the heathen. St. Anthony’s visit to Alexandria was to the gazing multitude like the visit of a messenger from the other world, and resulted in many conversions. His emaciated face, the glare of his eye, his spectral yet venerable form, his contempt of the world, and his few aphoristic sentences told more powerfully on that age and people than a most elaborate sermon. St. Symeon, standing on a column from year to year, fasting, praying, and exhorting the visitors to repentance, was to his generation a standing miracle and a sign that pointed them to heaven. Sometimes, in seasons of public calamity, such hermits saved whole cities and provinces from the imperial wrath, by their effectual intercessions. When Theodosius, in 387, was about to destroy Antioch for a sedition, the hermit Macedonius met the two imperial commissaries, who reverently dismounted and kissed his hands and feet; he reminded them and the emperor of their own weakness, set before them the value of men as immortal images of God, in comparison with the perishable statues of the emperor, and thus saved the city from demolition.\footnote{In Theodoret: Hist. relig. c. (vita) 13.} The heroism of the anchoretic life, in the voluntary renunciation of lawful pleasures and the patient endurance of self-inflicted pains, is worthy of admiration in its way, and not rarely almost incredible.

But this moral heroism—and these are the weak points of it—oversteps not only the present standard of Christianity, but all sound measure; it has no support either in the theory or the practice of Christ and the apostolic church; and it has far more resemblance to heathen than to biblical precedents. Many of the most eminent saints of the desert differ only in their Christian confession, and in some Bible phrases learnt by rote, from Buddhist fakirs and Mohammedan dervises. Their highest virtuousness consisted in bodily exercises of their own devising, which, without love, at best profit nothing at all, very often only gratify spiritual vanity, and entirely obscure the gospel way of salvation.

To illustrate this by a few examples, we may choose any of the most celebrated eastern anchorets of the fourth and fifth centuries, as reported by the most credible contemporaries.

The holy Scriptures instruct us to pray and to labor; and to pray not only mechanically with the lips, as the heathen do, but with all the heart. But Paul the Simple said daily three hundred prayers, counting them with pebbles, which he carried in his bosom (a sort of rosary); when he heard of a virgin who prayed seven hundred times a day, he was troubled, and told his distress to Macarius, who well answered him: “Either thou prayest not with thy heart, if thy conscience reproves thee, or thou couldst pray oftener. I have for six years prayed only a hundred times a day, without being obliged to condemn myself for neglect.” Christ ate and drank like other men, expressly distinguishing himself thereby from John, the representative of the old covenant; and Paul recommends to us to use the gifts of God temperately, with cheerful and childlike gratitude.\footnote{Comp. Matt. xi. 18, 19; 1 Tim. iv. 3-5.} But the renowned anchoret and presbyter Isidore of Alexandria (whom Athanasius ordained) touched no meat, never ate enough, and, as Palladius relates, often burst into tears at table for shame, that he, who was destined to eat angels’ food in paradise, should have to eat material stuff like the irrational brutes. Macarius the elder, or the Great, for a long time ate only once a week, and slept standing and leaning on a staff. The equally celebrated younger Macarius lived three years on four or five ounces of bread a day, and seven years on raw herbs and pulse. Ptolemy spent three years alone in an unwatered desert, and quenched his thirst with the dew, which he collected in December and January, and preserved in earthen vessels; but he fell at last into skepticism, madness, and debauchery.\footnote{Comp. Hist. Laus. c. 33 and 95.} Sozomen tells of a certain Batthaeus, that by reason of his extreme abstinence, worms crawled out of his teeth; of Alas, that to his eightieth year he never ate bread; of Heliodorus, that he spent many nights without sleep, and fasted without interruption seven days.\footnote{Hist. Eccles. lib. vi. cap. 34.} Symeon, a Christian Diogenes, spent six and thirty years praying, fasting, and preaching, on the top of a pillar thirty or forty feet high, ate only once a week, and in fast times not at all. Such heroism of abstinence was possible, however, only in the torrid climate of the East, and is not to be met with in the West.

Anchoretism almost always carries a certain cynic roughness and coarseness, which, indeed, in the light of that age, may be leniently judged, but certainly have no affinity with the morality of the Bible, and offend not only good taste, but all sound moral feeling. The ascetic holiness, at least according to the Egyptian idea, is incompatible with cleanliness and decency, and delights in filth. It reverses the maxim of sound evangelical morality and modern Christian civilization, that cleanliness is next to godliness. Saints Anthony and Hilarion, as their admirers, Athanasius the Great and Jerome the Learned, tell us, scorned to comb or cut their hair (save once a year, at Easter), or to wash their hands or feet. Other hermits went almost naked in the wilderness, like the Indian gymnosophists.\footnote{These latter themselves were not absolutely naked, but wore a covering over the middle, as Augustine, in the passage above cited, De Civit. Dei, l. xiv. c. 17, and later tourists tell us. On the contrary, there were monks who were very scrupulous on this point. It is said of Ammon, that he never saw himself naked. The monks in Tabennae, according to the rule of Pachomius, had to sleep always in their clothes.} The younger Macarius, according to the account of his disciple Palladius, once lay six months naked in the morass of the Scetic desert, and thus exposed himself to the incessant attacks of the gnats of Africa, “whose sting can pierce even the hide of a wild boar.” He wished to punish himself for his arbitrary revenge on a gnat, and was there so badly stung by gnats and wasps, that he was thought to be smitten with leprosy, and was recognized only by his voice.\footnote{Comp. Hist Lausiaca, c. 20, and Tillemont, tom. viii. p. 633.} St. Symeon the Stylite, according to Theodoret, suffered himself to be incessantly tormented for a long time by twenty enormous bugs, and concealed an abscess full of worms, to exercise himself in patience and meekness. In Mesopotamia there was a peculiar class of anchorets, who lived on grass, spending the greater part of the day in prayer and singing, and then turning out like beasts upon the mountain.\footnote{The \textgreek{βοσκοί} or pabulatores. Comp. Sozom. H. E. l. vi. 33. Ephraim Syrus delivered a special eulogy on them, cited in Tillemont, Mem. tom. viii. p. 292 sq.} Theodoret relates of the much lauded Akepsismas, in Cyprus, that he spent sixty years in the same cell, without seeing or speaking to any one, and looked so wild and shaggy, that he was once actually taken for a wolf by a shepherd, who assailed him with stones, till he discovered his error, and then worshipped the hermit as a saint.\footnote{Hist. Rel. cap. (vita) xv. (Opera omnia, ed Par. iii. 843 sqq.).} It was but a step from this kind of moral sublimity to beastly degradation. Many of these saints were no more than low sluggards or gloomy misanthropes, who would rather company with wild beasts, with lions, wolves, and hyenas, than with immortal men, and above all shunned the face of a woman more carefully than they did the devil. Sulpitius Severus saw an anchoret in the Thebaid, who daily shared his evening meal with a female wolf; and upon her discontinuing her visits for some days by way of penance for a theft she had committed, he besought her to come again, and comforted her with a double portion of bread.\footnote{Dial. i. c. 8. Severus sees in this a wonderful example of the power of Christ over wild beasts.} The same writer tells of a hermit who lived fifty years secluded from all human society, in the clefts of Mount Sinai, entirely destitute of clothing, and all overgrown with thick hair, avoiding every visitor, because, as he said, intercourse with men interrupted the visits of the angels; whence arose the report that he held intercourse with angels.\footnote{L. c. i. c 11.}

It is no recommendation to these ascetic eccentricities that while they are without Scripture authority, they are fully equalled and even surpassed by the strange modes of self-torture practised by ancient and modern Hindoo devotees, for the supposed benefit of their souls and the gratification of their vanity in the presence of admiring spectators. Some bury themselves—we are told by ancient and modern travellers—in pits with only small breathing holes at the top, while others disdaining to touch the vile earth, live in iron cages suspended from trees. Some wear heavy iron collars or fetters, or drag a heavy chain fastened by one end round their privy parts, to give ostentatious proof of their chastity. Others keep their fists hard shut, until their finger nails grow through the palms of their hands. Some stand perpetually on one leg; others keep their faces turned over one shoulder, until they cannot turn them back again. Some lie on wooden beds, bristling all over with iron spikes; others are fastened for life to the trunk of a tree by a chain. Some suspend themselves for half an hour at a time, feet uppermost, or with a hook thrust through their naked back, over a hot fire. Alexander von Humboldt, at Astracan, where some Hindoos had settled, found a Yogi in the vestibule of the temple naked, shrivelled up, and overgrown with hair like a wild beast, who in this position had withstood for twenty years the severe winters of that climate. A Jesuit missionary describes one of the class called Tapasonias, that he had his body enclosed in an iron cage, with his head and feet outside, so that he could walk, but neither sit nor lie down; at night his pious attendants attached a hundred lighted lamps to the outside of the cage, so that their master could exhibit himself walking as the mock light of the world.\footnote{See Ruffner, l.c. i. 49 sqq., and Wuttke, l.c. p. 369 sqq.}

In general, the hermit life confounds the fleeing from the outward world with the mortification of the inward world of the corrupt heart. It mistakes the duty of love; not rarely, under its mask of humility and the utmost self-denial, cherishes spiritual pride and jealousy; and exposes itself to all the dangers of solitude, even to savage barbarism, beastly grossness, or despair and suicide. Anthony, the father of anchorets, well understood this, and warned his followers against overvaluing solitude, reminding them of the proverb of the Preacher, iv. 10: “Woe to him that is alone when he falleth; for he hath not another to help him up.”

The cloister life was less exposed to these errors. It approached the life of society and civilization. Yet, on the other hand, it produced no such heroic phenomena, and had dangers peculiar to itself. Chrysostom gives us the bright side of it from his own experience. “Before the rising of the sun,” says he of the monks of Antioch, “they rise, hale and sober, sing as with one mouth hymns to the praise of God, then bow the knee in prayer, under the direction of the abbot, read the holy Scriptures, and go to their labors; pray again at nine, twelve, and three o’clock; after a good day’s work, enjoy a simple meal of bread and salt, perhaps with oil, and sometimes with pulse; sing a thanksgiving hymn, and lay themselves on their pallets of straw without care, grief, or murmur. When one dies, they say: ’He is perfected;’ and all pray God for a like end, that they also may come to the eternal sabbath-rest and to the vision of Christ.” Men like Chrysostom, Basil, Gregory, Jerome, Nilus, and Isidore, united theological studies with the ascetic exercises of solitude, and thus gained a copious knowledge of Scripture and a large spiritual experience.

But most of the monks either could not even read, or had too little intellectual culture to devote themselves with advantage to contemplation and study, and only brooded over gloomy feelings, or sank, in spite of the unsensual tendency of the ascetic principle, into the coarsest anthropomorphism and image worship. When the religious enthusiasm faltered or ceased, the cloister life, like the hermit life, became the most spiritless and tedious routine, or hypocritically practised secret vices. For the monks carried with them into their solitude their most dangerous enemy in their hearts, and there often endured much fiercer conflicts with flesh and blood, than amidst the society of men.

The temptations of sensuality, pride, and ambition externalized and personified themselves to the anchorets and monks in hellish shapes, which appeared in visions and dreams, now in pleasing and seductive, now in threatening and terrible forms and colors, according to the state of mind at the time. The monastic imagination peopled the deserts and solitudes with the very worst society, with swarms of winged demons and all kinds of hellish monsters.\footnote{According to a sensuous and local conception of Eph. vi. 12\textgreek{· Τὰ πνευματικὰ τῆς πονηρίας ἐν τοῖς ἐπουρανίοις} ; “die bösen Geister unter dem Himmel” (evil spirits under heaven), as Luther translates; while the Vulgate gives it literally, but somewhat obscurely: “Spiritualia nequitiae in coelestibus;” and the English Bible quite too freely: “Spiritual wickedness in high places.” In any case \textgreek{πνευματικά} is to be taken in a much wider sense than \textgreek{πνεύματα} or \textgreek{δαιμόνια}; and \textgreek{ἐπουράνια}, also, is not fully identical with the cloud heaven or the atmosphere, and besides admits a different construction, so that many put a comma after \textgreek{πονηρίας}. The monastic satanology and demonology, we may remark, was universally received in the ancient church and throughout the middle ages. And it is well known that Luther retained from his monastic life a sensuous, materialistic idea of the devil and of his influence on men.} It substituted thus a new kind of polytheism for the heathen gods, which were generally supposed to be evil spirits. The monastic demonology and demonomachy is a strange mixture of gross superstitions and deep spiritual experiences. It forms the romantic shady side of the otherwise so tedious monotony of the secluded life, and contains much material for the history of ethics, psychology, and pathology.

Especially besetting were the temptations of sensuality, and irresistible without the utmost exertion and constant watchfulness. The same saints, who could not conceive of true chastity without celibacy, were disturbed, according to their own confession, by unchaste dreams, which at least defiled the imagination.\footnote{Athanasius says of St. Anthony, that the devil sometimes appeared to him in the form of a woman; Jerome relates of St. Hilarion, that in bed his imagination was often beset with visions of naked women. Jerome himself acknowledges, in a letter to a virgin (!), Epist. xxii. (ed. Vallars. t. i. p. 91, 92), \emph{de Custodia Virginitatis,} ad Eustochium: “O quoties ego ipse in eremo constitutus et in illa vasta solitudine, quae exusta solis ardoribus horridum monachis praebebat habitaculum, putavi me Romanis interesse deliciis .... Ille igitur ego, qui ob gehennae metum tali me carcere ipse damnaveram, scorpionum tantum socius et ferarum, saepe choris intereram puellarum. Pallebant ora jejuniis, et mens desideriis aestuabat in frigido corpore, et ante hominem suum jam in carne praemortuum, sola libidinum incendia bulliebant. Itaque omni auxilio destitutus, ad Jesu jacebam pedes, rigabam lacrymis, crine tergebam et repugnantem carnem hebdomadarum inedia subjugabam.” St. Ephraim warns against listening to the enemy, who whispers to the monk: \textgreek{Οὐ δυνατὸν παύσασθει ἀπό σου, ἐὰν μὴ πληροφορήσῃς ἐπιθυμίαν σου}.} Excessive asceticism sometimes turned into unnatural vice; sometimes ended in madness, despair, and suicide. Pachomius tells us, so early as his day, that many monks cast themselves down precipices, others ripped themselves up, and others put themselves to death in other ways.\footnote{Vita Pach. § 61. Comp. Nilus, Epist. l. ii. p. 140: \textgreek{Τινὲς}... \textgreek{ἑαυτοὺς ἔσφαξαν μαχαίρᾳ} etc. Even among the fanatical Circumcelliones, Donatist medicant monks in Africa, suicide was not uncommon.}

A characteristic trait of monasticism in all its forms is a morbid aversion to female society and a rude contempt of married life. No wonder, then, that in Egypt and the whole East, the land of monasticism, women and domestic life never attained their proper dignity, and to this day remain at a very low stage of culture. Among the rules of Basil is a prohibition of speaking with a woman, touching one, or even looking on one, except in unavoidable cases. Monasticism not seldom sundered the sacred bond between husband and wife, commonly with mutual consent, as in the cases of Ammon and Nilus, but often even without it. Indeed, a law of Justinian seems to give either party an unconditional right of desertion, while yet the word of God declares the marriage bond indissoluble. The Council of Gangra found it necessary to oppose the notion that marriage is inconsistent with salvation, and to exhort wives to remain with their husbands. In the same way monasticism came into conflict with love of kindred, and with the relation of parents to children; misinterpreting the Lord’s command to leave all for His sake. Nilus demanded of the monks the entire suppression of the sense of blood relationship. St. Anthony forsook his younger sister, and saw her only once after the separation. His disciple, Prior, when he became a monk, vowed never to see his kindred again, and would not even speak with his sister without closing his eyes. Something of the same sort is recorded of Pachomius. Ambrose and Jerome, in full earnest, enjoined upon virgins the cloister life, even against the will of their parents. When Hilary of Poictiers heard that his daughter wished to marry, he is said to have prayed God to take her to himself by death. One Mucius, without any provocation, caused his own son to be cruelly abused, and at last, at the command of the abbot himself, cast him into the water, whence he was rescued by a brother of the cloister.\footnote{Tillem. vii. 430. The abbot thereupon, as Tillemont relates, was informed by a revelation, ”que Muce avait egalé par son obeissance celle d’Abraham,” and soon after made him his successor.}

Even in the most favorable case monasticism falls short of harmonious moral development, and of that symmetry of virtue which meets us in perfection in Christ, and next to him in the apostles. It lacks the finer and gentler traits of character, which are ordinarily brought out only in the school of daily family life and under the social ordinances of God. Its morality is rather negative than positive. There is more virtue in the temperate and thankful enjoyment of the gifts of God, than in total abstinence; in charitable and well-seasoned speech, than in total silence; in connubial chastity, than in celibacy; in self-denying practical labor for the church. than in solitary asceticism, which only pleases self and profits no one else.

Catholicism, whether Greek or Roman, cannot dispense with the monastic life. It knows only moral extremes, nothing of the healthful mean. In addition to this, popery needs the monastic orders, as an absolute monarchy needs large standing armies both for conquest and defence. But evangelical Protestantism, rejecting all distinction of a twofold morality, assigning to all men the same great duty under the law of God, placing the essence of religion not in outward exercises, but in the heart, not in separation from the world and from society, but in purifying and sanctifying the world by the free spirit of the gospel, is death to the great monastic institution.

\xsection{Position of Monks in the Church}

§ 33. Position of Monks in the Church.

As to the social position of monasticism in the system of ecclesiastical life: it was at first, in East and West, even so late as the council of Chalcedon, regarded as a lay institution; but the monks were distinguished as religiosi from the seculares, and formed thus a middle grade between the ordinary laity and the clergy. They constituted the spiritual nobility, but not the ruling class; the aristocracy, but not the hierarchy of the church. “A monk,” says Jerome, “has not the office of a teacher, but of a penitent, who endures suffering either for himself or for the world.” Many monks considered ecclesiastical office incompatible with their effort after perfection. It was a proverb, traced to Pachomius: “A monk should especially shun women and bishops, for neither will let him have peace.”\footnote{Omnino monachum fugere debere mulieres et episcopos.} Ammonius, who accompanied Athanasius to Rome, cut off his own ear, and threatened to cut out his own tongue, when it was proposed to make him a bishop.\footnote{Sozom. iv. 30.} Martin of Tours thought his miraculous power deserted him on his transition from the cloister to the bishopric. Others, on the contrary, were ambitious for the episcopal chair, or were promoted to it against their will, as early as the fourth century. The abbots of monasteries were usually ordained priests, and administered the sacraments among the brethren, but were subject to the bishop of the diocese. Subsequently the cloisters managed, through special papal grants, to make themselves independent of the episcopal jurisdiction. From the tenth century the clerical character was attached to the monks. In a certain sense, they stood, from the beginning, even above the clergy; considered themselves preëminently conversi and religiosi, and their life vita religiosa; looked down with contempt upon the secular clergy; and often encroached on their province in troublesome ways. On the other hand, the cloisters began, as early as the fourth century, to be most fruitful seminaries of clergy, and furnished, especially in the East, by far the greater number of bishops. The sixth novel of Justinian provides that the bishops shall be chosen from the clergy, or from the monastery.

In dress, the monks at first adhered to the costume of the country, but chose the simplest and coarsest material. Subsequently, they adopted the tonsure and a distinctive uniform.

\xsection{Influence and Effect of Monasticism}

§ 34. Influence and Effect of Monasticism.

The influence of monasticism upon the world, from Anthony and Benedict to Luther and Loyola, is deeply marked in all branches of the history of the church. Here, too, we must distinguish light and shade. The operation of the monastic institution has been to some extent of diametrically opposite kinds, and has accordingly elicited the most diverse judgments. “It is impossible,” says Dean Milman,\footnote{Hist. of (ancient) Christianity, Am. ed., p. 432.} “to survey monachism in its general influence, from the earliest period of its inworking into Christianity, without being astonished and perplexed with its diametrically opposite effects. Here it is the undoubted parent of the blindest ignorance and the most ferocious bigotry, sometimes of the most debasing licentiousness; there the guardian of learning, the author of civilization, the propagator of humble and peaceful religion.” The apparent contradiction is easily solved. It is not monasticism, as such, which has proved a blessing to the church and the world; for the monasticism of India, which for three thousand years has pushed the practice of mortification to all the excesses of delirium, never saved a single soul, nor produced a single benefit to the race. It was Christianity in monasticism which has done all the good, and used this abnormal mode of life as a means for carrying forward its mission of love and peace. In proportion as monasticism was animated and controlled by the spirit of Christianity, it proved a blessing; while separated from it, it degenerated and became at fruitful source of evil.

At the time of its origin, when we can view it from the most favorable point, the monastic life formed a healthful and necessary counterpart to the essentially corrupt and doomed social life of the Graeco-Roman empire, and the preparatory school of a new Christian civilization among the Romanic and Germanic nations of the middle age. Like the hierarchy and the papacy, it belongs with the disciplinary institutions, which the spirit of Christianity uses as means to a higher end, and, after attaining that end, casts aside. For it ever remains the great problem of Christianity to pervade like leaven and sanctify all human society in the family and the state, in science and art, and in all public life. The old Roman world, which was based on heathenism, was, if the moral portraitures of Salvianus and other writers of the fourth and fifth centuries are even half true, past all such transformation; and the Christian morality therefore assumed at the outset an attitude of downright hostility toward it, till she should grow strong enough to venture upon her regenerating mission among the new and, though barbarous, yet plastic and germinal nations of the middle age, and plant in them the seed of a higher civilization.

Monasticism promoted the downfall of heathenism and the victory of Christianity in the Roman empire and among the barbarians. It stood as a warning against the worldliness, frivolity, and immorality of the great cities, and a mighty call to repentance and conversion. It offered a quiet refuge to souls weary of the world, and led its earnest disciples into the sanctuary of undisturbed communion with God. It was to invalids a hospital for the cure of moral diseases, and at the same time, to healthy and vigorous enthusiasts an arena for the exercise of heroic virtue.\footnote{Chateaubriand commends the monastic institution mainly under the first view. “If there are refuges for the health of the body, ah ! permit religion to have such also for the health of the soul, which is still more subject to sickness, and the infirmities of which are so much more sad, so much more tedious and difficult to cure!” Montalembert (l.c. i. 25) objects to this view as poetic and touching but false, and represents monasticism as an arena for the healthiest and strongest souls which the world has ever produced, and quotes the passage of Chrysostom: “Come and see the tents of the soldiers of Christ; come and see their order of battle; they fight every day, and every day they defeat and immolate the passions which assail us.”} It recalled the original unity and equality of the human race, by placing rich and poor, high and low upon the same level. It conduced to the abolition, or at least the mitigation of slavery.\footnote{1 The abbot Isidore of Pelusium wrote to a slaveholder, Ep. l. i. 142 (cited by Neander): “I did not think that the man who loves Christ, and knows the grace which makes us all free, would still hold slaves.”} It showed hospitality to the wayfaring, and liberality to the poor and needy. It was an excellent school of meditation, self-discipline, and spiritual exercise. It sent forth most of those catholic, missionaries, who, inured to all hardship, planted the standard of the cross among the barbarian tribes of Northern and Western Europe, and afterward in Eastern Asia and South America. It was a prolific seminary of the clergy, and gave the church many of her most eminent bishops and popes, as Gregory I. and Gregory VII. It produced saints like Anthony and Bernard, and trained divines like Chrysostom and Jerome, and the long succession of schoolmen and mystics of the middle ages. Some of the profoundest theological discussions, like the tracts of Anselm, and the Summa of Thomas Aquinas, and not a few of the best books of devotion, like the “Imitation of Christ,” by Thomas a Kempis, have proceeded from the solemn quietude of cloister life. Sacred hymns, unsurpassed for sweetness, like the Jesu dulcis memoria, or tender emotion, like the Stabat mater dolorosa, or terrific grandeur, like the Dies irae, dies illa, were conceived and sung by mediaeval monks for all ages to come. In patristic and antiquarian learning the Benedictines, so lately as the seventeenth century, have done extraordinary service. Finally, monasticism, at least in the West, promoted the cultivation of the soil and the education of the people, and by its industrious transcriptions of the Bible, the works of the church fathers, and the ancient classics, earned for itself, before the Reformation, much of the credit of the modern civilization of Europe. The traveller in France, Italy, Spain, Germany, England, and even in the northern regions of Scotland and Sweden, encounters innumerable traces of useful monastic labors in the ruins of abbeys, of chapter houses, of convents, of priories and hermitages, from which once proceeded educational and missionary influences upon the surrounding hills and forests. These offices, however, to the progress of arts and letters were only accessory, often involuntary, and altogether foreign to the intention of the founders of monastic life and institutions, who looked exclusively to the religious and moral education of the soul. In seeking first the kingdom of heaven, these other things were added to them.

But on the other hand, monasticism withdrew from society many useful forces; diffused an indifference for the family life, the civil and military service of the state, and all public practical operations; turned the channels of religion from the world into the desert, and so hastened the decline of Egypt, Syria, Palestine, and the whole Roman empire. It nourished religious fanaticism, often raised storms of popular agitation, and rushed passionately into the controversies of theological parties; generally, it is true, on the side of orthodoxy, but often, as at the Ephesian “council of robbers,” in favor of heresy, and especially in behalf of the crudest superstition. For the simple, divine way of salvation in the gospel, it substituted an arbitrary, eccentric, ostentatious, and pretentious sanctity. It darkened the all-sufficient merits of Christ by the glitter of the over-meritorious works of man. It measured virtue by the quantity of outward exercises instead of the quality of the inward disposition, and disseminated self-righteousness and an anxious, legal, and mechanical religion. It favored the idolatrous veneration of Mary and of saints, the worship of images and relics, and all sorts of superstitious and pious fraud. It circulated a mass of visions and miracles, which, if true, far surpassed the miracles of Christ and the apostles and set all the laws of nature and reason at defiance. The Nicene age is full of the most absurd monks’ fables, and is in this respect not a whit behind the darkest of the middle ages.\footnote{The monkish miracles, with which the \emph{Vitae Patrum} of the Jesuit Rosweyde and the \emph{Acta Sanctorum} swarm, often contradict all the laws of nature and of reason, and would be hardly worthy of mention, but that they come from such fathers as Jerome, Rufinus, Severus, Palladius, and Theodoret, and go to characterize the Nicene age. We are far from rejecting all and every one as falsehood and deception, and accepting the judgment of Isaac Taylor (Ancient Christianity, ii. 106): “The Nicene miracles are of a kind which shocks every sentiment of gravity, of decency, and of piety:—in their obvious features they are childish, horrid, blasphemous, and foul.” Much more cautious is the opinion of Robertson (Hist. of the Christian Church, i. 312) and other Protestant historians, who suppose that, together with the innocent illusions of a heated imagination and the fabrications of intentional fraud, there must have been also much that was real, though in the nature of the case an exact sifting is impossible. But many of these stories are too much even for Roman credulity, and are either entirely omitted or at least greatly reduced and modified by critical historians. We read not only of innumerable visions, prophecies, healings of the sick and the possessed, but also of raising of the dead (as in the life of Martin of Tours), of the growth of a dry stick into a fruitful tree, and of a monk’s passing unseared, in absolute obedience to his abbot, through a furnace of fire as through a cooling bath. (Comp. Sulp. Sever. Dial. i. c. 12 and 13.) Even wild beasts play a large part, and are transformed into rational servants of the Egyptian saints of the desert. At the funeral of Paul of Thebes, according to Jerome, two lions voluntarily performed the office of sexton. Pachomius walked unharmed over serpents and scorpions, and crossed the Nile on crocodiles, which, of their own accord, presented their backs. The younger Macarius, or (according to other statements of the Historia Lausiaca; comp. the investigation of Tillemont, tom. viii. p. 811 sqq.) the monk Marcus stood on so good terms with the beasts, that a hyena (according to Rufinus, V. P. ii. 4, it was a lioness) brought her young one to him in his cell, that he might open its eyes; which he did by prayer and application of spittle; and the next day she offered him, for gratitude, a large sheepskin; the saint at first declined the gift, and reproved the beast for the double crime of murder and theft, by which she had obtained the skin; but when the hyena showed repentance, and with a nod promised amendment, Macarius took the skin, and afterward bequeathed it to the great bishop Athanasius. Severus (Dial. i. c. 9) gives a very similar account of an unknown anchoret, but, like Rufinus, substitutes for the hyena of Palladius a lioness with five whelps, and makes the saint receive the present of the skin without scruple or reproof. Shortly before (c. 8), he speaks, however, of a wolf, which once robbed a friendly hermit, whose evening meal she was accustomed to share, showed deep repentance for it, and with bowed head begged forgiveness of the saint. Perhaps Palladius or his Latin translator has combined these two anecdotes.} Monasticism lowered the standard of general morality in proportion as it set itself above it and claimed a corresponding higher merit; and it exerted in general a demoralizing influence on the people, who came to consider themselves the profanum vulgus mundi, and to live accordingly. Hence the frequent lamentations, not only of Salvian, but of Chrysostom and of Augustine, over the indifference and laxness of the Christianity of the day; hence to this day the mournful state of things in the southern countries of Europe and America, where monasticism is most prevalent, and sets the extreme of ascetic sanctity in contrast with the profane laity, but where there exists no healthful middle class of morality, no blooming family life, no moral vigor in the masses. In the sixteenth century the monks were the bitterest enemies of the Reformation and of all true progress. And yet the greatest of the reformers was a pupil of the convent, and a child of the monastic system, as the boldest and most free of the apostles had been the strictest of the Pharisees.

\xsection{Paul of Thebes and St. Anthony}

§ 35. Paul of Thebes and St. Anthony.

I. Athanasius: Vita S. Antonii (in Greek, Opera, ed. Ben. ii. 793–866). The same in Latin, by Evagrius, in the fourth century. Jerome: Catal. c. 88 (a very brief notice of Anthony); Vita S. Pauli Theb. (Opera, ed. Vallars, ii. p. 1–12). Sozom: H. E. l. i. cap. 13 and 14. Socrat.: H. E. iv. 23, 25.

II. Acta Sanctorum, sub Jan. 17 (tom. ii. p. 107 sqq.). Tillemont: Mem. tom. vii. p. 101–144 (St. Antoine, premier père des solitaires d’Egypte). Butler (R.C.): Lives of the Saints, sub Jan. 17. Möhler (R.C.): Athanasius der Grosse, p. 382–402. Neander: K. G. iii. 446 sqq. (Torrey’s Engl. ed. ii. 229–234). Böhringer: Die Kirche Christi in Biographien, i. 2, p. 122–151. H. Ruffner: l.c. vol. i. p. 247–302 (a condensed translation from Athanasius, with additions). K. Hase: K. Gesch. § 64 (a masterly miniature portrait).

The first known Christian hermit, as distinct from the earlier ascetics, is the fabulous Paul of Thebes, in Upper Egypt. In the twenty-second year of his age, during the Decian persecution, a.d. 250, he retired to a distant cave, grew fond of the solitude, and lived there, according to the legend, ninety years, in a grotto near a spring and a palm tree, which furnished him food, shade, and clothing,\footnote{Pliny counts thirty-nine different sorts of palm trees, of which the best grow in Egypt, are ever green, have thick foliage, and bear a fruit, from which in some places bread is made.} until his death in 340. In his later years a raven is said to have brought him daily half a loaf, as the ravens ministered to Elijah. But no one knew of this wonderful saint, till Anthony, who under a higher impulse visited and buried him, made him known to the world. After knocking in vain for more than an hour at the door of the hermit, who would receive the visits of beasts and reject those of men, he was admitted at last with a smiling face, and greeted with a holy kiss. Paul had sufficient curiosity left to ask the question, whether there were any more idolaters in the world, whether new houses were built in ancient cities and by whom the world was governed? During this interesting conversation, a large raven came gently flying and deposited a double portion of bread for the saint and his guest. “The Lord,” said Paul, “ever kind and merciful, has sent us a dinner. It is now sixty years since I have daily received half a loaf, but since thou hast come, Christ has doubled the supply for his soldiers.” After thanking the Giver, they sat down by the fountain; but now the question arose who should break the bread; the one urging the custom of hospitality, the other pleading the right of his friend as the elder. This question of monkish etiquette, which may have a moral significance, consumed nearly the whole day, and was settled at last by the compromise that both should seize the loaf at opposite ends, pull till it broke, and keep what remained in their hands. A drink from the fountain, and thanksgiving to God closed the meal. The day afterward Anthony returned to his cell, and told his two disciples: “Woe to me, a sinner, who have falsely pretended to be a monk. I have seen Elijah and John in the desert; I have seen St. Paul in paradise.” Soon afterward he paid St. Paul a second visit, but found him dead in his cave, with head erect and hands lifted up to heaven. He wrapped up the corpse, singing psalms and hymns, and buried him without a spade; for two lions came of their own accord, or rather from supernatural impulse, from the interior parts of the desert, laid down at his feet, wagging their tails, and moaning distressingly, and scratched a grave in the sand large enough for the body of the departed saint of the desert! Anthony returned with the coat of Paul, made of palm leaves, and wore it on the solemn days of Easter and Pentecost.

The learned Jerome wrote the life of Paul, some thirty years afterward, as it appears, on the authority of Anathas and Macarius, two disciples of Anthony. But he remarks, in the prologue, that many incredible things are said of him, which are not worthy of repetition. If he believed his story of the grave-digging lions, it is hard to imagine what was more credible and less worthy of repetition.

In this Paul we have an example, of a canonized saint, who lived ninety years unseen and unknown in the wilderness, beyond all fellowship with the visible church, without Bible, public worship, or sacraments, and so died, yet is supposed to have attained the highest grade of piety. How does this consist with the common doctrine of the Catholic church respecting the necessity and the operation of the means of grace? Augustine, blinded by the ascetic spirit of his age, says even, that anchorets, on their level of perfection, may dispense with the Bible. Certain it is, that this kind of perfection stands not in the Bible, but outside of it.

The proper founder of the hermit life, the one chiefly instrumental in giving it its prevalence, was St. Anthony of Egypt. He is the most celebrated, the most original, and the most venerable representative of this abnormal and eccentric sanctity, the “patriarch of the monks,” and the “childless father of an innumerable seed.”\footnote{Jeromesays of Anthony, in his Vita Pauli Theb. (c. i.): “Non tam ipse auto omnes (eremitas) fuit, quam ab eo omnium incitata sunt studia.”}

Anthony sprang from a Christian and honorable Coptic family, and was born about 251, at Coma, on the borders of the Thebaid. Naturally quiet, contemplative, and reflective, he avoided the society of playmates, and despised all higher learning. He understood only his Coptic vernacular, and remained all his life ignorant of Grecian literature and secular science.\footnote{According to the common opinion, which was also Augustine’s, Anthonycould not even read. But Tillemont (tom. vii. 107 and 666), Butler, and others think that this igorance related only to the Greek alphabet, not to the Egyptian. Athanasius, p. 795, expresses himself somewhat indistinctly; that, from dread of society, he would not \textgreek{μαθεῖν γράμματα} (letters? or the arts?), but speaks afterward of his regard for reading.} But he diligently attended divine worship with his parents, and so carefully heard the Scripture lessons, that he retained them in memory.\footnote{Augustinesays of him, De doctr. Christ. § 4, that, without being able to read from only hearing the Bible, he knew it by heart. The life of Athanasius shows, indeed, that a number of Scripture passages were very familiar to him. But of a connected and deep knowledge of Scripture in him, or in these anchorets generally, we find no trace.} Memory was his library. He afterward made faithful, but only too literal use of single passages of Scripture, and began his discourse to the hermits with the very uncatholic-sounding declaration: “The holy Scriptures give us instruction enough.” In his eighteenth year, about 270, the death of his parents devolved on him the care of a younger sister and a considerable estate. Six months afterward he heard in the church, just as he was meditating on the apostles’ implicit following of Jesus, the word of the Lord to the rich young ruler: “If thou wilt be perfect, go and sell that thou hast and give to the poor, and thou shalt have treasure in heaven; and come and follow me.”\footnote{Matt. xix. 21.} This word was a voice of God, which determined his life. He divided his real estate, consisting of three hundred acres of fertile land, among the inhabitants of the village, and sold his personal property for the benefit of the poor, excepting a moderate reserve for the support of his sister. But when, soon afterward, he heard in the church the exhortation, “Take no thought for the morrow,”\footnote{Matt. vi. 34.} he distributed the remnant to the poor, and intrusted his sister to a society of pious virgins.\footnote{\textgreek{Εἰς παρθενῶνα}, says Athanasius; i.e., not “un monastere de verges,” as Tillemont translates, for nunneries did not yet exist; but a society of female ascetics within the congregation; from which, however, a regular cloister might of course very easily grow.} He visited her only once after—a fact characteristic of the ascetic depreciation of natural ties.

He then forsook the hamlet, and led an ascetic life in the neighborhood, praying constantly, according to the exhortation: “Pray without ceasing;” and also laboring, according to the maxim: “If any will not work, neither should he eat.” What he did not need for his slender support, he gave to the poor. He visited the neighboring ascetics, who were then already very plentiful in Egypt, to learn humbly and thankfully their several eminent virtues; from one, earnestness in prayer; from another, watchfulness; from a third, excellence in fasting; from a fourth, meekness; from all, love to Christ and to fellow men. Thus he made himself universally beloved, and came to be reverenced as a friend of God.

But to reach a still higher level of ascetic holiness, he retreated, after the year 285, further and further from the bosom and vicinity of the church, into solitude, and thus became the founder of an anchoretism strictly so called. At first he lived in a sepulchre; then for twenty years in the ruins of a castle; and last on Mount Colzim, some seven hours from the Red Sea, a three days’ journey east of the Nile, where an old cloister still preserves his name and memory.

In this solitude he prosecuted his ascetic practices with ever-increasing rigor. Their monotony was broken only by basket making, occasional visits, and battles with the devil. In fasting he attained a rare abstemiousness. His food consisted of bread and salt, sometimes dates; his drink, of water. Flesh and wine he never touched. He ate only once a day, generally after sunset, and, like the presbyter Isidore, was ashamed that an immortal spirit should need earthly nourishment. Often he fasted from two to five days. Friends, and wandering Saracens, who always had a certain reverence for the saints of the desert, brought him bread from time to time. But in the last years of his life, to render himself entirely independent of others, and to afford hospitality to travellers, he cultivated a small garden on the mountain, near a spring shaded by palms.\footnote{Jerome, in his Vita Hilarionis, c. 31, gives an incidental description of this last residence of Anthony, according to which it was not so desolate as from Athanasius one would infer. He speaks even of palms, fruit trees, and vines in this garden, the fruit of which any one would have enjoyed.} Sometimes the wild beasts of the forest destroyed his modest harvest, till he drove them away forever with the expostulation: “Why do you injure me, who have never done you the slightest harm? Away with you all, in the name of the Lord, and never come into my neighborhood again.” He slept on bare ground, or at best on a pallet of straw; but often he watched the whole night through in prayer. The anointing of the body with oil he despised, and in later years never washed his feet; as if filthiness were an essential element of ascetic perfection. His whole wardrobe consisted of a hair shirt, a sheepskin, and a girdle. But notwithstanding all, he had a winning friendliness and cheerfulness in his face.

Conflicts with the devil and his hosts of demons were, as with other solitary saints, a prominent part of Anthony’s experience, and continued through all his life. The devil appeared to him in visions and dreams, or even in daylight, in all possible forms, now as a friend, now as a fascinating woman, now as a dragon, tempting him by reminding him of his former wealth, of his noble family, of the care due to his sister, by promises of wealth, honor, and renown, by exhibitions of the difficulty of virtue and the facility of vice, by unchaste thoughts and images, by terrible threatening of the dangers and punishments of the ascetic life. Once he struck the hermit so violently, Athanasius says, that a friend, who brought him bread, found him on the ground apparently dead. At another time he broke through the wall of his cave and filled the room with roaring lions, howling wolves, growling bears, fierce hyenas, crawling serpents and scorpions; but Anthony turned manfully toward the monsters, till a supernatural light broke in from the roof and dispersed them. His sermon, which he delivered to the hermits at their request, treats principally of these wars with demons, and gives also the key to the interpretation of them: “Fear not Satan and his angels. Christ has broken their power. The best weapon against them is faith and piety .... The presence of evil spirits reveals itself in perplexity, despondency, hatred of the ascetics, evil desires, fear of death .... They take the form answering to the spiritual state they find in us at the time.\footnote{Athanas. c. 42: \textgreek{Ἐλθόντες γὰρ} (\textgreek{οἱ ἐχθροὶ}) \textgreek{ὁποίους ἀν εὕρωσιν ἡμᾶς}\emph{,} \textgreek{τοιοῦτοι καὶ αὐτοὶ γίνονται}, etc.—an important psychological observation.} They are the reflex of our thoughts and fantasies. If thou art carnally minded, thou art their prey; but if thou rejoicest in the Lord and occupiest thyself with divine things, they are powerless .... The devil is afraid of fasting, of prayer, of humility and good works. His illusions soon vanish, when one arms himself with the sign of the cross.”

Only in exceptional cases did Anthony leave his solitude; and then he made a powerful impression on both Christians and heathens with his hairy dress and his emaciated, ghostlike form. In the year 311, during the persecution under Maximinus, he appeared in Alexandria in the hope of himself gaining the martyr’s crown. He visited the confessors in the mines and prisons, encouraged them before the tribunal, accompanied them to the scaffold; but no one ventured to lay hands on the saint of the wilderness. In the year 351, when a hundred years old, he showed himself for the second and last time in the metropolis of Egypt, to bear witness for the orthodox faith of his friend Athanasius against Arianism, and in a few days converted more heathens and heretics than had otherwise been gained in a whole year. He declared the Arian denial of the divinity of Christ worse than the venom of the serpent, and no better than heathenism which worshipped the creature instead of the Creator. He would have nothing to do with heretics, and warned his disciples against intercourse with them. Athanasius attended him to the gate of the city, where he cast out an evil spirit from a girl. An invitation to stay longer in Alexandria he declined, saying: “As a fish out of water, so a monk out of his solitude dies.” Imitating his example, the monks afterward forsook the wilderness in swarms whenever orthodoxy was in danger, and went in long processions with wax tapers and responsive singing through the streets, or appeared at the councils, to contend for the orthodox faith with all the energy of fanaticism, often even with physical force.

Though Anthony shunned the society of men, yet he was frequently visited in his solitude and resorted to for consolation and aid by Christians and heathens, by ascetics, sick, and needy, as a heaven-descended physician of Egypt for body and soul. He enjoined prayer, labor, and care of the poor, exhorted those at strife to the love of God, and healed the sick and demoniac with his prayer. Athanasius relates several miracles performed by him, the truth of which we leave undecided though they are far less incredible and absurd than many other monkish stories of that age. Anthony, his biographer assures us, never boasted when his prayer was heard, nor murmured when it was not, but in either case thanked God. He cautioned monks against overrating the gift of miracles, since it is not our work, but the grace of the Lord; and he reminds them of the word: “Rejoice not, that the spirits are subject unto you; but rather rejoice, because your names are written in heaven.” To Martianus, an officer, who urgently besought him to heal his possessed daughter, he said: “Man, why dost thou call on me? I am a man, as thou art. If thou believest, pray to God, and he will hear thee.” Martianus prayed, and on his return found his daughter whole.

Anthony distinguished himself above most of his countless disciples and successors, by his fresh originality of mind. Though uneducated and limited, he had sound sense and ready mother wit. Many of his striking answers and felicitous sentences have come down to us. When some heathen philosophers once visited him, he asked them: “Why do you give yourselves so much trouble to see a fool?” They explained, perhaps ironically, that they took him rather for a wise man. He replied: “If you take me for a fool, your labor is lost; but if I am a wise man, you should imitate me, and be Christians, as I am.” At another time, when taunted with his ignorance, he asked: “Which is older and better, mind or learning?” The mind, was the answer. “Then,” said the hermit, “the mind can do without learning.” “My book,” he remarked on a similar occasion, “is the whole creation, which lies open before me, and in which I can read the word of God as often as I will.” The blind church-teacher, Didymus, whom he met in Alexandria, he comforted with the words: “Trouble not thyself for the loss of the outward eye, with which even flies see; but rejoice in the possession of the spiritual eye, with which also angels behold the face of God, and receive his light.”\footnote{This is not told indeed by Athanasius, but by Rufinus, Jerome, and Socrates (Hist. Eccl. iv. 25). Comp. Tillemont, l.c. p. 129.} Even the emperor Constantine, with his sons, wrote to him as a spiritual father, and begged an answer from him. The hermit at first would not so much as receive the letter, since, in any case, being unable to write, he could not answer it, and cared as little for the great of this world as Diogenes for Alexander. When told that the emperor was a Christian, he dictated the answer: “Happy thou, that thou worshippest Christ. Be not proud of thy earthly power. Think of the future judgment, and know that Christ is the only true and eternal king. Practise justice and love for men, and care for the poor.” To his disciples he said on this occasion: “Wonder not that the emperor writes to me, for he is a man. Wonder much more that God has written the law for man, and has spoken to us by his own Son.”

During the last years of his life the patriarch of monasticism withdrew as much as possible from the sight of visitors, but allowed two disciples to live with him, and to take care of him in his infirm old age. When he felt his end approaching, he commanded them not to embalm his body, according to the Egyptian custom, but to bury it in the earth, and to keep the spot of his interment secret. One of his two sheepskins he bequeathed to the bishop Serapion, the other, with his underclothing, to Athanasius, who had once given it to him new, and now received it back worn out. What became of the robe woven from palm leaves, which, according to Jerome, he had inherited from Paul of Thebes, and wore at Easter and Pentecost, Athanasius does not tell us. After this disposition of his property, Anthony said to his disciples: “Children, farewell; for Anthony goes away, and will be no more with you.” With these words he stretched out his feet and expired with a smiling face, in the year 356, a hundred and five years old. His grave remained for centuries unknown. His last will was thus a protest against the worship of saints and relics, which, however, it nevertheless greatly helped to promote. Under Justinian, in 561, his bones, as the Bollandists and Butler minutely relate, were miraculously discovered, brought to Alexandria, then to Constantinople, and at last to Vienne in South France, and in the eleventh century, during the raging of an epidemic disease, the so-called “holy fire,” or “St. Anthony’s fire,” they are said to have performed great wonders.

Athanasius, the greatest man of the Nicene age, concludes his biography of his friend with this sketch of his character: “From this short narrative you may judge how great a man Anthony was, who persevered in the ascetic life from youth to the highest age. In his advanced age he never allowed himself better food, nor change of raiment, nor did he even wash his feet. Yet he continued healthy in all his parts. His eyesight was clear to the end, and his teeth sound, though by long use worn to mere stumps. He retained also the perfect use of his hands and feet, and was more robust and vigorous than those who are accustomed to change of food and clothing and to washing. His fame spread from his remote dwelling on the lone mountain over the whole Roman empire. What gave him his renown, was not learning nor worldly wisdom, nor human art, but alone his piety toward God .... And let all the brethren know, that the Lord will not only take holy monks to heaven, but give them celebrity in all the earth, however deep they may bury themselves in the wilderness.”

The whole Nicene age venerated in Anthony a model saint.\footnote{Comp. the proofs in Tillemont, l.c. p. 137 sq.} This fact brings out most characteristically the vast difference between the ancient and the modern, the old Catholic and the evangelical Protestant conception of the nature of the Christian religion. The specifically Christian element in the life of Anthony, especially as measured by the Pauline standard, is very small. Nevertheless we can but admire the needy magnificence, the simple, rude grandeur of this hermit sanctity even in its aberration. Anthony concealed under his sheepskin a childlike humility, an amiable simplicity, a rare energy of will, and a glowing love to God, which maintained itself for almost ninety years in the absence of all the comforts and pleasures of natural life, and triumphed over all the temptations of the flesh. By piety alone, without the help of education or learning, he became one of the most remarkable and influential men in the history of the ancient church. Even heathen contemporaries could not withhold from him their reverence, and the celebrated philosopher Synesius, afterward a bishop, before his conversion reckoned Anthony among those rare men, in whom flashes of thought take the place of reasonings, and natural power of mind makes schooling needless.\footnote{Dion, fol. 51, ed. Petav., cited in Tillemont and Neander.}

\xsection{Spread of Anchoretism. Hilarion}

§ 36. Spread of Anchoretism. Hilarion.

The example of Anthony acted like magic upon his generation, and his biography by Athanasius, which was soon translated also into Latin, was a tract for the times. Chrysostom recommended it to all as instructive and edifying reading.\footnote{Hom. viii. in Matth. tom. vii. 128 (ed. Montfaucon).} Even Augustine, the most evangelical of the fathers, was powerfully affected by the reading of it in his decisive religious struggle, and was decided by it in his entire renunciation of the world.\footnote{· Comp. Aug.: Confess. l. viii. c. 6 and 28.}

In a short time, still in the lifetime of Anthony, the deserts of Egypt, from Nitria, south of Alexandria, and the wilderness of Scetis, to Libya and the Thebaid, were peopled with anchorets and studded with cells. A mania for monasticism possessed Christendom, and seized the people of all classes like an epidemic. As martyrdom had formerly been, so now monasticism was, the quickest and surest way to renown upon earth and to eternal reward in heaven. This prospect, with which Athanasius concludes his life of Anthony, abundantly recompensed all self-denial and mightily stimulated pious ambition. The consistent recluse must continually increase his seclusion. No desert was too scorching, no rock too forbidding, no cliff too steep, no cave too dismal for the feet of these world-hating and man-shunning enthusiasts. Nothing was more common than to see from two to five hundred monks under the same abbot. It has been supposed, that in Egypt the number of anchorets and cenobites equalled the population of the cities.\footnote{“Quanti populi,” says Rufinus (Vitae Patr. ii c. 7), “habentur in urbibus, tantae paene habentur in desertis multitudines monachorum.” Gibbon adds the sarcastic remark: “Posterity might repeat the saying, which had formerly been applied to sacred animals of the same country, That in Egypt it was less difficult to find a god than a man.” Montalembert (Monks of the West, vol. i. p. 314) says of the increase of monks: “Nothing in the wonderful history of these hermits in Egypt is so incredible as their number. But the most weighty authorities agreed in establishing it (S. Augustine, De morib. Eccles. i. 31). It was a kind of emigration of towns to the desert, of civilization to simplicity, of noise to silence, of corruption to innocence. The current once begun, floods of men, of women, and of children threw themselves into it, and flowed thither during a century with irresistible force.”} The natural contrast between the desert and the fertile valley of the Nile, was reflected in the moral contrast between the monastic life and the world.

The elder Macarius\footnote{There were several (five or seven) anchorets of this name, who are often confounded. The most celebrated are Macarius the elder, or the Great († 390), to whom the Homilies probably belong; and Macarius the younger, of Alexandria († 404), the teacher of Palladius, who spent a long time with him, and set him as high as the other. Comp. Tillemont’s extended account, tom. viii. p. 574-650, and the notes, p. 811 sqq.} introduced the hermit life in the frightful desert of Scetis; Amun or Ammon,\footnote{On Ammon, or, in Egyptian, Amus and Amun, comp. Tillemont, viii. p. 153-166, and the notes, p. 672-674.} on the Nitrian mountain. The latter was married, but persuaded his bride, immediately after the nuptials, to live with him in the strictest abstinence. Before the end of the fourth century there were in Nitria alone, according to Sozomen, five thousand monks, who lived mostly in separate cells or laurae, and never spoke with one another except on Saturday and Sunday, when they assembled for common worship.

From Egypt the solitary life spread to the neighboring countries.

Hilarion, whose life Jerome has written graphically and at large,\footnote{Opera, tom. ii. p. 13-40.} established it in the wilderness of Gaza, in Palestine and Syria. This saint attained among the anchorets of the fourth century an eminence second only to Anthony. He was the son of pagan parents, and grew up “as a rose among thorns.” He went to school in Alexandria, diligently attended church, and avoided the circus, the gladiatorial shows, and the theatre. He afterward lived two months with St. Anthony, and became his most celebrated disciple. After the death of his parents, he distributed his inheritance among his brothers and the poor, and reserved nothing, fearing the example of Ananias and Sapphira, and remembering the word of Christ: “Whosoever he be of you, that forsaketh not all that he hath, he cannot be my disciple.”\footnote{Lu. xiv. 33.} He then retired into the wilderness of Gaza, which was inhabited only by robbers and assassins; battled, like Anthony, with obscene dreams and other temptations of the devil; and so reduced his body—the “ass,” which ought to have not barley, but chaff—with fastings and night watchings, that, while yet a youth of twenty years, he looked almost like a skeleton. He never ate before sunset. Prayers, psalm singing, Bible recitations, and basket weaving were his employment. His cell was only five feet high, lower than his own stature, and more like a sepulchre than a dwelling. He slept on the ground. He cut his hair only once a year, at Easter. The fame of his sanctity gradually attracted hosts of admirers (once, ten thousand), so that he had to change his residence several times, and retired to Sicily, then to Dalmatia, and at last to the island of Cyprus, where he died in 371, in his eightieth year. His legacy, a book of the Gospels and a rude mantle, he made to his friend Hesychius, who took his corpse home to Palestine, and deposited it in the cloister of Majumas. The Cyprians consoled themselves over their loss, with the thought that they possessed the spirit of the saint. Jerome ascribes to him all manner of visions and miraculous cures.

\xsection{St. Symeon and the Pillar Saints}

§ 37. St. Symeon and the Pillar Saints.

Respecting St. Symeon, or Simeon Stylites, we have accounts from three contemporaries and eye witnesses, Anthony, Cosmas, and especially Theodoret (Hist. Relig. c. 26). The latter composed his narrative sixteen years before the death the saint.

Evagrius: H. E. i. c. 13. The Acta Sanctorum and Butler, sub Jan. 5. Uhlemann: Symeon, der erste Säulenheilige in Syrien. Leipz. 1846. (Comp. also the fine poem of A. Tennyson: St. Symeon Stylites, a monologue in which S. relates his own experience.)

It is unnecessary to recount the lives of other such anchorets; since the same features, even to unimportant details, repeat themselves in all.\footnote{A peculiar, romantic, but not fully historical interest attaches to the biography of the imprisoned and fortunately escaping monk Malchus, with his nominal wife, which is preserved to us by Jerome.} But in the fifth century a new and quite original path\footnote{Original at least in the Christian church. Gieseler refers to a heathen precedent; the \textgreek{Φαλλοβατεῖς}in Syria, mentioned by Lucian, De Dea Syria, c. 28 and 29.} was broken by Symeon, the father of the Stylites or pillar saints, who spent long years, day and night, summer and winter, rain and sunshine, frost and heat, standing on high, unsheltered pillars, in prayer and penances, and made the way to heaven for themselves so passing hard, that one knows not whether to wonder at their unexampled self-denial, or to pity their ignorance of the gospel salvation. On this giddy height the anchoretic asceticism reached its completion.

St. Symeon the Stylite, originally a shepherd on the borders of Syria and Cilicia, when a boy of thirteen years, was powerfully affected by the beatitudes, which he heard read in the church, and betook himself to a cloister. He lay several days, without eating or drinking, before the threshold, and begged to be admitted as the meanest servant of the house. He accustomed himself to eat only once a week, on Sunday. During Lent he even went through the whole forty days without any food; a fact almost incredible even for a tropical climate.\footnote{Butler, l.c., however, relates something similar of a contemporary Benedictine monk, Dom Claude Leante: “In 1731, when he was about fifty-one years of age, he had fasted eleven years without taking any food the whole forty days, except what he daily took at mass; and what added to the wonder is, that during Lent he did not properly sleep, but only dozed. He could not bear the open air; and toward the end of Lent he was excessively pale and wasted. This fact is attested by his brethren and superiors, in a relation printed at Sens, in 1731.”} The first attempt of this kind brought him to the verge of death; but his constitution conformed itself, and when Theodoret visited him, he had solemnized six and twenty Lent seasons by total abstinence, and thus surpassed Moses, Elias, and even Christ, who never fasted so but once. Another of his extraordinary inflections was to lace his body so tightly that the cord pressed through to the bones, and could be cut off only with the most terrible pains. This occasioned his dismissal from the cloister. He afterward spent some time as a hermit upon a mountain, with an iron chain upon his feet, and was visited there by admiring and curious throngs. When this failed to satisfy him, he invented, in 423, a new sort of holiness, and lived, some two days’ journey (forty miles) east of Antioch, for six and thirty years, until his death, upon a pillar, which at the last was nearly forty cubits high;\footnote{The first pillar, which he himself erected, and on which he lived four years, was six cubits (\textgreek{πήχεων}) high, the second twelve, the third twenty-two, and the fourth, which the people erected for him, and on which he spent twenty years, was thirty-six, according to Theodoret; others say forty. The top was only three feet in diameter. It probably had a railing, however, on which he could lean in sleep or exhaustion. So at least these pillars are drawn in pictures. Food was carried up to the pillar saints by their disciples on a ladder.} for the pillar was raised in proportion as he approached heaven and perfection. Here he could never lie nor sit, but only stand, or lean upon a post (probably a banister), or devoutly bow; in which last posture he almost touched his feet with his head—so flexible had his back been made by fasting. A spectator once counted in one day no less than twelve hundred and forty-four such genuflexions of the saint before the Almighty, and then gave up counting. He wore a covering of the skins of beasts, and a chain about his neck. Even the holy sacrament he took upon his pillar. There St. Symeon stood many long and weary days, and weeks, and months, and years, exposed to the scorching sun, the drenching rain, the crackling frost, the howling storm, living a life of daily death and martyrdom, groaning under the load of sin, never attaining to the true comfort and peace of soul which is derived from a child-like trust in Christ’s infinite merits, earnestly striving after a superhuman holiness, and looking to a glorious reward in heaven, and immortal fame on earth. Alfred Tennyson makes him graphically describe his experience in a monologue to God:

Yet Symeon was not only concerned about his own salvation. People streamed from afar to witness this standing wonder of the age. He spoke to all classes with the same friendliness, mildness, and love; only women he never suffered to come within the wall which surrounded his pillar. From this original pulpit, as a mediator between heaven and earth, he preached repentance twice a day to the astonished spectators, settled controversies, vindicated the orthodox faith, extorted laws even from an emperor, healed the sick wrought miracles, and converted thousands of heathen Ishmaelites, Iberians, Armenians, and Persians to Christianity, or at least to the Christian name. All this the celebrated Theodoret relates as an eyewitness during the lifetime of the saint. He terms him the great wonder of the world,\footnote{\textgreek{Τὸ μέγα θαῦμα τῆς οἰκουμένης}. Hist. Relig. c. 26, at the beginning.} and compares him to a candle on a candlestick, and to the sun itself, which sheds its rays on every side. He asks the objector to this mode of life to consider that God often uses very striking means to arouse the negligent, as the history of the prophets shows;\footnote{Referring to Isa xx. 2; Jer. i. 17; xxviii. 12; Hos i. 2; iii. 1; Ezek. iv. 4; xii. 5.} and concludes his narrative with the remark: “Should the saint live longer, he may do yet greater wonders, for he is a universal ornament and honor of religion.”

He died in 459, in the sixty-ninth year of his age, of a long-concealed and loathsome ulcer on his leg; and his body was brought in solemn procession to the metropolitan church Of Antioch.

Even before his death, Symeon enjoyed the unbounded admiration of Christians and heathens, of the common people, of the kings of Persia, and of the emperors Theodosius II., Leo, and Marcian, who begged his blessing and his counsel. No wonder, that, with all his renowned humility, he had to struggle with the temptations of spiritual pride. Once an angel appeared to him in a vision, with a chariot of fire, to convey him, like Elijah, to heaven, because the blessed spirits longed for him. He was already stepping into the chariot with his right foot, which on this occasion he sprained (as Jacob his thigh), when the phantom of Satan was chased away by the sign of the cross. Perhaps this incident, which the Acta Sanctorum gives, was afterward invented, to account for his sore, and to illustrate the danger of self-conceit. Hence also the pious monk Nilus, with good reason, reminded the ostentatious pillar saints of the proverb: “He that exalteth himself shall be abased.”\footnote{Ep. ii. 114; cited in Gieseler, ii. 2, p. 246, note 47 (Edinb. Engl. ed. ii. p. 13, note 47), and in Neander.}

Of the later stylites the most distinguished were Daniel († 490), in the vicinity of Constantinople, and Symeon the younger († 592), in Syria. The latter is said to have spent sixty-eight years on a pillar. In the East this form of sanctity perpetuated itself, though only in exceptional cases, down to the twelfth century. The West, so far as we know, affords but one example of a stylite, who, according to Gregory of Tours, lived a long time on a pillar near Treves, but came down at the command of the bishop, and entered a neighboring cloister.

\xsection{Pachomius and the Cloister life}

§ 38. Pachomius and the Cloister life.

On St. Pachomius we have a biography composed soon after his death by a monk of Tabennae, and scattered accounts in Palladius, Jerome (Regula Pachomii, Latine reddita, Opp. Hieron. ed. Vallarsi, tom. ii. p. 50 sqq.), Rufinus, Sozomen, \&c. Comp. Tillemont, tom. vii. p. 167–235, and the Vit. Sanct. sub Maj. 14.

Though the strictly solitary life long continued in use, and to this day appears here and there in the Greek and Roman churches, yet from the middle of the fourth century monasticism began to assume in general the form of the cloister life, as incurring less risk, being available for both sexes, and being profitable to the church. Anthony himself gave warning, as we have already observed, against the danger of entire isolation, by referring to the proverb: “Woe to him that is alone.” To many of the most eminent ascetics anchoretism was a stepping stone to the coenobite life; to others it was the goal of coenobitism, and the last and highest round on the ladder of perfection.

The founder of this social monachism was Pachomius, a contemporary of Anthony, like him an Egyptian, and little below him in renown among the ancients. He was born about 292, of heathen parents, in the Upper Thebaid, served as a soldier in the army of the tyrant Maximin on the expedition against Constantine and Licinius, and was, with his comrades, so kindly treated by the Christians at Thebes, that he was won to the Christian faith, and, after his discharge from the military service, received baptism. Then, in 313, he visited the aged hermit Palemon, to learn from him the way to perfection. The saint showed him the difficulties of the anchorite life: “Many,” said he, “have come hither from disgust with the world, and had no perseverance. Remember, my son, my food consists only of bread and salt; I drink no wine, take no oil, spend half the night awake, singing psalms and meditating on the Scriptures, and sometimes pass the whole night without sleep.” Pachomius was astounded, but not discouraged, and spent several years with this man as a pupil.

In the year 325 he was directed by an angel, in a vision, to establish on the island of Tabennae, in the Nile, in Upper Egypt, a society of monks, which in a short time became so strong that even before his death (348) it numbered eight or nine cloisters in the Thebaid, and three thousand (according to some, seven thousand), and, a century later, fifty thousand members. The mode of life was fixed by a strict rule of Pachomius, which, according to a later legend, an angel communicated to him, and which Jerome translated into Latin. The formal reception into the society was preceded by a three-years’ probation. Rigid vows were not yet enjoined. With spiritual exercises manual labor was united, agriculture, boat building, basketmaking, mat and coverlet weaving, by which the monks not only earned their own living, but also supported the poor and the sick. They were divided, according to the grade of their ascetic piety, into four and twenty classes, named by the letters of the Greek alphabet. They lived three in a cell. They ate in common, but in strict silence, and with the face covered. They made known their wants by signs. The sick were treated with special care. On Saturday and Sunday they partook of the communion. Pachomius, as abbot, or archimandrite, took the oversight of the whole; each cloister having a separate superior and a steward.

Pachomius also established a cloister of nuns for his sister, whom he never admitted to his presence when she would visit him, sending her word that she should be content to know that he was still alive. In like manner, the sister of Anthony and the wife of Ammon became centres of female cloister life, which spread with great rapidity.

Pachomius, after his conversion never ate a full meal, and for fifteen years slept sitting on a stone. Tradition ascribes to him all sorts of miracles, even the gift of tongues and perfect dominion over nature, so that he trod without harm on serpents and scorpions, and crossed the Nile on the backs of crocodiles!\footnote{Möhler remarks on this (Vermischte Schriften, ii. p. 183): “Thus antiquity expresses its faith, that for man perfectly reconciled with God there is no enemy in nature. There is more than poetry here; there is expressed at least the high opinion his own and future generations had of Pachomius.” The last qualifying remark suggests a doubt even in the mind of this famous modern champion of Romanism as to the real historical character of the wonderful tales of this monastic saint.} Soon after Pachomius, fifty monasteries arose on the Nitrian mountain, in no respect inferior to those in the Thebaid. They maintained seven bakeries for the benefit of the anchorets in the neighboring Libyan desert, and gave attention also, at least in later days, to theological studies; as the valuable manuscripts recently discovered there evince.

From Egypt the cloister life spread with the rapidity of the irresistible spirit of the age, over the entire Christian East. The most eminent fathers of the Greek church were either themselves monks for a time, or at all events friends and patrons of monasticism. Ephraim propagated it in Mesopotamia; Eustathius of Sebaste in Armenia and Paphlagonia; Basil the Great in Pontus and Cappadocia. The latter provided his monasteries and nunneries with clergy, and gave them an improved rule, which, before his death (379), was accepted by some eighty thousand monks, and translated by Rufinus into Latin. He sought to unite the virtues of the anchorite and coenobite life, and to make the institution useful to the church by promoting the education of youth, and also (as Athanasius designed before him) by combating Arianism among the people.\footnote{Gregory Nazianzen, in his eulogy on Basil (Orat. xx. of the old order, Orat. xliii. in the new Par. ed.), gives him the honor of endeavoring to unite the theoretical and the practical modes of life in monasticism, \textgreek{ἲνα μήτε τὸ φιλόσοφον ἀκοινώνητον ᾗ, μήτε τὸ πρακτικὸν ἀφιλόσοφον}.} He and his friend Gregory Nazianzen were the first to unite scientific theological studies with the ascetic exercises of solitude. Chrysostom wrote three books in praise and vindication of the monastic life, and exhibits it in general in its noblest aspect.

In the beginning of the fifth century, Eastern monasticism was most worthily represented by the elder Nilus of Sinai, a pupil and venerator of Chrysostom, and a copious ascetic writer, who retired with his son from a high civil office in Constantinople to Mount Sinai, while his wife, with a daughter, travelled to an Egyptian cloister;\footnote{Comp. Neander, iii. 487 (Torrey’s translation, vol. ii. p. 250 sqq.), who esteems Nilus highly; and the article of Gass in Herzog’s Theol. Encykl. vol. x. p, 355 sqq. His works are in the Bibl. Max. vet. Patr. tom. vii., and in Migne’s Patrol. Gr. t. 79.} and by the abbot Isidore, of Pelusium, on the principal eastern mouth of the Nile, from whom we have two thousand epistles.\footnote{Comp. on him Tillemont, xv., and H. A. Niemeyer: “De Isid. Pel. vita, scripet doctrina,” Hal. 1825. His Epistles are in the 7th volume of the Bibliotheca Maxima, and in Migne’s Patrol. Graeca, tom. 58, Paris, 1860.} The writings of these two men show a rich spiritual experience, and an extended and fertile field of labor and usefulness in their age and generation.

\xsection{Fanatical and Heretical Monastic Societies in The East}

§ 39. Fanatical and Heretical Monastic Societies in The East.

Acta Concil. Gangrenensis, in Mansi, ii. 1095 sqq. Epiphan.: Haer. 70, 75 and 80. Socr.: H. E. ii. 43. Sozom.: iv. 24. Theodor.: H. E. iv. 9, 10; Fab. haer. iv. 10, 11. Comp. Neander: iii. p. 468 sqq. (ed. Torrey, ii. 238 sqq.).

Monasticism generally adhered closely to the orthodox faith of the church. The friendship between Athanasius, the father of orthodoxy, and Anthony, the father of monachism, is on this point a classical fact. But Nestorianism also, and Eutychianism, Monophysitism, Pelagianism, and other heresies, proceeded from monks, and found in monks their most vigorous advocates. And the monastic enthusiasm ran also into ascetic heresies of its own, which we must notice here.

1. The Eustathians, so named from Eustathius, bishop of Sebaste and friend of Basil, founder of monasticism in Armenia, Pontus, and Paphlagonia. This sect asserted that marriage debarred from salvation and incapacitated for the clerical office. For this and other extravagances it was condemned by a council at Gangra in Paphlagonia (between 360 and 370), and gradually died out.

2. The Audians held similar principles. Their founder, Audius, or Udo, a layman of Syria, charged the clergy of his day with immorality, especially avarice and extravagance. After much persecution, which he bore patiently, he forsook the church, with his friends, among whom were some bishops and priests, and, about 330, founded a rigid monastic sect in Scythia, which subsisted perhaps a hundred years. They were Quartodecimans in the practice of Easter, observing it on the 14th of Nisan, according to Jewish fashion. Epiphanius speaks favorably of their exemplary but severely ascetic life.

3. The Euchites or Messalians,\footnote{From \texthebrew{ןילִצְלִמַ }\emph{=} \textgreek{Εὐχίται–ϊ, –ϊ} from \textgreek{εὐχη–ΐ,–ϊ} \emph{prayer}.} also called Enthusiasts, were roaming mendicant monks in Mesopotamia and Syria (dating from 360), who conceived the Christian life as an unintermitted prayer, despised all physical labor, the moral law, and the sacraments, and boasted themselves perfect. They taught, that every man brings an evil demon with him into the world, which can only be driven away by prayer; then the Holy Ghost comes into the soul, liberates it from all the bonds of sense, and raises it above the need of instruction and the means of grace. The gospel history they declared a mere allegory. But they concealed their pantheistic mysticism and antinomianism under external conformity to the Catholic church. When their principles, toward the end of the fourth century, became known, the persecution of both the ecclesiastical and the civil authority fell upon them. Yet they perpetuated themselves to the seventh century, and reappeared in the Euchites and Bogomiles of the middle age.

\xsection{Monasticism in the West. Athanasius, Ambrose, Augustine, Martin of Tours}

§ 40. Monasticism in the West. Athanasius, Ambrose, Augustine, Martin of Tours.

I. Ambrosius: De Virginibus ad Marcellinam sororem suam libri tres, written about 377 (in the Benedictine edition of Ambr. Opera, tom. ii. p. 145–183). Augustinus (a.d. 400): De Opere Monachorum liber unus (in the Bened. ed., tom. vi. p. 476–504). Sulpitius Severus (about a.d. 403): Dialogi tres (de virtutibus monachorum orientalium et de virtutibus B. Martini); and De Vita Beati Martini (both in the Bibliotheca Maxima vet. Patrum, tom. vi. p. 349 sqq., and better in Gallandi’s Bibliotheca vet. Patrum, tom. viii. p. 392 sqq.).

II. J. Mabillon: Observat. de monachis in occidente ante Benedictum (Praef. in Acta Sanct. Ord. Bened.). R. H. Milman: Hist. of Latin Christianity, Lond. 1854, vol. i. ch. vi. p. 409–426: “Western Monasticism.” Count de Montalembert: The Monks of the West, Engl. translation, vol. i. p. 379 sqq.

In the Latin church, in virtue partly of the climate, partly of the national character,\footnote{Sulpitius Severus, in the first of his three dialogues, gives several amusing instances of the difference between the Gallic and Egyptian stomach, and was greatly astonished when the first Egyptian anchoret whom he visited placed before him and his four companions a half loaf of barley bread and a handful of herbs for a dinner, though they tasted very good after the wearisome journey. “Edacitas,” says he, “in Graecis gula est, in Gallia natura.” (Dial. i. c. 8, in Gallandi, t. viii. p. 405.)} the monastic life took a much milder form, but assumed greater variety, and found a larger field of usefulness than in the Greek. It produced no pillar saints, nor other such excesses of ascetic heroism, but was more practical instead, and an important instrument for the cultivation of the soil and the diffusion of Christianity and civilization among the barbarians.\footnote{“The monastic stream,” says Montalembert, l.c., “which had been born in the deserts of Egypt, divided itself into two great arms. The one spread in the East, at first inundated everything, then concentrated and lost itself there. The other escaped into the West, and spread itself by a thousand channels over an entire world, which had to be covered and fertilized.”} Exclusive contemplation was exchanged for alternate contemplation and labor. “A working monk,” says Cassian, “is plagued by one devil, an inactive monk by a host.” Yet it must not be forgotten that the most eminent representatives of the Eastern monasticism recommended manual labor and studies; and that the Eastern monks took a very lively, often rude and stormy part in theological controversies. And on the other hand, there were Western monks who, like Martin of Tours, regarded labor as disturbing contemplation.

Athanasius, the guest, the disciple, and subsequently the biographer and eulogist of St. Anthony, brought the first intelligence of monasticism to the West, and astounded the civilized and effeminate Romans with two live representatives of the semi-barbarous desert-sanctity of Egypt, who accompanied him in his exile in 340. The one, Ammonius, was so abstracted from the world that he disdained to visit any of the wonders of the great city, except the tombs of St. Peter and St. Paul; while the other, Isidore, attracted attention by his amiable simplicity. The phenomenon excited at first disgust and contempt, but soon admiration and imitation, especially among women, and among the decimated ranks of the ancient Roman nobility. The impression of the first visit was afterward strengthened by two other visits of Athanasius to Rome, and especially by his biography of Anthony, which immediately acquired the popularity and authority of a monastic gospel. Many went to Egypt and Palestine, to devote themselves there to the new mode of life; and for the sake of such, Jerome afterward translated the rule of Pachomius into Latin. Others founded cloisters in the neighborhood of Rome, or on the ruins of the ancient temples and the forum, and the frugal number of the heathen vestals was soon cast into the shade by whole hosts of Christian virgins. From Rome, monasticism gradually spread over all Italy and the isles of the Mediterranean, even to the rugged rocks of the Gorgon and the Capraja, where the hermits, in voluntary exile from the world, took the place of the criminals and political victims whom the justice or tyranny and jealousy of the emperors had been accustomed to banish thither.

Ambrose, whose sister, Marcellina, was among the first Roman nuns, established a monastery in Milan,\footnote{Augustine, Conf. vii. 6: “Erat monasterium Mediolani plenum bonis fratribus extra urbis moenia, sub Ambrosio nutritore.”} one of the first in Italy, and with the warmest zeal encouraged celibacy even against the will of parents; insomuch that the mothers of Milan kept their daughters out of the way of his preaching; whilst from other quarters, even from Mauritania, virgins flocked to him to be consecrated to the solitary life.\footnote{Ambr.: De virginibus, lib. iii., addressed to his sister Marcellina, about 377. Comp. Tillem. x. 102-105, and Schröckh, viii. 355 sqq.} The coasts and small islands of Italy were gradually studded with cloisters.\footnote{Ambr.: Hexaëmeron, l. iii. c. 5. Hieron.: Ep. ad Oceanum de morte Fabiolae, Ep. 77 ed. Vall. (84 ed. Ben., al. 30).}

Augustine, whose evangelical principles of the free grace of God as the only ground of salvation and peace were essentially inconsistent with the more Pelagian theory of the monastic life, nevertheless went with the then reigning spirit of the church in this respect, and led, with his clergy, a monk-like life in voluntary poverty and celibacy,\footnote{He himself speaks of a monasterium clericorum in his episcopal residence, and his biographer, Possidius, says of him, Vita, c. 5: “Factus ergo presbyter monasterium inter ecclesiam mox instituit, et cum Dei servis vivere coepit secundum modum, et regulam sub sanctis apostlis constitutam, maxime ut nemo quidquam proprium haberet, sed eis essent omnia communia.”} after the pattern, as he thought, of the primitive church of Jerusalem; but with all his zealous commendation he could obtain favor for monasticism in North Africa only among the liberated slaves and the lower classes.\footnote{De opera monach. c. 22. Still later, Salvian (De gubern. Dei, viii. 4) speaks of the hatred of the Africans for monasticism.} He viewed it in its noblest aspect, as a life of undivided surrender to God, and undisturbed occupation with spiritual and eternal things. But he acknowledged also its abuses; he distinctly condemned the vagrant, begging monks, like the Circumcelliones and Gyrovagi, and wrote a book (De opere monachorum) against the monastic aversion to labor.

Monasticism was planted in Gaul by Martin of Tours, whose life and miracles were described in fluent, pleasing language by his disciple, Sulpitius Severus,\footnote{In his Vita Martini, and also in three letters respecting him, and in three very eloquently and elegantly written dialogues, the first of which relates to the oriental monks, the two others to the miracles of Martin (translated, with some omissions, in Ruffner’s Fathers of the Desert, vol. ii. p. 68-178). He tells us (Dial. i. c. 23) that the book traders of Rome sold his Vita Martini more rapidly than any other book, and made great profit on it. The Acts of the Saints were read as romances in those days.} a few years after his death. This celebrated saint, the patron of fields, was born in Pannonia (Hungary), of pagan parents. He was educated in Italy, and served three years, against his will, as a soldier under Constantius and Julian the Apostate. Even at that time he showed an uncommon degree of temperance, humility, and love. He often cleaned his servant’s shoes, and once cut his only cloak in two with his sword, to clothe a naked beggar with half; and the next night he saw Christ in a dream with the half cloak, and plainly heard him say to the angels: “Behold, Martin, who is yet only a catechumen, hath clothed me.”\footnote{The biographer here refers, of course, to Matt. xxv. 40} He was baptized in his eighteenth year; converted his mother; lived as a hermit in Italy; afterward built a monastery in the vicinity of Poictiers (the first in France); destroyed many idol temples, and won great renown as a saint and a worker of miracles. About the year 370 he was unanimously elected by the people, against his wish, bishop of Tours on the Loire, but in his episcopal office maintained his strict monastic mode of life, and established a monastery beyond the Loire, where he was soon surrounded with eighty monks. He had little education, but a natural eloquence, much spiritual experience, and unwearied zeal. Sulpitius Severus places him above all the Eastern monks of whom he knew, and declares his merit to be beyond all expression. “Not an hour passed,” says he,\footnote{Toward the close of his biography, c. 26, 27 (Gallandi, tom. viii. 399).} “in which Martin did not pray .... No one ever saw him angry, or gloomy, or merry. Ever the same, with a countenance full of heavenly serenity, he seemed to be raised above the infirmities of man. There was nothing in his mouth but Christ; nothing in his heart but piety, peace, and sympathy. He used to weep for the sins of his enemies, who reviled him with poisoned tongues when he was absent and did them no harm .... Yet he had very few persecutors, except among the bishops.” The biographer ascribes to him wondrous conflicts with the devil, whom he imagined he saw bodily and tangibly present in all possible shapes. He tells also of visions, miraculous cures, and even, what no oriental anchoret could boast, three instances of restoration of the dead to life, two before and one after his accession to the bishopric;\footnote{Comp. Dial. ii. 5 (in Gallandi Bibl. tom. viii. p. 412).} and he assures us that he has omitted the greater part of the miracles which had come to his ears, lest he should weary the reader; but he several times intimates that these were by no means universally credited, even by monks of the same cloister. His piety was characterized by a union of monastic humility with clerical arrogance. At a supper at the court of the tyrannical emperor Maximus in Trier, he handed the goblet of wine, after he himself had drunk of it, first to his presbyter, thus giving him precedence of the emperor.\footnote{Vita M. c. 20 (in Gallandi, viii. 397).} The empress on this occasion showed him an idolatrous veneration, even preparing the meal, laying the cloth, and standing as a servant before him, like Martha before the Lord.\footnote{Dial. ii. 7, which probably relates to the same banquet, since Martin declined other invitations to the imperial table. Severus gives us to understand that this was the only time Martin allowed a woman so near him, or received her service. He commended a nun for declining even his official visit as bishop, and Severus remarks thereupon: “O glorious virgin, who would not even suffer herself to be seen by Martin! O blessed Martin, who took not this refusal for an insult, but commended its virtue, and rejoiced to find in that region so rare an example!” (Dial, ii. c. 12, Gall, viii. 414.)} More to the bishop’s honor was his protest against the execution of the Priscillianists in Treves. Martin died in 397 or 400: his funeral was attended by two thousand monks, besides many nuns and a great multitude of people; and his grave became one of the most frequented centres of pilgrimage in France.

In Southern Gaul, monasticism spread with equal rapidity. John Cassian, an ascetic writer and a Semipelagian († 432), founded two cloisters in Massilia (Marseilles), where literary studies also were carried on; and Honoratus (after 426, bishop of Arles) established the cloister of St. Honoratus on the island of Lerina.

\xsection{St. Jerome as a Monk}

§ 41. St. Jerome as a Monk.

S. Eus. Hieronymi: Opera omnia, ed. Erasmus (assisted by Oecolampadius), Bas. 1516–’20, 9 vols. fol.; ed. (Bened.) Martianay, Par. 1693–1706, 5 vols. fol. (incomplete); ed. Vallarsi and Maffei, Veron. 1734–’42, 11 vols. fol., also Venet. 1766 (best edition). Comp. especially the 150 Epistles, often separately edited (the chronological order of which Vallarsi, in tom. i. of his edition, has finally established).

For extended works on the life of Jerome see Du Pin (Nouvelle Biblioth. des auteurs Eccles. tom. iii. p. 100–140); Tillemont (tom. xii. 1–356); Martianay (La vie de St. Jerôme, Par. 1706); Joh. Stilting (in the Acta Sanctorum, Sept. tom. viii. p. 418–688, Antw. 1762); Butler (sub Sept. 30); Vallarsi (in Op. Hieron., tom. xi. p. 1–240); Schröckh (viii. 359 sqq., and especially xi. 3–254); Engelstoft (Hieron. Stridonensis, interpres, criticus, exegeta, apologeta, historicus, doctor, monachus, Havn. 1798); D. v. Cölln (in Ersch and Gruber’s Encycl. sect. ii. vol. 8); Collombet (Histoire de S. Jérôme, Lyons, 1844); and O. Zöckler (Hieronymus, sein Leben und Wirken. Gotha, 1865).

The most zealous promoter of the monastic life among the church fathers was Jerome, the connecting link between Eastern and Western learning and religion. His life belongs almost with equal right to the history of theology and the history of monasticism. Hence the church art generally represents him as a penitent in a reading or writing posture, with a lion and a skull, to denote the union of the literary and anchoretic modes of life. He was the first learned divine who not only recommended but actually embraced the monastic mode of life, and his example exerted a great influence in making monasticism available for the promotion of learning. To rare talents and attainments,\footnote{As he himself boasts in his second apology to Rufinus: “Ego philosophus(?), rhetor, grammaticus, dialecticus, hebraeus, graecus, latinus, trilinguis.” The celebrated Erasmus, the first editor of his works, and a very competent judge in matters of literary talent and merit, places Jeromeabove all the fathers, even St. Augustine(with whose doctrines of free grace and predestination he could not sympathize), and often gives eloquent expression to his admiration for him. In a letter to Pope Leo X. (Ep. ii. 1, quoted in Vallarsi’s ed. of Jerome’s works, tom. xi. 290), he says: “Divus Hieronymus sic apud Latinos est theologorum princeps, ut hunc prope solum habeamus theologi dignum nomine. Non quod caeteros damnem, sed quod illustres alioqui, si cum hoc conferantur, ob huius eminentiam velut obscurentur. Denique tot egregiis est cumulatus dotibus, ut vix ullum habeat et ipsa docta Graecia, quem cum hoc viro quest componere. Quantum in illo Romanae facundiae! quanta linguarum peritia! quanta omnis antiquitatis omnium historiarum notitia! quam fida memoria! quam felix rerum omnium mixtum! quam absoluta mysticarum literarum cognitio! super omnia, quis ardor ille, quam admirabilis divini pectoris afflatus? ut una et plurimum delectet eloquentia, et doceat eruditione, et rapiat sanctimonia.”} indefatigable activity of mind, ardent faith, immortal merit in the translation and interpretation of the Bible, and earnest zeal for ascetic piety, he united so great vanity and ambition, such irritability and bitterness of temper, such vehemence of uncontrolled passion, such an intolerant and persecuting spirit, and such inconstancy of conduct, that we find ourselves alternately attracted and repelled by his character, and now filled with admiration for his greatness, now with contempt or pity for his weakness.

Sophronius Eusebius Hieronymus was born at Stridon,\footnote{Hence called \emph{Stridonensis}; also in distinction from the contemporary but little known Greek Jerome, who was probably a presbyter in Jerusalem.} on the borders of Dalmatia, not far from Aquileia, between the years 331 and 342.\footnote{Martianay, Stilting, Cave, Schröckh, Hagenbach, and others, place his birth, according to Prosper, Chron. ad ann. 331, in the year 331; Baronius, Du Pin, and Tillemont, with greater probability, in the year 342. The last infers from various circumstances, that Jeromelived, not ninety-one years, as Prosper states, but only seventy-eight. Vallarsi (t. xi. 8) places his birth still later, in the year 346. His death is placed in the year 419 or 420.} He was the son of wealthy Christian parents, and was educated in Rome under the direction of the celebrated heathen grammarian Donatus, and the rhetorician Victorinus. He read with great diligence and profit the classic poets, orators, and philosophers, and collected a considerable library. On Sundays he visited, with Bonosus and other young friends, the subterranean graves of the martyrs, which made an indelible impression upon him. Yet he was not exempt from the temptations of a great and corrupt city, and he lost his chastity, as he himself afterward repeatedly acknowledged with pain.

About the year 370, whether before or after his literary tour to Treves and Aquileia is uncertain, but at all events in his later youth, he received baptism at Rome and resolved thenceforth to devote himself wholly, in rigid abstinence, to the service of the Lord. In the first zeal of his conversion he renounced his love for the classics, and applied himself to the study of the hitherto distasteful Bible. In a morbid ascetic frame, he had, a few years later, that celebrated dream, in which he was summoned before the judgment seat of Christ, and as a heathen Ciceronian,\footnote{“Mentiris,” said the Lord to him, when Jeromecalled himself a Christian, Ciceronianus es, non Christianus, ubi enim thesaurus tuus ibi et cor tuum.” Ep. xxii. ad Eustochium, “De custodia virginitatis ”(tom. i. p. 113\emph{).} C. A. Heumann has written a special treatise, De ecstasi Hieronymi anti-Ciceroniana. Comp. also Schröckh, vol. vii. p. 35 sqq., and Ozanam: ” Civilisation au 5e Siècle,” i. 301.} so severely reprimanded and scourged, that even the angels interceded for him from sympathy with his youth, and he himself solemnly vowed never again to take worldly books into his hands. When he woke, he still felt the stripes, which, as he thought, not his heated fancy, but the Lord himself had inflicted upon him. Hence he warns his female friend Eustochium, to whom several years afterward (a.d. 384) he recounted this experience, to avoid all profane reading: “What have light and darkness, Christ and Belial (2 Cor. vi. 14), the Psalms and Horace, the Gospels and Virgil, the Apostles and Cicero, to do with one another? ... We cannot drink the cup of the Lord and the cup of the demons at the same time.”\footnote{Ep. xxii. ed. Vall. i. 112).} But proper as this warning may be against overrating classical scholarship, Jerome himself, in his version of the Bible and his commentaries, affords the best evidence of the inestimable value of linguistic and antiquarian knowledge, when devoted to the service of religion. That oath, also, at least in later life, he did not strictly keep. On the contrary, he made the monks copy the dialogues of Cicero, and explained Virgil at Bethlehem, and his writings abound in recollections and quotations of the classic authors. When Rufinus of Aquileia, at first his warm friend, but afterward a bitter enemy, cast up to him this inconsistency and breach of a solemn vow, he resorted to the evasion that he could not obliterate from his memory what he had formerly read; as if it were not so sinful to cite a heathen author as to read him. With more reason he asserted, that all was a mere dream, and a dream vow was not binding. He referred him to the prophets, “who teach that dreams are vain, and not worthy of faith.” Yet was this dream afterward made frequent use of, as Erasmus laments, to cover monastic obscurantism.

After his baptism, Jerome divided his life between the East and the West, between ascetic discipline and literary labor. He removed from Rome to Antioch with a few friends and his library, visited the most celebrated anchorets, attended the exegetical lectures of the younger Apollinaris in Antioch, and then (374) spent some time as an ascetic in the dreary Syrian desert of Chalcis. Here, like so many other hermits, he underwent a grevious struggle with sensuality, which he described ten years after with indelicate minuteness in a long letter to his virgin friend Eustochium.\footnote{Ep. xxii. (i. p. 91, ed. Vallars.)} In spite of his starved and emaciated body, his fancy tormented him with wild images of Roman banquets and dances of women; showing that the monastic seclusion from the world was by no means proof against the temptations of the flesh and the devil. Helpless he cast himself at the feet of Jesus, wet them with tears of repentance, and subdued the resisting flesh by a week of fasting and by the dry study of Hebrew grammar (which, according to a letter to Rusticus,\footnote{Ep. cxxv., ed. Vallars. (al. 95 or 4.)} he was at that time learning from a converted Jew), until he found peace, and thought himself transported to the choirs of the angels in heaven. In this period probably falls the dream mentioned above, and the composition of several ascetic writings, full of heated eulogy of the monastic life.\footnote{De laude vitae solitariae, Ep. xiv. (tom. i. 28-36) ad Heliodorum. The Roman lady Fabiola learned this letter by heart, and Du Pin calls it a masterpiece of eloquence (Nouv. Bibl. des auteurs eccl. iii. 102), but it is almost too declamatory and turgid. He himself afterward acknowledged it overdrawn.} His biographies of distinguished anchorets, however, are very pleasantly and temperately written.\footnote{Gibbon says of them: “The stories of Paul, Hilarion, and Malchus are admirably told; and the only defect of these pleasing compositions is the want of truth and common sense.”} He commends monastic seclusion even against the will of parents; interpreting the word of the Lord about forsaking father and mother, as if monasticism and Christianity were the same. “Though thy mother”—he writes, in 373, to his friend Heliodorus, who had left him in the midst of his journey to the Syrian desert—“with flowing hair and rent garments, should show thee the breasts which have nourished thee; though thy father should lie upon the threshold; yet depart thou, treading over thy father, and fly with dry eyes to the standard of the cross. This is the only religion of its kind, in this matter to be cruel .... The love of God and the fear of hell easily, rend the bonds of the household asunder. The holy Scripture indeed enjoins obedience to parents; but he who loves them more than Christ, loses his soul .... O desert, where the flowers of Christ are blooming!. O solitude, where the stones for the new Jerusalem are prepared! O retreat, which rejoices in the friendship of God! What doest thou in the world, my brother, with thy soul greater than the world? How long wilt thou remain in the shadow of roofs, and in the smoky dungeon of cities? Believe me, I see here more of the light.”\footnote{Ep. xiv. (t. i. 29 sq.) Similar descriptions of the attractions of monastic life we meet with in the ascetic writings of Gregory, Basil, Ambrose, Chrysostom, Cassian, Nilus, and Isidor. “So great grace,” says the venerable monk Nilus of Mount Sinai, in the beginning of the fifth century (Ep. lib. i Ep. 1, as quoted by Neander, Am. ed. ii. 250), “so great grace his God bestowed on the monks, even in anticipation of the future world, that they wish for no honors from men, and feel no longing after the greatness of this world; but, on the contrary, often seek rather to remain concealed from men: while, on the other hand, many of the great, who possess all the glory of the world, either of their own accord, or compelled by misfortune, take refuge with the lowly monks, and, delivered from fatal dangers, obtain at once a temporal and an eternal salvation.”} The eloquent appeal, however, failed of the desired effect; Heliodorus entered the teaching order and became a bishop.

The active and restless spirit of Jerome soon brought him again upon the public stage, and involved him in all the doctrinal and ecclesiastical controversies of those controversial times. He received the ordination of presbyter from the bishop Paulinus in Antioch, without taking charge of a congregation. He preferred the itinerant life of a monk and a student to a fixed office, and about 380 journeyed to Constantinople, where he heard the anti-Arian sermons of the celebrated Gregory Nazianzen, and translated the Chronicle of Eusebius and the homilies of Origen on Jeremiah and Ezekiel. In 382, on account of the Meletian schism, he returned to Rome with Paulinus and Epiphanius. Here he came into close connection with the bishop, Damasus, as his theological adviser and ecclesiastical secretary,\footnote{As we infer from a remark of Jeromein Ep. cxxiii. c. 10, written a. 409 (ed. Vallars. i. p. 901): “Ante annos plurimos, quum in chartis ecclesiasticis” (i.e. probably in ecclesiastical documents; though Schröckh, viii. p. 122, refers it to the Holy Scriptures, appealing to a work of Bonamici unknown to me), “juvarem Damasum, Romanae urbis episcopum, et orientis atque occidentis synodicis consultationibus responderem,” etc. The latter words, which Schröckh does not quote, favor the common interpretation.} and was led by him into new exegetical labors, particularly the revision of the Latin version of the Bible, which he completed at a later day in the East.

At the same time he labored in Rome with the greatest zeal, by mouth and pen, in the cause of monasticism, which had hitherto gained very little foothold there, and met with violent opposition even among the clergy. He had his eye mainly upon the most wealthy and honorable classes of the decayed Roman society, and tried to induce the descendants of the Scipios, the Gracchi, the Marcelli, the Camilli, the Anicii to turn their sumptuous villas into monastic retreats, and to lead a life of self-sacrifice and charity. He met with great success. “The old patrician races, which founded Rome, which had governed her during all her period of splendor and liberty, and which overcame and conquered the world, had expiated for four centuries, under the atrocious yoke of the Caesars, all that was most hard and selfish in the glory of their fathers. Cruelly humiliated, disgraced, and decimated during that long servitude, by the masters whom degenerate Rome had given herself, they found at last in Christian life, such as was practised by the monks, the dignity of sacrifice and the emancipation of the soul. These sons of the old Romans threw themselves into it with the magnanimous fire and persevering energy which had gained for their ancestors the empire of the world. ’Formerly,’ says St. Jerome, ’according to the testimony of the apostles, there were few rich, few noble, few powerful among the Christians. Now it is no longer so. Not only among the Christians, but among the monks are to be found a multitude of the wise, the noble, and the rich.’... The monastic institution offered them a field of battle where the struggles and victories of their ancestors could be renewed and surpassed for a loftier cause, and over enemies more redoubtable. The great men whose memory hovered still over degenerate Rome had contended only with men, and subjugated only their bodies; their descendants undertook to strive with devils, and to conquer souls .... God called them to be the ancestors of a new people, gave them a new empire to found, and permitted them to bury and transfigure the glory of their forefathers in the bosom of the spiritual regeneration of the world.”\footnote{Montalembert, himself the scion of an old noble family in France, l.c. i. p. 388 sq. Comp. Hieron., Epist. lxvi. ad Pammachium, de obit. Paulinae (ed. Vallars. i. 391 sqq.).}

Most of these distinguished patrician converts of Jerome were women—such widows as Marcella, Albinia, Furia, Salvina, Fabiola, Melania, and the most illustrious of all, Paula, and her family; or virgins, as Eustochium, Apella, Marcellina, Asella, Felicitas, and Demetrias. He gathered them as a select circle around him; he expounded to them the Holy Scriptures, in which some of these Roman ladies were very well read; he answered their questions of conscience; he incited them to celibate life, lavish beneficence, and enthusiastic asceticism; and flattered their spiritual vanity by extravagant praises. He was the oracle, biographer, admirer, and eulogist of these holy women, who constituted the spiritual nobility of Catholic Rome. Even the senator Pammachius, son in-law to Paula and heir to her fortune, gave his goods to the poor, exchanged the purple for the cowl, exposed himself to the mockery of his colleagues, and became, in the flattering language of Jerome, the general in chief of Roman monks, the first of monks in the first of cities.\footnote{In one of his Epist. ad Pammach.: “Primus inter monachos in prima urbe ... archistrategos monachorum.”} Jerome considered second marriage incompatible with genuine holiness; even depreciated first marriage, except so far as it was a nursery of brides of Christ; warned Eustochium against all intercourse with married women; and hesitated not to call the mother of a bride of Christ, like Paula, a “mother-in-law of God.”\footnote{Ep. xxii. ad Eustochium, “de custodia virginitatis.” Even Rufinus was shocked at the profane, nay, almost blasphemous expression, \emph{socrus Dei,} and asked him from what \emph{heathen} poet he had stolen it.}

His intimacy with these distinguished women, whom he admired more, perhaps, than they admired him, together with his unsparing attacks upon the immoralities of the Roman clergy and of the higher classes, drew upon him much unjust censure and groundless calumny, which he met rather with indignant scorn and satire than with quiet dignity and Christian meekness. After the death of his patron Damasus, a.d. 384, he left Rome, and in August, 385, with his brother Paulinian, a few monks, Paula, and her daughter Eustochium, made a pilgrimage “from Babylon to Jerusalem, that not Nebuchadnezzar, but Jesus, should reign over him.” With religious devotion and inquiring mind he wandered through the holy places of Palestine, spent some time in Alexandria, where he heard the lectures of the celebrated Didymus; visited the cells of the Nitrian mountain; and finally, with his two female friends, in 386, settled in the birthplace of the Redeemer, to lament there, as he says, the sins of his youth, and to secure himself against others.

In Bethlehem he presided over a monastery till his death, built a hospital for all strangers except heretics, prosecuted his literary studies without cessation, wrote several commentaries, and finished his improved Latin version of the Bible—the noblest monument of his life—but entangled himself in violent literary controversies, not only with opponents of the church orthodoxy like Helvidius (against whom he had appeared before, in 384), Jovinian, Vigilantius, and Pelagius, but also with his long-tried friend Rufinus, and even with Augustine.\footnote{His controversy with Augustineon the interpretation of Gal. ii. 14 is not unimportant as an index of the moral character of the two most illustrous Latin fathers of the church. Jeromesaw in the account of the collision between Paul and Peter, in Antioch, an artifice of pastoral prudence, and supposed that Paul did not there reprove the senior apostle in earnest, but only for effect, to reclaim the Jews from their wrong notions respecting the validity of the ceremonial law. Augustine’s delicate sense of truth was justly offended by this exegesis, which, to save the dignity of Peter, ascribed falsehood to Paul, and he expressed his opinion to Jerome, who, however, very loftily made him feel his smaller grammatical knowledge. But they afterward became reconciled. Comp. on this dispute the letters on both sides, in Hieron. Opera, ed. Vall. tom. i. 632 sqq., and the treatise of Möhler, in his ”Vermischte Schriften,” vol. i. p. 1-18.} Palladius says, his jealousy could tolerate no saint beside himself, and drove many pious monks away from Bethlehem. He complained of the crowds of monks whom his fame attracted to Bethlehem.\footnote{“Tantis de toto orbe confluentibus obruimur turbis monachorum.”} The remains of the Roman nobility, too, ruined by the sack of Rome, fled to him for food and shelter. At the last his repose was disturbed by incursions of the barbarian Huns and the heretical Pelagians. He died in 419 or 420, of fever, at a great age. His remains were afterward brought to the Roman basilica of Maria Maggiore, but were exhibited also and superstitiously venerated in several copies in Florence, Prague, Clugny, Paris, and the Escurial.\footnote{The Jesuit Stilting, the author of the Vita Hieron. in the Acta Sanctorum, devotes nearly thirty folio pages to accounts of the veneration paid to him and his relics after his death.}

The Roman church has long since assigned him one of the first places among her standard teachers and canonical saints. Yet even some impartial Catholic historians venture to admit and disapprove his glaring inconsistencies and violent passions. The Protestant love of truth inclines to the judgment, that Jerome was indeed an accomplished and most serviceable scholar and a zealous enthusiast for all which his age counted holy, but lacking in calm self-control and proper depth of mind and character, and that he reflected, with the virtues, the failings also of his age and of the monastic system. It must be said to his credit, however, that with all his enthusiastic zeal and admiration for monasticism, he saw with a keen eye and exposed with unsparing hand the false monks and nuns, and painted in lively colors the dangers of melancholy, hypochondria, the hypocrisy and spiritual pride, to which the institution was exposed.\footnote{Most Roman Catholic biographers, as Martianay, Vallarsi, Stilting, Dolci, and even the Anglican Cave, are unqualified eulogists of Jerome. See also the “Selecta Veterum testimonia de Hieronymo ejusque scriptis,” in Vallarsi’s edition, tom. xi. pp. 282-300. Tillemont, however, who on account of his Jansenist proclivity sympathizes more with Augustine, makes a move toward a more enlightened judgment, for which Stilting sharply reproves him. Montalembert (l.c. i. 402) praises him as a man of genius, inspired by zeal and subdued by penitence, of ardent faith and immense resources of knowledge; yet he incidentally speaks also of his “almost savage impetuosity of temper,” and “that inexhaustible vehemence which sometimes degenerated into emphasis and affectation.” Dr. John H. Newman, in his opinion before his transition from Puseyism to Romanism, exhibits the conflict in which the moral feeling is here involved with the authority of the Roman Church: “I do not scruple to say, that, were he not a saint, there are things in his writings and views from which I should shrink; but as the case stands, I shrink rather from putting myself in opposition to something like a judgment of the catholic(?) world in favor of his saintly perfection.” (Church of the Fathers, 263, cited by Robertson.) Luther also here boldly broke through tradition, but, forgetful of the great value of the Vulgate even to his German version of the Bible, went to the opposite extreme of unjust derogation, expressing several times a distinct antipathy to this church father, and charging him with knowing not how to write at all of Christ, but only of fasts, virginity, and useless monkish exercises. Le Clerc exposed his defects with thorough ability, but unfairly, in his ”Quaestiones Hieronymianae“ (Amstel. 1700, over 500 pages). Mosheim and Schröckh are more mild, but the latter considers it doubtful whether Jeromedid Christianity more good than harm. Among later Protestant historians opinion has become somewhat more favorable, though rather to his learning than to his moral character, which betrays in his letters and controversial writings too many unquestionable weaknesses.}

\xsection{St. Paula}

§ 42. St. Paula.

Hieronymus: Epitaphium Paulae matris, ad Eustochium virginem, Ep. cviii. (ed. Vallarsi, Opera, tom. i. p. 684 sqq.; ed. Bened. Ep. lxxxvi). Also the Acta Sanctorum, and Butler’s Lives of Saints, sub Jan. 26.

Of Jerome’s many female disciples, the most distinguished is St. Paula, the model of a Roman Catholic nun. With his accustomed extravagance, he opens his eulogy after her death, in. 404, with these words: “If all the members of my body were turned into tongues, and all my joints were to utter human voices, I should be unable to say anything worthy of the holy and venerable Paula.”

She was born in 347, of the renowned stock of the Scipios and Gracchi and Paulus Aemilius,\footnote{Her father professed to trace his genealogy to Agamemnon, and her husband to Aeneas.} and was already a widow of six and thirty years, and the mother of five children, when, under the influence of Jerome, she renounced all the wealth and honors of the world, and betook herself to the most rigorous ascetic life. Rumor circulated suspicion, which her spiritual guide, however, in a letter to Asella, answered with indignant rhetoric: “Was there, then, no other matron in Rome, who could have conquered my heart, but that one, who was always mourning and fasting, who abounded in dirt,\footnote{This want of cleanliness, the inseparable companion of ancient ascetic holiness, is bad enough in monks, but still more intolerable and revolting in nuns.} who had become almost blind with weeping, who spent whole nights in prayer, whose song was the Psalms, whose conversation was the gospel, whose joy was abstemiousness, whose life was fasting? Could no other have pleased me, but that one, whom I have never seen eat? Nay, verily, after I had begun to revere her as her chastity deserved, should all virtues have at once forsaken me?” He afterward boasts of her, that she knew the Scriptures almost entirely by memory; she even learned Hebrew, that she might sing the psalter with him in the original; and continually addressed exegetical questions to him, which he himself could answer only in part.

Repressing the sacred feelings of a mother, she left her daughter Ruffina and her little son Toxotius, in spite of their prayers and tears, in the city, of Rome,\footnote{“Nesciebat se matrem,” says Jerome, “ut Christi probaret ancillam.” Revealing the conflict of monastic sanctity with the natural virtues which God has enjoined. Montalembert, also, quotes this objectionable passage with apparent approbation.} met Jerome in Antioch, and made a pilgrimage to Palestine and Egypt. With glowing devotion, she knelt before the rediscovered cross, as if the Lord were still hanging upon it; she kissed the stone of the resurrection which the angel rolled away; licked with thirsty tongue the pretended tomb of Jesus, and shed tears of joy as she entered the stable and beheld the manger of Bethlehem. In Egypt she penetrated into the desert of Nitria, prostrated herself at the feet of the hermits, and then returned to the holy land and settled permanently in the birthplace of the Saviour. She founded there a monastery for Jerome, whom she supported, and three nunneries, in which she spent twenty years as abbess, until 404.

She denied herself flesh and wine, performed, with her daughter Eustochium, the meanest services, and even in sickness slept on the bare ground in a hair shirt, or spent the whole night in prayer. “I must,” said she, “disfigure my face, which I have often, against the command of God, adorned with paint; torment the body, which has participated in many idolatries; and atone for long laughing by constant weeping.” Her liberality knew no bounds. She wished to die in beggary, and to be buried in a shroud which did not belong to her. She left to her daughter (she died in 419) a multitude of debts, which she had contracted at a high rate of interest for benevolent purposes.\footnote{Jeromesays, Eustochium hoped to pay the debts of her mother—probably by the help of others. Fuller justly remarks: “Liberality should have banks, as well as a stream.”}

Her obsequies, which lasted a week, were attended by the bishops of Jerusalem and other cities of Palestine, besides clergy, monks, nuns, and laymen innumerable. Jerome apostrophizes her: “Farewell, Paula, and help with prayer the old age of thy adorer!”

\xsection{Benedict of Nursia}

§ 43. Benedict of Nursia.

Gregorius M.: Dialogorum, l. iv. (composed about 594; lib. ii. contains the biography of St. Benedict according to the communications of four abbots and disciples of the saint, Constantine, Honoratus, Valentinian, and Simplicius, but full of surprising miracles). Mabillon and other writers of the Benedictine congregation of St. Maurus: Acta Sanctorum ordinis S. Benedicti in saeculorum classes distributa, fol. Par. 1668–1701, 9 vols. (to the year 1100), and Annales ordinis S. Bened. Par. 1703–’39, 6 vols. fol. (to 1157). Dom (Domnus) Jos. De Mège: Vie de St. Benoit, Par. 1690. The Acta Sanctorum, and Butler, sub Mart. 21. Montalembert: The Monks of the West, vol. ii. book iv.

Benedict of Nursia, the founder of the celebrated order which bears his name, gave to the Western monasticism a fixed and permanent form, and thus carried it far above the Eastern with its imperfect attempts at organization, and made it exceedingly profitable to the practical, and, incidentally, also to the literary interests of the Catholic Church. He holds, therefore, the dignity of patriarch of the Western monks. He has furnished a remarkable instance of the incalculable influence which a simple but judicious moral rule of life may exercise on many centuries.

Benedict was born of the illustrious house of Anicius, at Nursia (now Norcia) in Umbria, about the year 480, at the time when the political and social state of Europe was distracted and dismembered, and literature, morals, and religion seemed to be doomed to irremediable ruin. He studied in Rome, but so early as his fifteenth year he fled from the corrupt society of his fellow students, and spent three years in seclusion in a dark, narrow, and inaccessible grotto at Subiaco.\footnote{In Latin \emph{Sublaqueum}, or \emph{Sublacum}, in the States of the Church, over thirty English miles (Butler says “near forty,” Montalembert, ii. 7, \emph{“}fifty miles”) east of Rome, on the Teverone. Butler describes the place as “a barren, hideous chain of rocks, with a river and lake in the valley.”} A neighboring monk, Romanus, furnished him from time to time his scanty food, letting it down by a cord, with a little bell, the sound of which announced to him the loaf of bread. He there passed through the usual anchoretic battles with demons, and by prayer and ascetic exercises attained a rare power over nature. At one time, Pope Gregory tells us, the allurements of voluptuousness so strongly tempted his imagination that he was on the point of leaving his retreat in pursuit of a beautiful woman of previous acquaintance; but summoning up his courage, he took off his vestment of skins and rolled himself naked on thorns and briers, near his cave, until the impure fire of sensual passion was forever extinguished. Seven centuries later, St. Francis of Assisi planted on that spiritual battle field two rose trees, which grew and survived the Benedictine thorns and briers. He gradually became known, and was at first taken for a wild beast by the surrounding shepherds, but afterward reverenced as a saint.

After this period of hermit life he began his labors in behalf of the monastery proper. In that mountainous region he established in succession twelve cloisters, each with twelve monks and a superior, himself holding the oversight of all. The persecution of an unworthy priest caused him, however, to leave Subiaco and retire to a wild but picturesque mountain district in the Neapolitan province, upon the boundaries of Samnium and Campania. There he destroyed the remnants of idolatry, converted many of the pagan inhabitants to Christianity by his preaching and miracles, and in the year 529, under many difficulties, founded upon the ruins of a temple of Apollo the renowned cloister of Monte Cassino,\footnote{\emph{Monasterium Cassinense.} It was destroyed, indeed, by the Lombards, as early as 583, as Benedictis said to have predicted it would be, but was rebuilt in 731, consecrated in 748, again destroyed by the Saracens in 857, rebuilt about 950, and more completely, after many other calamities, in 1649, consecrated for the third time by BenedictXIII. in 1727, enriched and increased under the patronage of the emperors and popes, but in modern times despoiled of its enormous income (which at the end of the sixteenth century was reckoned at 500,000 ducats), and has stood through all vicissitudes to this day. In the days of its splendor, when the abbot was first baron of the kingdom of Naples, and commanded over four hundred towns and villages, it numbered several hundred monks, but in 1843 only twenty. It has a considerable library. Montalembert (l.c. ii. 19) calls Monte Cassino “the most powerful and celebrated monastery in the Catholic universe; celebrated especially because there Benedictwrote his rule and formed the type which was to serve as a model to innumerable communities submitted to that sovereign code.” He also quotes the poetic description from Dante’s \emph{Paradiso.} Dom Luigi Tosti published at Naples, in 1842, a full history of this convent, in three volumes.} the alma mater and capital of his order. Here he labored fourteen years, till his death. Although never ordained to the priesthood, his life there was rather that of a missionary and apostle than of a solitary. He cultivated the soil, fed the poor, healed the sick, preached to the neighboring population, directed the young monks, who in increasing numbers flocked to him, and organized the monastic life upon a fixed method or rule, which he himself conscientiously observed. His power over the hearts, and the veneration in which he was held, is illustrated by the visit of Totila, in 542, the barbarian king, the victor of the Romans and master of Italy, who threw himself on his face before the saint, accepted his reproof and exhortations, asked his blessing, and left a better man, but fell after ten years’ reign, as Benedict had predicted, in a great battle with the Graeco-Roman army under Narses. Benedict died, after partaking of the holy communion, praying, in standing posture, at the foot of the altar, on the 21st of March, 543, and was buried by the side of his sister, Scholastica, who had established, a nunnery near Monte Cassino and died a few weeks before him. They met only once a year, on the side of the mountain, for prayer and pious conversation. On the day of his departure, two monks saw in a vision a shining pathway of stars leading from Monte Cassino to heaven, and heard a voice, that by this road Benedict, the well beloved of God, had ascended to heaven.

His credulous biographer, Pope Gregory I., in the second book of his Dialogues, ascribes to him miraculous prophecies and healings, and even a raising of the dead.\footnote{Gregor. Dial. ii. 37.} With reference to his want of secular culture and his spiritual knowledge, he calls him a learned ignorant and an unlettered sage.\footnote{“Scienter nesciens, et sapienter indoctus.”} At all events he possessed the genius of a lawgiver, and holds the first place among the founders of monastic orders, though his person and life are much less interesting than those of a Bernard of Clairvaux, a Francis of Assisi, and an Ignatius of Loyola.\footnote{Butler, l.c., compares him even with Moses and Elijah. “Being chosen by God, like another Moses, to conduct faithful souls into the true promised land, the kingdom of heaven, he was enriched with eminent supernatural gifts, even those of miracles and prophecy. He seemed, like another Eliseus, endued by God with an extraordinary power, commanding all nature, and, like the ancient prophets, foreseeing future events. He often raised the sinking courage of his monks, and baffled the various artifices of the devil with the sign of the cross, rendered the heaviest stone light, in building his monastery, by a short prayer, and, in presence of a multitude of people, raised to life a novice who had been crushed by the fall of a wall at Monte Cassino.” Montalembert omits the more extraordinary miracles, except the deliverance of Placidus from the whirlpool, which he relates in the language of Bossuet, ii. 15.}

\xsection{The Rule of St. Benedict}

§ 44. The Rule of St. Benedict.

The Regula Benedicti has been frequently edited and annotated, best by Holstenius: Codex reg. Monast. tom. i. p. 111–135; by Dom Marténe: Commentarius in regulam S. Benedicti literalis, moralis, historicus, Par. 1690, in 4to.; by Dom Calmet, Par. 1734, 2 vols.; and by Dom Charles Brandes (Benedictine of Einsiedeln), in 3 vols., Einsiedeln and New York, 1857. Gieseler gives the most important articles in his Ch. H. Bd. i. AbtheiI. 2, § 119. Comp. also Montalembert, l.c. ii. 39 sqq.

The rule of St. Benedict, on which his fame rests, forms an epoch in the history of monasticism. In a short time it superseded all contemporary and older rules of the kind, and became the immortal code of the most illustrious branch of the monastic army, and the basis of the whole Roman Catholic cloister life.\footnote{The Catholic church has recognized three other rules besides that of St. Benedict, viz.: 1. That of St. Basil, which is still retained by the Oriental monks; 2. That of St. Augustine, which is adopted by the regular canons, the order of the preaching brothers or Dominicans, and several military orders; 3. The rule of St. Francis of Assisi, and his mendicant order, in the thirteenth century.} It consists of a preface or prologue, and a series of moral, social, liturgical, and penal ordinances, in seventy-three chapters. It shows a true knowledge of human nature, the practical wisdom of Rome, and adaptation to Western customs; it combines simplicity with completeness, strictness with gentleness, humility with courage, and gives the whole cloister life a fixed unity and compact organization, which, like the episcopate, possessed an unlimited versatility and power of expansion. It made every cloister an ecclesiola in ecclesia, reflecting the relation of the bishop to his charge, the monarchical principle of authority on the democratic basis of the equality of the brethren, though claiming a higher degree of perfection than could be realized in the great secular church. For the rude and undisciplined world of the middle age, the Benedictine rule furnished a wholesome course of training and a constant stimulus to the obedience, self-control, order, and industry which were indispensable to the regeneration and healthy growth of social life.\footnote{Pope Gregory believed the rule of St. Benedicteven to be directly inspired, and Bossuet \emph{(Panégyric de Saint Benoit),} in evident exaggeration, calls it “an epitome of Christianity, a learned and mysterious abridgment of all doctrines of the gospel, all the institutions of the holy fathers, and all the counsels of perfection.” Montalembert speaks in a similar strain of French declamatory eloquence. Monasticism knows very little of the gospel of freedom, and resolves Christianity into a new law of obedience.}

The spirit of the rule may be judged from the following sentences of the prologus, which contains pious exhortations: “Having thus,” he says, “my brethren, asked of the Lord who shall dwell in his tabernacle, we have heard the precepts prescribed to such a one. If we fulfil these conditions, we shall be heirs of the kingdom of heaven. Let us then prepare our hearts and bodies to fight under a holy obedience to these precepts; and if it is not always possible for nature to obey, let us ask the Lord that he would deign to give us the succor of his grace. Would we avoid the pains of hell and attain eternal life, while there is still time, while we are still in this mortal body, and while the light of this life is bestowed upon us for that purpose, let us run and strive so as to reap an eternal reward. We must then form a school of divine servitude, in which, we trust, nothing too heavy or rigorous will be established. But if, in conformity with right and justice, we should exercise a little severity for the amendment of vices or the preservation of charity, beware of fleeing under the impulse of terror from the way of salvation, which cannot but have a hard beginning. When a man has walked for some time in obedience and faith, his heart will expand, and he will run with the unspeakable sweetness of love in the way of God’s commandments. May he grant that, never straying from the instruction of the Master, and persevering in his doctrine in the monastery until death, we may share by patience in the sufferings of Christ, and be worthy to share together his kingdom.”\footnote{We have availed ourselves, in this extract from the preface, of the translation of Montalembert, ii. 44 sq.} The leading provisions of this rule are as follows:

At the head of each society stands an abbot, who is elected by the monks, and, with their consent, appoints a provost (praepositus), and, when the number of the brethren requires, deans over the several divisions (decaniae), as assistants. He governs, in Christ’s stead, by authority and example, and is to his cloister, what the bishop is to his diocese. In the more weighty matters he takes the congregation of the brethren into consultation; in ordinary affairs only the older members. The formal entrance into the cloister must be preceded by a probation of novitiate of one year (subsequently it was made three years), that no one might prematurely or rashly take the solemn step. If the novice repented his resolution, he could leave the cloister without hindrance; if he adhered to it, he was, at the close of his probation, subjected to an examination in presence of the abbot and the monks, and then, appealing to the saints, whose relics were in the cloister, he laid upon the altar of the chapel the irrevocable vow, written or at least subscribed by his own hand, and therewith cut off from himself forever all return to the world.

From this important arrangement the cloister received its stability and the whole monastic institution derived additional earnestness, solidity, and permanence.

The vow was threefold, comprising stabilitas, perpetual adherence to the monastic order; conversio morum, especially voluntary poverty and chastity, which were always regarded as the very essence of monastic piety under all its forms; and obedientia coram Deo et sanctis ejus, absolute obedience to the abbot, as the representative of God and Christ. This obedience is the cardinal virtue of a monk.\footnote{Cap. 5: “Primus humilitatis gradus est obedientia sine mora. Haec convenit iis, qui nihil sibi Christo carius aliquid existimant; propter servitium sanctum, quod professi sunt, seu propter metum gehennae, vel gloriam vitae aeternae, mox ut aliquid imperatum a majore fuerit, ac si divinitus imperetur, moram pati nesciunt in faciendo.”}

The life of the cloister consisted of a judicious alternation of spiritual and bodily exercises. This is the great excellence of the rule of Benedict, who proceeded here upon the true principle, that idleness is the mortal enemy of the soul and the workshop of the devil.\footnote{Cap. 48: “Otiositas inimica est animae; et ideo certis temporibus occupari debent fratres in labore manuum, certis iterum horis in lectione divina.”} Seven hours were to be devoted to prayer, singing of psalms, and meditation;\footnote{\emph{The horaecanonicae} are the \emph{Nocturnae vigiliae, Matutinae, Prima}, \emph{Tertia, Sexta, Nona, Vespera, and Completorium,} and are taken (c. 16) from a literal interpretation of Ps. cxix. 164: “Seven times a day do I praise thee,” and v. 62: “At midnight I will rise to give thanks unto thee.” The Psalter was the liturgy and hymn book of the convent. It was so divided among the seven services of the day, that the whole psalter should be chanted once a week.} from two to three hours, especially on Sunday, to religious reading; and from six to seven hours to manual labor in doors or in the field, or, instead of this, to the training of children, who were committed to the cloister by their parents (oblati).\footnote{Cap. 59: “Si quis forte de nobilibus offert filium suum Deo in monasterio, si ipse puer minori aetate est, parentes ejus faciant petitionem,” etc.}

Here was a starting point for the afterward celebrated cloister schools, and for that attention to literary pursuits, which, though entirely foreign to the uneducated Benedict and his immediate successors, afterward became one of the chief ornaments of his order, and in many cloisters took the place of manual labor.

In other respects the mode of life was to be simple, without extreme rigor, and confined to strictly necessary things. Clothing consisted of a tunic with a black cowl (whence the name: Black Friars); the material to be determined by the climate and season. On the two weekly fast days, and from the middle of September to Easter, one meal was to suffice for the day. Each monk is allowed daily a pound of bread and pulse, and, according to the Italian custom, half a flagon (hemina) of wine; though he is advised to abstain from the wine, if he can do so without injury to his health. Flesh is permitted only to the weak and sick,\footnote{Cap. 40: “Carnium quadrupedum ab omnibus abstinetur comestio, praeter omnino debiles et aegrotos.” Even birds are excluded, which were at that time only delicacies for princes and nobles, as Mabillon shows from the contemporary testimony of Gregory of Tours.} who were to be treated with special care. During the meal some edifying piece was read, and silence enjoined. The individual monk knows no personal property, not even his simple dress as such; and the fruits of his labor go into the common treasury. He should avoid all contact with the world, as dangerous to the soul, and therefore every cloister should be so arranged, as to be able to carry on even the arts and trades necessary for supplying its wants.\footnote{Cap. 66: “Monasterium, si possit fieri, ita debet construi, ut omnia necessaria, id est, aqua, molendinum, hortus, pistrinum, vel artes diversae intra monasterium exerceantur, ut non sit necessitas monachis vagandi foras, quia omnino non expedit animabus eorum.”} Hospitality and other works of love are especially commended.

The penalties for transgression of the rule are, first, private admonition, then exclusion from the fellowship of prayer, next exclusion from fraternal intercourse, and finally expulsion from the cloister, after which, however, restoration is possible, even to the third time.

\xsection{The Benedictines. Cassiodorus}

§ 45. The Benedictines. Cassiodorus.

Benedict had no presentiment of the vast historical importance, which this rule, originally designed simply for the cloister of Monte Cassino, was destined to attain. He probably never aspired beyond the regeneration and salvation of his own soul and that of his brother monks, and all the talk of later Catholic historians about his far-reaching plans of a political and social regeneration of Europe, and the preservation and promotion of literature and art, find no support whatever in his life or in his rule. But he humbly planted a seed, which Providence blessed a hundredfold. By his rule he became, without his own will or knowledge, the founder of an order, which, until in the thirteenth century the Dominicans and Franciscans pressed it partially into the background, spread with great rapidity over the whole of Europe, maintained a clear supremacy, formed the model for all other monastic orders, and gave to the Catholic church an imposing array of missionaries, authors, artists, bishops, archbishops, cardinals, and popes, as Gregory the Great and Gregory VII. In less than a century after the death of Benedict, the conquests of the barbarians in Italy, Gaul, Spain were reconquered for civilization, and the vast territories of Great Britain, Germany, and Scandinavia incorporated into Christendom, or opened to missionary labor; and in this progress of history the monastic institution, regulated and organized by Benedict’s rule, bears an honorable share.

Benedict himself established a second cloister in the vicinity of Terracina, and two of his favorite disciples, Placidus and St. Maurus,\footnote{This Maurus, the founder of the abbacy of Glanfeuil (St. Maur sur Loire), is the patron saint of a branch of the Benedictines, the celebrated Maurians in France (dating from 1618), who so highly distinguished themselves in the seventeenth and early part of the eighteenth centuries, by their thorough archaeological and historical researches, and their superior editions of the Fathers. The most eminent of the Maurians are D. (Dom, equivalent to Domnus, Sir) Menard, d’Achery, Godin, Mabillon, le Nourry, Martianay, Ruinart, Martene, Montfaucon, Massuet, Garnier, and de la Rue, and in our time Dom Pitra, editor of a valuable collection of patristic fragments, at the cloister of Solesme.} introduced the “holy rule,” the one into Sicily, the other into France. Pope Gregory the Great, himself at one time a Benedictine monk, enhanced its prestige, and converted the Anglo-Saxons to the Roman Christian faith, by Benedictine monks. Gradually the rule found so general acceptance both in old and in new institutions, that in the time of Charlemagne it became a question, whether there were any monks at all, who were not Benedictines. The order, it is true, has degenerated from time to time, through the increase of its wealth and the decay of its discipline, but its fostering care of religion, of humane studies, and of the general civilization of Europe, from the tilling of the soil to the noblest learning, has given it an honorable place in history and won immortal praise. He who is familiar with the imposing and venerable tomes of the Benedictine editions of the Fathers, their thoroughly learned prefaces, biographies, antiquarian dissertations, and indexes, can never think of the order of the Benedictines without sincere regard and gratitude.

The patronage of learning, however, as we have already said, was not within the design of the founder or his rule. The joining of this to the cloister life is duel if we leave out of view the learned monk Jerome, to Cassiodorus, who in 538 retired from the honors and cares of high civil office, in the Gothic monarchy of Italy,\footnote{He was the last of the Roman consuls—an office which Justinian abolished—and was successively the minister of Odoacer, Theodoric, and Athalaric, who made him prefect of the praetorium} to a monastery founded by himself at Vivarium\footnote{Or \emph{Vivaria,} so called from the numerous vivaria or fish ponds in that region.} (Viviers), in Calabria in Lower Italy. Here he spent nearly thirty years as monk and abbot, collected a large library, encouraged the monks to copy and to study the Holy Scriptures, the works of the church fathers, and even the ancient classics, and wrote for them several literary and theological text-books, especially his treatise De institutione divinarum literarum, a kind of elementary encyclopaedia, which was the code of monastic education for many generations. Vivarium at one time almost rivalled Monte Cassino, and Cassiodorus won the honorary title of the restorer of knowledge in the sixth century.\footnote{Comp. Mabillon, Ann. Bened. l. v. c. 24, 27; F. de Ste. Marthe, Vie de Cassiodore, 1684.}

The Benedictines, already accustomed to regular work, soon followed this example. Thus that very mode of life, which in its founder, Anthony, despised all learning, became in the course of its development an asylum of culture in the rough and stormy times of the migration and the crusades, and a conservator of the literary treasures of antiquity for the use of modern times.

\xsection{Opposition to Monasticism. Jovinian}

§ 46. Opposition to Monasticism. Jovinian.

I. Chrysostomus: Πρὸς τοὺς πολεμοῦτας τοῖς ἐπὶ τὸ μονάζειν ἐνάγουσιν(a vindication of monasticism against its opponents, in three books). Hieronymus: Ep. 61, ad Vigilantium (ed. Vallars. tom. i. p. 345 sqq.); Ep. 109, ad Riparium (i. 719 sqq.); Adv. Helvidium (a.d. 383); Adv. Jovinianum (a.d. 392); Adv. Vigilantium (a.d. 406). All these three tracts are in Opera Hieron. tom. ii. p. 206–402. Augustinus: De haeres. cap. 82 (on Jovinian), and c. 84 (on Helvidius and the Helvidians). Epiphanius: Haeres. 75 (on Aerius).

II. Chr. W. F. Walch: Ketzerhistorie (1766), part iii. p. 585 (on Helvidius and the Antidikomarianites); p. 635 sqq. (on Jovinian); and p. 673 sqq. (on Vigilantius). Vogel: De Vigilantio haeretico orthodoxo, Gött. 1756. G. B. Lindner: De Joviniano et Vigilantio purioris doctrinae antesignanis, Lips. 1839. W. S. Gilly: Vigilantius and his Times, Lond. 1844. Comp. also Neander: Der heil. Joh. Chrysostomus, 3d ed. 1848, vol. i. p. 53 sqq.; and Kirchengesch, iii. p. 508 sqq. (Torrey’s translation, ii. p. 265 sqq.). Baur: Die christliche Kirche von 4–6ten Jahrh. 1859, p. 311 sqq.

Although monasticism was a mighty movement of the age, engaging either the cooperation or the admiration of the whole church, yet it was not exempt from opposition. And opposition sprang from very different quarters: now from zealous defenders of heathenism, like Julian and Libanius, who hated and bitterly reviled the monks for their fanatical opposition to temples and idol-worship; now from Christian statesmen and emperors, like Valens, who were enlisted against it by its withdrawing so much force from the civil and military service of the state, and, in the time of peril from the barbarians, encouraging idleness and passive contemplation instead of active, heroic virtue; now from friends of worldly indulgence, who found themselves unpleasantly disturbed and rebuked by the religious earnestness and zeal of the ascetic life; lastly, however, also from a liberal, almost protestant, conception of Christian morality, which set itself at the same time against the worship of Mary and the saints, and other abuses. This last form of opposition, however, existed mostly in isolated cases, was rather negative than positive in its character, lacked the spirit of wisdom and moderation, and hence almost entirely disappeared in the fifth century, only to be revived long after, in more mature and comprehensive form, when monasticism had fulfilled its mission for the world.

To this class of opponents belong Helvidius, Jovinian, Vigilantius, and Aerius. The first three are known to us through the passionate replies of Jerome, the last through the Panarion of Epiphanius. They figure in Catholic church history among the heretics, while they have received from many Protestant historians a place among the “witnesses of the truth” and the forerunners of the Reformation.

We begin with Jovinian, the most important among them, who is sometimes compared, for instance, even by Neander, to Luther, because, like Luther, he was carried by his own experience into reaction against the ascetic tendency and the doctrines connected with it. He wrote in Rome, before the year 390 a work, now lost, attacking monasticism in its ethical principles. He was at that time himself a monk, and probably remained so in a free way until his death. At all events he never married, and according to Augustine’s account, he abstained “for the present distress,”\footnote{1 Cor. vii. 26.} and from aversion to the encumbrances of the married state. Jerome pressed him with the alternative of marrying and proving the equality of celibacy with married life, or giving up his opposition to his own condition.\footnote{Adv. Jovin. lib. i. c. 40 (Opera, ii. 304): “Et tamen iste formosus monachus, crassus, nitidus, dealbatus, et quasi sponsus semper incedens, aut uxorem ducat ut aequalem virginitatem nuptiis probet; aut, si non duxerit, frustra contra nos verbis agit, cum opere nobiscum sit.”} Jerome gives a very unfavorable picture of his character, evidently colored by vehement bitterness. He calls Jovinian a servant of corruption, a barbarous writer, a Christian Epicurean, who, after having once lived in strict asceticism, now preferred earth to heaven, vice to virtue, his belly to Christ, and always strode along as an elegantly dressed bridegroom. Augustine is much more lenient, only reproaching Jovinian with having misled many Roman nuns into marriage by holding before them the examples of pious women in the Bible. Jovinian was probably provoked to question and oppose monasticism, as Gieseler supposes, by Jerome’s extravagant praising of it, and by the feeling against it, which the death of Blesilla (384) in Rome confirmed. And he at first found extensive sympathy. But he was excommunicated and banished with his adherents at a council about the year 390, by Siricius, bishop of Rome, who was zealously opposed to the marriage of priests. He then betook himself to Milan, where the two monks Sarmatio and Barbatian held forth views like his own; but he was treated there after the same fashion by the bishop, Ambrose, who held a council against him. From this time he and his party disappear from history, and before the year 406 he died in exile.\footnote{Augustinesays, De haer. c. 82: “Cito ista haeresis oppressa et extincta est;” and Jeromewrites of Jovinian, in 406, Adv. Vigilant. c. 1, that, after having been condemned by the authority of the Roman church, he dissipated his mind in the enjoyment of his lusts.}

According to Jerome, Jovinian held these four points (1) Virgins, widows, and married persons, who have once been baptized into Christ, have equal merit, other things in their conduct being equal. (2) Those, who are once with full faith born again by baptism, cannot be overcome (subverti) by the devil. (3) There is no difference between abstaining from food and enjoying it with thanksgiving. (4) All, who keep the baptismal covenant, will receive an equal reward in heaven.

He insisted chiefly on the first point; so that Jerome devotes the whole first book of his refutation to this point, while he disposes of all the other heads in the second. In favor of the moral equality of married and single life, he appealed to Gen. ii. 24, where God himself institutes marriage before the fall; to Matt. xix. 5, where Christ sanctions it; to the patriarchs before and after the flood; to Moses and the prophets, Zacharias and Elizabeth, and the apostles, particularly Peter, who lived in wedlock; also to Paul, who himself exhorted to marriage,\footnote{1 Cor. vii. 36, 39.} required the bishop or the deacon to be the husband of one wife,\footnote{1 Tim. iii. 2, 12.} and advised young widows to marry and bear children.\footnote{1 Tim. v. 14; comp. 1 Tim. ii. 15; Heb. xiii. 4.} He declared the prohibition of marriage and of divinely provided food a Manichaean error. To answer these arguments, Jerome indulges in utterly unwarranted inferences, and speaks of marriage in a tone of contempt, which gave offence even to his friends.\footnote{From 1 Cor. vii. 1, for example (“It is good for a man not to touch a woman”), he argues, without qualification, l. i. c. 7 (Opera, ii. 246): “Si bonum est mulierem non tangere, \emph{malum} \emph{est ergo tangere,} nihil enim bono contrarium est, nisi malum; si autem malum est, et ignoscitur, ideo conceditur, ne malo quid deterius fiat .... Tolle fornicationem, et non dicet [apostolus], \emph{unusquisque} \emph{uxorem suam habeat.}“Immediately after this (ii. 247) he argues, from the exhortation of Paul to pray without ceasing, 1 Thess. v. 17: “Si semper orandum est, nunquam ergo conjugio serviendum, quoniam quotiescunque uxori debitum reddo, orare non possum.” Such sophistries and misinterpretations evidently proceed upon the lowest sensual idea of marriage, and called forth some opposition even at that age. He himself afterward felt that he had gone too far, and in his Ep. 48 (ed. Vallars. or Ep. 30, ed. Bened.) ad Pammachium, endeavored to save himself by distinguishing between the gymnastic (polemically rhetorical) and the dogmatic mode of writing.} Augustine was moved by it to present the advantages of the married life in a special work, De bono conjugali, though without yielding the ascetic estimate of celibacy.\footnote{De bono conj. c. 8: “Duo bona sunt connubium et continentia, quorum alterum est melius.”}

Jovinian’s second point has an apparent affinity with the Augustinian and Calvinistic doctrine of the perseverantia sanctorum. It is not referred by him, however, to the eternal and unchangeable counsel of God, but simply based on 1 Jno. iii. 9, and v. 18, and is connected with his abstract conception of the opposite moral states. He limits the impossibility of relapse to the truly regenerate, who “plena fide in baptismate renati sunt,” and makes a distinction between the mere baptism of water and the baptism of the Spirit, which involves also a distinction between the actual and the ideal church.

His third point is aimed against the ascetic exaltation of fasting, with reference to Rom. xiv. 20, and 1 Tim. iv. 3. God, he holds, has created all animals for the service of man; Christ attended the marriage feast at Cana as a guest, sat at table with Zaccheus, with publicans and sinners, and was called by the Pharisees a glutton and a wine-bibber; and the apostle says: To the pure all things are pure, and nothing to be refused, if it be received with thanksgiving.

He went still further, however, and, with the Stoics, denied all gradations of moral merit and demerit, consequently also all gradations of reward and punishment. He overlooked the process of development in both good and evil. He went back of all outward relations to the inner mind, and lost all subordinate differences of degree in the great contrast between true Christians and men of the world, between regenerate and unregenerate; whereas, the friends of monasticism taught a higher and lower morality, and distinguished the ascetics, as a special class, from the mass of ordinary Christians. As Christ, says he, dwells in believers, without difference of degree, so also believers are in Christ without difference of degree or stages of development. There are only two classes of men, righteous and wicked, sheep and goats, five wise virgins and five foolish, good trees with good fruit and bad trees with bad fruit. He appealed also to the parable of the laborers in the vineyard, who all received equal wages. Jerome answered him with such things as the parable of the sower and the different kinds of ground, the parable of the different numbers of talents with corresponding rewards, the many mansions in the Father’s house (by which Jovinian singularly understood the different churches on earth), the comparison of the resurrection bodies with the stars, which differ in glory, and the passage: “He which soweth sparingly, shall reap also sparingly; and he which soweth bountifully, shall reap also bountifully.”\footnote{2 Cor. ix. 6.}

\xsection{Helvidius, Vigilantius, and Aerius}

§ 47. Helvidius, Vigilantius, and Aerius.

See especially the tracts of Jerome quoted in the preceding section.

Helvidius, whether a layman or a priest at Rome it is uncertain, a pupil, according to the statement of Gennadius, of the Arian bishop Auxentius of Milan, wrote a work, before the year 383, in refutation of the perpetual virginity of the mother of the Lord—a leading point with the current glorification of celibacy. He considered the married state equal in honor and glory to that of virginity. Of his fortunes we know nothing. Augustine speaks of Helvidians, who are probably identical with the Antidicomarianites of Epiphanius. Jerome calls Helvidius, indeed, a rough and uneducated man,\footnote{At the very beginning of his work against him, he styles him “hominem rusticum et vix primis quoque imbutum literis.”} but proves by quotations of his arguments, that he had at least some knowledge of the Scriptures, and a certain ingenuity. He appealed in the first place to Matt. i. 18, 24, 25, as implying that Joseph knew his wife not before, but after, the birth of the Lord; then to the designation of Jesus as the “first born” son of Mary, in Matt. i. 25, and Luke ii. 7; then to the many passages, which speak of the brothers and sisters of Jesus; and finally to the authority of Tertullian and Victorinus. Jerome replies, that the “till” by no means always fixes a point after which any action must begin or cease;\footnote{Comp. Matt. xxviii. 20.} that, according to Ex. xxxiv. 19, 20; Num. xviii. 15 sqq., the “first born” does not necessarily imply the birth of other children afterward, but denotes every one, who first opens the womb; that the “brothers” of Jesus may have been either sons of Joseph by a former marriage, or, according to the wide Hebrew use of the term, cousins; and that the authorities cited were more than balanced by the testimony of Ignatius, Polycarp(?), and Irenaeus. “Had Helvidius read these,” says he, “he would doubtless have produced something more skilful.”

This whole question, it is well known, is still a problem in exegesis. The perpetua virginitas of Mary has less support from Scripture than the opposite theory. But it is so essential to the whole ascetic system, that it became from this time an article of the Catholic faith, and the denial of it was anathematized as blasphemous heresy. A considerable number of Protestant divines,\footnote{Luther, for instance (who even calls Helvidiusa “gross fool”), and Zuingle, among the Reformers; Olshausen and J. P. Lange, among the later theologians.} however, agree on this point with the Catholic doctrine, and think it incompatible with the dignity of Mary, that, after the birth of the Son of God and Saviour of the world, she should have borne ordinary children of men.

Vigilantius, originally from Gaul,\footnote{Respecting his descent, compare the diffuse treatise of the tedious but thorough Walch, l.c. p. 675-677.} a presbyter of Barcelona in Spain, a man of pious but vehement zeal, and of literary talent, wrote in the beginning of the fifth century against the ascetic spirit of the age and the superstition connected with it. Jerome’s reply, dictated hastily in a single night at Bethlehem in the year 406, contains more of personal abuse and low witticism, than of solid argument. “There have been,” he says, “monsters on earth, centaurs, syrens, leviathans, behemoths .... Gaul alone has bred no monsters, but has ever abounded in brave and noble men,—when, of a sudden, there has arisen one Vigilantius, who should rather be called Dormitantius,\footnote{This cheap pun he repeats, Epist. 109, ad Ripar. (Opera, i. p. 719), where he says that Vigilantius(Wakeful) was so called \textgreek{κατ  ἀντίφρασιν}, and should rather be called \emph{Dormitantius} (Sleepy). The fact is, that Vigilantiuswas wide-awake to a sense of certain superstitions of the age} contending in an impure spirit against the Spirit of Christ, and forbidding to honor the graves of the martyrs; he rejects the Vigils—only at Easter should we sing hallelujah; he declares abstemiousness to be heresy, and chastity a nursery of licentiousness (pudicitiam, libidinis seminarium) .... This innkeeper of Calagurris\footnote{In South Gaul; now Casères in Gascogne. As the business of innkeeper is incompatible with the spiritual office, it has been supposed that the father of Vigilantiuswas a \emph{caupo Calagurritanus.} Comp. Rössler’s Bibliothek der Kirchenväter, part ix. p. 880 sq., note 100; and Walch, l.c} mingles water with the wine, and would, according to ancient art, combine his poison with the genuine faith. He opposes virginity, hates chastity, cries against the fastings of the saints, and would only amidst jovial feastings amuse himself with the Psalms of David. It is terrible to bear, that even bishops are companions of his wantonness, if those deserve this name, who ordain only married persons deacons, and trust not the chastity of the single.”\footnote{Adv. Vigil.c. 1 and 2 (Opera, tom. ii. p. 387 sqq.).} Vigilantius thinks it better for a man to use his money wisely, and apply it gradually to benevolent objects at home, than to lavish it all at once upon the poor or give it to the monks of Jerusalem. He went further, however, than his two predecessors, and bent his main efforts against the worship of saints and relics, which was then gaining ascendency and was fostered by monasticism. He considered it superstition and idolatry. He called the Christians, who worshipped the “wretched bones” of dead men, ash-gatherers and idolaters.\footnote{“Cinerarios et idolatras, qui mortuorum ossa venerantur.” Hieron. Ep. 109, ad Riparium (tom. i. p. 719).} He expressed himself sceptically respecting the miracles of the martyrs, contested the practice of invoking them and of intercession for the dead, as useless, and declared himself against the Vigils, or public worship in the night, as tending to disorder and licentiousness. This last point Jerome admits as a fact, but not as an argument, because the abuse should not abolish the right use.

The presbyter Aerius of Sebaste, about 360, belongs also among the partial opponents of monasticism. For, though himself an ascetic, he contended against the fast laws and the injunction of fasts at certain times, considering them an encroachment upon Christian freedom. Epiphanius also ascribes to him three other heretical views: denial of the superiority of bishops to presbyters, opposition to the usual Easter festival, and opposition to prayers for the dead.\footnote{Epiph. Haer. 75. Comp. also Walch, l.c. iii. 321-338. Bellarmine, on account of this external resemblance, styles Protestantism the Aerian heresy.} He was hotly persecuted by the hierarchy, and was obliged to live, with his adherents, in open fields and in caves.

\xchapter{The Hierarchy and Polity of the Church}

THE HIERARCHY AND POLITY OF THE CHURCH.

Comp. in part the literature in vol. i. § 105 and 110 (to which should be added now, P. A. de Lagarde: Constitutiones Apostolorum, Lips. and Lond., 1862); also Gibbon, ch. xx.; Milman: Hist. of Ancient Christianity, book iv. c. 1 (Amer. ed. p. 438 sqq.), and the corresponding sections in Bingham, Schroeckh, Plank, Neander, Gieseler, Baur, etc. (see the particular literature below).

\xsection{Schools of the Clergy}

§ 48. Schools of the Clergy.

Having in a former section observed the elevation of the church to the position of the state religion of the Roman empire, and the influence of this great change upon the condition of the clergy and upon public morality, we turn now to the internal organization and the development of the hierarchy under its new circumstances. The step of progress which we here find distinguishing the organization of this third period from the episcopal system of the second and the apostolic supervision of the first, is the rise of the patriarchal constitution and of the system of ecumenical councils closely connected with it. But we must first glance at the character and influence of the teaching order in general.

The work of preparation for the clerical office was, on the one hand, materially facilitated by the union of the church with the state, putting her in possession of the treasures, the schools, the learning, and the literature of classic heathendom, and throwing the education of the rising generation into her hands. The numerous doctrinal controversies kept the spirit of investigation awake, and among the fathers and bishops of the fourth and fifth centuries we meet with the greatest theologians of the ancient church. These gave their weighty voices for the great value of a thorough education to the clerical office, and imparted much wholesome instruction respecting the studies proper to this purpose.\footnote{E.g. Chrysostom: \emph{De sacerdotio;} Augustine: \emph{De doctrina Christiana}; Jerome: in several letters; Gregory the Great: \emph{Regula pastoralis}.} The African church, by a decree of the council of Carthage, in 397, required of candidates a trial of their knowledge and orthodoxy. A law of Justinian, of the year 541, established a similar test in the East.

But on the other hand, a regular and general system of clerical education was still entirely wanting. The steady decay of the classic literature, the gradual cessation of philosophical and artistic production, the growth of monastic prejudice against secular learning and culture, the great want of ministers in the suddenly expanded field of the church, the uneasy state of the empire, and the barbarian invasions, were so many hinderances to thorough theological preparation. Many candidates trusted to the magical virtue of ordination. Others, without inward call, were attracted to the holy office by the wealth and power of the church. Others had no time or opportunity for preparation, and passed, at the instance of the popular voice or of circumstances, immediately from the service of the state to that of the church, even to the episcopal office; though several councils prescribed a previous test of their capacity in the lower degrees of reader, deacon, and presbyter. Often, however, this irregularity turned to the advantage of the church, and gave her a highly gifted man, like Ambrose, whom the acclamation of the people called to the episcopal see of Milan even before he was baptized. Gregory Nazianzen laments that many priests and bishops came in fresh from the counting house, sunburnt from the plow, from the oar, from the army, or even from the theatre, so that the most holy order of all was in danger of becoming the most ridiculous. “Only he can be a physician,” says he, “who knows the nature of diseases; he, a painter, who has gone through much practice in mixing colors and in drawing forms; but a clergyman may be found with perfect ease, not thoroughly wrought, of course, but fresh made, sown and full blown in a moment, as the legend says of the giants.\footnote{\textgreek{Ὡς ὁ μύθος ποιεῖ τοὺς γιγάντας}.} We form the saints in a day, and enjoin them to be wise, though they possess no wisdom at all, and bring nothing to their spiritual office, except at best a good will.”\footnote{Greg. Orat. xliii. c. 26 (Opera omnia, ed. Bened., Paris, 1842, tom. i. p. 791 sq.), and similar passages in his other orations, and his Carmen de se ipse et advers. Episc. Comp. Ullmann: Greg. v. Naz. p. 511 sqq.} If such complaints were raised so early as the end of the Nicene age, while the theological activity of the Greek church was in its bloom, there was far more reason for them after the middle of the fifth century and in the sixth, especially in the Latin church, where, even among the most eminent clergymen, a knowledge of the original languages of the Holy Scriptures was a rare exception.

The opportunities which this period offered for literary and theological preparation for the ministry, were the following:

1. The East had four or five theological schools, which, however, were far from supplying its wants.

The oldest and most celebrated was the catechetical school of Alexandria. Favored by the great literary treasures, the extensive commercial relations, and the ecclesiastical importance of the Egyptian metropolis, as well as by a succession of distinguished teachers, it flourished from the middle of the second century to the end of the fourth, when, amidst the Origenistic, Nestorian, and Monophysite confusion, it withered and died. Its last ornament was the blind, but learned and pious Didymus (340–395).

From the Alexandrian school proceeded the smaller institution of Caesarea in Palestine, which was founded by Origen, after his banishment from Alexandria, and received a new but temporary impulse in the beginning of the fourth century from his admirer, the presbyter Pamphilus, and from his friend Eusebius. It possessed the theological library which Eusebius used in the preparation of his learned works.

Far more important was the theological school of Antioch, founded about 290 by the presbyters Dorotheus and Lucian. It developed in the course of the fourth century a severe grammatico-historical exegesis, counter to the Origenistic allegorical method of the Alexandrians; now in connection with the church doctrine, as in Chrysostom; now in a rationalizing spirit, as in Theodore of Mopsuestia and Nestorius.

The seminary at Edessa, a daughter of the Antiochian school, was started by the learned deacon, Ephraim Syrus († 378), furnished ministers for Mesopotamia and Persia, and stood for about a hundred years.

The Nestorians, at the close of the fifth century, founded a seminary at Nisibis in Mesopotamia, which was organized into several classes and based upon a definite plan of instruction.

The West had no such institutions for theological instruction, but supplied itself chiefly from cloisters and private schools of the bishops. Cassiodorus endeavored to engage Pope Agapetus in founding a learned institution in Rome, but was discouraged by the warlike disquietude of Italy. Jerome spent some time at the Alexandrian school under the direction of Didymus.

2. Many priests and bishops, as we have already observed, emanated from the monasteries, where they enjoyed the advantages of retirement from the world, undisturbed meditation, the intercourse of kindred earnest minds, and a large spiritual experience; but, on the other hand, easily sank into a monkish narrowness, and rarely attained that social culture and comprehensive knowledge of the world and of men, which is necessary, especially in large cities, for a wide field of labor.

3. In the West there were smaller diocesan seminaries, under the direction of the bishops, who trained their own clergy, both in theory and in practice, as they passed through the subordinate classes of reader, sub-deacon, and deacon.

Augustine set a good example of this sort, having at Hippo a “monasterium clericorum,” which sent forth many good presbyters and bishops for the various dioceses of North Africa. Similar clerical monasteries or episcopal seminaries arose gradually in the southern countries of Europe, and are very common in the Roman Catholic church to this day.

4. Several of the most learned and able fathers of the fourth century received their general scientific education in heathen schools, under the setting sun of the classic culture, and then studied theology either in ascetic retirement or under some distinguished church teacher, or by the private reading of the Scriptures and the earlier church literature.

Thus Basil the Great and Gregory Nazianzen were in the high school of Athens at the same time with the prince Julian the Apostate; Chrysostom attended the lectures of the celebrated rhetorician Libanius in Antioch; Augustine studied at Carthage, Rome, and Milan; and Jerome was introduced to the study of the classics by the grammarian Donatus of Rome. The great and invaluable service of these fathers in the development and defence of the church doctrine, in pulpit eloquence, and especially in the translation and exposition of the Holy Scriptures, is the best evidence of the high value of a classical education. And the church has always, with good reason, acknowledged it.

\xsection{Clergy and Laity. Elections}

§ 49. Clergy and Laity. Elections.

The clergy, according to the precedent of the Old Testament, came to be more and more rigidly distinguished, as a peculiar order, from the body of the laity. The ordination, which was solemnized by the laying on of hands and prayer, with the addition at a later period of an anointing with oil and balsam, marked the formal entrance into the special priesthood, as baptism initiated into the universal priesthood; and, like baptism, it bore an indefeasible character (character indelebilis). By degrees the priestly office assumed the additional distinction of celibacy and of external marks, such as tonsure, and sacerdotal vestments worn at first only during official service, then in every-day life. The idea of the universal priesthood of believers retreated in proportion, though it never passed entirely out of sight, but was from time to time asserted even in this age. Augustine, for example, says, that as all are called Christians on account of their baptism, so all believers are priests, because they are members of the one High Priest.\footnote{De Civit. Dei, lib. xx. cap. 10: ”\emph{Erunt sacerdotes Dei et Christi et regnabunt cum eo mille annos (Apoc. xx. 6): non utique de solis episcopis et presbyteris dictum est, qui proprie jam vocantur in Ecclesia sacerdotes; sed sicut omnes Christianos dicimus propter mysticum chrisma, sic omnes sacerdotes, quoniam membra sunt unius sacerdotis. De quibus apostolus Petrus: Plebs, inquit, sancta regale sacerdotium (1 Pet. ii. 9)}.” Comp. Ambrosiaster ad Eph. iv. 11; Jerome ad Tit. i. 7 and Pope Leo I., Sermon. iv. 1.}

The progress of the hierarchical principle also encroached gradually upon the rights of the people in the election of their pastors.\footnote{According to Clemens Romanus, ad Corinth. c. 44, the consent of the whole congregation in the choice of their officers was the apostolic and post-apostolic custom; and the Epistles of Cyprian, especially Ep. 68, show that the same rule continued in the middle of the third century. Comp. vol. i. § 105.} But in this period it did not as yet entirely suppress them. The lower clergy were chosen by the bishops, the bishops by their colleagues in the province and by the clergy. The fourth canon of Nice, probably at the instance of the Meletian schism, directed that a bishop should be instituted and consecrated by all, or at least by three, of the bishops of the province. This was not aimed, however, against the rights of the people, but against election by only one bishop—the act of Meletius. For the consent of the people in the choice of presbyters, and especially of bishops, long remained, at least in outward form, in memory of the custom of the apostles and the primitive church. There was either a formal vote,\footnote{\textgreek{Ζήτησις, ψήφισμα, ψῆγος}, scrutinium.} particularly when there were three or more candidates before the people, or the people were thrice required to signify their confirmation or rejection by the formula: “Worthy,” or “unworthy.”\footnote{\textgreek{Ἄξιος}, dignus, or \textgreek{ἀνάξιος}, indignus. Constitut. Apost. viii. 4; Concil. Aurelat. ii. (A. D. 452) c. 54; Gregor. Naz. Orat. xxi. According to a letter of Peter of Alexandria, in Theodor. Hist. Eccl. iv. 22, the bishop in the East was elected\textgreek{ἐπισκόπων συνόδῳ, ψήφῳ κληρικῶν, αἰτήσει λαῶν}. He himself was elected archbishop of Alexandria and successor of Athanasius (a.d.373), according to the desire of the latter, “by the unanimous consent of the clergy and of the chief men of the city” (iv. cap. 20), and, after his expulsion, he objected to his wicked successor Lucius, among other things, that “he had purchased the episcopal office with gold, as though it had been a secular dignity, ... and had not been elected by a \emph{synod of bishops, by the votes of the clergy, or by the request of the people, according to the regulations of the church}“ (iv. c. 22).} The influence of the people in this period appears most prominently in the election of bishops. The Roman bishop Leo, in spite of his papal absolutism, asserted the thoroughly democratic principle, long since abandoned by his successors: “He who is to preside over all, should be elected by all.”\footnote{Epist. x. c. 4 (opera, ed. Baller. i. 637): “Expectarentur certe vota civium, testimonia populorum, quaereretur honoratorum arbitrium, electio clericorum .... In the same epistle, cap. 6: \emph{Qui praefuturus est omnibus, ab} \emph{omnibus eligatur.”}} Oftentimes the popular will decided before the provincial bishops and the clergy assembled and the regular election could be held. Ambrose of Milan and Nectarius of Constantinople were appointed to the bishopric even before they were baptized; the former by the people, the latter by the emperor Theodosius; though in palpable violation of the eightieth apostolic canon and the second Nicene.\footnote{Paulinus, Vita Ambros.; Sozomen, H. E. l. iv. c. 24, and vii. 8. This historian excuses the irregularity by a special interposition of Providence.} Martin of Tours owed his elevation likewise to the popular voice, while some bishops objected to it on account of his small and wasted form.\footnote{Sulpitius Severus, Vita Mart. c. 7: “Incredibilis multitudo non solum ex eo oppido [Tours], sed etiam ex vicinis urbibus \emph{ad suffragia ferenda} convenerat,” etc.} Chrysostom was called from Antioch to Constantinople by the emperor Arcadius, in consequence of a unanimous vote of the clergy and people.\footnote{Socrates, H. E. vi. 2:\textgreek{Ψηφίσματι κοινῷ ὁμοῦ πάντων κλήρου τε φημὶ καὶ λαοῦ}..} Sometimes the people acted under outside considerations and the management of demagogues, and demanded unworthy or ignorant men for the highest offices. Thus there were frequent disturbances and collisions, and even bloody conflicts, as in the election of Damasus in Rome. In short, all the selfish passions and corrupting influences, which had spoiled the freedom of the popular political elections in the Grecian and Roman republics, and which appear also in the republics of modern times, intruded upon the elections of the church. And the clergy likewise often suffered themselves to be guided by impure motives. Chrysostom laments that presbyters, in the choice of a bishop, instead of looking only at spiritual fitness, were led by regard for noble birth, or great wealth, or consanguinity and friendship.\footnote{De sacerdotio, lib. iii. c. 15. Further on in the same chapter he says even, that many are elected on account of their badness, to prevent the mischief they would otherwise do: \textgreek{Οἱ δὲ, διὰ, πονηρίαν, [εἰς τὴν τοῦ κλήρου καταλέγονται τάξιν, καὶ ἵνα μὴ, παροφθέντες , μεγάλα ἐργάσωνται κακά}. Quite parallel is the testimony of Gregory Nazianzen in his Carmen,\textgreek{εἰς ἑαυτὸν καὶ περὶ ἐπισκόπων}\emph{,} or De se ipso et de episcopis, ver. 330 sqq. (Opera, ed. Bened. Par. tom. ii. p. 796), and elsewhere.} The bishops themselves sometimes did no better. Nectarius, who was suddenly transferred, in 381, by the emperor Theodosius, from the praetorship to the bishopric of Constantinople, even before he was baptized,\footnote{Sozomenus, Hist. Eccl. vii. c. 8. Sozomen sees in this election a special interposition of God.} wished to ordain his physician Martyrius deacon, and when the latter refused, on the ground of incapacity, he replied: “Did not I, who am now a priest, formerly live much more immorally than thou, as thou thyself well knowest, since thou wast often an accomplice of my many iniquities?” Martyrius, however, persisted in his refusal, because he had continued to live in sin long after his baptism, while Nectarius had become a new man since his.\footnote{Sozomenus, vii. c. 10. Otherwise he, as well as Socrates, H. E. v. c. 8, and Theodoret, H. E. v. c. 8, speaks very favorably of the character of Nectarius.}

The emperor also, after the middle of the fourth century, exercised a decisive influence in the election of metropolitans and patriarchs, and often abused it in a despotic and arbitrary way.

Thus every mode of appointment was evidently exposed to abuse, and could furnish no security against unworthy candidates, if the electors, whoever they might be, were destitute of moral earnestness and the gift of spiritual discernment.

Toward the end of the period before us the republican element in the election of bishops entirely disappeared. The Greek church after the eighth century vested the franchise exclusively in the bishops.\footnote{The seventh ecumenical council, at Nice, 787, in its third canon, on the basis of a wrong interpretation of the fourth canon of the first council of Nice, expressly prohibited the people and the secular power from any share in the election of bishops. Also the eighth general council prescribes that the bishop should be chosen only by the college of bishops.} The Latin church, after the eleventh century, vested it in the clergy of the cathedral church, without allowing any participation to the people. But in the West, especially in Spain and France, instead of the people, the temporal prince exerted an important influence, in spite of the frequent protest of the church.

Even the election of pope, after the downfall of the West Roman empire, came largely under control of the secular authorities of Rome; first, of the Ostrogothic kings; then, of the exarchs of Ravenna in the name of the Byzantine emperor; and, after Charlemagne, of the emperor of Germany; till, in 1059, through the influence of Hildebrand (afterward Gregory VII.), it was lodged exclusively with the college of cardinals, which was filled by the pope himself. Yet the papal absolutism of the middle age, like the modern Napoleonic military despotism in the state, found it well, under favorable prospects, to enlist the democratic principle for the advancement of its own interests.

\xsection{Marriage and Celibacy of the Clergy}

§ 50. Marriage and Celibacy of the Clergy.

The progress and influence of monasticism, the general exaltation of the ascetic life above the social, and of celibacy above the married state, together with the increasing sharpness of the distinction between clergy and laity, all tended powerfully toward the celibacy of the clergy. What the apostle Paul, expressly discriminating a divine command from a human counsel, left to each one’s choice, and advised, in view of the oppressed condition of the Christians in the apostolic age, as a safer and less anxious state only for those who felt called to it by a special gift of grace, now, though the stress of circumstances was past, was made, at least in the Latin church, an inexorable law. What had been a voluntary, and therefore an honorable exception, now became the rule, and the former rule became the exception. Connubial intercourse appeared incompatible with the dignity and purity of the priestly office and of priestly functions, especially with the service of the altar. The clergy, as the model order, could not remain below the moral ideal of monasticism, extolled by all the fathers of the church, and must exhibit the same unconditional and undivided devotion to the church within the bosom of society, which monasticism exhibited without it. While placed by their calling in unavoidable contact with the world, they must vie with the monks at least in the virtue of sexual purity, and thereby increase their influence over the people. Moreover, the celibate life secured to the clergy greater independence toward the state and civil society, and thus favored the interests of the hierarchy. But, on the other hand, it estranged them more and more from the sympathies and domestic relations of the people, and tempted them to the illicit indulgence of appetite, which, perhaps, did more injury to the cause of Christian morality and to the true influence of the clergy, than the advantage of forced celibacy could compensate.

In the practice of clerical celibacy, however, the Greek and the Latin churches diverged in the fourth century, and are to this day divided. The Greek church stopped halfway, and limited the injunction of celibacy to the higher clergy, who were accordingly chosen generally from the monasteries or from the ranks of widower-presbyters; while the Latin church extended the law to the lower clergy, and at the same time carried forward the hierarchical principle to absolute papacy. The Greek church differs from the Latin, not by any higher standard of marriage, but only by a closer adherence to earlier usage and by less consistent application of the ascetic principle. It is in theory as remote from the evangelical Protestant church as the Latin is, and approaches it only in practice. It sets virginity far above marriage, and regards marriage only in its aspect of negative utility. In the single marriage of a priest it sees in a measure a necessary evil, at best only a conditional good, a wholesome concession to the flesh for the prevention of immorality,\footnote{1 Cor. vii. 9.} and requires of its highest office bearers total abstinence from all matrimonial intercourse. It wavers, therefore, between a partial permission and a partial condemnation of priestly marriage.

In the East, one marriage was always allowed to the clergy, and at first even to bishops, and celibacy was left optional. Yet certain restrictions were early introduced, such as the prohibition of marriage after ordination (except in deacons and subdeacons), as well as of second marriage after baptism; the apostolic direction, that a bishop should be the husband of one wife,\footnote{1 Tim. iii. 2, 12; Lit. i. 6.} being taken as a prohibition of successive polygamy, and at the same time as an allowance of one marriage. Besides second marriage, the marrying of a concubine, a widow, a harlot, a slave, and an actress, was forbidden to the clergy. With these restrictions, the “Apostolic Constitutions” and “Canons” expressly permitted the marriage of priests contracted before ordination, and the continuance of it after ordination.\footnote{Lib. vi. cap. 17 (ed. Ueltzen, p. 144):\textgreek{επίσκοπον καὶ πρεβύτερον καὶ διάκονον} [thus including the bishop] \textgreek{εἴπομεν μονογάμους καθίστασθαι}... \textgreek{μὴ ἐξεῖναι δὲ αὐτοῖς μετὰ χειροτονίαν ἀγάμοις οὗσιν ἔτι ἐπὶ γάμον ἔρχεσθαι}, etc. Can. Apost. can. 17 (p. 241): \textgreek{Ὁ δυσὶ γάμοις συμπλακεὶς μετὰ τὸ βάπτισμα}... \textgreek{οὐ δύναται εῖναι ἐπίσκοπος ἢ πρεσβύτερος ἢ διάκονος ἢ ὅλως τοῦ καταλόγου τοῦ ἱερατικοῦ}. Comp. can. 18 and can. 5.} The synod of Ancyra, in 314, permitted deacons to marry even after ordination, in case they had made a condition to that effect beforehand; otherwise they were to remain single or lose their office.\footnote{Can. 10. Comp. Dr. Hefele, Conciliengeschichte, i. p. 198.} The Synod of New Caesarea, which was held at about the same time, certainly before 325, does not go beyond this, decreeing: “If a presbyter (not a deacon) marry (that is, after ordination), he shall be expelled from the clergy; and if he practise lewdness, or become an adulterer, he shall be utterly thrust out and held to penance.”\footnote{Can. 1. In Harduin, tom. v. p. 1499; Hefele, Conciliengesch. i. 211 sq. This canon passed even into the Corpus juris can. c. 9, dist. 28.} At the general council of Nice, 325, it was proposed indeed, probably by the Western bishop Hosius,\footnote{Hosius of Cordova, who was present at the council of Elvira in Spain, in 305, where a similar proposition was made and carried (can. 33). In the opinion above given, Theiner, Gieseler, Robertson, and Hefele agree.} to forbid entirely the marriage of priests; but the motion met with strong opposition, and was rejected. A venerable Egyptian bishop, Paphnutius, though himself a strict ascetic from his youth up, and a confessor who in the last persecution had lost an eye and been crippled in the knee, asserted with impressiveness and success, that too great rigor would injure the church and promote licentiousness and that marriage and connubial intercourse were honorable and spotless things.\footnote{See the account in Socrates, H. E. i. c. 11, where that proposition to prohibit priestly marriage is called an innovation, a \textgreek{νόμος νεαρός}\emph{;} in Sozomen, H. E. i. c. 23; and in Gelasius, Hist. Conc. Nic. ii. 32. The statement is thus sufficiently accredited, and agrees entirely with the ancient practice of the Oriental church and the directions of the Apostolic Constitutions and Canons. The third canon of the council of Nice goes not against it, since it forbids only the immorality of mulieres \emph{subintroductae} (comp. vol. i. § 95). The doubts of several Roman divines (Baronius, Bellarmine, Valesius), who would fain trace the celibacy of the clergy to an apostolic origin, arise evidently from dogmatic bias, and are sufficiently refuted by Hefele, a Roman Catholic historian, in his Conciliengeschichte, vol. i. p. 417 sqq.} The council of Gangra in Paphlagonia (according to some, not till the year 380) condemned, among several ascetic extravagances of the bishop Eustathius of Sebaste and his followers, contempt for married priests and refusal to take part in their ministry.\footnote{Comp. Hefele, l.c. i. 753 sqq.} The so-called Apostolic Canons, which, like the Constitutions, arose by a gradual growth in the East, even forbid the clergy, on pain of deposition and excommunication, to put away their wives under the pretext of religion.\footnote{Can. 5 (ed. Ueltzen, p. 239): \textgreek{Ἐπίσκοπος ἢ πρεσβύτερος ἢ διάκονος τὴν ἑαυτοῦ  γυναῖκα μὴ ἐκβαλλέτω προφάσει εὐλαβείας; ἐὰν δὲ ἐκβαλῆ, ἀφοριξέσθω, ἐπιμένων δὲ καθαιρείσθω}. Comp. Const. Apost. vi. 17.} Perhaps this canon likewise was occasioned by the hyper-asceticism of Eustathius.

Accordingly we not unfrequently find in the Oriental church, so late as the fourth and fifth centuries, not only priests, but even bishops living in wedlock. One example is the father of the celebrated Gregory Nazianzen, who while bishop had two sons, Gregory and the younger Caesarius, and a daughter. Others are Gregory of Nyssa, who, however, wrote an enthusiastic eulogy of the unmarried life, and lamented his loss of the crown of virginity; and Synesius († about 430), who, when elected bishop of Ptolemais in Pentapolis, expressly stipulated for the continuance of his marriage connection.\footnote{Declaring: “God, the law, and the consecrated hand of Theophilus (bishop of Alexandria), have given me a wife. I say now beforehand, and I protest, that I will neither ever part from her, nor live with her in secret as if in an unlawful connection; for the one is utterly contrary to religion, the other to the laws; but I desire to receive many and good children from her” (Epist. 105 ed. Basil., cited in the original Greek in Gieseler). Comp. on the instances of married bishops, Bingham, Christ. Antiq. b. iv. ch. 5; J. A. Theiner and A. Theiner, Die Einführung der erzwungenen Ehelosigkeit der christl. Geistlichen u. ihre Folgen (Altenburg, 1828), vol. i. p. 263 sqq., and Gieseler, vol. i. div. 2, § 97, notes at the close. The marriage of Gregory of Nyssa with Theosebia is disputed by some Roman Catholic writers, but seems well supported by Greg. Naz. Ep. 95, and Greg Nyss. De virg. 3.} Socrates, whose Church History reaches down to the year 439, says of the practice of his time, that in Thessalia matrimonial intercourse after ordination had been forbidden under penalty of deposition from the time of Heliodorus of Trica, who in his youth had been an amatory writer; but that in the East the clergy and bishops voluntarily abstained from intercourse with their wives, without being required by any law to do so; for many, he adds, have had children during their episcopate by their lawful wives.\footnote{Hist. Eccl. v. cap. 22\textgreek{· Τῶν ἐν ἀνατολῇ πάντων γνώμῃ} (i.e. from principle or voluntarily—according to the reading of the Florentine codex) \textgreek{ἀπεχομένων, καὶ τῶν ἐπισκόπων, εἰ καὶ βούλοίντο, οὐ μὴν ἀνάγκῃ νόμου τοῦτο ποιούντων. Πολλοὶ γὰρ αὐτῶν ἐν τῷ καιρῷ τῆς ἐπισκοπῆς καὶ παῖδας ἐκ τῆς νομίμης γαμετῆς πεποιήκασιν}} There were Greek divines, however, like Epiphanius, who agreed with the Roman theory. Justinian I. was utterly opposed to the marriage of priests, declared the children of such connection illegitimate, and forbade the election of a married man to the episcopal office (a.d. 528). Nevertheless, down to the end of the seventh century, many bishops in Africa, Libya, and elsewhere, continued to live in the married state, as is expressly said in the twelfth canon of the Trullan council; but this gave offence and was forbidden. From that time the marriage of bishops gradually disappears, while marriage among the lower clergy continues to be the rule.

This Trullan council, which was the sixth ecumenical\footnote{More precisely, the \emph{second} Trullan council, held in the Trullan hall of the imperial palace in Constantinople; also called \emph{Concilium Quinisextum}, \textgreek{σύνοδος πενθέκτη}, being considered a supplement to the fifth and sixth general councils. Comp. respecting it Hefele, iii. 298 sqq.} (a.d. 692), closes the legislation of the Eastern church on the subject of clerical marriage. Here—to anticipate somewhat—the continuance of a first marriage contracted before ordination was prohibited in the case of bishops on pain of deposition, but, in accordance with the Apostolic Constitutions and Canons, allowed in the case of presbyters and deacons (contrary to the Roman practice), with the Old Testament restriction, that they abstain from sexual intercourse during the season of official service, because he who administers holy things must be pure.\footnote{1 Can. 3, 4, and especially 12, 13, and 48. In the latter canon bishops are directed, after ordination, to commit their wives to a somewhat remote cloister, though to provide for their support.} The same relation is thus condemned in the one case as immoral, in the other approved and encouraged as moral; the bishop is deposed if he retains his lawful wife and does not, immediately after being ordained, send her to a distant cloister; while the presbyter or deacon is threatened with deposition and even excommunication for doing the opposite and putting his wife away.

The Western church, starting from the perverted and almost Manichaean ascetic principle, that the married state is incompatible with clerical dignity and holiness, instituted a vigorous effort at the end of the fourth century, to make celibacy, which had hitherto been left to the option of individuals, the universal law of the priesthood; thus placing itself in direct contradiction to the Levitical law, to which in other respects it made so much account of conforming. The law, however, though repeatedly enacted, could not for a long time be consistently enforced. The canon, already mentioned, of the Spanish council of Elvira in 305, was only provincial. The first prohibition of clerical marriage, which laid claim to universal ecclesiastical authority, at least in the West, proceeded in 385 from the Roman church in the form of a decretal letter of the bishop Siricius to Himerius, bishop of Tarragona in Spain, who had referred several questions of discipline to the Roman bishop for decision. It is significant of the connection between the celibacy of the clergy and the interest of the hierarchy, that the first properly papal decree, which was issued in the tone of supreme authority, imposed such an unscriptural, unnatural, and morally dangerous restriction. Siricius contested the appeal of dissenting parties to the Mosaic law, on the ground that the Christian priesthood has to stand not merely for a time, but perpetually, in the service of the sanctuary, and that it is not hereditary, like the Jewish; and he ordained that second marriage and marriage with a widow should incapacitate for ordination, and that continuance in the married state after ordination should be punished with deposition.\footnote{Epist. ad Himerium Episc. Tarraconensem (in Harduin, Acta Conc. i. 849-850), c 7: “Hi vero, qui illiciti privilegii excusatione nituntur, ut sibi asserant veteri hoc lege concessum: noverint se ab omni ecclesiastico honore, quo indigne usi sunt, apostolicae sedis auctoritate dejectos .... Si quilibet episcopus, presbyter atque diaconus, quod non optamus, deinceps fuerit talis inventus, jam nunc sibi omnem per nos indulgentiae aditum intelligat obseratum: quia ferro necesse est excidantur vulnera, quae fomentorum non senserint medicinam.” The exegesis of Siricius is utterly arbitrary in limiting the demand of holiness (Lev. xx. 7) to the priests and to abstinence from matrimonial intercourse, and in referring the words of Paul respecting walking in the flesh, Rom. viii. 8, 9, to the married life, as if marriage were thus incompatible with the idea of holiness. Comp. also the striking remarks of Greenwood, Catheda Petri, vol. i. p. 265 sq., and Milman, Hist. of Latin Christianity, i. 119 (Amer. ed.), on Siricius.} And with this punishment he threatened not bishops only, but also presbyters and deacons. Leo the Great subsequently, extended the requirement of celibacy even to the subdiaconate. The most eminent Latin church fathers, Ambrose, Jerome, and even Augustine—though the last with more moderation—advocated the celibacy of priests. Augustine, with Eusebius of Vercella before him (370), united their clergy in a cloister life, and gave them a monastic stamp; and Martin of Tours, who was a monk from the first, carried his monastic life into his episcopal office. The councils of Italy, Africa, Spain, and Gaul followed the lead of Rome. The synod of Clermont, for example (a.d. 535), declared in its twelfth canon: “No one ordained deacon or priest may continue matrimonial intercourse. He is become the brother of her who was his wife. But since some, inflamed with lust, have rejected the girdle of the warfare [of Christ], and returned to marriage intercourse, it is ordered that such must lose their office forever.” Other councils, like that of Tours, 461, were content with forbidding clergymen, who begat children after ordination, to administer the sacrifice of the mass, and with confining the law of celibacy ad altiorem gradum.\footnote{Comp. Hefele, ii. 568, and Gieseler, l.c. (§ 97, note 7).}

But the very fact of the frequent repetition of these enactments, and the necessity of mitigating the penalties of transgression, show the great difficulty of carrying this unnatural restriction into general effect. In the British and Irish church, isolated as it was from the Roman, the marriage of priests continued to prevail down to the Anglo-Saxon period.

But with the disappearance of legitimate marriage in the priesthood, the already prevalent vice of the cohabitation of unmarried ecclesiastics with pious widows and virgins “secretly brought in,”\footnote{The so-called \emph{sorores,} or \emph{mulieres subintroductae}, or\textgreek{παρθένοι συνείσακτοι}. Comp. on the origin of this practice, vol. i. § 95.} became more and more common. This spiritual marriage, which had begun as a bold ascetic venture, ended only too often in the flesh, and prostituted the honor of the church.

The Nicene council of 325 met the abuse in its third canon with this decree: “The great council utterly forbids, and it shall not be allowed either to a bishop, or a priest, or a deacon, or any other clergyman, to have with him a συνθείσακτος, unless she be his mother, or sister, or aunt, or some such person, who is beyond all suspicion.”\footnote{By a misinterpretation of the term \textgreek{συνείσακτος}, the sense of which is fixed in the usage of the early church, Baronius and Bellarmine erroneously find in this canon a universal law of celibacy, and accordingly deny the above-mentioned statement respecting Paphnutius. Comp. Hefele, i. 364.} This canon forms the basis of the whole subsequent legislation of the church de cohabitatione clericorum et mulierum. It had to be repeatedly renewed and strengthened; showing plainly that it was often disobeyed. The council of Toledo in Spain, a.d. 527 or 531, ordered in its third canon: “No clergyman, from the subdeacon upward, shall live with a female, be she free woman, freed woman, or slave. Only a mother, or a sister, or other near relative shall keep his house. If he have no near relative, his housekeeper must live in a separate house, and shall under no pretext enter his dwelling. Whosoever acts contrary to this, shall not only be deprived of his spiritual office and have the doors of his church closed, but shall also be excluded from all fellowship of Catholics.” The Concilium Agathense in South Gaul, a.d. 506, at which thirty-five bishops met, decreed in the tenth and eleventh canons: “A clergyman shall neither visit nor receive into his house females not of his kin; only with his mother, or sister, or daughter, or niece may he live. Female slaves, also, and freed women, must be kept away from the house of a clergyman.” Similar laws, with penalties more or less severe, were passed by the council of Hippo, 393, of Angers, 453, of Tours, 461, of Lerida in Spain, 524, of Clermont, 535, of Braga, 563, of Orleans, 538, of Tours, 567.\footnote{Comp. the relevant canons of these and other councils in the second and third volumes of Hefele’s Conciliengeschichte.} The emperor Justinian, in the twenty-third Novelle, prohibited the bishop having any woman at all in his house, but the Trullan council of 692 returned simply to the Nicene law.\footnote{Can. 5: “No clergyman shall have a female in his house, but those allowed in the old canon (Nicaen. c. 3). Even eunuchs are to observe this.”} The Western councils also made attempts to abolish the exceptions allowed in the Nicene canon, and forbade clergymen all intercourse with women, except in presence of a companion.

This rigorism, however, which sheds an unwelcome light upon the actual state of things that made it necessary, did not better the matter, but rather led to such a moral apathy, that the Latin church in the middle age had everywhere to contend with the open concubinage of the clergy, and the whole energy of Gregory VII. was needed to restore in a measure the old laws of celibacy, without being sufficient to prevent the secret and, to morality, far more dangerous violations of it.\footnote{“Throughout the whole period,” says Milman (Hist. of Latin Christianity, i. 123), “from Pope Siricius to the Reformation, as must appear in the course of our history, the law [of clerical celibacy] was defied, infringed, eluded. It never obtained anything approaching to general observance, though its violation was at times more open, at times more clandestine.”} The later ecclesiastical legislation respecting the mulieres subintroductae is more lenient, and, without limiting the intercourse of clergymen to near kindred, generally excludes only concubines and those women “de quibus possit haberi suspicio.”\footnote{So the Concilium Tridentinum, sess. xxv. de reform. cap. 14. Comp. also the article Subintroductae, in the 10th volume of Wetzer and Welte’s Cath. Church Lexicon.}

\xsection{Moral Character of the Clergy in general}

§ 51. Moral Character of the Clergy in general.

Augustine gives us the key to the true view of the clergy of the Roman empire in both light and shade, when he says of the spiritual office: “There is in this life, and especially in this day, nothing easier, more delightful, more acceptable to men, than the office of bishop, or presbyter, or deacon, if the charge be administered superficially and to the pleasure of men; but nothing in the eye of God more wretched, mournful, and damnable. So also there is in this life, and especially in this day, nothing more difficult, more laborious) more hazardous than the office of bishop, or presbyter, or deacon; but nothing in the eye of God more blessed, if the battle be fought in the manner enjoined by our Captain.”\footnote{Epist. 21 ad Valerium Nihil esse in hac vita et maxime hoc tempore facilius et laetitius et hominibus acceptabilius episcopi aut presbyteri aut diaconi officio, si perfunctorie atque adulatorie res agatur: sed nihil apud Deum miserius et tristius et damnabilius. Item nihil esse in hac vita et maxime hoc tempore difficilius, laboriosius, periculosius episcopi aut presbyteri aut diaconi officio, sed apud Deum nihil beatius, si eo modo militetur, quo noster imperator jubet.” This epistle was written soon after his ordination to the priesthood, a.d.391. See Opera, ed. Bened. tom. ii p. 25.} We cannot wonder, on the one hand that, in the better condition of the church and the enlarged field of her labor, a multitude of light-minded and unworthy men crowded into the sacred office, and on the other, that just the most earnest and worthy bishops of the day, an Ambrose, an Augustine, a Gregory Nazianzen, and a Chrysostom, trembled before the responsibility of the office, and had to be forced into it in a measure against their will, by the call of the church.

Gregory Nazianzen fled into the wilderness when his father, without his knowledge, suddenly consecrated him priest in the presence of the congregation (361). He afterward vindicated this flight in his beautiful apology, in which he depicts the ideal of a Christian priest and theologian. The priest must, above all, he says, be a model of a Christian, offer himself a holy sacrifice to God, and be a living temple of the living God. Then he must possess a deep knowledge, of souls, and, as a spiritual physician, heal all classes of men of various diseases of sin, restore, preserve, and protect the divine image in them, bring Christ into their hearts by the Holy Ghost, and make them partakers of the divine nature and of eternal salvation. He must, moreover, have at command the sacred philosophy or divine science of the world and of the worlds, of matter and spirit, of good and evil angels, of the all-ruling Providence, of our creation and regeneration, of the divine covenants, of the first and second appearing of Christ, of his incarnation, passion, and resurrection, of the end of all things and the universal judgment, and above all, of the mystery of the blessed Trinity; and he must be able to teach and elucidate these doctrines of faith in popular discourse. Gregory, sets forth Jesus as the perfect type of the priest, and next to him he presents in an eloquent picture the apostle Paul, who lived only for Christ, and under all circumstances and amid all trials by sea and land, among Jews and heathen, in hunger and thirst, in cold and nakedness, in freedom and bonds, attested the divine power of the gospel for the salvation of the world. This ideal, however, Gregory found but seldom realized. He gives on the whole a very unfavorable account of the bishops, and even of the most celebrated councils of his day, charging them with ignorance unworthy means of promotion, ambition, flattery, pride, luxury, and worldly mindedness. He says even: “Our danger now is, that the holiest of all offices will become the most ridiculous; for the highest clerical places are gained not so much by virtue, as by iniquity; no longer the most worthy, but the most powerful, take the episcopal chair.”\footnote{Orat. xliii. c. 46 (Opera, ed. Bened. tom. i. p. 791), in the Latin translation: “Nunc autem periculum est, ne ordo omnium sanctissimus, sit quoque omnium maxime ridiculus. Non enim virtute magis, quam maleficio et scelere, sacerdotium paratur; nec digniorum, sed potentiorum, throni sunt.” In the following chapter, however, he represents his friend Basil as a model of all virtues.} Though his descriptions, especially in the satirical poem “to himself and on the bishops,” composed probably after his resignation in Constantinople (a.d. 381), may be in many points exaggerated, yet they were in general drawn from life and from experience.\footnote{Comp. Ullmann: Gregor von Nazianz, Erste Beilage, p. 509-521, where the views of this church father on the clerical office and the clergy of his time are presented at large in his own words. Also Gieseler, i., ii. § 103, gives copious extracts from the writings of Gregory on the vices of the clergy.}

Jerome also, in his epistles, unsparingly attacks the clergy of his time, especially the Roman, accusing them of avarice and legacy hunting, and drawing a sarcastic picture of a clerical fop, who, with his fine scented clothes, was more like a bridegroom than a clergyman.\footnote{Hieron. ad Eustochium, and especially ad Nepotianum, de vita clericorum et monachorum (Opera, ed. Vall. tom. i. p. 252 sqq.). Yet neither does he spare the monks, but says, ad Nepot.: “Nonnulli sunt ditiores monachi quam fuerant seculares et clerici qui possident opes sub Christo paupere, quas sub locuplete et fallaci Diabolo non habuerant.”} Of the rural clergy’, however, the heathen Ammianus Marcellinus bears a testimony, which is certainly reliable, to their simplicity, contentment, and virtue.\footnote{Lib. xxvii. c. 3, sub ann. 367.}

Chrysostom, in his celebrated treatise on the priesthood,\footnote{\textgreek{Περὶ ἱερωσύνης}, or De Sacerdotio libri sex. The work has been often published separately, and several times translated into modern languages (into German, for example, by Hasselbach, 1820, and Ritter, 1821; into English by Hollier, 1740, Bunce, 1759; Hohler, 1837; Marsh, 1844; and best by B. Harris Cowper, London, 1866). Comp. the list of twenty-three different separate editions and translations in Lomler: Joh. Chrysost. Opera praestantissima Gr. et Lat. Rudolph. 1840, p. viii, ix.} written probably, before his ordination (somewhere between the years 375 and 381), or while he was deacon (between 381 and 386), portrayed the theoretical and practical qualifications, the exalted duties, responsibilities, and honors of this office, with youthful enthusiasm, in the best spirit of his age. He requires of the priest, that he be in every respect better than the monk, though, standing in the world, he have greater dangers and difficulties to contend with.\footnote{De Sacerdotio, lib. vi. cap. 2-8.} He sets up as the highest object of the preacher, the great principle stated by, Paul, that in all his discourses he should seek to please God alone, not men. “He must not indeed despise the approving demonstrations of men; but as little must he court them, nor trouble himself when his hearers withhold them. True and imperturbable comfort in his labors he finds only in the consciousness of having his discourse framed and wrought out to the approval of God.”\footnote{\textgreek{Πρὸς ἀρέσκειαν τοῦ Θεοῦ}, lib. v. c. 7.} Nevertheless the book as a whole is unsatisfactory. A comparison of it with the “Reformed Pastor” of Baxter, which is far deeper and richer in all that pertains to subjective experimental Christianity and the proper care of souls, would result emphatically in favor of the English Protestant church of the seventeenth century.\footnote{Comp. also the remarks of B. H. Cowper in the introduction to his English translation, Lond. 1866, p. xiii.}

We must here particularly notice a point which reflects great discredit on the moral sense of many of the fathers, and shows that they had not wholly freed themselves from the chains of heathen ethics. The occasion of this work of Chrysostom was a ruse, by which he had evaded election to the bishopric, and thrust it upon his friend Basil.\footnote{Not Basil the Great (as Socrates supposes), for he was much older, and died in 379; but probably (as Montfaucon conjectures) the bishop of Raphanea in Syria, near Antioch, whose name appears among the bishops of the council of Constantinople, in 381.} To justify this conduct, he endeavors at large, in the fifth chapter of the first book, to prove that artifice might be lawful and useful; that is, when used as a means to a good end. “Manifold is the potency of deception, only it must not be employed with knavish intent. And this should be hardly called deception, but rather a sort of accommodation (οἰκονομία), wisdom, art, or sagacity, by which one can find many ways of escape in an exigency, and amend the errors of the soul.” He appeals to biblical examples, like Jonathan and the daughter of Saul, who by deceiving their father rescued their friend and husband; and, unwarrantably, even to Paul, who became to the Jews a Jew, to the Gentiles a Gentile, and circumcised Timothy, though in the Epistle to the Galatians he pronounced circumcision useless. Chrysostom, however, had evidently learned this, loose and pernicious principle respecting the obligation of truthfulness, not from the Holy Scriptures, but from the Grecian sophists.\footnote{Even the purest moral philosopher of antiquity, Plato, vindicates falsehood, and recommends it to physicians and rulers as a means to a good end, a help to the healing of the sick or to the advantage of the people. Comp. De republ. iii. p. 266, ed. Bipont.: \textgreek{Εἰ γὰρ ὀρθῶς ἐλέχγομεν ἄρτι, καὶ τῷ ὄντι θεοῖς μὲν ἄχρηστον ψεῦδος ἀνθρώποις δὲ χρήςιμον, ὡς ἐν φαρμάκου εἴδει, δῆλον ὅτι τὸ γε τοιοῦτον ἱατροῖς δοτέον, ἰδιώταις δὲ οὐχ ἁπτέον. Δῆλον, ἔφη. Τοῖς ἄρχουσι δὴ τῆς πόλεως , εἴπερ τισὶν ἄλλοις, προσήκει ψεύδεσθαι ἢ πολεμίων ἢ πολιτῶν ἕνεκα, ἐπ  ὠφελείᾳτῆς πόλεως· τοῖς δὲ ἄλλοις πᾶσιν οὐχ ἁπτέον τοῦ τοιούτου}. . The Jewish philosophizing theologian, Philo, had a similar view, in his work: Quod Deus sit immutabilis, p. 302.} Besides, he by no means stood alone in the church in this matter, but had his predecessors in the Alexandrian fathers,\footnote{Clemens Alex., Strom. vi. p. 802, and Origen, Strom. vi. (in Hieron. Apol. i. Adv. Ruf. c. 18), where he adduces the just cited passage of Plato in defence of a doubtful accommodation at the expense of truth. See the relevant passages in Gieseler, i. § 63, note 7.} and his followers in Cassian, Jerome, and other eminent Catholic divines.

Jerome made a doubtful distinction between γυμναστικῶς scribere and δογματικῶς scribere, and, with Origen, explained the severe censure of Paul on Peter in Antioch, for example, as a mere stroke of pastoral policy, or an accommodation to the weakness of the Jewish Christians at the expense of truth.\footnote{Epist. 48 (ed. Vall., or Ep. 30 ed. Bened., Ep. 50 in older editions), ad Pammachium, pro libris contra Jovinianum, and Comm. ad Gal. ii. 11 sqq. Also Johannes Cassianus, a pupil of Chrysostom, defends the lawfulness of falsehood and deception in certain cases, Coll. xvii. 8 and 17.} But Augustine’s delicate Christian sense of truth revolted at this construction, and replied that such an interpretation undermined the whole authority of Holy Scripture; that an apostle could never lie, even for a good object; that, in extremity, one should rather suppose a false reading, or wrong translation, or suspect his own apprehension; but that in Antioch Paul spoke the truth and justly censured Peter openly for his inconsistency, or for a practical (not a theoretical) error, and thus deserves the praise of righteous boldness, as Peter on the other hand, by his meek submission to the censure, merits the praise of holy humility.\footnote{Comp. the somewhat sharp correspondence of the two fathers in Hieron. Epist. 101-105, 110, 112, 115, 134, 141, in Vallarsi’s ed. (tom. i. 625 sqq.), or in August. Epist 67, 68, 72-75, 81, 82 (in the Bened. ed. of Aug. tom. ii. 161 sqq.); August.: De mendacio, and Contra mendacium; also the treatise of Möhler mentioned above, 41, on this controversy, so instructive in regard to the patristic ethics and exegesis.}

Thus in Jerome and Augustine we have the representatives of two opposite ethical views: one, unduly subjective, judging all moral acts merely by their motive and object, and sanctioning, for example, tyrannicide, or suicide to escape disgrace, or breach of faith with heretics (as the later Jesuitical casuistry does with the utmost profusion of sophistical subtlety); the other, objective, proceeding on eternal, immutable principles and the irreconcilable opposition of good and evil, and freely enough making prudence subservient to truth, but never truth subservient to prudence.

Meantime, in the Greek church also, as early as the fourth century, the Augustinian view here and there made its way; and Basil the Great, in his shorter monastic Rule,\footnote{Regul. brev. interrogate 76, cited by Neander in his monograph on Chrysostom(3d ed.) i. p. 97. Neander there adduces still another similar testimony against the lawfulness of the lie, by the contemporaneous Egyptian monk, John of Lycopolis, from Pallad. Hist. Lausiaca.} rejected even accommodation (οἰκονομία) for a good end, because Christ ascribes the lie, without distinction of kinds, exclusively to Satan.\footnote{John, viii. 44.} In this respect, therefore, Chrysostom did not stand at the head of his age, but represented without doubt the prevailing view of the Eastern church.

The legislation of the councils with reference to the clergy, shows in general the earnestness and rigor with which the church guarded the moral purity and dignity of her servants. The canonical age was, on the average, after the analogy of the Old Testament, the five-and-twentieth year for the diaconate, the thirtieth for the priesthood and episcopate. Catechumens, neophytes, persons baptized at the point of death, penitents, energumens (such as were possessed of a devil), actors, dancers, soldiers, curials (court, state, and municipal officials),\footnote{The ground on which even civil officers were excluded, is stated by the Roman council of 402, which ordained in the tenth canon: “One who is clothed with a civil office cannot, on account of the sins almost necessarily connected with it, become a clergyman without previous penance.” Comp. Mansi, iii. 1133, and Hefele; ii. 75.} slaves, eunuchs, bigamists, and all who led a scandalous life after baptism, were debarred from ordination. The frequenting of taverns and theatres, dancing and gambling, usury and the pursuit of secular business were forbidden to clergymen. But on the other hand, the frequent repetition of warnings against even the lowest and most common sins, such as licentiousness, drunkenness, fighting, and buffoonery, and the threatening of corporal punishment for certain misdemeanors, yield an unfavorable conclusion in regard to the moral standing of the sacred order.\footnote{Comp. the decrees of councils in Hefele, ii. 574, 638, 686, 687, 753, 760, \&c. Even the Can. Apost. 27, 65, and 72, are directed against common crimes in the clergy, such as battery, murder, and theft, which therefore must have already appeared, for legislation always has regard to the actual state of things. The Pastoral Epistles of Paul contain no exhortations or prohibitions of this kind.} Even at the councils the clerical dignity was not seldom desecrated by outbreaks of coarse passion; insomuch that the council of Ephesus, in 449, is notorious as the “council of robbers.”

In looking at this picture, however, we must not forget that in this, period of the sinking empire of Rome the task of the clergy was exceedingly difficult, and amidst the nominal conversion of the whole population of the empire, their number and education could not keep pace with the sudden and extraordinary expansion of their field of labor. After all, the clerical office was the great repository of intellectual and moral force for the world. It stayed the flood of corruption; rebuked the vices of the times; fearlessly opposed tyrannical cruelty; founded institutions of charity and public benefit; prolonged the existence of the Roman empire; rescued the literary treasures of antiquity; carried the gospel to the barbarians, and undertook to educate and civilize their rude and vigorous hordes. Out of the mass of mediocrities tower the great church teachers of the fourth and fifth centuries, combining all the learning, the talent, and the piety of the time, and through their immortal writings mightily moulding the succeeding ages of the world.

\xsection{The Lower Clergy}

§ 52. The Lower Clergy.

As the authority and influence of the bishops, after the accession of Constantine, increased, the lower clergy became more and more dependent upon them. The episcopate and the presbyterate were now rigidly distinguished. And yet the memory of their primitive identity lingered. Jerome, at the end of the fourth century, reminds the bishops that they owe their elevation above the presbyters, not so much to Divine institution as to ecclesiastical usage; for before the outbreak of controversies in the church there was no distinction between the two, except that presbyter is a term of age, and bishop a term of official dignity; but when men, at the instigation of Satan, erected parties and sects, and, instead of simply following Christ, named themselves of Paul, or Apollos, or Cephas, all agreed to put one of the presbyters at the head of the rest, that by his universal supervision of the churches, he might kill the seeds of division.\footnote{Hieron. Comm. ad Tit. i. 7: “Idem est ergo presbyter qui episcopus, et antequam diaboli instinctu studia in religione fierent ... communi presbyterorum consilio ecclesiae gubernabantur,” etc. Comp. Epist. ad Evangelum presbyterum Ep. 146, ed. Vall. Opera, i. 1074 sqq.; Ep. 101, ed. Bened.), and Epist. ad Oceanum (Ep. 69, ed. Vall., Ep. 82, ed. Bened.). In the latter epistle he remarks: “Apud veteres iidem episcopi et presbyteri fuerunt, quia illud nomen dignitatis est, hoc aetatis.”} The great commentators of the Greek church agree with Jerome in maintaining the original identity of bishops and presbyters in the New Testament.\footnote{Chrysostom, Hom. i. in Ep. ad Philipp. (Phil. i. 1, on the words \textgreek{συν ἐπισκόποις}\emph{,} which imply a number of bishops, i.e. presbyters in one and the same congregation), observes: \textgreek{τοὺς πρεσβυτέρου· οὕτως εκάλεσε· τότε γὰρ τέως ἐκοινώνουν τοῖς ὀνόμασι}.. Of the same opinion are Theodoret, ad Phil. i. 1, and ad Tim. iii. 1; Ambrosiaster, ad Eph. iv. 11; and the author of the pseudo-Augustinian Questiones V. et N.T., qu. 101. Comp. on this whole subject of the original identity of \textgreek{ἐπίσκοπος}and \textgreek{πρεσβύτερος}, my History of the Apostolic Church, § 132 (Engl. translation, p. 522-531), and Rich. Rothe: Anfänge der christlichen Kirche, i. p. 207-217.}

In the episcopal or cathedral churches the Presbyters still formed the council of the bishop. In town and country congregations, where no bishop officiated, they were more independent. Preaching, administration of the sacraments, and care of souls were their functions. In. North Africa they were for a long time not allowed to preach in the presence of the bishop; until Augustine was relieved by his bishop of this restriction. The seniores plebis in the African church of the fourth and fifth centuries were not clergymen, but civil personages and other prominent members of the congregation.\footnote{Optatus of Mileve calls them, indeed, \emph{ecclesiasticos viros}; not, however, in the sense of \emph{clerici}, from whom, on the contrary, he distinguishes them, but in the broad sense of catholic Christians as distinguished from heathens and heretics. Comp. on these \emph{seniores plebis,} or\emph{lay elders,} as they are called, the discussion of Dr. Rothe: Die Anfänge der christl. Kirche u. ihrer Verfassung, vol. i. p. 227 sqq.}

In the fourth century arose the office of archpresbyter, whose duty it was to preside over the worship, and sometimes to take the place of the bishop in his absence or incapacity.

The Deacons, also called Levites, retained the same functions which they had held in the preceding period. In the West, they alone, not the lectors, were allowed to read in public worship the lessons from the Gospels; which, containing the words of the Lord, were placed above the Epistles, or the words of the apostles. They were also permitted to baptize and to preach. After the pattern of the church in Jerusalem, the number of deacons, even in large congregations, was limited to seven; though not rigidly, for the cathedral of Constantinople had, under Justinian I., besides sixty presbyters, a hundred deacons, forty deaconesses, ninety subdeacons, a hundred and ten lectors, twenty-five precentors, and a hundred janitors—a total of five hundred and twenty-five officers. Though subordinate to the presbyters, the deacons frequently stood in close relations with the bishop, and exerted a greater influence. Hence they not rarely looked upon ordination to the presbyterate as a degradation. After the beginning of the fourth century an archdeacon stood at the head of the college, the most confidential adviser of the bishop, his representative and legate, and not seldom his successor in office. Thus Athanasius first appears as archdeacon of Alexandria at the council of Nice, clothed with important influence; and upon the death of the latter he succeeds to the patriarchal chair of Alexandria.

The office of Deaconess, which, under the strict separation of the sexes in ancient times, and especially in Greece, was necessary to the completion of the diaconate, and which originated in the apostolic age,\footnote{Comp. Rom. xii. 1, 12, and my Hist. of the Apost. Church, § 135, p. 535 sqq.} continued in the Eastern church down to the twelfth century. It was frequently occupied by the widows of clergymen or the wives of bishops, who were obliged to demit the married state before entering upon their sacred office. Its functions were the care of the female poor, sick, and imprisoned, assisting in the baptism of adult women, and, in the country churches of the East, perhaps also of the West, the preparation of women for baptism by private instruction.\footnote{Comp. Pelagius ad Rom. xvi. 1. Neander (iii. p. 314, note; Torrey’s transl. ii. p. 158) infers from a canon of the fourth council of Carthage, that the latter custom prevailed also in the West, since it is there required of “viduae quae ad ministerium baptizandarum mulierum eliguntur,” “ut possint apto et sano sermone docere imperitas et rusticas mulieres.”} Formerly, from regard to the apostolic precept in 1 Tim. v. 9, the deaconesses were required to be sixty years of age.\footnote{Comp. Codex Theodos. 1. xvi., Tit. ii. lex 27: “Nulla nisi emensis 60 annis secundum praeceptum apostoli ad diaconissarum consortium transferatur.”} The general council of Chalcedon, however, in 451, reduced the canonical age to forty years, and in the fifteenth canon ordered: “No female shall be consecrated deaconess before she is forty years old, and not then without careful probation. If, however, after having received consecration, and having been some time in the service, she marry, despising the grace of God, she with her husband shall be anathematized.” The usual ordination prayer in the consecration of deaconesses, according to the Apostolic Constitutions, runs thus: “Eternal God, Father of our Lord Jesus Christ, Creator of man and woman, who didst fill Miriam and Deborah and Hannah and Huldah with the Spirit, and didst not disdain to suffer thine only-begotten Son to be born of a woman; who also in the tabernacle and the temple didst appoint women keepers of thine holy gates: look down now upon this thine handmaid, who is designated to the office of deacon, and grant her the Holy Ghost, and cleanse her from all filthiness of the flesh and of the spirit, that she may worthily execute the work intrusted to her, to thine honor and to the praise of thine Anointed; to whom with thee and the Holy Ghost be honor and adoration forever. Amen.”\footnote{Const. Apost. lib. viii. cap. 20. We have given the prayer in full. Neander (iii. p. 322, note) omits some passages. The custom of ordaining deaconesses is placed by this prayer and by the canon quoted from the council of Chalcedon beyond dispute. The 19th canon of the council of Nice, however, appears to conflict with this, in reckoning deaconesses among the laity, who have no consecration (\textgreek{χειροθεσία}). Some therefore suppose that the ordination of deaconesses did not arise till after the Nicaenum (325), though the Apostolic Constitutions contradict this; while others (as Baronius, and recently Hefele, Concilien-Gesch. 1855, vol. i. p. 414) would resolve the contradiction by distinguishing between the proper\textgreek{χειροθεσία} and the simple benediction. But the consecration of the deaconesses was certainly accompanied with imposition of hands in presence of the whole clergy; since the Apost. Const., 1. viii. c. 19, expressly say to the bishop: \textgreek{Ἐπιθήσεις αὐτῃ τὰς χεῖρας, παρεστῶτος τοῦ πρεσβυτερίου καὶ τῶν διακόνων καὶ τῶν διακονισσῶν}. The contradiction lies, however, in that Nicene canon itself; for (according to the Greek Codices) the \emph{deaconesses} are immediately before counted among the clergy, if we do not, with the Latin translation, read \emph{deacons} instead. Neander helps himself by a distinction between proper deaconesses and widows \emph{abusivè so} called.}

The noblest type of an apostolic deaconess, which has come down to us from this period, is Olympias, the friend of Chrysostom, and the recipient of seventeen beautiful epistles from him.\footnote{They are found in Montfaucon’s Bened. edition of Chrysostom, tom. iii. p. 524-604, and in Lomler’s edition of Joann. Chrysost. Opera praestantissima, 1840, p. 168-252. These seventeen epistles to Olympias are, in the judgment of Photius as quoted by Montfaucon (Op. iii. 524), of the epistles of Chrysostom, “longissimae, elegantissimae, omniumque utilissimae.” Compare also Montfaucon’s prefatory remarks on Olympias.} She sprang from a respectable heathen family, but received a Christian education; was beautiful and wealthy; married in her seventeenth year (a.d. 384) the prefect of Constantinople, Nebridius; but in twenty months after was left a widow, and remained so in spite of the efforts of the emperor Theodosius to unite her with one of his own kindred. She became a deaconess; lived in rigid asceticism; devoted her goods to the poor; and found her greatest pleasure in doing good. When Chrysostom came to Constantinople, he became her pastor, and guided her lavish benefaction by wise counsel. She continued faithful to him in his misfortune; survived him by several years, and died in 420, lamented by all the poor and needy in the city and in the country around.

In the West, on the contrary, the office of deaconess was first shorn of its clerical character by a prohibition of ordination passed by the Gallic councils in the fifth and sixth centuries;\footnote{A mere benediction was appointed in place of ordination. The first synod of Orange (Arausicana i.), in 441, directed in the 26th canon: “Diaconae omnimodis non ordinandae [thus they had previously been ordained in Gaul also, and reckoned with the clergy]; si quae jam sunt, benedictioni, quae populo impenditur, capita submittant.” Likewise was the ordination of deaconesses forbidden by the council of Epaon in Burgundy, in 517, can. 21, and by the second council at Orleans, in 533, can. 17 and 18.} and at last it was wholly abolished. The second synod of Orleans, in 533, ordained in its eighteenth canon: “No woman shall henceforth receive the benedictio diaconalis [which had been substituted for ordinatio], on account of the weakness of this sex.” The reason betrays the want of good deaconesses, and suggests the connection of this abolition of an apostolic institution with the introduction of the celibacy of the priesthood, which seemed to be endangered by every sort of female society. The adoption of the care of the poor and sick by the state, and the cessation of adult baptisms and of the custom of immersion, also made female assistance less needful. In modern times, the Catholic church, it is true, has special societies or orders of women, like the Sisters of Mercy, for the care of the sick and poor, the training of children, and other objects of practical charity; and in the bosom of Protestantism also similar benevolent associations have arisen, under the name of Deaconess Institutes, or Sisters’ Houses, though in the more free evangelical spirit, and without the bond of a vow.\footnote{The Deaconess House (Hutterhaus) at Kaiserswerth on the Rhine, founded in 1836; Bethany in Berlin, 1847; and similar evangelical hospitals in Dresden, 1842, Strasburg, 1842, Paris (institution des deaconess des églises evangéliques de France), 1841, London (institution of Nursing Sisters), 1840, New York (St. Luke’s Hopital), Pittsburg, 1849, Smyrna, Jerusalem, etc.} But, though quite kindred in their object, these associations are not to be identified with the office of deaconess in the apostolic age and in the ancient church. That was a regular, standing office in every Christian congregation, corresponding to the office of deacon; and has never since the twelfth century been revived, though the local work of charity has never ceased.

To the ordinary clergy there were added in this period sundry extraordinary church offices, rendered necessary by the multiplication of religious functions in large cities and dioceses:

1. Stewards.\footnote{\textgreek{Οἰκόνομοι}. Besides these there were also \textgreek{κειμηλιάρχαι}, sacellarii, thesaurarii.} These officers administered the church property under the supervision of the bishop, and were chosen in part from the clergy, in part from such of the laity as were versed in law. In Constantinople the “great steward” was a person of considerable rank, though not a clergyman. The council of Chalcedon enjoined upon every episcopal diocese the appointment of such officers, and the selection of them from the clergy, “that the economy of the church might not be irresponsible, and thereby the church property be exposed to waste and the clerical dignity be brought into ill repute.”\footnote{Conc. Chalced. can. 26. This canon also occurs twice in the Corp. jur. can. c. 21, C. xvi. q. 7, and c. 4, Dist. lxxix.} For conducting the litigation of the church, sometimes a special advocate, called the e[kdiko”, or defensor, was appointed.

2. Secretaries,\footnote{·\textgreek{Ταχυγράφοι}, notarii, excerptores.} for drawing the protocols in public ecclesiastical transactions (gesta ecclesiastica). They were usually clergymen, or such as had prepared themselves for the service of the church.

3. Nurses or Parabolani,\footnote{\emph{Parabolani}, probably from \textgreek{παραβάλλειν τὴν ζωήν}, to risk life; because in contagious diseases they often exposed themselves to the danger of death.} especially in connection with the larger church hospitals. Their office was akin to that of the deacons, but had more reference to the bodily assistance than to the spiritual care of the sick. In Alexandria, by the fifth century, these officers formed a great guild of six hundred members, and were not rarely misemployed as a standing army of episcopal domination.\footnote{A perversion of a benevolent association to turbulent purposes similar to that of the firemen’s companies in the large cities of the United States.} Hence, upon a complaint of the citizens of Alexandria against them, to the emperor Theodosius II., their number were reduced to five hundred. In the West they were never introduced.

4. Buriers of the Dead\footnote{78 \textgreek{Κοπιάται}, copiattae, fossores, fossarii.} likewise belonged among these ordines minores of the church. Under Theodosius II. there were more than a thousand of them in Constantinople.

\xsection{The Bishops}

§ 53. The Bishops.

The bishops now stood with sovereign power at the head of the clergy and of their dioceses. They had come to be universally regarded as the vehicles and propagators of the gifts of the Holy Ghost, and the teachers and lawgivers of the church in all matters of faith and discipline. The specific distinction between them and the presbyters was carried into everything; while yet it is worthy of remark, that Jerome, Chrysostom, and Theodoret, just the most eminent exegetes of the ancient church, expressly acknowledged the original identity of the two offices in the New Testament, and consequently derive the proper episcopate, not from divine institution, but only from church usage.\footnote{See the passages quoted in § 52, and the works there referred to. The modern Romish divine, Perrone, in his Praelectiones Theologicae, t. ix. § 93, denies that the doctrine of the superiority of bishops over presbyters by \emph{divine right,} is an article of the Catholic faith. But the council of Trent, sess. xxiii. can. 6, condemns all who deny the divine institution of the three orders.}

The traditional participation of the people in the election, which attested the popular origin of the episcopal office, still continued, but gradually sank to a mere formality, and at last became entirely extinct. The bishops filled their own vacancies, and elected and ordained the clergy. Besides ordination, as the medium for communicating the official gifts, they also claimed from the presbyters in the West, after the fifth century, the exclusive prerogatives of confirming the baptized and consecrating the chrism or holy ointment used in baptism.\footnote{Innocent I., Ep. ad Decent.: “Ut sine chrismate et episcopi jussione neque presbyter neque diaconus jus habeant baptizandi.”} In the East, on the contrary, confirmation (the chrism) is performed also by the presbyters, and, according to the ancient custom, immediately follows baptism.

To this spiritual preëminence of the bishops was now added, from the time of Constantine, a civil importance. Through the union of the church with the state, the bishops became at the same time state officials of weight, and enjoyed the various privileges which accrued to the church from this connection.\footnote{Comp. above, ch. iii. § 14-16.} They had thenceforth an independent and legally valid jurisdiction; they held supervision of the church estates, which were sometimes very considerable, and they had partial charge even of the city, property; they superintended the morals of the people, and even of the emperor; and they exerted influence upon the public legislation. They were exempt from civil jurisdiction, and could neither be brought as witnesses before a court nor be compelled to take an oath. Their dioceses grew larger, and their power and revenues increased. Dominus beatissimus(μακαριώτατος), sanctissimus(ἁγιώτατος), or reverendissimus, Beatitudo or Sanctitas tua, and similar high-sounding titles, passed into universal use. Kneeling, kissing of the hand, and like tokens of reverence, came to be shown them by all classes, up to the emperor himself. Chrysostom, at the end of the fourth century, says: “The heads of the empire (hyparchs) and the governors of provinces (toparchs) enjoy no such honor as the rulers of the church. They are first at court, in the society of ladies, in the houses of the great. No one has precedence of them.”

To this position corresponded the episcopal insignia, which from the fourth century became common: the ring, as the symbol of the espousal of the bishop to the church; the crosier or shepherd’s staff (also called crook, because it was generally curved at the top); and the pallium,\footnote{2 \textgreek{Ἱερὰ στολή, ὡμοφόριον}, superhumerale, pallium, also ephod (\texthebrew{רוֹבאֵ}, \textgreek{ἐπωμίς}). The ephod (Ex. xxviii. 6-11; and xxxix. 2-5), in connection with the square breastplate belonging to it (\texthebrew{ןשֶׁח}, comp. Ex. xxviii. 15-30; xxxix. 8-21), was the principal official vestment of the Jewish high-priest, and no doubt served as the precedent for the archiepiscopal pallium, but exceeded the latter in costliness. It consisted of two shoulder pieces (like the pallium and the chasubles), which hung over the upper part of the body before and behind, and were skilfully wrought of fine linen in three colors, fastened by golden rings and chains, and richly ornamented with gold thread, and twelve precious stones, on which the names of the twelve tribes were graven. Whether the sacred oracle, Urim and Thummim (LXX.: \textgreek{δήλωσις καὶ ἀλήθεια}, Ex. xxviii. 30), was identical with the twelve precious stones in the breastplate, the learned are not agreed. Comp. Winer, Bibl. Reallex., and W. Smith, Dictionary of the Bible, sub \emph{Urim and Thummim.}}, a shoulder cloth, after the example of the ephod of the Jewish high-priest, and perhaps of the sacerdotal mantle worn by the Roman emperors as pontifices maximi. The pallium is a seamless cloth hanging over the shoulders, formerly of white linen, in the West subsequently of white lamb’s wool, with four red or black crosses wrought in it with silk. According to the present usage of the Roman church the wool is taken from the lambs of St. Agnes, which are every year solemnly blessed and sacrificed by the pope in memory of this pure virgin. Hence the later symbolical meaning of the pallium, as denoting the bishop’s following of Christ, the good Shepherd, with the lost and reclaimed sheep upon his shoulders. Alexandrian tradition traced this vestment to the evangelist Mark; but Gregory Nazianzen expressly says that it was first given by Constantine the Great to the bishop Macarius of Jerusalem.\footnote{Orat. xlvii. So Theodoret, Hist. eccl. ii. 27, at the beginning. Macarius is said to have worn the gilded vestment in the administration of baptism.} In the East it was worn by all bishops, in the West by archbishops only, on whom, from the time of Gregory I., it was conferred by the pope on their accession to office. At first the investiture was gratuitous, but afterward came to involve a considerable fee, according to the revenues of the archbishopric.

As the bishop united in himself all the rights and privileges of the clerical office, so he was expected to show himself a model in the discharge of its duties and a follower of the great Archbishop and Archshepherd of the sheep. He was expected to exhibit in a high degree the ascetic virtues, especially that of virginity, which, according to Catholic ethics, belongs to the idea of moral perfection. Many a bishop, like Athanasius, Basil, Ambrose, Augustine, Chrysostom, Martin of Tours, lived in rigid abstinence and poverty, and devoted his income to religious and charitable objects.

But this very power and this temporal advantage of the episcopate became also a lure for avarice and ambition, and a temptation to the lordly and secular spirit. For even under the episcopal mantle the human heart still beat, with all those weaknesses and passions, which can only be overcome by the continual influence of Divine grace. There were metropolitans and patriarchs, especially in Alexandria, Constantinople, and Rome, who, while yet hardly past the age of persecution, forgot the servant form of the Son of God and the poverty of his apostles and martyrs, and rivalled the most exalted civil officials, nay, the emperor himself, in worldly pomp and luxury. Not seldom were the most disgraceful intrigues employed to gain the holy office. No wonder, says Ammianus, that for so splendid a prize as the bishopric of Rome, men strive with the utmost passion and persistence, when rich presents from ladies and a more than imperial sumptuousness invite them.\footnote{Amm. Marcell. xxvii. c. 3, sub anno 367: “ut dotentur oblationibus matronarum procedantque vehiculis insidentes, circumspecte vestiti, epulas curantes profusas, adeo ut eorum convivia regales superent mensas.” But then with this pomp of the Roman prelates he contrasts the poverty of the worthy country bishops.} The Roman prefect, Praetextatus, declared jestingly to the bishop Damasus, who had obtained the office through a bloody battle of parties, that for such a price he would at once turn Christian himself.\footnote{Besides Ammianus, Jerome also states this, in his book against John of Jerusalem (Opera, tom. ii. p. 415, ed. Vallars.): “Miserabilis ille Praetextatus, qui designatus consul est mortuus, homo sacrilegus et idolorum cultor, solebat ludens beato papae Damaso dicere: ’Facite me Romanae urbis episcopum, et ero protinus Christianus.’ “} Such an example could not but shed its evil influence on the lower clergy of the great cities. Jerome sketches a sarcastic description of the Roman priests, who squandered all their care on dress and perfumery, curled their hair with crisping pins, wore sparkling rings, paid far too great attention to women, and looked more like bridegrooms than like clergymen.\footnote{Epist. ad Eustochium de virginitate servanda.} And in the Greek church it was little better. Gregory Nazianzen, himself a bishop, and for a long time patriarch of Constantinople, frequently mourns the ambition, the official jealousies, and the luxury of the hierarchy, and utters the wish that the bishops might be distinguished only by a higher grade of virtue.

\xsection{Organization of the Hierarchy: Country Bishop, City Bishops, and Metropolitans}

§ 54. Organization of the Hierarchy: Country Bishop, City Bishops, and Metropolitans.

The episcopate, notwithstanding the unity of the office and its rights, admitted the different grades of country bishop, ordinary city bishop, metropolitan, and patriarch. Such a distinction had already established itself on the basis of free religious sentiment in the church; so that the incumbents of the apostolic sees, like Jerusalem, Antioch, Ephesus, Corinth, and Rome, stood at the head of the hierarchy. But this gradation now assumed a political character, and became both modified and confirmed by attachment to the municipal division of the Roman empire.

Constantine the Great divided the whole empire into four praefectures (the Oriental, the Illyrian, the Italian, and the Gallic); the praefectures into vicariates, dioceses, or proconsulates, fourteen or fifteen in all;\footnote{The dioceses or vicariates were as follows: \par I. The Praefectura Orientalisconsisted of the five dioceses of \emph{Oriens,} with Antioch as its political and ecclesiastical capital; \emph{Aegyptus}, with Alexandria; \emph{Asia} \emph{proconsularis,} with Ephesus; \emph{Pontus}, with Caesarea in Cappadocia; \emph{Thracia}, with Heraklea, afterward Constantinople. \par II. The Praefectura Illyrica, with Thessalonica as its capital, had only the two dioceses of \emph{Macedonia} and \emph{Dacia}. \par III. The Praefectura Italicaembraced \emph{Roma} (i.e. South Italy and the islands of the Mediterranean, or the so-called Suburban provinces); \emph{Italia}, or the Vicariate of Italy, with its centre at Mediolanum (Milan); \emph{Illyricum} \emph{occidentale,} with its capital at Sirmium; and \emph{Africa occidentalis,} with Carthage. \par IV. The Praefectura Gallicaembraced the dioceses of \emph{Gallia}, with Treveri (Trier) and Lugdunum (Lyons); \emph{Hispania}, with Hispalis (Sevilla); and \emph{Britannnia}, with Eboracum (York).} and each diocese again into several provinces.\footnote{Thus the diocese of the Orient, for example, had five provinces, Egypt nine, Pontus thirteen, Gaul seventeen, Spain seven. Comp. Wiltsch, Kirchl. Geogr. u. Statistik, i. p. 67 sqq., where the provinces are all quoted, as is not necessary for our purpose here.} The praefectures were governed by Praefecti Praetorio, the dioceses by Vicarii, the provinces by Rectores, with various titles—commonly Praesides.

It was natural, that after the union of church and state the ecclesiastical organization and the political should, so far as seemed proper, and hence of course with manifold exceptions, accommodate themselves to one another. In the East this principle of conformity was more palpably and rigidly carried out than in the West. The council of Nice in the fourth century proceeds upon it, and the second and fourth ecumenical councils confirm it. The political influence made itself most distinctly felt in the elevation of Constantinople to a patriarchal see. The Roman bishop Leo, however, protested against the reference of his own power to political considerations, and planted it exclusively upon the primacy of Peter; though evidently the Roman see owed its importance to the favorable cooperation of both these influences. The power of the patriarchs extended over one or more municipal dioceses; while the metropolitans presided over single provinces. The word diocese (διοίκησις) passed from the political into the ecclesiastical terminology, and denoted at first a patriarchal district, comprising several provinces (thus the expression occurs continually in the Greek acts of councils), but afterward came to be applied in the West to each episcopal district. The circuit of a metropolitan was called in the East an eparchy (ἐπαρχία), in the West provincia. An ordinary bishopric was called in the East a parish (παροικία), while in the Latin church the term (parochia) was usually applied to a mere pastoral charge.

The lowest rank in the episcopal hierarchy was occupied by the country bishops,\footnote{\textgreek{Χωρεπίσκοποι}. The principal statements respecting them are: Epist. Synodi Antioch., a.d.270, in Euseb. H. E. vii. 36 (where they are called \textgreek{ἐπίσκοποι τῶν ὁμόρων ἀγρῶν}); Concil. Ancyr., a.d.315, can. 13 (where they are forbidden to ordain presbyters and deacons); Concil. Antioch., a.d.341, can. 10 (same prohibition); Conc. Laodic., between 320 and 372, can. 57 (where the erection of new country bishoprics is forbidden); and Conc. Sardic., a.d.343, can. 6 (where they are wholly abolished).} the presiding officers of those rural congregations, which were not supplied with presbyters from neighboring cities. In North Africa, with its multitude of small dioceses, these country bishops were very numerous, and stood on an equal footing with the others. But in the East they became more and more subordinate to the neighboring city bishops; until at last, partly on account of their own incompetence, chiefly for the sake of the rising hierarchy, they were wholly extinguished. Often they were utterly unfit for their office; at least Basil of Caesarea, who had fifty country bishops in his metropolitan district, reproached them with frequently receiving men totally unworthy into the clerical ranks. And moreover, they stood in the way of the aspirations of the city bishops; for the greater the number of bishops, the smaller the diocese and the power of each, though probably the better the collective influence of all upon the church. The council of Sardica, in 343, doubtless had both considerations in view, when, on motion of Hosius, the president, it decreed: “It is not permitted, that, in a village or small town, for which a single priest is sufficient, a bishop should be stationed, lest the episcopal dignity and authority suffer scandal;\footnote{Can. 6: ... \textgreek{ἲνα μὴ κατευτελίξηται τὸ τοῦ ἐπισκόπου ὄνομα καὶ ἡ αὐθεντία}; or, in the Latin version: “Ne vilescat nomen episcopi et auctoritas.” Comp. Hefele, i. p. 556. The differences between the Greek and Latin text in the first part of this canon have no influence on the prohibition of the appointment of country bishops.} but the bishops of the eparchy (province) shall appoint bishops only for those places where bishops have already been, or where the town is so populous that it is considered worthy to be a bishopric.” The place of these chorepiscopi was thenceforth supplied either by visitators (περιοδεῦται), who in the name of the bishop visited the country congregations from time to time, and performed the necessary functions, or by resident presbyters (parochi), under the immediate supervision of the city bishop.

Among the city bishops towered the bishops of the capital cities of the various provinces. They were styled in the East metropolitans, in the West usually archbishops.\footnote{\textgreek{Μητροπολίτης}, metropolitanus, and the kindred title \textgreek{ἔξαρχος} (applied to the most powerful metropolitans); \textgreek{ἀρχιεπίσκοπος}, archiepiscopus, and primas.} They had the oversight of the other bishops of the province; ordained them, in connection with two or three assistants; summoned provincial synods, which, according to the fifth canon of the council of Nice and the direction of other councils, were to be held twice a year; and presided in such synods. They promoted union among the different churches by the reciprocal communication of synodal acts, and confirmed the organism of the hierarchy.

This metropolitan constitution, which had gradually arisen out of the necessities of the church, became legally established in the East in the fourth century, and passed thence to the Graeco-Russian church. The council of Nice, at that early day, ordered in the fourth canon, that every new bishop should be ordained by all, or at least by three, of the bishops of the eparchy (the municipal province), under the direction and with the sanction of the metropolitan.\footnote{This canon has been recently discovered also in a Coptic translation, and published by Pitra, in the Spiclegium Solesmense, i. 526 sq.} Still clearer is the ninth canon of the council of Antioch, in 341: “The bishops of each eparchy (province) should know, that upon the bishop of the metropolis (the municipal capital) also devolves a care for the whole eparchy, because in the metropolis all, who have business, gather together from all quarters. Hence it has been found good, that he should also have a precedence in honor,\footnote{\textgreek{Καὶ τῇ τιμῇ προηγεῖσθαι αυτόν}.} and that the other bishops should do nothing without him—according to the old and still binding canon of our fathers—except that which pertains to the supervision and jurisdiction of their parishes (i.e. dioceses in the modern terminology), and the provinces belonging to them; as in fact they ordain presbyters and deacons, and decide all judicial matters. Otherwise they ought to do nothing without the bishop of the metropolis, and he nothing without the consent of the other bishops.” This council, in the nineteenth canon, forbade a bishop being ordained without the presence of the metropolitan and the presence or concurrence of the majority of the bishops of the province.

In Africa a similar system had existed from the time of Cyprian, before the church and the state were united. Every province had a Primas; the oldest bishop being usually chosen to this office. The bishop of Carthage, however, was not only primate of Africa proconsularis, but at the same time, corresponding to the proconsul of Carthage, the ecclesiastical head of Numidia and Mauretania, and had power to summon a general council of Africa.\footnote{Cyprian, Epist. 45, says of his province of Carthage: “Latius fusa est nostra provincia; habet enim Numidiam et Mauretaniam sibi cohaerentes.”}

\xsection{The Patriarchs}

§ 55. The Patriarchs.

Mich. Le Quien (French Dominican, † 1788): Oriens Christianus, in quatuor patriarchatus digestus, quo exhibentur ecclesiae, patriarchae caeterique preasules totius Orientis. Opus posthumum, Par. 1740, 3 vols. fol. (a thorough description of the oriental dioceses from the beginning to 1732). P. Jos. Cautelius (Jesuit): Metropolitanarum urbium historia civilis et ecclesiastic in qua Romanae Sedis dignitas et imperatorum et regum in eam merits explicantur, Par. 1685 (important for ecclesiastical statistics of the West, and the extension of the Roman patriarchate). Bingham (Anglican): Antiquities, l. ii. c. 17. Joh. El. Theod. Wiltsch (Evangel.): Handbuch der Kirchl. Geographie u. Statistik, Berl. 1846, vol. i. p. 56 sqq. Friedr. Maassen (R.C.): Der Primat des Bischofs von Rom. u. die alten Patriarchalkirchen, Bonn, 1853. Thomas Greenwood: Cathedra Petri, a Political History of the Latin Patriarchate, Lond. 1859 sqq. (vol. i. p. 158–489). Comp. my review of this work in the Am. Theol. Rev., New York, 1864, p. 9 sqq.

Still above the metropolitans stood the five Patriarchs,\footnote{\textgreek{Πατριάρχης}; patriarcha; sometimes also, after the political terminology, \textgreek{ἔξαρχος}. The name patriarch, originally applied to the progenitors of Israel (Heb. vii. 4, to Abraham; Acts vii. 8 sq., to the twelve sons of Jacob; ii. 29, to David, as founder of the Davidic Messianic house), was at first in the Eastern church an honorary title for bishops in general (so in Gregory Nazianzen, and Gregory of Nyssa), but after the council of Constantinople (381), and still more after that of Chalcedon (451), it came to be used in an official sense and restricted to the five most eminent metropolitans. In the West, several metropolitans, especially the bishop of Aquileia, bore this title \emph{honoris causa.} The bishop of Rome declined that particular term, as placing him on a level with other patriarchs, and preferred the name \emph{papa}. “Patriarch” bespeaks an oligarchical church government; “pope,” a monarchical.} the oligarchical summit, so to speak, the five towers in the edifice of the Catholic hierarchy of the Graeco-Roman empire.

These patriarchs, in the official sense of the word as already fixed at the time of the fourth ecumenical council, were the bishops of the four great capitals of the empire, Rome, Alexandria, Antioch, and Constantinople; to whom was added, by way of honorary distinction, the bishop of Jerusalem, as president of the oldest Christian congregation, though the proper continuity of that office had been broken by the destruction of the holy city. They had oversight of one or more dioceses; at least of two or more provinces or eparchies.\footnote{According to the political division of the empire after Constantine. Comp. § 54} They ordained the metropolitans; rendered the final decision in church controversies; conducted the ecumenical councils; published the decrees of the councils and the church laws of the emperors; and united in themselves the supreme legislative and executive power of the hierarchy. They bore the same relation to the metropolitans of single provinces, as the ecumenical councils to the provincial. They did not, however, form a college; each acted for himself. Yet in important matters they consulted with one another, and had the right also to keep resident legates (apocrisiarii) at the imperial court at Constantinople.

In prerogative they were equal, but in the extent of their dioceses and in influence they differed, and had a system of rank among themselves. Before the founding of Constantinople, and down to the Nicene council, Rome maintained the first rank, Alexandria the second, and Antioch the third, in both ecclesiastical and political importance. After the end of the fourth century this order was modified by the insertion of Constantinople as the second capital, between Rome and Alexandria, and the addition of Jerusalem as the fifth and smallest patriarchate.

The patriarch of Jerusalem presided only over the three meagre provinces of Palestine;\footnote{Comp. Wiltsch, i. p. 206 sqq. The statement of Ziegler, which Wiltsch quotes and seems to approve, that the fifth ecumenical council, of 553, added to the patriarchal circuit of Jerusalem the metropolitans of Berytus in Phenicia, and Ruba in Syria, appears to be an error. Ruba nowhere appears in the acts of the council, and Berytus belonged to Phoenicia prima, consequently to the patriarchate of Antioch. Le Quien knows nothing of such an enlargement of the patriarchate of Hierosolyma.} the patriarch of Antioch, over the greater part of the political diocese of the Orient, which comprised fifteen provinces, Syria, Phenicia, Cilicia, Arabia, Mesopotamia, \&c.;\footnote{Wiltsch, i. 189 sqq.} the patriarch of Alexandria, over the whole diocese of Egypt with its nine rich provinces, Aegyptus prima and secunda, the lower and upper Thebaid, lower and upper Libya, \&c.;\footnote{Ibid. i. 177 sqq.} the patriarch of Constantinople, over three dioceses, Pontus, Asia Minor, and Thrace, with eight and twenty provinces, and at the same time over the bishoprics among the barbarians;\footnote{Ibid. p. 143 sqq.} the patriarch of Rome gradually extended his influence over the entire West, two prefectures, the Italian and the Gallic, with all their dioceses and provinces.\footnote{Comp. § 57, below.}

The patriarchal system had reference primarily only to the imperial church, but indirectly affected also the barbarians, who received Christianity from the empire. Yet even within the empire, several metropolitans, especially the bishop of Cyprus in the Eastern church, and the bishops of Milan, Aquileia, and Ravenna in the Western, during this period maintained their autocracy with reference to the patriarchs to whose dioceses they geographically belonged. In the fifth century, the patriarchs of Antioch attempted to subject the island of Cyprus, where Paul first had preached the gospel, to their jurisdiction; but the ecumenical council of Ephesus, in 431, confirmed to the church of Cyprus its ancient right to ordain its own bishops.\footnote{Comp. Wiltsch, i. p. 232 sq., and ii. 469.} The North African bishops also, with all respect for the Roman see, long maintained Cyprian’s spirit of independence, and in a council at Hippo Regius, in 393, protested against such titles as princeps sacerdotum, summus sacerdos, assumed by the patriarchs, and were willing only to allow the title of primae sedis episcopus.\footnote{Cod. can. eccl. Afr. can. 39, cited by Neander, iii. p. 335 (Germ. ed.).}

When, in consequence of the Christological controversies, the Nestorians and Monophysites split off from the orthodox church, they established independent schismatic patriarchates, which continue to this day, showing that the patriarchal constitution answers most nearly to the oriental type of Christianity. The orthodox Greek church, as well as the schismatic sects of the East, has substantially remained true to the patriarchal system down to the present time; while the Latin church endeavored to establish the principle of monarchical centralization so early as Leo the Great, and in the course of the middle age produced the absolute papacy.

\xsection{Synodical Legislation on the Patriarchal Power and Jurisdiction}

§ 56. Synodical Legislation on the Patriarchal Power and Jurisdiction.

To follow now the ecclesiastical legislation respecting this patriarchal oligarchy in chronological order:

The germs of it already lay in the ante-Nicene period, when the bishops of Antioch, Alexandria, and Rome, partly in virtue of the age and apostolic origin of their churches, partly, on account of the political prominence of those three cities as the three capitals of the Roman empire, steadily asserted a position of preëminence. The apostolic origin of the churches of Rome and Antioch is evident from the New Testament: Alexandria traced its Christianity, at least indirectly through the evangelist Mark, to Peter, and was politically more important than Antioch; while Rome from the first had precedence of both in church and in state. This preëminence of the oldest and most powerful metropolitans acquired formal legislative validity and firm establishment through the ecumenical councils of the fourth and fifth centuries.

The first ecumenical council of Nice, in 325, as yet knew nothing of five patriarchs, but only the three metropolitans above named, confirming them in their traditional rights.\footnote{Accordingly Pope Nicolas, in 866, in a letter to the Bulgarian prince Bogoris, would acknowledge only the bishops of Rome, Alexandria, and Antioch as patriarchs in the proper sense, because they presided over apostolic churches; whereas Constantinople was not of apostolic founding, and was not even mentioned by the most venerable of all councils, the Nicene; Jerusalem was named indeed by these councils, but only under the name of Aelia.} In the much-canvassed sixth canon, probably on occasion of the Meletian schism in Egypt, and the attacks connected with it on the rights of the bishop of Alexandria, that council declared as follows:

“The ancient custom, which has obtained in Egypt, Libya, and the Pentapolis, shall continue in force, viz.: that the bishop of Alexandria have rule over all these [provinces], since this also is customary with the bishop of Rome [that is, not in Egypt, but with reference to his own diocese]. Likewise also at Antioch and in the other eparchies, the churches shall retain their prerogatives. Now, it is perfectly clear, that, if any one has been made bishop without the consent of the metropolitan, the great council does not allow him to be bishop.”\footnote{In the oldest Latin Cod. canonum (in Mansi, vi. 1186) this canon is preceded by the important words: \emph{Ecclesia Romana semper habuit primatum}. These are, however, manifestly spurious, being originally no part of the canon itself, but a superscription, which gave an expression to the Roman inference from the Nicene canon. Comp. Gieseler, i. 2, § 93, note 1; and Hefele, Hist. of Councils, i. 384 sqq.}

The Nicene fathers passed this canon not as introducing anything new, but merely as confirming an existing relation on the basis of church tradition; and that, with special reference to Alexandria, on account of the troubles existing there. Rome was named only for illustration; and Antioch and all the other eparchies or provinces were secured their admitted rights.\footnote{So Greenwood also views the matter, Cathedra Petri, 1859, vol. i. p. 181: “It was manifestly not the object of this canon to confer any new jurisdiction upon the church of Alexandria, but simply to confirm its customary prerogative. By way of illustration, it places that prerogative, whatever it was, upon the same level with that of the two other eparchal churches of Rome and Antioch. Moreover, the words of the canon disclose no other ground of claim but custom; and the customs of each eparchia are restricted to the territorial limits of the diocese or eparchia itself. And though, within those limits, the several customary rights and prerogatives may have differed, yet beyond them no jurisdiction of any kind could, by virtue of this canon, have any existence at all.”} The bishoprics of Alexandria, Rome, and Antioch were placed substantially on equal footing, yet in such tone, that Antioch, as the third capital of the Roman empire, already stands as a stepping stone to the ordinary metropolitans. By the “other eparchies” of the canon are to be understood either all provinces, and therefore all metropolitan districts, or more probably, as in the second canon of the first council of Constantinople, only the three eparchates of Caesarea in Cappadocia, Ephesus, and Asia Minor, and Heraclea in Thrace, which, after Constantine’s division of the East, possessed similar prerogatives, but were subsequently overshadowed and absorbed by Constantinople. In any case, however, this addition proves that at that time the rights and dignity of the patriarchs were not yet strictly distinguished from those of the other metropolitans. The bishops of Rome, Alexandria, and Antioch here appear in relation to the other bishops simply as primi inter pares, or as metropolitans of the first rank, in whom the highest political eminence was joined with the highest ecclesiastical. Next to them, in the second rank, come the bishops of Ephesus in the Asiatic diocese of the empire, of Neo-Caesarea in the Pontic, and of Heraclea in the Thracian; while Constantinople, which was not founded till five years later, is wholly unnoticed in the Nicene council, and Jerusalem is mentioned only under the name of Aelia.

Between the first and second ecumenical councils arose the new patriarchate of Constantinople, or New Rome, built by Constantine in 330, and elevated to the rank of the imperial residence. The bishop of this city was not only the successor of the bishop of the ancient Byzantium, hitherto under the jurisdiction of the metropolitan of Heraclea, but, through the favor of the imperial court and the bishops who were always numerously assembled there, it placed itself in a few decennia among the first metropolitans of the East, and in the fifth century became the most powerful rival of the bishop of old Rome.

This new patriarchate was first officially recognized at the first ecumenical council, held at Constantinople in 381, and was conceded “the precedence in honor, next to the bishop of Rome,” the second place among all bishops; and that, on the purely political consideration, that New Rome was the residence of the emperor.\footnote{Conc. Constant. i. can 3: \textgreek{Τὸν μέντοι Κώσταντινουπόλεως ἐπίσκοπον ἔχειν τὰ πρεσβεῖα τῆς τιμῆς, μετὰ τὸν τῆσ Ῥώμης ἐπίσκοπον, διὰ τὸ εἶναι αὐτὴν νέαν Ῥώμην}. This canon is quoted also by Socrates, v. 8, and Sozomen, vii. 9, and confirmed by the council of Chalcedon (see below); so that it must be from pure dogmatical bias, that Baronius (Annal. ad ann. 381, n. 35, 36) questions its genuineness} At the same time the imperial city and the diocese of Thrace (whose ecclesiastical metropolis hitherto had been Heraclea) were assigned as its district.\footnote{The latter is not, indeed, expressly said in the above canon, which seems to speak only of an honorary precedence. But the canon was so understood by the bishops of Constantinople, and by the historians Socrates (v. 8) and Theodoret (Epist. 86, ad Flavianum), and so interpreted by the Chalcedonian council (can. 28). The relation of the bishop of Constantinople to the metropolitan of Heraclea, however, remained for a long time uncertain, and at the council ad Quercum, 403, in the affair of Chrysostom, Paul of Heraclea took the presidency, though the patriarch Theophilus of Alexandria was present. Comp. Le Quien, tom. i. p. 18; and Wiltsch, i. p. 139.}

Many Greeks took this as a formal assertion of the equality of the bishop of Constantinople with the bishop of Rome, understanding “next” or “after” (metav) as referring only to time, not to rank. But it is more natural to regard this as conceding a primacy of honor, which the Roman see could claim on different grounds. The popes, as the subsequent protest of Leo shows, were not satisfied with this, because they were unwilling to be placed in the same category with the Constantinopolitan fledgling, and at the same time assumed a supremacy of jurisdiction over the whole church. On the other hand, this decree was unwelcome also to the patriarch of Alexandria, because this see had hitherto held the second rank, and was now required to take the third. Hence the canon was not subscribed by Timotheus of Alexandria, and was regarded in Egypt as void. Afterward, however, the emperors prevailed with the Alexandrian patriarchs to yield this point.

After the council of 381, the bishop of Constantinople indulged in manifold encroachments on the rights of the metropolitans of Ephesus and Caesarea in Cappadocia, and even on the rights of the other patriarchs. In this extension of his authority he was favored by the fact that, in spite of the prohibition of the council of Sardica, the bishops of all the districts of the East continually resided in Constantinople, in order to present all kinds of interests to the emperor. These concerns of distant bishops were generally referred by the emperor to the bishop of Constantinople and his council, the σύνοδος ἐνδημοῦσα, as it was called, that is, a council of the bishops resident (ἐνδημούντων) in Constantinople, under his presidency. In this way his trespasses even upon the bounds of other patriarchs obtained the right of custom by consent of parties, if not the sanction of church legislation. Nectarius, who was not elected till after that council, claimed the presidency at a council in 394, over the two patriarchs who were present, Theophilus of Alexandria and Flavian of Antioch; decided the matter almost alone; and thus was the first to exercise the primacy over the entire East. Under his successor, Chrysostom, the compass of the see extended itself still farther, and, according to Theodoret,\footnote{H. E. lib. v. cap. 28.} stretched over the capital, over all Thrace with its six provinces, over all Asia (Asia proconsularis) with eleven provinces, and over Pontus, which likewise embraced eleven provinces; thus covering twenty-eight provinces in all. In the year 400, Chrysostom went “by request to Ephesus,” to ordain there Heraclides of Ephesus, and at the same time to institute six bishops in the places of others deposed for simony.\footnote{According to Sozomen it was thirteen, according to Theophilus of Alexandria at the council ad Quercum seventeen bishops, whom he instituted; and this act was charged against him as an unheard-of crime. See Wiltsch, i. 141.} His second successor, Atticus, about the year 421, procured from the younger Theodosius a law, that no bishop should be ordained in the neighboring dioceses without the consent of the bishop of Constantinople.\footnote{Socrates, H. E. l. vii. 28, where such a law is incidentally mentioned. The inhabitants of Cyzicus in the Hellespont, however, transgressed the law, on the presumption that it was merely a personal privilege of Atticus.} This power still needed the solemn sanction of a general council, before it could have a firm legal foundation. It received this sanction at Chalcedon.

The fourth ecumenical council, held at Chalcedon in 451 confirmed and extended the power of the bishop of Constantinople, by ordaining in the celebrated twenty-eighth canon:

“Following throughout the decrees of the holy fathers, and being “acquainted with the recently read canon of the hundred and fifty bishops [i.e. the third canon of the second ecumenical council of 381], we also have determined and decreed the same in reference to the prerogatives of the most holy church of Constantinople or New Rome. For with reason did the fathers confer prerogatives (τὰ πρεσβεῖα) on the throne [the episcopal chair] of ancient Rome, on account of her character as the imperial city (διὰ τὸ βασιλεύειν); and, moved by the same consideration, the hundred and fifty bishops recognized the same prerogatives (τὰ ἴσα ςπρεσβεῖα) also in the most holy throne of New Rome; with good reason judging, that the city, which is honored with the imperial dignity and the senate [i.e. where the emperor and senate reside], and enjoys the same [municipal] privileges as the ancient imperial Rome, should also be equally elevated in ecclesiastical respects, and be the second after he(δευτέραν μετ  ἐκείνην.].ς

“And [we decree] that of the dioceses of Pontus, Asia [Asia proconsularis], and Thrace, only the metropolitans, but in such districts of those dioceses as are occupied by barbarians, also the [ordinary] bishops, be ordained by the most holy throne of the most holy church at Constantinople; while of course every metropolitan in those dioceses ordains the new bishops of a province in concurrence with the existing bishops of that province, as is directed in the divine (θείοις) canons. But the metropolitans of those dioceses, as already said, shall be ordained by the archbishop (ἀρχιεπισκόπου) of Constantinople, after they shall have been unanimously elected in the usual way, and he [the archbishop of Constantinople] shall have been informed of it.”

We have divided this celebrated Chalcedonian canon into two parts, though in the Greek text the parts are (by Καὶ ωὝστε) closely connected. The first part assigns to the bishop of Constantinople the second rank among the patriarchs, and is simply a repetition and confirmation of the third canon of the council of Constantinople; the second part goes farther, and sanctions the supremacy, already actually exercised by Chrysostom and his successors, of the patriarch of Constantinople, not only over the diocese of Thrace, but also over the dioceses of Asia Minor and Pontus, and gives him the exclusive right to ordain both the metropolitans of these three dioceses, and all the bishops of the barbarians\footnote{Among the barbarian tribes, over whom the bishops of Constantinople exercised an ecclesiastical jurisdiction, were the Huns on the Bosphorus, whose king, Gorda, received baptism in the time of Justinian; the Herulians, who received the Christian faith in 527; the Abasgians and Alanians on the Euxine sea, who about the same time received priests from Constantinople. Comp. Wiltsch, i. 144 and 145.} within those bounds. This gave him a larger district than any other patriarch of the East. Subsequently an edict of the emperor Justinian, in 530, added to him the special prerogative of receiving appeals from the other patriarchs, and thus of governing the whole Orient.

The council of Chalcedon in this decree only followed consistently the oriental principle of politico-ecclesiastical division. Its intention was to make the new political capital also the ecclesiastical capital of the East, to advance its bishop over the bishops of Alexandria and Antioch, and to make him as nearly as possible equal to the bishop of Rome. Thus was imposed a wholesome check on the ambition of the Alexandrian patriarch, who in various ways, as the affair of Theophilus and Dioscurus shows, had abused his power to the prejudice of the church.

But thus, at the same time, was roused the jealousy of the bishop of Rome, to whom a rival in Constantinople, with equal prerogatives, was far more dangerous than a rival in Alexandria or Antioch. Especially offensive must it have been to him, that the council of Chalcedon said not a word of the primacy of Peter, and based the power of the Roman bishop, like that of the Constantinopolitan, on political grounds; which was indeed not erroneous, yet only half of the truth, and in that respect unfair.

Just here, therefore, is the point, where the Eastern church entered into a conflict with the Western, which continues to this day. The papal delegates protested against the twenty-eighth canon of the Chalcedonian council, on the spot, in the sixteenth and last session of the council; but in vain, though their protest was admitted to record. They appealed to the sixth canon of the Nicene council, according to the enlarged Latin version, which, in the later addition, “Ecclesia Romana semper habuit primatum,” seems to assign the Roman bishop a position above all the patriarchs, and drops Constantinople from notice; whereupon the canon was read to them in its original form from the Greek Acts, without that addition, together with the first three canons of the second ecumenical council with their express acknowledgment of the patriarch of Constantinople in the second rank.\footnote{This correction of the Roman legates is so little to the taste of the Roman Catholic historians, especially the ultramontane, that the Ballerini, in their edition of the works of Leo the Great, tom. iii. p. xxxvii. sqq., and even Hefele, Conciliengesch. i. p. 385, and ii. p. 522, have without proof declared the relevant passage in the Greek Acts of the council of Chalcedon a later interpolation. Hefele, who can but concede the departure of the Latin version from the original text of the sixth canon of Nice, thinks, however, that the Greek text was not read in Chalcedon, because even this bore \emph{against} the elevation of Constantinople, and therefore \emph{in} \emph{favor} of the Roman legates. But the Roman legates, as also Leo in his protest against the 28th decree of Chalcedon, laid chief stress upon the Roman addition, \emph{Ecclesia Romana semper habuit primatum}, and considered the equalization of any other patriarch with the bishop of Rome incompatible with it. Since the legates, as is conceded, appealed to the Nicene canon, the Greeks had first to meet this appeal, before they passed to the canons of the council of Constantinople. Only the two together formed a sufficient answer to the Roman protest.} After the debate on this point, the imperial commissioners thus summed up the result: “From the whole discussion, and from what has been brought forward on either side, we acknowledge that the primacy over all (πρὸ πάντων τὰ πρωτεῖα) and the most eminent rank (καὶ τὴν ἐξείρετον τιμήν) are to continue with the archbishop of old Rome; but that also the archbishop of New Rome should enjoy the same precedence of honor (τὰ πρεσβεῖα τῆς τιμῆς), and have the right to ordain the metropolitans in the dioceses of Asia, Pontus, and Thrace,” \&c. Now they called upon the council to declare whether this was its opinion; whereupon the bishops gave their full, emphatic consent, and begged to be dismissed. The commissioners then closed the transactions with the words: “What we a little while ago proposed, the whole council hath ratified;” that is, the prerogative granted to the church of Constantinople is confirmed by the council in spite of the protest of the legates of Rome.\footnote{Mansi, vii. p. 446-454; Harduin, ii, 639-643; Hefele, ii. 524, 525.}

After the council, the Roman bishop, Leo, himself protested in three letters of the 22d May, 452; the first of which was addressed to the emperor Marcian, the second to the empress Pulcheria, the third to Anatolius, patriarch of Constantinople.\footnote{Leo, Epist. 104, 105, and 106 (al, Ep. 78-80). Comp. Hefele, l.c. ii. 530 sqq.} He expressed his satisfaction with the doctrinal results of the council, but declared the elevation of the bishop of Constantinople to the patriarchal dignity to be a work of pride and ambition—the humble, modest pope!—to be an attack upon the rights of other Eastern metropolitans—the invader of the same rights in Gaul!—especially upon the rights of the Roman see guaranteed by the council of Nice—on the authority of a Roman interpolation—and to be destructive of the peace of the church—which the popes have always sacredly kept! He would hear nothing of political considerations as the source of the authority of his chair, but pointed rather to Divine institution and the primacy of Peter. Leo speaks here with great reverence of the first ecumenical council, under the false impression that that council in its sixth canon acknowledged the primacy of Rome; but with singular indifference of the second ecumenical council, on account of its third canon, which was confirmed at Chalcedon. He charges Anatolius with using for his own ambition a council, which had been called simply for the extermination of heresy and the establishment of the faith. But the canons of the Nicene council, inspired by the Holy Ghost, could be superseded by no synod, however great; and all that came in conflict with them was void. He exhorted Anatolius to give up his ambition, and reminded him of the words: Tene quod habes, ne alius accipiat coronam tuam.\footnote{Rev. iii. 11.}

But this protest could not change the decree of the council nor the position of the Greek church in the matter, although, under the influence of the emperor, Anatolius wrote an humble letter to Leo. The bishops of Constantinople asserted their rank, and were sustained by the Byzantine emperors. The twenty-eighth canon of the Chalcedonian council was expressly confirmed by Justinian I., in the 131st Novelle (c. 1), and solemnly renewed by the Trullan council (can. 36), but was omitted in the Latin collections of canons by Prisca, Dionysius, Exiguus, and Isidore. The loud contradiction of Rome gradually died away; yet she has never formally acknowledged this canon, except during the Latin empire and the Latin patriarchate at Constantinople, when the fourth Lateran council, under Innocent III., in 1215, conceded that the patriarch of Constantinople should hold the next rank after the patriarch of Rome, before those of Alexandria and Antioch.\footnote{Harduin, tom. vii. 23; Schröckh, xvii. 43; and Hefele, ii. 544.}

Finally, the bishop of Jerusalem, after long contests with the metropolitan of Caesarea and the patriarch of Antioch, succeeded in advancing himself to the patriarchal dignity; but his distinction remained chiefly a matter of honor, far below the other patriarchates in extent of real power. Had not the ancient Jerusalem, in the year 70, been left with only a part of the city wall and three gates to mark it, it would doubtless, being the seat of the oldest Christian congregation, have held, as in the time of James, a central position in the hierarchy. Yet as it was, a reflection of the original dignity of the mother city fell upon the new settlement of Aelia Capitolina, which, after Adrian, rose upon the venerable ruins. The pilgrimage of the empress Helena, and the magnificent church edifices of her son on the holy places, gave Jerusalem a new importance as the centre of devout pilgrimage from all quarters of Christendom. Its bishop was subordinate, indeed, to the metropolitan of Caesarea, but presided with him (probably secundo loco) at the Palestinian councils.\footnote{Comp. Eusebius, himself the metropolitan of Caesarea, H. E. v. 23. He gives the succession of the bishops of Jerusalem, as well as of Rome, Alexandria, and Antioch, while he omits those of Caesarea.} The council of Nice gave him an honorary precedence among the bishops, though without affecting his dependence on the metropolitan of Caesarea. At least this seems to be the meaning of the short and some. what obscure seventh canon: “Since it is custom and old tradition, that the bishop of Aelia (Jerusalem) should be honored, he shall also enjoy the succession of honor,\footnote{\textgreek{Ἀκολουθία τῆς τιμῆς}; which is variously interpreted. Comp. Hefele, i. 389 sq.} while the metropolis (Caesarea) preserves the dignity allotted to her.” The legal relation of the two remained for a long time uncertain, till the fourth ecumenical council, at its seventh session, confirmed the bishop of Jerusalem in his patriarchal rank, and assigned to him the three provinces of Palestine as a diocese, without opposition.

\xsection{The Rival Patriarchs of Old and New Rome}

§ 57. The Rival Patriarchs of Old and New Rome.

Thus at the close of the fourth century we see the Catholic church of the Graeco-Roman empire under the oligarchy of five coordinate and independent patriarchs, four in the East and one in the West. But the analogy of the political constitution, and the tendency toward a visible, tangible representation of the unity of the church, which had lain at the bottom of the development of the hierarchy from the very beginnings of the episcopate, pressed beyond oligarchy to monarchy; especially in the West. Now that the empire was geographically and politically severed into East and West, which, after the death of Theodosius, in 395, had their several emperors, and were never permanently reunited, we can but expect in like manner a double head in the hierarchy. This we find in the two patriarchs of old Rome and New Rome; the one representing the Western or Latin church, the other the Eastern or Greek. Their power and their relation to each other we must now more carefully observe.

The organization of the church in the East being so largely influenced by the political constitution, the bishop of the imperial capital could not fail to become the most powerful of the four oriental patriarchs. By the second and fourth ecumenical councils, as we have already seen, his actual preëminence was ratified by ecclesiastical sanction, and he was designated to the foremost dignity.\footnote{\textgreek{Τὰ πρεσβεῖα τῆς τιμῆς}... \textgreek{διὰ τὸ εἷναι αὐτὴν} [i.e. Constantinople] \textgreek{νέαν Ῥώμην}. Comp. § 56.} From Justinian I. he further received supreme appellate jurisdiction, and the honorary title of ecumenical patriarch, which he still continues to bear.\footnote{The title \textgreek{οἰκομενικὸς πατριάρχης}\emph{, universalisepiscopus}, had before been used in flattery by oriental patriarchs, and the later Roman bishops bore it, in spite of the protest of Gregory I., without scruple. The statement of popes Gregory I. and Leo IX., that the council of Chalcedon conferred on the Roman bishop Leo the title \emph{of universal episcopus,} and that he rejected it, is erroneous. No trace of it can be found either in the Acts of the councils or in the epistles of Leo. In the Acts, Leo is styled \textgreek{ὁ ἁγιώτατος καὶ μακαριώτατος ἀρχιεπίσκοπος τῆς μεγάλης καὶ πρεσβυτέρασ Ῥώμης}; which, however, in the Latin Acts sent by Leo to the Gallican bishops, was thus enlarged: “Sanctus et beatissimus \emph{Papa, caput universalis ecclesiae}, Leo\emph{.”} The papal legates at Chalcedon subscribed themselves: Vicarii apostolici \emph{universalis} \emph{ecclesiae papae,} which the Greeks translated: \textgreek{τῆς οἰκουμενικῆς ἐκκλησίας ἐπισκόπου}. Hence probably arose the error of Gregory I. The popes wished to be \emph{papae} universalis ecclesiae, not \emph{episcopi} or \emph{patriarchae} universales; no doubt because the latter designation put them on a level with the Eastern patriarchs. Comp. Gieseler, i. 2, p. 192, not. 20, and p. 228, not. 72; and Hefele, ii. 525 sq.} He ordained the other patriarchs, not seldom decided their deposition or institution by his influence, and used every occasion to interfere in their affairs, and assert his supreme authority, though the popes and their delegates at the imperial court incessantly protested. The patriarchates of Jerusalem, Antioch, and Alexandria were distracted and weakened in the course of the fifth and sixth centuries by the tedious monophysite controversies, and subsequently, after the year 622, were reduced to but a shadow by the Mohammedan conquests. The patriarchate of Constantinople, on the contrary, made important advances southwest and north; till, in its flourishing period, between the eighth and tenth centuries, it embraced, besides its original diocese, Calabria, Sicily, and all the provinces of Illyricum, the Bulgarians, and Russia. Though often visited with destructive earthquakes and conflagrations, and besieged by Persians, Arabians, Hungarians, Russians, Latins, and Turks, Constantinople maintained itself to the middle of the fifteenth century as the seat of the Byzantine empire and centre of the Greek church. The patriarch of Constantinople, however, remained virtually only primus inter pares, and has never exercised a papal supremacy over his colleagues in the East, like that of the pope over the metropolitans of the West; still less has he arrogated, like his rival in ancient Rome, the sole dominion of the entire church. Toward the bishop of Rome he claimed only equality of rights and coordinate dignity.

In this long contest between the two leading patriarchs of Christendom, the patriarch of Rome at last carried the day. The monarchical tendency of the hierarchy was much stronger in the West than in the East, and was urging a universal monarchy in the church.

The patriarch of Constantinople enjoyed indeed the favor of the emperor, and all the benefit of the imperial residence. New Rome was most beautifully and most advantageously situated for a metropolis of government, of commerce, and of culture, on the bridge between two continents; and it formed a powerful bulwark against the barbarian conquests. It was never desecrated by an idol temple, but was founded a Christian city. It fostered the sciences and arts, at a time when the West was whelmed by the wild waves of barbarism; it preserved the knowledge of the Greek language and literature through the middle ages; and after the invasion of the Turks it kindled by its fugitive scholars the enthusiasm of classic studies in the Latin church, till Greece rose from the dead with the New Testament in her hand, and held the torch for the Reformation.

But the Roman patriarch had yet greater advantages. In him were united, as even the Greek historian Theodoret concedes,\footnote{Epist. 113, to Pope Leo I.} all the outward and the inward, the political and the spiritual conditions of the highest eminence.

In the first place, his authority rested on an ecclesiastical and spiritual basis, reaching back, as public opinion granted, through an unbroken succession, to Peter the apostle; while Constantinople was in no sense an apostolica sedes, but had a purely political origin, though, by transfer, and in a measure by usurpation, it had possessed itself of the metropolitan rights of Ephesus\footnote{That the apostle Andrew brought the gospel to the ancient Byzantium, is an entirely unreliable legend of later times.} Hence the popes after Leo appealed almost exclusively to the divine origin of their dignity, and to the primacy of the prince of the apostles over the whole church.

Then, too, considered even in a political point of view, old Rome had a far longer and grander imperial tradition to show, and was identified in memory with the bloom of the empire; while New Rome marked the beginning of its decline. When the Western empire fell into the hands of the barbarians, the Roman bishop was the only surviving heir of this imperial past, or, in the well-known dictum of Hobbes, “the ghost of the deceased Roman empire, sitting crowned upon the grave thereof.”

Again, the very remoteness of Rome from the imperial court was favorable to the development of a hierarchy independent of all political influence and intrigue; while the bishop of Constantinople had to purchase the political advantages of the residence at the cost of ecclesiastical freedom. The tradition of the donatio Constantini, though a fabrication of the eighth century, has thus much truth: that the transfer of the imperial residence to the East broke the way for the temporal power and the political independence of the papacy.

Further, amidst the great trinitarian and christological controversies of the Nicene and post-Nicene age, the popes maintained the powerful prestige of almost undeviating ecumenical orthodoxy and doctrinal stability;\footnote{One exception is the brief pontificate of the Arian, Felix II, whom the emperor Constantius, in 355, forcibly enthroned during the exile of Liberius, and who is regarded by some as an illegitimate anti-pope. The accounts respecting him are, however, very conflicting, and so are the opinions of even Roman Catholic historians. Liberius also, in 357, lapsed for a short time into Arianism that he might be recalled from exile. Another and later exception is Pope Honorius, whom even the sixth ecumenical council of Constantinople, 681, anathematized for Monothelite heresy.} while the see of Constantinople, with its Grecian spirit of theological restlessness and disputation, was sullied with the Arian, the Nestorian, the Monophysite, and other heresies, and was in general, even in matters of faith, dependent on the changing humors of the court. Hence even contending parties in the East were accustomed to seek counsel and protection from the Roman chair, and oftentimes gave that see the coveted opportunity to put the weight of its decision into the scale. This occasional practice then formed a welcome basis for a theory of jurisdiction. The Roma locuta est assumed the character of a supreme and final judgment. Rome learned much and forgot nothing. She knew how to turn every circumstances with consummate administrative tact, to her own advantage.

Finally, though the Greek church, down to the fourth ecumenical council, was unquestionably the main theatre of church history and the chief seat of theological learning, yet, according to the universal law of history, “Westward the star of empire takes its way,” the Latin church, and consequently the Roman patriarchate, already had the future to itself. While the Eastern patriarchates were facilitating by internal quarrels and disorder the conquests of the false prophet, Rome was boldly and victoriously striking westward, and winning the barbarian tribes of Europe to the religion of the cross.

\xsection{The Latin Patriarch}

§ 58. The Latin Patriarch.

These advantages of the patriarch of Rome over the patriarch of Constantinople are at the same time the leading causes of the rise of the papacy, which we must now more closely pursue.

The papacy is undeniably the result of a long process of history. Centuries were employed in building it, and centuries have already been engaged upon its partial destruction. Lust of honor and of power, and even open fraud,\footnote{Recall the interpolations of papistic passages in the works of Cyprian; the Roman enlargement of the sixth canon of Nice; the citation of the Sardican canon under the name and the authority of the Nicene council; and the later notorious pseudo-Isidorian decretals. The popes, to be sure, were not the original authors of these falsifications, but they used them freely and repeatedly for their purposes.} have contributed to its development; for human nature lies hidden under episcopal robes, with its steadfast inclination to abuse the power intrusted to it; and the greater the power, the stronger is the temptation, and the worse the abuse. But behind and above these human impulses lay the needs of the church and the plans of Providence, and these are the proper basis for explaining the rise, as well as the subsequent decay, of the papal dominion over the countries and nations of Europe.

That Providence which moves the helm of the history of world and church according to an eternal plan, not only prepares in silence and in a secrecy unknown even to themselves the suitable persons for a given work, but also lays in the depths of the past the foundations of mighty institutions, that they may appear thoroughly furnished as soon as the time may demand them. Thus the origin and gradual growth of the Latin patriarchate at Rome looked forward to the middle age, and formed part of the necessary, external outfit of the church for her disciplinary mission among the heathen barbarians. The vigorous hordes who destroyed the West-Roman empire were to be themselves built upon the ruins of the old civilization, and trained by an awe-inspiring ecclesiastical authority and a firm hierarchical organization, to Christianity and freedom, till, having come of age, they should need the legal schoolmaster no longer, and should cast away his cords from them. The Catholic hierarchy, with its pyramid-like culmination in the papacy, served among the Romanic and Germanic peoples. until the time of the Reformation, a purpose similar to that of the Jewish theocracy and the old Roman empire respectively in the inward and outward preparation for Christianity. The full exhibition of this pedagogic purpose belongs to the history of the middle age; but the foundation for it we find already being laid in the period before us.

The Roman bishop claims, that the four dignities of bishop, metropolitan, patriarch, and pope or primate of the whole church, are united in himself. The first three offices must be granted him in all historical justice; the last is denied him by the Greek church, and by the Evangelical, and by all non-Catholic sects.

His bishopric is the city of Rome, with its cathedral church of St. John Lateran, which bears over its main entrance the inscription: Omnium urbis et orbis ecclesiarum mater et caput; thus remarkably outranking even the church of St. Peter—as if Peter after all were not the first and highest apostle, and had to yield at last to the superiority of John, the representative of the ideal church of the future. Tradition says that the emperor Constantine erected this basilica by the side of the old Lateran palace, which had come down from heathen times, and gave the palace to Pope Sylvester; and it remained the residence of the popes and the place of assembly for their councils (the Lateran councils) till after the exile of Avignon, when they took up their abode in the Vatican beside the ancient church of St. Peter.

As metropolitan or archbishop, the bishop of Rome had immediate jurisdiction over the seven suffragan bishops, afterward called cardinal bishops, of the vicinity: Ostia, Portus, Silva candida, Sabina, Praeneste, Tusculum, and Albanum.

As patriarch, he rightfully stood on equal footing with the four patriarchs of the East, but had a much larger district and the primacy of honor. The name is here of no account, since the fact stands fast. The Roman bishops called themselves not patriarchs, but popes, that they might rise the sooner above their colleagues; for the one name denotes oligarchical power, the other, monarchical. But in the Eastern church and among modern Catholic historians the designation is also quite currently applied to Rome.

The Roman patriarchal circuit primarily embraced the ten suburban provinces, as they were called, which were under the political jurisdiction of the Roman deputy, the Vicarius Urbis; including the greater part of Central Italy, all Upper Italy, and the islands of Sicily, Sardinia, and Corsica.\footnote{Concil. Nicaean. of 325, can. 6, in the Latin version of Rufinus (Hist. Eccl. x. 6): “Et ut apud Alexandria et in urbe Roma vetusta consuetudo servetur, ut vel ille Ægypti, vel hic \emph{suburbicariarum ecclesiarum} sollicitudinem gerat.” The words \emph{suburb. eccl}. are wanting in the Greek original, and are a Latin definition of the patriarchal diocese of Rome at the end of the fourth century. Since the seventeenth century they have given rise to a long controversy among the learned. The jurist Gothofredus and his friend Salmasius limited the \emph{regiones suburbicariae} to the small province of the \emph{Praefectus} Urbis, i.e. to the city of Rome with the immediate vicinity to the hundredth milestone; while the Jesuit Sirmond extended it to the much greater official district of the \emph{Vicarius} Urbis, viz., the ten provinces of Campania, Tuscia with Umbria, Picenum suburbicarium, Valeria, Samnium, Apulia with Calabria, Lucania and Brutii, Sicilia, Sardinia, and Corsica. The comparison of the Roman bishop with the Alexandrian in the sixth canon of the Nicene council favors the latter view; since even the Alexandrian diocese likewise stretched over several provinces. The \emph{Prisca}, however—a Latin collection of canons from the middle of the fifth century—has perhaps hit the truth of the matter, in saying, in its translation of the canon in question: “Antiqui moris est ut urbis Romae episcopus habeat principatum, ut \emph{suburbicaria loca} [i.e. here, no doubt, the smaller province of the Praefectus] et \emph{omnem provinciam} \emph{suam} [i.e. the larger district of the Vicarius, or a still wider, indefinite extent] sollicitudine sua gubernet.” Comp. Mansi, Coll. Conc. vi. 1127, and Hefele, i. 380 sqq.} In its wider sense, however, it extended gradually over the entire west of the Roman empire, thus covering Italy, Gaul, Spain, Illyria, southeastern Britannia, and northwestern Africa.\footnote{According to the political division of the empire, the Roman patriarchate embraced in the fifth century three praefectures, which were divided into eight political dioceses and sixty-nine provinces. These are, (1) the praefecture of Italy, with the three dioceses of Italy, Illyricum, and Africa; (2) the praefectum Galliarum, with the dioceses of Gaul, Spain, and Britain; (3) the praefecture of Illyricum (not to be confounded with the \emph{province} of Illyria, which belonged to the praefecture of Italy), which, after 879, was separated indeed from the Western empire, as Illyricum orientale, but remained ecclesiastically connected with Rome, and embraced the two dioceses of Macedonia and Dacia. Comp. Wiltsch, l.c. i. 67 sqq.; Maassen, p. 125; and Hefele, i. 383.}

The bishop of Rome was from the beginning the only Latin patriarch, in the official sense of the word. He stood thus alone, in the first place, for the ecclesiastical reason, that Rome was the only sedes apostolica in the West, while in the Greek church three patriarchates and several other episcopal sees, such as Ephesus, Thessalonica, and Corinth, shared the honor of apostolic foundation. Then again, he stood politically alone, since Rome was the sole metropolis of the West, while in the East there were three capitals of the empire, Constantinople, Alexandria, and Antioch. Hence Augustine, writing from the religious point of view, once calls Pope Innocent I. the “ruler of the Western church;”\footnote{Contra Julianum, lib. i. cap. 6.} and the emperor Justinian, on the ground of political distribution, in his 109th Novelle, where he speaks of the ecclesiastical division of the whole world, mentions only five known patriarchates, and therefore only one patriarchate of the West. The decrees of the ecumenical councils, also, know no other Western patriarchate than the Roman, and this was the sole medium through which the Eastern church corresponded with the Western. In the great theological controversies of the fourth and fifth centuries the Roman bishop appears uniformly as the representative and the organ of all Latin Christendom.

It was, moreover, the highest interest of all orthodox churches in the West, amidst the political confusion and in conflict with the Arian Goths, Vandals, and Suevi, to bind themselves closely to a common centre, and to secure the powerful protection of a central authority. This centre they could not but find in the primitive apostolic church of the metropolis of the world. The Roman bishops were consulted in almost all important questions of doctrine or of discipline. After the end of the fourth century they issued to the Western bishops in reply, pastoral epistles and decretal letters,\footnote{\emph{Epistola decretales}; an expression, which, according to Gieseler and others, occurs first about 500, in the so-called decretum Gelasii de libris recipiendis et non recipiendis.} in which they decided the question at first in the tone of paternal counsel, then in the tone of apostolic authority, making that which had hitherto been left to free opinion, a fixed statute. The first extant decretal is the Epistola of Pope Siricius to the spanish bishop Himerius, a.d. 385, which contains, characteristically, a legal enforcement of priestly celibacy, thus of an evidently unapostolic institution; but in this Siricius appeals to “generalia decreta,” which his predecessor Liberius had already issued. In like manner the Roman bishops repeatedly caused the assembling of general or patriarchal councils of the West (synodos occidentales), like the synod of Axles in 314. After the sixth and seventh centuries they also conferred the pallium on the archbishops of Salona, Ravenna, Messina, Syracuse, Palermo, Arles, Autun, Sevilla, Nicopolis (in Epirus), Canterbury, and other metropolitans, in token of their superior jurisdiction.\footnote{See the information concerning the conferring of the pallium in Wiltsch, i. 68 sq.}

\xsection{Conflicts and Conquests of the Latin Patriarchate}

§ 59. Conflicts and Conquests of the Latin Patriarchate.

But this patriarchal power was not from the beginning and to a uniform extent acknowledged in the entire West. Not until the latter part of the sixth century did it reach the height we have above described.\footnote{This is conceded by Hefele, i. 383 sq.: “It is, however, not to be mistaken, that the bishop of Rome did not \emph{everywhere,} in all the West, exercise \emph{full} patriarchal rights; that, to wit, in several provinces, simple bishops were ordained without his coöperation.” And not only simple bishops, but also metropolitans. See the text.} It was not a divine institution, unchangeably fixed from the beginning for all times, like a Biblical article of faith; but the result of a long process of history, a human ecclesiastical institution under providential direction. In proof of which we have the following incontestable facts:

In the first place, even in Italy, several metropolitans maintained, down to the close of our period, their own supreme headship, independent of Roman and all other jurisdiction.\footnote{A\textgreek{ὐτοκέφαλοι}, also \textgreek{ἀκέφαλοι}, as in the East especially the archbishops of Cyprus and Bulgaria were called, and some other metropolitans, who were subject to no patriarch.} The archbishops of Milan, who traced their church to the apostle Barnabas, came into no contact with the pope till the latter part of the sixth century, and were ordained without him or his pallium. Gregory I., in 593, during the ravages of the Longobards, was the first who endeavored to exercise patriarchal rights there: he reinstated an excommunicated presbyter, who had appealed to him.\footnote{Comp. Wiltsch, i. 234.} The metropolitans of Aquileia, who derived their church from the evangelist Mark, and whose city was elevated by Constantine the Great to be the capital of Venetia and Istria, vied with Milan, and even with Rome, calling themselves “patriarchs,” and refusing submission to the papal jurisdiction even under Gregory the Great.\footnote{Comp. Gregory I., Epist. l. iv. 49; and Wiltsch, i. 236 sq. To the metropolis of Aquileia belonged the bishopric of Verona, Tridentum (the Trent, since become so famous), Aemona, Altinum, Torcellum, Pola, Celina, Sabiona, Forum Julii, Bellunum, Concordia, Feltria, Tarvisium, and Vicentia.} The bishop of Ravenna likewise, after 408, when the emperor Honorius selected that city for his residence, became a powerful metropolitan, with jurisdiction over fourteen bishoprics. Nevertheless he received the pallium from Gregory the Great, and examples occur of ordination by the Roman bishop.\footnote{Baron. Ann. ad ann. 433; Wiltsch, i. 69, 87.}

The North African bishops and councils in the beginning of the fifth century, with all traditional reverence for the apostolic see, repeatedly protested, in the spirit of Cyprian, against encroachments of Rome, and even prohibited all appeal in church controversies from their own to a transmarine or foreign tribunal, upon pain of excommunication.\footnote{Comp. the relevant Acts of councils in Gieseler, i. 2, p. 221 sqq., and an extended description of this case of appeal in Greenwood, Cath. Petri, i. p. 299-310, and in Hefele, Concilien-Gesch. ii. 107 sqq., 120, 123 sq.} The occasion of this was an appeal to Rome by the presbyter Apiarius, who had been deposed for sundry offences by Bishop Urbanus, of Sicca, a disciple and friend of Augustine, and whose restoration was twice attempted, by Pope Zosimus in 418, and by Pope Coelestine in 424. From this we see that the popes gladly undertook to interfere for a palpably unworthy priest, and thus sacrificed the interests of local discipline, only to make their own superior authority felt. The Africans referred to the genuine Nicene canon (for which Zosimus had substituted the Sardican appendix respecting the appellate jurisdiction of Rome, of which the Nicene council knew nothing), and reminded the pope, that the gift of the Holy Ghost, needful for passing a just judgment, was not lacking to any province, and that he could as well inspire a whole province as a single bishop. The last document in the case of this appeal of Apiarius is a letter of the (twentieth) council of Carthage, in 424, to Pope Coelestine I., to the following purport:\footnote{Mansi, iii. 839 sq.} “Apiarius asked a new trial, and gross misdeeds of his were thereby brought to light. The papal legate, Faustinus, has, in the face of this, in a very harsh manner demanded the reception of this man into the fellowship of the Africans, because he has appealed to the pope and been received into fellowship by him. But this very thing ought not to have been done. At last has Apiarius himself acknowledged all his crimes. The pope may hereafter no longer so readily give audience to those who come from Africa to Rome, like Apiarius, nor receive the excommunicated into church communion, be they bishops or priests, as the council of Nice (can. 5) has ordained, in whose direction bishops are included. The assumption of appeal to Rome is a trespass on the rights of the African church, and what has been [by Zosimus and his legates] brought forward as a Nicene ordinance for it, is not Nicene, and is not to be found in the genuine copies of the Nicene Acts, which have been received from Constantinople and Alexandria. Let the pope, therefore, in future send no more judges to Africa, and since Apiarius has now been excluded for his offences, the pope will surely not expect the African church to submit longer to the annoyances of the legate Faustinus. May God the Lord long preserve the pope, and may the pope pray for the Africans.” In the Pelagian controversy the weak Zosimus, who, in opposition to the judgment of his predecessor Innocent, had at first expressed himself favorably to the heretics, was even compelled by the Africans to yield. The North African church maintained this position under the lead of the greatest of the Latin fathers, St. Augustine, who in other respects contributed more than any other theologian or bishop to the erection of the Catholic system. She first made submission to the Roman jurisdiction, in the sense of her weakness, under the shocks of the Vandals. Leo (440–461) was the first pope who could boast of having extended the diocese of Rome beyond Europe into another quarter of the globe.\footnote{Epist. 87; Mansi, vi. 120.} He and Gregory the Great wrote to the African bishops entirely in the tone of paternal authority without provoking reply.

In Spain the popes found from the first a more favorable field. The orthodox bishops there were so pressed in the fifth century by the Arian Vandals, Suevi, Alani, and soon after by the Goths, that they sought counsel and protection with the bishop of Rome, which, for his own sake, he was always glad to give. So early as 385, Siricius, as we have before observed, issued a decretal letter to a Spanish bishop. The epistles of Leo to Bishop Turibius of Asturica, and the bishops of Gaul and Spain,\footnote{Ep. 93 and 95; Mansi, vi. 131 and 132.} are instances of the same authoritative style. Simplicius (467–483) appointed the bishop Zeno of Sevilla papal vicar,\footnote{Mansi, vii. 972.} and Gregory the Great, with a paternal letter, conferred the pallium on Leander, bishop of Sevilla.\footnote{Greg. Ep. i. 41; Mansi, ix. 1059. Comp. Wiltsch, i. 71.}

In Gaul, Leo succeeded in asserting the Roman jurisdiction, though not without opposition, in the affair of the archbishop Hilary of Arles, or Arelate. The affair has been differently represented from the Gallican and the ultramontane points of view.\footnote{This difference shows itself in the two editions of the works of Leo the Great, respectively: that of the French Pasquier Quesnel, a Gallican and Jansenist (exiled 1681, died at Brussels 1719), which also contains the works, and a vindication, of Hilaryof Arles (Par. 1675, in 2 vols.), and was condemned in 1676 by the Congregation of the Index, without their even reading it; and that of the two brothers Ballerini, which appeared in opposition to the former (Ven. 1755-1757, 3 vols.), and represents the Italian ultramontane side. Comp. further on this contest of Hilarius Arelatensis (not to be confounded with Hilarius Pictaviensis, Hilarius Narbonensis, and others of the same name) with Pope Leo, the Vita Hilarii of Honoratus Massiliensis, of about the year 490 (printed in Mansi, vi. 461 sqq., and in the Acta Sanct. ad d. 5 Maji); the article by Perthel, in Illgen’s Zeitschrift for Hist. Theol. 1843; Greenwood, l.c. i. p. 350-356; Milman, Lat. Christianity, i. p. 269-276 (Amer. ed.); and the article “Hilarius” in Wetzer and Welte’s Kirchenlexic vol. v. p. 181 sqq.} Hilary (born 403, died 449), first a rigid monk, then, against his will, elevated to the bishopric, an eloquent preacher, an energetic prelate, and the first champion of the freedom of the Gallican church against the pretensions of Rome, but himself not free from hierarchical ambition, deposed Celidonius, the bishop of Besançon, at a council in that city (synodus Vesontionensis), because he had married a widow before his ordination, and had presided as judge at a criminal trial and pronounced sentence of death; which things, according to the ecclesiastical law, incapacitated him for the episcopal office. This was unquestionably an encroachment on the province of Vienne, to which Besaçon belonged. Pope Zosimus had, indeed, in 417, twenty-eight years before, appointed the bishop of Arles, which was a capital of seven provinces, to be papal vicar in Gaul, and had granted him metropolitan rights in the provinces Viennensis, and Narbonensis prima and secunda, though with the reservation of causae majores.\footnote{“Nisi magnitudo causae etiam nostrum exquirat examen.” Gieseler, i. 2, p. 218; Greenwood, i. p. 299.} The metropolitans of Vienne, Narbonne, and Marseilles, however, did not accept this arrangement, and the succeeding popes found it best to recognize again the old metropolitans.\footnote{Comp. Bonifacii I Epist. 12 ad Hilarium Narbon. (not Arelatensen), a.d.422, in Gieseler, p. 219. Boniface here speaks in favor of the Nicene principle, that each metropolitan should rule simply over one province. Greenwood overlooks this change, and hence fully justifies Hilaryon the ground of the appointment of Zosimus. But even though this appointment had stood, the deposition of a bishop was still a \emph{causa} \emph{major}, which Hilary, as vicar of the pope, should have laid before him for ratification.} Celidonius appealed to Leo against that act of Hilary. Leo, in 445, assembled a Roman council (concilium sacerdotum), and reinstated him, as the accusation of Hilary, who himself journeyed on foot in the winter to Rome, and protested most vehemently against the appeal, could not be proven to the satisfaction of the pope. In fact, he directly or indirectly caused Hilary to be imprisoned, and, when he escaped and fled back to Gaul, cut him off from the communion of the Roman church, and deprived him of all prerogatives in the diocese of Vienne, which had been only temporarily conferred on the bishop of Arles, and were by a better judgment (sententia meliore) taken away. He accused him of assaults on the rights of other Gallican metropolitans, and above all of insubordination toward the principality of the most blessed Peter; and he goes so far as to say: “Whoso disputes the primacy of the apostle Peter, can in no way lessen the apostle’s dignity, but, puffed up by the spirit of his own pride, he destroys himself in hell.”\footnote{Leo, Epist. 10 (al. 89) ad Episc. provinciae Viennensis. What an awful perversion this of the true Christian stand-point!} Only out of special grace did he leave Hilary in his bishopric. Not satisfied with this, he applied to the secular arm for help, and procured from the weak Western emperor, Valentinian III., an edict to Aetius, the magister militum of Gaul, in which it is asserted, almost in the words of Leo, that the whole world (universitas; in Greek, οἰκουμένη) acknowledges the Roman see as director and governor; that neither Hilary nor any bishop might oppose its commands; that neither Gallican nor other bishops should, contrary to the ancient custom, do anything without the authority of the venerable pope of the eternal city; and that all decrees of the pope have the force of law.

The letter of Leo to the Gallican churches, and the edict of the emperor, give us the first example of a defensive and offensive alliance of the central spiritual and temporal powers in the pursuit of an unlimited sovereignty. The edict, however, could of course have power, at most, only in the West, to which the authority of Valentinian was limited. In fact, even Hilary and his successors maintained, in spite of Leo, the prerogatives they had formerly received from Pope Zosimus, and were confirmed in them by later popes.\footnote{The popes Vigil, 539-555, Pelagius, 555-559, and Gregory the Great conferred on the archbishop of Arles, besides the pallium, also the papal vicariate (vices). Comp. Wiltsch, i. 71 sq.} Beyond this the issue of the contest is unknown. Hilary of Arles died in 449, universally esteemed and loved, without, so far as we know, having become formally reconciled with Rome;\footnote{At all events, no reconciliation can be certainly proved. Hilarydid, indeed, according to the account of his disciple and biographer, who some forty years after his death encircled him with the halo, take some steps toward reconciliation, and sent two priests as delegates with a letter to the Roman prefect, Auxiliaris. The latter endeavored to act the mediator, but gave the delegates to understand, that Hilary, by his vehement boldness, had too deeply wounded the delicate ears of the Romans. In Leo’s letter a new trespass is charged upon Hilary, on the rights of the bishop Projectus, \emph{after} the deposition of Celidonius. And Hilarydied soon after this contest (449). Waterland ascribed to him the Athanasian Creed, though without good reason.} though, notwithstanding this, he figures in a remarkable manner in the Roman calendar, by the side of his papal antagonist Leo, as a canonical saint. Undoubtedly Leo proceeded in this controversy far too rigorously and intemperately against Hilary; yet it was important that he should hold fast the right of appeal as a guarantee of the freedom of bishops against the encroachments of metropolitans. The papal despotism often proved itself a wholesome check upon the despotism of subordinate prelates.

With Northern Gaul the Roman bishops came into less frequent contact; yet in this region also there occur, in the fourth and fifth centuries, examples of the successful assertion of their jurisdiction.

The early British church held from the first a very isolated position, and was driven back, by the invasion of the pagan Anglo-Saxons, about the middle of the fifth century, into the mountains of Wales, Cornwallis, Cumberland, and the still more secluded islands. Not till the conversion of the Anglo-Saxons under Gregory the Great did a regular connection begin between England and Rome.

Finally, the Roman bishops succeeded also in extending their patriarchal power eastward, over the praefecture of East Illyria. Illyria belonged originally to the Western empire, remained true to the Nicene faith through the Arian controversies, and for the vindication of that faith attached itself closely to Rome. When Gratian, in 379, incorporated Illyricum Orientale with the Eastern empire, its bishops nevertheless refused to give up their former ecclesiastical connection. Damasus conferred on the metropolitan Acholius, of Thessalonica, as papal vicar, patriarchal rights in the new praefecture. The patriarch of Constantinople endeavored, indeed, repeatedly, to bring this ground into his diocese, but in vain. Justinian, in 535, formed of it a new diocese, with an independent patriarch at Prima Justiniana (or Achrida, his native city); but this arbitrary innovation had no vitality, and Gregory I. recovered active intercourse with the Illyrian bishops. Not until the eighth century, under the emperor Leo the Isaurian, was East Illyria finally severed from the Roman diocese and incorporated with the patriarchate of Constantinople.\footnote{Comp. Gieseler, i. 2, p. 21 5 sqq.; and Wiltsch, i. 72 sqq., 431 sqq.}

\xsection{The Papacy}

§ 60. The Papacy.

Literature, as in § 55, and vol. i. § 110.

At last the Roman bishop, on the ground of his divine institution, and as successor of Peter, the prince of the apostles, advanced his claim to be primate of the entire church, and visible representative of Christ, who is the invisible supreme head of the Christian world. This is the strict and exclusive sense of the title, Pope.\footnote{The name \emph{papa}—according to some an abbreviation of \emph{pater patrum}, but more probably, like the kindred \emph{abbas}, \textgreek{πάππας}, or \textgreek{πάπας}, \emph{pa-pa}, simply an imitation of the first prattling of children, thus equivalent to \emph{father}—was, in the West, for a long time the honorary title of every bishop, as a spiritual father; but, after the fifth century, it became the special distinction of the patriarchs, and still later was assigned exclusively to the Roman bishop, and to him in an eminent sense, as father of the whole church. Comp. Du Cange, Glossar. s. verb. \emph{Papa} and \emph{Pater Patrum}; and Hoffmann, Lexic. univers. iv. p. 561. In the same exclusive sense the Italian and Spanish \emph{papa}, the French \emph{pape}, the English \emph{pope}, and the German \emph{Papst}or \emph{Pabst,} are used. In the Greek and Russian churches, on the contrary, all priests are called \emph{Popes} (from \textgreek{πάπας,} \emph{papa}). The titles \emph{apostolicus, vicarius Christi, summus} \emph{pontifex, sedes apostolica,} were for a considerable time given to various bishops and their sees, but subsequently claimed exclusively by the bishops of Rome.}

Properly speaking, this claim has never been fully realized, and remains to this day an apple of discord in the history of the church. Greek Christendom has never acknowledged it, and Latin, only under manifold protests, which at last conquered in the Reformation, and deprived the papacy forever of the best part of its domain. The fundamental fallacy of the Roman system is, that it identifies papacy and church, and therefore, to be consistent, must unchurch not only Protestantism, but also the entire Oriental church from its origin down. By the “una sancta catholica apostolica ecclesia” of the Niceno-Constantinopolitan creed is to be understood the whole body of Catholic Christians, of which the ecclesia Romana, like the churches of Alexandria, Antioch, Jerusalem, and Constantinople, is only one of the most prominent branches. The idea of the papacy, and its claims to the universal dominion of the church, were distinctly put forward, it is true, so early as the period before us, but could not make themselves good beyond the limits of the West. Consequently the papacy, as a historical fact, or so far as it has been acknowledged, is properly nothing more than the Latin patriarchate run to absolute monarchy.

By its advocates the papacy is based not merely upon church usage, like the metropolitan and patriarchal power, but upon divine right; upon the peculiar position which Christ assigned to Peter in the well-known words: “Thou art Peter, and on this rock will I build my church.”\footnote{Matt. xvi. 18: \textgreek{Σὺ εἷ Πέτρος, καὶ ἐπὶ ταύτῃ τῇ πέτρᾳ} [mark the change of the gender from the masculine to the feminine, from the person to the thing or the truth confessed—a change which disappears in the English and German versions] \textgreek{οἰκοδομήσω μου τὴν ἐκκλησίαν, καὶ πύλαι ᾅδου οὐ κατισχύσουσιν αὐτῆς}. Comp. the commentators, especially Meyer, Lange, Alford, Wordsworth, \emph{ad loc.,} and my Hist. of the Apost. Church, § 90 and 94 (N. Y. ed. p. 350 sqq., and 374 sqq.).} This passage was at all times taken as an immovable exegetical rock for the papacy. The popes themselves appealed to it, times without number, as the great proof of the divine institution of a visible and infallible central authority in the church. According to this view, the primacy is before the apostolate, the head before the body, instead of the reverse.

But, in the first place, this preëminence of Peter did not in the least affect the independence of the other apostles. Paul especially, according to the clear testimony of his epistles and the book of Acts, stood entirely upon his own authority, and even on one occasion, at Antioch, took strong ground against Peter. Then again, the personal position of Peter by no means yields the primacy to the Roman bishop, without the twofold evidence, first that Peter was actually in Rome, and then that he transferred his prerogatives to the bishop of that city. The former fact rests upon a universal tradition of the early church, which at that time no one doubted, but is in part weakened and neutralized by the absence of any clear Scripture evidence, and by the much more certain fact, given in the New Testament itself, that Paul labored in Rome, and that in no position of inferiority or subordination to any higher authority than that of Christ himself. The second assumption, of the transfer of the primacy to the Roman bishops, is susceptible of neither historical nor exegetical demonstration, and is merely an inference from the principle that the successor in office inherits all the official prerogatives of his predecessor. But even granting both these intermediate links in the chain of the papal theory, the double question yet remains open: first, whether the Roman bishop be the only successor of Peter, or share this honor with the bishops of Jerusalem and Antioch, in which places also Peter confessedly resided; and secondly, whether the primacy involve at the same time a supremacy of jurisdiction over the whole church, or be only an honorary primacy among patriarchs of equal authority and rank. The former was the Roman view; the latter was the Greek.

An African bishop, Cyprian († 258), was the first to give to that passage of the 16th of Matthew, innocently as it were, and with no suspicion of the future use and abuse of his view, a papistic interpretation, and to bring out clearly the idea of a perpetual cathedra Petri. The same Cyprian, however, whether consistently or not, was at the same time equally animated with the consciousness of episcopal equality and independence, afterward actually came out in bold opposition to Pope Stephen in a doctrinal controversy on the validity of heretical baptism, and persisted in this protest to his death.\footnote{Comp. vol. i. § 110.}

\xsection{Opinions of the Fathers}

§ 61. Opinions of the Fathers.

A complete collection of the patristic utterances on the primacy of Peter and his successors, though from the Roman point of view, may be found in the work of Rev. Jos. Berington and Rev. John Kirk: “The Faith of Catholics confirmed by Scripture and attested by the Fathers of the first five centuries of the Church,” 3d ed., London, 1846, vol. ii. p. 1–112. Comp. the works quoted sub § 55, and a curious article of Prof. Ferd. Piper, on Rome, the eternal city, in the Evang. Jahrbuch for 1864, p. 17–120, where the opinions of the fathers on the claims of the urbs aeterna and its many fortunes are brought out.

We now pursue the development of this idea in the church fathers of the fourth and fifth centuries. In general they agree in attaching to Peter a certain primacy over the other apostles, and in considering him the foundation of the church in virtue of his confession of the divinity of Christ; while they hold Christ to be, in the highest sense, the divine ground and rock of the church. And herein lies a solution of their apparent self-contradiction in referring the petra in Matt. xvi. 18, now to the person of Peter, now to his confession, now to Christ. Then, as the bishops in general were regarded as successors of the apostles, the fathers saw in the Roman bishops, on the ground of the ancient tradition of the martyrdom of Peter in Rome, the successor of Peter and the heir of the primacy. But respecting the nature and prerogatives of this primacy their views were very indefinite and various. It is remarkable that the reference of the rock to Christ, which Augustine especially defended with great earnestness, was acknowledged even by the greatest pope of the middle ages, Gregory VII., in the famous inscription he sent with a crown to the emperor Rudolph: “Petra [i.e., Christ] dedit Petro [i.e., to the apostle], Petrus [the pope] diadema Rudolpho.”\footnote{Baronius, Annal. ad ann. 1080, vol. xi. p. 704.} It is worthy of notice, that the post-Nicene, as well as the ante-Nicene fathers, with all their reverence for the Roman see, regarded the heathenish title of Rome, urbs aeterna, as blasphemous, with reference to the passage of the woman sitting upon a scarlet-colored beast, full of names of blasphemy, Rev. xvii. 3.\footnote{Hieronymus, Adv. Jovin. lib. ii. c. 38 (Opera, t. ii. p. 382), where he addresses Rome: “Ad te loquar, quae scriptam in fronte blasphemiam Christi confessione delesti.” Prosper: “Eterna cum dicitur quae temporalis est, utique nomen est blasphemiae.” Comp. Piper, l.c. p. 46.} The prevailing opinion seems to have been, that Rome and the Roman empire would fall before the advent of Antichrist and the second coming of the Lord.\footnote{So Chrysostomad 2 Thess. ii. 7; Hieronymus, Ep. cxxi. qu. 11 (tom. i. p. 880 sq.); Augustine, De Civit. Dei, lib. xx. cap. 19.}

1. The views of the Latin fathers.

The Cyprianic idea was developed primarily in North Africa, where it was first clearly pronounced.

Optatus, bishop of Milevi, the otherwise unknown author of an anti-Donatist work about a.d. 384, is, like Cyprian, thoroughly possessed with the idea of the visible unity of the church; declares it without qualification the highest good, and sees its plastic expression and its surest safeguard in the immovable cathedra Petri, the prince of the apostles, the keeper of the keys of the kingdom of heaven, who, in spite of his denial of Christ, continued in that relation to the other apostles, that the unity of the church might appear in outward fact as an unchangeable thing, invulnerable to human offence. All these prerogatives have passed to the bishops of Rome, as the successors of this apostle.\footnote{De schismate Donatistarum, lib. ii. cap. 2, 3, and l. vii. 3. The work was composed while Siricius was bishop of Rome, hence about 384.}

Ambrose of Milan († 397) speaks indeed in very high terms of the Roman church, and concedes to its bishops a religious magistracy like the political power of the emperors of pagan Rome;\footnote{Ambr. Sermo ii. in festo Petri et Pauli: “In urbe Romae, quae principatum et caput obtinet nationum: scilicet ut ubi caput superstitionis erat, illic caput quiesceret sanctitatis, et ubi gentilium principes habitabant, illic ecclesiarum principes morerentur.” In Ps. 40: “Ipse est Petrus cui dixit: Tu es Petrus ... ubi ergo Patrus, ibi ecclesia; ubi ecclesia, ibi mulla mors, sed vita eterna.” Comp. the poetic passage in his Morning Hymn, in the citation from Augustinefurther on. But in another passage he likewise refers the rock to Christ, in Luc. ix. 20: “Petra est Christus,” etc.} yet he calls the primacy of Peter only a “primacy of confession, not of honor; of faith, not of rank,”\footnote{De incarnat. Domini, c. 4: “Primatum confessionis utique, non honoris, primatum fidei, non ordinis.”} and places the apostle Paul on an equality with Peter.\footnote{De Spiritu S. ii. 12: “Nec Paulus inferior Petro, quamvis ille ecclesiae fundamentum.” Sermo ii. in festo P. et P., just before the above-quoted passage: “Ergo beati Petrus et Paulus eminent inter universos apostolos, et peculiari quadam praerogativa praecellunt. Verum inter ipsos, quis cui praeponatur, incertum est. Puto enim illos aequales esse meritis, qui aequales sunt passione.” Augustine, too, once calls Paul, not Peter, \emph{caput et princeps apostolorum,} and in another place that he \emph{tanti apostolatus meruit principatum}.} Of any dependence of Ambrose, or of the bishops of Milan in general during the first six centuries, on the jurisdiction of Rome, no trace is to be found.

Jerome († 419), the most learned commentator among the Latin fathers, vacillates in his explanation of the petra; now, like Augustine, referring it to Christ,\footnote{Hieron. in Amos, vi. 12: “Petra Christus est qui donavit apostolis suis, ut ipsi quoque petrae vocentur.” And in another place: “Ecclesia Catholica super Petram \emph{Christum} stabili radici fundata est.”} now to Peter and his confession.\footnote{Adv. Jovin. l. i. cap. 26 (in Vallars. ed., tom. ii. 279), in reply to Jovinian’s appeal to Peter in favor of marriage: “At dicis: super Petrum fundatur ecclesia; licet id ipsum in alio loco super omnes apostolos fiat, et cuncti claves regni coelorum accipiant, et ex aequo super eos fortitudo ecclesiae solidetur, tamen propterea inter duodecim unus eligitur, ut capite constituto, schismatis tollatur occasio.” So Epist. xv. ad Damasum papam (ed. Vall. i. 37).} In his commentary on Matt. xvi., he combines the two interpretations thus: “As Christ gave light to the apostles, so that they were called, after him, the light of the world, and as they received other designations from the Lord; so Simon, because he believed on the rock, Christ, received the name Peter, and in accordance with the figure of the rock, it is justly said to him: ’I will build my church upon thee (super te),’ ” He recognizes in the Roman bishop the successor of Peter, but advocates elsewhere the equal rights of the bishops,\footnote{Comp. Epist. 146, ed. Vall. i. 1076 (or Ep. 101 ed. Bened., al. 85) ad Evangelum: “Ubicunque fuerit episcopus, sive Romae, sive Eugubii, sive Constantinopoli, sive Rhegii, sive Alexandriae, sive Tanis [an intentional collocation of the most powerful and most obscure bishoprics], ejusdem est meriti, ejusdem est et sacerdotii. Potentia divitiarum et paupertatis humilitas vel sublimiorem vel inferiorem episcopum non facit. Caeterum omnes apostolorum successores sunt.”} and in fact derives even the episcopal office, not from direct divine institution, but from the usage of the church and from the presidency in the presbyterium.\footnote{Comp. § 52, above. J. Craigie Robertson, Hist. of the Christian Church to 590 (Lond. 1854), p. 286, note, finds a remarkable negative evidence against the papal claims in St. Jerome’s Ep. 125, “where submission to one head is enforced on monks by the instinctive habits of beasts, bees, and cranes, the contentions of Esau and Jacob, of Romulus and Remus, the oneness of an emperor in his dominions, of a judge in his province, of a master in his house, of a pilot in a ship, of a general in an army, of a bishop, the archpresbyter, and the archdeacon in a church; but there is no mention of the one universal bishop.”} He can therefore be cited as a witness, at most, for a primacy of honor, not for a supremacy of jurisdiction. Beyond this even the strongest passage of his writings, in a letter to his friend, Pope Damasus (a.d. 376), does not go: “Away with the ambition of the Roman head; I speak with the successor of the fisherman and disciple of the cross. Following no other head than Christ, I am joined in the communion of faith with thy holiness, that is, with the chair of Peter. On that rock I know the church to be built.”\footnote{Ep. xv. (alias 57) ad Damasum papam (ed. Vall. l. 37 sq.): “Facessat invidia: Romani culminis recedat ambitio, cum successore piscatoris et discipulo crucis loquor. Ego nullum primum, nisi Christum sequens, Beatitudini tuae, id est cathedrae Petri, communione consocior. Super illam petram aedificatam ecclesiam scio. Quicunque extra hanc domum agnum comederit, profanus est. Si quis in Noe arca non fuerit, peribit regnante diluvio.”} Subsequently this father, who himself had an eye on the papal chair, fell out with the Roman clergy, and retired to the ascetic and literary solitude of Bethlehem, where he served the church by his pen far better than he would have done as the successor of Damasus.

Augustine († 430), the greatest theological authority of the Latin church, at first referred the words, “On this rock I will build my church,” to the person of Peter, but afterward expressly retracted this interpretation, and considered the petra to be Christ, on the ground of a distinction between petra (ἐπὶ ταύτῃ τῇ πέτρᾳ) and Petrus (σὺ εἷ Πέτρος); a distinction which Jerome also makes, though with the intimation that it is not properly applicable to the Hebrew and Syriac Cephas.\footnote{Hier. Com. in Ep. ad Galat. ii. 11, 12 (ed. Vallars. tom. vii. col. 409): “Non quod aliud significat \emph{Petrus,} aliud\emph{Cephas,} sed quo quam nos Latine et Graece \emph{petram} vocemus, hanc Hebraei et Syri, propter linguae inter se viciniam, \emph{Cephan}, nuncupent.”} “I have somewhere said of St. Peter” thus Augustine corrects himself in his Retractations at the close of his life\footnote{Retract. l. i. c. 21.}—“that the church is built upon him as the rock; a thought which is sung by many in the verses of St. Ambrose:

’Hoc ipsa petra ecclesiae

Canente, culpam diluit.’\footnote{In the Ambrosian Morning Hymn: “Aeterne rerum conditor.”}

(The Rock of the church himself

In the cock-crowing atones his guilt.)

But I know that I have since frequently said, that the word of the Lord, ’Thou art Petrus, and on this petra I will build my church,’ must be understood of him, whom Peter confessed as Son of the living God; and Peter, so named after this rock, represents the person of the church, which is founded on this rock and has received the keys of the kingdom of heaven. For it was not said to him: ’Thou art a rock’ (petra), but, ’Thou art Peter’ (Petrus); and the rock was Christ, through confession of whom Simon received the name of Peter. Yet the reader may decide which of the two interpretations is the more probable.” In the same strain he says, in another place: “Peter, in virtue of the primacy of his apostolate, stands, by a figurative generalization, for the church .... When it was said to him, ’I will give unto thee the keys of the kingdom of heaven,’ \&c., he represented the whole church, which in this world is assailed by various temptations, as if by floods and storms, yet does not fall, because it is founded upon a rock, from which Peter received his name. For the rock is not so named from Peter, but Peter from the rock (non enim a Petro petra, sed Petrus a petra), even as Christ is not so called after the Christian, but the Christian after Christ. For the reason why the Lord says, ’On this rock I will build my church’ is that Peter had said: ’Thou art the Christ, the Son of the living God,’ On this rock, which then hast confessed, says he will build my church. For Christ was the rock (petra enim erat Christus), upon which also Peter himself was built; for other foundation can no man lay, than that is laid, which is Jesus Christ. Thus the church, which is built upon Christ, has received from him, in the person of Peter, the keys of heaven; that is, the power of binding and loosing sins.”\footnote{Tract. in Evang. Joannis, 124, § 5. The original is quoted among others by Dr. Gieseler, i. 2, p. 210 (4th ed.), but with a few unessential omissions.} This Augustinian interpretation of the petra has since been revived by some Protestant theologians in the cause of anti-Romanism.\footnote{Especially by Calov in the Lutheran church, and quite recently by Dr. Wordsworth in the Church of England (Commentary on Matt. xvi. 18). But Dr. Alford decidedly protests against it, with most of the modern commentators.} Augustine, it is true, unquestionably understood by the church the visible Catholic church, descended from the apostles, especially from Peter, through the succession of bishops; and according to the usage of his time he called the Roman church by eminence the sedes apostolica.\footnote{De utilit. credendi, § 35, he traces the development of the church “ab apostolica sede per successiones apostolorum;” and Epist. 43, he incidentally speaks of the “Romana ecclesia in qua semper apostolicae cathedrae viguit principatus.” Greenwood, i. 296 sq., thus resolves the apparent contradiction in Augustine: “In common with the age in which he lived, he (St. Augustine) was himself possessed with the idea of a visible representative unity, and considered that unity as equally the subject of divine precept and institution with the church-spiritual itself. The spiritual unity might therefore stand upon the \emph{faith} of Peter, while the outward and visible oneness was inherent in his person; so that while the church derived her esoteric and spiritual character from the faith which Peter had confessed, she received her external or executive powers from Peter through ’the succession of bishops’ sitting in Peter’s chair. Practically, indeed, there was little to choose between the two theories.” Comp. also the thorough exhibition of the Augustinian theory of the Catholic church and her attributes by Dr. Rothe, in his work Die Anfänge der christlichen Kirche, i. p. 679-711.} But on the other hand, like Cyprian and Jerome, he lays stress upon the essential unity of the episcopate, and insists that the keys of the kingdom of heaven were committed not to a single man, but to the whole church, which Peter was only set to represent.\footnote{De diversis Serm. 108: Has enim claves non homo unus, sed unitas accepit ecclesiae. Hinc ergo Petri excellentia praedicatur, quia ipsius universitatis et unitatis figuram gessit quando ei dictum est: \emph{tibi} trado, quod \emph{omnibus} traditum est, etc.} With this view agrees the independent position of the North African church in the time of Augustine toward Rome, as we have already observed it in the case of the appeal of Apiarius, and as it appears in the Pelagian controversy, of which Augustine was the leader. This father, therefore, can at all events be cited only as a witness to the limited authority of the Roman chair. And it should also, in justice, be observed, that in his numerous writings he very rarely speaks of that authority at all, and then for the most part incidentally; showing that he attached far less importance to this matter than the Roman divines.\footnote{Bellarmine, in Praef. in Libr. de Pontif., calls this article even \emph{rem summam fidei Christiana}!}

The later Latin fathers of the fourth and fifth centuries prefer the reference of the petra to Peter and his confession, and transfer his prerogatives to the Roman bishops as his successors, but produce no new arguments. Among them we mention Maximus of Turin (about 450), who, however, like Ambrose, places Paul on a level with Peter;\footnote{Hom. v., on the feast of Peter and Paul. To the one, says he, the keys of knowledge were committed, to the other the keys of power.” Eminent inter universos apostlos et peculiari quadam praerogativa praecellunt. Verum inter ipsos quis cui praeponatur, incertum est.” The same sentence in Ambrose, De Spir. S. ii. 12.} then Orosius, and several popes; above all Leo, of whom we shall speak more fully in the following section.

2. As to the Greek fathers: Eusebius, Cyril of Jerusalem, Basil, the two Gregories, Ephraim, Syrus, Asterius, Cyril of Alexandria, Chrysostom, and Theodoret refer the petra now to the confession, now to the person, of Peter; sometimes to both. They speak of this apostle uniformly in very lofty terms, at times in rhetorical extravagance, calling him the “coryphaeus of the choir of apostles,” the prince of the apostles,” the “tongue of the apostles,” the “bearer of the keys,” the “keeper of the kingdom of heaven,” the “pillar,” the “rock,” the “firm foundation of the church.” But, in the first place, they understand by all this simply an honorary primacy of Peter, to whom that power was but first committed, which the Lord afterward conferred on all the apostles alike; and, in the second place, they by no means favor an exclusive transfer of this prerogative to the bishop of Rome, but claim it also for the bishops of Antioch, where Peter, according to Gal. ii., sojourned a long time, and where, according to tradition, he was bishop, and appointed a successor.

So Chrysostom, for instance, calls Ignatius of Antioch a “successor of Peter, on whom, after Peter, the government of the church devolved,”\footnote{In S. Ignat. Martyr., n. 4.} and in another place says still more distinctly: “Since I have named Peter, I am reminded of another Peter [Flavian, bishop of Antioch], our common father and teacher, who has inherited as well the virtues as the chair of Peter. Yea, for this is the privilege of this city of ours [Antioch], to have first (ἐν ἀρχῇ) had the coryphaeus of the apostles for its teacher. For it was proper that the city, where the Christian name originated, should receive the first of the apostles for its pastor. But after we had him for our teacher, we, did not retain him, but transferred him to imperial Rome.”\footnote{Hom. ii. in Principium Actorum, n. 6, tom. iii. p. 70 (ed. Montfaucon). The last sentence (\textgreek{ἀλλὰ προσεχωρήσαμεν τῇ βασιλίδι Ρώμῃ}) is by some regarded as a later interpolation in favor of the papacy. But it contains no concession of superiority. Chrysostomimmediately goes on to say: “We have indeed not retained the body of Peter, but we have retained the faith of Peter; and while we retain his faith, we have himself.”}

Theodoret also, who, like Chrysostom, proceeded from the Antiochian school, says of the “great city of Antioch,” that it has the “throne of Peter.”\footnote{Epist. 86.} In a letter to Pope Leo he speaks, it is true, in very extravagant terms of Peter and his successors at Rome, in whom all the conditions, external and internal, of the highest eminence and control in the church are combined.\footnote{Epist. 113. Comp. Bennington and Kirk, l.c. p. 91-93. In the Epist. 116, to Renatus, one of the three papal legates at Ephesus, where he entreats his intercession with Leo, he ascribes to the Roman see the control of the church of the world (\textgreek{τῶν κατὰ την οἰκουμένην ἐκκλησιῶν τὴν ἡγεμονίαν}), but certainly in the oriental sense of an honorary supervision.} But in the same epistle he remarks, that the “thrice blessed and divine double star of Peter and Paul rose in the East and shed its rays in every direction;” in connection with which it must be remembered that he was at that time seeking protection in Leo against the Eutychian robber-council of Ephesus (449), which had unjustly deposed both himself and Flavian of Constantinople.

His bitter antagonist also, the arrogant and overbearing Cyril of Alexandria, descended some years before, in his battle against Nestorius, to unworthy flattery, and called Pope Coelestine “the archbishop of the whole [Roman] world.”\footnote{\textgreek{Ἀρχιεπίσκοπον πάσης τῆς οἰκουμένης} [i. e., of the Roman empire, according to the well-known \emph{usus loquendi,} even of the N. T., Comp. Luke ii. 1], \textgreek{πατέρα τε καὶ πατριάρχην Κελεστῖνον τὸν τῆς μεγαλοπόλεως Ρώμης}. Encom. in S. Mar. Deip. (tom. v. p. 384). Comp. his Ep. ix. ad Coelest.} The same prelates, under other circumstances, repelled with proud indignation the encroachments of Rome on their jurisdiction.

\xsection{The Decrees of Councils on the Papal Authority}

§ 62. The Decrees of Councils on the Papal Authority.

Much more important than the opinions of individual fathers are the formal decrees of the councils.

First mention here belongs to the council of Sardica in Illyria (now Sofia in Bulgaria) in 343,\footnote{That this is the true date appears from the recently discovered Festival Epistles of Athanasius, published in Syriac by Cureton (London, 1848), in an English translation by Williams (Oxford, 1854), and in German by Larsow (Leipzig, 1852). Mansi puts the council in the year 344, but most writers, including Gieseler, Neander, Milman, and Greenwood, following the erroneous statement of Socrates (ii. 20) and Sozomen (iii. 12), place it in the year 347. Comp. on the subject Larsow, Die Festbriefe des Athanasius, p. 31; and Hefele, Conciliengesch. i. p. 513 sqq.} during the Arian controversy. This council is the most favorable of all to the Roman claims. In the interest of the deposed Athanasius and of the Nicene orthodoxy it decreed:

(1) That a deposed bishop, who feels he has a good cause, may apply, out of reverence to the memory of the apostle Peter, to the Roman bishop Julius, and shall leave it with him either to ratify the deposition or to summon a new council.

(2) That the vacant bishopric shall not be filled till the decision of Rome be received.

(3) That the Roman bishop, in such a case of appeal, may, according to his best judgment, either institute a new trial by the bishops of a neighboring province, or send delegates to the spot with full power to decide the matter with the bishops.\footnote{Can. 3, 4, and 5 (in the Latin translation, can. 3, 4, and 7), in Mansi, iii. 23 sq., and in Hefele, i. 539 sqq., where the Greek and the Latin Dionysian text is given with learned explanations. The Greek and Latin texts differ in some points.}

Thus was plainly committed to the Roman bishops an appellate and revisory jurisdiction in the case of a condemned or deposed bishop even of the East. But in the first place this authority is not here acknowledged as a right already existing in practice. It is conferred as a new power, and that merely as an honorary right, and as pertaining only to the bishop Julius in person.\footnote{So the much discussed \emph{canones} are explained not only by Protestant historians, but also by Catholic of the Gallican school, like Peter de Marca, Quesnel, Du-Pin, Richer, Febronius. This interpretation agrees best with the whole connection; with the express mention of Julius (which is lacking indeed, in the Latin translation of Prisca and in Isidore, but stands distinctly in the Greek and Dionysian texts: \textgreek{Ἰουλίῳ τῷ ἐπισκόπῳ Ῥώμης}, Julio Romano episcopo); with the words, ” Si vobis placet” (can. 3), whereby the appeal in question is made dependent first on the decree of this council; and finally, with the words, “Sancti Petri apostoli memoriam honoremus,” which represent the Roman bishop’s right of review as an honorary matter. What Hefele urges against these arguments (i. 548 sq.), seems to me very insufficient.} Otherwise, either this bishop would not be expressly named, or his successors would be named with him. Furthermore, the canons limit the appeal to the case of a bishop deposed by his comprovincials, and say nothing of other cases. Finally, the council of Sardica was not a general council, but only a local synod of the West, and could therefore establish no law for the whole church. For the Eastern bishops withdrew at the very beginning, and held an opposition council in the neighboring town of Philippopolis; and the city of Sardica, too, with the praefecture of Illyricum, at that time belonged to the Western empire and the Roman patriarchate: it was not detached from them till 379. The council was intended, indeed, to be ecumenical; but it consisted at first of only a hundred and seventy bishops, and after the recession of the seventy-six Orientals, it had only ninety-four; and even by the two hundred signatures of absent bishops, mostly Egyptian, to whom the acts were sent for their approval, the East, and even the Latin Africa, with its three hundred bishoprics, were very feebly represented. It was not sanctioned by the emperor Constantius, and has by no subsequent authority been declared ecumenical.\footnote{Baronius, Natalis Alexander, and Mansi have endeavored indeed to establish for the council an ecumenical character, but in opposition to the weightiest ancient and modern authorities of the Catholic church. Comp. Hefele, i. 596 sqq,} Accordingly its decrees soon fell into oblivion, and in the further course of the Arian controversy, and even throughout the Nestorian, where the bishops of Alexandria, and not those of Rome, were evidently at the head of the orthodox sentiment, they were utterly unnoticed.\footnote{It is also to be observed, that the synodal letters, as well as the orthodox ecclesiastical writers of this and the succeeding age, which take notice of this council, like Socrates, Sozomen, Theodoret, and Basil, make no mention of those decrees concerning Rome.} The general councils of 381, 451, and 680 knew nothing of such a supreme appellate tribunal, but unanimously enacted, that all ecclesiastical matters, without exception, should first be decided in the provincial councils, with the right of appeal—not to the bishop of Rome, but to the patriarch of the proper diocese. Rome alone did not forget the Sardican decrees, but built on this single precedent a universal right. Pope Zosimus, in the case of the deposed presbyter Apiarius of Sicca (a.d. 417–418), made the significant mistake of taking the Sardican decrees for Nicene, and thus giving them greater weight than they really possessed; but he was referred by the Africans to the genuine text of the Nicene canon. The later popes, however, transcended the Sardican decrees, withdrawing from the provincial council, according to the pseudo-Isidorian Decretals, the right of deposing a bishop, which had been allowed by Sardica, and vesting it, as a causa major, exclusively in themselves.

Finally, in regard to the four great ecumenical councils, the first of Nice, the first of Constantinople, that of Ephesus, and that of Chalcedon: we have already presented their position on this question in connection with their legislation on the patriarchal system.\footnote{Comp. § 56.} We have seen that they accord to the bishop of Rome a precedence of honor among the five officially coequal patriarchs, and thus acknowledge him primus inter pares, but, by that very concession, disallow his claims to supremacy of jurisdiction, and to monarchical authority over the entire church. The whole patriarchal system, in fact, was not monarchy, but oligarchy. Hence the protest of the Roman delegates and of Pope Leo against the decrees of the council of Chalcedon in 451, which coincided with that of Constantinople in 381. This protest was insufficient to annul the decree, and in the East it made no lasting impression; for the subsequent incidental concessions of Greek patriarchs and emperors, like that of the usurper Phocas in 606, and even of the sixth ecumenical council of Constantinople in 680, to the see of Rome, have no general significance, but are distinctly traceable to special circumstances and prejudices.

It is, therefore, an undeniable historical fact, that the greatest dogmatic and legislative authorities of the ancient church bear as decidedly against the specific papal claims of the Roman bishopric, is in favor of its patriarchal rights and an honorary primacy in the patriarchal oligarchy. The subsequent separation of the Greek church from the Latin proves to this day, that she was never willing to sacrifice her independence to Rome, or to depart from the decrees of her own greatest councils.

Here lies the difference, however, between the Greek and the Protestant opposition to the universal monarchy of the papacy. The Greek church protested against it from the basis of the oligarchical patriarchal hierarchy of the fifth century; in an age, therefore, and upon a principle of church organization, which preceded the grand agency of the papacy in the history of the world. The evangelical church protests against it on the basis of a freer conception of Christianity, seeing in the papacy an institution, which indeed formed the legitimate development of the patriarchal system, and was necessary for the training of the Romanic and Germanic Nations of the middle ages, but which has virtually fulfilled its mission and outlived itself. The Greek church never had a papacy; the evangelical historically implies one. The papacy stands between the age of the patriarchal hierarchy and the age of the Reformation, like the Mosaic theocracy between the patriarchal period and the advent of Christianity. Protestantism rejects at once the papal monarchy and the patriarchal oligarchy, and thus can justify the former as well as the latter for a certain time and a certain stage in the progress of the Christian world.

\xsection{Leo the Great. a.d. 440-461}

§ 63. Leo the Great. a.d. 440–461.

I. St. Leo Magnus: Opera omnia (sermones et epistolae), ed. Paschas. Quesnel., Par. 1675, 2 vols. 4to. (Gallican, and defending Hilary against Leo, hence condemned by the Roman Index); and ed. Petr. et Hieron. Ballerini (two very learned brothers and presbyters, who wrote at the request of Pope Benedict XIV.), Venet. 1753–1757, 3 vols. fol. (Vol. i. contains 96 Sermons and 173 Epistles, the two other volumes doubtful writings and learned dissertations.) This edition is reprinted in Migne’s Patrologiae Cursus completus, vol. 54–57, Par. 1846.

II. Acta Sanctorum: sub Apr. 11 (Apr. tom. ii. p. 14–30, brief and unsatisfactory). Tillemont: Mem. t. xv. p. 414–832 (very full). Butler: Lives of the Saints, sub Apr. 11. W. A. Arendt (R.C.): Leo der Grosse u. seine Zeit, Mainz, 1835 (apologetic and panegyric). Edw. Perthel: P. Leo’s I. Leben u. Lehren, Jena, 1843 (Protestant). Fr. Boehringer: Die Kirche Christi u. ihre Zeugen, Zürich, 1846, vol. i. div. 4, p. 170–309. Ph. Jaffé: Regesta Pontif. Rom., Berol. 1851, p. 34 sqq. Comp. also Greenwood: Cathedra Petri, Lond. 1859, vol. i. bk. ii. chap. iv.-vi. (The Leonine Period); and H. H. Milman: Hist. of Latin Christianity, Lond. and New York, 1860, vol. i. bk. ii. ch. iv.

In most of the earlier bishops of Rome the person is eclipsed by the office. The spirit of the age and public opinion rule the bishops, not the bishops them. In the preceding period, Victor in the controversy on Easter, Callistus in that on the restoration of the lapsed, and Stephen in that on heretical baptism, were the first to come out with hierarchical arrogance; but they were somewhat premature, and found vigorous resistance in Irenaeus, Hippolytus, and Cyprian, though on all three questions the Roman view at last carried the day.

In the period before us, Damasus, who subjected Illyria to the Roman jurisdiction, and established the authority of the Vulgate, and Siricius, who issued the first genuine decretal letter, trod in the steps of those predecessors. Innocent I. (402–417) took a step beyond, and in the Pelagian controversy ventured the bold assertion, that in the whole Christian world nothing should be decided without the cognizance of the Roman see, and that, especially in matters of faith, all bishops must turn to St. Peter.\footnote{Ep. ad Conc. Cartha. and Ep. ad Concil. Milev., both in 416. In reference to this decision, which went against Pelagius, Augustineuttered the word so often quoted by Roman divines: ”\emph{Causa finita est}; utinam aliquando finiatur error.” But when Zosimus, the successor of Innocent, took the part of Pelagius, Augustineand the African church boldly opposed him, and made use of the Cyprianic right of protest.“Circumstances alter cases.”}

But the first pope, in the proper sense of the word, is Leo I., who justly bears the title of “the Great” in the history of the Latin hierarchy. In him the idea of the Papacy, as it were, became flesh and blood. He conceived it in great energy and clearness, and carried it out with the Roman spirit of dominion, so far as the circumstances of the time at all allowed. He marks the same relative epoch in the development of the papacy, as Cyprian in the history of the episcopate. He had even a higher idea of the prerogatives of the see of Rome than Gregory the Great, who, though he reigned a hundred and fifty years later, represents rather the patriarchal idea than the papal. Leo was at the same time the first important theologian in the chair of Rome, surpassing in acuteness and depth of thought all his predecessors, and all his successors down to Gregory I. Benedict XIV. placed him (a.d. 1744) in the small class of doctores ecclesiae, or authoritative teachers of the catholic faith. He battled with the Manichaean, the Priscillianist, the Pelagian, and other heresies, and won an immortal name as the finisher of the orthodox doctrine of the person of Christ.

The time and place of the birth and earlier life of Leo are unknown. His letters, which are the chief source of information, commence not before the year 442. Probably a Roman\footnote{As Quesnel and most of his successors infer from Prosper’s Chronicle, and a passage in Leo’s Ep. 31, c. 4, where he assigns among the reasons for not attending the council at Ephesus in 449, that he could not “deserere \emph{patriam} et sedem apostolicam.” \emph{Patria}, however, may as well mean Italy, or at least the diocese of Rome, including the ten suburbican provinces. In the Liber pontificalis he is called “natione \emph{Tuscus},\emph{“} but in two manuscript copies, “natione \emph{Romanus}.” Canisius, in the Acta Sanctorum, adopts the former view. Butler reconciles the difficulty by supposing that he was descended of a noble Tuscan family, but born at Rome.}—if not one by birth, he was certainly a Roman in the proud dignity of his spirit and bearing, the high order of his legislative and administrative talent, and the strength and energy of his will—he distinguished himself first under Coelestine (423–432) and Sixtus III. (432–440) as archdeacon and legate of the Roman church. After the death of the latter, and while himself absent in Gaul, he was elected pope by the united voice of clergy, senate, and people, and continued in that office one-and-twenty years (440–461). His feelings at the assumption of this high office, he himself thus describes in one of his sermons: “Lord, I have beard your voice calling me, and I was afraid: I considered the work which was enjoined on me, and I trembled. For what proportion is there between the burden assigned to me and my weakness, this elevation and my nothingness? What is more to be feared than exaltation without merit, the exercise of the most holy functions being intrusted to one who is buried in sin? Oh, you have laid upon me this heavy burden, bear it with me, I beseech you be you my guide and my support.”

During the time of his pontificate he was almost the only great man in the Roman empire, developed extraordinary activity, and took a leading part in all the affairs of the church. His private life is entirely unknown, and we have no reason to question the purity of his motives or of his morals. His official zeal, and all his time and strength, were devoted to the interests of Christianity. But with him the interests of Christianity were identical with the universal dominion of the Roman church.

He was animated with the unwavering conviction that the Lord himself had committed to him, as the successor of Peter, the care of the whole church.\footnote{Ep. v. ad Episcopos Metrop. per Illyricum constitutos, c. 2 (ed. Ball. i. 617, in Migne’s Patristic Libr. vol. liv. p. 515): “Quia per \emph{omnes} ecclesias cura nostra distenditur, \emph{exigente} hoc a nobis \emph{Domino}, qui apostolicae dignitatis beatissimo apostolo Petro primatum fidei suae remuneratione \emph{commisit, universalem} ecclesiam in fundamenti ipsius [Quesnel proposes \emph{istius} for \emph{ipsius}] soliditate constituens, necessitatem sollicitudinis quam habemus, cum his qui nobis collegii caritate juncti sunt, sociamus.”} He anticipated all the dogmatical arguments by which the power of the papacy was subsequently established. He refers the petra, on which the church is built, to Peter and his confession. Though Christ himself—to sum up his views on the subject—is in the highest sense the rock and foundation, besides which no other can be laid, yet, by transfer of his authority, the Lord made Peter the rock in virtue of his great confession, and built on him the indestructible temple of his church. In Peter the fundamental relation of Christ to his church comes, as it were, to concrete form and reality in history. To him specially and individually the Lord intrusted the keys of the kingdom of heaven; to the other apostles only in their general and corporate capacity. For the faith of Peter the Lord specially prayed in the hour of his passion, as if the standing of the other apostles would be the firmer, if the mind of their leader remained unconquered. On Peter rests the steadfastness of the whole apostolic college in the faith. To him the Lord, after his resurrection, committed the care of his sheep and lambs. Peter is therefore the pastor and prince of the whole church, through whom Christ exercises his universal dominion on earth. This primacy, however, is not limited to the apostolic age, but, like the faith of Peter, and like the church herself, it perpetuates itself; and it perpetuates itself through the bishops of Rome, who are related to Peter as Peter was related to Christ. As Christ in Peter, so Peter in his successors lives and speaks and perpetually executes the commission: “Feed my sheep.” It was by special direction of divine providence, that Peter labored and died in Rome, and sleeps with thousands of blessed martyrs in holy ground. The centre of worldly empire alone can be the centre of the kingdom of God. Yet the political position of Rome would be of no importance without the religious considerations. By Peter was Rome, which had been the centre of all error and superstition, transformed into the metropolis of the Christian world, and invested with a spiritual dominion far wider than her former earthly empire. Hence the bishopric of Constantinople, not being a sedes apostolica, but resting its dignity on a political basis alone, can never rival the Roman, whose primacy is rooted both in divine and human right. Antioch also, where Peter only transiently resided, and Alexandria, where he planted the church through his disciple Mark, stand only in a secondary relation to Rome, where his bones repose, and where that was completed, which in the East was only laid out. The Roman bishop is, therefore, the primus omnium episcoporum, and on him devolves the plenitudo potestatis, the solicitudo omnium pastorum, and communis cura universalis ecclesiae.\footnote{These views Leo repeatedly expresses in his sermons on the festival of St. Peter and on the anniversary of his own elevation, as well as in his official letters to the African, Illyrian, and South Gallic bishops, to Dioscurus of Alexandria, to the patriarch Anatolius of Constantinople, to the emperor Marcian and the empress Pulcheria. Particular proof passages are unnecessary. Comp. especially Ep. x., xi., xii., xiv., civ.-cvi. (ed. Baller.), and Perthel, l.c. p. 226-241, where the chief passages are given in full.}

Leo thus made out of a primacy of grace and of personal fitness a primacy of right and of succession. Of his person, indeed, he speaks in his sermons with great humility, but only thereby the more to exalt his official character. He tells the Romans, that the true celebration of the anniversary of his accession is, to recognize, honor, and obey, in his lowly person, Peter himself, who still cares for shepherd and flock, and whose dignity is not lacking even to his unworthy heir.\footnote{“Cujus dignitas etiam in indigno haerede non deficit,” Sermo iii. in Natal, ordin. c. 4 (vol. i. p. 13, ed. Ball.).“Etsi necessarium est trepidare de merito, religiosum est tamen gaudere de dono: quoniam qui mihi oneris est auctor, ipse est administrationis adjutor.” Serm. ii. c. 1.} Here, therefore, we already have that characteristic combination of humility and arrogance, which has stereotyped itself in the expressions: “Servant of the servants of God,” “vicar of Christ,” and even “God upon earth.” In this double consciousness of his personal unworthiness and his official exaltation, Leo annually celebrated the day of his elevation to the chair of Peter. While Peter himself passes over his prerogative in silence, and expressly warns against hierarchical assumption,\footnote{Pet. v. 3.} Leo cannot speak frequently and emphatically enough of his authority. While Peter in Antioch meekly submits to the rebuke of the junior apostle Paul,\footnote{Gal. ii. 11.} Leo pronounces resistance to his authority to be impious pride and the sure way to hell.\footnote{Ep. x. c. 2 (ed. Ball. i. p. 634; ed. Migne, vol. 54, p. 630), to the Gallican bishops in the matter of Hilary: “Cui (sc. Petro) quisquis principatum aestimat denegandum, illius quidem nullo modo potest minuere dignitatem; sed \emph{inflatus spiritu superbiae suae semetipsum in inferna demergit}.” Comp. Ep. clxiv. 3; clvii. 3.} Obedience to the pope is thus necessary to salvation. Whosoever, says he, is not with the apostolic see, that is, with the head of the body, whence all gifts of grace descend throughout the body, is not in the body of the church, and has no part in her grace. This is the fearful but legitimate logic of the papal principle, which confines the kingdom of God to the narrow lines of a particular organization, and makes the universal spiritual reign of Christ dependent on a temporal form and a human organ. But in its very first application this papal ban proved itself a brutum fulmen, when in spite of it the Gallican archbishop Hilary, against whom it was directed, died universally esteemed and loved, and then was canonized. This very impracticability of that principle, which would exclude all Greek and Protestant Christians from the kingdom of heaven, is a refutation of the principle itself.

In carrying his idea of the papacy into effect, Leo displayed the cunning tact, the diplomatic address, and the iron consistency which characterize the greatest popes of the middle age. The circumstances in general were in his favor: the East rent by dogmatic controversies; Africa devastated by the barbarians; the West weak in a weak emperor; nowhere a powerful and pure bishop or divine, like Athanasius, Augustine, or Jerome, in the former generation; the overthrow of the Western empire at hand; a new age breaking, with new peoples, for whose childhood the papacy was just the needful school; the most numerous and last important general council convened; and the system of ecumenical orthodoxy ready to be closed with the decision concerning the relation of the two natures in Christ.

Leo first took advantage of the distractions of the North African church under the Arian Vandals, and wrote to its bishops in the tone of an acknowledged over-shepherd. Under the stress of the times, and in the absence of a towering, character like Cyprian and Augustine, the Africans submitted to his authority (443). He banished the remnants of the Manichaeans and Pelagians from Italy, and threatened the bishops with his anger, if they should not purge their churches of the heresy. In East Illyrian which was important to Rome as the ecclesiastical outpost toward Constantinople, he succeeded in regaining and establishing the supremacy, which had been acquired by Damasus, but had afterward slipped away. Anastasius of Thessalonica applied to him to be confirmed in his office. Leo granted the prayer in 444, extending the jurisdiction of Anastasius over all the Illyrian bishops, but reserving to them a right of appeal in important cases, which ought to be decided by the pope according to divine revelation. And a case to his purpose soon presented itself, in which Leo brought his vicar to feel that he was called indeed to a participation of his care, but not to a plentitude of power (plenitudo potestatis). In the affairs of the Spanish church also Leo had an opportunity to make his influence felt, when Turibius, bishop of Astorga, besought his intervention against the Priscillianists. He refuted these heretics point by point, and on the basis of his exposition the Spaniards drew up an orthodox regula fidei with eighteen anathemas against the Priscillianist error.

But in Gaul he met, as we have already, seen, with a strenuous antagonist in Hilary of Arles, and, though he called the secular power to his aid, and procured from the emperor Valentinian an edict entirely favorable to his claims, he attained but a partial victory.\footnote{Comp. above, § 59.} Still less successful was his effort to establish his primacy in the East, and to prevent his rival at Constantinople from being elevated, by the famous twenty-eighth canon of Chalcedon, to official equality with himself.\footnote{See the particulars in § 36, above, near the close} His earnest protest against that decree produced no lasting effect. But otherwise he had the most powerful influence in the second stage of the Christological controversy. He neutralized the tyranny of Dioscurus of Alexandria and the results of the shameful robber-council of Ephesus (449), furnished the chief occasion of the fourth ecumenical council, presided over it by his legates (which the Roman bishop had done at neither of the three councils before), and gave the turn to the final solution of its doctrinal problem by that celebrated letter to Flavian of Constantinople, the main points of which were incorporated in the new symbol. Yet he owed this influence by no means to his office alone, but most of all to his deep insight of the question, and to the masterly tact with which he held the Catholic orthodox mean between the Alexandrian and Antiochian, Eutychian and Nestorian extremes. The particulars of his connection with this important dogma belong, however, to the history of doctrine.

Besides thus shaping the polity and doctrine of the church, Leo did immortal service to the city of Rome, in twice rescuing it from destruction.\footnote{Comp. Pertbel, l.c. p. 90 sqq., and p. 104 sqq.} When Attila, king of the Huns, the “scourge of God,” after destroying Aquileia, was seriously threatening the capital of the world (A. D. 452), Leo, with only two companions, crozier in hand, trusting in the help of God, ventured into the hostile camp, and by his venerable form, his remonstrances, and his gifts, changed the wild heathen’s purpose. The later legend, which Raphael’s pencil has employed, adorned the fact with a visible appearance of Peter and Paul, accompanying the bishop, and, with drawn sword, threatening Attila with destruction unless he should desist.\footnote{Leo himself says nothing of his mission to Attila. Prosper, in Chron. ad ann. 452, mentions it briefly, and Canisius, in the Vita Leonis (in the \emph{Acta Sanctorum}, for the month of April, tom. ii. p. 18), with later exaggerations.} A similar case occurred several years after (455), when the Vandal king Genseric, invited out of revenge by the empress Eudoxia, pushed his ravages to Rome. Leo obtained from him the promise that at least he would spare the city the infliction of murder and fire; but the barbarians subjected it to a fourteen days’ pillage, the enormous spoils of which they transported to Carthage; and afterward the pope did everything to alleviate the consequent destitution and suffering, and to restore the churches.\footnote{Comp. Leo’s 84th Sermon, which was preached soon after the departure of the Vandals, and Prosper, Chron ad ann. 455}

Leo died in 461, and was buried in the church of St. Peter. The day and circumstances of his death are unknown.\footnote{The Roman calendar places his name on the 11th of April. But different writers fix his death on June 28, Oct. 30 (Quesnel), Nov. 4 (Pagi), Nov. 10 (Butler). Butler quotes the concession of Bower, the apostate Jesuit, who, in his Lives of the Popes, says of Leo, that “he was without doubt a man of extraordinary parts, far superior to all who had governed that church before him, and scarce equalled by any since.”}

The literary works of Leo consist of ninety-six sermons and one hundred and seventy-three epistles, including epistles of others to him. They are earnest, forcible, full of thought, churchly, abounding in bold antitheses and allegorical freaks of exegesis, and sometimes heavy, turgid, and obscure in style. His collection of sermons is the first we have from a Roman bishops In his inaugural discourse he declared preaching to be his sacred duty. The sermons are short and simple, and were delivered mostly on high festivals and on the anniversaries of his own elevation.\footnote{Sermones de natali. Canisius (in Acta Sanct., l.c. p. 17) calls Leo Christianum Demosthenem.} Other works ascribed to him, such as that on the calling of all nations,\footnote{De vocatione omnium gentium—a work praised highly even by Erasmus, Luther, Bullinger, and Grotius. Quesnel has only proved the possibility of Leo’s being the author. Comp. Perthel, l.c. p. 127 sqq. The Sacramentarium Leonis, or a collection of liturgical prayers for all the festival days of the year, contains some of his prayers, but also many which are of a later date.} which takes a middle ground on the doctrine of predestination, with the view to reconcile the Semipelagians and Augustinians, are of doubtful genuineness.

\xsection{The Papacy from Leo I to Gregory I. a.d. 461-590}

§ 64. The Papacy from Leo I to Gregory I. a.d. 461–590.

The first Leo and the first Gregory are the two greatest bishops of Rome in the first six centuries. Between them no important personage appears on the chair of Peter; and in the course of that intervening century the idea and the power of the papacy make no material advance. In truth, they went farther in Leo’s mind than they did in Gregory’s. Leo thought and acted as an absolute monarch; Gregory as first among the patriarchs; but both under the full conviction that they were the successors of Peter.

After the death of Leo, the archdeacon Hilary, who had represented him at the council of Ephesus, was elected to his place, and ruled (461–468) upon his principles, asserting the strict orthodoxy in the East and the authority of the primacy in Gaul.

His successor, Simplicius (468–483), saw the final dissolution of the empire under Romulus Augustulus (476), but, as he takes not the slightest notice of it in his epistles, he seems to have ascribed to it but little importance. The papal power had been rather favored than hindered in its growth by the imbecility of the latest emperors. Now, to a certain extent, it stepped into the imperial vacancy, and the successor of Peter became, in the mind of the Western nations, sole heir of the old Roman imperial succession.

On the fall of the empire the pope became the political subject of the barbarian and heretical (for they were Arian) kings; but these princes, as most of the heathen emperors had done, allowed him, either from policy, or from ignorance or indifference, entire freedom in ecclesiastical affairs. In Italy the Catholics had by far the ascendency in numbers and in culture. And the Arianism of the new rulers was rather an outward profession than an inward conviction. Odoacer, who first assumed the kingdom of Italy (476–493), was tolerant toward the orthodox faith, yet attempted to control the papal election in 483 in the interest of the state, and prohibited, under penalty of the anathema, the alienation of church property by any bishop. Twenty years later a Roman council protested against this intervention of a layman, and pronounced the above prohibition null and void, but itself passed a similar decree against the alienation of church estates.\footnote{This was the fifth (al. fourth) council under, Symmachus, held in Nov. 502, therefore later than the \emph{synodus palmaris}. Comp. Hefele, ii. p. 625 sq.}

Pope Felix II., or, according to another reckoning,III. (483–492), continued the war of his predecessor against the Monophysitism of the East, rejected the Henoticon of the emperor Zeno, as an unwarrantable intrusion of a layman in matters of faith, and ventured even the excommunication of the bishop Acacius of Constantinople. Acacius replied with a counter anathema, with the support of the other Eastern patriarchs; and the schism between the two churches lasted over thirty years, to the pontificate of Hormisdas.

Gelasius I. (492–496) clearly announced the principle, that the priestly power is above the kingly and the imperial, and that from the decisions of the chair of Peter there is no appeal. Yet from this pope we have, on the other hand, a remarkable testimony against what he pronounces the “sacrilege” of withholding the cup from the laity, the communio sub una specie.

Anastasius II. (496–498) indulged in a milder tone toward Constantinople, and incurred the suspicion of consent to its heresy.\footnote{. Dante puts him in hell, and Baronius ascribes his sudden death to an evident judgment of God.}

His sudden death was followed by a contested papal election, which led to bloody encounters. The Ostrogothic king Theodoric (the Dietrich of Bern in the Niebelungenlied), the conqueror and master of Italy (493–526), and, like Odoacer, an Arian, was called into consultation in this contest, and gave his voice for Symmachus against Laurentius, because Symmachus had received the majority of votes, and had been consecrated first. But the party of Laurentius, not satisfied with this, raised against Symmachus the reproach of gross iniquities, even of adultery and of squandering the church estates. The bloody scenes were renewed, priests were murdered, cloisters were burned, and nuns were insulted. Theodoric, being again called upon by the senate for a decision, summoned a council at Rome, to which Symmachus gave his consent; and a synod, convoked by a heretical king, must decide upon the pope! In the course of the controversy several councils were held in rapid succession, the chronology of which is disputed.\footnote{Comp. Hefele, ii. p. 615 sqq.} The most important was the synodus palmaris,\footnote{So named from the building in Rome, in which it was held: “A porticu beati Petri Apostoli, quae appellatur ad Palmaria,” as Anastasius says. In the histories of councils it is erroneously given as Synodus III. Many historians, Gieseler among them, place it in the year 503.} the fourth council under Symmachus, held in October, 501. It acquitted this pope without investigation, on the presumption that it did not behove the council to pass judgment respecting the successor of St. Peter. In his vindication of this council—for the opposition was not satisfied with it—the deacon Ennodius, afterward bishop of Pavia († 521), gave the first clear expression to the absolutism upon which Leo had already acted: that the Roman bishop is above every human tribunal, and is responsible only to God himself.\footnote{Libellus apologeticus pro Synodo IV. Romana, in Mansi, viii. 274. This vindication was solemnly adopted by the sixth Roman council under Symmachus, in 503, and made equivalent to a decree of council.} Nevertheless, even in the middle age, popes were deposed and set up by emperors and general councils. This is one of the points of dispute between the absolute papal system and the constitutional episcopal system in the Roman church, which was left unsettled even by the council of Trent.

Under Hormisdas (514–523) the Monophysite party in the Greek church was destroyed by the energetic zeal of the orthodox emperor Justin, and in 519 the union of that church with Rome was restored, after a schism of five-and-thirty years.

Theodoric offered no hinderance to the transactions and embassies, and allowed his most distinguished subject to assert his ecclesiastical supremacy over Constantinople. This semi-barbarous and heretical prince was tolerant in general, and very liberal toward the Catholic church; even rising to the principle, which has waited till the modern age for its recognition, that the power of the prince should be restricted to civil government, and should permit no trespass on the conscience of its subjects.” No one,” says he, “shall be forced to believe against his will.” Yet, toward the close of his reign, on mere political suspicion, he ordered the execution of the celebrated philosopher Boethius, with whom the old Roman literature far more worthily closes, than the Roman empire with Augustulus; and on the same ground he caused the death of the senator Symmachus and the incarceration of Pope John I. (523–526).

Almost the last act of his reign was the nomination of the worthy Felix III. (IV.) to the papal chair, after a protracted struggle of contending parties. With the appointment he issued the order that hereafter, as heretofore, the pope should be elected by clergy and people, but should be confirmed by the temporal prince before assuming his office; and with this understanding the clergy and the city gave their consent to the nomination.

Yet, in spite of this arrangement, in the election of Boniface II. (530–532) and John II. (532–535) the same disgraceful quarrelling and briberies occurred;—a sort of chronic disease in the history of the papacy.

Soon after the death of Theodoric (526) the Gothic empire fell to pieces through internal distraction and imperial weakness. Italy was conquered by Belisarius (535), and, with Africa, again incorporated with the East Roman empire, which renewed under Justinian its ancient splendor, and enjoyed a transient after-summer. And yet this powerful, orthodox emperor was a slave to the intriguing, heretical Theodora, whom he had raised from the theatre to the throne; and Belisarius likewise, his victorious general, was completely under the power of his wife Antonina.

With the conquest of Italy the popes fell into a perilous and unworthy dependence on the emperor at Constantinople, who reverenced, indeed, the Roman chair, but not less that of Constantinople, and in reality sought to use both as tools of his own state-church despotism. Agapetus (535–536) offered fearless resistance to the arbitrary course of Justinian, and successfully protested against the elevation of the Eutychian Anthimus to the patriarchal see of Constantinople. But, by the intrigues of the Monophysite empress, his successor, Pope Silverius (a son of Hormisdas, 536–538), was deposed on the charge of treasonable correspondence with the Goths, and banished to the island of Pandataria, whither the worst heathen emperors used to send the victims of their tyranny, and where in 540 he died—whether a natural or a violent death, we do not know.

Vigilius, a pliant creature of Theodora, ascended the papal chair under the military protection of Belisarius (538–554). The empress had promised him this office and a sum of money, on condition that he nullify the decrees of the council of Chalcedon, and pronounce Anthimus and his friends orthodox. The ambitious and doubled-tongued prelate accepted the condition, and accomplished the deposition, and perhaps the death, of Silverius. In his pontificate occurred the violent controversy of the three chapters and the second general council of Constantinople (553). His administration was an unprincipled vacillation between the dignity and duties of his office and subservience to an alien theological and political influence; between repeated condemnation of the three chapters in behalf of a Eutychianizing spirit, and repeated retraction of that condemnation. In Constantinople, where he resided several years at the instance of the emperor, he suffered much personal persecution, but without the spirit of martyrdom, and without its glory. For example, at least according to Western accounts, he was violently torn from the altar, upon which he was holding with both hands so firmly that the posts of the canopy fell in above him; he was dragged through the streets with a rope around his neck, and cast into a common prison; because he would not submit to the will of Justinian and his council. Yet he yielded at last, through fear of deposition. He obtained permission to return to Rome, but died in Sicily, of the stone, on his way thither (554).

Pelagius I. (554–560), by order of Justinian, whose favor he had previously gained as papal legate at Constantinople, was made successor of Vigilius, but found only two bishops ready to consecrate him. His close connection with the East, and his approval of the fifth ecumenical council, which was regarded as a partial concession to the Eutychian Christology, and, so far, an impeachment of the authority of the council of Chalcedon, alienated many Western bishops, even in Italy, and induced a temporary suspension of their connection with Rome. He issued a letter to the whole Christian world, in which he declared his entire agreement with the first four general councils, and then vindicated the fifth as in no way departing from the Chalcedonian dogma. But only by the military aid of Narses could he secure subjection; and the most refractory bishops, those of Aquileia and Milan, he sent as prisoners to Constantinople.

In these two Justinian-made popes we see how much the power of the Roman hierarchy was indebted to its remoteness from the Byzantine despotism, and how much it was injured by contact with it.

With the descent of the Arian Longobards into Italy, after 668, the popes again became more independent of the Byzantine court. They continued under tribute indeed to the ex-archs in Ravenna, as the representatives of the Greek emperors (from 554), and were obliged to have their election confirmed and their inauguration superintended by them. But the feeble hold of these officials in Italy, and the pressure of the Arian barbarians upon them, greatly favored the popes, who, being the richest proprietors, enjoyed also great political consideration in Italy, and applied their influence to the maintenance of law and order amidst the reigning confusion.

In other respects the administrations of John III. (560–573), Benedict I. (574–578), and Pelagius II. (578–590), are among the darkest and the most sterile in the annals of the papacy.

But with Gregory I. (590–604) a new period begins. Next to Leo I. he was the greatest of the ancient bishops of Rome, and he marks the transition of the patriarchal system into the strict papacy of the middle ages. For several reasons we prefer to place him at the head of the succeeding period. He came, it is true, with more modest claims than Leo, who surpassed him in boldness, energy, and consistency. He even solemnly protested, as his predecessor Pelagius II. had done, against the title of universal bishop, which the Constantinopolitan patriarch, John Jejunator, adopted at a council in 587;\footnote{Even Justinian repeatedly applied to the patriarch of Constantinople officially the title \textgreek{οἰκομενικὸς πατριάρχης}, \emph{universalis patriarcha}.} he declared it an antichristian assumption, in terms which quite remind us of the patriarchal equality, and seem to form a step in recession from the ground of Leo. But when we take his operations in general into view, and remember the rigid consistency of the papacy, which never forgets, we are almost justified in thinking, that this protest was directed not so much against the title itself, as against the bearer of it, and proceeded more from jealousy of a rival at Constantinople, than from sincere humility.\footnote{Bellarmine disposes of this apparent testimony of one of the greatest and best popes against the system of popery, which has frequently been urged since Calvin by Protestant controversialists, by assuming that the term \emph{episcopus universalis} is used in two very different senses.“Respondeo,” he says (in his great controversial work, De controversiis christianae fidei, etc., de Romano pontifice, lib. ii. cap. 31), “duobus modis posse intelligi nomen universalis episcopi. Uno modo, ut ille, qui dicitur universalis, intelligatur esse solus episcopus omnium urbium Christianarum, ita ut caeteri non sint episcopi, sed vicarii tantum illius, qui dicitur episcopus universalis, et hoc modo nomen hoc est vere profanum, sacrilegum et antichristianum.... Altero modo dici potest episcopus universalis, qui habet curam totius ecclesiae, sed generalem, ita ut non excludat particulares episcopos. Et hoc modo nomen hoc posse tribui Romano pontifici ex mente Gregorii probatur.”} From the same motive the Roman bishops avoided the title of patriarch, as placing them on a level with the Eastern patriarchs, and preferred the title of pope, from a sense of the specific dignity of the chair of Peter. Gregory is said to have been the first to use the humble-proud title: “Servant of the servants of God.” His successors, notwithstanding his protest, called themselves “the universal bishops” of Christendom. What he had condemned in his oriental colleagues as antichristian arrogance, the later popes considered but the appropriate expression of their official position in the church universal.

\xsection{The Synodical System. The Ecumenical Councils}

§ 65. The Synodical System. The Ecumenical Councils.

I. The principal sources are the Acts of the Councils, the best and most complete collections of which are those of the Jesuit Sirmond (Rom. 1608–1612, 4 vols. fol.); the so-called Collectio regia (Paris, 1644, 37 vols. fol.; a copy of it in the Astor Libr., New York); but especially those of the Jesuit Hardouin († 1729): Collectio maxima Conciliorum generalium et provincialium (Par. 1715 sqq., 12 vols. fol.), coming down to 1714, and very available through its five copious indexes (tom. i. and ii. embrace the first six centuries; a copy of it, from Van Ess’s library, in the Union Theol. Sem. Library, at New York); and the Italian Joannes Dominicus Mansi (archbishop of Lucca, died 1769): Sacrorum Conciliorum nova et amplissima collection, Florence, 1759-’98, in 31 (30) vols. fol. This is the most complete and the best collection down to the fifteenth century, but unfinished, and therefore without general indexes; tom. i. contains the Councils from the beginning of Christianity to a.d. 304; tom. ii.-ix. include our period to a.d. 590 (I quote from an excellent copy of this rare collection in the Union Theol. Sem. Libr., at New York, 30 t. James Darling, in his Cyclop. Bibliographica, p. 740–756, gives the list of the contents of an earlier edition of the Councils by Nic. Coleti, Venet., 1728, in 23 vols., with a supplement of Mansi, in 6 vols. 1748-’52, which goes down to 1727, while the new edition of Mansi only reaches to 1509. Brunet, in the “Manuel Du Libraire,” quotes the edition of Mansi, Florence, 1759–1798, with the remark: “Cette collection, dont le dernier volume s’arrête à l’année 1509, est peu commune à Paris ou elle revenait à 600 fr.” Strictly speaking its stops in the middle of the 15th century, except in a few documents which reach further.) Useful abstracts are the Summa Conciliorum of Barth. Caranza, in many editions; and in the German language, the Bibliothek der Kirchenversammlungen (4th and 5th centuries), by Fuchs, Leipz., 1780–1784, 4 vols.

II. Chr. Wilh. Franz Walch (Luth.): Entwurf einer vollstaendigen Historie der Kirchenversammlungen, Leipz., 1759. Edw. H. Landon (Anglic.): A manual of Councils of the Holy Catholick Church, comprising the substance of the most remarkable and important canons, alphabetically arranged, 12mo. London, 1846. C. J. Hefele (R.C.): Conciliengeshichte, Freiburg, 1855–1863, 5 vols. (a very valuable work, not yet finished; vol. v. comes down to a.d. 1250). Comp. my Essay on Oekumenische Concilien, in Dorner’s Annals of Ger. Theol. vol. viii. 326–346.

Above the patriarchs, even above the patriarch of Rome, stood the ecumenical or general councils,\footnote{The name \textgreek{σύνοδος οἰκουμενική}(concilium universale, s. generale) occurs first in the sixth canon of the council of Constantinople in 381. The \textgreek{οἰκουμένη} (sc. \textgreek{γῆ}) is, properly, the whole inhabited earth; then, in a narrower sense, the earth inhabited by \emph{Greeks}, in distinction from the barbarian countries; finally, with the Romans, the \emph{orbis Romanus}, the political limits of which coincided with those of the ancient Graeco-Latin church. But as the bishops of the barbarians outside the empire were admitted, the ecumenical councils represented the entire Catholic Christian world.} the highest representatives, of the unity and authority of the old Catholic church. They referred originally to the Roman empire, but afterward included the adjacent barbarian countries, so far as those countries were represented in them by bishops. They rise up like lofty peaks or majestic pyramids from the plan of ancient church history, and mark the ultimate authoritative settlement of the general questions of doctrine and discipline which agitated Christendom in the Graeco-Roman empire.

The synodical system in general had its rise in the apostolic council at Jerusalem,\footnote{Acts xv., and Gal. ii. Comp. my History of the Apostolic Church, §§ 67-69 (Engl. ed., p. 245-257). Mansi, l.c. tom. i. p. 22 (De quadruplici Synodo Apostolorum), and other Roman Catholic writers, speak of \emph{four} Apostolic Synods: Acts i. 13 sqq., for the election of an apostle; ch. vi. for the election of deacons; ch. xv. for the settlement of the question of the binding authority of the law of Moses; and ch. xxi. for a similar object. But we should distinguish between a private conference and consultation, and a public synod.} and completed its development, under its Catholic form, in the course of the first five centuries. Like the episcopate, it presented a hierarchical gradation of orders. There was, first, the diocesan or district council, in which the bishop of a diocese (in the later sense of the word) presided over his clergy; then the provincial council, consisting of the metropolitan or archbishop and the bishops of his ecclesiastical province; next, the patriarchal council, embracing all the bishops of a patriarchal district (or a diocese in the old sense of the term); then the national council, inaccurately styled also general, representing either the entire Greek or the entire Latin church (like the later Lateran councils and the council of Trent); and finally, at the summit stood the ecumenical council, for the whole Christian world. There was besides these a peculiar and abnormal kind of synod, styled σύνοδος ἐνδημοῦσα, frequently held by the bishop of Constantinople with the provincial bishops resident (ἐνδημοῦντες) on the spot.\footnote{It is usually supposed there were only four or five different kinds of council. But Hefele reckons eight (i. p. 3 and 4) adding to those above named the irregular\textgreek{σύνοδοι ἐνδημοῦσαι}, also the synods of the bishops of \emph{two or more} provinces finally the \emph{concilia} \emph{mixta}, consisting of the \emph{secular} and spiritual dignitaries province, as separate classes.}

In the earlier centuries the councils assembled without fixed regularity, at the instance of present necessities, like the Montanist and the Easter controversies in the latter part of the second century. Firmilian of Cappadocia, in his letter to Cyprian, first mentions, that at his time, in the middle of the third century, the churches of Asia Minor held regular annual synods, consisting of bishops and presbyters. From that time we find an increasing number of such assemblies in Egypt, Syria, Greece, Northern Africa, Italy, Spain, and Gaul. The council of Nicaea, a.d. 325, ordained, in the fifth canon, that the provincial councils should meet twice a year: during the fast season before Easter, and in the fall.\footnote{A similar order, with different times, appears still earlier in the 37th of the apostolic canons, where it is said (in the ed. of Ueltzen, p. 244):\textgreek{Δεύτεροντοῦ ἔτους σύνοδος γενέσθω τῶν ἐπισκόπων}.} In regard to the other synods no direction was given.

The Ecumenical councils were not stated, but extraordinary assemblies, occasioned by the great theological controversies of the ancient church. They could not arise until after the conversion of the Roman emperor and the ascendancy of Christianity as the religion of the state. They were the highest, and the last, manifestation of the power of the Greek church, which in general took the lead in the first age of Christianity, and was the chief seat of all theological activity. Hence in that church, as well as in others, they are still held in the highest veneration, and kept alive in the popular mind by pictures in the churches. The Greek and Russian Christians have annually commemorated the seven ecumenical councils, since the year 842, on the first Sunday in Lent, as the festival of the triumph of orthodoxy\footnote{This Sunday, the celebration of which was ordered by the empress Theodora in 842, is called among the Greeks the \textgreek{κυριακήτῆς ὀρθοδοξίας}. On that day the ancient councils are dramatically reproduced in the public worship.} and they live in the hope that an eighth ecumenical council shall yet heal the divisions and infirmities of the Christian world. Through their symbols of faith those councils, especially of Nice and of Chalcedon, still live in the Western church, both Roman Catholic and Evangelical Protestant.

Strictly speaking, none of these councils properly represented the entire Christian world. Apart from the fact that the laity, and even the lower clergy, were excluded from them, the assembled bishops themselves formed but a small part of the Catholic episcopate. The province of North Africa alone numbered many more bishops than were present at either the second, the third, or the fifth general council.\footnote{The schismatical Donatists alone held a council at Carthage in 308, of two hundred and seventy bishops (Comp. Wiltsch, Kirchl. Geogr. u. Statistik, i. p. 53 and 54); while the second ecumenical council numbered only a hundred and fifty, the third a hundred and sixty (a hundred and ninety-eight), and the fifth a hundred and sixty-four.} The councils bore a prevailingly oriental character, were occupied with Greek controversies, used the Greek language, sat in Constantinople or in its vicinity, and consisted almost wholly of Greek members. The Latin church was usually represented only by a couple of delegates of the Roman bishop; though these delegates, it is true, acted more or less in the name of the entire West. Even the five hundred and twenty, or the six hundred and thirty members of the council of Chalcedon, excepting the two representatives of Leo I., and two African fugitives accidentally present, were all from the East. The council of Constantinople in 381 contained not a single Latin bishop, and only a hundred and fifty Greek, and was raised to the ecumenical rank by the consent of the Latin church toward the middle of the following century. On the other hand, the council of Ephesus, in 449, was designed by emperor and pope to be an ecumenical council; but instead of this it has been branded in history as the synod of robbers, for its violent sanction of the Eutychian heresy. The council of Sardica, in 343, was likewise intended to be a general council, but immediately after its assembling assumed a sectional character, through the secession and counter-organization of the Eastern bishops.

It is, therefore, not the number of bishops present, nor even the regularity of the summons alone, which determines the ecumenical character of a council, but the result, the importance and correctness of the decisions, and, above all, the consent of the orthodox Christian world.\footnote{Schröckh says (vol. viii. p. 201), unjustly, that this general consent belongs among the “empty conceits.” Of course the unanimity must be limited to \emph{orthodox} Christendom.}

The number of the councils thus raised by the public opinion of the Greek and Latin churches to the ecumenical dignity, is seven. The succession begins with the first council of Nicaea, in the year 325, which settled the doctrine of the divinity of Christ, and condemned the Arian heresy. It closes with the second council of Nice, in 787, which sanctioned the use of images in the church. The first four of these councils command high theological regard in the orthodox Evangelical churches, while the last three are less important and far more rarely mentioned.

The ecumenical councils have not only an ecclesiastical significance, but bear also a political or state-church character. The very name refers to the οἰκουμένη, the orbis Romanus, the empire. Such synods were rendered possible only by that great transformation, which is marked by the accession of Constantine. That emperor caused the assembling of the first ecumenical council, though the idea was probably suggested to him by friends among the bishops; at least Rufinus says, he summoned the council “ex sacerdotum sententia.” At all events the Christian Graeco-Roman emperor is indispensable to an ecumenical council in the ancient sense of the term; its temporal head and its legislative strength.

According to the rigid hierarchical or papistic theory, as carried out in the middle ages, and still asserted by Roman divines, the pope alone, as universal head of the church, can summon, conduct, and confirm a universal council. But the history of the first seven, or, as the Roman reckoning is, eight, ecumenical councils, from 325 to 867, assigns this threefold power to the Byzantine emperors. This is placed beyond all contradiction, by the still extant edicts of the emperors, the acts of the councils, the accounts of all the Greek historians, and the contemporary Latin sources. Upon this Byzantine precedent, and upon the example of the kings of Israel, the Russian Czars and the Protestant princes of Germany, Scandinavia, and England—be it justly or unjustly—build their claim to a similar and still more extended supervision of the church in their dominions.

In the first place, the call of the ecumenical councils emanated from the emperors.\footnote{This is conceded even by the Roman Catholic church historian Hefele (i. p. 7), in opposition to, Bellarmine and other Romish divines.“The first eight general councils,” says he, “were appointed and convoked by the \emph{emperors}; all the subsequent councils, on the contrary [i.e. all the \emph{Roman} Catholic general councils], by the popes; but even in those first councils there appears a certain \emph{participation of the popes} in their convocation, more or less prominent in particular instances.” The latter assertion is too sweeping, and can by no means be verified in the history of the first two of these councils, nor of the fifth.} They fixed the place and time of the assembly, summoned the metropolitans and more distinguished bishops of the empire by an edict, provided the means of transit, and paid the cost of travel and the other expenses out of the public treasury. In the case of the council of Nicaea and the first of Constantinople the call was issued without previous advice or consent from the bishop of Rome.\footnote{As regards the council of Nicaea: according to Eusebius and all the ancient authorities, it was called by Constantinealone; and not till three centuries later, at the council of 680, was it claimed that Pope Sylvester had any share in the convocation. As to the council of Constantinople in 381: the Roman theory, that Pope Damasus summoned it in conjunction with Theodosius, rests on a confusion of this council with another and an unimportant one of 382. Comp. the notes of Valesius to Theodoret, Hist. Ecel. v. 9; and Hefele (who here himself corrects his earlier view), vol. i. p. 8, and vol. ii. p. 36.} In the council of Chalcedon, in 451, the papal influence is for the first time decidedly prominent; but even there it appears in virtual subordination to the higher authority of the council, which did not suffer itself to be disturbed by the protest of Leo against its twenty-eighth canon in reference to the rank of the patriarch of Constantinople. Not only ecumenical, but also provincial councils were not rarely called together by Western princes; as the council of Arles in 314 by Constantine, the council of Orleans in 549 by Childebert, and—to anticipate an instance—the synod of Frankfort in 794 by Charlemagne. Another remarkable fact has been already mentioned: that in the beginning of the sixth century several Orthodox synods at Rome, for the purpose of deciding the contested election of Symmachus, were called by a secular prince, and he the heretical Theodoric; yet they were regarded as valid.

In the second place, the emperors, directly or indirectly, took an active part in all but two of the ecumenical councils summoned by them, and held the presidency. Constantine the Great, Marcian, and his wife Pulcheria, Constantine Progonatus, Irene, and Basil the Macedonian, attended in person; but generally the emperors, like the Roman bishops (who were never present themselves), were represented by delegates or commissioners, clothed with full authority for the occasion. These deputies opened the sessions by reading the imperial edict (in Latin and Greek) and other documents. They presided in conjunction with the patriarchs, conducted the entire course of the transactions, preserved order and security, closed the council, and signed the acts either at the head or at the foot of the signatures of the bishops. In this prominent position they sometimes exercised, when they had a theological interest or opinion of their own, no small influence on the discussions and decisions, though they had no votum; as the presiding officers of deliberative and legislative bodies generally have no vote, except when the decision of a question depends upon their voice.

To this presidency of the emperor or of his commissioners the acts of the councils and the Greek historians often refer. Even Pope Stephen V. (a.d. 817) writes, that Constantine the Great presided in the council of Nice. According to Eusebius, he introduced the principal matters of business with a solemn discourse, constantly attended the sessions, and took the place of honor in the assembly. His presence among the bishops at the banquet, which he gave them at the close of the council, seemed to that panegyrical historian a type of Christ among his saints!\footnote{Euseb., Vita Const. iii. 15: \textgreek{Χριστοῦ βασιλείας ἔδοξενἄντις φαντασιοῦσθαιεἰκόνα, ὄναρτεἷναι ἀλλ’ οὐχ ὕπαρ τὸ γινόμενον}.} This prominence of Constantine in the most celebrated and the most important of all the councils is the more remarkable, since at that time he had not yet even been baptized. When Marcian and Pulcheria appeared with their court at the council of Chalcedon, to confirm its decrees, they were greeted by the assembled bishops in the bombastic style of the East, as defenders of the faith, as pillars of orthodoxy, as enemies and persecutors of heretics; the emperor as a second Constantine, a new Paul, a new David; the empress as a second Helena; with other high-sounding predicates.\footnote{Mansi, vii. 170 sqq. The emperor is called there not simply divine, which would be idolatrous enough, but \emph{most divine}, \textgreek{ὁ θειότατος· καὶεὐσεβέστατος ἡμῶνδεσπότης}, divinissimus et piissimus noster imperator ad sanctam synodum dixit, etc. And these adulatory epithets occur repeatedly in the acts of this council.} The second and fifth general councils were the only ones at which the emperor was not represented, and in them the presidency was in the hands of the patriarchs of Constantinople.

But together with the imperial commissioners, or in their absence, the different patriarchs or their representatives, especially the legates of the Roman bishop, the most powerful of the patriarchs, took part in the presiding office. This was the case at the third and fourth, and the sixth, seventh, and eighth universal councils.

For the emperor’s connection with the council had reference rather to the conduct of business and to the external affairs of the synod, than to its theological and religious discussions. This distinction appears in the well-known dictum of Constantine respecting a double episcopate, which we have already noticed. And at the Nicene council the emperor acted accordingly. He paid the bishops greater reverence than his heathen predecessors had shown the Roman senators. He wished to be a servant, not a judge, of the successors of the apostles, who are constituted priests and gods on earth. After his opening address, he “resigned the word” to the (clerical) officers of the council,\footnote{Eusebius, Vita Const. iii. 13: \textgreek{Ὁμὲνδὴταῦτεἰπὼν  Ρωμαίᾳ γλώττῃ}[which was still the official language], \textgreek{ὑφερμηνεύοντος ἑτέρου, παρεδίδουτὸνλόγοντοῖς τῆς συνόδουπροέδροις.} Yet, according to the immediately following words of Eusebius, the emperor continued to take lively interest in the proceedings, hearing, speaking, and exhorting to harmony. Eusebius’whole account of this synod is brief and unsatisfactory.} by whom probably Alexander, bishop of Alexandria, Eustathius of Antioch, and Hosius of Cordova—the latter as special friend of the emperor, and as representative of the Western churches and perhaps of the bishop of Rome—are to be understood. The same distinction between a secular and spiritual presidency meets us in Theodosius II., who sent the comes Candidian as his deputy to the third general council, with full power over the entire business proceedings, but none over theological matters themselves; “for”—wrote he to the council-, “it is not proper that one who does not belong to the catalogue of most holy bishops, should meddle in ecclesiastical discussions.” Yet Cyril of Alexandria presided at this council, and conducted the business, at first alone, afterward in conjunction with the papal legates; while Candidian supported the Nestorian opposition, which held a council of its own under the patriarch John of Antioch.

Finally, from the emperors proceeded the ratification of the councils. Partly by their signatures, partly by special edicts, they gave the decrees of the council legal validity; they raised them to laws of the realm; they took pains to have them observed, and punished the disobedient with deposition and banishment. This was done by Constantine the Great for the decrees of Nice; by Theodosius the Great for those of Constantinople; by Marcian for those of Chalcedon. The second ecumenical council expressly prayed the emperor for such sanction, since he was present neither in person nor by commission. The papal confirmation, on the contrary, was not considered necessary, until after the fourth general council, in 451.\footnote{To wit, in a letter of the council to Leo (Ep. 89, in the Epistles of Leo, ed. Baller., tom. i. p. 1099), and in a letter of Marcian to Leo (Ep. 110, tom. i. p. 1182 sq.).} And notwithstanding this, Justinian broke through the decrees of the fifth council, of 553, without the consent, and in fact despite the intimated refusal of Pope Vigilius. In the middle ages, however, the case was reversed. The influence of the pope on the councils increased, and that of the emperor declined; or rather, the German emperor never claimed so preëminent a position in the church as the Byzantine. Yet the relation of the pope to a general council, the question which of the two is above the other, is still a point of controversy between the curialist or ultramontane and the episcopal or Gallican schools.

Apart from this predominance of the emperor and his commissioners, the character of the ecumenical councils was thoroughly hierarchical. In the apostolic council at Jerusalem, the elders and the brethren took part with the apostles, and the decision went forth in the name of the whole congregation.\footnote{Acts xv. 22: \textgreek{Τότεἔδοξετοῖς ἀποστόλοις καὶτοῖς πρεσβυτέροις σὺνὃλῃ τῇ ἐκκλησίᾳ}; and v. 23: \textgreek{Οἱἀπόστολοικαὶοἰπρεσβύτεροικαὶοἰἀδελφοὶτοῖς ... ἀδελφοῖς. κ.τ.λ.} Comp. my Hist. of the Apostolic Church, § 69, and § 128.} But this republican or democratic element, so to call it, had long since given way before the spirit of aristocracy. The bishops alone, as the successors and heirs of the apostles, the ecclesia docens, were members of the councils. Hence, in the fifth canon of Nice, even a provincial synod is termed “the general assembly of the bishops of the province.” The presbyters and deacons took part, indeed, in the deliberations, and Athanasius, though at the time only a deacon, exerted probably more influence on the council of Nice by his zeal and his gifts, than most of the bishops; but they had no votum decisivum, except when, like the Roman legates, they represented their bishops. The laity were entirely excluded.

Yet it must be remembered, that the bishops of that day were elected by the popular voice. So far as that went, they really represented the Christian people, and were not seldom called to account by the people for their acts, though they voted in their own name as successors of the apostles. Eusebius felt bound to justify, his vote at Nice before his diocese in Caesarea, and the Egyptian bishops at Chalcedon feared an uproar in their congregations.

Furthermore, the councils, in an age of absolute despotism, sanctioned the principle of common public deliberation, as the best means of arriving at truth and settling controversy. They revived the spectacle of the Roman senate in ecclesiastical form, and were the forerunners of representative government and parliamentary legislation.

In matters of discipline the majority decided; but in matters of faith unanimity was required, though, if necessary, it was forced by the excision of the dissentient minority. In the midst of the assembly an open copy of the Gospels lay upon a desk or table, as, a symbol of the presence of Christ, whose infallible word is the rule of all doctrine. Subsequently the ecclesiastical canons and the relics of the saints were laid in similar state. The bishops—at least according to later usage—sat in a circle, in the order of the dates of their ordination or the rank of their sees; behind them, the priests; before or beside them, the deacons. The meetings were opened and closed with religious solemnities in liturgical style. In the ancient councils the various subjects were discussed in open synod, and the Acts of the councils contain long discourses and debates. But in the council of Trent the subjects of action were wrought up in separate committees, and only laid before the whole synod for ratification. The vote was always taken by heads, till the council of Constance, when it was taken by nations, to avoid the preponderance of the Italian prelates.

The jurisdiction of the ecumenical councils covered the entire legislation of the church, all matters of Christian faith and practice (fidei et morum), and all matters of organization arid worship. The doctrinal decrees were called dogmata or symbola; the disciplinary, canones. At the same time, the councils exercised, when occasion required, the highest judicial authority, in excommunicating bishops and patriarchs.

The authority of these councils in the decision of all points of controversy was supreme and final.

Their doctrinal decisions were early invested with infallibility; the promises of the Lord respecting the indestructibleness of his church, his own perpetual presence with the ministry, and the guidance of the Spirit of truth, being applied in the full sense to those councils, as representing the whole church. After the example of the apostolic council, the usual formula for a decree was: Visum est Sprirtui Sancto et nobis.\footnote{\textgreek{Ἒδοξετῷπνεύματιἁγίῳ καὶἡμῖν}, Acts xv. 28. The provincial councils, too, had already used this phrase; e.g. the Concil. Carthaginiense, of 252 (in the Opera Cypriani): “Placuit nobis, \emph{Sancto Spiritu suggerente}, et Domino per visiones multas et manifestas admonente.” So the council of Arles, in 314: “Placuit ergo, \emph{presente Spiritu Sancto} et angelis ejus.”} Constantine the Great, in a circular letter to the churches, styles the decrees of the Nicene council a divine command;\footnote{\textgreek{Θείανἐντολήν,} and \textgreek{θείανβούλησιν}, in Euseb., Vita Const. iii. 20. Comp. his Ep. ad Eccl. Alexandr., in Socrates, H. E. i. 9 where he uses similar expressions.} a phrase, however, in reference to which the abuse of the word divine, in the language of the Byzantine despots, must not be forgotten. Athanasius says, with reference to the doctrine of the divinity of Christ: “What God has spoken by the council of Nice, abides forever.”\footnote{Isidore of Pelusium also styles the Nicene council divinely inspired, \textgreek{θεόθενἐμπνευςθεῖσα} (Ep. 1. iv. Ep. 99). So Basil the Great, Ep. 114 (in the Benedictine edition of his Opera omnia, tom. iii. p. 207), where he says that the 318 fathers of Nice have not spoken without the \textgreek{ἐνέργειατοῦἁγίουπνεύματος}(non sine Spiritus Sancti afflatu).} The council of Chalcedon pronounced the decrees of the Nicene fathers unalterable statutes, since God himself had spoken through them.\footnote{Act. i., in Mansi, vi. p. 672. We quote from the Latin translation: “Nullo autem modo patimur a quibusdam concuti definitam fidem, sive fidei symbolum, a sanctis patribus nostris qui apud Nicaeam convenerunt illis temporibus: nec permittimus aut nobis, aut aliis, mutare aliquod verbum ex his quae ibidem continentur, aut unam syllabam praeterire, memores dicentis: \emph{Ne transferas terminos aeternos, quos posuerunt patres tui} (Prov. xxii. 8; Matt. x. 20). Non enim erant ipsi loquentes, sed ipse Spiritus Dei et Patris qui procedit ex ipso.”} The council of Ephesus, in the sentence of deposition against Nestorius, uses the formula: “The Lord Jesus Christ, whom he has blasphemed, determines through this most holy council.”\footnote{\textgreek{Ὁβλασφημηθεὶς παῤαὐτοῦκύριος Ἰης. Χριστὸς ωὝρισεδιὰτῆς παρούσης ἁγιωτάτης συνόδου}.} Pope Leo speaks of an “irretractabilis consensus” of the council of Chalcedon upon the doctrine of the person of Christ. Pope Gregory the Great even placed the first four councils, which refuted and destroyed respectively the heresies and impieties of Arius, Macedonius, Nestorius, and Eutyches, on a level with the four canonical Gospels.\footnote{Lib. i. Ep. 25 (ad Joannem episcopum Constant., et caeteros patriarchas, in Migne’s edition of Gr. Opera, tom. iii. p. 478, or in the Bened. ed. iii. 515): “Praeterea, quia corde creditur ad justitiam, ore autem confessio fit ad salutem, sicut sancti evangelii quatuor libros, sic quatuor concilia suscipere et venerari me fateor. Nicaenum scilicet in quo perversum Arii dogma destruitur; Constantinopolitanum quoque, in quo Eunomii et Macedonii error convincitur; Ephesinum etiam primum, in quo Nestorii impietas judicatur; Chalcedonense vero, in quo Eutychetii [Eutychis] Dioscorique pravitas reprobatur, tota devotione complector, integerrima approbatione custodio: quia in his velut in quadrato lapide, sanctae fidei structura consurgit, et cujuslibet vitae atque actionis existat, quisquis eorum soliditatem non tenet, etiam si lapis esse cernitur, tamen extra aedificium jacet. Quintum quoque concilium pariter veneror, in quo et epistola, quae Ibae dicitur, erroris plena, reprobatur,” etc.} In like manner Justinian puts the dogmas of the first four councils on the same footing with the Holy Scriptures, and their canons by the side of laws of the realm.\footnote{Justin. Novell. cxxxi.“Quatuor synodorum dogmata sicut sanctas scripturas accipimus, et regulas sicut leges observamus.”} The remaining three general councils have neither a theological importance, nor therefore an authority, equal to that of those first four, which laid the foundations of ecumenical orthodoxy. Otherwise Gregory would have mentioned also the fifth council, of 553, in the passage to which we have just referred. And even among the first four there is a difference of rank; the councils of Nice and Chalcedon standing highest in the character of their results.

Not so with the rules of discipline prescribed in the canones. These were never considered universally binding, like the symbols of faith; since matters of organization and usage, pertaining rather to the external form of the church, are more or less subject to the vicissitude of time. The fifteenth canon of the council of Nice, which prohibited and declared invalid the transfer of the clergy from one place to another,\footnote{Conc. Nic. can. 15: \textgreek{ὛΩστεἀπὸπόλεως εἰς πόλινμὴμεταβαίνεινμήτεἐπίσκοπονμήτεπρεσβύτερονμήτεδιάκονον}. This prohibition arose from the theory of the relation between a clergyman and his congregation, as a mystical marriage, and was designed to restrain clerical ambition. It appears in the Can. Apost. 13, 14, but was often violated. At the Nicene council itself there were several bishops, like Eusebius of Nicomedia, and Eustathius of Antioch, who had exchanged their first bishopric for another and a better.} Gregory Nazianzen, fifty-seven years later (382), reckons among statutes long dead.\footnote{\textgreek{Νόμους πάλαιτεθνηκότας}, Carm. de vita sua, v. 1810.} Gregory himself repeatedly changed his location, and Chrysostom was called from Antioch to Constantinople. Leo I. spoke with strong disrespect of the third canon of the second ecumenical council, for assigning to the bishop of Constantinople the first rank after the bishop of Rome; and for the same reason be protested against the twenty-eighth canon of the fourth ecumenical council.\footnote{Epist. 106 (al. 80) ad Anatolium, and Epist. 105 ad Pulcheriam. Comp. above, § 57. Even Gregory I., so late as 600, writes in reference to the \emph{canones} of the Constantinopolitan council of 381: “Romana autem ecclesia eosdem canones vel gesta Synodi illius hactenus non habet, nec accepit; in hoc autem eam accepit, quod est per eam contra Macedonium definitum.” Lib. vii. Ep. 34, ad Eulogium episcopum Alexandr. (tom. iii. p. 882, ed. Bened., and in Migne’s ed., iii. 893.)} Indeed the Roman church has made no point of adopting all the disciplinary laws enacted by those synods.

Augustine, the ablest and the most devout of the fathers, conceived, in the best vein of his age, a philosophical view of this authority of the councils, which strikes a wise and wholesome mean between the extremes of veneration and disparagement, and approaches the free spirit of evangelical Protestantism. He justly subordinates these councils to the Holy Scriptures, which are the highest and the perfect rule of faith, and supposes that the decrees of a council may be, not indeed set aside and repealed, yet enlarged and completed by, the deeper research of a later day. They embody, for the general need, the results already duly prepared by preceding theological controversies, and give the consciousness of the church, on the subject in question, the clearest and most precise expression possible at the time. But this consciousness itself is subject to development. While the Holy Scriptures present the truth unequivocally and infallibly, and allow no room for doubt, the judgment of bishops may be corrected and enriched with new truths from the word of God, by the wiser judgment of other bishops; the judgment of the provincial council by that of a general; and the views of one general council by those of a later.\footnote{De Baptismo contra Donatistas, I. ii. 3 (in the Benedictine edition of August. Opera, tom. ix. p. 98): “Quis autem nesciat, sanctam Scripturam canonicam, tam Veteris quam Novi Testamenti, certis suis terminis contineri, eamque omnibus posterioribus Episcoporum literis ita praeponi, ut de illa omnino dubitari et disceptari non possit, utrum verum vel utrum rectum sit, quidquid in ea scriptum esse constiterit; Episcoporum autem literas quae post confirmatum canonem vel scripta sunt vel scribuntur, et per sermonem forte sapientiorem cujuslibet in ea re peritioris, et per aliorum Episcoporum graviorem auctoritatem doctioremque prudentiam, et per concilia licere \emph{reprehendi}, si quid in eis forte \emph{a veritate deviatum est}; et \emph{ipsa concilia}, quae per singulas regiones vel provincias fiunt, \emph{plenariorum} \emph{conciliorum auctoritate}, quae fiunt ex universo orbe Christiano, sine ullis ambagibus \emph{cedere}; \emph{ipsaque plenaria saepe priora posterioribus emendari}, quum aliquo experimento rerum aperitur quod clausum erat et cognoscitur quod latebat; sine ullo typho sacrilegae superbiae, sine ulla inflata cervice arrogantiae, sine ulla contentione lividae invidiae, cum sancta humilitate, cum pace catholica, cum caritate christiana.” Comp. the passage Contra Maximinum Arianum, ii. cap. 14, § 3 (in the Bened. ed., tom. viii. p. 704), where he will have even the decision of the Nicene council concerning the \emph{homousion} measured by the higher standard of the Scriptures.} In this Augustine presumed, that all the transactions of a council were conducted in the spirit of Christian humility, harmony, and love; but had he attended the council of Ephesus, in 431, to which he was summoned about the time of his death, he would, to his grief, have found the very opposite spirit reigning there. Augustine, therefore, manifestly acknowledges a gradual advancement of the church doctrine, which reaches its corresponding expression from time to time through the general councils; but a progress within the truth, without positive error. For in a certain sense, as against heretics, he made the authority of Holy Scripture dependent on the authority of the catholic church, in his famous dictum against the Manichaean heretics: “I would not believe the gospel, did not the authority of the catholic church compel me.”\footnote{Contra Epistolam Manichaei, lib. i. c. 5 (in the Bened. ed., tom. viii. p. 154): “Ego vero evangelio non crederem, nisi me ecclesiae catholicae commoveret auctoritas.”} In like manner Vincentius Lerinensis teaches, that the church doctrine passes indeed through various stages of growth in knowledge, and becomes more and more clearly defined in opposition to ever-rising errors, but can never become altered or dismembered.\footnote{Commonitorium, c. 23 (in Migne’s Curs. Patrol. tom. 50, p. 667): “Sed forsitan dicit aliquis: Nullusne ergo in ecclesia Christi profectus habebitur religionis? Habeatur plane et maximus .... Sed ita tamen ut vere profectus sit ille fidei, non permutatio. Si quidem ad profectum pertinet ut in semetipsum unaquaeque res amplificetur; ad permutationem vero, ut aliquid ex alio in aliud transvertatur. Crescat igitur oportet et multum vehementerque proficiat tam singulorum quam omnium, tam unius hominis, quam totius ecclesiae, aetatum ac seculorum gradibus, intelligentia, scientia, sapientia, sed in suo dutaxat genere, in eodem scilicet dogmate, eodem sensu, eademque sententia.”}

The Protestant church makes the authority of the general councils, and of all ecclesiastical tradition, depend on the degree of its conformity to the Holy Scriptures; while the Greek and Roman churches make Scripture and tradition coordinate. The Protestant church justly holds the first four general councils in high, though not servile, veneration, and has received their statements of doctrine into her confessions of faith, because she perceives in them, though compassed with human imperfection, the clearest and most suitable expression of the teaching of the Scriptures respecting the Trinity and the divine-human person of Christ. Beyond these statements the judgment of the church (which must be carefully distinguished from theological speculation) has not to this day materially advanced;—the highest tribute to the wisdom and importance of those councils. But this is not saying that the Nicene and the later Athanasian creeds are the non plus ultra of all the church’s knowledge of the articles therein defined. Rather is it the duty of theology and of the church, while prizing and holding fast those earlier attainments, to study the same problems ever anew, to penetrate further and further these sacred fundamental mysteries of Christianity, and to bring to light new treasures from the inexhaustible mines of the Word of God, under the guidance of the same Holy Spirit, who lives and works in the church at this day as mightily as he did in the fifth century and the fourth. Christology, for example, by the development of the doctrine of the two states of Christ in the Lutheran church, and of the three offices of Christ in the Reformed, has been substantially enriched; the old Catholic doctrine, which was fixed with unerring tact at the council of Chalcedon, being directly concerned only with the two natures of Christ, as against the dualism of Nestorius and the monophysitism of Eutyches.

With this provision for further and deeper soundings of Scripture truth, Protestantism feels itself one with the ancient Greek and Latin church in the bond of ecumenical orthodoxy. But toward the disciplinary canons of the ecumenical councils its position is still more free and independent than that of the Roman church. Those canons are based upon an essentially unprotestant, that is, hierarchical and sacrificial conception of church order and worship, which the Lutheran and Anglican reformation in part, and the Zwinglian and Calvinistic almost entirely renounced. Yet this is not to say that much may not still be learned, in the sphere of discipline, from those councils, and that perhaps many an ancient custom or institution is not worthy to be revived in the spirit of evangelical freedom.

The moral character of those councils was substantially parallel with that of earlier and later ecclesiastical assemblies, and cannot therefore be made a criterion of their historical importance and their dogmatic authority. They faithfully reflect both the light and the shade of the ancient church. They bear the heavenly treasure in earthen vessels. If even among the inspired apostles at the council of Jerusalem there was much debate,\footnote{Acts xv. 6: \textgreek{Πολλῆς συζητήσεως γενομένης}; which Luther indeed renders quite too strongly: ” After they had wrangled long.” The English versions from Tyndale to King James translate: ” much disputing.”} and soon after, among Peter, Paul, and Barnabas, a violent, though only temporary collision, we must of course expect much worse of the bishops of the Nicene and the succeeding age, and of a church already interwoven with a morally degenerate state. Together with abundant talents, attainments, and virtues, there were gathered also at the councils ignorance, intrigues, and partisan passions, which had already been excited on all sides by long controversies preceding and now met and arrayed themselves, as hostile armies, for open combat. For those great councils, all occasioned by controversies on the most important and the most difficult problems of theology, are, in fact, to the history of doctrine, what decisive battles are to the history of war. Just because religion is the deepest and holiest interest of man, are religious passions wont to be the most violent and bitter; especially in a time when all classes, from imperial court to market stall, take the liveliest interest in theological speculation, and are drawn into the common vortex of excitement. Hence the notorious rabies theologorum was more active in the fourth and fifth centuries than it has been in any other period of history, excepting, perhaps, in the great revolution of the sixteenth century, and the confessional polemics of the seventeenth.

We have on this point the testimony of contemporaries and of the acts of the councils themselves. St. Gregory Nazianzen, who, in the judgment of Socrates, was the most devout and eloquent man of his age,\footnote{Hist. Eccl. lib. v. cap. 7.} and who himself, as bishop of Constantinople, presided for a time over the second ecumenical council, had so bitter an observation and experience as even to lose, though without sufficient reason, all confidence in councils, and to call them in his poems “assemblies of cranes and geese.” “To tell the truth” thus in 382 (a year after the second ecumenical council, and doubtless including that assembly in his allusion) he answered Procopius, who in the name of the emperor summoned him in vain to a synod—“to tell the truth, I am inclined to shun every collection of bishops, because I have never yet seen that a synod came to a good end, or abated evils instead of increasing them. For in those assemblies (and I do not think I express myself too strongly here) indescribable contentiousness and ambition prevail, and it is easier for one to incur the reproach of wishing to set himself up as judge of the wickedness of others, than to attain any success in putting the wickedness away. Therefore I have withdrawn myself, and have found rest to my soul only in solitude.”\footnote{Ep. ad Procop. 55, old order (al. 130). Similar representations occur in Ep. 76, 84; Carm. de vita sua, v. 1680-1688; Carm. x. v. 92; Carm. Adv. Episc. v. 154. Comp. Ullmann, Gregor. von Naz., p. 246 sqq., and p. 270. It is remarkable that Gibbon makes no use of these passages to support his summary judgment of the general councils at the end of his twentieth chapter, where he says: “The progress of time and superstition erased the memory of the weakness, the passion, the ignorance, which disgraced these ecclesiastical synods; and the Catholic world has unanimously submitted to the \emph{infallible} decrees of the general councils.”} It is true, the contemplative Gregory had an aversion to all public life, and in such views yielded unduly to his personal inclinations. And in any case he is inconsistent; for he elsewhere speaks with great respect of the council of Nice, and was, next to Athanasius, the leading advocate of the Nicene creed. Yet there remains enough in his many unfavorable pictures of the bishops and synods of his time, to dispel all illusions of their immaculate purity. Beausobre correctly observes, that either Gregory the Great must be a slanderer, or the bishops of his day were very remiss. In the fifth century it was no better, but rather worse. At the third general council, at Ephesus, 431, all accounts agree that shameful intrigue, uncharitable lust of condemnation, and coarse violence of conduct were almost as prevalent as in the notorious robber-council of Ephesus in 449; though with the important difference, that the former synod was contending for truth, the latter for error. Even at Chalcedon, the introduction of the renowned expositor and historian Theodoret provoked a scene, which almost involuntarily reminds us of the modern brawls of Greek and Roman monks at the holy sepulchre under the restraining supervision of the Turkish police. His Egyptian opponents shouted with all their might: “The faith is gone! Away with him, this teacher of Nestorius!” His friends replied with equal violence: “They forced us [at the robber-council] by blows to subscribe; away with the Manichaeans, the enemies of Flavian, the enemies of the faith! Away with the murderer Dioscurus? Who does not know his wicked deeds? The Egyptian bishops cried again: Away with the Jew, the adversary of God, and call him not bishop!” To which the oriental bishops answered: “Away with the rioters, away with the murderers! The orthodox man belongs to the council!” At last the imperial commissioners interfered, and put an end to what they justly called an unworthy and useless uproar.\footnote{\textgreek{Ἐκβοήσεις δημοτικαί.} See Harduin, tom. ii. p. 71 sqq., and Mansi, tom. vi. p. 590 sq. Comp. also Hefele, ii. p. 406 sq.}

In all these outbreaks of human passion, however, we must not forget that the Lord was sitting in the ship of the church, directing her safely through the billows and storms. The Spirit of truth, who was not to depart from her, always triumphed over error at last, and even glorified himself through the weaknesses of his instruments. Upon this unmistakable guidance from above, only set out by the contrast of human imperfections, our reverence for the councils must be based. Soli Deo gloria; or, in the language of Chrysostom: Δόξα τῷ θεῷ πάντων ἕνεκεν!

\xsection{List of the Ecumenical Councils of the Ancient Church}

§ 66. List of the Ecumenical Councils of the Ancient Church,

We only add, by way of a general view, a list of all the ecumenical councils of the Graeco-Roman church, with a brief account of their character and work.

1. The Concilium Nicaenum I., a.d. 325; held at Nicaea in Bithynia, a lively commercial town near the imperial residence of Nicomedia, and easily accessible by land and sea. It consisted of three hundred and eighteen bishops,\footnote{This is the usual estimate, resting on the authority of Athanasius, Basil (Ep. 114; Opera, t. iii. p 207, ed. Bened.), Socrates, Sozomen, and Theodoret; whence the council is sometimes called the Assembly of the Three Hundred and Eighteen. Other data reduce the number to three hundred, or to two hundred and seventy, or two hundred and fifty, or two hundred and eighteen; while later tradition swells it to two thousand or more.} besides a large number of priests, deacons, and acolytes, mostly from the East, and was called by Constantine the Great, for the settlement of the Arian controversy. Having become, by decisive victories in 323, master of the whole Roman empire, he desired to complete the restoration of unity and peace with the help of the dignitaries of the church. The result of this council was the establishment (by anticipation) of the doctrine of the true divinity of Christ, the identity of essence between the Son and the Father. The fundamental importance of this dogma, the number, learning, piety and wisdom of the bishops, many of whom still bore the marks of the Diocletian persecution, the personal presence of the first Christian emperor, of Eusebius, “the father of church history,” and of Athanasius, “the father of orthodoxy” (though at that time only archdeacon), as well as the remarkable character of this epoch, combined in giving to this first general synod a peculiar weight and authority. It is styled emphatically “the great and holy council,” holds the highest place among all the councils, especially with the Greeks,\footnote{For some time the Egyptian and Syrian churches commemorated the council of Nicaea by an annual festival.} and still lives in the Nicene Creed, which is second in authority only to the ever venerable Apostles’ Creed. This symbol was, however, not finally settled and completed in its present form (excepting the still later Latin insertion of filioque), until the second general council. Besides this the fathers assembled at Nicaea issued a number of canons, usually reckoned twenty on various questions of discipline; the most important being those on the rights of metropolitans, the time of Easter, and the validity of heretical baptism.

2. The Concilium Constantinopolitanum I., a.d. 381 summoned by Theodosius the Great, and held at the imperial city, which had not even name in history till five years after the former council. This council, however, was exclusively oriental, and comprised only a hundred and fifty bishops, as the emperor had summoned none but the adherents of the Nicene party, which had become very much reduced under the previous reign. The emperor did not attend it. Meletius of Antioch was president till his death; then Gregory Nazianzen; and, after his resignation, the newly elected patriarch Nectarius of Constantinople. The council enlarged the Nicene confession by an article on the divinity and personality of the Holy Ghost, in opposition to the Macedonians or Pneumatomachists (hence the title Symbolum Nicaeno-Constantinopolitanum), and issued seven more canons, of which the Latin versions, however, give only the first four, leaving the genuineness of the other three, as many think, in doubt.

3. The Concilium Ephesinum, a.d. 431; called by Theodosius II., in connection with the Western co-emperor Valentinian III., and held under the direction of the ambitious and violent Cyril of Alexandria. This council consisted of, at first, a hundred and sixty bishops, afterward a hundred and ninety-eight,\footnote{The opposition council, which John of Antioch, on his subsequent arrival, held in the same city in the cause of Nestorius and under the protection of the imperial commissioner Candidian, numbered forty-three members, and excommunicated Cyril, as Cyril had excommunicated Nestorius.} including, for the first time, papal delegates from Rome, who were instructed not to mix in the debates, but to sit as judges over the opinions of the rest. It condemned the error of Nestorius on the relation of the two natures in Christ, without, stating clearly the correct doctrine. It produced, therefore, but a negative result, and is the least important of the first four councils, as it stands lowest also in moral character. It is entirely rejected by the Nestorian or Chaldaic Christians. Its six canons relate exclusively to Nestorian and Pelagian affairs, and are wholly omitted by Dionysius Exiguus in his collection.

4. The Concilium Chalcedonense, a.d. 451; summoned by the emperor Marcian, at the instance of the Roman bishop Leo; held at Chalcedon in Bithynia, opposite Constantinople; and composed of five hundred and twenty (some say six hundred and thirty) bishops.\footnote{The synod itself, in a letter to Leo, states the number as only five hundred and twenty; Leo, on the contrary (Ep. 102), speaks of about six hundred members; and the usual opinion (Tillemont, Memoires, t. xv. p. 641) raises the whole number of members, including deputies, to six hundred and thirty.} Among these were three delegates of the bishop of Rome, two bishops of Africa, and the rest all Greeks and Orientals. The fourth general council fixed the orthodox doctrine of the person of Christ in opposition to Eutychianism and Nestorianism, and enacted thirty canons (according to some manuscripts only twenty-seven or twenty-eight), of which the twenty-eighth was resisted by the Roman legates and Leo I. This was the most numerous, and next to the Nicene, the most important of all the general councils, but is repudiated by all the Monophysite sects of the Eastern church.

5. The Concilium Constantinopolitanum II. was assembled a full century later, by the emperor Tustinian, a.d. 553, without consent of the pope, for the adjustment of the tedious Monophysite controversy. It was presided over by the patriarch Eutychius of Constantinople, consisted of only one hundred and sixty-four bishops, and issued fourteen anathemas against the three chapters,\footnote{Tria capitula, \textgreek{Κεφάλεια}.} so called, or the christological views of three departed bishops and divines, Theodore of Mopsueste, Theodoret of Cyros, and Ibas of Edessa, who were charged with leaning toward the Nestorian heresy. The fifth council was not recognized, however, by many Western bishops, even after the vacillating Pope Vigilius gave in his assent to it, and it induced a temporary schism between Upper Italy and the Roman see. As to importance, it stands far below the four previous councils. Its Acts, in Greek, with the exception of the fourteen anathemas, are lost.

Besides these, there are two later councils, which have attained among the Greeks and Latins an undisputed ecumenical authority: the Third Council of Constantinople, under Constantine Progonatus, a.d. 680, which condemned Monothelitism (and Pope Honorius, † 638),\footnote{The condemnation of a departed pope as a heretic by an ecumenical council is so inconsistent with the claims of papal infallibility, that Romish historians have tried their utmost to dispute the fact, or to weaken its force by sophistical pleading.} and consummated the old Catholic christology; and the Second Council of Nicaea, under the empress Irene, a.d. 787, which sanctioned the image-worship of the Catholic church, but has no dogmatical importance.

Thus Nicaea—now the miserable Turkish hamlet Is-nik\footnote{\textgreek{Εἰς Νίκαιαν}. \emph{Nice} and \emph{Nicene} are properly misnomers, but sanctioned by the use of Gibbon and other great English writers.}—has the honor of both opening and closing the succession of acknowledged ecumenical councils.

From this time forth the Greeks and Latins part, and ecumenical councils are no longer to be named. The Greeks considered the second Trullan\footnote{\emph{Trullum} was a saloon with a cupola in the imperial palace of Constantinople.} (or the fourth Constantinopolitan) council of 692, which enacted no symbol of faith, but canons only, not an independent eighth council, but an appendix to the fifth and sixth ecumenical councils (hence, called the Quinisexta sc. synodus); against which view the Latin church has always protested. The Latin church, on the other hand, elevates the fourth council of Constantinople, a.d. 869,\footnote{The Latins call it the fourth because they reject the fourth Constantinopolitan (the second Trullan) council of 692, because of its canons, and the fifth of 754 because it condemned the worship of images, which was subsequently sanctioned by the second council of Nicaea in 787.} which deposed the patriarch Photius, the champion of the Greek church in her contest with the Latin, to the dignity of an eighth ecumenical council; but this council was annulled for the Greek church by the subsequent restoration of Photius. The Roman church also, in pursuance of her claims to exclusive catholicity, adds to the seven or eight Greek councils twelve or more Latin general councils, down to the Vatican (1870); but to all these the Greek and Protestant churches can concede only a sectional character. Three hundred and thirty-six years elapsed between the last undisputed Graeco-Latin ecumenical council of the ancient church (a.d. 787), and the first Latin ecumenical council of the mediaeval church (1123). The authority of the papal see had to be established in the intervening centuries.\footnote{On the number of the ecumenical councils till that of Trent the Roman divines themselves are not agreed. The Gallicans reckon twenty-one, Bellarmine eighteen, Hefele only sixteen. The undisputed ones, besides the eight already mentioned Graeco-Latin councils, are these eight Latin: the first Lateran (Roman) council, a.d.1123; the second Lateran, a.d.1139; the third Lateran, a.d.1179; the fourth Lateran, a.d.1215); the first of Lyons, a.d.1245; the second of Lyons, a.d.1274; that of Florence, a.d.1439; (the fifth Lateran, 1512-1517, is disputed;) and that of Trent, a.d.1545-1563. The ecumenical character of the three reformatory councils of Pisa, Constance, and Basle, in the beginning of the fifteenth century, and of the fifth Lateran council, a.d.1512-1517, is questioned among the Roman divines, and is differently viewed upon ultramontane and upon Gallican principles. Hefele considers them \emph{partially} ecumenical; that is, so far as they were ratified by the pope. [But in the Revised edition of his \emph{Conciliengeschichte}, 1873 sqq., he reckons twenty ecumenical councils, including the Vatican, 1870. See Appendix, p. 1032.]}

\xsection{Books of Ecclesiastical Law}

§ 67. Books of Ecclesiastical Law.

I. Bibiliotheca juris canonici veteris, ed. Voellus (theologian of the Sorbonne) and Justellus (Justeau, counsellor and secretary to the French king), Par. 1661, 2 vols. fol. (Vol. i. contains the canons of the universal church, Greek and Latin, the ecclesiastical canons of Dionysius Exiguus, or of the old Roman church, the canons of the African church, etc. See a list of contents in Darling’s Cyclop. Bibliographica, p. 1702 sq.)

II. See the literature in vol. ii. § 56 (p. 183). The brothers Ballerini: De antiquis tum editis tum ineditis collectionibus et collectoribus canonum ad Gratianum usque in ed. Opp. Leon M. Ven., 1753 sqq. The treatises of Quesnel, Marca, Constant, Drey, Theiner, etc., on the history of the collections of canons. Comp. Ferd. Walther: Lehrbuch des Kirchenrechts, p. 109 sqq., 8th ed., 1839.

The universal councils, through their disciplinary enactments or canons, were the main fountain of ecclesiastical law. To their canons were added the decrees of the most important provincial councils of the fourth century, at Ancyra (314), Neo-Caesarea (314), Antioch (341), Sardica (343), Gangra (365), and Laodicea (between 343 and 381); and in a third series, the orders of eminent bishops, popes, and emperors. From these sources arose, after the beginning of the fifth century, or at all events before the council of Chalcedon, various collections of the church laws in the East, in North Africa, in Italy, Gaul, and Spain; which, however, had only provincial authority, and in many respects did not agree among themselves. A codex canonum ecclesiae universae did not exist. The earlier collections because eclipsed by two, which, the one in the West, the other in the East, attained the highest consideration.

The most important Latin collection comes from the Roman, though by descent Scythian, abbot Dionysius Exiguus,\footnote{It is uncertain whether he obtained the surname Exiguus from his small stature or his monastic humility.} who also, notwithstanding the chronological error at the base of his reckoning, immortalized himself by the introduction of the Christian calendar, the “Dionysian Era.” It was a great thought of this “little” monk to view Christ as the turning point of ages, and to introduce this view into chronology. About the year 500 Dionysius translated for the bishop Stephen of Salona a collection of canons from Greek into Latin, which is still extant, with its prefatory address to Stephen.\footnote{It may be found in the above-cited Bibliotheca, vol. i., and in all good collections of councils. He says in the preface that, confusione priscae translationis (the Prisca or Itala) offensus, he has undertaken a new translation of the Greek canons.} It contains, first, the, fifty so-called Apostolic Canons, which pretend to have been collected by Clement of Rome, but in truth were a gradual production of the third and fourth centuries;\footnote{“Canones, qui dicuntur apostolorum, ... quibus plurimi consensum non praebuere facilem;” implying that Dionysius himself, with many others, doubted their apostolic origin. In a later collection of canons by Dionysius, of which only the preface remains, he entirely omitted the apostolic canons, with the remark: “Quos non admisit universitas, ego quoque in hoc opere praetermisi.” On the pseudo-apostolic Canons and Constitutions, comp. vol. i. § 113 (p. 440-442), and the well-known critical work of the Roman Catholic theologian Drey.} then the canons of the most important councils of the fourth and fifth centuries, including those of Sardica and Africa; and lastly, the papal decretal letters from Siricius (385) to Anastasius II. (498). The Codex Dionysii was gradually enlarged by additions, genuine and spurious, and through the favor of the popes, attained the authority of law almost throughout the West. Yet there were other collections also in use, particularly in Spain and North Africa.

Some fifty years after Dionysius, John Scholasticus, previously an advocate, then presbyter at Antioch, and after 564 patriarch of Constantinople, published a collection of canons in Greek,\footnote{\textgreek{Σύνταγμακανόνων}, Concordia canonum, in the Bibliotheca of Justellus, tom. ii.} which surpassed the former in completeness and convenience of arrangement, and for this reason, as well as the eminence of the author, soon rose to universal authority in the Greek church. In it he gives eighty-five Apostolic Canons, and the ordinances of the councils of Ancyra (314) and Nicaea (325), down to that of Chalcedon (451), in fifty titles, according to the order of subjects. The second Trullan council (Quinisextum, of 692), which passes with the Greeks for ecumenical, adopted the eighty-five Apostolic Canons, while it rejected the Apostolic Constitutions, because, though, like the canons, of apostolic origin, they had been early adulterated. Thus arose the difference between the Greek and Latin churches in reference to the number of the so-called Apostolic canons; the Latin church retaining only the fifty of the Dionysian collection.

The same John, while patriarch of Constantinople, compiled from the Novelles of Justinian a collection of the ecclesiastical state-laws or νόμοι, as they were called in distinction from the synodal church-laws or κανόνες. Practical wants then led to a union of the two, under the title of Nomocanon.

These books of ecclesiastical law served to complete and confirm the hierarchical organization, to regulate the life of the clergy, and to promote order and discipline; but they tended also to fix upon the church an outward legalism, and to embarrass the spirit of progress.

\xchapter{Church Discipline and Schisms}

CHURCH DISCIPLINE AND SCHISMS.

\xsection{Decline of Discipline}

§ 68. Decline of Discipline.

The principal sources are the books of ecclesiastical law and the acts of councils. Comp. the literature at § 67, and at vol. i. § 114.

The union of the church with the state shed, in general, an injurious influence upon the discipline of the church; and that, in two opposite directions.

On the one hand it increased the stringency of discipline and led to a penal code for spiritual offences. The state gave her help to the church, lent the power of law to acts of suspension and excommunication, and accompanied those acts with civil penalties. Hence the innumerable depositions and banishments of bishops during the theological controversies of the Nicene and the following age, especially under the influence of the Byzantine despotism and the religious intolerance and bigotry of the times. Even the penalty of death was decreed, at least against the Priscillianists, though under the protest of nobler divines, who clave to the spiritual character of the church and of her weapons.\footnote{Comp. § 27, above.} Heresy was regarded as the most grievous and unpardonable crime against society, and was treated accordingly by the ruling party, without respect of creed.

But on the other hand discipline became weakened. With the increasing stringency against heretics, firmness against practical errors diminished. Hatred of heresy and laxity of morals, zeal for purity of doctrine and indifference to purity of life, which ought to exclude each other, do really often stand in union. Think of the history of Pharisaism at the time of Christ, of orthodox Lutheranism in its opposition to Spener and the Pietistic movement, and of prelatical Anglicanism in its conflict with Methodism and the evangelical party. Even in the Johannean age this was the case in the church of Ephesus, which prefigured in this respect both the light and shade of the later Eastern church.\footnote{Rev. ii. 1-7. Comp. my Hist. of the Apostolic Church, p. 429.} The earnest, but stiff, mechanical penitential discipline, with its four grades of penance, which had developed itself during the Dioclesian persecution,\footnote{Comp. vol. i. § 114 (p. 444 sq.).} continued in force, it is true, as to the letter, and was repeatedly reaffirmed by the councils of the fourth century. But the great change of circumstances rendered the practical execution of it more and more difficult, by the very multiplication and high position, of those on whom it ought to be enforced. In that mighty revolution under Constantine the church lost her virginity, and allied herself with the mass of heathendom, which had not yet experienced an inward change. Not seldom did the emperors themselves, and other persons of authority, who ought to have led the way with a good example, render themselves, with all their zeal for theoretical orthodoxy, most worthy of suspension and excommunication by their scandalous conduct, while they were surrounded by weak or worldly bishops, who cared more for the favor of their earthly masters, than for the honor of their heavenly Lord and the dignity of the church. Even Eusebius, otherwise one of the better bishops of his time, had no word of rebuke for the gross crimes of Constantine, but only the most extravagant eulogies for his merits.

In the Greek church the discipline gradually decayed, to the great disadvantage of public morality, and every one was allowed to partake of the communion according to his conscience. The bishops alone reserved the right of debarring the vicious from the table of the Lord. The patriarch Nectarius of Constantinople, about 390, abolished the office of penitential priest (presbyter poenitentiarius), who was set over the execution of the penitential discipline. The occasion of this act was furnished by a scandalous occurrence: the violation of a lady of rank in the church by a worthless deacon, when she came to submit herself to public penance. The example of Nectarius was soon followed by the other oriental bishops.\footnote{Sozomen, vii. 16; Socrates, v. 19. This fact has been employed by the Roman church against the Protestant, in the controversy on the sacrament of penance. Nectarius certainly did abolish the institution of penitential priest, and the \emph{public} church penance. But for or against private penances no inference can be drawn from the statement of these historians.}

Socrates and Sozomen, who inclined to the severity of the Novatians, date the decline of discipline and of the former purity of morals from this act. But the real cause lay further back, in the connection of the church with the temporal power. Had the state been pervaded with the religious earnestness and zeal of Christianity, like the Genevan republic, for example, under the reformation of Calvin, the discipline of the church would have rather gained than lost by the alliance. But the vast Roman state could not so easily and quickly lay aside its heathen traditions and customs; it perpetuated them under Christian names. The great mass of the people received, at best, only John’s baptism of repentance, not Christ’s baptism of the Holy Ghost and of fire.

Yet even under these new conditions the original moral earnestness of the church continued, from time to time, to make itself known. Bishops were not wanting to confront even the emperors, as Nathan stood before David after his fall, in fearless rebuke. Chrysostom rigidly insisted, that the deacon should exclude all unworthy persons from the holy communion, though by his vehement reproof of the immoralities of the imperial court, he brought upon himself at last deposition and exile.” Though a captain,” says he to those who administer the communion, “or a governor, nay, even one adorned with the imperial crown, approach [the table of the Lord] unworthily, prevent him; you have greater authority than he .... Beware lest you excite the Lord to wrath, and give a sword instead of food. And if a new Judas should approach the communion, prevent him. Fear God, not man. If you fear man, he will treat you with scorn; if you fear God, you will appear venerable even to men.”\footnote{Hom. 82 (al. 83) in Matt., toward the close (in Montfaucon’s edition of Chrys., tom. vii. p. 789 sq.). Comp. his exposition of 1Cor. xi. 27, 28, in Hom. 27 and 28, in 1Corinth. (English translation in the Oxford Library of the Fathers, etc., p. 379 sqq., and 383 sqq.).} Synesius excommunicated the worthless governor of Pentapolis, Andronicus, for his cruel oppression of the poor and contempt of the exhortations of the bishop, and the discipline attained the desired effect. The most noted example of church discipline is the encounter between Ambrose and Theodosius I. in Milan about the year 390. The bishop refused the powerful and orthodox emperor the communion, and thrust him back from the threshold of the church, because in a tempest of rage he had caused seven thousand persons in Thessalonica., regardless of rank, sex, or guilt, to be hewn down by his soldiers in horrible cruelty on account of a riot. Eight months afterward Ambrose gave him absolution at his request, after he had submitted to the public penance of the church and promised in future not to execute a death penalty until thirty days after the pronouncing of it, that he might have time to revoke it if necessary, and to exercise mercy.\footnote{This occurrence is related by Ambrosehimself, in 395, in his funeral discourse on Theodosius (de obitu Theod. c. 34, in the Bened. ed. of his works, tom. ii. p. 1207), in these words: “Deflevit in ecclesia publice peccatum suum, quod ei aliorum fraude obrepserat; gemitu et lacrymis oravit veniam. Quod privati erubescunt, non erubuit imperator, publice agere poenitentiam; neque ullus postea dies fuit quo non illum doleret errorem. Quid, quod praeclaram adeptus victoriam; tamen quia hostes in acie prostrati sunt abstinuit a consortio sacramentorum, donec Domini circa se gratiam filiorum experiretur adventu.” Also by his biographer Paulinus (de vita Ambros. c. 24), by Augustine(De Civit. Dei, v. 26), by the historians Theodoret (v. 17), Sozomen (vii. 25), and Rufinus (xi. 18).} Here Ambrose certainly vindicated—though perhaps not without admixture of hierarchical loftiness—the dignity and rights of the church against the state, and the claims of Christian temperance and mercy against gross military power.” Thus,” says a modern historians “did the church prove, in a time of unlimited arbitrary power, the refuge of popular freedom, and saints assume the part of tribunes of the people.”\footnote{Hase, Church History, § 117 (p. 161, 7th ed.)}

\xsection{The Donatist Schism. External History}

§ 69. The Donatist Schism. External History.

I. Sources. Augustine: Works against the Donatists (Contra epistolam Parmeniani, libri iii.; De baptismo, contra Donatistas, libri vii.; Contra literas Petiliani, libri iii.; De Unitate Ecclesiae, lib. unus; Contra Cresconium, grammaticum Donat., libri iv.; Breviculus Collationis cum Donatistis; Contra Gaudentium, etc.), in the 9th vol. of his Opera, ed. Bened. (Paris, 1688). Optatus Milevitanus (about 370): De schismate Donatistarum. L. E. Du Pin: Monumenta vett. ad Donatist. Hist. pertinentia, Par. 1700. Excerpta et Scripta vetera ad Donatistarum Historiam pertinentia, at the close of the ninth volume of the Bened. ed. of Augustine’s works.

II. Literature. Valesius: De schism. Donat. (appended to his ed. of Eusebius). Walch: Historie der Ketzereien, etc., vol. iv. Neander: Allg. K. G. ii. 1, p. 366 sqq. (Torrey’s Engl. translation, ii. p. 182 sqq.). A. Roux: De Augustine adversario Donat. Lugd. Bat. 1838. F. Ribbeck: Donatus u. Augustinus, oder der erste entscheidende Kampf zwischen Separatismus u. Kirche., Elberf. 1858. (The author was for a short time a Baptist, and then returned to the Prussian established church, and wrote this work against separatism.)

Donatism was by far the most important schism in the church of the period before us. For a whole century it divided the North African churches into two hostile camps. Like the schisms of the former period,\footnote{Comp. vol. i. § 115, p. 447 sqq.} it arose from the conflict of the more rigid and the more indulgent theories of discipline in reference to the restoration of the lapsed. But through the intervention of the Christianized state, it assumed at the same time an ecclesiastico-political character. The rigoristic penitential discipline had been represented in the previous period especially by the Montanists and Novatians, who were still living; while the milder principle and practice had found its most powerful support in the Roman church, and, since the time of Constantine, had generally prevailed.

The beginnings of the Donatist schism appear in the Dioclesian persecution, which revived that controversy concerning church discipline and martyrdom. The rigoristic party, favored by Secundus of Tigisis, at that time primate of Numidia, and led by the bishop Donatus of Casae Nigrae, rushed to the martyr’s crown with fanatical contempt of death, and saw in flight from danger, or in the delivering up of the sacred books, only cowardice and treachery, which should forever exclude from the fellowship of the church. The moderate party, at whose head stood the bishop Mensurius and his archdeacon and successor Caecilian, advocated the claims of prudence and discretion, and cast suspicion on the motives of the forward confessors and martyrs. So early as the year 305 a schism was imminent, in the matter of an episcopal election for the city of Cita. But no formal outbreak occurred until after the cessation of the persecution in 311; and then the difficulty arose in connection with the hasty election of Caecilian to the bishopric of Carthage. The Donatists refused to acknowledge him, because in his ordination the Numidian bishops were slighted, and the service was performed by the bishop Felix of Aptungis, or Aptunga, whom they declared to be a traditor, that is, one who had delivered up the sacred writings to the heathen persecutors. In Carthage itself he had many opponents, among whom were the elders of the congregation (seniores plebis), and particularly a wealthy and superstitious widow, Lucilla, who was accustomed to kiss certain relics before her daily communion, and seemed to prefer them to the spiritual power of the sacrament. Secundus of Tigisis and seventy Numidian bishops, mostly of the rigoristic school, assembled at Carthage deposed and excommunicated Caecilian, who refused to appear, and elected the lector Majorinus, a favorite of Lucilla, in his place. After his death, in 315, Majorinus was succeeded by Donatus, a gifted man, of fiery energy and eloquence, revered by his admirers as a wonder worker, and styled The Great. From this man, and not from the Donatus mentioned above, the name of the party was derived.\footnote{“Pars Donati, Donatistae, Donatiani.” Previously they were commonly called “Pars Majorini.” Optatus of Mileve seems, indeed, to know of only one Donatus. But the Donatists expressly distinguish Donatus Magnus of Carthage from Donatus a Casis Nigris. Likewise Augustine, Contra Cresconium Donat, ii. 1; though he himself had formerly confounded the two.}

Each party endeavored to gain churches abroad to its side, and thus the schism spread. The Donatists appealed to the emperor Constantine—the first instance of such appeal, and a step which they afterward had to repent. The emperor, who was at that time in Gaul, referred the matter to the Roman bishop Melchiades (Miltiades) and five Gallican bishops, before whom the accused Caecilian and ten African bishops from each side were directed to appear. The decision went in favor of Caecilian, and he was now, except in Africa, universally regarded as the legitimate bishop of Carthage. The Donatists remonstrated. A second investigation, which Constantine intrusted to the council of Arles (Arelate) in 314, led to the same result. When the Donatists hereupon appealed from this ecclesiastical tribunal to the judgment of the emperor himself, he likewise declared against them at Milan in 316, and soon afterward issued penal laws against them, threatening them with the banishment of their bishops and the confiscation of their churches.

Persecution made them enemies of the state whose help they had invoked, and fed the flame of their fanaticism. They made violent resistance to the imperial commissioner, Ursacius, and declared that no power on earth could induce them to hold church fellowship with the “rascal” (nebulo) Caecilian. Constantine perceived the fruitlessness of the forcible restriction of religion, and, by an edict in 321, granted the Donatists full liberty of faith and worship. He remained faithful to this policy of toleration, and exhorted the Catholics to patience and indulgence. At a council in 330 the Donatists numbered two hundred and seventy bishops.

Constans, the successor of Constantine, resorted again to violent measures; but neither threats nor promises made any impression on the party. It came to blood. The Circumcellions, a sort of Donatist mendicant monks, who wandered about the country among the cottages of the peasantry,\footnote{“Cellas circumientes rusticorum.” Hence the name \emph{Circumcelliones}. But they called themselves \emph{Milites Christi Agonistici}. Their date and origin are uncertain. According to Optatus of Mileve, they first appeared under Constans, in 347.} carried on plunder, arson, and murder, in conjunction with mutinous peasants and slaves, and in crazy zeal for the martyr’s crown, as genuine soldiers of Christ, rushed into fire and water, and threw themselves down from rocks. Yet there were Donatists who disapproved this revolutionary frenzy. The insurrection was suppressed by military force; several leaders of the Donatists were executed, others were banished, and their churches were closed or confiscated. Donatus the Great died in exile. He was succeeded by one Parmenianus.

Under Julian the Apostate the Donatists again obtained, with all other heretics and schismatics, freedom of religion, and returned to the possession of their churches, which they painted anew, to redeem them from their profanation by the Catholics. But under the subsequent emperors their condition grew worse, both from persecutions without and dissensions within. The quarrel between the two parties extended into all the affairs of daily life; the Donatist bishop Faustinus of Hippo, for example, allowing none of the members of his church to bake bread for the Catholic inhabitants.

\xsection{Augustine and the Donatists. Their Persecution and Extinction}

§ 70. Augustine and the Donatists. Their Persecution and Extinction.

At the end of the fourth century, and in the beginning of the fifth, the great Augustine, of Hippo, where there was also a strong congregation of the schismatics, made a powerful effort, by instruction and persuasion, to reconcile the Donatists with the Catholic church. He wrote several works on the subject, and set the whole African church in motion against them. They feared his superior dialectics, and avoided him wherever they could. The matter, however, was brought, by order of the emperor in 411, to a three days’ arbitration at Carthage, attended by two hundred and eighty-six Catholic bishops and two hundred and seventy-nine Donatist.\footnote{Augustinegives an account of the debate in his Breviculus Collationis cum Donatists (Opera, tom. ix. p. 545-580).}

Augustine, who, in two beautiful sermons before the beginning of the disputation, exhorted to love, forbearance and meekness, was the chief speaker on the part of the Catholics Petilian, on the part of the schismatics. Marcellinus, the imperial tribune and notary, and a friend of Augustine, presided, and was to pass the decisive judgment. This arrangement was obviously partial, and secured the triumph of the Catholics. The discussions related to two points: (1) Whether the Catholic bishops Caecilian and Felix of Aptunga were traditors; (2) Whether the church lose her nature and attributes by fellowship with heinous sinners. The balance of skill and argument was on the side of Augustine, though the Donatists brought much that was forcible against compulsion in religion, and against the confusion of the temporal and the spiritual powers. The imperial commissioner, as might be expected, decided in favor of the Catholics. The separatists nevertheless persisted in their view, but their appeal to the emperor continued unsuccessful.

More stringent civil laws were now enacted against them, banishing the Donatist clergy from their country, imposing fines on the laity, and confiscating the churches. In 415 they were even forbidden to hold religious assemblies, upon pain of death.

Augustine himself, who had previously consented only to spiritual measures against heretics, now advocated force, to bring them into the fellowship of the church, out of which there was no salvation. He appealed to the command in the parable of the supper, Luke, xiv. 23, to “compel them to come in;” where, however, the “compel” (ἀνάγκασον) is evidently but a vivid hyperbole for the holy zeal in the conversion of the heathen, which we find, for example, in the apostle Paul.\footnote{On Augustine’s view Comp. § 27, toward the close.}

New eruptions of fanaticism ensued. A bishop Gaudentius threatened, that if the attempt were made to deprive him of his church by force, he, would burn himself with his congregation in it, and vindicated this intended suicide by the example of Rhazis, in the second book of Maccabees (ch. xiv.).

The conquest of Africa by the Arian Vandals in 428 devastated the African church, and put an end to the controversy, as the French Revolution swept both Jesuitism and Jansenism away. Yet a remnant of the Donatists, as we learn from the letters of Gregory I., perpetuated itself into the seventh century, still proving in their ruins the power of a mistaken puritanic zeal and the responsibility and guilt of state-church persecution. In the seventh century the entire African church sank under the Saracenic conquest.

\xsection{Internal History of the Donatist Schism. Dogma of the Church}

§ 71. Internal History of the Donatist Schism. Dogma of the Church.

The Donatist controversy was a conflict between separatism and catholicism; between ecclesiastical purism and ecclesiastical eclecticism; between the idea of the church as an exclusive community of regenerate saints and the idea of the church as the general Christendom of state and people. It revolved around the doctrine of the essence of the Christian church, and, in particular, of the predicate of holiness. It resulted in the completion by Augustine of the catholic dogma of the church, which had been partly developed by Cyprian in his conflict with a similar schism.\footnote{Comp. vol. i § 111, 115, and 131.}

The Donatists, like Tertullian in his Montanistic writings, started from an ideal and spiritualistic conception of the church as a fellowship of saints, which in a sinful world could only be imperfectly realized. They laid chief stress on the predicate of the subjective holiness or personal worthiness of the several members, and made the catholicity of the church and the efficacy of the sacraments dependent upon that. The true church, therefore, is not so much a school of holiness, as a society of those who are already holy; or at least of those who appear so; for that there are hypocrites not even the Donatists could deny, and as little could they in earnest claim infallibility in their own discernment of men. By the toleration of those who are openly sinful, the church loses, her holiness, and ceases to be church. Unholy priests are incapable of administering sacraments; for how can regeneration proceed from the unregenerate, holiness from the unholy? No one can give what he does not himself possess. He who would receive faith from a faithless man, receives not faith but guilt.\footnote{Aug. Contra literas Petil. l. i. cap. 5 (tom. ix. p. 208): “Qui fidem a perfido sumserit, non fidem percipit, sed reatum; omnis enim res origine et radice consistit, et si caput non habet aliquid, nihil est.”} It was on this ground, in fact, that they rejected the election of Caecilian: that he had been ordained bishop by an unworthy person. On this ground they refused to recognize the Catholic baptism as baptism at all. On this point they had some support in Cyprian, who likewise rejected the validity of heretical baptism, though not from the separatist, but from the catholic point of view, and who came into collision, upon this question, with Stephen of Rome.\footnote{Comp. vol. i. § 104, p. 404 sqq.}

Hence, like the Montanists and Novatians, they insisted on rigorous church discipline, and demanded the excommunication of all unworthy members, especially of such as had denied their faith or given up the Holy Scriptures under persecution. They resisted, moreover, all interference of the civil power in church affairs; though they themselves at first had solicited the help of Constantine. In the great imperial church, embracing the people in a mass, they saw a secularized Babylon, against which they set themselves off, in separatistic arrogance, as the only true and pure church. In support of their views, they appealed to the passages of the Old Testament, which speak of the external holiness of the people of God, and to the procedure of Paul with respect to the fornicator at Corinth.

In opposition to this subjective and spiritualistic theory of the church, Augustine, as champion of the Catholics, developed the objective, realistic theory, which has since been repeatedly reasserted, though with various modifications, not only in the Roman church, but also in the Protestant, against separatistic and schismatic sects. He lays chief stress on the catholicity of the church, and derives the holiness of individual members and the validity of ecclesiastical functions from it. He finds the essence of the church, not in the personal character of the several Christians, but in the union of the whole church with Christ. Taking the historical point of view, he goes back to the founding of the church, which may be seen in the New Testament, which has spread over all the world, and which is connected through the unbroken succession of bishops with the apostles and with Christ. This alone can be the true church. It is impossible that she should all at once disappear from the earth, or should exist only in the African sect of the Donatists.\footnote{Augustine, ad Catholicos Epistola contra Donatistas, usually quoted under the shorter title, De unitate Ecclesiae, c. 12 (Bened. ed. tom. ix. p. 360): “Quomodo coeptum sit ab Jerusalem, et deinde processum in Judaeam et Samariam, et inde in totam terram, ubi adhuc crescit ecclesia, donec usque in finem etiam reliquas gentes, ubi adhuc non est, obtineat, scripturis sanctis testibus consequenter ostenditur; quisquis aliud evangelizaverit, anathema sit. Aliud autem evangelizat, qui periisse dicit de caetero mundo ecclesiam et in parte Donati in sola Africa remansisse dicit. Ergo anathema sit. Aut legat mihi hoc in scripturis sanctis, et non sit anathema.”} What is all that they may say of their little heap, in comparison with the great catholic Christendom of all lands? Thus even numerical preponderance here enters as an argument; though under other circumstances it may prove too much, and would place the primitive church at a clear disadvantage in comparison with the prevailing Jewish and heathen masses, and the Evangelical church in its controversy with the Roman Catholic.

From the objective character of the church as a divine institution flows, according to the catholic view, the efficacy of all her functions, the Sacraments in particular. When Petilian, at the Collatio cum Donatistis, said: “He who receives the faith from a faithless priest, receives not faith, but guilt,” Augustine answered: “But Christ is not unfaithful (perfidus), from whom I receive faith (fidem), not guilt (reatum). Christ, therefore, is properly the functionary, and the priest is simply his organ.” “My origin,” said Augustine on the same occasion, “is Christ, my root is Christ, my head is Christ. The seed, of which I was born, is the word of God, which I must obey even though the preacher himself practise not what he preaches. I believe not in the minister by whom I am baptized, but in Christ, who alone justifies the sinner and can forgive guilt.”\footnote{Contra literas Petiliani, l. i. c. 7 (Opera, tom. ix. p. 209): “Origo mea Christus est, radix mea Christus est, caput meum Christus est.” ... In the same place: “Me innocentem non facit, nisi qui mortuus est propter delicta nostra et resurrexit propter justificationem nostram. Non enim in ministrum, per quem baptizor, credo; sed in cum qui justificat impium, ut deputetur mihi fides in justitiam.”}

Lastly, in regard to church discipline, the opponents of the Donatists agreed with them in considering it wholesome and necessary, but would keep it within the limits fixed for it by the circumstances of the time and the fallibility of men. A perfect separation of sinners from saints is impracticable before the final judgment. Many things must be patiently borne, that greater evil may be averted, and that those still capable of improvement may be improved, especially where the offender has too many adherents.” Man,” says Augustine, “should punish in the spirit of love, until either the discipline and correction come from above, or the tares are pulled up in the universal harvest.”\footnote{Aug. Contra Epistolam Parmeniani, l. iii. c. 2, § 10-15 (Opera, tom. ix. p. 62-66).} In support of this view appeal was made to the Lord’s parables of the tares among the wheat, and of the net which gathered together of every kind (Matt. xiii.). These two parables were the chief exegetical battle ground of the two parties. The Donatists understood by the field, not the church, but the world, according to the Saviour’s own exposition of the parable of the tares;\footnote{Breviculus Collat. c. Don. Dies tert. c. 8, § 10 (Opera, ix. p. 559): “Zizania inter triticum non in ecclesia, sed in ipso mundo permixta dixerunt, quoniam Dominus ait, \emph{Ager est} \emph{mundus}“ (Matt. xiii. 38). As to the exegetical merits of the controversy see Trench’s “Notes on the Parables,” p. 83 sqq. (9th Lond. edition, 1863), and Lange’s Commentary on Matt. xiii. (Amer. ed. by Schaff, p. 244 sqq.).} the Catholics replied that it was the kingdom of heaven or the church to which the parable referred as a whole, and pressed especially the warning of the Saviour not to gather up the tares before the final harvest, lest they root up also the wheat with them. The Donatists, moreover, made a distinction between unknown offenders, to whom alone the parable of the net referred, and notorious sinners. But this did not gain them much; for if the church compromises her character for holiness by contact with unworthy persons at all, it matters not whether they be openly unworthy before men or not, and no church whatever would be left on earth.

On the other hand, however, Augustine, who, no more than the Donatists, could relinquish the predicate of holiness for the church, found himself compelled to distinguish between a true and a mixed, or merely apparent body of Christ; forasmuch as hypocrites, even in this world, are not in and with Christ, but only appear to be.\footnote{Corpus Christi \emph{verum} atque \emph{permixtum}, or \emph{verum} atque \emph{simulatum}. Comp. De doctr. Christ. iii. 32, as quoted below in full.} And yet he repelled the Donatist charge of making two churches. In his view it is one and the same church, which is now mixed with the ungodly, and will hereafter be pure, as it is the same Christ who once died, and now lives forever, and the same believers, who are now mortal and will one day put on immortality’.\footnote{Breviculus Collationis cum Donatistis, Dies tertius, cap. 10, § 19 and 20 (Opera, ix. 564): “Deinde calumniantes, quod duas ecclesias Catholici dixerint, unam quae nunc habet permixtos malos, aliam quae post resurrectionem eos non esset habitura: veluti non iidem futuri essent sancti cum Christo regnaturi, qui nunc pro ejus nomine cum juste vivunt tolerant malos .... De duabus etiam ecclesiis calumniam eorum Catholici refutarunt, identidem expressius ostendentes, quid dixerint, id est, non eam ecclesiam, quos nunc habet permixtos malos, alienam se dixisse a regno Dei, ubi non erunt mali commixti, sed eandem ipsam unam et sanctam ecclesiam nunc esse aliter tunc autem aliter futuram, nunc habere malos mixtos, tunc non habituram ... sicut non ideo duo Christi, quia prior mortuus postea non moriturus.”}

With some modification we may find here the germ of the subsequent Protestant distinction of the visible and invisible church; which regards the invisible, not as another church, but as the ecclesiola in ecclesia (or ecclesiis), as the smaller communion of true believers among professors, and thus as the true substance of the visible church, and as contained within its limits, like the soul in the body, or the kernel in the shell. Here the moderate Donatist and scholarly theologian, Tychoius,\footnote{Or Tichonius, as Augustinespells the name. Although himself a Donatist, he wrote against them, “qui contra Donatistas invictissime scripsit, cum fuerit Donatista” (says Aug. De doctr. Christ. l. iii. c. 30, § 42). He was opposed to rebaptism and acknowledged the validity of the Catholic sacraments; but he was equally opposed to the secularism of the Catholic church and its mixture with the state, and adhered to the strict discipline of the Donatists. Of his works only one remains, viz., Liber regularum, or de septem regulis, a sort of Biblical hermeneutics, or a guide for the proper understanding of the mysteries of the Bible. It was edited by Gallandi, in his Bibliotheca Veterum Patrum, tom. viii. p. 107-129. Augustinenotices these rules at length in his work De doctrina Christiana, lib. iii. c. 30 sqq. (Opera, ed. Bened. tom. iii. p. 57 sqq.). Tychonius seems to have died before the close of the fourth century. Comp. on him Tillemont, Memoires, tom. vi. p. 81 sq., and an article of A. Vogel, in Herzog’s Real-Encyclopaedie, vol. xvi. p. 534-536.} approached Augustine; calling the church a twofold body of Christ,\footnote{“Corpus Domini bipartitum.” This was the second of his rules for the true understanding of the Scriptures.} of which the one part embraces the true Christians, the other the apparent.\footnote{Augustineobjects only to his mode of expression, De doctr. Christ. iii. 32 (tom. iii. 58): “Secunda [regula Tichonii] est \emph{de Domini corpore bipartito}; non enim revera Domini corpus est, quod cum illo non erit in aeternum; sed dicendum fuit de \emph{Domini corpore vero} atque \emph{permixto}, aut \emph{vero} atque \emph{simulato}, vel quid aliud; quia non solum in aeternum, verum etiam nunc hypocrites non cum illo esse dicendi sunt, quamvis in ejus esse videantur ecclesia, unde poterat ista regula et sic appellari, ut diceretur de permixta ecclesia.” Comp. also Dr. Baur, K. G. vom 4-6 Jahrh., p. 224.} In this, as also in acknowledging the validity of the Catholic baptism, Tychonius departed from the Donatists; while he adhered to their views on discipline and opposed the Catholic mixture of the church and the world. But neither he nor Augustine pursued this distinction to any clearer development. Both were involved, at bottom, in the confusion of Christianity with the church, and of the church with a particular outward organization.

\xsection{The Roman Schism of Damasus and Ursinus}

§ 72. The Roman Schism of Damasus and Ursinus.

Rufinus: Hist. Eccl. ii. 10. Hieronymus: Chron. ad ann. 366. Socrates: H. E. iv. 29 (all in favor of Damasus). Faustinus et Marcellinus (two presbyters of Ursinus): Libellus precum ad Imper. Theodos. in Bibl. Patr. Lugd. v. 637 (in favor of Ursinus). With these Christian accounts of the Roman schism may be compared the impartial statement of the heathen historian Ammianus Marcellinus, xxvii. c. 3, ad ann. 367.

The church schism between Damasus and Ursinus (or Ursicinus) in Rome, had nothing to do with the question of discipline, but proceeded partly from the Arian controversy, partly from personal ambition.\footnote{Ammianus Marc., l.c., intimates the latter: “Damasus et Ursinus supra humanum modum ad rapiendam episcopatus sedem ardentes scissis studiis asperrimo conflictabantur,” etc.} For such were the power and splendor of the court of the successor of the Galilean fisherman, even at that time, that the distinguished pagan senator, Praetextatus, said to Pope Damasus: “Make me a bishop of Rome, and I will be a Christian to-morrow.”\footnote{This is related even by St. Jerome(Comp. above § 53, p. 267, note), and goes to confirm the statements of Ammianus.} The schism presents a mournful example of the violent character of the episcopal elections at Rome. These elections were as important events for the Romans as the elections of the emperors by the Praetorian soldiers had formerly been. They enlisted and aroused all the passions of the clergy and the people.

The schism originated in the deposition and banishment of the bishop Tiberius, for his orthodoxy, and the election of the Arian Felix\footnote{Athanasius (Historia Arianorum ad Monachos, § 75, Opera ed. Bened. i. p. 389), and Socrates (H. E. ii. 37), decidedly condemn him as an Arian. Nevertheless this heretic and anti-pope has been smuggled into the Roman catalogue of saints and martyrs. Gregory XIII instituted an investigation into the matter, which was terminated by the sudden discovery of his remains, with the inscription: “Pope and Martyr.”} as pope in opposition by the arbitrary will of the emperor Constantius (a.d. 355). Liberius, having in his exile subscribed the Arian creed of Sirmium, \footnote{According to Baronius, ad a. 357, the jealousy of Felix was the Delilah, who robbed the catholic Samson (Liberius) of his strength.} was in 358 reinstated, and Felix retired, and is said to have subsequently repented his defection to Arianism. The parties, however, continued.

After the death of Liberius in 366, Damasus was, by the party of Felix, and Ursinus by the party of Liberius, elected successor of Peter. It came to repeated bloody encounters; even the altar of the Prince of Peace was desecrated, and in a church whither Ursinus had betaken himself, a hundred and thirty-seven men lost their lives in one day.\footnote{Ammian. Marc. l. xxvii. c. 3: “Constat in basilica Sicinini (Sicinii), ubi ritus Christiani est conventiculum, uno die cxxxvii. reperta cadavera peremtorum.” Then he speaks of the pomp and luxury of the Roman bishopric, on account of which it was the object of so passionate covetousness and ambition, and contrasts with it the simplicity and self-denial of the rural clergy. The account is confirmed by Augustine, Brevic. Coll.c. Donat. c. 16, and Hieron. in Chron. an. 367. Socrates, iv. 29, speaks generally of several fights, in which many lives were lost.} Other provinces also were drawn into the quarrel. It was years before Damasus at last, with the aid of the, emperor, obtained undisputed possession of his office, and Ursinus was banished. The statements of the two parties are so conflicting in regard to the priority and legitimacy of election in the two cases, and the authorship of the bloody scenes, that we cannot further determine on which side lay the greater blame. Damasus, who reigned from 367 to 384) is indeed depicted as in other respects a violent man,\footnote{His opponents also charged him with too great familiarity with Roman ladies. The same accusation, however, was made against his friend Jerome, on account of his zeal for the spread of the ascetic life among the Roman matrons.} but he was a man of learning and literary taste, and did good service by his patronage of Jerome’s Latin version of the Bible, and by the introduction of the Latin Psalter into the church song.\footnote{Comp. on Damasus his works, edited by Merenda, Rome, 1754, several epistles of Jerome, Tillemont, tom. viii. 386, and Butler’s Lives of the Saints, sub Dec. 11th.}

\xsection{The Meletian Schism at Antioch}

§ 73. The Meletian Schism at Antioch.

Hieronymus: Chron. ad ann. 864. Chrysostomus: Homilia in S. Patrem nostrum Meletium, archiepiscopum magnae Antiochiae (delivered a.d. 386 or 387, in Montfaucon’s ed. of Chrysost. Opera, tom. ii. p. 518–523). Sozomen: H. E. iv. 28; vii. 10, 11. Theodor.: H. E. V. 3, 35. Socrates: H. E. iii. 9; v. 9, 17. Comp. Walch: Ketzerhistorie, part iv. p. 410 sqq.

The Meletian schism at Antioch\footnote{Not to be confounded with the Meletian schism at Alexandria, which arose in the previous period. Comp. vol. i. § 115 (p. 451).} was interwoven with the Arian controversies, and lasted through more than half a century.

In 361 the majority of the Antiochian church elected as bishop Meletius, who had formerly been an Arian, and was ordained by this party, but after his election professed the Nicene orthodoxy. He was a man of rich persuasive eloquence, and of a sweet and amiable disposition, which endeared him to the Catholics and Arians. But his doctrinal indecision offended the extremists of both parties. When he professed the Nicene faith, the Arians deposed him in council, sent him into exile, and transferred his bishopric to Euzoius, who had formerly been banished with Arius.\footnote{Sozom. H. E. iv. c. 28..} The Catholics disowned Euzoius, but split among themselves; the majority adhered to the exiled Meletius, while the old and more strictly orthodox party, who had hitherto been known as the Eustathians, and with whom Athanasius communicated, would not recognize a bishop of Arian consecration, though Catholic in belief, and elected Paulinus, a presbyter of high character, who was ordained counter-bishop by Lucifer of Calaris.\footnote{This Lucifer was an orthodox fanatic, who afterward himself fell into conflict with Athanasius in Alexandria, and formed a sect of his own, the Luciferians, On rigid principles of church purity. Comp. Socr. iii. 9; Sozom. iii. 15; and Walch, Ketzerhist. iii. 338 sqq}

The doctrinal difference between the Meletians and the old Nicenes consisted chiefly in this: that the latter acknowledged three hypostases in the divine trinity, the former only three prosopa; the one laying the stress on the triplicity of the divine essence, the other on its unity.

The orthodox orientals declared for Meletius, the occidentals and Egyptians for Paulinus, as legitimate bishop of Antioch. Meletius, on returning from exile under the protection of Gratian, proposed to Paulinus that they should unite their flocks, and that the survivor of them should superintend the church alone; but Paulinus declined, since the canons forbade him to take as a colleague one who had been ordained by Arians.\footnote{Theodoret, H. E. lib. iii. 3. He highly applauds the magnanimous proposal of Meletius.} Then the military authorities put Meletius in possession of the cathedral, which had been in the hands of Euzoius. Meletius presided, as senior bishop, in the second ecumenical council (381), but died a few days after the opening of it—a saint outside the communion of Rome. His funeral was imposing: lights were borne before the embalmed corpse, and psalms sung in divers languages, and these honors were repeated in all the cities through which it passed on its transportation to Antioch, beside the grave of St. Babylas.\footnote{Sozom. vii. c. 10. The historian says that the singing of psalms on such occasions was quite contrary to Roman custom.} The Antiochians engraved his likeness on their rings, their cups, and the walls of their bedrooms. So St. Chrysostom informs us in his eloquent eulogy on Meletius.\footnote{Chrysostomsays in the beginning of this oration, that five years had elapsed since Meletius had gone to Jesus. He died in 381, consequently the oration must have been pronounced in 386 or 387.} Flavian was elected his successor, although Paulinus was still alive. This gave rise to fresh troubles, and excited the indignation of the bishop of Rome. Chrysostom labored for the reconciliation of Rome and Alexandria to Flavian. But the party of Paulinus, after his death in 389, elected Evarius as successor († 392), and the schism continued down to the year 413 or 415, when the bishop Alexander succeeded in reconciling the old orthodox remnant with the successor of Meletius. The two parties celebrated their union by a splendid festival, and proceeded together in one majestic stream to the church.\footnote{Theodoret, H. E. l. v. c. 35. Dr. J. R. Kurtz, in his large work on Church History (Handbuch der Kirchengesch. vol. i. part ii. § 181, p. 129) erroneously speaks of a resignation of Alexander, by which he, from love of peace, induced his congregation to acknowledge the Meletian bishop Flavian. But Flavian had died several years before (in 404), and Alexander was himself the second successor of Flavian, the profligate Porphyrius intervening. Theodoret knows nothing of a resignation. Kurtz must be used with considerable caution, as he is frequently inaccurate, and relies too much on secondary authorities.}

Thus a long and tedious schism was brought to a close, and the church of Antioch was permitted at last to enjoy that peace which the Athanasian synod of Alexandria in 362 had desired for it in vain.\footnote{See the Epist. Synodica Conc. Alex. in Mansi’s Councils, tom. iii. p. 345 sqq.}

\xchapter{Public Worship and Religious Customs and Ceremonies}

PUBLIC WORSHIP AND RELIGIOUS CUSTOMS AND CEREMONIES.

I. The ancient Liturgies; the Acts of Councils; and the ecclesiastical writers of the period.

II. The archaeological and liturgical works of Martene, Mamachi, Bona, Muratori, Pelicia, Asseman, Renaudot, Binterim, and Staudenmeier, of the Roman Catholic church; and Bingham, Augusti, Siegel, Alt, Piper, Neale, and Daniel, of the Protestant.

\xsection{The Revolution in Cultus}

§ 74. The Revolution in Cultus.

The change in the legal and social position of Christianity with reference to the temporal power, produced a mighty effect upon its cultus. Hitherto the Christian worship had been confined to a comparatively small number of upright confessors, most of whom belonged to the poorer classes of society. Now it came forth from its secrecy in private houses, deserts, and catacombs, to the light of day, and must adapt itself to the higher classes and to the great mass of the people, who had been bred in the traditions of heathenism. The development of the hierarchy and the enrichment of public worship go hand in hand. A republican and democratic constitution demands simple manners and customs; aristocracy and monarchy surround themselves with a formal etiquette and a brilliant court-life. The universal priesthood is closely connected with a simple cultus; the episcopal hierarchy, with a rich, imposing ceremonial.

In the Nicene age the church laid aside her lowly servant-form, and put on a splendid imperial garb. She exchanged the primitive simplicity of her cultus for a richly colored multiplicity. She drew all the fine arts into the service of the sanctuary, and began her sublime creations of Christian architecture, sculpture, painting, poetry, and music. In place of the pagan temple and altar arose everywhere the stately church and the chapel in honor of Christ, of the Virgin Mary, of martyrs and saints. The kindred ideas of priesthood, sacrifice, and altar became more fully developed and more firmly fixed, as the outward hierarchy grew. The mass, or daily repetition of the atoning sacrifice of Christ by the hand of the priest, became the mysterious centre of the whole system of worship. The number of church festivals was increased; processions, and pilgrimages, and a multitude of significant and superstitious customs and ceremonies were introduced. The public worship of God assumed, if we may so speak, a dramatic, theatrical character, which made it attractive and imposing to the mass of the people, who were as yet incapable, for the most part, of worshipping God in spirit and in truth. It was addressed rather to the eye and the ear, to feeling and imagination, than to intelligence and will. In short, we already find in the Nicene age almost all the essential features of the sacerdotal, mysterious, ceremonial, symbolical cultus of the Greek and Roman churches of the present day.

This enrichment and embellishment of the cultus was, on one hand, a real advance, and unquestionably had a disciplinary and educational power, like the hierarchical organization, for the training of the popular masses. But the gain in outward appearance and splendor was balanced by many a loss in simplicity and spirituality. While the senses and the imagination were entertained and charmed, the heart not rarely returned cold and hungry. Not a few pagan habits and ceremonies, concealed under new names, crept into the church, or were baptized only with water, not with the fire and Spirit of the gospel. It is well known with what peculiar tenacity a people cleave to religious usages; and it could not be expected that they should break off in an instant from the traditions of centuries. Nor, in fact, are things which may have descended from heathenism, to be by any means sweepingly condemned. Both the Jewish cultus and the heathen are based upon those universal religious wants which Christianity must satisfy, and which Christianity alone can truly meet. Finally, the church has adopted hardly a single existing form or ceremony of religion, without at the same time breathing into it a new spirit, and investing it with a high moral import. But the limit of such appropriation it is very hard to fix, and the old nature of Judaism and heathenism which has its point of attachment in the natural heart of man, continually betrayed its tenacious presence. This is conceded and lamented by the most earnest of the church fathers of the Nicene and post-Nicene age, the very persons who are in other respects most deeply involved in the Catholic ideas of cultus.

In the Christian martyr-worship and saint-worship, which now spread with giant strides over the whole Christian world, we cannot possibly mistake the succession of the pagan worship of gods and heroes, with its noisy popular festivities. Augustine puts into the mouth of a heathen the question: “Wherefore must we forsake gods, which the Christians themselves worship with us?” He deplores the frequent revels and amusements at the tombs of the martyrs; though he thinks that allowance should be made for these weaknesses out of regard to the ancient custom. Leo the Great speaks of Christians in Rome who first worshipped the rising sun, doing homage to the pagan Apollo, before repairing to the basilica of St. Peter. Theodoret defends the Christian practices at the graves of the martyrs by pointing to the pagan libations, propitiations, gods, and demigods. Since Hercules, Aesculapitis, Bacchus, the Dioscuri, and many other objects of pagan worship were mere deified men, the Christians, he thinks, cannot be blamed for honoring their martyrs—not making them gods but venerating them as witnesses and servants of the only, true God. Chrysostom mourns over the theatrical customs, such as loud clapping in applause, which the Christians at Antioch and Constantinople brought with them into the church. In the Christmas festival, which from the fourth century spread from Rome over the entire church, the holy commemoration of the birth of the Redeemer is associated—to this day, even in Protestant lands—with the wanton merriments of the pagan Saturnalia. And even in the celebration of Sunday, as it was introduced by Constantine, and still continues on the whole continent of Europe, the cultus of the old sun-god Apollo mingles, with the remembrance of the resurrection of Christ; and the widespread profanation of the Lord’s Day, especially on the continent of Europe, demonstrates the great influence which heathenism still exerts upon Roman and Greek Catholic, and even upon Protestant, Christendom.

\xsection{The Civil and Religious Sunday}

§ 75. The Civil and Religious Sunday.

Geo. Holden: The Christian Sabbath. Lond. 1825 (see ch. v.). John T. Baylee: History of the Sabbath. Lond. 1857 (see chs. x.-xiii.). James Aug. Hessey: Sunday, its Origin, History, and present Obligation; Bampton Lectures preached before the University of Oxford. Lond. 1860 (Patristic and high-Anglican). James Gilfillan: The Sabbath viewed in the Light of Reason, Revelation, and History, with Sketches of its Literature. Edinb. and New York, 1862 (The Puritan and Anglo-American view). Robert Cox: The Literature on the Sabbath Question. Edinb. 1865, 2 vols. (Latitudinarian, but very full and learned).

The observance of Sunday originated in the time of the apostles, and ever since forms the basis of public worship, with its ennobling, sanctifying, and cheering influences, in all Christian lands.

The Christian Sabbath is, on the one hand, the continuation and the regeneration of the Jewish Sabbath, based upon God’s resting from the creation and upon the fourth commandment of the decalogue, which, as to its substance, is not of merely national application, like the ceremonial and civil law, but of universal import and perpetual validity for mankind. It is, on the other hand, a new creation of the gospel, a memorial of the resurrection of Christ and of the work of redemption completed and divinely sealed thereby. It rests, we may say, upon the threefold basis of the original creation, the Jewish legislation, and the Christian redemption, and is rooted in the physical, the moral, and the religious wants of our nature. It has a legal and an evangelical aspect. Like the law in general, the institution of the Christian Sabbath is a wholesome restraint upon the people, and a schoolmaster to lead them to Christ. But it is also strictly evangelical: it was originally made for the benefit of man, like the family, with which it goes back beyond the fall to the paradise of innocence, as the second institution of God on earth; it was “a delight” to the pious of the old dispensation (Isa. lviii. 13), and now, under the new, it is fraught with the glorious memories and blessings of Christ’s resurrection and the outpouring of the Holy Spirit. The Christian Sabbath is the ancient Sabbath baptized with fire and the Holy Ghost, regenerated, spiritualized, and glorified. It is the connecting link of creation and redemption, of paradise lost, and paradise regained, and a pledge and preparation for the saints’ everlasting rest in heaven.\footnote{For a fuller exposition of the Author’s views on the Christian Sabbath, \emph{see} his Essay on the Anglo-American Sabbath (English and German), New York, 1863.}

The ancient church viewed the Sunday mainly, we may say, one-sidedly and exclusively, from its Christian aspect as a new institution, and not in any way as a continuation of the Jewish Sabbath. It observed it as the day of the commemoration of the resurrection or of the now spiritual creation, and hence as a day of sacred joy and thanksgiving, standing in bold contrast to the days of humiliation and fasting, as the Easter festival contrasts with Good Friday.

So long as Christianity was not recognized and protected by the state, the observance of Sunday was purely religious, a strictly voluntary service, but exposed to continual interruption from the bustle of the world and a hostile community. The pagan Romans paid no more regard to the Christian Sunday than to the Jewish Sabbath.

In this matter, as in others, the accession of Constantine marks the beginning of a new era, and did good service to the church and to the cause of public order and morality. Constantine is the founder, in part at least, of the civil observance of Sunday, by which alone the religious observance of it in the church could be made universal and could be properly secured. In the year 321 he issued a law prohibiting manual labor in the cities and all judicial transactions, at a later period also military exercises, on Sunday.\footnote{Lex Constantini a. 321 (Cod. Just. l. iii., Tit. 12, 3): Imperator Constantinus Aug. Helpidio: “Omnes judices, urbanaeque plebes et cunctarum artium officia venerabili die Solis quiescant. Ruri tamen positi agrorum culturae libere licenterque inserviant, quoniam frequenter evenit, ut non aptius alio die frumenta sulcis aut vineae scrobibus mandentur, ne occasione momenti pereat commoditas coelesti provisione concessa. Dat. Non. Mart. Crispo ii. et Constantino ii. Coss.” In English: “On the venerable Day of the Sun let the magistrates and people residing in cities rest, and let all workshops be closed. In the country, however, persons engaged in agriculture may freely and lawfully continue their pursuits; because it often happens that another day is not so suitable for grain-sowing or for vine-planting; lest by neglecting the proper moment for such operations the bounty of heaven should be lost. (Given the 7th day of March, Crispus and Constantinebeing consuls each of them for the second time.)” The prohibition of military exercises is mentioned by Eusebius, Vita Const. IV. 19, 20, and seems to refer to a somewhat later period. In this point Constantinewas in advance of modern Christian princes, who prefer Sunday for parades.} He exempted the liberation of slaves, which as an act of Christian humanity and charity, might, with special propriety, take place on that day.\footnote{Cod. Theod. l. ii. tit. 8, 1: “Sicut indignissimum videbatur, diem Solis ... altercantibus jurgiis et noxiis partium contentionibus occupari, ita gratum et jocundum est, eo die, quae sunt maxime votiva, compleri; atque ideo emancipandi et manumittendi die festo cuncti licentiam habeant.”} But the Sunday law of Constantine must not be overrated. He enjoined the observance, or rather forbade the public desecration of Sunday, not under the name of Sabbatum or Dies Domini, but under its old astrological and heathen title, Dies Solis, familiar to all his subjects, so that the law was as applicable to the worshippers of Hercules, Apollo, and Mithras, as to the Christians. There is no reference whatever in his law either to the fourth commandment or to the resurrection of Christ. Besides he expressly exempted the country districts, where paganism still prevailed, from the prohibition of labor, and thus avoided every appearance of injustice. Christians and pagans had been accustomed to festival rests. Constantine made these rests to synchronize, and gave the preference to Sunday, on which day Christians from the beginning celebrated the resurrection of their Lord and Saviour. This and no more was implied in the famous enactment of 321. It was only a step in the right direction, but probably the only one which Constantine could prudently or safely take at that period of transition from the rule of paganism to that of Christianity.

For the army, however, he went beyond the limits of negative and protective legislation, to which the state ought to confine itself in matters of religion, and enjoined a certain positive observance of Sunday, in requiring the Christian soldiers to attend Christian worship, and the heathen soldiers, in the open field, at a given signal, with eyes and hands raised towards heaven, to recite the following, certainly very indefinite, form of prayer: “Thee alone we acknowledge as God, thee we reverence as king, to thee we call as our helper. To thee we owe our victories, by thee have we obtained the mastery of our enemies. To thee we give thanks for benefits already received, from thee we hope for benefits to come. We all fall at thy feet, and fervently beg that thou wouldest preserve to us our emperor Constantine and his divinely beloved sons in long life healthful and victorious.”\footnote{Euseb. Vit. Const. iv. 20.}

Constantine’s successors pursued the Sunday legislation which he had initiated, and gave a legal sanction and civil significance also to other holy days of the church, which have no Scriptural authority, so that the special reverence due to the Lord’s Day was obscured in proportion as the number of rival claims increased. Thus Theodosius I. increased the number of judicial holidays to one hundred and twenty-four. The Valentinians, I. and II., prohibited the exaction of taxes and the collection of moneys on Sunday, and enforced the previously enacted prohibition of lawsuits. Theodosius the Great, in 386, and still more stringently the younger Theodosius, in 425, forbade theatrical performances, and Leo and Anthemius, in 460, prohibited other secular amusements, on the Lord’s Day.\footnote{Cod. Theod. xv. 5, 2, a. 386: “Nullus Solis die populo spectaculum praebeat.” If the emperor’s birthday fell on Sunday, the acknowledgment of it, which was accompanied by games, was to be postponed.} Such laws, however, were probably never rigidly executed. A council of Carthage, in 401, laments the people’s passion for theatrical and other entertainments on Sunday. The same abuse, it is well known, very generally prevails to this day upon the continent of Europe in both Protestant and Roman Catholic countries, and Christian princes and magistrates only too frequently give it the sanction of their example.

Ecclesiastical legislation in like manner prohibited needless mechanical and agricultural labor, and the attending of theatres and other public places of amusement, also hunting and weddings, on Sunday and on feast days. Besides such negative legislation, to which the state must confine itself, the church at the same time enjoined positive observances for the sacred day, especially the regular attendance of public worship, frequent communion, and the payment of free-will offerings (tithes). Many a council here confounded the legal and the evangelical principles, thinking themselves able to enforce by the threatening of penalties what has moral value only as a voluntary act. The Council of Eliberis, in 305, decreed the suspension from communion of any person living in a town who shall absent himself for three Lord’s Days from church. In the same legalistic spirit, the council of Sardica,\footnote{Can. xi. appealing to former ordinances, Comp. Can. Apost. xiii. and xiv. (xiv. and xv.), and the council of Elvira, can. xxi. Hefele: Conciliengesch. i. p. 570.} in 343, and the Trullan council\footnote{Can. lxxx.} of 692, threatened with deposition the clergy who should unnecessarily omit public worship three Sundays in succession, and prescribed temporary excommunication for similar neglect among the laity. But, on the other hand, the councils, while they turned the Lord’s Day itself into a legal ordinance handed down from the apostles, pronounced with all decision against the Jewish Sabbatism. The Apostolic Canons and the council of Gangra (the latter, about 450, in opposition to the Gnostic Manichaean asceticism of the Eustathians) condemn fasting on Sunday.\footnote{Can. Apost. liii. (alias Iii.): “Si quis episcopus aut presbyter aut diaconus in diebus festis non sumit carnem aut vinum, deponatur.” Comp. can. lxvi. (lxv.) and Const. Apost. v. 20. The council of Gangra says in the 18th canon: “If any one, for pretended ascetic reasons, fast on Sunday, let him be anathema.” The same council condemns those who despise the house of God and frequent schismatical assemblies.} In the Greek church this prohibition is still in force, because Sunday, commemorating the resurrection of Christ, is a day of spiritual joy. On the same symbolical ground kneeling in prayer was forbidden on Sunday and through the whole time of Easter until Pentecost. The general council of Nicaeea, in 325, issued on this point in the twentieth canon the following decision: “Whereas some bow the knee on Sunday and on the days of Pentecost [i.e., during the seven weeks after Easter], the holy council, that everything may everywhere be uniform, decrees that prayers be offered to God in a standing posture.” The Trullan council, in 692, ordained in the ninetieth canon: “From Saturday evening to Sunday evening let no one bow the knee.” The Roman church in general still adheres to this practice.\footnote{Comp. the Corpus juris can. c. 13, Dist. 3 de consecr. Roman Catholics, however, always kneel in the reception and adoration of the sacrament.} The New Testament gives no law for such secondary matters; the apostle Paul, on the contrary, just in the season of Easter and Pentecost, before his imprisonment, following an inward dictate, repeatedly knelt in prayer.\footnote{Acts xx. 36; xxi. 5.} The council of Orleans, in 538, says in the twenty-eighth canon: “It is Jewish superstition, that one may not ride or walk on Sunday, nor do anything to adorn the house or the person. But occupations in the field are forbidden, that people may come to the church and give themselves to prayer.”\footnote{Comp. the brief scattered decrees of the councils on the sanctification of Sunday, in Hefele, l.c. i. 414, 753, 760, 761, 794; ii 69, 647, 756; Neale’s Feasts and Fasts; and Gilfillan: The Sabbath, \&c., p. 390.}

As to the private opinions of the principal fathers on this subject, they all favor the sanctification of the Lord’s Day, but treat it as a peculiarly Christian institution, and draw a strong, indeed a too strong, line of distinction between it and the Jewish Sabbath; forgetting that they are one in essence and aim, though different in form and spirit, and that the fourth commandment as to its substance—viz., the keeping holy of one day out of seven—is an integral part of the decalogue or the moral law, and hence of perpetual obligation.\footnote{See the principal patristic passages on the Lord’s Day in Hessey, Sunday, etc., p. 90 ff. and p. 388 ff. Hessey says, p. 114: “In no clearly genuine passage that I can discover in any writer of these two [the fourth and fifth] centuries, or in any public document, ecclesiastical or civil, is the fourth commandment referred to as the ground of the obligation to observe the Lord’s Day.” The Reformers of the sixteenth century, likewise, in their zeal against legalism and for Christian freedom, entertained rather lax views on the Sabbath law. It was left for Puritanism in England, at the close of Queen Elizabeth’s reign, to bring out the perpetuity of the fourth commandment and the legal and general moral feature in the Christian Sabbath. The book of Dr. Bownd, first published in 1595, under the title, “The Doctrine of the Sabbath,” produced an entire revolution on the subject in the English mind, which is visible to this day in the strict observance of the Lord’s Day in England, Scotland, the British Provinces, and the United States. Comp. on Dr. Bownd’s book my Essay above quoted, p. 16 ff., Gilfillan, p. 69 ff., and Hessey, p. 276 ff.} Eusebius calls Sunday, but not the Sabbath, “the first and chief of days and a day of salvation,” and commends Constantine for commanding that “all should assemble together every week, and keep that which is called the Lord’s Day as a festival, to refresh even their bodies and to stir up their minds by divine precepts and instruction.”\footnote{De Laud. Const. c. 9 arid 17.} Athanasius speaks very highly of the Lord’s Day, as the perpetual memorial of the resurrection, but assumes that the old Sabbath has deceased.\footnote{In the treatise: De sabbatis et de circumcisione, which is among the doubtful works of Athanasius.} Macarius, a presbyter of Upper Egypt (350), spiritualizes the Sabbath as a type and shadow of the true Sabbath given by the Lord to the soul—the true and eternal Sabbath, which is freedom from sin.\footnote{Hom. 35.} Hilary represents the whole of this life as a preparation for the eternal Sabbath of the next. Epiphanius speaks of Sunday as an institution of the apostles, but falsely attributes the same origin to the observance of Wednesday and Friday as half fasts. Ambrose frequently mentions Sunday as an evangelical festival, and contrasts it with the defunct legal Sabbath. Jerome makes the same distinction. He relates of the Egyptian coenobites that they “devote themselves on the Lord’s Day to nothing but prayer and reading the Scriptures.” But he mentions also without censure, that the pious Paula and her companions, after returning from church on Sundays, “applied themselves to their allotted works and made garments for themselves and others.” Augustine likewise directly derives Sunday from the resurrection, and not from the fourth commandment. Fasting on that day of spiritual joy he regards, like Ambrose, as a grave scandal and heretical practice. The Apostolical Constitutions in this respect go even still further, and declare: “He that fasts on the Lord’s Day is guilty of sin.” But they still prescribe the celebration of the Jewish Sabbath on Saturday in addition to the Christian Sunday. Chrysostom warns Christians against sabbatizing with the Jews, but earnestly commends the due celebration of the Lord’s Day. Leo the Great, in a beautiful passage—the finest of all the patristic utterances on this subject—lauds the Lord’s Day as the day of the primitive creation, of the Christian redemption, of the meeting of the risen Saviour with the assembled disciples, of the outpouring of the Holy Spirit, of the principal Divine blessings bestowed upon the world.\footnote{Leon. Epist. ix. ad Dioscurum Alex. Episc. c. 1 (Opp. ed. Ballerini, tom. i. col. 630): “Dies resurrectionis Dominicae ... quae tantis divinarum dispositionum mysteriis est consecrata, ut quicquid est a Domino insignius constitutum, in huius piei dignitate sit gestum. In hac mundus sumpsit exordium. In hac per resurrectionem Christi et mors interitum, et vita accepit initium. In hac apostoli a Domino praedicandi omnibus gentibus evangelii tubam sumunt, et inferendum universo mundo sacramentum regenerationis accipiunt. In hac, sicut beatus Joannes evangelista testatur (Joann. xx. 22), congregatis in unum discipulis, januis clausis, cum ad eos Dominus introisset, insufflavit, et dixit: ’\emph{Accipite Spiritum Sanctum; quorum remiseritis peccata, remittuntur eis, et quorum detinueritis, detenta erunt.}’In hac denique promissus a Domino apostolis Spiritus Sanctus advenit: ut coelesti quadam regula insinuatum et traditum noverimus, in illa die celebranda nobis esse mysteria sacerdotalium benedictionum, in qua collata sunt omnia dona gratiarum.”} But he likewise brings it in no connection with the fourth commandment, and with the other fathers leaves out of view the proper foundation of the day in the eternal moral law of God.

Besides Sunday, the Jewish Sabbath also was distinguished in the Eastern church by the absence of fasting and by standing in prayer. The Western church, on the contrary, especially the Roman, in protest against Judaism, observed the seventh day of the week as a fast day, like Friday. This difference between the two churches was permanently fixed by the fifty-fifth canon of the Trullan council of 692: “In Rome fasting is practised on all the Saturdays of Quadragesima [the forty days’ fast before Easter]. This is contrary to the sixty-sixth apostolic canon, and must no longer be done. Whoever does it, if a clergyman, shall be deposed; if a layman, excommunicated.”

Wednesday and Friday also continued to be observed in many countries as days commemorative of the passion of Christ (dies stationum), with half-fasting. The Latin church, however, gradually substituted fasting on Saturday for fasting on Wednesday.

Finally, as to the daily devotions: the number of the canonical hours was enlarged from three to seven (according to Ps. cxix. 164: “Seven times in a day will I praise thee But they were strictly kept only in the cloisters, under the technical names of matina (about three o’clock), prima (about six), tertia (nine), sexta (noon), nona (three in the afternoon), vesper (six), completorium (nine), and mesonyctium or vigilia (midnight). Usually two nocturnal prayers were united. The devotions consisted of prayer, singing, Scripture reading, especially in the Psalms, and readings from the histories of the martyrs and the homilies of the fathers. In the churches ordinarily only morning and evening worship was held. The high festivals were introduced by a night service, the vigils.

\xsection{The Church Year}

§ 76. The Church Year.

R. Hospinian: Festa Christian. (Tiguri, 1593) Genev. 1675. M. A. Nickel (R.C.): Die heil. Zeiten u. Feste nach ihrer Entstehung u. Feier in der Kath. Kirche, Mainz, 1825 sqq. 6 vols. Pillwitz: Geschichte der heil. Zeiten. Dresden, 1842. E. Ranke: Das kirchliche Pericopensystem aus den aeltesten Urkunden dargelegt. Berlin, 1847. Fr. Strauss (late court preacher and professor in Berlin): Das evangelische Kirchenjahr. Berl. 1850. Lisco: Das christliche Kirchenjahr. Berl. (1840) 4th ed. 1850. Bobertag: Das evangelische Kirchenjahr, \&c. Breslau, 1857. Comp. also Augusti: Handbuch der Christlichen Archaeologie, vol. i. (1836), pp. 457–595.

After the, fourth century, the Christian year, with a cycle of regularly recurring annual religious festivals, comes forth in all its main outlines, though with many fluctuations and variations in particulars, and forms thenceforth, so to speak, the skeleton of the Catholic cultus.

The idea of a religious year, in distinction from the natural and from the civil year, appears also in Judaism, and to some extent in the heathen world. It has its origin in the natural necessity of keeping alive and bringing to bear upon the people by public festivals the memory of great and good men and of prominent events. The Jewish ecclesiastical year was, like the whole Mosaic cultus, symbolical and typical. The Sabbath commemorated the creation and the typical redemption, and pointed forward to the resurrection and the true redemption, and thus to the Christian Sunday. The passover pointed to Easter, and the feast of harvest to the Christian Pentecost. The Jewish observance of these festivals originally bore an earnest, dignified, and significant character, but in the hands of Pharisaism it degenerated very largely into slavish Sabbatism and heartless ceremony, and provoked the denunciation of Christ and the apostles. The heathen festivals of the gods ran to the opposite extreme of excessive sensual indulgence and public vice.\footnote{Philo, in his Tract de Cherubim (in Augusti, l.c. p. 481 sq.), paints this difference between the Jewish and heathen festivals in strong colors; and the picture was often used by the church fathers against the degenerate pagan character of the Christian festivals.}

The peculiarity of the Christian year is, that it centres in the person and work of Jesus Christ, and is intended to minister to His glory. In its original idea it is a yearly representation of the leading events of the gospel history; a celebration of the birth, passion, and resurrection of Christ, and of the outpouring of the Holy Spirit, to revive gratitude and devotion. This is the festival part, the semestre Domini. The other half, not festal, the semestre ecclesiae, is devoted to the exhibition of the life of the Christian church, its founding, its growth, and its consummation, both is a whole, and in its individual members, from the regeneration to the resurrection of the dead. The church year is, so to speak, a chronological confession of faith; a moving panorama of the great events of salvation; a dramatic exhibition of the gospel for the Christian people. It secures to every important article of faith its place in the cultus of the church, and conduces to wholeness and soundness of Christian doctrine, as against all unbalanced and erratic ideas.\footnote{This last thought is well drawn out by W. Archer Butler in one of his sermons: “It is the chief advantage of that religious course of festivals by which the church fosters the piety of her children, that they tend to preserve a due proportion and equilibrium in our religious views. We have all a tendency to adopt particular views of the Christian truths, to insulate certain doctrines from their natural accompaniments, and to call our favorite fragment the gospel. We hold a few texts so near our eyes that they hide all the rest of the Bible. The church festival system spreads the gospel history in all its fulness across the whole surface of the sacred year. It is a sort of chronological creed, and forces us, whether we will or no, by the very revolution of times and seasons, to give its proper place and dignity to every separate article. ’Day unto day uttereth speech,’ and the tone of each holy anniversary is distinct and decisive. Thus the festival year is a bulwark of orthodoxy as real as our confession of faith.” History shows, however (especially that of Germany and France), that neither the church year nor creeds can prevent a fearful apostasy to rationalism and infidelity.} It serves to interweave religion with the, life of the people by continually recalling to the popular mind the most important events upon which our salvation rests, and by connecting them with the vicissitudes of the natural and the civil year. Yet, on the other hand, the gradual overloading of the church year, and the multiplication of saints’ days, greatly encouraged superstition and idleness, crowded the Sabbath and the leading festivals into the background, and subordinated the merits of Christ to the patronage of saints. The purification and simplification aimed at by the Reformation became an absolute necessity.

The order of the church year is founded in part upon the history of Jesus and of the apostolic church; in part, especially in respect to Easter and Pentecost, upon the Jewish sacred year; and in part upon the natural succession of seasons; for the life of nature in general forms the groundwork of the higher life of the spirit, and there is an evident symbolical correspondence between Easter and spring, Pentecost and the beginning of harvest, Christmas and the winter solstice, the nativity of John the Baptist and the summer solstice.

The Christian church year, however, developed itself spontaneously from the demands of the Christian worship and public life, after the precedent of the Old Testament cultus, with no positive direction from Christ or the apostles. The New Testament contains no certain traces of annual festivals; but so early as the second century we meet with the general observance of Easter and Pentecost, founded on the Jewish passover and feast of harvest, and answering to Friday and Sunday in the weekly cycle. Easter was a season of sorrow, in remembrance of the passion; Pentecost was a time of joy, in memory of the resurrection of the Redeemer and the outpouring of the Holy Ghost.\footnote{Comp. vol. i. § 99} These two festivals form the heart of the church year. Less important was the feast of the Epiphany, or manifestation of Christ as Messiah. In the fourth century the Christmas festival was added to the two former leading feasts, and partially took the place of the earlier feast of Epiphany, which now came to be devoted particularly to the manifestation of Christ among the Gentiles. And further, in Easter the πάσχα σταυρώσιμονand ἀναστάσιμονcame to be more strictly distinguished, the latter being reckoned a season of joy.

From this time, therefore, we have three great festival cycles, each including a season of preparation before the feast and an after-season appropriate: Christmas, Easter, and Pentecost. The lesser feasts of Epiphany and Ascension arranged themselves under these.\footnote{. There was no unanimity, however, in this period, in the number of the feasts. Chrysostom, for example, counts seven principal feasts, corresponding to the seven days of the week: Christmas, Epiphany, Passion, Easter, Ascension, Pentecost, and the Feast of the Resurrection of the Dead. The last, however, is not a strictly ecclesiastical feast, and the later Greeks reckon only six principal festivals, answering to the six days of creation, followed by the eternal Sabbath of the church triumphant in heaven. Comp. Augusti, i. p. 530,} All bear originally a christological character, representing the three stages of the redeeming work of Christ: the beginning, the prosecution, and the consummation. All are for the glorification of God in Christ.

The trinitarian conception and arrangement of the festal half of the church year is of much later origin, cotemporary with the introduction of the festival of the Trinity (on the Sunday after Pentecost). The feast of Trinity dates from the ninth or tenth century, and was first authoritatively established in the Latin church by Pope John XXII., in 1334, as a comprehensive closing celebration of the revelation of God the Father, who sent His Son (Christmas), of the Son, who died for us and rose again (Easter), and of the Holy Ghost, who renews and sanctifies us (Pentecost).\footnote{The assertion that the festum Trinitatis descends from the time of Gregory the Great, has poor foundation in his words: “Ut de Trinitate specialia cantaremus; for these refer to the praise of the holy Trinity in the general public worship of God. The first clear traces of this festival appear in the time of Charlemagne and in the tenth century, when Bishop Stephen of Liege vindicated it. Yet so late as 1150 it was counted by the abbot Potho at Treves among the \emph{novae} celebritates. Many considered it improper to celebrate a special feast of the Trinity, while there was no distinct celebration of the unity of God. The Roman church year reached its culmination and mysterious close in the feast of Corpus Christi (the body of Christ), which was introduced under Pope Clement the Fifth, in 1311, and was celebrated on Thursday of Trinity week (feria quinta proxima post octavam Pentecostes) in honor of the mystery of transubstantiation.} The Greek church knows nothing of this festival to this day, though she herself, in the Nicene age, was devoted with special earnestness and zeal to the development of the doctrine of the Trinity. The reason of this probably is, that there was no particular historical fact to give occasion for such celebration, and that the mystery of the holy Trinity, revealed in Christ, is properly the object of adoration in all the church festivals and in the whole Christian cultus.

But with these three great feast-cycles the ancient church was not satisfied. So early as the Nicene age it surrounded them with feasts of Mary, of the apostles, of martyrs, and of saints, which were at first only local commemorations, but gradually assumed the character of universal feasts of triumph. By degrees every day of the church year became sacred to the memory of a particular martyr or saint, and in every case was either really or by supposition the day of the death of the saint, which was significantly called his heavenly birth-day.\footnote{Hence called \emph{Natales, natalitia, nativitas}, \textgreek{γενέθλια}, of the martyrs. The Greek church also has its saint for every day of the year, but varies in many particulars from the Roman calendar.} This multiplication of festivals has at bottom the true thought, that the whole life of the Christian should be one unbroken spiritual festivity. But the Romish calendar of saints anticipates an ideal condition, and corrupts the truth by exaggeration, as the Pharisees made the word of God “of none effect” by their additions. It obliterates the necessary distinction between Sunday and the six days of labor, to the prejudice of the former, and plays into the hands of idleness. And finally, it rests in great part upon uncertain legends and fantastic myths, which in some cases even eclipse the miracles of the gospel history, and nourish the grossest superstition.

The Greek oriental church year differs from the Roman in this general characteristic: that it adheres more closely to the Jewish ceremonies and customs, while the Roman attaches itself to the natural year and common life. The former begins in the middle of September (Tisri), with the first Sunday after the feast of the Holy Cross; the latter, with the beginning of Advent, four weeks before Christmas. Originally Easter was the beginning of the church year, both in the East and in the West; and the Apostolic Constitutions and Eusebius call the month of Easter the “first month” (corresponding to the month Nisan, which opened the sacred year of the Jews, while the first of Tisri, about the middle of our September, opened their civil year). In the Greek church also the lectiones continuae of the Holy Scriptures, after the example of the Jewish Parashioth and Haphthoroth, became prominent and the church year came to be divided according to the four Evangelists; while in the Latin church, since the sixth century, only select sections from the Gospels, and Epistles, called pericopes, have been read. Another peculiarity of the Western church year, descending from the fourth century, is the division into four portions, of three months each, called Quatember,\footnote{Quatuor tempora.} separated from each other by a three days’ fast. Pope Leo I. delivered several sermons on the quarterly Quatember fast,\footnote{Sermones de jejunio quatuor temporum.} and urges especially on that occasion charity to the poor. Instead of this the Greek church has a division according to the four Gospels, which are read entire in course; Matthew next after Pentecost, Luke beginning on the fourteenth of September, Mark at the Easter fast, and John on the first Sunday after Easter.

So early as the fourth century the observance of the festivals was enjoined under ecclesiastical penalties, and was regarded as an established divine ordinance. But the most eminent church teachers, a Chrysostom, a Jerome, and an Augustine, expressly insist, that the observance of the Christian festivals must never be a work of legal constraint, but always an act of evangelical freedom; and Socrates, the historian, says, that Christ and the apostles have given no laws and prescribed no penalties concerning it.\footnote{Comp. the passages in Augusti, l.c. i. p. 474 sqq.}

The abuse of the festivals soon fastened itself on the just use of them and the sensual excesses of the pagan feasts, in spite of the earnest warnings of several fathers, swept in like a wild flood upon the church. Gregory Nazianzen feels called upon, with reference particularly to the feast of Epiphany, to caution his people against public parade, splendor of dress, banquetings, and drinking revels, and says: “Such things we will leave to the Greeks, who worship their gods with the belly; but we, who adore the eternal Word, will find our only satisfaction in the word and the divine law, and in the contemplation of the holy object of our feast.”\footnote{Orat. 38 in Theoph., cited at large by Augusti, p. 483 sq. Comp. Augustine, Ep. 22, 3; 29, 9, according to which “comessationes et ebrietates in honorem etiam beatissimorum martyrum” were of almost daily occurrence in the African church, and were leniently judged, lest the transition of the heathen should be discouraged.} On the other hand, however, the Catholic church, especially after Pope Gregory I. (the “pater caerimoniarum”), with a good, but mistaken intention, favored the christianizing of heathen forms of cultus and popular festivals, and thereby contributed unconsciously to the paganizing of Christianity in the Middle Age. The calendar saints took the place of the ancient deities, and Rome became a second time a pantheon. Against this new heathenism, with its sweeping abuses, pure Christianity was obliged with all earnestness and emphasis to protest.

Note. – The Reformation of the sixteenth century sought to restore the entire cultus, and with it the Catholic church year, to its primitive Biblical simplicity; but with different degrees of consistency. The Lutheran, the Anglican, and the German Reformed churches—the latter with the greater freedom—retained the chief festivals, Christmas, Easter, and Pentecost, together with the system of pericopes, and in some cases also the days of Mary and the apostles (though these are passing more and more out of use); while the strictly Calvinistic churches, particularly the Presbyterians and Congregationalists, rejected all the yearly festivals as human institutions, but, on the other hand, introduced a proportionally stricter observance of the weekly day of rest instituted by God Himself. The Scotch General Assembly of August 6th, 1575, resolved: “That all days which heretofore have been kept holy, besides the Sabbath-days, such as Yule day [Christmas], saints’ days, and such others, may be abolished, and a civil penalty be appointed against the keepers thereof by ceremonies, banqueting, fasting, and such other vanities.” At first, the most of the Reformers, even Luther and Bucer, were for the abolition of all feast days, except Sunday; but the genius and long habits of the people were against such a radical reform. After the end of the sixteenth and beginning of the seventeenth century the strict observance of Sunday developed itself in Great Britain and North America; while the Protestantism of the continent of Europe is much looser in this respect, and not essentially different from Catholicism. It is remarkable, that the strictest observance of Sunday is found just in those countries where the yearly feasts have entirely lost place in the popular mind: Scotland and New England. In the United States, however, for some years past, the Christmas and Easter festivals have regained ground without interfering at all with the strict observance of the Lord’s day, and promise to become regular American institutions. Good Friday and Pentecost will follow. On Good Friday of the year 1864 the leading ministers of the different evangelical churches in New York (the Episcopalian, Presbyterian, Dutch and German Reformed, Lutheran, Congregational, Methodist, and Baptist) freely united in the celebration of the atoning death of their common Saviour and in humiliation and prayer to the great edification of the people. It is acknowledged more and more that the observance of the great facts of the evangelical history to the honor of Christ is a common inheritance of primitive Christianity and inseparable from Christian worship.” These festivals” (says Prof. Dr. Henry B. Smith in his admirable opening sermon of the Presbyterian General Assembly, N. S., of 1864, on Christian Union and Ecclesiastical Re-union), “antedate, not only our (Protestant) divisions, but also the corruptions of the Papacy; they exalt the Lord and not man; they involve a public and solemn recognition of essential Christian facts, and are thus a standing protest against infidelity; they bring out the historic side of the Christian faith, and connect us with its whole history; and all in the different denominations could unite in their observance without sacrificing any article of their creed or discipline.” There is no danger that American Protestantism will transgress the limits of primitive evangelical simplicity in this respect, and ever return to the papal Mariolatry and Hagiolatry. The Protestant churches have established also many new annual festivals, such as the feasts of the Reformation, of Harvest-home, and of the Dead in Germany; and in America, the frequent days of fasting and prayer, besides the annual Thanksgiving-day, which originated in Puritan New England, and has been gradually adopted in almost all the states of the Union, and quite recently by the general government itself, as a national institution. With the pericopes, or Scripture lessons, the Reformed church everywhere deals much more freely than the Lutheran, and properly reserves the right to expound the whole word of Scripture in any convenient order according to its choice. The Gospels and Epistles may be read as a regular part of the Sabbath service; but the minister should be free to select his text from any portion of the Canonical Scriptures; only it is always advisable to follow a system and to go, if possible, every year through the whole plan and order of salvation in judicious adaptation to the church year and the wants of the people.

\xsection{The Christmas Cycle}

§ 77. The Christmas Cycle.

Besides the general literature given in the previous section, there are many special treatises on the origin of the Christmas festival, by Bynaeus, Kindler, Ittig, Vogel, Wernsdorf, Jablonsky, Planck, Hagenbach, P. Cassel, \&c. Comp. Augusti: Archaeol. i. 533.

The Christmas festival\footnote{\emph{Natalis}, or \emph{natalitia Domini} or \emph{Christi},\textgreek{ἡμέρα γενέθλιος, γενέθλια τοῦ Χριστοῦ}.} is the celebration of the incarnation of the Son of God. It is occupied, therefore, with the event which forms the centre and turning-point of the history of the world. It is of all the festivals the one most thoroughly interwoven with the popular and family life, and stands at the head of the great feasts in the Western church year. It continues to be, in the entire Catholic world and in the greater part of Protestant Christendom, the grand jubilee of children, on which innumerable gifts celebrate the infinite love of God in the gift of his only-begotten Son. It kindles in mid-winter a holy fire of love and gratitude, and preaches in the longest night the rising of the Sun of life and the glory of the Lord. It denotes the advent of the true golden age, of the freedom and equality of all the redeemed before God and in God. No one can measure the joy and blessing which from year to year flow forth upon all ages of life from the contemplation of the holy child Jesus in his heavenly innocence and divine humility.

Notwithstanding this deep significance and wide popularity, the festival of the birth of the Lord is of comparatively late institution. This may doubtless be accounted for in the following manner: In the first place, no corresponding festival was presented by the Old Testament, as in the case of Easter and Pentecost. In the second place, the day and month of the birth of Christ are nowhere stated in the gospel history, and cannot be certainly determined. Again: the church lingered first of all about the death and resurrection of Christ, the completed fact of redemption, and made this the centre of the weekly worship and the church year. Finally: the earlier feast of Epiphany afforded a substitute. The artistic religious impulse, however, which produced the whole church year, must sooner or later have called into existence a festival which forms the groundwork of all other annual festivals in honor of Christ. For, as Chrysostom, some ten years, after the introduction of this anniversary in Antioch, justly said, without the birth of Christ there were also no baptism, passion, resurrection, or ascension, and no outpouring of the Holy Ghost; hence no feast of Epiphany, of Easter, or of Pentecost.

The feast of Epiphany had spread from the East to the West. The feast of Christmas took the opposite course. We find it first in Rome, in the time of the bishop Liberius, who on the twenty-fifth of December, 360, consecrated Marcella, the sister of St. Ambrose, nun or bride of Christ, and addressed her with the words: “Thou seest what multitudes are come to the birth-festival of thy bridegroom.”\footnote{Ambrose, De virgin. iii. 1: “Vides quantus ad natalem Sponsi tui populus convenerit, ut nemo impastus recedit?} This passage implies that the festival was already existing and familiar. Christmas was introduced in Antioch about the year 380; in Alexandria, where the feast of Epiphany was celebrated as the nativity of Christ, not till about 430. Chrysostom, who delivered the Christmas homily in Antioch on the 25th of December, 386,\footnote{Opp. ii. 384.} already calls it, notwithstanding its recent introduction (some ten years before), the fundamental feast, or the root, from which all other Christian festivals grow forth.

The Christmas festival was probably the Christian transformation or regeneration of a series of kindred heathen festivals—the Saturnalia, Sigillaria, Juvenalia, and Brumalia—which were kept in Rome in the month of December, in commemoration of the golden age of universal freedom and equality, and in honor of the unconquered sun, and which were great holidays, especially for slaves and children.\footnote{The Satumalia were the feast of Saturn or Kronos, in representation of the golden days of his reign, when all labor ceased, prisoners were set free, slaves went about in gentlemen’s clothes and in the hat (the mark of a freeman), and all classes gave themselves up to mirth and rejoicing. The Sigillaria were a festival of images and puppets at the close of the Saturnalia on the 21st and 22d of December, when miniature images of the gods, wax tapers, and all sorts of articles of beauty and luxury were distributed to children and among kinsfolk. The Brumalia, from bruma (brevissima, the shortest day), had reference to the winter solstice, and the return of the Sol invictus.} This connection accounts for many customs of the Christmas season, like the giving of presents to children and to the poor, the lighting of wax tapers, perhaps also the erection of Christmas trees, and gives them a Christian import; while it also betrays the origin of the many excesses in which the unbelieving world indulges at this season, in wanton perversion of the true Christmas mirth, but which, of course, no more forbid right use, than the abuses of the Bible or of any other gift of God. Had the Christmas festival arisen in the period of the persecution, its derivation from these pagan festivals would be refuted by the then reigning abhorrence of everything heathen; but in the Nicene age this rigidness of opposition between the church and the world was in a great measure softened by the general conversion of the heathen. Besides, there lurked in those pagan festivals themselves, in spite of all their sensual abuses, a deep meaning and an adaptation to a real want; they might be called unconscious prophecies of the Christmas feast. Finally, the church fathers themselves\footnote{Chrysostom, Gregory of Nyassa, Leo the Great, and others.} confirm the symbolical reference of the feast of the birth of Christ, the Sun of righteousness, the Light of the world, to the birth-festival of the unconquered sun,\footnote{\emph{Dies} or \emph{natales invicti Solis}. This is the feast of the Persian sun-god Mithras, which was formally introduced in Rome under Domitian and Trajan.} which on the twenty-fifth of December, after the winter solstice, breaks the growing power of darkness, and begins anew his heroic career. It was at the same time, moreover, the prevailing opinion of the church in the fourth and fifth centuries, that Christ was actually born on the twenty-fifth of December; and Chrysostom appeals, in behalf of this view, to the date of the registration under Quirinius (Cyrenius), preserved in the Roman archives. But no certainly respecting the birthday of Christ can be reached from existing data.\footnote{In the early church, the 6th of January, the day of the Epiphany festival, was regarded by some as the birth-day of Christ. Among Biblical chronologists, Jerome, Baronius, Lamy, Usher, Petavius, Bengel, and Seyffarth, decide for the 25th of December, while Scaliger, Hug Wieseler, and Ellicott (Hist. Lectures on the Life of our Lord Jesus Christ, p. 70, note 3, Am. ed.), place the birth of Christ in the month of February. The passage in Luke, ii. 8, is frequently cited against the common view, because, according to the Talmudic writers, the flocks in Palestine were brought in at the beginning of November, and not driven to pasture again till toward March. Yet this rule, certainly, admitted many exceptions, according to the locality and the season. Comp. the extended discussion in Wieseler: Chronologische Synopse, p. 132 ff., and Seyffarth, Chronologia Sacra.}

Around the feast of Christmas other festivals gradually gathered, which compose, with it, the Christmas Cycle. The celebration of the twenty-fifth of December was preceded by the Christmas Vigils, or Christmas Night, which was spent with the greater solemnity, because Christ was certainly born in the night.\footnote{Luke ii. 8.}

After Gregory the Great the four Sundays before Christmas began to be devoted to the preparation for the coming of our Lord in the flesh and for his second coming to the final judgment. Hence they were called Advent Sundays. With the beginning of Advent the church year in the West began. The Greek church reckons six Advent Sundays, and begins them with the fourteenth of November. This Advent season was designed to represent and reproduce in the consciousness of the church at once the darkness and the yearning and hope of the long ages before Christ. Subsequently all noisy amusements and also weddings were forbidden during this season. The pericopes are selected with reference to the awakening of repentance and of desire after the Redeemer.

From the fourth century Christmas was followed by the memorial days of St. Stephen, the first Christian martyr (Dec. 26), of the apostle and evangelist John (Dec. 27), and of the Innocents of Bethlehem (Dec. 28), in immediate succession; representing a threefold martyrdom: martyrdom in will and in fact (Stephen), in will without the fact (John), and in fact without the will, an unconscious martyrdom of infanthe innocence. But Christian martyrdom in general was regarded by the early church as a heavenly birth and a fruit of the earthly birth of Christ. Hence the ancient festival hymn for the day of St. Stephen, the leader of the noble army of martyrs: “Yesterday was Christ born upon earth, that to-day Stephen might be born in heaven.”\footnote{“Heri natus est Christus in terris, ut hodie Stephanus nasceretur in coelis.” The connection is, however, a purely ideal one; for at first the death-day of Stephen was in August; afterward, on account of the discovery of his relics, it was transferred to January.} The close connection of the feast of John the, Evangelist with that of the birth of Christ arises from the confidential relation of the beloved disciple to the Lord, and from the fundamental thought of his Gospel: “The Word was made flesh.” The innocent infant-martyrs of Bethlehem, “the blossoms of martyrdom, the rosebuds torn off by the hurricane of persecution, the offering of first-fruits to Christ, the tender flock of sacrificial lambs,” are at the same time the representatives of the innumerable host of children in heaven.\footnote{Comp. the beautiful hymn of the Spanish poet Prudentius, of the fifth century: “Salvete flores martyrum.” German versions by Nickel, Königsfeld, Bässler, Hagenbach, \&c. A good English version in “The Words of the Hymnal Noted, ” Lond, p. 45: \par “All hail! ye Infant-Martyr flowers, \par Cut off in life’s first dawning hours: \par As rosebuds, snapt in tempest strife, \par When Herod sought your Saviour’s life,” \&c.} More than half of the human race are said to die in infancy, and yet to children the word emphatically applies: “Theirs is the kingdom of heaven.” The mystery of infant martyrdom is constantly repeated. How many children are apparently only born to suffer, and to die; but in truth the pains of their earthly birth are soon absorbed by the joys of their heavenly birth, and their temporary cross is rewarded by an eternal crown.

Eight days after Christmas the church celebrated, though not till after the sixth or seventh century, the Circumcision and the Naming of Jesus. Of still later origin is the Christian New Year’s festival, which falls on the same day as the Circumcision. The pagan Romans solemnized the turn of the year, like the Saturnalia, with revels. The church teachers, in reaction, made the New Year a day of penance and prayer. Thus Augustine, in a sermon: “Separate yourselves from the heathen, and at the change of the year do the opposite of what they do. They give each other gifts; give ye alms instead. They sing worldly songs; read ye the word of God. They throng the theatre come ye to the church. They drink themselves drunken; do ye fast.”

The feast of Epiphany\footnote{τὰ ἐπιφάνεια, or \textgreek{ἐπιφανία, Χριστ φανία}, also \textgreek{θεοφανία}. Comp. vol i. § 99.} on the contrary, on the sixth of January, is older, as we have already observed, than Christmas itself, and is mentioned by Clement of Alexandria. It refers in general to the manifestation of Christ in the world, and originally bore the twofold character of a celebration of the birth and the baptism of Jesus. After the introduction of Christmas, it lost its reference to the birth. The Eastern church commemorated on this day especially the baptism of Christ, or the manifestation of His Messiahship, and together with this the first manifestation of His miraculous power at the marriage at Cana. The Westem church, more Gentile-Christian in its origin, gave this festival, after the fourth century, a special reference to the adoration of the infant Jesus by the wise men from the east,\footnote{Matt. ii. 1-11.} under the name of the feast of the Three Kings, and transformed it into a festival of Gentile missions; considering the wise men as the representatives of the nobler heathen world.\footnote{Augustine, Sermo 203: “Hodierno die manifestatus redemptor omnium gentium,” \&c. The transformation of the Persian magi or priest-philosophers into three kings (Caspar, Melchior, and Balthasar) by the mediaeval legend was a hasty inference from the triplicity of the gifts and from Ps. lxxii. 10, 11. The legend brings us at last to the cathedral at Cologne, where the bodies of the three saint-kings are to this day exhibited and worshipped.} Thus at the same time the original connection of the feast with the birth of Christ was preserved. Epiphany forms the close of the Christmas Cycle. It was an early custom to announce the term of the Easter observance on the day of Epiphany by the so-called Epistolae paschales, or gravmmata pascavlia. This was done especially by the bishop of Alexandria, where astronomy most flourished, and the occasion was improved for edifying instructions and for the discussion of important religious questions of the day.

\xsection{The Easter Cycle}

§ 78. The Easter Cycle.

Easter is the oldest and greatest annual festival of the church. As to its essential idea and observance, it was born with the Christian Sunday on the morning of the resurrection.\footnote{The late Dr. Fried. Strauss of Berlin, an eminent writer on the church year (Das evangelische Kirchenjahr, p. 218), says: ”Das heilige Osterfest ist das christliche Fest schlechthin. Es ist nicht blos Hauptfest, sondern das Fest, das einmal im Jahre vollstandig auftritt, aber in allen andern Festen von irgend einer Seite wiederkehrt, und eben dadurch diese zu Festen macht. Nannte man doch jeden Festtag, ja sogar jeden Sonntag aus diesem Grunde \emph{dies paschalis}. Daher musste es auch das ursprüngliche Fest in dem umfassendsten Sinne des Wortes sein. Man kann nicht sagen, in welcher christlichen Zeit es entstanden sei; es ist mit der Kirche entstanden, und die Kirche ist mit ihm entstanden.”} Like the passover with the Jews, it originally marked the beginning of the church year. It revolves entirely about the person and the work of Christ, being devoted to the great saving fact of his passion and resurrection. We have already spoken of the origin and character of this festival,\footnote{Vol. i. § 99 (p. 373 ff.).} and shall confine ourselves here to the alterations and enlargements which it underwent after the Nicene age.

The Easter festival proper was preceded by a forty days’ season of repentance and fasting, called Quadragesima, at least as early as the year 325; for the council of Nice presupposes the existence of this season.\footnote{In its fifth canon, where it orders that provincial councils be held twice a year, before \emph{Quadragesima} (\textgreek{πρὸ τῆς τεσσαρακοστῆς}), and in the autumn.} This fast was an imitation of the forty days’ fasting of Jesus in the wilderness, which itself was put in typical connection with the forty days’ fasting of Moses\footnote{Ex. xxxiv. 28.} and Elijah,\footnote{1 Kings xix. 8.} and the forty years’ wandering of Israel through the desert. At first a free-will act, it gradually assumed the character of a fixed custom and ordinance of the church. Respecting the length of the season much difference prevailed, until Gregory I. (590–604) fixed the Wednesday of the sixth week before Easter, Ash Wednesday as it is called,\footnote{\emph{Dies cinerum, caput jejunii}, or \emph{quadragesimae}.} as the beginning of it. On this day the priests and the people sprinkled themselves with dust and ashes, in token of their perishableness and their repentance, with the words: “Remember, O man, that dust thou art, and unto dust thou must return; repent, that thou mayest inherit eternal life.” During Quadragesima criminal trials and criminal punishments, weddings, and sensual amusements were forbidden; solemn, earnest silence was imposed upon public and private life; and works of devotion, penances and charity were multiplied. Yet much hypocrisy was practised in the fasting; the rich compensating with exquisite dainties the absence of forbidden meats. Chrysostom and Augustine are found already lamenting this abuse. During the days preceding the beginning of Lent, the populace gave themselves up to unrestrained merriment, and this abuse afterward became legitimized in all Catholic countries, especially in Italy (flourishing most in Rome, Venice, and Cologne), in the Carnival.\footnote{From \emph{caro} and \emph{vale}; flesh taking its departure for a time in a jubilee of revelling. According to others, it is the converse: dies quo caro \emph{valet}; \emph{i.e}., the day on which it is still allowed to eat flesh and to indulge the flesh. The Carnival, or Shrove-tide, embraces the time from the feast of Epiphany to Ash Wednesday, or, commonly, only the last three or the last eight days preceding Lent. It is celebrated in every city of Italy; in Rome, especially, with masquerades, races, dramatic plays, farces, jokes, and other forms of wild merriment and frantic joy, yet with good humor; replacing the old Roman feasts of Saturnalia, Lupercalia, and Floralia}

The six Sundays of Lent are called Quadragesima prima, secunda, and so on to sexta. They are also named after the initial words of the introit in the mass for the day: Invocabit (Ps. xci. 15), Reminiscere, (Ps. xxv. 6), Oculi (Ps. xxxiv. 15), Laetare (Is. lxvi. 10), Judica (Ps. xliii. 1), Palmarum (from Matt. xxi. 8). The three Sundays preceding Quadragesima are called respectively Estomihi (from Ps. xxxi. 2) or Quinquagesima (i.e., Dominica quinquagesimae diei, viz., before Easter), Sexagesima, and Septuagesima; which are, however, inaccurate designations. These three Sundays were regarded as preparatory to the Lenten season proper. In the larger cities it became customary to preach daily during the Quadragesimal fast; and the usage of daily Lenten sermons (Quadragesimales, or sermones Quadragesimales) has maintained itself in the Roman church to this day.

The Quadragesimal fast culminates in the Great, or Silent, or Holy Week,\footnote{\emph{Septimana sancta, magna, muta; hebdomas nigra}, or \emph{paschalis}; \textgreek{ἑβδομὰς μεγάλη}, \emph{Passion Week}.} which is especially devoted to the commemoration of the passion and death of Jesus, and is distinguished by daily public worship, rigid fasting, and deep silence. This week, again, has its prominent days. First Palm Sunday,\footnote{\emph{Dominica palmarum}; \textgreek{ἑορτὴ τῶν βαίων}.} which has been, in the East since the fourth century, in the West since the sixth, observed in memory of the entry of Jesus into Jerusalem for His enthronement on the cross. Next follows Maundy Thursday,\footnote{Feria quinta paschae, dies natalis eucharistiae, dies viridium; \textgreek{ἡμεγάληπέμπτη}. The English name, Maundy Thursday, is derived from \emph{maunds}or baskets, in which on that day the king of England distributed alms to certain poor at Whitehall, \emph{Maund} is connected with the Latin \emph{mendicare}, and French \emph{mendier}, to beg.} in commemoration of the institution of the Holy Supper, which on this day was observed in the evening, and was usually connected with a love feast, and also with feet-washing. The Friday of the Holy Week is distinguished from all others as Good Friday,\footnote{\emph{Dies dominicae passionis}; \textgreek{παρασκευή, πάσχα σταυρώσιμον, ἡμέρα τοῦ σταυροῦ.} In German \emph{Char-Freitag} either from the Greek \textgreek{χάρις}, or, more probably, from the Latin \emph{carus}, \emph{beloved, dear}, comp. the English \emph{Good} Friday. Other etymologists derive it from \emph{carena} (\emph{carême}), \emph{i.e., fasting}, or from \emph{kar} (\emph{küren, to} \emph{choose}), \emph{i.e., the chosen day}; others still from \emph{karo-parare, i.e., preparation-day}.} the day of the Saviour’s death; the day of the deepest penance and fasting of the year, stripped of all Sunday splendor and liturgical pomp, veiled in the deepest silence and holy sorrow; the communion omitted (which had taken place the evening before), altars unclothed, crucifixes veiled, lights extinguished, the story of the passion read, and, instead of the church hymns, nothing sung but penitential psalms. Finally the Great Sabbath,\footnote{\textgreek{Μέγα} or \textgreek{ἅγιον σάββατον}; \emph{sabbatum magnum}, or \emph{sanctum}.} the day of the Lord’s repose in the grave and descent into Hades; the favorite day in all the year for the administration of baptism, which symbolizes participation in the death of Christ.\footnote{Rom. vi. 4-6.} The Great Sabbath was generally spent as a fast day, even in the Greek church, which usually did not fast on Saturday.

In the evening of the Great Sabbath began the Easter Vigils,\footnote{\emph{Vigiliae paschales}; \textgreek{παννυχίδες.}} which continued, with Scripture reading, singing, and prayer, to the dawn of Easter morning, and formed the solemn transition from the πάσχα σταυρώσιμον to the πάσχα ἀναστάσιμον, and from the deep sorrow of penitence over the death of Jesus to the joy of faith in the resurrection of the Prince of life. All Christians, and even many pagans, poured into the church with lights, to watch there for the morning of the resurrection. On this night the cities were splendidly illuminated, and transfigured in a sea of fire; about midnight a solemn procession surrounded the church, and then triumphally entered again into the “holy gates,” to celebrate Easter. According to an ancient tradition, it was expected that on Easter night Christ would come again to judge the world.\footnote{Comp. Lactantius: Inst. divin. vii. c. 19; and Hieronymus ad Matt. xxv. 6 (t. vii. 203, ed. Vallarsi): “Unde traditionem apostolicam permansisse, ut in die vigiliarum Paschae ante noctis dimidium populos dimittere non liceat, \emph{expectantes adventum Christi}.”}

The Easter festival itself\footnote{\emph{Festum dominicae resurrectionis}; \textgreek{ἑορτὴἀναστάσιμος, κυριακὴμεγάλη}.} began with the jubilant salutation, still practized in the Russian church: “The Lord is risen !” and the response: “He is truly risen!\footnote{“Dominus resurrexit.”—“Vere resurrexit.”} Then the holy kiss of brotherhood scaled the newly fastened bond of love in Christ. It was the grandest and most joyful of the feasts. It lasted a whole week, and closed with the following Sunday, called the Easter Octave,\footnote{\emph{Octava paschae, pascha clausum}; \textgreek{ἀντίπασχα}. \emph{Octave}is applied in genera to the whole eight-days’ observance of the great church festivals; then especially to the eighth or last day of the feast.} or White Sunday,\footnote{\emph{Dominica in albis}. Also \emph{Quasimodogeniti}, from the Introit for public worship, 1 Pet. ii. 2 (“Quasimodo geniti infantes,”” As new-born babes,” \&c.). Among the Greeks it was called \textgreek{καινὴ κυριακή}.} when the baptized appeared in white garments, and were solemnly incorporated into the church.

\xsection{The Time of the Easter Festival}

§ 79. The Time of the Easter Festival.

Comp. the Literature in vol. i. at § 99; also L. Ideler: Handbuch der Chronologie. Berlin, 1826. Vol. ii. F. Piper: Geschichte des Osterfestes. Berlin, 1845. Hefele: Conciliengeschichte. Freiburg, 1855. Vol. i. p. 286 ff.

The time of the Easter festival became, after the second century, the subject of long and violent controversies and practical confusions, which remind us of the later Eucharistic disputes, and give evidence that human passion and folly have sought to pervert the great facts and institutions of the New Testament from holy bonds of unity into torches of discord, and to turn the sweetest honey into poison, but, with all their efforts, have not been able to destroy the beneficent power of those gifts of God.

These Paschal controversies descended into the present period, and ended with the victory of the Roman and Alexandrian practice of keeping Easter, not, like Christmas and the Jewish Passover, on a fixed day of the month, whatever day of the week it might be, but on a Sunday, as the day of the resurrection of our Lord. Easter thus became, with all the feasts depending on it, a movable feast; and then the different reckonings of the calendar led to many inconveniences and confusions. The exact determination of Easter Sunday is made from the first full moon after the vernal equinox; so that the day may fall on any Sunday between the 22d day of March and the 25th of April.

The council of Arles in 314 had already decreed, in its first canon, that the Christian Passover be celebrated “uno die et uno tempore per omnem orbem,” and that the bishops of Rome should fix the time. But as this order was not universally obeyed, the fathers of Nicaea proposed to settle the matter, and this was the second main object of the first ecumenical council in 325. The result of the transactions on this point, the particulars of which are not known to us, does not appear in the canons (probably out of consideration for the numerous Quartodecimanians), but is doubtless preserved in the two circular letters of the council itself and the emperor Constantine.\footnote{Socrates: Hist. Eccl. i. 9; Theodoret: H. E. i. 10; Eusebius: Vita Const ii. 17. Comp. Hefele, l.c. i. p. 309 sqq.} The feast of the resurrection was thenceforth required to be celebrated everywhere on a Sunday, and never on the day of the Jewish passover, but always after the fourteenth of Nisan, on the Sunday after the first vernal full moon. The leading motive for this regulation was opposition to Judaism, which had dishonored the passover by the crucifixion of the Lord.” We would,” says the circular letter of Constantine in reference to the council of Nice, “we would have nothing in common with that most hostile people, the Jews; for we have received from the Redeemer another way of honoring God [the order of the days of the week], and harmoniously adopting this method, we would withdraw ourselves from the evil fellowship of the Jews. For what they pompously assert, is really utterly absurd: that we cannot keep this feast at all without their instruction .... It is our duty to have nothing in common with the murderers of our Lord.” This bitter tone against Judaism runs through the whole letter.

At Nicaea, therefore, the Roman and Alexandrian usage with respect to Easter triumphed, and the Judaizing practice of the Quartodecimanians, who always celebrated Easter on the fourteenth of Nisan, became thenceforth a heresy. Yet that practice continued in many parts of the East, and in the time of Epiphanius, about a.d. 400, there were many, Quartodecimanians, who, as he says, were orthodox, indeed, in doctrine, but in ritual were addicted to Jewish fables, and built upon the principle: “Cursed is every one who does not keep his passover on the fourteenth of Nisan.”\footnote{Epiphanius, Haer. l.c. 1. Comp. Ex. xii. 15.} They kept the day with the Communion and with fasting till three o’clock. Yet they were divided into several parties among themselves. A peculiar offshoot of the Quartodecimanians was the rigidly ascetic Audians, who likewise held that the passover must be kept at the very same time (not after the same manner) with the Jews, on the fourteenth of Nisan, and for their authority appealed to their edition of the Apostolic Constitutions.

And even in the orthodox church these measures did not secure entire uniformity. For the council of Nicaea, probably from prudence, passed by the question of the Roman and Alexandrian computation of Easter. At least the Acts contain no reference to it.\footnote{Hefele thinks, however (i. p. 313 f.), from an expression of Cyril of Alexandria and Leo I., that the Nicaenum (1) gave the Alexandrian reckoning the preference over the Roman; (2) committed to Alexandria the reckoning, to Rome the announcing, of the Easter term; but that this order was not duly observed.} At all events this difference remained: that Rome, afterward as before, fixed the vernal equinox, the terminus a quo of the Easter full moon, on the 18th of March, while Alexandria placed it correctly on the 21st. It thus occurred, that the Latins, the very year after the Nicene council, and again in the years 330, 333, 340, 341, 343, varied from the Alexandrians in the time of keeping Easter. On this account the council of Sardica, as we learn from the recently discovered Paschal Epistles of Athanasius, took the Easter question again in hand, and brought about, by mutual concessions, a compromise for the ensuing fifty years, but without permanent result. In 387 the difference of the Egyptian and the Roman Easter amounted to fully five weeks. Later attempts also to adjust the matter were in vain, until the monk Dionysius Exiguus, the author of our Christian calendar, succeeded in harmonizing the computation of Easter on the basis of the true Alexandrian reckoning; except that the Gallican and British Christians adhered still longer to the old custom, and thus fell into conflict with the Anglo-Saxon. The introduction of the improved Gregorian calendar in the Western church in 1582 again produced discrepancy; the Eastern and Russian church adhered to the Julian calendar, and is consequently now about twelve days behind us. According to the Gregorian calendar, which does not divide the months with astronomical exactness, it sometimes happens that the Paschal full moon is put a couple of hours too early, and the Christian Easter, as was the case in 1825, coincides with the Jewish Passover, against the express order of the council of Nicaea.

\xsection{The Cycle of Pentecost}

§ 80. The Cycle of Pentecost.

The whole period of seven weeks from Easter to Pentecost bore a joyous, festal character. It was called Quinquagesima, or Pentecost in the wider sense,\footnote{\textgreek{Πεντεκοστή}. Comp. the author’s Hist. of the Apost. Ch. § 54.} and was the memorial of the exaltation of Christ at the right hand of the Father, His repeated appearances during the mysterious forty days, and His heavenly headship and eternal presence in the church. It was regarded as a continuous Sunday, and distinguished by the absence of all fasting and by standing in prayer. Quinquagesima formed a marked contrast with the Quadragesima which preceded. The deeper the sorrow of repentance had been in view of the suffering and dying Saviour, the higher now rose the joy of faith in the risen and eternally living Redeemer. This joy, of course, must keep itself clear of worldly amusements, and be sanctified by devotion, prayer, singing, and thanksgiving; and the theatres, therefore, remained closed through the fifty days. But the multitude of nominal Christians soon forgot their religious impressions, and sought to compensate their previous fasting with wanton merry-making.

The seven Sundays after Easter are called in the Latin church, respectively, Quasimodo-geniti, Misericordia Domini, Jubilate, Cantate, Rogate, (or, Vocem jucunditatis), Exaudi, and Pentecoste. In the Eastern church the Acts of the Apostles are read at this season.

Of the fifty festival days, the fortieth and the fiftieth were particularly prominent. The fortieth day after Easter, always a Thursday, was after the fourth century dedicated to the exaltation of Christ at the right hand of God, and hence named Ascension Day.\footnote{\emph{Dies ascensionis}; \textgreek{ἑορτὴ τῆς ἀναλήψεως.}} The fiftieth day, or the feast of Pentecost in the stricter sense,\footnote{\emph{Dies pentecostes}; \textgreek{πεντεκοστή, ἡμέρα τοῦ Πνεύματος}.} was the kernel and culminating point of this festival season, as Easter day was of the Easter cycle. It was the feast of the Holy Ghost, who on this day was poured out upon the assembled disciples with the whole fulness of the accomplished redemption; and it was at the same time the birth-day of the Christian church. Hence this festival also was particularly prized for baptisms and ordinations. Pentecost corresponded to the Jewish feast of that name, which was primarily the feast of first-fruits, and afterward became also the feast of the giving of the law on Sinai, and in this twofold import was fulfilled in the outpouring of the Holy Ghost and the founding of the Christian church.” Both revelations of the divine law,” writes Jerome to Fabiola, “took place on the fiftieth day after the passover; the one on Sinai, the other on Zion; there the mountain was shaken, here the temple; there, amid flames and lightnings, the tempest roared and the thunder rolled, here, also with mighty wind, appeared tongues of fire; there the sound of the trumpet pealed forth the words of the law, here the cornet of the gospel sounded through the mouth of the apostles.”

The celebration of Pentecost lasted, at least ultimately, three days or a whole week, closing with the Pentecostal Octave, which in the Greek church (so early as Chrysostom) was called The Feast of all Saints and Martyrs,\footnote{\textgreek{κυριακὴ τῶν ἁγίων πάντων μαρτυρησάντων}. The Western church kept a similar feast on the first of November, but not till the eighth century.} because the martyrs are the seed and the beauty of the church. The Latin church, on the contrary, though not till the tenth century, dedicated the Sunday after Pentecost to the Holy Trinity, and in the later times of the Middle Age, further added to the festival part of the church year the feast of Corpus Christi, in celebration of the mystery of transubstantiation, on the Thursday after Trinity. It thus invested the close of the church year with a purely dogmatic import. Protestantism has retained the feast of Trinity, in opposition to the Antitrinitarians; but has, of course, rejected the feast of Corpus Christi.

In the early church, Pentecost was the last great festival of the Christian year. Hence the Sundays following it, till Advent, were counted from Whitsunday.\footnote{So in the Roman church even after the introduction of the Trinity festival. The Protestants, on the contrary, as far as they retained the ecclesiastical calendar (Lutherans, Anglicans, \&c.), make the first Sunday \emph{after} Pentecost the basis, and count the First, Second, Third Sunday \emph{after Trinity}, instead of the First, Second, etc., Sunday after \emph{Whitsunday}.} The number of the Sundays in the second half of the church year therefore varies between twenty-seven and twenty-two, according to the time of Easter. In this part of the year we find even in the old lectionaries and sacramentaries some subordinate, feasts in memory of great men of the church; such as the feast of St. Peter and St. Paul, the founders of the church (June 29); the feast of the chief martyr, Laurentius, the representative of the church militant (August, 10); the feast of the archangel Michael, the representative of the church triumphant (September 29).

\xsection{The Exaltation of the Virgin Mariology}

§ 81. The Exaltation of the Virgin Mariology.

Canisius (R.C.): De Maria Virgine libri quinque. Ingolst. 1577. Lamberertini (R.C.): Comment. dum De J. Christi, matrisque ejus festis. Patav. 1751. Perrone (R.C.): De Immaculata B. V. Mariae conceptu. Rom. 1848. (In defence of the new papal dogma of the sinless conception of Mary.) F. W. Genthe: Die Jungfrau Maria, ihre Evangelien u. ihre Wunder. Halle, 1852. Comp. also the elaborate article, “Maria, Mutter des Herrn,” by Steitz, in Herzog’s Protest. Real-Encycl. (vol. ix. p. 74 ff.), and the article, “Maria, die heil. Jungfrau,” by Reithmayr (R.C.) in Wetzer u. Welte’s Kathol. Kirchenlex. (vi. 835 ff.); also the Eirenicon-controversy between Pusey and J. H. Newman, 1866.

Into these festival cycles a multitude of subordinate feasts found their way, at the head of which stand the festivals of the holy Virgin Mary, honored as queen of the army of saints.

The worship of Mary was originally only a reflection of the worship of Christ, and the feasts of Mary were designed to contribute to the glorifying of Christ. The system arose from the inner connection of the Virgin with the holy mystery of the Incarnation of the Son of God; though certainly, with this leading religious and theological interest other motives combined. As mother of the Saviour of the world, the Virgin Mary unquestionably holds forever a peculiar position among all women, and in the history of redemption. Even in heaven she must stand peculiarly near to Him whom on earth she bore nine months under her bosom, and whom she followed with true motherly care to the cross. It is perfectly natural, nay, essential, to sound religious feeling, to associate with Mary the fairest traits of maidenly and maternal character, and to revere her as the highest model of female purity, love, and piety. From her example issues a silent blessing upon all generations, and her name and memory are, and ever will be, inseparable from the holiest mysteries and benefits of faith. For this reason her name is even wrought into the Apostles’ Creed, in the simple and chaste words: “Conceived by the Holy Ghost, born of the Virgin Mary.”

The Catholic church, however, both Latin and Greek, did not stop with this. After the middle of the fourth century it overstepped the wholesome Biblical limit, and transformed the mother of the Lord”\footnote{\textgreek{Ἡ μήτηρ τοῦ κυρίου}, Luke i. 43.} into a mother of God, the humble handmaid of the Lord”\footnote{\textgreek{Ἡ δούλη κυρίου}, Luke i. 38.} into a queen of heaven, the “highly favored”\footnote{\textgreek{Κεχαριτωμένη} (pass. part.), Luke i. 28.} into a dispenser of favors, the “blessed among women”\footnote{\textgreek{Εὐλογημένη ἐν γυναιξίν}, Luke i. 28.} into an intercessor above all women, nay, we may almost say, the redeemed daughter of fallen Adam, who is nowhere in Holy Scripture excepted from the universal sinfulness, into a sinlessly holy co-redeemer. At first she was acquitted only of actual sin, afterward even of original; though the doctrine of the immaculate conception of the Virgin was long contested, and was not established as an article of faith in the Roman church till 1854. Thus the veneration of Mary gradually degenerated into the worship of Mary; and this took so deep hold upon the popular religious life in the Middle Age, that, in spite of all scholastic distinctions between latria, and dulia, and hyrerdulia, Mariolatry practically prevailed over the worship of Christ. Hence in the innumerable Madonnas of Catholic art the human mother is the principal figure, and the divine child accessory. The Romish devotions scarcely utter a Pater Noster without an Ave Maria, and turn even more frequently and naturally to the compassionate, tender-hearted mother for her intercessions, than to the eternal Son of God, thinking that in this indirect way the desired gift is more sure to be obtained. To this day the worship of Mary is one of the principal points of separation between the Graeco-Roman Catholicism and Evangelical Protestantism. It is one of the strongest expressions of the fundamental Romish error of unduly exalting the human factors or instruments of redemption, and obstructing, or rendering needless, the immediate access of believers to Christ, by thrusting in subordinate mediators. Nor can we but agree with nearly all unbiased historians in regarding the worship of Mary as an echo of ancient heathenism. It brings plainly to mind the worship of Ceres, of Isis, and of other ancient mothers of the gods; as the worship of saints and angels recalls the hero-worship of Greece and Rome. Polytheism was so deeply rooted among the people, that it reproduced itself in Christian forms. The popular religious want had accustomed itself even to female deities, and very naturally betook itself first of all to Mary, the highly favored and blessed mother of the divine-human Redeemer, as the worthiest object of adoration.

Let us trace now the main features in the historical development of the Catholic Mariology and Mariolatry.

The New Testament contains no intimation of any worship or festival celebration of Mary. On the one hand, Mary, is rightly called by Elizabeth, under the influence of the Holy Ghost, “the mother of the Lord”\footnote{Luke i. 43:\textgreek{Ἡ μήτηρ τοῦ κυρίου μου}.}—but nowhere “the mother of God,” which is at least not entirely synonymous—and is saluted by her, as well as by the angel Gabriel, as “blessed among women;”\footnote{Luke i. 28: \textgreek{Χαῖρε, κεχαριτωμένη·ὁκύριος μετὰσοῦ, εὐλογημένησὺἐνγυναιξίν}. So Elizabeth, Luke i. 42: \textgreek{Εὐλογημένη σὺ ἐν γυναιξί, καὶ εὐλογημένος ὁ καρπὸς τῆς κοιλίας σου}.} nay, she herself prophesies in her inspired song, which has since resounded through all ages of the church, that “henceforth all generations shall call me blessed.”\footnote{Luke i. 48: Ἀ\textgreek{πὸ τοῦ νῦν μακαριοῦσί με πᾶσαι αἱ γενεαί}.} Through all the youth of Jesus she appears as a devout virgin, full of childlike innocence, purity, and humility; and the few traces we have of her later life, especially the touching scene at the cross,\footnote{John xix. 25-27.} confirm this impression. But, on the other hand, it is equally unquestionable, that she is nowhere in the New Testament excepted from the universal sinfulness and the universal need of redemption, and represented as immaculately holy, or as in any way an object of divine veneration. On the contrary, true to the genuine female character, she modestly stands back throughout the gospel history, and in the Acts and the Epistles she is mentioned barely once, and then simply as the “mother of Jesus;”\footnote{Acts i. 14.} even her birth and her death are unknown. Her glory fades in holy humility before the higher glory of her Son. In truth, there are plain indications that the Lord, with prophetic reference to the future apotheosis of His mother according to the flesh, from the first gave warning against it. At the wedding in Cana He administered to her, though leniently and respectfully, a rebuke for premature zeal mingled perhaps with maternal vanity.\footnote{John ii. 4: \textgreek{Τί ἐμοὶ καὶ σοἰ, γύναι}; Comp. the commentators on the passage. The expression ”\emph{woman}“ is entirely respectful, comp. John xix. 21; xx. 13, 15. But the ”\emph{What have I to do with thee}?” is, like the Hebrew\texthebrew{ךְלָיָ ילּ־המַ }(Josh. xxii, 24; 2 Sam.xvi. 10; xix. 22; 1 Kings xvii 18; 2 Kings iii. 13; 2 Chron. xxxv. 21), a rebuke and censure of undue interference; comp. Matt. viii. 29; Luke viii. 28; Mark i. 24 (also the classics). Meyer, the best grammatical expositor, observes on \textgreek{γύναι}: “That Jesus did not say \textgreek{μῆτερ}, flowed involuntarily from the sense of His higher wonder-working position, whence He repelled the interference of feminine weakness, which here met Him even in His mother.”} On a subsequent occasion he put her on a level with other female disciples, and made the carnal consanguinity subordinate to the spiritual kinship of the doing of the will of God.\footnote{Matt. xii. 46-50.} The well-meant and in itself quite innocent benediction of an unknown woman upon His mother He did not indeed censure, but He corrected it with a benediction upon all who hear the word of God and keep it, and thus forestalled the deification of Mary by confining the ascription within the bounds of moderation.\footnote{Luke xi 27, 28. The \textgreek{μενοῦνγε} is emphatic, \emph{utique}, but also corrective, \emph{imo vero}; so here, and Rom. ii. 20; x. 18. Luther inexactly translates simply, \emph{ja}; the English Bible more correctly, \emph{yea rather}. Meyer \emph{ad loc}.: “Jesus does not forbid the congratulation of His mother, but He applies the predicate \textgreek{μακάριος} as the woman had done, to an outward relation, but to an \emph{ethical} category, in which \emph{any} one \emph{might} stand, so that the congratulation of His mother \emph{as} \emph{mother} is thereby corrected.” Van Oosterzee strikingly remarks in his Commentary on Luke (in Lange’s \emph{Bibelwerk}): ” The congratulating woman is the prototype of all those, who in all times have honored the mother of the Lord above her Son, and been guilty of Mariolatry. If the Lord even here disapproves this honoring of His mother, where it moves in so modest limits, what judgment would He pass upon the new dogma of Pio Nono, on which a whole new Mariology is built?”}

In striking contrast with this healthful and sober representation of Mary in the canonical Gospels are the numerous apocryphal Gospels of the third and fourth centuries, which decorated the life of Mary with fantastic fables and wonders of every kind, and thus furnished a pseudo-historical foundation for an unscriptural Mariology and Mariolatry.\footnote{Here belongs, above all, the Protevangelium Jacobi Minoris, which dates from the third or fourth century; then the Evangelium de nativitate S. Mariae; the Historia de nativitate Mariae et de infantia Salvatoris; the Evangelium infantiae Servatoris; the Evang. Josephi fabri lignarii. Comp. Thilo’s Cod. Apocryphus N. Ti. Lips. 1832, and the convenient digest of this apocryphal history in R. Hofmann’s Leben Jesu nach den Apocryphen. Leipz. 1851, pp. 5-117.} The Catholic church, it is true, condemned this apocryphal literature so early as the Decrees of Gelasius;\footnote{Decret de libris apocr. Coll Cone. ap. Harduin, tom. ii. p. 941. Comp. Pope Innocent I., Ep. ad Exuperium Tolosanum, c. 7, where the Protevang. Jacobi is rejected and condemned.} yet many of the fabulous elements of it—such as the names of the parents of Mary, Joachim (instead of Eli, as in Luke iii. 23) and Anna,\footnote{Epiphanius also, Haer. 78, no. 17, gives the parents of Jesus these names. To reconcile this with Luke iii: 23, the Roman theologians suppose, that Eli, or Heli, is an abbreviation of Heliakim, and that this is the same with Joakim, or Joachim.} the birth of Mary in a cave, her education in the temple, and her mock marriage with the aged Joseph\footnote{According to the apocryphal Historia Josephi he was already ninety years old; according to Epiphanius at least eighty; and was blessed with children by a former marriage. According to Origen, also, and Eusebius, and Gregory of Nyssa, Joseph was an aged widower. Jerome, on the contrary, makes him, like Mary, a pure \emph{coelebs}, and says of him: “Mariae quam putatus est habuisse, custos potius fait quam maritus;” consequently he must “virginem mansisse cum Maria, qui pater Domini meruit adpellari.” Contr. Helvid. c. 19.}—passed into the Catholic tradition.

The development of the orthodox Catholic Mariology and Mariolatry originated as early as the second century in an allegorical interpretation of the history of the fall, and in the assumption of an antithetic relation of Eve and Mary, according to which the mother of Christ occupies the same position in the history of redemption as the wife of Adam in the history of sin and death.\footnote{Rom. v. 12 ff.; 1 Cor. xv. 22. But Paul ignores here Eve and Mary altogether.} This idea, so fruitful of many errors, is ingenious, but unscriptural, and an apocryphal substitute for the true Pauline doctrine of an antitypical parallel between the first and second Adam.\footnote{In later times in the Latin church even the \emph{Ave} with which Gabriel saluted the Virgin, was received as the converse of the name of \emph{Eva}; though the Greek \textgreek{χαῖρε} Luke i. 28, admits no such far-fetched accommodation. In like manner the bruising of the serpent’s head, Gen. iii. 15, was applied to Mary instead of Christ, because the Vulgate wrongly translates the Hebrew \texthebrew{שׁאר ךָפְיּשׁיְאוּה}\emph{, ipsa}conteret caput tuum;“while the LXX. rightly refers the \texthebrew{אוּה }to \texthebrew{ﬠרַזֶ }as masc., \textgreek{αὐτός} and likewise all Protestant versions of the Bible.} It tends to substitute Mary for Christ. Justin Martyr, Irenaeus, and Tertullian, are the first who present Mary as the counterpart of Eve, as a “mother of all living” in the higher, spiritual sense, and teach that she became through her obedience the mediate or instrumental cause of the blessings of redemption to the human race, as Eve by her disobedience was the fountain of sin and death.\footnote{Irenaeus: Adv. haer. lib. iii. c. 22, § 4: “Consequenter autem et Maria virgo obediens invenitur, dicens: ’\emph{Ecce ancilla tua, Domine, fiat mihi secundum verbum tuum}’(Luke i. 38); Eva vero disobediens: non obedivit enim, quum adhuc esset virgo. Quemadmodum illa virum quidem habens Adam, virgo tamen adhuc existens ... inobediens facta, et sibi et universo generihumano causa facta est mortis: \emph{sic et Maria} habens praedestinatum virum, et tamen virgo obediens, et \emph{sibi et universo generi humano causa facta est salutis} .... Sic autem et Evae inobedientiea nodus solutionem accepit per obedientiam Mariae. Quod enim allgavit virgo Eva per incredulitatem, hoc virgo Maria solvit per fidem.” Comp. v. 19, § 1. Similar statements occur in Justin M. (Dial.c.Tryph. 100), Tertullian(De carne Christi, c. 17), Epiphanius (Haer. 78, 18), Ephraem (Opp. ii. 318; iii. 607), Jerome(Ep. xxii. ad Eustoch. 21: “Mors per Evam, vita per Mariam ”). Even St. Augustinecarries this parallel between the first and second Eve as far as any of the fathers, in a sermon De Adam et Eva et sancta Maria, not heretofore quoted, published from Vatican Manuscripts in Angelo Mai’s Nova Patrum Bibliotheca, tom. i. Rom. 1852, pp. 1-4. Here, after a most exaggerated invective against woman (whom he calls latrocinium vitae, suavis mors, blanda percussio, interfectio lenis, pernicies delicata, malum libens, sapida jugulatio, omnium calamitas rerum—and all that in a sermon!), goes on thus to draw a contrast between Eve and Mary: “O mulier ista exsecranda, dum decepit! o iterum beata colenda, dum salvat! Plus enim contulit gratiae, quam doloris. Licet ipsa docuerit mortem, ipsa tamen genuit dominum salvatorem. Inventa est ergo mors per mulierem, vita per virginem .... Ergo malum per feminam, immo et per feminam bonum: quia si per Evam cecidimus, magis stamus per Mariam: per Evam sumus servituti addicti effeti per Mariam liberi: Eva nobis sustulit diuturnitatem, aeternitatem nobis Maria condonavit: Eva nos damnari fecit per arboris pomum, absolvit Maria per arboris sacramentum, quia et Christus in ligno pependit ut fructus” (c. 3, pp. 2 and 3). And in conclusion: “Haec mater est humani generis, auctor illa salutis. Eva nos educavit, roboravit et Maria: per Evam cotidie crescimus, regnamus in aeternum per Mariam: per Evam deducti ad terram, ad coelum elevati per Mariam” (c. 4, p. 4). Comp. Aug Sermo 232, c. 2.} Irenaeus calls her also the “advocate of the virgin Eve,” which, at a later day, is understood in the sense of intercessor.\footnote{Adv. haer. v. cap. 19, § 1: “Quemadmodum illa [Eva] seducta est ut effugeret Deum ... sic haec [Maria] suasa est obedire Deo, uti virginis Evae virgo Maria fieret advocata [probably a translation of \textgreek{συνήγορος} or \textgreek{παράκλητος}]. Et quemadmodum adstrictum est morti genus humanum per virginem, \emph{salvatur per virginem}, aequa lance disposita, virginalis inobedientia per virginalem obedientiam.” p 415} On this account this father stands as the oldest leading authority in the Catholic Mariology; though with only partial justice; for he was still widely removed from the notion of the sinlessness of Mary, and expressly declares the answer of Christ in John ii. 4, to be a reproof of her premature haste.\footnote{Adv. haer. iii. cap. 16, § 7 (not. c. 18, as Gieseler, i. 2, p. 277, wrongly cited it): ”... Dominus \emph{repellus ejus intempestivam festinationem}, dixit: \emph{’Quid mihi et tibi est mulier}?’” So even Chrysostom, Hom. 21 in Joh n. 1.} In the same way Tertullian, Origen, Basil the Great, and even Chrysostom, with all their high estimate of the mother of our Lord, ascribe to her on one or two occasions (John ii. 3; Matt. xiii. 47) maternal vanity, also doubt and anxiety, and make this the sword (Luke ii. 35) which, under the cross, passed through her soul.\footnote{Tertullian, De carne Christi, c. 7; Origen, in Luc. Hom. 17; Basil, Ep. 260; Chrysostom, Hom. 44 in Matt, and Hom. 21 in Joh ; Cyril Alex. In Joann. l. xii.} In addition to this typological antithesis of Mary and Eve, the rise of monasticism supplied the development of Mariology a further motive in the enhanced estimate of virginity, without which no true holiness could be conceived. Hence the virginity of Mary, which is unquestioned for the part of her life before the birth of Christ, came to be extended to her whole life, and her marriage with the aged Joseph to be regarded as a mere protectorate, and, therefore, only a nominal marriage. The passage, Matt. i. 25, which, according to its obvious literal meaning (the ἕωςand πρωτότοκος\footnote{The reading \textgreek{πρωτότοκος} in Matt. i. 25 is somewhat doubtful, but it is certainly genuine in Luke ii. 7.}), seems to favor the opposite view, was overlooked or otherwise explained; and the brothers of Jesus,\footnote{They are always called \textgreek{ἀδελφοί} (four in number, James, Joseph or Joses, Simon, and Jude) and \textgreek{αδελφαί} (at least two), Matt. xii. 46, 47; xiii. 55, 56; Mark iii. 31, 32; vi. 3; John vii. 3, 5, 10; Acts i. 14, etc., but nowhere \textgreek{ἀνεψιοί}; Mark a term well known to the N. T. vocabulary (Col. iv. 10), or \textgreek{συγγενεῖς}, kinsmen (Mark vi. 4; Luke i. 36, 58; ii. 44; John viii. 26; Acts x. 24), or \textgreek{υἱοὶτῆς ἀδελφῆς}, sister’s sons (Acts xxiii. 26). This speaks strongly against the cousin-theory.} who appear fourteen or fifteen times in the gospel history and always in close connection with His mother, were regarded not as sons of Mary subsequently born, but either as sons of Joseph by a former marriage (the view of Epiphanius), or, agreeably to the wider Hebrew use of the term ἀcousins of Jesus (Jerome).\footnote{Comp. on this whole complicated question of the brothers of Christ and the connected question of James, the author’s treatise on Jakobus und die Brüder des Herrn, Berlin, 1842, his Hist. of the Apostolic Church, 2d ed. § 95 (p. 383 of the Leipzig ed.; p. 378 of the English), and his article on the Brethren of Christ in the Bibliotheca Sacra of Andover for Oct. 1864} It was felt—and this feeling is shared by many devout Protestants—to be irreconcilable with her dignity and the dignity of Christ, that ordinary children should afterward proceed from the same womb out of which the Saviour of the world was born. The name perpetua virgo, ἀεὶ παρθένος, was thenceforth a peculiar and inalienable predicate of Mary. After the fourth century it was taken not merely in a moral sense, but in the physical also, as meaning that Mary conceived and produced the Lord clauso utero.\footnote{Tertullian(De carne Christi, c. 23: Virgo quantum a viro; \emph{non virgo quantum a partu}), Clement of Alex. (Strom. vii. p. 889), and even Epiphanius (Haer. lxxviii. § 19, where it is said of Christ:\textgreek{Οὗτός ἐστινἀληθῶς ἀνοίγωνμήτρανμητρός}), were still of another opinion on this point. Ambroseof Milan is the first, within my knowledge, to propound this miraculous view (Epist. 42 ad Siricium). He appeals to Ezek. xliv. 1-3, taking the east gate of the temple, which must remain closed because Jehovah passed through it, to refer typically to Mary.“Quos est haec porta, nisi Maria? Ideo clausa, quia virgo. Porta igitur Maria, per quam Christus intravit in hunc mundum.” De inst. Virg. c. 8 (Op. ii. 262). So Ambrosealso in his hymn, ” A solis ortus cardine,“and Jerome, Adv. Pelag. l. ii. 4. The resurrection of Jesus from the closed tomb and the entrance of the risen Jesus through the closed doors, also, was often used as an analogy. The fathers assume that the stone which sealed the Saviour’s tomb, was not rolled away till after the resurrection, and they draw a parallel between the sealed tomb from which He rose to everlasting life, and the closed gate of the Virgin’s womb from which He was born to earthly life. Jerome, \emph{Comment. in Matth}. xxvii. 60: ” Potest novum sepulchrum Mariae virginalem uterum demonstrare.” Gregory the Great: ” Ut ex clauso Virginis utero natus, sic ex clauso sepulchro resurrexit in quo nemo conditus fuerat, et postquam resurrexisset, se per clausas fores in conspectum apostolorum induxit.” Subsequently the catholic view, consistently, removed every other incident of an ordinary birth, such as pain and the flow of blood. While Jeromestill would have Jesus born under all ” naturae contumeliis,“John Damascenus says (De orth. fide, iv. 14): ” Since this birth was not preceded by any [carnal] pleasure, it could also have been followed by no pangs.” Here, too, a passage of prophecy must serve as a proof: Is. lxvi. 7: ” Before she travailed, she brought forth,”\&c.} This, of course, required the supposition of a miracle, like the passage of the risen Jesus through the closed doors. Mary, therefore, in the Catholic view, stands entirely alone in the history of the world in this respect, as in others: that she was a married virgin, a wife never touched by her husband.\footnote{Augustine(De s. virg. c. 6): “Sola Maria et spiritu et corpore mater et virgo.”}

Epiphanius, in his seventy-eighth Heresy, combats the advocates of the opposite view in Arabia toward the end of the fourth century (367), as heretics under the title of Antidikomarianites, opposers of the dignity of Mary, i.e., of her perpetual virginity. But, on the other hand, he condemns, in the seventy-ninth Heresy, the contemporaneous sect of the Collyridians in Arabia, a set of fanatical women, who, as priestesses, rendered divine worship to Mary, and, perhaps in imitation of the worship of Ceres, offered little cakes (κολλυρίδες) to her; he claims adoration for God and Christ alone. Jerome wrote, about 383, with indignation and bitterness against Helvidius and Jovinian, who, citing Scripture passages and earlier church teachers, like Tertullian, maintained that Mary bore children to Joseph after the birth of Christ. He saw in this doctrine a desecration of the temple of the Holy Ghost, and he even compares Helvidius to Erostratus, the destroyer of the temple at Ephesus.\footnote{Helvidiusadduces the principal exegetical arguments for his view; the passages on the Lord’s brothers, and especially Matt. i. 25, pressing the words \textgreek{ἐγίνωσκε} and \textgreek{ἕως}. Jeromeremarks, on the contrary, that the \emph{knowing} by no means necessarily denotes nuptial intercourse, and that \emph{till} does not always fix a limit; e.g., Matt. xxviii. 20 and 1 Cor. xv. 25. In like manner Helvidiuslaid stress on the expression \textgreek{πρωτότοκος}, used of Christ, Matt. i. 25; Luke ii. 7; to which Jeromerightly replies that, according to the law, every son who first opens the womb is called the \emph{first-born}, Ex. xxxiv. 19, 20; Num. xviii. 15 ff., whether followed by other children or not. The “brothers of Jesus” he explains to be cousins, sons of Alpheus and the sister of the Virgin Mary, who likewise was called Mary (as he wrongly infers from John xix. 25). The main argument of Jerome, however, is the ascetic one: the overvaluation of celibacy. Joseph was probably only “custos,” not “maritus Mariae” cap. 19), and their marriage only nominal. He would not indeed deny that there are pious souls among married women and widows, but they are such as have abstained or ceased from living in conjugal intercourse (cap. 21). Helvidius, conversely, ascribed equal moral dignity to the married and the single state. So Jovinian. Comp. § 43.} The bishop Bonosus of Sardica was condemned for the same view by the Illyrican bishops, and the Roman bishop Siricius approved the sentence, a.d. 392.

Augustine went a step farther. In an incidental remark against Pelagius, he agreed with him in excepting Mary, “propter honorem Domini,” from actual (but not from original) sin.\footnote{De Nat. et grat. contra Pelag. c. 36, § 42: ”\emph{Excepta sancto virgine Maria, de qua propter honorem Domini nullam prorsus, cum de peccatis agitur, haberi volo guaestionem ... hac ergo} \emph{virgine excepta}, si omnes illos sanctos et sanctas (whom Peligius takes for sinless] ... congregare possemus et interrogare, utrum essent sine peccato, quid fuisse responsuros putamus: utrum hoc quod iste [Pelagius] dicit, an quod Joannes apostolus” [1 John i. 8]? In other places, however, Augustinesays, that the flesh of Mary came “de peccati propagine” (De Gen. ad Lit. x. c. 18), and that, in virtue of her descent from Adam, she was subject to death also as the consequence of sin (“Maria ex Adam mortua propter peocatum,” Enarrat. in Ps. 34, vs. 12). This was also the view of Anselm of Canterbury († 1109), in his Cur Deus homo, ii. 16, where he says of Christ that he assumed sinless manhood “de massa peccatrice, id est de humano genere, quod totum infectum errat peccato,” and of Mary: “Virgo ipsa, unde assumptus est, est in iniquitatibus concepta, et in peccatis concepit eam mater ejus, et cum originali peccato nata est, quoniam et ipsa in Adam peccavit in quo omnes peccaverunt.” Jerometaught the universal sinfulness without any exception, Adv. Pelag. ii. 4.} This exception he is willing to make from the sinfulness of the race, but no other. He taught the sinless birth and life of Mary, but not her immaculate conception. He no doubt assumed, as afterward Bernard of Clairvaux and Thomas Aquinas, a sanctificatio in utero, like that of Jeremiah (Jer. i. 5) and John the Baptist (Luke i. 15), whereby, as those two men were fitted for their prophetic office, she in a still higher degree was sanctified by a special operation of the Holy Ghost before her birth, and prepared to be a pure receptacle for the divine Logos. The reasoning of Augustine backward from the holiness of Christ to the holiness of His mother was an important turn, which was afterward pursued to further results. The same reasoning leads as easily to the doctrine of the immaculate conception of Mary, though also, just as well, to a sinless mother of Mary herself, and thus upward to the beginning, of the race, to another Eve who never fell. Augustine’s opponent, Pelagius, with his monastic, ascetic idea of holiness and his superficial doctrine of sin, remarkably outstripped him on this point, ascribing to Mary perfect sinlessness. But, it should be remembered, that his denial of original sin to all men, and his excepting of sundry saints of the Old Testament besides Mary, such as Abel, Enoch, Abraham, Isaac, Melchizedek, Samuel, Elijah, Daniel, from actual sin,\footnote{Augustine, De Nat. et grat. cap. 36.} so that πάντεςin Rom. v. 12, in his view, means only a majority, weaken the honor he thus appears to confer upon the mother of the Lord. The Augustinian view long continued to prevail; but at last Pelagius won the victory on this point in the Roman church.\footnote{The doctrine of the Immaculate Conception of Mary was, for the first time after Pelagius, plainly brought forward in 1140 at Lyons, but was opposed by Bernard of Clairvaux (Ep. 174), and thence continued an avowed issue between the Franciscans and Dominicans, till it gained the victory in the papal bull of 1854.}

Notwithstanding this exalted representation of Mary, there appear no clear traces of a proper worship of Mary, as distinct from the worship of saints in general, until the Nestorian controversy of 430. This dispute formed an important turning-point not only in Christology, but in Mariology also. The leading interest in it was, without doubt, the connection of the virgin with the mystery of the incarnation. The perfect union of the divine and human natures seemed to demand that Mary might be called in some sense the mother of God, θεοτόκος, Deipara; for that which was born of her was not merely the man Jesus, but the God-Man Jesus Christ.\footnote{The expression \textgreek{θεοτόκος} does not occur in the Scriptures, and is at best easily misunderstood. The nearest to it is the expression of Elizabeth:\textgreek{Ἡ μήτηρ τοῦ κυρίου μου}, Luke i. 43, and the words of the angel Gabriel: \textgreek{Τὸ γεννώμμενον [ἐκ σοῦ}, \emph{de te}, al. \emph{in te}, is not sufficiently attested, and is a later explanatory addition] \textgreek{ἅγιο νκληθήσεται υἱὸς Θεοῦ}, Luke i. 35. But with what right the distinguished Roman Catholic professor Reithmayr, in the Catholic Encyclop. above quoted, vol. vi. p. 844, puts into the mouth of Elizabeth the expression, “mother \emph{of God} my Lord;” I cannot see; for there is no such variation in the reading of Luke i. 43.} The church, however, did, of course, not intend by that to assert that she was the mother of the uncreated divine essence—for this would be palpably absurd and blasphemous—nor that she herself was divine, but only that she was the human point of entrance or the mysterious channel for the eternal divine Logos. Athanasius and the Alexandrian church teachers of the Nicene age, who pressed the unity of the divine and the human in Christ to the verge of monophysitism, had already used this expression frequently and without scruple,\footnote{The earliest witnesses for \textgreek{θεοτόκος} are Origen (according to Socrates, H. E. vii. 32), Eusebius (Vita Const iii. 43), Cyril of Jerus. (Catech. x. 146), Athanasius (Orat. iii. c. Arian. c. 14, 83), Didymus (De Trinit. i. 31, 94; ii 4, 133), and GregoryNaz. (Orat. li. 738). But it should be remembered that Hesychius, presbyter in Jerusalem († 343) calls David, as an ancestor of Christ, \textgreek{θεοπάτωρ} (Photius, Cod. 275), and that in many apocrypha James is called \textgreek{ἀδελφόθεος} (Gieseler, i. ii. 134). It is also worthy of note that Augustine(† 430), with all his reverence for Mary, never calls her \emph{mater} \emph{Dei} or \emph{Deipara}; on the contrary, he seems to guard against it, Tract. viii. in Ev. Joann. c. 9. “Secundum quod Deus erat [Christus] matrem non habebat.”} and Gregory Nazianzen even declares every one impious who denies its validity.\footnote{Orat. li. 738: \textgreek{Εἴ τιοὐθεοτόκον τὴν Μαρίαν ὑπολαμβάνει, χωρίς ἐστι τῆς θεότητος .}} Nestorius, on the contrary, and the Antiochian school, who were more devoted to the distinction of the two natures in Christ, took offence at the predicate θεοτόκος, saw in it a relapse into the heathen mythology, if not a blasphemy against the eternal and unchangeable Godhead, and preferred the expression Χριστοτόκος, mater Christi. Upon this broke out the violent controversy between him and the bishop Cyril of Alexandria, which ended in the condemnation of Nestorianism at Ephesus in 431.

Thenceforth the θεοτόκοςwas a test of orthodox Christology, and the rejection of it amounted to the beginning or the end of all heresy. The overthrow of Nestorianism was at the same time the victory of Mary-worship. With the honor of the Son, the honor also of the Mother was secured. The opponents of Nestorius, especially Proclus, his successor in Constantinople († 447), and Cyril of Alexandria († 444), could scarcely find predicates enough to express the transcendent glory of the mother of God. She was the crown of virginity, the indestructible temple of God, the dwelling place of the Holy Trinity, the paradise of the second Adam, the bridge from God to man, the loom of the incarnation, the sceptre of orthodoxy; through her the Trinity is glorified and adored, the devil and demons are put to flight, the nations converted, and the fallen creature raised to heaven.\footnote{Comp. Cyril’s Encom. in S. M. Deiparam and Homil. Ephes., and the Orationes of Proclus in Gallandi, vol. ix. Similar extravagant laudation had already been used by Ephraim Syrus († 378) in his work, De laudibus Dei genetricis, and in the collection of prayers which bore his name, but are in part doubtless of later origin, in the 3d volume of his works, pp. 524-552, ed. Benedetti and S. Assemani.} The people were all on the side of the Ephesian decision, and gave vent to their joy in boundless enthusiasm, amidst bonfires, processions, and illuminations.

With this the worship of Mary, the mother of God, the queen of heaven, seemed to be solemnly established for all time. But soon a reaction appeared in favor of Nestorianism, and the church found it necessary to condemn the opposite extreme of Eutychianism or Monophysitism. This was the office of the council of Chalcedon in 451: to give expression to the element of truth in Nestorianism, the duality of nature in the one divine-human person of Christ. Nevertheless the θεοτόκοςwas expressly retained, though it originated in a rather monophysite view.\footnote{\textgreek{Εκ Μαρίας τῆς παρθένον, τῆς θεοτόκου}.}

\xsection{Mariolatry}

§ 82. Mariolatry.

Thus much respecting the doctrine of Mary. Now the corresponding practice. From this Mariology follows Mariolatry. If Mary is, in the strict sense of the word, the mother of God, it seems to follow as a logical consequence, that she herself is divine, and therefore an object of divine worship. This was not, indeed, the meaning and purpose of the ancient church; as, in fact, it never asserted that Mary was the mother of the essential, eternal divinity of the Logos. She was, and continues to be, a created being, a human mother, even according to the Roman and Greek doctrine. But according to the once prevailing conception of her peculiar relation to deity, a certain degree of divine homage to Mary, and some invocation of her powerful intercession with God, seemed unavoidable, and soon became a universal practice.

The first instance of the formal invocation of Mary occurs in the prayers of Ephraim Syrus († 379), addressed to Mary and the saints, and attributed by the tradition of the Syrian church, though perhaps in part incorrectly, to that author. The first more certain example appears in Gregory Nazianzen († 389), who, in his eulogy on Cyprian, relates of Justina that she besought the virgin Mary to protect her threatened virginity, and at the same time disfigured her beauty by ascetic self-tortures, and thus fortunately escaped the amours of a youthful lover (Cyprian before his conversion).\footnote{\textgreek{Τὴν παρθένον Μαρίαν ἱκετεύουσα βοηθῆναι} (Virginem Mariam supplex obsecrans) \textgreek{παρθένῳ κινδυνευούσῃ}. Orat. xviii de St. Cypriano, tom. i. p. 279, ed. Paris. The earlier and authentic accounts respecting Cyprian know nothing of any such courtship of Cyprian and intercession of Mary.} But, on the other hand, the numerous writings of Athanasius, Basil, Chrysostom, and Augustine, furnish no example of an invocation of Mary. Epiphanius even condemned the adoration of Mary, and calls the practice of making offerings to her by the Collyridian women, blasphemous and dangerous to the soul.\footnote{Adv. Haer. Collyrid.: \textgreek{Ἐν τιμῇ ἔστω Μαρία, ὁ δὲ Πατὴρ ... προσκύνείσθω, τὴν Μαρίαν μηδεὶς προσκυνείτω.}} The entire silence of history respecting the worship of the Virgin down to the end of the fourth century, proves clearly that it was foreign to the original spirit of Christianity, and belongs among the many innovations of the post-Nicene age.

In the beginning of the fifth century, however, the worship of saints appeared in full bloom, and then Mary, by reason of her singular relation to the Lord, was soon placed at the head, as the most blessed queen of the heavenly host. To her was accorded the hyperdulia (ὑπερδουλεία)—to anticipate here the later scholastic distinction sanctioned by the council of Trent—that is, the highest degree of veneration, in distinction from mere dulia (δουλεία), which belongs to all saints and angels, and from latria (λατρεία), which, properly speaking, is due to God alone. From that time numerous churches and altars were dedicated to the holy Mother of God, the perpetual Virgin; among them also the church at Ephesus in which the anti-Nestorian council of 431 had sat. Justinian I., in a law, implored her intercession with God for the restoration of the Roman empire, and on the dedication of the costly altar of the church of St. Sophia he expected all blessings for church and empire from her powerful prayers. His general, Narses, like the knights in the Middle Age, was unwilling to go into battle till he had secured her protection. Pope Boniface IV. in 608 turned the Pantheon in Rome into a temple of Mary ad martyres: the pagan Olympus into a Christian heaven of gods. Subsequently even her images (made after an original pretending to have come from Luke) were divinely worshipped, and, in the prolific legends of the superstitious Middle Age, performed countless miracles, before some of which the miracles of the gospel history grow dim. She became almost coördinate with Christ, a joint redeemer, invested with most of His own attributes and acts of grace. The popular belief ascribed to her, as to Christ, a sinless conception, a sinless birth, resurrection and ascension to heaven, and a participation of all power in heaven and on earth. She became the centre of devotion, cultus, and art, the popular symbol of power, of glory, and of the final victory of catholicism over all heresies.\footnote{The Greek church even goes so far as to substitute, in the collects, the name of Mary for the name of Jesus, and to offer petitions in the name of the Theotokos.} The Greek and Roman churches vied throughout the Middle Age (and do so still) in the apotheosis of the human mother with the divine-human child Jesus in her arms, till the Reformation freed a large part of Latin Christendom from this unscriptural semi-idolatry and concentrated the affection and adoration of believers upon the crucified and risen Saviour of the world, the only Mediator between God and man.

A word more: respecting the favorite prayer to Mary, the angelic greeting, or the Ave Maria, which in the Catholic devotion runs parallel to the Pater Noster. It takes its name from the initial words of the salutation of Gabriel to the holy Virgin at the annunciation of the birth of Christ. It consists of three parts:

(1) The salutation of the angel (Luke i. 28): Ave Maria, gratiae plena, Dominus tecum!

(2) The words of Elizabeth (Luke i. 42): Benedicta tu in mulieribus\footnote{96 .These words, according to the \emph{textus receptus,} had been already spoken also by the angel, Luke i. 28: \textgreek{Εὐλογημένησὺ ἐν γυναιξίν}, though they are wanting here in important manuscripts, and are omitted by Tischendorf and Meyer asa later addition, from i. 42.}, et benedictus fructus ventris tui, Jesus.

(3) The later unscriptural addition, which contains the prayer proper, and is offensive to the Protestant and all sound Christian feeling: Sancta Maria, mater Dei, ora pro nobis peccatoribus, nunc et in hora mortis. Amen.

Formerly this third part, which gave the formula the character of a prayer, was traced back to the anti-Nestorian council of Ephesus in 431, which sanctioned the expression mater Dei, or Dei genitrix (θεοτόκος).But Roman archaeologists\footnote{Mast, for example, in Wetzer und Welte’s Kathol. Kirchenlexikon, vol. i. p, 563} now concede that it is a much later addition, made in the beginning of the sixteenth century (1508), and that the closing words, nunc et in hora mortis, were added even after that time by the Franciscans. But even the first two parts did not come into general use as a standing formula of prayer until the thirteenth century.\footnote{Peter Damiani (who died a.d. 1072) first mentions, as a solitary case, that a clergyman daily prayed the words: “Ave Maria, gratia plena! Dominus tecum, benedicta tu in mulieribus.” The first order on the subject was issued by Odo, bishop of Paris, after 1196 (comp. Mansi, xxii. 681): “Exhortentur populum semper presbyteri ad dicendam orationem dominicam et credo in Deum et \emph{salutationem beatae Virginis}.”} From that date the Ave Maria stands in the Roman church upon a level with the Lord’s Prayer and the Apostles’ Creed, and with them forms the basis of the rosary.

\xsection{The Festivals of Mary}

§ 83. The Festivals of Mary.

This mythical and fantastic, and, we must add, almost pagan and idolatrous Mariology impressed itself on the public cultus in a series of festivals, celebrating the most important facts and fictions of the life of the Virgin, and in some degree running parallel with the festivals of the birth, resurrection, and ascension of Christ.

1. The Annunciation of Mary\footnote{\textgreek{Ημέρα ἀσπασμοῦ,} or \textgreek{Χαριτισμοῦ, εύαγγελισμοῦ, ἐνσαρκώσεως}; \emph{festum annunciationis}, s. \emph{incarnationis}, \emph{conceptionis Domini.}} commemorates the announcement of the birth of Christ by the archangel Gabriel,\footnote{Luke i. 26-39.} and at the same time the conception of Christ; for in the view of the ancient church Mary conceived the Logos (Verbum) through the ear by the word of the angel. Hence the festival had its place on the 25th of March, exactly nine months before Christmas; though in some parts of the church, as Spain and Milan, it was celebrated in December, till the Roman practice conquered. The first trace of it occurs in Proclus, the opponent and successor of Nestorius in Constantinople after 430; then it appears more plainly in several councils and homilies of the seventh century.

2. The Purification of Mary\footnote{\emph{Festum purificationis Mariae,} or\emph{praesentationis Domini, Simeonis et Hannae occursus;} \textgreek{ὑπαπάντη}, or \textgreek{ὑπάντη}, or \textgreek{ὑπάντησις τοῦ Κυρίου} (\emph{the meeting of the Lord with Simeon and Anna in the temple}).} or Candlemas, in memory of the ceremonial purification of the Virgin,\footnote{Comp. Luke ii. 22; Lev. xi 2-7. The apparent incongruity of Mary’s need of purification with the prevalent Roman Catholic doctrine of her absolute purity and freedom from the ordinary accompaniments of parturition (even, according to Paschasius Radbert, from the flow of blood) gave rise to all kinds of artificial explanations. Augustinederived it from the consuetudo legis rather than the necessitas expiandi purgandique peccati, and places it on a par with the baptism of Christ. (Quaest. in Heptateuchum, l. iii. c. 40.)} forty days after the birth of Jesus, therefore on the 2d of February (reckoning from the 25th of December); and at the same time in memory of the presentation of Jesus in the temple and his meeting of Simeon and Anna.\footnote{Luke ii. 22-38.} This, like the preceding, was thus originally as much a festival of Christ as of Mary, especially in the Greek church. It is supposed to have been introduced by Pope Gelasius in 494, though by some said not to have arisen till 542 under Justinian I., in consequence of a great earthquake and a destructive pestilence. Perhaps it was a Christian transformation of the old Roman lustrations or expiatory sacrifices (Februa, Februalia), which from the time of Numa took place in February, the month of purification or expiation.\footnote{Februarius, from Februo, the purifying god; like Januarius, from the god Janus. Februare = purgare to purge. February was originally the last month.} To heathen origin is due also the use of lighted tapers, with which the people on this festival marched, singing, out of the church through the city. Hence the name Candlemas.\footnote{\emph{Festum candelarum} sive \emph{luminum}.}

3. The Ascension, or Assumption rather, of Mary\footnote{\textgreek{κοίμησις,} or \textgreek{ἀνάληψις τῆς ἁγίας Θεοτόκου}, \emph{festum assumptionis}.} is celebrated on the 15th of August. The festival was introduced by the Greek emperor Mauritius (582–602); some say, under Pope Gelasius († 496). In Rome, after the ninth century, it is one of the principal feasts, and, like the others, is distinguished with vigil and octave.

It rests, however, on a purely apocryphal foundation.

The entire silence of the apostles and the primitive church teachers respecting the departure of Mary stirred idle curiosity to all sorts of inventions, until a translation like Enoch’s and Elijah’s was attributed to her. In the time of Origen some were inferring from Luke ii. 35, that she had suffered martyrdom. Epiphanius will not decide whether she died and was buried, or not. Two apocryphal Greek writings de transitu Mariae, of the end of the fourth or beginning of the fifth century, and afterward pseudo-Dionysius the Areopagite and Gregory of Tours († 595), for the first time contain the legend that the soul of the mother of God was transported to the heavenly paradise by Christ and His angels in presence of all the apostles, and on the following morning\footnote{According to later representations, as in the three discourses of John Damascenus on this subject, her body rested, like the body of the Lord, \emph{three days} uncorrupted in the grave.} her body also was translated thither on a cloud and there united with the soul. Subsequently the legend was still further embellished, and, besides the apostles, the angels and patriarchs also, even Adam and Eve, were made witnesses of the wonderful spectacle.

Still the resurrection and ascension of Mary are in the Roman church only a matter of “devout and probable opinion,” not an article of faith;\footnote{The Greek council of Jerusalem in 1672, which was summoned against the Calvinists, officially proclaimed it, and thus almost raised it to the authority of a dogma.} and a distinction is made between the ascensio of Christ (by virtue of His divine nature) and the assumptio of Mary (by the power of grace and merit).

But since Mary, according to the most recent Roman dogma, was free even from original sin, and since death is a consequence of sin, it should strictly follow that she did not die at all, and rise again, but, like Enoch and Elijah, was carried alive to heaven.

In the Middle Age—to anticipate briefly—yet other festivals of Mary arose: the Nativity of Mary,\footnote{\emph{Nativitas, natalis B. M. V.;} \textgreek{γενέθλιον}, \&c.} after a.d. 650; the Presentation of Mary,\footnote{\emph{Festum presentationis.}} after the ninth century, founded on the apocryphal tradition of the eleven years’ ascetic discipline of Mary in the temple at Jerusalem; the Visitation of Mary\footnote{\emph{Festum visitationis.}} in memory of her visit to Elizabeth; a festival first mentioned in France in 1247, and limited to the western church; and the festival of the Immaculate Conception,\footnote{\emph{Festum immaculatae} \emph{conceptionis B. M.V.}} which arose with the doctrine of the sinless conception of Mary, and is interwoven with the history of that dogma down to its official and final promulgation by Pope Pius IX. in 1854.

\xsection{The Worship of Martyrs and Saints}

§ 84. The Worship of Martyrs and Saints.

I. Sources: The Memorial Discourses of Basil the Great on the martyr Mamas (a shepherd in Cappadocia, † about 276), and on the forty martyrs (soldiers, who are said to have suffered in Armenia under Licinius in 320); of Gregory Naz. on Cyprian († 248), on Athanasius († 372), and on Basil († 379); of Gregory Of Nyssa on Ephraim Syrus († 378), and on the megalomartyr Theodorus; of Chrysostom on Bernice and Prosdoce, on the Holy Martyrs, on the Egyptian Martyrs, on Meletius of Antioch; several homilies of Ambrose, Augustine, Leo the Great, Peter Chrysologus Caesarius, \&c.; Jerome against Vigilantius. The most important passages of the fathers on the veneration of saints are conveniently collected in: The Faith of Catholics on certain points of controversy, confirmed by Scripture and attested by the Fathers. By Berington and Kirk, revised by Waterworth.” 3d ed. 1846, vol. iii. pp. 322–416.

II. The later Literature: (1) On the Roman Catholic side: The Acta Sanctorum of the Bollandists, thus far 58 vols. fol. (1643–1858, coming down to the 22d of October). Theod. Ruinart: Acta primorum martyrum sincera et selecta. Par. 1689 (confined to the first four centuries). Laderchio: S. patriarcharum et prophetarum, confessorum, cultus perpetuus, etc. Rom. 1730. (2) On the Protestant side: J. Dallaeus: Adversus Latinorum de cultus religiosi objecto traditionem. Genev. 1664. Isaac Taylor: Ancient Christianity. 4th ed. Lond. 1844, vel. ii. p. 173 ff. (“Christianized demonolatry in the fourth century.”)

The system of saint-worship, including both Hagiology and Hagiolatry, developed itself at the same time with the worship of Mary; for the latter is only the culmination of the former.

The New Testament is equally ignorant of both. The expression ἅγιοι, sancti, saints, is used by the apostles not of a particular class, a spiritual aristocracy of the church, but of all baptized and converted Christians without distinction; because they are separated from the world, consecrated to the service of God, washed from the guilt of sin by the blood of Christ, and, notwithstanding all their remaining imperfections and sins, called to perfect holiness. The apostles address their epistles to “the saints” i.e., the Christian believers, “at Rome, Corinth, Ephesus,” \&c.\footnote{Comp. Acts ix. 13, 32, 41; xxvi. 10; Rom. i. 7; xii. 13; xv. 25, 26; 1 Cor. i. 2; vi. 1; Eph. i. 1, 15, 18; iv. 12; Phil. i. 1; iv. 21, 22; Rev. xiii. 7, 10, \&c.}

After the entrance of the heathen masses into the church the title came to be restricted to bishops and councils and to departed heroes of the Christian faith, especially the martyrs of the first three centuries. When, on the cessation of persecution, the martyr’s crown, at least within the limits of the Roman empire, was no longer attainable, extraordinary ascetic piety, great service to the church, and subsequently also the power of miracles, were required as indispensable conditions of reception into the Catholic calendar of saints. The anchorets especially, who, though not persecuted from without, voluntarily crucified their flesh and overcame evil spirits, seemed to stand equal to the martyrs in holiness and in claims to veneration. A tribunal of canonization did not yet exist. The popular voice commonly decided the matter, and passed for the voice of God. Some saints were venerated only in the regions where they lived and died; others enjoyed a national homage; others, a universal.

The veneration of the saints increased with the decrease of martyrdom, and with the remoteness of the objects of reverence. “Distance lends enchantment to the view;” but “familiarity” is apt “to breed contempt.” The sins and faults of the heroes of faith were lost in the bright haze of the past, while their virtues shone the more, and furnished to a pious and superstitious fancy the richest material for legendary poesy.

Almost all the catholic saints belong to the higher degrees of the clergy or to the monastic life. And the monks were the chief promoters of the worship of saints. At the head of the heavenly chorus stands Mary, crowned as queen by the side of her divine Son; then come the apostles and evangelists, who died a violent death, the protomartyr Stephen, and the martyrs of the first three centuries; the patriarchs and prophets also of the Old Covenant down to John the Baptist; and finally eminent hermits and monks, missionaries, theologians, and bishops, and those, in general, who distinguished themselves above their contemporaries in virtue or in public service. The measure of ascetic self-denial was the measure of Christian virtue. Though many of the greatest saints of the Bible, from the patriarch Abraham to Peter, the prince of the apostles, lived in marriage, the Romish ethics, from the time of Ambrose and Jerome, can allow no genuine holiness within the bonds of matrimony, and receives only virgines and some few vidui and viduae into its spiritual nobility.\footnote{To reconcile this perverted view with the Bible, the Roman tradition arbitrarily assumes that Peter separated from his wife after his conversion; whereas Paul, so late as the year 57, expressly presupposes the opposite, and claims for himself the right to take with him a sister as a wife on his missionary tours (\textgreek{ἀδελφὴν γυναῖκα περιάγειν}), like the other apostles, and the brothers of the Lord, and Cephas. 1 Cor. ix. 5. Married saints, like St. Elisabeth of Hungary and St. Louis of France, are rare exceptions.} In this again the close connection of saint-worship with monasticism is apparent.

To the saints, about the same period, were added angels as objects of worship. To angels there was ascribed in the church from the beginning a peculiar concern with the fortunes of the militant church, and a certain oversight of all lands and nations. But Ambrose is the first who expressly exhorts to the invocation of our patron angels, and represents it as a duty.\footnote{De viduis c. 9: “Obsecrandi sunt Angeli pro nobis, qui nobis ad praesidium dati sunt.” Origen had previously \emph{commended} the invocation of angels.} In favor of the guardianship and interest of angels appeal was rightly made to several passages of the Old and New Testaments: Dan. x. 13, 20, 21; xii. 1; Matt. xviii. 10; Luke xv. 7; Heb. i. 14; Acts xii. 15. But in Col. ii. 18, and Rev. xix. 10; xxii. 8, 9, the worship of angels is distinctly rebuked.

Out of the old Biblical notion of guardian angels arose also the idea of patron saints for particular countries, cities, churches, and classes, and against particular evils and dangers. Peter and Paul and Laurentius became the patrons of Rome; James, the patron of Spain; Andrew, of Greece; John, of theologians; Luke, of painters; subsequently Phocas, of seamen; Ivo, of jurists; Anthony, a protector against pestilence; Apollonia, against tooth-aches; \&c.

These different orders of saints and angels form a heavenly hierarchy, reflected in the ecclesiastical hierarchy on earth. Dionysius the Areopagite, a fantastical Christian Platonist of the fifth-century, exhibited the whole relation of man to God on the basis of the hierarchy; dividing the hierarchy into two branches, heavenly and earthly, and each of these again into several degrees, of which every higher one was the mediator of salvation to the one below it.

These are the outlines of the saint-worship of our period. Now to the exposition and estimate of it, and then the proofs.

The worship of saints proceeded originally, without doubt, from a pure and truly Christian source, to wit: a very deep and lively sense of the communion of saints, which extends over death and the grave, and embraces even the blessed in heaven. It was closely connected with love to Christ, and with gratitude for everything great and good which he has done through his instruments for the welfare of posterity. The church fulfilled a simple and natural duty of gratitude, when, in the consciousness of unbroken fellowship with the church triumphant, she honored the memory of the martyrs and confessors, who had offered their life for their faith, and had achieved victory for it over all its enemies. She performed a duty of fidelity to her own children, when she held up for admiration and imitation the noble virtues and services of their fathers. She honored and glorified Christ Himself when she surrounded Him with an innumerable company of followers, contemplated the reflection of His glory in them, and sang to His praise in the Ambrosian Te Deum:

In the first three centuries the veneration of the martyrs in general restricted itself to the thankful remembrance of their virtues and the celebration of the day of their death as the day of their heavenly birth.\footnote{\emph{Natalitia,} \textgreek{γενέθλια}.} This celebration usually took place at their graves. So the church of Smyrna annually commemorated its bishop Polycarp, and valued his bones more than gold and gems, though with the express distinction: “Christ we worship as the Son of God; the martyrs we love and honor as disciples and successors of the Lord, on account of their insurpassable love to their King and Master, as also, we wish to be their companions and fellow disciples.”\footnote{In the Epistle of the church of Smyrna De Martyr. Polycarpi, cap. 17 (Patres-Apost. ed. Dressel, p. 404): \textgreek{Τοῦτον μὲν γὰρ υἱὸν ὄντα τοῦ Θεοῦ προσκυνοῦμεν· τοὺς δὲ μάρτυρας , ὡς μαθητὰς καὶ μιμητὰς τοῦ κυρίου ἀγαπῶμεν ἀξίως , κ. τ. λ}} Here we find this veneration as yet in its innocent simplicity.

But in the Nicene age it advanced to a formal invocation of the saints as our patrons (patroni) and intercessors (intercessores, mediatores) before the throne of grace, and degenerated into a form of refined polytheism and idolatry. The saints came into the place of the demigods, Penates and Lares, the patrons of the domestic hearth and of the country. As once temples and altars to the heroes, so now churches and chapels\footnote{Memoriae, \textgreek{μαρτύρια}.} came to be built over the graves of the martyrs, and consecrated to their names (or more precisely to God through them). People laid in them, as they used to do in the temple of Aesculapius, the sick that they might be healed, and hung in them, as in the temples of the gods, sacred gifts of silver and gold. Their graves were, as Chrysostom says, move splendidly adorned and more frequently visited than the palaces of kings. Banquets were held there in their honor, which recall the heathen sacrificial feasts for the welfare of the manes. Their relics were preserved with scrupulous care, and believed to possess miraculous virtue. Earlier, it was the custom to pray for the martyrs (as if they were not yet perfect) and to thank God for their fellowship and their pious example. Now such intercessions for them were considered unbecoming, and their intercession was invoked for the living.\footnote{Augustine, Serm. 159, 1 (al. 17): “Injuria est pro martyre orare, cujus nos debemus orationibus commendari.” Serm. 284, 5: “Pro martyribus non orat [ecclesia], sed eorum potius orationibus se commendat.” Serm. 285, 5: “Pro aliis fidelibus defunctis oratur [to wit, for the souls in purgatory still needing purification]; \emph{pro martyribus non oratur}; tam enim perfecti exierunt, ut non sint suscepti nostri, sed \emph{advocati}.” Yet Augustineadds the qualification: “Neque hoc in se, sed in illo cui capiti perfecta membra cohaeserunt. Ille est enim vere \emph{advocatus unus}, qui interpellat pro nobis, sedens ad dexteram Patris: sed advocatus unus, sicut et pastor unus.” When the grateful intercessions for the departed saints and martyrs were exchanged for the invocation of their intercession, the old formula: “Annue nobis, Domine, ut animae famuli tui Leonis haec prosit oblatio,” was changed into the later: “Annue nobis, quaesumus, Domine, ut intercessione beati Leonis haec nobis prosit oblatio.” But instead of praying for the saints, the Catholic church now prays for the souls in purgatory.}

This invocation of the dead was accompanied with the presumption that they take the deepest interest in all the fortunes of the kingdom of God on earth, and express it in prayers and intercessions.\footnote{Ambrose, De viduis, c. 9, calls the martyrs “nostri praesules et speculatores (spectatores) vitae actuumque nostrorum.”} This was supposed to be warranted by some passages of Scripture, like Luke xv. 10, which speaks of the angels (not the saints) rejoicing over the conversion of a sinner, and Rev. viii. 3, 4, which represents an angel as laying the prayers of all the saints on the golden altar before the throne of God. But the New Testament expressly rebukes the worship of the angels (Col. ii. 18; Rev. xix. 10; xxii. 8, 9), and furnishes not a single example of an actual invocation of dead men; and it nowhere directs us to address our prayers to any creature. Mere inferences from certain premises, however plausible, are, in such weighty matters, not enough. The intercession of the saints for us was drawn as a probable inference from the duty of all Christians to pray for others, and the invocation of the saints for their intercession was supported by the unquestioned right to apply to living saints for their prayers, of which even the apostles availed themselves in their epistles.

But here rises the insolvable question: How can departed saints hear at once the prayers of so many Christians on earth, unless they either partake of divine omnipresence or divine omniscience? And is it not idolatrous to clothe creatures with attributes which belong exclusively to Godhead? Or, if the departed saints first learn from the omniscient God our prayers, and then bring them again before God with their powerful intercessions, to what purpose this circuitous way? Why not at once address God immediately, who alone is able, and who is always ready, to hear His children for the sake of Christ?

Augustine felt this difficulty, and concedes his inability to solve it. He leaves it undecided, whether the saints (as Jerome and others actually supposed) are present in so many places at once, or their knowledge comes through the omniscience of God, or finally it comes through the ministry of angels.\footnote{De cura pro mortuis (a.d. 421), c. 16. In another place he decidedly rejects the first hypothesis, because otherwise he himself would be always surrounded by his pious mother, and because in Isa. lxiii. 16 it is said: “Abraham is ignorant of us.”} He already makes the distinction between λατρεία, or adoration due to God alone, and the invocatio (δουλεία) of the saints, and firmly repels the charge of idolatry, which the Manichaean Faustus brought against the catholic Christians when he said: “Ye have changed the idols into martyrs, whom ye worship with the like prayers, and ye appease the shades of the dead with wine and flesh.” Augustine asserts that the church indeed celebrates the memory of the martyrs with religious solemnity, to be stirred up to imitate them, united with their merits, and supported by their prayers,\footnote{“Et ad excitandam imitationem, et ut mentis eorum consocietur, atque orationibus adjuvetur.” Contra Faustum l. 20, n. 21.} but it offers sacrifice and dedicates altars to God alone. Our martyrs, says he, are not gods; we build no temples to our martyrs, as to gods; but we consecrate to them only memorial places, as to departed men, whose spirits live with God; we build altars not to sacrifice to the martyrs, but to sacrifice with them to the one God, who is both ours and theirs.\footnote{De Civit. Dei, xxii. 10: “Nobis Martyres non sunt dii: quia unum eundemque Deum et nostrum scimus et Martyrum. Nec tamen miraculis, quae per Memorias nostrorum Martyrum fiunt, ullo modo comparanda sunt miracula, quae facta per templa perhibentur illorum. Verum si qua similia videntur, sicut a Moyse magi Pharaonis, sic eorum dii victi sunt a Martyribus nostris .... Martyribus nostris non templa sicut diis, sed Memorias sicut hominibus mortuis, quorum apud Deum vivunt spiritus, fabricamus; nec ibi erigimus altaria, in quibus sacrificemus Martyribus, sed uni Deo et Martyrum et nostro sacrificium [corpus Christi] immolamus.”}

But in spite of all these distinctions and cautions, which must be expected from a man like Augustine, and acknowledged to be a wholesome restraint against excesses, we cannot but see in the martyr-worship, as it was actually practised, a new form of the hero-worship of the pagans. Nor can we wonder in the least. For the great mass of the Christian people came, in fact, fresh from polytheism, without thorough conversion, and could not divest themselves of their old notions and customs at a stroke. The despotic form of government, the servile subjection of the people, the idolatrous homage which was paid to the Byzantine emperors and their statues, the predicates divina, sacra, coelestia, which were applied to the utterances of their will, favored the worship of saints. The heathen emperor Julian sarcastically reproached the Christians with reintroducing polytheism into monotheism, but, on account of the difference of the objects, revolted from the Christian worship of martyrs and relics, as from the “stench of graves and dead men’s bones.” The Manichaean taunt we have already mentioned. The Spanish presbyter Vigilantius, in the fifth century, called the worshippers of martyrs and relics, ashes-worshippers and idolaters,\footnote{\emph{Cinerarios} and \emph{idololatrae}.} and taught that, according to the Scriptures, the living only should pray with and for each other. Even some orthodox church teachers admitted the affinity of the saint-worship with heathenism, though with the view of showing that all that is good in the heathen worship reappears far better in the Christian. Eusebius cites a passage from Plato on the worship of heroes, demi-gods, and their graves, and then applies it to the veneration of friends of God and champions of true religion; so that the Christians did well to visit their graves, to honor their memory there, and to offer their prayers.\footnote{In his Praeparat. Evangelica, xiii. cap. 11, p, 663. Comp. Demostr. Evang. iii. § 3, p. 107.} The Greeks, Theodoret thinks, have the least reason to be offended at what takes place at the graves of the martyrs; for the libations and expiations, the demi-gods and deified men, originated with themselves. Hercules, Aesculapius, Bacchus, the Dioscuri, and the like, are deified men; consequently it cannot be a reproach to the Christians that they—not deify, but—honor their martyrs as witnesses and servants of God. The ancients saw nothing censurable in such worship of the dead. The saints, our helpers and patrons, are far more worthy of such honor. The temples of the gods are destroyed, the philosophers, orators, and emperors are forgotten, but the martyrs are universally known. The feasts of the gods are now replaced by the festivals of Peter, Paul, Marcellus, Leontius, Antonins, Mauricius, and other martyrs, not with pagan pomp and sensual pleasures, but with Christian soberness and decency.\footnote{Theodoret, Graec. affect. curatio. Disp. viii. (Ed. Schulz, iv. p. 902 sq.)}

Yet even this last distinction which Theodoret asserts, sometimes disappeared. Augustine laments that in the African church banqueting and revelling were daily practised in honor of the martyrs,\footnote{“Commessationes et ebrietates in honorem etiam beatissimorum Martyrum.” Ep. 22 and 29.} but thinks that this weakness must be for the time indulged from regard to the ancient customs of the pagans.

In connection with the new hero-worship a new mythology also arose, which filled up the gaps of the history of the saints, and sometimes even transformed the pagan myths of gods and heroes into Christian legends.\footnote{Thus, e.g., the fate of the Attic king’s son Hippolytus, who was dragged to death by horses on the sea shore, was transferred to the Christian martyr Hippolytus, of the beginning of the third century. The martyr Phocas, a gardener at Sinope in Pontus, became the patron of all mariners, and took the place of Castor and Pollux. At the daily meals on shipboard, Phocas had his portion set out among the rest, as an invisible guest, and the proceeds of the sale of these portions was finally distributed among the poor as a thank-offering for the prosperous voyage.} The superstitious imagination, visions, and dreams, and pious fraud famished abundant contributions to the Christian legendary poesy.

The worship of the saints found eloquent vindication and encouragement not only, in poets like Prudentius (about 405) and Paulinus of Nola (died 431), to whom greater freedom is allowed, but even in all the prominent theologians and preachers of the Nicene and post-Nicene age. It was as popular as monkery, and was as enthusiastically commended by the leaders of the church in the East and West.

The two institutions, moreover, are closely connected and favor each other. The monks were most zealous friends of saint-worship in their own cause. The church of the fifth century already went almost as far in it as the Middle Age, at all events quite as far as the council of Trent; for this council does not prescribe the invocation of the saints, but confines itself to approving it as “good and useful” (not as necessary) on the ground of their reigning with Christ in heaven and their intercession for us, and expressly remarks that Christ is our only, Redeemer and Saviour.\footnote{Conc. Trid. Sess. xxv.: “Sanctos una cum Christo regnantes orationes suas pro hominibus Deo offere;\emph{bonum} atque \emph{utile} esse suppliciter eos invocare et ob beneficia impetranda a Deo per Filium eius Jesum Christum, qui solus noster redemptor et salvator est, ad eorum orationes, opem auxiliumque confugere.”} This moderate and prudent statement of the doctrine, however, has not yet removed the excesses which the Roman Catholic people still practise in the worship of the saints, their images, and their relics. The Greek church goes even further in theory than the Roman; for the confession of Peter Mogilas (which was subscribed by the four Greek patriarchs in 1643, and again sanctioned by the council of Jerusalem in 1672), declares it duty and propriety (χρέος) to implore the intercession (μεσιτεία) of Mary and the saints with God for us.

We now cite, for proof and further illustration, the most important passages from the church fathers of our period on this point. In the numerous memorial discourses of the fathers, the martyrs are loaded with eulogies, addressed as present, and besought for their protection. The universal tone of those productions is offensive to the Protestant taste, and can hardly be reconciled with evangelical ideas of the exclusive and all-sufficient mediation of Christ and of justification by pure grace without the merit of works. But it must not be forgotten that in these discourses very much is to be put to the account of the degenerate, extravagant, and fulsome rhetoric of that time. The best church fathers, too, never separated the merits of the saints from the merits of Christ, but considered the former as flowing out of the latter.

We begin with the Greek fathers. Basil the Great calls the forty soldiers who are said to have suffered martyrdom under Licinius in Sebaste about 320, not only a “holy choir,” an “invincible phalanx,” but also “common patrons of the human family, helpers of our prayers and most mighty intercessors with God.”\footnote{Basil. M. Hom. 19, in XL. Martyres, § 8: \textgreek{Ὢ χορὸς ἅγιος, ὢ σύνταγμα ἱερόν, ὢ συναπισμὸς ἀῥῥαγής, ὢ κοινοὶ φύλακες τοῦ γένους τῶν ἀνθρώπων} (Ocommunes generis humani custodes), \textgreek{ἀγαθοὶ κοινωνοὶ φροντίδων, δεήσεως συνεργοὶ, πρεσβευταὶδυνατώτατοι} (legati apud Deum potentissimi), \textgreek{ἀστέρες τῆς οἰκουμένης, ἄνθη τῶν ἐκκλησιῶν, ὑμᾶς ούχ ἡ γῆ κατέκρυψεν, ἀλλ οὐρανὸς ὑπεδέξατο}.}

Ephraim Syrus addresses the departed saints, in general, in such words as these: “Remember me, ye heirs of God, ye brethren of Christ, pray to the Saviour for me, that I through Christ may be delivered from him who assaults me from day to day;” and the mother of a martyr: “O holy, true, and blessed mother, plead for me with the saints, and pray: ’Ye triumphant martyrs of Christ, pray for Ephraim, the least, the miserable,’ that I may find grace, and through the grace of Christ may be saved.”

Gregory of Nyssa asks of St. Theodore, whom he thinks invisibly present at his memorial feast, intercessions for his country, for peace, for the preservation of orthodoxy, and begs him to arouse the apostles Peter and Paul and John to prayer for the church planted by them (as if they needed such an admonition!). He relates with satisfaction that the people streamed to the burial place of this saint in such multitudes that the place looked like an ant hill. In his Life of St. Ephraim, he tells of a pilgrim who lost himself among the barbarian posterity of Ishmael, but by the prayer, “St. Ephraim, help me!”\footnote{\textgreek{Αγιε Εφραὶμ, βαήθειμοί}.} and the protection of the saint, happily found his way home. He himself thus addresses him at the close: “Thou who standest at the holy altar, and with angels servest the life-giving and most holy Trinity, remember us all, and implore for us the forgiveness of sins and the enjoyment of the eternal kingdom.”\footnote{Ἀ\textgreek{ιτούμενος ἡμῖν ἁμαρτημάτων ἄφεσιν, αίωνίου τὲ βασιλείας ἀπόλαυσιν.} De vita Ephraem. p. 616 (tom. iii.).}

Gregory Nazianzen is convinced that the departed Cyprian guides and protects his church in Carthage more powerfully by his intercessions than he formerly did by his teachings, because he now stands so much nearer the Deity; he addresses him as present, and implores his favor and protection.\footnote{\textgreek{Σὺ δὲ ἡμᾶς ἐποπτεύσις ἄνωθεν ἵλεως, καὶ τὸν ἡμέτερον διεξάγοις λόγον καὶ βίον, κ. τ.λ.} Orat. 18 in laud. Cypr. p. 286.} In his eulogy on Athanasius, who was but a little while dead, he prays: “Look graciously down upon us, and dispose this people to be perfect worshippers of the perfect Trinity; and when the times are quiet, preserve us—when they are troubled, remove us, and take us to thee in thy fellowship.”

Even Chrysostom did not rise above the spirit of the time. He too is an eloquent and enthusiastic advocate of the worship of the saints and their relics. At the close of his memorial discourse on Sts. Bernice and Prosdoce—two saints who have not even a place in the Roman calendar—he exhorts his hearers not only on their memorial days but also on other days to implore these saints to be our protectors: “For they have great boldness not merely during their life but also after death, yea, much greater after death.\footnote{\textgreek{Παρακαλῶμεν αὐτὰς, ἀξιῶμεν γενέσθαι προστάτιδας ἡμῶν; πολλὴν γὰρ ἔχουσιν παῤῥησίαν οὐχὶ ζῶσαι μόνον, ἀλλὰ καὶ τελευτήσασαι· καὶ πολλῶ μᾶλλον τελευτήσασαι.} Opp. tom. ii. 770.} For they now bear the stigmata of Christ [the marks of martyrdom], and when they show these, they can persuade the King to anything.” He relates that once, when the harvest was endangered by excessive rain, the whole population of Constantinople flocked to the church of the Apostles, and there elected the apostles Peter and Andrew, Paul and Timothy, patrons and intercessors before the throne of grace.\footnote{Contra ludos et theatra, n. 1, tom. vi. 318.} Christ, says he on Heb. i. 14, redeems us as Lord and Master, the angels redeem us as ministers.

Asterius of Amasia calls the martyr Phocas, the patron of mariners, “a pillar and foundation of the churches of God in the world, the most renowned of the martyrs, who draws men of all countries in hosts to his church in Sinope, and who now, since his death, distributes more abundant nourishment than Joseph in Egypt.”

Among the Latin fathers, Ambrose of Milan is one of the first and most decided promoters of the worship of saints. We cite a passage or two. “May Peter, who so successfully weeps for himself, weep also for us, and turn upon us the friendly look of Christ.\footnote{Hexaem. l. v. cap. 25, § 90: “Fleat pro nobis Petrus, qui pro se bene flevit, et in nos pia Christi ora convertat. Approperet Jesu Domini passio, quae quotidie delicta nostra condonat et munus remissionis operatur.”} The angels, who are appointed to guard us, must be invoked for us; the martyrs, to whose intercession we have claim by the pledge of their bodies, must be invoked. They who have washed away their sins by their own blood, may pray for our sins. For they are martyrs of God, our high priests, spectators of our life and our acts. We need not blush to use them as intercessors for our weakness; for they also knew the infirmity of the body when they gained the victory over it.”\footnote{De viduis, c. 9: “Obsecrandi sunt Angeli pro nobis, qui nobis ad praesidium dati sunt; martyres obsecrandi, quorum videmur nobis quoddam corporis pignore patrocinium vindicare. Possunt pro peccatis rogare nostris, qui proprio sanguine etiam si qua habuerunt peccata laverunt. Isti enim sunt Dei martyres, nostri prae sules, speculatores vitae actuumque nostrorum,” etc. Ambrosegoes farther than the council of Trent, which does not command the invocation of the saints, but only commends it, and represents it not as duty, but only as privilege. See the passage already cited, p. 437.}

Jerome disputes the opinion of Vigilantius, that we should pray for one another in this life only, and that the dead do not hear our prayers, and ascribes to departed saints a sort of omnipresence, because, according to Rev. xiv. 4, they follow the Lamb whithersoever he goeth.\footnote{Adv. Vigilant. n. 6: “Si agnus ubique, ergo et hi, qui cum agno sunt, ubique esse credendi sunt.” So the heathen also attributed ubiquity to their demons. Hesiodus, Opera et dies, v. 121 sqq.} He thinks that their prayers are much more effectual in heaven than they were upon earth. If Moses implored the forgiveness of God for six hundred thousand men, and Stephen, the first martyr, prayed for his murderers after the example of Christ, should they cease to pray, and to be heard, when they are with Christ?

Augustine infers from the interest which the rich man in hell still had in the fate of his five surviving brothers (Luke xvi. 27), that the pious dead in heaven must have even far more interest in the kindred and friends whom they have left behind.\footnote{Epist. 259, n. 5.} He also calls the saints our intercessors, yet under Christ, the proper and highest Intercessor, as Peter and the other apostles are shepherds under the great chief Shepherd.\footnote{Sermo 285, n. 5.} In a memorial discourse on Stephen, he imagines that martyr, and St. Paul who stoned him, to be present, and begs them for their intercessions with the Lord with whom they reign.\footnote{Sermo 317, n. 5: “Ambo modo sermonem nostrum auditis; ambo pro nobis orate ... orationibus suis commendent nos.”} He attributes miraculous effects, even the raising of the dead, to the intercessions of Stephen.\footnote{Serm. 324.} But, on the other hand, he declares, as we have already observed, his inability to solve the difficult question of the way in which the dead can be made acquainted with our wishes and prayers. At all events, in Augustine’s practical religion the worship of the saints occupies a subordinate place. In his “Confessions” and “Soliloquies” he always addresses himself directly to God, not to Mary nor to martyrs.

The Spanish poet Prudentius flees with prayers and confessions of sin to St. Laurentius, and considers himself unworthy to be heard by Christ Himself.\footnote{Hymn. ii. in hon. S. Laurent. vss. 570-584: \par “Indignus agnosco et scio, \par Quem Christus ipse exaudiat; \par —- Sed per patronos martyres \par Potest medelam consequi.”}

The poems of Paulinus of Nola are full of direct prayers for the intercessions of the saints, especially of St. Felix, in whose honor he erected a basilica, and annually composed an ode, and whom he calls his patron, his father, his lord. He relates that the people came in great crowds around the wonder-working relics of this saint on his memorial day, and could not look on them enough.

Leo the Great, in his sermons, lays great stress on the powerful intercession of the apostles Peter and Paul, and of the Roman martyr Laurentius.\footnote{“Cuius oratione,” says he of the latter, “et patrocinio adjuvari nos sine cessatione confidimus.” Serm. 85 in Natal. S. Laurent c. 4.}

Pope Gregory the Great, at the close of our period, went much farther.

According to this we cannot wonder that the Virgin Mary and the saints are interwoven also in the prayers of the liturgies,\footnote{E.g., the Liturgies of St. James, St. Mark, St. Basil, St. Chrysostom, the Coptic Liturgy of St. Cyril, and the Roman Liturgy.} and that their merits and intercession stand by the side of the merits of Christ as a ground of the acceptance of our prayers.

\xsection{Festivals of the Saints}

§ 85. Festivals of the Saints.

The system of saint-worship, like that of the worship of Mary, became embodied in a series of religious festivals, of which many had only a local character, some a provincial, some a universal. To each saint a day of the year, the day of his death, or his heavenly birthday, was dedicated, and it was celebrated with a memorial oration and exercises of divine worship, but in many cases desecrated by unrestrained amusements of the people, like the feasts of the heathen gods and heroes.

The most important saints’ days which come down from the early church, and bear a universal character, are the following:

1. The feast of the two chief apostles Peter and Paul,\footnote{\emph{Natalis apostolorum Petri et Pauli}.} on the twenty-ninth of June, the day of their martyrdom. It is with the Latins and the Greeks the most important of the feasts of the apostles, and, as the homilies for the day by Gregory Nazianzen, Chrysostom, Ambrose, Augustine, and Leo the Great show, was generally introduced as early as the fourth century

2. Besides this, the Roman church has observed since the fifth century a special feast in honor of the prince of the apostles and for the glorification of the papal office: the feast of the See of Peter\footnote{\emph{Festum cathedrae Petri.}} on the twenty-second of February, the day on which, according to tradition, he took possession of the Roman bishopric. With this there was also an Antiochan St. Peter’s day on the eighteenth of January, in memory of the supposed episcopal reign of this apostle in Antioch. The Catholic liturgists dispute which of these two feasts is the older. After Leo the Great, the bishops used to keep the Natales. Subsequently the feast of the Chains of Peter\footnote{\emph{Festum catenarum Petri}, commonly \emph{Petri ad vincula}, on the first of August. According to the legend, the Herodian Peter’s-chain, which the empress Eudoxia, wife of Theodosius II., discovered on a pilgrimage in Jerusalem, and sent as a precious relic to Rome, miraculously united with the Neronian Peter’s-chain at Rome on the first contact, so that the two have since formed only one holy and inseparable chain!} was introduced in memory of the chains which Peter wore, according to Acts xii. 6, under Herod at Jerusalem, and, according to the Roman legend, in the prison at Rome under Nero.

3. The feast of John, the apostle and evangelist, on the twenty-seventh of December, has already been mentioned in connection with the Christmas cycle.\footnote{Comp. § 77, volume 3}

4. Likewise the feast of the protomartyr Stephen, on the twenty-sixth of December, after the fourth century.\footnote{· Ibid.}

5. The feast of John the Baptist, the last representative of the saints before Christ. This was, contrary to the general rule, a feast of his birth, not his martyrdom, and, with reference to the birth festival of the Lord on the twenty-fifth of December, was celebrated six months earlier, on the twenty-fourth of June, the summer solstice. This was intended to signify at once his relation to Christ and his well-known word: “He must increase, but I must decrease.” He represented the decreasing sun of the ancient covenant; Christ, the rising sun of the new.\footnote{Comp. John iii. 30. This interpretation is given by Augustine, Serm. 12 in Nat. Dom.: “In nativitate Christi \emph{dies crescit,} in Johannis nativitate \emph{decrescit}. Profectum plane facit dies, quum mundi Salvator oritur; defectum patitur, quum ultimus prophetarum generatur.”} In order to celebrate more especially the martyrdom of the Baptist, a feast of the Beheading of John,\footnote{\emph{Festum decollationis S. Johannis B.}} on the twenty-ninth of August, was afterward introduced; but this never became so important and popular as the feast of his birth.

6. To be just to all the heroes of the faith, the Greek church, after the fourth century, celebrated a feast of All Saints on the Sunday after Pentecost (the Latin festival of the Trinity).\footnote{This Sunday is therefore called by the Greeks the \emph{Martyrs’} \emph{and Saints’} Sun\emph{day,} \textgreek{ἡ κυριακὴ τῶν ἁγίων πάντων,} or \textgreek{τῶν ἁγίων καὶ μαρτύρων}. We have a homily of Chrysostomon it: \textgreek{Ἐγκώμιονεἰς τοῦς ἁγίους πάντας τοῦς ἐνὅλῳ τῷ κόσμῳ μαρτυρήσαντες}, or De martyribus totius orbis. Hom. lxxiv. Opera, tom. ii. 711 sqq.} The Latin church, after 610, kept a similar feast, the Festum Omnium Sanctorum, on the first of November; but this did not come into general use till after the ninth century.

7. The feast of the Archangel Michael,\footnote{\emph{Festum S. Michaelis, archangeli.}} the leader of the hosts of angels, and the representative of the church triumphant,\footnote{Rev. xii. 7-9; comp. Jude, vs. 9.} on the twenty-ninth of September. This owes its origin to some miraculous appearances of Michael in the Catholic legends.\footnote{Comp. Augusti, Archaeologie, i. p. 585. Michael, e.g., in a pestilence in Rome in the seventh century, is said to have appeared as a deliverer on the Tomb of Hadrian (Moles Hadriani, or Mausoleo di Adriano), so that the place received the name of Angel’s Castle (Castello di S. Angelo). It lies, as is well known, at the great bridge of the Tiber, and is used as a fortress.} The worship of the angels developed itself simultaneously with the worship of Mary and the saints, and churches also were dedicated to angels, and called after their names. Thus Constantine the Great built a church to the archangel Michael on the right bank of the Black Sea, where the angel, according to the legend, appeared to some ship-wrecked persons and rescued them from death. Justinian I. built as many as six churches to him. Yet the feast of Michael, which some trace back to Pope Gelasius I., a.d. 493, seems not to have become general till after the ninth century.

\xsection{The Christian Calendar. The Legends of the Saints. The Acta Sanctorum}

§ 86. The Christian Calendar. The Legends of the Saints. The Acta Sanctorum.

This is the place for some observations on the origin and character of the Christian calendar with reference to its ecclesiastical elements, the catalogue of saints and their festivals.

The Christian calendar, as to its contents, dates from the fourth and later centuries; as to its form, it comes down from classical antiquity, chiefly from the Romans, whose numerous calendars contained, together with astronomical and astrological notes, tables also of civil and religious festivals and public sports. Two calendars of Christian Rome still extant, one of the year 354, the other of the year 448,\footnote{The latter is found in the Acta Sanct. Jun. tom. vii. p. 176 sqq.} show the transition. The former contains for the first time the Christian week beginning with Sunday, together with the week of heathen Rome; the other contains Christian feast days and holidays, though as yet very few, viz., four festivals of Christ and six martyr days. The oldest purely Christian calendar is a Gothic one, which originated probably in Thrace in the fourth century. The fragment still extant\footnote{Printed in Angelo Mai, Script. vet. nova collect. tom. v. P. 1, pp. 66-68. Comp. Krafft, Kirchengeschichte der germanischen Völker. Vol. i. Div. 1, pp. 385-387.} contains thirty-eight days for November and the close of October, among which seven days are called by the names of saints (two from the Bible, three from the church universal, and two from the Gothic church).

There are, however, still earlier lists of saints’ days, according to the date of the holiday; the oldest is a Roman one of the middle of the fourth century, which contains the memorial days of twelve bishops of Rome and twenty-four martyrs, together with the festival of the birth of Christ and the festival of Peter on the twenty-second of February.

Such tables are the groundwork of the calendar and the martyrologies. At first each community or province had its own catalogue of feasts, hence also its own calendar. Such local registers were sometimes called diptycha\footnote{From \textgreek{δίπτυχος}, folded double.} (δίπτυχα), because they were recorded on tables with two leaves; yet they commonly contained, besides the names of the martyrs, the names also of the earlier bishops and still living benefactors or persons, of whom the priests were to make mention by name in the prayer before the consecration of the elements in the eucharist. The spread of the worship of a martyr, which usually started from the place of his martyrdom, promoted the interchange of names. The great influence of Rome gave to the Roman festival-list and calendar the chief currency in the West.

Gradually the whole calendar was filled up with the names of saints. As the number of the martyrs exceeded the number of days in the year, the commemoration of several must fall upon the same day, or the canonical hours of cloister devotion must be given up. The oriental calendar is richer in saints from the Old Testament than the occidental.\footnote{The Roman Catholic saint-calendars have passed, without material change, to the Protestant church in Germany and other countries. Recently Prof. Piper in Berlin has attempted a thorough evangelical reform of the calendar by rejecting the doubtful or specifically Roman saints, and adding the names of the forerunners of the Reformation and the Reformers and distinguished men of the Protestant churches to the list under their birthdays. To this reform also his \emph{Evangelischer Kalender} is devoted, which has appeared annually since 1850, and contains brief, popular sketches of the Catholic and Protestant saints received into the improved calendar. Most English and American calendars entirely omit this list of saints.}

With the calendars are connected the Martyrologia, or Acta Martyrum, Acta Sanctorum, called by the Greeks Menologia and Menaea.\footnote{From \textgreek{μήν}month; hence, month-register. The Greek Menologies, \textgreek{μηνολόγια}are simply the lists of the martyrs in monthly order, with short biographical notices. The Menaea, \textgreek{μηναῖα}, are intended for the public worship and comprise twelve folio volumes, corresponding to the twelve months, with the \emph{officia} of the saints for every day, and the proper legends and hymns.} There were at first only “Diptycha” and “Calendaria martyrum,” i.e., lists of the names of the martyrs commemorated by the particular church in the order of the days of their death on the successive days of the year, with or without statements of the place and manner of their passion. This simple skeleton became gradually animated with biographical sketches, coming down from different times and various authors, containing a confused mixture of history and fable, truth and fiction, piety and superstition, and needing to be used with great critical caution. As these biographies of the saints were read on their annual days in the church and in the cloisters for the edification of the people, they were called Legenda.

The first Acts of the Martyrs come down from the second and third centuries, in part from eye-witnesses, as, for example, the martyrdom of Polycarp (a.d. 167), and of the martyrs of Lyons and Vienne in South Gaul; but most of them originated, at least in their present form, in the post-Constantinian age. Eusebius wrote a general martyrology, which is lost. The earliest Latin martyrology is ascribed to Jerome, but at all events contains many later additions; this father, however, furnished valuable contributions to such works in his “Lives of eminent Monks” and his “Catalogue of celebrated Church Teachers.” Pope Gelasius thought good to prohibit or to restrict the church reading of the Acts of the Saints, because the names of the authors were unknown, and superfluous and incongruous additions by heretics or uneducated persons (idiotis) might be introduced. Gregory the Great speaks of a martyrology in use in Rome and elsewhere, which is perhaps the same afterward ascribed to Jerome and widely spread. The presentMartyrologium Romanum, which embraces the saints of all countries, is an expansion of this, and was edited by Baronius with a learned commentary at the command of Gregory XIII. and Sixtus V. in 1586, and afterward enlarged by the Jesuit Heribert Rosweyd.

Rosweyd († 1629) also sketched, toward the close of the sixteenth century, the plan for the celebrated “Acta Sanctorum, quotquot toto orbe coluntur,” which Dr. John van Bolland († 1665) and his companions and continuators, called Bollandists (Henschen, † 1681; Papenbroek, † 1714; Sollier, † 1740; Stiltinck, † 1762, and others of inferior merit), published at Antwerp in fifty-three folio volumes, between the years 1643 and 1794 (including the two volumes of the second series), under the direction of the Jesuits, and with the richest and rarest literary aids.\footnote{When Rosweyd’s prospectus, which contemplated only 17 volumes, was shown to Cardinal Bellarmine, he asked: “What is the man’s age?” “Perhaps forty.” “Does he expect to live two hundred years?” More than 250 years have passed since, and still the work is unfinished. The relation of the principal authors is indicated in the following verse: \par “Quod Rosweydus praepararat, Quod Bollandus inchoarat, Quod Henschenius formarat, Perfecit (?) Papenbroekius.”} This work contains, in the order of the days of the year, the biography of every saint in the Catholic calendar, as composed by the Bollandists, down to the fifteenth of October, together with all the acts of canonization, papal bulls, and other ancient documents belonging thereto, with learned treatises and notes; and that not in the style of popular legends, but in the tone of thorough historical investigation and free criticism, so far as a general accordance with the Roman Catholic system of faith would allow.\footnote{The work was even violently persecuted at times in the Romish Church. Papenbroek, for proving that the prophet Elijah was not the founder of the Carmelite order, was stigmatized as a heretic, and the Acta condemned by the Spanish Inquisition, but the condemnation was removed by papal interference in 1715. The Bollandists took holy revenge of the Carmelites by a most elaborate biography and vindication of St. Theresa, the glory of that order, in the fifty-fourth volume (the first of the new series), 1845, sub Oct. 15th, pp. 109-776.} It was interrupted in 1773 by the abolition of the order of the Jesuits, then again in 1794, after a brief resumption of labor and the publication of two more volumes (the fifty-second and fifty-third), by the French Revolution and invasion of the Netherlands and the partial destruction of the literary, material; but since 1845 (or properly since 1837) it has been resumed at Brussels under the auspices of the same order, though not with the same historical learning and critical acumen, and proceeds tediously toward completion.\footnote{The names connected with the new (third) series are Joseph van der Moere, Joseph van Hecke, Bossue, Buch, Tinnebroek, etc. By 1858 five new folio volumes had appeared at Brussels (to the twenty-second of October), so that the whole work now embraces fifty-eight volumes, which cost from two thousand four hundred to three thousand francs. The present Bollandist library is in the convent of St. Michael in Brussels and embraces in three rooms every known biography of a saint, hundreds of the rarest missals and breviaries, hymnals and martyrologies, sacramentaries and rituals. A not very correct reprint of the Antwerp original has appeared at Venice since 1734. A new edition by Jo. Carnandet is now coming out at Paris and Rome, 1863 sqq. Complete copies have become very rare. I have seen and used at different times three copies, one in the Theol. Seminary Library at Andover, and two at New York (in the Astor Library, and in the Union Theol. Sem. Library). Comp. the Prooemium de ratione universa operis, in the Acta Sanctorum, vol. vi. for Oct. (published 1845). R. P. Dom Pitra: Etudes sur la Collection des Actes des Saintes, par les RR. PP. Jusuites Bollandistes. Par. 1850. Also an article on the Bollandists by J. M. Neale in his Essays on Liturgiology and Church History, Lond. 1863, p. 89 ff.} This colossal and amazing work of more than two centuries of pious industry and monkish learning will always remain a rich mine for the system of martyr and saint-worship and the history of Christian life.

\xsection{Worship of Relics. Dogma of the Resurrection. Miracles of Relics}

§ 87. Worship of Relics. Dogma of the Resurrection. Miracles of Relics.

Comp. the Literature at § 84. Also J. Mabillon (R.C.): Observationes de sanctorum reliquiis (Praef. ad Acta s. Bened. Ordinis). Par. 1669. Barrington and Kirk (R.C.): The Faith of Catholics, \&c. Lond. 1846. Vol. iii. pp. 250–307. On the Protestant side, J. H. Jung: Disquisitio antiquaria de reliqu. et profanis et sacris earumque cultu, ed. 4. Hannov. 1783.

The veneration of martyrs and saints had respect, in the first instance, to their immortal spirits in heaven, but came to be extended, also, in a lower degree, to their earthly remains or relics.\footnote{Reliquiae, and reliqua, \textgreek{λείψανα}.} By these are to be understood, first, their bodies, or rather parts of them, bones, blood, ashes; then all which was in any way closely connected with their persons, clothes, staff, furniture, and especially the instruments of their martyrdom. After the time of Ambrose the cross of Christ also, which, with the superscription and the nails, are said to have been miraculously discovered by the empress Helena in 326,\footnote{The legend of the “invention of the cross” (inventio s. crucis), which is celebrated in the Greek and Latin churches by a special festival, is at best faintly implied in Eusebius in a letter of Constantineto the bishop Macarius of Jerusalem (Vita Const. iii. 30—a passage which Gieseler overlooked—though in iii. 25, where it should be expected, it is entirely unnoticed, as Gieseler correctly observes), and does not appear till several decennia later, first in Cyril of Jerusalem (whose Epist. ad Constantium of 351, however, is considered by Gieseler and others, on critical and theological grounds, a much later production), then, with good agreement as to the main fact, in Ambrose, Chrysostom, Paulinus of Nola, Socrates, Sozomen, Theodoret, and other fathers. With all these witnesses the fact is still hardly credible, and has against it particularly the following considerations: (1) The place of the crucifixion was desecrated under the emperor Hadrian by heathen temples and statues, besides being filled up and defaced beyond recognition. (2) There is no clear testimony of a \emph{contemporary}. (3) The pilgrim from Bordeaux, who visited Jerusalem in 333, and in a still extant \emph{itinerarium} (Vetera Rom. itineraria, ed. P. Wesseling, p. 593) enumerates all the sacred things of the holy city, knows nothing of the holy cross or its Invention (comp. Gieseler, i. 2, p. 279, note 37; Edinb. ed. vol. ii. p. 36). This miracle contributed very much to the increase of the superstitious use of crosses and crucifixes. Cyril of Jerusalem remarks that about 380 the splinters of the holy cross filled the whole world, and yet, according to the account of the devout but credulous Paulinus of Nola (Epist. 31, al. 11), the original remained in Jerusalem undiminished,—a continual miracle! Besides Gieseler, comp. particularly the minute investigation of this legend by Isaac Taylor, The Invention of the Cross and the Miracles therewith connected, in “Ancient Christianity,” vol. ii. pp. 277-315.} was included, and subsequently His crown of thorns and His coat, which are preserved, the former, according to the legend, in Paris, and the latter in Treves.\footnote{Comp. Gildemeister: Der heil. Rock von Trier, 2d ed. 1845—a controversial work called forth by the Ronge excitement in German Catholicism in 1844.} Relics of the body of Christ cannot be thought of, since He arose without seeing corruption, and ascended to heaven, where, above the reach of idolatry and superstition, He is enthroned at the right hand of the Father. His true relics are the Holy Supper and His living presence in the church to the end of the world.

The worship of relics, like the worship of Mary and the saints, began in a sound religious feeling of reverence, of love, and of gratitude, but has swollen to an avalanche, and rushed into all kinds of superstitious and idolatrous excess. “The most glorious thing that the mind conceives,” says Goethe, “is always set upon by a throng of more and more foreign matter.”

As Israel could not sustain the pure elevation of its divinely revealed religion, but lusted after the flesh pots of Egypt and coquetted with sensuous heathenism so it fared also with the ancient church.

The worship of relics cannot be derived from Judaism; for the Levitical law strictly prohibited the contact of bodies and bones of the dead as defiling.\footnote{Num. xix. 11 ff.; xxxi. 19. The touching of a corpse or a dead bone, or a grave, made one unclean seven days, and was to be expiated by washing, upon pain of death. The tent, also, in which a person had died, and all open vessels in it, were unclean. Comp. Josephus, c. Apion. ii. 26; Antiqu. iii. 11, 3. The Talmudists made the laws still more stringent on this point.} Yet the isolated instance of the bones of the prophet Elisha quickening by their contact a dead man who was cast into his tomb,\footnote{2 Kings xiii. 21 (Sept.): \textgreek{ἥψατο τῶν ὀστῶν Ἑλισαιέ, καὶ ἔζησε καὶ ἔστη ἐπὶ τοὺς πόδας}Comp. the apocryphal book Jesus Sirach (Ecclesiasticus) xlviii. 13, 14; xlix. 12.} was quoted in behalf of the miraculous power of relics; though it should be observed that even this miracle did not lead the Israelites to do homage to the bones of the prophet nor abolish the law of the uncleanness of a corpse.

The heathen abhorred corpses, and burnt them to ashes, except in Egypt, where embalming was the custom and was imitated by the Christians on the death of martyrs, though St. Anthony protested against it. There are examples, however, of the preservation of the bones of distinguished heroes like Theseus, and of the erection of temples over their graves.\footnote{Plutarch, in his Life of Theseus, c. 86.}

The Christian relic worship was primarily a natural consequence of the worship of the saints, and was closely connected with the Christian doctrine of the resurrection of the body, which was an essential article of the apostolic tradition, and is incorporated in almost all the ancient creeds. For according to the gospel the body is not an evil substance, as the Platonists, Gnostics, Manichaeans held, but a creature of God; it is redeemed by Christ; it becomes by the regeneration an organ and temple of the Holy Ghost; and it rests as a living seed in the grave, to be raised again at the last day, and changed into the likeness of the glorious body of Christ. The bodies of the righteous “grow green” in their graves, to burst forth in glorious bloom on the morning of the resurrection. The first Christians from the beginning set great store by this comforting doctrine, at which the heathen, like Celsus and Julian, scoffed. Hence they abhorred also the heathen custom of burning, and adopted the Jewish custom of burial with solemn religious ceremonies, which, however, varied in different times and countries.

But in the closer definition of the dogma of the resurrection two different tendencies appeared: a spiritualistic, represented by the Alexandrians, particularly by Origen and still later by the two Gregories; the other more realistic, favored by the Apostles’ Creed,\footnote{In the phrase \textgreek{ἀνάστασις τῆς σαρκός}, instead of \textgreek{τοῦ σώματος}, resurrectio \emph{carnis}, instead of \emph{corporis.} The Nicene creed uses the expression \textgreek{ἀνάστασις νεκρῶν}, resurrectio \emph{mortuorum}. In the German version of the Apostles’ Creed the easily mistaken term \emph{Fleisch, flesh,} is retained; but the English churches say more correctly: resurrection of the \emph{body.}} advocated by Tertullian, but pressed by some church teachers, like Epiphanius and Jerome, in a grossly materialistic manner, without regard to the σῶμα πευματικόν of Paul and the declaration that “flesh and blood cannot inherit the kingdom of God.”\footnote{Jerome, on the ground of his false translation of Job xix. 26, teaches even the restoration of all bones, veins, nerves, teeth, and hair (because the Bible speaks of gnashing of teeth among the damned, and of the hairs our heads being all numbered!). “Habent dentes,” says he of the resurrection bodies, “ventrem, genitalia, et tamen nec cibis nec uxoribus indigent.” Augustineis more cautious, and endeavors to avoid gross, carnal conceptions. Comp. the passages in Hagenbach’s Dogmengeschichte, i. § 140 (Engl. ed., New York, i. p. 370 ff.).} The latter theory was far the more consonant with the prevailing spirit of our period, entirely supplanted the other, and gave the mortal remains of the saints a higher value, and the worship of them a firmer foundation.

Roman Catholic historians and apologists find a justification of the worship and the healing virtue of relics in three facts of the New Testament: the healing of the woman with the issue of blood by the touch of Jesus’ garment;\footnote{Matt. ix. 20.} the healing of the sick by the shadow of Peter;\footnote{Acts v. 14, 15.}and the same by handkerchiefs from Paul.\footnote{Acts xix. 11, 12.}

These examples, as well as the miracle wrought by the bones of Elisha, were cited by Origen, Cyril of Jerusalem, Ambrose, Chrysostom, and other fathers, to vindicate similar and greater miracles in their time. They certainly mark the extreme limit of the miraculous, beyond which it passes into the magical. But in all these cases the living and present person was the vehicle of the healing power; in the second case Luke records merely the popular belief, not the actual healing; and finally neither Christ nor the apostles themselves chose that method, nor in any way sanctioned the superstitions on which it was based.\footnote{On the contrary, the account of the healing of sick by the handkerchiefs of Paul is immediately followed by an account of the magical abuse of the name of Jesus, as a warning, Acts xix. 13 ff.} At all events, the New Testament and the literature of the apostolic fathers know nothing of an idolatrous veneration of the cross of Christ or the bones and chattels of the apostles. The living words and acts of Christ and the apostles so completely absorbed attention that we have no authentic accounts of the bodily appearance, the incidental externals, and transient possessions of the founders of the church. Paul would know Christ after the spirit, not after the flesh. Even the burial places of most of the apostles and evangelists are unknown. The traditions of their martyrdom and their remains date from a much later time, and can claim no historical credibility.

The first clear traces of the worship of relics appear in the second century in the church of Antioch, where the bones of the bishop and martyr Ignatius († 107) were preserved as a priceless treasure;\footnote{\textgreek{θησαυρὸς ἀτίμητος}. Martyr. S. Ignat. cap. vii. (Patrum Apostolic. Opera, ed. Dressel, p. 214). The genuineness of the Martyr-Acts of Ignatius, however, is disputed by many.} and in Smyrna, where the half-burnt bones of Polycarp († 167) were considered “more precious than the richest jewels and more tried than gold.”\footnote{Τ\textgreek{ὰ τιμιώτερα λίθων πολυτελῶν καὶ δοκιμώτερα ὑπὲρ χρυσίον ὀστᾶ αὐτοῦ}, Epist. Eccl. Smyrn. de Martyr. S. Polyc. c. 18 (ed. Dressel, p. 404), and in Euseb. H. E. iv. 15.} We read similar things in the Acts of the martyrs Perpetua and Cyprian. The author of the Apostolic Constitutions\footnote{Const. Apost. lib. vi. c. 30. The sixth book dates from the end of the third century.} exhorts that the relics of the saints, who are with the God of the living and not of the dead, be held in honor, and appeals to the miracle of the bones of Elisha, to the veneration which Joseph showed for the remains of Jacob, and to the bringing of the bones of Joseph by Moses and Joshua into the promised land.\footnote{Comp. Gen. l. 1, 2, 25, 26; Ex. xiii. 19; Jos. xxiv. 32; Acts vii. 16.} Eusebius states that the episcopal throne of James of Jerusalem was preserved to his time, and was held in great honor.\footnote{Hist. Ecel. vii. 19 and 32.}

Such pious fondness for relics, however, if it is confined within proper limits, is very natural and innocent, and appears even in the Puritans of New England, where the rock in Plymouth, the landing place of the Pilgrim Fathers in 1620, has the attraction of a place of pilgrimage, and the chair of the first governor of Massachusetts is scrupulously preserved, and is used at the inauguration of every new president of Harvard University.

But toward the middle of the fourth century the veneration of relics simultaneously with the worship of the saints, assumed a decidedly superstitious and idolatrous character. The earthly remains of the martyrs were discovered commonly by visions and revelations, often not till centuries after their death, then borne in solemn processions to the churches and chapels erected to their memory, and deposited under the altar;\footnote{With reference to Rev. vi. 9: “I saw under the altar (\textgreek{ὑποκάτω τοῦ θυσιαστηρίου}) the souls of them that were slain for the word of God,” \&c.} and this event was annually celebrated by a festival.\footnote{\emph{Festum translationis}.} The legend of the discovery of the holy cross gave rise to two church festivals: The Feast of the Invention of the Cross\footnote{\emph{Festum inventionis s. crucis}.} on the third of May, which has been observed in the Latin church since the fifth or sixth century; and The Feast of the Elevation of the Cross,\footnote{\emph{Festum exaltationis s. crucis}, \textgreek{σταυροφανεία}.} on the fourteenth of September, which has been observed in the East and the West, according to some since the consecration of the church of the Holy Sepulchre in 335, according to others only since the reconquest of the holy cross by the emperor Heraclius in 628. The relics were from time to time displayed to the veneration of the believing multitude, carried about in processions, preserved in golden and silver boxes, worn on the neck as amulets against disease and danger of every kind, and considered as possessing miraculous virtue, or more strictly, as instruments through which the saints in heaven, in virtue of their connection with Christ, wrought miracles of healing and even of raising the dead. Their number soon reached the incredible, even from one and the same original; there were, for example, countless splinters of the pretended cross of Christ from Jerusalem, while the cross itself is said to have remained, by a continued miracle, whole and undiminished! Veneration of the cross and crucifix knew no bounds, but can, by no means, be taken as a true measure of the worship of the Crucified; on the contrary, with the great mass the outward form came into the place of the spiritual intent, and the wooden and silver Christ was very often a poor substitute for the living Christ in the heart.\footnote{What Luther says of the “juggleries and idolatries” of the cross under the later papacy, which “would rather bear the cross of Christ in silver, than in heart and life,” applies, though, of course, with many noble exceptions, even to the period before us. Dr. Herzog, in his Theol. Encyclopaedia, vol. viii. p. 60 f., makes the not unjust remark: “The more the cross came into use in manifold forms and signs, the more the truly evangelical faith in Christ, the Crucified, disappeared. The more the cross of Christ was outwardly exhibited, the more it became inwardly an offence and folly to men. The Roman Catholic church in this respect resembles those Christians, who talk so much of their spiritual experiences, make so much ado about them that they at last talk themselves out, and produce glittering nonsense.”}

Relics became a regular article of trade, but gave occasion, also, for very many frauds, which even such credulous and superstitious relic-worshippers as St. Martin of Tours\footnote{Sulpit. Severus, Vita beati Mart. c. 11.} and Gregory the Great\footnote{Epist. lib. iv. Ep. 30. Gregory here relates that some Greek monks came to Rome to dig up bones near St. Paul’s church to sell, as they themselves confessed, for holy relics in the East (confessi sunt, quod illa ossa ad Graeciam essent tamquam Sanctorum reliquias portaturi).} lamented. Theodosius I., as early as 386, prohibited this trade; and so did many councils; but without success. On this account the bishops found themselves compelled to prove the genuineness of the relics by historical tradition, or visions, or miracles.

At first, an opposition arose to this worship of dead men’s bones. St. Anthony, the father of monasticism († 356), put in his dying protest against it, directing that his body should be buried in an unknown place. Athanasius relates this with approbation,\footnote{In his Vita Antoini, Opera Athan. ii. 502.} and he caused several relics which had been given to him to be fastened up, that they might be out of the reach of idolatry.\footnote{Rufinus, Hist. Ecel. ii. 28.} But the opposition soon ceased, or became confined to inferior or heretical authors, like Vigilantius and Eunomius, or to heathen opponents like Porphyry and Julian. Julian charges the Christians, on this point, with apostasy from their own Master, and sarcastically reminds them of His denunciation of the Pharisees, who were like whited sepulchres, beautiful without, but within full of dead men’s bones and all uncleanness.\footnote{Cyrillus Alex. Adv. Jul. l. x. tom. vi. p. 356.} This opposition, of course, made no impression, and was attributed to sheer impiety. Even heretics and schismatics, with few exceptions, embraced this form of superstition, though the Catholic church denied the genuineness of their relics and the miraculous virtue of them

The most and the best of the church teachers of our period, Hilary, the two Gregories, Basil, Chrysostom, Isidore of Pelusium, Theodoret, Ambrose, Jerome, Augustine, and Leo, even those who combated the worship of images on this point, were carried along by the spirit of the time, and gave the weight of their countenance to the worship of relics, which thus became an essential constituent of the Greek and Roman Catholic religion. They went quite as far as the council of Trent,\footnote{Sessio xxv. De Invocat. Sanct., etc.} which expresses itself more cautiously, on the worship of relics as well as of saints, than the church fathers of the Nicene age. With the good intent to promote popular piety by sensible stimulants and tangible supports, they became promoters of dangerous errors and gross superstition.

To cite some of the most important testimonies:

Gregory Nazianzen thinks the bodies of the saints can as well perform miracles, as their spirits, and that the smallest parts of the body or of the symbols of their passion are as efficacious as the whole body.\footnote{Adv. Julian. t. i. Orat. iii. p. 76 sq.}

Chrysostom values the dust and ashes of the martyrs more highly than gold or jewels, and ascribes to them the power of healing diseases and putting death to flight.\footnote{Opera, tom. ii. p. 828.} In his festal discourse on the translation of the relics of the Egyptian martyrs from Alexandria to Constantinople, he extols the bodies of the saints in eloquent strains as the best ramparts of the city against all visible enemies and invisible demons, mightier than walls, moats, weapons, and armies.\footnote{Hom. in MM. Aegypt. tom. ii. p. 834 sq.}

“Let others,” says Ambrose, “heap up silver and gold; we gather the nails wherewith the martyrs were pierced, and their victorious blood, and the wood of their cross.”\footnote{Exhort. virgin. 1.} He himself relates at large, in a letter to his sister, the miraculous discovery of the bones of the twin brothers Gervasius and Protasius, two otherwise wholly unknown and long-forgotten martyrs of the persecution under Nero or Domitian.\footnote{Epist. xxii. Sorori suae, Op. ii. pp. 874-878. Comp. Paulinus, Vit. Ambros. p. iv.; Paulinus Nol. Ep. xii. ad Severum; and Augustinein sundry places (see below).} This is one of the most notorious relic miracles of the early church. It is attested by the most weighty authorities, by Ambrose and his younger contemporaries, his secretary and biographer Paulinus, the bishop Paulinus of Nola, and Augustine, who was then in Milan; it decided the victory of the Nicene orthodoxy over the Arian opposition of the empress Justina; yet is it very difficult to be believed, and seems at least in part to rest on pious frauds.\footnote{Clericus, Mosheim, and Isaac Taylor (vol. ii. p. 242 ff.) do not hesitate to charge St. Ambrose, the author of the Te Deum, with fraud in this story. The latter, however, endeavors to save the character of Ambroseby distinguishing between himself and the spirit of his age. “Ambrose,” says he (ii. 270), “occupies a high position among the Fathers; and there was a vigor and dignity in his character, as well as a vivid intelligence, which must command respect; but in proportion as we assign praise to the man, individually, we condemn the system which could so far vitiate a noble mind, and impel one so lofty in temper to act a part which heathen philosophers would utterly have abhorred.”}

The story is, that when Ambrose, in 386, wished to consecrate the basilica at Milan, he was led by a higher intimation in a vision to cause the ground before the doors of Sts. Felix and Nahor to be dug up, and there he found two corpses of uncommon size, the heads severed from the bodies (for they died by the sword), the bones perfectly preserved, together with a great quantity of fresh blood.\footnote{“Invenimus mirae magnitudinis viros duos, ut prisca aetas ferebat, ossa omnia integra, sanguinis plurimum! ” Did Ambrosereally believe that men in the first century (prisca aetas) were of greater bodily stature than his contemporaries in the fourth? But especially absurd is the mass of fresh blood, which then was exported throughout Christendom as a panacea. According to Romish tradition, the blood of many saints, as of Januarius in Naples, becomes liquid every year. Taylor the miraculously healed Severus, by trade a butcher, had something to do with this blood.} These were the saints in question. They were exposed for two days to the wondering multitude, then borne in solemn procession to the basilica of Ambrose, performing on the way the healing of a blind man, Severus by name, a butcher by trade, and afterward sexton of this church. This, however, was not the only miracle which the bones performed. “The age of miracles returned,” says Ambrose. “How many pieces of linen, how many portions of dress, were cast upon the holy relics and were recovered with the power of healing from that touch.\footnote{“Et tactu ipso medicabilia reposcuntur.”} It is a source of joy to all to touch but the extremest portion of the linen that covers them; and whoso touches is healed. We give thee thanks, O Lord Jesus, that thou hast stirred up the energies of the holy martyrs at this time, wherein thy church has need of stronger defence. Let all learn what combatants I seek, who are able to contend for us, but who do not assail us, who minister good to all, harm to none.” In his homily De inventione SS. Gervasii et Protasii, he vindicates the miracle of the healing of the blind man against the doubts of the Arians, and speaks of it as a universally acknowledged and undeniable fact: The healed man, Severus, is well known, and publicly testifies that he received his sight by the contact of the covering of the holy relics.

Jerome calls Vigilantius, for his opposition to the idolatrous veneration of ashes and bones, a wretched man, whose condition cannot be sufficiently pitied, a Samaritan and Jew, who considered the dead unclean; but he protects himself against the charge of superstition. We honor the relics of the martyrs, says he, that we may adore the God of the martyrs; we honor the servants, in order thereby to honor the Master, who has said: “He that receiveth you, receiveth me.”\footnote{Ep. cix. ad Riparium.} The saints are not dead; for the God of Abraham, Isaac, and Jacob is not a God of the dead, but of the living. Neither are they enclosed in Abraham’s bosom as in a prison till the day of Judgment, but they follow the Lamb whithersoever he goeth.\footnote{Adv. Vigil.c. 6.}

Augustine believed in the above-mentioned miraculous discovery of the bodies of Gervasius and Protasius, and the healing of the blind man by contact with them, because he himself was then in Milan, in 386, at the time of his conversion,\footnote{Cum illic—Mediolani—essemus.} and was an eye-witness, not indeed of the discovery of the bones—for this he nowhere says—but of the miracles, and of the great stir among the people.\footnote{He speaks of this four times clearly and plainly, Confess. ix. 7; De Civit Dei, xxii. 8; Serm. 286 in Natali MM. Protasii et Gervasii; Retract. i. 13, § 7.}

He gave credit likewise to the many miraculous cures which the bones of the first martyr Stephen are said to have performed in various parts of Africa in his time.\footnote{Serm. 317 and 318 de Martyr. Steph. Is. Taylor (l.c. ii. pp. 316-350) has thoroughly investigated the legend of the relics of the proto-martyr, and comes to the conclusion that it likewise rests on pious frauds which Augustinehonestly believed.} These relics were discovered in 415, nearly four centuries after the stoning of Stephen, in an obscure hamlet near Jerusalem, through a vision of Gamaliel, by a priest of Lucian; and some years afterward portions of them were transported to Uzali, not far from Utica, in North Africa, and to Spain and Gaul, and everywhere caused the greatest ado in the superstitious populace.

But Augustine laments, on the other hand, the trade in real and fictitious relics, which was driven in his day,\footnote{De opere Monachorum, c. 28: “Tam multos hypocrites sub habitu monachomm [hostis] usquequoque dispersit, circumeuntes provincias, nusquam missos, nusquam fixos, nusquam stantes, nusquam sedentes. Alii membra martyrum, el tamen martyrum, venditant.” Augustinerejects the pretended miracles of the Donatists, and calls them wonderlings (mirabiliarii), who are either deceivers or deceived (In Joann. evang. Tract. xiii. § 17).} and holds the miracles to be really superfluous, now that the world is converted to Christianity, so that he who still demands miracles, is himself a miracle.\footnote{De Civit. Dei, xxii. c. 8: “Cur, inquiunt, nunc illa miracula, quae praedicatis facta esse, non fiunt? Possem quidem dicere, necessaria fuisse priusquam crederet mundus, ad hoc ut crederet mundus. Quisvis adhuc prodigia ut credat inquirit, magnum est ipse prodigium, qui mundo credente non credit.” Comp. De util. cred. c. 25, § 47; c. 50, § 98; De vera relig. c. 25, § 47.} Though he adds, that to that day miracles were performed in the name of Jesus by the sacraments or by the saints, but not with the same lustre, nor with the same significance and authority for the whole Christian world.\footnote{Ibid.: “Nam etiam nunc fiunt miracula in ejus nomine, Sive per sacramenta ejus, Sive per orationes vel memorias sanctorum ejus; sed non eadem claritate illustrantur, ut tanta quanta illa gloria diffamentur .... Nam plerumque etiam ibi [in the place where these miracles were wrought] \emph{paucissimi} sciunt, ignorantibus caeteris, maxime si magna sit civitas; et quando alibi aliisque narrantur, \emph{non tanta ea commendat auctoritas, ut sine difficultate vel dubitatione credantur}, quamvis Christianis fidelibus a fidelibus indicentur.” Then follows the account of the famous \emph{miraculum Protasii et Gervasii}, and of several cures in Carthage and Hippo. Those in Hippo were wrought by the relics of St. Stephen, and formally confirmed.} Thus he himself furnishes a warrant and an entering wedge for critical doubt in our estimate of those phenomena.\footnote{Comp. Fr. Nitzsch(jun.): Augustinus’ Lehre vom Wunder, Berlin, 1865, especially pp. 82-45. (A very full and satisfactory treatise.)}

\xsection{Observations on the Miracles of the Nicene Age}

§ 88. Observations on the Miracles of the Nicene Age.

Comp. on the affirmative side especially John H. Newman (now R.C., then Romanizing Anglican): Essay on Miracles, in the 1st vol. of the English translation of Fleury’s Ecclesiastical History, Oxford, 1842; on the negative, Isaac Taylor (Independent): Ancient Christianity, Lond. 4th ed. 1844. Vol. ii. pp. 233–365. Dr. Newman previously took the negative side on the question of the genuineness of the church miracles in a contribution to the Encyclopaedia Metropolitana, 1830.

In the face of such witnesses as Ambrose and Augustine, who must be accounted in any event the noblest and most honorable men of the early church, it is venturesome absolutely to deny all the relic-miracles, and to ascribe them to illusion and pious fraud. But, on the other hand, we should not be bribed or blinded by the character and authority of such witnesses, since experience sufficiently proves that even the best and most enlightened men cannot wholly divest themselves of superstition and of the prejudices of their age\footnote{Recall, e.g., Luther and the apparitions of the devil, the Magnalia of Cotton Mather, the old Puritans and their trials for witchcraft, as well as the modem superstitions of spiritual rappings and table-turnings by which many eminent and intelligent persons have been carried along.} Hence, too, we should not ascribe to this whole question of the credibility of the Nicene miracles an undue dogmatic weight, nor make the much wider issue between Catholicism and Protestantism dependent on it.\footnote{As is done by many Roman Catholic historians and apologists in the cause of Catholicism, and by Isaac Taylor in the interest of Protestantism. The latter says in his oft-quoted work, vol. ii. p. 239: “The question before us [on the genuineness of the Nicene miracles] is therefore in the strictest sense \emph{conclusive} as to the modem controversy concerning church principles and the authority of tradition. If the miracles of the fourth century, and those which follow in the same track, were real, then Protestantism is altogether indefensible, and ought to be denounced as an impiety of the most flagrant kind. But if these miracles were wicked frauds; and if they were the first series of a system of impious delusion—then, not only is the modern Papacy to be condemned, but the church of the fourth century must be condemned with it; and for the same reasons; and the Reformation is to be adhered to as the emancipation of Christendom from the thraldom of him who is the ’father of lies.’ ” Taylor accordingly sees in the old Catholic miracles sheer lying wonders of Satan, and signs of the apostasy of the church predicted in the Epistles of St. Paul. From the same point of view he treats also the phenomena of asceticism and monasticism, putting them with the unchristian hatred of the creature and the ascription of nature to the devil, which characterized the Gnostics. But he thus involves not only the Nicene age, but the ante-Nicene also, up to Irenaeus and Ignatius, in this apostasy, and virtually gives up the unbroken continuity of true Christianity. He is, moreover, not consistent in making the church fathers, on the one hand, the chief originators of monkish asceticism and false miracles, while, on the other hand, he sincerely reveres them and eloquently lauds them for their Christian earnestness and their immortal services. Comp. his beautiful concession in vol. i. p. 37 (cited in the 1st vol. of this Hist § 46, note 2).} In every age, as in every man, light and shade in fact are mingled, that no flesh should exalt itself above measure. Even the most important periods of church history, among which the Nicene age, with all its faults, must be numbered, have the heavenly treasure in earthen vessels, and reflect the spotless glory of the Redeemer in broken colors.

The most notorious and the most striking of the miracles of the fourth century are Constantine’s vision of the cross (a.d. 312), the finding of the holy cross (a.d. 326), the frustration of Julian’s building of the temple (a.d. 363) the discovery of the relics of Protasius and Gervasius (a.d. 386), and subsequently (a.d. 415) of the bones of St. Stephen, with a countless multitude of miraculous cures in its train. Respecting the most important we have already spoken at large in the proper places.

We here offer some general remarks on this difficult subject.

The possibility of miracles in general he only can deny who does not believe in a living God and Almighty Maker of heaven and earth. The laws of nature are organs of the free will of God; not chains by which He has bound Himself forever, but elastic threads which He can extend and contract at His pleasure. The actual occurrence of miracles is certain to every believer from Holy Scripture, and there is no passage in the New Testament to limit it to the apostolic age. The reasons which made miracles necessary as outward proofs of the divine mission of Christ and the apostles for the unbelieving Jews of their time, may reappear from time to time in the unbelieving heathen and the skeptical Christian world; while spiritual miracles are continually taking place in regeneration and conversion. In itself, it is by no means unworthy and incredible that God should sometimes condescend to the weakness of the uneducated mass, and should actually vouchsafe that which was implored through the mediation of saints and their relics.

But the following weighty considerations rise against the miracles of the Nicene and post-Nicene age; not warranting, indeed, the rejection of all, yet making us at least very cautious and doubtful of receiving them in particular:

1. These miracles have a much lower moral tone than those of the Bible, while in some cases they far exceed them in outward pomp, and make a stronger appeal to our faculty of belief. Many of the monkish miracles are not so much supernatural and above reason, as they are unnatural and against reason, attributing even to wild beasts of the desert, panthers and hyenas, with which the misanthropic hermits lived on confidential terms, moral feelings and states, repentance and conversion\footnote{Comp. the examples quoted in § 34, p. 177 f.} of which no trace appears in the New Testament.\footnote{The speaking serpent in Paradise (Gen. iii.), and the speaking ass of Balaam (Num. xxii. 22-33; comp. 2 Pet. ii. 16), can hardly be cited as analogies, since in those cases the irrational beast is merely the organ of a moral power foreign to him.}

2. They serve not to confirm the Christian faith in general, but for the most part to support the ascetic life, the magical virtue of the sacrament, the veneration of saints and relics, and other superstitious practices, which are evidently of later origin, and are more or less offensive to the healthy evangelical mind.\footnote{Is. Taylor, l.c. vol. ii. p. 235, says of the miracles of the Nicene age: “These alleged miracles were, \emph{almost in} \emph{every instance,} wrought expressly in support of those very practices and opinions which stand forward as the points of contrast, distinguishing Romanism from Protestantism ... the supernatural properties of the eucharistic elements, the invocation of saints, or direct praying to them, and the efficacy of their relics; and the reverence or worship due to certain visible and palpable religious symbols.” Historical questions, however, should be investigated and decided with all possible freedom from confessional prejudices.}

3. The further they are removed from the apostolic age, the more numerous they are, and in the fourth century alone there are more miracles than in all the three preceding centuries together, while the reason for them, as against the power of the heathen world, was less.

4. The church fathers, with all the worthiness of their character in other respects, confessedly lacked a highly cultivated sense of truth, and allowed a certain justification of falsehood ad majorem Dei gloriam, or fraus pia, under the misnomer of policy or accommodation;\footnote{So especially Jerome, Epist ad Pammachium (Lib. apologeticus pro libris contra Jovinianum Ep. xlviii. c. 12, ed. Vallarsi tom i. 222, or Ep. xix. in the Benedictine ed.): “Plura esse genera dicendi: et inter caetera, aliud esse \textgreek{γυμναστικῶς}scribere, aliud \textgreek{δογματικῶς.} In priori vagam esse disputationem; et adversario respondentem, nunc haec nunc illa proponere, argumentari ut libet, aliud loqui, aliud agere, panem, ut dicitur, ostendere, lapidem tenere. In sequenti autem aparta frons et, ut ita dicam, ingenuitas necessaria est. Aliud est quaerere, aliud definire. In altero pugnandum, in altero docendum est.” He then appeals to the Greek and Roman classics, the ancient fathers in their polemical writings, and even St. Paul in his arguments from the Old Testament. Of interest in this connection is his controversy with Augustineon the conduct of Paul toward Peter, Gal. ii. 11, which Jeromewould attribute to mere policy or accommodation. Even Chrysostomutters loose principles on the duty of veracity (De sacerdot. i. 5), and his pupil Cassian still more, appealing to the example of Rahab (Coll. xvii. 8, 17, etc.). Comp. Gieseler, i. ii. p. 307 (§ 102, note 17). The corrupt principle that “the end sanctifies the means,” is much older than Jesuitism, which is commonly made responsible for it. Christianity had at that time not yet wholly overcome the spirit of falsehood in ancient heathenism.} with the solitary exception of Augustine, who, in advance of his age, rightly condemned falsehood in every form.

5. Several church fathers like Augustine, Martin of Tours, and Gregory I., themselves concede that in their time extensive frauds with the relics of saints were already practised; and this is confirmed by the fact that there were not rarely numerous copies of the same relics, all of which claimed to be genuine.

6. The Nicene miracles met with doubt and contradiction even among contemporaries, and Sulpitius Severus makes the important admission that the miracles of St. Martin were better known and more firmly believed in foreign countries than in his own.\footnote{Dialog. i. 18.}

7. Church fathers, like Chrysostom and Augustine, contradict themselves in a measure, in sometimes paying homage to the prevailing faith in miracles, especially in their discourses on the festivals of the martyrs, and in soberer moments, and in the calm exposition of the Scriptures, maintaining that miracles, at least in the Biblical sense, had long since ceased.\footnote{This argument is prominently employed by James Craigie Robertson (moderate Anglican): History of the Christian Church to Gregory the Great, Lond. 1854, p. 334. “On the subject of miracles,” says he, “there is a remarkable inconsistency in the statements of writers belonging to the end of the fourth and beginning of the fifth centuries. St. Chrysostomspeaks of it as a notorious and long-settled fact that miracles had ceased (v. Newman, in Fleury, vol. i. p. xxxix). Yet at that very time, St. Martin, St. Ambrose, and the monks of Egypt and the East are said to have been in full thaumaturgical activity; and Sozomen (viii. 5) tells a story of a change of the eucharistic bread into a stone as having happened at Constantinople, while Chrysostomhimself was bishop. So again, St. Augustinesays that miracles such as those of Scripture were no longer done, yet he immediately goes on to reckon up a number of miracles which had lately taken place, apparently without exciting much sensation, and among them \emph{seventy} formally attested ones, wrought at Hippo alone, within two years, by the relics of St. Stephen (De Civit. Dei, xxii. 8. 1, 20). On the whole, while I would not deny that miracles may have been wrought after the times of the apostles and their associates, I can find very little satisfaction in the particular instances which are given.” On Augustine’s theory of miracles, comp. above, § 87 (p. 459 f.), and the treatise of Nitzsch jun. there quoted.}

We must moreover remember that the rejection of the Nicene miracles by no means justifies the inference of intentional deception in every case, nor destroys the claim of the great church teachers to our respect. On the contrary, between the proper miracle and fraud there lie many intermediate steps of self-deception, clairvoyance, magnetic phenomena and cures, and unusual states of the human soul, which is full of deep mysteries, and stands nearer the invisible spirit-world than the everyday mind of the multitude suspects. Constantine’s vision of the cross, for example, may be traced to a prophetic dream;\footnote{Comp. above, § 2 (p.25).} and the frustration of the building of the Jewish temple under Julian, to a special providence, or a historical judgment of God.\footnote{Comp. above, § 4 (p. 55).} The mytho-poetic faculty, too, which freely and unconsciously produces miracles among children, may have been at work among credulous monks in the dreary deserts and magnified an ordinary event into a miracle. In judging of this obscure portion of the history of the church we must, in general, guard ourselves as well against shallow naturalism and skepticism, as against superstitious mysticism, remembering that

\xsection{Processions and Pilgrimages}

§ 89. Processions and Pilgrimages.

Early Latin dissertations on pilgrimages by J. Gretser, Mamachi, Lazari, J. H. Heidegger, etc. J. Marx (R.C.): Das Wallfahren in der katholischen Kirche, historisch-kritisch dargestellt. Trier, 1842. Comp. the relevant sections in the church archaeologies of Bingham, Augusti, Binterim, \&c.

Solemn religious processions on high festivals and special occasions had been already customary among the Jews,\footnote{As in the siege of Jericho, Jos. vi. 3 ff.; at the dedication of Solomon’s temple, 1 Kings viii. 1ff; on the entrance of Jesus into Jerusalem, Matt. xxi. 8 ff.} and even among the heathen. They arise from the love of human nature for show and display, which manifests itself in all countries in military parades, large funerals, and national festivities.

The oppressed condition of the church until the time of Constantine made such public demonstrations impossible or unadvisable.

In the fourth century, however, we find them in the East and in the West, among orthodox and heretics,\footnote{The Arians, for example. Comp. Sozom., H. E. viii. 8, where weekly singing processions of the Arians are spoken of.} on days of fasting and prayer, on festivals of thanksgiving, at the burial of the dead, the induction of bishops, the removal of relics, the consecration of churches, and especially in times of public calamity. The two chief classes are thanksgiving and penitential processions. The latter were also called cross-processions, litanies.\footnote{Litaniae (\textgreek{λιτανεῖαι}), supplicationes, rogationes, \textgreek{ἐξομολογήσεις}, stations, collectae.}

The processions moved from church to church, and consisted of the clergy, the monks, and the people, alternately saying or singing prayers, psalms, and litanies. In the middle of the line commonly walked the bishop as leader, in surplice, stole, and pluvial, with the mitre on his head, the crozier in his left hand, and with his right hand blessing the people. A copy of the Bible, crucifixes, banners, images and relics, burning tapers or torches, added solemn state to the procession.\footnote{The antiquity of all these accessory ceremonies cannot be exactly fixed.}

Regular annual processions occurred on Candlemas, and on Palm Sunday. To these was added, after the thirteenth century, the procession on Corpus Christi, in which the sacrament of the altar is carried about and worshipped.

Pilgrimages are founded in the natural desire to see with one’s own eyes sacred or celebrated places, for the gratification of curiosity, the increase of devotion, and the proving of gratitude.\footnote{“Die Stätte, die ein guter Mensch betrat, \par Ist eingeweiht; nach hundert Jahren klingt \par Sein Wort und seine That dem Enkel wieder.”} These also were in use before the Christian era. The Jews went up annually to Jerusalem at their high festivals as afterward the Mohammedans went to Mecca. The heathen also built altars over the graves of their heroes and made pilgrimages thither.\footnote{“Religiosa cupiditas est,” says Paulinus of Nola, Ep. 3 6, “loca videre, in quibus Christus ingressus et passus est et resurrexit et unde ascendit.”} To the Christians those places were most interesting and holy of all, where the Redeemer was born, suffered, died, and rose again for the salvation of the world.

Christian pilgrimages to the Holy Land appear in isolated cases even in the second century, and received a mighty impulse from the example of the superstitiously pious empress Helena, the mother of Constantine the Great. In 326, at the age of seventy-nine, she made a pilgrimage to Jerusalem, was baptized in the Jordan, discovered the holy cross, removed the pagan abominations and built Christian churches on Calvary and Olivet, and at Bethany.\footnote{Euseb., Vita Const iii. 41 sq., and De locis Ebr. a. v. Bethabara.} In this she was liberally supported by her son, in whose arms she died at Nicomedia in 327. The influence of these famous pilgrims’ churches extended through the whole middle age, to the crusades, and reaches even to most recent times.\footnote{Recall the Crimean war of 1854-’56.}

The example of Helena was followed by innumerable pilgrims who thought that by such journeys they made the salvation of their souls more sure. They brought back with them splinters from the pretended holy cross, waters from the Jordan, earth from Jerusalem and Bethlehem, and other genuine and spurious relics, to which miraculous virtue was ascribed.\footnote{Thus Augustine, De Civit. Dei, xxii. 8, is already found citing examples of the supernatural virtue of the \emph{terra sancta} of Jerusalem.}

Several of the most enlightened church fathers, who approved pilgrimages in themselves, felt it necessary to oppose a superstitious estimate of them, and to remind the people that religion might be practised in any place. Gregory of Nyssa shows that pilgrimages are nowhere enjoined in the Scriptures, and are especially unsuitable and dangerous for women, and draws a very unfavorable picture of the immorality prevailing at places of such resort. “Change of place,” says he, “brings God no nearer. Where thou art, God will come to thee, if the dwelling of thy soul is prepared for him.”\footnote{Epist. ad Ambrosium et Basilissam.} Jerome describes with great admiration the devout pilgrimage of his friend Paula to the East, and says that he himself, in his Bethlehem, had adored the manger and birthplace of the Redeemer;\footnote{Adv. Ruffinum ultima Responsio, c. 22 (Opp. ed. Vall. tom. ii. p. 551), where he boastfully recounts his literary journeys, and says: “Protinus concito gradu Bethlehem meam reversus sum, ubi adoravi praesepe et incunabula Salvatoris.” Comp. his Vita Paulae, for her daughter Eustochium, where he describes the pilgrim-stations then in use.} but he also very justly declares that Britain is as near heaven as Jerusalem, and that not a journey to Jerusalem, but a holy living there, is the laudable thing.\footnote{Epist. lviii. ad Paulinum (Opp. ed. Vallarsi, tom. i. p. 318; in the Bened. ed. it is Ep. 49; in the older editions, Ep. 13): “Non Jerusolymis fuisse, sed Jerusolymis bene vixisse, laudandum est.” In the same epistle, p. 319, he commends the blessed monk Hilarion, that, though a Palestinian, he had been only a day in Jerusalem, “ut nec contemnere loca sancta propter viciniam, nec rursus Dominum loco claudere videretur.”}

Next to Jerusalem, Bethlehem, and other localities of the Holy Land, Rome was a preëminent place of resort for pilgrims from the West and East, who longed to tread the threshold of the princes of the apostles (limina apostolorum). Chrysostom regretted that want of time and health prevented him from kissing the chains of Peter and Paul, which made devils tremble and angels rejoice.

In Africa, Hippo became a place of pilgrimage on account of the bones of St. Stephen; in Campania, the grave of St. Felix, at Nola; in Gaul, the grave of St. Martin of Tours († 397). The last was especially renowned, and was the scene of innumerable miracles.\footnote{The Huguenots in the sixteenth century burnt the bones of St. Martin, as objects of idolatry, and scattered their ashes to the winds.} Even the memory of Job drew many pilgrims to Arabia to see the ash heap, and to kiss the earth, where the man of God endured so much.\footnote{So Chrysostomrelates, Hom. v. de statuis, § 1, tom. ii. f 59; \textgreek{ἱνα τὴν κορπίαν ἐκείνην ἴδωσι καὶ θεασάμενοι καταφιλήσωσι τὴν γῆν.}}

In the Roman and Greek churches the practice of pilgrimage to holy places has maintained itself to the present day. Protestantism has divested the visiting of remarkable places, consecrated by great men or great events, of all meritoriousness and superstitious accessories, and has reduced it to a matter of commendable gratitude and devout curiosity. Within these limits even the evangelical Christian cannot view without emotion and edification the sacred spots of Palestine, the catacombs of Rome, the simple slabs over Luther and Melanchthon in the castle-church of Wittenberg, the monuments of the English martyrs in Oxford, or the rocky landing-place of the Puritanic pilgrim fathers in Massachusetts. He feels himself nearer to the spirit of the great dead; but he knows that this spirit continues not in their dust, but lives immortally with God and the saints in heaven.

\xsection{Public Worship of the Lord's Day. Scripture-Reading and Preaching}

§ 90. Public Worship of the Lord’s Day. Scripture-Reading and Preaching.

J.A. Schmidt: De primitive ecclesiae lectionibus. Helmst. 1697. E. Ranke: Das kirchliche Perikopensystem aus den Ältesten Urkunden der röm. Liturgie. Berlin, 1847. H. T. Tzschirner: De claris Eccles. vet. oratoribus Comment. i.-ix. Lips. 1817 sqq. K. W. F. Paniel: Pragmatische Geschichte der christl. Beredtsamkeit. Leipz. 1839 ff.

The order and particular parts of the ordinary public worship of God remain the same as they were in the previous period. But the strict separation of the service of the Catechumens,\footnote{Missa catechumenorum, \textgreek{λειτουργίατῶνκατηχουμένων}.} consisting of prayer, scripture reading, and preaching, from the service of the faithful,\footnote{Missa fidelium, \textgreek{λειτουργίατῶνπιστῶν}.} consisting of the communion, lost its significance upon the universal prevalence of Christianity and the union of church and state. Since the fifth century the inhabitants of the Roman empire were now considered as Christians at least in name and confession and could attend even those parts of the worship which were formerly guarded by secrecy against the profanation of pagans. The Greek term liturgy, and the Latin term mass, which is derived from the customary formula of dismission,\footnote{\emph{Missa} is equivalent to \emph{missio, dismissio}, and meant originally the dismission of the congregation after the service by the customary formula: \emph{Ite, missa est} (ecclesia). After the first part of the service the catechumens were thus dismissed by the deacon, after the second part the faithful. But with the fusion of the two parts in one, the formula of dismission was used only at the close, and then it came to signify also the service itself, more especially the eucharistic sacrifice. In the Greek church the corresponding formula of dismission was: \textgreek{ἀπολύεσθε ἐν εἰρήνῃ}, \emph{i.e.}, \emph{ite in pace} (Apost. Const. lib. viii. c. 15). Ambrosius is the first who uses \emph{missa, missam facere} (Ep. 20), for the eucharistic sacrifice. Other derivations of the word, from the Greek \textgreek{μύησις}or the Hebrew verb \texthebrew{השָׂﬠָ}, \emph{to act}, etc., are too far fetched, and cut off by the fact that the word is used only in the Latin church. Comp. vol. i. § 101, p. 383 ff.} was applied, since the close of the fourth century (398), to the communion service or the celebration of the eucharistic sacrifice. This was the divine service in the proper sense of the term, to which all other parts were subordinate. We shall speak of it more fully hereafter.\footnote{Comp. below, §§ 96 and 97.} We have to do at present with those parts which were introductory to the communion and belong to the service of the catechumens as well as to that of the communicants.

The reading of a portion of the Holy Scriptures continued to be an essential constituent of divine service. Upon the close of the church canon, after the Council of Carthage in 397, and other synods, the reading of uncanonical books (such as writings of the apostolic fathers) was forbidden, with the exception of the legends of the martyrs on their memorial days.

There was as yet no obligatory system of pericopes, like that of the later Greek and Roman churches. The lectio continua, or the reading and exposition of whole books of the Bible, remained in practice till the fifth century, and the selection of books for the different parts and services of the church year was left to the judgment of the bishop. At high festivals, however, such portions were read as bore special reference to the subject of the celebration. By degrees, after the example of the Jewish synagogue,\footnote{The Jews, perhaps from the time of Ezra, divided the Old Testament into sections, larger or smaller, called \emph{Parashioth} (\texthebrew{תוֹיּשִׁרְפַּ}), to wit, the Pentateuch into54 \emph{Parashioth,} and the Prophets (\emph{i.e.}, the later historical books and the prophets proper) into as many \emph{Haphtharoth;} and these sections were read in course on the different Sabbaths. This division is much older than the division into verses.} a more complete yearly course of selections from the New Testament for liturgical use was arranged, and the selections were called lessons or pericopes.\footnote{\emph{Lectiones}, \textgreek{ἀναγνώσματα, ἀναγνώσεις, περικοπαί}.} In the Latin church this was done in the fifth century; in the Greek, in the eighth. The lessons\footnote{\emph{Lectiones}, \textgreek{ἀναγνώσματα, ἀναγνώσεις, περικοπαί}.} were taken from the Gospels and from the Epistles, or the Apostle (in part also from the Prophets), and were therefore called the Gospel and the Epistle for the particular Sunday or festival. Some churches, however, had three, or even four lessons, a Gospel, an Epistle, and a section from the Old Testament and from the Acts. Many manuscripts of the New Testament contained only the pericopes or lessons for public worship,\footnote{Hence called \emph{Lectionaria,} sc. volumina, or \emph{Lectionarii}, sc. libri; also \emph{Evangelia cum Epistolis, Comes} (manual of the clergy); in Greek, \textgreek{ἀναγνωστικά, εὐαγγελιστάρια, ἐκλογάδια.}} and many of these again, only the Gospel pericopes.\footnote{Hence \emph{Evangelistaria,} or \emph{Evangelistarim,} in distinction from the \emph{Epistolaria, Epistolare,} or \emph{Apostolus.}} The Alexandrian deacon Euthalius, about 460, divided the Gospel and the Apostle, excepting the Revelation, into fifty-seven portions each, for the Sundays and feast days of the year; but they were not generally received, and the Eastern church still adhered for a long time to the lectio continua. Among the Latin lectionaries still extant, the Lectionarium Gallicanum, dating from the sixth or seventh century, and edited by Mabillon, and the so-called Comes (i.e., Clergyman’s Companion) or Liber Comitis, were in especial repute. The latter is traced by tradition to the learned Jerome, and forms the groundwork of the Roman lectionary and the entire Western System of pericopes, which has passed from the Latin church into the Anglican and the Lutheran, but has undergone many changes in the course of time.\footnote{The high antiquity of the \emph{Comes} appears at any rate in its beginning with the Christmas Vigils instead of the Advent Sunday, and its lack of the festival of the Trinity and most of the saints’ days. There are different recensions of it, the oldest edited by Pamelius, another by Baluze, a third (made by Alcuin at the command of Charlemagne) by Thomasi. E. Ranke, l.c., has made it out probable that Jeromecomposed the \emph{Comes} under commission from Pope Damasus, and is consequently the original author of the Western pericope system.} This selection of Scripture portions was in general better fitted to the church year, but had the disadvantage of withholding large parts of the holy Scriptures from the people.

The lessons were read from the ambo or reading desk by the lector, with suitable formulas of introduction; usually the Epistle first, and then the Gospel; closing with the doxology or the singing of a psalm. Sometimes the deacon read the Gospel from the altar, to give it special distinction as the word of the Lord Himself.

The church fathers earnestly enjoined, besides this, diligent private reading of the Scriptures; especially Chrysostom, who attributed all corruption in the church to the want of knowledge of the Scriptures. Yet he already found himself compelled to combat the assumption that the Bible is a book only for clergy and monks, and not for the people; an assumption which led in the middle age to the notorious papal prohibitions of the Scriptures in the popular tongues. Strictly speaking, the Bible has been made what it was originally intended to be, really a universal book of the people, only by the invention of the art of printing, by the spirit of the Reformation, and by the Bible Societies of modern times. For in the ancient church, and in the middle age, the manuscripts of the Bible were so rare and so dear, and the art of reading was so limited, that the great mass were almost entirely dependent on the fragmentary reading of the Scriptures in public worship. This fact must be well considered, to forestall too unfavorable a judgment of that early age.

The reading of the Scripture was followed by the sermon, based either on the pericope just read, or on a whole book, in consecutive portions. We have from the greatest pulpit orators of antiquity, from Athanasius, Gregory Nazianzen, Basil the Great, Chrysostom, Ambrose, Augustine, connected homilies on Genesis, the Prophets, the Psalms, the Gospels, and the Epistles. But on high festivals a text was always selected suitable and usual for the occasion.\footnote{Comp. Augustine’s Expos. in Joh. in praef.: “Meminit sanctitas vestra, evangelium secundum Johannem ex ordine lectionum nos solere tractare. Sed quia nunc interposita est solemnitas sanctorum dierum, quibus certas ex evangelio lectiones oportet recitari, quae ita sunt annuae, ut aliae esse non possint, ordo ille quem susceperamus, ex necessitate paululum intermissus est, non omissus.”} There was therefore in the ancient church no forced conformity to the pericopes; the advantages of a system of Scripture lessons and a consecutive exposition of entire books of Scripture were combined. The reading of the pericopes belongs properly to the altar-service, and must keep its connection with the church year; preaching belongs to the pulpit, and may extend to the whole compass of the divine word.

Pulpit eloquence in the fourth and fifth centuries reached a high point in the Greek church, and is most worthily represented by Gregory Nazianzen and Chrysostom. But it also often degenerated there into artificial rhetoric, declamatory bombast, and theatrical acting. Hence the abuse of frequent clapping and acclamations of applause among the people.\footnote{\textgreek{Κρότος}, acclamatio, applausus. Chrysostomand Augustineoften denounced this theatrical disorder, but in vain.} As at this day, so in that, many went to church not to worship God, but to hear a celebrated speaker, and left as soon as the sermon was done. The sermon, they said, we can hear only in the church, but we can pray as well at home. Chrysostom often raised his voice against this in Antioch and in Constantinople. The discourses of the most favorite preachers were often written down by stenographers and multiplied by manuscripts, sometimes with their permission, sometimes without.

In the Western church the sermon was much less developed, consisted in most cases of a simple practical exhortation, and took the background of the eucharistic sacrifice. Hence it was a frequent thing there for the people to leave the church at the beginning of the sermon; so that many bishops, who had no idea of the free nature of religion and of worship, compelled the people to hear by closing the doors.

The sermon was in general freely delivered from the bishop’s chair or from the railing of the choir (the cancelli), sometimes from the reading-desk. The duty of preaching devolved upon the bishops; and even popes, like Leo I. and Gregory I., frequently preached before the Roman congregation. Preaching was also performed by the presbyters and deacons. Leo I. restricts the right of preaching and teaching to the ordained clergy;\footnote{Ep. 62 ad Maxim.: “Praeter eos qui sunt Domini sacerdotes nullus sibi jus docendi et praedicandi audeat vindicare, sive sit ille monachus, sive sit laicus, qui alicujus scientiae nomine glorietur.”} yet monks and hermits preached not rarely in the streets, from pillars (like St. Symeon), roofs, or trees; and even laymen, like the emperor Constantine and some of his successors, wrote and delivered (though not in church) religious discourses to the faithful people.\footnote{Euseb. Vita Const. iv. 29, 32, 55, and Constantine’s Oratio ad Sanctos, in the appendix}

\xsection{The Sacraments in General}

§ 91. The Sacraments in General.

G. L. Hahn: Die Lehre von den Sacramenten in ihrer geschichtlichen Entwicklung innerhalb der abendländischen Kirche bis zum Concil von Trient. Breslau, 1864 (147 pp.). Comp. also the article Sacramente by G. E. Steitz in Herzog’s Real-Encyklopädie, vol. xiii. pp. 226–286; and Const. von Schätzler: Die Lehre von der Wirksamkeit der Sacramente ex opere operato. Munich, 1860.

The use of the word sacramentum in the church still continued for a long time very indefinite. It embraced every mystical and sacred thing (omne mysticum sacrumque signum). Tertullian, Ambrose, Hilary, Leo, Chrysostom, and other fathers, apply it even to mysterious doctrines and facts, like the Trinity, the divinity of Christ, the incarnation, the crucifixion, and the resurrection. But after the fifth century it denotes chiefly sacred forms of worship, which were instituted by Christ and by which divine blessings are mystically represented, sealed, and applied to men. This catholic theological conception has substantially passed into the evangelical churches, though with important changes as to the number and operation of the sacraments.\footnote{. The word \emph{sacramentum} bears among the fathers the following senses: (1) The \emph{oath} in general, as in the Roman profane writers; and particularly the \emph{soldier’s oath}. (2) The \emph{baptismal vow}, by which the candidate bound himself to the perpetual service of Christ, as \emph{miles Christi}, against sin, the world, and the devil. (3) The \emph{baptismal confession}, which was regarded as a spiritual oath. (4) \emph{Baptism} itself, which, therefore, was often styled\emph{sacramentum fidei}, s. \emph{salutis}, also \emph{pignus} \emph{salutis}. (5) It became almost synonymous with \emph{mystery}, by reason of an inaccurate translation of the Greek \textgreek{μυστήριον}, in the Vulgate (comp. Eph. v. 32), and was accordingly applied to facts, truths, and precepts of the gospel which were concealed from those not Christiana, and to the Christian revelation in general. (6) The \emph{eucharist}, and other holy ordinances and usages of the church. (7) After the twelfth century the \emph{seven} well-known \emph{sacraments} of the Catholic church. Comp. the proofs in Hahn, l.c. pp. 5-10, where yet other less usual senses of the word are adduced.}

Augustine was the first to substitute a clear doctrine of the nature of the sacraments for a vague notion and rhetorical exaggerations. He defines a sacrament to be a visible sign of an invisible grace or divine blessing.\footnote{Signum visibile, or forma visibilis gratiae invisibilis. Augustinecalls the sacraments also verba visibilia, signacula corporalia, signa rerum spiritualium, signacula rerum divinarum visibilia, etc. See Hahn, l.c. p. 11 ff. The definition is not adequate. At least a third mark must be added, not distinctly mentioned by Augustine, viz., the divina institutio, or, more precisely, a mandatum Christi. This is the point of difference between the Catholic and Protestant conceptions of the sacrament. The Roman and Greek churches take the divine institution in a much broader sense, while Protestantism understands by it an express command of Christ in the New Testament, and consequently limits the number of sacraments to baptism and the Lord’s Supper, since for the other five sacraments the Catholic church can show no such command. Yet confirmation, ordination, and marriage have practically acquired a sacramental import in Protestantism, especially in the Lutheran and Anglican churches.} Two constituents, therefore, belong to such a holy act: the outward symbol or sensible element (the signum, also sacramentum in the stricter sense), which is visible to the eye, and the inward grace or divine virtue (the res or virtus sacramenti), which is an object of faith.\footnote{Augustine, De catechiz. rudibus, § 50: “Sacramenta signacula quidem rerum divinarum esse visibilia, sed res ipsas invisibiles in eis honoari.” Serm. ad pop. 292 (tom. v. p. 770): “Dicuntur sacramenta, quia in eis aliud videtur, aliud intelligitur. Quod videtur, speciem habet corporalem; quod intelligitur, fructum habet spiritalem.”} The two, the sign and the thing signified, are united by the word of consecration.\footnote{Augustine, In Joann Evang. tract. 80: “Detrahe verbum, et quid est aqua [the baptismal water] nisi aqua? \emph{Accedit verbum ad elementum et fit sacramemtum}, etiam ipsum tamquam visibile verbum.”} From the general spirit of Augustine’s doctrine, and several of his expressions, we must infer that he considered divine institution by Christ to be also a mark of such holy ordinance.\footnote{Comp. Epist. 82, §§ 14 and 15; Ep. 138, § 7; De vera relig. c. 16, § 33; and Hahn p. 154.} But subsequently this important point retired from the consciousness of the church, and admitted the widening of the idea, and the increase of the number, of the sacraments.

Augustine was also the first to frame a distinct doctrine of the operation of the sacraments. In his view the sacraments work grace or condemnation, blessing or curse, according to the condition of the receiver.\footnote{Comp. the proof passages in Hahn, p. 279 ff. Thus Augustinesays, e.g., De bapt. contra Donat. 1. iii. c. 10 (tom. ix. p. 76): “Sacramento suo divina virtus adsistit sive ad salutem bene utentium, sive ad perniciem male utentium.” De unit. eccl.c.21 (tom. ix. p. 256): “Facile potestis intelligere et in bonis esse et in malis sacramenta divina, sed in illis ad salutem, in malis ad damnationem.”} They operate, therefore, not immediately and magically, but mediately and ethically, not ex opere operato, in the later scholastic language, but through the medium of the active faith of the receiver. They certainly have, as divine institutions, an objective meaning in themselves, like the life-principle of a seed, and do not depend on the subjective condition of the one who administers them (as the Donatists taught); but they reach with blessing only those who seize the blessing, or take it from the ordinance, in faith; they bring curse to those who unworthily administer or receive them. Faith is necessary not as the efficient cause, but as the subjective condition, of the saving operation of the offered grace.\footnote{Hence the later formula: Fides non facit ut \emph{sit} sacramentum, sed ut \emph{prosit}. Faith does not produce the sacramental blessing, but subjectively receives and appropriates it.} Augustine also makes a distinction between a transient and a permanent effect of the sacrament, and thereby prepares the way for the later scholastic doctrine of the character indelebilis. Baptism and ordination impress an indelible character, and therefore cannot be repeated. He is fond of comparing baptism with the badge of the imperial service,\footnote{Stigma militare, character militaris. To this the expression \emph{character indelebilis} certainly attaches itself easily, though the doctrine concerning it cannot be traced with certainty back of the thirteenth century. Comp. Hahn, l.c. p. 298 ff., where it is referred to the time of Pope Innocent III.} which the soldier always retains either to his honor or to his shame. Hence the Catholic doctrine is: Once baptized, always baptized; once a priest, always a priest. Nevertheless a baptized person, or an ordained person, can be excommunicated and eternally lost. The popular opinion in the church already inclined strongly toward the superstitious view of the magical operation of the sacrament, which has since found scholastic expression in the opus operatum theory.

The church fathers with one accord assert a relative (not absolute) necessity of the sacraments to salvation.\footnote{Even Augustine, De peccat. merit. et remiss. lib. i. c. 24, § 34: “Praeter baptismum et participationem mensae dominicae non solum ad regnum Dei, sed nec ad salutem et vitam aeternam posse quemquam hominem pervenire.” This would, strictly considered, exclude all Quakers and unbaptized infants from salvation; but Augustineadmits as an exception the possibility of a conversion of the heart without baptism. See below. The scholastics distinguished more accurately a threefold necessity: (1) absolute: \emph{simpliciter necessarium}; (2) teleological: \emph{in ordine ad finem}; (3) hypothetical or relative\emph{: necessarium ex suppositione, quae est necessitas consequentiae}. To the sacraments belongs only the last sort of necessity, because now, under existing circumstances, God will not ordinarily save any one without these means which he has appointed. Comp. Hahn, 1. c. p. 26 ff. According to Thomas Aquinas only three sacraments are perfectly necessary, viz., baptism and penance for the individual, and ordination for the whole church.} They saw in them, especially in baptism and the eucharist, the divinely appointed means of appropriating the forgiveness of sins and the grace of God. Yet with this view they firmly held that not the want of the sacraments, but only the contempt of them, was damning.\footnote{“Non defectus, sed contemptus sacramenti damnat.” Comp. Augustine, De bapt. contra Donat. 1. iv. c. 25, §32: “Conversio cordis potest quidem inesse non percepto baptismo, sed contemto non potest. Neque enim ullo modo dicenda est conversio cordis ad Deum, cum Dei sacramentum contemnitur.”} In favor of this they appealed to Moses, Jeremiah, John the Baptist, the thief on the cross,—who all, however, belonged to the Old Testament economy—and to many Christian martyrs, who sealed their faith in Christ with their blood, before they had opportunity to be baptized and to commune. The Virgin Mary also, and the apostles, belong in some sense to this class, who, since Christ himself did not baptize, received not the Christian baptism of water, but instead were on the day of Pentecost baptized with Spirit and with fire. Thus Cornelius also received through Peter the gift of the Holy Ghost before baptism; but nevertheless submitted himself afterwards to the outward Sacrament. In agreement with this view, sincere repentance and true faith, and above all the blood-baptism of martyrdom,\footnote{Baptismus sanguinis.} were regarded as a kind of compensation for the sacraments.

The number of the sacraments remained yet for a long time indefinite; though among the church fathers of our period baptism and the Lord’s Supper were regarded either as the only Sacraments, or as the prominent ones.

Augustine considered it in general an excellence of the New Testament over the Old, that the number of the sacraments was diminished, but their import enhanced,\footnote{Contra Faust. xix. 13: “Prima sacramenta praenunciativa erant Christi venturi; quae cum suo adventu Christus implevisset ablata sunt et alia sunt instituta, \emph{virtute majora, numero pauciora}.”} and calls baptism and the Supper, with reference to the water and the blood which flowed from the side of the Lord, the genuine or chief sacraments, on which the church subsists.\footnote{De symb. ad Catech. c. 6: “Quomodo Eva facts est ex latere Adam, ita ecclesia formatur ex latere Christi. Percussum est ejus latus et statim manavit sanguis et \emph{aqua}, quae sunt ecclesiae \emph{genuina} \emph{sacramenta}.” De ordine baptismi, c. 5 (Bibl. max. tom. xiv. p. 11): “Profluxerunt ex ejus latere \emph{sanguis} et\emph{aqua, duo sanctae ecclesiae praecipua sacramenta}.” Serm. 218: “Sacraments, quibus formatur ecclesia.” Comp. Chrysostom, Homil. 85 in Joh: \textgreek{ἐξ ἀμφοτέρων ἡ ἐκκλησία συνέστηκε.} Tertulliancalled baptism and the eucharist “sacramenta propria,” Adv. Marc. i. 14.} But he includes under the wider conception of the sacrament other mysterious and holy usages, which were commended in the Scriptures,\footnote{“Et si quid aliud in divinis literis commendatur,” or: “omne mysticum sacrumque signum.”} naming expressly confirmation,\footnote{“Sacrimentum chrismatis, ” Contr. Lit. Petiliani ii. 104. So even Cyprian, Ep. 72.} marriage,\footnote{“Sacramentum nuptiarum,” De nuptiis et concupisc. i. 2.} and ordination.\footnote{“Sacramentum dandi baptismum,” De bapt ad Donat. i. 2; Epist. Parm ii. 13.} Thus he already recognizes to some extent five Christian sacraments, to which the Roman church has since added penance and extreme unction.

Cyril of Jerusalem, in his Mystagogic Catechism, and Ambrose of Milan, in the six books De Sacramentis ascribed to him, mention only three sacraments: baptism, confirmation, and the Lord’s supper; and Gregory of Nyssa likewise mentions three, but puts ordination in the place of confirmation. For in the Eastern church confirmation, or the laying on of hands, was less prominent, and formed a part of the sacrament of baptism; while in the Western church it gradually established itself in the rank of an independent sacrament.

The unknown Greek author of the pseudo-Dionysian writings of the sixth century enumerates six sacraments (μυστήρια):\footnote{De hierarch. eccles. c. 2 sq.}(1.) baptism, or illumination; (2.) the eucharist, or the consecration of consecrations; (3.) the consecration with anointing oil, or confirmation; (4.) the consecration of priests; (5.) the consecration of monks; (6.) the consecration of the dead, or extreme unction. Here marriage and penance are wanting; in place of them appears the consecration of monks, which however was afterwards excluded from the number of the sacraments.

In the North African, the Milanese, and the Gallican churches the washing of feet also long maintained the place of a distinct sacrament.\footnote{According to the testimony of Ambrose, Augustine, and the Missale Gallicum vetus. Comp. Hahn, l.c. p. 84 f.} Ambrose asserted its sacramental character against the church of Rome, and even declared it to be as necessary as baptism, because it was instituted by Christ, and delivered men from original sin, as baptism from the actual sin of transgression;—a view which rightly found but little acceptance.

This uncertainty as to the number of the sacraments continued till the twelfth century.\footnote{Beda Venerabilis († 735), Ratramnus of Corbie († 868), Ratherius of Verona († 974), in enumerating the sacraments, name only baptism and the Lord’s Supper; and even Alexander of Hales († 1245) expressly says (Summa p. iv. Qu. 8, Membr. 2, art. 1): “Christus duo sacraments instituit per se ipsum, sacramentum baptismi et sacramentum eucharistiae.” Damiani († 1072), on the other hand, mentions twelve sacraments, viz., baptism, confirmation, anointing of the sick, consecration of bishops, consecration of kings, consecration of churches, penance, consecration of canons, monks, hermits, and nuns and marriage. Opp. tom. ii. 372 (ed. C. Cajet.). Bernard of Clairvaux († 1151) names ten sacraments. Confirmation was usually reckoned among the sacraments. Comp. Hahn, l.c. 88 ff.} Yet the usage of the church from the fifth century downward, in the East and in the West, appears to have inclined silently to the number seven, which was commended by its mystical sacredness. This is shown at least by the agreement of the Greek and Roman churches in this point, and even of the Nestorians and Monophysites, who split off in the fifth century from the orthodox Greek church.\footnote{No plain trace, however, of such a definite number appears in the earliest monuments of the faith of these Oriental sects, or even in the orthodox theologian John Damascenus.}

In the West, the number seven was first introduced, as is usually supposed, by the bishop Otto of Bamberg (1124), more correctly by Peter Lombard (d. 1164), the “Master of Sentences;” rationally and rhetorically justified by Thomas Aquinas and other scholastics (as recently by Möhler) from the seven chief religious wants of human life and human society;\footnote{Usually: Birth = baptism; growth = confirmation; nourishment = the Supper; healing of Sickness = penance; perfect restoration = extreme unction; propagation of society = marriage; government of society = orders. Others compare the sacraments with the four cardinal natural virtues: prudence, courage, justice, and temperance, and the three theological virtues: faith, love, and hope; but vary in their assignments of the several sacraments to the several virtues respectively. All these comparisons are, of course, more or less arbitrary and fanciful.} and finally publicly sanctioned by the council of Florence in 1439 with the concurrence of the Greek church, and established by the council of Trent with an anathema against all who think otherwise.\footnote{The Council of Trent pronounces the anathema upon all who deny the number of seven sacraments and its institution by Christ, Sess. vii. de sacr. can. 1: “Si quis dixerit, sacramenta novae legis non fuisse omnia a Christo instituta, aut esse plum vel pauciora quam septem, anathema sit.” In default of a historical proof of the seven sacraments from the writings of the church fathers, Roman divines, like Brenner and Perrone, find themselves compelled to resort to the \emph{disciplina} \emph{arcani}; but this related only to the \emph{celebration} of the sacraments, and disappeared in the fourth century upon the universal adoption of Christianity. Comp. also the treatise of G. L. Hahn: Doctrinae Romanae de numero sacramentario septenario rationes historicae. Vratisl. 1859.} The Reformation returned, in this point as in others, to the New Testament; retained none but baptism and the Lord’s Supper as proper sacraments, instituted and enjoined by Christ himself; entirely rejected extreme unction (and at first confirmation); consigned penance to the province of the inward life, and confirmation, marriage, and orders to the more general province of sacred acts and usages, to which a more or less sacramental character may be ascribed, but by no means an equality in other respects with baptism and the holy Supper.\footnote{A more particular discussion of the differences between the Roman and the Protestant doctrines of the sacraments belongs to symbolism and polemics.}

\xsection{Baptism}

§ 92. Baptism.

For the Literature, see vol. i. § 37, p. 122; especially Höfling (Lutheran): Das Sacrament der Taufe. W. Wall (Anglican): The History of Infant Baptism (1705), new ed. Oxf. 1844, 4vols. C. A. G. v. Zezschwitz: System der christlich kirchlichen Katechetik. Vol. i. Leipz. 1863.

On heretical baptism in particular, See Mattes (R.C.): Ueber die Ketzertaufe, in the Tübingen “Theol. Quartalschrift,” for 1849, pp. 571–637, and 1850, pp. 24–69; and G. E. Steitz, art. Ketzertaufe in Herzog’s Theol. Encyclop. vol. vii. pp. 524–541 (partly in opposition to Mattes). Concerning the form of baptism, on the Baptist side, T. J. Conant: The Meaning and Use of Baptizein philologically and historically investigated. New York, 1861.

The views of the ante-Nicene fathers concerning baptism and baptismal regeneration were in this period more copiously embellished in rhetorical style by Basil the Great and the two Gregories, who wrote special treatises on this sacrament, and were more clearly and logically developed by Augustine. The patristic and Roman Catholic view on regeneration, however, differs considerably from the one which now prevails among most Protestant denominations, especially those of the more Puritanic type, in that it signifies not so much a subjective change of heart, which is more properly called conversion, but a change in the objective condition and relation of the sinner, namely, his translation from the kingdom of Satan into the kingdom of Christ. Some modern divines make a distinction between baptismal and moral regeneration, in order to reconcile the doctrine of the fathers with the fact that the evidences of a new life are wholly wanting in so many who are baptized. But we cannot enter here into a discussion of the difficulties of this doctrine, and must confine ourselves to a historical statement.

Gregory Nazianzen sees in baptism all blessings of Christianity combined, especially the forgiveness of sins, the new birth, and the restoration of the divine image. To children it is a seal (σφραγίς) of grace and a consecration to the service of God. According to Gregory of Nyssa, the child by baptism is instated in the paradise from which Adam was thrust out. The Greek fathers had no clear conception of original sin. According to the Pelagian Julian of Eclanum, Chrysostom taught: We baptize children, though they are not stained with sin, in order that holiness, righteousness, sonship, inheritance, and brotherhood may be imparted to them through Christ.\footnote{The passage is not found in the writings of Chrysostom. Augustine, however, does not dispute the citation, but tries to explain it away (contra Julian. i. c. 6, § 21).}

Augustine brought the operation of baptism into connection with his more complete doctrine of original sin. Baptism delivers from the guilt of original sin, and takes away the sinful character of the concupiscence of the flesh,\footnote{De nupt. et concup. i. 28: “Dimittitur concupiscentia carnis in baptismo, non ut non sit, sed ut in peccatum non imputetur.”} while for the adult it at the same time effects the forgiveness of all actual transgressions before baptism. Like Ambrose and other fathers, Augustine taught the necessity of baptism for entrance into the kingdom of heaven, on the ground of John iii. 5, and deduced therefrom, in logical consistency, the terrible doctrine of the damnation of all unbaptized children, though he assigned them the mildest grade of perdition.\footnote{“Parvulos in damnatione omnium mitissima futuros.” Comp. De peccat. mer. i. 20, 21, 28; Ep. 186, 27. To the heathen he also assigned a milder and more tolerable condemnation, Contr. Julian. iv. 23.}

The council of Carthage, in 318, did the same, and in its second canon rejected the notion of a happy middle state for unbaptized children. It is remarkable, however, that this addition to the second canon does not appear in all copies of the Acts of the council, and was perhaps out of some horror omitted.\footnote{Comp. Neander, l.c. i. p. 424, and especially Hefele, Conciliengeschichte, ii. p. 103. The passage in question, which is lacking both in Isidore and in Dionysius, runs thus: “Whoever says that there is, in the kingdom of heaven or elsewhere, a certain middle place, where children who die without baptism live happy (beate vivant), while yet they cannot without baptism enter into the kingdom of heaven, \emph{i.e.}, into eternal life, let him be anathema.”}

In Augustine we already find all the germs of the scholastic and Catholic doctrine of baptism, though they hardly agree properly with his doctrine of predestination, the absolute sovereignty of divine grace and the perseverance of saints. According to this view, baptism is the sacrament of regeneration, which is, negatively, the means of the forgiveness of sin, that is, both of original sin and of actual sins committed before baptism (not after it), and positively, the foundation of the new spiritual life of faith through the impartation of the gratia operans and co-operans. The subjective condition of this effect is the worthy receiving, that is, penitent faith. Since in the child there is no actual sin, the effect of baptism in this case is limited to the remission of the guilt of original sin; and since the child cannot yet itself believe, the Christian church (represented by the parents and the sponsors) here appears in its behalf, as Augustine likewise supposed, and assumes the responsibility of the education of the baptized child to Christian majority.\footnote{The scholastics were not entirely agreed whether baptism imparts positive grace to all, or only to adults. Peter Lombard was of the latter opinion; but most divines extended the positive effect of baptism even to children, though under various modifications. Comp. the full exposition of the scholastic doctrine of baptism (which does not belong here) in Hahn, l.c. p. 333 ff.}

As to infant baptism: there was in this period a general conviction of its propriety and of its apostolic origin. Even the Pelagians were no exception; though infant baptism does not properly fit into their system; for they denied original sin, and baptism, as a rite of purification, always has reference to the forgiveness of sins. They attributed to infant baptism an improving effect. Coelestius maintained that children by baptism gained entrance to the higher stage of salvation, the kingdom of God, to which, with merely natural powers, they could not attain. He therefore supposed a middle condition of lower salvation for unbaptized children, which in the above quoted second canon of the council of Carthage—if it be genuine—is condemned. Pelagius said more cautiously: Whither unbaptized children go, I know not; whither they do not go, I know.

But, notwithstanding this general admission of infant baptism, the practice of it was by no means universal. Forced baptism, which is contrary to the nature of Christianity and the sacrament, was as yet unknown. Many Christian parents postponed the baptism of their children, sometimes from indifference, sometimes from fear that they might by their later life forfeit the grace of baptism, and thereby make their condition the worse. Thus Gregory Nazianzen and Augustine, though they had eminently pious mothers, were not baptized till their conversion in their manhood. But they afterward regretted this. Gregory admonishes a mother: “Let not sin gain the mastery in thy child; let him be consecrated even in swaddling bands. Thou art afraid of the divine seal on account of the weakness of nature. What weakness of faith! Hannah dedicated her Samuel to the Lord even before his birth; and immediately after his birth trained him for the priesthood. Instead of fearing human weakness, trust in God.”

Many adult catechumens and proselytes likewise, partly from light-mindedness and love of the world, partly from pious prudence and superstitious fear of impairing the magical virtue of baptism, postponed their baptism until some misfortune or severe sickness drove them to the ordinance. The most celebrated example of this is the emperor Constantine, who was not baptized till he was on his bed of death. The postponement of baptism in that day was equivalent to the postponement of repentance and conversion so frequent in ours. This custom was resisted by the most eminent church teachers, but did not give way till the fifth century, when it gradually disappeared before the universal introduction of infant baptism.

Heretical baptism was now generally regarded as valid, if performed in the name of the triune God. The Roman view prevailed over the Cyprianic, at least in the Western church; except among the Donatists, who entirely rejected heretical baptism (as well as the catholic baptism), and made the efficacy of the sacrament depend not only on the ecclesiastical position, but also on the personal piety of the officiating priest.

Augustine, in his anti-Donatistic writings, defends the validity of heretical baptism by the following course of argument: Baptism is an institution of Christ, in the administration of which the minister is only an agent; the grace or virtue of the sacrament is entirely dependent on Christ, and not on the moral character of the administering agent; the unbeliever receives not the power, but the form of the sacrament, which indeed is of no use to the baptized as long as he is outside of the saving catholic communion, but becomes available as soon as he enters it on profession of faith; baptism, wherever performed, imparts an indelible character, or, as he calls it, a “character dominicus,” “regius.” He compares it often to the “nota militaris,” which marks the soldier once for all, whether it was branded on his body by the legitimate captain or by a rebel, and binds him to the service, and exposes him to punishment for disobedience.

Proselyted heretics were, however, always confirmed by the laying on of hands, when received into the catholic church. They were treated like penitents. Leo the Great says of them, that they have received only the form of baptism without the power of sanctification.\footnote{Epist. 129 ad Nicet. c. 7: “Qui baptismum ab haereticis acceperunt ... sola invocatione Spiritus S. per impositionem manuum confirmandi sunt, quia formam tantum baptismi sine sanctificationis virtute sumpserunt.”}

The most eminent Greek fathers of the Nicene age, on the other hand, adhered to the position of Cyprian and Firmilian. Athanasius, Gregory Nazianzen, Basil, and Cyril of Jerusalem regarded, besides the proper form, the true trinitarian faith on the part of the baptizing community, as an essential condition of the validity of baptism. The 45th of the so-called Apostolic Canons threatens those with excommunication who received converted heretics without rebaptism. But a milder view gradually obtained even in the East, which settled at last upon a compromise.

The ecumenical council of Constantinople in 381, in its seventh canon (which, however, is wanting in the Latin versions, and is perhaps later), recognizes the baptism of the Arians, the Sabbatians (a sort of Novatians, so called from their leader Sabbatius), the Quartodecimanians, the Apollinarians, but rejected the baptism of the Eunomians, “who baptize with only one immersion,” the Sabellians, “who teach the Son-Fatherhood (υἱοπατορία),” the Montanists (probably because they did not at that time use the orthodox baptismal formula), and all other heretics. These had first to be exercised, then instructed, and then baptized, being treated therefore as heathen proselytes.\footnote{Comp. Hefele, Conciliengeschichte, ii. 26; Mattes, Ueber die Ketzertaufe, in the Tübingen Quartalschrift, 1849, p. 580.} The Trullan council of 692, in its 95th canon repeated this canon, and added the Nestorians, the Eutychians, and the followers of Dioscurus and Severus to the list of those heretics who may be received into the church on a mere recantation of their error. These decisions lack principle and consistency.

The catechetical instruction which preceded the baptism of proselytes and adults, and followed the baptism of children, ended with a public examination (scrutinium) before the congregation. The Creed—in the East the Nicene, in the West the Apostles’—was committed to memory and professed by the candidates or the god-parents of the children.

The favorite times for baptism for adults were Easter and Pentecost, and in the East also Epiphany. In the fourth century, when the mass of the population of the Roman empire went over from heathenism to Christianity, the baptisteries were thronged with proselytes on those high festivals, and the baptism of such masses had often a very imposing and solemn character. Children were usually incorporated into the church by baptism soon after their birth.

Immersion continued to be the usual form of baptism, especially in the East; and the threefold immersion in the name of the Trinity. Yet Gregory the Great permitted also the single immersion, which was customary in Spain as a testimony against the Arian polytheism.\footnote{Greg. Ep. i. 43, to Bishop Leander of Seville: “Dum in tribus subsistentiis una substantia est, reprehensibile esse nullatenus potest infantem in baptismate vel ter vel semel mergere: quando et in tribus mersionibus personarum trinitas, et in una potest personarum singularitas designari. Sed quia nunc usque ab haereticis infans in baptismate tertio mergebatur, fiendum \emph{apud vos} non esse censeo, ne dum mersiones numerant, divinitatem dividant.” From this we see, at the same time, that even in infant baptism, and among heretics, immersion was the custom. Yet in the nature of the case, sprinkling, at least of weak or sick children, as in the \emph{baptismus clinicorum}, especially in northern climates, came early into use.}

With baptism, several preparatory and accompanying ceremonies, some of them as early as the second and third centuries, were connected; which were significant, but overshadowed and obscured the original simplicity of the sacrament. These were exorcism, or the expulsion of the devil;\footnote{Comp. vol. i. p. 399.} breathing upon the candidates,\footnote{Insufflare, \textgreek{ἐμφυσᾷν}.} as a sign of the communication of the Holy Ghost, according to John xx. 22; the touching of the ears,\footnote{Sacramentum apertionis.} with the exclamation: Ephphatha!—from Mark vii. 34, for the opening of the spiritual understanding; the sign of the cross made upon the forehead and breast, as the mark of the soldier of Christ; and, at least in Africa, the giving of salt, as the emblem of the divine word, according to Mark ix. 50; Matt. v. 13 Col. iv. 6. Proselytes generally took also a new name, according to Rev. ii. 17.

In the act of baptism itself, the candidate first, with his face toward the west, renounced Satan and all his pomp and service,\footnote{This was the \textgreek{ἀποταγή}, or abrenunciatio diaboli, with the words: \textgreek{Ἀποτάσσομαί σοι, Σατανᾶ, καὶ πάσῃ τῇ πομπῇ σου καὶ πάσῃ τῇ λατρείᾳ σου.} The Apostolic Constitutions add \textgreek{τοῖς ἔργοις}. In Tertullian: “Renunciare diabolo et pompae et angelis ejus.”} then, facing the east, he vowed fidelity to Christ,\footnote{Συντάσσομαί σοι, Χριστέ.} and confessed his faith in the triune God, either by rehearsing the Creed, or in answer to questions.\footnote{\textgreek{Ὁμολόγησις}, professio.} Thereupon followed the threefold or the single immersion in the name of the triune God, with the calling of the name of the candidate, the deacons and deaconesses assisting. After the second anointing with the consecrated oil (confirmation), the veil was removed, with which the heads of catechumens, in token of their spiritual minority, were covered during divine worship, and the baptized person was clothed in white garments, representing the state of regeneration, purity, and freedom. In the Western church the baptized person received at the same time a mixture of milk and honey, as a symbol of childlike innocence and as a fore-taste of the communion.

\xsection{Confirmation}

§ 93. Confirmation.

Comp. the Literature of Baptism, especially Höfling, and Zezschwitz: Der Katechumenat (first vol. of his System der Katechetik). Leipzig, 1863.

Confirmation, in the first centuries, was closely connected with the act of baptism as the completion of that act, especially in adults. After the cessation of proselyte baptism and the increase of infant baptism, it gradually came to be regarded as an independent sacrament. Even by Augustine, Leo I., and others, it is expressly called sacramentum.\footnote{Aug. Contra liter. Petil. l. ii. c. 104 (tom. ix. p. 199); Leo, Epist. 156, c. 5. Confirmation is called \emph{confirmatio} from its nature; \emph{sigillum} or \emph{consignatio}, from its design; \emph{chrisma} or\emph{unctio}, from its matter; and \emph{impositio manuum}, from its form.} This independence was promoted by the hierarchical interest, especially in the Latin church, where the performance of this rite is an episcopal function.

The catholic theory of confirmation is, that it seals and completes the grace of baptism, and at the same time forms in some sense a subjective complement to infant baptism, in which the baptized person, now grown to years of discretion, renews the vows made by his parents or sponsors in his name at his baptism, and makes himself personally responsible for them. The latter, however, is more properly a later Protestant (Lutheran and Anglican) view. Baptism, according to the doctrine of the ancient church, admits the man into the rank of the soldiers of Christ; confirmation endows him with strength and courage for the spiritual warfare.

The outward form of confirmation consists in the anointing of the forehead, the nose, the ear, and the breast with the consecrated oil, or a mixture of balsam,\footnote{\textgreek{Χρίσμα}. This was afterward, in the Latin church, the second anointing, in distinction from that which took place at baptism. The Greek church, however, which always conjoins confirmation with baptism, stopped with one anointing. Comp. Hahn, l.c. p. 91 f.} which symbolizes the consecration of the whole man to the spiritual priesthood; and in the laying on of the hands of the clergyman,\footnote{\emph{Impositio manuum}. This, however, subsequently became less prominent than the anointing; hence confirmation is also called simply \emph{chrisma}, or \emph{sacramentum chrismatis, unctionis}.} which signifies and effects the communication of the Holy Ghost for the general Christian calling.\footnote{The formula now used in the Roman church in the act of confirmation, which is not older, however, than the twelfth century, runs: “Signo te signo crucis et confirmo te chrismate salutis, in nomine Patris et Filii et Spiritus Sancti.”} The anointing takes precedence of the imposition of hands, in agreement with the Old Testament sacerdotal view; while in the Protestant church, wherever confirmation continues, it is entirely abandoned, and only the imposition of hands is retained.

In other respects considerable diversity prevailed in the different parts of the ancient church in regard to the usage of confirmation and the time of performing it.

In the Greek church every priest may administer confirmation or holy unction, and that immediately after baptism; but in the Latin church after the time of Jerome (as now in the Anglican) this function, like the power of ordination, was considered a prerogative of the bishops, who made periodical tours in their dioceses to confirm the baptized. Thus the two acts were often far apart in time.

\xsection{Ordination}

§ 94. Ordination.

J. Morinus (R.C.): Comment. Hist. so dogm. de sacris Eccles. ordinationibus. Par. 1655, etc. Fr. Halierius (R.C.): De sacris electionibus et ordinationibus. Rom. 1749. 3 vols. fol. G. L. Hahn: l.c. p. 96 and p. 354 ff. Comp. the relevant sections in the archaeological works of Bingham, Augusti, Binterim, etc.

The ordination of clergymen\footnote{\textgreek{Χειροτονία, καθιέρσις}\emph{, ordinatio}and in the case of bishops, \emph{consecratio}.} was as early as the fourth or fifth century admitted into the number of sacraments. Augustine first calls it a sacrament, but with the remark that in his time the church unanimously acknowledged the sacramental character of this usage.\footnote{De bono conjug. c. 18 (tom. vi. p. 242), c. 24 (p. 247); Contr. Epist. Parmen. l. ii. c. 12 (tom. ix. pp. 29, 30). Comp. Leo M. Epist. xii. c. 9; Gregor. M. Expos. in i. Regg. l. vi. c. 3. These and other passages in Hahn, p. 97.}

Ordination is the solemn consecration to the special priesthood, as baptism is the introduction to the universal priesthood; and it is the medium of communicating the gifts for the ministerial office. It confers the capacity and authority of administering the sacraments and governing the body of believers, and secures to the church order, care, and steady growth to the end of time. A ruling power is as necessary in the church as in the state. In the Jewish church there was a hereditary priestly caste; in the Christian this is exchanged for an unbroken succession of voluntary priests from all classes, but mostly from the middle and lower classes of the people.

Like baptism and confirmation, ordination imparts, according to the later scholastic doctrine, a character indelebilis, and cannot therefore be repeated.\footnote{Already intimated by Augustine, De bapt. c. Donat. ii. 2: “Sicut baptizatus, si ab unitate recesserit, sacramentum dandi non amittit, sic etiam ordinatus, Si ab unitate recesserit, sacramentum dandi baptismum [\emph{i.e.}, ordination] non amittit.”} But this of course does not exclude the possibility of suspension and excommunication in case of gross immorality or gross error. The council of Nice, in 325, acknowledged even the validity of the ordination of the schismatic Novatians.

Corresponding to the three ordines majores there were three ordinations: to the diaconate, to the presbyterate, and to the episcopate.\footnote{On the character of the ordination of the sub-deacons, as well as of diaconissae and presbyters, there were afterward diverse views. Usually this was considered ordination only in an improper sense.} Many of the most eminent bishops, however, like Cyprian and Ambrose, received the three rites in quick succession, and officiated only as bishops.

Different from ordination is installation, or induction into a particular congregation or diocese, which may be repeated as often as the minister is transferred.

Ordination was performed by laying on of hands and prayer, closing with the communion. To these were gradually added other preparatory and attendant practices; such as the tonsure\footnote{After the fifth century, but under various forms, tonsura Petri, etc. It was first applied to penitents, then to monks, and finally to the clergy.} the anointing with the chrism (only in the Latin church after Gregory the Great), investing with the insignia of the office (the holy books, and in the case of bishops the ring and staff), the kiss of brotherhood, etc. Only bishops can ordain, though presbyters assist. The ordination or consecraion of a bishop generally requires, for greater solemnity, the presence of three bishops.

No one can receive priestly orders without a fixed field of labor which yields him support.\footnote{Hence the old rules: “Ne quis vage ordinetur,” and, “Nemo ordinatur sine titulo.” Comp. Acts xiv. 23; Tit. i. 5; 1 Pet. v. 1.} In the course of time further restrictions, derived in part from the Old Testament, in regard to age, education, physical and moral constitution, freedom from the bonds of marriage, etc., were established by ecclesiastical legislation.

The favorite times for ordination were Pentecost and the quarterly Quatember terms\footnote{Quatuor tempora. Comp. the old verse: “Post crux (Holyrood day, 14th September), post cineres (Ash Wednesday), post spiritus (Pentecost) atque Luciae (13th December), Sit tibi in auguria quarta sequens feria.”} (i.e., the beginning of Quadragesima, the weeks after Pentecost, after the fourteenth of September, and after the thirteenth of December), which were observed, after Gelasius or Leo the Great, as ordinary penitential seasons of the church. The candidates were obliged to prepare themselves for consecration by prayer and fasting.

\xsection{The Sacrament of the Eucharist}

§ 95. The Sacrament of the Eucharist.

Comp. the Literature in vol. i. § 38 and § 102, the corresponding sections in the Doctrine Histories and Archaeologies, and the treatises of G. E. Steitz on the historical development of the doctrine of the Lord’s Supper in the Greek church, in Dorner’s “Jahrbücher für Deutsche Theologie,” for 1864 and 1865. In part also the liturgical works of Neale, Daniel, etc., cited below (§ 98), and Philip Freeman: The Principles of Divine Service. Lond. Part i. 1855, Part ii. 1862. (The author, in the introduction to the second part, states as his object: “To unravel, by means of an historical survey of the ancient belief concerning the Holy Eucharist, viewed as a mystery, and of the later departures from it, the manifold confusions which have grown up around the subject, more especially since the fatal epoch of the eleventh century.” But the book treats not so much of the doctrine of the Eucharist, as of the ceremony of it, and the eucharistic sacrifice, with special reference to the Anglican church.)

The Eucharist is both a sacrament wherein God conveys to us a certain blessing, and a sacrifice which man offers to God. As a sacrament, or the communion, it stands at the head of all sacred rites; as a sacrifice it stands alone. The celebration of it under this twofold character forms the holy of holies of the Christian cultus in the ancient church, and in the greater part of Christendom at this day.\footnote{Freeman, l.c. Introduction to Part ii. (1857), p. 2, says of the Eucharist, not without justice, from a historical and theological point of view: “It was confessedly through long ages of the church, and is by the vast majority of the Christian world at this hour, conceived to be ... no less than the highest line of contact and region of commingling between heaven and earth known to us, or provided for us;—a borderland of mystery, where, by gradations baffling sight and thought, the material truly blends with the spiritual, and the visible shades off into the unseen; a thing, therefore, which of all events or gifts in this world most nearly answers to the highest aspirations and deepest yearnings of our wonderfully compounded being; while in some ages and climes of the church it has been elevated into something yet more awful and mysterious.”}

We consider first the doctrine of the Eucharist as a sacrament, then the doctrine of the Eucharist as a sacrifice, and finally the celebration of the eucharistic communion and eucharistic sacrifice.

The doctrine of the sacrament of the Eucharist was not a subject of theological controversy and ecclesiastical action till the time of Paschasius Radbert, in the ninth century; whereas since then this feast of the Saviour’s dying love has been the innocent cause of the most bitter disputes, especially in the age of the Reformation, between Papists and Protestants, and among Lutherans, Zwinglians, and Calvinists. Hence the doctrine of the ancient church on this point lacks the clearness and definiteness which the Nicene dogma of the Trinity, the Chalcedonian Christology, and the Augustinian anthropology and soteriology acquired from the controversies preceding them. In the doctrine of baptism also we have a much better right to speak of a consensus patrum, than in the doctrine of the holy Supper.

In general, this period, following the representatives of the mystic theory in the previous one, was already very strongly inclined toward the doctrine of transubstantiation and toward the Greek and Roman sacrifice of the mass, which are inseparable in so far as a real sacrifice requires the real presence of the victim. But the kind and mode of this presence are not yet particularly defined, and admit very different views: Christ may be conceived as really present either in and with the elements (consubstantiation, impanation), or under the illusive appearance of the changed elements (transubstantiation), or only dynamically and spiritually.

In the previous period we distinguish three views: the mystic view of Ignatius, Justin Martyr, and Irenaeus; the symbolical view of Tertullian and Cyprian; and the allegorical or spiritualistic view of Clement of Alexandria and Origen. In the present the first view, which best answered the mystic and superstitious tendency of the time, preponderated, but the second also was represented by considerable authorities.\footnote{Rückert divides the fathers into 2 classes: the \emph{Metabolical}, and the \emph{Symbolical}. The symbolical view he assigns to Tertullian, Clement, Origen, Euseb., Athan., and Augustine. But to this designation there are many objections. Of the Synecdochian (Lutheran) interpretation of the words of institution the ancient church knew nothing.” So says Kahnis, Luth. Dogmatik, ii. p. 221.}

I. The realistic and mystic view is represented by several fathers and the early liturgies, whose testimony we shall further cite below. They speak in enthusiastic and extravagant terms of the sacrament and sacrifice of the altar. They teach a real presence of the body and blood of Christ, which is included in the very idea of a real sacrifice, and they see in the mystical union of it with the sensible elements a sort of repetition of the incarnation of the Logos. With the act of consecration a change accordingly takes place in the elements, whereby they become vehicles and organs of the life of Christ, although by no means necessarily changed into another substance. To denote this change very strong expressions are used, like μεταβολή, μεταβάλλειν, μεταβάλλεσθαι, μεταστοιχειοῦσθαι, μεταποιεῖσθαι, mutatio, translatio, transfiguratio, transformatio;\footnote{But not yet the technical term \emph{transubstantiatio}, which was introduced by Paschasius Radbertus toward the middle of the ninth century, and the corresponding Greek term \textgreek{μετουσίωσις}, which is still later.} illustrated by the miraculous transformation of water into wine, the assimilation of food, and the pervasive power of leaven.

Cyril of Jerusalem goes farther in this direction than any of the fathers. He plainly teaches some sort of supernatural connection between the body of Christ and the elements, though not necessarily a transubstantiation of the latter. Let us hear the principal passages.\footnote{Comp. especially his five mystagogical discourses, addresses to the newly baptized. Cyril’s doctrine is discussed at large in Rückert, Das Abendmahl, sein Wesen u. seine Geschichte, p. 415 ff. Comp. also Neander, Dogmengesch. i. p. 426, and, in part against Rückert, Kahnis, Die Luth. Dogmatik, ii. p. 211 f} “Then follows,” he says in describing the celebration of the Eucharist, “the invocation of God, for the sending of his Spirit to make the bread the body of Christ, the wine the blood of Christ. For what the Holy Ghost touches is sanctified and transformed.” “Under the type of the bread\footnote{\textgreek{Ἐν τύπῳ ἄρτου}, which may mean either under the emblem of the bread (still existing as such), or under the outward form, \emph{sub specie panis.} More naturally the former.} is given to thee the body, under the type of the wine is given to thee the blood, that thou mayest be a partaker of the body and blood of Christ, and be of one body and blood with him.”\footnote{\textgreek{Σύσσωμος καὶσύναιμος αὐτοῦ.}} “After the invocation of the Holy Ghost the bread of the Eucharist is no longer bread, but the body of Christ.” “Consider, therefore, the bread and the wine not as empty elements, for they are, according to the declaration of the Lord, the body and blood of Christ.” In support of this change Cyril refers at one time to the wedding feast at Cana, which indicates, the Roman theory of change of substance; but at another to the consecration of the chrism, wherein the substance is unchanged. He was not clear and consistent with himself. His opinion probably was, that the eucharistic elements lost by consecration not so much their earthly substance, as their earthly purpose.

Gregory of Nyssa, though in general a very faithful disciple of the spiritualistic Origen, is on this point entirely realistic. He calls the Eucharist a food of immortality, and speaks of a miraculous transformation of the nature of the elements into the glorified body of Christ by virtue of the priestly blessing.\footnote{Orat. catech. magna, c. 37. Comp. Neander, l.c. i. p. 428, and Kahnis, ii. 213.}

Chrysostom likewise, though only incidentally in his homilies, and not in the strain of sober logic and theology, but of glowing rhetoric, speaks several times of a union of our whole nature with the body of Christ in the Eucharist, and even of a manducatio oralis.\footnote{Of an \textgreek{ἐμπῆξαιτοὺς ὀδόντας τῇ σαρκὶκαὶσυμπλακῆναι} Comp. the passages from Chrysostomin Ebrard and Rückert, l.c., and Kahnis, ii. p. 215 ff.}

Of the Latin fathers, Hilary,\footnote{De Trinit. viii. 13 sq. Comp. Rückert, l.c. p. 460 ff.} Ambrose,\footnote{De Mysteriis, c. 8 and 9, where a \emph{mutatio} of the \emph{species elementorum} by the word of Christ is spoken of, and the changing of Moses’ rod into a serpent, and of the Nile into blood, is cited in illustration. The genuineness of this small work, however, is doubtful. Rückert considers Ambrosethe pillar of the mediaeval doctrine of the Supper, which he finds in his work De mysteries, and De initiandis.} and Gaudentius († 410) come nearest to the later dogma of transubstantiation. The latter says: “The Creator and Lord of nature, who produces bread from the earth, prepares out of bread his own body, makes of wine his own blood.”\footnote{Serm. p. 42: “Ipse naturarum creator et dominus, qui producit de terra panem, de pane rursus, quia et potest et promisit, efficit proprium corpus, et qui de aqua vinum fecit, facit et de vino sanguinem.” But, on the other hand, Gaudentius (bishop of Brixia) calls the supper a \emph{figure} of the passion of Christ, and the bread the \emph{figure} (figura) of the body of Christ (p. 43). Comp. Rückert, l.c. 477 f.}

But closely as these and similar expressions verge upon the Roman doctrine of transubstantiation, they seem to contain at most a dynamic, not a substantial, change of the elements into the body and the blood of Christ. For, in the first place, it must be remembered there is a great difference between the half-poetic, enthusiastic, glowing language of devotion, in which the fathers, and especially the liturgies, speak of the eucharistic sacrifice, and the clear, calm, and cool language of logic and doctrinal definition. In the second place, the same fathers apply the same or quite similar terms to the baptismal water and the chrism of confirmation, without intending to teach a proper change of the substance of these material elements into the Holy Ghost. On the other hand, they not rarely use, concerning the bread and wine, τύπος, ἀντίτυπα, figura, signum, and like expressions, which denote rather a symbolical than a metabolical relation of them to the body and blood of the Lord. Finally, the favorite comparison of the mysterious transformation with the incarnation of the Logos, which, in fact, was not an annihilation of the human nature, but an assumption of it into unity with the divine, is of itself in favor of the continuance of the substance of the elements; else it would abet the Eutychian heresy.

II. The symbolical view, though on a realistic basis, is represented first by Eusebius, who calls the Supper a commemoration of Christ by the symbols of his body and blood, and takes the flesh and blood of Christ in the sixth chapter of John to mean the words of Christ, which are spirit and life, the true food of the soul, to believers.\footnote{Demonstr. evang. 1, c. 10; Theol. eccl. iii. c. 12, and the fragment of a tract, De paschate, published by Angelo Mai in Scriptorum veterum nova collectio, vol. i. p. 247. Comp. Neander, l.c. i. 430, and especially Steitz, second article (1865), pp. 97-106.} Here appears the influence of his venerated Origen, whose views in regard to the sacramental aspect of the Eucharist he substantially repeats.

But it is striking that even Athanasius, “the father of orthodoxy,” recognized only a spiritual participation, a self-communication of the nourishing divine virtue of the Logos, in the symbols of the bread and wine, and incidentally evinces a doctrine of the Eucharist wholly foreign to the Catholic, and very like the older Alexandrian or Origenistic, and the Calvinistic, though by no means identical with the latter.\footnote{To this result H. Voigt comes, after the most thorough investigation, in his learned monograph on the doctrine of Athanasius, Bremen, 1861, pp. 170-181, and since that time also Steitz, in his second article, already quoted, pp. 109-127. Möhler finds in the passage Ad Scrap. iv. 19 (the principal eucharistic declaration of Athanasius then known), the Roman Catholic doctrine of the Supper (Athanasius der Gr p. 560 ff.), but by a manifestly strained interpretation, and in contradiction with passages in the more recently known Festival Letters of Athanasius, which confirm the exposition of Voigt.} By the flesh and blood in the mysterious discourse of Jesus in the sixth chapter of John, which he refers to the Lord’s Supper, he understands not the earthly, human, but the heavenly, divine manifestation of Jesus, a spiritual nutriment coming down from above, which the Logos through the Holy Ghost communicates to believers (but not to a Judas, nor to the unbelieving).\footnote{So in the main passage, the fourth Epistle to Serapion (Ad Scrap. iv. 19), which properly treats of the sin against the Holy Ghost (c. 8-23), and has been variously interpreted in the interest of different confessions, but now receives new light from several passages in the recently discovered Syriac Festival Letters of Athanasius, translated by Larsow, Leipzig, 1852, pp. 59, 78 sqq., 153 sqq., and especially p. 101.} With this view accords his extending of the participation of the eucharistic food to believers in heaven, and even to the angels, who, on account of their incorporeal nature, are incapable of a corporeal participation of Christ.\footnote{In the Festival Letters in Larsow, p. 101, Athanasius says: “And not only, my brethren, is this bread [of the Eucharist] a food of the righteous, and not only are the saints who dwell on earth nourished with such bread and blood, but also in heaven we eat such food; for even to the higher spirits and the angels the Lord is nutriment, and He is the delight of all the powers of heaven; to all He is all, and over every one He yearns in His love of man.”}

Gregory Nazianzen sees in the Eucharist a type of the incarnation, and calls the consecrated elements symbols and antitypes of the great mysteries, but ascribes to them a saving virtue.\footnote{Orat. xvii. 12; viii. 17; iv. 52. Comp. Ullmann’s Gregor. v. Naz. pp. 483-488; Neander, l.c. i. p. 431; and Steitz in Dorner’s Jahrbücher for 1865, pp. 133-141. Steitz makes Gregory an advocate of the symbolical theory.}

St. Basil, likewise, in explaining the words of Christ, “I live by the Father” (John vi. 57), against, the Arians who inferred from it that Christ was a creature, incidentally gives a spiritual meaning to the fruition of the eucharistic elements. “We eat the flesh of Christ,” he says, “and drink His blood, if we, through His incarnation and human life, become partakers of the Logos and of wisdom.”\footnote{Epist. viii. c. 4 (or Ep. 141 in the older editions): \textgreek{Τρώγομενγὰραὐτοῦτὴνσάρκακαὶπίνομεναὐτοῦτὸαἷμα, κοινωνοὶγινόμενοιδιὰτῆς ἐνανθρωπήσεως καὶτῆς αἰσθητῆς ζωῆς τοῦλόγουκαὶτῆς σοφίας. Σάρκαγὰρκαιαἷμαπασᾶναυτοῦτὴνμυστικὴνἐπιδημίαν}[\emph{i. e.,} a spiritual incarnation, or His internal coming to the soul, as distinct from His historical incarnation] \textgreek{ὠνόμασεκαὶτὴνἐκπρακτικῆς καὶφυσικῆς καὶθεολογικῆς συνεστῶσανδιδασκαλίαν, δἰἧς τρέφεταιψυχὴκαὶπρὸς τῶνὄντωνθεωρίανπαρασκευάζεται. Καὶτοῦτὸἐστιτὸἐκτοῦῥητοῦἴσως δηλούμενον}. This passage, overlooked by Klose, Ebrard, and Kahnis, but noticed by Rückert and more fully by Steitz (l.c. p. 127 ff.), in favor of the symbolical view, is the principal one in Basil on the Eucharist, and must regulate the interpretation of the less important allusions in his other writings.}

Macarius the Elder, a gifted representative of the earlier Greek mysticism († 390), belongs to the same Symbolical school; he calls bread and wine the antitype of the body and blood of Christ, and seems to know only a spiritual eating of the flesh of the Lord.\footnote{Hom. xxvii. 17, and other passages. Steitz (l.c. p. 142 ff.) enters more fully into the views of this monk of the Egyptian desert.}

Theodoret, who was acknowledged orthodox by the council of Chalcedon, teaches indeed a transformation (μεταβάλλειν) of the eucharistic elements by virtue of the priestly consecration, and an adoration of them, which certainly sounds quite Romish, but in the same connection expressly rejects the idea of an absorption of the elements in the body of the Lord, as an error akin to the Monophysite. “The mystical emblems of the body and blood of Christ,” says he, “continue in their original essence and form, they are visible and tangible as they were before [the consecration];\footnote{Dial. ii. Opera ed. Hal. tom. iv. p. 126, where the orthodox man says against the Eranist: \textgreek{Τὰμυστικὰσύμβολα... μένειἐπὶτῆς προτέρας οὐσίας καὶτοῦσχήματος καὶτοῦεἴδους, καὶὁρατάἐστικαὶἁπτὰ, οἶακαὶπρότερονἧν.}} but the contemplation of the spirit and of faith sees in them that which they have become, and they are adored also as that which they are to believers.”\footnote{\textgreek{Προσκυνεῖταιὡς εκεῖναὄνταἅπερπιστεύεται.}These words certainly prove that the consecrated elements are regarded as being not only subjectively, but in some sense objectively and really what the believer takes them for, namely, the body and blood of Christ. But with this they also retained, according to Theodoret, their natural reality and their symbolical character.}

Similar language occurs in an epistle to the monk Caesarius ascribed to Chrysostom, but perhaps not genuine;\footnote{Ep. ad Caesarium monach. (in Chrys. Opera, tom. iii. Pars altera, p. 897 of the new Paris ed. of Montfaucon after the Benedictine): “Sicut enim antequam sanctificetur panis, panem nominamus: divina autem illum sanctificante gratia, mediante sacerdote, liberatus est quidem ab appellatione panis; dignus autem habitus dominici corporis appellatione, \emph{etiamsi natura panis in ipso permansit}, et non duo corpora, sed unum corpus Filii praedicamus.” This epistle is extant in full only in an old Latin version.} in Ephraim of Antioch, cited by Photius; and even in the Roman bishop Gelasius at the end of the fifth century (492–496).

The latter says expressly, in his work against Eutyches and Nestorius: “The sacrament of the body and blood of Christ, which we receive, is a divine thing, because by it we are made partakers of the divine-nature. Yet the substance or nature of the bread and wine does not cease. And assuredly the image and the similitude of the body and blood of Christ are celebrated in the performance of the mysteries.”\footnote{De duabus naturis in Christo Adv. Eutychen et Nestorium (in the Bibl. Max. Patrum, tom. viii. p. 703) ... “et tamen esse non desinit \emph{substantia} vel \emph{natura} panis et vini. Et certe \emph{imago} et \emph{similitudo} corporis et sanguinis Christi in actione mysteriorum celebrantur.” Many Roman divines, through dogmatic prejudice, doubt the genuineness of this epistle. Comp. the Bibl. Max. tom. viii. pp. 699-700.}

It is remarkable that Augustine, in other respects so decidedly catholic in the doctrine of the church and of baptism, and in the cardinal points of the Latin orthodoxy, follows the older African theologians, Tertullian and Cyprian, in a symbolical theory of the Supper, which however includes a real spiritual participation of the Lord by faith, and in this respect stands nearest to the Calvinistic or Orthodox Reformed doctrine, while in minor points he differs from it as much as from transubstantiation and consubstantiation.\footnote{From his immense dogmatic authority, Augustinehas been an apple of contention among the different confessions in all controversies on the doctrine of the Supper. Albertinus (De euchar. pp. 602-742) and Rückert (l.c. p. 353 ff.) have successfully proved that he is no witness for the Roman doctrine; but they go too far when they make him a mere symbolist. That he as little favors the Lutheran doctrine, Kahnis (Vom Abendmahl, p. 221, and in the second part of his Luth. Dogmatik, p. 207) frankly concedes.} He was the first to make a clear distinction between the outward sign and the inward grace, which are equally essential to the conception of the sacrament. He maintains the figurative character of the words of institution, and of the discourse of Jesus, on the eating and drinking of his flesh and blood in the sixth chapter of John; with Tertullian, he calls the bread and wine “figurae” “or “signa corporis et sanguinis Christi” (but certainly not mere figures), and insists on a distinction between “that which is visibly received in the sacrament, and that which is spiritually eaten and drunk,” or between a carnal, visible manducation of the sacrament, and a spiritual eating of the flesh of Christ and drinking of his blood.\footnote{In Psalm. iii. 1: “Convivium, in quo corporis et sanguinis sui \emph{figuram} discipulis commendavit.” Contra Adamant. xii. 3 (”\emph{signum corporis sui}“); Contra advers. legis et prophet. ii. c. 9; Epist. 23; De Doctr. Christ. iii. 10, 16, 19; De Civit. Dei, xxi. c. 20, 25; De peccat. mer. ac rem. ii. 20 ”\emph{quamvis non sit corpus Christi}, sanctum est tamen, quoniam sacramentum est”).} The latter he limits to the elect and the believing, though, in opposition to the subjectivism of the Donatists, he asserts that the sacrament (in its objective import) is the body of Christ even for unworthy receivers. He says of Judas, that he only ate the bread of the Lord, while the other apostles “ate the Lord who was the bread.” In another place: The sacramentum “is given to some unto life, to others unto destruction;” but the res sacramenti, i.e., “the thing itself of which it is the sacramentum, is given to every one who is partaker of it, unto life.” “He who does not abide in Christ, undoubtedly neither eats His flesh nor drinks His blood, though he eats and drinks the sacramentum (i.e., the outward sign) of so great a thing to his condemnation.” Augustine at all events lays chief stress on the spiritual participation. “Why preparest thou the teeth and the belly? Believe, and thou hast eaten.”\footnote{Tract. in Joh. 25: “Quid paras dentes et ventrem? \emph{Crede}, et manducasti.” Comp. Tract. 26: “Qui non manet in Christo, \emph{nec manducat carnem ejus, nec bibit ejus sanguinem} licet premat dentibus sacramentum corporis et sanguinis Christi.”} He claims for the sacrament religious reverence, but not a superstitious dread, as if it were a miracle of magical effect.\footnote{De Trinit. iii. 10: “Honorem tamquam religiosa possunt habere, stuporem tamquam mira non possunt.”} He also expressly rejects the hypothesis of the ubiquity of Christ’s body, which had already come into use in support of the materializing view, and has since been further developed by Lutheran divines in support of the theory of consubstantiation. “The body with which Christ rose,” says he, “He took to heaven, which must be in a place .... We must guard against such a conception of His divinity as destroys the reality of His flesh. For when the flesh of the Lord was upon earth, it was certainly not in heaven; and now that it is in heaven, it is not upon earth.” “I believe that the body of the Lord is in heaven, as it was upon earth when he ascended to heaven.”\footnote{Ep. 146: “Ego Domini corpus ita in coelo esse credo, ut erat in terra, quando ascendit in coelum.” Comp. similar passages in Tract. in Joh. 13; Ep. 187; Serm. 264.} Yet this great church teacher at the same time holds fast the real presence of Christ in the Supper. He says of the martyrs: “They have drunk the blood of Christ, and have shed their own blood for Christ.” He was also inclined, with the Oriental fathers, to ascribe a saving virtue to the consecrated elements.

Augustine’s pupil, Facundus, taught that the sacramental bread “is not properly the body of Christ, but contains the mystery of the body.” Fulgentius of Ruspe held the same symbolical view; and even at a much later period we can trace it through the mighty influence of Augustine’s writings in Isidore of Sevilla, Beda Venerabilis, among the divines of the Carolingian age, in Ratramnus, and Berengar of Tours, until it broke forth in a modified form with greater force than ever in the sixteenth century, and took permanent foothold in the Reformed churches.

Pope Leo I. is sometimes likewise numbered with the symbolists, but without good reason. He calls the communion a “spiritual food,”\footnote{“Spiritualis alimonia.” This expression, however, as the connection of the passage in Serm. lix. 2 clearly shows, by no means excludes an operation of the sacrament on the body; for “spiritual” is often equivalent to “supernatural.” Even Ignatius called the bread of the Supper “a medicine of immortality, and all antidote of death” (\textgreek{φάρμακον ἀθανασίας, ἀντίδοτος τοῦ μὴ ἀποθανεῖν, ἀλλὰ ζῇν ἐν Χριστῷ διὰ παντός}Ad Ephes. c. 20; though this passage is wanting in the shorter Syriac recension.} as Athanasius had done before, but supposes a sort of assimilation of the flesh and blood of Christ by the believing participation. “What we believe, that we receive with the mouth .... The participation of the body and blood of Christ causes that we pass into that which we receive, and bear Christ in us in Spirit and body.” Voluntary abstinence from the wine in the Supper was as yet considered by this pope a sin.\footnote{Comp. the relevant passages from the writings of Leo in Perthel, Papst Leo 1. Leben u. Lehren, p. 216 ff., and in Rückert, l.c. p. 479 ff. Leo’s doctrine of the Supper is not so clearly defined as his doctrine of baptism, and has little that is peculiar. But he certainly had a higher than a purely symbolic view of the sacrament and of the sacrifice of the Eucharist.}

III. The old liturgies, whose testimony on this point is as important as that of the church fathers, presuppose the actual presence of Christ in the Supper, but speak throughout in the stately language of sentiment, and nowhere attempt an explanation of the nature and mode of this presence, and of its relation to the still visible forms of bread and wine. They use concerning the consecrated elements such terms as: The holy body, The dear blood, of our Lord Jesus Christ, The sanctified oblation, The heavenly, spotless, glorious, awful, divine gifts, The awful, unbloody, holy sacrifice, \&c. In the act of consecration the liturgies pray for the sending down of the Holy Ghost, that he may “sanctify and perfect”\footnote{In the liturgy of St. Mark (in Neale’s ed.: The Liturgies of S. Mark, S. James, S. Clement, S. Chrysostom, S. Basil, Lond. 1859, p. 26): \textgreek{Ἵνα αὐτὰ ἁγιάσῃ καὶ τελειώσῃ ... καὶ ποιήσῃ τὸν μὲν ἄρτον σῶμα}, to which the congregation answers: \textgreek{Ἀμήν}.} the bread and wine, or that he may sanctify and make “them the body and blood of Christ,\footnote{In the liturgy of St. James (in Neale, p. 64): \textgreek{Ἵνα ... ἁγιάσῃ καὶ ποιήσῃ τὸν μὲν ἄρτον τοῦτον σῶμα ἅγιον τοῦ Χριστοῦ σου, κ. τ. λ.}} or bless and make.”\footnote{The liturgy of St. Chrysostom(Neale, p. 137) uses the terms \textgreek{εὐλόγησον} and \textgreek{ποίησον}.}

IV. As to the adoration of the consecrated elements: This follows with logical necessity from the doctrine of transubstantiation, and is the sure touchstone of it. No trace of such adoration appears, however, in the ancient liturgies, and the whole patristic literature yields only four passages from which this practice can be inferred; plainly showing that the doctrine of transubstantiation was not yet fixed in the consciousness of the church.

Chrysostom says: “The wise men adored Christ in the manger; we see him not in the manger, but on the altar, and should pay him still greater homage.”\footnote{Hom. 24 in I Cor.} Theodoret, in the passage already cited, likewise uses the term προσκύνεῖν, but at the same time expressly asserts the continuance of the substance of the elements. Ambrose speaks once of the flesh of Christ “which we to-day adore in the mysteries,”\footnote{De Spir. S. iii. II: “Quam [carnem Christi] hodie in mysteriis adoramus, et quam apostoli in Domino Jesu adoraverant.”} and Augustine, of an adoration preceding the participation of the flesh of Christ.\footnote{In Psalm. 98, n. 9: “Ipsam carnem nobis manducandam ad salutem dedit; nemo autem illam carnem manducat nisi prius adoraverit ... et non modo non peccemus adorando, sed peccemus non adorando.”}

In all these passages we must, no doubt, take the term proskunei’n and adorare in the wider sense, and distinguish the bowing of the knee, which was so frequent, especially in the East, as a mere mark of respect, from proper adoration. The old liturgies contain no direction for any such act of adoration as became prevalent in the Latin church, with the elevation of the host, after the triumph of the doctrine of transubstantiation in the twelfth century.\footnote{So says also the Roman liturgist Muratori, De rebus liturgicis, c. xix. p. 227: “Uti omnes inter Catholicos eruditi fatentur, post \emph{Berengarii haeresiam} ritus in Catholica Romana ecclesia invaluit, scilicet post consecrationem elevare hostiam et calicem, ut a populo adoretur corpus et sanguis Domini.” Freeman, Principles of Div. Service, Introduction to Part ii. p. 169, asserts: “The Church throughout the world, down to the period of the unhappy change of doctrine in the Western church in the eleventh and twelfth centuries, never worshipped either the consecrated elements on account of their being the body and blood of Christ, or the presence of that body and blood; nor again, either Christ Himself as supernaturally present by consecration, or the presence of His divinity; neither have the churches of God to this hour, with the exception of those of the Roman obedience, any such custom.”}

\xsection{The Sacrifice of the Eucharist}

§ 96. The Sacrifice of the Eucharist.

Besides the works already cited on the holy Supper, comp. Höfling: Die Lehre der ältesten Kirche vom Opfer im Leben u. Cultus der Kirche. Erlangen, 1851. The articles: Messe, Messopfer, in Wetzer u. Welte: Kirchenlexicon der kathol. Theologie, vol. vii. (1851), p. 83 ff. G. E. Steitz: Art. Messe u. Messopfer in Herzog’s Protest. Real-Encyklopädie, vol. ix. (1858), pp. 375–408. Phil. Freeman: The Principles of Divine Service. Part ii. Oxf. and Lond. 1862. This last work sets out with a very full consideration of the Mosaic sacrificial cultus, and (in the Pref. p. vi.) unjustly declares all the earlier English and German works of Mede, Outram, Patrick, Magee, Bähr, Hengstenberg, and Kurtz, on this subject, entirely unsatisfactory and defective.

The Catholic church, both Greek and Latin, sees in the Eucharist not only a sacramentum, in which God communicates a grace to believers, but at the same time, and in fact mainly, a sacrificium, in which believers really offer to God that which is represented by the sensible elements. For this view also the church fathers laid the foundation, and it must be conceded they stand in general far more on the Greek and Roman Catholic than on the Protestant side of this question. The importance of the subject demands a preliminary explanation of the idea of sacrifice, and a clear discrimination of its original Christian form from its later perversion by tradition.

The idea of sacrifice is the centre of all ancient religions, both the heathen and the Jewish. In Christianity it is fulfilled. For by His one perfect sacrifice on the cross Christ has entirely blotted out the guilt of man, and reconciled him with the righteous God. On the ground of this sacrifice of the eternal High Priest, believers have access to the throne of grace, and may expect their prayers and intercessions to be heard. With this perfect and eternally availing sacrifice the Eucharist stands in indissoluble connection. It is indeed originally a sacrament and the main thing in it is that which we receive from God, not that which we give to God. The latter is only a consequence of the former; for we can give to God nothing which we have not first received from him. But the Eucharist is the sacramentum of a sacrificium, the thankful celebration of the sacrificial death of Christ on the cross, and the believing participation or the renewed appropriation of the fruits of this sacrifice. In other words, it is a feast on a sacrifice. “As oft as ye do eat this bread and drink this cup, ye do show the Lord’s death till He come.”

The Eucharist is moreover, as the name itself implies, on the part of the church a living and reasonable thank-offering, wherein she presents herself anew, in Christ and on the ground of his sacrifice, to God with prayers and intercessions. For only in Christ are our offerings acceptable to God, and only through the continual showing forth and presenting of His merit can we expect our prayers and intercessions to be heard.

In this view certainly, in a deep symbolical and ethical sense, Christ is offered to God the Father in every believing prayer, and above all in the holy Supper; i.e. as the sole ground of our reconciliation and acceptance. This is the deep truth which lies at the bottom of the Catholic mass, and gives it still such power over the religious mind.\footnote{Freeman states the result of his investigation of the Biblical sacrificial cultus and of the doctrine of the old Catholic church on the eucharistic sacrifice, as follows, on p. 280: “It is enough for us that the holy Eucharist is all that the ancient types foreshowed that it would be; that in it we present ’memorially,’ yet truly and with prevailing power, by the consecrating Hands of our Great High Priest, the wondrous Sacrifice once for all offered by Him at the Eucharistic Institution, consummated on the Cross, and ever since presented and pleaded by Him, Risen and Ascended, in Heaven; that our material Gifts are identified with that awful Reality, and as such are borne in upon the Incense of His Intercession, and in His Holy Hands, into the True Holiest Place: that we ourselves, therewith, are home in thither likewise, and abide in a deep mystery in the heavenly places in Christ Jesus; that thus we have all manner of acceptance,—sonship, kingship, and priesthood unto God; an our whole life, in all its complex action, being sanctified and purified for such access, and abiding continually in a heavenly sphere of acceptableness and privilege.—Enough for us, again, that on the sacramental side of the mystery, we have been thus privileged to give to God His own Gift of Himself to dwell in us, and we in Him;—-that we thereby possess an evermore renewedly dedicated being—strengthened with all might, and evermore made one with Him. Profoundly reverencing Christ’s peculiar Presence in us and around us in the celebration of such awful mysteries, we nevertheless take as the watchword of our deeply mysterious Eucharistic worship, ’Sursum corda,’ and ’Our life is hid with Christ in God.’ “}

But this idea in process of time became adulterated with foreign elements, and transformed into the Graeco-Roman doctrine of the sacrifice of the mass. According to this doctrine the Eucharist is an unbloody repetition of the atoning sacrifice of Christ by the priesthood for the salvation of the living and the dead; so that the body of Christ is truly and literally offered every day and every hour, and upon innumerable altars at the same time. The term mass, which properly denoted the dismissal of the congregation (missio, dismissio) at the close of the general public worship, became, after the end of the fourth century, the name for the worship of the faithful,\footnote{The \emph{missa fidelium}, in distinction from the \emph{missa catechumenorum}. Comp. 90 above.} which consisted in the celebration of the eucharistic sacrifice and the communion. The corresponding terms of the Orientals are λειτουργία, θυσία, προσφορά.

In the sacrifice of the mass the whole mysterious fulness and glory of the Catholic worship is concentrated. Here the idea of the priesthood reaches its dizzy summit; and here the devotion and awe of the spectators rises to the highest pitch of adoration. For to the devout Catholic nothing can be greater or more solemn than an act of worship in which the eternal Son of God is veritably offered to God upon the altar by the visible hand of the priest for the sins of the world. But though the Catholic worship here rises far above the vain sacrifices of heathendom and the merely typical sacrifices of Judaism, yet that old sacrificial service, which was interwoven with the whole popular life of the Jewish and Graeco-Roman world, exerted a controlling influence on the Roman Catholic service of the Eucharist, especially after the nominal conversion of the whole Roman heathendom, and obscured the original simplicity and purity of that service almost beyond recognition. The sacramentum became entirely eclipsed by the sacrificium, and the sacrificium became grossly materialized, and was exalted at the expense of the sacrifice on the cross. The endless succession of necessary repetitions detracts from the sacrifice of Christ.

The Biblical support of the sacrifice of the mass is weak, and may be reduced to an unduly literal interpretation or a downright perversion of some such passages as Mal. i. 10 f.; 1 Cor. x. 21; Heb. v. 6; vii. 1 f.; xiii. 10. The Epistle to the Hebrews especially is often misapplied, though it teaches with great emphasis the very opposite, viz., the abolition of the Old Testament sacrificial system by the Christian worship, the eternal validity of the sacrifice of our only High Priest on the right hand of the Father, and the impossibility of a repetition of it (comp. x. 14; vii. 23, 24).

We pass now to the more particular history. The ante-Nicene fathers uniformly conceived the Eucharist as a thank-offering of the church; the congregation offering the consecrated elements of bread and wine, and in them itself, to God.\footnote{Comp. vol. i. § 102, p. 389 ff.} This view is in itself perfectly innocent, but readily leads to the doctrine of the sacrifice of the mass, as soon as the elements become identified with the body and blood of Christ, and the presence of the body comes to be materialistically taken. The germs of the Roman doctrine appear in Cyprian about the middle of the third century, in connection with his high-churchly doctrine of the clerical priesthood. Sacerdotium and sacrificium are with him correlative ideas, and a Judaizing conception of the former favored a like Judaizing conception of the latter. The priest officiates in the Eucharist in the place of Christ,\footnote{“Vice Christi vere fungitur.”} and performs an actual sacrifice in the church.\footnote{“Sacrificium verum et plenum offert in ecclesia Patri.”} Yet Cyprian does not distinctly say that Christ is the subject of the spiritual sacrifice; rather is the mystical body of Christ, the Church, offered to God, and married with Christ.\footnote{Epist. 63 ad Caecil.c. 14. Augustine’s view is similar: the church offering herself to God in and with Christ as her Head.}

The doctrine of the sacrifice of the mass is much further developed in the Nicene and post-Nicene fathers, though amidst many obscurities and rhetorical extravagances, and with much wavering between symbolical and grossly realistic conceptions, until in all essential points it is brought to its settlement by Gregory the Great at the close of the sixth century. These points are the following:

1. The eucharistic sacrifice is the most solemn mystery of the church, and fills the faithful with a holy awe. Hence the predicates θυσία φοβερὰ, φρικτὴ, ἀναίμακτος, sacrificium tremendum, which are frequently applied to it, especially in the Oriental liturgies and homilies. Thus it is said in the liturgy of St. James: “We offer to Thee, O Lord, this awful and unbloody sacrifice.” The more surprising is it that the people should have been indifferent to so solemn an act, and that Chrysostom should lament: “In vain is the daily sacrifice, in vain stand we at the altar; there is no one to take part.”\footnote{Hom. iii. in Ep. ad Ephes. (new Par. Bened. ed. tom. xi. p. 26): \textgreek{Εἰκῇ θυσία καθημερινὴ, εἰκῇ παρεστήκαμεν τῷ θυσιαστηρίῳ, οὐδεὶς ὁ μετέχων}\emph{, i.e.}, Frustra est quotidianum sacrificium, frustra adstamus altari: nemo est qui participet.}

2. It is not a new sacrifice added to that of the cross, but a daily, unbloody repetition and perpetual application of that one only sacrifice. Augustine represents it, on the one hand, as a sacramentum memoriae a symbolical commemoration of the sacrificial death of Christ; to which of course there is no objection.\footnote{Contr. Faust. Manich. l. xx. 18: “Unde jam Christiani, \emph{peracti ejusdem sacrificii memoriam celebrant, sacrosancta oblatione et participatione corporis et sanguinis Christi.}“ Comp. l. xx. 21. This agrees with Augustine’s symbolical conception of the consecrated elements as signal imagines, similitudines corporis et sanguinis Christi. Steitz, l.c. p. 379, would make him altogether a symbolist, but does not succeed; comp. the preceding section, and Neander, Dogmengesch. i. p. 432.} But, on the other hand, he calls the celebration of the communion verissimum sacrificium of the body of Christ. The church, he says, offers (immolat) to God the sacrifice of thanks in the body of Christ, from the days of the apostles through the sure succession of the bishops down to our time. But the church at the same time offers, with Christ, herself, as the body of Christ, to God. As all are one body, so also all are together the same sacrifice.\footnote{De Civit. Dei, x. 20: “Per hoc [homo Jesus Christus] et sacerdos est ipse offerens, ipse et oblatio. Cujus rei sacramentum quotidianum esse voluit ecclesiae sacrificium, quae cum ipsius capitis corpus sit, se ipsam per ipsum offere discit.” And the faithful in heaven form with us one sacrifice, since they with us are one civitas Dei.} According to Chrysostom the same Christ, and the whole Christ, is everywhere offered. It is not a different sacrifice from that which the High Priest formerly offered, but we offer always the same sacrifice, or rather, we perform a memorial of this sacrifice.\footnote{Hom. xvii. in Ep. ad Hebr. tom. xii. pp. 241 and 242: \textgreek{Τοῦτο γὰρ ποιεῖτε φησὶν, εἰς τὴν ἐμὴν ἀνάμνησιν. Οὐκ ἄλλην θυσίαν, καθάπερ ὁ ἀρχιερεὺς τότε, ἀλλὰ τὴν αὐτὴν ἀεὶ ποιοῦμεν; μᾶλλον δὲ ἀνάμνησιν ἐργαζόμεθα θυσίας}.} This last clause would decidedly favor a symbolical conception, if Chrysostom in other places had not used such strong expressions as this: “When thou seest the Lord slain, and lying there, and the priest standing at the sacrifice,” or: “Christ lies slain upon the altar.”\footnote{De sacerd. iii. c. 4 (tom. i. 467): \textgreek{Ὅτανἰδῇς τὸνΚύριοντεθυμένονκαὶκείμενον, καὶτὸνἱερέαἐφεστῶτατῷθύματι, καὶἐπευχόμενον, κ. τ. λ}. Homil. xv. ad Popul. Antioch. c. 5 (tom. ii. p. 187): \textgreek{ἜνθαὁΧριστὸς κεῖταιτεθυμένος}. Comp. Hom. in tom. ii. p. 394, where it is said of the sacrifice of the Eucharist: \textgreek{Θυσίᾳ προσέρχῃ φρικτῇκαὶἀγιᾴ; ἐσφαγμένος πρόκειταιὁΧριστός.}}

3. The sacrifice is the anti-type of the Mosaic sacrifice, and is related to it as substance to typical shadows. It is also especially foreshadowed by Melchizedek’s unbloody offering of bread and wine. The sacrifice of Melchizedek is therefore made of great account by Hilary, Jerome, Augustine, Chrysostom, and other church fathers, on the strength of the well-known parallel in the seventh chapter of the Epistle to the Hebrews.

4. The subject of the sacrifice is the body of Jesus Christ, which is as truly present on the altar of the church, as it once was on the altar of the cross, and which now offers itself to God through his priest. Hence the frequent language of the liturgies: “Thou art he who offerest, and who art offered, O Christ, our God.” Augustine, however, connects with this, as we have already said, the true and important moral idea of the self-sacrifice of the whole redeemed church to God. The prayers of the liturgies do the same.\footnote{Freeman regards this as the main thing in the old liturgies. “In all liturgies,” says he, l.c. p. 190, “the Church has manifestly two distinct though closely connected objects in view. The first is, \emph{to offer herself in Christ to} \emph{God;} or rather, in strictness and as the highest conception of her aim, to \emph{procure that she may be offered by Christ Himself, and as in Christ, to the} \emph{Father.} And the second object, as the crowning and completing feature of the rite, and woven up with the other in one unbroken chain of service, is \emph{to obtain communion} \emph{through Christ with God;} or, more precisely again, that \emph{Christ} \emph{may Himself give her, through Himself, such communion}.”}

5. The offering of the sacrifice is the exclusive prerogative of the Christian priest. Later Roman divines take the words: “This do (ποιεῖτε) in remembrance of me,” as equivalent to: “This offer,” and limit this command to the apostles and their successors in office, whereas it is evidently an exhortation to all believers to the commemoration of the atoning death, the communio sacramenti, and not to the immolatio sacrificii.

6. The sacrifice is efficacious for the whole body of the church, including its departed members, in procuring the gifts which are implored in the prayers of the service.

All the old liturgies proceed under a conviction of the unbroken communion of saints, and contain commemorations and intercessions for the departed fathers and brethren, who are conceived to be, not in purgatory, but in communion with God and in a condition of progressive holiness and blessedness, looking forward in pious longing to the great day of consummation.

These prayers for an increase of bliss, which appeared afterwards very inappropriate, form the transition from the original simple commemoration of the departed saints, including the patriarchs, prophets and apostles, to intercessions for the suffering souls in purgatory, as used in the Roman church ever since the sixth century.\footnote{Neale has collected in an appendix to his English edition of the old liturgies (The Liturgies of S. Mark, S. James, etc., Lond. 1859, p. 216 ff.) the finest liturgical prayers of the ancient church for the departed saints, and deduces from them the positions, ”(1) that prayers for the dead, and more especially the oblation of the blessed Eucharist for them, have been from the beginning the practice of the Universal Church. (2) And this without any idea of a purgatory of pain, or of any state from which the departed soul has to be delivered as from one of misery.” The second point needs qualification.} In the liturgy of Chrysostom, still in use in the Greek and Russian church, the commemoration of the departed reads. “And further we offer to thee this reasonable service on behalf of those who have departed in the faith, our ancestors, Fathers, Patriarchs, Prophets, Apostles, Preachers, Evangelists, Martyrs, Confessors, Virgins, and every just spirit made perfect in the faith .... Especially the most holy, undefiled, excellently laudable, glorious Lady, the Mother of God and Ever-Virgin Mary .... the holy John the Prophet, Forerunner and Baptist, the holy, glorious and all-celebrated Apostles, and all thy Saints, through whose prayers look upon us, O God. And remember all those that are departed in the hope of the resurrection to eternal life, and give them rest where the light of Thy countenance shines upon them.”

Cyril of Jerusalem, in his fifth and last mystagogic Catechesis, which is devoted to the consideration of the eucharistic sacrifice and the liturgical service of God, gives the following description of the eucharistic intercessions for the departed: “When the spiritual sacrifice, the unbloody service of God, is performed, we pray to God over this atoning sacrifice for the universal peace of the church, for the welfare of the world, for the emperor, for soldiers and prisoners, for the sick and afflicted, for all the poor and needy. Then we commemorate also those who sleep, the patriarchs, prophets, apostles, martyrs, that God through their prayers and their intercessions may receive our prayer; and in general we pray for all who have gone from us, since we believe that it is of the greatest help to those souls for whom the prayer is offered, while the holy sacrifice, exciting a holy awe, lies before us.”\footnote{\textgreek{Τῆς ἁγίας καὶφρικωδεστάτης προκειμένης θυσίας}, Catech. xxiii. 8.}

This is clearly an approach to the later idea of purgatory in the Latin church. Even St. Augustine, with Tertullian, teaches plainly, as an old tradition, that the eucharistic sacrifice, the intercessions or suffragia and alms, of the living are of benefit to the departed believers, so that the Lord deals more mercifully with them than their sins deserve.\footnote{Serm. 172, 2 (Opp. tom. v. 1196): “Orationibus sanctae ecclesiae, et sacrificio salutari, et eleemosynis, quae pro eorum spiritibus erogantur, non est dubitandum mortuos adjuvari, ut cum eis misericordius agatur a Domino.” He expressly limits this effect, however, to those who have departed in \emph{the faith.}} His noble mother, Monica, when dying, told him he might bury her body where he pleased, and should give himself no concern for it, only she begged of him that he would remember her soul at the altar of the Lord.\footnote{Confess. l. ix. 27: “Tantum illud vos rogo, ut ad Domini altare memineritis mei, ubi fueritis.” Tertullianconsiders it the duty of a devout widow to pray for the soul of her husband, and to offer a sacrifice on the anniversary of his death; De monogam. c. 10. Comp. De corona, c. 2: “Oblationes pro defunctis pro natalitiis annua die facimus.”}

With this is connected the idea of a repentance and purification in the intermediate state between death and resurrection, which likewise Augustine derives from Matt. xii. 32, and 1 Cor. iii. 15, yet mainly as a mere opinion.\footnote{De Civit. Dei, xxi. 24, and elsewhere. The passages of Augustineand the other fathers in favor of the doctrine of purgatory are collected in the much-cited work of Berington and Kirk: The Faith of Catholics, etc., vol. iii. pp. 140-207.} From these and similar passages, and under the influence of previous Jewish and heathen ideas and customs, arose, after Gregory the Great, the Roman doctrine of the purgatorial fire for imperfect believers who still need to be purified from the dross of their sins before they are fit for heaven, and the institution of special masses for the dead, in which the perversion of the thankful remembrance of the one eternally availing sacrifice of Christ reaches its height, and the idea of the communion utterly disappears.\footnote{There are silent masses, missae solitariae, at which usually no one is present but the priest, with the attendant boys, who offers to God at a certain tariff the magically produced body of Christ for the deliverance of a soul from purgatory. This institution has also a heathen precedent in the old Roman custom of offering sacrifices to the Manes of beloved dead. On Gregory’s doctrine of the mass, which belongs in the next period, comp. the monograph of Lau, p. 484f. The horrible abuse of these masses for the dead, and their close connection with superstitious impostures of purgatory and of indulgence, explain the moral anger of the Reformers at the mass, and the strong declarations against it in several symbolical books, especially in the Smalcald Articles by Luther (ii. 2, where the mass is called \emph{draconis cauda}), and in the Heidelberg Catechism (the 80th question, which, by the way, is wanting entirely in the first edition of 1563, and was first inserted in the second edition by express command of the Elector Friedrich III., and in the third edition was enriched with the epithet “damnable idolatry”).}

In general, in the celebration of the Lord’s Supper the sacrament continually retired behind the sacrifice. In the Roman churches in all countries one may see and hear splendid masses at the high altar, where the congregation of the faithful, instead of taking part in the communion, are mere spectators of the sacrificial act of the priest. The communion is frequently despatched at a side altar at an early hour in the morning.

\xsection{The Celebration of the Eucharist}

§ 97. The Celebration of the Eucharist.

Comp. the Liturgical Literature cited in the next section, especially the works of Daniel, Neale, and Freeman.

The celebration of the eucharistic sacrifice and of the communion was the centre and summit of the public worship of the Lord’s day, and all other parts of worship served as preparation and accompaniment. The old liturgies are essentially, and almost exclusively, eucharistic prayers and exercises; they contain nothing besides, except some baptismal formulas and prayers for the catechumens. The word liturgy (λειτουργία), which properly embraces all parts of the worship of God, denotes in the narrower sense a celebration of the eucharist or the mass.

Here lies a cardinal difference between the Catholic and Evangelical cultus: in the former the sacrifice of the mass, in the latter the sermon, is the centre.

With all variations in particulars, especially in the introductory portions, the old Catholic liturgies agree in the essential points, particularly in the prayers which immediately precede and follow the consecration of the elements. They all (excepting some Syriac copies of certain Nestorian and Monophysite formularies) repeat the solemn Words of Institution from the Gospels,\footnote{Though in various forms. See below.} understanding them not merely in a declaratory but in an operative sense; they all contain the acts of Consecration, Intercession, and Communion; all (except the Roman) invoke the Holy Ghost upon the elements to sanctify them, and make them actual vehicles of the body and blood of Christ; all conceive the Eucharist primarily as a sacrifice, and then, on the basis of the sacrifice, as a communion.

The eucharistic action in the narrower sense is called the Anaphora, or the canon missae, and begins after the close of the service of the catechumens (which consisted principally of reading and preaching, and extended to the Offertory, i.e., the preparation of the bread and wine, and the placing of it on the altar). It is introduced with the Ἄνω τὰς καρδίας, or Sursum corda, of the priest: the exhortation to the faithful to lift up their hearts in devotion, and take part in the prayers; to which the congregation answers: Habemus ad Dominum, “We lift them up unto the Lord.” Then follows the exhortation: “Let us give thanks to the Lord,” with the response: “It is meet and right.”\footnote{Or, according to the Liturgia S. Jacobi: \textgreek{Ἄνω σχῶμεν τὸν νοῦν καὶ τὰς καρδίας}\emph{,} with the response: \textgreek{Ἄξιον καὶ δίκαιον}. In the Lit. S. Clem.: Priest: \textgreek{Ἄνω τὸν νοῦν}. \emph{All} (\textgreek{πάντες}): \textgreek{Ἔχομεν πρὸς τὸν Κύριον}.—\textgreek{Εὐχαριστήσωμεν τῷ Κυρίῳ}. Resp.:\textgreek{Ἄξιον καὶ δίκαιον}. In the Lit. S. Chrys. (still in use in the orthodox Greek and Russian church): \par \textgreek{Ὁ ἱερεύς· Ἄνω σχῶμεν τὰς καρδίας.} \par \textgreek{Ὁ χορός · Ἔχομεν πρὸ τὸν Κύριον.} \par \textgreek{Ὁ ἱερεύς· Εὐχαριστήσωμεν τῷ Κυρίῳ.} \par \textgreek{Ὁ χορός · Ἄξιον καὶ δίκαιον ἐστὶ προσκύνεῖν Πατέρα}, \textgreek{Υἱόν}, \textgreek{καὶ} \par \textgreek{ἂγιον Πνεῦμα, Τριάδα ὁμοούσιον καὶ ἀχώριστον}.}

The first principal act of the Anaphora is the great prayer of thanksgiving, the εὐλογίαor εὐχαριστία, after the example of the Saviour in the institution of the Supper. In this prayer the priest thanks God for all the gifts of creation and of redemption, and the choir generally concludes the thanksgiving with the so-called Trisagion or Seraphic Hymn (Is. vi. 3), and the triumphal Hosanna (Matt. xx. 9): “Holy, Holy, Holy Lord of Sabaoth; heaven and earth are full of Thy glory. Hosanna in the highest: blessed is He that cometh in the name of the Lord: Hosanna in the highest.”

Then follows the consecration and oblation of the elements, by the commemoration of the great facts in the life of Christ, by the rehearsing of the Words of Institution from the Gospels or from Paul, and by the invocation of the Holy Ghost, who brings to pass the mysterious change of the bread and wine into the sacramental body and blood of Christ.\footnote{Hence it is said, for example, in the Syriac version of the Liturgy of St. James: “How dreadful is this hour, in which the Holy Ghost hastens to come down from the heights of heaven, and broods over the Eucharist, and sanctifies it. In holy silence and fear stand and pray.”} This invocation of the Holy Ghost\footnote{\textgreek{Ἐπίκλησις Πνεύματος ἁγίου}, invocatio Spiritus Sancti.} appears in all the Oriental liturgies, but is wanting in the Latin church, which ascribes the consecration exclusively to the virtue of Christ’s Words of Institution. The form of the Words of Institution is different in the different liturgies.\footnote{They are collected by Neale, in his English edition of the Primitive Liturgies, pp. 175-215, from 67 ancient liturgies in alphabetical order. Freeman says, rather too strongly, l.c. p. 364: “No two churches in the world have even the same Words of Institution.”} The elevation of the consecrated elements was introduced in the Latin church, though not till after the Berengarian controversies in the eleventh century, to give the people occasion to show, by the adoration of the host, their faith in the real presence of Christ in the sacrament.

To add an example: The prayer of consecration and oblation in one of the oldest and most important of the liturgies, that of St. James, runs thus: After the Words of Institution the priest proceeds:

“Priest: We sinners, remembering His life-giving passion, His saving cross, His death, and His resurrection from the dead on the third day, His ascension to heaven, and His sitting at the right hand of Thee His God and Father, and His glorious and terrible second appearing, when He shall come in glory to judge the quick and the dead, and to render to every man according to his works,—offer to Thee, O Lord, this awful and unbloody sacrifice;\footnote{\textgreek{Προσφέρομέν σοι, Δέσποτα, τὴν φοβερὰν ταύτην καὶ ἀναίμακτον θυσίαν}. The term \textgreek{φοβερά}denotes holy awe, and is previously applied also to the second coming of Christ: \textgreek{Τῆς δευτέρας ἐνδόξου καὶ φοβερᾶς αὐτοῦ παρουσίας}, sc. \textgreek{μεμνημένοι.} The Liturgy of St. Chrysostomhas instead:\textgreek{Προσφέρομέν σοι τὴν λοφικὴν ταύτην καὶ ἀναίμακτον λατρείαν}(doubtless with reference to the\textgreek{λογικὴ λατρεία} in Rom. xii 1).} beseeching Thee that Thou wouldst deal with us not after our sins nor reward us according to our iniquities, but according to Thy goodness and unspeakable love to men wouldst blot out the handwriting which is against us Thy suppliants, and wouldst vouchsafe to us Thy heavenly and eternal gifts, which eye hath not seen, nor ear heard, neither hath it entered into the heart of man what Thou, O God, hast prepared for them that love Thee. And reject not Thy people, O loving Lord, for my sake and on account of my sins.

He repeats thrice: For Thy people and Thy Church prayeth to Thee.

People: Have mercy upon us, O Lord God, almighty Father!

Priest: Have mercy upon us, almighty God!

Have mercy upon us, O God, our Redeemer I

Have mercy upon us, O God, according to Thy great mercy, and send upon us, and upon these gifts here present, Thy most holy Spirit, Lord, Giver of life, who with Thee the God and Father, and with Thine only begotten Son, sitteth and reigneth upon one throne, and is of the same essence and co-eternal,\footnote{\textgreek{Ἐξαπόστειλον ἐφ  ἡμᾶς καὶ ἐπὶ τὰ προκείμενα δῶρα ταῦτα τὸ Πνεῦμά σου τὸ πανάγιον,} [\textgreek{εἶτα κλίνας τὸν αὐχένα λέγει·}] \textgreek{τὸ κύριον καὶ ζωοποιὸν, τὸ σύνθρονον σοὶ τῷ Θεῷ καὶ Πατρὶ, καὶ τῷ μονογενεῖ σου Υίῷ, τὸ συμβασιλεῦον, τὸ ὁμοούσιόν τε καὶ συναίδιον}. The \textgreek{ὁμοούσιον}, as well as the Nicene Creed in the preceding part of the Liturgy of St. James, indicates clearly a post-Nicene origin.} who spoke in the law and in the prophets, and in Thy new covenant, who descended in the form of a dove upon our Lord Jesus Christ in the river Jordan, and rested upon Him, who came down upon Thy holy apostles in the form of tongues of fire in the upper room of Thy holy and glorious Zion on the day of Pentecost: send down, O Lord, the same Holy Ghost upon us and upon these holy gifts here present, that with His holy and good and glorious presence He may sanctify this bread and make it the holy body of Thy Christ.\footnote{\textgreek{ἽΙνα ... }α\textgreek{γιάσῃ καὶ ποιήσῃ τὸν μὲν ἂρτον τοῦτον σῶμα ἂγιον τοῦ Χριστοῦ σου.}}

People: Amen.

Priest: And this cup the dear blood of Thy Christ.

People: Amen.

Priest (in a low voice): That they may avail to those who receive them, for the forgiveness of sins and for eternal life, for the sanctification of soul and body, for the bringing forth of good works, for the strengthening of Thy holy Catholic church which Thou hast built upon the rock of faith, that the gates of hell may not prevail against her; delivering her from all error and all scandal, and from the ungodly, and preserving her unto the consummation of all things.”

After the act of consecration come the intercessions, sometimes very long, for the church, for all classes, for the living, and for the dead from righteous Abel to Mary, the apostles, the martyrs, and the saints in Paradise; and finally the Lord’s Prayer. To the several intercessions, and the Lord’s Prayer, the people or the choir responds Amen. With this closes the act of eucharistic sacrifice.

Now follows the communion, or the participation of the consecrated elements. It is introduced with the words: “Holy things for holy persons,”\footnote{\textgreek{Τὰ ἅγια τοῖς ἁγίοις}, Sancta Sanctis. It is a warning to the unworthy not to approach the table of the Lord.} and the Kyrie eleison, or (as in the Clementine liturgy) the Gloria in Excelsis: “Glory be to God on high, peace on earth, and good will to men.\footnote{According to the usual reading \textgreek{ἐν ἀνθρώποις εὐδοκία}. But the older and better attested reading is \textgreek{εὐδοκίας}, which alters the sense and makes the angelic hymn bimembris: “Glory to God in the highest, and on earth peace among men of His good pleasure (\emph{i.e.}, the chosen people of God).”} Hosanna to the Son of David! Blessed is he that cometh in the name of the Lord: God is the Lord, and he hath appeared among us.” The bishop and the clergy communicate first, and then the people. The formula of distribution in the Clementine liturgy is simply: “The body of Christ;” “The blood of Christ, the cup of life,”\footnote{\textgreek{Σῶμα Χριστοῦ} —\textgreek{Αἷμα Χριστοῦ, ποτήριον ζωῆς}.} to which the receiver answers “Amen.” In other liturgies it is longer.\footnote{In the Liturgy of St. Mark::\textgreek{Σῶμα ἅγιον} —\textgreek{Αἷμα τίμιον τοῦ Κυρίου καὶ Θεοῦ καὶΣωτῆρος ἡμῶν}. In the Mozarabic Liturgy the communicating priest prays: “Corpus et sanguis Domini nostri Jesu Christi custodiat corpus et animam meam (tuam) in vitam aeternam.” Resp.: ”\emph{Amen}.” So in the Roman Liturgy, from which it passed into the Anglican.}

The holy act closes with prayers of thanksgiving, psalms, and the benediction.

The Eucharist was celebrated daily, or at least every Sunday. The people were exhorted to frequent communion, especially on the high festivals. In North Africa some communed every day, others every Sunday, others still less frequently.\footnote{Augustine, Epist. 118 ad Janua c. 2: “Alii quotidie communicant corpori et sanguini Dominico; alii certis diebus accipiunt; alibi nullus dies intermittitur quo non offeratur; alii sabbato tantum et dominico; alibi tantum dominico.”} Augustine leaves this to the needs of every believer, but says in one place: “The Eucharist is our daily bread.” The daily communion was connected with the current mystical interpretation of the fourth petition in the Lord’s Prayer. Basil communed four times in the week. Gennadius of Massilia commends at least weekly communion. In the East it seems to have been the custom, after the fourth century, to commune only once a year, or on great occasions. Chrysostom often complains of the indifference of those who come to church only to hear the sermon, or who attend the eucharistic sacrifice, but do not commune. One of his allusions to this neglect we have already quoted. Some later councils threatened all laymen with excommunication, who did not commune at least on Christmas, Easter, and Pentecost.

In the Oriental and North African churches prevailed the incongruous custom of infant communion, which seemed to follow from infant baptism, and was advocated by Augustine and Innocent I. on the authority of John vi. 53. In the Greek church this custom continues to this day, but in the Latin, after the ninth century, it was disputed or forbidden, because the apostle (1 Cor. xi. 28, 29) requires self-examination as the condition of worthy participation.\footnote{Comp. P. Zorn: Historia eucharistiae infantum, Berl. 1736; and the article by Kling in Herzog’s Encykl. vii. 549 ff.}

With this custom appear the first instances, and they exceptional, of a communio sub una specie; after a little girl in Carthage in the time of Cyprian had been made drunk by receiving the wine. But the withholding of the cup from the laity, which transgresses the express command of the Lord: “Drink ye all of it,” and is associated with a superstitious horror of profaning the blood of the Lord by spilling, and with the development of the power of the priesthood, dates only from the twelfth and thirteenth centuries and was then justified by the scholastic doctrine of concomitance.

In the Greek church it was customary to dip the bread in the wine, and deliver both elements in a spoon.

The customs of house-communion and after-communion for the sick and for prisoners, of distributing the unconsecrated remainder of the bread among the non-communicants, and of sending the consecrated elements, or their substitutes,\footnote{These substitutes for the consecrated elements were called \textgreek{ἀντίδωρα} (\emph{i.e.}, \textgreek{ἀντὶ τῶν δώρων εὐχαριστικῶν}), and \emph{eulogiae} (from the benediction at the close of the service).} to distant bishops or churches at Easter as a token of fellowship, are very old.

The Greek church used leavened bread, the Latin, unleavened. This difference ultimately led to intricate controversies.

The mixing of the wine with water was considered essential, and was explained in various mystical ways; chiefly by reference to the blood and water which flowed from the side of Jesus on the cross.

\xsection{The Liturgies. Their Origin and Contents}

§ 98. The Liturgies. Their Origin and Contents.

J. Goar. (a learned Dominican, † 1653): Εὐχολόγιον, sive Rituale Graecorum, etc. Gr. et Lat. Par. 1647 (another ed. at Venice, 1740). Jos. Aloys. Assemani (R.C.): Codex Liturgicus ecclesiae universae, ... in quo continentur libri rituales, missales, pontificales, officia, dypticha, etc., ecclesiarum Occidentis et Orientis (published under the auspices of Pope Boniface XIV.). Rom. 1749–’66, 13 vols. Euseb. Renaudot (R.C.): Liturgiarum orientalium collectio. Par. 1716 (reprinted 1847), 2 vols. L. A. Muratori (R.C., † 1750): Liturgia Romana vetus. Venet. 1748, 2 vols. (contains the three Roman sacramentaries of Leo, Gelasius, and Gregory I., also the Missale Gothicum, and a learned introductory dissertation, De rebus liturgicis). W. Palmer (Anglican): Origines Liturgicae. Lond. 1832 (and 1845), 2 vols. (with special reference to the Anglican liturgy). Ths. Brett: A Collection of the Principal Liturgies used in the Christian Church in the celebration of the Eucharist, particularly the ancient (translated into English), with a Dissertation upon them. Lond. 1838 (pp. 465). W. Trollope (Anglican): The Greek Liturgy of St. James. Edinb. 1848. H. A. Daniel (Lutheran, the most learned German liturgist): Codex Liturgicus ecclesiae universae in epitomem redactus. Lips. 1847 sqq. 4 vols. (vol. i. contains the Roman, vol. iv. the Oriental Liturgies). Fr. J. Mone (R.C.): Lateinische u. Griechische Messen aus dem 2ten his 6ten Jahrhundert. Frankf. a. M. 1850 (with valuable treatises on the Gallican, African, and Roman Mass). J. M. Neale († 1866, the most learned Anglican ritualist and liturgist, who studied the Eastern liturgies daily for thirty years, and almost knew them by heart); Tetralogia liturgica; sive S. Chrysostom, S. Jacobi, S. Marci divinae missae: quibus accedit ordo Mozarabicus. Lond. 1849. The Same: The Liturgies of S. Mark, S. James, S. Clement, S. Chrysostom, S. Basil, or according to the use of the churches of Alexandria, Jerusalem, Constantinople. Lond. 1859 f. (in the Greek original, and the same liturgies in an English translation, with an introduction and appendices, also at Lond. 1859). Comp. also Neale’s History of the Holy Eastern Church. Lond. 1850; Gen. Introd. vol. second; and his Essays on Liturgiology and Church History. Lond. 1863. (The latter, dedicated to the metropolitan Philaret of Moscow, is a collection of various learned treatises of the author from the “Christian Remembrancer” on the Roman and Gallican Breviary, the Church Collects, the Mozarabic and Ambrosian Liturgies, Liturgical Quotations, etc.) The already cited work, of kindred spirit, by the English Episcopal divine, Freeman, likewise treats much of the old Liturgies, with a predilection for the Western, while Neale has an especial reverence for the Eastern ritual. (Comp. also Bunsen: Christianity and Mankind, Lond. 1854, vol. vii., which contains Reliquiae Liturgicae; the Irvingite work: Readings upon the Liturgy and other Divine Offices of the Church. Lond. 1848–’54; Höfling: Liturgisches Urkundenbuch. Leipz. 1854.)

Liturgy\footnote{\textgreek{Λειτουργία,} from \textgreek{λεῖτος}, \emph{i.e.}, belonging to the \textgreek{λεώς} \emph{or} \textgreek{λαός}public, and \textgreek{ἔργον} = \textgreek{ἔργον τοῦ λεώ} \emph{or} \textgreek{τοῦ λαοῦ}, public work, office, function. In Athens the term was applied especially to the directing of public spectacles, festive dances, and the distribution of food to the people on festal occasions. Paul, in Rom. xiii. 6, calls secular magistrates \textgreek{λειτουργοὶ Θεοῦ.}} means, in ecclesiastical language,\footnote{Comp. Luke i. 23, where the priestly service of Zacharias is called \textgreek{λειτουργία}; Heb. viii. 2, 6; ix. 21; x. 11, where the word is applied to the High-Priesthood of Christ; Acts xiii. 2; Rom. xv. 16; Rom. xv. 27; 2 Cor. ix. 12, where religious fasting, missionary service, and common beneficences are called \textgreek{λειτουργία}or \textgreek{λειτουργεῖν}. . The restriction of the word to divine worship or sacerdotal action occurs as early as Eusebius, Vita Const. iv. 37, bishops being there called \textgreek{λειτουργοί.} The limitation of the word to the service of the Lord’s Supper is connected with the development of the doctrine of the eucharistic sacrifice.} the order and administration of public worship in general, and the celebration of the Eucharist in particular; then, the book or collection of the prayers used in this celebration. The Latin church calls the public eucharistic service Mass, and the liturgical books, sacramentarium, rituale, missale, also libri mysteriorum, or simply libelli.

The Jewish worship consisted more of acts than of words, but it included also fixed prayers and psalms (as Ps. 113–118) and the Amen of the congregation (Comp. 1 Cor. xiv. 16). The pagan Greeks and Romans had, in connection with their sacrifices, some fixed prayers and formulas of consecration, which, however, were not written) but perpetuated by oral tradition. The Indian literature, on the contrary, has liturgical books, and even the Koran contains prescribed forms of prayer.

The New Testament gives us neither a liturgy nor a ritual, but the main elements for both. The Lord’s Prayer, and the Words of the Institution of baptism and of the Holy Supper, are the living germs from which the best prayers and baptismal and eucharistic formulas of the church, whether oral or written, have grown. From the confession of Peter and the formula of baptism gradually arose in the Western church the Apostles’ Creed, which besides its doctrinal import, has also a liturgical office, as a public profession of candidates for baptism and of the faithful. In the Eastern church the Nicene creed is used instead. The Song of the angelic host is the ground-work of the Gloria in Excelsis. The Apocalypse is one sublime liturgic vision. With these belong also the Psalms, which have passed as a legitimate inheritance to the Christian church, and have afforded at all times the richest material for public edification.

In the ante-Nicene age we find as yet no traces of liturgical books. In each church, of course, a fixed order of worship gradually formed itself, which in apostolic congregations ran back to a more or less apostolic origin, but became enlarged and altered in time, and, until the fourth century, was perpetuated only by oral tradition. For the celebration of the sacraments, especially of the Eucharist, belonged to the Disciplina arcani, and was concealed, as the most holy thing of the church, from the gaze of Jews and heathens, and even of catechumens, for fear of profanation; through a misunderstanding of the warning of the Lord against casting pearls before swine, and after the example of the Samothracian and Eleusinian mysteries.\footnote{Comp. Tertullian, Apolog. c. 7; Origen, Homil. 9 in Levit. toward the end; Cyril of Jerusalem, Praefat. ad Catech. § 7, etc.} On the downfall of heathenism in the Roman empire the Disciplina arcani gradually disappeared, and the administration of the sacraments became a public act, open to all.

Hence also we now find, from the fourth and fifth centuries onward, a great number of written liturgies, and that not only in the orthodox catholic church, but also among the schismatics (as among the Nestorians, and the Monophysites). These liturgies bear in most cases apostolic names, but in their present form can no more be of apostolic origin than the so-called Apostolic Constitutions and Canons, nor nearly so much as the Apostles’ Creed. They contrast too strongly with the simplicity of the original Christian worship, so far as we can infer it from the New Testament and from the writings of the apologists and the ante-Nicene fathers. They contain also theological terms, such as ὁμοούσιος(concerning the Son of God), θεοτόκος(concerning the Virgin Mary), and some of them the whole Nicene Creed with the additions of the second oecumenical council of 381, also allusions to the worship of martyrs and saints, and to monasticism, which point unmistakably to the Nicene and post-Nicene age. Yet they are based on a common liturgical tradition, which in its essential elements reaches back to an earlier time, perhaps in some points to the apostolic age, or even comes down from the Jewish worship through the channel of the Jewish Christian congregations. Otherwise their affinity, which in many respects reminds one of the affinity of the Synoptical Gospels cannot be satisfactorily explained. These old catholic liturgies differ from one another in the wording, the number, the length, and the order of the prayers, and in other unessential points, but agree in the most important parts of the service of the Eucharist. They are too different to be derived from a common original, and yet too similar to have arisen each entirely by itself.\footnote{Trollope says, in the Introduction to his edition of the Liturgia Jacobi: “Nothing short of the reverence due to the authority of an apostle, could have preserved intact, through successive ages, that strict uniformity of rite and striking identity of sentiment, which pervade these venerable compositions; but there is, at the same time, a sufficient diversity both of expression and arrangement, to mark them as the productions of different authors, each writing without any immediate communication with the others, but all influenced by the same prevailing motives of action and the same constant habit of thought.” Neale goes further, and, in a special article on Liturgical Quotations (Essays on Liturgiology and Church History, Lond. 1863, p. 411 ff.), endeavors to prove that Paul several times quotes the primitive liturgy, viz., in those passages in which he introduces certain statements with a \textgreek{γέγραπται}, or \textgreek{λέγει}, or \textgreek{πιστὸς ὁ λόγος}, while the statements are not to be found in the Old Testament: 1 Cor. ii. 9; xv. 45; Eph. v. 14; 1 Tim. i. 15; iii. 1; iv. 1, 9; 2 Tim. ii. 11-13, 19; Tit. iii. 8. But the only plausible instance is 1 Cor. ii. 9: \textgreek{Καθὼς γέγραπται· ἅ ὀφθαλμὸς οὐκ εἷδε, καὶ οὖς οὐκ ἤκουσε, καὶ ἐπὶ καρδίαν ἀνθρώπου οὐκ ἀνέβη, ἅ ἡτοίμασεν ὁ Θεὸς τοῖς ἀγαπῶσιν αὐτόν}, which, it is true, occur word for word (though in the form of prayer, therefore with \textgreek{ἡτοίμασας}, and \textgreek{ἀγαπῶσί σε} instead of \textgreek{ἀγαπῶσιν αὐτόν}) the Anaphora of the Liturgia Jacobi, while the parallel commonly cited from Is. lxiv. 4 is hardly suitable. But if there had been such a primitive written apostolic liturgy, there would have undoubtedly been other and clearer traces of it. The passages adduced may as well have been quotations from primitive Christian hymns and psalms, though such are very nearly akin to liturgical prayers.}

All the old liturgies combine action and prayer, and presuppose, according to the Jewish custom, the participation of the people, who frequently respond to the prayers of the priest, and thereby testify their own priestly character. These responses are sometimes a simple Amen, sometimes Kyrie eleison, sometimes a sort of dialogue with the priest:

Priest: The Lord be with you!

People: And with thy spirit!

Priest: Lift—up your hearts!

People: We lift them up unto the Lord.

Priest: Let us give thanks!

People: It is meet and right.

Some parts of the liturgy, as the Creed, the Seraphic Hymn, the Lord’s Prayer, were said or sung by the priest and congregation together. Originally the whole congregation of the faithful\footnote{In the Clementine Liturgy, \emph{all}, \textgreek{πάντες}; in the Liturgy of St. James, \emph{the} \emph{People,} \textgreek{ὁ λαός}.} was intended to respond; but with the advance of the hierarchical principle the democratic and popular element fell away, and the deacons or the choir assumed the responses of the congregation, especially where the liturgical language was not intelligible to the people.\footnote{In the Liturgies of St. Basil and St. Chrysostom, which have displaced the older Greek liturgies, the \textgreek{διάκονος}or \textgreek{χορός}usually responds. In the Roman mass the people fall still further out of view, but accompany the priest with silent prayers.}

Several of the oldest liturgies, like those of St. Clement and St. James, have long since gone out of use, and have only a historical interest. Others, like those of St. Basil and St. Chrysostom, and the Roman, are still used, with various changes and additions made at various times, in the Greek and Latin churches. Many of their most valuable parts have passed, through the medium of the Latin mass-books, into the liturgies and agenda of the Anglican, the Lutheran, and some of the Reformed churches.

But in general they breathe an entirely different atmosphere from the Protestant liturgies, even the Anglican not excepted. For in them all the eucharistic sacrifice is the centre around which all the prayers and services revolve. This act of sacrifice for the quick and the dead is a complete service, the sermon being entirely unessential, and in fact usually dispensed with. In Protestantism, on the contrary, the Lord’s Supper is almost exclusively Communion, and the sermon is the chief matter in every ordinary service.

Between the Oriental and Occidental liturgies there are the following characteristic differences:

1. The Eastern retain the ante-Nicene division of public worship into two parts: the λειτουργία κατηχουμένων, Missa Catechumenorum, which is mainly didactic, and the λειτουργία τῶν πιστῶν, Missa Fidelium, which contains the celebration of the Eucharist proper. This division lost its primitive import upon the union of church and state, and the universal introduction of infant baptism. The Latin liturgies connect the two parts in one whole.

2. The Eastern liturgies contain, after the Words of Institution, an express Invocation of the Holy Ghost, without which the sanctification of the elements is not fully effected. Traces of this appear in the Gallican liturgies. But in the Roman liturgy this invocation is entirely wanting, and the sanctification of the elements is considered as effected by the priest’s rehearsal of the Words of Institution. This has remained a point of dispute between the Greek and the Roman churches. Gregory the Great asserts that the apostles used nothing in the consecration but the Words of Institution and the Lord’s Prayer.\footnote{Epist. ad Joann. Episc. Syriac.} But whence could he know this in the sixth century, since the New Testament gives us no information on the subject? An invocatio Spiritus Sancti upon the elements is nowhere mentioned; only a thanksgiving of the Lord, preceding the Words of Institution, and forming also, it may be, an act of consecration, though neither in the sense of the Greek nor of the Roman church. The Words of Institution: “This is my body,” \&c., are more-over addressed not to God, but to the disciples, and express, so to speak, the result of the Lord’s benediction.\footnote{On this disputed point Neale agrees with the Oriental church, Freeman with the Latin. Comp. Neale, Tetralogia Liturgica, Praefat. p. xv. sqq., and his English edition of the Primitive Liturgies of S. Mark, S. James, etc., p. 23. In the latter place he says of the \textgreek{ἐπίκλησις Πνεύματος ἁγίου}: “By the Invocation of the Holy Ghost, according to the doctrine of the Eastern church, and not by the words of institution, the bread and wine are ’changed,’ ’transmuted,’ ’transelemented,’ ’transubstantiated’ into our Lord’s Body and Blood. This has always been a point of contention between the two churches—the time at which the change takes place. Originally, there is no doubt that the Invocation of the Holy Ghostformed a part of all liturgies. The Petrine has entirely lost it: the Ephesine (Gallican and Mozarabic) more or less retains it: as do also those mixtures of the Ephesine and Petrine—the Ambrosian and Patriarchine or Aquileian. To use the words of the authorized Russian Catechism: ’Why is this (the Invocation) so essential? Because at the moment of this act, the bread and wine are changed or transubstantiated into the very Body of Christand into the very Blood of Christ. How are we to understand the word Transubstantiation? In the exposition of the faith by the Eastern Patriarchs, it is said that the word is not to be taken to define the manner in which the bread and wine are changed into the Body and Blood of our Lord; for this none can understand but God; but only this much is signified, that the bread, truly, really, and substantially becomes the very true Body of the Lord, and the wine the very Blood of the Lord.’ ” Freeman, on the contrary in his Principles of Div. Serv. vol. ii. Part ii. p. 196 f, asserts: “The Eastern church cannot maintain the position which, as represented by her doctors of the last four hundred years, and alleging the authority of St. Cyril, she has taken up, that there is no consecration till there has followed (1) a prayer of oblation and (2) one of Invocation of the Holy Ghost. In truth, the view refutes itself, for it disqualifies the oblation for the very purpose for which it is avowedly placed there, namely to make offering of the already consecrated Gifts, \emph{i.e.}, of the Body and Blood of Christ; thus reducing it to a level with the oblation at the beginning of the office. The only view that can be taken of these very ancient prayers, is that they are to be conceived of as offered simultaneously with the recitation of the Institution.”} 3. The Oriental liturgy allowed, more like the Protestant church, the use of the various vernaculars, Greek, Syriac, Armenian, Coptic, \&c.; while the Roman mass, in its desire for uniformity, sacrifices all vernacular tongues to the Latin, and so makes itself unintelligible to the people.

4. The Oriental liturgy is, so to speak, a symbolic drama of the history of redemption, repeated with little alteration every Sunday. The preceding vespers represent the creation, the fall, and the earnest expectation of Christ; the principal service on Sunday morning exhibits the life of Christ from his birth to his ascension; and the prayers and lessons are accompanied by corresponding symbolical acts of the priests and deacon: lighting and extinguishing candles, opening and closing doors, kissing the altar and the gospel, crossing the forehead, mouth, and breast, swinging the censer, frequent change of liturgical vestments, processions, genuflexions, and prostrations. The whole orthodox Greek and Russian worship has a strongly marked Oriental character, and exceeds the Roman in splendor and pomp of symbolical ceremonial.\footnote{On the mystical meaning of the Oriental cultus comp. the Commentary of Symeon of Thessalonica († 1429) on the Liturgy of St. Chrysostom, and Neale’s Introduction to his English edition of the Oriental Liturgies, pp. xxvii.-xxxvi.}

The Roman mass is also a dramatic commemoration and representation of the history of redemption, especially of the passion and atoning death of Christ, but has a more didactic character, and sets forth not so much the objective history, as the subjective application of redemption from the Confiteor to the Postcommunio. It affords less room for symbolical action, but more for word and song, and follows more closely the course of the church year with varying collects and prefaces for the high festivals,\footnote{The Collectsbelong strictly only to the Latin church, which has produced many hundred such short prayers. The word comes either from the fact that the prayer collects the sense of the Epistle and Gospel for the day in the form of prayer; or that the priest collects therein the wishes and petitions of the people. The collect is a short liturgical prayer, consisting of one petition, closing with the form of mediation through the merits of Christ, and sometimes with a doxology to the Trinity. Comp. a treatise of Neale on The Collects of the Church, in Essays on Liturgiology and Church History, p. 46 ff, and William Bright: Ancient Collects and Prayers, selected from various rituals, Oxford and London, 1860.} thus gaining variety. In this it stands the nearer to the Protestant worship, which, however, entirely casts off symbolical veils, and makes the sermon the centre.

Every Oriental liturgy has two main divisions. The first embraces the prayers and acts before the Anaphora or Oblation (canon Missae) to the Sursum corda; the second, the Anaphora to the close.

The first division again falls into the Mass of the Catechumens, and the Mass of the Faithful, to the Sursum corda. To it belong the Prefatory Prayer, the Introit, Ingressa, or Antiphon, the Little Entrance, the Trisagion, the Scripture Lessons, the Prayers after the Gospel, and the Expulsion of the Catechumens; then the Prayers of the Faithful, the Great Entrance, the Offertory, the Kiss of Peace, the Creed.

The Anaphora comprises the great Eucharistic Prayer of Thanksgiving, the Commemoration of the life of Jesus, the Words of Institution, the Oblation of the Elements, the Invocation of the Holy Ghost, the Great Intercession for Quick and Dead, the Lord’s Prayer, and finally the Communion with its proper prayers and acts, the Thanksgiving, and the Dismissal.\footnote{It is a curious fact, that in the Protestant Episcopal Trinity chapel of New York, with the full approval of the bishop, Horatio Potter, and the assistance of the choir, on the second of March, 1865, the anniversary of the accession of the Russian Czar, Alexander II., the full liturgy or mass of the orthodox Graeco-Russian church was celebrated before a numerous assembly by a recently arrived Graeco-Russian monk and priest (or deacon), Agapius Honcharenko. This is the first instance of an Oriental service in the United States (for the Russian fleet which was in the harbor of New York in 1863 held its worship exclusively upon the ships), and probably also the first instance of the celebration of the unbloody sacrifice of the mass and the mystery of transubstantiation in a Protestant church and with the sanction of Protestant clergy. The liturgy of St. Chrysostom, in the Slavonic translation, was intoned by the priest; the short responses, such as Hospode, Pomelue (Kyrie, Eleison), were grandly sung by the choir in the Slavonic language, and the Beatitudes, the Nicene Creed (of course, without the “Filioque,” which is condemned by the Greek church as a heretical innovation), and the Gloria in Excelsis in English There were wanting only the many genuflexions and prostrations, the trine immersion, and infant communion, to complete the illusion of a marriage of the two churches. Some secular journals gave the matter the significance of a political demonstration in favor of Russia! One of the religious papers saw in it an exhibition of the unity and catholicity of the church, and a resemblance to the miracle of Pentecost, in that Greeks, Slavonians, and Americans heard in their own tongues the wonderful works of God! But most of the Episcopal and other Protestant papers exposed the doctrinal inconsistency, since the Greek liturgy coincides in au important points with the Roman mass. Unfortunately for the philo-Russian movement, the Russo-Greek monk Agapius soon afterward publicly declared himself an opponent of the holy orthodox oriental church, an d charged it with serious error. The present Greek church, which regards even the archbishop of Canterbury and, the pope of Rome as unbaptized (because unimmersed) heretics and schismatics, could, of course, never consent to such an anomalous service as was held in Trinity chapel for the first, and in all probability for the last time.}

\xsection{The Oriental Liturgies}

§ 99. The Oriental Liturgies.

There are, in all, probably more than a hundred ancient liturgies, if we reckon revisals, modifications, and translations. But according to modern investigations they may all be reduced to five or six families, which may be named after the churches in which they originated and were used, Jerusalem (or Antioch), Alexandria, Constantinople, Ephesus, and Rome.\footnote{Neale now (The Liturgies of S. Mark, etc., 1859, p. vii) divides the primitive liturgies into five families: (1) That of St. James, or of Jerusalem; (2) that of St. Mark, or of Alexandria.; (3) that of St. Thaddaeus, or of the East; (4) that of St. Peter, or of Rome; (5) that of St. John, or of Ephesus. Formerly (Hist. of the Holy Eastern Church) he counted the Clementine Liturgy separately; but since Daniel has demonstrated the affinity of it with the Jerusalem (or, as he calls it, the Antiochian) family, he has put it down as a branch of that family.} Most of them belong to the Orientalchurch; for this church was in general much more productive, and favored greater variety, than the Western, which sought uniformity in organization and worship. And among the Oriental liturgies the Greek are the oldest and most important.

1. The liturgy of St. Clement. This is found in the eighth book of the Apostolic Constitutions, and, with them, is erroneously ascribed to the Roman bishop Clement.\footnote{It is given in Cotelier’s edition of the Patres Apostolici, in the various editions of the pseudo-Apostolic Constitutions, and in the liturgical collections of Daniel, Neale, and others.} It is the oldest complete order of divine service, and was probably composed in the East in the beginning of the fourth century.\footnote{Neale considers the liturgy the oldest part of the Apostolic Constitutions, places its composition in the second or third century, and ascribes its chief elements to the apostle Paul, with whose spirit and ideas it in many respects coincides.} It agrees most with the liturgy of St. James and of Cyril of Jerusalem, and may for this reason be considered a branch of the Jerusalem family. We know not in what churches, or whether at all, it was used. It was a sort of normal liturgy, and is chiefly valuable for showing the difference between the Nicene or ante-Nicene form of worship and the later additions and alterations.

The Clementine liturgy rigidly separates the service of the catechumens from that of the faithful.\footnote{Before the Sursum corda, or beginning of the Eucharist proper, the deacon says: “No catechumens, no hearers, no unbelievers, no heretics may remain here (\textgreek{μή τις τῶν κατηχουμένων, μή τις τῶν ἀκροωμένων, μὴ τις τῶν ἀπίστων, μή τις τῶν ἑτεροδόξων}). Depart, ye who have spoken the former prayer. Mothers, take your children,” etc. This arrangement is traced to James, the brother of John, the son of Zebedee.} It contains the simplest form for the distribution of the sacred elements: “The body of Christ,” and “The blood of Christ, the cup of life,” with the “Amen” of the congregation to each. In the commemoration of the departed it mentions no particular names of saints, not even the mother of God, who first found a place in public worship after the council of Ephesus in 431; and it omits several prefatory prayers of the priest. Finally it lacks the Nicene creed, and even the Lord’s Prayer, which is added to all other eucharistic prayers, and, according to the principles of some canonists, is absolutely necessary.\footnote{. The absence of the Lord’s Prayer in the Clementine Liturgy is sufficient to refute the view of Bunsen, that this prayer was originally the Prayer of Consecration in all liturgies.}

2. The liturgy of St. James. This is ascribed by tradition to James, the brother to the Lord, and bishop of Jerusalem.\footnote{Neale even supposes, as already observed, that St. Paul quotes from the Liturgia Jacobi, and not \emph{vice versa}, especially in I Cor. ii. 9} It, of course, cannot have been composed by him, even considering only the Nicene creed and the expressions ὁμοούσιοςand θεοτόκος, which occur in it, and which belong to the Nicene and post-Nicene theology. The following passage also bespeaks a much later origin: “Let us remember the most holy, immaculate, most glorious, blessed Mother of God and perpetual Virgin Mary, with all saints, that we through their prayers and intercessions may obtain mercy.” The first express mention of its use meets us in Proclus of Constantinople about the middle of the fifth century. But it is, as to substance, at all events one of the oldest liturgies, and must have been in use as early as the fourth century; for the liturgical quotations in Cyril of Jerusalem (in his fifth Mystagogic Catechesis), who died in 386, verbally agree with it. It was intended for the church of Jerusalem, which is mentioned in the beginning of the prayer for the church universal, as “the glorious Zion, the mother of all churches.”\footnote{\textgreek{Ὑπὲρ τῆς ἐνδόξου Σιὼν, τῆς μητρὸς πασῶν τῶν ἐκκλησιῶν· καὶ ὑπὲρ τῆς κατὰ πᾶσαν τὴν οἰκουμένην ἁγίας σου καθολικῆς καὶ ἀποστολικῆς ἐκκλησίας .}The intercessions for Jerusalem, and for the holy places which God glorified by the appearance of Christ and the outpouring of the Holy Ghost (\textgreek{ὑπὲρ τῶν ἁγίων σου τόπων, οὓς ἐδοξασας τῇ θεοφάνείᾳ τοῦ Χριστοῦ σου, κ.τ.λ.}), appears in no other liturgy.}

In contents and diction it is the most important of the ancient liturgies, and the fruitful mother of many, among which the liturgies of St. Basil and St. Chrysostom must be separately named.\footnote{Neale arranges the Jerusalem family in three divisions, as follows: \par “1. Sicilian S. James, as said in that island before the Saracen conquest, and partly assimilated to the Petrine Liturgy. \par 2. S. Cyril: where used uncertain, but assumilated to the Alexandrian form. \par 3. Syriac S. James, the source of the largest number of extant Liturgies. They are these: [1] \emph{Lesser S. James} [2] \emph{S. Clement}; [3] \emph{S. Mark}; [4] \emph{S. Dionysius}; [5] \emph{S. Xystus}; [6] \emph{S. Ignatius}; [7] \emph{S. Peter I}; [8] \emph{S. Peter II}; [9] \emph{S. Julius}; [10] \emph{S. John Evangelist}; [11] \emph{S}. \emph{Basil}; [12] (\emph{S}.) \emph{Dioscorus}; [13] \emph{S. John} Chrysostom\emph{I}; [14] \emph{All Apostles}; [15] \emph{S. Marutas}; [16] \emph{S. Eustathius}; [17] \emph{Philoxenus I}; [18] \emph{Matthew the Shepherd}; [19] \emph{James Bardaeus}; [20]. \emph{James of Botra}; [21] \emph{James of Edessa}; [22] \emph{Moses Bar-Cephas}; [23] \emph{Thomas of} \emph{Heraclea}; [24] \emph{Holy Doctors}; [25] \emph{Philoxenus II}; [26] \emph{S. John} Chrysostom\emph{II}; [27] \emph{Abu’lfaraj}; [28] \emph{John of Dara}; [291 \emph{S. Celestine}; [30] \emph{John Bar-Susan}; [31] \emph{Eleatar of} \emph{Babylon}; [32] \emph{John the Scribe}; [33] \emph{John Maro}; [34] \emph{Dionysius of Cardon}; [35] \emph{Michael of} \emph{Antioch}; [36] John Bar-Vahib; [37] \emph{John Bar-Maaden}; [38] \emph{Dionysius of Diarbekr}; [39] \emph{Philoxenus of Bagdad}. All these, from Syriac S. James inclusive, are Monophysite Liturgies} It spread over the whole patriarchate of Antioch, even to Cyprus, Sicily, and Calabria, but was supplanted in the orthodox East, after the Mohammedan conquest, by the Byzantine liturgy. Only once in a year, on the 23d of October, the festival of St. James, it is yet used at Jerusalem and on some islands of Greece.\footnote{There are only two manuscripts, with the fragment of a third, from which the ancient text of the Greek Liturgia Jacobi is derived. The first printed edition appeared at Rome in 1526; then one at Paris in 1560. Besides these we have the copies in the Bibliotheca Patrum, the Codex Apocryphus Novi Testamenti, the Codex Liturgicus of Assemani, the Codex Liturgicus of Daniel, and the later separate editions of Trollope (Edinburgh, 1848), and Neale (twice, in his Tetralogia Liturgica, 1849, and improved, in his Primitive Liturgies, 1860).}

The Syriac liturgy of James is a free translation from the Greek; it gives the Invocation of the Holy Spirit in a larger form, the other prayers in a shorter; and it betrays a later date. It is the source of thirty-nine Monophysite liturgies, which are in use still among the schismatic Syrians or Jacobites.\footnote{See the names of them in the preceding quotation from Neale.}

3. The liturgy of St. Mark, or the Alexandrian liturgy. This is ascribed to the well-known Evangelist, who was also, according to tradition, the founder of the church and catechetical school in the Egyptian capital. Such origin involves, of course, a shocking anachronism, since the liturgy contains the Nicaeno-Constantinopolitan creed of 381. In its present form it comes probably from Cyril, bishop of Alexandria († 444), who was claimed by the orthodox, as well as the Monophysites, as an advocate of their doctrine of the person of Christ.\footnote{Daniel (iv. 137 sqq.) likewise considers Cyril the probable author, and endeavors to separate the apostolical and the later elements. Neale, in the preface to his edition of the Greek text, thinks: “The general form and arrangement of the Liturgy of S. Mark may safely be attributed to the Evangelist himself, and to his immediate followers, S. Amianus, S. Abilius, and S. Cerdo. With the exception of certain manifestly interpolated passages, it had probably assumed its present appearance by the end of the second century.”} It agrees, at any rate, exactly with the liturgy which bears Cyril’s name.

It is distinguished from the other liturgies by the position of the great intercessory prayer for quick and dead before the Words of Institution and Invocation of the Holy Ghost, instead of after them. It was originally composed in Greek, and afterwards translated into Coptic and Arabic. It was used in Egypt till the twelfth century, and then supplanted by the Byzantine. The Copts still retained it. The Ethiopian canon is an offshoot from it. There are three Coptic and ten Ethiopian liturgies, which belong to the same family.\footnote{There is only one important manuscript of the Greek Liturgy of St. Mark, the Codex Rossanensis, printed in Renaudot’s Collectio, and more recently by Daniel and Neale.}

4. The liturgy of Edessa or Mesopotamia, or of All Apostles. This is traced to the apostles Thaddaeus (Adaeus) and Maris, and is confined to the Nestorians. From it afterwards proceeded the Nestorian liturgies: (1) of Theodore the Interpreter; (2) of Nestorius; (3) Narses the Leper; (4) of Barsumas; (5) of Malabar, or St. Thomas. The liturgy of the Thomas-Christians of Malabar has been much adulterated by the revisers of Diamper.\footnote{The printed edition is a revision by the Portuguese archbishop of Goa, Alexis of Menuze, and the council of Diamper (1599), who understood nothing of the Oriental liturgies. Neale says: “The Malabar Liturgy I have never been able to see in the original; and an \emph{unadulterated} copy of the original does not seem to exist.” He gives a translation of this liturgy in Primitive Liturgies, p. 128 ff.}

5. The liturgy of St. Basil and that of St. Chrysostom form together the Byzantine or ConstantinopolItan liturgy, and passed at the same time into the Graeco-Russian church. Both descend from the liturgy of St. James and give that ritual in an abridged form. They are living books, not dead like the liturgies of Clement and of James.

The liturgy of bishop Basil of Neo-Caesarea († 379) is read in the orthodox Greek, and Russian church, during Lent (except on Palm Sunday), on the eve of Epiphany, Easter and Christmas, and on the feast of St. Basil (1st of January). From it proceeded the Armenian liturgy.

The liturgy of St. Chrysostom († 407) is used on all other Sundays. It is an abridgment and improvement of that of St. Basil, and, through the influence of the distinguished patriarchs of Constantinople, it has since the sixth century dislodged the liturgies of St. James and St. Mark. The original text can hardly be ascertained, as the extant copies differ greatly from one another.

The present Greek and Russian ritual, which surpasses even the Roman in pomp, cannot possibly have come down in all its details from the age of Chrysostom. Chrysostom is indeed supposed, as Proclus says, to have shortened in many respects the worship in Constantinople on account of the weakness of human nature; but the liturgy which bears his name is still in the seventh century called “the Liturgy of the Holy Apostles,” and appears to have received his name not before the eighth.

\xsection{The Occidental Liturgies}

§ 100. The Occidental Liturgies.

The liturgies of the Western church may be divided into three classes: (1) the Ephesian family, which is traced to a Johannean origin, and embraces the Mozarabic and the Gallican liturgies; (2) the Roman liturgy, which, of course, like the papacy itself, must come down from St. Peter; (3) the Ambrosian and Aquileian, which is a mixture of the other two. We have therefore here less diversity than in the East. The tendency of the Latin church everywhere pressed strongly toward uniformity, and the Roman liturgy at last excluded all others.

1. The Old Gallican liturgy,\footnote{Edited by Mabillon: De liturgia Gallicana, libri iii. Par. 1729; and recently in much more complete form, from older MSS. by Francis Joseph Mone (archive-director in Carlsruhe): Lateinische u. griechische Messen aus dem 2ten his 6ten Jahrlhundert, Frankf. a. M. 1850. This is one of the most important liturgical discoveries. Mono gives fragments of eleven mass-formularies from a codex rescriptus of the former cloister of Reichenau, which are older than those previously known, but hardly reach back, as he thinks, to the century (the time of the persecution at Lyons, a.d.177). Comp. against this, Denzinger, in the Tübingen Quartalschrift, 1850, p. 500 ff. Neale agrees with Mone: Essays on Liturgiology, p. 137.} in many of its features, points back, like the beginnings of Christianity in South Gaul, to an Asiatic, Ephesian, and so far we may say Johannean origin, and took its later form in the fifth century. Among its composers, or rather the revisers, Hilary, of Poictiers is particularly named. In the time of Charlemagne it was superseded by the Roman. Gallicanism, which in church organization and polity boldly asserted its rights, suffered itself easily to be Romanized in its worship.

The Old British liturgy was without doubt identical with the Gallican, but after the conversion of the Anglo-Saxons it was likewise supplanted by the Roman.

2. The Old Spanish or (though incorrectly so called) Gothic, also named Mozarabic liturgy.\footnote{Called “Gothic,” because its development and bloom falls in the time of the Gothic rule in Spain; “Mozarabic” it came to be called after the conquest of Spain by the Arabs. Mozarab, Muzarab, Mostarab, is a kind of term of contempt for the Spanish Christians under the Arabic dominion, in distinction from the Arabs of pure blood. The word comes not from \emph{mixti} and \emph{Arabes}, nor from \emph{Muza}, the Maurian chieftain who subjugated Spain, but from a participle of the tenth conjugation of the Arabic verb \emph{araba}; therefore something like “arabizing Arab,” or Arab by adoption, in distinction from Arabs of the pure blood. Comp. the distinction between Hellenist and Hebrew.} This is in many respects allied to the Gallic, and probably came through the latter from a similar Eastern Source. It appears to have existed before the incursion of the West Goths in 409; for it shows no trace of the influence of the Arian heresy, or of the ritual system of Constantinople.\footnote{Pinius (in a dissertation prefixed to the 32d vol. of the Acta Sanctorum) supposes that the Spanish liturgy came from the Goths, therefore from Constantinople; but Neale (Essays on Liturgiology, p. 130 ff.) endeavors to prove that it was contemporaneous with the introduction of Christianity in Spain, but afterward, by Leander of Seville (about 589), was conformed in some points to the Oriental ceremonial.} Its present form is attributed to Isidore of Seville and the fourth council of Toledo in 633. It maintained itself in Spain down to the thirteenth century and was then superseded by the Roman liturgy.\footnote{The Spanish cardinal Ximenes edited from defective manuscripts the first printed edition at Toledo, 1500, which, however, is in a measure conformed to the Roman order. He also founded in the cathedral of Toledo a chapel (ad Corpus Christi), where the so renovated Mozarabic service is still continued daily. A similar chapel was founded in Salamanca for the same purpose. Neale, in his Tetralogia Liturgica, gives the Ordo Mozarabicus for comparison with the Liturgies of Chrysostom, James, and Mark. The latest edition is that in the 85th volume of Migne’s Patrologie, Paris, 1850, with a learned preface.}

It has, like the Gallican, besides the Gospels and Epistles, lessons also from the Old Testament;\footnote{On the Mozarabic pericopes comp. an article by Ernst Ranke in Herzog’s Encykop. vol. x. pp. 79-82. He attributes to them great intrinsic value and historical importance. “They even seem important,” says he,“for the general history of the ancient church. With the unmistakable affinity they bear to the Greek on the one hand, and to the Gallican on the other, they evince by themselves an intercourse between the Eastern and Western regions of the church, which, begun or at least aimed at by Paul, further established by Irenaeus, still under lively prosecution in the time of Jerome, afterward ruptured in the most violent manner, is without doubt one of the most noteworthy currents in the life of the church.”} it differs from the Roman liturgy in the order of festivals; and it contains, before the proper sacrificial action, a homiletic exhortation. The formula Sancta Sanctis, before the communion) the fraction of the host into nine parts (in memory of the nine mysteries of the life of Christ), the daily communion, the distribution of the cup by the deacon, remind us of the oriental ritual. The Mozarabic chant has much resemblance to the Gregorian, but exhibits besides a certain independent national character.\footnote{Neale has made the discovery, that the Mozarabic litanies were originally metrical, and attempts to restore the measure, l c. p. 143 ff.}

3. The African liturgy is known to us only through fragmentary quotations in Tertullian, Cyprian, and Augustine, from which we gather that it belonged to the Roman family.

4. The liturgy Of St. Ambrose.\footnote{Missale Ambrosianum, Mediol. 1768; a later edition under authority of the archbishop and cardinal Gaisruck, Mediol. 1850. Comp. an article by Neale: The Ambrosian Liturgy, in his Essays on Liturgiology, p. 171 ff. Neale considers the Ambrosian liturgy, like the Gallican and Mozarabic, a branch of the Ephesian family. “All three have been moulded by contact with the Petrine family; but the Ambrosian, as it might be expected, most of all.” He places it, however, far below the two others.} This is attributed to the renowned bishop of Milan († 397), and even to St. Barnabas. It is certain, that Ambrose introduced the responsive singing of psalms and hymns, and composed several prayers, prefaces, and hymns. His successor, Simplicius (a.d. 397–400), is supposed to have made several additions to the ritual. Many elements date from the reign of the Gothic kings (a.d. 493–568), and the Lombard kings (a.d. 568–739).

The Ambrosian liturgy is still used in the diocese of Milan; and after sundry vain attempts to substitute the Roman, it was confirmed by Alexander VI. in 1497 by a special bull, as the Ritus Ambrosianus. Excepting some Oriental peculiarities, it coincides substantially with the Roman liturgy, but has neither the pregnant brevity of the Roman, nor the richness and fullness of the Mozarabic. The prayers for the oblation of the sacrificial gifts differ from the Roman; the Apostles’ Creed is not recited till after the oblation; some saints of the diocese are received into the canonical lists of the saints; the distribution of the host takes place before the Paternoster, with formulas of its own, \&c.

The liturgy which was used for a long time in the patriarchate of Aquileia, is allied to the Ambrosian, and likewise stands midway between the Roman and the Oriental Gallican liturgies.

5. The Roman liturgy is ascribed by tradition, in its main features, to the Apostle Peter, but cannot be historically traced beyond the middle of the fifth century. It has without doubt slowly grown to its present form. The oldest written records of it appear in three sacramentaries, which bear the names of the three Popes, Leo, Gelasius, and Gregory.

(a) The Sacramentarium Leonianum, falsely ascribed to Pope Leo I. († 461), probably dates from the end of the fifth century, and is a planless collection of liturgical formularies. It was first edited in 1735 from a codex of Verona.\footnote{Hence called also Sacram. Veronense.}

(b) The Sacramentarium Gelasianum, which was first printed at Rome in 1680, passes for the work of the Roman bishop Gelasius († 492–496), who certainly did compose a Sacramentarium. Many saints’ days are wanting in it, which have been in use since the seventh century.

(c) The Sacramentarium Gregorianum, edited by Muratori and others. Gregory I. (590–604) is reputed to be the proper father of the Roman Ordo et Canon Missae, which, with various additions and modifications at later periods, gradually attained almost exclusive prevalence in the Latin church, and was sanctioned by the Council of Trent.

The collection of the various parts of the Roman liturgy\footnote{Sacramentarium, antiphonarium, lectionarium (containing the lessons from the Old Testament, the Acts, the Epistles, and the Apocalypse), evangelarium (the lessons from the Gospels), ordo Romanus.} in one book is called Missale Romanum, and the directions for the priests are called Rubricae.\footnote{From their being written or printed in red.}

\xsection{Liturgical Vestments}

§ 101. Liturgical Vestments.

Besides the liturgical works already cited, Comp. John England (late R.C. bishop of Charleston, S. C., d. 1842): An Historical Explanation of the Vestments, Ceremonies, etc., appertaining to the holy Sacrifice of the Mass (an Introduction to the American Engl. edition of the Roman Missal). Philad. 1843. Fr. Bock. (R.C.): Geschichte der liturgischen Gewänder des Mittelalters. Bonn, 1856, 2 vols. C. Jos. Hefele: Beiträge zur Kirchengeschichte, Archäologie und Liturgik. Vol ii. Tüb. 1864, p. 150 ff.

The stately outward solemnity of public worship, and the strict separation of the hierarchy from the body of the laity, required corresponding liturgical vesture, after the example of the Jewish priesthood and cultus,\footnote{To which in general the Greek and Roman system of vestments is very closely allied. On the Jewish sacred vestments, see Ex. xxviii. 1-53; xxxix. 1-31, etc.} symbolical of the grades of the clergy and of the different parts of the worship.

In the Greek church the liturgical vestments and ornaments are the sticharion,\footnote{\textgreek{Στοιχάριον, στιχάριον}(by Goar always translated, \emph{dalmatica}), a long coat corresponding to the broidered coat (\texthebrew{תנֶתׄכְּ}, χιτών,tunica, Ex. xxviii. 39) of the Jewish priest, and the alba and dalmatica of the Latin church.} and the orarion, or horarion\footnote{\textgreek{Ὡράριον}(from \textgreek{ὡρα}, hour of prayer), or \textgreek{ὠράριον}, corresponding to the Latin stola.} for the deacon; the sticharion, the phelonion,\footnote{\textgreek{Φελώνιον, φαιλώνιον}, a wide mantle, corresponding to the casula.} the zone,\footnote{\textgreek{Ζώνη}, girdle, cingulum, balteus, corresponding to the \texthebrew{ַאבְצֵט}} the epitrachelion,\footnote{\textgreek{Ἐπιτραχήλιον}, collarium, a double orarion, a scapulary or cape.} and the epimanikia\footnote{\textgreek{Ἐπιμανίκια,}on the arms, corresponding to the manipulus.} for the priest; the saccos,\footnote{\textgreek{Σάκκος,}a short coat with rich embroidery, without sleeves, and with little bells.} the omophorion\footnote{\textgreek{Ὠμοφόριον} corresponding to the Latin pallium (and so translated by Goar) but broader, and fastened about the neck with a button.} the epigonation,\footnote{\textgreek{Ἐπιγονάτιον,}also \textgreek{ὑπογονάτιον}, a quadrangular shield, reaching from the \textgreek{ζώνη}to the knee, and signifying, according to Simeon Metaphrastes, the victory over death and the devil.} and the crozier\footnote{\textgreek{Ράβδος,} sceptrum.} for the bishop. The mitre is not used by the Greeks.

The vestments in the Latin church are the amict or humeral\footnote{The linen cloth which the priest, before celebrating, threw about his neck and shoulders, with the prayer: “Impone, Domine, capiti meo galeam salutis ad expugnandos diabolicos excursus.” It is nowhere mentioned before the eighth century. It answers to the Jewish ephod.} the alb (white cope or surplice),\footnote{Alba vestis, tunica, camisia, the white linen robe which hangs from the neck to the feet. From the alb arose, by shortening, the surplice (superpelliceum, rochetturn; French: surplis; German: Chorrock), which is the ordinary official dress of the lower clergy.} the cincture,\footnote{Cingulum, balteus, zona, a linen girdle for gathering up the alb.} the maniple,\footnote{Manipulus, sudarium, fano, mappula, originally a napkin, hung upon the left arm of deacons and priests, afterward only of bishops, after the \emph{Confiteor}.} the orarium or stole\footnote{The stola is a linen vestment hanging from both shoulders. The pope wears the stole always; the priest, only when officiating. The council of Laodicea after 347 prohibited the wearing of it by subdeacons and the lower clergy.} for the priest; the chasuble,\footnote{Casula, planeta, the mass-vestment, covering the whole body, but without sleeves, with a cross behind and before embroidered in gold or fine silk. From the casula arose the pluviale, a festive mantle with a hood (casula calcullata), used in processions and on other state occasions.} the dalmatic,\footnote{So called from the place of its origin. It is an overgarment of costly material, similar to the casula, and worn under it.} the pectoral\footnote{The pectorale, crux pectoralis, is the breast-cross of bishops and archbishops, and answers probably to the breastplate of the Jewish high-priest.} and the mitre\footnote{The mitra, tiara, infula, birretum, is the episcopal head dress, after the type of the Jewish \texthebrew{תפֶנֶצְמִ }(LXX.: \textgreek{κίδαρις}, Vulgate: tiara, mitra), originally single, after the eleventh century with two points, supposed to denote the two Testaments.} for the bishop; the pallium for the archbishop. To these are to be added the episcopal ring and staff or crozier.

These clerical vestments almost all appear to have been more or less in use before the seventh century, though only in public worship; it is impossible exactly to determine the age of each. The use of priestly vestments itself originated in fact in the Old Testament, and undoubtedly passed into the church through the medium of the Jewish Christianity, but of course with many modifications. Constantine the Great presented the bishop Macarius of Jerusalem a splendid stole wrought with gold for use at baptism.

The Catholic ritualists of course give to the various mass-vestments a symbolical interpretation, which is in part derived from the undeniable meaning of the Jewish priestly garments,\footnote{On the Jewish sacerdotal vesture and its symbolical meaning, Comp. Braun: Vestitus sacerdotum Hebraeorum, Amstel. 1698; Lundius: Die jüdischen Heiligthümer, pp. 418-445; Baehr: Symbolik des mosaischen Cultus, vol. ii. pp. 61-165.} but in part is arbitrary, and hence variable. The amict, for example, denotes the collecting of the mind from distraction; the alb, the righteousness and holiness of the priests; the maniple, the fruits of good works; the stole, the official power of the priest; the mitre, the clerical chieftainship; the ring, the marriage of the bishop with the church; the staff his oversight of the flock.

The color of the liturgical garments was for several centuries white; as in the Jewish sacerdotal vesture the white color, the symbol of light and salvation, prevailed. But gradually five ecclesiastical colors established themselves. The material varied, except that for the amict and the alb linen (as in the Old Testament) was prescribed. According to the present Roman custom the sacred vestments, like other sacred utensils and the holy water, must be blessed by the bishop or a clergyman even appointed for the purpose. The Greeks bless them even before each use of them. The Roman Missal, and other liturgical books, give particular directions in the rubrics for the use of the mass vestments.

In everyday life, for the first five or six centuries the clergy universally wore the ordinary citizens’ dress; then gradually, after the precedent of the Jewish priests\footnote{The prevailing color of the ordinary Jewish priestly costume was white; that of the Christian clerical costume, on the contrary, is black.} and Christian monks, exchanged it for a suitable official costume, to make manifest their elevation above the laity. So late as the year 428, the Roman bishop Celestine censured some Gallic priests for having, through misinterpretation of Luke xii. 35, exchanged the universally used under-garment (tunica) and over-garment (toga) for the Oriental monastic dress, and rightly reminded them that the clergy should distinguish themselves from other people not so much by outward costume, as by purity of doctrine and of life.\footnote{“Discernendi a caeteris sumus doctrina, non veste, conversatione, non habitu, mentis puritate, non cultu.” Comp. Thomassin, Vetus ac nova ecclesiae disciplina, P. i. lib. ii. cap. 43.} Later popes and councils, however, enacted various laws and penalties respecting these externals, and the council of Trent prescribed an official dress befitting the dignity of the priesthood.\footnote{Sess. xiv. cap. 6 de reform.: “Oportet clericos vestes proprio congruentes ordini semper deferre, ut per decentiam habitus extrinseci morum honestatem intrinsecam ostendant.”}

\xchapter{Christian Art}

CHRISTIAN ART.

\xsection{Religion and Art}

§ 102. Religion and Art.

Man is a being intellectual, or thinking and knowing, moral, or willing and acting, and aesthetic, or feeling and enjoying. To these three cardinal faculties corresponds the old trilogy of the true, the good, and the beautiful, and the three provinces of science, or knowledge of the truth, virtue, or practice of the good, and art, or the representation of the beautiful, the harmony of the ideal and the real. These three elements are of equally divine origin and destiny.

Religion is not so much a separate province besides these three, as the elevation and sanctification of all to the glory of God. It represents the idea of holiness, or of union with God, who is the original of all that is true, good, and beautiful. Christianity, as perfect religion, is also perfect humanity. It hates only sin; and this belongs not originally to human nature, but has invaded it from without. It is a leaven which pervades the whole lump. It aims at a harmonious unfolding of all the gifts and powers of the soul. It would redeem and regenerate the whole man, and bring him into blessed fellowship with God. It enlightens the understanding, sanctifies the will, gives peace to the heart, and consecrates even the body a temple of the Holy Ghost. The ancient word: “Homo sum, nihil humani a me alienum puto,” is fully true only of the Christian. “All things are yours,” says the Apostle. All things are of God, and for God. Of these truths we must never lose sight, notwithstanding the manifold abuses or imperfect and premature applications of them.

Hence there is a Christian art, as well as a Christian science, a spiritual eloquence, a Christian virtue. Feeling and imagination are as much in need of redemption, and capable of sanctification, as reason and will.

The proper and highest mission of art lies in the worship of God. We are to worship God “in the beauty of holiness.” All science culminates in theology and theosophy, all art becomes perfect in cultus. Holy Scripture gives it this position, and brings it into the closest connection with religion, from the first chapter of Genesis to the last chapter of the Revelation, from the paradise of innocence to the new glorified earth. This is especially true of the two most spiritual and noble arts, of poetry and music, which proclaim the praise of God—in all the great epochs of the history of his kingdom from the beginning to the consummation. A considerable part of the Bible: the Psalms, the book of Job, the song of Solomon, the parables, the Revelation, and many portions of the historical, prophetical, and didactic books, are poetical, and that in the purest and highest sense of the word. Christianity was introduced into the world with the song of the heavenly host, and the consummation of the church will be also the consummation of poetry and song in the service of the heavenly sanctuary.

Art has always, and in all civilized nations, stood in intimate connection with worship. Among the heathen it ministered to idolatry. Hence the aversion or suspicion of the early Christians towards it. But the same is true of the philosophy of the Greeks, and the law of the Romans; yet philosophy and law are not in themselves objectionable. All depends on the spirit which animates these gifts, and the purpose which they are made to serve.

The great revolution in the outward condition of the church under Constantine dissipated the prejudices against art and the hindrances to its employment in the service of the church. There now arose a Christian art which has beautified and enriched the worship of God, and created immortal monuments of architecture, painting, poetry, and melody, for the edification of all ages; although, as the cultus of the early church in general perpetuated many elements of Judaism and heathenism, so the history of Christian art exhibits many impurities and superstitions which provoke and justify protest. Artists have corrupted art, as theologians theology, and priests the church. But the remedy for these imperfections is not the abolition of art and the banishment of it from the church, but the renovation and ever purer shaping of it by the spirit and in the service of Christianity, which is the religion of truth, of beauty, and of holiness.

From this time, therefore, church history also must bring the various arts, in their relation to Christian worship, into the field of its review. Henceforth there is a history of Christian architecture, sculpture, painting, and above all of Christian poetry and music.

\xsection{Church Architecture}

§ 103. Church Architecture.

On the history of Architecture in general, comp. the works of Kugler, Kinkel, Schnaase, and others, on the plastic arts; also Kreuser: Der christliche Kirchenbau, seine Geschichte, Symbolik u. Bildnerei, Bonn, 1851. 2 vols., and the English works of Knight, Brown, Close, J. Ferguson (A Hist. of Architecture, Lond. 1865, 3 vols.), etc.

Architecture is required to provide the suitable outward theatre for the public worship of God, to build houses of God among men, where he may hold fellowship with his people, and bless them with heavenly gifts. This is the highest office and glory of the art of building. Architecture is a handmaid of devotion. A beautiful church is a sermon in stone, and its spire a finger pointing to heaven. Under the old covenant there was no more important or splendid building than the temple at Jerusalem, which was erected by divine command and after the pattern of the tabernacle of the wilderness. And yet this was only a significant emblem and shadow of what was to come.

Christianity is, indeed, not bound to place, and may everywhere worship the omnipresent God. The apostles and martyrs held the most solemn worship in modest private dwellings, and even in deserts and subterranean catacombs, and during the whole period of persecution there were few church buildings properly so called. The cause of this want, however, lay not in conscientious objection, but in the oppressed condition of the Christians. No sooner did they enjoy external and internal peace, than they built special places of devotion, which in a normal, orderly condition of the church are as necessary to public worship as special sacred times. The first certain traces of proper church buildings, in distinction from private places, appear in the second half of the third century, during the three-and-forty years’ rest between the persecution of Decius and that of Diocletian.\footnote{Euseb. Hist. Eccl. viii. 1.} But these were destroyed in the latter persecution.

The period of church building properly begins with Constantine the Great. After Christianity was acknowledged by the state, and empowered to hold property, it raised houses of worship in all parts of the Roman empire. There was probably more building of this kind in the fourth century than there has been in an period since, excepting perhaps the nineteenth century in the United States, where, every ten years, hundreds of churches and chapels are erected, while in the great cities of Europe the multiplication of churches by no means keeps pace with the increase of population.\footnote{The cities of New York, Brooklyn, and Philadelphia, for instance, have more churches than the much older cities of Berlin, Vienna, and Paris. New York has some three hundred, Berlin and Paris each hardly fifty. This is a noble triumph of the voluntary principle in religion.} Constantine and his mother Helena led the way with a good example. The emperor adorned not only his new residential city, but also the holy Places in Palestine, and the African city Constantine, with basilicas, partly at his own expense, partly from the public treasury. His successors on the throne, excepting Julian, as well as bishops and wealthy laymen, vied with each other in building, beautifying, and enriching churches. This was considered a work pleasing to God and meritorious. Ambition and self-righteousness mingled themselves here, as they almost everywhere do, with zeal for the glory of God. Chrysostom even laments that many a time the poor are forgotten in the church buildings, and suggests that it is not enough to adorn the altar, the walls, and the floor, but that we must, above all, offer the soul a living sacrifice to the Lord.\footnote{Homil. lxxxi in Matth. § 2, and l. § 3.} Jerome also rebukes those who haughtily pride themselves in the costly gifts which they offer to God, and directs them to help needy fellow-Christians rather, since not the house of stone, but the soul of the believer is the true temple of Christ.

The fourth century saw in the city of Rome above forty great churches.\footnote{Optatus of Mileve, De schism. Donat. ii. 4: “Inter quadraginta et quod excurris basilicas.”} In Constantinople the Church of the Apostles and the church of St. Sophia, built by Constantine, excelled in magnificence and beauty, and in the fifth century were considerably enlarged and beautified by Justinian. Sometimes heathen temples or other public buildings were transformed for Christian worship. The Emperor Phocas (602–610), for example, gave to the Roman bishop Boniface IV, the Pantheon, built by Agrippa under Augustus, and renowned for its immense and magnificent dome (now called chiesa della rotonda), and it was thenceforth consecrated to the virgin Mary and the martyrs.

But generally the heathen temples, from their small size and their frequent round form, were not adapted for the Christian worship, as this is held within the building, and requires large room for the congregation, that the preaching and the Scripture-reading may be heard; while the heathen sacrifices were performed before the portico, and the multitude looked on without the sanctuary. The sanctuary of Pandrosos, on the Acropolis at Athens, holds but few persons, and even the Parthenon is not so capacious as an ordinary church. The Pantheon in Rome is an exception, and is much larger than most temples. The small round pagan temples were most easily convertible into Christian baptisteries and burial chapels. Far more frequently, doubtless, was the material of forsaken or destroyed temples applied to the building of churches.

\xsection{The Consecration of Churches}

§ 104. The Consecration of Churches.

New churches were consecrated with great solemnity by prayer, singing, the communion, eulogies of present bishops, and the depositing of relics of saints.\footnote{I This last was, according to Ambrose, Epist. 54, the custom in Rome, and certainly wherever such relics were to be had.} This service set them apart from all profane uses, and designated them exclusively for the service and praise of God and the edification of his people. The dedication of Solomon’s temple,\footnote{2 Chron. c. 5-7.} as well as the purification of the temple after its desecration by the heathen Syrians,\footnote{Macc. iv. 44 ff.} furnished the biblical authority for this custom. In times of persecution the consecration must have been performed in silence. But now these occasions became festivals attended by multitudes. Many bishops, like Theodoret, even invited the pagans to attend them. The first description of such a festivity is given us by Eusebius: the consecration of the church of the Redeemer at the Holy Sepulchre,\footnote{Vita Constant. iv. 43-46.} and of a church at Tyre.\footnote{Hist. Eccl. x. 2-4. Eusebius speaks here in general of the consecration of churches after the cessation of persecution, and then, c. 4, gives an oratio panegyriea, delivered probably by himself, in which he describes the church at Tyre in a minute, but pompous way.}

After the Jewish precedent,\footnote{\textgreek{Τὰ ἐγκαίνια}, in memory of the purification of the temple under the Maccabees, I Macc. iv. 59; John x. 22.} it was usual to celebrate the anniversary of the consecration.\footnote{Sozomen, H. E. ii. 25 (26). Gregory the Great ordered: “Solemnitates ecclesiarum dedicationum per singulos annos sunt celebrandae.”}

Churches were dedicated either to the holy Trinity, or to one of the three divine Persons, especially Christ, or to the Virgin Mary, or to apostles, especially Peter, Paul, and John, or to distinguished martyrs and saints.

The idea of dedication, of course, by no means necessarily involves the superstitious notion of the omnipresent God being inclosed in a definite place. On the contrary, Solomon had long before said at the dedication of the temple at Jerusalem: “Behold, the heaven and heaven of heavens cannot contain thee; how much less this house that I have builded.” When Athanasius was once censured for assembling the congregation on Easter, for want of room, in a newly built but not yet consecrated church, he appealed to the injunction of the Lord, that we enter into our closet to pray, as consecrating every place. Chrysostom urged that every house should be a church, and every head of a family a spiritual shepherd, remembering the account which he must give even for his children and servants.\footnote{Hom. vi. in Gen., § 2 \textgreek{Ἐκκλησίαν ποίησόν σου τὴν οἰκίαν καὶ γὰρ καὶ ἐπεύθυνος εἶ καὶ τῆς τῶν παιδίων καὶ τῆς οἰκετῶν σωτηρίας .}} Not walls and roof, but faith and life, constitute the church,\footnote{Serm. in Eutrop.:\textgreek{Ἡ ἐκκλησία οὐ τεῦχος καὶ ὄροφος , ἀλλὰ πίστις καὶ βίος .}} and the advantage of prayer in the church comes not so much from a special holiness of the place, as from the Christian fellowship, the bond of love, and the prayer of the priests.\footnote{De incomprehensibili: \textgreek{Ἐνταῦθα ἐστί τι πλέον, οἷον ἡ ὁμόνοια καὶ ἡ συμφωνία καὶ τῆς ἀγάπης ὁ σύνδεσμος καὶ αἱ τῶν ἱερέων εὐχαί.}} Augustine gives to his congregation the excellent admonition: “It is your duty to put your talent to usury; every one must be bishop in his own house; he must see that his wife, his son, his daughter, his servant, since he is bought with so great a price, continues in the true faith. The apostle’s doctrine has placed the master over the servant, and has bound the servant to obedience to the master, but Christ has paid a ransom for both.”\footnote{Serm. 94.}

\xsection{Interior Arrangement of Churches}

§ 105. Interior Arrangement of Churches.

The interior arrangement of the Christian churches in part imitated the temple at Jerusalem, in part proceeded directly, from the Christian spirit. It exhibits, therefore, like the whole catholic system, a mixture of Judaism and Christianity. At the bottom of it lay the ideas of the priesthood and of sacrifice, and of fellowship with God administered thereby.

Accordingly, in every large church after Constantine there were three main divisions, which answered, on the one hand, to the divisions of Solomon’s temple, on the other, to the three classes of attendants, the catechumens, the faithful, and the priests, or the three stages of approach to God. The evangelical idea of immediate access of the whole believing congregation to the throne of grace, does not yet appear. The priesthood everywhere comes between.

1. The portico: In this again must be distinguished:

(a) The inner portico, a covered hall which belonged to the church itself, and was called πρόναος, or commonly, from its long, narrow shape,νάρθηξ, ferula, i.e., literally, staff, rod.\footnote{Sometimes the narthex again was divided into two rooms, the \emph{upper} place for the kneelers (locus substratorum), \emph{i.e}., catechumens who might participate, kneeling, in the prayers after the sermon (hence genuflectentes, \textgreek{γονυκλίνοντες}and the \emph{lower} place, bordering on the outer portico, for mere hearers, Jews, and pagans (locus audientium).} The name paradise also occurs, because on one side of the wall of the portico Adam and Eve in paradise were frequently painted,—probably to signify that the fallen posterity of Adam find again their lost paradise in the church of Christ. The inner court was the place for all the unbaptized, for catechumens, pagans, and Jews, and for members of the church condemned to light penance, who might hear the preaching and the reading of the Scriptures, but must withdraw before the administration of the Holy Supper.

(b) The outer portico, αὐλή, atrium, also locus lugentium or hiemantium, which was open, and not in any way enclosed within the sacred walls, hence not a part of the house of God properly so called. Here those under heavy penance, the “weepers”\footnote{Flentes, hiemantes.} as they were called, must tarry, exposed to all weather, and apply with tears to those entering for their Christian intercessions.

In this outer portico, or atrium, stood the laver,\footnote{\textgreek{Κρήνη}, cantharus, phiala.} in which, after the primitive Jewish and heathen custom, maintained to this day in the Roman church, the worshipper, in token of inward purification, must wash every time he entered the church.\footnote{In Num. xix. 2 ff.; xxxi. 19 ff. (comp. Heb. ix. 13) the sprinkling-water, or “water of separation” (i.e., water of purification, LXX.: \textgreek{ὕδωρ ῤαντισμοῦ}), already appears, prepared from the ashes of the burned red heifer and water, and used for the cleansing of those made unclean by contact with a corpse. The later Jews were very strict in this; no one could appear in the temple or synagogue, or perform any act of worship, prayer, or sacrifice, without being washed, 1 Sam. xvi. 6; 2 Chron. xxx. 17. Therefore synagogues were built by preference in the neighborhood of streams. The Pharisees were very paltry and pedantic in the matter of these washings; comp. Matt. xv. 2; Mark vii. 3; Luke xi. 38. The same custom of symbolical purification before worship we find among the ancient Egyptians, Persians, Brahmans (who ascribed to the water of the Ganges saving virtue), Greeks, and Romans, and among the Mohammedans. At the entrance of every Turkish mosque stands a large font for this purpose.}

After about the ninth century, when churches were no longer built with spacious porticoes, this laver was transferred to the church itself, and fixed at the doors in the form of a holywater basin, supposed to be an imitation of the brazen sea in the priest’s court of Solomon’s temple.\footnote{1 Kings vii. 23-26; 2 Chron. iv. 2-5.} This symbolical usage could easily gather upon itself superstitious notions of the magical virtue of the holy water. Even in the pseudo-Apostolic Constitutions the consecrated water is called “a means of warding off diseases, frightening away evil spirits, a medicine for body and soul, and for purification from sins;” and though these expressions related primarily to the sacramental water of baptism as the bath of regeneration, yet they were easily applied by the people to consecrated water in general. In the Roman Catholic church the consecration of the water\footnote{Benedictio fontis.} is performed on Easter Sunday evening; in the Greco-Russian church, three times in the year.

2. The temple proper,\footnote{\textgreek{Ναός}.} the holy place,\footnote{\textgreek{Ἱερόν..}} or the nave of the church,\footnote{\textgreek{Ναῦς}, navis ecclesiae. Many derive this expression from a confusion of the Greek \textgreek{ναός}with \textgreek{ναῦς} and navis. Not till the ninth and tenth centuries is \emph{navis} used in this way. The more exact equivalent in English would be \emph{long-room}, or hall.} as it were the ark of the new covenant. This part extended from the doors of entrance to the steps of the altar, had sometimes two or four side-naves, according to the size of the church, and was designed for the body of the laity, the faithful and baptized. The men sat on the right towards the south (in the men’s nave), the women on the left towards the north (in the women’s nave), or, in Eastern countries, where the sexes were more strictly separated, in the galleries above.\footnote{Called \textgreek{ἐπερῶα}, the elevated galleries on the side walls. Besides this the women’s places were protected by wooden lattices from all curious or lascivious glances of the men. Chrysostomsays, Homil. 74 in Matth.: “Formerly these lattices certainly did not exist; for in Christ there is neither male nor female (Gal. iii. 28), and in the time of the apostles men and women were together with one accord. But then men were still men, and women were women; now women have sunk to the level of prostitutes, and men are like horses in rutting.” A sad commentary on the moral and religious condition of that time!} The monks and nuns, and the higher civil officers, especially the emperors with their families, usually had special seats of honor in semicircular niches on both sides of the altar.

About the middle of the main nave was the pulpit or the ambo,\footnote{\textgreek{Ἄμβων} from \textgreek{ἀναβαίνω}, pulpitum, suggestus. Hence the English pulpit, while the corresponding German \emph{Kanzel}is derived from \emph{cancelli}.} or subsequently two desks, at the left the Gospel-desk, at the right the Epistle-desk, where the lector or deacon read the Scripture lessons. The sermon was not always delivered from the pulpit, but more frequently either from the steps of the altar (hence the phrase: “speaking from the rails”), or from the seat of the bishop behind the altar-table.\footnote{\textgreek{Βῆμα}, exedra.}

Between the reading-desks and the altar was the odeum,\footnote{\textgreek{Ὠδεῖον}.Subsequently the singers were usually placed in the galleries of upper-church.} the place for the singers, and at the right and left the seats for the lower clergy (anagnosts or readers, exorcists, acolytes). This part of the nave lay somewhat higher than the floor of the church, though not so high as the altar-choir, and hence was also called the lower choir, and the gradual, because steps (gradus) led up to it. In the Eastern church the choir and nave are scarcely separated, and they form together the ναός, or temple hall; in the Western the choir and the sanctuary are put together under the name cancelli or chancel.

3. The most holy place,\footnote{\textgreek{Τὰ ἅγια τῶν ἁγίων, τὰ ἄδυτα, ἱερατεῖον,}, sacrarium, sanctuarium.} or the choir proper;\footnote{\textgreek{Χορός, βῆμα,} (ascensus).} called also in distinction from the lower choir, the high choir,\footnote{Hence the terms \emph{high mass, high altar}.} for the priests, and for the offering of the sacrifice of the Eucharist. No layman, excepting the emperor (in the east), might enter it. It was semi-circular or conchoidal\footnote{Hence called also \textgreek{κόγχη}, \emph{shell}.} in form, and was situated at the eastern end of the church, opposite the entrance doors, because the light, to which Christians should turn themselves, comes from the east.\footnote{Thus so early as this was the line of east and west established as the \emph{sacred} (or church-building) \emph{line}. Yet there were exceptions. Socrates, H. E. v. 22, notes it as peculiar in the church of Antioch, that the altar here stood not in the eastern end, but in the western (\textgreek{οὐ γὰρ πρὸς ἀνατολὰς τὸ θυσιαστήριον, ἀλλὰ πρὸς δύσιν ὁρᾷ}).} It was separated from the other part of the church by rails or a lattice,\footnote{\textgreek{Ἀμφίθυρα, κιγκλίδες},\emph{cancelli}, whence the name \emph{chancel}.} and by a curtain, or by sacred doors called in the Greek church the picture-wall, iconostas, on account of the sacred paintings on it.\footnote{Eusebius mentions, in his description of the church of the bishop Paulinus in Tyre, H. E. x. 4, an elegantly wrought lattice, and Athanasius mentions the curtains. Indeed, the pictures placed upon these curtains date back even to the fourth century, since Epiphanius, Ep. ad Joann. Hierosolymit., inveighed against a painted curtain in a village of Palestine. The lattice has perpetuated itself to this day in the picture wall or iconostas (\textgreek{εἰκονόστασις}) in the Russo-Greek church. It bears, on the right the picture of Christ, and on the left that of the Virgin Mary, and is pierced with three doors; the middle one, called the Emperor’s gate (dweri Zarskija), because only the emperor, besides the chief priest, may pass through it to take the holy Supper, is decorated and distinguished with the utmost splendor; oftentimes a golden sun with a thousand rays appears, which suddenly separates during the worship, and discloses the altar; or a Mount Zion with innumerable temples and battlements; or a network of golden garlands of flowers and fruits, among which especially clusters of grapes, probably with reference to the sacramental wine, frequently occur.} While in the Eastern churches this screen is still used, it in time gave place in the West to a low balustrade.

In the middle of the sanctuary stood the altar,\footnote{Altare, mensa sacra, \textgreek{θυσιαστήριον, ἁγία τράπεζα}The \emph{altar-cloth}, \emph{palla, pallia}, covers the whole upper face of the altar. This must not be confounded with the \emph{corporale} (\textgreek{εἴλητον,}from \textgreek{εἰλέω}, involvo), \emph{i}.\emph{e.,} a white linen cloth, with which the oblations prepared upon the altar are covered.} generally a table, or sometimes a chest with a lid; at first of wood, then, after the beginning of the sixth century, of stone or marble, or even of silver and gold, with a wall behind it, and an overshadowing, dome-shaped canopy,\footnote{\textgreek{Πυργος}, tower; \textgreek{κιβώριον}(of doubtful origin), ciborium, umbraculum. Subsequently the ciborium gave place to the steeple-shaped \emph{tabernaculum} for the preservation of the body of Christ. With the ciborium the dove-shaped form of the receptacle for the body of Christ (hence called \textgreek{περιστήριον}) also gradually disappeared.} above which a cross was usually fixed. The altar was hollow, and served as the receptacle for the relics of the martyrs; it was placed, where this was possible, exactly over the grave of a martyr, probably with reference to the passage in the Revelation: “I saw under the altar the souls of them that were slain for the word of God, and for the testimony which they held.”\footnote{Rev. vi. 9. In the Greek and Roman churches every altar must contain some relics, be they never so unimportant.} Often a subterranean chapel or crypt\footnote{\textgreek{Κρυπταί}, memoriae, confessiones, testimonia.} was built under the church, in order to have the church exactly upon the burial place of the saint, and at the same time to keep alive the memory of the primitive worship in underground vaults in the times of persecution.

The altar held therefore the twofold office of a tomb (though at the same time the monument of a new, higher life), and a place of sacrifice. It was manifestly the most holy place in the entire church, to which everything else had regard; whereas in Protestantism the pulpit and the word of God come into the foreground, and altar and sacrament stand back. Hence the altar was adorned also in the richest manner with costly cloths, with the cross, or at a later period the crucifix, with burning tapers, symbolical of Christ the light of the world,\footnote{This usage also no doubt came from Judaism into the Christian church; for in the temple at Jerusalem, and in the tabernacle before it, a lamp was perpetually burning according to divine command, Exod. xxvii. 30f. Probably lamps were in earlier use in the church. But tapers also were already in use in the time of Chrysostom, especially for lighting the altar, while lamps were rather employed in chapels and before images of saints.} and previously consecrated for ecclesiastical use,\footnote{In the Roman church the second of February, or the fortieth day after Christmas, when Mary presented the Lord in the temple, and when the aged Simeon prophetically called the child Jesus “a light to lighten the Gentiles,” is appointed for this consecration, and is hence called \emph{Candlemas of Mary}, a contraction of the two names, \emph{Purification of} \emph{Mary} and \emph{Candlemas}.} with a splendid copy of the Holy Scriptures, or the mass-book, but above all with the tabernacle, or little house for preserving the consecrated host, on which in the middle ages the German stone-cutters and sculptors displayed wonderful art.

Side altars did not come into use until Gregory the Great. Ignatius,\footnote{He even expressly (Ep. ad Philad. c. 4) likens the unity of the church in the episcopate to the unity of the altar: \textgreek{Ἔν θυσιαστήριον, ὡς εἷς ἐπίσκοπος}.} Athanasius, Gregory Nazianzen, and Augustine know of only one altar in the church. The Greek church has no more to this day. The introduction of such side altars, which however belong not to the altar space, but to the nave of the church, is connected with the progress of the worship of martyrs and relics.

At the left of the altar war, the table of prothesis,\footnote{\textgreek{Πρόθεσις.} oblationarium, still used in the Greek church.} on which the elements for the holy Supper were prepared, and which is still used in the Greek church; at the right the sacristy,\footnote{\textgreek{Σκευοφυλάκτιον, διακονικόν}, sacristia, sacrorum custodia, salutatorium, etc.} where the priests robed themselves, and retired for silent prayer. Behind the altar on the circular wall (and under the painting of Christ enthroned, if there was one) stood the bishop’s chair,\footnote{\textgreek{Θρόνος}, cathedra.} overlooking the whole church. On both sides of it, in a semicircle, were the seats of the presbyters. None but the clergy were allowed to receive the holy Supper within the altar rails.\footnote{Before Ambrosethe emperors were permitted to take their seats within the altar-space. But Ambrose, with the approval of Theodosius, abolished this custom, and assigned to the emperors a special place at the head of the congregation, just outside the rails. Sozomen, H. E. vii. 25.}

\xsection{Architectural Style. The Basilicas}

§ 106. Architectural Style. The Basilicas.

Comp. the works on the Basilicas by P. Sarnelli (Antica Basilicografia. Neapoli, 1686), Ciampini (Rom. 1693), Guttensohn \& Knapp (Monumenta di Rel. christ., ossia raccolta delle antiche chiese di Roma. Rom. 1822 sqq. 3 vols.; also in German, München, 1843), Bunsen (Die Basiliken des christlichen Roms. München, 1843, a commentary on the preceding), Von Quast (Berl. 1845), and Zestermann (Die antiken und die christlichen Basiliken. Leipz. 1847).

The history of church building, from the simple basilicas of the fourth century to the perfect Gothic cathedrals of the thirteenth and fourteenth, exhibits, like the history of the other Christian arts and the sciences, a gradual subjection and transformation of previous Jewish and heathen forms by the Christian principle. The church succeeded to the inheritance of all nations, but could only by degrees purge this inheritance of its sinful adulterations, pervade it with her spirit, and subject it to her aims; for she fulfils her mission through human freedom, not in spite of it, and does not magically transform nations, but legitimately educates them.

The history of Western architecture is the richer. The East contented itself with the Byzantine style, and adhered more strictly to the forms of the round temples, baptisteries, and mausoleums; while the West, starting from the Roman basilica, developed various styles.

The style of the earliest Christian churches was not copied from the heathen temples, because, apart from their connection with idolatry, which was itself highly offensive to the Christian sentiment, they were in form and arrangement, as we have already remarked, entirely unsuitable to Christian worship. The primitive Christian architecture followed the basilicas, and hence the churches built in this style were themselves called basilicas. The connection of the Christian and heathen basilicas, which has been hitherto recognized, and has been maintained by celebrated connoisseurs,\footnote{Bunsen, Schnaase, Kugler, Kinkel, Quast, \&e.} has been denied by some modem investigators,\footnote{Zestermann (1847) and Krauser (1851).} who have claimed for the Christian an entirely independent origin. And it is perfectly true, as concerns the interior arrangement and symbolical import of the building, that these can be ascribed to the Christian mind alone. Nor have any forensic or mercantile basilicas, to our knowledge, been transformed into Christian churches.\footnote{The passage quoted for this view from Ausonius in his address of thanks to the emperor Gratian, his pupil, c. 2: “Forum et basilica olim negotiis plena, nunc votis, votisque pro tua salute susceptis,” implies only, according to the connection, that now all houses and public places are full of good wishes for the emperor.} But in external architectural form there is without question an affinity, and there appears no reason why the church should not have employed this classic form.

The basilicas,\footnote{\textgreek{Στοαὶ βασιλικαί}The name comes from that of the highest civil magistrate, the \textgreek{ἄρχων βασιλεύς}, who held court in these buildings. In the church this designation was very naturally transferred to Christ, as the supreme King and Judge. Though of Greek origin, the basilicas first reached their full development in Rome, and, properly speaking, arose from the \emph{forum Romanum}. They were strictly \emph{fora} for the people, but roofed, and so protected from rain and heat. The city of Rome had ten of them: the Bas. Julia, Ulpia, Porcia, Marciana, \&c. Zestermann, however, denies the connection of the Roman basilica with the Athenian \textgreek{στοὰ βασίλειος ,}from the later times of Roman luxury, when the name basilicus was applied to everything grand and costly.} or royal halls, were public judicial and mercantile buildings, of simple, but beautiful structure, in the form of a long rectangle, consisting of a main hall, or main nave, two, often four, side naves,\footnote{Basilicas with a single nave are very rare. The pagan basilica of Trier is an instance, and the small church of St. Balbina in Rome, said to have been built by Gregory I. in the beginning of the seventh century.} which were separated by colonnades from the central space, and were somewhat lower. Here the people assembled for business and amusement. At the end of the hall opposite the entrance, stood a semicircular, somewhat elevated niche (apsis, tribune), arched over with a half-dome, where were the seats of the judges and advocates, and where judicial business was transacted. Under the floor of the tribunal was sometimes a cellar-like place of confinement for accused criminals.

In the history of architecture, too, there is a Nemesis. As the cross became changed from a sign of weakness to a sign of honor and victory, so must the basilica in which Christ and innumerable martyrs were condemned to death, become a place for the worship of the crucified One. The judicial tribune became the altar; the seat of the praetor behind it became the bishop’s chair; the benches of the jurymen became the seats of presbyters; the hall of business and trade became a place of devotion for the faithful people; the subterranean jail became a crypt or burial place, the superterrene birth-place, of a Christian martyr. To these were added other changes, especially the introduction of a cross-nave between the apse and the main nave, giving to the basilica the symbolical form of the once despised, but now glorious cross, and forming, so to speak, a recumbent crucifix. The cross with equal arms is called the Greek; that with unequal arms, in which the transept is shorter than the main nave from the entrance to the altar, the Latin. Towers, which express the heavenward spirit of the Christian religion, were not introduced till the ninth century, and were then built primarily for bells.

This style found rapid acceptance in the course of the fourth century with East and West; most of all in Rome, where a considerable number of basilicas, some in their ancient venerable simplicity, some with later alterations, are still preserved. The church of St. Maria Maggiore on the Esquiline hill affords the best view of an ancient basilica; the oldest principal church of Rome—S. Giovanni in Laterano (so named from the Roman patrician family of the Laterans), dedicated to the Evangelist John and to John the Baptist; the church of St. Paul, outside the city on the way to Ostia, which was burnt in 1823, but afterwards rebuilt splendidly in the same style, and consecrated by the pope in December, 1854; also S. Clemente, S. Agnese, and S. Lorenzo, outside the walls—are examples. The old church of St. Peter (Basilica Vaticana), which was built on the spot of this apostle’s martyrdom, the Neronian circus, and was torn down in the fifteenth century (the last remnant did not fall till 1606), surpassed all other churches of Rome in splendor and wealth, and was rebuilt, not in the same style, but, as is well known, in the Italian style of the sixteenth century.

Next to Rome, Ravenna is rich in old church buildings, among which the great basilica of S. Apollinare in Classe (in the port town, three miles from the main city, and built about the middle of the sixth century) is the most notable. The transept, as in all the churches of this city, is wanting.

In the East Roman empire there appeared even under Constantine sundry departures and transitions toward the Byzantine style. The oldest buildings there, which follow more or less the style of the Roman basilica, are the church at Tyre, begun in 313, destroyed in the middle ages, but known to us from the description of the historian Eusebius;\footnote{In the panegyric addressed to Paulinus, bishop of Tyre, Hist. Eccl. x. c. 4.} the original St. Sophia of Constantine in Constantinople; and the churches in the Holy Land, built likewise by him and his mother Helena, at, Mamre or Hebron, at Bethlehem over the birth-spot of Christ, on the Mount of Olives in memory of the ascension, and over the holy sepulchre on Mount Calvary. Justinian also sometimes built basilicas, for variety, together with his splendid Byzantine churches; and of these the church of St. Mary in Jerusalem was the finest, and was destined to imitate the temple of Solomon, but it was utterly blotted out by the Mohammedans.\footnote{Comp. the more minute descriptions of these churches in the above-mentioned illustrated work of Guttensohn and Knapp: Monumenta di religione christ., etc., 1822-’27, and the explanatory text by Bunsen: Die Basiliken des christl. Roms. München, 1848. Also Gottfried Kinkel: Geschichte der bildenden Künsten bei der christlichen Völkern, i. p. 61 sqq., and Ferd. von Quast: Die Basilika der Alten.}

\xsection{The Byzantine Style}

§ 107. The Byzantine Style.

Procopius: De aedificiis Justiniani. L.i.c. 1–3. Car. Dufresne Dom. du Cange: Constantinopolis Christiana. Venet. 1729. Salzenberg und Kortüm: Altchristliche Baudenkmale Constantinopels vom V. bis XII Jahrh. (40 magnificent copperplates and illustrations). Berlin, 1854.

The second style which meets us in this period, is the Byzantine, which in the West modified the basilica style, in the East soon superseded it, and in the Russo-Greek church has maintained itself to this day. It dates from the sixth century, from the reign of the scholarly and art-loving emperor Justinian I. (527–565), which was the flourishing period of Constantinople and of the centralized ecclesiastico-political despotism, in many respects akin to the age of Louis XIV. of France.

The characteristic feature of this style is the hemispherical dome, which, like the vault of heaven with its glory, spanned the centre of the Greek or the Latin cross, supported by massive columns (instead of slender pillars like the basilicas), and by its height and its prominence ruling the other parts of the building. This dome corresponds on the one hand to the centralizing principle of the Byzantine empire,\footnote{Kurtz, in his large Handbuch der K. Gesch., 3d ed. i. 372, well says: “The Byzantine state, in that maturity of it which Constantineintroduced and Justinian completed, was, in polity, as astonishing, gorgeous, majestic a centralized edifice, as the church of St. Sophia in architecture. The imperial power, as absolute autocracy, was the all-ruling, all-moving centre of the whole state life. The main dome, over-topping all, the full expression of the majesty of the centre, towards which all parts of the building strove, to which all were subservient, in the splendor of which all basked, was the court and the residence; on it the provinces and the authorities set over them leaned, as the subordinate side-domes or half-domes on the main one.”} but at the same time, and far more clearly than the flat basilica, to that upward striving of the Christian spirit from the earth towards the height of heaven, which afterwards more plainly expressed itself in the pointed arches and the towers of the Germanic cathedral. “While in the basilica style everything looks towards the end of the building where the altar and episcopal throne are set, and by this prevailing connection the upward direction is denied a free expression, in the dome structure everything concentrates itself about the spacious centre of the building over which, drawing the eye irresistibly upward, rises to an awe-inspiring height the majestic central dome. The basilica presents in the apse a figure of the horizon from which the sun of righteousness arises in his glory; the Byzantine building unfolds in the dome a figure of the whole vault of heaven in sublime, imposing majesty, but detracts thereby from the prominence of the altar, and leaves for it only a place of subordinate import.”

The dome is not, indeed, absolutely new. The Pantheon in Rome, whose imposing dome has a diameter of a hundred and thirty-two feet, dates from the age of Augustus, b.c. 26. But here the dome rises on a circular wall, and so strikes root in the earth, altogether in character with the heathen religion. The Byzantine dome rests on few columns connected by arches, and, like the vault of heaven, freely spans the central space of the church in airy height, without shutting up that space by walls.

Around the main central dome\footnote{\textgreek{Θόλος.}} stand four smaller domes in a square, and upon each dome rises a lofty gilded cross, which in the earlier churches stands upon a crescent, hung with all sorts of chains, and fastened by these to the dome.

The noblest and most complete building of this kind is the renowned church of St. Sophia at Constantinople, which was erected in lavish Asiatic splendor by the emperor Justinian after a plan by the architects Anthemius of Tralles and Isidore of Miletus (a.d. 532–537), and consecrated to the Redeemer,\footnote{The Wisdom, the Logos, of God; called in Proverbs and the Book of Wisdom \textgreek{σοφία}. Hence the name of the church. There is still standing in Constantinople a small church of St Sophia, which was likewise erected by Justinian.} but was transformed after the Turkish conquest into a Mohammedan mosque (Aja Sofia). It is two hundred and twenty-eight feet broad, and two hundred and fifty-two feet long; the dome, supported by four gigantic columns, rises a hundred and sixty-nine feet high over the altar, is a hundred and eight feet in diameter, and floats so freely and airily above the great central space, that, in the language of the Byzantine court biographer Procopius, it seems not to rest on terra firma, but to hang from heaven by golden chains.\footnote{In 557, the 32d year of Justinian, the eastern part of the dome fell in, and destroyed the altar, together with the tabernacle and the ambo, but was restored in 561. A similar misfortune befell it by an earthquake in the twelfth century, and again in 1346. The Turks let the grand structure gradually decay, till finally, by command of the Sultan, a.d. 1847-’49, a thorough restoration was undertaken under the direction of an Italian architect, Fossati. This brought to light the magnificence of the Mosaic pictures which Mohammedan picture-hatred and Turkish barbarism had in part destroyed, in part plastered over. The Sultan now caused them to be covered with plates of glass, cemented with lime; so that they are secure for a time, till the pile shall come again into the service of Christianity.} The most costly material was used in the building; the Phrygian marble with rose-colored and white veins, the dark red marble of the Nile, the green of Laconia, the black and white spotted of the Bosphorus, the gold-colored Libyan. And when the dome reflected the brilliance of the lighted silver chandeliers, and sent it back doubled from above, it might well remind one of the vault of heaven with its manifold starry glories, and account for the proud satisfaction with which Justinian on the day of the consecration, treading in solemn procession the finished building, exclaimed: “I have outdone thee, O Solomon!”\footnote{\textgreek{Νενίκηκά σε Σολομών}. Comp. the descriptions in Evagrius: Hist. Eccl. l iv. cap. 31; Procopius: De aedific. i. 1; and the poem of Paul Silentiarius: \textgreek{Εκφρασις ναοῦ τῆς Σοφίας}(a metrical translation of it in the above cited work of Salzenberg and Kortüm).} The church of St. Sophia stood thenceforth the grand model of the new Greek architecture, not only for the Christian East and the Russian church, but even for the Mohammedans in the building of their mosques.

In the West the city of Ravenna, on the Adriatic coast, after Honorius, (a.d. 404) the seat of the Western empire, or of the eparchate, and the last refuge of the old Roman magnificence and art, affords beautiful monuments of the Byzantine style; especially in the church of St. Vitale, which was erected by the bishop Maximian in 547.\footnote{Comp. on these Byzantine churches Kinkel, l. c., i. p. 100 sqq. and p. 121 sqq., and the splendid work of Salzenberg and Kortüm, Altchristliche Baudenkmale Konstantinopels, etc.}

In the West the ground plan of the basilica was usually retained, with pillars and entablature, until the ninth century, and the dome and vaultings of the Byzantine style were united with it. Out of this union arose what is called the Romanesque or the round-arch style, which prevailed from the tenth to the thirteenth century, and was then, from the thirteenth to the fifteenth, followed by the Germanic or pointed-arch style, with its gigantic masterpieces, the Gothic cathedrals. From the fifteenth century eclecticism and confusion prevailed in architecture, till the modern attempts to reproduce the ancient style. The Oriental church, on the contrary, has never gone beyond the Byzantine, its productivity almost entirely ceasing with the age of Justinian. But it is possible that the Graeco Russian church will in the future develop something new.

\xsection{Baptisteries, Grave-Chapels, and Crypts}

§ 108. Baptisteries, Grave-Chapels, and Crypts.

Baptisteries or Photisteries,\footnote{\textgreek{Φωτιστήρια,} places of enlightening; because the baptized were, according to Heb. vi. 4, called “enlightened.”} chapels designed exclusively for the administration of baptism, are a form of church building by themselves. In the first centuries baptism was performed on streams in the open air, or in private houses. But after the public exercise of Christian worship became lawful, in the fourth century special buildings for this holy ordinance began to appear, either entirely separate, or connected with the main church (at the side of the western main entrance) by a covered passage; and they were generally, dedicated to John the Baptist. The need of them arose partly from the still prevalent custom of immersion, partly from the fact that the number of candidates often amounted to hundreds and thousands; since baptism was at that time administered) as a rule, only three or four times a year, on the eve of the great festivals (Easter, Pentecost, Epiphany, and Christmas), and at episcopal sees, while the church proper was filled with the praying congregation.

These baptismal chapels were not oblong, like the basilicas, but round (like most of the Roman temples), and commonly covered with a dome. They had in the centre, like the bathing and swimming houses of the Roman watering places, a large baptismal basin,\footnote{\textgreek{Κολομβήθρα}, piscina fons baptismalis.} into which several steps descended. Around this stood a colonnade and a circular or polygonal gallery for spectators; and before the main entrance there was a spacious vestibule in the form of an entirely walled rectangle or oval. Generally the baptisteries had two divisions for the two sexes. The interior was sumptuously ornamented; especially the font, on which was frequently represented the symbolical figure of a hart panting for the brook, or a lamb, or the baptism of Christ by John. The earliest baptistery, of the Constantinian church of St. Peter in Rome, whose living flood was supplied from a fountain of the Vatican hill, was adorned with beautiful mosaic, the green, gold, and purple of which were reflected in the water. The most celebrated existing baptistery is that of the Lateran church at Rome, the original plan of which is ascribed to Constantine, but has undergone changes in the process of time.\footnote{In it, according to tradition, the emperor received baptism from pope Silvester I But this must be an error; for Constantinedid not receive baptism until he was on his death-bed in Nicomedia. Comp. § 2, above.}

After the sixth century, when the baptism of adults had become rare, it became customary to place a baptismal basin in the porch of the church, or in the church itself, at the left of the entrance, and, after baptism came to be administered no longer by the bishop alone, but by every pastor, each parish church contained such an arrangement. Still baptisteries also continued in use, and even in the later middle ages new ones were occasionally erected.

Finally, after the time of Constantine it became customary to erect small houses of worship or memorial chapels upon the burial-places of the martyrs, and to dedicate them to their memory.\footnote{Hence the name \textgreek{μαρτύρια,} \emph{martyrum memoriae, confessiones}. The clergy who officiated in them were called \textgreek{κληρικοὶ μαρτυρίων,} \emph{martyrarii}. The name \emph{capellae} occurs first in the seventh and eighth centuries, and is commonly derived from the \emph{cappa} (a clerical vestment covering the head and body) of St. Martin of Tours, which was preserved and carried about as a precious relic and as a national palladium of France.} These served more especially for private edification.

The subterranean chapels, or crypts, were connected with the churches built over them, and brought to mind the worship of the catacombs in the times of persecution. These crypts always produce a most earnest, solemn impression, and many of them are of considerable archaeological interest.

\xsection{Crosses and Crucifixes}

§ 109. Crosses and Crucifixes.

Jac. Gretser. (R.C.): De cruce Christi. 2 vols. Ingolst. 1608. Just. Lipsius: De cruce Christi. Antw. 1694. Fr. Münter: Die Sinnbilder u. Kunstvorstellungen der alten Christen. Altona, 1825. C. J. Hefele (R.C.): Alter u. älteste Form der Crucifixe (in the 2d vol. of his Beiträge zur Kirchengesch., Archäologie u. Liturgik. Tübingen, 1864, p. 265 sqq.).

The cross, as the symbol of redemption, and the signing of the cross upon the forehead, the eyes, the mouth, the breast, and even upon parts of clothing, were in universal use in this period, as they had been even in the second century, both in private Christian life and in public worship. They were also in many ways abused in the service of superstition; and the nickname cross-worshippers,\footnote{Religiosi crucis.} which the heathen applied to the Christians in the time of Tertullian,\footnote{Tert. Apolog. c. 16.} was in many cases not entirely unwarranted. Besides simple wooden crosses, now that the church had risen to the kingdom, there were many crosses of silver and gold, or sumptuously set with pearls and gems.\footnote{The cross occurs in three forms: the \emph{crux decussata} x(called St. Andrew’s cross, because this apostle is said to have died upon such an one); the \emph{crux commissa} T; and the \emph{crux immissa}, either with equal arms +(the Greek cross), or with unequal † (the Roman).}

The conspicuous part which, according to the statements of Eusebius, the cross played in the life of Constantine, is well known: forming the instrument of his conversion; borne by fifty men, leading him to his victories over Maxentius and Licinius; inscribed upon his banners, upon the weapons of his soldiers in his palace, and upon public places, and lying in the right hand of his own statue. Shortly afterwards Julian accused the Christians of worshipping the wood of the cross. “The sign of universal detestation,” says Chrysostom,\footnote{In the homily on the divinity of Christ, § 9, tom. i. 571.} “the sign of extreme penalty, is now become the object of universal desire and love. We see it everywhere triumphant; we find it on houses, on roofs, and on walls, in cities and hamlets, on the markets, along the roads, and in the deserts, on the mountains and in the valleys, on the sea, on ships, on books and weapons, on garments, in marriage chambers, at banquets, upon gold and silver vessels, in pearls, in painting upon walls, on beds, on the bodies of very sick animals, on the bodies of the possessed [—to drive away the disease and the demon—], at the dances of the merry, and in the brotherhoods of ascetics.” Besides this, it was usual to mark the cross on windows and floors, and to wear it upon the forehead.\footnote{\textgreek{Ἐκτυποῦν τόν σταυρὸν ἐν τῷ μετώπῳ}, effingere crucem in fronte, postare in fronte, which cannot always be understood as merely making the sign with the finger on the forehead. Comp. Neander, iii. 547, note.} According to Augustine this sign was to remind believers that their calling is to follow Christ in true humility, through suffering, into glory.

We might speak in the same way of the use of other Christian emblems from the sphere of nature; the representation of Christ by a good Shepherd, a lamb, a fish, and the like, which we have already observed in the period preceding.\footnote{Vol. i. § 100 (p. 377 sqq.).}

Towards the end of the present period we for the first time meet with crucifixes; that is, crosses not bare, but with the figure of the crucified Saviour upon them. The transition to the crucifix we find in the fifth century in the figure of a lamb, or even a bust of Christ, attached to the cross, sometimes at the top, sometimes at the bottom.\footnote{Crosses of this sort, colored red, with a white lamb, are thus described by Paulinus of Nola in the beginning of the fifth century, Epist. 32: \par “Sub cruce sanguinea niveo stat Christus in agno.”} Afterwards the whole figure of Christ was fastened to the cross, and the earlier forms gave place to this. The Trullan council of Constantinople (the Quinisextum), a.d. 692, directed in the 82d canon: “Hereafter, instead of the lamb, the human figure of Christ shall be set up on the images.”\footnote{\textgreek{Κατὰ τὸν ἀνθρώπινον χαρακτῆρα.}Hefele (l.c. 266 sq.) proves that crucifixes did not make their first appearance with this council, but that some existed before. The Venerable Bede, for example (Opp. ed. Giles, tom. iv. . p. 376), relates that a crucifix, bearing on one side the Crucified, on the other the serpent lifted up by Moses, was brought from Rome to the British cloister of Weremouth in 686. Gregory of Tours, also († 595), De gloria martyrum, lib. i. c. 23, describes a crucifix in the church of St. Genesius in Narbonne, which presented the Crucified One almost entirely naked (pictura, quae Dominum nostrum quasi praecinctum linteo indicat crucifixum). But this crucifix gave offence, and was veiled, by order of the bishop, with a curtain, and only at times exposed to the people.} But subsequently the orthodox church of the East prohibited all plastic images, crucifixes among them, and it tolerates only pictures of Christ and the saints. The earlier Latin crucifixes offend the taste and disturb devotion; but the Catholic art in its flourishing period succeeded in combining, in the figure of the suffering and dying Redeemer, the expression of the deepest and holiest anguish with that of supreme dignity. In the middle age there was frequently added to the crucifix a group of Mary, John, a soldier, and the penitent Magdalene, who on her knees embraced the post of the cross.

\xsection{Images of Christ}

§ 110. Images of Christ.

Fr. Kugler: Handbuch der Geschichte der Malerei seit Constantin dem Berlin, 1847, 2 vols.; and other works on the history of painting. Also C. Grüneisen: Die bildliche Darstellung der Gottheit. Stuttgart 1828. On the Iconoclastic controversies, comp. Maimbourg (R.C.): Histoire de l’hérésie de l’Iconoclastes. Par. 1679 sqq. 2 vols. Dallaeus (Calvinist): De imaginibus. Lugd. Bat. 1642. Fr. Spanheim: Historia imaginum restituta. Lugd. Bat. 1686. P. E. Jablonski († 1757): De origine imaginum Christi Domini, in Opuscul. ed. Water, Lugd. Bat. 1804, tom. iii. Walch: Ketzergesch., vols. x. and xi. J. Marx: Der Bildersturm der byzantinischen Kaiser. Trier, 1839. W. Grimm: Die Sage vom Ursprunge der Christusbilder. Berlin, 1843, L. Glückselig: Christus-Archäologie, Prag, 1863. Hefele: Beiträge zur Kirchengeschichte, vol. ii. Tüb. 1861 (Christusbilder, p. 254 sqq.). Comp. the liter. in Hase’s Leben Jesu, p. 79 (5th ed. 1865).

While the temple of Solomon left to the Christian mind no doubt concerning the lawfulness and usefulness of church architecture, the second commandment seemed directly to forbid a Christian painting or sculpture. “The primitive church,” says even a modern Roman Catholic historian,\footnote{Hefele, 1. c. p. 254.} “had no images, of Christ, since most Christians at that time still adhered to the commandment of Moses (Ex. xx. 4); the more, that regard as well to the Gentile Christians as to the Jewish forbade all use of images. To the latter the exhibition and veneration of images would, of course, be an abomination, and to the newly converted heathen it might be a temptation to relapse into idolatry. In addition, the church was obliged, for her own honor, to abstain from images, particularly from any representation of the Lord, lest she should be regarded by unbelievers as merely a new kind and special sort of heathenism and creature-worship. And further, the early Christians had in their idea of the bodily form of the Lord no temptation, not the slightest incentive, to make likenesses of Christ. The oppressed church conceived its Master only under the form of a servant, despised and uncomely, as Isaiah, liii. 2, 3, describes the Servant of the Lord.”

The first representations of Christ are of heretical and pagan origin. The Gnostic sect of the Carpocratians worshipped crowned pictures of Christ, together with images of Pythagoras, Plato, Aristotle, and other sages, and asserted that Pilate had caused a portrait of Christ to be made.\footnote{Irenaeus, Adv. haer. 1, 25, § 6: “Imagines quasdam quidem depictas, quasdam autem et de reliqua materia fabricatas habent, dicentes formam Christi factam a Pilato illo in tempore, quo fuit Jesus cum hominibus. Et has coronant et proponunt eas cum imaginibus mundi philosophorum, videlicet cum imagine Pythagorae et Platonis et Aristotelis et reliquorum; et reliquam observationem circa eas, similiter ut gentes, faciunt.” Comp. Epiphanius, Adv. haer. xxvi. no. 6; August., De haer. c. 7.} In the same spirit of pantheistic hero-worship the emperor Alexander Severus (a.d. 222–235) set up in his domestic chapel for his adoration the images of Abraham, Orpheus, Apollonius, and Christ.

After Constantine, the first step towards images in the orthodox church was a change in the conception of the outward form of Christ. The persecuted church had filled its eye with the humble and suffering servant-form of Jesus, and found therein consolation and strength in her tribulation. The victorious church saw the same Lord in heavenly glory on the right hand of the Father, ruling over his enemies. The one conceived Christ in his state of humiliation (but not in his state of exaltation), as even repulsive, or at least “having no form nor comeliness;” taking too literally the description of the suffering servant of God in Is. lii. 14 and liii. 2, 3.\footnote{So Justin Martyr, Dial. c. Tryph.; Clement, Alex., in several places of the Paedagogus and the Stromata; Tertullian, De carne Christi, c. 9, and Adv. Jud. c. 14; and Origen, Contra Cels. vi. c. 75. Celsus made this low conception of the form of the founder of their religion one of his reproaches against the Christians.} The other beheld in him the ideal of human beauty, “fairer than the children of men,” with “grace poured into his lips;” after the Messianic interpretation of Ps. xlv. 3.\footnote{So Chrysostom, Homil. 27 (al. 28) in Matth. (tom. vii. p. 371, in the new Paris ed.): \textgreek{Οὐδὲ γὰρ θαυματουργῶν ἦν θαυμαστὸς μόνον, ἀλλὰ καὶ φαινόμενος ἁπλῶς πολλῆς ἔγεμε χάριτος· καὶ τοῦτο ὁ προφήτης} (Ps. xlv.) \textgreek{δηλῶν ἔλεγεν· ὡραῖος κάλλει παρὰ τοὺς υἱοὺς τῶν ἀνθρώπων.}The passage in Isaiah (liii. 2) be refers to the ignominy which Christ suffered on the cross. So also Jerome, who likewise refers Ps. xlv. to the personal appearance of Jesus, and says of him: “Absque passionibus crucis universis [hominibus] pulchrior est .... Nisi enim babuisset et in vultu quiddam oculisque sidereum, numquam cum statim secuti fuissent apostoli, nec qui ad comprehendendum cum venerant. corruissent (Jno. xviii.).” Hieron. Ep. 65, c. 8.}

This alone, however, did not warrant images of Christ. For, in the first place, authentic accounts of the personal appearance of Jesus were lacking; and furthermore it seemed incompetent to human art duly to set forth Him in Whom the whole fulness of the Godhead and of perfect sinless humanity dwelt in unity.

On this point two opposite tendencies developed themselves, giving occasion in time to the violent and protracted image controversies, until, at the seventh ecumenical council at Nice in 787, the use and adoration of images carried the day in the church.

1. On the one side, the prejudices of the ante-Nicene period against images in painting or sculpture continued alive, through fear of approach to pagan idolatry, or of lowering Christianity into the province of sense. But generally the hostility was directed only against images of Christ; and from it, as Neander justly observes,\footnote{Kirchengesch., vol. iii. p. 550 (Germ. ed.).} we are by no means to infer the rejection of all representations of religious subjects; for images of Christ encounter objections peculiar to themselves.

The church historian Eusebius declared himself in the strongest manner against images of Christ in a letter to the empress Constantia (the widow of Licinius and sister of Constantine), who had asked him for such an image. Christ, says he, has laid aside His earthly servant-form, and Paul exhorts us to cleave no longer to the sensible;\footnote{Comp. 2 Cor. v. 16.} and the transcendent glory of His heavenly body cannot be conceived nor represented by man; besides, the second commandment forbids the making to ourselves any likeness of anything in heaven or in earth. He had taken away from a lady an image of Christ and of Paul, lest it should seem as if Christians, like the idolaters, carried their God about in images. Believers ought rather to fix their mental eye, above all, upon the divinity of Christ, and, for this purpose, to purify their hearts; since only the pure in heart shall see God.\footnote{In Harduin, Collect. concil. tom. iv. p. 406. A fragment of this letter of Eusebius is preserved in the acts of the council of the Iconoclasts at Constantinople in 754, and in the sixth act of the second council of Nice in 787.} The same Eusebius, however, relates of Constantine, without the slightest disapproval, that, in his Christian zeal, he caused the public monuments in the forum of the new imperial city to be adorned with symbolical representations of Christ, to wit, with figures of the good Shepherd and of Daniel in the lion’s den.\footnote{Vita Const. iii. c. 49.} He likewise tells us, that the woman of the issue of blood, after her miraculous cure (Matt. ix. 20), and out of gratitude for it, erected before her dwelling in Caesarea Philippi (Paneas) two brazen statues, the figure of a kneeling woman, and of a venerable man (Christ) extending his hand to help her, and that he had seen these statues with his own eyes at Paneas.\footnote{Hist. Eccl. lib. vii. cap. 18. According to Philostorgius (vii. 3), it was for a long time unknown whom the statues at Paneas represented, until a medicinal plant was discovered at their feet, and then they were transferred to the sacristy. The emperor Juliandestroyed them, and substituted his own statue, which was riven by lightning (Sozom. v. 21). Probably that statue of Christ was a monument of Hadrian or some other emperor, to whom the Phoenicians did obeisance in the form of a kneeling woman. Similar representations are to be seen upon coins, particularly of the time of Hadrian.} In the same place he speaks also of pictures (probably Carpocratian) of Christ and the apostles Peter and Paul, which he had seen, and observes that these cannot be wondered at in those who were formerly heathen, and who had been accustomed to testify their gratitude towards their benefactors in this way.

The narrow fanatic Epiphanius of Cyprus († 403) also seems to have been an opponent of images. For when he saw the picture of Christ or a saint\footnote{“Imaginem quasi Christi vel sancti cujusdam.”} on the altar-curtain in Anablatha, a village of Palestine, he tore away the curtain, because it was contrary to the Scriptures to hang up the picture of a man in the church, and he advised the officers to use the cloth for winding the corpse of some poor person.\footnote{Epiph. Ep. ad Joann Hierosolym., which Jeromehas preserved in a Latin translation. The Iconoclastic council at Constantinople in 754 cited several works of Epiphanius against images, the genuineness of which, however, is suspicious.} This arbitrary conduct, however, excited great indignation, and Epiphanius found himself obliged to restore the injury to the village church by another curtain.

2. The prevalent spirit of the age already very decidedly favored this material representation as a powerful help to virtue and devotion, especially for the uneducated classes, whence the use of images, in fact, mainly proceeded.

Plastic representation, it is true, was never popular in the East. The Greek church tolerates no statues, and forbids even crucifixes. In the West, too, in this period, sculpture occurs almost exclusively in bas relief and high relief, particularly on sarcophagi, and in carvings of ivory and gold in church decorations. Sculpture, from its more finite nature, lies farther from Christianity than the other arts.

Painting, on the contrary, was almost universally drawn into the service of religion; and that, not primarily from the artistic impulse which developed itself afterwards, but from the practical necessity of having objects of devout reverence in concrete form before the eye, as a substitute for the sacred books, which were accessible to the educated alone. Akin to this is the universal pleasure of children in pictures.

The church-teachers approved and defended this demand, though they themselves did not so directly need such helps. In fact, later tradition traced it back to apostolic times, and saw in the Evangelist Luke the first sacred painter. Whereof only so much is true: that he has sketched in his Gospel and in the Acts of the Apostles vivid and faithful pictures of the Lord, His mother, and His disciples, which are surely of infinitely greater value than all pictures in color and statues in marble.\footnote{Jerome, in his biographical sketch of Luke, De viris illustr. c. 7, is silent concerning this tradition (which did not arise till the seventh century or later), and speaks of Luke merely as medicus, according to Col. iv. 4.}

Basil the Great († 379) says “I confess the appearance of the Son of God in the flesh, and the holy Mary as the mother of God, who bore Him according to the flesh. And I receive also the holy apostles and prophets and martyrs. Their likenesses I revere and kiss with homage, for they are handed down from the holy apostles, and are not forbidden, but on the contrary painted in all our churches.”\footnote{Epist. 205. Comp. his Oratio in Barlaam, Opp. i. 515, and similar expressions in Gregory Naz., Orat. 19 (al. 18).} His brother, Gregory of Nyssa, also, in his memorial discourse on the martyr Theodore, speaks in praise of sacred painting, which “is wont to speak silently from the walls, and thus to do much good.” The bishop Paulinus of Nola, who caused biblical pictures to be exhibited annually at the festival seasons in the church of St. Felix, thought that by them the scenes of the Bible were made clear to the uneducated rustic, as they could not otherwise be; impressed themselves on his memory, awakened in him holy thoughts and feelings, and restrained him from all kinds of vice.\footnote{Paulinus, Carmen ix. et x. de S. Felicis natali.} The bishop Leontius of Neapolis in Cyprus, who at the close of the sixth century wrote an apology for Christianity against the Jews, and in it noticed the charge of idolatry, asserts that the law of Moses is directed not unconditionally against the use of religious images, but only against the idolatrous worship of them; since the tabernacle and the temple themselves contained cherubim and other figures; and he advocates images, especially for their beneficent influences. “In almost all the world,” says he, “profligate men, murderers, robbers, debauchees, idolaters, are daily moved to contrition by a look at the cross of Christ, and led to renounce the world, and practise every virtue.”\footnote{See the fragments of this apology in the 4th act of the second council of Nicaea, and Neander, iii. 560 (2d Germ. ed.), who adds the unprejudiced remark: “We cannot doubt that what Leontius here says, though rhetorically exaggerated, is nevertheless drawn from life, and is founded on impressions actually produced by the contemplation of images in certain states of feeling.”} And Leontius already appeals to the miraculous fact, that blood flowed from many of the images.\footnote{\textgreek{Πολλάκις αἱμάτων ῥύσεις ἐξ εἰκόνων γεγόνασι.}}

Owing to the difficulty, already noticed, of worthily representing Christ Himself, the first subjects were such scenes from the Old Testament as formed a typical prophecy of the history of the Redeemer. Thus the first step from the field of nature, whence the earliest symbols of Christ—the lamb, the fish, the shepherd—were drawn, was into the field of pre-Christian revelation, and thence it was another step into the province of gospel history itself. The favorite pictures of this kind were, the offering-up of Isaac—the pre-figuration of the great sacrifice on the cross; the miracle of Moses drawing forth water from the rock with his rod—which was interpreted either, according to 1 Cor. x. 4, of Christ Himself, or, more especially—and frequently, of the birth of Christ from the womb of the Virgin; the suffering Job—a type of Christ in His deepest humiliation; Daniel in the lion’s den—the symbol of the Redeemer subduing the devil and death in the underworld; the miraculous deliverance of the prophet Jonah from the whale’s belly—foreshadowing the resurrection;\footnote{Comp. Matt. xii. 39, 40.} and the translation of Elijah—foreshadowing the ascension of Christ.

About the middle of the fifth century, just when the doctrine of the person of Christ reached its formal settlement, the first representations of Christ Himself appeared, even said by tradition to be faithful portraits of the original.\footnote{The image-hating Nestorians ascribed the origin of iconolatry to their hated, in opponent, Cyril of Alexandria, and put it into connection with the Monophysite heresy (Assem., Bibl. orient. iii. 2, p. 401).} From that time the difficulty of representing the God-Man was removed by an actual representation, and the recognition of the images of Christ, especially of the Madonna with the Child, became even a test of orthodoxy, as against the Nestorian heresy of an abstract separation of the two natures in Christ. In the sixth century, according to the testimony of Gregory of Tours, pictures of Christ were hung not only in churches but in almost every private house.\footnote{De gloria martyrum, lib. i. c. 22.}

Among these representations of Christ there are two distinct types received in the church:

(1) The Salvator picture, with the expression of calm serenity and dignity, and of heavenly gentleness, without the faintest mark of grief. According to the legend, this was a portrait, miraculously imprinted on a cloth, which Christ Himself presented to Abgarus, king of Edessa, at his request.\footnote{4 First mentioned by the Armenian historian Moses of Chorene in the fifth century, partly on the basis of the spurious correspondence, mentioned by Eusebius (H. E i. 13), between Christ and Abgarus Uchomo of Edessa. The Abgarus likeness is said to have come, in the tenth century, into the church of St. Sophia at Constantinople, thence to Rome, where it is still shown in the church of St. Sylvester. But Genoa also pretends to possess the original. The two do not look much alike, and are of course only copies. Mr. Glückselig (Christus-Archaeologie, Prag, 1863) has recently made an attempt to restore from many copies an Edessenum redivivum.} The original is of course lost, or rather never existed, and is simply a mythical name for the Byzantine type of the likeness of Christ which appeared after the fifth century, and formed the basis of all the various representations of Christ until Raphael and Michael Angelo. These pictures present the countenance of the Lord in the bloom of youthful vigor and beauty, with a free, high forehead, clear, beaming eyes, long, straight nose, hair parted in the middle, and a somewhat reddish beard.

(2) The Ecce Homo picture of the suffering Saviour with the crown of thorns. This is traced back by tradition to St. Veronica, who accompanied the Saviour on the way to Golgotha, and gave Him her veil to wipe the sweat from His face; whereupon the Lord miraculously imprinted on the cloth the image of His thorn-crowned head.\footnote{This Veronica likeness is said to have come to Rome about a.d. 700, where it is preserved among the relics in St. Peter’s, but is shown only to noble personages. According to the common view, advocated especially by Mabillon and Papebroch, the name Veronica arose from the simple error of contracting the two words vera icon (\textgreek{εἰκών}), the \emph{true image}. W. Grimm considers the whole Veronica story a Latin version of the Greek Abgarus legend.}

The Abgarus likeness and the Veronica both lay claim to a miraculous origin, and profess to be εἰκόνες ἀχειροποίηται, pictures not made with human hands. Besides these, however, tradition tells of pictures of Christ taken in a natural way by Luke and by Nicodemus. The Salvator picture in the Lateran chapel Sancta Sanctorum in Rome, which is attributed to Luke, belongs to the Edessene or Byzantine type.

With so different pretended portraits of the Lord we cannot wonder at the variations of the pictures of Christ, which the Iconoclasts used as an argument against images. In truth, every nation formed a likeness of its own, according to its existing ideals of art and virtue.

Great influence was exerted upon the representations of Christ by the apocryphal description of his person in the Latin epistle of Publius Lentulus (a supposed friend of Pilate) to the Roman senate, delineating Christ as a man of slender form, noble countenance, dark hair parted in the middle, fair forehead, clear eyes, faultless mouth and nose, and reddish beard.\footnote{The letter of Lentulus has been rightly known in its present form only since the eleventh century. Comp. Gabler: De \textgreek{αὐθεντίᾳ} Epistolae Publii Lentuli ad Senatum R. de J. C. scriptae. Jenae, 1819, and 1822 (2 dissertations).} An older, and in some points different, description is that of John of Damascus, or some other writer of the eighth century, who says: “Christ was of stately form, with beautiful eyes, large nose, curling hair, somewhat bent, in the prime of life, with black beard, and sallow complexion, like his mother.”\footnote{Epist. ad Theoph. imper. de venerandis imag. (of somewhat doubtful origin), in Joh. Damasc. Opera, tom. i. p. 631, ed. Le Quien. A third description of the personal appearance of Christ, but containing nothing new, occurs in the fourteenth century, in Nicephorus Callisti, Hist. Eccl. lib. i. cap. 40.}

No figure of Christ, in color, or bronze, or marble, can reach the ideal of perfect beauty which came forth into actual reality in the Son of God and Son of man. The highest creations of art are here but feeble reflections of the original in heaven, yet prove the mighty influence which the living Christ continually exerts even upon the imagination and sentiment of the great painters and sculptors, and which He will exert to the end of the world.

\xsection{Images of Madonna and Saints}

§ 111. Images of Madonna and Saints.

Besides the images of Christ, representations were also made of prominent characters in sacred history, especially of the blessed Virgin with the Child, of the wise men of the east, as three kings worshipping before the manger,\footnote{Into the representation of the child Jesus in the manger the ox and ass were almost always brought, with reference to Is. i. 3: “The ox knoweth his owner, and the ass his master’s crib: but Israel doth not know, my people doth not consider.”} of the four Evangelists, the twelve Apostles, particularly Peter and Paul,\footnote{Usually Christ in the middle, and the leading apostles on either side. Augustine, De consensu Evangelist. i. 16: “Christus simul cum Petro et Paulo in pictis parietibus.”} of many martyrs and saints of the times of persecution, and honored bishops and monks of a later day.\footnote{Especially the pillar-saint, Symeon. The Antiochians had the picture of their deceased bishop Meletius on their seal rings, bowls, cups, and on the walls of their apartments. Comp. Chrysostom, Homil. in Miletium.}

According to a tradition of the eighth century or later, the Evangelist Luke painted not only Christ, but Mary also, and the two leading apostles. Still later legends ascribe to him even seven Madonnas, several of which, it is pretended, still exist; one, for example, in the Borghese chapel in the church of Maria Maggiore at Rome. The Madonnas early betray the effort to represent the Virgin as the ideal of female beauty, purity, and loveliness, and as resembling her divine Son.\footnote{The earliest pictures of the Madonna with the child are found in the Roman catacombs, and are traced in part by the Cavaliere de Rossi (Imagini Scelte, 1863) to the third and second centuries.} Peter is usually represented with a round head, crisped hair and beard; Paul, with a long face, bald crown, and pointed beard; both, frequently, carrying rolls in their hands, or the first the cross and the keys (of the kingdom of heaven), the second, the sword (of the word and the Spirit).

Such representations of Christ, of the saints, and of biblical events, are found in the catacombs and other places of burial, on sarcophagi and tombstones, in private houses, on cups and seal rings, and (in spite of the prohibition of the council of Elvira in 305)\footnote{Conc. Eliberin. or Illiberitin. can. 36: “Placuit picturas in ecclesia esse non debere, ne quod colitur aut adoratur, in parietibus depingatur.” This prohibition seems to have been confined, however, to pictures of Christ Himself; else we must suppose that martyrs and saints are accounted objects of \emph{cultus} and \emph{adoratio}.} on the walls of churches, especially behind the altar.

Manuscripts of the Bible also, liturgical books, private houses, and even the vestments of officials in the large cities of the Byzantine empire were ornamented with biblical pictures. Bishop Asterius of Amasea in Pontus, in the second half of the fourth century, protested against the wearing of these “God-pleasing garments,”\footnote{\textgreek{Ἱμαντια κεχαρισμένα τῷ Θεῷ..}} and advised that it were better with the proceeds of them to honor the living images of God, and support the poor; instead of wearing the palsied on the clothes, to visit the sick; and instead of carrying with one the image of the sinful woman kneeling and embracing the feet of Jesus, rather to lament one’s own sins with tears of contrition.

The custom of prostration\footnote{\textgreek{Προσκύνησις.}} before the picture, in token of reverence for the saint represented by it, first appears in the Greek church in the sixth century. And then, that the unintelligent people should in many cases confound the image with the object represented, attribute to the outward, material thing a magical power of miracles, and connect with the image sundry superstitious notions—must be expected. Even Augustine laments that among the rude Christian masses there are many image-worshippers,\footnote{De moribus ecclesiae cath. i. 75: “Novi multos esse picturarum adulatores.” The Manichaeans charged the entire catholic church with image-worship.} but counts such in the great number of those nominal Christians, to whom the essence of the Gospel is unknown.

As works of art, these primitive Christian paintings and sculptures are, in general, of very little value; of much less value than the church edifices. They are rather earnest and elevated, than beautiful and harmonious. For they proceeded originally not from taste, but from practical want, and, at least in the Greek empire, were produced chiefly by monks. It perfectly befitted the spirit of Christianity, to begin with earnestness and sublimity, rather than, as heathenism, with sensuous beauty. Hence also its repugnance to the nude, and its modest draping of voluptuous forms; only hands, feet, and face were allowed to appear.

The Christian taste, it is well known, afterwards changed, and, on the principle that to the pure all things are pure, it represented even Christ on the cross, and the holy Child at His mother’s breast or in His mothers arms, without covering.

Furthermore, in the time of Constantine the ancient classical painting and sculpture had grievously degenerated; and even in their best days they reached no adequate expression of the Christian principle.

In this view, the loss of so many of those old works of art, which, as the sheer apparatus of idolatry, were unsparingly destroyed by the iconoclastic storms of the succeeding period, is not much to be regretted. It was in. the later middle ages, when church architecture had already reached its height, that Christian art succeeded in unfolding an unprecedented bloom of painting and sculpture, and in far surpassing, on the field of painting at least, the masterpieces of the ancient Greeks. Sculpture, which can present man only in his finite limitation, without the flush of life or the beaming eye, like a shadowy form from the realm of the dead, probably attained among the ancient Greeks the summit of perfection, above which even Canova and Thorwaldsen do not rise. But painting, which can represent man in his organic connection with the world about him, and, to a certain degree, in his unlimited depth of soul and spirit, as expressed in the countenance and the eye, has waited for the influence of the Christian principle to fulfil its perfect mission, and in the Christs of Leonardo da Vinci, Fra Beato Angelico, Correggio, and Albrecht Dürer, and the Madonnas of Raphael, has furnished the noblest works which thus far adorn the history of the art.

\xsection{Consecrated Gifts}

§ 112. Consecrated Gifts.

It remains to mention in this connection yet another form of decoration for churches, which had already been customary among heathen and Jews: consecrated gifts. Thus the temple of Delphi, for example, had become exceedingly rich through such presents of weapons, silver and golden vessels, statues, \&c. In almost every temple of Neptune hung votive tablets, consecrated to the god in thankfulness for deliverance from shipwreck by him.\footnote{Comp. Horace, Ars poet. v. 20.} A similar custom seems to have existed among the Jews; for I Sam. xxi. implies that David had deposited the sword of the Philistine Goliath in the sanctuary. In the court of the priests a multitude of swords, lances, costly vessels, and other valuable things, were to be seen.

Constantine embellished the altar space in the church of Jerusalem with rich gifts of gold, silver, and precious stones. Sozomen tells us\footnote{H. E. iv. 25.} that Cyril, bishop of Jerusalem, in a time of famine, sold the treasures and sacred gifts of the church, and that afterwards some one recognized in the dress of an actress the vestment he once presented to the church.

A peculiar variety of such gifts, namely, memorials of miraculous cures,\footnote{\textgreek{Ἐκτυπώματα.}} appeared in the fifth century; at least they are first mentioned by Theodoret, who said of them in his eighth discourse on the martyrs: “That those who ask with the confidence of faith, receive what they ask, is plainly proved by their sacred gifts in testimony of their healing. Some offer feet, others hands, of gold or silver, and these gifts show their deliverance from those evils, as tokens of which they have been offered by the restored.” With the worship of saints this custom gained strongly, and became in the middle age quite universal. Whoever recovered from a sickness, considered himself bound first to testify by a gift his gratitude to the saint whose aid he had invoked in his distress. Parents, whose children fortunately survived the teething-fever, offered to St. Apollonia (all whose teeth, according to the legend, had been broken out with pincers by a hangman’s servant) gifts of jawbones in wax. In like manner St. Julian, for happily accomplished journeys, and St. Hubert, for safe return from the perils of the chase, were very richly endowed; but the Virgin Mary more than all. Almost every church or chapel which has a miracle-working image of the mother of God, possesses even now a multitude of golden and silver acknowledgments of fortunate returns and recoveries.

\xsection{Church Poetry and Music}

§ 113. Church Poetry and Music.

J. Rambach: Anthologie christl. Gesänge aus allen Jahrh. der christl. Kirche. Altona, 1817–’33. H. A. Daniel: Thesaurus hymnologicus. Hal. 1841–’56, 5 vols. Edélestand du Méril: Poésies populaires latines antérieures au douzième siècle. Paris, 1843. C. Fortlage: Gesänge der christl. Vorzeit. Berlin, 1844. G. A. Königsfeld u. A. W. v. Schlegel: Altchristliche Hymnen u. Gesaenge lateinisch u. Deutsch. Bonn, 1847. Second collection by Königsfeld, Bonn, 1865. E. E. Koch: Geschichte des Kirchenlieds u. Kirchengesangs der christl., insbesondere der deutschen evangel. Kirche. 2d ed. Stuttgart, 1852 f. 4 vols. (i. 10–30). F. J. Mone: Latein. Hymnen des Mittelalters (from MSS.), Freiburg, 1853–’55. (Vol. i., hymns of God and angels; ii., h. of Mary; iii., h. of saints.) Bässler: Auswahl Alt-christl. Lieder vom 2–15ten Jahrh. Berlin, 1858. R. Ch. Trench: Sacred Latin Poetry, chiefly lyrical, selected and arranged for use; with Notes and Introduction (1849), 2d ed. improved, Lond. and Cambr. 1864. The valuable hymnological works of Dr. J. M. Neale (of Sackville College, Oxford): The Ecclesiastical Latin Poetry of the Middle Ages (in Henry Thompson’s History of Roman Literature, Lond. and Glasgow., 1852, p. 213 ff.); Mediaeval Hymns and Sequences, Lond. 1851; Sequentiae ex Missalibus, 1852; Hymns of the Eastern Church, 1862, several articles in the Ecclesiologist; and a Latin dissertation, De Sequentiis, in the Essays on Liturgiology, etc., p. 359 sqq. (Comp. also J. Chandler: The Hymns of the Primitive Church, now first collected, translated, and arranged, Lond. 1837.)

Poetry, and its twin sister music, are the most sublime and spiritual arts, and are much more akin to the genius of Christianity, and minister far more copiously to the purposes of devotion and edification than architecture, painting, and sculpture. They employ word and tone, and can speak thereby more directly to the spirit than the plastic arts by stone and color, and give more adequate expression to the whole wealth of the world of thought and feeling. In the Old Testament, as is well known, they were essential parts of divine worship; and so they have been in all ages and almost all branches of the Christian church.

Of the various species of religious poetry, the hymn is the earliest and most important. It has a rich history, in which the deepest experiences of Christian life are stored. But it attains full bloom in the Evangelical church of the German and English tongue, where it, like the Bible, becomes for the first time truly the possession of the people, instead of being restricted to priest or choir.

The hymn, in the narrower sense, belongs to lyrical poetry, or the poetry of feeling, in distinction from the epic and dramatic. It differs also from the other forms of the lyric (ode, elegy, sonnet, cantata, \&c.) in its devotional nature, its popular form, and its adaptation to singing. The hymn is a popular spiritual song, presenting a healthful Christian sentiment in a noble, simple, and universally intelligible form, and adapted to be read and sung with edification by the whole congregation of the faithful. It must therefore contain nothing inconsistent with Scripture, with the doctrines of the church, with general Christian experience, or with the spirit of devotion. Every believing Christian can join in the Gloria in Excelsis or the Te Deum. The classic hymns, which are, indeed, comparatively few, stand above confessional differences, and resolve the discords of human opinions in heavenly harmony. They resemble in this the Psalms, from which all branches of the militant church draw daily nourishment and comfort. They exhibit the bloom of the Christian life in the Sabbath dress of beauty and holy rapture. They resound in all pious hearts, and have, like the daily rising sun and the yearly returning spring, an indestructible freshness and power. In truth, their benign virtue increases with increasing age, like that of healing herbs, which is the richer the longer they are bruised. They are true benefactors of the struggling church, ministering angels sent forth to minister to them who shall be heirs of salvation. Next to the Holy Scripture, a good hymn-book is the richest fountain of edification.

The book of Psalms is the oldest Christian hymn-book, inherited by the church from the ancient covenant. The appearance of the Messiah upon earth was the beginning of Christian poetry, and was greeted by the immortal songs of Mary, of Elizabeth, of Simeon, and of the heavenly host. Religion and poetry are married, therefore, in the gospel. In the Epistles traces also appear of primitive Christian songs, in rhythmical quotations which are not demonstrably taken from the Old Testament.\footnote{\emph{E.g}., Eph. v. 14, where either the Holy Spirit moving in the apostolic poesy, or (as I venture to suggest) the previously mentioned \emph{Light} personified, is introduced (\textgreek{διὸ λέγει}) speaking in three strophes: \par \textgreek{Ἔγειρε ὁ καθεύδων,} \par \textgreek{Καὶ ἀνάστα ἐκ τῶν νεκρῶν·} \par \textgreek{Καὶ ἐπιφαύσει σοι ὁ Χριστός.} \par Comp. Rev. iv. 8; 1 Tim. iii. 16; 2 Tim. ii. 11; and my History of the Apostolic Church, § 141.} We know from the letter of the elder Pliny to Trajan, that the Christians, in the beginning of the second century, praised Christ as their God in songs; and from a later source, that there was a multitude of such songs.\footnote{2 Comp. Euseb. H. E. v. 28.}

Notwithstanding this, we have no complete religious song remaining from the period of persecution, except the song of Clement of Alexandria to the divine Logos—which, however, cannot be called a hymn, and was probably never intended for public use—the Morning Song\footnote{\textgreek{ὝΥμνος ἑωθινός}, beginning: \textgreek{Δόξα ἐν ὑψίστοις Θεῷ}, in Const. Apost. vii. 47 (al, 48), and in Daniel’s Thesaur. hymnol. iii. p. 4.} and the Evening Song\footnote{\textgreek{ὝΥμνος ἐσπερινός}, which begins: \textgreek{Φῶς ἱλαρὸν ἁγίας δόξης}, see Daniel, iii 5.} in the Apostolic Constitutions, especially the former, the so-called Gloria in Excelsis, which, as an expansion of the doxology of the heavenly hosts, still rings in all parts of the Christian world. Next in order comes the Te Deum, in its original Eastern form, or the καθ  ἑκάστην ἡμέραν, which is older than Ambrose. The Ter Sanctus, and several ancient liturgical prayers, also may be regarded as poems. For the hymn is, in fact, nothing else than a prayer in the festive garb of poetical inspiration, and the best liturgical prayers are poetical creations. Measure and rhyme are by no means essential.

Upon these fruitful biblical and primitive Christian models arose the hymnology of the ancient catholic church, which forms the first stage in the history of hymnology, and upon which the mediaeval, and then the evangelical Protestant stage, with their several epochs, follow.

\xsection{The Poetry of the Oriental Church}

§ 114. The Poetry of the Oriental Church.

Comp. the third volume of Daniel’s Thesaurus hymnologicus (the Greek section prepared by B. Vormbaum); the works of J. M. Neale, quoted sub § 113; an article on Greek Hymnology in the Christian Remembrancer, for April, 1859, London; also the liturgical works quoted § 98.

We should expect that the Greek church, which was in advance in all branches of Christian doctrine and culture, and received from ancient Greece so rich a heritage of poetry, would give the key also in church song. This is true to a very limited extent. The Gloria in excelsis and the Te Deum are unquestionably the most valuable jewels of sacred poetry which have come down from the early church, and they are both, the first wholly, the second in part of Eastern origin, and going back perhaps to the third or second century.\footnote{That the so-called Hymnus angelicus, based on Luke ii. 14, is of Greek origin, and was used as a morning hymn, is abundantly proven by Daniel, Thesaurus hymnol. tom. ii. p. 267 sqq. It is found in slightly varying forms in the Apostolic Constitutions, l. vii. 47 (al. 48), in the famous Alexandrian Codex of the Bible, and other places. Of the so called Ambrosian hymn or Te Deum, parts at least are Greek, Comp. Daniel, l. c. p. 276 sqq.} But, excepting these hymns in rhythmic prose, the Greek church of the first six centuries produced nothing in this field which has had permanent value or general use.\footnote{We cannot agree with the anonymous author of the article in the “Christian Remembrancer” for April, 1859, p. 282, who places Cosmas of Maiuma as high as Adam of S. Victor, John of Damascus as high as Notker, Andrew of Crete as high as S. Bernard, and thinks Theophanes and Theodore of the Studium in no wise inferior to the best of Sequence writers of the eleventh and twelfth centuries.} It long adhered almost exclusively to the Psalms of David, who, as Chrysostom says, was first, middle, and last in the assemblies of the Christians, and it had, in opposition to heretical predilections, even a decided aversion to the public use of uninspired songs. Like the Gnostics before them, the Arians and the Apollinarians employed religious poetry and music as a popular means of commending and propagating their errors, and thereby, although the abuse never forbids the right use, brought discredit upon these arts. The council of Laodicea, about a.d. 360, prohibited even the ecclesiastical use of all uninspired or “private hymns,”\footnote{Can. 59: \textgreek{Οὐ δεῖ ἰδιωτικοὺς ψαλμοὺς λέγεσθαι ἐν τῇ ἐκκλησίᾳ.} By this must doubtless be understood not only heretical, but, as the connection shows, all extrabiblical hymns composed by men, in distinction from the \textgreek{κανονικὰ βιβλία τῆς καινῆς καὶ παλαιᾶς διαθήκης .}} and the council of Chalcedon, in 451, confirmed this decree.

Yet there were exceptions. Chrysostom thought that the perverting influence of the Arian hymnology in Constantinople could be most effectually counteracted by the positive antidote of solemn antiphonies and doxologies in processions. Gregory Nazianzen composed orthodox hymns in the ancient measure; but from their speculative theological character and their want of popular spirit, these hymns never passed into the use of the church. The same may be said of the productions of Sophronius of Jerusalem, who glorified the high festivals in Anacreontic stanzas; of Synesius of Ptolemais (about a.d. 410), who composed philosophical hymns; of Nonnus of Panopolis in Egypt, who wrote a paraphrase of the Gospel of John in hexameters; of Eudoxia, the wife of the emperor Theodosius II.; and of Paul Silentiarius, a statesman under Justinian I., from whom we have several epigrams and an interesting poetical description of the church of St. Sophia, written for its consecration. Anatolius, bishop of Constantinople († 458), is properly the only poet of this period who realized to any extent the idea of the church hymn, and whose songs were adapted to popular use.\footnote{Neale, in his Hymns of the Eastern Church, p. 3 sqq., gives several of them in free metrical reproduction. See below.}

The Syrian church was the first of all the Oriental churches to produce and admit into public worship a popular orthodox poetry, in opposition to the heretical poetry of the Gnostic Bardesanes (about a.d. 170) and his son Harmonius. Ephraim Syrus († 378) led the way with a large number of successful hymns in the Syrian language, and found in Isaac, presbyter of Antioch, in the middle of the fifth century, and especially in Jacob, bishop of Sarug in Mesopotamia († 521), worthy successors.\footnote{On the Syrian hymnology there are several special treatises, by Augusti: De hymnis Syrortim sacris, 1814; Hahn: Bardesanes Gnosticus, Syrorum primus hymnologus, 1819; Zingerle: Die heil. Muse der Syrer, 1833 (with German translations from Ephraim). Comp. also Jos. Six. Assemani: Bibl. orient. i. 80 sqq. (with Latin versions), and Daniel’s Thes. hymnol. tom. iii. 1855, pp. 139-268. The Syrian hymns for Daniel’s Thesaurus were prepared by L. Splieth, who gives them with the German version of Zingerle. An English version by H. Burgess: Select metrical Hymns and Homilies of Ephraem S., Lond. 1853, 2 vols.}

After the fifth century the Greek church lost its prejudices against poetry, and produced a great but slightly known abundance of sacred songs for public worship.

In the history of the Greek church poetry, as well as the Latin, we may distinguish three epochs: (1) that of formation, while it was slowly throwing off classical metres, and inventing its peculiar style, down to about 650; (2) that of perfection, down to 820; (3) that of decline and decay, to 1400 or to the fall of Constantinople. The first period, beautiful as are some of the odes of Gregory of Nazianzen and Sophronius of Jerusalem, has impressed scarcely any traces on the Greek office books. The flourishing period of Greek poetry coincides with the period of the image controversies, and the most eminent poets were at the same time advocates of images; pre-eminent among them being John of Damascus, who has the double honor of being the greatest theologian and the greatest poet of the Greek church.

The flower of Greek poetry belongs, therefore, in a later division of our history. Yet, since we find at least the rise of it in the fifth century, we shall give here a brief description of its peculiar character.

The earliest poets of the Greek church, especially Gregory Nazianzen, in the fourth, and Sophronius of Jerusalem in the seventh century, employed the classical metres, which are entirely unsuitable to Christian ideas and church song, and therefore gradually fell out of use.\footnote{See some odes of Gregory, Euthymius and Sophronius in Daniel’s Thes. tom. iii. p. 5 sqq. He gives also the hymn of Clement of Alex. (\textgreek{ὕμνος τοῦ σωτῆρος Χριστοῦ}), the \textgreek{ὔμνος ἑωθινός ,}and \textgreek{ὕμνος ἑσπερινὸς}, of the third century.} Rhyme found no entrance into the Greek church. In its stead the metrical or harmonic prose was adopted from the Hebrew poetry and the earliest Christian hymns of Mary, Zacharias, Simeon, and the angelic host. Anatolius of Constantinople († 458) was the first to renounce the tyranny of the classic metre and strike out a new path. The essential points in the peculiar system of the Greek versification are the following:\footnote{See the details in Neale’s works, whom we mainly follow as regards the Eastern hymnology, and in the article above alluded to in the “Christian Remembrancer” (probably also by Neale).}

The first stanza, which forms the model of the succeeding ones, is called in technical language Hirmos, because it draws the others after it. The succeeding stanzas are called Troparia (stanzas), and are divided, for chanting, by commas, without regard to the sense. A number of troparia, from three to twenty or more, forms an Ode, and this corresponds to the Latin Sequence, which was introduced about the same time by the monk Notker in St. Gall. Each ode is founded on a hirmos and ends with a troparion in praise of the Holy Virgin.\footnote{Hence this last troparion is called \emph{Theotokion}, from \textgreek{θεοτόκος}, the constant predicate of the Virgin Mary. The \emph{Stauro-theotokion} celebrates Mary at the cross.} The odes are commonly arranged (probably after the example of such Psalms as the 25th, 112th, and 119th) in acrostic, sometimes in alphabetic, order. Nine odes form a Canon.\footnote{\textgreek{Κανών.}Neale says (Hymns of the East. Ch. Introd. p. xxix.): “A canon consists of Nine Odes—each Ode containing any number of troparia from three to beyond twenty. The reason for the number nine is this: that there are nine Scriptural canticles employed at Lauds (\textgreek{εἰς τὸν Ὄρθρον}), on the model of which those in every Canon are formed. The first: that of Moses after the passage of the Red Sea—the second, that of Moses in Deuteronomy (ch. xxxiii.)- the third, that of Hannah—the fourth, that of Habakkuk—the fifth, that of Isaiah (ch. xxvi. 9-20)—the sixth, that of Jonah—the seventh, that of the Three Children (verses 3-34, our “Song” in the Bible Version)—the eighth, Benedicite—the ninth, \emph{Magnificat} and Benedict\emph{us}.”} The older odes on the great events of the incarnation, the resurrection, and the ascension, are sometimes sublime; but the later long canons, in glorification of unknown martyrs are extremely prosaic and tedious and full of elements foreign to the gospel. Even the best hymnological productions of the East lack the healthful simplicity, naturalness, fervor, and depth of the Latin and of the Evangelical Protestant hymn.

The principal church poets of the East are Anatolius († 458), Andrew of Crete (660–732), Germanus I. (634–734), John Of Damascus († about 780), Cosmas of Jerusalem, called the Melodist (780), Theophanes (759–818), Theodore of the Studium (826), Methodius I. (846), Joseph of the Studium (830), Metrophanes of Smyrna († 900), Leo VI. (886–917), and Euthymius († 920).

The Greek church poetry is contained in the liturgical books, especially in the twelve volumes of the Menaea, which correspond to the Latin Breviary, and consist, for the most part, of poetic or half-poetic odes in rhythmic prose.\footnote{Neale, l. c. p. xxxviii., says of the Oriental Breviary: “This is the staple of those three thousand pages—under whatever name the stanzas may be presented forming Canons and Odes; as Troparia, Idiomela, Stichera, Stichoi, Contakia, Cathismata, Theotokia, Triodia, Stauro-theotokia, Catavasiai—or whatever else. Nine-tenths of the Eastern Service-book is poetry.” Besides these we find poetical pieces also in the other liturgical books: the \emph{Paracletice} or the \emph{Great Octoechus}, in eight parts (for eight weeks and Sundays), the small \emph{Octoechus}, the \emph{Triodion} (for the Lent season), and the \emph{Pentecostarion} (for the Easter season). Neale (p. xli.) reckons that all these volumes together would form at least 5,000 closely-printed, double column quarto pages, of which 4,000 pages would be poetry. He adds an expression of surprise at the “marvellous ignorance in which English ecclesiastical scholars are content to remain of this huge treasure of divinity—the gradual completion of nine centuries at least.” Respecting the value of these poetical and theological treasures, however, few will agree with this learned and enthusiastic Anglican venerator of the Oriental church.} These treasures, on which nine centuries have wrought, have hitherto been almost exclusively confined to the Oriental church, and in fact yield but few grains of gold for general use. Neale has latterly made a happy effort to reproduce and make accessible in modern English metres, with very considerable abridgments, the most valuable hymns of the Greek church.\footnote{Neale, in his preface, says of his translations: “These are literally, I believe, the only English versions of any part of the treasures of Oriental Hymnology. There is scarcely a first or second-rate hymn of the Roman Breviary which has not been translated: of many we have six or eight versions. The eighteen quarto volumes of Greek church-poetry can only at present be known to the English reader by my little book.”}

We give a few specimens of Neale’s translations of hymns of St. Anatolius, patriarch of Constantinople, who attended the council of Chalcedon (451). The first is a Christmas hymn, commencing in Greek:

Μέγα καὶ παράδοξον θαῦμα.

Another specimen of a Christmas hymn by the same, commencing ἐν Βηθλεέμ:

One more on Christ calming the storm, ζοφερᾶς τρικυμίας, as reproduced by Neale:

\xsection{The Latin Hymn}

§ 115. The Latin Hymn.

More important than the Greek hymnology is the Latin from the fourth to the sixteenth century. Smaller in compass, it surpasses it in artless simplicity and truth, and in richness, vigor, and fulness of thought, and is much more akin to the Protestant spirit. With objective churchly character it combines deeper feeling and more subjective appropriation and experience of salvation, and hence more warmth and fervor than the Greek. It forms in these respects the transition to the Evangelical hymn, which gives the most beautiful and profound expression to the personal enjoyment of the Saviour and his redeeming grace. The best Latin hymns have come through the Roman Breviary into general use, and through translations and reproductions have become naturalized in Protestant churches. They treat for the most part of the great facts of salvation and the fundamental doctrines of Christianity. But many of them are devoted to the praises of Mary and the martyrs, and vitiated with superstitions.

In the Latin church, as in the Greek, heretics gave a wholesome impulse to poetical activity. The two patriarchs of Latin church poetry, Hilary and Ambrose, were the champions of orthodoxy against Arianism in the West.

The genius of Christianity exerted an influence, partly liberating, partly transforming, upon the Latin language and versification. Poetry in its youthful vigor is like an impetuous mountain torrent, which knows no bounds and breaks through all obstacles; but in its riper form it restrains itself and becomes truly free in self-limitation; it assumes a symmetrical, well-regulated motion and combines it with periodical rest. This is rhythm, which came to its perfection in the poetry of Greece and Rome. But the laws of metre were an undue restraint to the new Christian spirit which required a new form. The Latin poetry of the church has a language of its own, a grammar of its own, a prosody of its own, and a beauty of its own, and in freshness, vigor, and melody even surpasses the Latin poetry of the classics. It had to cast away all the helps of the mythological fables, but drew a purer and richer inspiration from the sacred history and poetry of the Bible, and the heroic age of Christianity. But it had first to pass through a state of barbarism like the Romanic languages of the South of Europe in their transition from the old Latin. We observe the Latin language under the influence of the youthful and hopeful religion of Christ, as at the breath of a second spring, putting forth fresh blossoms and flowers and clothing itself with a new garment of beauty, old words assuming new and deeper meanings, obsolete words reviving, new words forming. In all this there is much to offend a fastidious classical taste, yet the losses are richly compensated by the gains. Christianity at its triumph in the Roman empire found the classical Latin rapidly approaching its decay and dissolution; in the course of time it brought out of its ashes a new creation.

The classical system of prosody was gradually loosened, and accent substituted for quantity. Rhyme, unknown to the ancients as a system or rule, was introduced in the middle or at the end of the verse, giving the song a lyrical character, and thus a closer affinity with music. For the hymns were to be sung in the churches. This accented and rhymed poetry was at first, indeed, very imperfect, yet much better adapted to the freedom, depth, and warmth of the Christian spirit, than the stereotyped, stiff, and cold measure of the heathen classics.\footnote{Archbishop Trench (Sacred Latin Poetry, 2d ed. Introd. p. 9): “A struggle commenced from the first between the form and the spirit, between the old heathen form and the new Christian spirit—the latter seeking to release itself from the shackles and restraints which the former imposed upon it; and which were to it, not a help and a support, as the form should be, but a hindrance and a weakness—not liberty, but now rather a most galling bondage. The new wine went on fermenting in the old bottles, till it burst them asunder, though not itself to be spilt and lost in the process, but to be gathered into nobler chalices, vessels more fitted to contain it—new, even as that which was poured into them was new.” This process of liberation Trench illustrates in Prudentius, who still adheres in general to the laws of prosody, but indulges the largest license.} Quantity is a more or less arbitrary and artificial device; accent, or the emphasizing of one syllable in a polysyllabic word, is natural and popular, and commends itself to the ear. Ambrose and his followers, with happy instinct, chose for their hymns the Iambic dimeter, which is the least metrical and the most rhythmical of all the ancient metres. The tendency to euphonious rhyme went hand in hand with the accented rhythm, and this tendency appears occasionally in its crude beginnings in Hilary and Ambrose, but more fully in Damasus, the proper father of this improvement.

Rhyme is not the invention of either a barbaric or an overcivilized age, but appears more or less in almost all nations, languages, and grades of culture. Like rhythm it springs from the natural esthetic sense of proportion, euphony, limitation, and periodic return.\footnote{Comp. the excellent remarks of Trench, l. c. p. 26 sqq., on the import of rhyme. Milton, as is well known, blinded by his predilection for the ancient classics, calls rhyme (in the preface to “Paradise Lost”) “the invention of a barbarous age, to set off wretched matter and lame metre; a thing of itself to all judicious ears trivial and of no true musical delight.” Trench answers this biassed judgment by pointing to Milton’s own rhymed odes and sonnets,” the noblest lyrics which English literature possesses.”} It is found here and there, even in the oldest popular poetry of republican Rome, that of Ennius, for example.\footnote{“It is a curious thing,” says J. M. Neale (The Eccles. Lat. Poetry of the Middle Ages, p. 214), “that, in rejecting the foreign laws in which Latin had so long gloried, the Christian poets were in fact merely reviving in an inspired form, the early melodies of republican Rome;—the rhythmical ballads which were the delight of the men that warred with the Samnites, and the Volscians, and Hannibal.”} It occurs not rarely in the prose even of Cicero, and especially of St. Augustine, who delights in ingenious alliterations and verbal antitheses, like patet and latet, spes and res, fides and vides, bene and plene, oritur and moritur. Damasus of Rome introduced it into sacred poetry.\footnote{In his Hymnus de S. Agatha, see Daniel, Thes. hymnol. tom. i. p. 9, and Fortlage, Gesänge christl. Vorzeit, p. 365.} But it was in the sacred Latin poetry of the middle age that rhyme first assumed a regular form, and in Adam of St. Victor, Hildebert, St. Bernard, Bernard of Clugny, Thomas Aquinas, Bonaventura, Thomas a Celano, and Jacobus de Benedictis (author of the Stabat mater), it reached its perfection in the twelfth and thirteenth centuries; above all, in that incomparable giant hymn on the judgment, the tremendous power of which resides, first indeed in its earnest matter, but next in its inimitable mastery of the musical treatment of vowels. I mean, of course, the Dies irae of the Franciscan monk Thomas a Celano (about 1250), which excites new wonder on every reading, and to which no translation in any modern language can do full justice. In Adam of St. Victor, too, of the twelfth century, occur unsurpassable rhymes; e.g., the picture of the Evangelist John (in the poem: De, S. Joanne evangelista), which Olshausen has chosen for the motto of his commentary on the fourth Gospel, and which Trench declares the most beautiful stanza in the Latin church poetry:

The metre of the Latin hymns is various, and often hard to be defined. Gavanti\footnote{Thesaur. rit sacr., cited in the above-named hymnological work of Königsfeld and A. W. Schlegel, p. xxi., first collection.} supposes six principal kinds of verse:

1. Iambici dimetri(as: “Vexilla regis prodeunt”).

2. Iambici trimetri(ternarii vel senarii, as: “Autra deserti teneris sub annis”).

3. Trochaici dimetri(“Pange, lingua, gloriosi corporis mysterium,” a eucharistic hymn of Thomas Aquinas).

4. Sapphici, cum Adonicoin fine (as: “Ut queant axis resonare fibris”).

5. Trochaici(as: “Ave maris stella”).

6. Asclepiadici, cum Glyconicoin fine (as: “Sacris solemniis juncta sint gaudia”).

In the period before us the Iambic dimeter prevails; in Hilary and Ambrose without exception.

\xsection{The Latin Poets and Hymns}

§ 116. The Latin Poets and Hymns.

The poets of this period, Prudentius excepted, are all clergymen, and the best are eminent theologians whose lives and labors have their more appropriate place in other parts of this work.

Hilary, bishop of Poitiers (hence Pictaviensis, † 368), the Athanasius of the West in the Arian controversies, is, according to the testimony of Jerome,\footnote{Catal. vir. illustr. c. 100. Comp. also Isidore of Seville, De offic. Eccles. l. i., and Overthür, in the preface to his edition of the works of Hilary.} the first hymn writer of the Latin church. During his exile in Phrygia and in Constantinople, he became acquainted with the Arian hymns and was incited by them to compose, after his return, orthodox hymns for the use of the Western church. He thus laid the foundation of Latin hymnology. He composed the beautiful morning hymn: “Lucis largitor splendide;” the Pentecostal hymn: “Beata nobis gaudia;” and, perhaps, the Latin reproduction of the famous Gloria in excelsis. The authorship of many of the hymns ascribed to him is doubtful, especially those in which the regular rhyme already appears, as in the Epiphany hymn:

“Jesus refulsit omnium

Pius redemptor gentium.”

We give as a specimen a part of the first three stanzas of his morning hymn, which has been often translated into German and English:\footnote{The Latin has 8 stanzas. See Daniel, Thesaur. hymnol. tom. i. p. 1.}

Ambrose, the illustrious bishop of Milan, though some-what younger († 397), is still considered, on account of the number and value of his hymns, the proper father of Latin church song, and became the model for all successors. Such was his fame as a hymnographer that the words Ambrosianus and hymnus were at one time nearly synonymous. His genuine hymns are distinguished for strong faith, elevated but rude simplicity, noble dignity, deep unction, and a genuine churchly and liturgical spirit. The rhythm is still irregular, and of rhyme only imperfect beginnings appear; and in this respect they certainly fall far below the softer and richer melodies of the middle age, which are more engaging to ear and heart. They are an altar of unpolished and unhewn stone. They set forth the great objects of faith with apparent coldness that stands aloof from them in distant adoration; but the passion is there, though latent, and the fire of an austere enthusiasm burns beneath the surface. Many of them have, in addition to their poetical value, a historical and theological value as testimonies of orthodoxy against Arianism.\footnote{Trench sees in the Ambrosian hymns, not without reason (I. c. p. 86), “a rocklike firmness, the old Roman stoicism transmuted and glorified into that nobler Christian courage, which encountered and at length overcame the world.” Fortlage judged the same way before in a brilliant description of Latin hymns, l. c. p. 4 f., comp. Daniel, Cod. Lit. iii. p. 282 sq.}

Of the thirty to a hundred so-called Ambrosian hymns,\footnote{Daniel, ii. pp. 12-115.} however, only twelve, in the view of the Benedictine editors of his works, are genuine; the rest being more or less successful imitations by unknown authors. Neale reduces the number of the genuine Ambrosian hymns to ten, and excludes all which rhyme regularly, and those which are not metrical. Among the genuine are the morning hymn: “Aeterne rerum conditor;”\footnote{The genuineness of this hymn is put beyond question by two quotations of the contemporary and friend of Ambrose, Augustine, Confess. ix. 12, and Retract. i. 12, and by the affinity of it with a passage in the Hexaëmeron of Ambrose, xxiv. 88, where the same thoughts are expressed in prose. Not so certain is the genuineness of the other Ambrosian morning hymns: “Aeterna coeli gloria,” and “Splendor paternae gloriae.”} the evening hymn: “Deus creator omnium;”\footnote{The other evening hymn: “O lux beata Trinitas,” ascribed to him (in the Roman Breviary and in Daniel’s Thesaur. i. 36), is scarcely from Ambrose: it has already the rhyme in the form as we find it in the hymns of Fortunatus.} and the Advent or Christmas hymn: “Veni, Redemptor gentium.” This last is justly considered his best. It has been frequently reproduced in modern languages,\footnote{Especially in the beautiful German by John Frank: “Komm, Heidenheiland, Lösegeld,” which is a free recomposition rather than a translation. For another English version (abridged), see “The Voice of Christian Life in Song,” p. 97: \par “Redeemer of the nations, come; \par Pure offspring of the Virgin’s womb, \par Seed of the woman, promised long, \par Let ages swell Thine advent song.”} and we add this specimen of its matter and form with an English version:

By far the most celebrated hymn of the Milanese bishop, which alone would have made his name immortal, is the Ambrosian doxology, Te Deum laudamus. This, with the Gloria in excelsis, is, as already remarked, by far the most valuable legacy of the old Catholic church poetry; and will be prayed and sung with devotion in all parts of Christendom to the end of time. According to an old legend, Ambrose composed it on the baptism of St. Augustine, and conjointly with him; the two, without preconcert, as if from divine inspiration, alternately singing the words of it before the congregation. But his biographer Paulinus says nothing of this, and, according to later investigations, this sublime Christian psalm is, like the Gloria in excelsis, but a free reproduction and expansion of an older Greek hymn in prose, of which some constituents appear in the Apostolic Constitutions, and elsewhere.\footnote{“For instance, the beginning of a morning hymn, in the Codex Alexandrinus of the Bible, has been literally incorporated into the Te Deum: \par \textgreek{Καθ  ἑκάστην ἡμέραν εὐλογήσω σε,} Per singulas dies benedicimus te, \par \textgreek{Καὶ αἰνέσω τὸ ὄνομά σου εἰς τὸν αἰῶνα} Et laudamus nomen tuum in saeculum \par \textgreek{Καὶ εἰς τὸν αἰῶνα τοῦ αἰῶνος.} Et in saeculum saeculi. \par \textgreek{Καταξίωσον, κύριε, καὶ τὴν ἡμέραν ταύτην} Dignare, Domine, die isto \par \textgreek{Ἀναμαρτήτους φυλαχθήναι ἡμᾶς.} Sine peccato nos custodire. \par Comp. on this whole hymn the critical investigation of Daniel, l.c. vol. ii, p. 289}

Ambrose introduced also an improved mode of singing in Milan, making wise use of the Greek symphonies and antiphonies, and popular melodies. This Cantus Ambrosianus, or figural song, soon supplanted the former mode of reciting the Psalms and prayers in monotone with musical accent and little modulation of the voice, and spread into most of the Western churches as a congregational song. It afterwards degenerated, and was improved and simplified by Gregory the Great, and gave place to the so-called Cantus Romanus, or choralis.

Augustine, the greatest theologian among the church fathers († 430), whose soul was filled with the genuine essence of poetry, is said to have composed the resurrection hymn: “Cum rex gloriae Christus;” the hymn on the glory of paradise: “Ad perennis vitae fontem melis sitivit arida;” and others. But he probably only furnished in the lofty poetical intuitions and thoughts which are scattered through his prose works, especially in the Confessions, the materia carminis for later poets, like Peter Damiani, bishop of Ostia, in the eleventh century, who put into flowing verse Augustine’s meditations on the blessedness of heaven.\footnote{This beautiful hymn, “De gloria et gaudiis Paradisi,” is found in the appendix to the 6th volume of the Benedictine edition of the Opera Augustini, in Daniel’s Thesaurus, tom. i. p. 116, and in Trench’s Collection, p. 315 sqq., and elsewhere. Like all the new Jerusalem hymns it derives its inspiration from St. John’s description in the concluding chapters of the Apocalypse. There is an excellent German translation of it by Königsfeld and an English translation by Wackerbarth, given in part by Neale in his Mediaeval Hymns and Sequences, p. 59. The whole hymn is very fine, but not quite equal to the long poem of Bernard of Cluny (in the twelfth century), on the contempt of the world, which breathes the same sweet home-sickness to heaven, and which Neale (p. 58) justly regards as the most lovely, in the same way that the \emph{Dies irae}, is the most sublime, and the \emph{Stabat Mater} the most pathetic, of mediaeval hymns. The original has not less than 3,000 lines; Neale gives an admirable translation of the concluding part, commencing “Hic breve vivitur,” and a part of this translation: To thee, O dear, dear Country” (p. 55), is well worthy of a place in our hymn books. From these and similar mediaeval sources (as the “Urbs beata Jerusalem,” \&c.) is derived in part the famous English hymn: “ O mother dear, Jerusalem!” (in 31 stanzas), which is often ascribed to David Dickson, a Scotch clergyman of the seventeenth century, and which has in turn become the mother of many English hymns on the new Jerusalem. (Comp. on it the monographs of H. Bonar, Edinb. 1852, and of W. C. Prime: ” O Mother dear, Jerusalem,” New York, 1865.)—To Augustineis also ascribed the hymn: “O gens beata ccelitum,” a picture of the blessedness of the inhabitants of heaven, and: Quid, tyranne! quid miraris? ” an antidote for the tyranny of sin.} Damasus, bishop of Rome († 384), a friend of Jerome, likewise composed some few sacred songs, and is considered the author of the rhyme.\footnote{Jerome(De viris ill.c. 103) says of him: “Elegans in versibus componendis ingenium habet, multaque et brevia metro edidit.” Neale omits Damasus altogether. Daniel, Thes. i. pp. 8 and 9, gives only two of his hymns, a Hymnus de S. Andrea, and a Hymnus de S. Agatha, the latter with regular rhymes, commencing: \par “Martyris ecce dies Agathae Christus eam sibi qua sociat \par Virginis emicat eximiae, Et diadema duplex decorat.”}

Coelius Sedulius, a native of Scotland or Ireland, presbyter in the first half of the fifth century, composed the hymns: “Herodes, hostis impie,” and “A solis ortus cardine,” and some larger poems.

Marcus Aurelius Clemens Prudentius († 405), an advocate and imperial governor in Spain under Theodosius, devoted the last years of his life to religious contemplation and the writing of sacred poetry, and stands at the head of the more fiery and impassioned Spanish school. Bently calls him the Horace and Virgil of Christians, Neale, “the prince of primitive Christian poets.” Prudentius is undoubtedly the most gifted and fruitful of the old Catholic poets. He was master of the classic measure, but admirably understood how to clothe the new ideas and feelings of Christianity in a new dress. His poems have been repeatedly edited.\footnote{2 E.g., by Th. Obbarius, Tub. 1845; and by Alb. Dressel, Lips. 1860.} They are in some cases long didactic or epic productions in hexameters, of much historical value;\footnote{The \emph{Apotheosis}, a celebration of the divinity of Christ against its opponents (in 1,063 lines); the \emph{Harmatigenia}, on the origin of sin (in 966 lines); the \emph{Psychomachia}, on the warfare of good and evil in the soul (915 lines); \emph{Contra Symmachum}, on idolatry, \&c.} in others, collections of epic poems, as the Cathemerinon,\footnote{\textgreek{Καθημερινῶν} = Diurnorum (the \emph{Christian Day}, as we might call it, after the analogy of Keble’s \emph{Christian} \emph{Year}), hymns for the several hours of the day.} and Peristephanon.\footnote{\textgreek{Περὶ στεφάνων}, concerning the crowns, fourteen hymns on as many martyrs who have inherited the crown of eternal life. Many of them are intolerably tedious and in bad taste.} Extracts from the latter have passed into public use. The best known hymns of Prudentius are: “Salvete, flores martyrum,” in memory of the massacred innocents at Bethlehem,\footnote{3 \emph{De SS. Innocentibus}, from the twelfth book of the Cathemerinon, in Prudentii Carmina, ed. Obbarius, Tüb. 1845, p. 48, in Daniel, tom. i. p. 124, and in Trench, p. 121.} and his grand burial hymn: “Jam moesta quiesce querela,” which brings before us the ancient worship in deserts and in catacombs, and of which Herder says that no one can read it without feeling his heart moved by its touching tones.\footnote{It is the close of the tenth Cathemerinon, and was the usual burial hymn of the ancient church. It has been translated into German by Weiss, Knapp, Puchta, Königsfeld, Bässler, Schaff (in his Deutsches Gesangbuch, No. 468), and others. Trench, p. 281, calls it “the crowning glory of the poetry of Prudentius.” He never attained this grandeur on any other occasion. Neale, in his treatise on the Eccles. Latin Poetry, l.c. p. 22, gives translations of several parts of it, in the metre of the original, but without rhyme, commencing thus: \par “Each sorrowful mourner be silent! \par Fond mothers, give over your weeping! \par None grieve for those pledges as perished: \par This dying is life’s reparation.” \par Another translation by E. Caswall: “Cease, ye tearful mourners.”}

We must mention two more poets who form the transition from the ancient Catholic to mediaeval church poetry.

Venantius Fortunatus, an Italian by birth, a friend of queen Radegunde (who lived apart from her husband, and presided over a cloister), the fashionable poet of France, and at the time of his death (about 600), bishop of Poitiers, wrote eleven books of poems on various subjects, an epic on the life of St. Martin of Tours, and a theological work in vindication of the Augustinian doctrine of divine grace. He was the first to use the rhyme with a certain degree of mastery and regularity, although with considerable license still, so that many of his rhymes are mere alliterations of consonants or repetitions of vowels.\footnote{Such as prodeunt—mysterium, viscera—vestigia, fulgida—purpura, etc.} He first mastered the trochaic tetrameter, a measure which, with various modifications, subsequently became the glory of the mediaeval hymn. Prudentius had already used it once or twice, but Fortunatus first grouped it into stanzas. His best known compositions are the passion hymns: “Vexilla regis prodeunt,” and “Pange, lingua, gloriosi proelium (lauream) certaminis,” which, though not without some alterations, have passed into the Roman Breviary.\footnote{Daniel, Thes. i. p. 160 sqq., gives both forms: the original, and that of the Brev. Romanum.} The “Vexilla regis” is sung on Good Friday during the procession in which the consecrated host is carried to the altar. Both are used on the festivals of the Invention and the Elevation of the Cross.\footnote{Trench has omitted both in his Collection, and admitted instead of them some less valuable poems of Fortunatus, De cruce Christi, and De passione Domini, in hexameters.} The favorite Catholic hymn to Mary: “Ave maris stella,”\footnote{3 Daniel, i. p. 204.} is sometimes ascribed to him, but is of a much later date.

We give as specimens his two famous passion hymns, which were composed about 580.

Vexilla Regis Prodeunt.\footnote{274 The original text in Daniel, i. p. 160. The translation by Neale, from the Hymnal of the English Ecclesiological Society and Neale’s Mediaeval Hymns p. 6. it omits the second stanza, as does the Roman Breviary.}

Pange, Lingua, Gloriosi Proelium Certaminis.\footnote{See the original, which is not rhymed, in Daniel, i. p. 163 sqq., and in somewhat different form in the Roman Breviary. The masterly English translation in, the metre of the original is Neale’s, l.c. p. 237 sq., and in his Mediaeval Hymns and Sequences, p. 1. Another excellent English version by E. Caswell commences: Sing, my tongue, the Saviour’s glory; tell His triumph far and wide.”}

Far less important as a poet is Gregory I. (590–604), the last of the fathers and the first of the mediaeval popes. Many hymns of doubtful origin have been ascribed to him and received into the Breviary. The best is his Sunday hymn: “Primo dierum omnium.”\footnote{See Daniel’s Cod. i. p. 175 sqq. For au excellent English version of the hymn above alluded to, see Neale, l. c. p. 233.}

The hymns are the fairest flowers of the poetry of the ancient church. But besides them many epic and didactic poems arose, especially in Gaul and Spain, which counteracted the invading flood of barbarism, and contributed to preserve a connection with the treasures of the classic culture. Juvencus, a Spanish presbyter under Constantine, composed the first Christian epic, a Gospel history in four books (3,226 lines), on the model of Virgil, but as to poetic merit never rising above mediocrity. Far superior to him is Prudentius († 405); he wrote, besides the hymns already mentioned, several didactic, epic, and polemic poems. St. Pontius Paulinus, bishop of Nola († 431), who was led by the poet Ausonius to the mysteries of the Muses,\footnote{Ausonius yielded the palm to his pupil when he wrote of the verses of Paulinus: \par “Cedimus ingenio, quantum praecedimus aevo: \par Assurget Musae nostra camoena tuae.”} and a friend of Augustine and Jerome, is the author of some thirty poems full of devout spirit; the best are those on the festival of S. Felix, his patron. Prosper Aquitanus († 460), layman, and friend of Augustine, wrote a didactic poem against the Pelagians, and several epigrams; Avitus, bishop of Vienne († 523), an epic on the creation and the origin of evil; Arator, a court official under Justinian, afterwards a sub-deacon of the Roman church (about 544), a paraphrase, in heroic verse, of the Acts of the Apostles, in two books of about 1,800 lines. Claudianus Mamertus,\footnote{Not to be confounded with Claudius Claudianus, of Alexandria, the most gifted Latin poet at the end of the fourth and beginning of the fifth century. The Christian Idyls, Epistles, and Epigrams ascribed to him, were probably the work of Claudianus Mamertus, of Vienne (Comp. H. Thompson’s Manual of Rom. Lit. p. 204, and J. J. Brunet’s Manual du libraire, tom. iii. p. 1351 of the 5th ed. Par. 1862). For Claudius Claudianus was a heathen, according to the express testimony of Paulus Orosiusand of Augustine(De Civit. Dei, v. p. 26: “Poeta Claudianus, quamvis a Christi nomine alienus,” \&c.), and in one of his own epigrams, \emph{In Jacobum, magistrum equitum}, shows his contempt of the Christian religion.} Benedictus Paulinus, Elpidius, Orontius, and Draconti

\xchapter{Theological Controversies, and Development of the Ecumenical Orthodoxy}

THEOLOGICAL CONTROVERSIES, AND DEVELOPMENT OF THE ECUMENICAL ORTHODOXY.

\xsection{General Observations. Doctrinal Importance of the Period. Influence of the Ancient Philosophy}

§ 117. General Observations. Doctrinal Importance of the Period. Influence of the Ancient Philosophy.

The Nicene and Chalcedonian age is the period of the formation and ecclesiastical settlement of the ecumenical orthodoxy; that is, the doctrines of the holy trinity and of the incarnation and the divine-human person of Christ, in which the Greek, Latin, and evangelical churches to this day in their symbolical books agree, in opposition to the heresies of Arianism and Apollinarianism, Nestorianism and Eutychianism. Besides these trinitarian and christological doctrines, anthropology also, and soteriology, particularly the doctrines of sin and grace, in opposition to Pelagianism and Semi-Pelagianism were developed and brought to a relative settlement; only, however, in the Latin church, for the Greek took very little part in the Pelagian controversy.

The fundamental nature of these doctrines, the greatness of the church fathers who were occupied with them, and the importance of the result, give this period the first place after the apostolic in the history of theology. In no period, excepting the Reformation of the sixteenth century, have there been so momentous and earnest controversies in doctrine, and so lively an interest in them. The church was now in possession of the ancient philosophy and learning of the Roman empire, and applied them to the unfolding and vindication of the Christian truth. In the lead of these controversies stood church teachers of imposing talents and energetic piety, not mere book men, but venerable theological characters, men all of a piece, as great in acting and suffering as in thinking. To them theology was a sacred business of heart and life,\footnote{Or, as Gregory Nazianzen says of the true theologian, contemplation was a prelude to action, and action a prelude to contemplation, \textgreek{πρᾶξις}(a religious walk) \textgreek{ἐπίβασις θεωρίας}(actio gradus est ad contemplationem), Oratio xx. 12 (ed. Bened. Paris. tom. i. p. 383).} and upon them we may pass the judgment of Eusebius respecting Origen: “Their life was as their word, and their word was as their life.”

The theological controversies absorbed the intellectual activity of that time, and shook the foundations of the church and the empire. With the purest zeal for truth were mingled much of the odium and rabies theologorum, and the whole host of theological passions; which are the deepest and most bitter of passions, because religion is concerned with eternal interests.

The leading personages in these controversies were of course bishops and priests. By their side fought the monks, as a standing army, with fanatical zeal for the victory of orthodoxy, or not seldom in behalf even of heresy. Emperors and civil officers also mixed in the business of theology, but for the most part to the prejudice of its free, internal development; for they imparted to all theological questions a political character, and entangled them with the cabals of court and the secular interests of the day. In Constantinople, during the Arian controversy, all classes, even mechanics, bankers, frippers, market women, and runaway slaves took lively part in the questions of Homousion and sub-ordination, of the begotten and the unbegotten.\footnote{So Gregory of Nyssa (not Nazianzen, as J. H. Kurtz, wrongly quoting from Neander, has it in his large K. Gesch. i. ii. p. 99) relates from his own observation: Orat. de Deitate Filii et Spiritus S. (Opera ii. p. 898, ed. Paris. of 1615). He compares his contemporaries in this respect with the Athenians, who are always wishing to hear some new thing.}

The speculative mind of the Eastern church was combined with a deep religious earnestness and a certain mysticism, and at the same time with the Grecian curiosity and disputatiousness, which afterwards rather injured than promoted her inward life. Gregory Nazianzen, who lived in Constantinople in the midst of the Arian wars, describes the division and hostility which this polemic spirit introduced between parents and children, husbands and wives, old and young, masters and slaves, priests and people. “It has gone so far that the whole market resounds with the discourses of heretics, every banquet is corrupted by this babbling even to nausea, every merrymaking is transformed into a mourning, and every funeral solemnity is almost alleviated by this brawling as a still greater evil; even the chambers of women, the nurseries of simplicity, are disturbed thereby, and the flowers of modesty are crushed by this precocious practice of dispute.”\footnote{Orat. xxvii. 2 (Opera, tom. i. p. 488). Comp. Orat. xxxii. (tom. i. p. 581); Carmen de vita sua, vers. 1210 sqq. (tom. ii. p. 737 sq.).} Chrysostom, like Melanchthon at a later day, had much to suffer from the theological pugnacity of his times.

The history of the Nicene age shows clearly that the church of God carries the heavenly treasure in earthly vessels. The Reformation of the sixteenth century was likewise in fact an incessant war, in which impure personal and political motives of every kind had play, and even the best men often violated the apostolic injunction to speak the truth in love. But we must not forget that the passionate and intolerant dogmatism of that time was based upon deep moral earnestness and strong faith, and so far forth stands vastly above the tolerance of indifferentism, which lightly plays with the truth or not rarely strikes out in most vehement intolerance against the faith. (Remember the first French revolution.) The overruling of divine Providence in the midst of these wild conflicts is unmistakable, and the victory of the truth appears the greater for the violence of error. God uses all sorts of men for his instruments, and brings evil passions as well as good into his service. The Spirit of truth guided the church through the rush and the din of contending parties, and always triumphed over error in the end.

The ecumenical councils were the open battle-fields, upon which the victory of orthodoxy was decided. The doctrinal decrees of these councils contain the results of the most profound discussions respecting the Trinity and the person of Christ; and the Church to this day has not gone essentially beyond those decisions.

The Greek church wrought out Theology and Christology, while the Latin church devoted itself to Anthropology and Soteriology. The one, true to the genius of the Greek nationality, was predominantly speculative, dialectical, impulsive, and restless; the other, in keeping with the Roman character, was practical, traditional, uniform, consistent, and steady. The former followed the stimulation of Origen and the Alexandrian school; the latter received its impulse from Tertullian and Cyprian, and reached its theological height in Jerome and Augustine. The speculative inclination of the Greek church appeared even in its sermons, which not rarely treated of the number of worlds, the idea of matter, the different classes of higher spirits, the relation of the three hypostases in the Godhead, and similar abstruse questions. The Latin church also, however, had a deep spirit of investigation (as we see in Tertullian and Augustine), took an active part in the trinitarian and christological controversies of the East, and decided the victory of orthodoxy by the weight of its authority. The Greek church almost exhausted its productive force in those great struggles, proved indifferent to the deeper conception of sin and grace, as developed by Augustine, and after the council of Chalcedon degenerated theologically into scholastic formalism and idle refinements.

The fourth and fifth centuries are the flourishing, classical period of the patristic theology and of the Christian Graeco-Roman civilization. In the second half of the fifth century the West Roman empire, with these literary treasures, went down amidst the storms of the great migration, to take a new and higher sweep in the Germano-Roman form under Charlemagne. In the Eastern empire scholarship was better maintained, and a certain connection with antiquity was preserved through the medium of the Greek language. But as the Greek church had no middle age, so it has had no Protestant Reformation.

The prevailing philosophy of the fathers was the Platonic, so far as it was compatible with the Christian spirit. The speculative theologians of the East, especially those of the school of Origen, and in the West, Ambrose and pre-eminently Augustine, were moulded by the Platonic idealism.

A remarkable combination of Platonism with Christianity, to the injury of the latter, appears in the system of mystic symbolism in the pseudo-Dionysian books, which cannot have been composed before the fifth century, though they were falsely ascribed to the Areopagite of the book of Acts (xvii. 34), and proceeded from the later school of New-Platonism, as represented by Proclus of Athens († 485). The fundamental idea of these Dionysian writings (on the celestial hierarchy; on the ecclesiastical hierarchy; on the divine names; on mystic theology; together with ten letters) is a double hierarchy, one in heaven and one on earth, each consisting of three triads, which mediates between man and the ineffable, transcendent hyper-essential divinity. This idea is a remnant of the aristocratic spirit of ancient heathenism, and forms the connecting link with the hierarchical organization of the church, and explains the great importance and popularity which the pseudo-Dionysian system acquired, especially in the mystic theology of the middle ages.\footnote{Comp. Engelhardt: Die angeblichen Schriften des Areop. Dionysius übersetzt und erklärt, 1823, 2 Parts; Ritter: Geschichte der christl. Philosophie, Bd. ii. p. 515; Baur: Geschichte der Lehre von der Dreieinigkeit, ii. 207 f., and his Geschichte der Kirche, from the fourth to the sixth century, p. 59 ff.; Joh. Huber: Die Philosophie der Kirchenväter, pp. 327-341; and an article of K. Vogt, in Herzog’s Encycl. iii. p. 412 ff.}

In Synesius of Cyrene also the Platonism outweighs the Christianity. He was an enthusiastic pupil of Hypatia, the famous female philosopher at Alexandria, and in 410 was called to the bishopric of Ptolemais, the capital of Pentapolis. Before taking orders he frankly declared that he could not forsake his philosophical opinions, although he would in public accommodate himself to the popular belief. Theophilus of Alexandria, the same who was one of the chief persecutors of the admirers of Origen, the father of Christian Platonism, accepted this doubtful theory of accommodation. Synesius was made bishop, but often regretted that he exchanged his favorite studies for the responsible and onerous duties of the bishopric. In his hymns he fuses the Christian doctrine of the Trinity with the Platonic idea of God, and the Saviour with the divine Helios, whose daily setting and rising was to him a type of Christ’s descent into Hades and ascension to heaven. The desire of the soul to be freed from the chains of matter, takes the place of the sorrow for sin and the longing after salvation.\footnote{Comp. Clausen: De Synesio philosopho, Hafn. 1831; Huber: Philos. der Kirchenväter, pp. 315-321; Baur: Church Hist. from the fourth to the sixth century, p. 52 ff., and W. Möllerin Herzog’s Encycl. vol. xv. p. 335 ff.}

As soon as theology assumed a scholastic character and began to deal more in dialectic forms than in living ideas, the philosophy of Aristotle rose to favor and influence, and from John Philoponus, a.d. 550, throughout the middle age to the Protestant Reformation, kept the lead in the Catholic church. It was the philosophy of scholasticism, while mysticism sympathized rather with the Platonic system.

The influence of the two great philosophies upon theology was beneficial or injurious, according as the principle of Christianity was the governing or the governed factor. Both systems are theistic (at bottom monotheistic), and favorable to the spirit of earnest and profound speculation. Platonism, with its ideal, poetic views, stimulates, fertilizes, inspires and elevates the reason and imagination, but also easily leads into the errors of gnosticism and the twilight of mysticism. Aristotelianism, with its sober realism and sharp logical distinctions, is a good discipline for the understanding, a school of dialectic practice, and a help to logical, systematic, methodical treatment, but may also induce a barren formalism. The truth is, Christianity itself is the highest philosophy, as faith is the highest reason; and she makes successive philosophies, as well as the arts and the sciences, tributary to herself, on the Pauline principle that “all things are hers.”\footnote{Concerning the influence of philosophy on the church fathers, comp. Ritter’sGeschichte der christl. Philosophie; Ackermann, and Baur: Ueber das Christliche im Platonismus; Huber’sPhilosophie der Kirchenväter (Munich, 1859); NEANDER’S Dogmengeschichte, i p. 59 sqq.; Archer Butler’s Lectures on Ancient Philosophy; Shedd’sHistory of Christian Doctrine, vol. i. ch. 1 (Philosophical Influences in the Ancient Church); Alb. Stöckl: Geschichte der Philosophie des Mittelalters, Mainz, 1865, 2 Bde.}

\xsection{Sources of Theology. Scripture and Tradition}

§ 118. Sources of Theology. Scripture and Tradition.

Comp. the literature in vol. i. § 75 and § 76. Also: Eusebius: Hist. Eccl. iii. 3; vi. 25 (on the form of the canon in the Nicene age); Leander van Ess (R.C.): Chrysostomus oder Stimmen der Kirchenväter für’s Bibellesen. Darmstadt, 1824.

Vincentius Lirinensis († about 450): Commonitorium pro cathol. fidei antiquitate et universitate Adv. profanas omnium haer. novitates; frequent editions, e.g. by Baluzius (1663 and 1684), Gallandi, Coster, Kluepfel (with prolegom. and notes), Viennae, 1809, and by Herzog, Vratisl. 1839; also in connection with the Opera Hilarii Arelatensis, Rom. 1731, and the Opera Salviani, Par. 1669, and in Migne’s Patrolegis, vol. 50, p. 626 sqq.

The church view respecting the sources of Christian theology and the rule of faith and practice remains as it was in the previous period, except that it is further developed in particulars.\footnote{Comp. vol. i. § 75 and 76.} The divine Scriptures of the Old and New Testaments, as opposed to human writings; and the oral tradition or living faith of the catholic church from the apostles down, as opposed to the varying opinions of heretical sects together form the one infallible source and rule of faith. Both are vehicles of the same substance: the saving revelation of God in Christ; with this difference in form and office, that the church tradition determines the canon, furnishes the key to the true interpretation of the Scriptures, and guards them against heretical abuse. The relation of the two in the mind of the ancient church may be illustrated by the relation between the supreme law of a country (such as the Roman law, the Code Napoleon, the common law of England, the Constitution of the United States) and the courts which expound the law, and decide between conflicting interpretations. Athanasius, for example, “the father of orthodoxy,” always bases his conclusions upon Scripture, and appeals to the authority of tradition only in proof that he rightly understands and expounds the sacred books. The catholic faith, says he, is that which the Lord gave, the apostles preached, and the fathers have preserved; upon this the church is founded, and he who departs from this faith can no longer be called a Christian.\footnote{Ad Serap. Ep. i., cap. 28 (Opera, tom. i. pars ii. p. 676): \textgreek{Ἴδωμεν ... τὴν τῆς ἀρχῆς παράδοσιν καὶ διδασκαλίαν καὶ πίστιν τῆς καθολικῆς ἐκκλησίας ἥν ὁ μέν κύριος ἔδωκεν, οἱ δὲ ἀπόστολοι ἐκήρυξαν, καὶ οἱ πατέρες ἐφύλαξαν.} Voigt (Die Lehre des Athanasius, \&c. p. 13 ff.) makes Athanasius even the representative of the formal principle of Protestantism, the supreme authority, sufficiency, and self-interpreting character of the Scriptures; while Möhler endeavors to place him on the Roman side. Both are biassed, and violate history by their preconceptions.}

The sum of doctrinal tradition was contained in what is called the Apostles’ Creed, which at first bore various forms, but after the beginning of the fourth century assumed the Roman form now commonly used. In the Greek church its place was supplied after the year 325 by the Nicene Creed, which more fully expresses the doctrine of the deity of Christ. Neither of these symbols goes beyond the substance of the teaching of the apostles; neither contains any doctrine specifically Greek or Roman.

The old catholic doctrine of Scripture and tradition, therefore, nearly as it approaches the Roman, must not be entirely confounded with it. It makes the two identical as to substance, while the Roman church rests upon tradition for many doctrines and usages, like the doctrines of the seven sacraments, of the mass, of purgatory, of the papacy, and of the immaculate conception, which have no foundation in Scripture. Against this the evangelical church protests, and asserts the perfection and sufficiency of the Holy Scriptures as the record of divine revelation; while it does not deny the value of tradition, or of the consciousness of the church, in the interpretation of Scripture, and regulates public teaching by symbolical books. In the Protestant view tradition is not coordinate with Scripture, but subordinate to it, and its value depends on its agreement with the Scriptures. The Scriptures alone are the norma fidei; the church doctrine is only the norma doctrinae. Protestantism gives much more play to private judgment and free investigation in the interpretation of the Scriptures, than the Roman or even the Nicene church.\footnote{On this point compare the relevant sections in the works on Symbolic and Polemic Theology, and Schaff’s Principle of Protestantism, 1845.}

I. In respect to the Holy Scriptures:

At the end of the fourth century views still differed in regard to the extent of the canon, or the number of the books which should be acknowledged as divine and authoritative.

The Jewish canon, or the Hebrew Bible, was universally received, while the Apocrypha added to the Greek version of the Septuagint were only in a general way accounted as books suitable for church reading,\footnote{\textgreek{Βιβλία ἀναγινωσκόμενα} (libri ecclesiastici), in distinction from \textgreek{κανονικά}or \textgreek{κανονιζόμενα}on the one hand, and \textgreek{ἀπόκρυφα}on the other. So Athanasius.} and thus as a middle class between canonical and strictly apocryphal (pseudonymous) writings. And justly; for those books, while they have great historical value, and fill the gap between the Old Testament and the New, all originated after the cessation of prophecy, and they cannot therefore be regarded as inspired, nor are they ever cited by Christ or the apostles.\footnote{Heb. xi. 35 ff. probably alludes, indeed, to 2 Macc. vi. ff.; but between a historical allusion and a corroborative citation with the solemn \textgreek{ἡ γραφὴ λέγει}there is a wide difference.}

Of the New Testament, in the time of Eusebius, the four Gospels, the Acts, thirteen Epistles of Paul, the first Epistle of John, and the first Epistle of Peter, were universally recognized as canonical,\footnote{Hence called \textgreek{ὁμολογούμενα}.} while the Epistle to the Hebrews, the second and third Epistles of John, the second Epistle of Peter, the Epistle of James, and the Epistle of Jude were by many disputed as to their apostolic origin, and the book of Revelation was doubted by reason of its contents.\footnote{Hence called \textgreek{ἀντιλεγόμενα}, which, however, is by no means to be confounded with \textgreek{ἀπόκρυφα}and \textgreek{νόθα}. There are no apocrypha, properly speaking, in the New Testament. The apocryphal Gospels, Acts, and Apocalypses in every case differ greatly from the apostolic, and were never received into the canon. The idea of apocrypha in the Old Testament is innocent, and is applied to later Jewish writings, the origin of which is not accurately known, but the contents of which are useful and edifying.} This indecision in reference to the Old Testament Apocrypha prevailed still longer in the Eastern church; but by the middle of the fourth century the seven disputed books of the New Testament were universally acknowledged, and they are included in the lists of the canonical books given by Athanasius, Gregory Nazianzen, Amphilochius of Iconium, Cyril of Jerusalem, and Epiphanius; except that in some cases the Apocalypse is omitted.

In the Western church the canon of both Testaments was closed at the end of the fourth century through the authority of Jerome (who wavered, however, between critical doubts and the principle of tradition), and more especially of Augustine, who firmly followed the Alexandrian canon of the Septuagint, and the preponderant tradition in reference to the disputed Catholic Epistles and the Revelation; though he himself, in some places, inclines to consider the Old Testament Apocrypha as deutero-canonical books, bearing a subordinate authority. The council of Hippo in 393, and the third (according to another reckoning the sixth) council of Carthage in 397, under the influence of Augustine, who attended both, fixed the catholic canon of the Holy Scriptures, including the Apocrypha of the Old Testament, and prohibited the reading of other books in the churches, excepting the Acts of the Martyrs on their memorial days. These two African councils, with Augustine,\footnote{De doctr. Christ. l. ii. c. 8.} give forty-four books as the canonical books of the Old Testament, in the following order: Genesis, Exodus, Leviticus, Numbers, Deuteronomy, Joshua, Judges, Ruth, four books of Kings (the two of Samuel and the two of Kings), two books of Paralipomena (Chronicles), Job, the Psalms, five books of Solomon, the twelve minor Prophets, Isaiah, Jeremiah, Daniel, Ezekiel, Tobias, Judith, Esther, two books of Ezra, two books of Maccabees. The New Testament canon is the same as ours.

This decision of the transmarine church however, was subject to ratification; and the concurrence of the Roman see it received when Innocent I. and Gelasius I. (a.d. 414) repeated the same index of biblical books.

This canon remained undisturbed till the sixteenth century, and was sanctioned by the council of Trent at its fourth session.

Protestantism retained the New Testament canon of the Roman church,\footnote{The well-known doubts of Luther respecting some of the \emph{antilegomena}, especially the Epistle of James, the Epistle to the Hebrews, and the Revelation, are mere private opinions, which have latterly been reasserted by individual Lutheran divines, like Philippi and Kahnis, but have had no influence upon the church doctrine.} but, in accordance with the orthodox Jewish and the primitive Christian view, excluded the Apocrypha from the Old.\footnote{The more particular history of the canon belongs to historical and critical Introduction to the Bible. Besides the relevant sections in works of this sort, and in Lardner’s Credibility of the Gospel History, and Kirchhofer’sQuellensammlung (1844), comp. the following special treatises: Thiersch: Herstellung des historischen Standpunkts für die Kritik der N. T’tlichen Schriften, 1845; Credner. Zur Geschichte des Kanons, 1847; Oehler: Kanon des A. Ts. in Herzog’s Encyklopädie, vol. vii. pp. 243-270; Landerer: Kanon des Neuen Testaments, ibid. pp. 270-303; also an extended article: Canon of Scripture, in W. Smith’sDictionary of the Bible (London and Boston, 1860), vol. i. pp. 250-268.}

The most eminent of the church fathers speak in the strongest terms of the full inspiration and the infallible authority of the holy Scriptures, and commend the diligent reading of them even to the laity. Especially Chrysostom. The want of general education, however, and the enormous cost of books, left the people for the most part dependent on the mere hearing of the word of God in public worship; and the free private study of the Bible was repressed by the prevailing Spirit of the hierarchy. No prohibition, indeed, was yet laid upon the reading of the Bible; but the presumption that it was a book of the priests and monks already existed. It remained for a much later period, by the invention of printing, the free spirit of Protestantism, and the introduction of popular schools, to make the Bible properly a people’s book, as it was originally designed to be; and to disseminate it by Bible societies, which now print and circulate more copies of it in one year, than were made in the whole middle age, or even in the fifteen centuries before the Reformation.

The oldest manuscripts of the Bible now extant date no further back than the fourth century, are very few, and abound in unessential errors and omissions of every kind; and the problem of a critical restoration of the original text is not yet satisfactorily solved, nor can it be more than approximately solved in the absence of the original writings of the apostles.

The oldest and most important manuscripts in uncial letters are the Sinaitic (first discovered by Tischendorf in 1859, and published in 1862), the Vatican (in Rome, defective), the Alexandrian (in London); then the much mutilated codex of Ephraim Syrus in Paris, and the incomplete codex of Cambridge. From these and a few other uncial codices the oldest attainable text must be mainly gathered. Secondary sources are quotations in the fathers, the earliest versions, Stich as the Syriac Peshito and the Latin Vulgate, and the later manuscripts.\footnote{Full information on this subject may be found in the Introductions to the New Testament, and in the Prolegomena of the critical editions of the New Testament, among which the editions of Lachmann, Tischendorf, Tregelles, and Alford are the most important. Comp. particularly the eighth large edition of Tischendorf, begun in 1865, and diligently employing all existing critical helps.}

The faith which rests not upon the letter, but upon the living spirit of Christianity, is led into no error by the defects of the manuscripts and ancient and modern versions of the Bible, but only excited to new and deeper study.

The spread of the church among all the nations of the Roman empire, and even among the barbarians on its borders, brought with it the necessity of translating the Scriptures into various tongues. The most important of these versions, and the one most used, is the Latin Vulgate, which was made by the learned Jerome on the basis of the older Itala, and which afterwards, notwithstanding its many errors, was placed by the Roman church on a level with the original itself. The knowledge of Hebrew among the fathers was very rare; the Septuagint was considered sufficient, and even the knowledge of Greek diminished steadily in the Latin church after the invasion of the barbarians and the schism with the East, so that the Bible in its original languages became a sealed book, and remained such until the revival of learning in the fifteenth century.

In the interpretation of the Scriptures the system of allegorical exposition and imposition was in high repute, and often degenerated into the most arbitrary conceits, especially in the Alexandrian school, to which most of the great dogmatic theologians of the Nicene age belonged. In opposition to this system the Antiochian school, founded by Lucian († 311), and represented by Diodorus of Tarsus, Theodore of Mopsuestia, and best by John Chrysostom and Theodoret, advocated a soberer grammatical and historical exegesis, and made a sharper distinction between the human and the divine elements in the Scriptures. Theodore thereby incurred the suspicion and subsequently even the condemnation of the Greek church.

Among the Latin fathers a similar difference in the interpretation of Scripture appears between the discerning depth and lively play of Augustine and the grammatical and archaeological scholarship and dogmatical superficiality of Jerome.

II. The Holy Scriptures were universally accepted as the supreme authority and infallible rule of faith. But as the Scriptures themselves were variously interpreted, and were claimed by the heretics for their views, the fathers of our period, like Irenaeus and Tertullian before them, had recourse at the same time to Tradition, as preserved from the apostles through the unbroken succession of the bishops. With them the Scriptures are the supreme law; the combined wisdom and piety of the catholic church, the organic body of the faithful, is the judge which decides the true sense of the law. For to be understood the Bible must be explained, either by private judgment or by the universal faith of Christendom.

Strictly speaking, the Holy Ghost, who is the author, is also the only infallible interpreter of the Scriptures. But it was held that the Holy Ghost is given only to the orthodox church not to heretical and schismatic sects, and that he expresses himself through assembled orthodox bishops and universal councils in the clearest and most authoritative way. “The heretics,” says Hilary, “all cite the Scriptures, but without the sense of the Scriptures; for those who are outside the church can have no understanding of the, word of God.” They imagine they follow the Scriptures, while in truth they follow their own conceits, which they put into the Scriptures instead of drawing their thoughts from them.

Even Augustine, who of all the fathers stands nearest to evangelical Protestantism, on this point advocates the catholic principle in the celebrated maxim which he urges against the Manichaeans: “I would not believe the gospel, if I were not compelled by the authority of the universal church.” But he immediately adds: “God forbid that I should not believe the gospel.”\footnote{“Ego vero evangelio non crederem, nisi me Catholicae ecclesiae commoveret autoritas .... Sed absit ut ego Evangelio non credam. Illi enim credens, non invenio quomodo possim etiam tibi [Manichaeus] credere. Apostolorum enim nomina, quae ibi leguntur, non inter se continent nomen Manichaei.” Contra Epist. Manichaei, quam vocant fundamenti, cap. 6 (ed. Bened. tom. viii. p. 154). His object in this argument is to show that the Manichaeans have no right in the Scriptures, that the Catholic church is the legitimate owner and interpreter of the Bible. But it is an abuse to press this argument at once into the service of Rome as is so often done. Between the controversy of the old Catholic church with Manichaeism, and the controversy of Romanism with Protestantism, there is an immense difference.}

But there are different traditions; not to speak of various interpretations of the catholic tradition. Hence the need of a criterion of true and false tradition. The semi-Pelagian divine, Vincentius, a monk and priest in the South-Gallic cloister of Lirinum († 450), \footnote{Lerinum or Lirinum (now St. Honorat) is one of the group of small islands in the Mediterranean which formerly belonged to Roman Gaul, afterwards to France. In the fifth century it was a seminary of learned monks and priests for France, as Faustus Regiensis, Hilarius Arelatensis, Salvianus, and others.} otherwise little known, propounded the maxim which formed an epoch in this matter, and has since remained the standard in the Roman church: We must hold “what has been everywhere, always, and by all believed.”\footnote{Commonit. cap. 2 (in Migne’s\emph{Patrolog}. vol. 50, p. 640): “In ipsa item Catholica Ecclesia magnopere curandum est, ut id teneamus quod \emph{ubique}, quod \emph{semper}, quod \emph{ab} \emph{omnibus}creditum est.” The Commonitorium was composed, as we learn from the preface and from ch. 42, about three years after the ecumenical council of Ephesus, therefore about 434, under the false name of Peregrinus, as a help to the memory of the author that he might have the main points of ecclesiastical tradition constantly at hand against the heretics. Baronius calls it “opus certe aureum,” and Bellarmin “parvum mole et virtute maximum.” It consisted originally of two books, but the manuscript of the second book was stolen from the author, who then added a brief summary of both books at the close of the first (c. 41-43). Vossius, Cardinal Norisius (Historia Pelagiana, I. ii. c. 11), Natalis Alexander, Hefele, and Schmidt give this work a polemic aim against strict Augustinism, for which certainly the Greek church cannot be claimed, so that the three criteria of catholicity are wanting. There is pretty strong evidence in the book itself that Vincentius belonged to the semi-Pelagian school which arose in Marseille and Lirinum. He was probably also the author of the \emph{Vincentianae objectiones} against Augustine’s doctrine of predestination. Comp. on Vincentius, Tillemont’s Mémoires, tom. xv. pp. 143-147; the art. \emph{Vincentius v}. \emph{L}. by H. Schmidt in Herzog’s Encykl. vol. xvii. pp. 211-217; and an essay of C. J. Hefele (R.C.), in his Beiträge zur Kirchengeschichte, Archäologie und Liturgik, vol. i. p. 146 ff. (Tüb. 1864).} Here we have a threefold test of the ecclesiastical orthodoxy: Catholicity of place, of time, and of number; or ubiquity, antiquity, and universal consent;\footnote{As Vincentius expresses himself in the succeeding sentence: Universitas, antiquitas, consensio. Comp. c. 27.} in other words, an article of faith must be traced up to the apostles, and be found in all Christian countries, and among all believers. But this principle can be applied only to a few fundamental articles of revealed religion, not to any of the specifically Romish dogmas, and, to have any reasonable meaning, must be reduced to a mere principle of majority. In regard to the consensus omnium, which properly includes both the others, Vincentius himself makes this limitation, by defining the condition as a concurrence of the majority of the clergy.\footnote{“Consensio omnium vel certe \emph{paene} omnium sacerdotum pariter et magistrorum,” etc. Common. c. 2 (in Migne, p. 640).} To the voice of the people neither he nor the whole Roman system, in matters of faith, pays the slightest regard. In many important doctrines, however, there is not even a consensus patrum, as in the doctrine of free will, of predestination, of the atonement. A certain freedom of divergent private opinions is the indispensable condition of all progress of thought, and precedes the ecclesiastical settlement of every article of faith. Even Vincentius expressly asserts a steady advance of the church in the knowledge of the truth, though of course in harmony with the previous steps, as a man or a tree remains identical through the various stages of growth.\footnote{Cap 23 (in Migne, vol 50, p. 667 sqq.).}

Vincentius is thoroughly Catholic in the spirit and tendency of his work, and has not the most remote conception of the free Protestant study of the Scriptures. But on the other hand he would have as little toleration for new dogmas. He wished to make tradition not an independent source of knowledge and rule of faith by the side of the Holy Scriptures, but only to have it acknowledged as the true interpreter of Scripture, and as a bar to heretical abuse. The criterion of the antiquity of a doctrine, which he required, involves apostolicity, hence agreement with the spirit and substance of the New Testament. The church, says he, as the solicitous guardian of that which is intrusted to her, changes, diminishes, increases nothing. Her sole effort is to shape, or confirm, or preserve the old. Innovation is the business of heretics not of orthodox believers. The canon of Scripture is complete in itself, and more than sufficient.\footnote{Cap. 2: “Quum sit \emph{perfectus} Scripturarum Canon et sibi \emph{ad omnia satis superque} \emph{sufficiat},” etc. Cap. 29.} But since all heretics appeal to it, the authority of the church must be called in as the rule of interpretation, and in this we must follow universality, antiquity, and consent.\footnote{“Hoc facere curabant ... ut divinum canonem secundum universalis ecclesiae traditiones et juxta catholici dogmatis regulas interpretentur, in qua item catholica et apostolica ecclesia sequantur necesse est universitatem, antiquitatem, consensionem.” Commonit. cap. 27 (in Migne, vol. 50, p. 674). Comp. c. 2-4.} It is the custom of the Catholics, says he in the same work, to prove the true faith in two ways: first by the authority of the holy Scriptures, then by the tradition of the Catholic church; not because the canon alone is not of itself sufficient for all things, but on account of the many conflicting interpretations and perversions of the Scriptures.\footnote{Cap. 29 (in Migne, vol. 50, p. 677): “Non quia canon solus non sibi ad universa sufficiat, sed quia verba divina, pro suo plerique arbitratu interpretantes, varias opiniones erroresque concipiant,” etc.}

In the same spirit says pope Leo I.: “It is not permitted to depart even in one word from the doctrine of the Evangelists and the Apostles, nor to think otherwise concerning the Holy Scriptures, than the blessed apostles and our fathers learned and taught.”\footnote{Epist. 82 ad Episc. Marcianum Aug. (Opera, tom. i. p. 1044, ed. Ballerini, and in Migne, liv. p. 918): “Quum ab evangelica apostolicaque doctrina ne uno quidem verbo liceat dissidere, aut aliter de Scripturis divinis sapere quam beati apostoli et patres nostri didicerunt atque docuerunt,” etc.}

The catholic principle of tradition became more and more confirmed, as the authority of the fathers and councils increased and the learned study of the Holy Scriptures declined; and tradition gradually set itself in practice on a level with Scripture, and even above it. It fettered free investigation, and promoted a rigid, stationary and intolerant orthodoxy, which condemned men like Origen and Tertullian as heretics. But on the other hand the principle of tradition unquestionably exerted a wholesome conservative power, and saved the substance of the ancient church doctrine from the obscuring and confusing influence of the pagan barbarism which deluged Christendom.

\xsection{The Arian Controversy down to the Council of Nicaea, 318-325}

§ 119. The Arian Controversy down to the Council of Nicaea, 318–325.

I. – Trinitarian Controversies.

GENERAL LITERATURE OF THE ARIAN CONTROVERSIES.

I. Sources: On the orthodox side most of the fathers of the fourth century; especially the dogmatic and polemic works of Athanasius (Orationes c. Arianos; De decretis Nicaenae Synodi; De sententia Dionysii; Apologia c. Arianos; Apologia de fuga sua; Historia Arianorum, etc., all in tom. i. pars i. ii. of the Bened. ed.), Basil (Adv. Eunomium), Gregory Nazianzen (Orationes theologicae), Gregory Of Nyssa (Contra Eunom.), Epiphanius (Ancoratus), Hilary (De Trinitate), Ambrose (De Fide), Augustine (De Trinitate, and Contra Maximinimum Arianum), Rufinus, and the Greek church historians.

On the heretical side: The fragments of the writings of Arius (Qavleia, and two Epistolae to Eusebius of Nicomedia and Alexander of Alexandria), preserved in quotations in Athanasius, Epiphanius, Socrates, and Theodoret; comp. Fabricius: Biblioth. gr. viii. p. 309. Fragmenta Arianorum about 388 in Angelo Mai: Scriptorum veterum nova collect. Rom. 1828, vol. iii. The fragments of the Church History of the Arian Philostorgius, a.d. 350–425.

II. Works: Tillemont (R.C.): Mémoires, etc. tom. vi. pp. 239–825, ed. Paris. 1699, and ed. Ven. (the external history chiefly). Dionysius Petavius (Jesuit, † 1652): De theologicis dogmatibus, tom. ii., which treats of the divine Trinity in eight books; and in part toms. iv. and v. which treat in sixteen books of the Incarnation of the Word. This is still, though incomplete, the most learned work of the Roman church in the History of Doctrines; it first appeared at Paris, 1644–’50, in five volumes fol., then at Amsterdam, 1700 (in 6 vols.), and at Venice, 1757 (ed. Zacharia), and has been last edited by Passaglia and Schrader in Rome, 1857. J. M. Travasa (R.C.): Storia critica della vita di Ario. Ven. 1746. S. J. Maimburg: Histoire de l’Arianisme. Par. 1675. John Pearson (bishop of Chester, † 1686): An Exposition of the Creed (in the second article), 1689, 12th ed. Lond. 1741, and very often edited since by Dobson, Burton, Nichols, Chevalier, etc. George Bull (Anglican bishop of St. David’s, † 1710): Defensio fidei Nicaenae. Ox. 1685 (Opp. Lat. fol. ed. Grabe, Lond. 1703. Complete Works, ed. Burton, Oxf. 1827, and again in 1846, vol. 5th in two parts, and in English in the Anglo-Catholic Library, 1851). This classical work endeavors, with great learning, to exhibit the Nicene faith in all the ante-Nicene fathers, and so belongs more properly to the previous period. Dan. Waterland (archdeacon of Middlesex, † 1730, next to Bull the ablest Anglican defender of the Nicene faith): Vindication of Christ’s Divinity, 1719 ff., in Waterland’s Works, ed. Mildert, vols. i. ii. iii. Oxf. 1843. (Several acute and learned essays and sermons in defence of the orthodox doctrine of the Trinity against the high Arianism of Dr. Sam. Clarke and Dr. Whitby.) Chr. W. F. Walch: Vollständige Historic der Ketzereien, etc. 11 vols. Leipzig, 1762 ff. Vols. ii. and iii. (exceedingly thorough and exceedingly dry). Gibbon: History of the Decline and Fall of the Roman Empire, ch. xxi. A. Möhler (R.C.): Athanasius der Grosse u. die Kirche seiner Zeit. Mainz (1827); 2d ed. 1844 (Bk ii.-vi.). J. H. Newman (at the time the learned head of Puseyism, afterwards R.C.): The Arians of the Fourth Century. Lond. 1838; 2d ed. (unchanged), 1854. F. Chr. Baur: Die christl. Lehre von der Dreieinigkeit u. Menschwerdung in ihrer geschichtl. Entwicklung. 3 vols. Tübingen, 1841–’43. Vol. i. pp. 306–825 (to the council of Chalcedon). Comp. also Baur’s Kirchengesch. vom 4ten his 6ten Jahrh. Tüb. 1859, pp. 79–123. Js. A. Dorner: Entwicklungsgesch. der Lehre von der Person Christi. 1836, 2d ed. in 2 vols. Stuttg. 1845–’53. Vol. i. pp. 773–1080 (English transl. by W. L. Alexander and D. W. Simon, in Clark’s Foreign Theol. Library, Edinb. 1861). R. Wilberforce (at the time archdeacon of East Riding, afterwards R.C.): The Doctrine of the Incarnation of our Lord Jesus Christ. 4th ed. Lond. 1852. Bishop Kaye: Athanasius and the council of Nicaea. Lond. 1853. C. Jos. Hefele (R.C.): Conciliengeschichte. Freib. 1855 ff. Vol. i. p. 219 ff. Albert Prince de Broglie (R.C.): L’église et l’empire romain, au IV. siècle. Paris, 1856–’66, 6 vols. Vol. i. p. 331 sqq.; vol. ii. 1 sqq. W. W. Harvey: History and Theology of the Three Creeds. Lond. 1856, 2 vols. H. Voigt: Die Lehre des Athanasius von Alexandrien. Bremen, 1861. A. P. Stanley: Lectures on the History of the Eastern Church. 2d ed. 862 (reprinted in New York). Sects. ii.-vii. (more brilliant than solid). Comp. also the relevant sections in the general Church Histories of Fleury, Schröckh(vols. v. and vi.), Neander, Gieseler, and in the Doctrine Histories of Münscher-cölln, Baumgarten-Crusius, Hagenbach, Baur, Beck, Shedd.

The Arian controversy relates primarily to the deity of Christ, but in its course it touches also the deity of the Holy Ghost, and embraces therefore the whole mystery of the Holy Trinity and the incarnation of God, which is the very centre of the Christian revelation. The dogma of the Trinity came up not by itself in abstract form, but in inseparable connection with the doctrine of the deity of Christ and the Holy Ghost. If this latter doctrine is true, the Trinity follows by logical necessity, the biblical monotheism being presumed; in other words: If God is one, and if Christ and the Holy Ghost are distinct from the Father and yet participate in the divine substance, God must be triune. Though there are in the Holy Scriptures themselves few texts which directly prove the Trinity, and the name Trinity is wholly wanting in them, this doctrine is taught with all the greater force in a living form from Genesis to Revelation by the main facts of the revelation of God as Creator, Redeemer, and Sanctifier, besides being indirectly involved in the deity of Christ and the Holy Ghost.

The church always believed in this Trinity of revelation, and confessed its faith by baptism into the name of the Father, and of the Son, and of the Holy Ghost. This carried with it from the first the conviction, that this revelation of God must be grounded in a distinction immanent in the divine essence. But to bring this faith into clear and fixed knowledge, and to form the baptismal confession into doctrine, was the hard and earnest intellectual work of three centuries. In the Nicene age minds crashed against each other, and fought the decisive battles for and against the doctrines of the true deity of Christ, with which the divinity of Christianity stands or falls.

The controversies on this fundamental question agitated the Roman empire and the church of East and West for more than half a century, and gave occasion to the first two ecumenical councils of Nicaea and Constantinople. At last the orthodox doctrine triumphed, and in 381 was brought into the form in which it is to this day substantially held in all orthodox churches.

The external history of the Arian controversy, of which we first sketch the main features, falls into three stages:

1. From the outbreak of the controversy to the temporary victory of orthodoxy at the council of Nicaea; a.d. 318–325.

2. The Arian and semi-Arian reaction, and its prevalence to the death of Constantius; a.d. 325–361.

3. The final victory, and the completion of the Nicene creed; to the council of Constantinople, a.d. 381.

Arianism proceeded from the bosom of the Catholic church, was condemned as heresy at the council of Nicaea, but afterwards under various forms attained even ascendency for a time in the church, until at the second ecumenical council it was cast out forever. From that time it lost its importance as a politico-theological power, but continued as an uncatholic sect more than two hundred years among the Germanic nations, which were converted to Christianity under the Arian domination.

The roots of the Arian controversy are to be found partly in the contradictory elements of the christology of the great Origen, which reflect the crude condition of the Christian mind in the third century; partly in the antagonism between the Alexandrian and the Antiochian theology. Origen, on the one hand, attributed to Christ eternity and other divine attributes which logically lead to the orthodox doctrine of the identity of substance; so that he was vindicated even by Athanasius, the two Cappadocian Gregories, and Basil. But, on the other hand, in his zeal for the personal distinctions in the Godhead, he taught with equal clearness a separateness of essence between the Father and the Son\footnote{\textgreek{Ἑτερότης τῆς οὐσίας,} or \textgreek{τοῦ ὑποκειμένου.}De Orat. c. 15.} and the subordination of the Son, as a second or secondary God beneath the Father,\footnote{Hence be termed the Logos \textgreek{δεύτερος Θεός}, or \textgreek{Θεός}(without the article, comp. John i. 1), in distinction from the Father, who is absolute God, \textgreek{ὁ Θεός}, or \textgreek{αὐτόθεος}, Deus per se. He calls the Father also the root (\textgreek{ῥιζα}) and fountain (\textgreek{πηγή}) of the whole Godhead. Comp. vol. i. § 78. Redepenning: Origenes, ii. 304 sq., and Thomasius: Origenes, p. 118 sq.} and thus furnished a starting point for the Arian heresy. The eternal generation of the Son from the will of the Father was, with Origen, the communication of a divine but secondary substance, and this idea, in the hands of the less devout and profound Arius, who with his more rigid logic could admit no intermediate being between God and the creature, deteriorated to the notion of the primal creature.

But in general Arianism was much more akin to the spirit of the Antiochian school than to that of the Alexandrian. Arius himself traced his doctrine to Lucian of Antioch, who advocated the heretical views of Paul of Samosata on the Trinity, and was for a time excommunicated, but afterwards rose to great consideration, and died a martyr under Maximinus.

Alexander, bishop of Alexandria, made earnest of the Origenistic doctrine of the eternal generation of the Son (which was afterwards taught by Athanasius and the Nicene creed, but in a deeper sense, as denoting the generation of a person of the same substance from the substance of the Father, and not of a person of different substance from the will of the Father), and deduced from it the homo-ousia or consubstantiality of the Son with the Father.

Arius,\footnote{\textgreek{Ἄρειος}.} a presbyter of the same city after 313, who is represented as a tall, thin, learned, adroit, austere, and fascinating man, but proud, artful, restless, and disputatious, pressed and overstated the Origenistic view of the subordination, accused Alexander of Sabellianism, and taught that Christ, while he was indeed the creator of the world, was himself a creature of God, therefore not truly divine.\footnote{This, however, is manifestly contrary to Origen’s view, which made Christ an intermediate being between the uncreated Father and the creature, Contra Cels. iii. 34.}

The contest between these two views broke out about the year 318 or 320. Arius and his followers, for their denial of the true deity of Christ, were deposed and excommunicated by a council of a hundred Egyptian and Libyan bishops at Alexandria in 321. In spite of this he continued to hold religious assemblies of his numerous adherents, and when driven from Alexandria, agitated his doctrine in Palestine and Nicomedia, and diffused it in an entertaining work, half poetry, half prose: The Banquet (Θάλεια), of which a few fragments are preserved in Athanasius. Several bishops, especially Eusebius of Nicomedia and Eusebius of Caesarea, who either shared his view or at least considered it innocent, defended him. Alexander issued a number of circular letters to all the bishops against the apostates and Exukontians.\footnote{\textgreek{Οἱ ἐξ οὐκ ὄντων}. So he named the Arians, for their assertion that the Son of God was made \textgreek{ἐξ οὐκ ὄντων} out of nothing.} Bishop rose against bishop, and province against province. The controversy soon involved, through the importance of the subject and the zeal of the parties, the entire church, and transformed the whole Christian East into a theological battle-field.

Constantine, the first emperor who mingled in the religious affairs of Christendom, and who did this from a political, monarchical interest for the unity of the empire and of religion, was at first inclined to consider the contest a futile logomachy, and endeavored to reconcile the parties in diplomatic style by letters and by the personal mission of the aged bishop Hosius of Spain; but without effect. Questions of theological and religious principle are not to be adjusted, like political measures, by compromise, but must be fought through to their last results, and the truth must either conquer or (for the time) succumb. Then, in pursuance, as he thought, of a “divine inspiration,” and probably also with the advice of bishops who were in friendship with him,\footnote{At least Rufinus says, H. E. i. 1: “Ex sacerdotum sententia.” Probably Hosius and Eusebius of Caesarea had most influence with the emperor in this matter, as in others. But of any coöperation of the pope in the summoning of the council of Nicaea the earliest documents know nothing.} he summoned the first universal council, to represent the whole church of the empire, and to give a final decision upon the relation of Christ to God, and upon some minor questions of discipline, the time of Easter, and the Meletian schism in Egypt.

\xsection{The Council of Nicaea, 325}

§ 120. The Council of Nicaea, 325.

SOURCES.

(1) The twenty Canones, the doctrinal Symbol, and a Decree of the Council of Nicaea, and several Letters of bishop Alexander of Alexandria and the emperor Constantine (all collected in Greek and Latin in Mansi: Collect. sacrorum Conciliorum, tom. ii. fol. 635–704). Official minutes of the transactions themselves were not at that time made; only the decrees as adopted were set down in writing and subscribed by all (comp. Euseb. Vita Const. iii. 14). All later accounts of voluminous acts of the council are sheer fabrications (Comp. Hefele, i. p. 249 sqq.)

(2) Accounts of eye-witnesses, especially Eusebius, Vita Const. iii. 4–24 (superficial, rather Arianizing, and a panegyric of the emperor Constantine). The Church History of Eusebius, which should have closed with the council of Nice, comes down only to the year 324. Athanasius: De decretis Synodi Nic.; Orationes iv contra Arianos; Epist. ad Afros, and other historical and anti-Arian tracts in tom. i. and ii. of his Opera, ed. Bened. and the more important of them also in the first vol. of Thilo’s Bibliotheca Patrum Graec. dogmat. Lips. 1853. (Engl. transl. in the Oxford Library of the Fathers.)

(3) The later accounts of Epiphanius: Haer. 69; Socrates: H. E. i. 8 sqq.; Sozomen: H. E. i. 17 sqq.; Theodoret: H. E. i. 1–13; Rufinus: H. E. i. 1–6 (or lib. x., if his transl. of Eusebius be counted in). Gelasius Cyzicenus (about 476): Commentarius actorum Concilii Nicaeni (Greek and Latin in Mansi, tom. ii. fol. 759 sqq.; it professes to be founded on an old MS., but is filled with imaginary speeches). Comp. also the four Coptic fragments in Pitra: Spicilegium Solesmense, Par. 1852, vol. i. p. 509 sqq., and the Syriac fragments in Analecta Nicaena. Fragments relating to the Council of Nicaea. The Syriac text from am ancient MS. by H. Cowper, Lond. 1857.

LITERATURE.

Of the historians cited at § 119 must be here especially mentioned Tillemont (R.C.), Walch, Schröckh, Gibbon, Hefele (i. pp. 249–426), A. de Broglie (vol. ii. ch. iv. pp. 3–70), and Stanley. Besides them, Ittig: Historia concilii Nicaeni, Lips. 1712. Is. Boyle: A historical View of the Council of Nice, with a translation of Documents, New York, 1856 (in Crusé’s ed. of Euseb.’s Church History). Comp. also § 65 and 66 above, where this in connection with the other ecumenical councils has already been spoken of.

Nicaea, the very name of which speaks victory, was the second city of Bithynia, only twenty English miles from the imperial residence of Nicomedia, and easily accessible by sea and land from all parts of the empire. It is now a miserable Turkish village, Is-nik,\footnote{\emph{I.e}., \textgreek{Εἰς Νίκαιαν}, like Stambul, Is-tam-bul, from \textgreek{εἰς τὴν πόλιν.} Isnik now contains only some fifteen hundred inhabitants.} where nothing but a rude picture in the solitary church of St. Mary remains to the memory of the event which has given the place a name in the history of the world.

Hither, in the year 325, the twentieth of his reign (therefore the festive vicennalia), the emperor summoned the bishops of the empire by a letter of invitation, putting at their service the public conveyances, and liberally defraying from the public treasury the expenses of their residence in Nicaea and of their return. Each bishop was to bring with him two presbyters and three servants.\footnote{The imperial letter of convocation is not extant. Eusebius says, Vita Const. iii. 6, the emperor by very respectful letters invited the bishops of all countries to come with all speed to Nicaea (\textgreek{σπεύδειν ἁπανταχόθεν τοὺς ἐπισκόπους ·γράμμασι τιμητικοῖς προκαλούμενος}). Arius also was invited (Rufinus, H. E. i. 1). In an invitation of Constantineto the bishop of Syracuse to attend the council of Arles (as given by Eusebius, H. E. x. c. 5), the emperor directs him to bring with him two priests and three servants, and promises to defray the travelling expenses. The same was no doubt done at the council of Nice. Comp. Eus. V. Const. iii. 6 and 9.} They travelled partly in the public post carriages, partly on horses, mules, or asses, partly on foot. Many came to bring their private disputes before the emperor, who caused all their papers, without reading them, to be burned, and exhorted the parties to reconciliation and harmony.

The whole number of bishops assembled was at most three hundred and eighteen;\footnote{According to Athanasius (Ad Afros, c. 2, and elsewhere), Socrates (H. E. l. 8), Theodoret (H. E. i. 7), and the usual opinion. The spirit of mystic interpretation gave to the number 318, denoted in Greek by the letters TIH, a reference to the cross (T), and to the holy name Jesus (\textgreek{ἸΗσοῦς}). It was also (Ambrose, De fide, i. 18) put in connection with the three hundred and eighteen servants of Abraham, the father of the faithful (Gen. xiv. 14). Eusebius, however, gives only two hundred and fifty bishops (\textgreek{πεντήκοντα καὶ διακοσίων ἀριθμόν}), or a few over; but with an indefinite number of attendant priests, deacons, and acolyths (Vit. Const. iii. 8). The later Arabic accounts of more than two thousand bishops probably arose from confounding bishops and clergy in general. Perhaps the number of members increased towards the close, so that Eusebius with his 260, and Athanasius with his 318, may both be right. The extant Latin lists of the subscribers contain the names of no more than two hundred and twenty-four bishops and chorepiscopi, and many of these are mutilated and distorted by the mistakes of transcribers, and varied in the different copies. Comp. the list from an ancient Coptic cloister in Pitra’s Spicilegium Solesmense (Par. 1852.), tom. i. p. 516 sqq.; and Hefele, Conciliengesch. i. 284.} that is, about one sixth of all the bishops of the empire, who are estimated as at least eighteen hundred (one thousand for the Greek provinces, eight hundred for the Latin), and only half as many as were at the council of Chalcedon. Including the presbyters and deacons and other attendants the number may, have amounted to between fifteen hundred and two thousand. Most of the Eastern provinces were strongly represented; the Latin church, on the contrary, had only seven delegates: from Spain Hosius of Cordova, from France Nicasius of Dijon, from North Africa Caecilian of Carthage, from Pannonia Domnus of Strido, from Italy Eustorgius of Milan and Marcus of Calabria, from Rome the two presbyters Victor or Vitus and Vincentius as delegates of the aged pope Sylvester I. A Persian bishop John, also, and a Gothic bishop, Theophilus, the forerunner and teacher of the Gothic Bible translator Ulfilas, were present.

The formal sessions began, after preliminary disputations between Catholics, Arians, and philosophers, probably about Pentecost, or at farthest after the arrival of the emperor on the 14th of June. They closed on the 25th of July, the anniversary of the accession of Constantine; though the members did not disperse till the 25th of August.\footnote{On the various dates, comp. Hefele, l. c. i. p. 261 sqq. Broglie, ii. 26, puts the arrival of the emperor earlier, on the 4th or 5th of June.} They were held, it appears, part of the time in a church or some public building, part of the time in the emperor’s house.

The formal opening of the council was made by the stately entrance of the emperor, which Eusebius in his panegyrical flattery thus describes:\footnote{Vita Const. iii. 10. The above translation is somewhat abridged.} “After all the bishops had entered the central building of the royal palace, on the sides of which very many seats were prepared, each took his place with becoming modesty, and silently awaited the arrival of the emperor. The court officers entered one after another, though only such as professed faith in Christ. The moment the approach of the emperor was announced by a given signal, they all rose from their seats, and the emperor appeared like a heavenly messenger of God,\footnote{\textgreek{Οἷα Θεοῦ τις οὐράνιος · ἄγγελος.}} covered with gold and gems, a glorious presence, very tall and slender, full of beauty, strength, and majesty. With this external adornment he united the spiritual ornament of the fear of God, modesty, and humility, which could be seen in his downcast eyes, his blushing face, the motion of his body, and his walk. When he reached the golden throne prepared for him, he stopped, and sat not down till the bishops gave him the sign. And after him they all resumed their seats.”

How great the contrast between this position of the church and the time of her persecution but scarcely passed! What a revolution of opinion in bishops who had once feared the Roman emperor as the worst enemy of the church, and who now greeted the same emperor in his half barbarous attire as an angel of God from heaven, and gave him, though not yet even baptized, the honorary presidency of the highest assembly of the church!

After a brief salutatory address from the bishop on the right of the emperor, by which we are most probably to understand Eusebius of Caesarea, the emperor himself delivered with a gentle voice in the official Latin tongue the opening address, which was immediately after translated into Greek, and runs thus:\footnote{According to Eusebius, l. c. iii. c. 12. Sozomen, Socrates, and Rufinus also give the emperor’s speech, somewhat differently, but in substantial agreement with this.}

“It was my highest wish, my friends, that I might be permitted to enjoy your assembly. I must thank God that, in addition to all other blessings, he has shown me this highest one of all: to see you all gathered here in harmony and with one mind. May no malicious enemy rob us of this happiness, and after the tyranny of the enemy of Christ [Licinius and his army] is conquered by the help of the Redeemer, the wicked demon shall not persecute the divine law with new blasphemies. Discord in the church I consider more fearful and painful than any other war. As soon as I by the help of God had overcome my enemies, I believed that nothing more was now necessary than to give thanks to God in common joy with those whom I had liberated. But when I heard of your division, I was convinced that this matter should by no means be neglected, and in the desire to assist by my service, I have summoned you without delay. I shall, however, feel my desire fulfilled only when I see the minds of all united in that peaceful harmony which you, as the anointed of God, must preach to others. Delay not therefore, my friends, delay not, servants of God; put away all causes of strife, and loose all knots of discord by the laws of peace. Thus shall you accomplish the work most pleasing to God, and confer upon me, your fellow servant,\footnote{\textgreek{τῷ ὑμετέρῷ συνθεράποντι}.} an exceeding great joy.”

After this address he gave way to the (ecclesiastical) presidents of the council\footnote{\textgreek{Παρεδίδου τὸνλόγον τοῖς τῆς συνόδου προέδροις}, says Euseb. iii. 13. The question of the presidency in the ecumenical councils has already been spoken of in § 65.} and the business began. The emperor, however, constantly, took an active part, and exercised a considerable influence.

Among the fathers of the council, besides a great number of obscure mediocrities, there were several distinguished and venerable men. Eusebius of Caesarea was most eminent for learning; the young archdeacon Athanasius, who accompanied the bishop Alexander of Alexandria, for zeal, intellect, and eloquence. Some, as confessors, still bore in their body the marks of Christ from the times of persecution: Paphnutius of the Upper Thebaid, Potamon of Heraklea, whose right eye had been put out, and Paul of Neo-Caesarea, who had been tortured with red hot iron under Licinius, and crippled in both his hands. Others were distinguished for extraordinary ascetic holiness, and even for miraculous works; like Jacob of Nisibis, who had spent years as a hermit in forests and eaves, and lived like a wild beast on roots and leaves, and Spyridion (or St. Spiro) of Cyprus, the patron of the Ionian isles, who even after his ordination remained a simple shepherd. Of the Eastern bishops, Eusebius of Caesarea, and of the Western, Hosius, or Osius, of Cordova,\footnote{Athanasius always calls him the Great, \textgreek{ὁ μέγας}.} had the greatest influence with the emperor. These two probably sat by his side, and presided in the deliberations alternately with the bishops of Alexandria and Antioch.

In reference to the theological question the council was divided in the beginning into three parties.\footnote{The ancient and the Roman Catholic historians (and A. de Broglie, l.c. vol. ii. p. 21) generally assume only two parties, an orthodox majority and a heretical minority. But the position of Eusebius of Caesarea, the character of his confession, and the subsequent history of the controversy, prove the existence of a middle, semi-Arian party. Athanasius, too, who usually puts all shades of opponents together, accuses Eusebius of Caesarea and others repeatedly of insincerity in their subscription of the Nicene creed, and yet these were not proper Arians, but semi-Arians.}

The orthodox party, which held firmly to the deity of Christ, was at first in the minority, but in talent and influence the more weighty. At the head of it stood the bishop (or “pope”) Alexander of Alexandria, Eustathius of Antioch, Macarius of Jerusalem, Marcellus of Ancyra, Rosins of Cordova (the court bishop), and above all the Alexandrian archdeacon, Athanasius, who, though small and young, and, according to later practice not admissible to a voice or a seat in a council, evinced more zeal and insight than all, and gave promise already of being the future head of the orthodox party.

The Arians or Eusebians numbered perhaps twenty bishops, under the lead of the influential bishop Eusebius of Nicemedia (afterwards of Constantinople), who was allied with the imperial family, and of the presbyter Arius, who attended at the command of the emperor, and was often called upon to set forth his views.\footnote{Rufinus, i. 5: “Evocabatur frequenter Arius in concilium.”} To these also belonged Theognis of Nicaea, Maris of Chalcedon, and Menophantus of Ephesus; embracing in this remarkable way the bishops of the several seats of the orthodox ecumenical councils.

The majority, whose organ was the renowned historian Eusebius of Caesarea, took middle ground between the right and the left, but bore nearer the right, and finally went over to that side. Many of them had an orthodox instinct, but little discernment; others were disciples of Origen, or preferred simple biblical expression to a scholastic terminology; others had no firm convictions, but only uncertain opinions, and were therefore easily swayed by the arguments of the stronger party or by mere external considerations.

The Arians first proposed a creed, which however was rejected with tumultuous disapproval, and torn to pieces; whereupon all the eighteen signers of it, excepting Theonas and Secundus, both of Egypt, abandoned the cause of Arius.

Then the church historian Eusebius, in the name of the middle party, proposed an ancient Palestinian Confession, which was very similar to the Nicene, and acknowledged the divine nature of Christ in general biblical terms, but avoided the term in question, ὁμοούσιοςconsubstantialis, of the same essence. The emperor had already seen and approved this confession, and even the Arian minority were ready to accept it.

But this last circumstance itself was very suspicious to the extreme right. They wished a creed which no Arian could honestly subscribe, and especially insisted on inserting the expression homo-ousios, which the Arians hated and declared to be unscriptural, Sabellian, and materialistic.\footnote{Athanasius himself, however, laid little stress on the term, and rarely used it in his theological expositions; he cared more for the thing than the name. The word \textgreek{ὁμοούσιος}, from \textgreek{ὁμος}and \textgreek{οὐσία} was not an invention of the council of Nice, still less of Constantine, but had previously arisen in theological language, and occurs even in Origen and among the Gnostics, though of course it is no more to be found in the Bible than the word \emph{trinity}.} The emperor saw clearly that the Eusebian formula would not pass; and, as he had at heart, for the sake of peace, the most nearly unanimous decision which was possible, he gave his voice for the disputed word.

Then Hosius of Cordova appeared and announced that a confession was prepared which would now be read by the deacon (afterwards bishop) Hermogenes of Caesarea, the secretary of the synod. It is in substance the well-known Nicene creed with some additions and omissions of which we are to speak below. It is somewhat abrupt; the council not caring to do more than meet the immediate exigency. The direct concern was only to establish the doctrine of the true deity of the Son. The deity of the Holy Spirit, though inevitably involved, did not then come up as a subject of special discussion, and therefore the synod contented itself on this point with the sentence: “And (we believe) in the Holy Ghost.”\footnote{Dr. Shedd, therefore, is plainly incorrect in saying, Hist. of Chr. Doctrine, vol. i. p. 308: “The problem to be solved by the Nicene council was to exhibit the doctrine of the trinity in its \emph{completeness}; to bring into the creed statement the \emph{total} data of Scripture upon both the side of unity and trinity.” This was not done till the council of Constantinople in 381, and strictly not till the still later Symbolum Athanasianum.} The council of Constantinople enlarged the last article concerning the Holy Ghost. To the positive part of the Nicene confession is added a condemnation of the Arian heresy, which dropped out of the formula afterwards received.

Almost all the bishops subscribed the creed, Hosius at the head, and next him the two Roman presbyters in the name of their bishop. This is the first instance of such signing of a document in the Christian church. Eusebius of Caesarea also signed his name after a day’s deliberation, and vindicated this act in a letter to his diocese. Eusebius of Nicomedia and Theognis of Nicaea subscribed the creed without the condemnatory formula, and for this they were deposed and for a time banished, but finally consented to all the decrees of the council. The Arian historian Philostorgius, who however deserves little credit,\footnote{Even Gibbon (ch. xxi.) places very little dependence on this historian: “The credibility of Philostorgius is lessened, in the eyes of the orthodox, by his Arianism; and in those of rational critics [as if the orthodox were necessarily irrational and uncritical!] by his passion, his prejudice, and his ignorance.”} accuses them of insincerity in having substituted, by the advice of the emperor, for ὁμο-ούσιος(of the same essence) the semi-Arian word ὁμοι-ούσιος(of like essence). Only two Egyptian bishops, Theonas and Secundus, persistently refused to sign, and were banished with Arius to Illyria. The books of Arius were burned and his followers branded as enemies of Christianity.\footnote{Jerome(Adv. Lucifer, c. 20; Opera, ed. Vallars. tom. ii. p. 192 sqq.) asserts, on the authority of aged witnesses then still living, that Arius and his adherents were pardoned even before the close of the council. Socrates also says (H. E. i. c. 14) that Arius was recalled from banishment before Eusebius and Theognis, but under prohibition of return to Alexandria. This isolated statement, however, cannot well be harmonized with the subsequent recalling, and probably arose from some confusion.}

This is the first example of the civil punishment of heresy; and it is the beginning of a long succession of civil persecutions for all departures from the Catholic faith. Before the union of church and state ecclesiastical excommunication was the extreme penalty. Now banishment and afterwards even death were added, because all offences against the church were regarded as at the same time crimes against the state and civil society.

The two other points on which the council of Nicaea decided, the Easter question and the Meletian schism, have been already spoken of in their place. The council issued twenty canons in reference to discipline. The creed and the canons were written in a book, and again signed by the bishops. The council issued a letter to the Egyptian and Libyan bishops as to the decision of the three main points; the emperor also sent several edicts to the churches, in which he ascribed the decrees to divine inspiration, and set them forth as laws of the realm. On the twenty-ninth of July, the twentieth anniversary of his accession, he gave the members of the council a splendid banquet in his palace, which Eusebius (quite too susceptible to worldly splendor) describes as a figure of the reign of Christ on earth; he remunerated the bishops lavishly, and dismissed them with a suitable valedictory, and with letters of commendation to the authorities of all the provinces on their homeward way.

Thus ended the council of Nicaea. It is the first and most venerable of the ecumenical synods, and next to the apostolic council at Jerusalem the most important and the most illustrious of all the councils of Christendom. Athanasius calls it “a true monument and token of victory against every heresy;” Leo the Great, like Constantine, attributes its decrees to the inspiration of the Holy Ghost, and ascribes even to its canons perpetual validity; the Greek church annually observes (on the Sunday before Pentecost) a special feast in memory of it. There afterwards arose a multitude of apocryphal orations and legends in glorification of it, of which Gelasius of Cyzicus in the fifth century collected a whole volume.\footnote{Stanly (sic) interweaves several of these miraculous legends with graphical minuteness into the text of his narrative, thus giving it the interest of romance, at the expense of the dignity of historical statement. The simple Spyridion performed, on his journey to the Council, the amazing feat of restoring in the dark his two mules to life by annexing the white head to the chestnut mule, and the chestnut head to its white companion, and overtook the rival bishops who had cut off the heads of the mules with the intention to prevent the rustic bishop from reaching Nicaea and hurting the cause of orthodoxy by his ignorance! According to another version of this silly legend the decapitation of the mules is ascribed to malicious Arians.}

The council of Nicaea is the most important event of the fourth century, and its bloodless intellectual victory over a dangerous error is of far greater consequence to the progress of true civilization, than all the bloody victories of Constantine and his successors. It forms an epoch in the history of doctrine, summing up the results of all previous discussions on the deity of Christ and the incarnation, and at the same time regulating the further development of the Catholic orthodoxy for centuries. The Nicene creed, in the enlarged form which it received after the second ecumenical council, is the only one of all, the symbols of doctrine which, with the exception of the subsequently added filioque, is acknowledged alike by the Greek, the Latin, and the Evangelical churches, and to this, day, after a course of fifteen centuries, is prayed and sung from Sunday to Sunday in all countries of the civilized world. The Apostles’ Creed indeed, is much more generally used in the West, and by its greater simplicity and more popular form is much better adapted to catechetical and liturgical purposes; but it has taken no root in the Eastern church; still less the Athanasian Creed, which exceeds the Nicene in logical precision and completeness. Upon the bed of lava grows the sweet fruit of the vine. The wild passions and the weaknesses of men, which encompassed the Nicene council, are extinguished, but the faith in the eternal deity of Christ has remained, and so long as this faith lives, the council of Nicaea will be named with reverence and with gratitude.

\xsection{The Arian and Semi-Arian Reaction, a.d. 325-361}

§ 121. The Arian and Semi-Arian Reaction, a.d. 325–361.

The victory of the council of Nicaea over the views of the majority of the bishops was a victory only in appearance. It had, to be sure, erected a mighty fortress, in which the defenders of the essential deity of Christ might ever take refuge from the assaults of heresy; and in this view it was of the utmost importance, and secured the final triumph of the truth. But some of the bishops had subscribed the homoousion with reluctance, or from regard to the emperor, or at best with the reservation of a broad interpretation; and with a change of circumstances they would readily turn in opposition. The controversy now for the first time fairly broke loose, and Arianism entered the stage of its political development and power. An intermediate period of great excitement ensued, during which council was held against council, creed was set forth against creed, and anathema against anathema was hurled. The pagan Ammianus Marcellinus says of the councils under Constantius: “The highways were covered with galloping bishops;” and even Athanasius rebuked the restless flutter of the clergy, who journeyed the empire over to find the true faith, and provoked the ridicule and contempt of the unbelieving world. In intolerance and violence the Arians exceeded the orthodox, and contested elections of bishops not rarely came to bloody encounters. The interference of imperial politics only poured oil on the flame, and embarrassed the natural course of the theological development.

The personal history of Athanasius was interwoven with the doctrinal controversy; he threw himself wholly into the cause which he advocated. The question whether his deposition was legitimate or not, was almost identical with the question whether the Nicene Creed should prevail.

Eusebius of Nicomedia and Theognis of Nicaea threw all their influence against the adherents of the homoousion. Constantine himself was turned by Eusebius of Caesarea, who stood between Athanasius and Arius, by his sister Constantia and her father confessor, and by a vague confession of Arius, to think more favorably of Arius, and to recall him from exile. Nevertheless he afterwards, as before, thought himself in accordance with the orthodox view and the Nicene creed. The real gist of the controversy he had never understood. Athanasius, who after the death of Alexander in April, 328,\footnote{According to the Syriac preface to the Syriac Festival Letters of Athanasius, first edited by Cureton in 1848. It was previously supposed that Alexander died two years earlier. Comp. Hefele, i. p. 429.} became bishop of Alexandria and head of the Nicene party, refused to reinstate the heretic in his former position, and was condemned and deposed for false accusations by two Arian councils, one at Tyre under the presidency of the historian Eusebius, the other at Constantinople in the year 335 (or 336), and banished by the emperor to Treves in Gaul in 336, as a disturber of the peace of the church.

Soon after this Arius, having been formally acquitted of the charge of heresy by a council at Jerusalem (a.d. 335), was to have been solemnly received back into the fellowship of the church at Constantinople. But on the evening before the intended procession from the imperial palace to the church of the Apostles, he suddenly died (a.d. 336), at the age of over eighty years, of an attack like cholera, while attending to a call of nature. This death was regarded by many as a divine judgment; by others, it was attributed to poisoning by enemies; by others, to the excessive joy of Arius in his triumph.\footnote{Comp. Athanasius, De morte Arii Epist. ad Serapionem (Opera, tom. i. p. 340). He got his information from his priest Macarius, who was in Constantinople at the time.}

On the death of Constantine (337), who had shortly before received baptism from the Arian Eusebius of Nicomedia, Athanasius was recalled from his banishment (338) by Constantine II. († 340), and received by the people with great enthusiasm; “more joyously than ever an emperor.”\footnote{So says Gregory Nazianzen. The date of his return, according to the Festival Letters of Athanasius, was the 23d November, 338.} Some months afterwards (339) he held a council of nearly a hundred bishops in Alexandria for the vindication of the Nicene doctrine. But this was a temporary triumph.

In the East Arianism prevailed. Constantius, second son of Constantine the Great, and ruler in the East, together with his whole court, was attached to it with fanatical intolerance. Eusebius of Nicomedia was made bishop of Constantinople (338), and was the leader of the Arian and the more moderate, but less consistent semi-Arian parties in their common opposition to Athanasius and the orthodox West. Hence the name Eusebians.\footnote{\textgreek{Οἱ περὶ Εὐσέβιον}.} Athanasius was for a second time deposed, and took refuge with the bishop Julius of Rome (339 or 340), who in the autumn of 341 held a council of more than fifty bishops in defence of the exile and for the condemnation of his opponents. The whole Western church was in general more steadfast on the side of the Nicene orthodoxy, and honored in Athanasius a martyr of the true faith. On the contrary a synod at Antioch, held under the direction of the Eusebians on the occasion of the dedication of a church in 341,\footnote{Hence called the council \emph{in} \emph{encoeniis} (\textgreek{ἐγκαινίοις}) or \emph{in dedicatione}.} issued twenty-five canons, indeed, which were generally accepted as orthodox and valid, but at the same time confirmed the deposition of Athanasius, and set forth four creeds, which rejected Arianism, yet avoided the orthodox formula, particularly the vexed homoousion.\footnote{This apparent contradiction between orthodox canons and semi-Arian confessions has occasioned all kinds of hypotheses in reference to this Antiochian synod. Comp. on them, Hefele, i. p. 486 sqq.}

Thus the East and the West were in manifest conflict.

To heal this division, the two emperors, Constantius in the East and Constans in the West, summoned a general council at Sardica in Illyria, a.d. 343.\footnote{Not a.d.347, as formerly supposed. Comp. Hefele, i. 515 sqq.} Here the Nicene party and the Roman influence prevailed.\footnote{About a hundred and seventy bishops in all (according to Athanasius) were present at Sardica, ninety-four occidentals and seventy-six orientals or Eusebians. Sozomen and Socrates, on the contrary, estimate the number at three hundred. The signatures of the acts of the council are lost, excepting a defective list of fifty-nine names of bishops in Hilary.} Pope Julius was represented by two Italian priests. The Spanish bishop Hosius presided. The Nicene doctrine was here confirmed, and twelve canons were at the same time adopted, some of which are very important in reference to discipline and the authority of the Roman see. But the Arianizing Oriental bishops, dissatisfied with the admission of Athanasius, took no part in the proceedings, held an opposition council in the neighboring city of Philippopolis, and confirmed the decrees of the council of Antioch. The opposite councils, therefore, inflamed the discord of the church, instead of allaying it.

Constantius was compelled, indeed, by his brother to restore Athanasius to his office in 346; but after the death of Constans, a.d. 350, be summoned three successive synods in favor of a moderate Arianism; one at Sirmium in Pannonia (351), one at Arelate or Arles in Gaul (353), and one at Milan in Italy, (355); he forced the decrees of these councils on the Western church, deposed and banished bishops, like Liberius of Rome, Hosius of Cordova, Hilary of Poictiers, Lucifer of Calaris, who resisted them, and drove Athanasius from the cathedral of Alexandria during divine service with five thousand armed soldiers, and supplied his place with an uneducated and avaricious Arian, George of Cappadocia (356). In these violent measures the court bishops and Eusebia, the last wife of Constantius and a zealous Arian, had great influence. Even in their exile the faithful adherents of the Nicene faith were subjected to all manner of abuse and vexation. Hence Constantius was vehemently attacked by Athanasius, Hilary, and Lucifer, compared to Pharaoh, Saul, Ahab, Belshazzar, and called an inhuman beast, the forerunner of Antichrist, and even Antichrist himself.

Thus Arianism gained the ascendency in the whole Roman empire; though not in its original rigorous form, but in the milder form of homoi-ousianism or the doctrine of similarity of essence, as opposed on the one hand to the Nicene homo-ousianism (sameness of essence), and on the other hand to the Arian hetero-ousianism (difference of essence).

Even the papal chair was desecrated by heresy during this Arian interregnum; after the deposition of Liberius, the deacon Felix II., “by antichristian wickedness,” as Athanasius expresses it, was elected his successor.\footnote{Comp. above, § 72, \textbf{p. 371}.} Many Roman historians for this reason regard him as a mere anti-pope. But in the Roman church books this Felix is inserted, not only as a legitimate pope, but even as a saint, because, according to a much later legend, he was executed by Constantius, whom he called a heretic. His memory is celebrated on the twenty-ninth of July. His subsequent fortunes are very differently related. The Roman people desired the recall of Liberius, and he, weary of exile, was prevailed upon to apostatize by subscribing an Arian or at least Arianizing confession, and maintaining church fellowship with the Eusebians.\footnote{The apostasy of Liberius comes to us upon the clear testimony of the most orthodox fathers, Athanasius, Hilary, Jerome, Sozomen, \&c., and of three letters of Liberius himself, which Hilaryadmitted into his sixth fragment, and accompanied with some remarks. Jeromesays in his Chronicle: “Liberius, taedio victus exilii, in haereticam pravitatem subscribens Romam quasi victor intravit.” Comp. his Catal. script. eccl c. 97. He probably subscribed what is called the third Sirmian formula, that is, the collection of Semi-Arian decrees adopted at the third council of Sirmium in 358. Hefele (i. 673), from his Roman point of view, knows no way of saving him but by the hypothesis that he renounced the Nicene \emph{word} (\textgreek{ὁμοούσιος}), but not the Nicene \emph{faith}. But this, in the case of so current a party term as \textgreek{ὁμοούσιος}, which Liberius himself afterwards declared “the bulwark against all Arian heresy” (Socr. H. E. iv. 12), is entirely untenable.} On this condition he was restored to his papal dignity, and received with enthusiasm into Rome (358). He died in 366 in the orthodox faith, which he had denied through weakness, but not from conviction.

Even the almost centennarian bishop Hosius was induced by long imprisonment and the threats of the emperor, though not himself to compose (as Hilary states), yet to subscribe (as Athanasius and Sozomen say), the Arian formula of the second council of Sirmium, a.d. 357, but soon after repented his unfaithfulness, and condemned the Arian heresy shortly before his death.

The Nicene orthodoxy was thus apparently put down. But now the heretical majority, having overcome their common enemy, made ready their own dissolution by divisions among themselves. They separated into two factions. The right wing, the Eusebians or Semi-Arians, who were represented by Basil of Ancyra and Gregory of Laodicea, maintained that the Son was not indeed of the same essence (ὁμο-ούσιος), yet of like essence (ὁμοι-ούσιος), with the Father. To these belonged many who at heart agreed with the Nicene faith, but either harbored prejudices against Athanasius, or saw in the term ὁμο-ούσιος an approach to Sabellianism; for theological science had not yet duly fixed the distinction of substance (οὐσία) and person (ὑπόστασις), so that the homoousia might easily be confounded with unity of person. The left wing, or the decided Arians, under the lead of Eudoxius of Antioch, his deacon Aëtius,\footnote{He was hated among the orthodox and Semi-Arians, and called \textgreek{ἄθεος}. He was an accomplished dialectician, a physician and theological author in Antioch, and died about 370 in Constantinople.} and especially the bishop Eunomius of Cyzicus in Mysia\footnote{He was a pupil and friend of Aëtius, and popularized his doctrine. He died in 392. Concerning him, comp. Klose, Geschichte u. Lehre des Eunomius, Kiel, 1833, and Dorner, l.c. vol. i. p. 853 sqq. cDorner calls him a deacon; but through the mediation of the bishop Eudoxius of Constantinople (formerly of Antioch) he received the bishopric of Cyzicus or Cyzicum as early as 360, before he became the head of the Arian party. Theodoret, H. E. l. ii. c. 29.} (after whom they were called also Eunomians), taught that the Son was of a different essence (ἑτεροούσιος), and even unlike the Father (ἀνόμοιος), and created out of nothing (ἐκ οὐκ ὄντων). They received also, from their standard terms, the names of Heterousiasts, Anomaeans, and Exukontians.

A number of councils were occupied with this internal dissension of the anti-Nicene party: two at Sirmium (the second, a.d. 357; the third, a.d. 358), one at Antioch (358), one at Ancyra (358), the double council at Seleucia and Rimini (359), and one at Constantinople (360). But the division was not healed. The proposed compromise of entirely avoiding the word ούσια, and substituting ὅμοιος like, for ὁμοιούσιος of like essence, and ἀνόμοιος, unlike, satisfied neither party. Constantius vainly endeavored to suppress the quarrel by his imperio-episcopal power. His death in 361 opened the way for the second and permanent victory of the Nicene orthodoxy.

\xsection{The Final Victory of Orthodoxy, and the Council of Constantinople, 381}

§ 122. The Final Victory of Orthodoxy, and the Council of Constantinople, 381.

Julian the Apostate tolerated all Christian parties, in the hope that they would destroy one another. With this view he recalled the orthodox bishops from exile. Even Athanasius returned, but was soon banished again as an “enemy of the gods,” and recalled by Jovian. Now for a time the strife of the Christians among themselves was silenced in their common warfare against paganism revived. The Arian controversy took its own natural course. The truth regained free play, and the Nicene spirit was permitted to assert its intrinsic power. It gradually achieved the victory; first in the Latin church, which held several orthodox synods in Rome, Milan, and Gaul; then in Egypt and the East, through the wise and energetic administration of Athanasius, and through the eloquence and the writings of the three great Cappadocian bishops Basil, Gregory of Nazianzum, and Gregory of Nyssa.

After the death of Athanasius in 373, Arianism regained dominion for a time in Alexandria, and practised all kinds of violence upon the orthodox.

In Constantinople Gregory Nazianzen labored, from 379, with great success in a small congregation, which alone remained true to the orthodox faith during the Arian rule; and he delivered in a domestic chapel, which he significantly named Anastasia (the church of the Resurrection), those renowned discourses on the deity of Christ which won him the title of the Divine, and with it many persecutions.

The raging fanaticism of the Arian emperor Valens (364–378) against both Semi-Arians and Athanasians wrought an approach of the former party to the latter. His successor, Gratian, was orthodox, and recalled the banished bishops.

Thus the heretical party was already in reality intellectually and morally broken, when the emperor Theodosius I., or the Great, a Spaniard by birth, and educated in the Nicene faith, ascended the throne, and in his long and powerful reign (379–395) externally completed the triumph of orthodoxy in the Roman empire. Soon after his accession he issued, in 380, the celebrated edict, in which he required all his subjects to confess the orthodox faith, and threatened the heretics with punishment. After his entrance into Constantinople he raised Gregory Nazianzen to the patriarchal chair in place of Demophilus (who honestly refused to renounce his heretical conviction), and drove the Arians, after their forty years’ reign, out of all the churches of the capital.

To give these forcible measures the sanction of law, and to restore unity in the church of the whole empire, Theodosius called the second ecumenical council at Constantinople in May, 381. This council, after the exit of the thirty-six Semi-Arian Macedonians or Pneumatomachi, consisted of only a hundred and fifty bishops. The Latin church was not represented at all.\footnote{In the earliest Latin translation of the canons of this council, indeed, three Roman legates, Paschasinus, Lucentius, and Bonifacius, are recorded among the signers (in \emph{Mansi}, t. vi. p. 1176), but from an evident confusion of this council with the fourth ecumenical of 451, which these delegates attended. Comp. Hefele, ii. p. 3 and 393. The assertion of Baronius that in reality pope Damasus summoned the council, rests likewise on a mistake of the first council of Constantinople for the second in 382.} Meletius (who died soon after the opening), Gregory Nazianzen, and after his resignation Nectarius of Constantinople, successively presided. This preferment of the patriarch of Constantinople before the patriarch of Alexandria is explained by the third canon of the council, which assigns to the bishop of new Rome the first rank after the bishop of old Rome. The emperor attended the opening of the sessions, and showed the bishops all honor.

At this council no new symbol was framed, but the Nicene Creed, with some unessential changes and an important addition respecting the deity of the Holy Ghost against Macedonianism or Pneumatoinachism, was adopted.\footnote{This modification and enlargement of the Nicene Creed seems not to have originated with the second ecumenical council, but to have been current in substance about ten years earlier. For Epiphanius, in his Ancoratus, which was composed in 374, gives two similar creeds, which were then already in use in the East; the shorter one literally agrees with that of Constantinople (c. 119, ed. Migne, tom. iii. p. 231); the longer one (c. 120) is more lengthy on the Holy Ghost; both have the anathema. Hefele, ii. 10, overlooks the shorter and more important form.} In this improved form the Nicene Creed has been received, though in the Greek church without the later Latin addition: filioque.

In the seven genuine canons of this council the heresies of the Eunomians or Anomoeans, of the Arians or Eudoxians, of the Semi-Arians or Pneumatomachi, of the Sabellians, Marcellians, and Apollinarians, were condemned, and questions of discipline adjusted.

The emperor ratified the decrees of the council, and as early as July, 381, enacted the law that all churches should be given up to bishops who believed in the equal divinity of the Father, the Son, and the Holy Ghost, and who stood in church fellowship with certain designated orthodox bishops. The public worship of heretics was forbidden.

Thus Arianism and the kindred errors were forever destroyed in the Roman empire, though kindred opinions continually reappear as isolated cases and in other connections.\footnote{John Milton and Isaac Newton cannot properly be termed Arians. Their view of the relation of the Son to the Father was akin to that of Arius, but their spirit and their system of ideas were totally different. Bishop Bull’sgreat work, Defensio fidei Nicaenae, first published 1685, was directed against Socinian and Arian views which obtained in England, but purely with historical arguments drawn from the ante-Nicene fathers. Shortly afterwards the high Arian view was revived and ably defended with exegetical, patristic, and philosophical arguments by Whiston, Whitby, and especially by Dr. Samuel Clarke(died 1729), in his treatise on the “Scripture Doctrine of the Trinity” (1712), which gave rise to a protracted controversy, and to the strongest dialectical defence (though broken and irregular in method) of the Nicene doctrine in the English language by Dr. Waterland. This trinitarian controversy, one of the ablest and most important in the history of English theology, is very briefly and superficially touched in the great works of Dr. Baur (vol. iii. p. 685 ff.) and Dorner (vol. ii. p. 903 ff.); but the defect has been supplied by Prof. Patrick Fairbairnin an Appendix to the English translation of Dorner’s History of Christology, Divis. Secd. vol. iii. p. 337 ff.}

But among the different barbarian peoples of the West, especially in Gaul and Spain, who had received Christianity from the Roman empire during the ascendency of Arianism, this doctrine was perpetuated two centuries longer: among the Goths till 587; among the Suevi in Spain till 560; among the Vandals who conquered North Africa in 429 and cruelly persecuted the Catholics, till their expulsion by Belisarius in 530; among the Burgundians till their incorporation in the Frank empire in 534, and among the Longobards till the close of the sixth century. These barbarians, however, held Arianism rather through accident than from conviction, and scarcely knew the difference between it and the orthodox doctrine. Alaric, the first conqueror of Rome; Genseric, the conqueror of North Africa; Theodoric the Great, king of Italy and hero of the Niebelungen Lied, were Arians. The first Teutonic translation of the Bible came from the Arian missionary Ulfilas.

\xsection{The Theological Principles involved: Import of the Controversy}

§ 123. The Theological Principles involved: Import of the Controversy.

Here should be compared, of the works before mentioned, especially Petavius (tom. sec. De sanctissima Trinitate), and Möhler (Athanasius, third book), of the Romanists, and Baur, Dorner, and Voigt, of the Protestants.

We pass now to the internal history of the Arian conflict, the development of the antagonistic ideas; first marking some general points of view from which the subject must be conceived.

To the superficial and rationalistic eye this great struggle seems a metaphysical subtilty and a fruitless logomachy, revolving about a Greek iota. But it enters into the heart of Christianity, and must necessarily affect in a greater or less degree all other articles of faith. The different views of the contending parties concerning the relation of Christ to the Father involved the general question, whether Christianity is truly divine, the highest revelation, and an actual redemption, or merely a relative truth, which may be superseded by a more perfect revelation.

Thus the controversy is conceived even by Dr. Baur, who is characterized by a much deeper discernment of the philosophical and historical import of the conflicts in the history of Christian doctrine, than all other rationalistic historians. “The main question,” he says, “was, whether Christianity is the highest and absolute revelation of God, and such that by it in the Son of God the self-existent absolute being of God joins itself to man, and so communicates itself that man through the Son becomes truly one with God, and comes into such community of essence with God, as makes him absolutely certain of pardon and salvation. From this point of view Athanasius apprehended the gist of the controversy, always finally summing up all his objections to the Arian doctrine with the chief argument, that the whole substance of Christianity, all reality of redemption, everything which makes Christianity the perfect salvation, would be utterly null and meaningless, if he who is supposed to unite man with God in real unity of being, were not himself absolute God, or of one substance with the absolute God, but only a creature among creatures. The infinite chasm which separates creature from Creator, remains unfilled; there is nothing really mediatory between God and man, if between the two there be nothing more than some created and finite thing, or such a mediator and redeemer as the Arians conceive the Son of God in his essential distinction from God: not begotten from the essence of God and coeternal, but created out of nothing and arising in time. Just as the distinctive character of the Athanasian doctrine lies in its effort to conceive the relation of the Father and Son, and in it the relation of God and man, as unity and community of essence, the Arian doctrine on the contrary has the opposite aim of a separation by which, first Father and Son, and then God and man, are placed in the abstract opposition of infinite and finite. While, therefore, according to Athanasius, Christianity is the religion of the unity of God and man, according to Arius the essence of the Christian revelation can consist only in man’s becoming conscious of the difference which separates him, with all the finite, from the absolute being of God. What value, however, one must ask, has such a Christianity, when, instead of bringing man nearer to God, it only fixes the chasm between God and man?”\footnote{Die christliche Kirche vom 4-6ten Jahrhundert, 1859, p. 97 sq.}

Arianism was a religious political war against the spirit of the Christian revelation by the spirit of the world, which, after having persecuted the church three hundred years from without, sought under the Christian name to reduce her by degrading Christ to the category of the temporal and the created, and Christianity to the level of natural religion. It substituted for a truly divine Redeemer, a created demigod, an elevated Hercules. Arianism proceeded from human reason, Athanasianism from divine revelation; and each used the other source of knowledge as a subordinate and tributary factor. The former was deistic and rationalistic, the latter theistic and supernaturalistic, in spirit and effect. The one made reasonableness, the other agreement with Scripture, the criterion of truth. In the one the intellectual interest, in the other the moral and religious, was the motive principle. Yet Athanasius was at the same time a much deeper and abler thinker than Arius, who dealt in barren deductions of reason and dialectic formulas.\footnote{Baur, Newman (The Arians, p. 17), and others put Arianism into connection with the Aristotelian philosophy, Athanasianism with the Platonic; while Petavius, Ritter, to some extent also Voigt (I. c. p. 194), and others exactly reverse the relation, and derive the Arian idea of God from Platonism and Neo-Platonism. This contrariety of opinion itself proves that such a comparison is rather confusing than helpful. The empirical, rational, logical tendency of Arianism is, to be sure, more Aristotelian than Platonic; and so far Baur is right. But the Aristotelian logic and dialectics may be used equally well in the service of Catholic orthodoxy, as they were in fact in the mediaeval scholasticism; while, on the other hand, the Platonic idealism, which was to Justin, Origen, and Augustine, a bridge to faith, may lead into all kinds of Gnostic and mystic error. All depends on making revelation and faith, or philosophy and reason, the starting-point and the ruling power of the theological system. Comp. also the observations of Dr. Dorner against Dr. Baur, in his Entwicklungsgesch. der Christologie, vol. i. p. 859, note.}

In close connection with this stood another distinction. Arianism associated itself with the secular political power and the court party; it represented the imperio-papal principle, and the time of its prevalence under Constantius was an uninterrupted season of the most arbitrary and violent encroachments of the state upon the rights of the church. Athanasius, on the contrary, who was so often deposed by the emperor, and who uttered himself so boldly respecting Constantius, is the personal representative not only of orthodoxy, but also of the independence of the church with reference to the secular power, and in this respect a precursor of Gregory VII. in his contest with the German imperialism.

While Arianism bent to the changing politics of the court party, and fell into diverse schools and sects the moment it lost the imperial support, the Nicene faith, like its great champion Athanasius, remained under all outward changes of fortune true to itself, and made its mighty advance only by legitimate growth outward from within. Athanasius makes no distinction at all between the various shades of Arians and Semi-Arians, but throws them all into the same category of enemies of the catholic faith.\footnote{I cannot refrain from quoting the striking judgment of George Bancroft, once a Unitarian preacher, on the import of the Arian controversy and the vast influence of the Athanasian doctrine on the onward march of true Christian civilization. “In vain,” says he in his address on the \emph{Progress of the} \emph{Human Race}, delivered before the New York Historical Society in 1854, p. 25 f., “did restless pride, as that of Arius, seek to paganize Christianity and make it the ally of imperial despotism; to prefer a belief resting on authority and unsupported by an inward witness, over the clear revelation of which the millions might see and feel and know the divine glory; to substitute the conception, framed after the pattern of heathenism, of an agent, superhuman yet finite, for faith in the ever continuing presence of God with man; to wrong the greatness and sanctity of the Spirit of God by representing it as a birth of time. Against these attempts to subordinate the enfranchising virtue of truth to false worship and to arbitrary power reason asserted its supremacy, and the party of superstition was driven from the field. Then mooned Ashtaroth was eclipsed and Osiris was seen no more in Memphian grove; then might have been heard the crash of the falling temples of Polytheism; and instead of them, came that harmony which holds Heaven and Earth in happiest union. Amid the deep sorrows of humanity during the sad conflict which was protracted through centuries for the overthrow of the past and the reconstruction of society, the consciousness of an incarnate God carried peace into the bosom of mankind. That faith emancipated the slave, broke the bondage of woman, redeemed the captive, elevated the low, lifted up the oppressed, consoled the wretched, inspired alike the heroes of thought and the countless masses. The down-trodden nations clung to it as to the certainty of their future emancipation; and it so filled the heart of the greatest poet of the Middle Ages—perhaps the greatest poet of all time—that he had no prayer so earnest as to behold in the profound and clear substance of the eternal light, that circling of reflected glory which showed the image of man.”}

\xsection{Arianism}

§ 124. Arianism.

The doctrine of the Arians, or Eusebians, Aëtians, Eunomians, as they were called after their later leaders, or Exukontians, Heteroousiasts, and Anomoeans, as they were named from their characteristic terms, is in substance as follows:

The Father alone is God; therefore he alone is unbegotten, eternal, wise, good, and unchangeable, and he is separated by an infinite chasm from the world. He cannot create the world directly, but only through an agent, the Logos. The Son of God is pre-existent,\footnote{\textgreek{Πρὸ χρόνων καὶ αἰώνων}.} before all creatures, and above all creatures, a middle being between God and the world, the creator of the world, the perfect image of the Father, and the executor of his thoughts, and thus capable of being called in a metaphorical sense God, and Logos, and Wisdom.\footnote{\textgreek{Θεός, λόγος , σοφία}.} But on the other hand, he himself is a creature, that is to say, the first creation of God, through whom the Father called other creatures into existence; he was created out of nothing\footnote{\textgreek{Ποίημα, κτίσμα ἐξ οὐκ ὄντων.} Hence the name Exukontians} (not out of the essence of God) by the will of the Father before all conceivable time; he is therefore not eternal, but had a beginning, and there was a time when he was not.\footnote{\textgreek{Ἀρχὴν ἔχει} —\textgreek{οὐκ ἦν πρὶν γεννηθῇ, ἤτοι κτισθῇ} —\textgreek{ἦν ποτε ὅτε οὐκ ἦν.}}

Arianism thus rises far above Ebionism, Socinianism, deism, and rationalism, in maintaining the personal pre-existence of the Son before all worlds, which were his creation; but it agrees with those systems in lowering the Son to the sphere of the created, which of course includes the idea of temporalness and finiteness. It at first ascribed to him the predicate of unchangeableness also,\footnote{\textgreek{Ἀναλλοίωτος, ἄτρεπτος ὁ υἱός.}} but afterwards subjected him to the vicissitudes of created being.\footnote{\textgreek{Τρεπτὸς φύσει ὡς τὰ κρίσματα}.} This contradiction, however, is solved, if need be, by the distinction between moral and physical unchangeableness; the Son is in his nature (φύσει) changeable, but remains good (καλός) by a free act of his will. Arius, after having once robbed the Son of divine essence,\footnote{\textgreek{οὐσία}} could not consistently allow him any divine attribute in the strict sense of the word; he limited his duration, his power, and his knowledge, and expressly asserted that the Son does not perfectly know the Father, and therefore cannot perfectly reveal him. The Son is essentially distinct from the Father,\footnote{\textgreek{Ἑτεροούσιος τῷ πατρί.}} and—as Aëtius and Eunomius afterward more strongly expressed it—unlike the Father;\footnote{\textgreek{Ἀνόμοιος κατὰ οὐσίαν}. Hence the name \textgreek{Ἀνόμοιοι}, Anomoeans.} and this dissimilarity was by some extended to all moral and metaphysical attributes and conditions.\footnote{\textgreek{Ἀνόμοιος κατὰ πάντᾳ.}.} The dogma of the essential deity of Christ seemed to Arius to lead of necessity to Sabellianism or to the Gnostic dreams of emanation. As to the humanity of Christ, Arius ascribed to him only a human body, but not a rational soul, and on this point Apollinarius came to the same conclusion, though from orthodox premises, and with the intention of saving the unity of the divine personality of Christ.

The later development of Arianism brought out nothing really new, but rather revealed many inconsistencies and contradictions. Thus, for example, Eunomius, to whom clearness was the measure of truth, maintained that revelation has made everything clear, and man can perfectly know God; while Arius denied even to the Son the perfect knowledge of God or of himself. The negative and rationalistic element came forth in ever greater prominence, and the controversy became a metaphysical war, destitute of all deep religion, spirit. The eighteen formulas of faith which Arianism and Semi-Arianism produced between the councils of Nice and Constantinople, are leaves without blossoms, and branches without fruit. The natural course of the Arian heresy is downward, through the stage of Socinianism, into the rationalism which sees in Christ a mere man, the chief of his kind.

To pass now to the arguments used for and against this error:

1. The Arians drew their exegetical proofs from the passages of Scripture which seem to place Christ in any way in the category of that which is created,\footnote{Such as Prov. viii. 22-25 (Comp. Sir. i. 4; xxiv. 8f.), where personified Wisdom, \emph{i.e}., the Logos, says (according to the Septuagint): \textgreek{Κύριος ἔκτισέν με} [Heb. \texthebrew{קָנָנִי }\par Vulg. possedit me] \textgreek{ἀρχὴν ὁδῶν αὐτοῦ εἰς ἔργα αὐτοῦ· πρὸ τοῦ αἰῶνος ἐθεμελίωσέν με, κ.τ.λ.}This passage seemed clearly to prove the two propositions of Arius, that the Father created the Son, and that he created him for the purpose of creating the world through him (\textgreek{εἰς ἔργα αὐτοῦ}). Acts ii. 36: \textgreek{Ὅτι καὶ κύριον αὐτὸν καὶ Χριστὸν ἐποίησεν ὁ θεός.}Heb. i. 4: \textgreek{Κρείττων γενόμενος τῶν ἀγγέλων}. Heb. iii. 2: \textgreek{Πιστὸν ὄντα τῷ ποιήσαντι αὐτόν.} John i 14: \textgreek{Ὁ λόγος σάρξ ἐγένετο.} Phil. ii. 7-9. The last two passages are of course wholly inapposite, as they treat of the incarnation of the Son of God, not of his pre-temporal existence and essence. Heb. i. 4 refers to the exaltation of the God-Man. Most plausible of all is the famous passage: \textgreek{πρωτότοκος πάσης κτίσεως}, Col. i. 15, from which the Arians inferred that Christ himself is a \textgreek{κτίσις} of God, to wit, the first creature of all. But \textgreek{πρωτότοκος}is not equivalent to \textgreek{πρωτόκτιστος}or \textgreek{πρωτόπλαστος}: on the contrary, Christ is by this very term distinguished from the creation, and described as the Author, Upholder, and End of the creation. A creature cannot possibly be the source of life for all creatures. The meaning of the expression, therefore, is: born before every creature, \emph{i.e}., before anything was made. The text indicates the distinction between the eternal generation of the Son from the essence of the Father, and the temporal creation of the world out of nothing by the Son. Yet there is a difference between \textgreek{μονογενής}and \textgreek{πρωτότοκος ,} which Athanasius himself makes: the former referring to the relation of the Son to the Father, the latter, to his relation to the world.} or ascribe to the incarnate (not the pre-temporal, divine) Logos growth, lack of knowledge, weariness, sorrow, and other changing human affections and states of mind,\footnote{Such as Luke ii. 52; Heb. v. 8, 9; John xii. 27, 28; Matt. xxvi. 39; Mark xii. 52; \&c.} or teach a subordination of the Son to the Father.\footnote{\emph{E.g}., John xiv. 28: \textgreek{Ὁ πατήρ μείζων μού ἐστιν.}. This passage also refers not to the pre-existent state of Christ, but to the state of humiliation of the God-Man.}

Athanasius disposes of these arguments somewhat too easily, by referring the passages exclusively to the human side of the person of Jesus. When, for example, the Lord says he knows not the day, nor the hour of the judgment, this is due only to his human nature. For how should the Lord of heaven and earth, who made days and hours, not know them! He accuses the Arians of the Jewish conceit, that divine and human are incompatible. The Jews say How could Christ, if he were God, become man, and die on the cross? The Arians say: How can Christ, who was man, be at the same time God? We, says Athanasius, are Christians; we do not stone Christ when he asserts his eternal Godhead, nor are we offended in him when he speaks to us in the language of human poverty. But it is the peculiar doctrine of Holy Scripture to declare everywhere a double thing of Christ: that he, as Logos and image of the Father, was ever truly divine, and that he afterwards became man for our salvation. When Athanasius cannot refer such terms as “made,” “created,” “became,” to the human nature he takes them figuratively for “testified,” “constituted,” “demonstrated.”\footnote{The \textgreek{ἔκτισε}and \textgreek{ἐθεμελίωσε} in Prov. viii. 22 ff., on which the Arians laid special stress, and of which Athanasius treats quite at large in his second oration against the Arians, he refers not to the essence of the Logos (with whom the \textgreek{σοφία}was by both parties identified), but to the incarnation of the Logos and to the renovation of our race through him: appealing to Eph. ii. 10: “We are his workmanship, created in Christ Jesus unto good works.” As to the far more important passage in Col. i. 15, Athanasius gives substantially the correct interpretation in his Expositio fidei, cap. 3 (ed. Bened. tom. i. 101), where he says: \textgreek{πρωτότοκον εἰπὼν} [\textgreek{Παῦλος}] \textgreek{δηλοῖ μὴ εἶναι αὐτὸν κτίσμα, ἀλλὰ γέννημα τοῦ πατρός · ξένον γάρ ἐπὶ τῆς θεότητος αὐτοῦ τὸ λέγεσθαι κτίσμα. Τὰ γὰρ πάντα ἐκτίσθησαν ὑπὸ τοῦ πατρὸς διὰ τοῦ υἱοῦ, ὁ δὲ υἱὸς μόνος ἐκ τοῦ πατρὸς ἀϊδίως ἐγεννήθη· διὸ πρωτότοκός ἐστι πάσης κτίσεως ὁ Θεὸς λόγος ·, ἄτρεπτος ἐξ ἀτρέπτου.}}

As positive exegetical proofs against Arianism, Athanasius cites almost all the familiar proof-texts which ascribe to Christ divine names, divine attributes, divine works, and divine dignity, and which it is unnecessary here to mention in detail.

Of course his exegesis, as well as that of the fathers in general, when viewed from the level of the modern grammatical, historical, and critical method, contains a great deal of allegorizing caprice and fancy and sophistical subtilty. But it is in general far more profound and true than the heretical.

2. The theological arguments for Arianism were predominantly negative and rationalizing. The amount of them is, that the opposite view is unreasonable, is irreconcilable with strict monotheism and the dignity of God, and leads to Sabellian or Gnostic errors. It is true, Marcellus of Ancyra, one of the most zealous advocates of the Nicene homoousianism, fell into the Sabellian denial of the tri-personality,\footnote{Comp. on Marcellus of Ancyra below, § 126.} but most of the Nicene fathers steered with unerring tact between the Scylla of Sabellianism, and the Charybdis of Tritheism.

Athanasius met the theological objections of the Arians with overwhelming dialectical skill, and exposed the internal contradictions and philosophical absurdities of their positions. Arianism teaches two gods, an uncreated and a created, a supreme and a secondary god, and thus far relapses into heathen polytheism. It holds Christ to be a mere creature, and yet the creator of the world; as if a creature could be the source of life, the origin and the end of all creatures! It ascribes to Christ a pre-mundane existence, but denies him eternity, while yet time belongs to the idea of the world, and is created only therewith,\footnote{Mundus non factus est \emph{in} tempore, sed \emph{cum} tempore, says Augustine, although I cannot just now lay my hand on the passage. Time is the successional form of existence of all created things. Now Arius might indeed have said: Time arose with the Son as the first creature. This, however, he did not say, but put a time before the Son.} so that before the world there was nothing but eternity. It supposes a time before the creation of the pre-existent Christ; thus involving God himself in the notion of time; which contradicts the absolute being of God. It asserts the unchangeableness of God, but denies, with the eternal generation of the Son, also the eternal Fatherhood; thus assuming after all a very essential change in God.\footnote{Of less weight is the objection, which was raised by Alexander of Alexandria: Since the Son is the Logos, the Arian God must have been, until the creation of the Son, \textgreek{ἄλογος}, a being without reason.} Athanasius charges the Arians with dualism and heathenism, and he accuses them of destroying the whole doctrine of salvation. For if the Son is a creature, man remains still separated, as before, from God; no creature can redeem other creatures, and unite them with God. If Christ is not divine, much less can we be partakers of the divine nature and children of God.\footnote{Comp. the second Oration against the Arians, cap. 69 ff.}

\xsection{Semi-Arianism}

§ 125. Semi-Arianism.

The Semi-Arians,\footnote{\textgreek{Ἡμιάρειοι}.} or, as they are called, the Homoiousiasts,\footnote{\textgreek{Ὁμοιουσιαστοί}. The name \emph{Eusebians} is used of the Arians and Semi-Arians, who both for a time made common cause, as a political party under the lead of Eusebius of Nicomedia (not of Caesarea), against the Athanasians and Nicenes.} wavered in theory and conduct between the Nicene orthodoxy and the Arian heresy. Their doctrine makes the impression, not of an internal reconciliation of opposites which in fact were irreconcilable, but of diplomatic evasion, temporizing compromise, flat, half and half juste milieu. They had a strong footing in the subordination of most of the ante-Nicene fathers; but now the time for clear and definite decision had come.

Their doctrine is contained in the confession which was proposed to the council of Nicaea by Eusebius of Caesarea, but rejected, and in the symbols of the councils of Antioch and Sirmium from 340 to 360. Theologically they were best represented first by Eusebius of Caesarea, who adhered more closely to his admired Origen, and later by Cyril of Jerusalem, who approached nearer the orthodoxy of the Nicene party.

The signal term of Semi-Arianism is homoi-ousion, in distinction from homo-ousion and hetero-ousion. The system teaches that Christ if; not a creature, but co-eternal with the Father, though not of the same, but only of like essence, and subordinate to him. It agrees with the Nicene creed in asserting the eternal generation of the Son, and in denying that he was a created being; while, with Arianism, it denies the identity of essence. Hence it satisfied neither of the opposite parties, and was charged by both with logical incoherence. Athanasius and his friends held, against the Semi-Arians, that like attributes and relations might be spoken of, but not like essences or substances; these are either identical or different. It may be said of one man that he is like another, not in respect of substance, but in respect of his exterior and form. If the Son, as the Semi-Arians ad-mit, is of the essence of the Father, he must be also of the same essence. The Arians argued: There is no middle being between created and uncreated being; if God the Father alone is uncreated, everything out of him, including the Son, is created, and consequently of different essence, and unlike him.

Thus pressed from both sides, Semi-Arianism could not long withstand; and even before the council of Constantinople it passed over, in the main, to the camp of orthodoxy.\footnote{Bull judges Semi-Arianisn very contemptuously.“Semi-AriAnus,” says he (l. iv. 4, 8, vol. v. pars ii p. 779), “et semi-Deus, et semi-creatura perinde monstra et portenta sunt quae sani et pii omnes merito exhorrent. Filius Dei aut verus omnino Deus, aut mera creatura statuatur necesse est; aeternae veritatis axioma est inter Deum et creaturam, inter non factum et factum, medium esse nihil.” Quite similarly Waterland: A Defence of some Queries relating to Dr. Clarke’s Scheme of the Holy Trinity, Works, vol. i p. 404.}

\xsection{Revived Sabellianism. Marcellus and Photinus}

§ 126. Revived Sabellianism. Marcellus and Photinus.

I. Eusebius Caesar.: Two books contra Marcellum (κατὰ Μαρκέλλου), and three books De ecclesiastica theologia (after his Demonstratio evang.). Hilary: Fragmenta, 1–3. Basil the Great: Epist. 52. Epiphanius: Haeres. 72. Retberg: Marcelliana. Gött. 1794 (a collection of the Fragments of Marcellus).

II. Montfaucon: Diatribe de causa Marcelli Ancyr. (in Collect. nova Patr. tom. ii. Par. 1707). Klose: Geschichte u. Lehre des Marcellus u. Photinus. Hamb. 1837. Möhler: Athanasius der Gr. Buch iv. p. 318 sqq. (aiming to vindicate Marcellus, as Neander also does). Baur: l.c. vol. i. pp. 525–558. Dorner: l.c. i. pp. 864–882. (Both against the orthodoxy of Marcellus.) Hefele: Conciliengesch. i. 456 sq. et passim. Willenborg: Ueber die Orthodoxie des Marc. Münster, 1859

Before we pass to the exhibition of the orthodox doctrine, we must notice a trinitarian error which arose in the course of the controversy from an excess of zeal against the Arian subordination, and forms the opposite extreme.

Marcellus, bishop of Ancyra in Galatia, a friend of Athanasius, and one of the leaders of the Nicene party, in a large controversial work written soon after the council of Nicaea against Arianism and Semi-Arianism, so pushed the doctrine of the consubstantiality of Christ that he impaired the personal distinction of Father and Son, and, at least in phraseology, fell into a refined form of Sabellianism.\footnote{In his work \textgreek{περὶ ὑποταγῆς}, De subjectione Domini Christi, founded on 1 Cor. xv. 28.} To save the full divinity of Christ and his equality with the Father, he denied his hypostatical pre-existence. As to the orthodoxy of Marcellus, however, the East and the West were divided, and the diversity continues even among modem scholars. A Semi-Arian council in Constantinople, a.d. 335, deposed him, and intrusted Eusebius of Caesarea with the refutation of his work; while, on the contrary, pope Julius of Rome and the orthodox council of Sardica (343), blinded by his equivocal declarations, his former services, and his close connection with Athanasius, protected his orthodoxy and restored him to his bishopric. The counter-synod of Philippopolis, however, confirmed the condemnation. Finally even Athanasius, who elsewhere always speaks of him with great respect, is said to have declared against him.\footnote{Hilary, Fragm. ii. n. 21 (p. 1299, ed. Bened.), states that Athanasius as early as 349 renounced church fellowship with Marcellus.} The council of Constantinople, a.d. 381, declared even the baptism of the Marcellians and Photinians invalid.\footnote{These are meant by the \textgreek{οἱ ἀπὸ τῆς Γαλατῶν χώρας ἐρχόμενοι}in the 7th canon of the second ecumenical council. Marcellus and Photinus were both of Ancyra in Galatia. Comp. Hefele, Conciliengeschichte, vol. ii. p. 26.}

Marcellus wished to hold fast the true deity of Christ without falling under the charge of subordinatianism. He granted the Arians right in their assertion that the Nicene doctrine of the eternal generation of the Son involves the subordination of the Son, and is incompatible with his own eternity. For this reason he entirely gave up this doctrine, and referred the expressions: Son, image, firstborn, begotten, not to the eternal metaphysical relation, but to the incarnation. He thus made a rigid separation between Logos and Son, and this is the πρῶτον ψεῦδοςof this system. Before the incarnation there was, he taught, no Son of God, but only a Logos, and by that he understood,—at least so he is represented by Eusebius,—an impersonal power, a reason inherent in God, inseparable from him, eternal, unbegotten, after the analogy of reason in man. This Logos was silent (therefore without word) in God before the creation of the world, but then went forth out of God as the creative word and power, the δραστικὴ ἐνέργεια πράξεως of God (not as a hypostasis). This power is the principle of creation, and culminates in the incarnation, but after finishing the work of redemption returns again into the repose of God. The Son, after completing the work of redemption, resigns his kingdom to the Father, and rests again in God as in the beginning. The sonship, therefore, is only a temporary state, which begins with the human advent of Christ, and is at last promoted or glorified into Godhead. Marcellus reaches not a real God-Man, but only an extraordinary dynamical indwelling of the divine power in the man Jesus. In this respect the charge of Samosatenism, which the council of Constantinople in 335 brought against him, has a certain justice, though he started from premises entirely different from those of Paul of Samosata.\footnote{Dorner (l.c. 880 sq.) asserts of Marcellus, that his Sabellianism ran out to a sort of Ebionitism.} His doctrine of the Holy Spirit and of the Trinity is to a corresponding degree unsatisfactory. He speaks, indeed, of an extension of the indivisible divine monad into a triad, but in the Sabellian sense, and denies the three hypostases or persons.

Photinus, first a deacon at Ancyra, then bishop of Sirmium in Pannonia, went still further than his preceptor Marcellus. He likewise started with a strict distinction between the notion of Logos and Son,\footnote{He called God \textgreek{λογοπάτηρ}, because, in his view, God is both Father and Logos. Sabellius had used the expression \textgreek{υἱοπάτηρ}, to deny the personal distinction between the Father and the Son. Photinus had to say instead of this, \textgreek{λογοπάτηρ}, because, in his view, the \textgreek{λόγος}, not the \textgreek{υἱός}, is eternally in God.} rejected the idea of eternal generation, and made the divine in Christ an impersonal power of God. But while Marcellus, from the Sabellian point of view, identified the Son with the Logos as to essence, and transferred to him the divine predicates attaching to the Logos, Photinus, on the contrary, quite like Paul of Samosata, made Jesus rise on the basis of his human nature, by a course of moral improvement and moral merit, to the divine dignity, so that the divine in him is a thing of growth.

Hence Photinus was condemned as a heretic by several councils in the East and in the West, beginning with the Semi-Arian council at Antioch in 344. He died in exile in 366.\footnote{Comp. on Photinus, Athanas., De syn. 26; Epiph., Haer. 71; Hilary, De trinit. vii. 3-7, etc.; Baur, l.c. vol. i. p. 542 sqq.; Dorner, l.c. i. p. 881 sq.; and Hefele, l.c, i. p. 610 sqq.}

\xsection{The Nicene Doctrine of the Consubstantiality of the Son with the Father}

§ 127. The Nicene Doctrine of the Consubstantiality of the Son with the Father.

Comp. the literature in §§ 119 and 120, especially the four Orations of Athanasius against the Arians, and the other anti-Arian tracts of this “father of orthodoxy.”

The Nicene, Homo-ousian, or Athanasian doctrine was most clearly and powerfully represented in the East by Athanasius, in whom it became flesh and blood;\footnote{Particularly distinguished are his four Orations against the Arians, written in 356.} and next to him, by Alexander of Alexandria, Marcellus of Ancyra (who however strayed into Sabellianism), Basil, and the two Gregories of Cappadocia; and in the West by Ambrose and Hilary.

The central point of the Nicene doctrine in the contest with Arianism is the identity of essence or the consubstantiality of the Son with the Father, and is expressed in this article of the (original) Nicene Creed: ”[We believe] in one Lord Jesus Christ, the Son of God; who is begotten the only-begotten of the Father; that is, of the essence of the Father, God of God, and Light of Light, very God of very God, begotten, not made, being of one substance with the Father.”\footnote{\textgreek{Καὶ εἰς ἕνα Κύριον Ἰησοῦν Χριστὸν, τὸν υἱὸν τοῦ Θεοῦ· γεννηθέντα ἐκ τοῦ Πατρὸς μονογενῆ· τοῦτ  ἐστιν ἐκ τῆς οὐσίας τοῦ Πατρὸς, Θεὸν ἐκ Θεοῦ καὶ φῶς ἐκ φωτὸς, Θεὸν ἀληθινον ἐκ Θεοῦ ἀληθινοῦ· γεννηθέντα, οὐ ποιηθέντα, ὁμοούσιον τῷ Πατρί, κ.τ.λ.}}

The term ὁμοούσιος, consubstantial, is of course no more a biblical term,\footnote{Though John’s \textgreek{Θεὸς ἧν ὁ λόγος} (John i. 1), and Paul’s \textgreek{τὸ εἷναι ἴσα Θεῷ} (Phil. ii. 6), are akin to it. The latter passage, indeed, since i1sa is adverbial, denotes rather divine existence, than divine being or essence, which would be more correctly expressed by \textgreek{τὸ εἷναι ἴσον Θεῷ}?, or by \textgreek{ἰσόθεος}. But the latter would be equally in harmony with Paul’s theology. The Jews used the masc. \textgreek{ἴσος}, though in a polemical sense, when they drew from the way in which he called himself preeminently and exclusively the Son of God the logical inference, that he made himself equal with God, John v. 18: \textgreek{Ὅτι ... πατέρα ἴδιον ἔλεγε τὸν Θεὸν, ἴσον ἑαυτὸν ποιῶν τῷ Θεῷ}. The Vulgate translates: \emph{aequalem} se faciens Deo.} than trinity;\footnote{The word \textgreek{τριὰς}and trinitas, in this application to the Godhead, appears first in Theophilus of Antioch and Athenagoras in the second century, and in Tertullianin the third. Confessions of faith must be drawn up in language different from the Scriptures—else they mean nothing or everything—since they are an \emph{interpretation} of the Scriptures and intended to exclude false doctrines.} but it had already been used, though in a different sense, both by heathen writers\footnote{Bull, Def. fidei Nic., Works, vol. v. P. i. p. 70: ” \textgreek{Ὁμοούσιον} a probatis Graecis scriptoribus id dicitur, quod ejusdem cum altero substantiae, essentiae, sive naturae est.” He then cites some passages from profane writers. Thus Porphyry says, De abstinentia ab esu animalium, lib. i. n. 19: \textgreek{Εἴγε ὁμοούσιοι οἱ τῶν ζώων ψύχαὶ ἡμετέραις}, \emph{i.e}., siquidem animae animalium sunt ejusdem cum nostris essentiae. Aristotle (in a quotation in Origen) speaks of the consubstantiality of all stars, \textgreek{ὁμοούσια πάντα ἄστρα}, omnia astra sunt ejusdem essentiae sive naturae.} and by heretics,\footnote{First by the Gnostic Valentine, in Irenaeus, Adv. haer. l. i. cap. 1, § 1 and § 5(ed. Stieren, vol. i. 67 and 66). In the last passage it is said of man that he is \textgreek{ὑλικός}, and as such very like God, indeed, but not consubstantial, \textgreek{παραπλήσιον μὲν, ἀλλ  οὐχ ὁμοούσιον τῷ Θεῷ}. The Manichaeans called the human soul, in the sense of their emanation system, \textgreek{ὁμοούσιον τῷ. Θεῷ}. Agapius, in Photius (Bibl. Cod. 179), calls even the sun and the moon, in a pantheistic sense, \textgreek{ὁμοούσια Θεῷ}. The Sabellinas used the word of the trinity, but in opposition to the distinction of persons.} as well as by orthodox fathers.\footnote{Origen deduces from the figurative description \textgreek{ἀπαύγασμα}, Heb. i. 3, the \textgreek{ὁμοούσιον} of the Son. His disciples rejected the term, indeed, at the council at Antioch in 264, because the heretical Paul of Samosata gave it a perverted meaning, taking oujsiva for the common source from which the three divine persons first derived their being. But towards the end of the third century the word was introduced again into church use by Theognostus and Dionysius of Alexandria, as Athanasius, De Decr. Syn. Nic. c. 25 (ed. Bened. i. p. 230), demonstrates. Eusebius, Ep. ad Caesarienses c. 7 (in Socr. H. E. i. 8, and in Athan. Opera i. 241), says that some early bishops and authors, learned and celebrated (\textgreek{τῶν παλαιῶν τινὰς λογίους καὶ ἐπιφανεῖς ἐπιστόπους καὶ συγγραφεῖς} ), used \textgreek{ὁμοούσιον} of the Godhead of the Father and Son. Tertullian(Adv. Prax.) applied the corresponding Latin phrase \emph{unius} \emph{substantiae} to the persons of the holy Trinity.} It formed a bulwark against Arians and Semi-Arians, and an anchor which moored the church during the stormy time between the first and the second ecumenical councils.\footnote{Cunningham (Hist. Theology, i. p. 291) says of \textgreek{ὁμοούσιος}: ” The number of these individuals who held the substance of the Nicene doctrine, but objected to the phraseology in which it was expressed, was very small [?]—and the evil thereof, was very inconsiderable; while the advantage was invaluable that resulted from the possession and the use of a definite phraseology, which shut out all supporters of error, combined nearly all the maintainers of truth, and formed a rallying-point around which the whole orthodox church ultimately gathered, after the confusion and distinction occasioned by Arian cunning and Arian persecution had passed away.”} At first it had a negative meaning against heresy; denying, as Athanasius repeatedly says, that the Son is in any sense created or produced and changeable.\footnote{Athanas. Epist. de Decretis Syn. Nicaenae, cap. 20 (i. p. 226); c. 26 (p. 231); and elsewhere.} But afterwards the homoousion became a positive testword of orthodoxy, designating, in the sense of the Nicene council, clearly and unequivocally, the veritable and essential deity of Christ, in opposition to all sorts of apparent or half divinity, or mere similarity to God. The same divine, eternal, unchangeable essence, which is in an original way in the Father, is, from eternity, in a derived way, through generation, in the Son; just as the water of the fountain is in the stream, or the light of the sun is in the ray, and cannot be separated from it. Hence the Lord says: “I am in the Father, and the Father in Me; He that hath seen Me hath seen the Father; I and My Father are one.” This is the sense of the expression: “God of God,” “very God of very God.” Christ, in His divine nature, is as fully consubstantial with the Father, as, in His human nature, He is with man; flesh of our flesh, and bone of our bone; and yet, with all this, He is an independent person with respect to the Father, as He is with respect to other men. In this view Basil turns the term ὁμοούσιοςagainst the Sabellian denial of the personal distinctions in the Trinity, since it is not the same thing that is consubstantial with itself, but one thing that is consubstantial with another.\footnote{Basil. M. Epist. lii. 3 (tom. iii. 146): \textgreek{Αὔτη δὲ ἡ φωνὴ καὶ τὸ τοῦ Σαβελλίου κακὸν ἐπανορθοῦται· ἀναιρεῖ γὰρ τὴν ταυτότητα τῆς ὑποστάσεως καὶ εἰσάγει ταλείαν τῶν προσώπων τὴν ἔννοιαν}: (tollit enim hypostaseos identitatem perfectamque personarum notionem inducit) \textgreek{οὐ γὰρ αὐτὸ τί ἐστιν ἑαυτῷ ὁμοούσιον, ἀλλ  ἕτερον ἑτέρῳ} (non enim idem sibi ipsi consubstantiale est, sed alterun alteri).} Consubstantiality among men, indeed, is predicated of different individuals who partake of the same nature, and the term in this view might denote also unity of species in a tritheistic sense.

But in the case before us the personal distinction of the Son from the Father must not be pressed to a duality of substances of the same kind; the homoousion, on the contrary, must be understood as identity or numerical unity of substance, in distinction from mere generic unity. Otherwise it leads manifestly into dualism or tritheism. The Nicene doctrine refuses to swerve from the monotheistic basis, and stands between Sabellianism and tritheism; though it must be admitted that the usage of οὐσίαand ὑπόστασις;still wavered for a time, and the relation of the consubstantiality to the numerical unity of the divine essence did not come clearly out till a later day. Athanasius insists that the unity of the divine essence is indivisible, and that there is only one principle of Godhead.\footnote{Orat. iv. contra Arianos, c. 1 (tom. i. p. 617): \textgreek{ὝΩστε δύο μὲν εἷναι πατέρα καὶ υἱὸν, μονάδα δὲ θεότητος ἀδιαίρετον καὶ ἄσχιστον ... μία ἀρχὴ θεότητος καὶ οὐ δύο ἀρχαί, ὅθεν κυρίως καὶ μοναρχία ἐστίν.}} He frequently illustrates the relation) as Tertullian had done before him, by the relation between fire and brightness,\footnote{\emph{E.g}., Orat. iv. c. Arianos, c. 10 (p. 624): \textgreek{Ἔστω δὲ παράδειγμα ἀνθρώπινον τὸ πῦρ καὶ τὸ ἐξ αύτοῦ ἀπαύγασμα} (ignes et splendor ex eo ortus), \textgreek{Δύο μὲν τῷ εἶναι} [this is not accurate, and strictly taken would lead to two \textgreek{οὐσίαι}] \textgreek{καὶ ὁρᾶσθαι, ἕν δὲ τῷ ἐξ αὐτοῦ καὶ ἀδιαίρετον εἶναι τὸ ἀπαύγασμα αὐτοῦ}.} or between fountain and stream; though in these illustrations the proverbial insufficiency of all similitudes must never be forgotten. “We must not,” says he, “take the words in John xiv. 10: ’I am in the Father and the Father in Me’ as if the Father and the Son were two different interpenetrating and mutually complemental substances, like two bodies which fill one vessel. The Father is full and perfect, and the Son is the fulness of the Godhead.”\footnote{0rat. iii. c. Arian. c. 1 (p. 551): \textgreek{Πλήρης καὶ τέλειός ἐστιν ὁ πατὴρ, καὶ πλήρωμα θεότητός ἐστιν ὁ Υἱός .}} “We must not imagine,” says he in another place, “three divided substances\footnote{\textgreek{Τρεῖς ὑποστάσεις} [here, as often in the Nicene age, synonymous with \textgreek{οὐσίαι}] \textgreek{μεμερισμένας καθ  ἑαυτάς}. Athan. Expos. Fidei or \textgreek{Ἔκθεσις πίστεως}, cap. 2 (Opera, ed. Bened. i. p. 100).} in God, as among men, lest we, like the heathen, invent a multiplicity of gods; but as the stream which is born of the fountain, and not separated from it, though there are two forms and names. Neither is the Father the Son, nor the Son the Father; for the Father is the Father of the Son, and the Son is the Son of the Father. As the fountain is not the stream, nor the stream the fountain, but the two are one and the same water which flows from the fountain into the stream; so the Godhead pours itself, without division, from the Father into the Son. Hence the Lord says: I went forth from the Father, and come from the Father. Yet He is ever with the Father, He is in the bosom of the Father, and the bosom of the Father is never emptied of the Godhead of the Son.”\footnote{Expositio Fidei, cap. 2: \textgreek{Ὡς γὰρ οὐκ ἔστιν ἡ πηγὴ ποταμὸς , οὐδε ὁ ποταμὸς πηγὴ, ἀμφότερα δὲ ἓν καὶ ταὐτόν ἐστιν ὕδωρ τὸ ἐκ της πηγῆς μετεχευόμενον, οὕτως ἡ ἐκ τοῦ πατρὸς εἰς τὸν υἱὸν θεότης ἀῤῥεύστως καὶ ἀδιαιρέτως τυγχάνει, κ.τ.λ.}}

The Son is of the essence of the Father, not by division or diminution, but by simple and perfect self-communication. This divine self-communication of eternal love is represented by the figure of generation, suggested by the biblical terms Father and Son, the only-begotten Son, the firstborn.\footnote{\textgreek{Πατὴρ, υἱὸς, μονογενης υἱός}(frequent in John), \textgreek{πρωτότοκος πάσης κτίσεως}(Col. i. 15). Waterland (Works, i. p. 368) says of this point of the Nicene doctrine, “that an explicit profession of \emph{eternal generation} might have been dispensed with: provided only that the eternal existence of the \textgreek{λόγος}. as \emph{a real subsisting person}, \emph{in and of the Father}, which comes to the same thing, might be secured. This was the point; and this was all.”} The eternal generation is an internal process in the essence of God, and the Son is an immanent offspring of this essence; whereas creation is an act of the will of God, and the creature is exterior to the Creator, and of different substance. The Son, as man, is produced;\footnote{\textgreek{Γενητός}(not to be confounded with \textgreek{γεννητός}), \textgreek{ποιητός}, factus. Comp. John i. 14: \textgreek{Ὁ λόγος σὰρξ ἐγένετο}.} as God, he is unproduced or uncreated;\footnote{\textgreek{Ἀγένητος, οὐ ποιηθείς}, non-factus, increatus; not to be confounded with \textgreek{ἀγέννητος}, non-genitus, which belongs to the Father alone.} he is begotten\footnote{\textgreek{Γεννητός}, or, as in the Symb. Nic. \textgreek{γεννηθείς}, genitus.} from eternity of the unbegotten\footnote{\textgreek{Ἀγέννητος}, non-genitus. This terminology is very frequent in the writings of Athanasius, especially in the Orat. i. contra Arianos, and in his Epist. de decretis Syn. Nic.} Father. To this Athanasius refers the passage concerning the Only-begotten who is in the bosom of the Father.\footnote{John i. 18: \textgreek{Ὁ μονογενὴς υἱὸς , ὁ ὢν} (a perpetual or eternal relation, not \textgreek{ἧν}) \textgreek{εἰς}(motion, in distinction from \textgreek{ἐν}) \textgreek{τὸν κολπον τοῦ πατρός}. Comp. Athanas. Epist. de decr. S. N. c. 22 (tom. i. p. 227): \textgreek{Τί γὰρ ἄλλο τὸ ἐν κόλποις σημαίνει, ἣ τὴν γνησίαν ἐκ τοῦ πατρὸς τοῦ υἱοῦ γέννησιν;}}

Generation and creation are therefore entirely different ideas. Generation is an immanent, necessary, and perpetual process in the essence of God himself, the Father’s eternal communication of essence or self to the Son; creation, on the contrary, is an outwardly directed, free, single act of the will of God, bringing forth a different and temporal substance out of nothing. The eternal fatherhood and sonship in God is the perfect prototype of all similar relations on earth. But the divine generation differs from all human generation, not only in its absolute spirituality, but also in the fact that it does not produce a new essence of the same kind, but that the begotten is identical in essence with the begetter; for the divine essence is by reason of its simplicity, incapable of division, and by reason of its infinity, incapable of increase.\footnote{Bishop John Pearson, in his well-known work: An Exposition of the Creed (Art. ii. p. 209, ed. W. S. Dobson, New York, 1851), thus clearly and rightly exhibits the Nicene doctrine in this point: “In human generations the son is of the same nature with the father, and yet is not the same man; because though he has an essence of the same kind, yet he has not the same essence; the power of generation depending on the first prolifical benediction, \emph{increase and multiply}, it must be made by way of multiplication, and thus every son becomes another man. But the divine essence, being by reason of its simplicity not subject to division, and in respect of its infinity incapable of multiplication, is so communicated as not to be multiplied; insomuch that he who proceeds by that communication, has not only the same nature, but is also the same God. The Father God, and the Word God; Abraham man and Isaac man: but Abraham one man, Isaac another man; not so the Father one God and the Word another, but the Father and the Word both the same God.”} The generation, properly speaking, has no reference at all to the essence, but only to the hypostatical distinction. The Son is begotten not as God, but as Son, not as to his natura, but as to his ἰδιότης, his peculiar property and his relation to the Father. The divine essence neither begets, nor is begotten. The same is true of the processio of the Holy Ghost, which has reference not to the essence, but only to the person, of the Spirit. In human generation, moreover, the father is older than the son; but in the divine generation, which takes place not in time, but is eternal, there can be no such thing as priority or posteriority of one or the other hypostasis. To the question whether the Son existed before his generation, Cyril of Alexandria answered: “The generation of the Son did not precede his existence, but he existed eternally, and eternally existed by generation.” The Son is as necessary to the being of the Father, as the Father to the being of the Son.

The necessity thus asserted of the eternal generation does not, however, impair its freedom, but is intended only to deny its being arbitrary and accidental, and to secure its foundation in the essence of God himself. God, to be Father, must from eternity beget the Son, and so reproduce himself; yet he does this in obedience not to a foreign law, but to his own law and the impulse of his will. Athanasius, it is true, asserts on the one hand that God begets the Son not of his will,\footnote{\textgreek{Μὴ ἐκ βουλήσεως}.} but by his nature,\footnote{\textgreek{Φύσει.}} yet on the other hand he does not admit that God begets the Son without will,\footnote{\textgreek{Ἀβουλήτως} and \textgreek{ἀθελήτως}.} or of force or unconscious necessity. The generation, therefore, rightly understood, is an act at once of essence and of will. Augustine calls the Son “will of will.”\footnote{Voluntas de voluntate. De trinit. xv. 20.} In God freedom and necessity coincide.

The mode of the divine generation is and must be a mystery. Of course all human representations of it must be avoided, and the matter be conceived in a purely moral and spiritual way. The eternal generation, conceived as an intellectual process, is the eternal self-knowledge of God; reduced to ethical terms, it is his eternal and absolute love in its motion and working within himself.

In his argument for the consubstantiality of the Son, Athanasius, in his four orations against the Arians, besides adducing the proof from Scripture, which presides over and permeates all other arguments, sets out now in a practical method from the idea of redemption, now in a speculative, from the idea of God.

Christ has delivered us from the curse and power of sin, reconciled us with God, and made us partakers of the eternal, divine life; therefore he must himself be God. Or, negatively: If Christ were a creature, he could not redeem other creatures from sin and death. It is assumed that redemption is as much and as strictly a divine work, as creation.\footnote{Comp. particularly the second oration contra Arianos, c. 69 sqq.}

Starting from the idea of God, Athanasius argues: The relation of Father is not accidental, arising in time; else God would be changeable;\footnote{Orat. i. contra Arianos, c. 28 (p. 433): \textgreek{Διὰ τοῦτο ἀεὶ πατὴρ καὶ οὐκ ἐπιγέγονε} (accidit) \textgreek{τῷ Θεῷ τὸ πατὴρ, ἳνα μὴ καὶ τρεπτὸς εἶναι νομισθῇ. Εἰ γὰρ καλὸν τὸ εἶναι αὐτὸν πατέρα, οὐκ ἀεὶ δὲ ἦν πατὴρ, οὐκ ἀεὶ ἄρα τὸ καλὸν ἦν αὐτῷ}. Though to this it might be objected that by the incarnation of the Logos and the permanent reception of human nature into fellowship with the divine, a certain change has passed, after all, upon the deity.} it belongs as necessarily to the essence and character of God as the attributes of eternity, wisdom, goodness, and holiness; consequently he must have been Father from eternity, and this gives the eternal generation of the Son.\footnote{Orat. ii. c. Arianos, c. 1 sqq. (p. 469 sqq.); Orat. iii. c. 66 (p. 615), and elsewhere.} The divine fatherhood and sonship is the prototype of all analagous relations on earth. As there is no Son without Father, no more is there Father without Son. An unfruitful Father were like a dark light, or a dry fountain, a self-contradiction. The non-existence of creatures, on the contrary, detracts nothing from the perfection of the Creator, since he always has the power to create when he will.\footnote{This last argument, in the formally logical point of view, may not be perfectly valid; for there may as well be a distinction between an ideal and real fatherhood, as between an ideal and real creatorship; and, on the other hand, one might reason with as good right backwards from the notion of essential omnipotence to an eternal creation, and say with Hegel: Without the world God is not God. But from the speculative and ethical point of view a difference must unquestionably be admitted, and an element of truth be acknowledged in the argument of Athanasius. The Father needed the Son for his own self-consciousness, which is inconceivable without an object. God is essentially love, and this realizes itself in the relation of Father and Son, and in the fellowship of the Spirit: Ubi amor ibi trinitas.} The Son is of the Father’s own interior essence, while the creature is exterior to God and dependent on the act of his will.\footnote{Orat. i. c. 29 (p. 433): \textgreek{Τὸ ποίημα ἔξωθεν τοῦ ποιοῦντός ἐστιν ... ὁ δὲ υἱὸς ἴδιον τῆς οὐσίας γέννημά ἐστι· διὸ καὶ τὸ μὲν ποίημα οὐκ ἀνάγκη ἀεὶ εἶναι, ὅτε γὰρ βούλεται ὁ δημιουργὸς ἐργάζεται, τὸ δὲ γέννημα οὐ βουλήσει ὑπόκειται, ἀλλὰ της οὐσίας ἐστὶν ἰδιότης .}} God, furthermore, cannot be conceived without reason (ἄλογος), wisdom, power, and according to the Scriptures (as the Arians themselves concede) the Son is the Logos, the wisdom, the power, the Word of God, by which all things were made. As light rises from fire, and is inseparable from it, so the Word from God, the Wisdom from the Wise, and the Son from the Father.\footnote{Comp. the 4th Oration against the Arians, cap. 1 sqq. (p. 617 sqq.)} The Son, therefore, was in the beginning, that is, in the beginning of the eternal divine being, in the original beginning, or from eternity. He himself calls himself one with the Father, and Paul praises him as God blessed forever.\footnote{The \textgreek{Θεός}in the well-known passage, Rom. ix. 6, is thus repeatedly by Athanasus, \emph{e.g}., Orat. i. contra Arianos, c. 11; Orat. iv. c. 1, and by other fathers (Irenaeus, Tertullian, Cyprian, Origen, Chrysostom), as well as by the Reformers and most of the orthodox expositors, referred to Christ. This interpretation, too, is most suitable to the connection, and in perfect harmony with the Christology of Paul, who sets forth Christ as the image of God, the possessor of the fulness of the divine life and glory, the object of worship (Phil. ii. 6; Col. i. 15 ff.; ii. 9; 2 Cor. iv. 4; Eph. v. 5; 1 Tim. iii. 16; Tit. ii. 13); and who therefore, as well as John, i. 1, could call him in the predicative sense \textgreek{Θεός}, \emph{i.e}., of divine essence, in distinction from \textgreek{ὁ Θεός}with the article.}

Finally Christ cannot be a proper object of worship, as he is represented in Scripture and has always been regarded in the Church, without being strictly divine. To worship a creature is idolatry.

When we attentively peruse the warm, vigorous, eloquent, and discriminating controversial writings of Athanasius and his co-laborers, and compare with them the vague, barren, almost entirely negative assertions and superficial arguments of their opponents, we cannot escape the impression that, with all their exegetical and dialectical defects in particulars, they have on their side an overwhelming preponderance of positive truth, the authority of holy Scripture, the profounder speculations of reason, and the prevailing traditional faith of the early church.\footnote{We say the prevailing \emph{faith}; not denying that the theological \emph{knowledge} and \emph{statement} of the doctrine of the trinity had hitherto been in many respects indefinite and wavering. The learned bishop Bull, indeed, endeavored to prove, in opposition to the Jesuit Petavius, that the ante-Nicene fathers taught concerning the deity of the Son the very same things as the Nicene. Comp. the Preface to his Defensio fidei Nicaenae, ed. Burton, Oxf. 1827, vol. v. Pars. 1, p. ix.: “De summa rei, quam aliis persuadere volo, plane ipse, neque id temere, persuasus sum, nempe, quod de Filii divinitate contra Arium, idem re ipsa (quanquam aliis fortasse nonnunquam verbis, alioque loquendi modo) docuisse Patres ac doctores ecclesiae probatos ad unum omnes, qui ante tempora synodi Nicaenae, ab ipsa usque apostolorum aetate, floruerunt.” But this assertion can be maintained only by an artificial and forced interpretation of many passages, and goes upon a mechanical and lifeless view of history. Comp. also the observations of W. Cunningham, Historical Theology, vol. i. p. 269 ff.}

The spirit and tendency of the Nicene doctrine is edifying; it magnifies Christ and Christianity. The Arian error is cold and heartless, degrades Christ to the sphere of the creature, and endeavors to substitute a heathen deification of the creature for the true worship of God. For this reason also the faith in the true and essential deity of Christ has to this day an inexhaustible vitality, while the irrational Arian fiction of a half-deity, creating the world and yet himself created, long ago entirely outlived itself.\footnote{Dorner, l.c. i. p. 883, justly says: “Not only to the mind of our time, but to all sound reason, does it seem absurd, nay, superstitious, that an under-god, a finite, created being, should be the creator.”}

\xsection{The Doctrine of the Holy Spirit}

§ 128. The Doctrine of the Holy Spirit.

The decision of Nicaea related primarily only to the essential deity of Christ. But in the wider range of the Arian controversies the deity of the Holy Ghost, which stands and falls with the deity of the Son, was indirectly involved. The church always, indeed, connected faith in the Holy Spirit with faith in the Father and Son, but considered the doctrine concerning the Holy Spirit as only an appendix to the doctrine concerning the Father and the Son, until the logical progress brought it to lay equal emphasis on the deity and personality of the Holy Ghost, and to place him with the Father and Son as an element of equal claim in the Trinity.

The Arians made the Holy Ghost the first creature of the Son, and as subordinate to the Son as the Son to the Father. The Arian trinity was therefore not a trinity immanent and eternal, but arising in time and in descending grades, consisting of the uncreated God and two created demi-gods. The Semi-Arians here, as elsewhere, approached the orthodox doctrine, but rejected the consubstantiality, and asserted the creation, of the Spirit. Thus especially Macedonius, a moderate Semi-Arian, whom the Arian court-party had driven from the episcopal chair of Constantinople. From him the adherents of the false doctrine concerning the Holy Spirit, were, after 362, called Macedonians;\footnote{\textgreek{Μακεδονιανοί}.} also Pneumatomachi,\footnote{\textgreek{Πνευματόμαχοι}.} and Tropici.\footnote{\textgreek{Τροπικοί}. This name comes probably from their explaining as mere tropes (figurative expressions) or metaphors the passages of Scripture from which the orthodox derived the deity of the Holy Spirit. Comp. Athanas., Ad Serap. Ep. i. c. 2 (tom. i. Pars ii. p. 649).}

Even among the adherents of the Nicene orthodoxy an uncertainty still for a time prevailed respecting the doctrine of the third person of the Holy Trinity. Some held the Spirit to be an impersonal power or attribute of God; others, at farthest, would not go beyond the expressions of the Scriptures. Gregory Nazianzen, who for his own part believed and taught the consubstantiality of the Holy Ghost with the Father and the Son, so late as 380 made the remarkable concession:\footnote{Orat. xxxi. De Spiritu sancto, cap. 5 (Op. tom. i. p. 559, and in Thilo’s Bibliotheca P. Gr. dogm. vol. ii. p. 503).} “Of the wise among us, some consider the Holy Ghost an influence, others a creature, others God himself,\footnote{\textgreek{τῶν καθ  ἡμᾶς σοφῶν οἱ μὲν ἐνέργειαν τοῦτο} [\textgreek{τὸ πνεῦμα ἅγιον}] \textgreek{ὑπέλαβον, οἱ δὲ κτίσμα, οἱ δὲ Θεόν.}} and again others know not which way to decide, from reverence, as they say, for the Holy Scripture, which declares nothing exact in the case. For this reason they waver between worshipping and not worshipping the Holy Ghost,\footnote{Ou\textgreek{τε σέβουσιν, οὔτε ἀτιμάζουσι}.} and strike a middle course, which is in fact, however, a bad one.” Basil, in 370, still carefully avoided calling the Holy Ghost God, though with the view of gaining the weak. Hilary of Poictiers believed that the Spirit, who searches the deep things of God, must be divine, but could find no Scripture passage in which he is called God, and thought that he must be content with the existence of the Holy Ghost, which the Scripture teaches and the heart attests.\footnote{De trinitate, ii. 29; and xii. 55.}

But the church could not possibly satisfy itself with only two in one. The baptismal formula and the apostolic benediction, as well as the traditional trinitarian doxologies, put the Holy Ghost on an equality with the Father and the Son, and require a divine tri-personality resting upon a unity of essence. The divine triad tolerates in itself no inequality of essence, no mixture of Creator and creature. Athanasius well perceived this, and advocated with decision the consubstantiality of the Holy Spirit against the Pneumatomachi or Tropici.\footnote{In the four Epistles to Serapion, bishop of Tmuis, written in 362 (Ep. ad Serapionem Thmuitanum episcopum contra illos qui blasphemant et dicunt Spiritum S. rem creatam esse), in his Opera, ed. Bened. tom. i. Pars ii. pp. 647-714; also in Thilo’s Biblioth. Patr. Graec. dogmatica, vol. i. pp. 666-819.} Basil did the same,\footnote{De Spiritu Sancto ad S. Amphilochium Iconii episcopum (Opera, ed. Bened. tom. iii. and in Thilo’s Bibl. vol. ii. pp. 182-343).} and Gregory of Nazianzum,\footnote{Orat. xxxi. De Spiritu Sancto (Opera, tom. i. p. 556 sqq. and in Thilo’s Bibl. vol. ii. pp. 497-537).} Gregory of Nyssa,\footnote{Orat. catech. c. 2. Comp. Rupp, Gregor v. Nyasa, p. 169 sq.} Didymus,\footnote{De Spiritu S., translated by Jerome.} and Ambrose.\footnote{De Spiritu S. libri 3.}

This doctrine conquered at the councils of Alexandria, a.d. 362, of Rome, 375, and finally of Constantinople, 381, and became an essential constituent of the ecumenical orthodoxy.

Accordingly the Creed of Constantinople supplemented the Nicene with the important addition: “And in the Holy Ghost, who is Lord and Giver of life, who with the Father is worshipped and glorified, who spake by the prophets.”\footnote{Similar additions had already been previously made to the Nicene Creed. Thus Epiphanius in his Ancoratus, c. 120, which was written in 374, gives the Nicene Creed as then already in general use with the following passage on the Holy Spirit: \textgreek{Καὶ εἰς τὸ ἅγιον πνεῦμα πιστεύομεν, τὸ λαλῆσαν ἐν νόμῳ, καὶ κηρύξαν ἐν τοῖς προφήταις καὶ καταβὰν ἐπὶ τὸν Ἰορδάνην, λαλοῦν ἐν ἀποστόλοις , οἰκοῦν ἐν ἁγίοις· οὕτως δὲ πιστεύομεν ἐν αὐτῷ, ὅτι ἐστὶ πνεῦμα ἅγιον, πνεῦμα Θεοῦ, πνεῦμα τέλειον, πνεῦμα παράκλητον, ἄκτιστον, ἐκ τοῦ πατρὸς ἐκπορευόμενον, καὶ ἐκ τοῦ υἱοῦ λαμβανόμενον καὶ πιστευόμενον}. His shorter Creed, Anc. c. 119 (in Migne’s ed. tom. iii. 231), even literally agrees with that of Constantinople, but in both he adds the anathema of the original Nicene Creed.}

This declares the consubstantiality of the Holy Ghost, not indeed in words, yet in fact, and challenges for him divine dignity and worship.

The exegetical proofs employed by the Nicene fathers for the deity of the Holy Ghost are chiefly the following. The Holy Ghost is nowhere in Scripture reckoned among creatures or angels, but is placed in God himself, co-eternal with God, as that which searches the depths of Godhead (1 Cor. ii. 11, 12). He fills the universe, and is everywhere present (Ps. cxxxix. 7), while creatures, even angels, are in definite places. He was active even in the creation (Gen. i. 3), and filled Moses and the prophets. From him proceeds the divine work of regeneration and sanctification (John iii. 5; Rom. i. 4; viii. 11; 1 Cor. vi. 11; Tit. iii. 5–7; Eph. iii. 16; v. 17, 19, \&c). He is the source of all gifts in the church (1 Cor. xii). He dwells in believers, like the Father and the Son, and makes them partakers of the divine life. Blasphemy against the Holy Ghost is the extreme sin, which cannot be forgiven (Matt. xii. 31). Lying to the Holy Ghost is called lying to God (Acts v. 3, 4). In the formula of baptism (Matt. xxviii. 19), and likewise in the apostolic benediction (2 Cor. xiii. 13), the Holy Ghost is put on a level with the Father and the Son and yet distinguished from both; he must therefore be truly divine, yet at the same time a self-conscious person.\footnote{The well-known passage concerning the three witnesses in heaven, I John v. 7, is not cited by the Nicene fathers: a strong evidence that it was wanting in the manuscripts of the Bible at that time.} The Holy Ghost is the source of sanctification, and unites us with the divine life, and thus must himself be divine. The divine trinity tolerates in itself nothing created and changeable. As the Son is begotten of the Father from eternity, so the Spirit proceeds from the Father through the Son. (The procession of the Spirit from the Son, on the contrary, is a subsequent inference of the Latin church from the consubstantiality of the Son, and was unknown to the Nicene fathers.)

The distinction between generation and procession is not particularly defined. Augustine calls both ineffable and inexplicable.\footnote{“Ego distinguere nescio, non valeo, non sufficio, propterea quia sicut generatio ita processio inenarrabilis est.”} The doctrine of the Holy Ghost was not in any respect so accurately developed in this period, as the doctrine concerning Christ, and it shows many gaps.

\xsection{The Nicene and Constantinopolitan Creed}

§ 129. The Nicene and Constantinopolitan Creed.

We look now at the Creeds of Nicaea and Constantinople side by side, which sum up the result of these long controversies. We mark the differences by inclosing in brackets the parts of the former omitted by the latter, and italicizing the additions which the latter makes to the former.

The Nicene Creed of 325\footnote{It is found, together with the similar Eusebian (Palestinian) confession, in the well-known Epistle of Eusebius of Caesarea to his diocese (Epist. ad suae parochiae homines), which is given by Athanasius at the close of his Epist. de decretis Nicaenae Synodi (Opera, tom. i. p. 239, and in Thilo’s Bibl. vol. i. p. 84 sq.); also, though with some variation by Theodoret, H. E. i. 12, and Socrates, H. E. i. 8. Sozomen omitted it (H. E. i. 10) from respect to the disciplina arcani. The Symbolum Nicaenum is given also, with unessential variations, by Athanasius in his letter to the emperor Jovian c. 3, and by Gelasius Cyzic., Lib. Synod. de Concil. Nicaeno, ii. 36. On the unimportant variations in the text, Comp. Walch, Bibl. symbol. p. 75 sqq., and A. Rahn, Bibliothek der Symbole, 1842. Comp. also the parallel Creeds of the Nicene age in the Appendix to Pearson’s Exposition of the Creed.}

the Nicaeno-Constantinopolitan Creed of 381\footnote{Found in the Acts of the second ecumenical council in all the collections (Mansi, tom. iii. 566; Harduin, i. 814). It probably does not come directly from this council still less from the individual authorship of Gregory of Nyssa or Gregory of Nazianzum to whom it has sometimes been ascribed, but the additions by which it is distinguished from the Nicene, were already extant in substance under different forms (in the Symbolum Epiphanii, for example, and the Sym b. Basilii Magni), and took shape gradually in the course of the controversy. It is striking that it is not mentioned as distinct from the Nicene by Gregory Nazianzen in his Epist. 102 to Cledonius (tom. ii. 93 ed. Paris 1842), nor by the third ecumenical council at Ephesus. On the other hand, it was twice recited at the council of Chalcedon, twice adopted in the acts, and thus solemnly sanctioned. Comp. Hefele, ii. 11, 12.}

Πιστεύομεν εἰς ἕνα Θεὸν πατέρα παντοκράτορα, πάντων ὁρατῶν τε καὶ ἀοράτων ποιητήν

Πιστεύομεν εἰς ἕνα Θεὸν, πατέρα παντοκράτορα, ποιητὴν οὐρανοῦ καὶ γῆς, ὁρατῶν τε πάντων καὶ ἀοράτων.

Καὶ εἰς ἕνα κύριον Ἰησοῦν Χριστὸν, τὸν υἱὸν τοῦ Θεοῦ· γεννηθέντα ἐκ τοῦ πατρὸς [μονογενῆ· τοῦτ ἔστιν ἐκ τῆς οὐσίας τοῦ πατρὸς· Θεὸν ἐκ Θεοῦ καὶ\footnote{\textgreek{Καὶ} is wanting in Athanasius (De decretis, etc,).} ] φῶς ἐκ φωτὸς, Θεὸν ἀληθινὸν ἐκ Θεοῦ ἀληθινοῦ· γεννηθεντα, οὐ ποιηθέντα, ὁμοούσιον τῷ πατρὶ· δι  οὗ τὰ πάντα ἐγένετο [τά τε ἐν τῷ οὐρανῷ καὶ τὰ ἐν τῇ γῇ·] τὸν δι  ἡμᾶς τοὺς ἀνθρώπους καὶ διὰ τὴν ἡμετέραν σωτηρίαν κατελθόντα καὶ σαρκωθέντα, καὶ \footnote{\textgreek{Καὶ} is wanting in Athanasius; Socrates and Galerius have it.} ἐνανθρωπήσαντα· παθόντα\footnote{Gelasius adds \textgreek{ταφέντα}, buried.} καὶ ἀναστάντα τῇ τρίτῃ ἡμέρᾳ, ἀνελθόντα εἰς τοὺς\footnote{Without the article in Athanasius.} οὐρανοὺς,\footnote{Al. \textgreek{καί}.} ἐρχόμενον κρίναι ξῶντας καὶ νεκρους.

Καὶ εἰς ἕνα κύριον Ἰησοῦν Χριστὸν τὸν υἱὸν τοῦ Θεοῦ τὸν μονογενῆ· τὸν ἐκ τοῦ πατρός γεννηθέντα πρὸ πάντων τῶν αἰώνων· φῶς ἐκ φωτὸς, Θεὸν ἀληθινὸν ἐκ Θεοῦ ἀληθινοῦ, γεννηθέντα, οὐ ποιηθέντα, ὁμοούσιον τῷ πατρὶ· δι  οὗ τὰ πάντα ἐγένετο· τὸν δι  ἡμᾶς τοὺς ἀνθρώπους καὶ διὰ τὴν ἡμετέραν σωτηρίαν κατελθόντα ἐκ τῶν οὐρανῶν, καὶ σαρκωθέντα ἐκ πνεύματος ἁγίου καὶ Μαρίας τῆς παρθένου, καὶ ἐνανθρωπήσαντα· σταυρωθέντα τε ὑπέρ ἡμῶν ἐπὶ Ποντίου Πιλάτου, καὶ παθόντα, καὶ ταφέντα, καὶ ἀναστάντα τῇ τρίτῃ ἡμέρᾳ κατὰ τὰς γραφὰς, καὶ ἀνελθόντα εἰς τούς οὐρανοὺς, καὶ καθεζόμενον ἐκ δεξιῶν τοῦ πατρὸς, καὶ πάλιν ἐρχόμενον μετὰ δόξης κρίναι ζῶντας καὶ νεκροὺς· οὗ τῆς βασιλείας οὐκ ἔσται τέλος.

Καὶ εἰς τὸ ἅγιον πνεῦμα.

Καὶ εἰς τὸ πνεῦμα τὸ ἂγιον, τὸ κύριον, τὸ ζωοποιὸν, τὸ ἐκ τοῦ πα-τρὸς ἐκπορευόμενον, τὸ σὺν πατρὶ καὶ υἱῷ προσκυνούμενον καὶ συνδοξαζόμενον, τὸ λαλῆσαν διὰ τῶν προφητῶν. Εἰς μίαν ἁγίαν καθολικὴν καὶ ἀποστολικὴν ἐκ-κλησίαν· ὁμολογοῦμεν ἓν βάπτισμα εἰς ἄφεσιν ἁμαρτιῶν· προσδοκ-ῶμεν ἀνάστασιν νεκρῶν καὶ ζωὴν τοῦ μέλλοντος αἰῶνος. Ἀμήν.

Τοὺς δὲ λέγοντας, ὅτι \footnote{Athanasius omits \textgreek{ὅτι}.} ἦν ποτε ὅτε οὐκ ἦν· καὶ· πρὶν γεννηθῆναι οὐκ ἦν· καὶ ὅτι ἐξ οὐκ ὄντων ἐγένετο· ἢ ἐξ ἑτέρας ὑποστάσεως ἢ οὐσίας\footnote{Here hypostasis and essence are still used interchangeably; though Basil and Bull endeavor to prove a distinction. Comp. on the contrary, Petavius, De trinit. l. iv. c. 1 (p. 314 sqq.). Rufinus, i. 6, translates: “Ex alia subsistentia aut substantia.”} φάσκοντας εἶναι· ἢ κτιστὸν, ἢ τρεπτὸν, ἢ ἀλλοιωτὸν τὸν υἰὸν τοῦ Θεοῦ· ἀναθεματίζει ἡ ἁγία καθολικὴ καὶ ἀποστολικὴ\footnote{Athanasius omits \textgreek{ἁγία} and \textgreek{ἀποστολική}. Theodoret has both predicates, Socrates has \textgreek{ἀποστολική}, all read \textgreek{καθολική}.} ἐκκλησία. 

We believe in one God, the Father Almighty, Maker of all things visible, and invisible.

We believe in one God, the Father Almighty, Maker of heaven and earth, and of all things visible

“And in one Lord Jesus Christ, the Son of God, begotten of the Father [the only-begotten, i.e., of the essence of the Father, God of God, and] Light of Light, very God of very God, begotten, not made, being of one substance with the Father; by whom all things were made [in heaven and on earth]; who for us men, and for our salvation, came down and was incarnate and was made man; he suffered, and the third day he rose again, ascended into heaven; from thence he cometh to judge the quick and the dead.

“And in one Lord Jesus Christ, the only-begotten Son of God, begotten of the Father before all worlds ( aeons ), \footnote{This addition appears as early as the creeds of the council of Antioch in 341.} Light of Light, very God of very God, begotten, not made, being of one substance with the Father; by whom all things were made; who for us men, and for our salvation, came down from heaven , and was incarnate by the Holy Ghost of the Virgin Mary , and was made man; he was crucified for us under Pontius Pilate, and suffered , and was buried , and the third day he rose again, according to the Scriptures, and ascended into heaven , and sitteth on the right hand of the Father ; from thence he cometh again, with glory , to judge the quick and the dead ; whose kingdom shall have no end . \footnote{This addition likewise is found substantially in the Antiochian creeds of 341, and is directed against Marcellus of Ancyra, Sabellius, and Paul of Samosata, who taught that the union of the power of God (\textgreek{ἐνέργεια δραστική}) with the man Jesus will cease at the end of the world, so that the Son and His kingdom are not eternal Comp. Hefele, i. 438 and 507 sq.}

“And in the Holy Ghost.

“And in the Holy Ghost, who is Lord and Giver of life, who pro-ceedeth from the Father, who with the Father and the Son together is worshipped and glorified, who spake by the prophets. In one holy catholic and apostolic church, we acknowledge one baptism for the remission of sins; we look fo r the resurrection of the dead, and the life of the world to come. Amen . \footnote{Similar additions concerning the Holy Ghost, the catholic church, baptism and life everlasting are found in the older symbols of Cyril of Jerusalem, Basil, and the two Creeds of Epiphanius. See § 128 above, and Appendix to Pearson on the Creed, p. 594 ff.}

[“And those who say: there was a time when he was not; and: he was not before he was made; and: he was made out of nothing, or out of another substance or thing, or the Son of God is created, or changeable, or alterable; they are condemned by the holy catholic and apostolic Church.” ]

A careful comparison shows that the Constantinopolitan Creed is a considerable improvement on the Nicene, both in its omission of the anathema at the close, and in its addition of the articles concerning the Holy Ghost and concerning the church and the way of salvation. The addition: according to the Scriptures, is also important, as an acknowledgment of this divine and infallible guide to the truth. The whole is more complete and symmetrical than the Nicaenum, and in this respect is more like the Apostles’ Creed, which, in like manner, begins with the creation and ends with the resurrection and the life everlasting, and is disturbed by no polemical dissonance; but the Apostles’ Creed is much more simple in structure, and thus better adapted to the use of a congregation and of youth, than either of the others.

The Constantinopolitan Creed maintained itself for a time by the side of the Nicene, and after the council of Chalcedon in 451, where it was for the first time formally adopted, it gradually displaced the other. Since that time it has itself commonly borne the name of the Nicene Creed. Yet the original Nicene confession is still in use in some schismatic sects of the Eastern church.

The Latin church adopted the improved Nicene symbol from the Greek, but admitted, in the article on the Holy Ghost, the further addition of the well-known filioque, which was first inserted at a council of Toledo in 589, and subsequently gave rise to bitter disputes between the two

\xsection{The Nicene, Doctrine of the Trinity. The Trinitarian Terminology}

§ 130. The Nicene, Doctrine of the Trinity. The Trinitarian Terminology.

The doctrine of the essential deity and the personality of the Holy Ghost completed the Nicene doctrine of the Trinity ; and of this doctrine as a whole we can now take a closer view.

This fundamental and comprehensive dogma secured both the unity and the full life of the Christian conception of God; and in this respect it represents, as no other dogma does, the whole of Christianity. It forms a bulwark against heathen polytheism on the one hand, and Jewish deism and abstract monotheism on the other. It avoids the errors and combines the truth of these two opposite conceptions. Against the pagans, says Gregory of Nyssa, we hold the unity of essence; against the Jews, the distinction of hypostases. We do not reject all multiplicity, but only such as destroys the unity of the being, like the pagan polytheism; no more do we reject all unity, but only such unity as denies diversity and full vital action. The orthodox doctrine of the Trinity, furthermore, formed the true mean between Sabellianism and tritheism, both of which taught a divine triad, but at the expense, in the one case, of the personal distinctions, in the other, of the essential unity. It exerted a wholesome regulative influence on the other dogmas. It overcame all theories of emanation, established the Christian conception of creation by a strict distinction of that which proceeds from the essence of God, and is one with him, like the Son and the Spirit, from that which arises out of nothing by the free will of God, and is of different substance. It provided for an activity and motion of knowledge and love in the divine essence, without the Origenistic hypothesis of an eternal creation. And by the assertion of the true deity of the Redeemer and the Sanctifier, it secured the divine character of the work of redemption and sanctification.

The Nicene fathers did not pretend to have exhausted the mystery of the Trinity, and very well understood that all human knowledge, especially in this deepest, central dogma, proves itself but fragmentary. All speculation on divine things ends in a mystery, and reaches an inexplicable residue, before which the thinking mind must bow in humble devotion. “Man,” says Athanasius, “can perceive only the hem of the garment of the triune God; the cherubim cover the rest with their wings.” In his letter to the Monks, written about 358, he confesses that the further he examines, the more the mystery eludes his understanding,\footnote{Ep. ad Monachos (Opera tom. i. p. 343).} and he exclaims with the Psalmist: “Such knowledge is too wonderful for me; it is high, I cannot attain unto it.”\footnote{Ps. cxxxix 6.} Augustine says in one place: “If we be asked to define the Trinity, we can only say, it is not this or that.”\footnote{Enarrat. in PS. xxvi. 8. John Damascenus (Expos. fidei) almost reaches the Socratic confession, when he says: All we can know concerning the divine nature is, that it cannot be conceived. Of course, such concessions are to be understood \emph{cum grano salis}.} But though we cannot explain the how or why of our faith, still the Christian may know, and should know, what he believes, and what he does not believe, and should be persuaded of the facts and truths which form the matter of his faith.

The essential points of the orthodox doctrine of the Trinity are these:

1. There is only one divine essence or substance.\footnote{\textgreek{Οὐσία}, substantia, essentia, \textgreek{φύσις}, natura, \textgreek{τὸ ὄν, τὸ ὑποκείμενον}. Comp. Petavius, De Trinitate lib. iv. c. 1 (ed. Par. tom. ii. p. 311): “Christiani scriptores ... \textgreek{οὐσίαν}appellant non singularem individuamque, sed communem individuis substantiam.” The word \textgreek{ὑποκείμενον,}however, is sometimes taken as equivalent to provswpon.} Father, Son, and Spirit are one in essence, or consubstantial.\footnote{\textgreek{Ὁμοούσιοι.} On the import of this, comp. § 127, and in the text above.} They are in one another, inseparable, and cannot be conceived without each other. In this point the Nicene doctrine is thoroughly monotheistic or monarchian, in distinction from tritheism, which is but a new form of the polytheism of the pagans.

The terms essence (οὐσία) and nature (φύσις), in the philosophical sense, denote not an individual, a personality, but the genus or species; not unum in numero, but ens unum in multis. All men are of the same substance, partake of the same human nature, though as persons and individuals they are very different.\footnote{“We men,” says Athanasius, “consisting of body and soul are all \textgreek{μίας φύσεως καὶ οὐσίας ,} but many persons.”} The term homoousion, in its strict grammatical sense, differs from monoousion or toutoousion, as well as from heteroousion, and signifies not numerical identity, but equality of essence or community of nature among several beings. It is clearly used thus in the Chalcedonian symbol, where it is said that Christ is “consubstantial (homoousios) with the Father as touching the Godhead, and consubstantial with us [and yet individually, distinct from us] as touching the manhood.” The Nicene Creed does not expressly assert the singleness or numerical unity of the divine essence (unless it be in the first article: “We believe in one God”); and the main point with the Nicene fathers was to urge against Arianism the strict divinity and essential equality of the Son and Holy Ghost with the Father. If we press the difference of homoousion from monoousion, and overlook the many passages in which they assert with equal emphasis the monarchia or numerical unity of the Godhead, we must charge them with tritheism.\footnote{Cudworth (in his great work on the Intellectual System of the Universe, vol. ii p. 437 ff.) elaborately endeavors to show that Athanasius and the Nicene fathers actually taught three divine substances in the order of subordination. But he makes no account of the fact that the terminology and the distinction of \textgreek{οὐσία} and \textgreek{ὑποστασις}were at that time not yet clearly settled.}

But in the divine Trinity consubstantiality denotes not only sameness of kind, but at the same time numerical unity; not merely the unum in specie, but also the unum in numero. The, three persons are related to the divine substance not as three individuals to their species, as Abraham, Isaac, and Jacob, or Peter, John, and Paul, to human nature; they are only one God. The divine substance is absolutely indivisible by reason of its simplicity, and absolutely inextensible and untransferable by reason of its infinity; whereas a corporeal substance can be divided, and the human nature can be multiplied by generation. Three divine substances would limit and exclude each other, and therefore could not be infinite or absolute. The whole fulness of the one undivided essence of God, with all its attributes, is in all the persons of the Trinity, though in each in his own way: in the Father as original principle, in the Son by eternal generation, in the Spirit by, eternal procession. The church teaches not one divine essence and three persons, but one essence in three persons. Father, Son, and Spirit cannot be conceived as three separate individuals, but are in one another, and form a solidaric unity.\footnote{Comp. the passages from Athanasius and other fathers cited at § 126. “The Persons of the Trinity,” says R. Hooker (Eccles. Polity, B. v. ch. 56, voL ii. p. 315 in Keble’s edition), quite in the spirit of the Nicene orthodoxy, “are not three particular substances to whom one \emph{general} nature is common, but three that subsist by one substance \emph{which itself} \emph{is particular}: yet they all three have it and their several ways of having it are that which makes their personal distinction. The Father therefore is in the Son, and the Son in Him, they both in the Spirit and the Spirit in both them. So that the Father’s offspring, which is the Son, remaineth eternally in the Father; the Father eternally also in the Son, no way severed or divided by reason of the sole and single unity of their substance. The Son in the Father as light in that light out of which it floweth without separation; the Father in the Son as light in that light which it causeth and leaveth not. And because in this respect his eternal being is of the Father, which eternal being is his life, therefore he by the Father liveth.” In a similar strain, Cunningham says in his exposition of the Nicene doctrine of the Trinity (Hist. Theology, i. p. 285): “The unity of the divine nature as distinguished from the nature of a creature, might be only a specific and not a numerical unity, and this nature might be possessed by more than one divine being; but the Scriptures plainly ascribe a numerical unity to the Supreme Being, and, of course, preclude the idea that there are several different beings who are possessed of the one divine nature. This is virtually the same thing as teaching us that the one divine nature is possessed only by one essence or substance, from which the conclusion is clear, that if the Father be possessed of the divine nature, and if the Son, with a distinct personality, be also possessed of the divine nature, the Father and the Son must be of one and the same substance; or rather—for it can scarcely with propriety be called a conclusion or consequence—the doctrine of the consubstantiality of the Son with the Father is just an expression or embodiment of the one great truth, the different component parts of which are each established by scriptural authority, viz.: that the Father and the Son, having distinct personality in the unity of the Godhead, are both equally possessed of the divine, as distinguished from the created, nature. Before any creature existed, or had been produced by God out of nothing, the Son existed in the possession of the divine nature. If this be true, and if it be also true that God is in any sense one, then it is likewise true—for this is just according to the established meaning of words, the current mode of expressing it—that the Father and the Son are the same in substance as well as equal in power and glory.”}

Many passages of the Nicene fathers have unquestionably a tritheistic sound, but are neutralized by others which by themselves may bear a Sabellian construction so that their position must be regarded as midway between these two extremes. Subsequently John Philoponus, an Aristotelian and Monophysite in Alexandria about the middle of the sixth century, was charged with tritheism, because he made no distinction between φύσιςand ὑπόστασις, and reckoned in the Trinity three natures, substances, and deities, according to the number of persons.\footnote{On tritheism, and the doctrine of John Philoponus and John Ascusnages, which is known to us only in fragments, comp. especially Baur, Lehre von der Dreieinigkeit, etc., vol. ii. pp. 13-32. In the English Church the error of tritheism was revived by Dean Sherlockin his “Vindication of the Doctrine of the Holy and ever Blessed Trinity,” 1690. He maintained that, with the exception of a mutual consciousness of each other, which no created spirits can have, the three divine persons are “three distinct infinite minds” or “three intelligent beings.” He was opposed by South, Wallis, and others. See Patrick Fairbairn’s Appendix to the English translation of Dorner’s History of Christology, vol. iii. p. 354 ff. (Edinburgh, 1863).}

2. In this one divine essence there are three persons\footnote{\textgreek{Πρόσωπα,} \emph{personae}. This term occurs very often in the New Testament, now in the sense of \emph{person}, now of \emph{face} or \emph{countenance}, again of \emph{form} or external appearance. Etymologically (from \textgreek{πρός·}and \textgreek{ἡ ὤψ} , the eye, face), it means strictly \emph{face}; then in general, \emph{front}; also \emph{mask, visor, character} (of a drama); and finally, \emph{person}, in the grammatical sense. In like manner the Latin word \emph{persona} (from \emph{sonus}, sound) signifies the mask of the Roman actor, through which he made himself audible (\emph{personuit}); then the actor himself; then any assumed or real character; and finally an individual a reasonable being. Sabellianism used the word in the sense of face or character; tritheism in the grammatical sense. Owing to this ambiguity of the word, the term \emph{hypostasis} is to be preferred, though this too is somewhat inadequate. Comp. the Lexicons, and especially Petavius, De trinit., lib. iv. Dr. Shedd also prefers \emph{hypostasis}, and observes, vol. i. p. 371: ” This term (\emph{persona}), it is obvious to remark, though the more common one in English, and perhaps in Protestant trinitarianism generally, is not so well adapted to express the conception intended, as the Greek \textgreek{ὐπόστασις}. It has a Sabellian leaning, because it does not with sufficient plainness indicate the \emph{subsistence} in the Essence. The Father, Son, and Spirit are more than mere aspects or appearances of the Essence. The Latin \emph{persona} was the mask worn by the actor in the play, and was representative of his particular character for the particular time. Now, although those who employed these terms undoubtedly gave them as full and solid a meaning as they could, and were undoubtedly true trinitarians, yet the representation of the eternal and necessary hypostatical distinctions in the Godhead, by terms derived from transitory scenical exhibitions, was not the best for purposes of science, even though the poverty of human language should justify their employment for popular and illustrative statements.”} or, to use a better term, hypostases, \footnote{\textgreek{Ὑποστάσεις} \emph{subsistentiae}. Comp. Heb. i 3. (The other passages of the New Testament where the word is used, Heb. iii. 14; xi. 1; 2 Cor. ix. 4; xi. 17, do not belong here.) \textgreek{Ὑπόστασις}, and the corresponding Latin \emph{sub-stantia}, strictly \emph{foundation}, then \emph{essence, substance}, is originally pretty much synonymous with \textgreek{οὐσία}, \emph{essentia}, and is in fact as we have already said, frequently interchanged with it, even by Athanasius, and in the anathema at the close of the original Nicene Creed. But gradually (according to Petavius, after the council at Alexandria in 862) a distinction established itself in the church terminology, in which Gregory of Nyasa, particularly in his work: De differentia essentiae et hypostaseos (tom. iii. p. 32 sqq.) had an important influece. Comp. Petavius, l.c. p. 314 sqq.} that is, three different modes of subsistence\footnote{\textgreek{Τρόποι ὑπάρξεως}, an expression, however, capable of a Sabellian sense.} of the one same undivided and indivisible whole, which in the Scriptures are called the Father, the Son, and the Holy Ghost.\footnote{This question of the \emph{tri}-personality of God must not be confounded with the modern question of the \emph{personality} of God in general. The tri-personality was asserted by the Nicene fathers in opposition to abstract monarchianism and Sabellianism; the personality is asserted by Christian theism against pantheism, which makes a personal relation of the spirit of man to God impossible. Schleiermacher, who as a philosopher leaned decidedly to pantheism, admitted (in a note to his Reden über die Religion) that devotion and prayer always presume and require the personality of God. The philosophical objection, that personality necessarily includes limitation by other personalities, and so contradicts the notion of the absoluteness of God, is untenable; for we can as well conceive an absolute personality, as an absolute intelligence and an absolute will, to which, however, the power of self-limitation must be ascribed, not as a weakness, but as a perfection. The orthodox tri-personality does not conflict with this total personality, but gives it full organic life.} These distinctions are not merely different attributes, powers, or activities of the Godhead, still less merely subjective aspects under which it presents itself to the human mind; but each person expresses the whole fulness of the divine being with all its attributes, and the three persons stand in a relation of mutual knowledge and love. The Father communicates his very life to the Son, and the Spirit is the bond of union and communion between the two. The Son speaks, and as the God-Man, even prays, to the Father, thus standing over against him as a first person towards a second; and calls the Holy Ghost “another Comforter” whom he will send from the Father, thus speaking of him as of a third person.\footnote{John xiv. 16: \textgreek{Ἄλλον παράκλητον,}comp. v. 26; c. xv. 26: \textgreek{Ὁ παράκλητος·, ὅν ἐγὼ πεμψω ὑμῖν παρὰ πατρός ,} —a clear distinction of Spirit, Son, and Father.}

Here the orthodox doctrine forsook Sabellianism or modalism, which, it is true, made Father, Son, and Spirit strictly coordinate, but only as different denominations and forms of manifestation of the one God.

But, on the other hand, as we have already intimated, the term person must not be taken here in the sense current among men, as if the three persons were three different individuals, or three self-conscious and separately acting beings. The trinitarian idea of personality lies midway between that of a mere form of manifestation, or a personation, which would lead to Sabellianism, and the idea of an independent, limited human personality, which would result in tritheism. In other words, it avoids the monoousian or unitarian trinity of a threefold conception and aspect of one and the same being, and the triousian or tritheistic trinity of three distinct and separate beings.\footnote{Comp. Petavius, l.c., who discusses very fully the trinitarian terminology of the Nicene fathers. Also J. H. Newman, The Arians, etc. p. 208: “The word \emph{person}, which we venture to use in speaking of those three distinct manifestations of Himself, which it has pleased Almighty God to give us, is in its philosophical sense too wide for our meaning. Its essential signification, as applied to ourselves, is that of \emph{an individual intelligent} agent, answering to the Greek \textgreek{ὑπόστασις}, or reality. On the other hand, if we restrict it to its etymological sense of \emph{persona} or \textgreek{πρόσωπον}, \emph{i.e., character}, it evidently means less than Scripture doctrine, which we wish to ascertain by it; denoting merely certain outward expressions of the Supreme Being relatively to ourselves, which are of an accidental and variable nature. The statements of Revelation then lie between this internal and external view of the Divine Essence, between Tritheism, and what is popularly called Unitarianism.” Dr. Shedd, History of Christian Doctrine, vol. i. p. 365: ” The doctrine of a subsistence in the substance of the Godhead brings to view a species of existence that is so anomalous and unique, that the human mind derives little or no aid from those analogies which assist it in all other cases. The hypostasis is a real subsistence,—a solid essential form of existence, and not a mere emanation, or energy, or manifestation,—but it is intermediate between substance and attributes. It is not identical with the substance, for there are not three substances. It is not identical with attributes, for the three Persons each and equally possess all the divine attributes .... Hence the human mind is called upon to grasp the notion of a species of existence that is totally \emph{sui generis}, and not capable of illustration by any of the ordinary comparisons and analogies.”} In each person there is the same inseparable divine substance, united with the individual property and relation which distinguishes that person from the others. The word person is in reality only a make-shift, in the absence of a more adequate term. Our idea of God is more true and deep than our terminology, and the essence and character of God far transcends our highest ideas.\footnote{As Augustinesays, De trinitate, lib. vii. cap. 4 (§ 7, ed. Bened. Venet. tom. viii. foL 858): “Verius cogitatur Deus quam dicitur, et verius est quam cogitatur.”}

The Nicene fathers and Augustine endeavored, as Tertullian and Dionysius of Alexandria had already done, to illustrate the Trinity by analogies from created existence. Their figures were sun, ray, and light; fountain, stream, and flow; root, stem, and fruit; the colors of the rainbow;\footnote{Used by Basil and Gregory of Nyasa.} soul, thought, and spirit;\footnote{\textgreek{Ψυχὴ, ἐνθύμησις, πνεῦμα}, in Gregory Nazianzen.} memory, intelligence, and will;\footnote{Augustine, De trinit. x. c. 11 (§ 18), tom. viii. fol. 898: “Haec tria, memoria, intelligentia, voluntas, quoniam non sunt tres vitae, sed una vita, nec tres mentes, sed una mens: consequenter utique non tres substantiae sunt, sed una substantia.”} and the idea of love, which affords the best illustration, for God is love.\footnote{Augustine, ib. viii. 8 (f. 875): “Immo vero vides trinitatem, si caritatem vides; ” ix. 2 (f. 879): “Tria sunt, amans, et quod amatur, et amor.” And in another place: “Tres sunt, amans, amatus, et mutuus amor.”} Such figures are indeed confessedly insufficient as proofs, and, if pressed, might easily lead to utterly erroneous conceptions. For example: sun, ray, and light are not co-ordinate, but the two latter are merely qualities or emanations of the first. “Omne simile claudicat.”\footnote{This was clearly felt and confessed by the fathers themselves, who used these illustrations merely as helps to their understanding. Joh. Damascenus (De fide orthod. l. i. c. 8; Opera, tom. i. p. 137) says: “It is impossible for any image to be found in created things, representing in itself the nature of the Holy Trinity without any point of dissimilitude. For can a thing created, and compound, and changeable, and circumscribed, and corruptible, clearly express the superessential divine essence, which is exempt from all these defects?” Comp. Mosheim’s notes to Cudworth, vol. ii. 422 f. (Lond. ed. of 1845).} Analogies, however, here do the negative service of repelling the charge of unreasonableness from a doctrine which is in fact the highest reason, and which has been acknowledged in various forms by the greatest philosophers, from Plato to Schelling and Hegel, though often in an entirely unscriptural sense. A certain trinity undeniably runs through all created life, and is especially reflected in manifold ways in man, who is created after the image of God; in the relation of body, soul, and spirit; in the faculties of thought, feeling, and will; in the nature of self-consciousness;\footnote{The trinity of self-consciousness consists in a process of becoming objective to one’s self, and knowing one’s self in this objectivity, according to the logical law of thesis, antithesis, and synthesis, or in the unity of the subject thinking and the subject thought. This speculative argument has been developed by Leibnitz, Hegel, and other German philosophers, and is adopted also by Dr. Shedd, Hist. of Christian Doct. i. p. 366 ff., note. But this analogy properly leads at best only to a Sabellian tri-personality, not to the orthodox.} and in the nature of love.\footnote{The ethical induction of the Trinity from the idea of love was first attempted by Augustine, and has more recently been pursued by Sartorius, J. Müller, J. P. Lange, Martensen, Liebner, Schöberlein, and others. It is suggested by the moral essence of God, which is love, the relation of the Father to the Son, and the “fellowship” of the Holy Ghost, and it undoubtedly contains a deep element of truth; but, strictly taken, it yields only two different personalities and an impersonal relation, thus proving too much for the Father and the Son, and too little for the Holy Spirit.} 3. Each divine person has his property, as it were a characteristic individuality, expressed by the Greek word ἰδιότης,\footnote{Also \textgreek{ἴδιον.}Gregory of Nyssa calls these characteristic distinctions \textgreek{γνωριστικαὶ ἰφιότητες,}peculiar marks of recognition. The terms \textgreek{ἰδιότης ,} and \textgreek{ὑπόστασις}were sometimes used synonymously. The word \textgreek{ἰδιότης}, fem. (from \textgreek{ἴδιος}), peculiarity, is of course not to be confounded with \textgreek{ἰδιώτης}, masc., which likewise comes from \textgreek{ἴδιος}, but means a private man, then layman, then an imbecile, idiot,} and the Latin proprietas.\footnote{\emph{Proprietas personalis}; also \emph{character hypostaticus}.} This is not to be confounded with attribute; for the divine attributes, eternity, omnipresence, omnipotence, wisdom, holiness, love, etc., are inherent in the divine essence, and are the common possession of all the divine hypostases. The idiotes, on the contrary, is a peculiarity of the hypostasis, and therefore cannot be communicated or transferred from one to another.

To the first person fatherhood, or the being unbegotten, \footnote{\textgreek{Ἀγεννησία}, \emph{paternitas}.} is ascribed as his property; to the second, sonship, or the being begotten;\footnote{\textgreek{γεννησία, γέννησις,} \emph{generatio filiatio}.} to the Holy Ghost, procession.\footnote{\textgreek{Εκπόρευσις}, \emph{procesio}; also \textgreek{ἔκπεμψις}, \emph{missio}; both from John xv. 15 (\textgreek{πέμψω} —\textgreek{ἐκπορεύεται}) and similar passages, which relate, however, not to the eternal trinity of constitution, but to the historical trinity of manifestation. Gregory Nazianzen says: \textgreek{Ἴδιον πατρὸς μὲν ἡ ἀγεννησία, υἱοῦ δέ ἡ γέννησις , πνεύματος δὲ ἡ ἔκπεμψις .}} In other words: The Father is unbegotten, but begetting; the Son is uncreated, but begotten; the Holy Ghost proceeds from the Father (and, according to the Latin doctrine, also from the Son). But these distinctions relate, as we have said, only to the hypostases, and have no force with respect to the divine essence which is the same in all, and neither begets nor is begotten, nor proceeds, nor is sent.

4. The divine persons are in one another, mutually interpenetrate, and form a perpetual intercommunication and motion within the divine essence; as the Lord says: “I am in the Father, and the Father in me;” and “the Father that dwelleth in me, he doeth the works.”\footnote{John xiv. 10: \textgreek{Ὁ δὲ πατήρ ὁ ἐν ἐμοὶ μένων, αὐτὸς ποιεῖ τὰ ἔργα;} v.11: \textgreek{Ἐγω ἐν τῷ πατρὶ, καὶ ὁ πατήρ ἐν ἐμοί.}This also refers, strictly, not to the eternal relation, but to the indwelling of the Father in the historical, incarnate Christ.} This perfect indwelling and vital communion was afterwards designated (by John of Damascus and the scholastics) by such terms as ἐνύπαρξις, περιχώρησις,\footnote{From \textgreek{περιχωρέω} (with \textgreek{εἰς}), to circulate, go about, \emph{progredi, ambulare}. Comp. Petavius, De trinit., lib. iv. c. 16 (tom. ii. p. 453 sqq.), and De incarnatione, lib. iv. c. 14 (tom. iv. p. 473 sqq.). The thing itself is clearly taught even by the Nicene fathers, especially by Athanasius in his third Oration against the Arians, c. 3 sqq., and elsewhere, with reference to the relation of the Son to the Father, although he never, so far as I know, used the word \textgreek{περιχώρησις}. Gregory Nazianzen uses the verb \textgreek{περιχωρεῖν}(not the noun) of the vital interpenetration of the two natures in Christ. Gibbon, in his contemptuous account of the Nicene controversy (chapter xxi.) calls the \textgreek{περιχώρησις}or \emph{circumincessio} ” the deepest and darkest corner of the whole theological abyss,” but takes no pains even to explain this idea. The old Protestant theologians defined the \textgreek{περιχώρησις}as “immanentia, h. e. inexistentia mutua et singularissima, intima et perfectissima inhabitatio unius personae in alia.” Comp. Joh. Gerhard, Loci theologici, tom. i. p. 197 (ed. Cotta).} inexistentia, immanentia, inhabitatio, circulatio, permeatio, intercommunio, circumincessio.\footnote{From \emph{incedo}, denoting the perpetual internal motion of the Trinity, the circumfusio or mutua commeatio, et communicatio personarum inter se. Petavius (in the 2d and 4th vol. l. c.), Cudworth (Intellectual System of the Universe, vol. ii. p. 454, ed. of Harrison, Lond. 1845), and others use instead of this, \emph{circuminsessio}, from \emph{sedeo}, which rather expresses the \emph{repose} of the persons in one another, the inexistentia or mutua existentia personarum. This would correspond to the Greek \textgreek{ἐνύπαρξις}rather than to \textgreek{περιχώρησις}.}

5. The Nicene doctrine already contains, in substance, a distinction between two trinities: an immanent trinity of constitution,\footnote{Ad intra, \textgreek{τρόπος ὑπάρξεως .}} which existed from eternity, and an economic trinity of manifestation;\footnote{Ad extra, \textgreek{τρόπος·ἀποκαλύψεως}} though this distinction did not receive formal expression till a much later period. For the generation of the Son and the procession of the Spirit are, according to the doctrine, an eternal process. The perceptions and practical wants of the Christian mind start, strictly speaking, with the trinity of revelation in the threefold progressive work of the creation, the redemption, and the preservation of the world, but reason back thence to a trinity of being; for God has revealed himself as he is, and there can be no contradiction between his nature and his works. The eternal pre-existence of the Son and the Spirit is the background of the historical revelation by which they work our salvation. The Scriptures deal mainly with the trinity of revelation, and only hint at the trinity of essence, as in the prologue of the Gospel of John which asserts an eternal distinction between God and the Logos. The Nicene divines, however, agreeably to the metaphysical bent of the Greek mind, move somewhat too exclusively in the field of speculation and in the dark regions of the intrinsic and ante-mundane relations of the Godhead, and too little upon the practical ground of the facts of salvation.

6. The Nicene fathers still teach, like their predecessors, a certain subordinationism, which seems to conflict with the doctrine of consubstantiality. But we must distinguish between, a subordinatianism of essence (οὐσία) and a subordinatianism of hypostasis, of order and dignity.\footnote{\textgreek{Ὑποταγὴ τάξεως καὶ ἀξιώματος}.} The former was denied, the latter affirmed. The essence of the Godhead being but one, and being absolutely perfect, can admit of no degrees. Father, Son, and Spirit all have the same divine essence, yet not in a co-ordinate way, but in an order of subordination. The Father has the essence originally and of himself, from no other; he is the primal divine subject, to whom alone absoluteness belongs, and he is therefore called preeminently God,\footnote{\textgreek{Ὁ Θεός}, and \textgreek{αὐτόθεος}, in distinction from \textgreek{Θεός}. Waterland (Works, vol. i. p. 315) remarks on this: ” The title of \textgreek{ὁ Θεός}, being understood in the same sense with \textgreek{αὐτόθεος ,} was, as it ought to be, generally reserved to the Father, as the distinguishing personal character of the first Person of the Holy Trinity. And this amounts to no more than the acknowledgment of the Father’s prerogative, as Father. But as it might also signify any Person who is truly and essentially God, it might properly be applied to the Son too: and it is so applied sometimes, though not so often as it is to the Father.”} or the principle, the fountain, and the root of Godhead.\footnote{\textgreek{Ἡ πηγὴ, ἡ αἰτία, ἡ ῥίζα τῆς θεότητος}: \emph{fons, origo}, \emph{principium}.} The Son, on the contrary, has his essence by communication from the Father, therefore, in a secondary, derivative way. “The Father is greater than the Son.” The one is unbegotten, the other begotten; the Son is from the Father, but the Father is not from the Son; fatherhood is in the nature of the case primary, sonship secondary. The same subordination is still more applicable to the Holy Ghost. The Nicene fathers thought the idea of the divine unity best preserved by making the Father, notwithstanding the triad of persons, the monad from which Son and Spirit spring, and to which they return.

This subordination is most plainly expressed by Hilary of Poictiers, the champion of the Nicene doctrine in the West.\footnote{De trinit. iii. 12: “Et quis non Patrem potiorem confitebitur, ut ingenitum a genito, ut Patrem a Filio, ut eum qui miserit ab eo qui missus sit, ut volentem ab eo qui obediat? Et ipse nobis erit testis: Pater major me est. Haec ita ut sunt intelligenda sunt, sed cavendum est, ne apud imperitos gloriam Filii honor Patris infirmet.” In the same way Hilaryderives all the attributes of the Son from the Father. Comp. also Hilary, De Synodis, seu de fide Orientalium, pp. 1178 and 1182 (Opera, ed. Bened.), and the third and eighteenth canons of the Sirmian council of 357.} The familiar comparisons of fountain and stream, sun and light, which Athanasius, like Tertullian, so often uses, likewise lead to a dependence of the Son upon the Father\footnote{Comp. the relevant passages from Athanasius, Basil, and the Gregories, in Bull, Defensio, sect. iv. (Pars ii. p. 688 sqq.). Even John of Damascus, with whom the productive period of the Greek theology closes, still teaches the same subordination, De orthod. fide, i. 10: \textgreek{Πάντα ὅσα ἔχει ὁ υἱὸς καὶ τὸ πνεῦμα, ἐκ τοῦ πατρὸς ἔχει, καὶ αὐτὸ τὸ εἶναι.}} Even the Nicaeno-Constantinopolitan Creed favors it, in calling the Son God of God, Light of Light, very God of very God. For if a person has anything, or is anything, of another, he has not that, or is not that, of himself. Yet this expression may be more correctly understood, and is in fact sometimes used by the later Nicene fathers, as giving the Son and Spirit only their hypostases from the Father, while the essence of deity is common to all three persons, and is co-eternal in all.

Scriptural argument for this theory of subordination was found abundant in such passages as these: “As the Father hath life in himself (ἔχει ζωὴν ἐν ἑαυτῷ), so hath he given (ἔδωκε) to the Son to have life in himself; and hath given him authority to execute judgment also;”\footnote{John v. 26, 27.} “All things are delivered unto me (πάντα μοι παρεδόθη) of my Father;”\footnote{Matt. xi. 27; Comp. xxviii. 18.} “My father is greater than I.”\footnote{John xiv. 28. Cudworth (I. c. ii. 422) agrees with several of the Nicene fathers in referring this passage to the divinity of Christ, for the reason that the superiority of the eternal God over mortal man was no news at all. Mosheim in a learned note to Cudworth \emph{in} \emph{loco}, protests against both interpretations, and correctly so. For Christ speaks here of his entire divine-human person, but in the state of \emph{humiliation}.} But these and similar passages refer to the historical relation of the Father to the incarnate Logos in his estate of humiliation, or to the elevation of human nature to participation in the glory and power of the divine,\footnote{John xvii. 5; Pbil. ii. 9-11.} not to the eternal metaphysical relation of the Father to the Son.

In this point, as in the doctrine of the Holy Ghost, the Nicene system yet needed further development. The logical consistency of the doctrine of the consubstantiality of the Son, upon which the Nicene fathers laid chief stress, must in time overcome this decaying remnant of the ante-Nicene subordinationism.\footnote{All important scholars since Petavius admit the subordinatism in the Nicene doctrine of the trinity; \emph{e.g}., Bull, who in the fourth (not third, as Gibbon says) section of his famous Defensio fidei Nic. (Works, vol. v. Pars ii. pp. 685-796) treats quite at large of the subordination of the Son to the Father, and in behalf of the identity of the Nicene and ante-Nicene doctrine proves that all the orthodox fathers, before and after the council of Nice, “uno ore docuerunt naturam perfectionesque divinas Patri Filioque competere non callateraliter aut coördinate, sed subordinate; hoc est, Filium eandem quidem naturam divinam cum Patre communem habere, sed a Patre communicatam; ita scilicet ut Pater solus naturam illam divinam a se habeat, sive a nullo alio, Filius autem a Patre; proinde Pater divinitatis, quae in Filio est, origo se principium sit,” etc. So Waterland, who, in his vindication of the orthodox doctrine of the Trinity against Samuel Clarke, asserts such a supremacy of the Father as is consistent with the eternal and necessary existence, the consubstantiality, and the infinite perfection of the Son. Among modem historians Neander, Gieseler, Baur (Lehre von der Dreieinigkeit, etc. i. p. 468 ff.), and Dorner (Lehre von der Person Christi, i p. 929 ff.) arrive at the same result. But while Baur and Dorner (though from different points of view) recognize in this a defect of the Nicene doctrine, to be overcome by the subsequent development of the church dogma, the great Anglican divines, Cudworth (Intellectual System, vol. ii. p. 421 ff.), Pearson, Bull, Waterland (and among American divines Dr. Shedd) regard the Nicene subordinationism as the true, Scriptural, and final form of the trinitarian doctrine, and make no account of Augustine, who went beyond it. Kahnis (Der Kirchenglaube, ii. p. 66 ff.) thinks that the Scriptures go still further than the Nicene fathers in subordinating the Son and the Spirit to the Father.}

\xsection{The Post-Nicene Trinitarian Doctrine of Augustine}

§ 131. The Post-Nicene Trinitarian Doctrine of Augustine.

Augustine: De trinitate, libri xv., begun in 400, and finished about 415; and his anti-Arian works: Contra sermonem Arianorum; Collatio cum Maximino Arianorum episcopo; Contra Maximinum haereticum, libri ii. (all in his Opera omnia, ed. Bened. of Venice, 1733, in tom. viii. pp. 626–1004; and in Migne’s ed. Par. 1845, tom. viii. pp. 683–1098).

While the Greek church stopped with the Nicene statement of the doctrine of the Trinity, the Latin church carried the development onward under the guidance of the profound and devout speculative spirit of Augustine in the beginning of the fifth century, to the formation of the Athanasian Creed. Of all the fathers, next to Athanasius, Augustine performed the greatest service for this dogma, and by his discriminating speculation he exerted more influence upon the scholastic theology and that of the Reformation, than all the Nicene divines. The points in which he advanced upon the Nicene Creed, are the following:\footnote{The Augustinian doctrine of the trinity is discussed at length by Baur, Die christl. Lehre von der Dreieinigkeit, etc. vol. i. pp. 826-888. Augustinehad but an imperfect knowledge of the Greek language, and was therefore not accurately acquainted with the writings of the Nicene fathers, but was thrown the more upon his own thinking. Comp. his confession, De trinit. l. iii. cap. 1 (tom. viii. f. 793, ed. Bened. Venet., from which in this section I always quote, though giving the varying chapter-division of other editions).}

1. He eliminated the remnant of subordinationism, and brought out more clearly and sharply the consubstantiality of the three persons and the numerical unity of their essence.\footnote{De trinit. l. vii. cap. 6 (§ 11), tom. viii. f. 863: “Non major essentia est Pater et Filius et Spiritus Sanctus simul, quam solus Pater, aut solus Filius; sed tres simul illae substantiae [here equivalent to \textgreek{ὑποστάσει}“ ] sive personae, si ita dicendae sunt, aequales sunt singulis: quod animalis homo non percipit.” Ibid. (f. 863): “Ita dicat unam essentiam, ut non existimet aliud alio vel majus, vel melius, vel aliqua ex parte divisum.” Ibid. lib. viii. c. 1 (fol. 865): “Quod vero ad se dicuntur singuli, non dici pluraliter tres, sed unam ipsam trinitatem: sicut Deus Pater, Deus Filius, Deus Spiritus Sanctus; et bonus Pater, bonus Filius, bonus Spiritus Sanctus; et omnipotens Pater, omnipotens Filius, omnipotens Spiritus Sanctus; nec tamen tres Dii, aut tres boni, aut tres omnipotentes, sed unus Deus, bonus, omnipotens ipsa Trinitas.” Lib. xv. 17 (fol. 988): “Pater Deus, et Filius Deus, et Spiritus S. Deus, et simul unus Deus.” De Civit. Dei, xi. cap. 24: “Non tres Dii vel tres omnipotentes, sed unus Deus omnipotens.” So the Athanasian Creed, vers. 11.}

Yet he too admitted that the Father stood above the Son and the Spirit in this: that he alone is of no other, but is absolutely original and independent; while the Son is begotten of him, and the Spirit proceeds from him, and proceeds from him in a higher sense than from the Son.\footnote{De trinit. l. xv. c. 26 (§ 47, fol. 1000): ”\emph{Pater solus non est de alio}, ideo solus appellatur ingenitus, non quidem in Scripturis, sed in consuetudine disputantium ... Filius autem de Patre natus est: et Spiritus Sanctus de Patre \emph{principaliter}, et ipso sine ullo temporis intervallo dante, communiter de utroque procedit.”} We may speak of three men who have the same nature; but the persons in the Trinity are not three separately subsisting individuals. The divine substance is not an abstract generic nature common to all, but a concrete, living reality. One and the same God is Father, Son, and Spirit. All the works of the Trinity are joint works. Therefore one can speak as well of an incarnation of God, as of an incarnation of the Son, and the theophanies of the Old Testament, which are usually ascribed to the Logos, may also be ascribed to the Father and the Holy Ghost.

If the orthodox doctrine of the Trinity lies midway between Sabellianism and tritheism, Augustine bears rather to the Sabellian side. He shows this further in the analogies from the human spirit, in which he sees the mystery of the Trinity reflected, and by which he illustrates it with special delight and with fine psychological discernment, though with the humble impression that the analogies do not lift the veil, but only make it here and there a little more penetrable. He distinguishes in man being, which answers to the Father, knowledge or consciousness, which answers to the Son, and will, which answers to the Holy Ghost.\footnote{Confess. xiii. 11: “Dico haec tria: esse, nosse, velle. Sum enim, et novi, et volo; sum sciens, et volens; et scio esse me, et velle; et volo esse, et scire. In his igitur tribus quam sit inseparabilis vita, et una vita, et una mens, et una essentia, quam denique inseparabilis distinctio, et tamen distinctio, videat qui potest.” This comparison he repeats in a somewhat different form, De Civit. Dei, xi. 26.} A similar trinity he finds in the relation of mind, word, and love; again in the relation of memory, intelligence, and will or love, which differ, and yet are only one human nature (but of course also only one human person).\footnote{Mens, verbum, amor;—memoria, intelligentia, voluntas or caritas; for voluntas and caritas are with him essentially the same: “Quid enim est aliud caritas quam voluntas?” Again: amans, amatus, mutuus amor. On these, and similar analogies which we have already mentioned in § 130, comp. Augustine, De Civit. Dei, l. xi. c. 24; De trinit. xiv. and xv., and the criticism of Baur, l.c. i. p. 844 sqq.}

2. Augustine taught the procession of the Holy Ghost from the Son as well as from the Father, though from the Father mainly. This followed from the perfect essential unity of the hypostases, and was supported by some passages of Scripture which speak of the Son sending the Spirit.\footnote{John xv. 26: \textgreek{Ὁ παράκλητος, ὃν ἐγω πέμψω ὑμῖν παρὰ τοῦ πατρός ,}and xvi. 7: \textgreek{Πεμψω αὐτὸν πρὸς ὑμᾶς}; compared with John xiv. 26: \textgreek{Τὸ πνεῦμα τὸ ἅγιον, ὃ πέμψει ὁ πατήρ ἐν τῷ ὀνοματί μου.} Augustineappeals also to John xx. 22, where Christ breathes the Holy Ghost on his disciples, De trinit. iv. c. 20 (§ 29), fol. 829: “Nec possumus dicere quod Spiritus S. et a Filio non procedat, neque enim frusta idem Spiritus et Patris et Filii Spiritus dicitur. Nec video quid aliud significare voluerit, cum sufflans in faciem discipulorum ait: ’Accipite Spiritum S.’ ” Tract. 99 in Evang. Joh. § 9: “Spiritus S. non de Patre procedit in Filium, et de Filio procedit ad sanctificandam creatuam, sed simul de utroque procedit.” But after all, he makes the Spirit proceed \emph{mainly} from the Father: de patre \emph{principaliter}. De trinit. xv. c. 26 (§ 47). Augustinemoreover regards the procession of the Spirit from the Son as the gift of the Father which is implied in the communication of life to the Son. Comp. Tract 99 in Evang. Joh. § 8: “A quo habet Filius ut sit Deus (est enim de Deo Deus), ab illo habet utique ut etiam de illo procedat Spiritus Sanctus: ac per hoc Spiritus Sanctus ut etiam de Filio procedat sicut procedit de Patre, ab ipso habet Patre.”} He also represented the Holy Ghost as the love and fellowship between Father and Son, as the bond which unites the two, and which unites believers with God.\footnote{De trinit. xv. c. 17 (§ 27) fol. 987: “Spiritus S. secundum Scriptams sacras nec Patris solius est, nec Filii solius, sed amborum, et ideo communem, qua invicem se diligunt Pater et Filius, nobis insinuat caritatem.” Undoubtedly God is love; but this may be said in a special sense of the Holy Ghost. De trinit. xv. c. 17 (§ 29), fol. 988: “Ut scillicet in illa simplici summaque natura non sit aliud substantia et aliud caritas, sed substantia ipsa sit caritas et caritas ipsa sit substantia, sive in Patre, sive in Filio, sive in Spiritu S., et tamen proprie Spiritus S. caritas nuncupetur.”}

The Nicaeno-Constantinopolitan Creed affirms only the processio Spiritus a Patre, though not with an exclusive intent, but rather to oppose the Pneumatomachi, by giving the Spirit a relation to the Father as immediate as that of the Son. The Spirit is not created by the Son, but eternally proceeds directly from the Father, as the Son is from eternity begotten of the Father. Everything proceeds from the Father, is mediated by the Son, and completed by the Holy Ghost. Athanasius, Basil, and the Gregories give this view, without denying procession from the Son. Some Greek fathers, Epiphanius,\footnote{Ancor. § 9: \textgreek{Ἄρα Θεὸς ἐκ πατρὸς καὶ υἱοῦ τὸ πνεῦμα.}Yet he says not expressly: \textgreek{ἐκπορεύεται ἐκ τοῦ υἱου.}} Marcellus of Ancyra,\footnote{Though in a Sabellian sense.} and Cyril of Alexandria,\footnote{Who in his anathemas against Nestorius condemns also those who do not derive the Holy Ghost from Christ. Theodoret replied: If it be meant that the Spirit is of the same essence with Christ, and proceeds from the Father, we agree; but if it be intended that the Spirit has his existence through the Son, this is impious. Comp. Neander, Dogmengesch. i. p. 822.} derived the Spirit from the Father and the Son; while Theodore of Mopsuestia and Theodoret would admit no dependence of the Spirit on the Son.

Augustine’s view gradually met universal acceptance in the West. It was adopted by Boëthius, Leo the Great and others. \footnote{Comp. the passages in Hagenbach’s Dogmengeschichte, vol. i. p. 267 (in the Engl. ed. by H. B. Smith, New York, 1861), and in Perthel, Leo der G. p. 138 ff Leo says, \emph{e.g}., Serm. lxxv. 2: “Huius enim beata trinitatis incommutabilis deitas una est in substantia indivisa in opere, concors in voluntate, par in potential aequalis in gloria.”} It was even inserted in the Nicene Creed by the council of Toledo in 589 by the addition of filioque, together with an anathema against its opponents, by whom are meant, however, not the Greeks, but the Arians.

Here to this day lies the main difference in doctrine between the Greek and Latin churches, though the controversy over it did not break out till the middle of the ninth century under patriarch Photius, (867). \footnote{Comp. on this Controversy J. G. Walch: Historia Controversiae Graecorum Latinorumque de Processione Spir. S., Jen. 1751. Also John Mason Neale; A History of the Holy Eastern Church, Lond. 1850, vol. i. 1093. A. P. Stanley (Eastern Church, p. 142) calls this dispute which once raged so long and so violently, “an excellent specimen of the race of extinct controversies.”} Dr. Waterland briefly sums up the points of dispute thus:\footnote{Works, vol. iii. p. 237 f.} “The Greeks and Latins have had many and tedious disputes about the procession. One thing is observable, that though the ancients, appealed to by both parties, have often said that the Holy Ghost proceeds from the Father, without mentioning the Son, yet they never said that he proceeded from the Father alone; so that the modern Greeks have certainly innovated in that article in expression at least, if not in real sense and meaning. As to the Latins, they have this to plead, that none of the ancients ever condemned their doctrine; that many of them have expressly asserted it; that the oriental churches themselves rather condemn their taking upon them to add anything to a creed formed in a general council, than the doctrine itself; that those Greek churches that charge their doctrine as heresy, yet are forced to admit much the same thing, only in different words; and that Scripture itself is plain, that the Holy Ghost proceeds at least by the Son, if not from him; which yet amounts to the same thing.”

This doctrinal difference between the Greek and the Latin Church, however insignificant it may appear at first sight, is characteristic of both, and illustrates the contrast between the conservative and stationary theology of the East, after the great ecumenical councils, and the progressive and systematizing theology of the West. The wisdom of changing an ancient and generally received formula of faith may be questioned. It must be admitted, indeed, that the Nicene Creed has undergone several other changes which were embodied in the Constantinopolitan Creed, and adopted by the Greeks as well as the Latins. But in the case of the Filioque, the Eastern Church which made the Nicene Creed, was never consulted, and when the addition was first brought to the notice of the bishop of Rome by Charlemagne, he protested against the innovation. His successors acquiesced in it, and the Protestant churches accepted the Nicene Creed with the Filioque, though without investigation. The Greek Church has ever protested against it since the time of Photius, and will never adopt it. She makes a sharp distinction between the procession, which is an eternal and internal process in the Holy Trinity itself, and the mission, of the Spirit, which is an act of revelation in time. The Spirit eternally proceeds from the Father alone (though through the Son); but was sent by the Father and the Son on the day of Pentecost. Hence the present tense is used of the former (John 15:26), and the future of the latter (14:26; 15:26). The Greek Church is concerned for the dignity and sovereignty of the Father, as the only source and root of the Deity. The Latin Church is concerned for the dignity of the Son, as being of one substance with the Father, and infers the double procession from the double mission.

\xsection{The Athanasian Creed}

§ 132. The Athanasian Creed.

G. Joh. Voss (Reform): De tribus symbolis, diss. ii. 1642, and in his Opera Omnia, Amstel. 1701 (forming an epoch in critical investigation). Archbishop Usher: De symbolis. 1647. J. H. Heidegger (Ref.): De symbolo Athanasiano. Zür. 1680. Em. Tentzel (Luth.): Judicia eruditoram de Symb. Athan. studiose collecta. Goth. 1687. Montfaucon (R.C.): Diatribe in Symbolum Quicunque, in the Benedictine ed. of the Opera Athanasii, Par. 1698, tom. ii. pp. 719–735. Dan. Waterland (Anglican): A Critical History of the Athanasian Creed. Cambridge, 1724, sec. ed. 1728 (in Waterland’s Works, ed. Mildert, vol. iii. pp. 97–270, Oxf. 1843). Dom. M. Speroni (R.C.): De symbolo vulgo S. Athanasii. Dias. i. and ii. Patav. 1750–’51. E. Köllner (Luth.): Symbolik aller christl. Confessionen. Hamb. vol. i. 1837, pp. 53–92. W. W. Harvey (Angl.): The History and Theology of the Three Creeds. Lond. 1854, vol. ii. pp. 541–695. Ph. Schaff: The Athanasian Creed, in the Am. Theolog. Review, New York, 1866, pp. 584–625. (Comp. the earlier literature, in chronological order, in Waterland, l. c. p. 108 ff., and in Köllner).

[Comp. here the notes in Appendix, p. 1034, and the later and fuller treatment in Schaff: Creeds of Christendom, N. York, 4th ed., 1884, vol. i. 34–42; vol. ii. 66–72, with the facsimile of the oldest MS. of the Athan. Creed in the Utrecht Psalter, ii. 555 sq. The rediscovery of that MS. in 1873 occasioned a more thorough critical investigation of the whole subject, with the result that the Utrecht Psalter dates from the ninth century, and that there is no evidence that the pseudo-Athanasian Creed, in its present complete form, existed before the age of Charlemagne. The statements in this section which assume an earlier origin, must be modified accordingly. Added 1889.]

The post-Nicene or Augustinian doctrine of the Trinity reached its classic statement in the third and last of the ecumenical confessions, called the Symbolum Athanasianum, or, as it is also named from its initial words, the Symbolum Quicumque; beyond which the orthodox development of the doctrine in the Roman and Evangelical churches to this day has made no advance.\footnote{In striking contrast with this unquestionable historical eminence of this Creed is Baur’s slighting treatment of it in his work of three volumes on the history of the doctrine of the Trinity, where he disposes of it in a brief note, vol. ii. p. 33, as a vain attempt to vindicate by logical categories the harsh and irreconcilable antagonism of unity and triad.} This Creed is unsurpassed as a masterpiece of logical clearness, rigor, and precision; and so far as it is possible at all to state in limited dialectic form, and to protect against heresy, the inexhaustible depths of a mystery of faith into which the angels desire to look, this liturgical theological confession achieves the task. We give it here in full, anticipating the results of the Christological controversies; and we append parallel passages from Augustine and other older writers, which the unknown author has used, in some cases word for word, and has woven with great dexterity into an organic whole.\footnote{In the Latin text we follow chiefly the careful revision of Waterland, ch. ix. (Works, vol. iii. p. 221 ff.), who also adds the various readings of the best manuscripts, and several parallel passages from the church fathers previous to 430, as he pushes the composition back before the third ecumenical council (431). We have also compared the text of Montfaucon (in his edition of Athanasius) and of Walch (Christl. Concordienbuch, 1750). The numbering of verses differs after ver. 19. Waterland puts vers. 19 and 20 in one, also vers. 25 and 26, 89 and 40, 41 and 42, making only forty verses in all. So Montfaucon, p. 735 ff. Walch makes forty-four verses.}

1. Quicumque vult salvus esse, ante omnia opus est, ut teneat catholicam fidem.\footnote{Comp. Augustine, Contra Maximin. Arian. l. ii. c. 3 (Opera, tom. viii. f. 729, ed. Venet.): “Haec est fides nostra, quoniam haec est fides recta, quae etiam catholica nuncupatur.”}

1. Whosoever will be saved, before all things it is necessary that he hold the catholic [true Christian] faith

2. Quam nisi quisque integram inviolatamque\footnote{Some manuscripts: “inviolabilemque.”} servaverit, absque dubio\footnote{“Absque dubio is wanting in the Cod. reg. Paris., according to Waterland.} in aeternum peribit.

2. Which faith except every one do keep whole and undefiled, without doubt he shall perish everlastingly.

3. Fides autem catholica haec est, ut unum Deum in trinitate et trinitatem in unitate veneremur;\footnote{Gregory Naz. Orat. xxiii. p. 422: ... \textgreek{μονάδα ἐν τριάδι, καὶ τριάδα ἐν μονάδι προσκύνουμένην.}}

3. But this is the catholic faith: That we worship one God in trinity, and trinity in unity;

4. Neque confundentes personas; neque substantiam separantes.\footnote{A similar sentence occurs in two places in the Commonitorium of Vincentius of Lerinum († 450): “Ecclesia vero \emph{catholica unam} divinitatem \emph{in trinitatis} plenitudine et \emph{trinitatis} aequalitatem, \emph{in una} atque eadem majestate \emph{veneratur}, ut \emph{neque} singularitas \emph{substantiae personarum} confundat proprietatem, \emph{neque} item trinitatis distinctio unitatem \emph{separet} deitatis ” (cap. 18 and 22). See the comparative tables in Montfaucon in Opera Athan. tom. ii. p. 725 sq. From this and two other parallels Anthelmi (Disquisitio de Symb. Athan., Par. 1698) has inferred that Vincentius of Lerinum was the author of the Athanasian Creed. But such arguments point much more strongly to Augustine, who affords many more parallels, and from whom Vincentius drew.}

4. Neither confounding the persons; nor dividing the substance.

5. Alia est enim persona Patris: alia Filii: alia Spiritus Sancti.\footnote{Vincentius Lirl.c. cap. 19: ”\emph{Alia est persona Patris, alia Patris, alia Spiritus Sancti. Sed Patris et Filii et Spiritus S. non alia et alia, sed una eademque natura}.” A similar passage is quoted by Waterland from the Symbolum Pelagii.}

5. For there is one person of the Father: another of the Son: another of the Holy Ghost.

6. Sed Patris et Filii et Spiritus Sancti una est divinitas: aequalis gloria, coaeterna majestas.\footnote{Augustine, tom. viii. p. 744 (ed. Venet.): ”\emph{Patris et Filii et Spiritus Sancti unam virtutem, unam substantiam, unam deitatem, unam majestatem, unam gloriam}.”}

6. But the Godhead of the Father, and of the Son, and of the Holy Ghost is all one: the glory equal, the majesty coëternal.

7. Qualis Pater, talis Filius, talis (et) Spiritus Sanctus.\footnote{Faustini Fid. (cited by Waterland): “\emph{Qualis} est \emph{Pater} secundum substantiam, \emph{talem} genuit \emph{Filium},” etc.}

7. Such as the Father is, such is the Son, and such is the Holy Ghost.

8. Increatus Pater: increatus Filius: increatus (et) Spiritus Sanctus.

8. The Father is uncreated: the Son is uncreated: the Holy Ghost is uncreated.

9. Immensus Pater: immensus Filius: immensus Spiritus Sanctus.\footnote{So Augustine, except that he has \emph{magnus} for immensus. Comp. below. \emph{lmmensus} is differently translated in the different Greek copies: (\textgreek{ἀκατάληπτος. ἄπειρος ,}and \textgreek{ἄμετρος},—a proof that the original is Latin. Venantius Fortunatus, in his Expositio fidei Catholicae, asserts: “Non est mensurabuis in sua natura, quia illocaus est, incircumscriptus, ubique totus, ubique praesens, ubique potens.” The word is thus quite equivalent to omnipresent. The translation ” incomprehensible” in the Anglican Book of Common Prayer is inaccurate, and probably came from the Greek translation \textgreek{ἀκατάληπτος}.}

9. The Father is immeasurable: the Son is immeasurable: the Holy Ghost is immeasurable.

10. Aeternus Pater: aeternus Filius: aeternus (et) Spiritus Sanctus.\footnote{Augustine, Op. tom. v. p. 543: ”\emph{Aeternus Pater, coaeternus Filius, coaeternus Spiritus Sanctus}.”}

10. The Father is eternal: the Son eternal: the Holy Ghost eternal.

11. Et tamen non tree aeterni: sed unus aeternus.

11. And yet there are not three eternals; but one eternal.

12. Sicut non tres increati: nec tres immensi: sed unus increatus et unus immensus.

12. As also there are not three uncreated: nor three immeasurable: but one uncreated, and one immeasurable.

13. Similiter omnipotens Pater: omnipotens Filius: omnipotens (et) Spiritus Sanctus.

13. So likewise the Father is almighty: the Son almighty: and the Holy Ghost almighty,

14. Et tamen non tres omnipo-entes; sed unus omnipotens.\footnote{In quite parallel terms Augustine, De trinit, lib. v. cap. 8 (tom. viii. 837 sq.); At “Magnus \emph{Pater}, magnus \emph{Filius}, magnus \emph{Spiritus S}., non tamen tres magni, \emph{sed unus} magnus ... Et bonus \emph{Pater}, bonus \emph{Filius}, bonus \emph{Spiritus S}.; nec \emph{tres} boni, \emph{sed unus} bonus; de quo dictum est, ’Nemo bonus nisi unus Deus.’ ... Itaque \emph{omnipotens Pater, omnipotens Filius, omnipotens Spiritus S}.; nec \emph{tamen tres omnipotentes, sed unus omnipotens}, ’ex quo omnia, per quem omnia, in quo omnia, ipsi gloria’ (Rom. ix. 36).”}

14. And yet there are not three almighties: but one almighty.

15. Ita Deus Pater: Deus Filius: Deus (et) Spiritus Sanctus.\footnote{Comp. Augustine, De trinit. lib. viii. in Prooem. to cap. 1 “Sicut \emph{Deus Pater, Deus Filius, Deus Spiritus S}.; et bonus P., bonus F., bonus Sp. S.; et omnipotens P., omnipotens F., omnipotens Sp. S.; \emph{nec tamen tres} Dii, aut tres boni, aut tree omnipotentes; \emph{sed unus Deus}, bonus, omnipotens, ipse Trinitas.”—Serm. 215 (Opera, tom. v. p. 948): “Unus Pater Deus, unus Filius Deus, unus Spiritus S. Deus: nec tamen Pater et F. et Sp. S. \emph{tres Dii}, sed \emph{unus Deus}.” De trinit. x. c. 11 (§18); “Haec igitur tria, memoria, intelligentia, voluntas, quoniam non sunt tres vitae, sed una vita; nec tres mentes, sed una mens; consequenter utique nec tres substantiae sunt, sed una substantia.” Comp. also Ambrosius, De Spiritu S. iii. 111: “Ergo sanctus Pater, sanctus Filius, santus et Spiritus; sed non tres sancti; quia unus est Deus sanctus, unus est Dominus;” and similar places.}

15. So the Father is God: the Son is God: and the Holy Ghost is God.

16. Et tamen non tres Dii; sed unus est Deus.\footnote{Comp. the above passage from Augustine, and De trinit l.c. 5 (al. 8): “Et tamen hanc trinitatem non tres Deos, sed unum Deum.” A similar passage in Vigilius of Tapsus, De trinitate, and in a sermon of Caesarius of Arles, which is ascribed to Augustine(v. 399).}

16. And yet there are not three Gods; but one God.

17. Ita Dominus Pater: Dominus Filius: Dominus (et) Spiritus Sanctus.

17. So the Father is Lord: the Son Lord: and the Holy Ghost Lord.

18. Et tamen non tres Domini; sed unus est Dominus.\footnote{Augustine: “Non tamen sunt duo Dii et duo Domini secundum formam Dei, sed ambo cum Spiritu suo unus est Dominus ... sed simul omnes \emph{non tres Dominos} esse Deos, sed \emph{unum Dominum} Deum dico.” Contra Maximin. Arian. 1. ii. c. 2 and 8 (Opera, viii. f. 729).}

18. And yet not three Lords; but one Lord

19. Quia sicut singulatim unamquamque personam et

19. For like as we are compelled by the Christian verity to Deum et Dominum confiteri christiana veritate compellimur:\footnote{1 Others read: “Deum \emph{ac} Dominum.”} acknowledge every Person by himself to be God and Lord

20. Ita tres Deos, aut (tres)\footnote{Waterland omits \emph{tres}, Walch has it.} Dominos dicere catholica religione prohibemur.

20. So are we forbidden by the catholic religion to say, there are three Gods, or three Lords.

21. Pater a nullo est factus; nec creatus; nec genitus.

21. The Father is made of none; neither created; nor begotten.

22. Filius a Patre solo est: \footnote{\emph{Solo} is intended to distinguish the Son from the Holy Ghost, who is of the Father \emph{and of the Son}; thus containing already the Latin doctrine of the double procession. Hence some Greek copies strike out \emph{alone}, while others inconsistently retain it.} non factus; nec creatus; sed genitus.

22. The Son is of the Father alone: not made; nor created; but begotten.

23. Spiritus Sanctus a Patre et Filio: non factus; nec creatus; nec genitus (est); sed procedens.\footnote{This is manifestly the Latin doctrine of the processio, which would be still more plainly expressed if it were said: “sed ab \emph{utroque} procedens.” Comp. Augustine, De trinit. lib. xv. cap. 26 (§ 47): “Non igitur ab utroque est \emph{genitus, sed procedit} ab utroque amborum \emph{Spiritus}.” Most Greek copies (comp. in Moutfaucon in Athan. Opera, tom. ii. p. 728 sqq.) omit \emph{et Filio}, and read only \textgreek{ἀπὸ τοῦ πατρός}.}

23. The Holy Ghost is of the Father and the Son: not made; neither created; nor begotten; but proceeding.

24. Unus ergo Pater, non tres Patres: unus Filius, non tres Filii: unus Spiritus Sanctus, non tres Spiritus Sancti.\footnote{Augustine, Contra Maxim. ii. 3 (tom. viii. f. 729): “In Trinitate quae Deus est, \emph{unus} est \emph{Pater}, \emph{non} duo vel \emph{tres}; et \emph{unus Filius, non} duo vel \emph{tres}; et \emph{unus} amborum \emph{Spiritus}, \emph{non} duo vel \emph{tres}.”}

24. Thus there is one Father, not three Fathers: one Son, not three Sons: one Holy Ghost, not three Holy Ghosts.

25. Et in hac trinitate nihil prius, aut posterius: nihil maius, aut minus.\footnote{August. Serm. 215, tom. v. f. 948: ”\emph{In hac trinitate non est aliud alio majus aut minus, nulla operum separatio, nulla dissimilitude substantiae}.” Waterland quotes also a kindred passage from the Symb. Pelagii.}

25. And in this Trinity none is before or after another: none is greater or less than another.

26. Sed totae tres personae coaeternae sibi sunt et coaequales.

26. But the whole three Persons are co-eternal together, and co-equal

27. Ita, ut per omnia, sicut jam supra dictum est, et unitas in trinitate et trinitas in unitate veneranda sit.\footnote{So Waterland and the Anglican Liturgy. The Lutheran Book of Concord reverses the order, and reads: trinitas in unitate, et unitas in trinitate.}

27. So that in all things, as aforesaid, the Unity in Trinity, and the Trinity in Unity is to be worshipped.

28. Qui vult ergo salvos esse, ita de trinitate sentiat.

28. He therefore that will be saved, must thus think of the Trinity.

The origin of this remarkable production is veiled in mysterious darkness. Like the Apostles’ Creed, it is not so much the work of any one person, as the production of the spirit of the church. As the Apostles’ Creed represents the faith of the ante-Nicene period, and the Nicene Creed the faith of the Nicene, so the Athanasian Creed gives formal expression to the post-Nicene faith in the mystery of the Trinity and the incarnation of God. The old tradition which, since the eighth century, has attributed it to Athanasius as the great champion of the orthodox doctrine of the Trinity, has been long ago abandoned on all hands; for in the writings of Athanasius and his contemporaries, and even in the acts of the third and fourth ecumenical councils, no trace of it is to be found.\footnote{Ger. Vossius first demonstrated the spuriousness of the tradition in his decisive treatise of 1642. Even Roman divines, like Quesnel, Dupin, Pagi, Tillemont, Montfaucon, and Muratori, admit the spuriousness. Köllner adduces nineteen proofs against the Athanasian origin of the Creed, two or three of which are perfectly sufficient without the rest. Comp. the most important in my treatise, l.c. p. 592 ff.} It does not appear at all in the Greek church till the eleventh or twelfth century; and then it occurs in a few manuscripts which bear the manifest character of translations, vary from one another in several points, and omit or modify the clause on the procession of the Holy Ghost from the Father and the Son (v. 23).\footnote{Wherever the creed has come into use in the Greek churches, this verse has been omitted as a Latin interpolation.} It implies the entire post-Nicene or Augustinian development of the doctrine of the Trinity, and even the Christological discussions of the fifth century, though it does not contain the anti-Nestorian test-word θεοτόκος, mother of God. It takes several passages verbally from Augustine’s work on the Trinity which was not completed till the year 415, and from the Commonitorium of Vincentius of Lerinum, 434; works which evidently do not quote the passages from an already existing symbol, but contribute them as stones to the building. On the other hand it contains no allusion to the Monophysite and Monothelite controversies, and cannot be placed later than the year 570; for at that date Venantius Fortunatus of Poictiers wrote a short commentary on it.

It probably originated about the middle of the fifth century, in the school of Augustine, and in Gaul, where it makes its first appearance, and acquires its first ecclesiastical authority. But the precise author or compiler cannot be discovered, and the various views of scholars concerning him are mere opinions.\footnote{Comp. the catalogue of opinions in Waterland, vol. iii. p. 117; in Köllner; and in my own treatise. The majority of voices have spoken in favor of Vigilius of Tapsus in Africa, a.d.484; others for Vincentius of Lerinum, 434; Waterland for Hilaryof Arles, about 430; while others ascribe it indefinitely to the North African, or Gallic, or Spanish church in the sixth or seventh century. Harvey recently, but quite groundlessly, has dated the composition back to the year 401, and claims it for the bishop Victricius of Rouen (Hist. and Theol. of the Three Creeds, vol. ii. p. 583 f.). He thinks that Augustinequotes from it, but this father nowhere alludes to such a symbol; the author of the Creed, on the contrary, has taken several passages from Augustine, De Trinitate, as well as from Vincentius of Lerinum and other source. Comp. the notes to the Creed above, and my treatise, p. 596 ff.} From Gaul the authority of this symbol spread over the whole of Latin Christendom, and subsequently made its way into some portions of the Greek church in Europe. The various Protestant churches have either formally adopted the Athanasian Creed together with the Nicene and the Apostles’, or at all events agree, in their symbolical books, with its doctrine of the trinity and the person of Christ.\footnote{On this agreement of the symbolical books of the Evangelical churches with the Athanasianum, comp. my treatise, l.c. p. 610 ff. Luther considers this Creed the weightiest and grandest production of the church since the time of the Apostles. In the Church of England it is still sung or chanted in the cathedrals. The Protestant Episcopal church in the United States, on the contrary, has excluded it from the Book of Common Prayer.}

The Athanasian Creed presents, in short, sententious articles, and in bold antitheses, the church doctrine of the Trinity in opposition to Unitarianism and tritheism, and the doctrine of the incarnation and the divine-human person of Christ in opposition to Nestorianism and Eutychianism, and thus clearly and concisely sums up the results of the trinitarian and Christological controversies of the ancient church. It teaches the numerical unity of substance and the triad of persons in the Father and the Son and the Holy Ghost, with the perfect deity and perfect humanity of Christ in one indivisible person. In the former case we have one substance or nature in three persons; in the latter, two natures in one divine-human person.

On this faith eternal salvation is made to depend. By the damnatory clauses in its prologue and epilogue the Athanasianum has given offence even to those who agree with its contents. But the original Nicene Creed contained likewise an anathema, which afterwards dropped out of it; the anathema is to be referred to the heresies, and may not be applied to particular persons, whose judge is God alone; and finally, the whole intention is, not that salvation and perdition depend on the acceptance and rejection of any theological formulary or human conception and exhibition of the truth, but that faith in the revealed truth itself, in the living God, Father, Son, and Spirit, and in Jesus Christ the God-Man and the Saviour of the world, is the thing which saves, even where the understanding may be very defective, and that unbelief is the thing which condemns; according to the declaration of the Lord: “He that believeth and is baptized shall be saved; but he that believeth not shall be damned.” In particular actual cases Christian humility and charity of course require the greatest caution, and leave the judgment to the all-knowing and just God.

The Athanasian Creed closes the succession of ecumenical symbols; symbols which are acknowledged by the entire orthodox Christian world, except that Evangelical Protestantism ascribes to them not an absolute, but only a relative authority, and reserves the right of freely investigating and further developing all church doctrines from the inexhaustible fountain of the infallible word of God.

II. The Origenistic Controversies.

I. Epiphanius: Haeres. 64. Several Epistles of Epiphanius, Theophilus of Alex., and Jerome (in Jerome’s Epp. 51 and 87–100, ed. Vallarsi). The controversial works of Jerome and Rufinus on the orthodoxy of Origen (Rufini Praefatio ad Orig. περὶ ἀρχῶν; and Apologia s. invectivarum in Hieron.; Hieronymi Ep. 84 ad Pammachium et Oceanum de erroribus Origenis; Apologia Adv. Rufinum libri iii, written 402–403, etc.). Palladius: Vita Johannis Chrysostomi (in Chrysost. Opera, vol. xiii. ed. Montfaucon). Socrates: H. E. vi. 3–18. Sozomenus: H. E. viii. 2–20. Theodoret: H. E. v. 27 sqq. Photius: Biblioth. Cod. 59. Mansi:

II. Huetius: Origeniana (Opera Orig. vol. iv. ed. De la Rue). Doucin: Histoire des mouvements arrivés dans l’église au sujet d’Origène. Par. 1700. Walch: Historie der Ketzereien. Th. vii. p. 427 sqq. Schröckh: Kirchengeschichte, vol. x. 108 sqq. Comp. the monographs Of Redepenning and Thomasius on Origen; and Neander: Der heil. Joh. Chrysostomus. Berl. 1848, 3d ed. vol. ii. p. 121 sqq. Hefele (R.C.): Origenistenstreit, in the Kirchenlexicon of Wetzer and Welte, vol. vii. p. 847 sqq., and Conciliengeschichte, vol. ii. p. 76 sqq. O. Zöckler: Hieronymus. Gotha, 1865, p. 238 ff; 391 ff.

\xsection{The Orgenistic Controversy in Palestine. Epiphanius, Rufinus, and Jerome, a.d. 394-399}

§133. The Orgenistic Controversy in Palestine. Epiphanius, Rufinus, and Jerome, a.d. 394–399.

Between the Arian and the Nestorian controversies and in indirect connection with the former, come the vehement and petty personal quarrels over the orthodoxy of Origen, which brought no gain, indeed, to the development of the church doctrine, yet which have a bearing upon the history of theology, as showing the progress of orthodoxy under the twofold aspect of earnest zeal for the pure faith, and a narrow-minded intolerance towards all free speculation. The condemnation of Origen was a death blow to theological science in the Greek church, and left it to stiffen gradually into a mechanical traditionalism and formalism. We shall confine ourselves, if possible, to the points of general interest, and omit the extremely insipid and humiliating details of personal invective and calumny.

It is the privilege of great pioneering minds to set a mass of other minds in motion, to awaken passionate sympathy and antipathy, and to act with stimulating and moulding power even upon after generations. Their very errors are often more useful than the merely traditional orthodoxy of unthinking men, because they come from an honest search after truth, and provoke new investigation. One of these minds was Origen, the most learned and able divine of the ante-Nicene period, the Plato or the Schleiermacher of the Greek church. During his life-time his peculiar, and for the most part Platonizing, views already aroused contradiction, and to the advanced orthodoxy of a later time they could not but appear as dangerous heresies. Methodius of Tyre († 311) first attacked his doctrines of the creation and the resurrection; while Paulphilus († 309), from his prison, wrote an apology for Origen, which Eusebius afterwards completed. His name was drawn into the Arian controversies, and used and abused by both parties for their own ends. The question of the orthodoxy of the great departed became in this way a vital issue of the day, and rose in interest with the growing zeal for pure doctrine and the growing horror of all heresy.

Upon this question three parties arose: free, progressive disciples, blind adherents, and blind opponents.\footnote{Similar parties have arisen with reference to Luther, Schleiermacher, and other great theologians and philosophers.}

1. The true, independent followers of Origen drew from his writings much instruction and quickening, without committing themselves to his words, and, advancing with the demands of the time, attained a clearer knowledge of the specific doctrines of Christianity than Origen himself, without thereby losing esteem for his memory and his eminent services. Such men were Pamphilus, Eusebius of Caesarea, Didymus of Alexandria, and in a wider sense Athanasius, Basil the Great, Gregory of Nazianzum, and Gregory of Nyssa; and among the Latin fathers, Hilary, and at first Jerome, who afterwards joined the opponents. Gregory of Nyssa, and perhaps also Didymus, even adhered to Origen’s doctrine of the final salvation of all created intelligences.

2. The blind and slavish followers, incapable of comprehending the free spirit of Origen, clave to the letter, held all his immature and erratic views, laid greater stress on them than Origen himself, and pressed them to extremes. Such mechanical fidelity to a master is always apostasy to his spirit, which tends towards continual growth in knowledge. To this class belonged the Egyptian monks in the Nitrian mountains; four in particular: Dioscurus, Ammonius, Eusebius, and Enthymius, who are known by the name of the “tall brethren,”\footnote{\textgreek{Ἀδελφοὶ μακροί}, on account of their bodily size.} and were very learned.

3. The opponents of Origen, some from ignorance, others from narrowness and want of discrimination, shunned his speculations as a source of the most dangerous heresies, and in him condemned at the same time all free theological discussion, without which no progress in knowledge is possible, and without which even the Nicene dogma would never have come into existence. To these belonged a class of Egyptian monks in the Scetic desert, with Pachomius at their head, who, in opposition to the mysticism and spiritualism of the Origenistic monks of Nitria, urged grossly sensuous views of divine things, so as to receive the name of Anthropomorphites. The Roman church, in which Origen was scarcely known by name before the Arian disputes, shared in a general way the strong prejudice against him as an unsound and dangerous writer.

The leader in the crusade against the bones of Origen was the bishop Epiphanius of Salamis (Constantia) in Cyprus († 403), an honest, well-meaning, and by his contemporaries highly respected, but violent, coarse, contracted, and bigoted monastic saint and heresy hunter. He had inherited from the monks in the deserts of Egypt an ardent hatred of Origen as an arch-heretic, and for this hatred he gave documentary justification from the numerous writings of Origen in his Panarion, or chest of antidotes for eighty heresies, in which he branded him as the father of Arianism and many other errors.\footnote{Haeer. 64. Compare also his Epistle to bishop John of Jerusalem, written 394 and translated by Jeromeinto Latin (Ep. 51, ed. Vallarsi), where he enumerates eight heresies of Origen relating to the trinity, the doctrine of man, of angels, of the world, and the last things.} Not content with this, he also endeavored by journeying and oral discourse to destroy everywhere the influence of the long departed teacher of Alexandria, and considered himself as doing God and the church the greatest service thereby.

With this object the aged bishop journeyed in 394 to Palestine, where Origen was still held in the highest consideration, especially with John, bishop of Jerusalem, and with the learned monks Rufinus and Jerome, the former of whom was at that time in Jerusalem and the latter in Bethlehem. He delivered a blustering sermon in Jerusalem, excited laughter, and vehemently demanded the condemnation of Origen. John and Rufinus resisted; but Jerome, who had previously considered Origen the greatest church teacher after the apostles, and had learned much from his exegetical writings, without adopting his doctrinal errors, yielded to a solicitude for the fame of his own orthodoxy, passed over to the opposition, broke off church fellowship with John, and involved himself in a most violent literary contest with his former friend Rufinus; which belongs to the chronique scandaleuse of theology. The schism was terminated indeed by the mediation of the patriarch Theophilus in 397, but the dispute broke out afresh. Jerome condemned in Origen particularly his doctrine of pre-existence, of the final conversion of the devils, and of demons, and his spiritualistic sublimation of the resurrection of the body; while Rufinus, having returned to the West (398), translated several works of Origen into Latin, and accommodated them to orthodox taste. Both were in fact equally zealous to defend themselves against the charge of Origenism, and to fasten it upon each other, and this not by a critical analysis and calm investigation of the teachings of Origen, but by personal denunciations and miserable invectives.\footnote{2 Comp. the description of their conduct by Zöckler, Hieronymus, p. 396 ff.}

Rufinus was cited before pope Anastasius (398–402), who condemned Origen in a Roman synod; but he sent a satisfactory defense and found an asylum in Aquileia. He enjoyed the esteem of such men as Paulinus of Nola and Augustine, and died in Sicily (410).

\xsection{The Origenistic Controversy in Egypt and Constantinople. Theophilus and Chrysostom a.d. 399-407}

§ 134. The Origenistic Controversy in Egypt and Constantinople. Theophilus and Chrysostom a.d. 399–407.

Meanwhile a second act of this controversy was opened in Egypt, in which the unprincipled, ambitious, and intriguing bishop Theophilus of Alexandria plays the leading part. This bishop was at first an admirer of Origen, and despised the anthropomorphite monks, but afterwards, through a personal quarrel with Isidore and the “four tall brethren,” who refused to deliver the church funds into his hands, he became an opponent of Origen, attacked his errors in several documents (399–403),\footnote{In his Epistola Synodica ad episcopos Palaestinos et ad Cyprios, 400, and in three successive Epistolae Paschales, from 401-403, all translated by Jeromeand forming Epp. 92, 96, 98, and l00 of his Epistles, according to the order of Vallarsi. They enter more deeply into the topics of the controversy than Jerome’s own writings against Origen. Jerome(Ep. 99 ad Theophilum) pays him the compliment: “Rhetoricae eloquentiae jungis philosophos, et Demosthenem atque Platonem nobis consocias.”} and pronounced an anathema on his memory, in which he was supported by Epiphanius, Jerome, and the Roman bishop Anastasius. At the same time he indulged in the most violent measures against the Origenistic, monks, and banished them from Egypt. Most of these monks fled to Palestine; but some: fifty, among whom were the four tall brethren, went to Constantinople, and found there a cordial welcome with the bishop John Chrysostom in 401.

In this way that noble man became involved in the dispute. As an adherent of the Antiochian school, and as a practical theologian, he had no sympathy with the philosophical speculation of Origen, but he knew how to appreciate his merits in the exposition of the Scriptures, and was impelled by Christian love and justice to intercede with Theophilus in behalf of the persecuted monks, though he did not admit them to the holy communion till they proved their innocence.

Theophilus now set every instrument in motion to overthrow the long envied Chrysostom, and employed even Epiphanius, then almost an octogenarian, as a tool of his hierarchical plans. This old man journeyed in mid-winter in 402 to Constantinople, in the imagination that by his very presence he would be able to destroy the thousand-headed hydra of heresy, and he would neither hold church fellowship with Chrysostom, who assembled the whole clergy of the city to greet him, nor pray for the dying son of the emperor, until all Origenistic heretics should be banished from the capital, and he might publish the anathema from the altar. But he found that injustice was done to the Nitrian monks, and soon took ship again to Cyprus, saying to the bishops who accompanied him to the sea shore: “I leave to you the city, the palace, and hypocrisy; but I go, for I must make great haste.” He died on the ship in the summer of 403.

What the honest coarseness of Epiphanius failed to effect, was accomplished by the cunning of Theophilus, who now himself travelled to Constantinople, and immediately appeared as accuser and judge. He well knew how to use the dissatisfaction of the clergy, of the empress Eudoxia, and of the court with Chrysostom on account of his moral severity and his bold denunciations.\footnote{According to Socrates (H.E. vi. 4) another special reason for the disaffection was, that Chrysostomalways ate alone, and never accepted an invitation to a banquet, either on account of dyspepsia or habitual abstemiousness. But by the people he was greatly esteemed and loved as a man and as a preacher.} In Chrysostom’s own diocese, on an estate “at the oak”\footnote{\textgreek{Πρὸς τὴν δρῦν,} Synodus ad Quercum. The estate belonged to the imperial prefect Rufinus, and had a palace, a large church, and a monastery. Sozomen, viii. 17.} in Chalcedon, he held a secret council of thirty-six bishops against Chrysostom, and there procured, upon false charges of immorality, unchurchly conduct, and high treason, his deposition and banishment in 403.\footnote{Among the twenty-nine charges were these: that Chrysostomcalled the Saint Epiphanius a fool and demon; that he wrote a book full of abuse of the clergy; that he received visits from females without witnesses; that he bathed alone and ate alone! See Hefele, ii. p. 78 sqq.} Chrysostom was recalled indeed in three days in consequence of an earthquake and the dissatisfaction of the people, but was again condemned by a council in 404, and banished from the court, because, incensed by the erection of a silver statue of Eudoxia close to the church of St. Sophia, and by the theatrical performances connected with it, he had with unwise and unjust exaggeration opened a sermon on Mark vi. 17 ff., in commemoration of John the Baptist with the personal allusion: “Again Herodias rages, again she raves, again she dances, and again she demands the head of John [this was Chrysostom’s own name] upon a charger.”\footnote{\textgreek{Πάλιν Ἡρωδίας μαίνεται, πάλιν ταράσσεται, πάλιν ὀρχεῖται, πάλιν ἐπὶ πίνακι τὴν κεφαλὴν τοῦ Ἰωάννου ζητεῖ λαβεῖν.}Comp. Socr. H. E. vi. 18. Eudoxia was a young and beautiful woman, who despised her husband, and indulged her passions. She died four years after the birth of her son Theodosius the Younger, whose true father is said to have been the comes John. Comp, Gibbon, ch. xxxii.} From his exile in Cucusus and Arabissus he corresponded with all parts of the Christian world, took lively interest in the missions in Persia and Scythia, and appealed to a general council. His opponents procured from Arcadius an order for his transportation to the remote desert of Pityus. On the way thither he died at Comana in Pontus, a.d. 407, in the sixtieth year of his age, praising God for everything, even for his unmerited persecutions.\footnote{\textgreek{Δόξα τῷ Θεῷ πάντων ἕνεκεν,} were his last words, the motto of his life and work.}

Chrysostom was venerated by the people as a saint, and thirty years after his death, by order of Theodosius II. (438), his bones were brought back in triumph to Constantinople, and deposited in the imperial tomb. The emperor himself met the remains at Chalcedon, fell down before the coffin, and in the name of his guilty parents, Arcadius and Eudoxia, implored the forgiveness of the holy man. The age could not indeed understand and appreciate the bold spirit of Origen, but was still accessible to the narrow piety of Epiphanius and the noble virtues of Chrysostom.

In spite of this prevailing aversion of the time to free speculation, Origen always retained many readers and admirers, especially among the monks in Palestine, two of whom, Domitian and Theodorus Askidas, came to favor and influence at the court of Justinian I. But under this emperor the dispute on the orthodoxy of Origen was renewed about the middle of the sixth century in connection with the controversy on the Three Chapters, and ended with the condemnation of fifteen propositions of Origen at a council in 544.\footnote{It was only a \textgreek{σύνοδος ἐνδημοῦσα}, \emph{i.e}., a council of the bishops just then in Constantinople, and is not to be confounded with the fifth \emph{ecumenical} council at Constantinople in 553, which decided only the controversy of the Three Chapters. Comp. Mansi, Conc. tom. ix. fol. 395-399 (where the fifteen canons are given); Walch, Ketzerhistorie, vii. 660; and Gieseler, K. Gesch. i. ii. p. 368.} Since then no one has ventured until recent times to raise his voice for Origen, and many of his works have perished.

With Cyril of Alexandria the theological productivity of the Greek church, and with Theodoret the exegetical, became almost extinct. The Greeks thenceforth contented themselves for the most part with revisions and collections of the older treasures. A church which no longer advances, goes backwards, or falls in stagnation.

III. The Christological Controversies.

Among the works on the whole field of the Christological controversies should be compared especially the already cited works of Petavius (tom. iv. De incarnatione Verbi), Walch (Ketzerhistorie, vol. v.-ix.), Baur, and Dorner. The special literature will be given at the heads of the several sections.

\xsection{General View. Alexandrian and Antiochian Schools}

§ 135. General View. Alexandrian and Antiochian Schools.

The Trinity and Christology, the two hardest problems and most comprehensive dogmas of theology, are intimately connected. Hence the settlement of the one was immediately followed by the agitation and study of the other. The speculations on the Trinity had their very origin in the study of the person of Christ, and led back to it again. The point of union is the idea of the incarnation of God. But in the Arian controversy the Son of God was viewed mainly in his essential, pre-mundane relation to the Father; while in the Christological contest the incarnate historical Christ and the constitution of his divine-human person was the subject of dispute.

The notion of redemption, which forms the centre of Christian thinking, demands a Redeemer who unites in his person the nature of God and the nature of man, yet without confusion. In order to be a true Redeemer, the person must possess all divine attributes, and at the same time enter into all relations and conditions of mankind, to raise them to God. Four elements thus enter into the orthodox doctrine concerning Christ: He is true God; be is true man; he is one person; and the divine and human in him, with all the personal union and harmony, remain distinct.

The result of the Arian controversies was the general acknowledgment of the essential and eternal deity of Christ. Before the close of that controversy the true humanity of Christ at the same time came in again for treatment; the church having indeed always maintained it against the Gnostic Docetism, but now, against a partial denial by Apollinarianism, having to express it still more distinctly and lay stress on the reasonable soul. And now came into question, further, the relation between the divine and the human natures in Christ. Origen, who gave the impulse to the Arian controversy, had been also the first to provoke deeper speculation on the mystery of the person of Christ. But great obscurity and uncertainty had long prevailed in opinions on this great matter. The orthodox Christology is the result of powerful and passionate conflicts. It is remarkable that the notorious rabies theologorum has never in any doctrinal controversy so long and violently raged as in the controversies on the person of the Reconciler, and in later times on the love-feast of reconciliation.

The Alexandrian school of theology, with its characteristic speculative and mystical turn, favored a connection of the divine and human in the act of the incarnation so close, that it was in danger of losing the human in the divine, or at least of mixing it with the divine;\footnote{Even Athanasius is not wholly free from this leaning to the monophysite view, and speaks of an \textgreek{ἕνωσις φυσική} of the Logos with his flesh, and of one incarnate nature of the divine Logos, \textgreek{μία φύσις τοῦ Θεοῦ λόγου σεσαρκωμένη}, which with his flesh is to be worshipped; see his little tract De incarnatione Dei Verbi (\textgreek{περὶ τῆς σαρκώσεως τοῦ Θεοῦ λόγου}) in the 3d tom. of the Bened. ed. p. 1. But in the first place it must be considered that this tract (which is not to be confounded with his large work De incarnatione Verbi Dei, \textgreek{περὶ τῆς ἐνανθρωπήσεως τοῦ λόγου}, in the first tom. P. i. of the Bened. ed. pp. 47-97), is by many scholars (Montfaucon, Möhler, Hefele) denied to Athanasius, though on insufficient grounds; and further, that at that time \textgreek{φύσις,οὐσία}and \textgreek{ὑπόστασις}were often interchanged, and did not become sharply distinguished till towards the end of the Nicene age. “In the indefiniteness of the notions of \textgreek{φύσις} and \textgreek{ὑπόστασις},” says Neander (Dogmengeschichte, i. p. 340), “the Alexandrians were the more easily moved, for the sake of the one \textgreek{ὑπόστασις}, to concede also only one \textgreek{φύσις}in Christ, and set the \textgreek{ἕνωσις φυσική}against those who talked of two natures.” Comp. Petavius, De incarn. Verbi, lib. ii. c. 3 (tom. iv. p. 120, de vocabulis \textgreek{φύσεως}et \textgreek{ὑποστάσεως}); also the observations of Dorner, l. c. i. p. 1072, and of Hefele, Conciliengesch. ii. p. 128 f. The two Gregories speak, indeed, of \textgreek{δύο φύσεις}in Christ, yet at the same time of a \textgreek{σύγκρασις}, and \textgreek{ἀνάκρασις}, \emph{i.e., mingling} of the two.} while, conversely, the Antiochian or Syrian school, in which the sober intellect and reflection prevailed, inclined to the opposite extreme of an abstract separation of the two natures.\footnote{Theodore, bishop of Mopsuestia in Cilicia, the head of the Antiochian school, compares the union of the divine and human in Christ with the marriage union of man and woman, and says that one cannot conceive a complete nature without a complete person (\textgreek{ὑπόστασις}). Comp. Neander, l.c. i. p. 343; Dorner, ii. p. 39 ff.; Fritzsche: De Theodori Mopsvest. vita et scriptis, Halae, 1837, and an article by W. Möller in Herzog’s Encycl. vol. xv. P. 715 ff. Of the works of Theodore of Mopsuestia we have only fragments, chiefly in the acts of the fifth ecumenical council (in Mansi, Conc. tom. ii. fol. 203 sqq.), and a commentary on the twelve Prophets, which cardinal Angelo Mai discovered, and edited in 1854 at Rome in his Nova Bibliotheca SS. Patrum, tom. vii. Pars i. pp. 1-408, together with some fragments of commentaries on New Testament books, edited by Fritzsche, jun., Turici, 1847; and by Pitra in Spicileg. Solesm. tom. i. Par. 1852.} In both cases the mystery of the incarnation, the veritable and permanent union of the divine and human in the one person of Christ, which is essential to the idea of a Redeemer and Mediator, is more or less weakened or altered. In the former case the incarnation becomes a transmutation or mixture (σύγκρασις) of the divine and human; in the latter, a mere indwelling (ἐνοίκησις) of the Logos in the man, or a moral union (συνάφεια) of the two natures, or rather of the two persons.

It was now the problem of the church, in opposition to both these extremes, to assert the personal unity and the distinction of the two natures in Christ with equal solicitude and precision. This she did through the Christological controversies which agitated the Greek church for more than two hundred years with extraordinary violence. The Roman church, though in general much more calm, took an equally deep interest in this work by some of its more eminent leaders, and twice decided the victory of orthodoxy, at the fourth general council and at the sixth, by the powerful influence of the bishop of Rome.

We must distinguish in this long drama five acts:

1. The Apollinarian controversy, which comes in the close of the Nicene age, and is concerned with the full humanity of Christ, that is, the question whether Christ, with his human body and human soul (anima animans), assumed also a human spirit (νοῦς, πνεῦμα, anima rationalis).

2. The Nestorian controversy, down to the rejection of the doctrine of the double personality of Christ by the third ecumenical council of Ephesus, a.d. 431.

3. The Eutychian controversy, to the condemnation of the doctrine of one nature, or more exactly of the absorption of the human in the divine nature of Christ; to the fourth ecumenical council at Chalcedon, a.d. 451.

4. The Monophysite dispute; the partial reaction towards the Eutychian theory; down to the fifth general council at Constantinople a.d. 553.

5. The Monothelite controversy, a.d. 633–680, which terminated with the rejection of the doctrine of one will in Christ by the sixth general council at Constantinople in 680, and lies this side of our period.

\xsection{The Apollinarian Heresy, a.d. 362-381}

§ 136. The Apollinarian Heresy, a.d. 362–381.

Sources:

I. Apollinaris: Περὶ σαρκώσεως—ΠεριΠίστεως —Περὶ ἀναστάσεως—Κατὰ κεφαλειον,—and controversial works against Porphyry, and Eunomius, biblical commentaries, and epistles. Only fragments of these remain in the answers of Gregory of Nyassa and Theodoret, and in Angelo Mai: Nov. Biblioth. Patrum, tom. vii. (Rom. 1854), Pars secunda, pp. 82–91 (commentary on Ezekiel), in Leontinus Byzantinus, and in the Catenae, especially the Catena in Evang. Joh., ed. Corderius, 1630.

II. Against Apollinaris: Athanasius: Contra Apollinarium, libri ii. (Περὶ σαρκώσεωστοῦ Κυρίου ἡμῶν Ἰ. Χ. κατὰ Ἀπολλιναρίου, in Opera, tom. i. pars secunda, pp. 921–955, ed. Bened., and in Thilo’s Bibl. Patr. Gr. dogm., vol. i. pp. 862–937). This work was written about the year 372 against Apollinarianism in the wider sense, without naming, Apollinaris or his followers; so that the title above given is wanting in the oldest codices. Similar errors, though in like manner without direct reference to Apollinaris, and evading his most important tenet, were combated by Athanasius in the Epist. ad Epictetum episcopum Corinthi contra haereticos (Opp. i. ii. 900 sqq., and in Thilo, i. p. 820 sqq.), which is quoted even by Epiphanius. Gregory Of Nyssa: Λόγος ἀντιῤῥητικὸς πρὸς τὰ Ἀπολλιναρίου, first edited by L. A. Zacagni from the treasures of the Vatican library in the unfortunately incomplete Collectanea monumentorum veterum ecclesiae Graecae et Latinae, Romae, 1698, pp. 123–287, and then by Gallandi, Bibliotheca Vet. Patrum, tom. vi. pp. 517–577. Gregory Naz.: Epist. ad Nectarium, and Ep. i. and ii. ad Cledonium (or Orat. 46 and 51–52; comp. Ullmann’s Gregor v. Naz. p. 401 sqq.). Basilius M.: Epist. 265 (a.d. 377), in the new Bened. ed. of his Opera, Par. 1839, tom. iii. Pars ii. p. 591 sqq. Epiphanius: Haer. 77. Theodoret: Fabul. haer. iv. 8; v. 9; and Diolog. i.-Iii.

Literature.

Dion. Petavius: De incarnatione Verbi, lib. i. cap. 6 (in the fourth vol. of the Theologicorum dogmatum, pp. 24–34, ed. Par. 1650). Jac. Basnage: Dissert. de Hist. haer. Apollinar. Ultraj. 1687. C. W. F. Walch: l.c. iii. 119–229. Baur: l.c. vol. i. pp. 585–647. Dorner: l.c. i. pp. 974–1080. H. Voigt: Die Lehre des Athanasius, \&c. Bremen, 1861. Pp. 306–345.

Apollinaris,\footnote{The name is usually written \emph{Apollinaris}, even by Petavius, Baur, and Dorner, and by all English writers. We have no disposition to disturb the established usage in a matter of so little moment. But the Greek fathers always write \textgreek{Ἀπολλινάριος}, and hence \emph{Apollinarius} (as in Jerome, De viris illustr., c. 104) is more strictly correct.} bishop of Laodicea in Syria, was the first to apply the results of the trinitarian discussions of the Nicene age to Christology, and to introduce the long Christological controversies. He was the first to call the attention of the Church to the psychical and pneumatic side of the humanity of Christ, and by contradiction brought out the doctrine of a reasonable human soul in him more clearly and definitely than it had before been conceived.

Apollinaris, like his father (Apollinaris the Elder, who was a native of Alexandria, and a presbyter in Laodicea), was distinguished for piety, classical culture, a scholarly vindication of Christianity against Porphyry and the emperor Julian, and adhesion to the Nicene faith. He was highly esteemed, too, by Athanasius, who, perhaps through personal forbearance, never mentions him by name in his writings against his error.

But in his zeal for the true deity of Christ, and his fear of a double personality, he fell into the error of denying his integral humanity. Adopting the psychological trichotomy, he attributed to Christ a human body, and a human (animal) soul,\footnote{\textgreek{Σῶμα}.} but not a human spirit or reason;\footnote{\textgreek{Ψυχὴ ἄλογος}, the inward vitality which man has in common with animals.} putting the divine Logos in the place of the human spirit. In opposition to the idea of a mere connection of the Logos with the man Jesus, he wished to secure an organic unity of the two, and so a true incarnation; but he sought this at the expense of the most important constituent of man. He reaches only a θεὸς σαρκοφόρος, as Nestorianism only an ἄνθρωπος θεοφόρος, instead of the proper θεάνθρωπος. He appealed to the fact that the Scripture says, the word was made flesh—not spirit;\footnote{\textgreek{Νοῦς, πνεῦμα}, or the \textgreek{ψυχὴ λογική}, anima rationalis, the motive, self-active, free element, the \textgreek{αὐτοκίνητον}, the thinking and willing, immortal spirit, which distinguishes man from animals. Apollinaris followed the psychological trichotomy of Plato. \textgreek{Ὁ ἄνθρωπος}, says he in Gregory of Nyssa, \textgreek{εἷς ἐστιν ἐκ πνεύματος καὶ φυχῆς καί σώματος ,} for which he quotes 1 Thess. v. 23, and Gal. v. 17. But in another fragment he designates the whole spiritual principle in man by \textgreek{ψυχῆ}, and makes the place of it in Christ to be supplied by the Logos. Comp. the passages in Gieseler, vol. i. Div. ii. p. 73 (4th ed.). From this time the triple division of human nature was unjustly accounted heterodox.} God was manifest in the flesh, \&c.; to which Gregory Nazianzen justly replied that in these passages the term σάρξwas used by synecdoche for the whole human nature. In this way Apollinaris established so close a connection of the Logos with human flesh, that all the divine attributes were transferred to the human nature, and all the human attributes to the divine, and the two were merged in one nature in Christ. Hence he could speak of a crucifixion of the Logos, and a worship of his flesh. He made Christ a middle being between God and man, in whom, as it were, one part divine and two parts human were fused in the unity of a new nature.\footnote{He even ventured to adduce created analogies, such as the mule, midway between the horse and the ass; the grey color, a mixture of white and black; and spring in distinction from winter and summer. Christ says he, is \textgreek{οὔτε ἄνθρωπος ὅλος , οὔτε θεὸς, ἀλλὰ θεοῦ καὶ ἀνθρώπου μίξις}.}

Epiphanius expresses himself concerning the beginning of the controversy in these unusually lenient and respectful terms: “Some of our brethren, who are in high position, and who are held in great esteem with us and all the orthodox, have thought that the spirit (ὁ νοῦς) should be excluded from the manifestation of Christ in the flesh, and have preferred to hold that our Lord Christ assumed flesh and soul, but not our spirit, and therefore not a perfect man. The aged and venerable Apollinaris of Laodicea, dear even to the blessed father Athanasius, and in fact to all the orthodox has been the first to frame and promulgate this doctrine. At first, when some of his disciples communicated it to us, we were unwilling to believe that such a man would put this doctrine in circulation. We supposed that the disciples had not understood the deep thoughts of so learned and so discerning a man, and had themselves fabricated things which he did not teach,” \&c.

So early as 362, a council at Alexandria rejected this doctrine (though without naming the author), and asserted that Christ possessed a reasonable soul. But Apollinaris did not secede from the communion of the Church, and begin to form a sect of his own, till 375. He died in 390. His writings, except numerous fragments in the works of his opponents, are lost.

Apollinaris, therefore, taught the deity of Christ, but denied the completeness (τελειότης) of his humanity, and, taking his departure from the Nicene postulate of the homoousion ran into the Arian heresy, which likewise put the divine Logos in the place of the human spirit in Christ, but which asserted besides this the changeableness (τρεπτότης) of Christ; while Apollinaris, on the contrary, aimed to establish more firmly the unchangeableness of Christ, to beat the Arians with their own weapons, and provide a better vindication of the Nicene dogma. He held the union of full divinity with full humanity in one person, therefore, of two wholes in one whole, to be impossible.\footnote{The result of this construction he called \textgreek{ἀνθρωπόθεος}, a sort of monstrosity, which he put in the same category with the mythological figures of the minotaur, the well-known Cretan monster with human body and bull’s head, or the body of a bull and the head of a man. But the Apollinarian idea of the union of the Logos with a truncated human nature might be itself more justly compared with this monster.} He supposed the unity of the person of Christ, and at the same time his sinlessness, could be saved only by the excision of the human spirit; since sin has its seat, not in the will-less soul, nor in the body, but in the intelligent, free, and therefore changeable will or spirit of man. He also charged the Church doctrine of the full humanity of Christ with limiting the atoning suffering of Christ to the human nature, and so detracting from the atoning virtue of the work of Christ; for the death of a man could not destroy death. The divine nature must participate in the suffering throughout. His opponents, for this reason, charged him with making deity suffer and die. He made, however, a distinction between two sides of the Logos, the one allied to man and capable of suffering, and the other allied to God and exalted above all suffering. The relation of the divine pneumatic nature in Christ to the human psychical and bodily nature Apollinaris illustrated by the mingling of wine and water, the glowing fire in the iron, and the union of soul and body in man, which, though distinct, interpenetrate and form one thing.

His doctrine, however, in particulars, is variously represented, and there arose among his disciples a complex mass of opinions, some of them differing strongly from one another. According to one statement Apollinaris asserted that Christ brought even his human nature from heaven, and was from eternity ἔνσαρκος; according to another this was merely an opinion of his disciples, or an unwarranted inference of opponents from his assertion of an eternal determination to incarnation, and from his strong emphasizing of the union of the Logos with the flesh of Christ, which allowed that even the flesh might be worshipped without idolatry.\footnote{Dorner, who has treated this section of the history of Christology, as well as others, with great thoroughness, says, i. 977: “That the school of Apollinaris did not remain in all points consistent with itself, nor true to its founder, is certain; but it is less certain whether Apollinaris himself always taught the same thing.” Theodoret charges him with a change of opinion, which Dorner attributes to different stages of the development of his system.}

The Church could not possibly accept such a half Docetistic incarnation, such a mutilated and stunted humanity of Christ, despoiled of its royal head, and such a merely partial redemption as this inevitably involved. The incarnation of the Logos is his becoming completely man.\footnote{\textgreek{Ἐσάρκωσις}is at the same time \textgreek{ἐνανθρώπησις}. Christ was really \textgreek{ἄνθρωπος}, not merely \textgreek{ὡς ἄνθρωπος}, as Apollinaris taught on the strength of Phil. ii. 7.} It involves, therefore, his assumption of the entire undivided nature of man, spiritual and bodily, with the sole exception of sin, which in fact belongs not to the original nature of man, but has entered from without, as a foreign poison, through the deceit of the devil. Many things in the life of Jesus imply a reasonable soul: sadness, anguish, and prayer. The spirit is just the most essential and most noble constituent of man, the controlling principle,\footnote{\textgreek{Τὸ κυριώτατον}.} and it stands in the same need of redemption as the soul and the body. Had the Logos not assumed the human spirit, he would not have been true man at all, and could not have been our example. Nor could he have redeemed the spirit; and a half-redemption is no redemption at all. To be a full Redeemer, Christ must also be fully man, τέλειος ἄνθρωπος. This was the weighty doctrinal result of the Apollinarian controversy.

Athanasius, the two Gregories, Basil, and Epiphanius combated the Apollinarian error, but with a certain embarrassment, attacking it rather from behind and from the flank, than in front, and unprepared to answer duly its main point, that two integral persons cannot form one person. The later orthodox doctrine surmounted this difficulty by teaching the impersonality of the human nature of Christ, and by making the personality of Christ to reside wholly in the Logos.

The councils at Rome under Damasus, in 377 and 378, and likewise the second ecumenical council, in 381, condemned the Apollinarians.\footnote{Conc. Constant. i. can. 1, where, with the Arians, semi-Arians, Pneumatomachi, Sabellians, and Marcellians or Photinians, the Apollinarians also are anathematized.} Imperial decrees pursued them, in 388, 397, and 428. Some of them returned into the catholic church; others mingled with the Monophysites, for whose doctrine Apollinaris had, in some measure, prepared the way.

With the rejection of this error, however, the question of the proper relation of the divine and human natures in Christ was not yet solved, but rather for the first time fairly raised. Those church teachers proved the necessity of a reasonable human soul in Christ. But respecting the mode of the union of the two natures their views were confused and their expressions in some cases absolutely incorrect and misleading.\footnote{This is true even of Athanasius. Comp. the note on him in § 136, \textbf{p. 706 f.}} It was through the succeeding stages of the Christological controversies that the church first reached a clear insight into this great mystery: God manifest in the flesh.

\xsection{The Nestorian Controversy, a.d. 428-431}

§ 137. The Nestorian Controversy, a.d. 428–431.

Sources.

I. Nestorius: Ὁμιλίαι, Sermones; Anathematismi. Extracts from the Greek original in the Acts of the council of Ephesus; in a Latin translation in Marius Mercator, a North African layman who just then resided in Constantinople, Opera, ed. Garnerius, Par. 1673. Pars ii, and better ed. Baluzius, Par. 1684; also in Gallandi, Bibl. vet. P. P. viii. pp. 615–735, and in Migne’s Patrol. tom. 48. Nestorius’ own account (Evagr. H. E. i. 7) was used by his friend Irenaeus (comes, then bishop of Tyre till 448) in his Tragödia s. comm. de rebus in synodo Ephesina ac in Oriente toto gestis, which, however, is lost; the documents attached to it were revised in the 6th century in the Synodicon adversus tragödiam Irenaei, in Mansi, tom. v. fol. 731 sqq. In favor of Nestorius, or at least of his doctrine, Theodoret († 457) in his works against Cyril, and in three dialogues entitled Ἑρανιστής(Beggar). Comp. also the fragments of Theodore of Mopsuestia, († 429).

II. Against Nestorius: Cyril of Alex.: Ἀναθεματισμοὶ, Five Books κατὰ Νεστορίου, and several Epistles against Nest., and Theod., in vol. vi. of Aubert’s ed. of his Opera, Par. 1638 (in Migne’s ed. t. ix.). Socrates: vii. c. 29–35 (written after 431, but still before the death of Nestorius; comp. c. 84). Evagrius: H. E. i. 2–7. Liberatus (deacon of Carthage about 553): Breviarium causes Nestorianorum et Eutychianorum (ed. Gartnier, Par. 1675, and printed in Gallandi, Bibl. vet. Patr. tom. xii. pp. 121–161). Leontinus Byzant. (monachus): De sectis; and contra Nestorium et Eutychen (in Gallandi, Bibl. tom. xii. p. 625 sqq., and 658–700). A complete collection of all the acts of the Nestorian controversy in Mansi, tom. iv. fol. 567 sqq., and tom. v. vii. ix.

Later Literature.

Petavius: Theolog. dogmatum tom. iv. (de incarnations), lib. i. c. 7 sqq. Jo. Garnier: De haeresi et libris Nestorii (in his edition of the Opera Marii Mercator. Par. 1673, newly edited by Migne, Par. 1846). Gibbon: Decline and Fall of the R. E. ch. 41. P. E. Jablonski: De Nestorianismo. Berol. 1724. Gengler (R.C.): Ueber die Verdammung des Nestorius (Tübinger Quartalschrift, 1835, No. 2). Schröckh: K. Geschichte, vol. xviii. pp. 176–312. Walch: Ketzerhist. v. 289–936. Neander: K. Gesch. vol. iv. pp. 856–992. Gieseler, vol. i. Div. ii. pp. 131 ff. (4th ed.). Baur: Dreieinigkeit, vol. i. 693–777. Dorner: Christologie, vol. ii. pp. 60–98. Hefele (R.C.): Conciliengesch., vol. ii. pp. 134:ff. H. H. Milman: History of Latin Christianity, vol. i. ch. iii. pp. 195–252. (Stanley, in his History of the Eastern Church, has seen fit to ignore the Nestorian, and the other Christological controversies—the most important in the history of the Greek church!) Comp. also W. Möller: Article Nestorius, in Herzog’s Theol. Encykl. vol. x. (1858) pp. 288–296, and the relevant sections in the works on Doctrine History.

Apollinarianism, which sacrificed to the unity of the person the integrity of the natures, at least of the human nature, anticipated the Monophysite heresy, though in a peculiar way, and formed the precise counterpart to the Antiochian doctrine, which was developed about the same time, and somewhat later by Diodorus, bishop of Tarsus (died 394), and Theodore, bishop of Mopsuestia (393–428), and which held the divine and human in Christ so rigidly apart as to make Christ, though not professedly, yet virtually a double person.

From this school proceeded Nestorius, the head and martyr of the Christological heresy which bears his name. His doctrine differs from that of Theodore of Mopsuestia only in being less speculative and more practical, and still less solicitous for the unity of the person of Christ.\footnote{So Dorner also states the difference, vol. ii. p. 62 f.} He was originally a monk, then presbyter in Antioch, and after 428 patriarch of Constantinople. In Constantinople a second Chrysostom was expected in him, and a restorer of the honor of his great predecessor against the detraction of his Alexandrian rival. He was an honest man, of great eloquence, monastic piety, and the spirit of a zealot for orthodoxy, but impetuous, vain, imprudent, and wanting in sound, practical judgment. In his inaugural sermon he addressed Theodosius II. with these words: “Give me, O emperor, the earth purified of heretics, and I will give thee heaven for it; help me to fight the heretics, and I will help thee to fight the Persians.”\footnote{Socrates, H. E., vii. 29.}

He immediately instituted violent measures against Arians, Novatians, Quartodecimanians, and Macedonians, and incited the emperor to enact more stringent laws against heretics. The Pelagians alone, with whose doctrine of free will (but not of original sin) he sympathized, he treated indulgently, receiving to himself Julian of Eclanum, Coelestius, and other banished leaders of that party, interceding for them in 429 with the emperor and with the pope Celestine, though, on account of the very unfavorable reports concerning Pelagianism which were spread by the layman Marius Mercator, then living in Constantinople, his intercessions were of no avail. By reason of this partial contact of the two, Pelagianism was condemned by the council of Ephesus together with Nestorianism.

But now Nestorius himself fell out with the prevailing faith of the church in Constantinople. The occasion was his opposition to the certainly very bold and equivocal expression mother of God, which had been already sometimes applied to the virgin Mary by Origen, Alexander of Alexandria, Athanasius, Basil, and others, and which, after the Arian controversy, and with the growth of the worship of Mary, passed into the devotional language of the people.\footnote{\textgreek{Θεοτόκος}Deipara, genitrix Dei, mater Dei. On the earlier use of this word comp. Petavius: De incarnatione, lib. v. c. 15 (tom. iv. p. 47 1 sqq., Paris ed. of 1650). In the Bible the expression does not occur and only the approximate \textgreek{μήτηρ τοῦ κυρίου}, in Luke i. 43; \textgreek{μήτηρ Ἰησοῦ}, on the contrary, is frequent. Cyril appeals to Gal. iv. 4: ” God sent forth his Son, made of a woman.” To the Protestant mind \textgreek{θεοτόκος} is offensive on account of its undeniable connection with the Roman Catholic worship of Mary, which certainly reminds us of the pagan mothers of gods. Comp. §§ 82 and 83.}

It was of course not the sense, or monstrous nonsense, of this term, that the creature bore the Creator, or that the eternal Deity took its beginning from Mary; which would be the most absurd and the most wicked of all heresies, and a shocking blasphemy; but the expression was intended only to denote the indissoluble union of the divine and human natures in Christ, and the veritable incarnation of the Logos, who took the human nature from the body, of Mary, came forth God-Man from her womb, and as God-Man suffered on the cross. For Christ was borne as a person, and suffered as a person; and the personality in Christ resided in his divinity, not in his humanity. So, in fact, the reasonable soul of man, which is the centre of the human personality, participates in the suffering and the death-struggle of the body, though the soul itself does not and cannot die.

The Antiochian theology, however, could not conceive a human nature without a human personality, and this it strictly separated from the divine Logos. Therefore Theodore of Mopsuestia had already disputed the term theotokos with all earnestness. “Mary,” says he, “bore Jesus, not the Logos, for the Logos was, and continues to be, omnipresent, though he dwelt in Jesus in a special manner from the beginning. Therefore Mary is strictly the mother of Christ, not the mother of God. Only in a figure, per anaphoram, can she be called also the mother of God, because God was in a peculiar sense in Christ. Properly speaking, she gave birth to a man in whom the union with the Logos had begun, but was still so incomplete that he could not yet (till after his baptism) be called the Son of God.” He even declared it “insane” to say that God was born of the Virgin; “not God, but the temple in which God dwelt, was born of Mary.”

In a similar strain Nestorius, and his friend Anastasius, a priest whom he had brought with him from Antioch, argued from the pulpit against the theotokon. Nestorius claimed that he found the controversy already existing in Constantinople, because some were calling Mary mother of God (θεοτόκος), others, mother of Man (ἀνθρωποτόκος). He proposed the middle expression, mother of Christ (Χριστοτόκος), because Christ was at the same time God and man. He delivered several discourses on this disputed point. “You ask,” says he in his first sermon, “whether Mary may be called mother of God. Has God then a mother? If so, heathenism itself is excusable in assigning mothers to its gods; but then Paul is a liar, for he said of the deity of Christ that it was without father, without mother, and without descent.\footnote{Heb. vii. 3: \textgreek{ἀπάτωρ, ἀμήτωπ, ἄνευ γενεαλογίας}.} No, my dear sir, Mary did not bear God; ... the creature bore not the uncreated Creator, but the man who is the instrument of the Godhead; the Holy Ghost conceived not the Logos, but formed for him, out of the virgin, a temple which he might inhabit (John ii. 21). The incarnate God did not die, but quickened him in whom he was made flesh .... This garment, which he used, I honor on account of the God which was covered therein and inseparable therefrom; ... I separate the natures, but I unite the worship. Consider what this must mean. He who was formed in the womb of Mary, was not himself God, but God assumed him [assumsit, i.e., clothed himself with humanity], and on account of Him who assumed, he who was assumed is also called God.”\footnote{In the original in Mansi, iv. 1197; in a Latin translation in Marius Mercator, ed. Garnier, Migne, p. 757 ff. Comp. this and similar passages also in Hefele, ii. p. 137, and Gieseler, i. 2, 139.}

From this word the Nestorian controversy took its rise; but this word represented, at the same time, a theological idea and a mighty religious sentiment; it was intimately connected with the growing veneration of Mary; it therefore struck into the field of devotion, which lies much nearer the people than that of speculative theology; and thus it touched the most vehement passions. The word theotokos was the watchword of the orthodox party in the Nestorian controversy, as the term homoousios had been in the Arian; and opposition to this word meant denial of the mystery of the incarnation, or of the true union of the divine and human natures in Christ.

And unquestionably the Antiochian Christology, which was represented by Nestorius, did not make the Logos truly become man. It asserted indeed, rightly, the duality of the natures, and the continued distinction between them; it denied, with equal correctness, that God, as such, could either be born, or suffer and die; but it pressed the distinction of the two natures to double personality. It substituted for the idea of the incarnation the idea of an assumption of human nature, or rather of an entire man, into fellowship with the Logos,\footnote{\textgreek{Πρόσληψις}. Theodore of Mopsuestia says (Act. Conc. Ephes. in Mansi, iv. fol. 1349): \textgreek{Ὁ δεσπότης θεὸς λόγος ἄνθρωπον εἴληφε τέλειον}(hominem perfectum assumpsit), instead of \textgreek{φύσιν ἀνθρώπου εἴληφε,}or \textgreek{σάρξ ἐγένετο}.} and an indwelling of Godhead in Christ.\footnote{\textgreek{Ἐνοίκησις,} in distinction from \textgreek{ἐνσάρκωσις}.} Instead of God-Man,\footnote{\textgreek{Θεάνθρωπος.}} we have here the idea of a mere God-bearing man;\footnote{\textgreek{θεοφόρος,} also \textgreek{θεοδόχος}, from \textgreek{δέχεσθαι}, God-assuming.} and the person of Jesus of Nazareth is only the instrument or the temple,\footnote{Instrumentum, templum, \textgreek{ναὸς}, a favorite term with the Nestorians.} in which the divine Logos dwells. The two natures form not a personal unity,\footnote{\textgreek{Ἕνωσις καθ  ὑπόστασιν.}} but only a moral unity, an intimate friendship or conjunction.\footnote{\textgreek{Συνάφεια}, connection, affinity, intercourse, attachment in distinction from \textgreek{ἔνωσις}, true interior union. Cyril of Alexandria charges Nestorius, in his Epist. ad Coelestinum: -\textgreek{Φεύγει πανταχοῦ τὸ λέγειν, τὴν ἔνωσιν, ἀλλ  ὀνομάζει τὴν συνάφειαν, ωὝσπερ ἐστιν ὃ ἔξωθεν}.} They hold an outward, mechanical relation to each other,\footnote{\textgreek{Ἕνωσις σχετική,} a unity of relation (from \textgreek{σχέσις}, condition, relation) in distinction from a \textgreek{ἕνωσις φυσική}, or \textgreek{σύγκρασις}, physical unity or commixture.} in which each retains its peculiar attributes,\footnote{\textgreek{Ἰδιώματα.}} forbidding any sort of communicatio idiomatum. This union is, in the first place, a gracious condescension on the part of God,\footnote{\textgreek{Ἕνωσις κατὰ χάριν,} or \textgreek{κατ  εὐδοκίαν}.} whereby the Logos makes the man an object of the divine pleasure; and in the second place, an elevation of the man to higher dignity and to sonship with God.\footnote{\textgreek{Ἕνωσις κατ  αξίαν, καθ  υἱοθεσίαν.}.} By virtue of the condescension there arises, in the third place, a practical fellowship of operation,\footnote{\textgreek{Ἕνωσις κατ  ἐνέργειαν}} in which the humanity becomes the instrument and temple of the deity and the ἕνωσις σχετικήcuIminates. Theodore of Mopsuestia, the able founder of the Antiochian Christology, set forth the elevation of the man to sonship with God (starting from Luke ii. 53) under the aspect of a gradual moral process, and made it dependent on the progressive virtue and meritoriousness of Jesus, which were completed in the resurrection, and earned for him the unchangeableness of the divine life as a reward for his voluntary victory of virtue.

The Antiochian and Nestorian theory amounts therefore, at bottom, to a duality of person in Christ, though without clearly avowing it. It cannot conceive the reality of the two natures without a personal independence for each. With the theanthropic unity of the person of Christ it denies also the theanthropic unity of his work, especially of his sufferings and death; and in the same measure it enfeebles the reality of redemption.\footnote{Cyril charges upon Nestorius (Epist. ad Coelest.), that he does not say the Son of God died and rose again, but always only the man Jesus died and rose. Nestorius himself says, in his second homily (in Mar. Merc. 763 sq.): It may be said that the Son of God, in the \emph{wider} sense, died, but not that God died. Moreover, the Scriptures, in speaking of the birth, passion, and death, never say \emph{God}, but \emph{Christ}, or \emph{Jesus}, or the \emph{Lord},—all of them names which suit both natures. A born, dead, and buried God, cannot be worshipped. Pilate, says he in another sermon, did not crucify the Godhead, but the clothing of the Godhead, and Joseph of Arimathea did not shroud and bury the Logos (in Marius Merc. 789 sqq.).}

From this point of view Mary, of course, could be nothing more than mother of the man Jesus, and the predicate theotokos, strictly understood, must appear absurd or blasphemous. Nestorius would admit no more than that God passed through (transiit) the womb of Mary.

This very war upon the favorite shibboleth of orthodoxy provoked the bitterest opposition of the people and of the monks, whose sympathies were with the Alexandrian theology. They contradicted Nestorius in the pulpit, and insulted him on the street; while he, returning evil for evil, procured corporal punishments and imprisonment for the monks, and condemned the view of his antagonists at a local council in 429.\footnote{According to a partisan report of Basilius to the emperor Theodosius, Nestorius struck, with his own hand, a presumptuous monk who forbade the bishop, as an obstinate heretic, to approach the altar, and then made him over to the officers, who flogged him through the streets and then cast him out of the city.}

His chief antagonist in Constantinople was Proclus, bishop of Cyzicum, perhaps an unsuccessful rival of Nestorius for the patriarchate, and a man who carried the worship of Mary to an excess only surpassed by a modern Roman enthusiast for the dogma of the immaculate conception. In a bombastic sermon in honor of the Virgin\footnote{See Mansi, tom iv. 578; and the remarks of Walch, vol. v. 373 ff.} he praised her as “the spotless treasure-house of virginity; the spiritual paradise of the second Adam; the workshop, in which the two natures were annealed together; the bridal chamber in which the Word wedded the flesh; the living bush of nature, which was unharmed by the fire of the divine birth; the light cloud which bore him who sat between the Cherubim; the stainless fleece, bathed in the dews of Heaven, with which the Shepherd clothed his sheep; the handmaid and the mother, the Virgin and Heaven.”

Soon another antagonist, far more powerful, arose in the person of the patriarch Cyril of Alexandria, a learned, acute, energetic, but extremely passionate, haughty, ambitious, and disputatious prelate. Moved by interests both personal and doctrinal, he entered the field, and used every means to overthrow his rival in Constantinople, as his like-minded uncle and predecessor, Theophilus, had overthrown the noble Chrysostom in the Origenistic strife. The theological controversy was at the same time a contest of the two patriarchates. In personal character Cyril stands far below Nestorius, but he excelled him in knowledge of the world, shrewdness, theological learning and acuteness, and had the show of greater veneration for Christ and for Mary on his side; and in his opposition to the abstract separation of the divine and human he was in the right, though he himself pressed to the verge of the opposite error of mixing or confusing the two natures in Christ.\footnote{Comp. in particular his assertion of a \textgreek{ἕνωσις φυσική}in the third of his Anathematismi against Nestorius; Hefele (ii. 155), however, understands by this not a \textgreek{ἔνωσις εἰς μίαν φύσιν,} but only a real union in \emph{one being, one existence}.} In him we have a striking proof that the value of a doctrine cannot always be judged by the personal worth of its representatives. God uses for his purposes all sorts of instruments, good, bad, and indifferent.

Cyril first wrote to Nestorius; then to the emperor, the empress Eudokia, and the emperor’s sister Pulcheria, who took lively interest in church affairs; finally to the Roman bishop Celestine; and he warned bishops and churches east and west against the dangerous heresies of his rival. Celestine, moved by orthodox instinct, flattered by the appeal to his authority, and indignant at Nestorius for his friendly reception of the exiled Pelagians, condemned his doctrine at a Roman council, and deposed him from the patriarchal chair, unless he should retract within ten days (430).

As Nestorius persisted in his view, Cyril, despising the friendly mediation of the patriarch John of Antioch, hurled twelve anathemas, or formulas of condemnation, at the patriarch of Constantinople from a council at Alexandria by order of the pope (430).\footnote{Cyrilli Opera, tom. iii. 67; in Mansi, iv. fol. 1067 sqq.; in Gieseler, i. ii. p. 143 ff. (§ 88, not. 20); in Hefele, ii. 155 ff.}

Nestorius replied with twelve counter-anathemas, in which he accused his opponents of the heresy of Apollinaris.\footnote{In Marius Mercator, p. 909; Gieseler, i. ii. 145 f.; Hefele, ii. 158 ff.} Theodoret of Cyros, the learned expositor and church historian, also wrote against Cyril at the instance of John of Antioch.

The controversy had now become so general and critical, that it could be settled only by an ecumenical council.

\xsection{The Ecumenical Council of Ephesus, a.d. 431. The Compromise}

§ 138. The Ecumenical Council of Ephesus, a.d. 431. The Compromise.

For the Acts of the Council, see Mansi (tom. iv. fol. 567–1482, and a part of tom. v.), Harduin, and Fuchs, and an extended history of the council and the transactions connected with it in Walch, Schröckh, and Hefele (ii. pp. 162–271). We confine ourselves to the decisive points.

Theodosius II., in connection with his Western colleague, Valentinian III., summoned a universal council on Pentecost, a.d. 431, at Ephesus, where the worship of the Virgin mother of God had taken the place of the worship of the light and life dispensing virgin Diana. This is the third of the ecumenical councils, and is held, therefore, by all churches, in high regard. But in moral character this council stands far beneath that of Nicaea or of the first council of Constantinople. An uncharitable, violent, and passionate Spirit ruled the transactions. The doctrinal result, also, was mainly only negative; that is to say, condemnation of Nestorianism. The positive and ecumenical character of the council was really secured only by the subsequent transactions, and the union of the dominant party of the council with the protesting minority of Oriental bishops.\footnote{It is with reference to this council mainly that Dean Milman (Latin Christianity, i. 227) passes the following harsh and sweeping judgment on the ecumenical councils of the ancient church: “Nowhere is Christianity less attractive, and, if we look to the ordinary tone and character of the proceedings, less authoritative, than in the councils of the church. It is in general a fierce collision of two rival factions, neither of which will yield, each of which is solemnly pledged against conviction. Intrigue, injustice, violence, decisions on authority alone, and that the authority of a turbulent majority, decisions by wild acclamation rather than after sober inquiry, detract from the reverence, and impugn the judgments, at least of the later councils. The close is almost invariably a terrible anathema, in which it is impossible not to discern the tones of human hatred, of arrogant triumph, of rejoicing at the damnation imprecated against the humiliated adversary. Even the venerable council of Nicaea commenced with mutual accusals and recriminations, which were suppressed by the moderation of the emperor; and throughout the account of Eusebius there is an adulation of the imperial convert with something of the intoxication, it might be of pardonable vanity, at finding themselves the objects of royal favor, and partaking in royal banquets. But the more fatal error of that council was the solicitation, at least the acquiescence in the infliction, of a civil penalty, that of exile, against the recusant prelates. The degeneracy is rapid from the council of Nicaea to that of Ephesus, where each party came determined to use every means of haste, manoeuvre, court influence, bribery, to crush his adversary; where there was an encouragement of, if not an appeal to, the violence of the populace, to anticipate the decrees of the council; where each had his own tumultuous foreign rabble to back his quarrel; and neither would scruple at any means to obtain the ratification of their anathemas through persecution by the civil government.” This is but the dark side of the picture. In spite of all human passions and imperfections truth triumphed at last, and this alone accounts for the extraordinary effect of these ecumenical councils, and the authority they still enjoy in the whole Christian world.}

Nestorius came first to Ephesus with sixteen bishops, and with an armed escort, as if he were going into battle. He had the imperial influence on his side, but the majority of the bishops and the prevailing voice of the people in Ephesus, and also in Constantinople, were against him. The emperor himself could not be present in person, but sent the captain of his body-guard, the comes Candidian. Cyril appeared with a numerous retinue of fifty Egyptian bishops, besides monks, parabolani, slaves, and seamen, under the banner of St. Mark and of the holy Mother of God. On his side was the archbishop Memnon of Ephesus, with forty of his Asiatic suffragans and twelve bishops from Pamphilia; and the clergy, the monks, and the people of Asia Minor were of the same sentiment. The pope of Rome—for the first time at an ecumenical council—was represented by two bishops and a priest, who held with Cyril, but did not mix in the debates, as they affected to judge between the contending parties, and thus maintain the papal authority. This deputation, however, did not come in at the beginning.\footnote{St Augustinealso was one of the Western bishops who were summoned, the emperor having sent a special officer to him; but he had died shortly before, on the 28th of August, 430.} The patriarch John of Antioch, a friend of Nestorius, was detained on the long journey with his bishops.

Cyril refused to wait, and opened the council in the church of St. Mary with a hundred and sixty bishops\footnote{Before the sentence of deposition came to be subscribed, the number had increased to a hundred and ninety-eight. According to the Roman accounts Cyril presided in the name and under the commission of the pope; but in this case he should have yielded the presidency in the second and subsequent sessions, at which the papal legates were present; which he did not do.} sixteen days after Pentecost, on the 22d of June, in spite of the protest of the imperial commissioner. Nestorius was thrice cited to appear, but refused to come until all the bishops should be assembled. The council then proceeded without him to the examination of the point in dispute, and to the condemnation of Nestorius. The bishops unanimously cried: “Whosoever does not anathematize Nestorius, let himself be anathema; the true faith anathematizes him; the holy council anathematizes him. Whosoever holds fellowship with Nestorius, let him be anathema. We all anathematize the letter and the doctrines of Nestorius. We all anathematize Nestorius and his followers, and his ungodly faith, and his ungodly doctrine. We all anathematize Nestorius,” \&c.\footnote{In Mansi, tom. iv. p. 1170 sq.; Hefele, ii. 169.} Then a multitude of Christological expressions of the earlier fathers and several passages from the writings of Nestorius were read, and at the close of the first session, which lasted till late in the night, the following sentence of deposition was adopted and subscribed by about two hundred bishops: “The Lord Jesus Christ, who is blasphemed by him [Nestorius], determines through this holy council that Nestorius be excluded from the episcopal office, and from all sacerdotal fellowship.”\footnote{\textgreek{Ὁ βλασφημηθεὶς τοίνυν παρ  αὐτοῦ κύριος ·ἡμῶν Ἰησοῦς Χριστὸς ωὝρισε διὰ τῆς παρούσης ἁγιωτάτης συνόδου, ἀλλότριον εἶναι τὸν αὐτὸν Νεστόριον τοῦ ἐπισκοπικοῦ ἀξιώματος καὶ παντὸς συλλόγου ἱερατικοῦ}. Mansi, iv. fol. 1211; Hefele, ii. 172.}

The people of Ephesus hailed this result with universal jubilee, illuminated the city, and accompanied Cyril with torches and censers in state to his house.\footnote{So Cyril himself complacently relates in a letter to his friends in Egypt. See Mansi, tom. iv. 1241 sq.}

On the following day Nestorius was informed of the sentence of deposition in a laconic edict, in which he was called a new Judas. But he indignantly protested against the decree, and made complaint in an epistle to the emperor. The imperial commissioner declared the decrees invalid, because they were made by only a portion of the council, and he prevented as far as possible the publication of them.

A few days after, on the 26th or 27th of June, John of Antioch at last reached Ephesus, and immediately, with forty-two bishops of like sentiment, among whom was the celebrated Theodoret, held in his dwelling, under the protection of the imperial commissioner and a body-guard, a counter council or conciliabulum, yielding nothing to the haste and violence of the other, deposed Cyril of Alexandria and Memnon of Ephesus from all priestly functions, as heretics and authors of the whole disorder and declared the other bishops who voted with them excommunicate until they should anathematize the heretical propositions of Cyril.\footnote{The Acts of this counter council in Mansi, tom. iv. 1259 sqq. (Acta Conciliabuli). Comp. also Hefele, ii. 178 ff.}

Now followed a succession of mutual criminations, invectives, arts of church diplomacy and politics, intrigues, and violence, which give the saddest picture of the uncharitable and unspiritual Christianity of that time. But the true genius of Christianity is, of course, far elevated above its unworthy organs, and overrules even the worst human passions for the cause of truth and righteousness.

On the 10th of July, after the arrival of the papal legates, who bore themselves as judges, Cyril held a second session, and then five more sessions (making seven in all), now in the house of Memnon, now in St. Mary’s church, issuing a number of circular letters and six canons against the Nestorians and Pelagians.

Both parties applied to the weak emperor, who, without understanding the question, had hitherto leaned to the side of Nestorius, but by public demonstrations and solemn processions of the people and monks of Constantinople under the direction of the aged and venerated Dalmatius, was awed into the worship of the mother of God. He finally resolved to confirm both the deposition of Nestorius and that of Cyril and Memnon, and sent one of the highest civil officers, John, to Ephesus, to publish this sentence, and if possible to reconcile the contending parties. The deposed bishops were arrested. The council, that is the majority, applied again to the emperor and his colleague, deplored their lamentable condition, and desired the release of Cyril and Memnon, who had never been deposed by them, but on the contrary had always been held in high esteem as leaders of the orthodox doctrine. The Antiochians likewise took all pains to gain the emperor to their side, and transmitted to him a creed which sharply distinguished, indeed, the two natures in Christ, yet, for the sake of the unconfused union of the two (ἀσύγχυτος ἕωσις), conceded to Mary the disputed predicate theotokos.

The emperor now summoned eight spokesmen from each of the two parties to himself to Chalcedon. Among them were, on the one side, the papal deputies, on the other John of Antioch and Theodoret of Cyros, while Cyril and Memnon were obliged to remain at Ephesus in prison, and Nestorius at his own wish was assigned to his former cloister at Antioch, and on the 25th of October, 431, Maximian was nominated as his successor in Constantinople. After fruitless deliberations, the council of Ephesus was dissolved in October, 431, Cyril and Memnon set free, and the bishops of both parties commanded to go home.

The division lasted two years longer, till at last a sort of compromise was effected. John of Antioch sent the aged bishop Paul of Emisa a messenger to Alexandria with a creed which he had already, in a shorter form, laid before the emperor, and which broke the doctrinal antagonism by asserting the duality of the natures against Cyril, and the predicate mother of God against Nestorius.\footnote{In Mansi, tom. v. fol. 305; Hefele, ii. 246; and Gieseler, i. ii. p. 150} “We confess,” says this symbol, which was composed by Theodoret, “that our Lord Jesus Christ, the only begotten Son of God, is perfect God and perfect man, of a reasonable soul and body subsisting;\footnote{. \textgreek{Θεὸν τέλειον καὶ ἄνθρωπον τέλειον ἐκ ψυχῆς λογικῆς}(against Apollinaris), \textgreek{καὶ σώματος}.} as to his Godhead begotten of the Father before all time, but as to his manhood, born of the Virgin Mary in the end of the days for us and for our salvation; of the same essence with the Father as to his Godhead, and of the same substance with us as to his manhood;\footnote{\textgreek{Ὁμοούσιον τῷ πατρὶ κατὰ τὴν θεότητα, καὶ ὁμοούσιον ἡμῖν κατὰ τὴν ἀνθρωπότητα}. Here \emph{homoousios}, at least in the second clause, evidently does not imply numerical unity, but only generic unity.} for two natures are united with one another.\footnote{\textgreek{Δύο γὰρ φύσεων ἕνωσις·γέγονε}, in opposition to the \textgreek{μία φύσις}of Cyril.} Therefore we confess one Christ, one Lord, and one Son. By reason of this union, which yet is without confusion,\footnote{. \textgreek{Κατὰ ταύτην τὴν τῆς ἀσυγχύτου}(against Cyril) \textgreek{ἑνώσεως ἔννοιαν}.} we also confess that the holy Virgin is mother of God, because God the Logos was made flesh and man, and united with himself the temple [humanity] even from the conception; which temple he took from the Virgin. But concerning the words of the Gospel and Epistles respecting Christ, we know that theologians apply some which refer to the one person to the two natures in common, but separate others as referring to the two natures, and assign the expressions which become God to the Godhead of Christ, but the expressions of humiliation to his manhood.”\footnote{\textgreek{Καὶ τὰς μὲν θεοπρεπεῖς κατὰ τὴν θεότητα τοῦ Χριστοῦ, τὰς δὲ ταπεινὰς κατὰ τὴν ἀνθρωπότητα αὐτοῦ παραδιδόντας.}Gieseler says (i. ii. p. 152), Nestorius never asserted anything but what agrees with this confession which Cyril subscribed. But he pressed the distinction of the natures in Christ so far that it amounted, in substance, though not in expression, to two persons; he taught not a true becoming man, but the union of the Logos with a \textgreek{τέλειος ἄνθρωπος}, a human person therefore not nature; and he constantly denied the \emph{theotokos}, except in an improper sense. His doctrine was unquestionably much distorted by his cotemporaries; but so also was the doctrine of Cyril.} Cyril assented to this confession, and repeated it verbally, with some further doctrinal explanations, in his answer to the irenical letter of the patriarch of Antioch, but insisted on the condemnation and deposition of Nestorius as the indispensable condition of church fellowship. At the same time he knew how to gain the imperial court to the orthodox side by all kinds of presents, which, according to the Oriental custom of testifying submission to princes by presents, were not necessarily regarded as bribes. The Antiochians, satisfied with saving the doctrine of two natures, thought it best to sacrifice the person of Nestorius to the unity of the church, and to anathematize his “wicked and unholy innovations.”\footnote{\textgreek{Τὰς φαύλας αὐτοῦ καὶ βεβήλους καινοφωνίας .}} Thus in 433 union was effected, though not without much contradiction on both sides, nor without acts of imperial force.

The unhappy Nestorius was dragged from the stillness of his former cloister, the cloister of Euprepius before the gates of Antioch, in which he had enjoyed four years of repose, from one place of exile to another, first to Arabia, then to Egypt, and was compelled to drink to the dregs the bitter cup of persecution which he himself, in the days of his power, had forced upon the heretics. He endured his suffering with resignation and independence, wrote his life under the significant title of Tragedy,\footnote{Fragments in Evagrius, H. E. i. 7, and in the Synodicon adversus Tragoediam Irenaei, c. 6. That the book bore the name of \emph{Tragedy}, is stated by Ebedjesu, a Nestorian metropolitan. The imperial commissioner, Irenaeus, afterwards bishop of Tyre, a friend of Nestorius, composed a book concerning him and the ecclesiastical history of his time, likewise under the title of \emph{Tragedy}, fragments of which, in a Latin translation, are preserved in the so-called Synodicon, in Mansi, v. 731 sqq.} and died after 439, no one knows where nor when. Characteristic of the fanaticism of the times is the statement quoted by Evagrius, \footnote{Hist. Eccl. i. 6.} that Nestorius, after having his tongue gnawed by worms in punishment for his blasphemy, passed to the harder torments of eternity. The Monophysite Jacobites are accustomed from year to year to cast stones upon his supposed grave in Upper Egypt and have spread the tradition that it has never been moistened by the rain of heaven, which yet falls upon the evil and the good. The emperor, who had formerly favored him, but was now turned entirely against him, caused all his writings to be burned, and his followers to be named after Simon Magus, and stigmatized as Simonians.\footnote{For his sad fate and his upright character Nestorius, after having been long abhorred, has in modem times, since Luther, found much sympathy; while Cyril by his violent conduct has incurred much censure. Walch, l.c. v. p. 817 ff., has collected the earlier opinions. Gieseler and Neander take the part of Nestorius against Cyril, and think that he was unjustly condemned. So also Milman, who would rather meet the judgment of the Divine Redeemer loaded with the errors of Nestorius than with the barbarities of Cyril, but does not enter into the theological merits of the controversy. (History of Latin Christianity, i. 210.) Petavius, Baur, Hefele, and Ebrard, on the contrary, vindicate Cyril against Nestorius, not as to his personal conduct, which was anything but Christian, but in regard to the particular matter in question, viz., the defence of the unity of Christ against the division of his personality. Dorner (ii. 81 ff.) justly distributes right and wrong, truth and error, on both sides, and considers Nestorius and Cyril representatives of two equally one-sided conceptions, which complement each other. Cyril’s strength lay on the religious and speculative side of Christology, that of Nestorius on the ethical and practical. Kahnis gives a similar judgment, Dogmatik, ii. p. 8 6.}

The same orthodox zeal turned also upon the writings of Theodore of Mopsuestia, the long deceased teacher of Nestorius and father of his error. Bishop Rabulas of Edessa († 435) pronounced the anathema upon him and interdicted his writings; and though his successor Ibas (436–457) again interested himself in Theodore, and translated several of his writings into Syriac (the ecclesiastical tongue of the Persian church), yet the persecution soon broke out afresh, and the theological school of Edessa where the Antiochian theology had longest maintained its life, and whence the Persian clergy had proceeded, was dissolved by the emperor Zeno in 489. This was the end of Nestorianism in the Roman empire.

\xsection{The Nestorians}

§ 139. The Nestorians.

Jos. Sim. Assemani: De Syris Nestorianis, in his Bibliotheca Orientalis. Rom. 1719–1728, fol. tom. iii. P. ii. Ebedjesu (Nestorian metropolitan of Nisibis, † 1318): Liber Margaritae de veritate fidei (a defence of Nestorianism), in Ang. Mai’s Scrip. vet. nova collect. x. ii. 317. Gibbon: Chap. xlvii., near the end. E. Smith and H. G. O. Dwight: Researches in Armenia; with a visit to the Nestorian and Chaldean Christians of Oormiah and Salmas. 2 vols. Bost. 1833. Justin Perkins: A Residence of eight years in Persia. Andover, 1843. Wiltsch: Kirchliche Geographie u. Statistik. Berl. 1846, i. 214 ff. Geo. Percy Badger: The Nestorians and their Rituals. Illustrated (with colored plates), 2 vols. Lond. 1852. H. Newcomb: A Cyclopaedia of Missions. New York, 1856, p. 553 ff. Petermann: Article Nestorianer, in Herzog’s Theol. Encykl. vol. x. (1858), pp. 279–288.

While most of the heresies of antiquity, Arianism not excepted, have been utterly obliterated from history, and only raise their heads from time to time as individual opinions under peculiar modifications, the Christological heresies of the fifth century, Nestorianism and Monophysitism, continue in organized sects to this day. These schismatic churches of the East are the petrified remains or ruins of important chapters in the history of the ancient church. They are sunk in ignorance and superstition; but they are more accessible to Western Christianity than the orthodox Greek church, and offer to the Roman and Protestant churches an interesting field of missions, especially among the Nestorians and the Armenians.

The Nestorians differ from the orthodox Greek church in their repudiation of the council of Ephesus and of the worship of Mary as mother of God, of the use of images (though they retain the sign of the cross), of the doctrine of purgatory (though they have prayers for the dead), and of transubstantiation (though they hold the real presence of Christ in the eucharist), as well as in greater simplicity of worship. They are subject to a peculiar hierarchical organization with eight orders, from the catholicus or patriarch to the sub-deacon and reader. The five lower orders, up to the priests, may marry; in former times even the bishops, archbishops, and patriarchs had this privilege. Their fasts are numerous and strict. The feast-days begin with sunset, as among the Jews. The patriarch eats no flesh; he is chosen always from the same family; he is ordained by three metropolitans. Most of the ecclesiastical books are written in the Syriac language.

After Nestorianism was exterminated from the Roman empire, it found an asylum in the kingdom of Persia, whither several teachers of the theological school of Edessa fled. One of them, Barsumas, became bishop of Nisibis (435–489),\footnote{Not to be confounded with the contemporary Monophysite abbot Barsumas, a saint of the Jacobites.} founded a new theological seminary there, and confirmed the Persian Christians in their aversion to the Cyrillian council of Ephesus, and in their adhesion to the Antiochian and Nestorian theology. They were favored by the Persian kings, from Pherozes, or Firuz, onward (461–488), out of political opposition to Constantinople. At the council of Seleucia (498) they renounced all connection with the orthodox church of the empire. They called themselves, after their liturgical language, Chaldaean or Assyrian Christians, while they were called by their opponents Nestorians. They had a patriarch, who after the year 496 resided in the double city of Seleucia-Ctesiphon, and after 762 in Bagdad (the capital of the Saracenic empire), under the name of Yazelich (catholicus), and who, in the thirteenth century, had no less than twenty-five metropolitans under his supervision.

The Nestorian church flourished for several centuries, spread from Persia, with great missionary zeal, to India, Arabia, and even to China and Tartary, and did good service in scholarship and in the founding of schools and hospitals. Mohammed is supposed to owe his imperfect knowledge of Christianity to a Nestorian monk, Sergius; and from him the sect received many privileges, so that it obtained great consideration among the Arabians, and exerted an influence upon their culture, and thus upon the development of philosophy and science in general.\footnote{The observations of Alex. von Humboldt, in the 2d vol. of his Kosmos (Stuttg. and Tüb. 1847, p. 247 E), on the connection of Nestorianism with the culture and physical science of the Arabians, are worthy of note: “It was one of the wondrous arrangements in the system of things, that the Christian sect of the Nestorians, which has exerted a very important influence on the geographical extension of knowledge, was of service even to the Arabians before the latter found their way to learned and disputatious Alexandria; that Christian Nestorianism, in fact, under the protection of the arms of Islam, was able to penetrate far into Eastern Asia. The Arabians, in other words, gained their first acquaintance with Grecian literature through the Syrians, a kindred Semitic race; while the Syrians themselves, scarcely a century and a half before, had first received the knowledge of Grecian literature through the anathematized Nestorians. Physicians who had been educated in the institutions of the Greeks, and at the celebrated medical school founded by the Nestorian Christians at Edessa in Mesopotamia, were, so early as the times of Mohammed, living, befriended by him and by Abu-Bekr, in Mecca. \par “ The school of Edessa, a model of the Benedictine schools of Monte Casino and Salerno, awakened the scientific search for \emph{materia medica} in the mineral and vegetable kingdoms. When it was dissolved by Christian fanaticism under Zeno the Isaurian, the Nestorians scattered towards Persia, where they soon attained political importance, and established a new and thronged medical institute at Dschondisapur in Khuzistan. They succeeded in spreading their science and their faith to China towards the middle of the seventh century under the dynasty of Thang, five hundred and seventy-two years after Buddhism had penetrated thither from India. \par “The seed of Western culture, scattered in Persia by educated monks, and by the philosophers of the last Platonic school of Athens who were persecuted by Justinian, took beneficent root among the Arabians during their first Asiatic campaign. Feeble as the science of the Nestorian priests may have been, it could still, with its peculiar medical and pharmaceutic turn, act genially upon a race which had long lived in free converse with nature, and had preserved a more fresh sensibility to every sort of study of nature, than the people of Greek and Italian cities. What gives the Arabian epoch the universal importance which we must here insist upon, is in great part connected with the trait of national character just indicated. The Arabians, we repeat, are to be regarded as the proper founders of the \emph{physical sciences}, in the sense which we are now accustomed to attach to the word.”} Among the Tartars, in the eleventh century, it succeeded in converting to Christianity a king, the priest-king Presbyter John (Prester John) of the Kerait, and his successor of the same name.\footnote{On this fabulous priest-kingdom, which the popes endeavored by unsuccessful embassies to unite to the Roman church, and whose light was quenched by the tide of the conquests of Zengis Khan, comp. Mosheim: Historia Tartarorum Eccles. Helmst. 1741; Neander: Kirchengesch. vol. v. p. 84 ff. (9th part of the whole work, 2d. 1841); and Ritter: Erdkunde, part ii. vol. i. pp. 256, 283 (2d ed. 1832).} But of this we have only uncertain accounts, and at all events Nestorian Christianity has since left but slight traces in Tartary and in China.

Under the Mongol dynasty the Nestorians were cruelly persecuted. The terrible Tamerlane, the scourge and the destroyer of Asia, towards the end of the fourteenth century almost exterminated them. Yet they have maintained themselves on the wild mountains and in the valleys of Kurdistan and in Armenia under the Turkish dominion to this day, with a separate patriarch, who from 1559 till the seventeenth century resided at Mosul, but has since dwelt in an almost inaccessible valley on the borders of Turkey and Persia. They are very ignorant and poor, and have been much reduced by war, pestilence, and cholera.

A portion of the Nestorians, especially those in cities, united from time to time, under the name of Chaldaeans, with the Roman church, and have a patriarch of their own at Bagdad.

And on the other side, Protestant missionaries from America have made vigorous and successful efforts, since 1833, to evangelize and civilize the Nestorians by preaching, schools, translations of the Bible, and good books.\footnote{Dr. Justin Perkins, Asahel Grant, Rhea, Stoddard, Wright, and other missionaries of the American Board of Commisioners for Foreign Missions. The centre of their labors is Oormiah, a city of 25,000 inhabitants, of whom 1,000 are Nestorians. Comp. on this subject Newcomb, l.c. 556 ff., especially the letter of Dr. Perkins of 1854, p. 564 ff., on the present condition of this mission; also Joseph P. Thompson: Memoir of the Rev. David Tappan Stoddard, missionary to the Nestorians, Boston, 1858; and a pamphlet issued by the American B.C. F. M.: Historical Sketch of the Mission to the Nestorians by Justin Perkins, and of the Assyrian Mission by Rev. Thomas Laurie, New York, 1862. The American Board of Foreign Missions look upon the Nestorian and Armenian missions as a means and encouraging pledge of the conversion of the millions of Mohammedans among whom Providence has placed and preserved those ancient sects, as it would seem, for such an end.}

The Thomas-Christians in East India are a branch of the Nestorians, named from the apostle Thomas, who is supposed to have preached the gospel on the coast of Malabar. They honor the memory of Theodore and Nestorius in their Syriac liturgy, and adhere to the Nestorian patriarchs. In the sixteenth century they were, with reluctance, connected with the Roman church for sixty years (1599–1663) through the agency of Jesuit missionaries. But when the Portuguese power in India was shaken by the Dutch, they returned to their independent position, and since the expulsion of the Portuguese they have enjoyed the free exercise of their religion on the coast of Malabar. The number of the Thomas-Christians is said still to amount to seventy thousand souls, who form a province by themselves under the British empire, governed by priests and elders.

\xsection{The Eutychian Controversy. The Council of Robbers, a.d. 449}

§ 140. The Eutychian Controversy. The Council of Robbers, a.d. 449.

Comp. the Works at § 137.

Sources.

Acts of the council of Chalcedon, of the local council of Constantinople, and of the Robber Synod of Ephesus. The correspondence between Leo and Flavian, etc. For these acts, letters, and other documents, see Mansi, Conc. tom. v. vi. and vii. (Gelasius?): Breviculus historiae Eutychianistarum a. gesta de nomine Acacii (extending to 486, in Mansi, vii. 1060 sqq.). Liberatus: Breviarium causae Nest. et Eutych. Leontinus Byzant.: Contra Nest. et Eutych. The last part of the Synodicon adv. tragödiam Irenaei (in Mansi, v. 731 sqq.). Evagrius: H. E. i. 9 sqq. Theodoret: Ἐρανιστής(the Beggar) or Πολύμορφος(the Multiformed),—a refutation of the Egyptian Eutychian system of doctrines (which begged together so much from various old heresies, as to form a now one), in three dialogues, written in 447 (Opera, ed. Schulze, vol. iv.).

Literature.

Petavius: De incarnatione Verbi, lib. i. c. 14–18, and the succeeding books, particularly iii., iv., and v. (Theolog. dogmatum, tom. iv. p. 65 sqq. ed. Par. 1650). Tillemont: Mémoires, tom. xv. pp. 479–719. C. A. Salig: De Eutychianismo ante Eutychen. Wolfenb. 1723. Walch: Ketzerhist. vol. vi. 3–640. Schröckh: vol. xviii. 433–492. Neander: Kirchengesch. iv. pp. 942–992. Baur: Gesch. der Lehre von d. Dreieinigkeit, etc. i. 800–825. Dorner: Gesch. d. Lehre v. d. Pers. Chr. ii. 99–149. Hefele (R.C.): Conciliengesch. ii. pp. 295–545. W. Cunningham: Historical Theology, i. pp. 311–’15. Comp. also the Monographs of Arendt (1835) and Perthel (1848) on Leo I.

The result of the third universal council was rather negative than positive. The council condemned the Nestorian error, without fixing the true doctrine. The subsequent union of the Alexandrians and the Antiochians was only a superficial peace, to which each party had sacrificed somewhat of its convictions. Compromises are generally of short duration; principles and systems must develope themselves to their utmost consequences; heresies must ripen, and must be opened to the core. As the Antiochian theology begot Nestorianism, which stretched the distinction of the human and divine natures in Christ to double personality; so the Alexandrian theology begot the opposite error of Eutychianism or Monophysitism, which urged the personal unity of Christ at the expense of the distinction of natures, and made the divine Logos absorb the human nature. The latter error is as dangerous as the former. For if Christ is not true man, he cannot be our example, and his passion and death dissolve at last into mere figurative representations or docetistic show.

A large portion of the party of Cyril was dissatisfied with the union creed, and he was obliged to purge himself of inconsistency. He referred the duality of natures spoken of in the symbol to the abstract distinction of deity and humanity, while the two are so made one in the one Christ, that after the union all separation ceases, and only one nature is to be recognized in the incarnate Son. The Logos, as the proper subject of the one nature, has indeed all human, or rather divine-human, attributes, but without a human nature. Cyril’s theory of the incarnation approaches Patripassianism, but differs from It in making the Son a distinct hypostasis from the Father. It mixes the divine and human; but It mixes them only in Christ, and so is Christo-theistic, but not pantheistic.\footnote{Cyril’s true view is most clearly expressed in the following propositions (comp. Mansi, v. 320, and Niedner, p. 364): The\textgreek{ἐνσάρκωσις}was\textgreek{φυσικὴ ἕνωσις}or \emph{becoming} man, on the part of God, so that there is only\textgreek{μία σεσαρκωμένη φύσις τοῦ λόγου. Ὁ Ωεὸς λόγος, ἑνωθεὶς σαρκὶ καθ  ὑπόστασιν, ἐγενετο ἄνθρωπος·, συνήφθη ἀνθρώπῳ. Μία ἤδη νοεῖται φύσις μετὰ τὴν ἕνωσιν, ἡ αὐτοῦ τοῦ λόγου σεσαρκωμένη. Ἡ τοῦ κυρίου σάρξ ἐστιν ἰδία τοῦ Θεοῦ λόγου, οὐχ ἑτέρου τινὸς παρ  αὐτόν.} The\textgreek{ἕνωσις τῶν φύσεων}is not, indeed, exactly a \textgreek{σύγχυσις τῶν φύσεων}, but at all events excludes all \textgreek{διαίρεσις}, and demands an absolute co-existence and interpenetration of the \textgreek{λόγος}and the\textgreek{σάρξ}. The consequence of this incarnation is the existence of a new entity, a divine-human subject, which is in nothing \emph{only} God or \emph{only} man, but in everything is both in one, and whose attributes (proprietates, idiomata) are not, some divine and others human, but all divine-human.}

On the other side, the Orientals or Antiochians, under the lead of John, Ibas, and especially Theodoret, interpreted the union symbol in their sense of a distinction of the two natures continuing in the one Christ even after the incarnation, and actually obtained the victory for this moderate Nestorianism, by the help of the bishop of Rome, at the council of Chalcedon.

The new controversy was opened by the party of monophysite sentiment.

Cyril died in 444. His arch-deacon, Dioscurus (Διόσκορος), who had accompanied him to the council at Ephesus, succeeded him in the patriarchal chair of Alexandria (444–451), and surpassed him in all his had qualities, while he fell far behind him in intellect and in theological capacity.\footnote{Towards the memory of Cyril he behaved very recklessly. He confiscated his considerable estate (Cyril was of wealthy family), accused him of squandering the church funds in his war against Nestorius, and unseated several of his relatives. He was himself charged, at the council of Chalcedon, with embezzlement of the moneys of the church and of the poor.} He was a man of unbounded ambition and stormy passion, and shrank from no measures to accomplish his designs and to advance the Alexandrian see to the supremacy of the entire East; in which he soon succeeded at the Council of Robbers. He put himself at the head of the monophysite party, and everywhere stirred the fire of a war against the Antiochian Christology.

The theological representative, but by no means the author, of the monophysite heresy which bears his name, was Eutyches,\footnote{That is, \emph{the Fortunate}. His opponents said he should rather have been named \emph{Atyches, the Unfortunate}. He must not be confounded with the deacon Eutyches, who attended Cyril to the council of Ephesus. Leo the Great in his renowned letter to Flavian, calls him “very ignorant and unskilled,” multum imprudens et nimis imperious, and justly attributes his error rather to imperitia than to versutia. So also Petavius and Hefele (ii. p. 800).} an aged and respected, but not otherwise important presbyter and archimandrite (head of a cloister of three hundred monks) in Constantinople, who had lived many years in monastic seclusion, and had only once appeared in public, to raise his voice, in that procession, for the Cyrillian council of Ephesus and against Nestorius. His relation to the Alexandrian Christology is like that of Nestorius to the Antiochian; that is, he drew it to a head, brought it to popular expression, and adhered obstinately to it; but he is considerably inferior to Nestorius in talent and learning. His connection with this controversy is in a great measure accidental.

Eutyches, like Cyril, laid chief stress on the divine in Christ, and denied that two natures could be spoken of after the incarnation. In our Lord, after his birth, he worshipped only one nature, the nature of God become flesh and man.\footnote{\textgreek{Μίαν φύσιν προσκυνεῖν, καὶ ταύτην Θεοῦ σαρκωθέντος καὶ ἐνανθρωπήσαντος}, or as he declared before the synod at Constantinople: \textgreek{Ὁμολογῶ ἐκ δύο φύσεων γεγεννῆσθαι τὸν κύριονἡμῶν πρὸ τῆς ἑνήσεως · μετὰ δὲ τὴν ἕνωσιν μίαν φυσιν ὁμολογῶ.}Mansi, tom. vi. fol. 744. In behalf of his view he appealed to the Scriptures, to Athanasius and Cyril, and to the council of Ephesus in 431.} The impersonal human nature is assimilated and, as It were, deified by the personal Logos, so that his body is by no means of the same substance (ὁμοούσιον) with ours, but a divine body.\footnote{The other side imputed to Eutychianism the doctrine of a heavenly body, or of an apparent body, or of the transformation of the Logos into flesh. So Theodoret, Fab. haer. iv. 13. Eutyches said, Christ had a \textgreek{σῶμα ἀνθρώπου}, but not a \textgreek{σῶμα ἀνθρώπινον}, and he denied the consubstantiality of his \textgreek{σάρξ}with ours. Yet he expressly guarded himself against Docetism, and against all speculation: \textgreek{Φυσιολογεῖν ἐμαυτῷ οὐκ ἐπιτρέπω}. He was really neither a philosopher nor a theologian, but only insisted on some theological opinions and points of doctrine with great tenacity and obstinacy.} All human attributes are transferred to the one subject, the humanized Logos. Hence it may and must be said: God is born, God suffered, God was crucified and died. He asserted, therefore, on the one hand, the capability of suffering and death in the Logos-personality, and on the other hand, the deification of the human in Christ.

Theodoret, in three dialogues composed in 447, attacked this Egyptian Eutychian type of doctrine as a beggar’s basket of Docetistic, Gnostic, Apollinarian, and other heresies,\footnote{Hence the title of the dialogues: \textgreek{Ἐρανιστής}, Beggar, and \textgreek{Πολύμορφος}, the Multiform. Under this name the Eutychian speaker is introduced. Theodoret also wrote an \textgreek{ἀπολογία ὑπὲρ Διοδώρου καὶ Θεοδώρου}which is lost.} and advocated the qualified Antiochian Christology, i.e., the doctrine of the unfused union of two natures in one person. Dioscurus accused him to the patriarch Domnus in Antioch of dividing the one Lord Christ into two Sons of God; and Theodoret replied to this with moderation. Dioscurus, on his part, endeavored to stir up the court in Constantinople against the whole church of Eastern Asia. Domnus and Theodoret likewise betook themselves to the capital, to justify their doctrine. The controversy now broke forth with greater violence, and concentrated on the person of Eutyches in Constantinople.

At a local synod of the patriarch Flavian at Constantinople in 448\footnote{\textgreek{Σύνοδος ἐδημοῦσα}. Its acts are incorporated in the acts of the council of Chalcedon, in Mansi, vi. 649 sqq.} Eutyches was charged with his error by Eusebius, bishop of Dorylaeum in Phrygia, and upon his wilful refusal, after repeated challenges, to admit the dyophysitism after the incarnation, and the consubstantiality of Christ’s body with our own, he was deposed and put under the ban of the church. On his way home, he was publicly insulted by the populace. The council confessed its faith that “Christ, after the incarnation, consisted of two natures\footnote{609 \textgreek{Ἐκ δύο φύσεων}, or, as others more accurately said, \textgreek{ἐν δύο φύσεσι},—an unessential difference, which reappears in the Creed of the council of Chalcedon. Comp. Mansi, tom. vi. fol. 685, and Neander, iv. p. 988. The first form may be taken also in a monophysite sense.} in one hypostasis and in one person, one Christ, one Son, one Lord.”

Both parties endeavored to gain the public opinion, and addressed themselves to distant bishops, especially to Leo I. of Rome. Leo, in 449, confirmed the decision of the council in several epistles, especially in a letter to Flavian, which forms an epoch in the history, of Christology, and in which he gave a masterly, profound, and clear analysis of the orthodox doctrine of two natures in one person.\footnote{This Epistola Dogmatica ad Flavianum (Ep. 28 in Ballerini, 24 in Quesnel), which Leo transmitted, with letters to the emperor and the emperor’s sister, Pulcheria, and the Robber Synod, by his legates, was afterwards formally approved at the council of Chalcedon in 451, and invested with almost symbolical authority. It may be found in the Opera Leonis, ed. Baller. tom. i. pp. 801-838; in Mansi, tom. v. fol. 1359; and in Hefele (Latin and German), ii. 335-346. Comp. on It also Walch, vi. p. 182 ff., and Baur, i. 809 ff.} But Eutyches had powerful friends among the monks and at the court, and a special patron in Dioscurus of Alexandria, who induced the emperor Theodosius II. to convoke a general council.

This synod met at Ephesus, in August, 449, and consisted of one hundred and thirty-five bishops. It occupies a notorious place in the chronique scandaleuse of church history. Dioscurus presided, with brutal violence, protected by monks and an armed soldiery; while Flavian and his friends hardly dared open their lips, and Theodoret was entirely excluded. When an explanation from Eusebius of Dorylaeum, who had been the accuser of Eutyches at the council of Constantinople, was presented, many voices exclaimed: “Let Eusebius be burnt; let him be burnt alive. As he has cut Christ in two, so let him be cut in two.”\footnote{Conc. Chalced. Actio i. in Harduin, tom. ii. fol. 161.} The council affirmed the orthodoxy and sanctity of Eutyches, who defended himself in person; adopted the twelve anathematisms of Cyril; condemned dyophysitism as a heresy, and deposed and excommunicated its advocates, including Theodoret, Flavian, and Leo. The three Roman delegates (the bishops Julius and Renatus, and the deacon Hilarus) dared not even read before the council the epistle addressed to it by Leo,\footnote{This, moreover, made reference to the famous Epistola Dogmatica, addressed to Flavian, which was also intended to be read before the council. Comp. Hefele, ii. 352.} and departed secretly, that they might not be compelled to subscribe its decisions.\footnote{Leo at least asserts this in reference to the deacon Hilarus. The two other delegates appear to have returned home before the council broke up. Renatus does not appear at all in the Acta, but Theodoret praises him for his courage at the Synod of Robbers. With the three delegates Leo sent also a notary, Dulcitius.} Flavian was so grossly maltreated by furious monks that he died of his wounds a few days later, in banishment, having first appealed to a new council. In his stead the deacon Anatolius, a friend and agent of Dioscurus, was chosen patriarch of Constantinople. He, however, afterwards went over to the orthodox party, and effaced the infamy of his elevation by his exquisite Greek hymns.

The conduct of these unpriestly priests was throughout so arbitrary and tyrannical, that the second council of Ephesus has ever since been branded with the name of the “Council of Robbers.”\footnote{\textgreek{Σύνοδος λῃστρική}, latrocinium Ephesinum; first so called by pope Leo in a letter to Pulcheria, dated July 20th, 451 (Ep. 95, ed. Ballerini, alias Ep. 75). The official Acta of the Robber Synod were read before the council of Chalcedon, and included in its records. These of themselves show dark enough. But with them must be compared the testimony of the defeated party, which was also rendered at the council of Chalcedon; the contemporaneous correspondence of Leo; and the accounts of the old historians. Comp. the details in Tillemont, Walch, Schröckh, Neander, and Hefele.} “Nothing,” Neander justly observes,\footnote{Kirchengesch. iv. p. 969 (2d Germ. ed. 1847).} “could be more contradictory to the spirit of the gospel than the fanatical zeal of the dominant party in this council for dogmatical formulas, in which they fancied they had Christ, who is spirit and life, although in temper and act they denied Him.” Dioscurus, for example, dismissed a charge of unchastity and other vices against a bishop, with the remark: “If you have an accusation against his orthodoxy, we will receive it; but we have not come together to pass judgment concerning unchastity.”\footnote{At the third session of the council of Chalcedon, Dioscurus himself was accused of gross intemperance and other evil habits. Comp. Hefele, ii. p. 429.} Thus fanatical zeal for doctrinal formulas outweighed all interests of morality, as if, as Theodoret remarks, Christ had merely prescribed a system of doctrine, and had not given also rules of life.

\xsection{The Ecumenical Council of Chalcedon, A.D. 451}

§ 141. The Ecumenical Council of Chalcedon, A.D. 451.

Comp. the Acta Concilii, together with the previous and subsequent epistolary correspondence, in Mansi (tom. vii.), Harduin (tom. ii.), and Fuchs, and the sketches of Evagrius: H. E. l. ii. c. 4; among later historians: Walch; Schröckh; Neander; Hefele, l.c. The latter, ii. 392, gives the literature in detail.

Thus the party of Dioscurus, by means of the court of the weak Theodosius II., succeeded in subjugating the Eastern church, which now-looked to the Western for help.

Leo, who occupied the papal chair from 440 to 461, with an ability, a boldness, and an unction displayed by none of his predecessors, and by few of his successors, and who, moreover, on this occasion represented the whole Occidental church, protested in various letters against the Robber Synod, which had presumed to depose him; and he wisely improved the perplexed state of affairs to enhance the authority of the papal see. He wrote and acted with imposing dignity, energy, circumspection, and skill, and with a perfect mastery of the question in controversy;—manifestly the greatest mind and character of his age, and by far the most distinguished among the popes of the ancient Church. He urged the calling of a new council in free and orthodox Italy, but afterwards advised a postponement, ostensibly on account of the disquiet caused in the West by Attila’s ravages, but probably in the hope of reaching a satisfactory result, even without a council, by inducing the bishops to subscribe his Epistola Dogmatica. \footnote{Respecting this apparent inconsistency of Leo, see Hefele, who considers it at length, ii. 387 ff.}

At the same time a political change occurred, which, as was often the case in the East, brought with it a doctrinal revolution. Theodosius died, in July, 450, in consequence of a fall from his horse; he left no male heirs, and the distinguished general and senator Marcian became his successor, by marriage with his sister Pulcheria, \footnote{Who, however, stipulated as a condition of the marriage, that she still be allowed to keep her vow of perpetual virginity. Marcian was a widower, sixty years of age, and had the reputation of great ability and piety. Some authors place him, as emperor, by the side of Constantineand Theodosius, or even above them. Comp. Leo’s Letters, Baronius (Annales), Tillemont (Emper. iii. 284), and Gibbon (at the end of ch. xxxiv.). The last-named author says of Marcian: “The zeal which he displayed for the orthodox creed, as it was established by the council of Chalcedon, would alone have inspired the grateful eloquence of the Catholics. But the behavior of Marcian, in a private life, and afterwards on the throne, may support a more rational belief, that he was qualified to restore and invigorate an empire, which had been almost dissolved by the successive weakness of two hereditary monarchs .... His own example gave weight to the laws which he promulgated for the reformation of manners.”} who favored Pope Leo and the dyophysite doctrine. The remains of Flavian were honorably interred, and several of the deposed bishops were reinstated.

To restore the peace of the empire, the new monarch, in May, 451, in his own name and that of his Western colleague, convoked a general council; not, however, to meet in Italy, but at Nicaea, partly that he might the better control it partly that he might add to its authority by the memories of the first ecumenical council. The edict was addressed to the metropolitans, and reads as follows:

“That which concerns the true faith and the orthodox religion must be preferred to all other things. For the favor of God to us insures also the prosperity of our empire. Inasmuch, now, as doubts have arisen concerning the true faith, as appears from the letters of Leo, the most holy archbishop of Rome, we have determined that a holy council be convened at Nicaea in Bithynia, in order that by the consent of all the truth may be tested, and the true faith dispassionately and more explicitly declared, that in time to come no doubt nor division may have place concerning it. Therefore let your holiness, with a convenient number of wise and orthodox bishops from among your suffragans, repair to Nicaea, on the first of September ensuing. We ourselves also, unless hindered by wars will attend in person the venerable synod.”\footnote{This promise was in fact fulfilled, although only at one session, the sixth.}

Leo, though dissatisfied with the time and place of the council, yielded, sent the bishops Paschasinus and Lucentius, and the priest Boniface, as legates, who, in conjunction with the legates already in Constantinople, were to represent him at the synod, over which Paschasinus was to preside in his name.\footnote{Evagrius, H. E. ii. c. 4: “The bishops Paschasinus and Lucentius, and the presbyter Boniface, were the representatives of Leo, archpriest of the elder Rome.” Besides them bishop Julianof Cos, Leo’s legate at Constantinople, also frequently appears in the council, but he had his seat among the bishops, not the papal delegates.}

The bishops assembled at Nicaea, in September, 451, but, on account of their turbulent conduct, were soon summoned to Chalcedon, opposite Constantinople, that the imperial court and senate might attend in person, and repress, as far as possible, the violent outbreaks of the religious fanaticism of the two parties. Here, in the church of St. Euphemia, on a hill commanding a magnificent prospect, and only two stadia or twelve hundred paces from the Bosphorus, the fourth ecumenical council was opened on the 8th of October, and sat till the lst of November. In number of bishops it far exceeded all other councils of the ancient Church,\footnote{There are only imperfect registers of the subscriptions yet extant, and the statements respecting the number of members vary from 520 to 630.} and in doctrinal importance is second only to the council of Nicaea. But all the five or six hundred bishops, except the papal delegates and two Africans, were Greeks and Orientals. The papal delegates had, therefore, to represent the whole of Latin Christendom. The imperial commissioners,\footnote{\textgreek{Ἄρχοντες}, judices. There were six of them.} who conducted the external course of the proceedings, in the name of the emperor, with the senators present, sat in the middle of the church, before the screen of the sanctuary. On the left sat the Roman delegates, who, for the first time at an ecumenical council, conducted the internal proceedings, as spiritual presidents; next them sat Anatolius, of Constantinople, Maximus, of Antioch, and most of the bishops of the East;—all opponents of Eutychianism. On the right sat Dioscurus, of Alexandria (who, however, soon had to give up his place and sit in the middle), Juvenal, of Jerusalem, and the other bishops of Egypt, Illyricum, and Palestine;—the Eutychians.

The proceedings were, from the outset, very tumultuous, and the theological fanaticism of the two parties broke out at times in full blaze, till the laymen present were compelled to remind the bishops of their clerical dignity.\footnote{Such tumultuous outcries (\textgreek{ἐκβοήσεις δημοτικαί}), said the commissioners and senators, ill-beseemed bishops, and were of no advantage to either side.} When Theodoret, of Cyrus, was introduced, the Orientals greeted him with enthusiasm, while the Egyptians cried: “Cast out the Jew, the enemy of God, the blasphemer of Christ!” The others retorted, with equal passion: “Cast out the murderer Dioscurus! Who is there that knows not his crimes?” The feeling against Nestorius was so strong, that Theodoret could only quiet the council by resolving (in the eighth session) to utter the anathema against his old friend, and against all who did not call Mary “mother of God,” and who divided the one Christ into two sons. But the abhorrence of Eutyches and the Council of Robbers was still stronger, and was favored by the court. Under these influences most of the Egyptians soon went over to the left, and confessed their error, some excusing themselves by the violent measures brought to bear upon them at the Robber Synod. The records of that Synod, and of the previous one at Constantinople (in 448), with other official documents, were read by the secretaries, but were continually interrupted by incidental debates, acclamations, and imprecations, in utter opposition to all our modern conceptions of parliamentary decorum, though experience is continually presenting us with fresh examples of the uncontrollable vehemence of human passions in excited assemblies.

So early as the close of the first session the decisions of the Robber Synod had been annulled, the martyr Flavian declared orthodox, and Dioscurus of Alexandria, Juvenal of Jerusalem, and other chiefs of Eutychianism, deposed. The Orientals exclaimed: “Many years to the Senate! Holy God, holy mighty, holy immortal God, have mercy upon us. Many years to the emperors! The impious must always be overthrown! Dioscurus, the murderer [of Flavian], Christ has deposed! This is a righteous judgment, a righteous senate, a righteous council!”

Dioscurus was in a subsequent session three times cited in vain to defend himself against various charges of avarice, injustice, adultery, and other vices, and divested of all spiritual functions; while the five other deposed bishops acknowledged their error, and were readmitted into the council.

At the second session, on the 10th of October, Dioscurus having already departed, the Nicaeno-Constantinopolitan symbol, two letters of Cyril (but not his anathemas), and the famous Epistola Dogmatica of Leo to Flavian, were read before the council amid loud applause—the bishops exclaiming: “That is the faith of the fathers! That is the faith of the apostles! So we all believe! So the orthodox believe Anathema to him who believes otherwise! Through Leo, Peter has thus spoken. Even so did Cyril teach! That is the true faith.”\footnote{Mansi, tom. vi. 971: \textgreek{αὕτη ἡ πίστις τῶν πατέρων, αὕτη ἡ πίστις τῶν ἀποστόλων, παντες οὕτω πιστεύομεν, οἱ ὀρθόδοξοι οὕτω πιστεύουσιν, ἀνάθεμα τῷ μὴ οὕτω πιστεύοντι, κ.τ.λ.}}

At the fifth and most important session, on the 22d of October, the positive confession of faith was adopted, which embraces the Nicaeno-Constantinopolitan symbol, and then, passing on to the point in controversy, expresses itself as follows, almost in the words of Leo’s classical epistle:\footnote{Complete in Mansi, tom. vii. f. 111-118, The Creed is also given by Evagrius, ii. 4.}

“Following the holy fathers, we unanimously teach one and the same Son, our Lord Jesus Christ, complete as to his Godhead, and complete as to his manhood; truly God, and truly man, of a reasonable soul and human flesh subsisting; consubstantial with the Father as to his Godhead, and consubstantial also with us as to his manhood;\footnote{\textgreek{Ὁμοούσιος} is used in both clauses, though with a shade of difference: Christ’s \emph{homoousia} with the Father implies numerical unity or identity of substance (God being one in essence, \emph{monoousios}); Christ’s \emph{homoousia} with men means only generic unity or equality of nature. Compare the remarks in § 130, p. 672 f.} like unto us in all things, yet without sin;\footnote{\textgreek{Ὕ Ενα καὶ αὐτὸν υἱὸν τὸν κύριον ἡμῶν Ἰ. Χριστὸν τὸν αὐτὸν ἐν θεότητι καὶ τέλειον τὸν αὐτὸν ἐν ἀνθρωπότητι, θεὸν ἀληθῶς καὶ ἄθρωπον ἀληθῶς τὸν αὐτὸν, ἐκ ψυχῆς λογικῆς}[against Apollinaris] \textgreek{καὶ σώματος, ὁμοιούσιον τῷ Πατρὶ κατὰ τὴν θεότητα, καὶ ὀμοούσιον τὸν αὐτὸν ἠμῖν κατὰτὴν ἀνθρωπότητα, κατὰ πάντα ὅμοιον ἠμῖν χωρὶς ἁμαρτίας.}} as to his Godhead begotten of the Father before all worlds, but as to his manhood, in these last days born, for us men and for our salvation, of the Virgin Mary, the mother of God;\footnote{\textgreek{Τῆς θεοτόκου}, against Nestorius. This, however, is immediately after modified by the phrase \textgreek{κατὰ τὴν ἀνθρωπότητα}(in distinction from \textgreek{κατὰ τὴν θεότητα}). Mary was the mother not merely of the \emph{human} nature of Jesus, but of the theanthropic \emph{person} Jesus Christ; not, however, according to his eternal Godhead, but according to his humanity. In like manner, the subject of the passion was the theanthropic \emph{person}, yet not according to his divine impassible nature, but according to his human nature.} one and the same Christ, Son, Lord, Only-begotten, known in (of) two natures, \footnote{\textgreek{Ἐν δύο φύσεσιν} and the Latin translation, \emph{in duabus naturis}, is directed against Eutyches. The present Greek text reads, it is true, \textgreek{ἐκ δύο φύσεων}, which, however, signifies, and according to the connection, can only signify, essentially the same thing, but is also capable of being understood in an Eutychian and Monophysite sense, namely, that Christ has arisen from the confluence of two natures, and since the incarnation has only one nature. Understood in this sense, Dioscurus at the council was very willing to accept the formula \textgreek{ἐκ δύο φύσεων}. But for this very reason the Orientals, and also the Roman legates, protested with one voice against \textgreek{ἐκ}and insisted upon another formula with \textgreek{ἐν}which was adopted. Baur (l.c. i. p. 820 f.) and Dorner (ii. p. 129) assert that \textgreek{ἐκ} is the accurate and original expression, and is a concession to Monophysitism, that It also agrees better (?) with the verb \textgreek{γνωρίζομεν}(to recognize by certain tokens) but that it was from the very beginning changed by the Occidentals into \textgreek{ἐν}. But we prefer the view of Gieseler, Neander (iv. 988), Hefele (ii. 451 f), and Beck (Dogmengeschichte, p. 251), that \textgreek{ἐν δύο φύσεσιν} was the original reading of the symbol, and that It was afterwards altered in the interest of Monophysitism. This is proved by the whole course of the proceedings at the fifth session of the council of Chalcedon, where the expression \textgreek{ἐκ δύο φύσεων} was protested against, and is proved by the testimony of the abbot Euthymius, a cotemporary, and by that of Severus, Evagrius, and Leontius of Byzantium. Severus, the Monophysite patriarch of Antioch since 513, charges the fathers of Chalcedon with the inexcusable crime of having taught: \textgreek{ἐν δύο φύσεσιν ἀδιαιρέτοις γνωρίζεσθαι τὸν Θεόν}(see Mansi, vii. 839). Evagrius (H. E. ii. 5) maintains that both formulas amount to essentially the same thing, and reciprocally condition each other. Dorner also affirms the same. His words are: ” The Latin formula has ’to acknowledge Christ as Son \emph{in} two natures,’ the Greek has ’to recognize Christ as Son \emph{from} two natures,’ which is plainly the same thought. The Latin formula is only a free, but essentially faithful translation, only that its coloring expresses somewhat more definitely still Christ’s subsisting in two natures, and is therefore more literally conformable to the Roman type of doctrine” (l.c. ii. p. 129 f.).} without confusion, without conversion, without severance, and without division;\footnote{\textgreek{Ἀσύγχύτως, ἀτρέπτως}[against Eutyches], \textgreek{ἀδιαρέτως, ἀχωρίστως}[against Nestorius]—\textgreek{γνωριζόμενον}.} the distinction of the natures being in no wise abolished by their union, but the peculiarity of each nature being maintained, and both concurring in one person and hypostasis.\footnote{\textgreek{Εἰς ἓν πρόσωπον καὶ μίαν ὑπόστασιν}.} We confess not a Son divided and sundered into two persons, but one and the same Son, and Only-begotten, and God-Logos, our Lord Jesus Christ, even as the prophets had before proclaimed concerning him, and he himself hath taught us, and the symbol of the fathers hath handed down to us.

“Since now we have drawn up this decision with the most comprehensive exactness and circumspection, the holy and ecumenical synod \footnote{\textgreek{Ἡ ἁγία καὶ οἰκουμενικὴ σύνοδος}.} hath ordained, that no one shall presume to propose, orally, or in writing, another faith, or to entertain or teach it to others; and that those who shall dare to give another symbol or to teach another faith to converts from heathenism or Judaism, or any heresy, shall, if they be bishops or clergymen, be deposed from their bishopric and spiritual function, or if they be monks or laymen, shall be excommunicated.”

After the public reading of this confession, all the bishops exclaimed: “This is the faith of the fathers; this is the faith of the apostles; to this we all agree; thus we all think.

The symbol was solemnly ratified at the sixth session (Oct. 25th), in the presence of the emperor and the empress. The emperor thanked Christ for the restoration of the unity of faith, and threatened all with heavy punishment, who should thereafter stir up new controversies; whereupon the synod exclaimed: “Thou art both priest and king, victor in war, and teacher of the faith.”

At its subsequent sessions the synod was occupied with the appeal of Ibas, bishop of Edessa, who had been deposed by the Robber Synod, and was now restored; with other cases of discipline; with some personal matters; and with the enactment of twenty-eight canons, which do not concern us here.\footnote{Respecting the famous 28th canon of the council, which gives the bishop of Constantinople equal rights with the bishop of Rome, and places him next after him in rank, Comp. above § 56 (p. 279 ff.).}

The emperor, by several edicts, gave the force of law to the decisions of the council, and commanded that all Eutychians should be banished from the empire, and their writings burned.\footnote{Eutyches, who, in the very beginning of the controversy, said of himself, that he had lived seventy years a monk died probably soon after the meeting of the council. Dioscurus was banished to Gangra, in Paphlagonia, and lived tin 454. Comp. Schröckh, Th. xviii. p. 492.} Pope Leo confirmed the doctrinal confession of the council, but protested against the twenty-eighth canon, which placed the patriarch of Constantinople on an equality with him. Notwithstanding these ratifications and rejoicings, the peace of the Church was only apparent, and the long Monophysite troubles were at hand.\footnote{Dorner judges very unfavorably of the council of Chalcedon (ii. p. 83), and denies it all vocation, inward or outward, to render a positive decision of the great question in controversy; forgetting that the third ecumenical council, which condemned Nestorius, was, in Christian spirit and moral dignity, decidedly inferior to the fourth. “Notwithstanding its 630 bishops,” says he (ii. 130), “it is very far from being able to claim canonical authority. The fathers of this council exhibIt neither the harmony of an assembly animated by the Holy Ghost, nor that certainty of judgment, past wavering and inconsistency, nor that manly courage in maintaining a well-gained conviction, which is possible where, out of antitheses long striving for unity, a bright and clear persuasion, shared by the general body, has arisen.” Kahnis (Der Kirchenglaube, Bd. ii. 1864, p. 89) judges as follows: “The significance of the Chalcedonian symbol does not lie in the ecumenical character of this council, for ecumenical is an exceedingly elastic idea; nor in its results being a development of those of the council of Ephesus (431), for, while at Ephesus the doctrine of the unity, here that of the distinction, in Christ’s person, was the victorious side; nor in the spirit with which all the proceedings were conducted, for passions, intrigues, political views, tumultuous disorder, \&c., prevailed in it in abundant measure: but it lies rather in the unity of acknowledgment which it has received in the Church, even to our day, and in the inner unity of its definitions.”}

But before we proceed to these, we must enter into a more careful exposition of the Chalcedonian Christology, which has become the orthodox doctrine of Christendom.

\xsection{The Orthodox Christology--Analysis and Criticism}

§ 142. The Orthodox Christology—Analysis and Criticism.

The first council of Nicaea had established the eternal preexistent Godhead of Christ. The symbol of the fourth ecumenical council relates to the incarnate Logos, as he walked upon earth and sits on the right hand of the Father, and it is directed against errors which agree with the Nicene Creed as opposed to Arianism, but put the Godhead of Christ in a false relation to his humanity. It substantially completes the orthodox Christology of the ancient Church; for the definitions added by the Monophysite and Monothelite controversies are few and comparatively unessential.

The same doctrine, in its main features, and almost in its very words (though with less definite reference to Nestorianism and Eutychianism), was adopted in the second part of the pseudo-Athanasian Creed,\footnote{Comp. above § 132.} and in the sixteenth century passed into all the confessions of the Protestant churches.\footnote{Comp. my article cited in § 132 upon the Symbolum \emph{Quicunque}. One of the briefest and clearest Protestant definitions of the person of Christ in the sense of the Chalcedonian formula, is the one in the Westminster (Presbyterian) Shorter Catechism: “Dominus Jesus Christus est electorum Dei Redemptor unicus, qui eternus Dei filius cum esset factus est homo; adeoque fuit, est eritque \textgreek{θεάνθρωπος}e [in] naturis duabus distinctis persona unica in sempiternum or, as it is in English: ”\emph{The only Redeemer of God’s elect is the Lord Jesus Christ, who, being the eternal Son of God, became man, and so was, and continueth to be, God and Man, in two distinct natures, and one person forever}.” The Westminster Confession formulates this doctrine (ch. viii. sec 21) in very nearly the words of the Chalcedonian symbol: “The Son of God, the second person in the Trinity, being very and eternal God, of one substance and equal with the Father, did, when the fulness of time was come, take upon Him man’s nature, with all the essential properties and common infirmities thereof, yet without sin; being conceived by the power of the Holy Ghost, in the womb of the Virgin Mary, of her substance. So that two whole, perfect, and distinct natures,—the Godhead and the manhood,—were inseparably joined together in one person, without conversion, composition, or confusion. Which person is very God and very man, yet one Christ, the only Mediator between God and man.”} Like the Nicene doctrine of the Trinity, it is the common inheritance of Greek, Latin, and Evangelical Christendom; except that Protestantism, here as elsewhere, reserves the right of searching, to ever new depths, the inexhaustible stores of this mystery in the living Christ of the Gospels and the apostolic writings.\footnote{The Lutheran Church has framed the doctrine of a threefold \emph{communicatio} \emph{idiomatum}, and included It in the Formula Concordiae. The controversy between the Lutheran theologians of Giessen and Tübingen, in the seventeenth century, concerning the \textgreek{κτῆσις}(the possession), the \textgreek{ξρῆσις} (the use), the \textgreek{κρύψις}(the secret use), and the \textgreek{κένωσις}(the entire abdication) of the divine attributes by the incarnate Logos, led to no definite results, and was swallowed up in the thirty years’ war. It has been resumed in modified form by modern German divines.}

The person of Jesus Christ in the fulness of its theanthropic life cannot be exhaustively set forth by any formulas of human logic. Even the imperfect, finite personality of man has a mysterious background, that escapes the speculative comprehension; how much more then the perfect personality of Christ, in which the tremendous antitheses of Creator and creature, Infinite and finite, immutable, eternal Being and changing, temporal becoming, are harmoniously conjoined! The formulas of orthodoxy can neither beget the true faith, nor nourish it; they are not the bread and the water of life, but a standard for theological investigation and a rule of public teaching.\footnote{Comp. Cunningham (Historical Theology, vol. i. p. 319): “The chief use now to be made of an examination of these controversies [the Eutychian and Nestorian] is not so much to guard us against errors [?] which may be pressed upon us, and into which we may be tempted to fall, but rather to aid us in forming clear and definite conceptions of the truths regarding the person of Christ, which all profess to believe; in securing precision and accuracy of language in explaining them, and especially to assist us in realizing them; in habitually regarding as great and actual realities the leading features of the constitution of Christ’s person, which the word of God unfolds to us.”}

Such considerations suggest the true position and the just value of the Creed of Chalcedon, against both exaggeration and disparagement. That symbol does not aspire to comprehend the Christological mystery, but contents itself with setting forth the facts and establishing the boundaries of orthodox doctrine. It does not mean to preclude further theological discussion, but to guard against such erroneous conceptions as would mutilate either the divine or the human in Christ, or would place the two in a false relation. It is a light-house, to point out to the ship of Christological speculation the channel between Scylla and Charybdis, and to save it from stranding upon the reefs of Nestorian dyophysitism or of Eutychian monophysitism. It contents itself with settling, in clear outlines, the eternal result of the theanthropic process of incarnation, leaving the study of the process itself to scientific theology. The dogmatic letter of Leo, it is true, takes a step beyond this, towards a theological interpretation of the doctrine; but for this very reason it cannot have the same binding and normative force as the symbol itself.

As the Nicene doctrine of the Trinity stands midway between tritheism and Sabellianism, so the Chalcedonian formula strikes the true mean between Nestorianism and Eutychianism.

It accepts dyophysitism; and so far it unquestionably favored and satisfied the moderate Antiochian party rather than the Egyptian.\footnote{Accordingly in Leo’s \emph{Epistola Dogmatica} also, which was the basis of the Creed, Nestorius is not even mentioned, while Eutyches, on the other hand, is refuted at length. But in a later letter of Leo, addressed to the emperor, a.d.457 (Ep. 156, ed. Ballerini), he classes Nestorius and Eutyches together, as equally dangerous heretics. The Creed of Chalcedon is also regarded by Baur, Niedner, and Dorner as exhibiting a certain degree of preference for the Nestorian dyophysitism.} But at the same time it teaches with equal distinctness, in opposition to consistent Nestorianism, the inseparable unity of the person of Christ.

The following are the leading ideas of this symbol:

1. A true incarnation of the Logos, or of the second person in the Godhead.\footnote{\textgreek{Ἐνανθρώπησις Θεοῦ, ἐνσάρκωσις}, incarnatio,—in distinction from a mere \textgreek{συνάφεια}, conjunctio, or \textgreek{σχετική ἕνωσις}, of the divine and human, by \textgreek{πρόσληψις}(from, \textgreek{προσλαμβάνω}), \emph{assumptio}, of the human, and \textgreek{ἐνοίκησις}the divine; and on the other hand, from a \textgreek{φυσικὴ ἕνωσις}or \textgreek{κρᾶσις, σύγχυσις ,}or \textgreek{σάρκωσις}in the sense of transmutation. The diametrical opposite of the \textgreek{ἐνανθρώπησις Θεοῦ}is the heathen \textgreek{ἀποθέωσις ἀνθρώπου}.} The motive is the unfathomable love of God; the end, the redemption of the fallen race, and its reconciliation with God. This incarnation is neither a conversion of God into a man, nor a conversion of a man into God; neither a humanizing of the divine, nor a deification or apotheosis of the human; nor on the other hand is it a mere outward, transitory connection of the two factors; but an actual and abiding union of the two in one personal life.

It is primarily and pre-eminently a condescension and self-humiliation of the divine Logos to human nature, and at the same time a consequent assumption and exaltation of the human nature to inseparable and eternal communion with the divine person. The Logos assumes the body, soul, and spirit of man, and enters into all the circumstances and infirmities of human life on earth, with the single exception of sin, which indeed is not an essential or necessary element of humanity, but accidental to it. “The Lord of the universe,” as Leo puts the matter in his epistle, “took the form of a servant; the impassible God became a suffering man; the Immortal One submitted himself to the dominion of death; Majesty assumed into itself lowliness; Strength, weakness; Eternity, mortality.” The same, who is true God, is also true man, without either element being altered or annihilated by the other, or being degraded to a mere accident.

This mysterious union came to pass, in an incomprehensible way, through the power of the Holy Ghost, in the virgin womb of Mary. But whether the miraculous conception was only the beginning, or whether it at the same time completed the union, is not decided in the Creed of Chalcedon. According to his human nature at least Christ submitted himself to the laws of gradual development and moral conflict, without which, indeed, he could be no example at all for us.

2. The precise distinction between nature and person. Nature or substance is the totality of powers and qualities which constitute a being; person is the Ego, the self-conscious, self-asserting, and acting subject. There is no person without nature, but there may be nature without person (as in irrational beings).\footnote{Compare the weighty dissertation of Boethius: De duabus naturis et una persona Christi, adversus Eutychen et Nestorium (Opera, ed. Basil., 1546, pp. 948-957), in which he defines \emph{natura} (\textgreek{φύσις}or \textgreek{οὐσία}), \emph{substantia} (\textgreek{ὑπόστασις}), and \emph{persona} (\textgreek{πρόσωπον}).”\emph{Natura},” he says, “est cujuslibet substantia specificata proprietas; persona vero rationabilis naturae individua subsistentia.”} The Church doctrine distinguishes in the Holy Trinity three persons (though not in the ordinary human sense of the word) in one divine nature or substance which they have in common; in its Christology it teaches, conversely, two natures in one person (in the usual sense of person) which pervades both. Therefore it cannot be said: The Logos assumed a human person,\footnote{\textgreek{Τέλειον ἄνθρωπον εἴληφε}, as Theodore of Mopsuestia and the strict Nestorians expressed themselves.} or united himself with a definite human individual: for then the God-Man would consist of two persons; but he took upon himself the human nature, which is common to all men; and therefore he redeemed not a particular man, but all men, as partakers of the same nature or substance.\footnote{As Augustinesays: Deus Verbum non accepit personam hominis, sed naturam, et in eternam personam divinitatis accepit temporalem substantiam carnis. And again: “Deus naturam nostram, id est, animam rationalem carnemque hominis Christi suscepit.” (De corrept. et grat. §30, tom. x. f. 766.) Comp. Johannes Damascenus De fide orthod. iii. c. 6, II. The Anglican theologian, Richard Hooker, styled on account of his sober equipoise of intellect “the judicious Hooker,” sets forth this point of the Church doctrine as follows: “He took not angels but the seed of Abraham. It pleased not the Word or Wisdom of God to take to itself some one person amongst men, for then should that one have been advanced which was assumed, and no more, but Wisdom to the end she might save many built her house of that Nature which is common unto all, she made not \emph{this or that man} her habitation, but dwelt \emph{in us}. If the Son of God had taken to himself a man now made and already perfected, it would of necessity follow, that there are in Christ two persons, the one assuming, and the other assumed; whereas the Son of God did not assume a man’s \emph{person} into his own, but a man’s \emph{nature} to his own person; and therefore took \emph{semen}, the seed of Abraham, the very first original and element of our nature, before it was come to have any personal human subsistence. The flesh and the conjunction of the flesh with God began both at one instant; his making and taking to himself our flesh was but one act, so that in Christ there is no personal subsistence but one, and that from everlasting. By taking only the nature of man he still continueth one person, and changeth but the manner of his subsisting, which was before in the glory of the Son of God, and is now in the habIt of our flesh.” (Ecclesiastical Polity, book v. ch. 52, in Keble’s edition of Hooker’s works, vol. ii. p. 286 f.) In just the same manner Anastasius Sinaita and John of Damascus express themselves. Comp. Dorner, ii. p. 183 ff. Hooker’s allusion to Heb. ii. 16 (\textgreek{οὐ γὰρ δήπου ἀγγέλων ἐπιλαμβάνεται, ἀλλὰ σπέρματοσ Ἁβραὰμ ἐπιλαμβάνεται}), it may be remarked, rests upon a false interpretation, since \textgreek{ἐπιλαμβάνεσθαι}does not refer to the incarnation, but signifies: to take hold of in order to help or redeem (as in Sirach, iv. 11). Comp. \textgreek{βοηθῆσαι}, Heb. ii. 18.} The personal Logos did not become an individual ἄνθρωπος, but σάρξ, flesh, which includes the whole of human nature, body, soul, and spirit. The personal self-conscious Ego resides in the Logos. But into this point we shall enter more fully below.

3. The result of the incarnation, that infinite act of divine love, is the God-Man. Not a (Nestorian) double being, with two persons; nor a compound (Apollinarian or Monophysite) middle being a tertium quid, neither divine nor human; but one person, who is both divine and human. Christ has a rational human soul, and—according to a definition afterwards added—a human will,\footnote{The sixth ecumenical council, held at Constantinople, a.d.680, condemned monothelitism, and decided in favor of dyothelitism, or the doctrine of two wills (or volitions) in Christ, which are necessary to the ethical conflict and victory of his own life and to his office as an example for us. This council teaches (Mansi, tom xi. 637): \textgreek{Δύο φυσικὰς θελήσεις ἤτοι θελήματα ἐν αὐτῷ καὶ δύο φυσικὰς ἐνεργείας ἀδιαιρέτως, ἀτρέπτως , ἀμερίστως, ἀσυγχύτως} … \textgreek{κηρύττομεν.} These wills are not opposite to one another, but the human will is ever in harmony with the divine, and in all things obedient to it. “Not my will, but thine be done:” therein is found the distinction and the unity.} and is therefore in the full sense of the word the Son of man; while yet at the same time he is the eternal Son of God in one person, with one undivided self-consciousness.

4. The duality of the natures. This was the element of truth in Nestorianism, and on this the council of Chalcedon laid chief stress, because this council was principally concerned with the condemnation of Eutychianism or monophysitism, as that of Ephesus (431) had been with the condemnation of Nestorianism, or abstract dyophysitism. Both views, indeed, admitted the distinction of the natures, but Eutychianism denied it after the act of the incarnation, and (like Apollinarianism) made Christ a middle being, an amalgam, as it were, of the two natures, or, more accurately, one nature in which the human element is absorbed and deified.

Against this it is affirmed by the Creed of Chalcedon, that even after the incarnation, and to all eternity, the distinction of the natures continues, without confusion or conversion,\footnote{\textgreek{Ἀσυγχύτως}and \textgreek{ἀτρεπτως .}} yet, on the other hand, without separation or division,\footnote{\textgreek{Ἀδιαιρέτως} and \textgreek{ἀχωρίστως}.} so that the divine will remain ever divine, and the human, ever human,\footnote{“Tenet,” says Leo, in his epistle to Flavian, “sine defectu proprietatem suam utraque natura, et sicut formam servi Dei formam non adimit, ita formam Dei servi forma non minuit .... Agit utraque cum alterius communione quod Proprium est; Verbo scilicet operante quod Verbi est, et came exsequente quod carnis est. Unum horum coruscat miraculis, aliud succumbIt injuriis.”} and yet the two have continually one common life, and interpenetrate each other, like the persons of the Trinity.\footnote{Here belongs John of Damascus’ doctrine of the \textgreek{περιχώρησις}, \emph{Permeatio, circummeatio, circulatio, circumincessio, intercommunio}, or reciprocal indwelling and pervasion, which has relation not merely to the Trinity but also to Christology. The verb \textgreek{περιχωρεῖν}, is, so far as I know, first applied by Gregory of Nyasa (Contra Apollinarium) to the interpenetration and reciprocal pervasion of the two natures in Christ. On this rested also the doctrine of the exchange or communication of attributes, \textgreek{ἀντίδοσις , ἀντιμετάστασις, κοινωνία ἰδιωμάτων}, \emph{communicatio idiomatum}. The \textgreek{ἀντιμετάστασις τῶν ὀνομάτων}, also \textgreek{ἀντιμεθίστασις}, \emph{transmutatio proprietalum}, transmutation of attributes, is, strictly speaking, not identical with \textgreek{ἀντίδοσις}, but a deduction from it, and the rhetorical expression for it. The doctrine of the communicatio idiomatum, however, awaited a full development much later, in the Lutheran church, where great subtlety was employed in perfecting it. This Lutheran doctrine has never found access into the Reformed church, and least of all the ubiquitarian hypothesis invented as a prop to consubstantiation; although a certain measure of truth lies at the basis of this, if it is apprehended dynamically, and not materially.}

The continuance of the divine nature unaltered is involved in its unchangeableness, and was substantially conceded by all parties. The controversy, therefore, had reference only to the human nature.

And here the Scriptures are plainly not on the Eutychian side. The Christ of the Gospels by no means makes the impression of a person in whom the human nature had been absorbed, or extinguished, or even weakened by the divine; on the contrary, he appears from the nativity to the sepulchre as genuinely and truly human in the highest and fairest sense of the word. The body which he had of the substance of Mary, was born, grew, hungered and thirsted, slept and woke, suffered and died, and was buried, like any other human body. His rational soul felt joy and sorrow, thought, spoke, and acted after the manner of men. The only change which his human nature underwent, was its development to full manhood, mental and physical, in common with other men, according to the laws of growth, yet normally, without sin or inward schism; and its ennoblement and completion by its union with the divine.

5. The unity of the person.\footnote{The \textgreek{ἕνωσις καθ  ὑπόστασιν,}or \textgreek{ἕνωσις ὑποστατική}, unio hypostatica or personalis, unitas personae. The unio personalis is the status unionis, the result of the unitio or incarnatio.} This was the element of truth in Eutychianism and the later monophysitism, which, however, they urged at the expense of the human factor. There is only one and the self-same Christ, one Lord, one Redeemer. There is an unity in the distinction, as well as a distinction in the unity. “The same who is true God,” says Leo, “is also true man, and in this unity there is no deceit; for in it the lowliness of man and the majesty of God perfectly pervade one another .... Because the two natures make only one person, we read on the one hand: ’The Son of man came down from heaven’ (John iii. 13), while yet the Son of God took flesh from the Virgin; and on the other: ’The Son of God was crucified and buried’ (1 Cor. ii. 8), while yet he suffered not in his Godhead as co-eternal and consubstantial with the Father, but in the weakness of human nature.”

Here again the Chalcedonian formula has a firm and clear basis in Scripture. In the gospel history this personal unity everywhere unmistakably appears. The self-consciousness of Christ is not divided. It is one and the self-same theanthropic subject that speaks, acts, and suffers, that rises from the dead, ascends to heaven, sits at the right hand of God, and shall come again in glory to judge the quick and the dead.

The divine and the human are as far from forming a double personality in Christ, as the soul and the body in man, or as the regenerate and the natural life in the believer. As the human personality consists of such a union of the material and the spiritual natures that the spirit is the ruling principle and personal centre: so does the person of Christ consist in such a union of the human and the divine natures that the divine nature is the seat of self-consciousness, and pervades and animates the human.\footnote{Comp. the Athanasian Creed: “Sicut anima rationalis et caro unus est homo, ita Deus et homo unus est Christus.” In the same does Augustineexpress himself, and indeed this passage in the Creed, as well as several others, appears to be taken from him. Dr. Shedd (History of Christian Doctrine, i. p. 402) carries out vividly this analogy of the human personality with that of Christ, as follows: “This union of the two natures in one self-conscious Ego may be illustrated by reference to man’s personal constitution. An individual man is one person. But this one person consists of two natures,—a material nature and a mental nature. The personality, the self-consciousness, is the resultant of the \emph{union} of the two. Neither one of itself makes the person.” [This is not quite exact. Personality lies in the reasonable soul, which can maintain its self-conscious existence without the body, even as in Christ His personality resides in the divine nature, as Dr. Shedd himself clearly states on p. 406.] “Both body and soul are requisite in order to a complete individuality. The two natures do not make two individuals. The material nature, taken by itself, is not the man; and the mental part, taken by itself, is not the man. But only the union of the two is. Yet in this intimate union of two such diverse substances as matter and mind, body and soul, there is not the slightest alteration of the properties of each substance or nature. The body of a man is as truly and purely material as a piece of granite; and the immortal mind of a man is as truly and purely spiritual and immaterial as the Godhead itself. Neither the material part nor the mental part, taken by itself, and in separation, constitutes the personality; otherwise every human individual would be two persons in juxtaposition. There is therefore a material ’nature,’ but no material ’person,’ and there is a mental ’nature,’ but no mental ’person.’ The person is the \emph{union} of these two natures, and is not to be denominated either material or mental, but \emph{human}. In like manner the person of Christ takes its denomination of \emph{theanthropic}, or \emph{divine-human}, neither from the divine nature alone, nor the hurnan nature alone, but from the \emph{union} of both natures.”}

I may refer also to the familiar ancient analogy of the fire and the iron.

6. The whole work of Christ is to be referred to his person, and not to be attributed to the one or the other nature exclusively. It is the one divine-human Christ, who wrought miracles of almighty power,—by virtue of the divine nature dwelling in him,—and who suffered and was buried,—according to his passible, human nature. The person was the subject, the human nature the seat and the sensorium, of the passion. It is by this hypostatical union of the divine and the human natures in all the stages of the humiliation and exaltation of Christ, that his work and his merits acquire an infinite and at the same time a genuinely human and exemplary significance for us. Because the God-Man suffered, his death is the reconciliation of the world with God; and because he suffered as Man, he has left us an example, that we should follow his steps.\footnote{Here also the orthodox Protestant theology is quite in agreement with the old Catholic. We cite two examples from the two opposite wings of English Protestantism. The Episcopalian theologian, Richard Hooker, says, with evident reference to the above-quoted passage from the letter of Leo: “To Christ we ascribe both working of wonders and suffering of pains, we use concerning Him speeches as well of humility as of divine glory, but the one we apply unto that nature which He took of the Virgin Mary, the other to that which was in the beginning” (Eccles. Polity, book v. ch. 52, vol. ii. p, 291, Keble’s edition). The great Puritan theologian of the seventeenth century, John Owen, says, yet more explicitly: “In all that Christ did as the King, Priest, and Prophet of the church,—in all that He did and suffered, in all that He continueth to do for us, in or by virtue of whether nature soever it be done or wrought,—it is not to be considered as the act and work of this or that nature in Him alone, but it is the act and work of the whole person,—of Him that is both God and man in one person.” (Declaration of the Glorious mystery of the Person of Christ; chap. xviii., in Owen’s Works, vol. i. p. 234). Comp. also the admirable exposition of the article \emph{Passus est} in Bishop Pearson’s Exposition of the Creed (ed. Dobson, p. 283 ff.).}

7. The anhypostasia, impersonality, or, to speak more accurately, the enhypostasia, of the human nature of Christ. This is a difficult point, but a necessary link in the orthodox doctrine of the one God-Man; for otherwise we must have two persons in Christ, and, after the incarnation, a fourth person, and that a human, in the divine Trinity. The impersonality of Christ’s human nature, however, is not to be taken as absolute, but relative, as the following considerations will show.

The centre of personal life in the God-Man resides unquestionably in the Logos, who was from eternity the second person in the Godhead, and could not lose his personality. He united himself, as has been already observed, not with a human person, but with human nature. The divine nature is therefore the root and basis of the personality of Christ. Christ himself, moreover, always speaks and acts in the full consciousness of his divine origin and character; as having come from the Father, having been sent by him, and, even during his earthly life, living in heaven and in unbroken communion with the Father.\footnote{The Logos is, according to the scholastic terminology of the later Greek theologians, especially John of Damascus, \textgreek{ἰδιοσύστατος}, or \textgreek{ἰδιουπόστατος}, i.e., per se subsistens, and \textgreek{ἰδιοπεριόριστος}, proprio termino circumscriptus.“Haec et similia vocabula,” says the learned Petavius (Theol. Dogm. tom. iv. p. 430), “demonstrant hypostasin non aliena ope fultam ac sustentatam existere, sed per semet ipsam, ac proprio termino definitam.” Schleiermacher’s Christology therefore, on this point, forms the direct opposite, of the Chalcedonian; it makes the \emph{man} Jesus the bearer of the personality, that is, transfers the proper centre of gravity in the personality to the human individuality of Christ, and views the divine nature as the supreme revelation of God in Him, as an impersonal principle, as a vital power. In this view the proper idea of the incarnation is lost. The same thing is true of the Christology of Hase, Keim, Beyschlag (and R. Rothe).} And the human nature of Christ had no independent personality of its own, besides the divine; it had no existence at all before the incarnation, but began with this act, and was so incorporated with the preexistent Logos-personality as to find in this alone its own full self-consciousness, and to be permeated and controlled by it in every stage of its development. But the human nature forms a necessary element in the divine personality, and in this sense we may say with the older Protestant theologians, that Christ is a persona σύνθετος, which was divine and human at once.\footnote{The correct Greek expression is, therefore, not \textgreek{ἀνυποστασία}but \textgreek{ἐνυποστασία}. The human nature of Christ was \textgreek{ἀνυπόστατος}, impersonalis, before the incarnation, but became \textgreek{ἀνυπόστατος}by the incarnation, that is, \textgreek{ἐν αὐτῇ τῇ τοῦ Θεοῦ λόγου ὑποστάσει ὑποστᾶσα}, and also \textgreek{ἑτερουπόστατος}, and \textgreek{συνυπόστατος}(compersonata), \emph{i.e}., quod per se et proprie modo non subsistit, sed inest in alio per se subsistente et substantia cum eo copulatur. Christ did not assume a human person, but a humana natura, in qua ipse Deus homo nasceretur. The doctrine of the \emph{an}hypostasia, impersonalitas, or rather \emph{en}hypostasia, of the human nature of Christ, is already observed, in incipient form, in Cyril of Alexandria, and was afterwards more fully developed by John of Damascus (De orthodoxa fide, lib. iii.), who, however, did not, for all this, conceive Christ as a mere generic being typifying mankind, but as a concrete human individual. Comp. Petavius, De incarnatione, l. v. c. 5-8 (tom. iv. p. 421 sqq.); Dorner, l. c. ii. p. 262 ff.; and J. P. Lange, Christliche Dogmatik, Part ii. p. 713.}

Thus interpreted, the church doctrine of the enhypostasia presents no very great metaphysical or psychological difficulty. It is true we cannot, according to our modern way of thinking, conceive a complete human nature without personality. We make personality itself consist in intelligence and free will, so that without it the nature sinks to a mere abstraction of powers, qualities, and functions.\footnote{Even in the scholastic era this difficulty was felt. Peter the Lombard says (Sentent. iii. d. 5 d.): Non accepit Verbum Dei \emph{personam} hominis, sed \emph{naturam}, quia non erat ex carne illa una composita persona, quam Verbum accepit, sed accipiendo univIt et uniendo accepit. \emph{E}: A quibusdam opponitur, quod persona assumpsit personam. Persona enim est substantia naturalis individuae naturae, hoc autem est anima. Ergo si animam assumpsit et personam. Quod ideo non sequitur, quia anima non est persona, quando alii rei unita est personaliter, sed quando per se est. Illa autem anima nunquam fuIt quin esset alii rei conjuncta.} But the human nature of Jesus never was, in fact, alone; it was from the beginning inseparably united with another nature, which is personal, and which assumed the human into a unity of life with itself. The Logos-personality is in this case the light of self-consciousness, and the impelling power of will, and pervades as well the human nature as the divine.\footnote{The Puritan theologian, John Owen (Works, vol. i. p. 223), says of the human nature of Christ quite correctly, and in agreement with the Chalcedonian Christology: “In itself it is \textgreek{ἀνυπόστατος}—that which hath not a subsistence of its own, which should give it individuation and distinction from the same nature in any other person. But \emph{it hath its subsistence in the person of the Son, which thereby is its own}. The divine nature, as in that person, is its \emph{suppositum}.”}

8. Criticism and development. This Chalcedonian Christology has latterly been subjected to a rigorous criticism, and has been charged now with dualism, now with docetism, according as its distinction of two natures or its doctrine of the impersonality of the human nature has most struck the eye.\footnote{Dr. Baur (Geschichte der Trinitätslehre, Bd. i. p. 823 f.) imputes to the Creed of Chalcedon “untenable inconsistency, equivocal indefiniteness, and discordant incompleteness,” but ascribes to it the merit of insisting upon the human in Christ as having equal claims with the divine, and of thus leaving the possibility of two equally legitimate points of view. Dr. Dorner, who regards the Chalcedonian statement as premature and inadequate (Geschichte der Christologie, Bd. ii. pp. 83, 130), raises against it the double objection of leaning to docetism on the one hand and to dualism on the other. He sums up his judgment of the labors of the ancient church down to John of Damascus in the sphere of Christology in the following words (ii. 273): “If we review the result of the Christological speculation of the ancient church, it is undeniable that the satisfying and final result cannot be found in it, great as its traditional influence even to this day is. It mutilates the human nature, inasmuch as, in an Apollinarian way, it joins to the trunk of a human nature the head of the divine hypostasis, and thus sacrifices the integrity of the humanity to the unity of the person. Yet after all—and this is only the converse of the same fault—in its whole doctrine of the natures and the will, it gives the divine and the human only an outward connection, and only, as it were, pushes the two natures into each other, without modification even of their properties. We discover, it is true, endeavors after something better, which indicate that the Christological image hovering before the mind, has not yet, with all the apparent completeness of the theory, found its adequate expression. But these endeavors are unfruitful.” Dr. W. Beyschlag, in his essay before the German \emph{Evangelische Kirchentag} at Altenburg, hold in 1864, concurs with these remarks, and says of the Chalcedonian dogma: “Instead of starting from the living intuition of the God-filled humanity of Christ, it proceeded from the defective and abstract conception of two separate natures, to be, as it were, added together in Christ; introduced thereby an irremediable dualism into his personal life; and at the same time, by transferring the personality wholly to the divine nature, depressed the humanity which \emph{in thesi} it recognized, to a mere unsubstantial accident of the Godhead, at bottom only apparent and docetistic.” But Beyschlag denies the real personal pre-existence of Christ and consequently a proper incarnation, and has by this denial caused no small scandal among the believing party in Germany. Dorner holds firmly to the pre-existence and incarnation, but makes the latter a gradual ethical unification of the Logos and the human nature, consummated in the baptism and the exaltation of Christ.}

But these imputations neutralize each other, like the imputations of tritheism and modalism which may be made against the orthodox doctrine of the Trinity when either the tripersonality or the consubstantiality is taken alone. This, indeed, is the peculiar excellence of the creed of Chalcedon, that it exhibits so sure a tact and so wise a circumspection in uniting the colossal antitheses in Christ, and seeks to do justice alike to the distinction of the natures and to the unity of the person.\footnote{F. R. Hasse (Kirchengeschichte, i. p. 177): “By the Creed of Chalcedon justice has been done to both the Alexandrian and the Antiochian Christology; the antagonism of the two is adjusted, and in the dogma of the one \textgreek{θεάνθρωπος} done away.”} In Christ all contradictions are reconciled.

Within these limits there remains indeed ample scope for further Christological speculations on the possibility, reality, and mode of the incarnation; on its relation to the revelation of God and the development of man; on its relation to the immutability of God and the trinity of essence and the trinity of revelation:—questions which, in recent times especially, have been earnestly and profoundly discussed by the Protestant theologians of Germany.\footnote{Witness the Christological investigations of Schleiermacher, R. Rothe, Göschel, Dorner, Liebner, Lange, Thomasius, Martensen, Gess, Ebrard, Schöberlein, Plitt, Beyschlag, and others. A thorough criticism of the latest theories is given by Dorner, in his large work on Christology, Bd. ii. p. 1260 ff. (Eng. transl. Div. 2d, vol. iii. p. l00 ff.), and in several dissertations upon the immutability of God, found in his Jabrbücher für Deutsche Theologie, 1856 and 1858; also by Philippi, Kirchliche Glaubenslehre, iv. i. pp. 344-382; Plitt, Evangelische Glaubenslehre (1863), i. p. 360 ff.; and Woldemar Schmidt, Das Dogma vom Gottmenschen, mit Beziehung auf die neusten Lösungsversuche der Gegensätze, Leipzig, 1865. The English theology has contented itself with the traditional acceptance and vindication of the old Catholic doctrine of Christ’s person, without instituting any special investigations of its own, while the doctrine of the Trinity has been thoroughly reproduced and vindicated by Cudworth, Bull, and Waterland, without, however, being developed further. Dr. Shedd also considers the Chalcedonian symbol as the \emph{ne} \emph{plus ultra} of Christological knowledge, “beyond which it is probable the human mind is unable to go, in the endeavor to unfold the mystery of Christ’s complex person, which in some of its aspects is even more baffling than the mystery of the Trinity” (History of Christian Doctrine, i. p. 408). This is probably also the reason why this work, in surprising contrast with every other History of Doctrine, makes no mention whatever of the Monophysite, Monothelite, Adoptian, Scholastic, Lutheran, Socinian, Rationalistic, and later Evangelical controversies and theories respecting this central dogma of Christianity.}

The great want, in the present state of the Christological controversy, is, on the one hand, a closer discussion of the Pauline idea of the kenosis, the self-limitation, self-renunciation of the Logos, and on the other hand, a truly human portrait of Jesus in his earthly development from childhood to the fall maturity of manhood, without prejudice to his deity, but rather showing forth his absolute uniqueness and sinless perfection as a proof of his Godhead. Both these tasks can and should be so performed, that the enormous labor of deep and earnest thought in the ancient church be not condemned as a sheer waste of strength, but in substance confirmed, expanded, and perfected.

And even among believing Protestant scholars, who agree in the main views of the theanthropic glory of the person of Christ, opinions still diverge. Some restrict the kenosis to the laying aside of the divine form of existence, or divine dignity and glory;\footnote{Of the \textgreek{δόξα θεοῦ}, John xvii. 5; the \textgreek{μορφὴ Θεοῦ}, Phil. ii. 6 ff.} others strain it in different degrees, even to a partial or entire emptying of the divine essence out of himself, so that the inner trinitarian process between Father and Son, and the government of the world through the Son, were partially or wholly suspended during his earthly life.\footnote{Among these modem Kenotics, W. F. Gess goes the farthest in his Lehre von der Person Christi (Basel, 1856). Dorner opposes the theory of the Kenotics and calls them Theopaschites and Patripassians (ii. 126 ff.). There is, however, an essential distinction, inasmuch as the ancient Monophysite Theopaschitism reduces the human nature of Christ to a mere accident of his Godhead, while Thomasius, Gess, and the other German Kenotics or Kenosists acknowledge the full humanity of Christ, and lay great stress on it.} Some, again, view the incarnation as an instantaneous act, consummated in the miraculous conception and nativity; others as a gradual process, an ethical unification of the eternal Logos and the man Jesus in continuous development, so that the complete God-Man would be not so much the beginning as the consummation of the earthly life of Jesus.

But all these more recent inquiries, earnest, profound, and valuable as they are, have not as yet led to any important or generally accepted results, and cannot supersede the Chalcedonian Christology. The theology of the church will ever return anew to deeper and still deeper contemplation and adoration of the theanthropic person of Jesus Christ, which is, and ever will be, the sun of history, the miracle of miracles, the central mystery of godliness, and the inexhaustible fountain of salvation and life for the lost race of man.

\xsection{The Monophysite Controversies}

§ 143. The Monophysite Controversies.

I. The Acta in Mansi, tom. vii.-ix. The writings already cited of Liberatus and Leontinus Byzant. Evagrius: H. E. ii. v. Nicephorus: H. E. xvi. 25. Procopius († about 552): Ἀνέκδοτα, Hist. arcana (ed. Orelli, Lips. 1827). Facundus (bishop of Hermiane in Africa, but residing mostly in Constantinople): Pro defensione trium capitulorum, in 12 books (written a.d. 547, ed. Sirmond, Paris, 1629, and in Galland. xi. 665). Fulgentius Ferrandus (deacon in Carthage, † 551): Pro tribus capitulis (in Gall. tom. xi.). Anastasius Sinaita (bishop of Antioch, 564): Ὁδηγόςadv. Acephalos. Angelo Mai: Script vet. Bova collectio, tom. vii. A late, though unimportant, contribution to the history of Monophysitism (from 581 to 583) is the Church History of the Monophysite bishop John of Ephesus (of the sixth century): The Third Part of the Eccles. History of John, bishop of Ephesus, Oxford, 1853 (edited by W. Cureton from the Syrian literature of the Nitrian convent).

II. Petavius: De Incarnatione, lib. i. c. 16–18 (tom. iv. p. 74 sqq.). Walch: Bd, vi.-viii. Schröckh: Th. xviii. pp. 493–636. Neander: Kirchengeschichte, !v. 993–1038. Gieseler: i. ii. pp. 347–376 (4th ed.), and his Commentatio qua Monophysitarum veterum variae de Christi persona opiniones ... illustrantur (1835 and 1838). Baur: Geschichte der Trinitätslehre, Bd. ii. pp. 37–96. Dorner: Geschichte der Christologie, ii. pp. 150–193. Hefele (R.C.): Conciliengeschichte, ii. 545 ff. F. Rud. Hasse: Kirchengeschichte (1864), Bd. i. p. 177 ff. A. Ebrard: Handbuch der Kirchen- und Dogmengeschichte (1865), Bd. i. pp. 263–279.

The council of Chalcedon did not accomplish the intended pacification of the church, and in Palestine and Egypt it met with passionate opposition. Like the council of Nicaea, it must pass a fiery trial of conflict before it could be universally acknowledged in the church. “The metaphysical difficulty,” says Niedner, “and the religious importance of the problem, were obstacles to the acceptance of the ecumenical authority of the council.” Its opponents, it is true, rejected the Eutychian theory of an absorption of the human nature into the divine, but nevertheless held firmly to the doctrine of one nature in Christ; and on this account, from the time of the Chalcedonian council they were called Monophysites,\footnote{\textgreek{Μονοφυσίται}, from \textgreek{μόνη} or \textgreek{μία, φύσις}. They conceded the \textgreek{ἐκ δύο φύςεσιν}(as even Eutyches and Dioscurus had done), but denied the \textgreek{ἐν δύο φύσεσιν}, after the \textgreek{ἕνωσις .}} while they in return stigmatized the adherents of the council as Dyophysites and Nestorians. They conceded, indeed, a composite nature (μία φύσις σύνθετοςor μία φύσις διττή), but not two natures. They assumed a diversity of qualities without corresponding substances, and made the humanity in Christ a mere accident of the immutable divine substance.

Their main argument against Chalcedon was, that the doctrine of two natures necessarily led to that of two persons, or subjects, and thereby severed the one Christ into two Sons of God. They were entirely at one with the Nestorians in their use of the terms “nature” and “person,” and in rejecting the orthodox distinction between the two. They could not conceive of human nature without personality. From this the Nestorians reasoned that, because in Christ there are two natures, there must be also two independent hypostases; the Monophysites, that, because there is but one person in Christ, there can be only one nature. They regarded the nature as something common to all individuals of a species (κοίνον), yet as never existing simply as such, but only in individuals. According to them, therefore, over, φύσιςor οὐσία is in fact always an individual existence.\footnote{\textgreek{Ἰδικόν.}}

The liturgical shibboleth of the Monophysites was: God has been crucified. This they introduced into their public worship as an addition to the Trisagion: “Holy, God, holy Mighty, holy Immortal, who hast been crucified for us, have mercy upon us.”\footnote{\textgreek{Ἅγιος ὁ Θεὸς , ἄγιος ἴσχυρος , ἅγιος ἀθάνατος , ὁ σταυρωθεὶς δι  ἡμᾶς , ἐλέησον ἡμᾶς}. An extension of the seraphic ascription, Isa. vi. 3.} From this they were also called Theopaschites.\footnote{665 \textgreek{Θεοπασχῖται.}} This formula is in itself orthodox, and forms the requisite counterpart to θεοτόκος, provided we understand by God the Logos, and in thought supply: “according to the flesh” or “according to the human nature.” In this qualified sense it was afterwards in fact not only sanctioned by Justinian in a dogmatical decree, but also by the fifth ecumenical council, though not as an addition to the Trisagion. For the theanthropic person of Christ is the subject, as of the nativity, so also of the passion; his human nature is the seat and the organ (sensorium) of the passion. But as an addition to the Trisagion, which refers to the Godhead generally, and therefore to the Father, and the Holy Ghost, as well as the Son, the formula is at all events incongruous and equivocal. Theopaschitism is akin to the earlier Patripassianism, in subjecting the impassible divine essence, common to the Father and the Son, to the passion of the God-Man on the cross; yet not, like that, by confounding the Son with the Father, but by confounding person with nature in the Son.

Thus from the council of Chalcedon started those violent and complicated Monophysite controversies which convulsed the Oriental church, from patriarchs and emperors down to monks and peasants, for more than a hundred years, and which have left their mark even to our day. They brought theology little appreciable gain, and piety much harm; and they present a gloomy picture of the corruption of the church. The intense concern for practical religion, which animated Athanasius and the Nicene fathers, abated or went astray; theological speculation sank towards barren metaphysical refinements; and party watchwords and empty formulas were valued more than real truth. We content ourselves with but a summary of this wearisome, though not unimportant chapter of the history of doctrines, which has recently received new light from the researches of Gieseler, Baur, and Dorner.\footnote{The \emph{external} history of Monophysitism is related with wearisome minuteness by Walch in three large volumes (vi.-viii.) of his Entwurf einer vollständigen Historie der Ketzereien, etc., his auf die Zeiten der Reformation.}

The external history of the controversy is a history of outrages and intrigues, depositions and banishments, commotions, divisions, and attempted reunions. Immediately after the council of Chalcedon bloody fights of the monks and the rabble broke out, and Monophysite factions went off in schismatic churches. In Palestine Theodosius (451–453) thus set up in opposition to the patriarch Juvenal of Jerusalem; in Alexandria, Timotheus Aelurus\footnote{\textgreek{Αἴλουρος}, Cat.} and Peter Mongus\footnote{\textgreek{Μόγγος}, the Stammerer; literally, the Hoarse.} (454–460), in opposition to the newly-elected patriarch Protarius, who was murdered in a riot in Antioch; Peter the Fuller\footnote{Fullo, \textgreek{γναφεύς}. He introduced the formula: \textgreek{Θεὸς ἐσταυρώθη δι  ἡμᾶς}into the liturgy. He was in 485 again raised to the patriarchate.} (463–470). After thirty years’ confusion the Monophysites gained a temporary victory under the protection of the rude pretender to the empire, Basiliscus (475–477), who in an encyclical letter,\footnote{\textgreek{Ἐγκυκλιον}. This, however, excited so much opposition, that the usurper in 477 revoked it in an \textgreek{ἀντεγκύκλιον}.} enjoined on all bishops to condemn the council of Chalcedon (476). After his fall, Zeno (474–475 and 477–491), by advice of the patriarch Acacius of Constantinople, issued the famous formula of concord, the Henoticon, which proposed, by avoiding disputed expressions, and condemning both Eutychianism and Nestorianism alike, to reconcile the Monophysite and dyophysite views, and tacitly set aside the Chalcedonian formula (482). But this was soon followed by two more schisms, one among the Monophysites themselves, and one between the East and the West. Felix II., bishop of Rome, immediately rejected the Henoticon, and renounced communion with the East (484–519). The strict Monophysites were as ill content with the Henoticon, as the adherents of the council of Chalcedon; and while the former revolted from their patriarchs, and became Acephali,\footnote{\textgreek{Ἀκέφαλοι}, without head.} the latter attached themselves to Rome. It was not till the reign of the emperor Justin I. (518–527), that the authority of the council of Chalcedon was established under stress of a popular tumult, and peace with Rome was restored. The Monophysite bishops were now deposed, and fled for the most part to Alexandria, where their party was too powerful to be attacked.

The internal divisions of the Monophysites turned especially on the degree of essential difference between the humanity of Christ and ordinary human nature, and the degree, therefore, of their deviation from the orthodox doctrine of the full consubstantiality of the humanity of Christ with ours.\footnote{Petavius, \emph{l.c.} lib. i. c. 17, enumerates twelve factions of the Monophysites.} The most important of these parties were the Severians (from Severus, the patriarch of Antioch) or Phthartolaters (adorers of the corruptible),\footnote{\textgreek{Φθαρτολάτραι}(from \textgreek{φθαρτός}corruptible, and \textgreek{λάτρης}, servant, worshipper), corrupticolae.} who taught that the body of Christ before the resurrection was mortal and corruptible; and the Julianists (from bishop Julian of Halicarnassus, and his contemporary Xenajas of Hierapolis) or Aphthartodocetae,\footnote{\textgreek{Ἀφθαρτοδοκῆται}, also called Phantasiastae, because they appeared to acknowledge only a \emph{seeming} body of Christ. Gieseler, however, in the second part of the above-mentioned dissertation, has shown that the Julianist view was not strictly docetistic, but kindred with the view of Clement of Alexandria, Origen, Hilary, Gregory of Nyssa, and Apollinaris.} who affirmed the body of Christ to have been originally incorruptible, and who bordered on docetism. The former conceded to the Catholics, that Christ as to the flesh was consubstantial with us (κατὰ σάρκα ὁμοούσιος ἡμῖν). The latter argued from the commingling (σύγχυσις) of the two natures, that the corporeality of Christ became from the very beginning partaker of the incorruptibleness of the Logos, and was subject to corruptibleness merely κατ  οἰκονομίαν. They appealed in particular to Jesus’ walking on the sea. Both parties were agreed as to the incorruptibleness of the body of Christ after the resurrection. The word fqorav, it may be remarked, was sometimes used in the sense of frailty, sometimes in that of corruptibleness.

The solution of this not wholly idle question would seem to be, that the body of Christ before the resurrection was similar to that of Adam before the fall; that is, it contained the germ of immortality and incorruptibleness; but before its glorification it was subject to the influence of the elements, was destructible, and was actually put to death by external violence, but, through the indwelling power of the sinless spirit, was preserved from corruption, and raised again to imperishable life. A relative immortality thus became absolute.\footnote{Comp. the Augustinian distinction of immortalitas minor and immortalitas major.} So far we may without self-contradiction affirm both the identity of the body of Christ before and after his resurrection, and its glorification after resurrection.\footnote{As was done by Augustineand Leo the Great. The latter affirms, Sermo 69, De resurrectione Domini, c. 4: “Resurrectio Domini non finis carnis, sed commutatio fuit, nec virtutis augmento consumpta substantia est. Qualitas transiit, non natura defecit; et factum est corpus impassibile, immortale, incorruptibile ... nihil remansit in carne Christi infirmum, ut et ipsa sit per essentiam et non sit ipsa per gloriam.” Comp. moreover, respecting the Aphthartodocetic controversy of the Monophysites, the remarks of Dorner, ii. 159 ff. and of Ebrard, Kirchen- und Dogmengeschichte, i. 268 f.}

The Severians were subdivided again, in respect to the question of Christ’s omniscience, into Theodosians, and Themistians, or Agnoetae.\footnote{After their leader Themistius, deacon of Alexandria; also called by their opponents, Agnoetae, \textgreek{Ἀγνοηταί}, because they taught that Christ in his condition of humiliation was not omniscient, but shared our ignorance of many things (Comp. Luke ii. 52; Mark xiii. 32). This view leads necessarily to dyophysitism, and accordingly was rejected by the strict Monophysites.} The Julianists were subdivided into Ktistolatae,\footnote{\textgreek{Κτιστολάτραι}, or, from their founder, Gajanitae. These viewed the body of Christ as created, \textgreek{κτιστόν}.} and Aktistetae\footnote{680 \textgreek{Ἀκτιστηταί}. These said that the body of Christ in itself was created, but that by its union with the Logos it became increate, and therefore also incorruptible.} according as they asserted or denied that the body of Christ was a created body. The most consistent Monophysite was the rhetorician Stephanus Niobes (about 550), who declared every attempt to distinguish between the divine and the human in Christ inadmissible, since they had become absolutely one in him.\footnote{His adherents were condemned by the other Monophysites as Niobitae.} An abbot of Edessa, Bar Sudaili, extended this principle even to the creation, which be maintained would at last be wholly absorbed in God. John Philoponus (about 530) increased the confusion; starting with Monophysite principles, taking φύσιςin a concrete instead of an abstract sense, and identifying it with ὑπόστασις, he distinguished in God three individuals, and so became involved in tritheism. This view he sought to justify by the Aristotelian categories of genus, species, and individuum.\footnote{His followers were called Philoponiaci, Tritheistae. Philoponus, it may be remarked, was not the first promulgator of this error; but (as appears from Assem. Bibl. orient. tom. ii. p. 327; comp. Hefele, ii. 655) the Monophysite John Askusnages, who ascribed to Christ only one nature, but to each person in the Godhead a separate nature, and on this account was banished by the emperor and excommunicated by the patriarch of Constantinople. Among the more famous Tritheists we have also Stephen Gobarus, about 600.}

\xsection{The Three, Chapters, and the Fifth Ecumenical Council, A.D. 553}

§ 144. The Three, Chapters, and the Fifth Ecumenical Council, A.D. 553.

Comp., besides the literature already cited, H. Noris (R.C.): Historia Pelagiana et dissertatio de Synodo Quinta oecumen. in qua Origenis et Th. Mopsuesteni Pelagiani erroris auctorum justa damnatio, et Aquilejense schisma describitur, etc. Padua, 1673, fol., and Verona, 1729. John Garnier (R.C.): Dissert. de V. Synodo. Paris, 1675 (against Card. Noris). Hefele (R.C.): vol. ii. 775–899.—The Greek Acta of the 5th council, with the exception of the 14 anathemas and some fragments, have been lost; but there is extant an apparently contemporary Latin translation (in Mansi, tom. ix. 163 sqq.), respecting whose genuineness and completeness there has been much controversy (comp. Hefele, ii. p. 831 ff.).

The further fortunes of Monophysitism are connected with the emperor Justinian I. (527–565). This learned and unweariedly active ruler, ecclesiastically devout, but vain and ostentatious, aspired, during his long and in some respects brilliant reign of nearly thirty years, to the united renown of a lawgiver and theologian, a conqueror and a champion of the true faith. He used to spend whole nights in prayer and fasting, and in theological studies and discussions; he placed his throne under the special protection of the Blessed Virgin and the archangel Michael; in his famous Code, and especially in the Novelles, he confirmed and enlarged the privileges of the clergy; he adorned the capital and the provinces with costly temples and institutions of charity; and he regarded it as his especial mission to reconcile heretics, to unite all parties of the church, and to establish the genuine orthodoxy for all time to come. In all these undertakings he fancied himself the chief actor, though very commonly he was but the instrument of the empress, or of the court theologians and eunuchs; and his efforts to compel a general uniformity only increased the divisions in church and state.

Justinian was a great admirer of the decrees of Chalcedon, and ratified the four ecumenical councils in his Code of Roman law. But his famous wife Theodora, a beautiful, crafty, and unscrupulous woman, whom he—if we are to believe the report of Procopius\footnote{Historia Arcana. c. 9.}—raised from low rank, and even from a dissolute life, to the partnership of his throne, and who, as empress, displayed the greatest zeal for the church and for ascetic piety, was secretly devoted to the Monophysite view, and frustrated all his plans. She brought him to favor the liturgical formula of the Monophysites: “God was crucified for us, so that he sanctioned it in an ecclesiastical decree (533).\footnote{This addition remained in use among the Catholics in Syria till it was thrown out by the \emph{Concilium Quinisextum} (can. 81). Thenceforth it was confined to the Monophysites and Monothelites. The opinion gained ground among the Catholics, that the formula taught a quaternity, instead of a trinity. Gieseler, i. P. ii. p. 366 ff.}

Through her influence the Monophysite Anthimus was made patriarch of Constantinople (535), and the characterless Vigilius bishop of Rome (538), under the secret stipulation that he should favor the Monophysite doctrine. The former, however, was soon deposed as a Monophysite (536), and the latter did not keep his promise.\footnote{Hefele (ii. p. 552) thinks that Vigilius was never a Monophysite at heart, and that he only gave the promise in the interest of “his craving ambition.” The motive, however, of course cannot alter the fact, nor weaken the argument, furnished by his repeated recantations, against the claims of the papal see to infallibility.} Meanwhile the Origenistic controversies were renewed. The emperor was persuaded, on the one hand, to condemn the Origenistic errors in a letter to Mennas of Constantinople; on the other hand, to condemn by an edict the Antiochian teachers most odious to the Monophysites: Theodore of Mopsuestia (the teacher of Nestorius), Theodoret of Cyros, and Ibas of Edessa (friends of Nestorius); though the last two had been expressly declared orthodox by the council of Chalcedon. Theodore he condemned absolutely, but Theodoret only as respected his writings against Cyril and the third ecumenical council at Ephesus, and Ibas as respected his letter to the Persian bishop Maris, in which he complains of the outrages of Cyril’s party in Edessa, and denies the communicatio idiomatum. These are the so-called Three Chapters, or formulas of condemnation, or rather the persons and writings designated and condemned therein.\footnote{\textgreek{Τρία κεφάλαια}, tria capitula. “Chapters” are properly articles, or brief propositions, under which certain errors are summed up in the form of anathemas. The twelve anathemas of Cyril against Nestorius were also called icEI)dAaia. By the Three Chapters, however, are to be understood in this case: 1. The \emph{person} and \emph{writings} of Theodore of Mopsuestia; 2. the anti-Cyrillian \emph{writings} of Theodoret; 3. the \emph{letter} of Ibas to Maris. Hence the appellation \emph{impia capitula}, aO’CGi Ke4paAaia. This deviation from ordinary usage has occasioned much confusion.}

Thus was kindled the violent controversy of the Three Chapters, of which it has been said that it has filled more volumes than it was worth lines. The East yielded easily to craft and force; the West resisted.\footnote{Especially the African Fulgentius Ferrandus, Liberatus, and Facundus of Hermiane, who wrote in defence of the Three Chapters; also the Roman deacon Rusticus.} Pontianus of Carthage declared that neither the emperor nor any other man had a right to sit in judgment upon the dead. Vigilius of Rome, however, favored either party according to circumstances, and was excommunicated for awhile by the dyophysite Africans, under the lead of Facundus of Hermiane. He subscribed the condemnation of the Three Chapters in Constantinople, a.d. 548, but refused to subscribe the second edict of the, emperor against the Three Chapters (551), and afterwards defended them.

To put an end to this controversy, Justinian, without the concurrence of the pope, convoked at Constantinople, a.d. 553, the Fifth Ecumenical Council, which consisted of a hundred and sixty-four bishops, and held eight sessions, from the 5th of May to the 2d of June, under the presidency of the patriarch Eutychius of Constantinople. It anathematized the Three Chapters; that is, the person of Theodore of Mopsuestia, the anti-Cyrillian writings of Theodoret, and the letter of Ibas,\footnote{These anathemas are found in the concluding sentence of the council (Mansi, tom. ix. 376): “Praedicta igitur tria capitula anathematizamus, id est Theodorum impium Mopsuestenum, cum nefandis ejus conscriptis, et quae impie Theodoretus conscripsit, et impiam epistolam, quae dicitur Ibae.”} and sanctioned the formula “God was crucified,” or “One of the Trinity has suffered,” yet not as an addition to the Trisagion.\footnote{Collect viii. can. 10: ῎Ει τιsὀυκ ὁμολογεῖτὸν ἐσταυωμενον σάρκι κύριον ημῶν Ἰησοῦν Χριδτὸν εἶναι Θεὸν ἀληθινῶν καὶκύριον τῆsδόξηs, καὶ῞ενα τῆsαγίαsτριάδοs, ὁτοιτῦτοsἀνάθεμα ῎εστων. “Whoever does not acknowledge that our Lord Jesus Christ, who was crucified in the flesh, is true God and Lord of glory, and one of the Holy Trinity, let him be anathema.”} The dogmatic decrees of Justinian were thus sanctioned by the church. But no further mention appears to have been made of Origenism; and in truth none was necessary, since a local synod of 544 had already condemned it. Perhaps also Theodore Askidas, a friend of the Origenists, and one of the leaders of the council, prevented the ecumenical condemnation of Origen. But this is a disputed point, and is connected with the difficult question of the genuineness and completeness of the Acts of the council.\footnote{In the 11th anathema, it is true, the name of Origen is condemned along with other heretics (Arius, Eunomius, Macedonius, Apollinaris, Nestorius, Eutyches), but the connection is incongruous, and the name is regarded by Halloix, Garnier, Jacob Basnage, Walch, and others, as all interpolation. Noris and Hefele (ii. p. 874) maintain its genuineness. At all events the fifteen anathemas against Origen do not belong to it, but to an earlier Constantinopolitan synod, held in 544. Comp. Hefele, ii. p. 768 ff.}

Vigilius at first protested against the Council, which, in spite of repeated invitations, he had not attended, and by which he was suspended; but he afterwards signified his adherence, and was permitted, after seven years’ absence, to return to Rome, but died on the journey, at Syracuse, in 555. His fourfold change of opinion does poor service to the claim of papal infallibility. His successor, Pelagius I., immediately acknowledged the council. But upon this the churches in Northern Italy, Africa, and Illyria separated themselves from the Roman see, and remained in schism till Pope Gregory I. induced most of the Italian bishops to acknowledge the council.

The result of this controversy, therefore, was the condemnation of the Antiochian theology, and the partial victory of the Alexandrian monophysite doctrine, so far as it could be reconciled with the definitions of Chalcedon. But the Chalcedonian dyophysitism afterwards reacted, in the form of dyothelitism, and at the sixth ecumenical council, at Constantinople, a.d. 680 (called also Concilium Trullanum I.), under the influence of a letter of pope Agatho, which reminds us of the Epistola Dogmatica of Leo, it gained the victory over the Monothelite view, which so far involves the Monophysite, as the ethical conception of one will depends upon the physical conception of one nature.

But notwithstanding the concessions of the fifth ecumenical council, the Monophysites remained separated from the orthodox church, refusing to acknowledge in any manner the dyophysite council of Chalcedon. Another effort of Justinian to gain them, by sanctioning the Aphthartodocetic doctrine of the incorruptibleness of Christ’s body (564), threatened to involve the church in fresh troubles; but his death soon afterwards, in 565, put an end to these fruitless and despotic plans of union. His successor Justin II. in 565 issued an edict of toleration, which exhorted all Christians to glorify the Lord, without contending about persons and syllables. Since that time the history of the Monophysites has been distinct from that of the catholic church.

\xsection{The Monophysite Sects: Jacobites, Copts, Abyssinians, Armenians, Maronites}

§ 145. The Monophysite Sects: Jacobites, Copts, Abyssinians, Armenians, Maronites.

Euseb. Renaudot (R.C., † 1720): Historia patriarcharum Alexandrinorum Jacobitarum a D. Marco usque ad finem saec. xiii. Par. 1713. Also by the same: Liturgiarum orientalium collectio. Par. 1716, 2 vols. 4to. Jos. Sim. Assemani (R.C., † 1768): Bibliotheca orientalis. Rom. 1719 sqq., 4 vols. folio (vol. ii. treats De scriptoribus Syria Monophysitis). Michael le Quien (R.C., † 1733): Oriens Christianus. Par. 1740, 3 vols. folio (vols. 2 and 3). Veyssière De La Croze: Histoire du Christianisme d’Ethiope et d’Armenie. La Haye, 1739. Gibbon: Chapter xlvii. towards the end. Makrîzi (Mohammedan, an historian and jurist at Cairo, died 1441): Historia Coptorum Christianorum (Arabic and Latin), ed. H. T. Wetzer, Sulzbach, 1828; a better edition by F. Wüstenfeld, with translation and annotations, Göttingen, 1845. J. E. T. Wiltsch Kirchliche Statistik. Berlin, 1846, Bd. i. p. 225 ff. John Mason Neale (Anglican): The Patriarchate of Alexandria. London, 1847, 2 vols. Also: A History of the Holy Eastern Church. Lond. 1850, 2 vols. (vol. ii. contains among other things the Armenian and Copto-Jacobite Liturgy). E. Dulaurier: Histoire, dogmes, traditions, et liturgie de l’Eglise Armeniane. Par. 1859. Arthur Penrhyn Stanley: Lectures on the History of the Eastern Church. New York, 1862, Lect. i. p. 92 ff. Respecting the present condition of the Jacobites, Copts, Armenians, and Maronites, consult also works of Eastern travel, and the numerous accounts in missionary magazines and other religious periodicals.

The Monophysites, like their antagonists, the Nestorians, have maintained themselves in the East as separate sects under their own bishops and patriarchs, even to the present day; thus proving the tenacity of those Christological errors, which acknowledge the full Godhead and manhood of Christ, while those errors of the ancient church, which deny the Godhead, or the manhood (Ebionism, Gnosticism, Manichaeism, Arianism, etc.), as sects, have long since vanished. These Christological schismatics stand, as if enchanted, upon the same position which they assumed in the fifth century. The Nestorians reject the third ecumenical council, the Monophysites the fourth; the former hold the distinction of two natures in Christ even to abstract separation, the latter the fusion of the two natures in one with a stubbornness which has defied centuries, and forbids their return to the bosom of the orthodox Greek church. They are properly the ancient national churches of Egypt, Syria, and Armenia, in distinction from the orthodox Greek church, and the united or Roman church of the East.

The Monophysites are scattered upon the mountains and in the valleys and deserts of Syria, Armenia, Assyria, Egypt, and Abyssinia, and, like the orthodox Greeks of those countries, live mostly under Mohammedan, partly under Russian, rule. They supported the Arabs and Turks in weakening and at last conquering the Byzantine empire, and thus furthered the ultimate victory of Islam. In return, they were variously favored by the conquerors, and upheld in their separation from the Greek church. They have long since fallen into stagnation, ignorance, and superstition, and are to Christendom as a praying corpse to a living man. They are isolated fragments of the ancient church history, and curious petrifactions from the Christological battle-fields of the fifth and sixth centuries, coming to view amidst Mohammedan scenes. But Providence has preserved them, like the Jews, and doubtless not without design, through storms of war and persecution, unchanged until the present time. Their very hatred of the orthodox Greek church makes them more accessible both to Protestant and Roman missions, and to the influences of Western Christianity and Western civilization.

On the other hand, they are a door for Protestantism to the Arabs and the Turks; to the former through the Jacobites, to the latter through the Armenians. There is the more reason to hope for their conversion, because the Mohammedans despise the old Oriental churches, and must be won, if at all, by a purer type of Christianity. In this respect the American missions among the Armenians in the Turkish empire, are, like those among the Nestorians in Persia, of great prospective importance, as outposts of a religion which is destined sooner or later to regenerate the East.

With the exception of the Chalcedonian Christology, which they reject as Nestorian heresy, most of the doctrines, institutions, and rites of the Monophysite sects are common to them with the orthodox Greek church. They reject, or at least do not recognize, the filioque; they hold to the mass, or the Eucharistic sacrifice, with a kind of transubstantiation; leavened bread in the Lord’s Supper; baptismal regeneration by trine immersion; seven sacraments (yet not explicitly, since they either have no definite term for sacrament, or no settled conception of it); the patriarchal polity; monasticism; pilgrimages, and fasting; the requisition of a single marriage for priests and deacons (bishops are not allowed to marry);\footnote{Laymen are allowed to marry twice, but a third marriage is regarded as fornication.} the prohibition of the eating of blood or of things strangled.\footnote{Comp. Acts xv. 20. The Latin church saw in this ordinance of the apostolic council merely a temporary measure during the existence of Jewish Christianity.} On the other hand, they know nothing of purgatory and indulgences, and have a simpler worship than the Greeks and Romans. According to their doctrine, all men after death go into Hades, a place alike without sorrow or joy; after the general judgment they enter into heaven or are cast into hell; and meanwhile the intercessions and pious works of the living have an influence on the final destiny of the departed. Like the orthodox Greeks, they honor pictures and relics of the saints, but not in the same degree. Scripture and tradition are with them coordinate sources of revelation and rules of faith. The reading of the Bible is not forbidden, but is limited by the ignorance of the people themselves. They use in worship the ancient vernacular tongues, which, however, are now dead languages to them.

There are four branches of the Monophysites: the Syrian Jacobites; the Copts, including the Abyssinians; the Armemians; and the less ancient Maronites.

I. The Jacobites in Syria, Mesopotamia, and Babylonia. Their name comes down from their ecumenical\footnote{\emph{Ecumenical, i.e}., not restricted to any particular province.} metropolitan Jacob, surnamed Baradai, or Zanzalus.\footnote{From his beggarly clothing. Barâdai signifies in Arabic and Syriac horse blanket, of coarse cloth, and \textgreek{τζάνζαλον}is \emph{vile aliquid et tritum} (see Rödiger in Herzog’s Encycl. vi. 401).} This remarkable man, in the middle of the sixth century, devoted himself for seven and thirty years (511–578), with unwearied zeal to the interests of the persecuted Monophysites. “Light-footed as Asahel,”\footnote{2 Sam. ii. 18.} and in the garb of a beggar, he journeyed hither and thither amid the greatest dangers and privations; revived the patriarchate of Antioch; ordained bishops, priests, and deacons; organized churches; healed divisions; and thus saved the Monophysite body from impending extinction.

The patriarch bears the title of patriarch of Antioch, because the succession is traced back to Severus of Antioch; but he commonly resides in Diarbekir, or other towns or monasteries. Since the fourteenth century, the patriarch has always borne the name Ignatius, after the famous martyr and bishop of Antioch. The Jacobite monks are noted for gross superstition and rigorous asceticism. A part of the Jacobites have united with the church of Rome. Lately some Protestant missionaries from America have also found entrance among them.

II. The Copts,\footnote{From \textgreek{αἴγυπτος}, Guptos, and not, as some suppose, from the town Koptos, nor from an abbreviation of Jacobite. They are the most ancient, but Christian Egyptians, in distinction from the Pharaonic (Chem), those of the Old Testament (Mizrim), the Macedonian or Greek (\textgreek{αἴγ}.) and the modem Arab Egyptians (Misr).} in Egypt, are in nationality the genuine descendants of the ancient Egyptians, though with an admixture of Greek and Arab blood. Soon after the council of Chalcedon, they chose Timotheus Aelurus in opposition to the patriarch Proterius. After varying fortunes, they have, since 536, had their own patriarch of Alexandria, who, like most of the Egyptian dignitaries, commonly resides at Cairo. He accounts himself the true successor of the evangelist Mark, St. Athanasius, and Cyril. He is always chosen from among the monks, and, in rigid adherence to the traditionary nolo episcopari, he is elected against his will; he is obliged to lead a strict ascetic life, and at night is waked every quarter of an hour for a short prayer. He alone has the power to ordain, and he performs this function not by imposition of hands, but by breathing on and anointing the candidate. His jurisdiction extends over the churches of Egypt, Nubia, and Abyssinia, or Ethiopia. He chooses and anoints the Abuna (i.e., Our Father), or patriarch for Abyssinia. Under him are twelve bishops, some with real jurisdiction, some titular; and under these again other clergy, down to readers and exorcists. There are still extant two incomplete Coptic versions of the Scriptures, the Upper Egyptian or Thebaic, called also, after the Arabic name of the province, the Sahidic, i.e., Highland version; and the Lower Egyptian or Memphitic.\footnote{Of this latter H. Tattam and P. Bötticher (1852) have lately published considerable fragments.}

The Copts were much more numerous than the Catholics, whom they scoffingly nicknamed Melchites,\footnote{From the Hebrew \emph{melech}, king.} or Caesar-Christians. They lived with them on terms of deadly enmity, and facilitated the conquest of Egypt by the Saracens (641). But they were afterwards cruelly persecuted by these very Saracens,\footnote{So that even their Arabic historian Makr\&lt;zi was moved to compassion for them.} and dwindled from some two millions of souls to a hundred and fifty or two hundred thousand, of whom about ten thousand, or according to others from thirty to sixty thousand, live in Cairo, and the rest mostly in Upper Egypt. They now, in common with all other religious sects, enjoy toleration. They and the Abyssinians are distinguished from the other Monophysites by the Jewish and Mohammedan practice of circumcision, which is performed by lay persons (on both sexes), and in Egypt is grounded upon sanitary considerations. They still observe the Jewish law of meats. They are sunk in poverty, ignorance, and semi-barbarism. Even the clergy, who indeed are taken from the lowest class of the people, are a beggarly set, and understand nothing but how to read mass, and perform the various ceremonies. They do not even know the Coptic or old Egyptian, their own ancient ecclesiastical language. They live by farming, and their official fees. The literary treasures of their convents in the Coptic, Syriac, and Arabic languages, have been of late secured for the most part to the British Museum, by Tattam and other travellers.

Missions have lately been undertaken among them, especially by the Church Missionary Society of England (commencing in 1825), and the United Presbyterians of America, but with little success so far.\footnote{A detailed, but very unfavorable description of the Copts is given by Edward W. Lane in his “Manners and Customs of the Modern Egyptians,” 1833. Notwithstanding this they stand higher than the other Egyptians. A. P. Stanley (Hist. of the Eastern Church, p. 95) says of them: “The Copts are still, even in their degraded state, the most civilized of the natives: the intelligence of Egypt still lingers in the Coptic scribes, who are on this account used as clerks in the offices of their conquerors, or as registrars of the water-marks of the Nile.” Comp. also the occasional notices of the Copts in the Egyptological writings of Wilkinson, Bunsen, Lepsius, Brugsch, and others.}

The Abyssinian church is a daughter of the Coptic, and was founded in the fourth century, by two missionaries from Alexandria, Frumentius and Aedesius. It presents a strange mixture of barbarism, ignorance, superstition, and Christianity. Its Ethiopic Bible, which dates perhaps from the first missionaries, includes in the Old Testament the apocryphal book of Enoch. The Chronicles of Axuma (the former capital of the country), dating from the fourth century, receive almost the same honor as the Bible. The council of Chalcedon is accounted an assembly of fools and heretics. The Abyssinian church has retained even more Jewish elements than the Coptic. It observes the Jewish Sabbath together with the Christian Sunday; it forbids the use of the flesh of swine and other unclean beasts; it celebrates a yearly feast of general lustration or rebaptizing of the whole nation; it retains the model of a sacred ark, called the ark of Zion, to which gifts and prayers are offered, and which forms the central point of public worship. It believes in the magical virtue of outward ceremonies, especially immersion, as the true regeneration. Singularly enough it honors Pontius Pilate as a saint, because be washed his hands of innocent blood. The endless controversies respecting the natures of Christ, which have died out elsewhere still rage there. The Abyssinians honor saints and pictures, but not images; crosses, but not the crucifix. Every priest carries a cross in his hand, and presents it to every one whom he meets, to be kissed. The numerous churches are small and dome-shaped above, and covered with reeds and straw. On the floor lie a number of staves and crutches, on which the people support themselves during the long service, as, like all the Orientals, they are without benches. Slight as are its remains of Christianity, Abyssinia still stands, in agriculture, arts, laws, and social condition, far above the heathen countries of Africa—a proof that even a barbaric Christianity is better than none.

The influences of the West have penetrated even to Abyssinia. The missions of the Jesuits in the seventeenth and eighteenth centuries, and of the Protestants in the nineteenth, have been prosecuted amidst many dangers and much self-denial, yet hitherto with but little success.\footnote{Especially worthy of note are the labors of the Basle missionaries, Samuel Gobat (now Anglican bishop in Jerusalem), Kugler, Isenberg, Blumhardt, and Krapf since 1830. Comp. Gobatin the Basler Missionsmagazin for 1834, Part 1 and 2. Isenberg: Abyssinien und die evangelische Mission, Bonn, 1844, 2 Bde. and Isenbergand Krapf: Journals, 1843. Also Harris: Highlands of Ethiopia 1844. The imported fragments of an Abyssinian translation of the Bible, dating from the fourth or fifth century, have drawn the attention of Westem scholars. Prof. A. Dillmann (now in Giessen) has since 1854 published the Aethiopic Old Testament, a grammar, and a lexicon of the Aethiopic language. Of the older works on Abyssinia the principal are Ludolphus: Historia Aethiopica, Frankf. 1681; Geddes: Church History of Aethiopia, Lond. 1696, and La Croze: Histoire du Christianisme d’Ethiopie et d’Armenie, La Haye, 1739. They have all drawn their principal materials from the Jesuits, especially from the general history of Tellez, published 1660.}

III. The Armenians. These are the most numerous, interesting, and hopeful of the Monophysite sects, and now the most accessible to evangelical Protestantism. Their nationality reaches back into hoary antiquity, like Mount Ararat, at whose base lies their original home. They were converted to Christianity in the beginning of the fourth century, under King Tiridates, by Gregory the Enlightener, the first patriarch and ecclesiastical writer and the greatest saint of the Armenians.\footnote{\textgreek{Φωτιστής}, Illuminator. He was married and had several sons. He was urgently invited to the Nicene council, but sent his son Aristax in his stead, to whom he resigned his office, and then withdrew himself for the rest of his life into a mountain-cave. There are homilies of his still extant, which were first printed in 1737 in Constantinople.} They were provided by him with monasteries and seminaries, and afterwards by Mesrob\footnote{Called Mesrop, Mjesrob, Mjesrop, and Marchtoz. Comp. respecting this man and the origin of the Armenian version of the Bible, the chronicle of his pupil, Moses Chorenensis, and the article by Petermann in Herzog’s Encycl. Bd. ix. p. 370 ff.} with a version of the Scriptures, made from the Greek with the help of the Syriac Peschito; which at the same time marks the beginning of the Armenian literature, since Mesrob had first to invent his alphabet. The Armenian canon has four books found in no other Bible; in the Old Testament, the History of Joseph and Asenath, and the Testament of the twelve Patriarchs, and in the New, the Epistle of the Corinthians to Paul and a Third, but spurious, Epistle of Paul to the Corinthians. The next oldest work in the Armenian language is the history of their land and people, by Moses Chorenensis, a half century later.

The Armenians fell away from the church of the Greek Empire in 552, from which year they date their era. The Persians favored the separation on political grounds, but were themselves thoroughly hostile to Christianity, and endeavored to introduce the Zoroastrian religion into Armenia. The Armenian church, being left unrepresented at the council of Chalcedon through the accidental absence of its bishops, accepted in 491 the Henoticon of the emperor Zeno, and at the synod of Twin (Tevin or Tovin, the capital at that time), held a.d. 595, declared decidedly for the Monophysite doctrine. The Confessio Armenica, which in other respects closely resembles the Nicene Creed, is recited by the priest at every morning service. The Armenian church had for a long time only one patriarch or Catholicus, who at first resided in Sebaste, and afterwards in the monastery of Etschmiezin (Edschmiadsin), their holy city, at the foot of Mount Ararat, near Erivan (now belonging to Russia), and had forty-two archbishops under him. At his consecration the dead hand of Gregory the Enlightener is even yet always used, as the medium of tactual succession. Afterwards other patriarchal sees were established, at Jerusalem (in 1311), at Sis, in Cilicia (in 1440), and after the fall of the Greek empire in Constantinople (1461).\footnote{Respecting the patriarchal and metropolitan sees and the bishoprics of the Armenians, comp. Le Quien, tom. i., and Wiltsch, Kirchliche Geographie und Statistik, ii. p. 375 ff.} In 637 Armenia fell under Mohammedan dominion, and belongs now partly to Turkey and partly to Russia. But the varying fortunes and frequent oppressions of their country have driven many thousands of the Armenians abroad, and they are now scattered in other parts of Russia and Turkey, as well as in Persia, India, and Austria.

The Armenians of the diaspora are mostly successful traders and brokers, and have become a nation and a church of merchant princes, holding great influence in Turkey. Their dispersion, and love of trade, their lack of political independence, their tenacious adherence to ancient national customs and rites, the oppressions to which they are exposed in foreign countries, and the influence which they nevertheless exercise upon these countries, make their position in the Orient, especially in Turkey, similar to that of the Jews in the Christian world.

The whole number of the Armenians is very variously estimated, from two and a half up to fifteen millions.\footnote{Stanley (History of the Eastern Church, p. 92), supported by Neale and Haxthausen (Transcaucasia), estimates the number of the Armenians at over eight millions. But Dr. G. W. Wood, of New York, formerly a missionary among them, informs me that their total number probably does not exceed six millions, of whom about two and a half millions are probably in Turkey.}

The Armenian church, it may be remarked, has long been divided into two parts, which, although internally very similar, are inflexibly opposed to each other. The united Armenians, since the council of Florence, a.d. 1439, have been connected with the church of Rome. To them belongs the congregation of the Mechitarists, which was founded by the Abbot Mechitar († 1749), and possesses a famous monastery on the island of San Lazzaro near Venice, from which centre it has successfully labored since 1702 for Armenian literature and education in the interest of the Roman Catholic church.\footnote{Comp. C. F. Neumann: Geschichte der armenischen Literatur nach den Werken der Mechitaristen, Leipzig, 1836. The chief work of the Mechitarists is the history of Armenia, by P. Michael Tschamtschean († 1823), in three vols., Venice, 1784.} The schismatical Armenians hold firmly to their peculiar ancient doctrines and polity. They regard themselves as the orthodox, and call the united or Roman Armenians schismatics.

Since 1830, the Protestant Missionary, Tract, and Bible societies of England, Basle, and the United States, have labored among the Armenians especially among the Monophysite portion, with great success, The American Board of Commissioners for Foreign Missions,\footnote{This oldest and most extensive of American missionary societies was founded a.d.1810, and is principally supported by the Congregationalists and New School Presbyterians.} in particular, has distributed Bibles and religious books in the Armenian and Armeno-Turkish\footnote{The Armeno-Turkish is the Turkish language written in Armenian characters.} language, and founded flourishing churches and schools in Constantinople, Broosa, Nicomedia, Trebizond, Erzroom, Aintab, Kharpoot, Diarbekir, and elsewhere. Several of these churches have already endured the crucial test of persecution, and justify bright hopes for the future. As the Jewish Synagogues of the diaspora were witnesses for monotheism among idolaters, and preparatory schools of Christianity, so are these Protestant Armenian churches, as well as the Protestant Nestorian, outposts of evangelical civilization in the East, and perhaps the beginning of a resurrection of primitive Christianity in the lands of the Bible and harbingers of the future conversion of the Mohammedans.\footnote{Compare, respecting the Armenian mission of the American Board, the publications of this Society; Eli Smithand H. G. O. Dwight: Missionary Researches in Armenia, Boston, 1833; Dr. H. G. O. Dwight: Christianity revived in the East, New York, 1850; H. Newcomb: Cyclopaedia of Missions, pp. 124-154. The principal missionaries among the Armenians are H. G. O. Dwight, W. Goodell, C. Hamlin, G. W. Wood, F. Riggs, D. Ladd, P. O. Powers, W. G. Schauffler (a Würtemberger, but educated at the Theol. Seminary of Andover, Mass.), and Benj. Schneider (a German from Pennsylvania, but likewise a graduate of Andover).}

IV. The youngest sect of the Monophysites, and the solitary memorial of the Monothelite controversy, are the Maronites, so called from St. Maron, and the eminent monastery founded by him in Syria (400).\footnote{He is probably the same Maron whose life Theodoret wrote, and to whom Chrysostomaddressed a letter when in exile. He is not to be confounded with the later John Maron, of the seventh century, who, according to the legendary traditions of the catholic Maronites, acting as papal legate at Antioch, converted the whole of Lebanon to the Roman church, and became their first patriarch. T he name “Maronites” occurs first in the eighth century, and that as a name of heretics, in John of Damascus.} They inhabit the range of Lebanon, with its declivities and valleys, from Tripolis on the North to the neighborhood of Tyre and the lake of Gennesaret on the South, and amount at most to half a million. They have also small churches in Aleppo, Damascus, and other places. They are pure Syrians, and still use the Syriac language in their liturgy, but speak Arabic. They are subject to a patriarch, who commonly resides in the monastery of Kanobin on Mt. Lebanon. They were originally Monothelites, even after the doctrine of one will of Christ, which is the ethical complement of the doctrine of one nature, had been rejected at the sixth ecumenical Council (a.d. 680). But after the Crusades (1182), and especially after 1596, they began to go over to the Roman church, although retaining the communion under both kinds, their Syriac missal, the marriage of priests, and their traditional fast-days, with some saints of their own, especially St: Maron.

From these came, in the eighteenth century, the three celebrated Oriental scholars, the Assemani, Joseph Simon († 1768), his brother Joseph Aloysius, and their cousin Stephen Evodius. These were born on Mt. Lebanon, and educated at the Maronite college at Rome.

There are also Maronites in Syria, who abhor the Roman church.\footnote{Respecting the present condition of the Maronites, comp. also Robinson’s Palestine, Ritter’s Erdkunde, Bd. xvii. Abtheil. 1, and Rödiger’s article in Herzog’s Encycl. Bd. x. p. 176 ff. A few years ago (1860), the Maronites drew upon themselves the sympathies of Christendom by the cruelties which their old hereditary enemies, the Druses, perpetrated upon them.}

\xsection{Character of the Pelagian Controversy}

§ 146. Character of the Pelagian Controversy.

IV. The Anthropological Controversies.

WORKS ON THE PELAGIAN CONTROVERSY IN GENERAL.

Sources:

I. Pelagius: Expositiones in epistolas Paulinas (composed before 410); Epistola ad Demetriadem, in 30 chapters (written a.d. 413); Libellus fidei ad Innocentium I. (417, also falsely called Explanatio Symboli ad Damasum). These three works have been preserved complete, as supposed works of Jerome, and have been incorporated in the Opera of this father (tom. xi. ed. of Vallarsius). Of the other writings of Pelagius (De natura; De libero arbitrio; Capitula; Epist. ad Innocent. I., which accompanied the Libellus fidei), we have only fragments in the works of his opponents, especially Augustine. In like manner we have only fragments of the writings of Coelestius: Definitiones; Symbolum ad Zosimum; and of Julianus of Eclanum: Libri iv. ad Turbantium episcopum contra Augustini primum de nuptiis; Libri viii. ad Florum contra Augustini secundum de nuptiis. Large and literal extracts in the extended replies of Augustine to Julian

II. Augustinus: De peccatorum meritis et remissione (412); De spiritu et litera (418); De natura et gratia (415); De gestis Pelagii (417); De gratia Christi et de peccato originali (418); De nuptiis et concupiscentia (419); Contra duas Epistolas Pelagianorum (420); Contra Julianum, libri vi. (421); Opus imperfectum contra Julianum (429); De gratia et libero arbitrio (426 or 427); De correptione et gratia (427) De praedestinatione sanctorum (428 or 429); De dono perseverantivae (429); and other anti-Pelagian writings, which are collected in the 10th volume of his Opera, in two divisions, ed. Bened. Par. 1690, and again Venet. 1733. (it is the Venice Bened. edition from which I have quoted throughout in this section. In Migne’s edition of Aug., Par. 1841, the anti-Pelagian writings form likewise the tenth tomus of 1912 pages.) Hieronymus: Ep. 133 (in Vallarsi’s, and in Migne’s ed.; or, Ep. 43 in the Bened. ed.) ad Ctesiphontem (315); Dialogi contra Pelagianos, libri iii. (Opera, ed. Vallars. vol. ii. f. 693–806, and ed. Migne, ii. 495–590). P. Orosius: Apologeticus c. Pelag. libri iii. (Opera, ed. Haverkamp). Marius Mercator, a learned Latin monk in Constantinople (428–451): Commonitoria, 429, 431 (ed. Baluz. Paris, 1684, and Migne, Par. 1846). Collection of the Acta in Mansi, tom. iv.

Literature:

Gerh. Joh. Vossius: Hist. de controversiis, quas Pelagius ejusque reliquiae moverunt, libri vii. Lugd. Batav. 1618 (auct. ed. Amstel. 1655). Cardinal Henr. Norisius: Historia Pelagiana et dissert. de Synodo Quinta Oecumen. Batavii, 1673, fol. (and in Opera, Veron. 1729, i.). Garnier (Jesuit): Dissert. vii. quibus integra continentur Pelagianorum hist. (in his ed. of the Opera of Marius Mercator, i. 113). The Praefatio to the 10th vol. of the Benedictine edition of Augustine’s Opera. Corn. Jansenius († 1638): Augustinus, sive doctrina S. Augustini de humanae naturae sanitate, aegritudine, medicina, adv. Pelagianos et Massilienses. Lovan. 1640, fol. (He read Augustine twenty times, and revived his system in the Catholic church.) Tillemont: Mémoires, etc. Tom. xiii. pp. 1–1075, which is entirely devoted to the life of Augustine. Ch. Wilh. Fr. Walch: Ketzerhistorie. Leipz. 1770. Bd. iv. and v. Schröckh: Kirchengeschichte. Parts xiv. and xv. (1790). G. F. Wiggers (sen.): Versuch einer pragmatischen Darstellung des Augustinismus und Pelagianismus, in zwei Theilen. Hamburg, 1833. (The first part appeared 1821 in Berlin; the second, which treats of Semi-Pelagianism, in 1833 at Hamburg. The common title-page bears date 1833. The first part has also been translated into English by Prof. Emerson, Andover, 1840). J. L. Jacobi: Die Lehre des Pelagius. Leipzig, 1842. F. Böhringer: Die Kirche Christi in Biographien. Bd. i. Th. 3, pp. 444–626, Zürich, 1845. Gieseler: Kirchengeschichte. Bd. i. Abth. 2 pp. 106–131 (4th ed. 1845, entirely favorable to Pelagianism). Neander: Kirchengeschichte. Bd. iv. (2d ed. 1847, more Augustinian). Schaff: The Pelagian Controversy, in the Bibliotheca Sacra, Andover, May, 1848 (No. xviii.). Theod. Gangauf: Metaphysische Psychologie des heiligen Augustinus. Augsb. 1852. Thorough, but not completed. H. Hart Milman: History of Latin Christianity. New York, 1860, vol. i. ch. ii. pp. 164–194. Jul. Müller: Die christliche Lehre von der Sünde. Bresl. 1838, 5th ed. 1866, 2 vols. (An English translation by Pulsford, Edinburgh.) The same: Der Pelagianismus. Berlin, 1854. (A brief, but admirable essay.) Hefele: Conciliengeschichte. Bd. ii. 1856, p. 91 ff. W. Cunningham: Historical Theology. Edinburgh, 1863, vol. i, pp. 321–358. Fr. Wörter (R.C.): Der Pelagianismus nach seinem Ursprung und seiner Lehre. Freiburg, 1866. Nourrisson: La philosophie de S. Augustin. Par. 1866, 2 vols. (vol. i. 452 ff.; ii. 352 ff.). Comp. also the literature in § 178, and the relevant chapters in the Doctrine-Histories of Münscher, Baumgarten-Crusius, Hagenbach, Neander, Baur, Beck, Shedd.

While the Oriental Church was exhausting her energies in the Christological controversies, and, with the help of the West, was developing the ecumenical doctrine of the person of Christ, the Latin church was occupied with the great anthropological and soteriological questions of sin and grace, and was bringing to light great treasures of truth, without either help from the Eastern church or influence upon her. The third ecumenical council, it is true, condemned Pelagianism, but without careful investigation, and merely on account of its casual connection with Nestorianism. The Greek historians, Socrates, Sozomen, Theodoret, and Evagrius, although they treat of that period, take not the slightest notice of the Pelagian controversies. In this fact we see the predominantly practical character of the West, in contradistinction to the contemplative and speculative East. Yet the Christological and anthropologico-soteriological controversies are vitally connected, since Christ became man for the redemption of man. The person and the work of the Redeemer presuppose on the one hand man’s capability of redemption, and on the other his need of redemption. Manichaeism denies the former, Pelagianism the latter. In opposition to these two fundamental anthropological heresies, the church was called to develope the whole truth.

Before Augustine the anthropology of the church was exceedingly crude and indefinite. There was a general agreement as to the apostasy and the moral accountability of man, the terrible curse of sin, and the necessity of redeeming grace; but not as to the extent of native corruption, and the relation of human freedom to divine grace in the work of regeneration and conversion. The Greek, and particularly the Alexandrian fathers, in opposition to the dualism and fatalism of the Gnostic systems, which made evil a necessity of nature, laid great stress upon human freedom, and upon the indispensable cooperation of this freedom with divine grace; while the Latin fathers, especially Tertullian and Cyprian, Hilary and Ambrose, guided rather by their practical experience than by speculative principles, emphasized the hereditary sin and hereditary guilt of man, and the sovereignty of God’s grace, without, however, denying freedom and individual accountability.\footnote{On the anthropology of the ante-Nicene and Nicene fathers, comp. the relevant sections in the larger works on Doctrine History, and Wiggers, \emph{l.c.} vol. l. p. 407 ff.} The Greek church adhered to her undeveloped synergism,\footnote{From \textgreek{σὺν}, and \textgreek{ἔργον.} There are, it may be remarked, different forms of synergism. The synergism of Melanchthon subordinates the human activity to the divine, and assigns to grace the initiative in the work of conversion.} which coordinates the human will and divine grace as factors in the work of conversion; the Latin church, under the influence of Augustine, advanced to the system of a divine, monergism,\footnote{From \textgreek{μόνον} and \textgreek{ἔργον}.} which gives God all the glory, and makes freedom itself a result of grace; while Pelagianism, on the contrary, represented the principle of a human monergism, which ascribes the chief merit of conversion to man, and reduces grace to a mere external auxiliary. After Augustine’s death, however the intermediate system of Semi-Pelagianism, akin to the Greek synergism, became prevalent in the West.

Pelagius and Augustine, in whom these opposite forms of monergism were embodied, are representative men, even more strictly than Arius and Athanasius before them, or Nestorius and Cyril after them. The one, a Briton, more than once convulsed the world by his errors; the other, an African, more than once by his truths. They represented principles and tendencies, which, in various modifications, extend through the whole history of the church, and reappear in its successive epochs. The Gottschalk controversy in the ninth century, the Reformation, the synergistic controversy in the Lutheran church, the Arminian in the Reformed, and the Jansenistic in the Roman Catholic, only reproduce the same great contest in new and specific aspects. Each system reflects the personal character and experience of its author. Pelagius was an upright monk, who without inward conflicts won for himself, in the way of tranquil development, a legal piety which knew neither the depths of sin nor the heights of grace. Augustine, on the other hand, passed through sharp convulsions and bitter conflicts, till he was overtaken by the unmerited grace of God, and created anew to a life of faith and love. Pelagius had a singularly clear, though contracted mind, and an earnest moral purpose, but no enthusiasm for lofty ideals; and hence he found it not hard to realize his lower standard of holiness. Augustine had a bold and soaring intellect, and glowing heart, and only found peace after he had long been tossed by the waves of passion; he had tasted all the misery of sin, and then all the glory of redemption, and this experience qualified him to understand and set forth these antagonistic powers far better than his opponent, and with a strength and fulness surpassed only by the inspired apostle Paul. Indeed, Augustine, of all the fathers, most resembles, in experience and doctrine, this very apostle, and stands next to him in his influence upon the Reformers.

The Pelagian controversy turns upon the mighty antithesis of sin and grace. It embraces the whole cycle of doctrine respecting the ethical and religious relation of man to God, and includes, therefore, the doctrines of human freedom, of the primitive state, of the fall, of regeneration and conversion, of the eternal purpose of redemption, and of the nature and operation of the grace of God. It comes at last to the question, whether redemption is chiefly a work of God or of man; whether man needs to be born anew, or merely improved. The soul of the Pelagian system is human freedom; the soul of the Augustinian is divine grace. Pelagius starts from the natural man, and works up, by his own exertions, to righteousness and holiness. Augustine despairs of the moral sufficiency of man, and derives the now life and all power for good from the creative grace of God. The one system proceeds from the liberty of choice to legalistic piety; the other from the bondage of sin to the evangelical liberty of the children of God. To the former Christ is merely a teacher and example, and grace an external auxiliary to the development of the native powers of man; to the latter he is also Priest and King, and grace a creative principle, which begets, nourishes, and consummates a new life. The former makes regeneration and conversion a gradual process of the strengthening and perfecting of human virtue; the latter makes it a complete transformation, in which the old disappears and all becomes new. The one loves to admire the dignity and strength of man; the other loses itself in adoration of the glory and omnipotence of God. The one flatters natural pride, the other is a gospel for penitent publicans and sinners. Pelagianism begins with self-exaltation and ends with the sense of self-deception and impotency. Augustinianism casts man first into the dust of humiliation and despair, in order to lift him on the wings of grace to supernatural strength, and leads him through the hell of self-knowledge up to the heaven of the knowledge of God. The Pelagian system is clear, sober, and intelligible, but superficial; the Augustinian sounds the depths of knowledge and experience, and renders reverential homage to mystery. The former is grounded upon the philosophy of common sense, which is indispensable for ordinary life, but has no perception of divine things; the latter is grounded upon the philosophy of the regenerate reason, which breaks through the limits of nature, and penetrates the depths of divine revelation. The former starts with the proposition: Intellectus praecedit fidem; the latter with the opposite maxim: Fides praecedit intellectum. Both make use of the Scriptures; the one, however, conforming them to reason, the other subjecting reason to them. Pelagianism has an unmistakable affinity with rationalism, and supplies its practical side. To the natural will of the former system corresponds the natural reason of the latter; and as the natural will, according to Pelagianism, is competent to good, so is the natural reason, according to rationalism, competent to the knowledge of the truth. All rationalists are Pelagian in their anthropology; but Pelagius and Coelestius were not consistent, and declared their agreement with the traditional orthodoxy in all other doctrines, though without entering into their deeper meaning and connection. Even divine mysteries may be believed in a purely external, mechanical way, by inheritance from the past, as the history of theology, especially in the East, abundantly proves.

The true solution of the difficult question respecting the relation of divine grace to human freedom in the work of conversion, is not found in the denial of either factor; for this would either elevate man to the dignity of a self-redeemer, or degrade him to an irrational machine, and would ultimately issue either in fatalistic pantheism or in atheism; but it must be sought in such a reconciliation of the two factors as gives full weight both to the sovereignty of God and to the responsibility of man, yet assigns a preëminence to the divine agency corresponding to the infinite exaltation of the Creator and Redeemer above the sinful creature. And although Angustine’s solution of the problem is not altogether satisfactory, and although in his zeal against the Pelagian error he has inclined to the opposite extreme; yet in all essential points, he has the Scriptures, especially the Epistles of Paul, as well as Christian experience, and the profoundest speculation, on his side. Whoever reads the tenth volume of his works, which contains his Anti-Pelagian writings in more than fourteen hundred folio columns (in the Benedictine edition), will be moved to wonder at the extraordinary wealth of thought and experience treasured in them for all time; especially if he considers that Augustine, at the breaking out of the Pelagian controversy, was already fifty-seven years old, and had passed through the Manichaen and Donatist controversies. Such giants in theology could only arise in an age when this queen of the sciences drew into her service the whole mental activity of the time.

The Pelagian controversy was conducted with as great an expenditure of mental energy, and as much of moral and religious earnestness, but with less passion and fewer intrigues, than the Trinitarian and Christological conflicts in the East. In the foreground stood the mighty genius and pure zeal of Augustine, who never violated theological dignity, and, though of thoroughly energetic convictions, had a heart full of love. Yet even he yielded so far to the intolerant spirit of his time as to justify the repression of the Donatist and Pelagian errors by civil penalties.

\xsection{External History of the Pelagian Controversy, A.D. 411-431}

§ 147. External History of the Pelagian Controversy, A.D. 411–431.

Pelagius\footnote{His British name is said to have been \emph{Morgan}, that is, Of the sea, \emph{Marigena}, in Greek \textgreek{Πελάγιος.}} was a simple monk, born about the middle of the fourth century in Britain, the extremity of the then civilized world. He was a man of clear intellect, mild disposition, learned culture, and spotless character; even Augustine. with all his abhorrence of his doctrines, repeatedly speaks respectfully of the man.\footnote{Comp. the passages where Augustinespeaks of Pelagius, in Wiggers, l. c. i. p. 35 f. Yet Augustine, not without reason, accuses him of duplicity, on account of his conduct at the synod of Diospolis in Palestine. Wiggers (i. p. 40) says of him: “it must be admitted that Pelagius was not always sufficiently straightforward; that he did not always express his views without ambiguity; that, in fact, he sometimes in synods condemned opinions which were manifestly his own. This may have arisen, it is true, in great part from his love of peace and the slight value which he attached to theoretical opinions.”} He studied the Greek theology, especially that of the Antiochian school, and early showed great zeal for the improvement of himself and of the world. But his morality was not so much the rich, deep life of faith, as it was the external legalism, the ascetic self-discipline and self-righteousness of monkery. It was characteristic, that, even before the controversy, he took great offence at the well-known saying of Augustine: “Give what thou commandest, and command what thou wilt.”\footnote{“Da quod jubes, et jube quod vis,” Confess. l. x. c. 29, et passim. Augustinehimself relates the above-mentioned fact, De dono persev. c. 20 (or § 53, tom. x. f. 851): “Quae mea verba, Pelagius Romae, cum a quodam fratre et coëpiscopo meo fuissent eo praesente commemorata, ferre non potuit, et contradicens aliquanto commotius pene cum eo, qui illa commemoraverat, litigavit.”} He could not conceive, that the power to obey the commandment must come from the same source as the commandment itself. Faith, with him, was hardly more than a theoretical belief; the main thing in religion was moral action, the keeping of the commandments of God by one’s own strength. This is also shown in the introductory remarks of his letter to Demetrias, a noble Roman nun, of the gens Anicia, in which he describes a model virgin as a proof of the excellency of human nature: “As often as I have to speak concerning moral improvement and the leading of a holy life, I am accustomed first to set forth the power and quality of human nature, and to show what it can accomplish.\footnote{“Soleo prius humanae naturae vim qualitatemque monstrare, et quid efficere possit, ostendere.” Ep. ad Demetr. c. 2.} For never are we able to enter upon the path of the virtues, unless hope, as companion, draws us to them. For every longing after anything dies within us, so soon as we despair of attaining that thing.”

In the year 409, Pelagius, already advanced in life, was in Rome, and composed a brief commentary on the Epistles of Paul. This commentary, which has been preserved among the works of Jerome, displays a clear and sober exegetical talent.\footnote{It found its way among the works of Jerome(tom. xi. ed. Vallars., and in Migne’s edition, tom. xi. f. 643-902) before the breaking out of the controversy, but has received doctrinal emendations from Cassiodorus, at least in the Epistle to the Romans. The confounding of Pelagius with Jeromearose partly from his accommodation to the ecclesiastical terminology, partly from his actual agreement with the prevailing tendency of monasticism. It is remarkable that both wrote an ascetic letter to the nun Demetrias. Comp. Jerome, Ep. 130 (ed. Vallarsi, and Migne, or 97 in the Bened. ed.) ad Demetriadem de servanda Virginitate (written in 414). She had also correspondence with Augustine. Semler has published the letters of Augustine, Jerome, and Pelagius to Demetrias in a separate form (Halle, 1775). Some have also ascribed to Pelagius the ascetic Epistola ad Celantiam matronam de ratione pie vivendi, which, like his Ep. ad Demetriadem, has found its way into the Epistles of Jerome(Ep. 148 in Vallarsi’s ed. tom. i. 1095, and in Migne’s ed. tom. i. 1204). The monasticism of Pelagius, however, was much cooler, more sober, and more philosophical than that of the enthusiastic Jerome, inclined as he was to all manner of extravagances.} He labored quietly and peacefully for the improvement of the corrupt morals of Rome, and converted the advocate Coelestius, of distinguished, but otherwise unknown birth, to his monastic life, and to his views. It was from this man, younger, more skilful in argument, more ready for controversy, and more rigorously consistent than his teacher, that the controversy took its rise. Pelagius was the moral author, Coelestius the intellectual author, of the system represented by them.\footnote{To this extent Pelagius and Coelestius appear to sustain a relation to Pelagianism similar to that which Dr. Pusey and John Henry Newman did to Puseyism. Jerome(in his letter to Ctesiphon) says of Coelestius, that he was, although the disciple of Pelagius, yet teacher and leader of the whole array (magister et totius ductor exercitus). Augustinecalls Pelagius more dissembling and crafty, Coelestius more frank and open (De peccato orig. c. 12). Marius Mercator ascribes to Coelestius an incredibilis loquacitas. But Augustineand Julianof Eclanum also mutually reproach each other with a vagabunda loquacitas.} They did not mean actually to found a new system, but believed themselves in accordance with Scripture and established doctrine. They were more concerned with the ethical side of Christianity than with the dogmatic; but their endeavor after moral perfection was based upon certain views of the natural power of the will, and these views proved to be in conflict with anthropological principles which had been developed in the African church for the previous ten years under the influence of Augustine.

In the year 411, the two friends, thus united in sentiment, left Rome, to escape the dreaded Gothic King Alaric, and went to Africa. They passed through Hippo, intending to visit Augustine, but found that he was just then at Carthage, occupied with the Donatists. Pelagius wrote him a very courteous letter, which Augustine answered in a similar tone; intimating, however, the importance of holding the true doctrine concerning sin. “Pray for me,” he said, “that God may really make me that which you already take me to be.” Pelagius soon proceeded to Palestine. Coelestius applied for presbyters’ orders in Carthage, the very place where he had most reason to expect opposition. This inconsiderate step brought on the crisis. He gained many friends, it is true, by his talents and his ascetic zeal, but at the same time awakened suspicion by his novel opinions.

The deacon Paulinus of Milan, who was just then in Carthage, and who shortly afterwards at the request of Augustine wrote the life of Ambrose, warned the bishop Aurelius against Coelestius, and at a council held by Aurelius at Carthage in 412,\footnote{According to Mansi and the common view. The brothers Ballerini and Hefele (ii. 91) decide in favor of the year 411. The incomplete Acta of the council are found in Mansi, tom. iv. fol. 289 sqq., and in the Commonitorium Marii Mercatoris \emph{ibidem}, f. 293.} appeared as his accuser. Six or seven errors, he asserted he had found in the writings of Coelestius:

1. Adam was created mortal, and would have died, even if he had not sinned.

2. Adam’s fall injured himself alone, not the human race.

3. Children come into the world in the same condition in which Adam was before the fall.

4. The human race neither dies in consequence of Adam’s fall, nor rises again in consequence of Christ’s resurrection.

5. Unbaptized children, as well as others, are saved.\footnote{Marius Mercator, it is true, does not cite this proposition among the others, f. 292, but he brings it up subsequently, f. 296: “In ipsa autem accusatione capitulorum, quae eidem Pelagio tum objecta sunt, etiam haec continentur, cum aliis execrandis, quae Coelestius ejus discipulus sentiebat, id est, \emph{infantes etiamsi non baptizentur, habere vitam aeternam}.”}

6. The law, as well as the gospel, leads to the kingdom of heaven.

7. Even before Christ there were sinless men.

The principal propositions were the second and third, which are intimately connected, and which afterwards became the especial subject of controversy.

Coelestius returned evasive answers. He declared the propositions to be speculative questions of the schools, which did not concern the substance of the faith, and respecting which different opinions existed in the church. He refused to recant the errors charged upon him, and the synod excluded him from the communion of the church. He immediately went to Ephesus, and was there ordained presbyter.

Augustine had taken no part personally in these transactions. But as the Pelagian doctrines found many adherents even in Africa and in Sicily, he wrote several treatises in refutation of them so early as 412 and 415, expressing himself, however, with respect and forbearance.\footnote{De peccatorum meritis et remissione; De spiritu et liters; De natura et gratia; De perfectione justitiae hominis.}

\xsection{The Pelagian Controversy in Palestine}

§ 148. The Pelagian Controversy in Palestine.

Meanwhile, in 414, the controversy broke out in Palestine, where Pelagius was residing, and where he had aroused attention by a letter to the nun Demetrias. His opinions gained much wider currency there, especially among the Origenists; for the Oriental church had not been at all affected by the Augustinian views, and accepted the two ideas of freedom and grace, without attempting to define their precise relation to each other. But just then there happened to be in Palestine two Western theologians, Jerome and Orosius; and they instituted opposition to Pelagius.

Jerome, who lived a monk at Bethlehem, was at first decidedly favorable to the synergistic theory of the Greek fathers, but at the same time agreed with Ambrose and Augustine in the doctrine of the absolutely universal corruption of sin.\footnote{Compare, respecting his relation to Pelagianism, O. Zöckler: Hieronymus (1865), p. 310 ff. and p. 420 ff.} But from an enthusiastic admirer of Origen he had been changed to a bitter enemy. The doctrine of Pelagius concerning free will and the moral ability of human nature he attributed to the influence of Origen and Rufinus; and he took as a personal insult an attack of Pelagius on some of his writings.\footnote{Comp. Jerome: Praefat. libri i. in Jeremiam (Opera, ed. Vallarsi, tom. iv. 834 sq.), where he speaks very contemptuously of Pelagius: “Nuper indoctus calumniator erupit, qui commentarios meos in epistolam Pauli ad Ephesios reprehendendos putat.” Soon afterwards he designates Grunnius, i.e., Rufinus, as his praecursor, and thus connects him with the Origenistic heresies. Pelagius had also expressed himself unfavorably respecting his translation of the Old Testament from the Hebrew.} He therefore wrote against him, though from wounded pride and contempt he did not even mention his name; first in a letter answering inquiries of a certain Ctesiphon at Rome (415);\footnote{Epist. 133 ad Ctesiphont. Adv. Pelag. (Opera, i. 1025-1042).} then more at length in a dialogue of three books against the Pelagians, written towards the end of the year 415, and soon after the acquittal of Pelagius by the synod of Jerusalem.\footnote{Dialogus c. Pelag. (Opera, tom. ii. 693-806).} Yet in this treatise and elsewhere Jerome himself teaches the freedom of the will, and only a conditional predestination of divine foreknowledge, and thus, with all his personal bitterness against the Pelagians, stands on Semi-Pelagian ground, though Augustine eulogizes the dialogue.\footnote{Op. imperf. contra Jul. iv. 88, where he says of it: Mira et ut talem fidem decebat, venustate composuit. The judgment is just as to the form, but too favorable as to the contents of this dialogue. Comp. Zöckler, Hieronymus, p. 428.}

A young Spanish ecclesiastic, Paul Orosius, was at that time living with Jerome for the sake of more extended study, and had been sent to him by Augustine with letters relating to the Origenistic and Pelagian controversy.

At a diocesan synod, convoked by the bishop John of Jerusalem in June, 415,\footnote{The Acta of the Conventus Hierosolymitanus, according to a report of Orosius, in his Apologia pro libertate arbitrii, cap. 3 and 4, are found in Mansi, iv. 301 sqq.} this Orosius appeared against Pelagius, and gave information that a council at Carthage had condemned Coelestius, and that Augustine had written against his errors. Pelagius answered with evasion and disparagement: “What matters Augustine to me?” Orosius gave his opinion, that a man who presumed to speak contumeliously of the bishop to whom the whole North African church owed her restoration (alluding apparently to the settlement of the Donatist controversies), deserved to be excluded from the communion of the whole church. John, who was a great admirer of the condemned Origen, and made little account of the authority of Augustine, declared: “I am Augustine,”\footnote{“Augustinus ego sum.” To this Orosiusreplied not infelicitously: “Si Angustini personam sumis, Augustini sententiam sequere.” Mansi, iv. 308.} and undertook the defence of the accused. He permitted Pelagius, although only a monk and layman, to take his seat among the presbyters.\footnote{Orosiuswas much scandalized by the fact that a bishop should order “laicum in consessu presbyterorum, reum haereseos manifestae in medio catholicorum sedere.”} Nor did he find fault with Pelagius’ assertion, that man can easily keep the commandments of God, and become free from sin, after the latter had conceded, in a very indefinite manner, that for this the help of God is necessary. Pelagius had the advantage of understanding both languages, while John spoke only Greek, Orosius only Latin, and the interpreter often translated inaccurately. After much discussion it was resolved, that the matter should be laid before the Roman bishop, Innocent, since both parties in the controversy belonged to the Western church. Meanwhile these should refrain from all further attacks on each other.

A second Palestinian council resulted still more favorably to Pelagius. This consisted of fourteen bishops, and was held at Diospolis or Lydda, in December of the same year, under the presidency of Eulogius, bishop of Caesarea, to judge of an accusation preferred by two banished bishops of Gaul, Heros and Lazarus, acting in concert with Jerome.\footnote{The scattered accounts of the Concilium Diospolitanum are collected in Mansi, tom. iv. 311 sqq. Comp. Hefele, ii. p. 95 ff.} The charges were unskilfully drawn up, and Pelagius was able to avail himself of equivocations, and to condemn as folly, though not as heresy, the teachings of Coelestius, which were also his own. The synod, of which John of Jerusalem was a member, did not go below the surface of the question, nor in fact understand it, but acquitted the accused of all heresy. Jerome is justified in calling this a “miserable synod;”\footnote{“Quidquid in illa miserabili synodo Diospolitana dixisse se denegat, in hoc opere confitetur,” he wrote, a.d.419, in a letter to Augustine(Ep. 143, ed. Vallars. tom. i. 1067). Comp. Mansi, iv. 315.} although Augustine is also warranted in saying: “it was not heresy, that was there acquitted, but the man who denied the heresy.”\footnote{Comp. Augustine, De gestis Pelagii, c. 1 sqq. (tom. x. fol. 192 sqq.). Pope innocent I. (402-417) wrote a consoling letter to Jerome, and a letter of reproof to John of Jerusalem for his inaction. Epp. 136 and 137 in Jerome’s Epistles.}

Jerome’s polemical zeal against the Pelagians cost him dear. In the beginning of the year 416, a mob of Pelagianizing monks, ecclesiastics, and vagabonds broke into his monastery at Bethlehem, maltreated the inmates, set the building on fire, and compelled the aged scholar to take to flight. Bishop John of Jerusalem let this pass unpunished. No wonder that Jerome, even during the last years of his life, in several epistles indulges in occasional sallies of anger against Pelagius, whom he calls a second Catiline.

\xsection{Position of the Roman Church. Condemnation of Pelagianism}

§ 149. Position of the Roman Church. Condemnation of Pelagianism.

The question took another turn when it was brought before the Roman see. Two North African synods, in 416, one at Carthage and one at Mileve (now Mela), again condemned the Pelagian error, and communicated their sentence to pope Innocent.\footnote{See the proceedings of the Concilium Carthaginense in Mansi, iv. 321 sqq., and of the Concilium Milevitanurn, ibid. f. 326 sqq.} A third and more confidential letter was addressed to him by five North African bishops, of whom Augustine was one.\footnote{Mansi, iv. 337 sqq.} Pelagius also sent him a letter and a confession of faith, which, however, were not received in due time.

Innocent understood both the controversy and the interests of the Roman see. He commended the Africans for having addressed themselves to the church of St. Peter, before which it was seemly that all the affairs of Christendom should be brought; he expressed his full agreement with the condemnation of Pelagius, Coelestius, and their adherents; but he refrained from giving judgment respecting the synod of Diospolis.\footnote{The answers of Innocent are found in Mansi, tom. iii. f. 1071 sqq.}

But soon afterwards (in 417) Innocent died, and was succeeded by Zosimus, who was apparently of Oriental extraction (417–418).\footnote{The notices of his life, as well as the Epistolae and Decreta Zosimi papae, are collected in Mansi, iv. 345 sqq.} At this juncture, a letter from Pelagius to Innocent was received, in which he complained of having suffered wrong, and gave assurance of his orthodoxy. Coelestius appeared personally in Rome, and succeeded by his written and oral explanations in satisfying Zosimus. He, like Pelagius, demonstrated with great fulness his orthodoxy on points not at all in question, represented the actually controverted points as unimportant questions of the schools, and professed himself ready, if in error, to be corrected by the judgment of the Roman bishop.

Zosimus, who evidently had no independent theological opinion whatever, now issued (417) to the North African bishops an encyclical letter accompanied by the documentary evidence, censuring them for not having investigated the matter more thoroughly, and for having aspired, in foolish, overcurious controversies, to know more than the Holy Scriptures. At the same time he bore emphatic testimony to the orthodoxy of Pelagius and Coelestius, and described their chief opponents, Heros and Lazarus, as worthless characters, whom he had visited with excommunication and deposition. They in Rome, he says, could hardly refrain from tears, that such men, who so often mentioned the gratia Dei and the adjutorium divinum, should have been condemned as heretics. Finally he entreated the bishops to submit themselves to the authority of the Roman see.\footnote{See the two epistles of Zosimus ad Africanos episcopos, in Mansi, iv. 350 and 353.}

This temporary favor of the bishop of Rome towards the Pelagian heresy is a significant presage of the indulgence of later popes for Pelagianizing tendencies, and of the papal condemnation of Jansenism.

The Africans were too sure of their cause, to yield submission to so weak a judgment, which, moreover, was in manifest conflict with that of Innocent. In a council at Carthage, in 417 or 418, they protested, respectfully but decidedly, against the decision of Zosimus, and gave him to understand that he was allowing himself to be greatly deceived by the indefinite explanations of Coelestius. In a general African council held at Carthage in 418, the bishops, over two hundred in number, defined their opposition to the Pelagian errors, in eight (or nine) Canons, which are entirely conformable to the Augustinian view.\footnote{It is the 16th Carthaginian synod. Mansi gives the canons in full, tom. iii. 810-823 (Comp. iv. 377). So also Wiggers, i. 214 ff. Hefele, ii. pp. 102-106, gives only extracts of them.} They are in the following tenor:

1. Whosoever says, that Adam was created mortal, and would, even without sin, have died by natural necessity, let him be anathema.

2. Whoever rejects infant baptism, or denies original sin in children, so that the baptismal formula, “for the remission of sins,” would have to be taken not in a strict, but in a loose sense, let him be anathema.

3. Whoever says, that in the kingdom of heaven, or elsewhere, there is a certain middle place, where children dying without baptism live happy (beate vivant), while yet without baptism they cannot enter into the kingdom of heaven, i.e., into eternal life, let him be anathema.\footnote{It is significant, that the third canon, which denies the salvation of unbaptized children, is of doubtful authenticity, and is wanting in Isidore and Dionysius. Hence the difference in the number of the canons against the Pelagians, as to whether there are 8 or 9.}

The fourth canon condemns the doctrine that the justifying grace of God merely effects the forgiveness of sins already committed; and the remaining canons condemn other superficial views of the grace of God and the sinfulness of man.

At the same time the Africans succeeded in procuring from the emperor Honorius edicts against the Pelagians.

These things produced a change in the opinions of Zosimus, and about the middle of the year 418, he issued an encyclical letter to all the bishops of both East and West, pronouncing the anathema upon Pelagius and Coelestius (who had meanwhile left Rome), and declaring his concurrence with the decisions of the council of Carthage in the doctrines of the corruption of human nature, of baptism, and of grace. Whoever refused to subscribe the encyclical, was to be deposed, banished from his church, and deprived of his property.\footnote{Epistola tractoria, or tractatoria, of which only some fragments are extant. Comp. Mansi, iv. 370. This letter was written \emph{after} and not \emph{before} the African council of 418 and the promulgation of the sacrum rescriptum of Honorius against the Pelagians, as Tillemont (xiii. 738) and the Benedictines (in the Preface to the 10th volume of the Opera August. § 18) have proved, in opposition to Baronius, Noris, and Garnier.}

Eighteen bishops of Italy refused to subscribe, and were deposed. Several of these afterwards recanted, and were restored.

The most distinguished one of them, however, the bishop Julian, of Eclanum, a small place near Capua in Campania, remained steadfast till his death, and in banishment vindicated his principles with great ability and zeal against Augustine, to whom he attributed all the misfortunes of his party, and who elaborately confuted him.\footnote{In two large works: Contra Julianum, libri vi. (Opera, tom. x. f. 497-711), and in the Opus imperfectum contra secundam Juliani responsionem, in six books (tom. x. P. ii. f. 874-1386), before completing which he died (a.d.430).} Julian was the most learned, the most acute, and the most systematic of the Pelagians, and the most formidable opponent of Augustine; deserving respect for his talents, his uprightness of life, and his immovable fidelity to his convictions, but unquestionably censurable for excessive passion and overbearing pride.\footnote{Gennadius, in his Liber de scriptoribus ecclesiastics, calls Julianof Eclanum “vir acer ingenio, in divinis scripturis doctus, Graeca et Latina lingua scholasticus.” By Augustine, however, in the Opus imperf. contra Jul. l. iv. 50 (Opera, x. P. ii. fol. 1163), he is called “in disputatione loquacissimus, in contentione calumniosissimus, in professions fallacissimus,” because he maligned the Catholics, while giving himself out for a Catholic. He was married.}

Julian, Coelestius, and other leaders of the exiled Pelagians, were hospitably received in Constantinople, in 429, by the patriarch Nestorius, who sympathized with their doctrine of the moral competency of the will, though not with their denial of original sin, and who interceded for them with the emperor and with pope Celestine, but in vain. Theodosius, instructed by Marius Mercator in the merits of the case, commanded the heretics to leave the capital (429). Nestorius, in a still extant letter to Coelestius,\footnote{In Marius Mercator, in a Latin translation, ed. Garnier-Migne, p. 182.} accords to him the highest titles of honor, and comforts him with the examples of John the Baptist and the persecuted apostles. Theodore of Mopsuestia († 428), the author of the Nestorian Christology, wrote in 419 a book against the Augustinian anthropology, of which fragments only are left.\footnote{In Photius, Bibl. Cod. 177, and in the Latin translation of Marius Mercator, also in the works of Jerome, tom. ii. 807-814 (ed. Vall.). The book was written contra Hiramum, i.e., Hieronymum, and was entitled: \textgreek{Πρὸς τοὺς λέγοντας φύσει καὶ οὐ γνώμῃ πταίειν τοὺς ἀνθρώπους λόγοι πέντε}, against those who say that men sin by nature, and not by free will.}

Of the subsequent life of Pelagius and Coelestius we have no account. The time and place of their death are entirely unknown. Julian is said to have ended his life a schoolmaster in Sicily, a.d. 450, after having sacrificed all his property for the poor during a famine.

Pelagianism was thus, as early as about the year 430, externally vanquished. It never formed an ecclesiastical sect, but simply a theological school. It continued to have individual adherents in Italy till towards the middle of the fifth century, so that the Roman bishop, Leo the Great, found himself obliged to enjoin on the bishops by no means to receive any Pelagian to the communion of the church without an express recantation.

At the third ecumenical council in Ephesus, a.d. 431 (the year after Augustine’s death), Pelagius (or more properly Coelestius) was put in the same category with Nestorius. And indeed there is a certain affinity between them: both favor an abstract separation of the divine and the human, the one in the person of Christ, the other in the work of conversion, forbidding all organic unity of life. According to the epistle of the council to pope Celestine, the Western Acta against the Pelagians were read at Ephesus and approved, but we do not know in which session. We are also ignorant of the discussions attending this act. In the canons, Coelestius, it is true, is twice condemned together with Nestorius, but without statement of his teachings.\footnote{Can. i. and Can. iv. The latter reads: “If clergymen fall away and either secretly or publicly hold with Nestorius or Coelestius, the synod decrees that they also be deposed.” Dr. Shedd (ii. 191) observes with justice: “The condemnation of Pelagianism which was finally passed by the council of Ephesus, seems to have been owing more to a supposed connection of the views of Pelagius with those of Nestorius, than to a clear and conscientious conviction that his system was contrary to Scripture and the Christian experience.”}

The position of the Greek church upon this question is only negative; she has in name condemned Pelagianism, but has never received the positive doctrines of Augustine. She continued to teach synergistic or Semi-Pelagian views, without, however, entering into a deeper investigation of the relation of human freedom to divine grace.\footnote{Comp. Münscher, Dogmengeschichte, vol. iv. 238, and Neander, Dogmengeschichte, vol. i. p. 412.}

\xsection{The Pelagian System: Primitive State and Freedom of Man; the Fall}

§ 150. The Pelagian System: Primitive State and Freedom of Man; the Fall.

The peculiar anthropological doctrines, which Pelagius clearly apprehended and put in actual practice, which Coelestius dialectically developed, and bishop Julian most acutely defended, stand in close logical connection with each other, although they were not propounded in systematic form. They commend themselves at first sight by their simplicity, clearness, and plausibility, and faithfully express the superficial, self-satisfied morality of the natural man. They proceed from a merely empirical view of human nature, which, instead of going to the source of moral life, stops with its manifestations, and regards every person, and every act of the will, as standing by itself, in no organic connection with a great whole.

We may arrange the several doctrines of this system according to the great stages of the moral history of mankind.

I. The Primitive State of mankind, and the doctrine of Freedom.

The doctrine of the primitive state of man holds a subordinate position in the system of Pelagius, but the doctrine of freedom is central; because in his view the primitive state substantially coincides with the present, while freedom is the characteristic prerogative of man, as a moral being, in all stages of his development.

Adam, he taught, was created by God sinless, and entirely competent to all good, with an immortal spirit and a mortal body. He was endowed with reason and free will. With his reason he was to have dominion over irrational creatures; with his free will he was to serve God. Freedom is the supreme good, the honor and glory of man, the bonum naturae, that cannot be lost. It is the sole basis of the ethical relation of man to God, who would have no unwilling service. It consists according to Pelagius, essentially in the liberum arbitrium, or the possibilitas boni et mali; the freedom of choice, and the absolutely equal ability at every moment to do good or evil.\footnote{De gratia Christi et de pecc. origin. c. 18 (§ 19, tom. x. fol. 238) where Augustinecites the following passage from the treatise of Pelagius, De libero arbitrio: “Habemus possibilitatem utriusque partis a Deo insitam, velut quamdam, ut ita dicam, radicem fructiferam et fecundam, quae ex voluntate hominis diversa gignat et pariat, et quae possit ad proprii cultoris arbitrium, vel nitere flore virtutum, vel sentibus horrere vitiorum.” Against this Augustinecites the declaration of our Lord, Matt. vii. 18, that “a good tree cannot bear evil fruit, nor a corrupt tree good fruit,” that therefore there cannot be “una eademque radix bonortim et malorum.”} The ability to do evil belongs necessarily to freedom, because we cannot will good without at the same time being able to will evil. Without this power of contrary choice, the choice of good itself would lose its freedom, and therefore its moral value. Man is not a free, self-determining moral subject, until good and evil, life and death, have been given into his hand.\footnote{Ep. ad Demet. cap. 3: “In hoc enim gemini itineris ne, in hoc utriusque libertate partis, rationabilis animae decus positum est. Hinc, inquam, totus naturae nostrae honor consistit, hinc dignitas, hinc denique optimi quique laudem merentur, hinc praemium. Nec esset omnino virtus ulla in bono perseverantis, si is ad malum transire non potuisset. Volens namque Deus rationabilem creaturam voluntarii boni munere [al. munire] et liberi arbitrii potestate donare, utriusque partis possibilitatem homini inserendo, proprium ejus fecit esse quod velit, ut boni ac mali capax, naturaliter utrumque posset, et ad alterutrum voluntatem deflecteret. Neque enim aliter spontaneum habere poterat bonum, nisi aeque etiam ea creatura malum habere potusit. Utrumque nos posse voluit optimus Creator, sed unum facere, bonum scilicet, quod et imperavit; malique facultatem ad hoc tantum dedit, ut voluntatem ejus ex nostra voluntate faceremus. Quod ut ita sit, hoc quoque ipsum, quia etiam mala facere possumus, bonum est. Donum, inquam, quia boni partem meliorem facit. Facit enim ipsam voluntariam sui juris, non necessitate devinctam, sed judicio liberam.”}

This is the only conception of freedom which Pelagius has, and to this he and his followers continually revert. He views freedom in its form alone, and in its first stage, and there fixes and leaves it, in perpetual equipoise between good and evil, ready at any moment to turn either way. It is without past or future; absolutely independent of everything without or within; a vacuum, which may make itself a plenum, and then becomes a vacuum again; a perpetual tabula rasa, upon which man can write whatsoever he pleases; a restless choice, which, after every decision, reverts to indecision and oscillation. The human will is, as it were, the eternal Hercules at the cross-road, who takes first a step to the right, then a step to the left, and ever returns to his former position. Pelagius knows only the antithesis of free choice and constraint; no stages of development, no transitions. He isolates the will from its acts, and the acts from each other, and overlooks the organic connection between habit and act. Human liberty, like every other spiritual power, has its development; it must advance beyond its equilibrium, beyond the mere ability to sin or not to sin, and decide for the one or the other. When the will decides, it so far loses its indifference, and the oftener it acts, the more does it become fixed; good or evil becomes its habit, its second nature; and the will either becomes truly free by deciding for virtue, and by practising virtue, or it becomes the slave of vice.\footnote{Pelagius himself, it must be admitted, recognized to some extent the power of habit and its effect upon the will (Ep. ad Demetr. c. 8); but Coelestius and Juliancarried out his idea of the freedom of choice more consistently to the conception of a purely qualitative or formal power which admits of no growth or change by actual exercise, but remains always the same. Comp. Niedner (in the posthumous edition of his Lehrbuch der Kirchengeschichte, Berlin, 1866, p. 345 f.), who justly remarks, in opposition to Baur’s defense of the Pelagian conception of freedom: “Freedom in its first stage, as the power of choice, is a moral (as well as a natural) faculty, and hence capable of development either by way of deterioration into a sinful inclination, or by rising to a higher form of freedom. This is the point which Coelestius and Julianignored: they attached too little weight to the \emph{use} of freedom.”} “Whosoever committeth sin, is the servant of sin.” Goodness is its own reward, and wickedness is its own punishment. Liberty of choice is not a power, but a weakness, or rather a crude energy, waiting to assume some positive form, to reject evil and commit itself to good, and to become a moral self-control, in which the choice of evil, as in Christ, is a moral, though not a physical, impossibility. Its impulse towards exercise is also an impulse towards self-annihilation, or at least towards self-limitation. The right use of the freedom of choice leads to a state of holiness; the abuse of it, to a state of bondage under sin. The state of the will is affected by its acts, and settles towards a permanent character of good or evil. Every act goes to form a moral state or habit; and habit is in turn the parent of new acts. Perfect freedom is one with moral necessity, in which man no longer can do evil because he will not do it, and must do good because he wills to do it; in which the finite will is united with the divine in joyful obedience, and raised above the possibility of apostasy. This is the blessed freedom of the children of God in the state of glory. There is, indeed, a subordinate sphere of natural virtue and civil justice, in which even fallen man retains a certain freedom of choice, and is the artificer of his own character. But as respects his relation to God, he is in a state of alienation from God, and of bondage under sin; and from this he cannot rise by his own strength, by a bare resolution of his will, but only by a regenerating act of grace. received in humility and faith, and setting him free to practise Christian virtue. Then, when born again from above, the will of the new man co-operates with the grace of God, in the growth of the Christian life.\footnote{Comp. the thorough and acute criticism of the Pelagian conception of freedom by Julius Müller, Die christliche Lehre von der Sünde, Bd. ii. p. 49 ff. (3d ed. 1849).}

Physical death Pelagius regarded as a law of nature, which would have prevailed even without sin.\footnote{Coelestius in Marius Mercator. Common. ii. p. 133: “Adam mortalem factum, qui sive peccaret, sive non peccaret, moriturus fuisset.”} The passages of Scripture which represent death as the consequence of sin, he referred to moral corruption or eternal damnation.\footnote{The words of God to Adam, Gen. iii. 19: “Dust thou \emph{art}, and unto dust shalt thou return,” Julianinterpreted not as a curse, but as a consolation, and as an argument for the natural mortality of Adam, by straining the “Dust thou \emph{art}.” See August. Opus imperfectum contra Julian. l. vi. cap. 27 (x. fol. 1346 sqq.).} Yet be conceded that Adam, if he had not sinned, might by a special privilege have been exempted from death.

II. The Fall of Adam and its Consequences.

Pelagius, destitute of all idea of the organic wholeness of the race or of human nature, viewed Adam merely as an isolated individual; he gave him no representative place, and therefore his acts no bearing beyond himself.

In his view, the sin of the first man consisted in a single, isolated act of disobedience to the divine command. Julian compares it to the insignificant offence of a child, which allows itself to be misled by some sensual bait, but afterwards repents its fault. “Rude, inexperienced, thoughtless, having not yet learned to fear, nor seen an example of virtue,”\footnote{“Rudis, imperitus, incautus, sine experimento timoris, sine exemplo justitiae.”} Adam allowed himself to be enticed by the pleasant look of the forbidden fruit, and to be determined by the persuasion of the woman. This single and excusable act of transgression brought no consequences, either to the soul or the body of Adam, still less to his posterity who all stand or fall for themselves.

There is, therefore, according to this system, no original sin, and no hereditary guilt. Pelagius merely conceded, that Adam, by his disobedience, set a bad example, which exerts a more or less injurious influence upon his posterity. In this view he condemned at the synod of Diospolis (415) the assertion of Coelestius, that Adam’s sin injured himself alone, not the human race.\footnote{“Adae peccatum ipsi soli obfuisse, et non generi humano; et infantes qui nascuntur, in eo statu esse, in quo fuit Adam ante praevaricationem.” In Angustine’s De pecc. orig. c. 13 (f. 258).} He was also inclined to admit an increasing corruption of mankind, though he ascribed it solely to the habit of evil, which grows in power the longer it works and the farther it spreads.\footnote{Ep. ad Demet. cap. 8: “Longa consuetudo vitiorum, quae nos infecit a parvo paulatimque per multos corrupit annos, et ita postea obligatos sibi et addictos tenet, ut vim quodammodo videatur habere natura.” He also says of consuetudo, that it “aut vitia aut virtutes alit.”} Sin, however, is not born with man; it is not a product of nature, but of the will.\footnote{Coelestius, Symb. fragm. i.: In remissionem autem peccatorum baptizandos infantes non idcireo diximus, ut \emph{peccatum ex traduce} [or, peccatum naturm, pecca. tum naturale] firmare videamur, quod longe a catholico sensu alienum est; quia peccatum non cum homine nascitur, quod postmodum exercetur ab homine quia non naturm delictum, sed voluntatis ease demonstrator.} Man is born both without virtue and without vice, but with the capacity for either.\footnote{Pelagius, in the first book of the Pro libero arbitrio, cited in Augustine’s De pecc. orig. cap. 13 (§ 14, tom. x. f. 258): “Omne bonum ac malum, quo vel laudabiles vel vituperabiles sumus, non nobiscum \emph{oritur}, sed \emph{agitur} a nobis: capaces enim utriusque rei, non pleni nascimur, et ut sine virtute, ita et sine vitio procreamur; atque ante actionem propriae voluntatis id solum in homine eat, quod Deus condidit.” It is not, however, very congruous with this, that in another place he speaks of a natural or inborn holiness. Ad Demet. c. 4: “Est in animis nostris naturalis quaedam, ut its dixerim, \emph{sanctitas}.”} The universality of sin must be ascribed to the power of evil example and evil custom.

And there are exceptions to it. The “all” in Rom. v. 12 is to be taken relatively for the majority. Even before Christ there were men who lived free from sin, such as righteous Abel, Abraham, Isaac, the Virgin Mary, and many others.\footnote{Comp. Pelagius, Com. in Rom. v. 12, and in August. De natum et gratia, cap. 36 (§ 42, Opera, tom. x. fol. 144): “Deinde commemorat [Pelagius] eos, qui non modo non peccasse, verum etiam juste vixisse referuntur, Abel, Enoch, Melchisedech, Abraham, Isaac, Jacob, Jesu Nove, Phineas, Samuel, Nathan, Elias, Joseph, Elizaeus, Micheas, Daniel, Ananias, Agarias, Meisael, Ezechiel, Mardochaeus, Simeon, Joseph, cui despondata erat virgo Maria, Johannes. Adjungit etiam feminas, Debboram, Annam, Samuelis matrem, Judith, Esther, alteram Annam faliam Phanuel, Elizabeth, ipsam etiam Domini ac Salvatoris nostri matrem, quam dicit sine peccato confiteri necesse ease pietati.”} From the silence of the Scriptures respecting the sins of many righteous men, he inferred that such men were without sin.\footnote{“De illis, quorum justitiae meminit [Scriptura sacra] et peccatorum sine dubio meminisset, si qua eos peccasse sensisset.” In Aug. De Nat. et grat. c. 37 (§ 43; tom. x. fol. 145).} In reference to Mary, Pelagius is nearer the present Roman Catholic view than Augustine, who exempts her only from actual sin, not from original.\footnote{In the passage cited, Augustineagrees with Pelagius in reference to Mary ’propter honorem Domini,’ but only as respects actual sin, of which the connection shows him to be speaking; for in other passages he affirms the conception of Mary in sin. Comp. Enarratio in Psalmum xxxiv. vs. 13 (ed. Migne, tom. iv. 335): “Maria ex Adam mortua \emph{propter peccatum}, Adam mortuus propter peccatum, et caro Domini ex Maria mortua est propter \emph{delenda} peccata.” De Genesi ad literam, lib. x. c. 18 (§ 32), where he discusses the origin of Christ’s soul, and says: “Quid incoinquinatius illo utero Virginia, cujus caro \emph{etiamsi de peccati propagine venit}, non tamen de peccati propagine concepit ...? ” See above, § 80, \textbf{p. 418.}} Jerome, with all his reverence for the blessed Virgin, does not even make this exception but says, without qualification, that every creature is under the power of sin and in need of the mercy of God.\footnote{Adv. Pelag. l. ii. c. 4 (tom. ii. 744, ed. Vallarsi): ”\textgreek{Ἀναμάρτητον}, id est sine peccato esse [hominem posse] nego, id enim soli Deo competit, omnisque creatura peccato subjacet, et indiget misericordia Dei, dicente Scriptura: Misericordia Domini plena est terra.”}

With original sin, of course, hereditary guilt also disappears; and even apart from this connection, Pelagius views it as irreconcilable with the justice of God. From this position a necessary deduction is the salvation of unbaptized infants. Pelagius, however, made a distinction between vita aeterna or a lower degree of salvation, and the regnum coelorum of the baptized saints; and he affirmed the necessity of baptism for entrance into the kingdom of heaven.\footnote{August. De peccatorum meritis et remissione, lib. i. c. 21 (§ 30, tom. x. f. 17); De haeresibus, cap. 88.}

In this doctrine of the fall we meet with the same disintegrating view of humanity as before. Adam is isolated from his posterity; his disobedience is disjoined from other sins. He is simply an individual, like any other man, not the representative of the whole race. There are no creative starting-points; every man begins history anew. In this system Paul’s exhibitions of Adam and Christ as the representative ancestors of mankind have no meaning. If the act of the former has merely an individual significance, so also has that of the latter. If the sin of Adam cannot be imputed, neither can the merit of Christ. In both cases there is nothing left but the idea of example, the influence of which depends solely upon our own free will. But there is an undeniable solidarity between the sin of the first man and that of his posterity.

In like manner sin is here regarded almost exclusively as an isolated act of the will, while yet there is also such a thing as sinfulness; there are sinful states and sinful habits, which are consummated and strengthened by sins of act, and which in turn give birth to other sins of act.

There is a deep truth in the couplet of Schiller, which can easily be divested of its fatalistic intent:

Finally, the essence and root of sin is not sensuality, as Pelagius was inclined to assume (though he did not express himself very definitely on this point), but self-seeking, including pride and sensuality as the two main forms of sin. The sin of Satan was a pride that aimed at equality with God, rebellion against God; and in this the fall of Adam began, and was inwardly consummated before he ate of the forbidden fruit.

\xsection{The Pelagian System Continued: Doctrine, of Human Ability and Divine Grace}

§ 151. The Pelagian System Continued: Doctrine, of Human Ability and Divine Grace.

III. The Present Moral Condition of man is, according to the Pelagian system, in all respects the same as that of Adam before the fall. Every child is born with the same moral powers and capabilities with which the first man was created by God. For the freedom of choice, as we have already seen, is not lost by abuse, and is altogether the same in heathens, Jews, and Christians, except that in Christians it is aided by grace.\footnote{Pelagius, in Aug. De gratia Christi, c. 31 (x. 244): “Liberi arbitrii potestatem dicimus in omnibus esse generaliter, in Christianis, Judaeis atque gentilibus. In omnibus est liberum arbitrium aequaliter per natumam, sed in solis Christianis juvatur gratia.”} Pelagius was a creationist, holding that the body alone is derived from the parents, and that every soul is created directly by God, and is therefore sinless. The sin of the father, inasmuch as it consists in isolated acts of will, and does not inhere in the nature, has no influence upon the child. The only difference is, that, in the first place, Adam’s posterity are born children, and not, like him, created full-grown; and secondly, they have before them the bad example of his disobedience, which tempts them more or less to imitation, and to the influence of which by far the most—but not all—succumb.

Julian often appeals to the virtues of the heathen, such as valor, chastity, and temperance, in proof of the natural goodness of human nature.

He looked at the matter of moral action as such, and judged it accordingly. “If the chastity of the heathen,” he objects to Augustine’s view of the corrupt nature of heathen virtue, “were no chastity, then it might be said with the same propriety that the bodies of unbelievers are no bodies; that the eyes of the heathen could not see; that grain which grew in their fields was no grain.”

Augustine justly ascribed the value of a moral act to the inward disposition or the direction of the will, and judged it from the unity of the whole life and according to the standard of love to God, which is the soul of all true virtue, and is bestowed upon us only through grace. He did not deny altogether the existence of natural virtues, such as moderation, lenity, benevolence, generosity, which proceed from the Creator, and also constitute a certain merit among men; but he drew a broad line of distinction between them and the specific Christian graces, which alone are good in the proper sense of the word, and alone have value before God.

The Holy Scriptures, history, and Christian experience, by no means warrant such a favorable view of the natural moral condition of man as the Pelagian system teaches. On the contrary, they draw a most gloomy picture of fearful corruption and universal inclination to all evil, which can only be overcome by the intervention of divine grace. Yet Augustine also touches an extreme, when, on a false application of the passage of St. Paul: “Whatsoever is not of faith, is sin” (Rom. xiv. 23), he ascribes all the virtues of the heathen to ambition and love of honor, and so stigmatizes them as vices.\footnote{De Civit. Dei, v. 13-20 and xix. 25. In the latter place he calls the virtues, which do not come from true religion, vices. “Virtutes ... nisi ad Deum retulerit, etiam ipsa \emph{vitia sunt potius quam virtutes}.” From this is doubtless derived the sentence so often attributed to Augustine: “The virtues of the heathen are splendid vices,” which, however, in this form and generality, does not, to my knowledge, occur in his writings. More on this point, see below, § 156.} And in fact he is in this inconsistent with himself. For, according to his view, the nature which God created, remains, as to its substance, good; the divine image is not wholly lost, but only defaced; and even man’s sorrow in his loss reveals a remaining trace of good.\footnote{De Genesi ad Lit. viii. 14; ReTract. ii. 24. Comp. Wiggers, i. p. 120 ff.}

Pelagius distinguishes three elements in the idea of good: . Power, will, and act (posse, velle, and esse). The first appertains to man’s nature, the second to his free will, the third to his conduct. The power or ability to do good, the ethical constitution, is grace, and comes therefore from God, as an original endowment of the nature of man. It is the condition of volition and action, though it does not necessarily produce them. Willing and acting belong exclusively to man himself.\footnote{Pelagius, Pro libero arbitrio, cited in Augustine’s De gratia Christi, c. 4 (§ 5, tom. x. fol. 232): ”\emph{Posse} in natura, \emph{velle} in arbitrio, \emph{esse} in effectu locamus. Primum illud, id est \emph{posse}, ad Deum proprie pertinet, qui illud creatrrae suae contulit, duo vero reliqua, hoc est \emph{velle} et \emph{esse}, ad hominem referenda sunt quia de arbitrii fonte descendunt. Ergo in voluntate et opera bono laus hominis est: immo et hominis et Dei, qui ipsius voluntatis et operis possibilitatem dedit, quique ipsam possibilitatem gratiae suae adjuvat semper auxilio.”} The power of speech, of thought, of sight, is God’s, gift; but whether we shall really think, speak, or see, and whether we shall think, speak, or see well or ill, depends upon ourselves.\footnote{“Quod possumus videre oculis, nostrum non est: quod vero bene aut male videmus, hoc nostrum est .... Quod loqui possumus, Dei est: quod vero bene vel male loquimur, nostrum est.” Quoted in Augustine’s De gratia Christi, c. 15 and 16 (fol. 237 and 238). Augustinecites against these examples Ps. cxix. 37: “Averte oculos meos, ne videant vanitatem.”}

Here the nature of man is mechanically sundered from his will and act; and the one is referred exclusively to God, the others to man. Moral ability does not exist over and above the will and its acts, but in them, and is increased by exercise; and thus its growth depends upon man himself. On the other hand, the divine help is indispensable even to the willing and doing of good; for God works in us both to will and to do.\footnote{Phil. ii. 13. Augustineappeals to this passage, De gratia Christi, c. 5 (f. 232 sq.) with great emphasis, as if Paul with prophetic eye had had in view the error of Pelagius.} The Pelagian system is founded unconsciously upon the deistic conception of the world as a clock, made and wound up by God, and then running of itself, and needing at most some subsequent repairs. God, in this system, is not the omnipresent and everywhere working Upholder and Governor of the world, in whom the creation lives and moves and has its being, but a more or less passive spectator of the operation of the universe.\footnote{It is against this deistic view that the pregnant lines of Goethe are directed: “Was wär’ ein Gott der nur von aussen stiesse, Im Kreis das All am Finger laufen liesse;} Jerome therefore fairly accuses the Pelagians (without naming them) of denying the absolute dependence of man on God, and cites against them the declaration of Christ, John v. 17, concerning the uninterrupted activity of God.\footnote{Epistola ad Ctesiphontem. Dr. Neander (Church History, vol. ii. p. 604 ff. Torrey’s transl.) regards this difference of view concerning the relation of the Creator to the creature as the most original and fundamental difference between the Augustinian and Pelagian system, although it did not clearly come to view in the progress of the controversy.}

IV. The doctrine of the Grace of God.

The sufficiency of the natural reason and will of man would seem to make supernatural revelation and grace superfluous. But this Pelagius does not admit. Besides the natural grace, as we may call his concreated ability, he assumes also a supernatural grace, which through revelation enlightens the understanding, and assists man to will and to do what is good.\footnote{Pelagius, in Aug. De gratia Christi, c. 7 (§ 8, x. f. 233): ” ... Deus ... gratiae suae auxilium subministrat, ut quod per liberum homines facere jubentur arbitrium, \emph{facilius possent implere} per gratiam.”} This grace confers the negative benefit of the forgiveness of past sins, or justification, which Pelagius understands in the Protestant sense of declaring righteous, and not (like Augustine) in the Catholic sense of making righteous;\footnote{Pelag. Com. in Rom. iv. 6: “Ad hoc fides prima ad justitiam reputatur, ut de praeterito absolvatur et de paesenti justificatur, et ad futura fidei opera praeparatur.” Similarly Julianof Eclanum. Augustine, on the contrary, has the evangelical conception of faith and of grace, but not of justification, which he interprets subjectively as a progressive making righteous, like the Roman church. Comp. De gratia Christi, c. 47 (§ 52, x. f. 251): ”... gratiam Dei ... in qua nos sua, non nostrae justitiae \emph{justos facit}, ut ea sit vera nostra justitia quae nobis ab illo est.” In another passage, however, he seems to express the Protestant view. De spir. et Lit. c. 26 (§ 45, tom. x. 109): “Certe ita dictum est: justificabuntur, se si diceretur: justi \emph{habebuntur}, justi \emph{deputabuntur}, sicut dictum est de quodam: \emph{Ille autem volens se justificare} (Luc. x. 29), i.e., ut justus \emph{haberetur} et \emph{deputaretur}.”}

and the positive benefit of a strengthening of the will by the power of instruction and example. As we have been followers of Adam in sin, so should we become imitators of Christ in virtue. “In those not Christians,” says Pelagius, “good exists in a condition of nakedness and helplessness; but in Christians it acquires vigor through the assistance of Christ.”\footnote{In Aug. De gratia Chr. c. 31 (tom. x. fol. 244): “In illis nudum et inerme est conditionis bonum; in his vero qui ad Christum pertinent, Christi munitur auxilio.”} He distinguishes different stages of development in grace corresponding to the increasing corruption of mankind. At first, he says, men lived righteous by nature (justitia per naturam), then righteous under the law (justitia sub lege), and finally righteous under grace (justitia gratiae), or the gospel.\footnote{Aug. De pecc. orig. c. 26 (§ 30, tom. x. f. 266): “Non, sicut Pelagius et ejus discipuli, tempora dividamus dicentes\emph{: primum vixisse justos homines ex natura, deinde sub lege, tertio sub gratia}.”} When the inner law, or the conscience, no longer sufficed, the outward or Mosaic law came in; and when this failed, through the overmastering habit of sinning, it had to be assisted by the view and imitation of the virtue of Christ, as set forth in his example.\footnote{Cited from Pelagius, l. c.: “Postquam nimia, sicut disputant, peccandi consuetudo praevaluit cui sanandae lex parum valeret, Christus advenit et tanquam morbo desperatissimo non per discipulos, sed per se ipsum medicus ipse subvenit.”} Julian of Eclanum also makes kinds and degrees of the grace of God. The first gift of grace is our creation out of nothing; the second, our rational soul; the third, the written law; the fourth, the gospel, with all its benefits. In the gift of the Son of God grace is completed.\footnote{In Angustine’s Opus imperf. i. 94 (tom. x. f. 928)}

Grace is therefore a useful external help (adjutorium) to the development of the powers of nature, but is not absolutely necessary. Coelestius laid down the proposition, that grace is not given for single acts.\footnote{· “Gratiam Dei et adjutorium non ad singulos actus dari.”} Pelagius, it is true, condemned those who deny that the grace of God in Christ is necessary for every moment and every act; but this point was a concession wrung from him in the controversy, and does not follow logically from his premises.\footnote{Comp., respecting this, Augustine, De gratia Christi, cap. 2 (tom. x fol. 229 sq.).}

Grace moreover, according to Pelagius, is intended for all men (not, as Augustine taught, for the elect few only), but it must first be deserved. This, however, really destroys its freedom.\footnote{Comp. Rom. iv. 4, 5; Eph. ii. 8, 9. Shakespeare has far better understood the nature of grace than Pelagius, in the famous speech of Portia in the Merchant of Venice (Act IV. Sc, 1): The quality of mercy is not strained: It droppeth as the gentle rain from heaven Upon the place beneath; it is twice blessed, It blesseth him that gives and him that takes.”} “The heathen,” he says, “are liable to judgment and damnation, because they, notwithstanding their free will, by which they are able to attain unto faith and to deserve God’s grace, make an evil use of the freedom bestowed upon them; Christians, on the other hand, are worthy of reward, because they through good use of freedom deserve the grace of God, and keep his commandments.”\footnote{Pelagius in Aug. De gratis Chr. c. 31 (x. f. 245). The \emph{illi}, according to the connection, must refer to those not Christians, the \emph{hi} to Christians. Yet according to his principles we might in turn fairly subdivide each class since according to him there are good heathens and bad Christians. Against this Augustineurges: “Ubi est illud apostoli: Justificati gratis per gratiam ipsius (Rom. iii. 24)? Ubi est illud: Gratis salvi facti estis (Eph. ii. 8)?” He concludes with the just proposition: “Non est gratia, nisi gratuita.”}

Pelagianism, therefore, extends the idea of grace too far, making it include human nature itself and the Mosaic law; while, on the other hand, it unduly restricts the specifically Christian grace to the force of instruction and example. Christ is indeed the Supreme Teacher, and the Perfect Example, but He is also High-priest and King, and the Author of a new spiritual creation. Had He been merely a teacher, He would not have been specifically distinct from Moses and Socrates, and could not have redeemed mankind from the guilt and bondage of sin. Moreover, He does not merely influence believers from without, but lives and works in them through the Holy Ghost, as the principle of their spiritual life. Hence Augustine’s wish for his opponent: “Would that Pelagius might confess that grace which not merely promises us the excellence of future glory, but also brings forth in us the faith and hope of it; a grace, which not merely admonishes to all good, but also from within inclines us thereto; not merely reveals wisdom, but also inspires us with the love of wisdom.”\footnote{De gratia Christi, c. 10 (tom. x. f. 235).} This superficial conception of grace is inevitable, with the Pelagian conception of sin. If human nature is uncorrupted, and the natural will competent to all good, we need no Redeemer to create in us a new will and a new life, but merely an improver and ennobler; and salvation is essentially the work of man. The Pelagian system has really no place for the ideas of redemption, atonement, regeneration, and new creation. It substitutes for them our own moral effort to perfect our natural powers, and the mere addition of the grace of God as a valuable aid and support. It was only by a happy inconsistency, that Pelagius and his adherents traditionally held to the church doctrines of the Trinity and the person of Christ. Logically their system led to a rationalistic Christology.\footnote{Wiggers, l.c. vol. i. p. 467, judges similarly. Also Neander, in his Dogmengeschichte, Bd. i. p. 884: “The Pelagian principles would logically have led to rationalistic views, to an entire rejection of the supernatural element, and to the belief that mankind needs only to develop itself from within itself, without the revelation and self-impartation of God, in order to attain the good. But they do not develop their first principles so consistently as this, and what Biblical elements they incorporate with their system are unquestionably not taken in merely by way of accommodation, but through the persuasion that a supernatural revelation is necessary, in order to realize the destiny of mankind.” Comp. Cunningham, Hist. Theology, i. p. 829: “Modern Socinians and Rationalists are the only consistent Pelagians. When men reject what Pelagius rejected, they are bound in consistency to reject everything that is peculiar and distinctive in the Christian system as a remedial scheme.”}

Pelagianism is a fundamental anthropological heresy, denying man’s need of redemption, and answering to the Ebionistic Christology, which rejects the divinity of Christ. It is the opposite of Manichaeism, which denies man’s capability of redemption, and which corresponds to the Gnostic denial of the true humanity of Christ.\footnote{Comp. Augustine, Contra duas Epist. Pelagianorum l. ii. c. 2, where he describes Manichaeism and Pelagianism at length as the two opposite extremes, and opposes to them the Catholic doctrine.}

\xsection{The Augustinian System: The Primitive State of Man, and Free Will}

§ 152. The Augustinian System: The Primitive State of Man, and Free Will.

Augustine (354–430) had already in his Confessions, in the year 400, ten years before the commencement of the Pelagian controversy, set forth his, deep and rich experiences of human sin and divine grace. This classical autobiography, which every theological student should read, is of universal application, and in it every Christian may bewail his own wanderings, despair of himself, throw himself unconditionally into the arms of God, and lay hold upon unmerited grace.\footnote{An ingenious but somewhat far-fetched parallel is drawn by Dr. Kleinert between Augustineand Faust, as two antipodal representatives of mankind, in a brochure: Augustin und Goethe’s Faust, Berlin, 1866. A more obvious comparison is that of the Confessions of Augustinewith the Confessions of Rousseau, and with Goethe’s Wahrheit und Dichtung.} Augustine had in his own life passed through all the earlier stages of the history of the church, and had overcome in theory and in practice the heresy of Manichaeism, before its opposite, Pelagianism, appeared. By his theological refutation of this latter heresy, and by his clear development of the Biblical anthropology, he has won the noblest and most lasting renown. As in the events recorded in his Confessions he gives views of the evangelical doctrines of sin and of grace, so in the doctrines of his anti-Pelagian writings he sets forth his personal experience. He teaches nothing which he has not felt. In him the philosopher and the living Christian are everywhere fused. His loftiest metaphysical speculation passes unconsciously into adoration. The living aroma of personal experience imparts to his views a double interest, and an irresistible attraction for all earnest minds.\footnote{Dr. Baur, in his posthumous Vorlesungen über the Dogmengeschichte, published by his son (1866, Bd. i. p. ii. p. 26), makes the fine remark respecting him: “With Augustinehimself everything lies in the individuality of his nature, as it was shaped by the course of his life, by his experiences and circumstances.” He should have added, however, that in so magnificent a personality as Augustine’s, that which is most individual is also the most universal, and the most subjective is the most objective.}

Yet his system was not always precisely the same; it became perfect only through personal conflict and practical tests. Many of his earlier views—e.g., respecting the freedom of choice, and respecting faith as a work of man—he himself abandoned in his Retractations;\footnote{Retract. l. i. c. 9.} and hence he is by no means to be taken as an infallible guide. He holds, moreover, the evangelical doctrines of sin and grace not in the Protestant sense, but, like his faithful disciples, the Jansenists, in connection with the sacramental and strict churchly system of Catholicism; he taught the necessity of baptismal regeneration and the damnation of all unbaptized children, and identified justification in substance with sanctification, though he made sanctification throughout a work of free grace, and not of human merit. It remains the exclusive prerogative of the inspired apostles to stand above the circumstances of their time, and never, in combating one error, to fall into its opposite. Nevertheless, Augustine is the brightest star in the constellation of the church fathers, and diffuses his light through the darkest periods of the middle ages, and among Catholics and Protestants alike, even to this day.\footnote{Baur, l.c. p. 32 f.: “From the time that Augustinedirected the development of the Christian system to the two doctrines of sin and grace, this tendency always remained in the Occidental dogmatics the prevailing one, and so great and increasingly predominant in the course of time did the authority of Augustinebecome in the church, that even those who had departed from his genuine teachings, which many were unwilling to follow out with rigid consistency, yet believed themselves bound to appeal to his authority, which his writings easily gave them opportunity to do, since his system, as the result of periods of development so various, and antitheses so manifold, offers very different sides, from which it can be interpreted.”}

His anthropology may be exhibited under the three stages of the religious development of mankind, the status integritatis, the status corruptionis, and the status redemtionis.

I. The Primitive State of man, or the State of Innocence.

Augustine’s conception of paradise is vastly higher than the Pelagian, and involves a far deeper fall and a far more glorious manifestation of redeeming grace. The first state of man resembles the state of the blessed in heaven, though it differs from that final state as the undeveloped germ from the perfect fruit. According to Augustine man came from the hand of his Maker, his genuine masterpiece, without the slightest fault. He possessed freedom, to do good; reason, to know God; and the grace of God. But by this grace Augustine (not happy in the choice of his term) means only the general supernatural assistance indispensable to a creature, that he may persevere in good.\footnote{Grace, in this wider sense, as source of all good, Augustinemakes independent of sin, and ascribes the possession of it even to the good angels. Comp. De corrupt. et grat. § 32 (tom. x. 767, 768): “Dederat [Deus homini] adjutorium sine quo in ea [bona voluntate] non posset permanere si vellet; ut autem vellet, in ejus libero reliquit arbitrio. Posset ergo permanere si vellet: quia non deerat adjutorium per quod posset et sine quo non posset perseveranter bonum tenere quod vellet .... Si autem hoc adjutorium vel \emph{angelo} vel homini, cum primum facti sunt, defuisset, quoniam non talis natura facta erat, ut sine divino adjutorio posset manere si vellet, non utique sua culpa cecidissent: adjutorium quippe defuisset, sine quo manere non possent.” We see here plainly the germ of the scholastic and Roman Catholic doctrine of the justitia originalis, which was ascribed to the first man as a special endowment of divine grace or a supernatural accident, on the ground of the familiar distinction between the imago Dei (which belongs to the \emph{essence} of man and consists in reason and free will) and the similitude Dei (the actual conformity to the divine will).} The relation of man to God was that of joyful and perfect obedience. The relation of the body to the soul was the same. The flesh did not yet lust against the spirit; both were in perfect harmony, and the flesh was wholly subject to the spirit. “Tempted and assailed by no strife of himself against himself, Adam enjoyed in that place the blessedness of peace with himself.” To this inward state, the outward corresponded. The paradise was not only spiritual, but also visible and material, without heat or cold, without weariness or excitement, without sickness, pains, or defects of any kind. The Augustinian, like the old Protestant, delineations, of the perfection of Adam and the blissfulness of paradise often exceed the sober standard of Holy Scripture, and borrow their colors in part from the heavenly paradise of the future, which can never be lost.\footnote{Comp. several passages in the Opus imperf. i. 71; iii. 147; vi. 9, 17; Contra Jul. v. 5; De civitate Dei, xiii. 1, 13, 14, 21; xiv. 10, where he depicts the beatitudo and deliciae of Eden in poetic colors, and extends the perfection even to the animal and vegetable realms. Yet he is not everywhere consistent. His views became more exaggerated from his opposition to Pelagianism. In the treatise, De libero arbitrio, iii. c. 24, §§ 71, 72, which he completed a.d.395, he says, that the first human beings were neither wise nor foolish, but had at first only the capability to become one or the other. “Infans nec stultus nec sapiens dici potest, quamvis jam homo sit; ex quo apparet natumm hominis recipere aliquid medium, quod neque stultitiam neque sapientiam recte vocaris.” ...“Ita factus est homo, ut quamvis sapiens nondum esset, praeceptum tamen posset accipere.” On the other hand, in his much later Opus imperf. c. Julianum, l. v. c. 1 (tom. x. f. 1222) he ascribes to the first men excellentissima sapientia, appealing to Pythagoras, who is said to have declared him the wisest who first gave names to things.}

Yet Augustine admits that the original state of man was only relatively perfect, perfect in its kind; as a child may be a perfect child, while he is destined to become a man; or as the seed fulfils its idea as seed, though it has yet to become a tree. God alone is immutable and absolutely good; man is subject to development in time, and therefore to change. The primal gifts were bestowed on man simply as powers, to be developed in either one of two ways. Adam could go straight forward, develop himself harmoniously in untroubled unity with God, and thus gradually attain his final perfection; or he could fall away, engender evil ex nihilo by abuse of his free will, and develop himself through discords and contradictions. It was graciously made possible that his mind should become incapable of error, his will, of sin, his body, of death; and by a normal growth this possibility would have become actual. But this was mere possibility, involving, in the nature of the case, the opposite possibility of error, sin, and death.

Augustine makes the important distinction between the possibility of not sinning\footnote{Posse \emph{non peccare}, which at the same time implies the possibilitas peccandi Comp. Opus imperf. v. 60 (fol. 1278): “Prorsus ita factus est, ut peccandi possibilitatem haberet a necessario, peccatum vero a possibili,” i.e., the \emph{possibility} of sinning was necessary, but the sinning itself merely possible. The peccare posse, says Augustine, in the same connection, is natura, the peccare is culpa.} and the impossibility of sinning.\footnote{\emph{Non posse} peccare, or impossibilitas peccandi.} The former is conditional or potential freedom from sin, which may turn into its opposite, the bondage of sin. This belonged to man before the fall. The latter is the absolute freedom from sin or the perfected holiness, which belongs to God, to the holy angels who have acceptably passed their probation, and to the redeemed saints in heaven.

In like manner he distinguishes between absolute and relative immortality.\footnote{Between the non posse mori and the posse non mori, or between the immortalitas major and the immortalitas minor.} The former is the impossibility of dying, founded upon the impossibility of sinning; an attribute of God and of the saints after the resurrection. The latter is the bare pre-conformation for immortality, and implies the opposite possibility of death. This was the immortality of Adam before the fall, and if he had persevered, it would have passed into the impossibility of dying; but it was lost by sin.\footnote{Comp. Opus imperf. l. vi. cap. 30 (tom. x. fol 1360): “Illa vero immortalitas in qua sancti angeli vivunt, et in qua nos quoque victuri sumus, procul dubio major est. Non enim talis, in qua homo habeat quidem in potestate non mori, sicut non peccare, sed etiam possit et mori, quia potest peccare: sed talis est illa immortalitas, in qua omnis qui ibi est, vel erit, mori non poterit, quia nec peccare jam poterit.” De corrept. et grat. § 33 (x. f. 168): “Prima libertas voluntatis erat, posse non peccare, novissima erit multo major, non posse peccare: Prima immortalitas erat, posse non mori, novissima erit multo major, non posse mori: prima erat perseverantiae potestas, bonum posse non deserere; novissima erit felicitas perseverantiae, bonum non posse deserere.”}

Freedom, also, Augustine holds to be an original endowment of man; but he distinguishes different kinds of it, and different degrees of its development, which we must observe, or we should charge him with self-contradiction.\footnote{The distinctions in the Augustinian idea of freedom have been overlooked by Wiggers and most of the old historians, but, on the other hand, brought out with more or less clearness by Neander (in the Kirchengeschichte and in the Dogmengeschichte), by Ritter (Gesch. der christl. Philosophic, ii. p. 341 ff.), Jul. Müller (Die christl. Lehre von der Sünde, ii. 45 ff.), Joh. Huber (Philosophic der Kirchenväter, p. 296 ff.). Baur bases his acute criticism of the Augustinian system in part upon the false assumption that Augustine’s view of the liberum arbitrium was precisely the same as that of Pelagius. See below.}

By freedom Augustine understands, in the first place, simply spontaneity or self-activity, as opposed to action under external constraint or from animal instinct. Both sin and holiness are voluntary, that is, acts of the will, not motions of natural necessity.\footnote{Retract. i. c. 9, § 4: “Voluntas est qua et peccatur, et recte vivitur.”} This freedom belongs at all times and essentially to the human will, even in the sinful state (in which the will is, strictly speaking, self-willed); it is the necessary condition of guilt and punishment, of merit and reward. In this view no thinking man can deny freedom, without destroying the responsibility and the moral nature of man. An involuntary, will is as bald a self-contradiction as an unintelligent intelligence.\footnote{Here belong especially the first chapters of the treatises, De gratia et libero arbitrio (tom. x. fol. 717-721), of the Opus imperf. contra Julianum, and Contra duas epistolas Pelagianorum. In this sense even the strictest adherents of the Augustinian and Calvinistic system have always more or less explicitly conceded human freedom. Thus Cunningham, a Calvinist of the Free Church of Scotland, in his presentation of the Pelagian controversy (Hist. Theol. i. p. 325): ”Augustinecertainly did not deny man’s free will altogether, and in every sense of the word; and the most zealous defenders of the doctrines of grace and of Calvinistic principles have admitted that there is a free will or free agency, in some sense, which man has, and which is necessary to his being responsible for his transgressions of God’s law. It is laid down in our own [the Westminster) Confession, that ’God hath endued the will of man with that natural liberty, that it is neither forced, nor by any absolute necessity of nature determined to good or evil.’ ” Dr. Shedd, an American Presbyterian of the Old School, in his History of Christian Doctrine, ii. p. 66, where he, in Augustine’s view, expresses his own, says: “The guilt of sin consists in its unforced wilfulness; and this guilt is not in the least diminished by the fact that the will cannot overcome its own wilfulness. For this wicked wilfulness was not created in the will, but is the product of the will’s act of apostasy. The present impotence to holiness is not an original and primitive impotence. By creation Adam had plenary power, not indeed to \emph{originate} holiness, for no creature has this, but to \emph{preserve} and \emph{perpetuate} it. The present destitution of holiness, and impossibility of originating it, is due therefore to the creature’s apostatizing agency, and is a part of his condemnation.” Also, p. 80: “There is no author in the whole theological catalogue, who is more careful and earnest than Augustine, to assert that sin is \emph{self}-activity, and that its source is in the voluntary nature of man. Sin, according to him, is not a substance, but an agency; it is not the essence of any faculty in man, but only the action of a faculty.” Neither Dr. Cunningham nor Dr. Shedd, however, takes any account of the different forms and degrees of freedom in the Augustinian system.}

A second form of freedom is the liberum arbitrium, or freedom of choice. Here Augustine goes half-way with Pelagius; especially in his earlier writings, in opposition to Manichaeism, which denied all freedom, and made evil a natural necessity and an original substance. Like Pelagius he ascribes freedom of choice to the first man before the fall. God created man with the double capacity of sinning or not sinning, forbidding the former and commanding the latter. But Augustine differs from Pelagius in viewing Adam not as poised in entire indifference between good and evil, obedience and disobedience but as having a positive constitutional tendency to the good, yet involving, at the same time, a possibility of sinning.\footnote{This important distinction is overlooked by Baur, in his Kirchengeschichte vom 4-6ten Jahrhundert, p. 143. It takes off the edge from his sharp criticism of the Augustinian system, in which he charges it with inconsistency in starting from the same idea of freedom as Pelagius and yet opposing it.} Besides, Augustine, in the interest of grace and of true freedom, disparages the freedom of choice, and limits it to the beginning, the transient state of probation. This relative indecision cannot be at all predicated of God or the angels, of the saints or of sinners. It is an imperfection of the will, which the actual choosing of the good or the evil more or less surmounts. Adam, with the help of divine grace, without which be might will the good, indeed, but could not persevere in it, should have raised himself to the true freedom, the moral necessity of good; but by choosing the evil, he fell into the bondage of sin.\footnote{Comp. respecting this conception of freedom, the treatise, De libero arbitrio (in Opera, tom. i. f. 569 sqq.), which was begun a.d.388, and finished a.d.395, and belongs therefore to his earliest writings; also, De correptione et gratia (especially cap. 9-11), and the sixth book of the Opus imperf. c. Julianum. Also Contra duas epistolas Pelag. l. ii. c. 2 (tom. x. f. 432), where he opposes both the Manichaean denial of the liberum arbitrium and the Pelagian assertion of its continuance after the fall. “Manichaei negant, homini bono ex libero arbitrio fuisse initium mali; Pelagiani dicunt, etiam hominem malum sufficienter habere liberum arbitrium ad faciendum praeceptum bonum; catholica [fides] utrosque redarguit, et illis dicens: Fecit Deus hominem rectum, et istis dicens: Si vos Filius liberaverit, vere liberi eritis.”} Augustine, however, incidentally concedes, that the liberum arbitrium still so far exists even in fallen man, that he can choose, not indeed between sin and holiness, but between individual actions within the sphere of sinfulness and of justitia civilis.\footnote{Contra duas Epist. Pelag. ii. c. 5 (or § 9, tom. x. f. 436): “Peccato Adae arbitrium liberum de hominum natura periisse non dicimus, sed ad peccandum valere in hominibus subditis diabolo, ad bene autem pieque vivendum non valere, nisi ipsa voluntas hominis Dei gratia fuerit liberata, et ad omne bonum actionis, sermonis, cogitationis adjuta.” Also, De gratia et libero arbitrio, c. 15 (x. f. 184): “Semper est autem in nobis voluntas libera, sed non semper est bona. Aut enim a justitia libera est, quando servit peccato, et tunc est mala; aut a peceato libera est, quando servit justitiae, et tunc est bona. Gratia vero Dei semper est bona.” Dr. Baur, it is true (Die christl Kirche vom Anfang des 4ten bis Ende des 6ten Jahrhunderts, p. 140), is not wholly wrong when he, with reference to this passage, charges Augustinewith an equivocal play upon words, in retaining the term freedom, but changing its sense into its direct opposite. “Meaningless as it is,” says Baur, “to talk in this equivocal sense of freedom, we however see even from this what interest the idea of freedom still had for him, even after he had sacrificed it to the determinism of his system.” The Lutheran theolgians likewise restricted the liberum arbitrium of fallen man to the justitia civilis, in distinction from the justitia Dei, or spiritualis. Comp. Melanchthon, in the Confessio Augustana, art. xviii. The Formula Concordiae goes even beyond Augustine, and compares the natural man in spiritualibus et divinis rebus with a “status salis,” “truncus,” and “lapis” nay, makes him out yet worse off, inasmuch as he is not merely passive, but “voluntati divinae rebellis est et inimicus ” (pp. 661 and 662).}

Finally, Augustine speaks most frequently and most fondly of the highest freedom, the free self-decision or self-determination of the will towards the good and holy, the blessed freedom of the children of God; which still includes, it is true, in this earthly life, the possibility of sinning, but becomes in heaven the image of the divine freedom, a felix necessitas boni, and cannot, because it will not, sin.\footnote{De corrept. et gratia, § 32 (x. 768): “Quid erit liberius libero arbitrio, quando non poterit servire peccato? ... § 33: Prima libertas voluntatis erat, posse non peccare, novissima erit multo major, non posse peccare.”} it is the exact opposite of the dura necessitas mali in the state of sin. It is not a faculty possessed in common by all rational minds, but the highest stage of moral development, confined to true Christians. This freedom Augustine finds expressed in that word of our Lord: “If the Son shall make you free, ye shall be free indeed.” It does not dispense with grace, but is generated by it; the more grace, the more freedom. The will is free in proportion as it is healthy, and healthy in proportion as it moves in the element of its true life, in God, and obeys Him of its own spontaneous impulse. To serve God is the true freedom.\footnote{“Deo servire vera libertas est;” a profound and noble saying. This higher conception of freedom Augustinehad substantially expressed long before the Pelagian controversy, e.g., in the Confessions. Comp. also De civit. Dei l. xiv. c. 11: “Arbitriam igitur voluntatis tunc est vere liberum, quum vitiis peccatisque non servit. Tale datum est a Deo: quod amissum proprio vitio, nisi a quo dari potuit, reddi non potest. Unde veritas dicit: \emph{Si vos filisliberaverit, tunc vere liberi eritis}, Id ipsum est autem, ac si diceret: Si vos Filius salvos fecerit, tunc vere salvi eritis. inde quippe liberatur, unde salvatur.”}

\xsection{The Augustinian System: The Fall and its Consequences}

§ 153. The Augustinian System: The Fall and its Consequences.

To understand Augustine’s doctrine of the fall of man, we must remember, first of all, that he starts with the idea of the organic unity of the human race, and with the profound parallel of Paul between the first and the second Adam;\footnote{Rom. v. 12 ff.; 1 Cor. xv. 22.} that he views the first man not merely as an individual, but at the same time as the progenitor and representative of the whole race, standing to natural mankind in the same relation as that of Christ to redeemed and regenerate mankind. The history of the fall, recorded in a manner at once profound and childlike in the third chapter of Genesis, has, therefore, universal significance. In Adam human nature fell, and therefore all, who have inherited that nature from him, who were in him as the fruit in the germ, and who have grown up, as it were, one person with him.\footnote{De civit. Dei, l. xiii. c. 14: “Omnes enim fuimus in illo uno, quando omnes fuimus ille unus, qui per feminam lapsus est in peccatum, quae de illo facta est ante peccatum.” Compare other passages below.}

But Augustine did not stop with the very just idea of an organic connection of the human race, and of the sin of Adam with original sin; he also supposed a sort of pre-existence of all the posterity of Adam in himself, so that they actually and personally sinned in him, though not, indeed, with individual consciousness. Since we were, at the time of the fall, “in lumbis Adami,” the sin of Adam is “jure seminationis et germinationis,” our sin and guilt, and physical death is a penalty even upon infant children, as it was a penalty upon Adam. The posterity of Adam therefore suffer punishment not for the sin of another, but for the sin which they themselves committed in Adam. This view, as we shall see farther on, Augustine founds upon a false interpretation of Rom. v. 12.

I. The Fall. The original state of man included the possibility of sinning, and this was the imperfection of that state. This possibility became reality. Why it should have been realized, is incomprehensible; since evil never has, like good, a sufficient reason. It is irrationality itself. Augustine fixes an immense gulf between the primitive state and the state of sin. But when thought has accomplished this adventurous leap, it finds his system coherent throughout.

Adam did not fall without temptation from another. That angel, who, in his pride, had turned away from God to himself, tempted man, who, standing yet in his integrity, provoked his envy. He first approached the woman, the weaker and the more credulous. The essence of the sin of Adam consisted not in the eating of the fruit; for this was in itself neither wrong nor harmful; but in disobedience to the command of God. “Obedience was enjoined by that commandment, as the virtue which, in the rational creature, is, as it were, the mother and guardian of all virtues.” The principle, the root of sin, was pride, self-seeking, the craving of the will to forsake its author, and become its own. This pride preceded the outward act. Our first parents were sinful in heart, before they had yet fallen into open disobedience. “For man never yet proceeded to an evil work, unless incited to it by an evil will.” This pride even preceded the temptation of the serpent. “If man had not previously begun to take pleasure in himself, the serpent could have had no hold upon him.”

The fall of Adam appears the greater, and the more worthy of punishment, if we consider, first, the height he occupied, the divine image in which he was created; then, the simplicity of the commandment, and ease of obeying it, in the abundance of all manner of fruits in paradise; and finally, the sanction of the most terrible punishment from his Creator and greatest Benefactor.

Thus Augustine goes behind the appearance to the substance; below the surface to the deeper truth. He does not stop with the outward act, but looks chiefly at the disposition which lies at its root.

II. The Consequences of the primal sin, both for Adam and for his posterity, are, in Augustine’s view, comprehensive and terrible in proportion to the heinousness of the sin itself. And all these consequences are at the same time punishments from the righteous God, who has, by one and the same law, joined reward with obedience and penalty with sin. They are all comprehended under death, in its widest sense; as Paul says: “The wages of sin is death;” and in Gen. ii. 17 we are to understand by the threatened death, all evil both to body and to soul.

Augustine particularizes the consequences of sin under seven heads; the first four being negative, the others positive:

1. Loss of the freedom of choice,\footnote{Of course not in indifferent things of ordinary life, in which the greatest sinner is free to choose, but in reference to the great religous decision for or against God and divine things.} which consisted in a positive inclination and love to the good, with the implied possibility of sin. In place of this freedom has come the hard necessity of sinning, bondage to evil. “The will, which, aided by grace, would have become a source of good, became to Adam, in his apostasy from God, a source of evil.”

2. Obstruction of knowledge. Man was originally able to learn everything easily, without labor, and to understand everything aright. But now the mind is beclouded, and knowledge can be acquired and imparted only in the sweat of the face.

3. Loss of the grace of God, which enabled man to perform the good which his freedom willed, and to persevere therein. By not willing, man forfeited his ability, and now, though he would do good, he cannot.

4. Loss of paradise. The earth now lies under the curse of God: it brings forth thorns and thistles, and in the sweat of his face man must eat his bread.

5. Concupiscence, i.e., not sensuousness in itself, but the preponderance of the sensuous, the lusting of the flesh against the spirit. Thus God punishes sin with sin—a proposition which Julian considered blasphemy. Originally the body was as joyfully obedient to the spirit, as man to God. There was but one will in exercise. By the fall this beautiful harmony has been broken, and that antagonism has arisen which Paul describes in the seventh chapter of the Epistle to the Romans. (Augustine referred this passage to the regenerate state.) The rebellion of the spirit against God involved, as its natural punishment, the rebellion of the flesh against the spirit. Concupiscentia, therefore, is substantially the same as what Paul calls in the bad sense “flesh.” It is not the sensual constitution in itself, but its predominance over the higher, rational nature of man.\footnote{Not the “sentiendi vivacitas,” but the “ribido sentiendi, quae nos ad sentiendum, sive consentientes mente, sive repugnantes, appetitu voluptatis impellit.” C. Julianum, l. iv. c. 14 (§ 65, tom. x. f. 615). He illustrates the difference by a reference to Matt. v. 28. “Non ait Dominus: qui viderit mulierem, sed: qui\emph{viderit ad conpiscendum, jam maechatus est eam in corde suo}. ... Illud [videre] Deus condidit, instruendo corpus humanum; illud [videre ad concupiscendum] diabolus seminavit, persuadendo peccatum.”} It is true, however, that Augustine, in his longing after an unimpeded life in the spirit, was inclined to treat even lawful appetites, such as hunger and thirst, so far as they assume the form of craving desire, as at least remotely connected with the fall.\footnote{“Quis autem mente sobrius non mallet, si fieri posset, sine ulla mordaci voluptate carnali vel arida sumere alimenta, vel humida, sicut suminus haec aëria, quae de circumfusis auris respirando et spirando sorbemus et fundimus?” Contra Jul. iv. c. 14, § 68, f. 616.} Julian attributed the strength of animal desire to the animal element in the original nature of man. Augustine answered, that the superiority of man to the brute consists in the complete dominion of reason over the sensual nature, and that therefore his approach to the brute in this respect is a punishment from God. Concupiscence then is no more a merely corporeal thing than the biblical σάρξ, but has its seat in the soul, without which no lust arises. We must, therefore, suppose a conflict in the soul itself, a lower, earthly, self-seeking instinct, and a higher, god-like impulse.

This is the generic sense of concupiscentia: the struggle of the collective sensual and psychical desires against the god-like spirit. But Augustine frequently employs the word, as other corresponding terms are used, in the narrower sense of unlawful sexual desire. This appeared immediately after the fall, in the shame of our first parents, which was not for their nakedness itself, since this was nothing new to them, but for the lusting of the body; for something, therefore, in and of itself good (the body’s, own enjoyment, as it were), but now unlawfully rising, through the discord between body and soul. But would there then have been propagation without the fall? Unquestionably; but it would have left the dominion of reason over the sensual desire undisturbed. Propagation would have been the act of a pure will and chaste love, and would have had no more shame about it than the scattering of seed upon the maternal bosom of the earth. But now lust rules the spirit; and Augustine in his earlier years had had bitter experience of its tyranny. To this element of sin in the act of procreation he ascribes the pains of childbirth, which in fact appear in Genesis as a consequence of the fall, and as a curse from God. Had man remained pure, “the ripe fruit would have descended from the maternal womb without labor or pain of the woman, as the fruit descends from the tree.”\footnote{De civitate Dei, xiv. 26.}

6.Physical death, with its retinue of diseases and bodily pains. Adam was indeed created mortal, that is, capable of death, but not subject to death. By a natural development the possibility of dying would have been overcome by the power of immortality; the body would have been gradually spiritualized and clothed with glory, without a violent transition or even the weakness of old age. But now man is fallen under the bitter necessity of death. Because the spirit forsook God willingly, it must now forsake the body unwillingly. With profound discernment Augustine shows that not only the actual severance of soul and body, but the whole life of sinful man is a continual dying. Even with the pains of birth and the first cry of the child does death begin. The threatening of the Lord, therefore: “In the day ye eat thereof, ye shall die,” began at once to be fulfilled. For though our first parents lived many years afterwards, they immediately began to grow old and to die. Life is an unceasing march towards death, and “to no one is it granted, even for a little, to stand still, or to go more slowly, but all are constrained to go with equal pace, and no one is impelled differently from others. For he whose life has been shorter, saw therefore no shorter day than he whose life was longer. And he who uses more time to reach death, does not therefore go slower, but only makes a longer journey.”

7. The most important consequence of the fall of Adam is original sin and hereditary guilt in his whole posterity; and as this was also one of the chief points of controversy, it must be exhibited at length.

\xsection{The Augustinian System: Original Sin, and the Origin of the Human Soul}

§ 154. The Augustinian System: Original Sin, and the Origin of the Human Soul.

Original sin,\footnote{Peccatum originals, vitium hereditarium.} according to Augustine, is the native bent of the soul towards evil, with which all the posterity of Adam—excepting Christ, who was conceived by the Holy Ghost and born of a pure Virgin—come into the world, and out of which all actual sins of necessity proceed. It appears principally in concupiscence, or the war of the flesh against the spirit. Sin is not merely an individual act, but also a condition, a status and habitus, which continues, by procreation, from generation to generation. Original sin results necessarily, as has been already remarked, from the generic and representative character of Adam, in whom human nature itself, and so, potentially, all who should inherit that nature, fell.\footnote{De peccatorum meritis et remissione, l. iii. c. 7 (§ 14, tom. x. f. 78): “In Adam omnes tunc peccaverunt, quando in ejus natura illa insita vi, qua eos gignere poterat, adhuc omnes ille unus fuerunt.” De corrept. et gratia, § 28 (x. f 765): “Quia vero [Adam] per liberum arbitrium Deum deseruit, justum judicium Dei expertus est, ut cum tota sua stirpe, quae in illo adhuc posita tota cum illo peccaverat, damnaretur.” This view easily fell in with Augustine’s Platonico-Aristotelian realism, which regarded the general conceptions as the original types of individual things. But the root of it lay deeper in his Christian consciousness and profound conviction of the all-pervading power of sin.} The corruption of the root communicates itself to the trunk and the branches. But where sin is, there is always guilt and ill-desert in the eyes of a righteous God. The whole race, through the fall of its progenitor, has become a massa perditionis. This, of course, still admits different degrees both of sinfulness and of guilt.

Original sin and guilt are propagated by natural generation. The generic character planted in Adam unfolds itself in a succession of individuals, who organically grow one out of another. As sin, however, is not merely a thing of the body, but primarily and essentially of the spirit, the question arises, on which of the current theories as to the origin and propagagation of souls Augustine based his view.

This metaphysical problem enters theology in connection with the doctrine of original sin; this, therefore, is the place to say what is needful upon it.\footnote{“La première difficulté est,” says Leibnitz in the Theodicée, Partie i. 86, comment l’âme a pu être infectée du péché originel, qui est la racine des péchés actuels, sans qu’il y sit en de l’injustice en Dieu à l’y exposer.”} The Gnostic and pantheistic emanation-theory had long since been universally rejected as heretical. But three other views had found advocates in the church:

1. The Traducian\footnote{From tradux, propagator. The author of this theory is Tertullian, De anima, c. 27 (Opera, ed. Fr. Oehler, tom. ii. p. 599 sqq.): “Immo simul ambas [animam et corpus] et concipi et confici et perfici dicimus, sicut et promi, nec ullum intervenire momentum in concepta quo locus ordinetur. ... Igitur ex uno homine tota haec animarum redundantia.” Cap. 86 (p. 617): “Anima in utero seminata pariter cum came pariter cum ipsa sortitur.” Comp. c. 19 (anima velut surculus quidam ex matrice Adam in propaginem deducta); De resurr. carnis, c. 45; Adv. Valentin. c. 26 (tradux animae). With Tertullianthis theory was connected with a materializing view of the soul.} or Generation-theory teaches that the soul originates with the body from the act of procreation, and therefore through human agency. It is countenanced by several passages of Scripture, such as Gen. v. 3; Ps. li. 5; Rom. v. 12; 1 Cor. xv. 22; Eph. ii. 3; it is decidedly suitable to the doctrine of original sin; and hence, since Tertullian, it has been adopted by most Western theologians in support and explanation of that doctrine.\footnote{Jeromesays of the maxima pars occidentalium, that they teach: “Ut quomodo corpus ex corpore, sic anima nascatur ex anima, et simili cum brutis animalibus conditione subsistat.” Ep. 78 ad Marcell. Leo the Great declared it even to be catholica fides, that every man “in corporis et animae subetantiam fomari intra materna viscera.” Ep. 15 ad Turrib. Similarly among the Oriental fathers, Theodoret, Fab. haer. v. 9: \textgreek{ἡ ἐκκλησία τοῖς θείοις πείθομένη λόγοις,} —\textgreek{λέγει τὴν ψυχὴν συνδημιουργεῖσθαιτῷ σώματι}.}

2. The Creation-theory ascribes each individual soul to a direct creative act of God, and supposes it to be united with the body either at the moment of its generation, or afterwards. This view is held by several Eastern theologians and by Jerome, who appeals to the unceasing creative activity of God (John v. 17). It required the assumption that the Soul, which must proceed pure from the hand of the Creator, becomes sinful by its connection with the naturally generated body. Pelagius and his followers were creationists.\footnote{Jeromesays, appealing to John v. 17; Zech. xii. 1; Ps. xxxiii. 15: “Quotidie Deus fabricatur animas, cujus velle fecisse est, et conditor esse non cessat.” Pelagius, in his Confession of Faith, declares for the view that souls are made and given by God Himself.}

3. The theory of Pre-existence, which was originated by Plato and more fully developed by Origen, supposes that the soul, even before the origin of the body, existed and sinned in another world, and has been banished in the body as in a prison,\footnote{The \textgreek{σῶμα}interpreted as \textgreek{σῆμα}(sepulchre). Origen appeals to the groaning of the creation, Rom. viii. 19.} to expiate that personal Adamic guilt, and by an ascetic process to be restored to its original state. This is one of the Origenistic heresies, which were condemned under Justinian. Even Gregory of Nyssa, although, like Nemesius and Cyril of Alexandria, he supposed the soul to be created before the body, compares Origen’s theory to the heathen myths and fables. Origen himself allowed that the Bible does not directly teach the pre-existence of the soul, but maintained that several passages, such as the strife between Esau and Jacob in the womb, and the leaping of John the Baptist in the womb of Elizabeth at the salutation of Mary, imply it. The only truth in this theory is that every human soul has from eternity existed in the thought and purpose of God.\footnote{Lately the theory of pre-existence has found in America an advocate in Dr. Edward Beecher, in his book: The Conflict of Ages, Boston, 1853. Wordsworth has given it a poetic garb in his Ode on Immortality: “Our birth is but a sleep and a forgetting: The soul that rises with us, our life’s star, Hath had elsewhere its setting, And cometh from afar.”}

Augustine emphatically rejects the doctrine of pre-existence,\footnote{De civit. Dei, xi. 23. Ad Oros. c. Priscill. et Orig. c. 8. In his earlier work, De libero arbitrio (about 395), he spoke more favorably of Pre-existentianism.} without considering that his own theory of a generic pre-existence and apostasy of all men in Adam is really liable to similar objections. For he also hangs the whole fate of the human race on a transcendental act of freedom, lying beyond our temporal consciousness though, it is true, he places this act in the beginning of earthly history, and ascribes it to the one general ancestor, while Origen transfers it into a previous world, and views it as an act of each individual soul.\footnote{Comp. Baur, Vorlesungen über die Dogmengeschichte, Bd. i. Th. ii. p. 31: “What essentially distinguishes the Augustinian system from that of Origen, consists only [?] in this, that in place of the pretemporal fall of souls we have the Adamic apostasy, and that what in Origen bears yet a heathen impress, has in Augustineassumed a purely Old Testament [certainly, however, also a Pauline] form.”}

But between creationism and traducianism Augustine wavers, because the Scriptures do not expressly decide. He wishes to keep both the continuous creative activity of God and the organic union of body and soul.

Augustine regards this whole question as belonging to science and the schools, not to faith and the church, and makes a confession of ignorance which, in a man of his speculative genius, involves great self-denial. “Where the Scripture,” he says, “renders no certain testimony, human inquiry must beware of deciding one way or the other. If it were necessary to salvation to know anything concerning it, Scripture would have said more.”\footnote{De peccatorum mer. et remiss. l. ii. c. 36, § 59. He still remained thus undecided in his Retractations, lib. i. cap. 1, § 3 (Opera, tom. i. f. 4), where he honestly acknowledges: “Quod attinet ad ejus [animi] originem ... nec tunc sciebam, nec adhuc scio.” He frequently treats of this question, e.g., De anima et ejus origine De Genesi ad literam, x. 23; Epist. 190 ad Optatum; and Opus imperf. iv. 104. Comp. also Gangauf, l.c. p. 248 ff. and John Huber, Philosophie der Kirchenväter, p. 291ff. Huber gives the following terse presentation of the Augustinian doctrine: “In the problem of the origin of the soul Augustinearrived at no definite view. In his earlier writings he is as yet even unsettled as to the doctrine of pre-existence (De lib. arbitr. i. 12, 24; iii. 20 and 21), but afterwards he rejects it most decidedly, especially as presented by Origen, and at the same time criticizes his whole theory of the origin of the world (De civit. Dei, xi. 23). In like manner he declares against the theory of emanation, according to which the soul has flowed out of God (De Genes. ad. lit. vii. 2, 3), is of one \emph{nature} (Epist. 166 ad Hieron. § 3) and coeternal (De civ. Dei, x. 31). Between creationism and generationism, however, he can come to no decision, being kept in suspense not so much by scientific as by theological considerations. As to generationism, he remembers Tertullian, and fears being compelled, like him, to affirm the corporeality of the soul. He perceives, however, that this theory explains the transmission of original sin, and propounds the inquiry, whether perchance one soul may not spring from another, as one light is kindled from another without diminution of its flame (Ep. 190 ad Optatum, 4, 14-15). But for creationism the chief difficulty lies in this very doctrine of original sin. If the soul is created directly by God, it is pure and sinless, and the question arises, how it has deserved to be clothed with corrupt flesh and brought into the succession of original sin. God Himself appears there to be the cause of its sinfulness, inasmuch as he caused it to become guilty by uniting it with the body (De an. et ejus orig. i. 8, 9; ii. 9, 13). All the passages of Scripture relevant to this point agree only in this, that God is the Giver, Author, and Former of souls; but how he forms them—whether he creates them out of nothing or derives them from the parents, they do not declare (lb. iv. 11, 15).—His doctrine, that God created everything together as to the germ, might naturally have inclined him rather to generationism, yet he does not get over his indecision, and declares even in his Retractations (i. 1, 3), that he neither know previously nor knows now, whether succeeding souls were descended from the first one or newly created as individuals.}

The three theories of the origin of the soul, we may remark by way of concluding criticism, admit of a reconciliation. Each of them contains an element of truth, and is wrong only when exclusively held. Every human soul has an ideal pre-existence in the divine mind, the divine will, and we may add, in the divine life; and every human soul as well as every human body is the product of the united agency of God and the parents. Pre-existentianism errs in confounding an ideal with a concrete, self-conscious, individual pre-existence; traducianism, in ignoring the creative divine agency without which no being, least of all an immortal mind, can come into existence, and in favoring a materialistic conception of the soul; creationism, in denying the human agency, and thus placing the soul in a merely accidental relation to the body.

\xsection{Arguments for the Doctrine of Original Sin and Hereditary Guilt}

§ 155. Arguments for the Doctrine of Original Sin and Hereditary Guilt.

We now pass to the proofs by which Augustine established his doctrine of original sin and guilt, and to the objections urged by his opponents.

1. For Scriptural authority he appealed chiefly and repeatedly to the words in Rom. v. 12, ἐφ  ωὟ πάντες ἥμαρτον, which are erroneously translated by the Vulgate: in quo\footnote{822 Which presupposes \textgreek{ἐν ωὟ}.The whole verse reads in the Vulgate: “Propterea, sicut per unum hominem peccatum in hunc mundum intravit et per peccatum more, et ita in omnes homines mors pertransiit, \emph{in quo} omnes peccaverunt.” Comp. Augustine, De peccat. merit. et remissione, i. 8, 10; Op. imperf. ii. 63; Contra duas Ep. Pel. iv. 4; De nupt. et concup. ii. 5. Pelagius explained the passage (ad Rom. v. 12): “In eo, quod omnes peccaverunt, exemplo Adae peccant,” or per imitationem in contrast with per propagationem. Juliantranslated \textgreek{ἐφ  ωὟ} propter quod. Comp. Contra Jul. vi. 75; Op. imperf. ii. 66.} omnes peccaverunt. As Augustine had but slight knowledge of Greek, he commonly confined himself to the Latin Bible, and here he referred the in quo to Adam (the “one man” in the beginning of the verse, which is far too remote); but the Greek ἐφ  ωὟ must be taken as neuter and as a conjunction in the sense: on the ground that, or because, all have sinned.\footnote{\textgreek{Ἐφ  ωὟ} (= \textgreek{ἐφ  οἷς} ) is equivalent to \textgreek{ἐπὶ τούτῳ ὅτι}, on the ground that, presupposing that, propterea quod. So Meyer, \emph{in loco}, and others. R. Rothe (in an extremely acute exegetical monograph upon Rom. v. 12-21, Wittenberg, 1836) and Chr. Fr. Schmid (Bibl. Theol. ii. p. 126) explain \textgreek{ἐφ  ωὟ} by\textgreek{ἐπὶ τούτῶ ωὝστε}, i.e., under the more particular specification that, inasmuch as. Comp. the Commentaries.} The exegesis of Augustine, and his doctrine of a personal fall, as it were, of all men in Adam, are therefore doubtless untenable. On the other hand, Paul unquestionably teaches in this passage a causal connection between sin and death, and also a causal connection between the sin of Adam and the sinfulness of his posterity, therefore original sin. The proof of this is found in the whole parallel between Adam and Christ, and their representative relation to mankind (Comp. 1 Cor. xv. 45 ff.), and especially in the pavnte” h]marton, but not in the ejf j w| as translated by the Vulgate and Augustine. Other passages of Scripture to which Augustine appealed, as teaching original sin, were such as Gen. viii. 21; Ps. li. 7; John iii. 6; 1 Cor. ii. 14; Eph. ii. 3.

2. The practice of infant baptism in the church, with the customary formula, “for remission of sins,” and such accompanying ceremonies as exorcism, presupposes the dominion of sin and of demoniacal powers even in infancy. Since the child, before the awakening of self-consciousness, has committed. no actual sin, the effect of baptism must relate to the forgiveness of original sin and guilt.\footnote{Comp. De nuptiis et concup. i. c. 26 (tom. x. f. 291 sq.); De peccat. mer. et remiss. i. c. 26 (§ 39, tom. x. fol. 22); De gratia Christi, c. 82, 33 (x. 245 sq.), and other passages. The relation of the doctrine of original sin to the practice of infant baptism came very distinctly into view from the beginning of the controversy. Some have even concluded from a passage of Augustine(De pecc. mer. iii. 6), that the controversy began with infant baptism and not with original sin. Comp. Wiggers, i. p. 59.} This was a very important point from the beginning of the controversy, and one to which Augustine frequently reverted.

Here he had unquestionably a logical advantage over the Pelagians, who retained the traditional usage of infant baptism, but divested it of its proper import, made it signify a mere ennobling of a nature already good, and, to be consistent, should have limited baptism to adults for the forgiveness of actual sins.

The Pelagians, however, were justly offended by the revolting inference of the damnation of unbaptized infants, which is nowhere taught in the Holy Scriptures, and is repugnant to every unperverted religious instinct. Pelagius inclined to assign to unbaptized infants a middle state of half-blessedness, between the kingdom of heaven appointed to the baptized and the hell of the ungodly; though on this point he is not positive.\footnote{“Quo non eant scio, quo eant nescio,” says he of unbaptized children. He ascribed to them, it is true, salus or vita aeterna, but not the reguum coelorum. Aug. De pecc. mer. et remissione, i. 18; iii. 3. In the latter place Augustinesays, that it is absurd to affirm a “vita aeterna \emph{extra} regnum Dei.” In his book, De haeresibus, cap. 88, Augustinesays of the Pelagians that they assign to unbaptized children “aeternam et beatam quandam vitam extra regnum Dei,” and teach that children being born without original sin, are baptized for the purpose of being admitted “ad regnum Dei,” and transferred “de bono in melius.”} He evidently makes salvation depend, not so much upon the Christian redemption, as upon the natural moral character of individuals. Hence also baptism had no such importance in his view as in that of his antagonist.

Augustine, on the authority of Matt. xxv. 34, 46, and other Scriptures, justly denies a neutral middle state, and meets the difficulty by supposing different degrees of blessedness and damnation (which, in fact, must be admitted), corresponding to the different degrees of holiness and wickedness. But, constrained by the idea of original sin, and by the supposed necessity of baptism to salvation, he does not shrink from consigning unbaptized children to damnation itself,\footnote{De pecc. orig. c. 31 (§ 36, tom. x. f. 269): “Unde ergo recte \emph{infans} illa perditione punitur, nisi quia pertinet ad massam perditionis?” De nupt et concup. c. 22 (x. 292): “Remanet originale peccatum, per quod [parvuli] sub diaboli potestate captivi sunt, nisi inde lavacro regenerationis et Christi sanguine redimantur et transeant in regnum redemtoris sui.” De peccat. merit. et remissione, iii. cap. 4 (x. 74): “Manifestum est, eos [parvulos] ad damnationem, nisi hoc [incorporation with Christ through baptism] eis collatum fuerit, pertinere. Non autem damnari possent, si peccatum utique non haberent.”} though he softens to the utmost this frightful dogma, and reduces the damnation to the minimum of punishment or the privation of blessedness.\footnote{Contra Julianum, l v. c. 11 (§ 44, tom. x. f. 651): “Si enim quod de Sodomis sit [Matt. x. 15; xi. 24] et utique non solis intelligi voluit, alius alio tolerabilius in die judicii punietur quis dubitaverit \emph{parvulos non baptizatos}, qui solum habent originale peccatum, nec ullis propriis aggravantur, in \emph{damnatione omnium} \emph{levissima}futuros? ” Comp. De pecc. meritis et remissione, l. i. c. 16 (or § 21, tom. x. 12): “Potest proinde recte dici, parvulos sine baptismo de corpore exeuntes in damnatione omnium \emph{mitissima} futuros.”} He might have avoided the difficulty, without prejudice to his premises, by his doctrine of the election of grace, or by assuming an extraordinary application of the merits of Christ in death or in Hades. But the Catholic doctrine of the necessity of outward baptism to regeneration and entrance into the kingdom of God, forbade him a more liberal view respecting the endless destiny of that half of the human race which die in childhood.

We may recall, however, the noteworthy fact, that the third canon of the North-African council at Carthage in 418, which condemns the opinion that unbaptized children are saved, is in many manuscripts wanting, and is therefore of doubtful authenticity. The sternness of the Augustinian system here gave way before the greater power of Christian love. Even Augustine, De civitate Dei, speaking of the example of Melchisedec, ventures the conjecture, that God may have also among the heathen an elect people, true Israelites according to the spirit, whom He draws to Himself through the secret power of His spirit. Why, we may ask, is not this thought applicable above all to children, to whom we know the Saviour Himself, in a very special sense (and without reference to baptism) ascribes a right to the kingdom of heaven?

3. The testimony of Scripture and of the church is confirmed by experience. The inclination to evil awakes with the awaking of consciousness and voluntary activity. Even the suckling gives signs of self-will, spite, and disobedience. As moral development advances, the man feels this disposition to be really bad, and worthy of punishment, not a mere limitation or defect. Thus we find even the child subject to suffering, to sickness, and to death. It is contrary to the pure idea of God, that this condition should have been the original one. God must have created man faultless and inclined towards good. The conviction that human nature is not as it should be, in fact pervades all mankind. Augustine, in one place, cites a passage of the third book of Cicero’s Republic: “Nature has dealt with man not as a real mother, but as a step-mother, sending him into the world with a naked, frail, and feeble body, and with a soul anxious to avoid burdens, bowed down under all manner of apprehensions, averse to effort, and inclined to sensuality. Yet can we not mistake a certain divine fire of the spirit, which glimmers on in the heart as it were under ashes.” Cicero laid the blame of this on creative nature. “He thus saw clearly the fact, but not the cause, for he had no conception of original sin, because he had no knowledge of the Holy Scriptures.”

\xsection{Answers to Pelagian Objections}

§ 156. Answers to Pelagian Objections.

To these positive arguments must be added the direct answers to the objections brought against the Augustinian theory, sometimes with great acuteness, by the Pelagians, and especially by Julian of Eclanum, in the dialectic course of the controversy.

Julian sums up his argument against Augustine in five points, intended to disprove original sin from premises conceded by Augustine himself: If man is the creature of God, he must come from the hands of God good; if marriage is in itself good, it cannot generate evil; if baptism remits all sins and regenerates, the children of the baptized cannot inherit sin; if God is righteous, he cannot condemn children for the sins of others; if human nature is capable of perfect righteousness, it cannot be inherently defective.\footnote{Contra Julianum Pelagianum, l. ii. c. 9 (§ 31, tom. x. f. 545 sq.).}

We notice particularly the first four of these points; the fifth is substantially included in the first.

1. If original sin propagates itself in generation, if there is a tradux peccati and a malum naturale, then sin is substantial, and we are found in the Manichaean error, except that we make God, who is the Father of children, the author of sin, while Manichaeism refers sin to the devil, as the father of human nature.\footnote{Comp. as against this the 2d book De nuptiis et concup.; Contra Jul. l. i. and ii., and the Opus imperf., in the introduction, and lib. iv. cap. 38.}

This imputation was urged repeatedly and emphatically by the sharp and clear-sighted Julian. But according to Augustine all nature is, and ever remains, in itself good, so far as it is nature (in the sense of creature); evil is only corruption of nature, vice cleaving to it. Manichaeus makes evil a substance, Augustine, only an accident; the former views it as a positive and eternal principle, the latter derives it from the creature, and attributes to it a merely negative or privative existence; the one affirms it to be a necessity of nature, the other, a free act; the former locates it in matter, in the body, the latter, in the will.\footnote{“Non est ulla substantia vel natura, sed vitium.” De nupt. et concup. l. ii. c. 34 (§ 57, x. f. 332).“Non ortum est malum nisi in bono; nec tamen summo et immutabli, quod est natura Dei, sed facto de nihilo per sapientiam Dei.” Ibid. lib. ii. c. 29 (or § 50, tom. x. f 327). Comp. particularly also Contra duas Epist. Pelag.ii. c. 2, where he sharply discriminates his doctrine alike from Manichaeism and Pelagianism. These passages were overlooked by Baurand Milman, who think that there is good foundation for the charge of Manichaeism against Augustine’s doctrine of sin. Gibbon; (ch. xxxiii.) derived the orthodoxy of Augustinefrom the Manichaean school!} Augustine retorted on the Pelagians the charge of Manichaeism, for their locating the carnal lust of man in his original nature itself, and so precluding its cure. But in their view the concupiscentia carnis was not what it was to Augustine, but an innocent natural impulse, which becomes sin only when indulged to excess.

2. If evil is nothing substantial, we should expect that the baptized and regenerate, in whom its power is broken, would beget sinless children. If sin is propagated, righteousness should be propagated also.

But baptism, according to Augustine, removes only the guilt (reatus) of original sin, not the sin itself (concupiscentia). In procreation it is not the regenerate spirit that is the agent, but the nature which is still under the dominion of the concupiscentia. “Regenerate parents produce not as sons of God, but as children of the world.” All that are born need therefore regeneration through the same baptism, which washes away the curse of original sin. Augustine appeals to analogies; especially to the fact that from the seed of the good olive a wild olive grows, although the good and the wild greatly differ.\footnote{De peccat. mer. et remiss. ii. cap. 9 and c. 25; De nuptiis et concup. i. c. 18; Contra Julian. vi. c. 5.}

3. But if the production of children is not possible without fleshly lust, must not marriage be condemned?\footnote{Comp. against this especially the first book De nuptiis et concupiscentia (tom. x. f. 279 sqq.), written 418 or 419, in order to refute this objection. Juliananswered this in a work of four books, which gave Augustineoccasion to compose the second book De nuptiis et concup., and the six books Contra Julianum, a.d.421. Julianpublished an answer to this again, which Augustinein turn refuted in his Opus imperf., a.d.429, during the writing of which he died, a.d.430.}

No; marriage, and the consequent production of children, are, like nature, in themselves good. They belong to the mutual polarity of the sexes. The blessing: “Be fruitful and multiply,” and the declaration: “Therefore shall a man leave his father and his mother, and shall cleave unto his wife, and they shall be one flesh,” come down from paradise itself, and generation would have taken place even without sin, yet “sine ulla libidine,” as a “tranquilla motio et conjunctio vel commixtio membrorum.” Carnal concupiscence is subsequent and adventitious, existing now as an accident in the act of generation, and concealed by nature herself with shame; but it does not annul the blessing of marriage. It is only through sin that the sexual parts have become pudenda; in themselves they are honorable. Undoubtedly the regenerate are called to reduce concupiscence to the mere service of generation, that they may produce children, who shall be children of God, and therefore born again in Christ. Such desire Augustine, with reference to 1 Cor. vii. 3 ff., calls “a pardonable guilt.” But since, in the present state, the concupiscentia carnis is inseparable from marriage, it would have been really more consistent to give up the “bonum nuptiarum,” and to regard marriage as a necessary evil; as the monastic asceticism, favored by the spirit of the age, was strongly inclined to do. And in this respect there was no material difference between Augustine and Pelagius. The latter went fully as far, and even farther, in his praise of virginity, as the highest form of Christian virtue; his letter to the nun Demetrias is a picture of a perfect virgin who in her moral purity proves the excellency of human nature.

4. It contradicts the righteousness of God, to suppose one man punished for the sin of another. We are accountable only, for sins which are the acts of our own will. Julian appealed to the oft-quoted passage, Ezek. xviii. 2–4, where God forbids the use of the proverb in Israel: “The fathers have eaten sour grapes, and the children’s teeth are set on edge,” and where the principle is laid down: “The soul that sinneth, it shall die.”\footnote{Aug. Opus imperf. iii. 18, 19 (tom. x. 1067, 1069). Augustine’s answer is unsatisfactory.}

On the individualizing principle of Pelagius this objection is very, natural, and is irrefragable; but in the system of Augustine, where mankind appears as an organic whole, and Adam as the representative of human nature and as including all his posterity, it partially loses its force. Augustine thus makes all men sharers in the fall, so that they are, in fact, punished for what they themselves did in Adam. But this by no means fully solves the difficulty. He should have applied his organic view differently, and should have carried it farther. For if Adam must not be isolated from his descendants, neither must original sin be taken apart from actual sin. God does not punish the one without the other. He always looks upon the life of man as a whole; upon original sin as the fruitful mother of actual sins; and he condemns a man not for the guilt of another, but for making the deed of Adam his own, and repeating the fall by his own voluntary transgression. This every one does who lives beyond unconscious infancy. But Augustine, as we have already, seen, makes even infancy subject to punishment for original sin alone, and thus unquestionably trenches not only upon the righteousness of God, but also upon his love, which is the beginning and end of his ways, and the key to all his works.

To sum up the Augustinian doctrine of sin: This fearful power is universal; it rules the species, as well as individuals; it has its seat in the moral character of the will, reaches thence to the particular actions, and from them reacts again upon the will; and it subjects every man, without exception, to the punitive justice of God. Yet the corruption is not so great as to alter the substance of man, and make him incapable of redemption. The denial of man’s capacity for redemption is the Manichaean error, and the opposite extreme to the Pelagian denial of the need of redemption. “That is still good,” says Augustine, “which bewails lost good; for had not something good remained in our nature, there would be no grief over lost good for punishment.”\footnote{De Genesi ad literam, viii. 14.} Even in the hearts of the heathen the law of God is not wholly obliterated,\footnote{Rom. ii. 14.} and even in the life of the most abandoned men there are some good works. But these avail nothing to salvation. They are not truly good, because they proceed from the turbid source of selfishness. Faith is the root, and love the motive, of all truly good actions, and this love is shed abroad in our hearts by the Holy Ghost. “Whatsoever is not of faith, is sin.” Before the time of Christ, therefore, all virtues were either, like the virtues of the Old Testament saints, who hoped in the same Christ in whom we believe, consciously or unconsciously Christian; or else they prove, on closer inspection, to be comparative vices or seeming virtues, destitute of the pure motive and the right aim. Lust of renown and lust of dominion were the fundamental traits of the old Romans, which first gave birth to those virtues of self-devotion to freedom and country, so glorious in the eyes of men; but which afterwards, when with the destruction of Carthage all manner of moral corruption poured in, begot the Roman vices.\footnote{The sentence often ascribed to Augustine, that “all pagan virtues are but splendid vices,” is not Augustinian in form, but in substance. Comp. the quotation and remarks above, §151. Dr. Baurstates his view correctly and clearly when he says (Vorlesungen über die Dogmengeschichte, Bd. i. Part 2, p. 342): “If, as Augustinetaught, faith in Christ is the highest principle of willing and acting, nothing can be truly good, which has not its root in faith, which principle Augustinethus expressed, using the words of the apostle Paul, Rom. xiv. 23: ’Omne, quod non ex fide, peccatum.’ Augustinejudged therefore all good in the will and act of man after the absolute standard of Christian good, and accordingly could only regard the virtues of the heathen as seeming virtues, and ascribe to anything pre-Christian an inner value only so far as it had an inner reference to faith in Christ.” Comp. also Baur’sGeschichte der christl. Kirche vom 4-6ten Jahrhundert, p. 153 ff. Neanderrepresents Augustine’s doctrine on heathen virtue thus (Church History, vol. iv. 1161, 2d Germ. ed., or vol. ii. p. 620, in Torrey’s translation): ”Augustinevery justly distinguishes the patriotism of the ancients from that which is to be called ’virtue,’ in the genuinely Christian sense, and which depends on the disposition towards God (\emph{virtus} from \emph{virtus vera}); but then he goes so far as to overlook altogether what bears some relationship to the divine life in such occasional coruscations of the moral element of human nature, and to see in them nothing but a service done for evil spirits and for man’s glory. He contributed greatly, on this particular side, to promote in the Western church the partial and contracted way of judging the ancient pagan times, as opposed to the more liberal Alexandrian views of which we still find traces in many of the Orientals in this period, and to which Augustinehimself, in the earlier part of his life, as a Platonist, had been inclined. Still the vestiges of his earlier and loftier mode of thinking are to be discerned in his later writings, where he searches after and recognizes the scattered fragments of truth and goodness in the pagan literature, which he uniformly traces to the revelation of the Spirit, who is the original source of all that is true and good, to created minds; though this is inconsistent with his own theory respecting the total corruption of human nature, and with the \emph{particularism} of his doctrine of predestination.”}

This view of heathen or natural morality as a specious form of vice, though true to a large extent, is nevertheless an unjust extreme, which Augustine himself cannot consistently sustain. Even he was forced to admit important moral differences among the heathen: between, for example, a Fabricius, of incorruptible integrity, and the traitor Catiline; and though he merely defines this difference negatively, as a greater and less degree of sin and guilt, yet this itself involves the positive concession, that Fabricius stands nearer the position of Christian morality, and that there exists at least relative goodness among the heathen. Moreover, he cannot deny, that there were before Christ, not only among the Israelites, but also among the Gentiles, God-fearing souls, such as Melchisedec and Job, true Israelites, not according to the flesh, but according to the spirit, whom God by the secret workings of His Spirit drew to Himself even without baptism and the external means of grace.\footnote{Comp. De peccat. orig. c. 24 (§ 28, tom. x. f. 265), where he asserts that the grace and faith of Christ operated even unconsciously “sive in eis justis quos sancta Scriptura commemorat, sive in eis justis quos quidem illa non commemorat, sed tamen fuisse credendi sunt, vel ante diluvium, vel inde usque ad legem datam, vel ipsius legis tempore, non solum in filiis Israel, sicut fuerant prophetae, sed etiam\emph{extra eundem sicut, fuit Job}. Et ipsorum enim corda eadem mundabantur mediators fide, et diffundebatur in eis caritas per Spiritum Sanctum, qui ubi vult spirat, non merita sequens, sed etiam ipsa merita faciens.”} So the Alexandrian fathers saw scattered rays of the Logos in the dark night of heathenism; only they were far from discriminating so sharply between what was Christian and what was not Christian.

All human boasting is therefore excluded, man is sick, sick unto death out of Christ, but he is capable of health; and the worse the sickness, the greater is the physician, the more powerful is the remedy—redeeming grace.

\xsection{Augustine's Doctrine of Redeeming Grace}

§ 157. Augustine’s Doctrine of Redeeming Grace.

Augustine reaches his peculiar doctrine of redeeming grace in two ways. First he reasons upwards from below, by the law of contrast; that is, from his view of the utter incompetency of the unregenerated man to do good. The greater the corruption, the mightier must be the remedial principle. The doctrine of grace is thus only the positive counterpart of the doctrine of sin. In the second place he reasons downwards from above; that is, from his conception of the all-working, all-penetrating presence of God in natural life, and much more in the spiritual. While Pelagius deistically severs God and the world after the creation, and places man on an independent footing, Augustine, even before this controversy, was, through his speculative genius and the earnest experience of his life, deeply penetrated with a sense of the absolute dependence of the creature on the Creator, in whom we live, and move, and have our being. But Augustine’s impression of the immanence of God in the world has nothing pantheistic; it does not tempt him to deny the transcendence of God and his absolute independence of the world. Guided by the Holy Scriptures, he maintains the true mean between deism and pantheism. In the very beginning of his Confessions\footnote{Liber i. c. 2.} he says very beautifully: ’How shall I call on my God, on my God and Lord? Into myself must I call Him, if I call on Him; and what place is there in me, where my God may enter into me, the God, who created heaven and earth? O Lord my God, is there anything in me, that contains Thee? Do heaven and earth contain Thee, which Thou hast created, in which Thou didst create me? Or does all that is, contain Thee, because without Thee there had existed nothing that is? Because then I also am, do I supplicate Thee, that Thou wouldst come into me, I, who had not in any wise been, if Thou wert not in me? I yet live, I do not yet sink into the lower world, and yet Thou art there. If I made my bed in hell, behold, Thou art there. I were not, then, O my God, I utterly were not, if Thou wert not in me. Yea, still more, I were not, O my God, if I were not in Thee, from whom all, in whom all, through whom all is. Even so, Lord, even so.” In short, man is nothing without God, and everything in and through God. The undercurrent of this sentiment could not but carry this father onward to all the views he developed in opposition to the Pelagian heresy.

While Pelagius widened the idea of grace to indefiniteness, and reduced it to a medley of natural gifts, law, gospel, forgiveness of sins, enlightenment, and example, Augustine restricted grace to the specifically Christian sphere (and, therefore, called it gratia Christi), though admitting its operation previous to Christ among the saints of the Jewish dispensation; but within this sphere he gave it incomparably greater depth. With him grace is, first of all, a creative power of God in Christ transforming men from within. It produces first the negative effect of forgiveness of sins, removing the hindrance to communion with God; then the positive communication of a new principle of life. The two are combined in the idea of justification, which, as we have already remarked, Augustine holds, not in the Protestant sense of declaring righteous once for all, but in the Catholic sense of gradually making righteous; thus substantially identifying it with sanctification.\footnote{De spiritu et litera, c. 26 (tom. x. f. 109): “Quid est enim aliud, justificati, quam \emph{justi facti}, ab illo scilicet qui justificat impium, ut ex impio fiat justus?” Retract. ii. 33: “Justificamur gratia Dei, hoc est \emph{justi efficimur}.”} Yet, as he refers this whole process to divine grace, to the exclusion of all human merit, he stands on essentially Evangelical ground.\footnote{Comp. De gratia et libero arbitrio, c. 8 (§ 19), and many other places, where he ascribes fides, caritas, omnia bona opera, and vita aeterna to the free, unmerited grace of God.} As we inherit from the first Adam our sinful and mortal life, so the second Adam implants in us, from God, and in God, the germ of a sinless and immortal life. Positive grace operates, therefore, not merely from without upon our intelligence by instruction and admonition, as Pelagius taught, but also in the centre of our personality, imparting to the will the power to do the good which the instruction teaches, and to imitate the example of Christ.\footnote{“Non lege atque doctrina insonante forinsecus, sed interna et occulta, mirabili ac ineffabili potestate operatur Deus in cordibus hominum non solum veras revelationes, sed bonas etiam voluntates.” De grat. Christi, cap. 24 (x. f. 24).} Hence he frequently calls it the inspiration of a good will, or of love, which is the fulfilling of the law.\footnote{De corrept. et grat. cap. 2 (x. 751): “Inspiratio bonae voluntatis atque operis.” Without this grace men can “nullum prorsus sive cogitando, sive volendo et amando, sive agendo facere bonum.” Elsewhere he calls it also, “inspiratio dilectionis” and “caritatis.” C. duas Epist. Pel. iv., and De gratia Christi, 39.} “Him that wills not, grace comes to meet, that he may will; him that wills, she follows up, that he may not will in vain.”\footnote{“Nolentem praevenit, ut velit; volentem subsequitur, ne frustra velit.” Enchir. c. 82.} Faith itself is an effect of grace; indeed, its first and fundamental effect, which provides for all others, and manifests itself in love. He had formerly held faith to be a work of man (as, in fact, though not exclusively, the capacity of faith, or receptivity for the divine, may be said to be); but he was afterwards led, particularly by the words of Paul in 1 Cor. iv. 7: “What hast thou, that thou hast not received?” to change his view.\footnote{Comp. Retract i. c. 23; De dono perseverantiae, c. 20, and De praedest. c. 2.} In a word, grace is the breath and blood of the new man; from it proceeds all that is truly good and divine, and without it we can do nothing acceptable to God.

From this fundamental conception of grace arise the several properties which Augustine ascribes to it in opposition to Pelagius:

First, it is absolutely necessary to Christian virtue; not merely auxiliary, but indispensable, to its existence. It is necessary “for every good act, for every good thought, for every good word of man at every moment.” Without it the Christian life can neither begin, proceed, nor be consummated. It was necessary even under the old dispensation, which contained the gospel in the form of promise. The saints before Christ lived of His grace by anticipation. “They stood,” says Augustine, “not under the terrifying, convicting, punishing law, but under that grace which fills the heart with joy in what is good, which heals it, and makes it free.”\footnote{“Erant tamen et legis tempore homines Dei, non sub lege terrente, convincente, puniente, sed sub gratia delectante, sanante, liberante.” De grat. Christi et de peccato origin. l. ii. c. 25 (§ 29).}

It is, moreover, unmerited. Gratia would be no gratia if it were not gratuita, gratis data.\footnote{Comp. De gestis Pelagii, § 33 (x. 210); De pecc. orig. § 28 (x. 265): “Non Dei gratia erit ullo modo, nisi gratuita fuerit omni modo.” In many other passages he says: gratia gratis datur; gratia praecedit bona opera; gratia praecedit merita; gratia indignis datur.} As man without grace can do nothing good, he is, of course, incapable of deserving grace; for, to deserve grace, he must do something good. “What merits could we have, while as yet we did not love God? That the love with which we should love might be created, we have been loved, while as yet we had not that love. Never should we have found strength to love God, except as we received such a love from Him who had loved us before, and because He had loved us before. And, without such a love, what good could we do? Or, how could we not do good, with such a love?” “The Holy Spirit breathes where He will, and does not follow merits, but Himself produces the merits!\footnote{De pecc. orig. § 28 (x. 265): “Et ipsorum [prophetarum] corda eadem mundabantur mediatoris fide, et diffundebatur in eis caritas per Spiritum Sanctum, qui ubi vult spirat, non merita sequens, sed etiam ipsa merita faciens.”} Grace, therefore, is not bestowed on man because he already believes, but that he may believe; not because he has deserved it by good works, but that he may deserve good works.” Pelagius reverses the natural relation by making the cause the effect, and the effect the cause. The ground of our salvation can only be found in God Himself, if He is to remain immutable. Augustine appeals to examples of pardoned sinners, “where not only no good deserts, but even evil deserts, had preceded.” Thus the apostle Paul, “averse to the faith, which he wasted, and vehemently inflamed against it, was suddenly converted to that faith by the prevailing power of grace, and that in such wise that he was changed not only from an enemy to a friend, but from a persecutor to a sufferer of persecution for the sake of the faith he had once destroyed. For to him it was given by Christ, not only to believe on him, but also to suffer for his sake.” He also points to children, who without will, and therefore without voluntary merit preceding, are through holy baptism incorporated in the kingdom of grace.\footnote{De gratia et libero arbitrio, cap. 22 (§ 44, tom. x. f. 742). Parvuli, he says, have no will to receive grace, nay, often struggle with tears against being baptized, “quod eis ad magnum impietatis peccatum imputaretur, si jam libero uterentur arbitrio: et tamen haeret etiam in reluctantibus gratia, apertissime nullo bono merito praecedente, alioquin gratia jam non esset gratia.” He then calls attention to the fact that grace is sometimes bestowed on children of unbelievers, and is withheld from many children of believers.} His own experience, finally, afforded him an argument, to him irrefutable, for the free, undeserved compassion of God. And if in other passages he speaks of merits, he means good works which the Holy Ghost effects in man, and which God graciously rewards, so that eternal life is grace for grace. “If all thy merits are gifts of God, God crowns thy merits not as thy merits, but as the gifts of his grace.”\footnote{De grat. et lib. arbitrio, c. 6 (f. 726), where Augustine, from passages like James i. 17; John iii. 27; Eph. ii. 8, draws the conclusion: “Si ergo Dei dona sunt bona merita tua, non Deus coronat merita tua tamquam merita tua, sed tamquam dona sua.”}

Grace is irresistible in its effect; not, indeed, in the way of physical constraint imposed on the will, but as a moral power, which makes man willing, and which infallibly attains its end, the conversion and final perfection of its subject.\footnote{“Subventum est infirmitati voluntatis humanae, ut divina gratia \emph{indeclinabiliter} et \emph{insuperabiliter} [not \emph{inseparabiliter}, as the Jesuit edition of Louvain, 1577, reads] ageretur; et ideo, quamvis infirma, non tamen deficeret, neque adversitate aliqua vinceretur.” De corrept. et grat. § 38 (tom. x. p. 771).} This point is closely connected with Augustine’s whole doctrine of predestination, and consistently leads to it or follows from it. Hence the Pelagians repeatedly raised the charge that Augustine, under the name of grace, introduced a certain fatalism. But the irresistibility must manifestly not be extended to all the influences of grace; for the Bible often speaks of grieving, quenching, lying to, and blaspheming the Holy Ghost, and so implies that grace may be resisted; and it presents many living examples of such resistance. It cannot be denied, that Saul, Solomon, Ananias, and Sapphira, and even the traitor Judas, were under the influence of divine grace, and repelled it. Augustine, therefore, must make irresistible grace identical with the specific grace of regeneration in the elect, which at the same time imparts the donum perseverantiae.\footnote{It is in this sense that the Calvinistic theologians have always understood the Augustinian system, especially the Presbyterians. So, e.g., Dr. Cunningham(l.c. vol. ii. p. 352): ”Augustine, in asserting the invincibility or irresistibility of grace, did not mean—and those who in subsequent times have embraced this general system of doctrine as scriptural, did not intend to convey the idea—that man was compelled to do that which was good, or that he was forced to repent and believe against his will, whether he would or not, as the doctrine is commonly misrepresented, but merely that he was certainly and effectually made willing, by the renovation of his will through the power of God, \emph{whenever that power was put forth in a measure} Sufficient \emph{and} Adequate \emph{to produce the result}. Augustine, and those who have adopted his system, did not mean to deny that men may, in some sense and to some extent, resist the Spirit, the possibility of which is clearly indicated in Scripture; inasmuch as they have most commonly held that, to use the language of our [the Westminster] Confession, ’persons who are not elected and who finally perish, may have some common operations of the Spirit,’ which, of course, they resist and throw off.” Similarly Dr. Shedd(Hist. of Doct. vol. ii. 73), who, however, extends irresistible grace to all the regenerate. “Not all grace,” he says, “but the grace which actually regenerates, Augustinedenominates \emph{irresistible}. By this he meant, not that the human will is converted unwillingly or by compulsion, but that divine grace is able to overcome the utmost obstinacy of the human spirit. ... Divine grace is irresistible, not in the sense that no form of grace is resisted by the sinner; but when grace reaches that special degree which constitutes it \emph{regenerating}, it then overcomes the sinner’s opposition, and makes him willing in the day of God’s power.” This is Calvinistic, but not Augustinian, although given as Augustine’s view. For according to Augustineall the baptized are regenerate, and yet many are eternally lost. (Comp. Ep. 98, 2; De pecc. mer. et rem. i. 39, and the passages in Hagenbach’s Doctrine History, vol. i. p. 358 ff. in the Anglo-American edition.) The gratia irresistiblis must therefore be restricted to the narrower circle of the \emph{electi}. Augustine’s doctrine of baptism is far more Lutheran and Catholic than Calvinistic. According to Calvin, the regenerating effect of baptism is dependent on the \emph{decretum} \emph{divinum}, and the truly regenerate is also elect, and therefore can never finally fall from grace. Augustine, for the honor of the sacrament, assumes the possibility of a fruitless regeneration; Calvin, in the interest of election and regeneration, assumes the possibility of an ineffectual baptism.}

Grace, finally, works progressively or by degrees. It removes all the consequences of the fall; but it removes them in an order agreeable to the finite, gradually unfolding nature of the believer. Grace is a foster-mother, who for the greatest good of her charge, wisely and lovingly accommodates herself to his necessities as they change from time to time. Augustine gives different names to grace in these different steps of its development. In overcoming the resisting will, and imparting a living knowledge of sin and longing for redemption, grace is gratia praeveniens or praeparans. In creating faith and the free will to do good, and uniting the soul to Christ, it is gratia operans. Joining with the emancipated will to combat the remains of evil, and bringing forth good works as fruits of faith, it is gratia cooperans. Finally, in enabling the believer to persevere in faith to the end, and leading him at length, though not in this life, to the perfect state, in which he can no longer sin nor die, it is gratia perficiens.\footnote{Summing all the stages together, Augustinesays: “Et quis istam etsi parvam dare coeperat caritatem, nisi ille qui \emph{praeparat} voluntatem, et \emph{coöperando perficit}, quod \emph{operando incipit}? Quoniam ipse ut velimus operatur incipiens qui volentibus coöperatur perficiens. Propter quod ait Apostolus: Certus sum, quoniam qui operatur in vobis opus bonum, perficiet usque in diem Christi Jesu” (Phil. i. 6)). De grat. et lib. arbitr. c. 27, § 33 (tom. x. 735).} This includes the donum perseverantiae, which is the only certain token of election.\footnote{Augustinetreats of this in the Liber de dono persevemntiae, one of his latest writings, composed in 428 or 429 (tom. x. f. 821 sqq.).} “We call ourselves elect, or children of God, because we so call all those whom we see regenerate, visibly leading a holy life. But he alone is in truth what he is called, who perseveres in that from which he receives the name.” Therefore so long as a man yet lives, we can form no certain judgment of him in this respect. Perseverance till death, i.e., to the point where the danger of apostasy ceases, is emphatically a grace, “since it is much harder to possess this gift of grace than any other; though for him to whom nothing is hard, it is as easy to bestow the one as the other.”

And as to the relation of grace to freedom: Neither excludes the other, though they might appear to conflict. In Augustine’s system freedom, or self-determination to good, is the correlative in man of grace on the part of God. The more grace, the more freedom to do good, and the more joy in the good. The two are one in the idea of love, which is objective and subjective, passive and active, an apprehending and a being apprehended.\footnote{Comp. upon this especially the book De gratia et libero arbitrio, which Augustinewrote a.d.426, addressed to Valentinus and other monks of Adrumetum, to refute the false reasoning of those, “qui sic gratiam Dei defendunt, ut negent hominis liberum arbitrium” (c. 1, tom. x. f. 717).}

We may sum up the Augustinian anthropology under these three heads:

1. The Primitive State: Immediate, undeveloped unity of man with God; child-like innocence; germ and condition of everything subsequent; possibility of a sinless and a sinful development.

2. The State of Sin: Alienation from God; bondage; dominion of death; with longing after redemption.

3. The State of Redemption or of Grace: Higher, mediated unity with God; virtue approved through conflict; the blessed freedom of the children of God; here, indeed, yet clogged with the remains of sin and death, but hereafter absolutely perfect, without the possibility of apostasy.

\xsection{The Doctrine of Predestination}

§ 158. The Doctrine of Predestination.

I. Augustinus: De praedestinatione sanctorum ad Prosperum et Hilarium (written a.d. 428 or 429 against the Semi-Pelagians); De dono perseverantiae (written in the same year and against the same opponents); De gratia et libero arbitrio (written a.d. 426 or 427 ad Valentinum et Monachos Adrumetinos); De correptione et gratia (written to the same persons and in the same year).

II Corn. Jansenius: Augustinus. Lovan. 1640, tom. iii. Jac. Sirmond (Jesuit): Historia praedestinatiana. Par. 1648 (and in his Opera, tom. iv. p. 271). Carl Beck: Die Augustinische, Calvinistische und Lutherische Lehre von der Praedestination aus den Quellen dargestellt und mit besonderer Rücksicht auf Schleiermacher’s Erwählungslehre comparativ beurtheilt. “Studien und Kritiken,” 1847. J. B. Mozley: Augustinian Doctrine of Predestination. Lond. 1855.

Augustine did not stop with this doctrine of sin and grace. He pursued his anthropology and soteriology to their source in theology. His personal experience of the wonderful and undeserved grace of God, various passages of the Scriptures, especially the Epistle to the Romans, and the logical connection of thought, led him to the doctrine of the unconditional and eternal purpose of the omniscient and omnipotent God. In this he found the programme of the history of the fall and redemption of the human race. He ventured boldly, but reverentially, upon the brink of that abyss of speculation, where all human knowledge is lost in mystery and in adoration.

Predestination, in general, is a necessary attribute of the divine will, as foreknowledge is an attribute of the divine intelligence; though, strictly speaking, we cannot predicate of God either a before or an after, and with him all is eternal present. It is absolutely inconceivable that God created the world or man blindly, without a fixed plan, or that this plan can be disturbed or hindered in any way by his creatures. Besides, there prevails everywhere, even in the natural life of man, in the distribution of mental gifts and earthly blessings, and yet much more in the realm of grace, a higher guidance, which is wholly independent of our will or act. Who is not obliged, in his birth in this or that place, at this or that time, under these or those circumstances, in all the epochs of his existence, in all his opportunities of education, and above all in his regeneration and sanctification, to recognize and adore the providence and the free grace of God? The further we are advanced in the Christian life, the less are we inclined to attribute any merit to ourselves, and the more to thank God for all. The believer not only looks forward into eternal life, but also backward into the ante-mundane eternity, and finds in the eternal purpose of divine love the beginning and the firm anchorage of his salvation.\footnote{Rom. viii. 29; Eph. i. 4.}

So far we may say every reflecting Christian must believe in some sort of election by free grace; and, in fact, the Holy Scriptures are full of it. But up to the time of Augustine the doctrine had never been an object of any very profound inquiry, and had therefore never been accurately defined, but only very superficially and casually touched. The Greek fathers, and Tertullian, Ambrose, Jerome, and Pelagius, had only taught a conditional predestination, which they made dependent on the foreknowledge of the free acts of men. In this, as in his views of sin and grace, Augustine went far beyond the earlier divines, taught an unconditional election of grace, and restricted the purpose of redemption to a definite circle of the elect, who constitute the minority of the race.\footnote{Comp. the opinions of the pre-Augustinian fathers respecting grace, predestination, and the extent of redemption, as given in detail in Wiggers, i. p. 440 ff. He says, p. 448: “In reference to predestination, the fathers before Augustinewere entirely at variance with him, and in agreement with Pelagius. They, like Pelagius, founded predestination upon prescience, upon the fore-knowledge of God, as to who would make themselves worthy or unworthy of salvation. They assume, therefore, not the unconditional predestination of Augustine, but the conditional predestination of the Pelagians. The Massilians had, therefore, a full right to affirm (Aug. Ep. 225), that Augustine’s doctrine of predestination was opposed to the opinions of the fathers and the sense of the church (ecclesiastico sensui), and that no ecclesiastical author had ever yet explained the Epistle to the Romans as Augustinedid, or in such a way as to derive from it a grace that had no respect to the merits of the elect. And it was only by a doubtful inference (De dono pers. 19) that Augustineendeavored to prove that Cyprian, Ambrose, and Gregory Nazianzen had known and received his view of predestination, by appealing to the agreement between this doctrine and their theory of grace.” Pelagius says of predestination in his Commentary on Rom. viii. 29 and ix. 80: “Quos praevidit conformes esse in vita, voluit ut fierent conformes in gloria. .... Quos praescivit credituros, hos vocavit, vocatio autem volentes colligit, non invitos.”}

In Augustine’s system the doctrine of predestination is not, as in Calvin’s, the starting-point, but the consummation. It is a deduction from his views of sin and grace. It is therefore more practical than speculative. It is held in check by his sacramental views. If we may anticipate a much later terminology, it moves within the limits of infralapsarianism, but philosophically is less consistent than supralapsarianism. While the infralapsarian theory, starting with the consciousness of sin, excludes the fall—the most momentous event, except redemption, in the history of the world—from the divine purpose, and places it under the category of divine permission, making it dependent on the free will of the first man; the supralapsarian theory, starting with the conception of the absolute sovereignty of God, includes the fall of Adam in the eternal and unchangeable plan of God, though, of course, not as an end, or for its own sake (which would be blasphemy), but as a temporary means to an opposite end, or as the negative condition of a revelation of the divine justice in the reprobate, and of the divine grace in the elect. Augustine, therefore, strictly speaking, knows nothing of a double decree of election and reprobation, but recognizes simply a decree of election to salvation; though logical instinct does sometimes carry him to the verge of supralapsarianism. In both systems, however, the decree is eternal, unconditioned, and immutable; the difference is in the subject, which, according to one system, is man fallen, according to the other, man as such. It was a noble, inconsistency which kept Augustine from the more stringent and speculative system of supralapsarianism; his deep moral convictions revolted against making any allowance for sin by tracing its origin to the divine will; and by his peculiar view of the inseparable connection between Adam and the race, he could make every man as it were individually responsible for the fall of Adam. But the Pelagians, who denied this connection, charged him with teaching a kind of fatalism.

The first sin, according to Augustine’s theory, was an act of freedom, which could and should have been avoided. But once committed, it subjected the whole race, which was germinally in the loins of Adam, to the punitive justice of God. All men are only a mass of perdition,\footnote{Massa perditionis, a favorite expression of Augustine.} and deserve, both for their innate and their actual sin, temporal and eternal death. God is but just, if He leave a great portion, nay (if all heathen and unbaptized children are lost), the greatest portion, of mankind to their deserved fate. But He has resolved from eternity to reveal in some His grace, by rescuing them from the mass of perdition, and without their merit saving them.

This is the election of grace, or predestination. It is related to grace itself, as cause to effect, as preparation to execution.\footnote{De praedest. Sanct. c. 10 (or § 19, tom. x f. 803): “Inter gratiam et praedestinationem hoc tantum interest, quod praedestinatio eat gratiae praeparatio, gratia vero jam ipsa donatio. Quod itaque ait apostolus: Non ex operibus ne forte quis extollatur, ipsius enim sumus figmentum, creati in Christo Jesu in operibus bonis (Eph. ii. 9), gratia est; quod autem sequitur: Quae praeparavit Deus, ut in illis ambulamus, praedestinatio est, quae sine praescientia non potest esse.” Further on in the same chapter: “Gratia est ipsius praedestinationis effectus.”} It is the ultimate, unfathomable ground of salvation. It is distinguished from foreknowledge, as will from intelligence; it always implies intelligence, but is not always implied in it.\footnote{De praed. sanctorum, cap. 10: “Praedestinatio ... sine praescientia non potest esse; potest autem esse sine praedestinatione praescientia. Praedestinatione quippe Deus ea praescivit, quod fuerat ipse facturus ... praescire autem potens est etiam quae ipse non facit, sicut quaecumque peccata.” Comp. De dono perseverantiae, c. 18 (f847 sq.).} God determines and knows beforehand what He will do; the fall of man, and the individual sins of men, He knows perfectly even from eternity, but He does not determine or will them, He only permits them. There is thus a point, where prescience is independent of predestination, and where human freedom, as it were, is interposed. (Here lies the philosophical weakness, but, on the other hand, the ethical strength of the infralapsarian system, as compared with the supralapsarian). The predetermination has reference only to good, not to evil. It is equivalent to election, while predestination, in the supralapsarian scheme, includes the decretum electionis and the decretum reprobationis. Augustine, it is true, speaks also in some places of a predestination to perdition (in consequence of sin), but never of a predestination to sin.\footnote{De anima et ejus origine (written a.d.419), l. iv. c. 11 (or § 16, tom. x. f 395): “Ex uno homine omnes homines ire in condemnationem qui nascuntur ex Adam, nisi ita renascantur in Christo ... \emph{quos praedestinavit ad aeternam vitam} misericordissimus gratiae largitor: qui eat et illis \emph{quos praedestinavit ad aeternam mortem}, justissimus supplicii retributor.” Comp. Tract. in Joann. xlviii. 4: “ad sempiternum \emph{interitum} praedestinatos,” and similar passages.} The election of grace is conditioned by no foreseen merit, but is absolutely free. God does not predestinate His children on account of their faith, for their faith is itself a gift of grace; but He predestinates them to faith and to holiness.\footnote{De praed. sanct. c. 18 (§ 37, x. f. 815): “Elegit ergo nos Deus in Christo ante mundi constitutionem, praedestinans nos in adoptionem filiorum: non quia per nos sancti et immaculati futuri eramus, sed \emph{elegit praedestinavitque ut essemus}.” Augustinethen goes on to attack the Pelagian and Semi-Pelagian theory of a predestination conditioned upon the foreseen holiness of the creature. Cap. 19 (§ 38): “Nec \emph{quia} credidimus, sed \emph{ut} credamus, vocamur.”}

Thus also the imputation of teaching that a man may be elect, and yet live a godless life, is precluded.\footnote{This imputation of some monks of Adrumetum in Tunis is met by Augustineparticularly in his treatise De correptione et gratia (a.d.427), in which he shows that as gratia and the liberum arbitrium, so also correptio and gratia, admonition and grace, are by no means mutually exclusive, but rather mutually condition each other.} Sanctification is the infallible effect of election. Those who are thus predestinated as vessels of mercy, may fall for a while, like David and Peter, but cannot finally fall from grace. They must at last be saved by, the successive steps of vocation, justification, and glorification, as certainly as God is almighty and His promises Yea and Amen;\footnote{De corrept. et grat. c. 7 (§ 14): “Nemo eorum [electorum] perit, quia non fallitur Deus. Horum si quisquam perit, vitio humano vincitur Deus; sed nemo eorum perit, quia nulla re vincitur Deus.” Ibid. c. 9 (§ 23, f. 763): “Quicunque ergo in Dei providentissima dispositione praesciti, praedestinati, vocati, justificati, glorificati sunt, non dico etiam nondum renati, sed etiam nondum nati, jam filii Dei sunt, et omnino perire non possunt.” For this he appeals to Rom. viii. 31 ff.; John vi. 37, 39, etc.} while the vessels of wrath are lost through their own fault. To election necessarily belongs the gift of perseverance, the donum perseverantiae, which is attested by a happy death. Those who fall away, even though they have been baptized and regenerated, show thereby, that they never belonged to the number of the elect.\footnote{De corrept. et gratia, c. 9 (§ 23, x. f. 763): “Ab illo [Deo] datur etiam perseverantia in bono usque in finem; neque enim datur nisi eis qui non peribunt: quoniam qui non perseverant peribunt.” Ibid. c. 11 (§ 36, f. 770): “Qui autem cadunt et pereunt, in praedestinatorum numero non fuerunt.”} Hence we cannot certainly know in this life who are of the elect, and we must call all to repentance and offer to all salvation, though the vocation of grace only proves effectual to some.

Augustine, as, already remarked, deduced this doctrine from his view of sin. If all men are by nature utterly incompetent to good, if it is grace that works in us to will and to do good, if faith itself is an undeserved gift of grace: the ultimate ground of salvation can then be found only in the inscrutable counsel of God. He appealed to the wonderful leadings in the lives of individuals and of nations, some being called to the gospel and to baptism, while others die in darkness. Why precisely this or that one attains to faith and others do not, is, indeed, a mystery. We cannot, says he, in this life explain the readings of Providence; if we only believe that God is righteous, we shall hereafter attain to perfect knowledge.

He could cite many Scripture texts, especially the ninth chapter of the Epistle to the Romans, for his doctrine. But other texts, which teach the universal vocation to salvation, and make man responsible for his reception or rejection of the gospel, he could only explain by forced interpretations. Thus, for instance, be understands in 1 Tim. ii. 4 by the all men, whom God will have to be saved, all manner of men, rich and poor, learned and unlearned, or he wrests the sense into: All who are saved, are saved only by the will of God.\footnote{Opus imperf. iv. 124; De corrept. et gratia, i. 28; De praed. sanct. 8; Enchir. c. 103; Epist. 217, c. 6. Comp. Wiggers, l.c. pp. 365 and 463 ff.} When he finds no other way of meeting objections, be appeals to the inscrutable wisdom of God.

Augustine’s doctrine of predestination was the immediate occasion of a theological controversy which lasted almost a hundred years, developed almost every argument for and against the doctrine, and called forth a system holding middle ground, to which we now turn.

\xsection{Semi-Pelagianism}

§ 159. Semi-Pelagianism.

Comp. the Works at § 146.

Sources.

I. Joh. Cassianus († 432): Collationes Patrum xxiv, especially the xiii. In the Opera omnia, cum commentaries D. Alardi Gazaei (Gazet), Atrebati (Atrecht or Arras in France), 1628 and 1733; reprinted, with additions, in Migne’s Patrologia, tom. xlix. and l. (tom. i. pp. 478–1328), and also published several times separately. Vincentius Lirinsis († 450), Faustus Rhegiensis († 490–500), and other Semi-Pelagian writers, see Gallandi, Biblioth. tom. x., and Migne, Patrol. tom. l. and liii.

II. Augustinus: De gratia et libero arbitrio; De correptione et gratia; De praedestinatione Sanctorum; De dono perseverantiae (all in the 10th vol. of the Benedict. ed.). Prosper Aquitanus (a disciple and admirer of Augustine, † 460): Epistola ad Augustinum de reliquiis Pelagianae haereseos in Gallia (Aug. Ep. 225, and in Opera Aug. tom. x. 780), and De gratia et libero arbitrio (contra Collatorem). Hilarius: Ad Augustinum de eodem argumento (Ep. 226 among the Epp. Aug., and in tom. x. 783). Also the Augustinian writings of Avitus of Vienne, Caesarius of Arles, Fulgentius of Ruspe, and others. (Comp. Gallandi, Bibl. tom. xi.; Migne, Patrol. vol. li.)

The Acta of the Synod of Orange, a.d. 529, in Mansi, tom. viii. 711 sqq.

Literature.

Jac. Sirmond: Historia praedestinatiana. Par. 1648. Johann Geffken: Historia Semipelagianismi antiquissima (more properly antiquissimi). Gott. 1826 (only goes to the year 434). G. Fr. Wiggers: Versuch einer pragmatischen Darstellung des Semipelagianismus in seinem Kampfe gegen den Augustinismus his zur zweiten Synode zu Orange. Hamburg, 1833 (the second part of his already cited work upon Augustinianism and Pelagianism). A very thorough work, but unfortunately without index. Comp, also Walch, Schröckh, and the appropriate portions of the later works upon the history of the church and of doctrines.

Semi-Pelagianism is a somewhat vague and indefinite attempt at reconciliation, hovering midway between the sharply marked systems of Pelagius and Augustine, taking off the edge of each, and inclining now to the one, now to the other. The name was introduced during the scholastic age, but the system of doctrine, in all essential points, was formed in Southern France in the fifth century, during the latter years of Augustine’s life and soon after his death. It proceeded from the combined influence of the pre-Augustinian synergism and monastic legalism. Its leading idea is, that divine grace and the human will jointly accomplish the work of conversion and sanctification, and that ordinarily man must take the first step. It rejects the Pelagian doctrine of the moral roundness of man, but rejects also the Augustinian doctrine of the entire corruption and bondage of the natural man, and substitutes the idea of a diseased or crippled state of the voluntary power. It disowns the Pelagian conception of grace as a mere external auxiliary; but also, quite as decidedly, the Augustinian doctrines of the sovereignty, irresistibleness, and limitation of grace; and affirms the necessity and the internal operation of grace with and through human agency, a general atonement through Christ, and a predestination to salvation conditioned by the foreknowledge of faith. The union of the Pelagian and Augustinian elements thus attempted is not, however, an inward organic coalescence, but rather a mechanical and arbitrary combination, which really satisfies neither the one interest nor the other, but commonly leans to the Pelagian side.\footnote{Wiggers (ii. pp. 359-364) gives a comparative view of the three systems in parallel columns. Comp. also the criticism of Baur, Die christliche Kirche vom vierten bis zum sechsten Jahrhundert, p. 181 ff. The latter, with his wonted sharpness of criticism, judges very unfavorably of Semi-Pelagianism as a whole. “This halving and neutralizing,” he says, p. 199 ff., “this attempt at equal distribution of the two complementary elements, not only setting them apart, but also balancing them with one another, so that sometimes the one, sometimes the other, is predominant, and thus within this whole sphere everything is casual and arbitrary, varying and indefinite according to the diversity of circumstances and individuals, this is characteristic of Semi-Pelagianism throughout. If the two opposing theories cannot be inwardly reconciled, at least they must be combined in such a way as that a specific element must be taken from each; the Pelagian freedom and the Augustinian grace must be advanced to equal rank. But this method only gains an external juxtaposition of the two.”}

For this reason it admirably suited the legalistic and ascetic piety of the middle age, and indeed always remained within the pale of the Catholic church, and never produced a separate sect.

We glance now at the main features of the origin and progress of this school.

The Pelagian system had been vanquished by Augustine, and rejected and condemned as heresy by the church. This result, however, did not in itself necessarily imply the complete approval of the Augustinian system. Many, even opponents of Pelagius, recoiled from a position so wide of the older fathers as Augustine’s doctrines of the bondage of man and the absolute election of grace, and preferred a middle ground.

First the monks of the convent of Adrumetum in North Africa differed among themselves over the doctrine of predestination; some perverting it to carnal security, others plunging from it into anguish and desperation, and yet others feeling compelled to lay more stress than Augustine upon human freedom and responsibility. Augustine endeavored to allay the scruples of these monks by his two treatises, De gratia et libero arbitrio, and De correptione et gratia. The abbot Valentinus answered these in the name of the monks in a reverent and submissive tone.\footnote{His answer is found in the Epistles of Augustine, Ep. 216, and in Opera, tom. x. f. 746 (ed. Bened.).}

But simultaneously a more dangerous opposition to the doctrine of predestination arose in Southern Gaul, in the form of a regular theological school within the Catholic church. The members of this school were first called “remnants of the Pelagians,”\footnote{“Reliquiae Pelagianorum.” So Prosper calls them in his letter to Augustine. He saw in them disguised, and therefore only so much the more dangerous, Pelagians.} but commonly Massilians, from Massilia (Marseilles), their chief centre, and afterwards Semi-Pelagians. Augustine received an account of this from two learned and pious lay friends, Prosper, and Hilarius,\footnote{Not to be confounded with Hilarius, bishop of Arles, in distinction from whom he is called Hilarius Prosperi. Hilarycalls himself a layman (Aug. Ep. 226, § 9). Comp. the Benedictines in tom. x. f. 785; Wiggers, ii. 187).} who begged that he himself would take the pen against it. This was the occasion of his two works, De praedestinatione sanctorum, and De dono perseverentiae, with which he worthily closed his labors as an author. He deals with these disputants more gently than with the Pelagians, and addresses them as brethren. After his death (430) the discussion was continued principally in Gaul; for then North Africa was disquieted by the victorious invasion of the Vandals, which for several decades shut it out from the circle of theological and ecclesiastical activity.

At the head of the Semi-Pelagian party stood John Cassian, the founder and abbot of the monastery at Massilia, a man of thorough cultivation, rich experience, and unquestioned orthodoxy.\footnote{Wiggers treats thoroughly and at length of him, in the above cited monograph, vol. ii. pp. 7-136. He has been mistakenly supposed a Scythian. His name and his fluent Latinity indicate an occidental origin. Yet he was in part educated at Bethlehem and in Constantinople, and spent seven years among the anchorites in Egypt. He mentioned John Chrysostomeven in the evening of his life with grateful veneration. (De incarn. vii. 30 sq.) “What I have written,” he says, “John has taught me, and therefore account it not so much mine as his. For a brook rises from a spring, and what is ascribed to the pupil, must be reckoned wholly to the honor of the teacher.” On the life and writings of Cassian compare also Schönemann, Bibliotheca, vol. ii. (reprinted in Migne’s ed. vol. i.).} He was a grateful disciple of Chrysostom, who ordained him deacon, and apparently also presbyter. His Greek training and his predilection for monasticism were a favorable soil for his Semi-Pelagian theory. He labored awhile in Rome with Pelagius, and afterwards in Southern France, in the cause of monastic piety, which he efficiently promoted by exhortation and example. Monasticism sought in cloistered retreats a protection against the allurements of sin, the desolating incursions of the barbarians, and the wretchedness of an age of tumult and confusion. But the enthusiasm for the monastic life tended strongly to over-value external acts and ascetic discipline, and resisted the free evangelical bent of the Augustinian theology. Cassian wrote twelve books De coenobiorum institutis, in which be first describes the outward life of the monks, and then their inward conflicts and victories over the eight capital vices: intemperance, unchastity, avarice, anger, sadness, dulness, ambition, and pride. More important are his fourteen Collationes Patrum, conversations which Cassian and his friend Germanus had had with the most experienced ascetics in Egypt, during a seven years’ sojourn there.

In this work, especially in the thirteenth Colloquy,\footnote{De protectione Dei. In Migne’s edition of Cass. Opera, vol. i. pp. 397-964} he rejects decidedly the errors of Pelagius,\footnote{He calls the Pelagian doctrine of the native ability of man ”\emph{profanam} opinionem ” (Coll. xiii. 16, in Migne’s ed. tom. i. p. 942), and even says: “Pelagium paene omnes impietate [probably here equivalent to ” contempt of grace,” as Wiggers, ii. 20, explains it] et amentia vicisse ” (De incarn. Dom. v. 2, tom. ii. 101).} and affirms the universal sinfulness of men, the introduction of it by the fall of Adam, and the necessity, of divine grace to every individual act. But, with evident reference to Augustine, though without naming him, he combats the doctrines of election and of the irresistible and particular operation of grace, which were in conflict with the church tradition, especially, with the Oriental theology, and with his own earnest ascetic legalism.

In opposition to both systems he taught that the divine image and human freedom were not annihilated, but only weakened, by the fall; in other words, that man is sick, but not dead, that he cannot indeed help himself, but that he can desire the help of a physician, and either accept or refuse it when offered, and that he must cooperate with the grace of God in his salvation. The question, which of the two factors has the initiative, he answers, altogether empirically, to this effect: that sometimes, and indeed usually, the human will, as in the cases of the Prodigal Son, Zacchaeus, the Penitent Thief, and Cornelius, determines itself to conversion; sometimes grace anticipates it, and, as with Matthew and Paul, draws the resisting will—yet, even in this case, without constraint—to God.\footnote{“Nonnumquam,” says he, De institut. coeno b. xii. 18 (Opera, vol. ii. p. 456, ed. Migne), “etiam inviti trahimur ad salutem.” This is, however, according to Cassian, a rare exception. The general distinction between Semi-Pelagianism and the Melanchthonian synergism may be thus defined, that the former ascribes the initiative in the work of conversion to the human will, the latter to divine grace, which involves also a different estimate of the importance of the gratia praeveniens or praeparans.} Here, therefore, the gratia praeveniens is manifestly overlooked.

These are essentially Semi-Pelagian principles, though capable of various modifications and applications. The church, even the Roman church, has rightly emphasized the necessity of prevenient grace, but has not impeached Cassian, who is properly the father of the Semi-Pelagian theory. Leo the Great even commissioned him to write a work against Nestorianism,\footnote{De incarnations Christi, libri vii. in Migne’s ed. tom. ii. 9-272.} in which he found an excellent opportunity to establish his orthodoxy, and to clear himself of all connection with the kindred heresies of Pelagianism and Nestorianism, which were condemned together at Ephesus in 431. He died after 432, at an advanced age, and though not formally canonized, is honored as a saint by some dioceses. His works are very extensively read for practical edification.

Against the thirteenth Colloquy of Cassian, Prosper Aquitanus, an Augustinian divine and poet, who, probably on account of the desolations of the Vandals, had left his native Aquitania for the South of Gaul, and found comfort and repose in the doctrines of election amid the wars of his age, wrote a book upon grace and freedom,\footnote{Found in the works of Prosper, Paris, 1711 (tom. li. in Migne’s Patrol.), and also in the Appendix to the Opera Augustini (tom. x. 171-198, ed. Bened.), under the title Pro Augustino, liber contra Collatorem. Comp. Wiggers, ii. p. 138 ff.} about 432, in which he criticises twelve propositions of Cassian, and declares them all heretical, except the first. He also composed a long poem in defence of Augustine and his system,\footnote{Carmen de ingratis. He charges the Semi-Pelagians with ingratitude to Augustineand his great merits to the cause of religion.} and refuted the “Gallic slanders and Vincentian imputations,” which placed the doctrine of predestination in the most odious light.\footnote{These Responsiones Prosperi Aquitani ad capitula calumniantium Gallorum and Ad capitula objectionum Vincentianorum (of Vincentius Lirinensis) are also found in the Appendix to the 10th vol. of the Benedictine edition of the Opera Augustini, f. 198 sqq. and f. 207 sqq. Among the objections of Vincentius are, e.g., the following: 3. Quia Deus majorem partem generis humani ad hoc creet, ut illam perdat in aeternum. 4. Quia major pars generis humani ad hoc creetur a Deo, ut non Dei, sed diaboli faciat voluntatem. 10. Quia adulteria et corruptelae virginum sacrarum ideo contingant, quia illas Deus ad hoc praedestinavit ut caderent.}

But the Semi-Pelagian doctrine was the more popular, and made great progress in France. Its principal advocates after Cassian are the following: the presbyter-monk Vicentius of Lerinum, author of the Commonitorium, in which he developed the true catholic test of doctrine, the threefold consensus, in covert antagonism to the novel doctrines of Augustinianism (about 434);\footnote{Comp. above, § 118; also Wiggers, ii. p. 208 ff., and Baur, l.c. p. 185 ff, who likewise impute to the Commonitorium a Semi-Pelagian tendency. This is beyond doubt, if Vincentius was the author of the above-mentioned Objectiones Vincentisanae. Perhaps the second part of the Commonitorium, which, except the last chapters, has been lost, was specially directed against the Augustinian doctrine of predestination, and was on this account destroyed, while the first part acquired almost canonical authority in the Catholic church.} Faustus, bishop of Rhegium (Riez), who at the council of Arles (475) refuted the hyper-Augustinian presbyter Lucidus, and was commissioned by the council to write a work upon the grace of God and human freedom;\footnote{De gratia Dei et humanae mentis libero arbitrio (in the Biblioth. maxima Patrum, tom. viii.). This work is regarded as the ablest defence of Semi-Pelagianism written in that age. Comp. upon it Wiggers, ii. p. 224 ff.} Gennadius, presbyter at Marseilles (died after 495), who continued the biographical work of Jerome, De viris illustribus, down to 495, and attributed Augustine’s doctrine of predestination to his itch for writing;\footnote{De viris illustr. c. 38, where he speaks in other respects eulogistically of Augustine. He refers to the passage in Prov. x. 19: “In multiloquio non fugies peccatum.” Comp. respecting him Wiggers, ii. 350 ff. and Neander, Dogmengeschichte, i. p. 406. His works are found in Migne’s Patrol. vol. 58.} Arnobius the younger;\footnote{In his Commentarius in Psalmos, written about 460, especially upon Ps. cxxvii.: “Nisi Dominus aedificaverit domum.” Some, following Sirmond, consider him as the author of the next-mentioned treatise \emph{Praedestinatus}, but without good ground. Comp. Wiggers, ii. p. 348 f.} and the much discussed anonymous tract Praedestinatus (about 460), which, by gross exaggeration, and by an unwarranted imputation of logical results which Augustine had expressly forestalled, placed the doctrine of predestination in an odious light, and then refuted it.\footnote{“Praedestinatus, seu Praedestinatorum haeresis, et libri S. Augustino temere adscripti refutatio.” The haeresis Praedestinatorum is the last of ninety heresies, and consists in the assertion: “Dei praedestinatione peccata committi.” This work was first discovered by J. Sirmond and published at Paris in 1643 (also in Gallandi, Biblioth. tom. x. p. 359 sqq., and in Migne’s Patrol. tom. liii. p. 587 sqq., together with Sirmond’s Historia Praedestinatiana). It occasioned in the seventeenth century a lively controversy between the Jesuits and the Jansenists, as to whether there had existed a distinct sect of Praedestinarians. The author, however, merely feigned such a sect to exist, in order to avoid the appearance of attacking Augustine’s authority. See details in Wiggers, ii. p. 329 ff.; Neander, Dogmengeschichte, i. 399 ff.; and Baur, p. 190 ff. The latter says: “The treatise [more accurately the second book of it; the whole consists of three books] is ascribed to Augustine, but as the ascription is immediately after declared false, both assertions are evidently made with the purpose of condemning Augustine’s doctrine with its consequences (only not directly in his name), as one morally most worthy of reprobation.” Neander ascribes only the first and the third book, Baur also the second book, to a Semi-Pelagian.}

The author of the Praedestinatus says, that a treatise had fallen into his hands, which fraudulently bore upon its face the name of the Orthodox teacher Augustine, in order to smuggle in, under a Catholic name, a blasphemous dogma, pernicious to the faith. On this account he had undertaken to transcribe and to refute this work. The treatise itself consists of three books; the first, following Augustine’s book, De haeresibus, gives a description of ninety heresies from Simon Magus down to the time of the author, and brings up, as the last of them, the doctrine of a double predestination, as a doctrine which makes God the author of evil, and renders all the moral endeavors of men fruitless;\footnote{The first book has also been reprinted in the Corpus haereseolog. ed. F. Oehler, tom. i. Berol. 1856, pp. 233-268.} the second book is the pseudo-Augustinian treatise upon this ninetieth heresy, but is apparently merely a Semi-Pelagian caricature by the same author;\footnote{Just as the Capitula Gallorum and the Objectiones Vincentianae exaggerate Angustinianism, in order the more easily to refute it.} the third book contains the refutation of the thus travestied pseudo-Augustinian doctrine of predestination, employing the usual Semi-Pelagian arguments.

A counterpart to this treatise is found in the also anonymous work, De vocatione omnium gentium, which endeavors to commend Augustinianism by mitigation, in the same degree that the Praedestinatus endeavors to stultify it by exaggeration.\footnote{It is found among the works of Leo I. and also of Prosper Aquitanus, but deviates from the views of the latter. Comp. Quesnel’s learned Dissertationes de auctore libri de vocatione gentium, in the second part of his edition of Leo’s works, and also Wiggers, ii. p. 218 ff.} It has been ascribed to pope Leo I. († 461), of whom it would not be unworthy; but it cannot be supposed that the work of so distinguished a man could have remained anonymous. The author avoids even the term praedestinatio, and teaches expressly, that Christ died for all men and would have all to be saved; thus rejecting the Augustinian particularism. But, on the other hand, he also rejects the Semi-Pelagian principles, and asserts the utter inability of the natural man to do good. He unhesitatingly sets grace above the human will, and represents the whole life of faith, from beginning to end, as a work of unmerited grace. He develops the three thoughts, that God desires the salvation of all men; that no one is saved by his own merits, but by grace; and that the human understanding cannot fathom the depths of divine wisdom. We must trust in the righteousness of God. Every one of the damned suffers only the righteous punishment of his sins; while no saint can boast himself in his merits, since it is only of pure grace that he is saved. But how is it with the great multitude of infants that die every year without baptism, and without opportunity of coming to the knowledge of salvation? The author feels this difficulty, without, however, being able to solve it. He calls to his help the representative character of parents, and dilutes the Augustinian doctrine of original sin to the negative conception of a mere defect of good, which, of course, also reduces the idea of hereditary guilt and the damnation of unbaptized children. He distinguishes between a general grace which comes to man through the external revelation in nature, law, and gospel, and a special grace, which effects conversion and regeneration by an inward impartation of saving power, and which is only bestowed on those that are saved.

Semi-Pelagianism prevailed in Gaul for several decades. Under the lead of Faustus of Rhegium it gained the victory in two synods, at Arles in 472 and at Lyons in 475, where Augustine’s doctrine of predestination was condemned, though without mention of his name.

\xsection{Victory of Semi-Augustinianism. Council of Orange, A.D. 529}

§ 160. Victory of Semi-Augustinianism. Council of Orange, A.D. 529.

But these synods were only provincial, and were the cause of a schism. In North Africa and in Rome the Augustinian system of doctrine, though in a somewhat softened form, attained the ascendency. In the decree issued by pope Gelasius in 496 de libris recipiendis et non recipiendis (the beginning of an Index librorum prohibitorum), the writings of Augustine and Prosper Aquitanus are placed among books ecclesiastically sanctioned, those of Cassian and Faustus of Rhegium among the apocryphal or forbidden. Even in Gaul it found in the beginning of the sixth century very capable and distinguished advocates, especially in Avitus, archbishop of Vienne (490523), and Caesarius, archbishop of Arles (502–542). Associated with these was Fulgentius of Ruspe († 533), in the name of the sixty African bishops banished by the Vandals and then living in Sardinia.\footnote{He wrote De veritate praedestinationis et gratiae Dei, three libb. against Faustus. He uses in these the expression praedestinatio \emph{duplex}, but understands by the second praedestinatio the praedestination to damnation, not to sin, and censures those who affirmed a predestination to sin. Yet he expressly consigned to damnation all unbaptized children, even such as die in their mother’s womb. Comp. Wiggers, ii. p. 378.}

The controversy was stirred up anew by the Scythian monks, who in their zeal for the Monophysite theopaschitism, abhorred everything connected with Nestorianism, and urged first pope Hormisdas, and then with better success the exiled African bishops, to procure the condemnation of Semi-Pelagianism.

These transactions terminated at length in the triumph of a moderate Augustinianism, or of what might be called Semi-Augustinianism, in distinction from Semi-Pelagianism. At the synod of Orange (Arausio) in the year 529, at which Caesarius of Arles was leader, the Semi-Pelagian system, yet without mention of its adherents, was condemned in twenty-five chapters or canons, and the Augustinian doctrine of sin and grace was approved, without the doctrine of absolute or particularistic predestination.\footnote{Comp. the transactions of the Concilium Arausicanum, the twenty-five Capitula, and the Symbolum in the Opera Aug. ed. Bened. Appendix to tom. x. 157 sqq.; in Mansi, tom. viii. p. 712 sqq.; and in Hefele, ii. p. 704 ff. The Benedictine editors trace back the several Capitula to their sources in the works of Augustine, Prosper, and others.} A similar result was reached at a synod of Valence (Valencia), held the same year, but otherwise unknown.\footnote{The Acts of the synod of Valence, in the metropolitan province of Vienne, held in the same year or in 530, have been lost. Pagi, and the common view, place this synod \emph{after} the synod of Orange, Hefele, on the contrary (ii. 718), \emph{before} it. But we have no decisive data.}

The synod of Orange, for its Augustinian decisions in anthropology and soteriology, is of great importance. But as the chapters contain many repetitions (mostly from the Bible and the works of Augustine and his followers), it will suffice to give extracts containing in a positive form the most important propositions.

Chap. 1. The sin of Adam has not injured the body only, but also the soul of man.

2. The sin of Adam has brought sin and death upon all mankind.

3. Grace is not merely bestowed when we pray for it, but grace itself causes us to pray for it.

5. Even the beginning of faith, the disposition to believe, is effected by grace.

9. All good thoughts and works are God’s gift.

10. Even the regenerate and the saints need continually the divine help.

12. What God loves in us, is not our merit, but his own gift.

13. The free will weakened\footnote{“Arbitrium voluntatis in primo homine in \emph{infirmatum}“ (not “amissum”).} in Adam, can only be restored through the grace of baptism.

16. All good that we possess is God’s gift, and therefore no one should boast.

18. Unmerited grace precedes meritorious works.\footnote{There are then meritorious works. “Debetur merces bonis operibus, si fiant, sed gratia quae non debetur praecedit, ut fiant” Chap. 18 taken from Augustine’s Opus imperf. c. Jul. i. c. 133 and from the Sentences of Prosper Aquitanus, n. 297. But, on the other hand, Augustinealso says: “Merita nostra sunt Dei munera.”}

19. Even had man not fallen, he would have needed divine grace for salvation.

23. When man sins, he does his own will; when he does good, he executes the will of God, yet voluntarily.

25. The love of God is itself a gift of God.

To these chapters the synod added a Creed of anthropology and soteriology, which, in opposition to Semi-Pelagianism, contains the following five propositions:\footnote{In the Latin original, the Epilogus reads as follows (Aug. Opera, tom. x. Appendix, f. 159 sq.): \par “Ac sic secundum suprascriptas sanctarum scripturarum Bententias vel antiquerum patrum definitiones hoc, Deo propitiante, et praedicare debemus et credere, quod per peccatum primi hominis its inclinatum et attenuatum fuerit liberum arbitrium, ut nullus postea aut diligere Deum sicut oportuit, aut credere in Deum, aut operari propter Deum quod bonum eat, possit, nisi gratia cum et misericordia divina praevenerit. Unde Abel justo et Noe, et Abrahae, et Isaac, et Jacob, et omni antiquorum sanctorum multitudini illam praeclaram fidem, quam in ipsorum laude praedicat apostolus Paulus, non per bonum naturae, quod prius in Adam datum fuerat, sed per gratiam Dei credimus fuisse collatam. Quam gratiam etiam post adventum Domini omnibus qui baptizari desiderant, non in libero arbitrio haberi, sed Christi novimus simul et credimus largitate conferri, secundum illud quod jam supra dictum est, et quod praedicat Paulus apostolus: Vobis donatum est pro Christo non solum ut in eum credatis, sed etiam ut pro illopatiamini (Phil. i. 29); et illud: Deus qui caepit in vobis bonum opus, perficiet usque in diem Domini nostri Jesu Christi (Phil. i. 6); et illud: Gratia salvi facti estis per fidem, et hoc non ex vobis, Dei enim donum est (Ephes. ii. 8); et quod de se ipso ait apostolus: Misericordiam consecutus sum ut fidelis essem (1 Cor. vii. 29); non dixit quia eram, sed ut essem; et illud: Quid habes quod non accepisti? (1 Cor. iv. 7); et illud: Omne datum bonum et omne donum perfectum de sursum est, descendens a Patre luminum (Jac. i. 17); et illud: Nemo habet quidquam boni, nisi illi datum fuerit de super (Joann. iii. 23). Innumerabilia sunt sanctarum scripturarum testimonia quae possunt ad probandam gratiam proferri, sed brevitatis studio praetermissa sunt, quia et revera cui pauca non sufficiunt plura non proderunt. \par “ Hoc etiam secundum fidem catholicam credimus, quod accepta per baptismum gratia, omnes baptizati, Christo auxilante et coöperante, quae ad salutem animae pertinent, possint et debeant, si fideliter laborare voluerint, adimplere. \par “Aliquos vero ad malum divina potestate praedestinatos esse non serum non credimus, sed etiam si sunt, qui tantum malum credere velint cum omni detestatione illis anathema dicimus. \par Hoc etiam salubriter profitemur et credimus, quod in omni opere bono non nos incipimus et postea per Dei misericordiam adjuvamur, sed ipse nobis, nullis praecedentibus bonis meritis, et fidem et amorem sui prius inspirat, ut et baptismi sacramenta fideliter requiramus, et post baptismum cum ipsius adjutorio ea quae sibi sunt placita implere possimus. Unde manifestissime credendum est, quod et illius latronis, quem Dominus ad paradisi patriam revocavit, et Cornelii centurionis, ad quem angelus Domini missus est, et Zachaei, qui ipsum Dominum suscipere meruit, illa tam admirabilis fides non fuit de natum, sed divinae largitatis donum. \par “ Et quia definitionem antiquoram patrum nostamque, quae suprascripta est, non solum religiosis, sed etiam laicis medicamentum esse, et desideramus et cupimus: placuit ut eam etiam illustres ac magnifici viri, qui nobiscum ad praefatam festivitatem convenerunt, propria manu subscriberent.” \par Then follow the names of fourteen bishops (headed by Caesarius) and eight laymen (headed by Petrus Marcellinus Felix Liberius, vir clarissimus et illustris Praefectus Praetorii Galliarum atque Patricius).}

1. Through the fall free will has been so weakened, that without prevenient grace no one can love God, believe on Him, or do good for God’s sake, as he ought (sicut oportuit, implying that he may in a certain measure).

2. Through the grace of God all may, by the co-operation of God, perform what is necessary for their soul’s salvation.

3. It is by no means our faith, that any have been predestinated by God to sin (ad malum), but rather: if there are people who believe so vile a thing, we condemn them with utter abhorrence (cum omni detestatione).\footnote{This undoubtedly takes for granted, that Augustinedid not teach this; and in fact he taught only a predestination of the wicked to perdition, not a predestination to sin.}

4. In every good work the beginning proceeds not, from us, but God inspires in us faith and love to Him without merit precedent on our part, so that we desire baptism, and after baptism can, with His help, fulfil His will.

5. Because this doctrine of the fathers and the synod is also salutary for the laity, the distinguished men of the laity also, who have been present at this solemn assembly, shall subscribe these acts.

In pursuance of this requisition, besides the bishops, the Praefectus praetorio Liberius, and seven other viri illustres, signed the Acts. This recognition of the lay element, in view of the hierarchical bent of the age, is significant, and indicates an inward connection of evangelical doctrine with the idea of the universal priesthood. And they were two laymen, we must remember, Prosper and Hilarius, who first came forward in Gaul in energetic opposition to Semi-Pelagianism and in advocacy of the sovereignty of divine grace.

The decisions of the council were sent by Caesarius to Rome, and were confirmed by pope Boniface II. in 530. Boniface, in giving his approval, emphasized the declaration, that even the beginning of a good will and of faith is a gift of prevenient grace, while Semi-Pelagianism left open a way to Christ without grace from God. And beyond question, the church was fully warranted in affirming the pre-eminence of grace over freedom, and the necessity and importance of the gratia praeveniens.

Notwithstanding this rejection of the Semi-Pelagian teachings (not teachers), they made their way into the church again, and while Augustine was universally honored as a canonized saint and standard teacher, Cassian and Faustus of Rhegium remained in grateful remembrance as saints in France.\footnote{Comp. respecting the further history of anthropology Wiggers: Schicksale der augustinischen Anthropologie von der Verdammung des Semipelagianismus auf den Synoden zu Orange und Valence, 529, bis zur Reaction des Mönchs Gottschalk für den Augustinimus, in Niedner’s “Zeitschrift für Hist. Theologie,” 1854, p. 1 ff.}

At the close of this period Gregory the Great represents the moderated Augustinian system, with the gratia praeveniens, but without the gratia irresistibilis and without a particularistic decretum absolutum. Through him this milder Augustinianism exerted great influence upon the mediaeeval theology. Yet the strict Augustinianism always had its adherents, in such men as Bede, Alcuin, and Isidore of Seville, who taught a gemina praedestinatio, sive electorum ad salutem, sive reproborum ad mortem; it became prominent again in the Gottschalk controversy in the ninth century, was repressed by scholasticism and the prevailing legalism; was advocated by the precursors of the Reformation, especially by Wiclif and Huss; and in the Reformation of the sixteenth century, it gained a massive acknowledgment and an independent development in Calvinism, which, in fact, partially recast it, and gave it its most consistent form.

\xchapter{Church Fathers, and Theological Literature}

CHURCH FATHERS, AND THEOLOGICAL LITERATURE.

Comp. the general literature on the Fathers in vol. i. § 116, and the special literature in the several sections following.

I.—The Greek Fathers.

\xsection{Eusebius of Caesarea}

§ 161. Eusebius of C sarea.

I. Eusebius Pamphili: Opera omnia Gr. et Lat., curis variorum nempe II. Valesii, Fr. Vigeri, B. Montfaucon, Card. Angelo Maii edita; collegit et denuo recognovit J. P. Migne. Par. (Petit-Montrouge) 1857. 6 vols. (tom. xix.-xxiv. of Migne’s Patrologia Graeca). Of his several works his Church History has been oftenest edited, sometimes by itself, sometimes in connection with his Vita Constantini, and with the church histories of his successors; best by Henr. Valesius (Du Valois), Par. 1659–’73, 8 vols., and Cantabr. 1720, 3 vols., and again 1746 (with additions by G. Reading, best ed.); also (without the later historians) by E. Zimmermann, Francof. 1822; F. A. Heinichen, Lips. 1827–’8, 3 vols.; E. Burton, Oxon. 1838, 2 vols. (1845 and 1856 in 1 vol.); Schwegler, Tüb. 1852; also in various translations: In German by Stroth, Quedlinburg, 1776 ff., 2 vols.; by Closs, Stuttg. 1839; and several times in French and English; in English by Hanmer (1584), T. Shorting, and better by Chr. Fr. Cruse (an Amer. Episcopalian of German descent, died in New York, 1865): The Ecclesiastical History of Euseb. Pamph., etc., Now York, 1856 (10th ed.), and Lond. 1858 (in Bohn’s Eccles. Library). Comp. also the literary notices in Brunet, sub Euseb., and James Darling, Cyclop. Bibliograph. p. 1072 ff.

II. Biographies by Hieronymus (De viris illustr. c. 81, a brief sketch, with a list of his works), Valesius (De vita scriptisque Eusebii Caesar.), W. Cave (Lives of the most eminent Fathers of the Church, vol. ii. pp. 95–144, ed. H. Cary, Oxf. 1840), Heinichen, Stroth, Cruse, and others, in their editions of the Eccles. Hist. of Eusebius. F. C. Baur: Comparatur Eusebius Hist. eccl. parens cum parente Hist. Herodoto. Tub. 1834. Haenell: De Euseb. Caes. religionis christ. defensore. Gott. 1843. Sam. Lee: Introductory treatise in his Engl. edition of the Theophany of Eusebius, Cambr. 1843. Semisch: Art. Eusebius v. Caes. in Herzog’s Encycl. vol. iv. (1855), pp. 229–238. Lyman Coleman: Eusebius as an historian, in the Bibliotheca Sacra, Andover, 1858, pp. 78–96. (The biography by Acacius, his successor in the see of Caesarea, Socr. ii. 4, is lost.)

This third period is uncommonly rich in great teachers of the church, who happily united theological ability and practical piety, and who, by their development of the most important dogmas in conflict with mighty errors, earned the gratitude of posterity. They monopolized all the learning and eloquence of the declining Roman empire, and made it subservient to the cause of Christianity for the benefit of future generations. They are justly called fathers of the church; they belong to Christendom without distinction of denominations; and they still, especially Athanasius and Chrysostom among the Greek fathers, and Augustine and Jerome among the Latin, by their writings and their example, hold powerful sway, though with different degrees of authority according to the views entertained by the various churches concerning the supremacy of the Bible and the value of ecclesiastical tradition.

We begin the series of the most important Nicene and post-Nicene divines with Eusebius of Caesarea, the “father of church history,” the Christian Herodotus.

He was born about the year 260 or 270, probably in Palestine, and was educated at Antioch, and afterwards at Caesarea in Palestine, under the influence of the works of Origen. He formed an intimate friendship with the learned presbyter Pamphilus,\footnote{Hence the surname \textgreek{Εὐσέβιος} (\textgreek{ὁ φίλος}) \textgreek{τοῦ Παμφίλου}, Pamphhili, by which anciently he was most frequently distinguished from many other less noted men of the same name, e.g.: Eusebius of Nicomedia († 341), Eusebius of Vercelli († 371), Eusebius Emesenus, of Emesa or Emisa in Phoenicia († 360), and others. On this last comp. Opuscula quae supersunt Graeca, ed. Augusti, Elberfeld, 1829, somewhat hastily; corrected by Thilo, Ueber die Schriften des Euseb. von Alex. und des Euseb. von Emisa, Halle, 1832.} who had collected a considerable biblical and patristic library, and conducted a flourishing theological school which he had founded at Caesarea, till in 309 he died a martyr in the persecution under Diocletian.\footnote{Jeromeremarks of Pamphilus (De viris illustribus, c. 75): “Tanto bibliothecae divinae amore flagravit, ut maximam partem Origenis voluminum sua manu descripserit, quae usque hodie [a. 392] in Caesariensi bibliotheca habentur.”} Eusebius taught for a long time in this school; and after the death of his preceptor and friend, he travelled to Tyre and to Egypt, and was an eye-witness of the cruel scenes of the last great persecution of the Christians. He was imprisoned as a confessor, but soon released.

Twenty years later, when Eusebius, presiding at the council at Tyre (335 or 336), took sides against Athanasius, the bishop Potamon of Hieraclea, according to the account of Epiphanius, exclaimed in his face: “How dost thou, Eusebius, sit as judge of the innocent Athanasius? Who can bear it? Why! didst thou not sit with me in prison in the time of the tyrants? They plucked out my eye for my confession of the truth; thou camest forth unhurt; thou hast suffered nothing for thy confession; unscathed thou art here present. How didst thou escape from prison? On some other ground than because thou didst promise to do an unlawful thing [to sacrifice to idols]? or, perchance, didst thou actually do this? “But this insinuation of cowardice and infidelity to Christ arose probably from envy and party passion in a moment of excitement. With such a stain upon him, Eusebius would hardly have been intrusted by the ancient church with the episcopal staff.\footnote{So Valesius also views the matter, while Baronius puts faith in the rebuke.}

About the year 315, or earlier, Eusebius was chosen bishop of Caesarea,\footnote{Hence he is also called Eusebius Caesariensis or Palestinus.} where he labored till his death in 340. The patriarchate of Antioch, which was conferred upon him after the deposition of Eustathius in 331, he in honorable self-denial, and from preference for a more quiet literary life, declined.

He was drawn into the Arian controversies against his will, and played an eminent part at the council of Nicaea, where he held the post of honor at the right hand of the presiding emperor. In the perplexities of this movement he took middle ground, and endeavored to unite the opposite parties. This brought him, on the one hand, the peculiar favor of the emperor Constantine, but, on the other, from the leaders of the Nicene orthodoxy, the suspicion of a secret leaning to the Arian heresy.\footnote{So thought, among the ancients, Hilary, Jerome(who otherwise speaks favorably of Eusebius), Theodoret, and the second council of Nicaea (a.d.787), which unjustly condemned him even expressly, as an Arian heretic; and so have thought, among modems, Baronius, Petavius, Clericus, Tillemont, Gieseler; while the church historian Socrates, the Roman bishops Gelasius and Pelagius II., Valesius, G. Bull, Cave (who enters into a fall vindication, l.c. p. 135 sqq.), and Sam. Lee (and most Anglicans), have defended the orthodoxy of Eusebius, or at least mention him with very high respect. The Gallican church has even placed him in the catalogue of saints. Athanasius never expressly charges him with apostasy from the Nicene faith to Arianism or to Semi-Arianism, but frequently says that before 325 he held with Arius, and changed his opinion at Nicaea. This is the view of Möhler also (Athanasius der Grosse, p. 333 ff.), whom Dorner (History of Christology, i. 792) inaccurately reckons among the opponents of the orthodoxy of Eusebius. The testimonies of the ancients for and against Eusebius are collected in Migne’s edition of his works, tom. i. pp. 68-98. Among recent writers Dr. Samuel Lee has most fully investigated the orthodoxy of Eusebius in the Preliminary Dissertation to his translation of the Theophania from the Syriac, pp. xxiv.-xcix. He arrives at the conclusion (p. xcviii.), “that Eusebius was no Arian; and that the same reasoning must prove that he was no Semi-Arian; that he did in no degree partake of the error of Origen, ascribed to him so positively and so groundlessly by Photius.” But this is merely a negative result.} It is certain that, before the council of Nicaea, he sympathized with Arius; that in the council he proposed an orthodox but indefinite compromise-creed; that after the council he was not friendly with Athanasius and other defenders of orthodoxy; and that, in the synod of Tyre, which deposed Athanasius in 335, he took a leading part, and, according to Epiphanius, presided. In keeping with these facts is his silence respecting the Arian controversy (which broke out in 318) in an Ecclesiastical History which comes down to 324, and was probably not completed till 326, when the council of Nicaea would have formed its most fitting close. He would rather close his history with the victory of Constantine over Licinius than with the Creed over which theological parties contended, and with which he himself was implicated. But, on the other hand, it is also a fact that he subscribed the Nicene Creed, though reluctantly, and reserving his own interpretation of the homoousion; that he publicly recommended it to the people of his diocese; and that he never formally rejected it.

The only satisfactory solution of this apparent inconsistency is to be found in his own indecision and leaning to a doctrinal latitudinarianism, not unfrequent in historians who become familiar with a vast variety of opinions in different ages and countries. On the important point of the homoousion he never came to a firm and final conviction. He wavered between the older Origenistic subordinationism and the Nicene orthodoxy. He asserted clearly and strongly with Origen the eternity of the Son, and so far was decidedly opposed to Arianism, which made Christ a creature in time; but he recoiled from the homoousion, because it seemed to him to go beyond the Scriptures, and hence he made no use of the term, either in his book against Marcellus, or in his discourses against Sabellius. Religious sentiment compelled him to acknowledge the full deity of Christ; fear of Sabellianism restrained him. He avoided the strictly orthodox formulas, and moved rather in the less definite terms of former times. Theological acumen he constitutionally lacked. He was, in fact, not a man of controversy, but of moderation and peace. He stood upon the border between the ante-Nicene theology and the Nicene. His doctrine shows the color of each by turns, and reflects the unsettled problem of the church in the first stage of the Arian controversy.\footnote{The same view is taken substantially by Baur (Geschichte der Lehre von der Dreieinigkeit und Menschwerdung, i. 475 ff.), Domer (Entwicklungsgeschichte der Lehre von der Person Christi, i. 792 ff), Semisch (Art. Eusebius in Herzog’s Encyklopädie, vol iv. 233), and other modem German theologians.}

With his theological indecision is connected his weakness of character. He was an amiable and pliant court-theologian, and suffered himself to be blinded and carried away by the splendor of the first Christian emperor, his patron and friend. Constantine took him often into his counsels, invited him to his table, related to him his vision of the cross, showed him the famous labarum, listened standing to his occasional sermons, wrote him several letters, and intrusted to him the supervision of the copies of the Bible for the use of the churches in Constantinople.

At the celebration of the thirtieth anniversary of this emperor’s reign (336), Eusebius delivered a panegyric decked with the most pompous hyperbole, and after his death, in literal obedience to the maxim: “De mortuis nihil nisi bonum,” he glorified his virtues at the expense of veracity and with intentional omission of his faults. With all this, however, he had noble qualities of mind and heart, which in more quiet times would have been an ornament to any episcopal see. And it must be said, to his honor, that he never claimed the favor of the emperor for private ends.

The theological and literary value of Eusebius lies in the province of learning. He was an unwearied reader and collector, and probably surpassed all the other church fathers, hardly excepting even Origen and Jerome, in compass of knowledge and of acquaintance with Grecian literature both heathen and Christian; while in originality, vigor, sharpness, and copiousness of thought, he stands far below Origen, Athanasius, Basil, and the two Gregories. His scholarship goes much further in breadth than in depth, and is not controlled and systematized by a philosophical mind or a critical judgment.

Of his works, the historical are by far the most celebrated and the most valuable; to wit, his Ecclesiastical History, his Chronicle, his Life of Constantine, and a tract on the Martyrs of Palestine in the Diocletian persecution. The position of Eusebius, at the close of the period of persecution, and in the opening of the period of the imperial establishment of Christianity, and his employment of many ancient documents, some of which have since been lost, give these works a peculiar value. He is temperate, upon the whole, impartial, and truth-loving—rare virtues in an age of intense excitement and polemical zeal like that in which he lived. The fact that he was the first to work this important field of theological study, and for many centuries remained a model in it, justly entitles him to his honorable distinction of Father of Church History. Yet he is neither a critical student nor an elegant writer of history, but only a diligent and learned collector. His Ecclesiastical History, from the birth of Christ to the victory of Constantine over Licinius in 324, gives a colorless, defective, incoherent, fragmentary, yet interesting picture of the heroic youth of the church, and owes its incalculable value, not to the historic art of the author, but almost entirely to his copious and mostly literal extracts from foreign, and in some cases now extinct, sources. As concerns the first three centuries, too, it stands alone; for the successors of Eusebius begin their history where he leaves off.

His Chronicle consists of an outline-sketch of universal history down to 325, arranged by ages and nations (borrowed largely from the Chronography of Julius Africanus), and an abstract of this universal chronicle in tabular form. The Greek original is lost, with the exception of unconnected fragments by Syncellus; but the second part, containing the chronological tables, was translated and continued by Jerome to 378, and remained for centuries the source of the synchronistic knowledge of history, and the basis of historical works in Christendom.\footnote{The Greek title was: \textgreek{Χρονικῶν κανόνων παντοδαπὴ ἱστορία}(Hieron. De viris illustr. c. 81); the Latin is: Chronica Eusebii s. Canones historiae universae, Hieronymo interprete. See Vallarsi’s ed. of Jerome’s works, tom. viii. 1-820. Jeromealso calls it Temporum librum. It is now known also (since 1818) in an Armenian translation. Most complete edition by Angelo Mai, in Script vet. nova coll. tom. viii. Rom. 1833, republished in Migne’s edition of the complete works of Eusebius, tom. i. p. 100 sqq.} Jerome also translated, with several corrections and additions, a useful antiquarian work of Eusebius, the so-called Onomasticon, a description of the places mentioned in the Bible.\footnote{\textgreek{Περὶ τῶν τοπικῶν ὀνομάτων τῶν ἐν τῇ θεία γραφῇ}, De situ et nominibus locorum Hebraicorum, in Jerome’s works, tom. iii. 121-290. A new edition, Greek and Latin, by Larsow and Parthey, Berol. 1862.}

In his Life, and still more in his Eulogy, of Constantine, Eusebius has almost entirely forgotten the dignity of the historian in the zeal of the panegyrist. Nevertheless, this work is the chief source of the history of the reign of his imperial friend.\footnote{Socrates already observes (in the first book of his Church History) that Eusebius wrote the Life of Constantinemore as a panegyrical oration than as an accurate account of events. Baronius (Annal. ad an. 324, n. 5) compares the Vita Constantini, not unfitly, with the Cyropaedia of Xenophon, who, as Cicero says, “vitam Cyri non tam ad historiae fidem conscripsit, quam ad effigiem justi principis exhibendam.” This is the most charitable construction we can put upon this book, the tone of which is intolerably offensive to a manly and independent spirit acquainted with the crimes of Constantine. But we should remember that stronger men, such as Athanasius, Hilary, and Epiphanius have overrated Constantine, and called him, “most pious,” and “of blessed memory.” Burckhardt, in his work on Constantine, p. 346 and passim, speaks too contemptuously of Eusebius, without any reference to his good qualities and great merits.}

Next in importance to his historical works are his apologetic; namely, his Praeparatio evangelica,\footnote{Best edited by Thomas Gaisford, Oxon. 1843, 4 vols. 8vo. In Migne’s edition it forms tom. iii.} and his Demonstratio evangelica.\footnote{Likewise edited by Gaisford, Oxf. 1852, 2 vols. 8vo. In Migne’s edition tom. !v.} These were both written before 324, and are an arsenal of the apologetic material of the ancient church. The former proposes, in fifteen books, to give a documentary refutation of the heathen religious from Greek writings. The latter gives, in twenty books, of which only the first ten are preserved, the positive argument for the absolute truth of Christianity, from its nature, and from the fulfilment of the prophecies in the Old Testament. The Theophany, in five books, is a popular compend from these two works, and was probably written later, as Epiphanius wrote his Anacephalaeosis after the Panarion, for more general use.\footnote{Dr. Sam. Lee, however, is of the opposite opinion, see p. xxii. of the Preface to his translation.“It appears probable to me,” he says, “that this more popular and more useful work [the Theophania] was first composed and published, and that the other two [the Praeparatio, and the Demonstratio Evangelica]—illustrating, as they generally do, some particular points only—argued in order in our work—were reserved for the reading and occasional writing of our author during a considerable number of years, as well for the satisfaction of his own mind, as for the general reading of the learned. It appears probable to me, therefore, that this was one of the first productions of Eusebius, if not the first after the persecutions ceased.”} It is known in the Greek original from fragments only, published by Cardinal Mai,\footnote{In the fourth volume of the Novae Patrum Bibliothecae, Rom. 1847, pp. 108-156, reprinted in Migne’s edition of the works of Eusebius tom. v. 609 sqq.} and now complete in a Syriac version which was discovered in 1839 by Tattam, in a Nitrian monastery, and was edited by Samuel Lee at London in 1842.\footnote{Also in English, under the title: On the Theophania, or Divine Manifestation of our Saviour Jesus Christ, by Eusebius, translated into English, with Notes, from an ancient Syriac Version of the Greek original, now lost; to which is prefixed a Vindication of the orthodoxy, and prophetical views, of that distinguished writer, by Sam. Lee, D. D., Cambr. 1848. The MS. of this work is deposited in the British Museum; it was written at Edessa in the Estranghelo, or old church-handwriting of the Syrians, on very fine and well-prepared skin. Dr. Lee assigns it to the year 411 (I. c. p. xii.).} To this class also belongs his apologetic tract Against Hierocles.\footnote{In Migne’s edition, tom. iv. 195-868.}

Of much less importance are the two dogmatic works of Eusebius: Against Marcellus, and Upon the Church Theology (likewise against Marcellus), in favor of the hypostatical existence of the Son.\footnote{In Migne’s edition, tom vi. p. 107 sqq.}

His Commentaries on several books of the Bible (Isaiah, Psalms, Luke) pursue, without independence, and without knowledge of the Hebrew, the allegorical method of Origen.\footnote{Angelo Mai has published new fragments of Commentaries of Eusebius on the Psalms and on the Gospel of Luke in Novae Patrum Bibliothecae tom. iv. p. 77 sqq. and p. 160 sqq., and republished in Migne’s ed. vol. vi.}

To these are to be added, finally, some works in Biblical Introduction and Archaeology, the Onomasticon, already alluded to, a sort of sacred geography, and fragments of an enthusiastic Apology for Origen, a juvenile work which he and Pamphilus jointly produced before 309, and which, in the Origenistic controversy, was the target of the bitterest shots of Epiphanius and Jerome.\footnote{The sixth book was added by Eusebius alone after the death of his friend. The first book is still extant in the Latin version of Rufinus, and some extracts in Photius.}

\xsection{The Church Historians after Eusebius}

§ 162. The Church Historians after Eusebius.

I. The Church Histories of Socrates, Sozomen, Theodoret, Evagrius, Philostorgius, and Theodorus Lector have been edited, with the Eccles. Hist. of Eusebius, by Valesius, Par. 1659–’73, in 3 vols. (defective reprint, Frankf. a. M. 1672–’79); best ed., Cambridge, 1720, and again 1746, in 3 vols., with improvements and additions by Guil. Reading. Best English translation by Meredith, Hanmer, and Wye Saltonstall, Cambr. 1688, 1692, and London, 1709. New ed. in Bohn’s Ecclesiastical Library, Land. 1851, in 4 vols. small 8vo.

II. F. A. Holzhausen: De fontibus, quibus Socrates, Sozomenus, ac Theodoretus in scribenda historia sacra usi sunt. Gött. 1825. G. Dangers: De fontibus, indole et dignitate librorum Theod. Lectoris et Evagrii. Gött. 1841. J. G. Dowling: An Introduction to the Critical Study of Eccl. History. Lond. 1838, p. 84 ff. F. Chr. Baur: Die Epochen der kirchlichen Geschichtschreibung. Tüb. 1852, pp. 7–32. Comp. P. Schaff: History of the Apostolic Church, Gen. Introd. p. 52 f.

Eusebius, without intending it, founded a school of church historians, who continued the thread of his story from Constantine the Great to the close of the sixth century, and, like him, limited themselves to a simple, credulous narration of external facts, and a collection of valuable documents, without an inkling of the critical sifting, philosophical mastery, and artistic reproduction of material, which we find in Thucydides and Tacitus among the classics, and in many a modern historian. None of them touched the history of the first three centuries; Eusebius was supposed to have done here all that could be desired. The histories of Socrates, Sozomen, and Theodoret run nearly parallel, but without mutual acquaintance or dependence, and their contents are very similar.\footnote{The frequent supposition (of Valois with others) that Sozomen wrote to complete Socrates, and Theodoret to complete both, cannot be proved. The authors seem independent of one another. Theodoret says in the Prooemium: “Since Eusebius of Palestine, commencing his history with the holy apostles, has described the events of the church to the reign of the God-beloved Constantine, I have begun my history where he ended his.” He makes no mention of any other writers on the same subject. Nor does Sozomen, l. i. c. 1, where he alludes to his predecessors. Valesius charges Sozomen with plagiarism.} Evagrius carried the narrative down to the close of the sixth century. All of them combine ecclesiastical and political history, which after Constantine were inseparably interwoven in the East; and (with the exception of Philostorgius) all occupy essentially the same orthodox stand-point. They ignore the Western church, except where it comes in contact with the East.

These successors of Eusebius are:

Socrates, an attorney or scholasticus in Constantinople, born in 380. His work, in seven books, covers the period from 306 to 439, and is valuable for its numerous extracts from sources, and its calm, impartial representation. It has been charged with a leaning towards Novatianism. He had upon the whole a higher view of the duty of the historian than his contemporaries and successors; he judged more liberally of heretics and schismatics, and is less extravagant in the praise of emperors and bishops.\footnote{Separate edition by Hussey: Socratis scholastici Historia Eccl. Oxon. 1853, 3 vols. 8vo.}

Hermias Sozomen, a native of Palestine, a junior contemporary of Socrates, and likewise a scholasticus in Constantinople, wrote the history of the church, in nine books, from 323 to the death of Honorius in 423,\footnote{According to the usual, but incorrect statement, to the year 489.} and hence in its subjects keeps pace for the most part with Socrates, though, as it would appear, without the knowledge of his work, and with many additions on the history of the hermits and monks, for whom he had a great predilection.\footnote{He informs us (Book v. c. 15) that his grandfather, with his whole family, was converted to Christianity by a miracle of the monk Hilarion.}

Theodoret, bishop of Cyrus, was born at Antioch about 390, of an honorable and pious mother; educated in the cloister of St. Euprepius (perhaps with Nestorius); formed upon the writings of Diodorus of Tarsus and Theodore of Mopsuestia; made bishop of Cyros, or Cyrrhos, in Syria, after 420; and died in 457. He is known to us from the Christological controversies as the most scholarly advocate of the Antiochian dyophysitism or moderate Nestorianism; condemned at Ephesus in 431, deposed by the council of Robbers in 449, acquitted in 451 by the fourth ecumenical council on condition of his condemning Nestorius and all deniers of the theotokos, but again partially condemned at the fifth long after his death. He was, therefore, like Eusebius, an actor as well as an author of church history. As bishop, he led an exemplary life, his enemies themselves being judges, and was especially benevolent to the poor. He owned nothing valuable but books, and applied the revenues of his bishopric to the public good. He shared the superstitions and weaknesses of his age.

His Ecclesiastical History, in five books, composed about 450, reaches from 325 to 429. It is the most valuable continuation of Eusebius, and, though shorter, it furnishes an essential supplement to the works of Socrates and Sozomen.

His “Historia religiosa” consists of biographies of hermits and monks, written with great enthusiasm for ascetic holiness, and with many fabulous accessories, according to the taste of the day. His “Heretical Fables,”\footnote{\textgreek{Αἱρετικῆς κακομυθίας ἐπιτομή,} in five books; in Schulze’s edition of the Opera, tom. iv. p. 280 sqq. The fifth book presents a summary of the chief articles of the orthodox faith, a sort of dogmatical compend.} though superficial and marred by many errors, is of some importance for the history of Christian doctrine. It contains a severe condemnation of Nestorius, which we should hardly expect from Theodoret.\footnote{Book iv. ch. 12. Garnier, Cave, and Oudin regard this anti-Nestorian chapter as a later interpolation, though without good reason; Schulze (note in loco, tom. iv. p. 368) defends it as genuine. It should be remembered that Theodoret at the council of Chalcedon could only save himself from expulsion by anathematizing Nestorius.}

Theodoret was a very fruitful author. Besides these histories, he wrote valuable commentaries on most of the books of the Old Testament and on all the Epistles of Paul; dogmatic and polemic works against Cyril and the Alexandrian Christology, and against the heretics; an apology of Christianity against the Greek philosophy; and sermons and letters.\footnote{TheodoretiOpera omnia cum et studio Jac. Sirmondi, Par. 1642, 4 vols. fol., with an additional vol. v. by Gamier, 1684. Another edition by \emph{J. L. Schulze}, Halle, 1768-’74, 5 tom. in 10 vols., which has been republished by \emph{J. P. Migne}, Par. 1860, in 5 vols. (Patrologia Graeca, tom. lxxx.-lxxxiv.). The last volume in Schulze’s and Migne’s editions contains Garnier’s Auctarium ad opera Theod. and his Dissertations on the life and on the faith of Theodoret, and on the fifth ecumenical Synod. Comp. also Schröckh, Church History, vol. xviii.}

Evagrius (born about 536 in Syria, died after 594) was a lawyer in Antioch, and rendered the patriarch Gregory great service, particularly in an action for incest in 588. He was twice married, and the Antiochians celebrated his second wedding (592) with public plays. He is the last continuator of Eusebius and Theodoret, properly so called. He begins his Ecclesiastical History of six books with the council of Ephesus, 431, and closes it with the twelfth year of the reign of the emperor Maurice, 594. He is of special importance on the Nestorian and Eutychian controversies; gives accounts of bishops and monks, churches and public buildings, earthquakes and other calamities; and interweaves political history, such as the wars of Chosroes and the assaults of the barbarians.\footnote{VaIesius blames him “quod non tantam diligentiam adhibuit in conquirendis antiquitatis ecclesiasticae monumentis, quam in legendis profanis auctoribus.”} He was strictly orthodox, and a superstitious venerator of monks, saints, and relics.\footnote{The first edition was from a Parisian manuscript by Rob. Stephanus, Par. 1544. Valesius, in his complete edition, employed two more manuscripts. A new edition, according to the text of Valesius, appeared at Oxford in 1844.}

Theodorus Lector, reader in the church of Constantinople about 525, compiled an abstract from Socrates, Sozomen, and Theodoret, under the title of Historia tripartita, which is still extant in the manuscript;\footnote{Valesius intended to edit it, and contented himself with giving the variations, since the book furnished nothing new.} and composed a continuation of Socrates from 431 to 518, of which fragments only are preserved in John Damascenus, Nilus, and Nicephorus Callisti.\footnote{Collected in the edition of Valesius.}

Of Philostorgius, an Arian church historian (born in 368), nothing has come down to us but fragments in Photius; and these breathe so strong a partisan spirit, that the loss of the rest is not to be regretted. He described the period from the commencement of the Arian controversy to the reign of Valentinian III. a.d. 423.

The series of the Greek church historians closes with Nicephorus Callistus or Callisti (i.e., son of Callistus),\footnote{Not to be confounded with Nicephorus, patriarch of Constantinople, who was deposed during the image controversy and died 828. His works, among which is also a brief Chronographia ab Adamo ad Michaelis et Theophili tempora (828), form tom. c. in Migne’s Patrologia Graeca.} who lived at Constantinople in the fifteenth century. He was surprised that the voice of history had been silent since the sixth century, and resumed the long-neglected task where his predecessors had left it, but on a more extended plan of a general history of the catholic church from the beginning to the year 911. We have, however, only eighteen books to the death of emperor Phocas in 610, and a list of contents of five other books. He made large use of Eusebius and his successors, and added unreliable traditions of the later days of the Apostles, the history of Monophysitism, of monks and saints, of the barbarian irruptions, \&c. He, too, ignores the Pelagian controversy, and takes little notice of the Latin church after the fifth century.\footnote{First edition in Latin by \emph{John Lange}, Basil. 1658; in Greek and Latin by \emph{Front. Ducaeus}, Par. 1630, in 2 vols. There exists but one Greek manuscript copy of Nicephorus, as far as we know, which is in the possession of the imperial library of Vienna.}

In the Latin church—to anticipate thus much—Eusebius found only one imitator and continuator, the presbyter and monk Rufinus, of Aquileia (330–410). He was at first a friend of Jerome, afterwards a bitter enemy. He translated, with abridgments and insertions at his pleasure, the Ecclesiastical History of Eusebius, and continued it to Theodosius the Great (392). Yet his continuation has little value. He wrote also biographies of hermits; an exposition of the Apostles’ Creed; and translations of several works of Origen, with emendations of offensive portions.\footnote{His works are edited by \emph{Vallarsi}, Veron. 1745, vol. i. fol. (unfinished). The Ecclesiastical History has several times appeared separately, and was long a needed substitute for Eusebius in the West.}

Cassiodorus, consul and monk (died about 562), composed a useful abstract of the works of Socrates, Sozomen, and Theodoret, in twelve books, under the title of Historia tripartita, for the Latin church of the middle age.

The only properly original contributions to church history from among the Latin divines were those of Jerome († 419) in his biographical and literary Catalogue of Illustrious Men (written in 392), which Gennadius, a Semi-Pelagian presbyter of South Gaul, continued to the year 495. Sulpicius Severus († 420) wrote in good style a Sacred History, or History of the Old and New Testament, from the creation down to the year 400; and Paulus Orosius (about 415) an apologetic Universal History, which hardly, however, deserves the name of a history.

\xsection{Athanasius the Great}

§ 163. Athanasius the Great.

I. S. Athanasius: Opera omnia quae extant vel quae ejus nomine circumferuntur, etc., Gr. et lat., opera et studio monachorum ordinis S. Benedicti e congregatione S. Mauri (Jac. Lopin et B. de, Montfaucon). Paris, 1698. 3 tom. fol. (or rather 2 tomi, the first in two parts). This is the most elegant and correct edition, but must be completed by two volumes of the Collectio nova Patrum, ed. B. de Montfaucon. Par. 1706. 2 tom. fol. More complete, but not so handsome, is the edition of 1777, Patav., in 4 vols. fol. (Brunet says of the latter “Édition moins belle et moins chère quo cello de Paris, mais augmentée d’un 4e vol., lequel renferme les opuscules de S. Athan., tirés de la Collectio nova du P. Montfaucon et des Anecdota de Wolf, et de plus l’interpretatio Psalmorum.”) But now both these older editions need again to be completed by the Syrian Festal Letters of Athanasius, discovered by Dr. Tattam in a Nitrian monastery in 1843; edited by W. Cureton in Syriac and English at London in 1846 and 1848 (and in English by H. Burgess and H. Williams, Oxf. 1854, in the Libr. of the Fathers); in German, with notes by F. Larson, at Leipzig in 1852; and in Syriac and Latin by Card. Angelo Mai in the Nova Patr. Bibliotheca, Rom. 1853, tom. vi. pp. 1–168. A new and more salable, though less accurate, edition of the Opera omnia Athan. (a reprint of the Benedictine) appeared at Petit-Montrouge (Par.) in J.P. Migne’s Patrologia Gr. (tom. xxv.-xxviii.), 1857, in 4 vols.

The more important dogmatic works of Athanasius have been edited separately by J. C. Thilo, in the first volume of the Bibliotheca Patrum Graec. dogmatica, Lips. 1853; and in an English translation, with explanations and indexes, by J. H. Newman, Oxf. 1842–’44 (Library of the Fathers, vols. 8, 13, 19).

II. Gregorius Naz.: Oratio panegyrica in Magnum Athanasium (Orat. xxi.). Several Vitae Athan. in the 1st vol. of the Bened. ed. of his Opera. Acta Sanctorum for May 2d. G. Hermant: La Vie d’Athanase, etc. Par. 1679. 2 vols. Tillemont: Mémoires, vol. viii. pp. 2–258 (2d ed. Par. 1713). W. Cave: Lives of the most eminent Fathers of the first Four Centuries, vol. ii. pp. 145–364 (Oxf. ed. of 1840). Schröckh: Th. xii. pp. 101–270. J. A. Möhler: Athanasius der Grosse und die Kirche seiner Zeit, besonders im Kampfe mit dem Arianismus. Mainz, 1827. 2d (title) ed. 1844. Heinrich Voigt: Die Lehre des heiligen Athanasius von Alexandria oder die kirchliche Dogmatik des 4ten Jahrhunderts auf Grund der biblischen Lehre vom Logos. Bremen, 1861. A. P. Stanley: Lectures on the History of the Eastern Church. New York, 1862, lecture vii. (pp. 322–358).

Athanasius is the theological and ecclesiastical centre, as his senior contemporary Constantine is the political and secular, about which the Nicene age revolves. Both bear the title of the Great; the former with the better right, that his greatness was intellectual and moral, and proved itself in suffering, and through years of warfare against mighty, errors and against the imperial court. Athanasius contra mundum, et mundus contra Athanasium, is a well-known sentiment which strikingly expresses his fearless independence and immovable fidelity to his convictions. He seems to stand an unanswerable contradiction to the catholic maxim of authority: Quod sem per, quod ubique, quod ab omnibus creditum est, and proves that truth is by no means always on the side of the majority, but may often be very unpopular. The solitary Athanasius even in exile, and under the ban of council and emperor, was the bearer of the truth, and, as he was afterwards named, the “father of orthodoxy.”\footnote{\textgreek{Ο πατήρ τῆς ὀρθοδοξίας .}. So Epiphanius already calls him, Haer. 69, c. 2.}

On a martyrs’ day in 313 the bishop Alexander of Alexandria saw a troop of boys imitating the church services in innocent sport, Athanasius playing the part of bishop, and performing baptism by immersion.\footnote{So Rufinus relates, H. E. l. i. c. 14. Most Roman historians, Hermant, Tillemont, Butler, and the author of the Vita Athan. in the Bened. ed. (tom. i. p. iii.), reject this legend, partly on account of chronological difficulty, partly because it seemed incompatible with the dignity of such a saint. Möhler passes it in silence.} He caught in this a glimpse of future greatness; took the youth into his care; and appointed him his secretary, and afterwards his archdeacon. Athanasius studied the classics, the Holy Scriptures, and the church fathers, and meantime lived as an ascetic. He already sometimes visited St. Anthony in his solitude.

In the year 325 he accompanied his bishop to the council of Nicaea, and at once distinguished himself there by his zeal and ability in refuting Arianism and vindicating the eternal deity of Christ, and incurred the hatred of this heretical party, which raised so many storms about his life.

In the year 328\footnote{This is the true date, according to the summaries of the newly-discovered Festal Letters of Athanasius, and not “a few weeks (or months rather] after the close of the council,” as the editor of the English translation of the historical tracts of Athanasius (Oxford Library of the Fathers, 1843, Preface, p. xxi), and even Stanley (I. c. p. 325), still say. The older hypothesis rests on a misapprehension of the \textgreek{πέντε μῆνες}in a passage of Athanasius, Apol. pro fuga sua, tom. i. P. 1, p. 140, which Theodoret erroneously dates from the close of the council of Nicaea, instead of the readmission of the Meletians into the fellowship of the church (H. E. i. 26). Alexander died in 328, not in 326. See particulars in Larsow, l. c. p. 26, and § 121 above.} he was nominated to the episcopal succession of Alexandria, on the recommendation of the dying Alexander, and by the voice of the people, though not yet of canonical age, and at first disposed to avoid the election by flight; and thus he was raised to the highest ecclesiastical dignity of the East. For the bishop of Alexandria was at the same time metropolitan of Egypt, Libya, and Pentapolis.

But now immediately began the long series of contests with the Arian party, which had obtained influence at the court of Constantine, and had induced the emperor to recall Arius and his adherents from exile. Henceforth the personal fortunes of Athanasius are so inseparably interwoven with the history of the Arian controversy that Nicene and Athanasian are equivalent terms, and the different depositions and restorations of Athanasius denote so many depressions and victories of the Nicene orthodoxy. Five times did the craft and power of his opponents, upon the pretext of all sorts of personal and political offences, but in reality on account of his inexorable opposition to the Arian and Semi-Arian heresy, succeed in deposing and banishing him. The first exile he spent in Treves, the second chiefly in Rome, the third with the monks in the Egyptian desert; and he employed them in the written defence of his righteous cause. Then the Arian party, was distracted, first by internal division, and further by the death of the emperor Constantius (361), who was their chief support. The pagan Julian recalled the banished bishops of both parties, in the hope that they might destroy one another. Thus, Athanasius among them, who was the most downright opposite of the Christian-hating emperor, again received his bishopric. But when, by his energetic and wise administration, he rather restored harmony in his diocese, and sorely injured paganism, which he feared far less than Arianism, and thus frustrated the cunning plan of Julian, the emperor resorted to violence, and banished him as a dangerous disturber of the peace. For the fourth time Athanasius left Alexandria, but calmed his weeping friends with the prophetic words: “Be of good cheer; it is only a cloud, which will soon pass over.” By presence of mind he escaped from an imperial ship on the Nile, which had two hired assassins on board. After Julian’s death in 362 he was again recalled by Jovian. But the next emperor Valens, an Arian, issued in 367 an edict which again banished all the bishops who had been deposed under Constantius and restored by Julian. The aged Athanasius was obliged for the fifth time to leave his beloved flock, and kept himself concealed more than four months in the tomb of his father. Then Valens, boding ill from the enthusiastic adherence of the Alexandrians to their orthodox bishop, repealed the edict.

From this time Athanasius had peace, and still wrote, at a great age, with the vigor of youth, against Apollinarianism. In the year 373\footnote{Opinions concerning the year of his death waver between 371 and 373. As he was bishop forty-six years, and came to the see in 328 (not 326, as formerly supposed), he cannot have died before 372 or 373.} he died, after an administration of nearly forty-six years, but before the conclusion of the Arian war. He had secured by his testimony the final victory of orthodoxy, but, like Moses, was called away from the earthly scene before the goal was reached.

Athanasius, like many great men (from David and Paul to Napoleon and Schleiermacher), was very small of stature,\footnote{Juliancalled him contemptuously (Ep. 51) \textgreek{μηδε ἀνὴρ, ἀλλ  ἀνθρωπίσκος εὐτελής .}} somewhat stooping and emaciated by fasting and many troubles, but fair of countenance, with a piercing eye and a personal appearance of great power even over his enemies.\footnote{Comp. Gregory Naz. in his Eulogy.} His omnipresent activity, his rapid and his mysterious movements, his fearlessness, and his prophetic insight into the future, were attributed by his friends to divine assistance, by his enemies to a league with evil powers. Hence the belief in his magic art.\footnote{This belief embodied itself in the Arian form of the legend of St. George of Cappadocia, the Arian bishop elected in opposition to Athanasius, and killed by the populace in Alexandria, in his contest with the wizard Athanasius. In this way Arians revenged themselves on the memory of their great adversary. Afterwards the wizard became a dragon, whom George on his horse overcomes. According to others, George was a martyr under Diocletian.} His congregation in Alexandria and the people and monks of Egypt were attached to him through all the vicissitudes of his tempestuous life with equal fidelity and veneration. Gregory Nazianzen begins his enthusiastic panegyric with the words: “When I praise Athanasius, I praise virtue itself, because he combines all virtues in himself.” Constantine the Younger called him “the man of God;” Theodoret, “the great enlightener;” and John of Damascus, the corner-stone of the church of God.”

All this is, indeed, very hyperbolical, after the fashion of degenerate Grecian rhetoric. Athanasius was not free from the faults of his age. But he is, on the whole, one of the purest, most imposing, and most venerable personages in the history of the church; and this judgment will now be almost universally accepted.\footnote{The rationalistic historian Henke(Geschichte der christl. Kirche, 5th ed. 1818, i. p. 212) called him, indeed, a “haughty hard-head,” and the “author of many broils and of the unhappiness of many thousand men.” But the age of the rationalistic debasement of history, thank God, is past. Quite different is the judgment of Gibbon, who despised the faith of Athanasiis, yet could not withhold from him personally the tribute of his admiration. “We have seldom,” says he in ch. xxi. of his celebrated work, “an opportunity of observing, either in active or speculative life, what effect may be produced, or what obstacles may be surmounted by the force of a single mind, when it is inflexibly applied to the pursuit of a single object. The immortal name of Athanasius will never be separated from the Catholic doctrine of the Trinity, to whose defence he consecrated every moment and every faculty of his being .... Amidst the storms of persecution the archbishop of Alexandria was patient of labor, jealous of fame, careless of safety; and although his mind was tainted by the contagion of fanaticism, Athanasius displayed a superiority of charater and abilities which would have qualified him far better than the degenerate sons of Constantinefor the government of a great monarchy.” Dr. Baurthus characterizes Athanasius (Vorlesungen über die Dogmengeschichte, vol. i. ii. p. 41): “His talent for speculative dogmatic investigations, in which he knew how to lay hold, sharply and clearly, of the salient point of the dogma, was as great as the power with which he stood at the head of a party and managed a theological controversy. ... The devotion, with which he defended the cause of orthodoxy, and the importance of the dogma, which was the subject of dispute, have made his name one of the most venerable in the church. In modern times he has been frequently charged with a passionate love for theological controversy. But the most recent ecclesiastical and doctrinal historians are more and more unanimous in according to him a pure zeal for Christian truth, and a profound sense for the apprehension of the same. It is a strong testimony for the purity of his character that his congregation at Antioch adhered to him with tender affection to the last.” A. De Broglie(L’église et l’empire romain au IVesiècle, vol. ii. p. 25) finds the principal quality of the mind of Athanasius in “un rare mélange de droiture de sens et de subtilité de raisonnement. Dans la discussion la plus compliquée rien ne lui échappait, mais rien no l’ébranlait. Il démêlait toutes les nuances de la pensée de son adversaire, en pénétrait tous les détours; mais il ne perdait jamais de vue le point principal et le but du débat .... Unissant lea qualités des deux écoles, il discutait comme un Grec et concluait nettement comme un Latin. Cette combinaison originale, relevée par une indomptable fermeté de caractère, fait encore aujourd’hui le seul mérite qu’ à distance nous puissions pleinement apprécier dans sea écrits.”}

He was (and there are few such) a theological and churchly character in magnificent, antique style. He was a man of one mould and one idea, and in this respect one-sided; yet in the best sense, as the same is true of most great men who are borne along with a mighty and comprehensive thought, and subordinate all others to it. So Paul lived and labored for Christ crucified, Gregory VII. for the Roman hierarchy, Luther for the doctrine of justification by faith, Calvin for the idea of the sovereign grace of God. It was the passion and the life-work of Athanasius to vindicate the deity of Christ, which he rightly regarded as the corner-stone of the edifice of the Christian faith, and without which he could conceive no redemption. For this truth he spent all his time and strength; for this he suffered deposition and twenty years of exile; for this he would have been at any moment glad to pour out his blood. For his vindication of this truth he was much hated, much loved, always respected or feared. In the unwavering conviction that he had the right and the protection of God on his side, he constantly disdained to call in the secular power for his ecclesiastical ends, and to degrade himself to an imperial courtier, as his antagonists often did.

Against the Arians he was inflexible, because he believed they hazarded the essence of Christianity itself, and he allowed himself the most invidious and the most contemptuous terms. He calls them polytheists, atheists, Jews, Pharisees, Sadducees, Herodians, spies, worse persecutors than the heathen, liars, dogs, wolves, antichrists, and devils. But he confined himself to spiritual weapons, and never, like his successor Cyril a century later, used nor counselled measures of force. He suffered persecution, but did not practise it; he followed the maxim: Orthodoxy should persuade faith, not force it.

Towards the unessential errors of good men, like those of Marcellus of Ancyra, he was indulgent. Of Origen he spoke with esteem, and with gratitude for his services, while Epiphanius, and even Jerome, delighted to blacken his memory and burn his bones. To the suspicions of the orthodoxy of Basil, whom, by the way, be never personally knew, he gave no ear, but pronounced his liberality a justifiable condescension to the weak. When he found himself compelled to write against Apollinaris, whom he esteemed and loved, he confined himself to the refutation of his error, without the mention of his name. He was more concerned for theological ideas than for words and formulas; even upon the shibboleth homoousios he would not obstinately insist, provided only the great truth of the essential and eternal Godhead of Christ were not sacrificed. At his last appearance in public, as president of the council of Alexandria in 362, he acted as mediator and reconciler of the contending parties, who, notwithstanding all their discord in the use of the terms ousia and hypostasis, were one in the ground-work of their faith.

No one of all the Oriental fathers enjoyed so high consideration in the Western church as Athanasius. His personal sojourn in Rome and Treves, and his knowledge of the Latin tongue, contributed to this effect. He transplanted monasticism to the West. But it was his advocacy of the fundamental doctrine of Christianity that, more than all, gave him his Western reputation. Under his name the Symbolum Quicunque, of much later, and probably of French, origin, has found universal acceptance in the Latin church, and has maintained itself to this day in living use. His name is inseparable from the conflicts and the triumph of the doctrine of the holy Trinity.

As an author, Athanasius is distinguished for theological depth and discrimination, for dialectical skill, and sometimes for fulminating eloquence. He everywhere evinces a triumphant intellectual superiority over his antagonists, and shows himself a veritable malleus haereticorum. He pursues them into all their hiding-places, and refutes all their arguments and their sophisms, but never loses sight of the main point of the controversy, to which he ever returns with renewed force. His views are governed by a strict logical connection; but his stormy fortunes prevented him from composing a large systematic work. Almost all his writings are occasional, wrung from him by circumstances; not a few of them were hastily written in exile.

They may be divided as follows:

1. Apologetic works in defence of Christianity. Among these are the two able and enthusiastic kindred productions of his youth (composed before 325): “A Discourse against the Greeks,” and “On the Incarnation of the Divine Word,)”\footnote{\textgreek{Λόγος κατὰ Ἑλλήνων}(or Contra Gentes), and \textgreek{Περὶ τῆς ἐνανθρωπήσεως τοῦ λογου} in the first volume, Part 1, of the Bened. ed. pp. 1-97. The latter tract (De incarnatione Verbi Dei) against unbelievers is not to be confounded with the tract written much later (a.d.364), and by some considered spurious: De incarnations Dei Verbi et contra Arianos, tom. i. Pars ii. pp. 871-890.} which he already looked upon as the central idea of the Christian religion.

2. Dogmatic and Controversial works in defence of the Nicene faith; which are at the same time very important to the history of the Arian controversies. Of these the following are directed against Arianism: An Encyclical Letter to all Bishops (written in 341); On the Decrees of the Council of Nicaea (352); On the Opinion of Dionysius of Alexandria (352); An Epistle to the Bishops of Egypt and Libya (356); four Orations against the Arians (358); A Letter to Serapion on the Death of Arius (358 or 359); A History of the Arians to the Monks (between 358 and 360). To these are to be added four Epistles to Serapion on the Deity of the Holy Spirit (358), and two books Against Apollinaris, in defence of the full humanity of Christ (379).

3. Works in his own Personal Defence: An Apology against the Arians (350); an Apology to Constantius (356); an Apology concerning his Flight (De fuga, 357 or 358); and several letters.

4. Exegetical works; especially a Commentary on the Psalms, in which he everywhere finds types and prophecies of Christ and the church, according to the extravagant allegorizing method of the Alexandrian school; and a synopsis or compendium of the Bible. But the genuineness of these unimportant works is by many doubted.\footnote{Comp. the arguments on both sides in the Opera, tom. ii. p. 1004 sqq. and tom. iii. p. 124 sqq.}

5. Ascetic and Practical works. Chief among these are his “Life of St. Anthony,” composed about 365, or at all events after the death of Anthony,\footnote{Opera, tom. ii. (properly tom. i. Pars. ii.), pp. 785-866. Comp. above, § 35.} and his “Festal Letters,” which have but recently become known.\footnote{Comp. the cited editions of the Festal Letters by Cureton, Larsow, and Angelo Mai.} The Festal Letters give us a glimpse of his pastoral fidelity as bishop, and throw new light also on many of his doctrines, and on the condition of the church in his time. In these letters Athanasius, according to Alexandrian custom, announced annually, at Epiphany, to the clergy and congregations of Egypt, the time of the next Easter, and added edifying observations on passages of Scripture, and timely exhortations. These were read in the churches, during the Easter season, especially on Palm-Sunday. As Athanasius was bishop forty-five years, he would have written that number of Festal Letters, if he had not been several times prevented by flight or sickness. The letters were written in Greek, but soon translated into Syriac, and lay buried for centuries in the dust of a Nitrian cloister, till the research of Protestant Scholarship brought them again to the light.

\xsection{Basil the Great}

§ 164. Basil the Great.

I. S. Basilius Caes. Cappad. archiepisc.: Opera omnia quae exstant vel quae ejus nomine circumferuntur, Gr. et Lat. ed. Jul. Garnier, presbyter and monk of the Bened. order. Paris, 1721–’30. 3 vols. fol. Eadem ed. Parisina altera, emendata et aucta a Lud. de Sinner, Par. (Gaume Fratres) 1839, 8 tomi in 6 Partes (an elegant and convenient ed.). Reprinted also by Migne, Par. 1857, in 4 vols. (Patrol. Gr. tom xxix, xxxii.). The first edition of St. Basil was superintended by Erasmus with Froben in Basle, 1532. Comp. also Opera Bas. dogmatica selecta in Thilo’s Bibl. Patr. Gr. dogm. vol. ii. Lips. 1854 (under care of J. D. H. Goldhorn, and containing the Libri iii. adversus Eunomium, and Liber i. de Spiritu Sancto).

II. Ancient accounts and descriptions of Basil in the funeral discourses and eulogies of Gregory Naz. (Oratio xliii.), Gregory Nyss., Amphilochius, Ephraem Syrus. Garnier: Vita S. Basilii, in his edition of the Opera, tom. iii. pp. xxxviii.-ccliv. (in the new Paris ed. of 1839; or tom. i. in Migne’s reprint). Comp. also the Vitae in the Acta Sanctorum, sub Jan. 14, by Hermant, Tillemont (tom. ix.), Fabricius (Bibl. tom. ix.), Cave, Pfeiffer, Schröckh (Part xiii. pp. 8–220), Böhringer, W. Klose (Basilius der Grosse, Stralsund, 1835), and Fialon (Etude historique et littéraire sur S. Basile, Par. 1866).

The Asiatic province of Cappadocia produced in the fourth century the three distinguished church teachers, Basil and the two Gregories, who stand in strong contrast with the general character of their countrymen; for the Cappadocians are described as a cowardly, servile, and deceitful race.\footnote{Particularly in the Letters of Isidore of Pelusium, who flourished in the beginning of the fifth century. Gregory Nazianzen gives a more favorable picture of the Cappadocians, and boasts of their orthodoxy, which, however, might easily be united with the faults above mentioned, especially in the East.}

Basil was born about the year 329,\footnote{According to Garnier; Comp. his Vita Bas. c. 1, § 2. Fabricius puts the birth erroneously into the year 816.} at Caesarea, the capital of Cappadocia, in the bosom of a wealthy and pious family, whose ancestors had distinguished themselves as martyrs. The seed of piety had been planted in him by his grandmother, St. Macrina, and his mother, St. Emmelia. He had four brothers and five sisters, who all led a religious life; two of his brothers, Gregory, bishop of Nyssa, and Peter, bishop of Sebaste, and his sister, Macrina the Younger, are, like himself, among the saints of the Eastern church. He received his literary education at first from his father, who was a rhetorician; afterwards at school in Constantinople (347), where he enjoyed the instruction and personal esteem of the celebrated Libanius; and in Athens, where he spent several years, between 351 and 355,\footnote{On the time of his residence in Athens, see Tillemont and Garnier.} studying rhetoric, mathematics, and philosophy, in company with his intimate friend Gregory Nazianzen, and at the same time with prince Julian the Apostate.

Athens, partly through its ancient renown and its historical traditions, partly by excellent teachers of philosophy and eloquence, Sophists, as they were called in an honorable sense, among whom Himerius and Proaeresius were at that time specially conspicuous, was still drawing a multitude of students from all quarters of Greece, and even from the remote provinces of Asia. Every Sophist had his own school and party, which was attached to him with incredible zeal, and endeavored to gain every newly arriving student to its master. In these efforts, as well as in the frequent literary contests and debates of the various schools among themselves, there was not seldom much rude and wild behavior. To youth who were not yet firmly grounded in Christianity, residence in Athens, and occupation with the ancient classics, were full of temptation, and might easily kindle an enthusiasm for heathenism, which, however, had already lost its vitality, and was upheld solely by the artificial means of magic, theurgy, and an obscure mysticism.\footnote{On this Athenian student-life of that day see especially the 43d, ch. 14 sqq. (in older editions the 20th) Oration of Gregory Nazianzen, and Libanius, De vita sua, p. 13, ed. Reiske.}

Basil and Gregory remained steadfast, and no poetical or rhetorical glitter could fade the impressions of a pious training. Gregory says of their studies in Athens, in his forty-third Oration:\footnote{The Oratio funebris in laudem Basilii M. c. 21 (Opera, ed. Migne, ii. p. 523).} “We knew only two streets of the city, the first and the more excellent one to the churches, and to the ministers of the altar; the other, which, however, we did not so highly esteem, to the public schools and to the teachers of the sciences. The streets to the theatres, games, and places of unholy amusements, we left to others. Our holiness was our great concern; our sole aim was to be called and to be Christians. In this we placed our whole glory.”\footnote{Ἡ\textgreek{μῖν δὲ τὸ μέγα πρᾶγμα καὶ ὄνομα, Χριστιανούς καὶ εἶναι καὶ ὀνομάζεσθαι.}} In a later oration on classic studies Basil encourages them, but admonishes that they should be pursued with caution, and with constant regard to the great Christian purpose of eternal life, to which all earthly objects and attainments are as shadows and dreams to reality. In plucking the rose one should beware of the thorns, and, like the bee, should not only delight himself with the color and the fragrance, but also gain useful honey from the flower.\footnote{Oratio ad adolescentes, quomodo possint ex gentilium libris fructum capere? or more simply, De legendis libris gentilium (in Gamier’s ed. tom. ii. P. i. pp. 243-259). This famous oration, which helped to preserve at least some regard for classical studies in the middle age, has been several times edited separately; as by Hugo Grotius (with a new Latin translation and Prolegomena), 1623; Joh. Potter, 1694; J. H. Majus, 1714; \&c.}

The intimate friendship of Basil and Gregory, lasting from fresh, enthusiastic youth till death, resting on an identity of spiritual and moral aims, and sanctified by Christian piety, is a lovely and engaging chapter in the history of the fathers, and justifies a brief episode in a field not yet entered by any church historian.

With all the ascetic narrowness of the time, which fettered even these enlightened fathers, they still had minds susceptible to science and art and the beauties of nature. In the works of Basil and of the two Gregories occur pictures of nature such as we seek in vain in the heathen classics. The descriptions of natural scenery among the poets and philosophers of ancient Greece and Rome can be easily compressed within a few pages. Socrates, as we learn from Plato, was of the opinion that we can learn nothing from trees and fields, and hence he never took a walk; he was so bent upon self-knowledge, as the true aim of all learning, that he regarded the whole study of nature as useless, because it did not tend to make man either more intelligent or more virtuous. The deeper sense of the beauty of nature is awakened by the religion of revelation alone, which teaches us to see everywhere in creation the traces of the power, the wisdom, and the goodness of God. The book of Ruth, the book of Job, many Psalms, particularly the 104th, and the parables, are without parallel in Grecian or Roman literature. The renowned naturalist, Alexander von Humboldt, collected some of the most beautiful descriptions of nature from the fathers for his purposes.\footnote{In the second volume of his Kosmos, Stuttgart and Tübingen, 1847, p. 27 ff. Humboldt justly observes, p. 26: “The tendency of Christian sentiment was, to prove from the universal order and from the beauty of nature the greatness and goodness of the Creator. Such a tendency, to glorify the Deity from His works, occasioned a prepension to descriptions of nature.” The earliest and largest picture of this kind he finds in the apologetic writer, Minucius Felix. Then he draws several examples from Basil (for whom he confesses he had “long entertained a special predilection”), Epist. xiv. and Epist. ccxxiii. (tom. iii. ed. Gamier), from Gregory of Nyasa, and from Chrysostom.} They are an interesting proof of the transfiguring power of the spirit of Christianity even upon our views of nature.

A breath of sweet sadness runs through them, which is entirely foreign to classical antiquity. This is especially manifest in Gregory of Nyssa, the brother of Basil. “When I see,” says he, for example, “every rocky ridge, every valley, every plain, covered with new-grown grass; and then the variegated beauty of the trees, and at my feet the lilies doubly enriched by nature with sweet odors and gorgeous colors; when I view in the distance the sea, to which the changing cloud leads out—my soul is seized with sadness which is not without delight. And when in autumn fruits disappear, leaves fall, boughs stiffen, stripped of their beauteous dress—we sink with the perpetual and regular vicissitude into the harmony of wonder-working nature. He who looks through this with the thoughtful eye of the soul, feels the littleness of man in the greatness of the universe.”\footnote{From several fragments of Gregory of Nyasa combined and translated (into German) by Humboldt, l.c. p. 29 f.} Yet we find sunny pictures also, like the beautiful description of spring in an oration of Gregory Nazianzen on the martyr Mamas.\footnote{See Ullmann’s Gregor von Nazianz, p. 210 ff.}

A second characteristic of these representations of nature, and for the church historian the most important, is the reference of earthly beauty to an eternal and heavenly principle, and that glorification of God in the works of creation, which transplanted itself from the Psalms and the book of Job into the Christian church. In his homilies on the history of the Creation, Basil describes the mildness of the serene nights in Asia Minor, where the stars, “the eternal flowers of heaven, raised the spirit of man from the visible to the invisible.” In the oration just mentioned, after describing the spring in the most lovely and life-like colors, Gregory Nazianzen proceeds: “Everything praises God and glorifies Him with unutterable tones; for everything shall thanks be offered also to God by me, and thus shall the song of those creatures, whose song of praise I here utter, be also ours .... Indeed it is now [alluding to the Easter festival] the spring-time of the world, the spring-time of the spirit, spring-time for souls, spring-time for bodies, a visible spring, an invisible spring, in which we also shall there have part, if we here be rightly transformed, and enter as new men upon a new life.” Thus the earth becomes a vestibule of heaven, the beauty of the body is consecrated an image of the beauty of the spirit.

The Greek fathers placed the beauty of nature above the works of art, having a certain prejudice against art on account of the heathen abuses of it. “If thou seest a splendid building, and the view of its colonnades would transport thee, look quickly at the vault of the heavens and the open fields, on which the flocks are feeding on the shore of the sea. Who does not despise every creation of art, when in the silence of the heart he early wonders at the rising sun, as it pours its golden (crocus-yellow) light over the horizon? when, resting at a spring in the deep grass or under the dark shade of thick trees, he feeds his eye upon the dim vanishing distance?” So Chrysostom exclaims from his monastic solitude near Antioch, and Humboldt\footnote{L.. c. p. 30.} adds the ingenious remark: “It was as if eloquence had found its element, its freedom, again at the fountain of nature in the then wooded mountain regions of Syria and Asia Minor.”

In the rough times of the first introduction of Christianity among the Celtic and Germanic tribes, who had worshipped the dismal powers of nature in rude symbols, an opposition to intercourse with nature appeared, like that which we find in Tertullian to pagan art; and church assemblies of the twelfth and thirteenth centuries, at Tours (1163) and at Paris (1209), forbid the monks the sinful reading of books on nature, till the renowned scholastics, Albert, the Great († 1280), and the gifted Roger Bacon († 1294), penetrated the mysteries of nature and raised the study of it again to consideration and honor.

We now return to the life of Basil. After finishing his studies in Athens he appeared in his native city of Caesarea as a rhetorician. But he soon after (a.d. 360) took a journey to Syria, Palestine, and Egypt, to become acquainted with the monastic life; and he became more and more enthusiastic for it. He distributed his property to the poor, and withdrew to a lonely romantic district in Pontus, near the cloister in which his mother Emmelia, with his sister Macrina, and other pious and cultivated virgins, were living. “God has shown me,” he wrote to his friend Gregory, “a region which exactly suits my mode of life; it is, in truth, what in our happy jestings we often wished. What imagination showed us in the distance, that I now see before me. A high mountain, covered with thick forest, is watered towards the north by fresh perennial streams. At the foot of the mountain a wide plain spreads out, made fruitful by the vapors which moisten it. The surrounding forest, in which many varieties of trees crowd together, shuts me off like a strong castle. The wilderness is bounded by two deep ravines. On one side the stream, where it rushes foaming down from the mountain, forms a barrier hard to cross; on the other a broad ridge obstructs approach. My hut is so placed upon the summit, that I overlook the broad plain, as well as the whole course of the Iris, which is more beautiful and copious than the Strymon near Amphipolis. The river of my wilderness, more rapid than any other that I know, breaks upon the wall of projecting rock, and rolls foaming into the abyss: to the mountain traveller, a charming, wonderful sight; to the natives, profitable for its abundant fisheries. Shall I describe to you the fertilizing vapors which rise from the (moistened) earth, the cool air which rises from the (moving) mirror of the water? Shall I tell of the lovely singing of the birds and the richness of blooming plants? What delights me above all is the silent repose of the place. It is only now and then visited by huntsmen; for my wilderness nourishes deer and herds of wild goats, not your bears and your wolves. How would I exchange a place with him? Alcmaeon, after he had found the Echinades, wished to wander no further.”\footnote{Ep. xiv. \textgreek{Γρηγορίῳ ἑταίρῶ}(tom. iii. p. 132, ed. nova Paris. Garn.), elegantly reproduced in German by Humboldt, l.c. p. 28, with the observation: “In this simple description of landscape and of forest-life, sentiments are expressed which more intimately blend with those of modem times, than anything that has come down to us from Greek or Roman antiquity.”}

This romantic picture shows that the monastic life had its ideal and poetic side for cultivated minds. In this region Basil, free from all cares, distractions, and interruptions of worldly life, thought that he could best serve God. “What is more blessed than to imitate on earth the choir of angels, at break of day to rise to prayer, and praise the Creator with anthems and songs; then go to labor in the clear radiance of the sun, accompanied everywhere by prayer, seasoning work with praise, as if with salt? Silent solitude is the beginning of purification of the soul. For the mind, if it be not disturbed from without, and do not lose itself through the senses in the world, withdraws into itself, and rises to thoughts of God.” In the Scriptures he found, “as in a store of all medicines, the true remedy for his sickness.”

Nevertheless, he had also to find that flight from the city was not flight from his own self. “I have well forsaken,” says he in his second Epistle,\footnote{Addressed to his friend Gregory, Ep. ii. c. 1 (tom. iii. p. 100).} “my residence in the city as a source of a thousand evils, but I have not been able to forsake myself. l am like a man who, unaccustomed to the sea, becomes seasick, and gets out of the large ship, because it rocks more, into a small skiff, but still even there keeps the dizziness and nausea. So is it with me; for while I carry about with me the passions which dwell in me, I am everywhere tormented with the same restlessness, so that I really get not much help from this solitude.” In the sequel of the letter, and elsewhere, he endeavors, however, to show that seclusion from worldly business, celibacy, solitude, perpetual occupation with the Holy Scriptures, and with the life of godly men, prayer and contemplation, and a corresponding ascetic severity of outward life, are necessary for taming the wild passions, and for attaining the true quietness of the soul.

He succeeded in drawing his friend Gregory to himself. Together they prosecuted their prayer, studies, and manual labor; made extracts from the works of Origen, which we possess, under the name of Philocalia, as the joint work of the two friends; and wrote monastic rules which contributed largely to extend and regulate the coenobite life.

In the year 364 Basil was made presbyter against his will, and in 370, with the co-operation of Gregory and his father, was elected bishop of Caesarea and metropolitan of all Cappadocia. In this capacity he had fifty country bishops under him, and devoted himself thenceforth to the direction of the church and the fighting of Arianism, which had again come into power through the emperor Valens in the East. He endeavored to secure to the catholic faith the victory, first by close connection with the orthodox West, and then by a certain liberality in accepting as sufficient, in regard to the not yet symbolically settled doctrine of the Holy Ghost, that the Spirit should not be considered a creature. But the strict orthodox party, especially the monks, demanded the express acknowledgment of the divinity of the Holy Ghost, and violently opposed Basil. The Arians pressed him still more. The emperor wished to reduce Cappadocia to the heresy, and threatened the bishop by his prefects with confiscation, banishment, and death. Basil replied: “Nothing more? Not one of these things touches me. His property cannot be forfeited, who has none; banishment I know not, for I am restricted to no place, and am the guest of God, to whom the whole earth belongs; for martyrdom I am unfit, but death is a benefactor to me, for it sends me more quickly to God, to whom I live and move; I am also in great part already dead, and have been for a long time hastening to the grave.”

The emperor was about to banish him, when his son, six years of age, was suddenly taken sick, and the physicians gave up all hope. Then he sent for Basil, and his son recovered, though he died soon after. The imperial prefect also recovered from a sickness, and ascribed his recovery to the prayer of the bishop, towards whom he had previously behaved haughtily. Thus this danger was averted by special divine assistance.

But other difficulties, perplexities, and divisions, continually met him, to obstruct the attainment of his desire, the restoration of the peace of the church. These storms, and all sorts of hostilities, early wasted his body. He died in 379, two years before the final victory of the Nicene orthodoxy, with the words: “Into Thy hands, O Lord I commit my spirit; Thou hast redeemed me, O Lord, God of truth.”\footnote{With this prayer of David, Ps. xxxi. 5, Lutheralso took leave of the world.} He was borne to the grave by a deeply sorrowing multitude.

Basil was poor, and almost always sickly; he had only a single worn-out garment, and ate almost nothing but bread, salt, and herbs. The care of the poor and sick he took largely upon himself. He founded in the vicinity of Caesarea that magnificent hospital, Basilias, which we have already mentioned, chiefly for lepers, who were often entirely abandoned in those regions, and left to the saddest fate; he himself took in the sufferers, treated them as brethren, and, in spite of their revolting condition, was not afraid to kiss them.\footnote{Greg. Naz. Orat. xliii. 63, p. 817 sq.}

Basil is distinguished as a pulpit orator and as a theologian, and still more as a shepherd of souls and a church ruler; and in the history of monasticism he holds a conspicuous place.\footnote{K. Hase (§ 102) thus briefly and concisely characterizes him: “An admirer of Libanius and St. Anthony, as zealous for science as for monkery, greatest in church government.”} In classical culture he yields to none of his contemporaries, and is justly placed with the two Gregories among the very first writers among the Greek fathers. His style is pure, elegant, and vigorous. Photius thought that one who wished to become a panegyrist, need take neither Demosthenes nor Cicero for his model, but Basil only.

Of his works, his Five Books against Eunomius, written in 361, in defence of the deity of Christ, and his work on the Holy Ghost, written in 375, at the request of his friend Amphilochius, are important to the history of doctrine.\footnote{The former in tom. i., the latter in tom. iii., ed. Garnier. Both are incorporated in Thilo’s Bibliotheca Patr. Graec. dogm. tom. ii.} He at first, from fear of Sabellianism, recoiled from the strong doctrine of the homoousia; but the persecution of the Arians drove him to a decided confession. Of importance in the East is the Liturgy ascribed to him, which, with that of St. Chrysostom, is still in use, but has undoubtedly reached its present form by degrees. We have also from St. Basil nine Homilies on the history of the Creation, which are full of allegorical fancies, but enjoyed the highest esteem in the ancient church, and were extensively used by Ambrose and somewhat by Augustine, in similar works;\footnote{\textgreek{Εξαήμερον,} or Homiliae ix. in Hexaëmeron. Opera, i. pp. 1-125, ed. Garnier (new ed.). An extended analysis of these sermons is given by Schröckh, xiii. pp. 168-181.} Homilies on the Psalms; Homilies on various subjects; several ascetic and moral treatises;\footnote{Moralia, or short ethical rules, Constitutiones monasticae, \&c., in tom. ii.} and three hundred and sixty-five Epistles,\footnote{Including some spurious, some doubtful, and some from other persons. Tom. iii, pp. 97-681. The numbering of Garnier differs from those of former editors.} which furnish much information concerning his life and times.

\xsection{Gregory of Nyssa}

§ 165. Gregory of Nyssa.

I. S. Gregorius Nyssenus: Opera omnia, quae reperiri potuerunt, Gr. et Lat., nunc primum e mss. codd. edita, stud. Front. Ducaei (Fronto le Duc, a learned Jesuit). Paris, 1615, 2 vols. fol. To be added to this. Appendix Gregorii ex ed. Jac. Gretseri, Par. 1618, fol.; and the Antirrhetoricus adv. Apollinar., first edited by L. Al. Zacagni, Collectanea monum. vet. eccl. Graec. et Lat. Rom. 1698, and in Gallandi, Bibliotheca, tom. vi. Later editions of the Opera by Aeg. Morél, Par. 1638, 3 vols. fol. (“moins belle que cello de 1615, mais plus ample et plus commode ... peu correcte,” according to Brunet); by Migne, Petit-Montrouge (Par.), 1858, 3 vols.; and by Franc. Oehler, Halis Saxonum, 1865 sqq. (Tom. i. continens libros dogmaticos, but only in the Greek original.) Oehler has also commenced an edition of select treatises of Gregory of Nyasa in the original with a German version. The Benedictines of St. Maur had prepared the critical apparatus for an edition of Gregory, but it was scattered during the French Revolution. Angelo Mai, in the Nov. Patrum Biblioth. tom. iv. Pars i. pp. 1–53 (Rom. 1847), has edited a few writings of Gregory unknown before, viz., a sermon Adversus Arium et Sabellium, a sermon De Spiritu Sancto adv. Macedonianos, and a fragment De processione Spiritus S. a Filio (doubtful).

II. Lives in the Acta Sanctorum, and in Butler, sub Mart. 9. Tillemont: Mém. tom. ix. p. 561 sqq. Schröckh: Part xiv. pp. 1–147. Jul. Rupp: Gregors des Bischofs von Nyssa Leben und Meinungen. Leipz. 1834 (unsatisfactory). W. Möller: Gregorii Nyss. doctrina de hominis natura, etc. Halis, 1854, and article in Herzog’s Encykl. vol. v. p. 354 sqq.

Gregory of Nyssa was a younger brother of Basil, and the third son of his parents. Of his honorable descent he made no account. Blood, wealth, and splendor, says he, we should leave to the friends of the world; the Christian’s lineage is his affinity with the divine, his fatherland is virtue, his freedom is the sonship of God. He was weakly and timid, and born not so much for practical life, as for study and speculation. He formed his mind chiefly upon the writings of Origen, and under the direction of his brother, whom he calls his father and preceptor. Further than this his early life is unknown.

After spending a short time as a rhetorician he broke away from the world, retired into solitude in Pontus, and became enamored of the ascetic life.

Quite in the spirit of the then widely-spread tendency towards the monastic life, he, though himself married, commends virginity in a special work, as a higher grade of perfection, and depicts the happiness of one who is raised above the incumbrances and snares of marriage, and thus, as he thinks, restored to the original state of man in Paradise.\footnote{That he was married appears from his own concession, De virginitate, c. 3, where by Theosebia he means his wife (not, as some earlier Roman scholars, and Rupp, l. c. p. 25, suppose, his sister), and from Gregory Nazianzen’s letter of condolence, Ep. 95. He laments that his eulogy of \textgreek{παρθενία} can no longer bring him the desired fruit. Theosebia seems to have lived till 384. Gregory Nazianzen, in his short eulogy of her, says that she rivalled her brothers-in-law (Basil and Peter) who were in the priesthood.} “From all the evils of marriage,” he says, “virginity is free; it has no lost children, no lost husband to bemoan; it is always with its Bridegroom, and delights in its devout exercises, and, when death comes, it is not separated from him, but united with him forever.” The essence of spiritual virginity, however, in his opinion, by no means consists merely in the small matter of sensual abstinence, but in the purity of the whole life. Virginity is to him the true philosophy, the perfect freedom. The purpose of asceticism in general he considered to be not the affliction of the body—which is only a means—but the easiest possible motion of the spiritual functions.

His brother Basil, in 372, called him against his will from his learned ease into his own vicinity as bishop of Nyssa, an inconsiderable town of Cappadocia. He thought it better that the place should receive its honor from his brother, than that his brother should receive his honor from his place. And so it turned out. As Gregory labored zealously for the Nicene faith, he drew the hatred of the Arians, who succeeded in deposing him at a synod in 376, and driving him into exile. But two years later, when the emperor Valens died and Gratian revoked the sentences of banishment, Gregory recovered his bishopric.

Now other trials came upon him. His brothers and sisters died in rapid succession. He delivered a eulogy upon Basil, whom he greatly venerated, and he described the life and death of his beautiful and noble sister Macrina, who, after the death of her betrothed, that she might remain true to him, chose single life, and afterwards retired with her mother into seclusion, and exerted great influence over her brothers.

Into her mouth he put his theological instructions on the soul, death, resurrection, and final restoration.\footnote{In his dialogue, De anima et resurrectione (\textgreek{Περὶ ψυχῆς καὶ ἀναστάσεως μετὰ τῆς ἰδίας ἀδελφῆς Μακρίνης διάλογο;ς}), Opp. iii. 181 sqq. (ed. Morell. 1638), also separately edited by J. G. Krabinger, Lips. 1837, and more recently, together with his biography of his sister, by Franc. Oehler, with a German translation, Leipz. 1858. The last-mentioned edition is at the same time the first volume of a projected Select Library of the Fathers, presenting the original text with a new German translation. The dialogue was written after the death of his brother Basil, and occasioned by it.} She died in the arms of Gregory, with this prayer: “Thou, O God, hast taken from me the fear of death. Thou hast granted me, that the end of this life should be the beginning of true life. Thou givest our bodies in their time to the sleep of death, and awakest them again from sleep with the last trumpet .... Thou hast delivered us from the curse and from sin by Thyself becoming both for us; Thou hast bruised the head of the serpent, hast broken open the gates of hell, hast overcome him who had the power of death, and hast opened to us the way to, resurrection. For the ruin of the enemy and the security of our life, Thou hast put upon those who feared Thee a sign, the sign of Thy holy cross, O eternal God, to whom I am betrothed from the womb, whom my soul has loved with all its might, to whom I have dedicated, from my youth up till now, my flesh and my soul. Oh! send to me an angel of light, to lead me to the place of refreshment, where is the water of peace, in the bosom of the holy fathers. Thou who hast broken the flaming sword, and bringest back to Paradise the man who is crucified with Thee and flees to Thy mercy. Remember me also in Thy kingdom!... Forgive me what in word, deed, or thought, I have done amiss! Blameless and without spot may my soul be received into Thy hands, as a burnt-offering before Thee!”\footnote{Nyss. \textgreek{Περὶ τοῦ βίου τῆς μακαρίας Μακρίνης.}}

Gregory attended the ecumenical council of Constantinople, and undoubtedly, since he was one of the most eminent theologians of the time, exerted a powerful influence there, and according to a later, but erroneous, tradition, he composed the additions to the Nicene Creed which were there sanctioned.\footnote{In Niceph. Call. H. E. xiii. 13. These additions were in use several years before 881, and are found in Epiphanius, Anchorate, n. 120 (tom. ii. p. 122).} The council intrusted to him, as “one of the pillars of catholic orthodoxy,” a tour of visitation to Arabia and Jerusalem, where disturbances had broken out which threatened a schism. He found Palestine in a sad condition, and therefore dissuaded a Cappadocian abbot, who asked his advice about a pilgrimage of his monks to Jerusalem. “Change of place,” says he, “brings us no nearer God, but where thou art, God can come to thee, if only the inn of thy soul is ready .... It is better to go out of the body and to raise one’s self to the Lord, than to leave Cappadocia to journey to Palestine.” He did not succeed in making peace, and he returned to Cappadocia lamenting that there were in Jerusalem men “who showed a hatred towards their brethren, such as they ought to have only towards the devil, towards sin, and towards the avowed enemies of the Saviour.”

Of his later life we know very little. He was in Constantinople thrice afterwards, in 383, 385, and 394, and he died about the year 395.

The wealth of his intellectual life he deposited in his numerous writings, above all in his controversial doctrinal works: Against Eunomius; Against Apollinaris; On the Deity of the Son and the Holy Ghost; On the difference between ousia and hypostasis in God; and in his catechetical compend of the Christian faith.\footnote{The \textgreek{Λόγος κατηχητικὸς ὁ μέγας}stands worthily by the side of the similar work of Origen, De principiis. Separate edition, Gr. and Lat. with notes, by J. G. Krabinger, Munich, 1888.} The beautiful dialogue with his sister Macrina on the soul and the resurrection has been already mentioned. Besides these he wrote many Homilies, especially on the creation of the world, and of man,\footnote{The Hexaëmeron of Gregory is a supplement to his brother Basil’s Hexaëmeron, and discusses the more obscure metaphysical questions connected with this subject. His book on the Workmanship of Man, though written first, may be regarded as a continuation of the Hexaëmeron, and beautifully sets forth the spiritual and royal dignity and destination of man, for whom the world was prepared and adorned as his palace.} on the life of Moses, on the Psalms, on Ecclesiastes, on the Song of Solomon, on the Lord’s Prayer, on the Beatitudes; Eulogies on eminent martyrs and saints (St. Stephen, the Forty Martyrs, Gregory Thaumaturgus, Ephrem, Meletius, his brother Basil); various valuable ascetic tracts; and a biography of his sister Macrina, addressed to the monk Olympios.

Gregory was more a man of thought than of action. He had a fine metaphysical head, and did lasting service in the vindication of the mystery of the Trinity and the incarnation, and in the accurate distinction between essence and hypostasis. Of all the church teachers of the Nicene age he is the nearest to Origen. He not only follows his sometimes utterly extravagant allegorical method of interpretation, but even to a great extent falls in with his dogmatic views.\footnote{On his relation to Origen, Comp. the appendix of Rupp, l.c. pp. 243-262.} With him, as with Origen, human freedom plays a great part. Both are idealistic, and sometimes, without intending it or knowing it, fall into contradiction with the church doctrine, especially in eschatology. Gregory adopts, for example, the doctrine of the final restoration of all things. The plan of redemption is in his view absolutely universal, and embraces all spiritual beings. Good is the only positive reality; evil is the negative, the non-existent, and must finally abolish itself, because it is not of God. Unbelievers must indeed pass through a second death, in order to be purged from the filthiness of the flesh. But God does not give them up, for they are his property, spiritual natures allied to him. His love, which draws pure souls easily and without pain to itself, becomes a purifying fire to all who cleave to the earthly, till the impure element is driven off. As all comes forth from God, so must all return into him at last.

\xsection{Gregory Nazianzen}

§ 166. Gregory Nazianzen.

I. S. Gregorius Theologus, vulgo Nazianzenus: Opera omnia, Gr. et Lat. opera et studio monachorum S. Benedicti e congreg. S. Mauri (Clemencet). Paris, 1778, tom. i. (containing his orations). This magnificent edition (one of the finest of the Maurian editions of the fathers) was interrupted by the French Revolution, but afterwards resumed, and with a second volume (after papers left by the Maurians) completed by A. B. Caillau, Par. l837–’40, 2 vols. fol. Reprinted in Migne’s Patrolog. Graec. (tom. 35–38), Petit-Montrouge, 1857, in 4 vols. (on the separate editions of his Orationes and Carmina, see Brunet, Man. du libraire, tom. ii. 1728 sq.)

II. Biographical notices in Gregory’s Epistles and Poems, in Socrates, Sozomen, Theodoret, Rufinus, and Suidas (s. v. Γρηγόριος). Gregorius Presbyter (of uncertain origin, perhaps of Cappadocia in the tenth century): Βίος τοῦ Γρηγορίου(Greek and Latin in Migne’s ed. of the Opera, tom. i. 243–304). G. Hermant: La vie de S. Basile le Grand et celle de S. Gregoire de Nazianz. Par. 1679, 2 vols. Acta Sanctorum, tom. ii. Maji, p. 373 sqq. Bened. Editores: Vita Greg. ex iis potissimum scriptis adornata (in Migne’s ed. tom. i. pp. 147–242). Tillemont: Mémoires, tom. ix. pp. 305–560, 692–731. Le Clerc: Bibliothèque Universelle, tom. xviii. pp. 1–128. W. Cave: Lives of the Fathers, vol. iii. pp. 1–90 (ed. Oxf. 1840). Schröckh: Part xiii. pp. 275–466. Carl Ullmann: Gregorius von Nazianz, der Theologe. Ein Beitrag zur Kirchen- und Dogmengeschichte des 4ten Jahrhunderts. Darmstadt, 1825. (One of the best historical monographs by a theologian of kindred spirit.) Comp. also the articles of Hefele in Wetzer und Welte’s Kirchenlexikon, vol. iv. 736 ff., and Gass in Herzog’s Encykl. vol. v. 349.

Gregory Nazianzen, or Gregory the Theologian, is the third in the Cappadocian triad; inferior to his bosom friend Basil as a church ruler, and to his namesake of Nyssa as a speculative thinker, but superior to both as an orator. With them he exhibits the flower of Greek theology in close union with the Nicene faith, and was one of the champions of orthodoxy, though with a mind open to free speculation. His life, with its alternations of high station, monastic seclusion, love of severe studies, enthusiasm for poetry, nature, and friendship, possesses a romantic charm. He was “by inclination and fortune tossed between the silence of a contemplative life and the tumult of church administration, unsatisfied with either, neither a thinker nor a poet, but, according to his youthful desire, an orator, who, though often bombastic and dry, labored as powerfully for the victory of orthodoxy as for true practical Christianity.”\footnote{So K. Haseadmirably characterizes him, in his Lehrbuch, p. 138 (7th ed.). The judgment of Gibbon(Decline and Fall, ch. xxii.) is characteristic: “The title of Saint has been added to his name: but the tenderness of his heart, and the elegance of his genius, reflect a more pleasing lustre on the memory of Gregory Nazianzen.” The praise of “the tenderness of his heart” suggests to the skeptical historian another fling at the ancient church, by adding the note: “I can only be understood to mean, that such was his natural temper when it was not hardened, or inflamed, by religious zeal. From his retirement, he exhorts Nectarius to prosecute the heretics of Constantinople.”}

Gregory Nazianzen was born about 330, a year before the emperor Julian, either at Nazianzum, a market-town in the south-western part of Cappadocia, where his father was bishop, or in the neighboring village of Arianzus.\footnote{Respecting the time and place of his birth, views are divided. According to Suidas, Gregory was over ninety years old, and therefore, since he died in 389 or 390, must have been born about the year 800. This statement was accepted by Pagi and other Roman divines, to remove the scandal of his canonized father’s having begotten children after he became bishop; but it is irreconcilable with the fact that Gregory, according to his own testimony (Carmen de vita sua, v. 112 and 238, and Orat. v. c. 23), studied in Athens at the same time with Julianthe Apostate, therefore in 355, and left Athens at the age of thirty years. Comp. Tillemont, tom. ix. pp. 693-697; Schröckh, Part xiii. p. 276, and the admirable monograph of Ullmann, p. 548 sqq. (of which I have made special use in this section).}

In the formation of his religious character his mother Nonna, one of the noblest Christian women of antiquity, exerted a deep and wholesome influence. By her prayers and her holy life she brought about the conversion of her husband from the sect of the Hypsistarians, who, without positive faith, worshipped simply a supreme being; and she consecrated her son, as Hannah consecrated Samuel, even before his birth; to the service of God. “She was,” as Gregory describes her, “a wife according to the mind of Solomon; in all things subject to her husband according to the laws of marriage, not ashamed to be his teacher and his leader in true religion. She solved the difficult problem of uniting a higher culture, especially in knowledge of divine things and strict exercise of devotion, with the practical care of her household. If she was active in her house, she seemed to know nothing of the exercises of religion; if she occupied herself with God and his worship, she seemed to be a stranger to every earthly occupation: she was whole in everything. Experiences had instilled into her unbounded confidence in the effects of believing prayer; therefore she was most diligent in supplications, and by prayer overcame even the deepest feelings of grief over her own and others’ sufferings. She had by this means attained such control over her spirit, that in every sorrow she encountered, she never uttered a plaintive tone before she had thanked God.” He especially celebrates also her extraordinary liberality and self-denying love for the poor and the sick. But it seems to be not in perfect harmony with this, that he relates of her: “Towards heathen women she was so intolerant, that she never offered her mouth or hand to them in salutation.\footnote{Against the express injunction of love for enemies, Hatt. v. 44 ff. The command of 2 John v. 10, 11, which might be quoted in justification of Nonna, refers not to pagans, but to anti-Christian heretics.} She ate no salt with those who came from the unhallowed altars of idols. Pagan temples she did not look at, much less would she have stepped upon their ground; and she was as far from visiting the theatre.” Of course her piety moved entirely in the spirit of that time, bore the stamp of ascetic legalism rather than of evangelical freedom, and adhered rigidly to certain outward forms. Significant also is her great reverence for sacred things. “She did not venture to turn her back upon the holy table, or to spit upon the floor of the church.” Her death was worthy of a holy life. At a great age, in the church which her husband had built almost entirely with his own means, she died, holding fast with one hand to the altar and raising the other imploringly to heaven, with the words: “Be gracious to me, O Christ, my King!” Amidst universal sorrow, especially among the widows and orphans whose comfort and help she had been, she was laid to rest by the side of her husband near the graves of the martyrs. Her affectionate son says in one of the poems in which he extols her piety and her blessed end: “Bewail, O mortals, the mortal race; but when one dies, like Nonna, praying, then weep I not.”

Gregory was early instructed in the Holy Scriptures and in the rudiments of science. He soon conceived a special predilection for the study of oratory, and through the influence of his mother, strengthened by a dream,\footnote{There appeared to him two veiled virgins, of unearthly beauty, who called themselves \emph{Purity} and \emph{Chastity}, companions of Jesus Christ, and friends of those who renounced all earthly connections for the sake of leading a perfectly divine life. After exhorting the youth to join himself to them in spirit, they rose again to heaven. Carmen iv. v. 205-285.} he determined on the celibate life, that he might devote himself without distraction to the kingdom of God. Like the other church teachers of this period, he also gave this condition the preference, and extolled it in orations and poems, though without denying the usefulness and divine appointment of marriage. His father, and his friend Gregory of Nyssa were among the few bishops who lived in wedlock.

From his native town he went for his further education to Caesarea in Cappadocia, where he probably already made a preliminary acquaintance with Basil; then to Caesarea in Palestine, where there were at that time celebrated schools of eloquence; thence to Alexandria, where his revered Athanasius wore the supreme dignity of the church; and finally to Athens, which still maintained its ancient renown as the seat of Grecian science and art. Upon the voyage thither he survived a fearful storm, which threw him into the greatest mental anguish, especially because, though educated a Christian, he, according to a not unusual custom of that time, had not yet received holy baptism, which was to him the condition of salvation. His deliverance he ascribed partly to the intercession of his parents, who had intimation of his peril by presentiments and dreams, and he took it as a second consecration to the spiritual office.

In Athens be formed or strengthened the bond of that beautiful Christian friendship with Basil, of which we have already spoken in the life of Basil. They were, as Gregory says, as it were only one soul animating two bodies. He became acquainted also with the prince Julian, who was at that time studying there, but felt wholly repelled by him, and said of him with prophetic foresight: “What evil is the Roman empire here educating for itself!”\footnote{\textgreek{Οἷον κακὸν ἡ Ῥωμαίων τρέφει}.} He was afterwards a bitter antagonist of Julian, and wrote two invective discourses against him after his death, which are inspired, however, more by the fire of passion than by pure enthusiasm for Christianity, and which were intended to expose him to universal ignominy as a horrible monument of enmity to Christianity and of the retributive judgment of God.\footnote{These Invectivae, or \textgreek{λόγοιστηλιτευτικοί}, are, according to the old order, the 3d and 4th, according to the new the 4th and 5th, of Gregory’s Orations, tom. i. pp. 78-176, of the Benedictine edition.}

Friends wished him to settle in Athens as a teacher of eloquence, but he left there in his thirtieth year, and returned through Constantinople, where he took with him his brother Caesarius, a distinguished physician,\footnote{To this Caesarius, who was afterwards physician in ordinary to the emperor in Constantinople, many, following Photius, ascribe the still extant collection of theological and philosophical questions, Dialogi iv sive Quaestiones Theol. et philos. 145; but without sufficient ground. Comp. Fabricius, Bibl. Gr. viii. p. 435. He was a true Christian, but was not baptized till shortly before his death in 368. His mother Nonna followed the funeral procession in the white raiment of festive joy. He was afterwards, like his brother Gregory, his sister Gorgonia, and his mother, received into the number of the saints of the Catholic church.} to his native city and his parents’ house. At this time his baptism took place. With his whole soul he now threw himself into a strict ascetic life. He renounced innocent enjoyments, even to music, because they flatter the senses. “His food was bread and salt, his drink water, his bed the bare ground, his garment of coarse, rough cloth. Labor filled the day; praying, singing, and holy contemplation, a great part of the night. His earlier life, which was anything but loose, only not so very strict, seemed to him reprehensible; his former laughing now cost him many tears. Silence and quiet meditation were law and pleasure to him.”\footnote{Ullmann, l. c. p. 50.} Nothing but love to his parents restrained him from entire seclusion, and induced him, contrary to talent and inclination, to assist his father in the management of his household and his property.

But he soon followed his powerful bent toward the contemplative life of solitude, and spent a short time with Basil in a quiet district of Pontus in prayer, spiritual contemplations, and manual labors. “Who will transport me,” he afterwards wrote to his friend concerning this visit,\footnote{Epist. ix. p. 774, of the old order, or Ep. vi. of the new (ed. Bened. ii. p. 6).} “back to those former days, in which I revelled with thee in privations? For voluntary poverty is after all far more honorable than enforced enjoyment. Who will give me back those songs and vigils? who, those risings to God in prayer, that unearthly, incorporeal life, that fellowship and that spiritual harmony of brothers raised by thee to a God-like life? who, the ardent searching of the Holy Scriptures, and the light which, under the guidance of the Spirit, we found therein?” Then he mentions the lesser enjoyments of the beauties of surrounding nature.

On a visit to his parents’ house, Gregory against his will, and even without his previous knowledge, was ordained presbyter by his father before the assembled congregation on a feast day of the year 361. Such forced elections and ordinations, though very offensive to our taste, were at that time frequent, especially upon the urgent wish of the people, whose voice in many instances proved to be indeed the voice of God. Basil also, and Augustine, were ordained presbyters, Athanasius and Ambrose bishops, against their will. Gregory fled soon after, it is true, to his friend in Pontus, but out of regard to his aged parents and the pressing call of the church, he returned to Nazianzum towards Easter in 362, and delivered his first pulpit discourse, in which he justified himself in his conduct, and said: “It has its advantage to hold back a little from the call of God, as Moses, and after him Jeremiah, did on account of their age; but it has also its advantage to come forward readily, when God calls, like Aaron and Isaiah; provided both be done with a devout spirit, the one on account of inherent weakness, the other in reliance upon the strength of him who calls.” His enemies accused him of haughty contempt of the priestly office; but he gave as the most important reason of his flight, that he did not consider himself worthy to preside over a flock, and to undertake the care of immortal souls, especially in such stormy times.

Basil, who, as metropolitan, to strengthen the catholic interest against Arianism, set about the establishment of new bishoprics in the small towns of Cappadocia, intrusted to his young friend one such charge in Sasima, a poor market town at the junction of three highways, destitute of water, verdure, and society, frequented only by rude wagoners, and at the time an apple of discord between him and his opponent, the bishop Anthimus of Tyana. A very strange way of showing friendship, unjustifiable even by the supposition that Basil wished to exercise the humility and self-denial of Gregory.\footnote{Gibbon (ch. xxvii.) very unjustly attributes this action of Basil to hierarchical pride and to an intention to insult Gregory. Basil treated his own brother not much better; for Nyssa was likewise an insignificant place.} No wonder that, though a bishopric in itself was of no account to Gregory, this act deeply wounded his sense of honor, and produced a temporary alienation between him and Basil.\footnote{Gregory gave to the pangs of injured friendship a touching expression in the following lines from the poem on his own Life (De vita sua, vss. 476 sqq. tom. ii. p. 699, of the Bened.ed., or tom. iii. 1062, in Migne’s ed.): \par Τοιαῦτ Ἀθῆναι, καὶ πόνοι κοινο λόγων, \par ομΌογετσς ιακ ετ ενυσοιτσs ιβος, \par Νοῦs εις εν αμφοῖν, οὖ δύω, θυᾶμ Ἑλλάδος, \par Καὶ δεξιαὶ, κόσμον μὲν ὡς πόῤῥω βαλεῖν, \par Αὐτοὺς δὲ κοινὸν τῷθεῷ ζῆσαι βίον, \par Λόγους τε δοῦναι τῷ μόνῳ σοφῷ Λόγῳ. \par Διεσκέδασται πάντα, ῎εῤῥιπται χαμαὶ, \par Αὖραι φέπουσι τὰς παλαιὰςελπίδας. \par   \par “Talia Athenae, et communia studia, \par Ejusdem texti et mensae consors vita, \par Mena una, non duae in ambobus, res mira Graeciae, \par Dataeque dexterae, mundum ut procul rejiceremus, \par Deoque simul viveremus, \par Et literas soli sapienti Verbo dedicaremus. \par Dissipata haec sunt omnia, et humi projects, \par Venti auferunt spes nostras antiquas.” \par Gibbon (ch. xxvii.) quotes this passage with admiration, though with characteristic omission of vss. 479-481, which refer to their harmony in religion; and he aptly alludes to a parallel from Shakespeare, who had never read the poems of Gregory Nazianzen, but who gave to similar feelings a similar expression, in the Midsummer Night’s Dream, where Helena utters the same pathetic complaint to her friend Hermia: \par “Is all the counsel that we two have shared, \par The sister’s vows,” \&e.} At the combined request of his friend and his aged father, he suffered himself indeed to be consecrated to the new office; but it is very doubtful whether he ever went to Sasima.\footnote{Gibbon says: “He solemnly protests, that he never consummated his spiritual marriage with this disgusting bride.”} At all events we soon afterwards find him in his solitude, and then again, in 372, assistant of his father in Nazianzum. In a remarkable discourse delivered in the presence of his father in 372, he represented to the congregation his peculiar fluctuation between an innate love of the contemplative life of seclusion and the call of the Spirit to public labor.

“Come to my help,” said he to his hearers,\footnote{Orat. xii. 4; tom. i. 249 sq. (in Migne’s ed. tom. i. p. 847).} “for I am almost torn asunder by my inward longing and by the Spirit. The longing urges me to flight, to solitude in the mountains, to quietude of soul and body, to withdrawal of spirit from all sensuous things, and to retirement into myself, that I may commune undisturbed with God, and be wholly penetrated by the rays of His Spirit .... But the other, the Spirit, would lead me into the midst of life, to serve the common weal, and by furthering others to further myself, to spread light, and to present to God a people for His possession, a holy people, a royal priesthood (Tit. ii. 14; 1 Pet. ii. 9), and His image again purified in many. For as a whole garden is more than a plant, and the whole heaven with all its beauties is more glorious than a star, and the whole body more excellent than one member, so also before God the whole well-instructed church is better than one well-ordered person, and a man must in general look not only on his own things, but also on the things of others. So Christ did, who, though He might have remained in His own dignity and divine glory, not only humbled Himself to the form of a servant, but also, despising all shame, endured the death of the cross, that by His suffering He might blot out sin, and by His death destroy death.”

Thus he stood a faithful helper by the side of his venerable and universally beloved father, who reached the age of almost an hundred years, and had exercised the priestly office for forty-five; and on the death of his father, in 374, he delivered a masterly funeral oration, which Basil attended.\footnote{Orat. xviii. \textgreek{Ἐπιτάφιος εἰς τὸν πατέρα, παρόντος Βασιλείου} (ed. Bened. tom. i. pp. 330-362; in Migne’s ed. i. 981 sqq.).} “There is,” said he in this discourse, turning to his still living mother, “only one life, to behold the (divine) life; there is only one death—sin; for this is the corruption of the soul. But all else, for the sake of which many exert themselves, is a dream which decoys us from the true; it is a treacherous phantom of the soul. When we think so, O my mother, then we shall not boast of life, nor dread death. For whatsoever evil we yet endure, if we press out of it to true life, if we, delivered from every change, from every vortex, from all satiety, from all vassalage to evil, shall there be with eternal, no longer changeable things, as small lights circling around the great.”

A short time after he had been invested with the vacant bishopric, he retired again, in 375, to his beloved solitude, and this time be went to Seleucia in Isauria, to the vicinity of a church dedicated to St. Thecla.

There the painful intelligence reached him of the death of his beloved Basil, a.d. 379. On this occasion be wrote to Basil’s brother, Gregory of Nyssa: “Thus also was it reserved for me still in this unhappy life to hear of the death of Basil and the departure of this holy soul, which is gone out from us, only to go in to the Lord, after having already prepared itself for this through its whole life.” He was at that time bodily and mentally very much depressed. In a letter to the rhetorician Eudoxius he wrote: “You ask, how it fares with me. Very badly. I no longer have Basil; I no longer have Caesarius; my spiritual brother, and my bodily brother. I can say with David, my father and my mother have forsaken me. My body is sickly, age is coming over my head, cares become more and more complicated, duties overwhelm me, friends are unfaithful, the church is without capable pastors, good declines, evil stalks naked. The ship is going in the night, a light nowhere, Christ asleep. What is to be done? O, there is to me but one escape from this evil case: death. But the hereafter would be terrible to me, if I had to judge of it by the present state.”

But Providence had appointed him yet a great work and in exalted position in the Eastern capital of the empire. In the year 379 he was called to the pastoral charge by the orthodox church in Constantinople, which, under the oppressive reign of Arianism, was reduced to a feeble handful; and he was exhorted by several worthy bishops to accept the call. He made his appearance unexpectedly. With his insignificant form bowed by disease, his miserable dress, and his simple, secluded mode of life, he at first entirely disappointed the splendor-loving people of the capital, and was much mocked and persecuted.\footnote{Once the Arian populace even stormed his church by night, desecrated the altar, mixed the holy wine with blood, and Gregory but barely escaped the fury of common women and monks, who were armed with clubs and stones. The next day he was summoned before the court for the tumult, but so happily defended himself, that the occurrence heightened the triumph of his just cause. Probably from this circumstance he afterwards received the honorary title of \emph{confessor}. See Ullmann, p. 176.} But in spite of all he succeeded, by his powerful eloquence and faithful labor, in building up the little church in faith and in Christian life, and helped the Nicene doctrine again to victory. In memory of this success his little domestic chapel was afterwards changed into a magnificent church, and named Anastasia, the Church of the Resurrection.

People of all classes crowded to his discourses, which were mainly devoted to the vindication of the Godhead of Christ and to the Trinity, and at the same time earnestly inculcated a holy walk befitting the true faith. Even the famous Jerome, at that time already fifty years old, came from Syria to Constantinople to hear these discourses, and took private instruction of Gregory in the interpretation of Scripture. He gratefully calls him his preceptor and catechist.

The victory of the Nicene faith, which Gregory had thus inwardly promoted in the imperial city, was outwardly completed by the celebrated edict of the new emperor Theodosius, in February, 380. When the emperor, on the 24th of December of that year, entered Constantinople, he deposed the Arian bishop, Demophilus, with all his clergy, and transferred the cathedral church\footnote{Not the church of St. Sophia, as Tillemont assumes, but the church of the Apostles, as Ullmann, p. 223, supposes; for Gregory never names the former, but mentions the latter repeatedly, and that as the church in which he himself preached. Constantinebuilt both, but made the church of the Apostles the more magnificent, and chose it for his own burial place (Euseb. Vita Const. iv. 58-60); St. Sophia afterwards became under Justinian the most glorious monument of the later Greek architecture, and the cathedral of Constantinople.} to Gregory with the words: “This temple God by our hand intrusts to thee as a reward for thy pains.” The people tumultuously demanded him for bishop, but he decidedly refused. And in fact he was not yet released from his bishopric of Nazianzum or Sasima (though upon the latter he had never formally entered); he could be released only by a synod.

When Theodosius, for the formal settlement of the theological controversies, called the renowned ecumenical council in May, 381, Gregory was elected by this council itself bishop of Constantinople, and, amidst great festivities, was inducted into the office. In virtue of this dignity he held for a time the presidency of the council.

When the Egyptian and Macedonian bishops arrived, they disputed the validity of his election, because, according to the fifteenth canon of the council of Nice, he could not be transferred from his bishopric of Sasima to another; though their real reason was, that the election had been made without them, and that Gregory would probably be distasteful to them as a bold preacher of righteousness. This deeply wounded him. He was soon disgusted, too, with the operations of party passions in the council, and resigned with the following remarkable declaration:

“Whatever this assembly may hereafter determine concerning me, I would fain raise your mind beforehand to something far higher: I pray you now, be one, and join yourselves in love! Must we always be only derided as infallible, and be animated only by one thing, the spirit of strife? Give each other the hand fraternally. But I will be a second Jonah. I will give myself for the salvation of our ship (the church), though I am innocent of the storm. Let the lot fall upon me, and cast me into the sea. A hospitable fish of the deep will receive me. This shall be the beginning of your harmony. I reluctantly ascended the episcopal chair, and gladly I now come down. Even my weak body advises me this. One debt only have I to pay: death; this I owe to God. But, O my Trinity! for Thy sake only am I sad. Shalt Thou have an able man, bold and zealous to vindicate Thee? Farewell, and remember my labors and my pains.”

In the celebrated valedictory which be delivered before the assembled bishops, he gives account of his administration; depicts the former humiliation and the present triumph of the Nicene faith in Constantinople, and his own part in this great change, for which he begs repose as his only reward; exhorts his hearers to harmony and love; and then takes leave of Constantinople and in particular of his beloved church, with this address:

“And now, farewell, my Anastasia, who bearest a so holy name; thou hast exalted again our faith, which once was despised; thou, our common field of victory, thou new Shiloh, where we first established again the ark of the covenant, after it had been carried about for forty years on our wandering in the wilderness.”

Though this voluntary resignation of so high a post proceeded in part from sensitiveness and irritation, it is still an honorable testimony to the character of Gregory in contrast with the many clergy of his time who shrank from no intrigues and by-ways to get possession of such dignities. He left Constantinople in June, 381, and spent the remaining years of his life mostly in solitude on his paternal estate of Arianzus in the vicinity of Nazianzum, in religious exercises and literary pursuits. Yet he continued to operate through numerous epistles upon the affairs of the church, and took active interest in the welfare and sufferings of the men around him. The nearer death approached, the more he endeavored to prepare himself for it by contemplation and rigid ascetic practice, that he “might be, and might more and more become, in truth a pure mirror of God and of divine things; might already in hope enjoy the treasures of the future world; might walk with the angels; might already forsake the earth, while yet walking upon it; and might be transported into higher regions by the Spirit.” In his poems he describes himself, living solitary in the clefts of the rocks among the beasts, going about without shoes, content with one rough garment, and sleeping upon the ground covered with a sack. He died in 390 or 391; the particular circumstances of his death being now unknown. His bones were afterwards brought to Constantinople; and they are now shown at Rome and Venice.

Among the works of Gregory stand pre-eminent his five Theological Orations in defence of the Nicene doctrine against the Eunomians and Macedonians, which he delivered in Constantinople, and which won for him the honorary title of the Theologian (in the narrower sense, i.e., vindicator of the deity of the Logos).\footnote{Hence called also \textgreek{λόγοι θεολογικοί}, Orationes theologicae. They are Orat. xxvii.-xxxi. in the Bened. ed. tom. i. pp. 487-577 (in Migne, tom. ii. 9 sqq.), and in the Bibliotheca Patrum Graec. dogmatica of Thilo, vol. ii. pp. 366-537.} His other orations (forty-five in all) are devoted to the memory of distinguished martyrs, friends, and kindred, to the ecclesiastical festivals, and to public events or his own fortunes. Two of them are bitter attacks on Julian after his death.\footnote{Invectivae, Orat. iv. et v. in the Bened. ed. tom. i. 73-176 (in Migne’s ed. tom. i. pp. 531-722). His horror of Julianmisled him even to eulogize the Arian emperor Constantius, to whom his brother was physician.} They are not founded on particular texts, and have no strictly logical order and connection.

He is the greatest orator of the Greek church, with the exception perhaps of Chrysostom; but his oratory often degenerates into arts of persuasion, and is full of labored ornamentation and rhetorical extravagances, which are in the spirit of his age, but in violation of healthful, natural taste.

As a poet he holds a subordinate, though respectable place. He wrote poetry only in his later life, and wrote it not from native impulse, as the bird sings among the branches, but in the strain of moral reflection, upon his own life, or upon doctrinal and moral themes. Many of his orations are poetical, many of his poems are prosaic. Not one of his odes or hymns passed into use in the church. Yet some of his smaller pieces, apothegms, epigrams, and epitaphs, are very beautiful, and betray noble affections, deep feeling, and a high order of talent and cultivation.\footnote{His poems fill together with the Epistles the whole second tome of the magnificent Benedictine edition, so delightful to handle, which was published at Paris, 1842 (edente et curante D. A. B. Caillau), and vols. iii. and iv. of Migne’s reprint. They are divided by the Bened. editor into: I. Poëmata theologica (dogmatica, moralia); II. historica (a. autobiographical, quae spectant ipsum Gregorium, \textgreek{περὶ ἑαυτοῦ}, De seipso; and b. \textgreek{περὶ τῶν ἑτέρων}, quae spectant alios); III. epitaphia; IV. epigrammata; and V. a long tragedy, \emph{Christus patiens}, with Christ, the Holy Virgin, Joseph, Theologus, Mary Magdalene, Nicodemus, Nuntius, and Pilate as actors. This is the first attempt at a Christian drama. The order of the poems, as well as the Orations and Epistles, differs in the Benedictine from that of the older editions. See the comparative table in tom. ii. p. xv. sqq. One of the finest passages in his poems is his lamentation over the temporary suspension of his friendship with Basil, quoted above, \textbf{p. 914}.}

We have, finally, two hundred and forty-two (or 244) Epistles from Gregory, which are important to the history of the time, and in some cases very graceful and interesting.

\xsection{Didymus of Alexandria}

§ 167. Didymus of Alexandria.

I. Didymi Alexandrini Opera omnia: accedunt S. Amphilochii et Nectarii scripta quae supersunt Graece, accurante et denuo recognoscente J. P. Migne. Petit-Montrouge (Paris), 1858. (Tom. xxxix. of the Patrologia Graeca.)

II. Hieronymus: De viris illustr. c. 109, and Prooem. in Hoseam. Scattered accounts in Rufinus, Palladius, Socrates, Sozomen, and Theodoret. Tillemont: Mémoires, x. 164. Fabricius: Bibl. Gr. tom. ix. 269 sqq. ed. Harless (also in Migne’s ed. of the Opera, pp. 131–140). Schröckh: Church History, vii. 74–87. Guericke: De schola Alexandrina. Hal. 1824.

Didymus, the last great teacher of the Alexandrian catechetical school, and a faithful follower of Origen, was born probably at Alexandria about the year 309. Though he became in his fourth year entirely blind, and for this reason has been surnamed Caecus, yet by extraordinary industry he gained comprehensive and thorough knowledge in philosophy, rhetoric, and mathematics. He learned to write by means of wooden tablets in which the characters were engraved; and he became so familiar with the Holy Scriptures by listening to the church lessons, that he knew them almost all by heart.

Athanasius nominated him teacher in the theological school, where he zealously labored for nearly sixty years. Even men like Jerome, Rufinus, Palladius, and Isidore, sat at his feet with admiration. He was moreover an enthusiastic advocate of ascetic life, and stood in high esteem with the Egyptian anchorites; with St. Anthony in particular, who congratulated him, that, though blind to the perishable world of sense, he was endowed with the eye of an angel to behold the mysteries of God. He died at a great age, in universal favor, in 395.

Didymus was thoroughly orthodox in the doctrine of the Trinity, and a discerning opponent of the Arians, but at the same time a great venerator of Origen, and a participant of his peculiar views concerning the pre-existence of souls, and probably concerning final restoration. For this reason he was long after his death condemned with intolerant zeal by several general councils.\footnote{First at the fifth ecumenical council in 553. The sixth council in 680 stigmatized him as a defender of the abominable doctrine of Origen, who revived the heathen fables of the transmigration of souls; and the seventh repeated this in 787}

We have from him a book On the Holy Ghost, translated by Jerome into Latin, in which he advocates, with much discrimination, and in simple, biblical style, the consubstantiality of the Spirit with the Father, against the Semi-Arians and Pneumatomachi of his time;\footnote{Didymus wrote only \emph{one} book De Spiritu Sancto (see Jerome, De viris illustr. c. 135: librum unum de Sp. S. Didymi quem in Latinum transtuli). The division into three books is of later date.} and three books on the Trinity, in the Greek original.\footnote{Discovered and edited by Joh. Aloys. Mingarelli, at Bologna, 1769, with a Latin translation and learned treatises on the life, doctrine, and writings of Didymus. (Dr. Herzog, Encykl. iii. p. 384, confounds this edition with a preliminary advertisement by the brother Ferdinand Mingarelli: Veterum testimonia de Didymo Alex. coeco, ex quibus tres libri de Trinitate nuper detecti eidem asseruntur, Rom. 1764. The title of the work itself is: Didymus, De Trinitate libri tres, nunc primum ex Passioneiano codice Gr. editi, Latine conversi, ac notis illustrati a D. Joh. Aloys. Mingarellio, Bononiae 1769, fol.)} He wrote also a brief treatise against the Manichaeans. Of his numerous exegetical works we have a commentary on the Catholic Epistles,\footnote{The Latin version is found in the libraries of the church fathers. The original Greek has been edited by Dr. Fr. Lücke from Muscovite manuscripts in four academic dissertations: Quaestiones ac vindiciae Didymianae, sive Didymi Alex. enarratio in Epistolas Catholicas Latina, Graeco exemplari magnam partem e Graecis scholiis restituta, Gotting. 1829-’32. Reprinted in Migne’s edition of Opera Didymi, pp. 1731-1818.} and large fragments, in part uncertain, of commentaries on the Psalms, Job, Proverbs, and some Pauline Epistles.\footnote{In Migne’s ed. p. 1109 sqq.}

\xsection{Cyril of Jerusalem}

§ 168. Cyril of Jerusalem.

I. S. Cyrilus, archiepisc. Hierosolymitanus: Opera quae exstant omnia, \&c., cura et studio Ant. Aug. Touttaei (Touttée), presb. et monachi Bened. e congreg. S. Mauri. Paris, 1720. 1 vol. fol. (edited after Touttée’s death by the Benedictine D. Prud. Maranus. Comp. therewith Sal. Deyling: Cyrillus Hieros. a corruptelis Touttaei aliorumque purgatus. Lips. 1728). Reprint, Venice, 1763. A new ed. by Migne, Petit-Montrouge, 1857 (Patrol. Gr. tom. xxxiii., which contains also the writings of Apollinaris of Laodicea, Diodor of Tarsus, and others). The Catecheses of Cyril have also been several times edited separately, and translated into modern languages. Engl. transl. in the Oxford Library of the Fathers, vol. ii. Oxf. 1839.

II. Epiphanius: Haer. lx. 20; lxxiii. 23, 27, 37. Hieronymus: De viris illustr. c. 112. Socrates: H. E. ii. 40, 42, 45; iii. 20. Sozomen: iv. 5, 17, 20, 22, 25. Theodoret: H. E. ii. 26, 27; iii. 14; v. 8. The Dissertationes Cyrillianae de vita et scriptis S. Cyr. \&c. in the Benedictine edition of the Opera, and in Migne’s reprint, pp. 31–822. The Acta Sanctorum, and Butler, sub mense Martii 18. Tillemont: tom. viii. pp. 428–439, 779–787. Also the accounts in the well-known patristic works of Dupin, Ceillier, Cave, Fabricius. Schröckh: Part xii. pp. 369–476.

Cyrilus, presbyter and, after 350, bishop of Jerusalem, was extensively involved during his public life in the Arian controversies. His metropolitan, Acacius of Caesarea, an Arian, who had elevated him to the episcopal chair, fell out with him over the Nicene faith and on a question of jurisdiction, and deposed him at a council in 357. His deposition was confirmed by an Arian council at Constantinople in 360.

After the death of the emperor Constantius he was restored to his bishopric in 361, and in 363 his embittered adversary, Acacius, converted to the orthodox faith. When Julian encouraged the Jews to rebuild the temple, Cyril is said to have predicted the miscarriage of the undertaking from the prophecies of Daniel and of Christ, and he was justified by the result. Under the Arian emperor Valens he was again deposed and banished, with all the other orthodox bishops, till he finally, under Theodosius, was permitted to return to Jerusalem in 379, to devote himself undisturbed to the supervision and restoration of his sadly distracted church until his death.

He attended the ecumenical council in Constantinople in 381, which confirmed him in his office, and gave him the great praise of having suffered much from the Arians for the faith. He died in 386, with his title to office and his orthodoxy universally acknowledged, clear of all the suspicions which many had gathered from his friendship with Semi-Arian bishops during his first exile.\footnote{His sentiments on the holy Trinity are discussed at length in the third preliminary dissertation of the Bened. editor (in Migne’s ed. p. 167 sqq.).}

From Cyril we have an important theological work, complete, in the Greek original: his twenty-three Catecheses.\footnote{\textgreek{Κατηχήσεις φωτιζομένων}(or \textgreek{βαπτιζομένων}), catecheses illuminandorum. They are preceded by a procatechesis.} The work consists of connected religious lectures or homilies, which he delivered while presbyter about the year 347, in preparing a class of catechumens for baptism. It follows that form of the Apostles’ Creed or the Rule of Faith which was then in use in the churches of Palestine and which agrees in all essential points with the Roman; it supports the various articles with passages of Scripture, and defends them against the heretical perversions of his time. The last five, called the Mystagogic Catecheses,\footnote{\textgreek{Κατηχήσεις μυσταγωγικαί}. The name is connected with the mysterious practices of the \emph{disciplina arcani} of the early church. Comp. the conclusion of the first Mystagogic Catechesis, c. 11 (Migne, p. 1075). The mystagogic lectures are also separately numbered. The first is a general exhortation to the baptized on 1 Pet. v. 8; the second treats De baptismo; the third, De chrismate; the fourth, De corpore et sanguine Christi; the fifth, De sacra liturgia et communione.} are addressed to newly baptized persons, and are of importance in the doctrine of the sacraments and the history of liturgy. In these he explains the ceremonies then customary at baptism: Exorcism, the putting off of garments, anointing, the short confession, triple immersion, confirmation by the anointing oil; also the nature and ritual of the holy Supper, in which he sees a mystical vital union of believers with Christ, and concerning which he uses terms verging at least upon the doctrine of transubstantiation. In connection with this he gives us a full account of the earliest eucharistic liturgy, which coincides in all essential points with such other liturgical remains of the Eastern church, as the Apostolic Constitutions and the Liturgy of St. James.

The Catecheses of Cyril are the first example of a popular compend of religion; for the catechetical work of Gregory of Nyssa (λόγος κατηχητικὸς ὁ μέγας) is designed not so much for catechumens, as for catechists and those intending to become teachers.

Besides several homilies and tracts of very doubtful genuineness, a homily on the healing of the cripple at Bethesda\footnote{Homilia in paralyticum, John v. 2-16 (in Migne’s ed. pp. 1131-1158).} and a remarkable letter to the emperor Constantius of the year 351, are also ascribed to Cyril.\footnote{Ep. ad Constantium imper. De viso Hierosolymus lucidae crucis signo, pp. 1154-1178.} In the letter he relates to the emperor the miraculous appearance of a luminous cross extending from Golgotha to a point over the mount of Olives (mentioned also by Socrates, Sozomen, and others), and calls upon him to praise the “consubstantial Trinity.”\footnote{\textgreek{Τὴν ἁγίαν καὶ ὁμοούσιον Τριάδα, τὸν ἀληθινὸν Θεὸν ἡμῶν, ωὟ πρέπει πᾶσα δόξα εἰς τοὺς · αἰῶνος τῶν αἰώνων.}}

\xsection{Epiphanius}

§ 169. Epiphanius.

I. S. Epiphanius: Opera omnia, Gr. et Lat., Dionysius Petavius ex veteribus libris recensuit, Latine vertit et animadversionibus illustravit. Paris, 1622, 2 vols. fol. The same edition reprinted with additions at Cologne (or rather at Leipsic), 1682, and by J. P. Migne Petit-Montrouge, 1858, in 3 vols. (tom. xli.-xliii. of Migne’s Patrologia Graeca). The Πανάριονor Panaria of Epiphanius, together with his Anacephalaeosis, with the Latin version of both by Petavius, has also been separately edited by Fr. Oehler, as tom. ii. and iii. of his Corpus haereseologicum, Berol. 1859–’61. (Part second of tom. iii. contains the Animadversiones of Petavius, and A. Jahn’s Symbolae ad emendanda et illustranda S. Epiphanii Panaria.)

II. Hieronymus: De viris illustr. c. 114, and in several of his Epistles relating to the Origenistic controversies, Epp. 66 sqq. ed. Vallarsi. Socrates: Hist. Eccl. l. vi. c. 10–14. Sozomen: H. E. viii. 11–15. Old biographies, full of fables, see in Migne’s edition, tom. i., and in Petav. ii. 318 sqq. The Vita Epiph. in the Acta Sanctorum for May, tom. iii. die 12, pp. 36–49 (also reprinted in Migne’s ed. tom. i.). Tillemont: Mémoires, tom. x. pp. 484–521, and the notes, pp. 802–809. Fr. Arm. Gervaise: L’histoire et la vie De saint Epiphane. Par. 1738. Fabricius: Biblioth. Graeca ed. Harless, tom. viii. p. 255 sqq. (also reprinted in Migne’s ed. of Epiph. i. 1 sqq.). W. Cave: Lives of the Fathers, iii. 207–236 (new Oxf. ed.). Schröckh: Th. x. 3 ff. R. Adelb. Lipsius: Zur Quellenkritik des Epiphanies. Wien, 1865. (A critical analysis of the older history of heresies, in Epiph. haer. 13–57, with special reference to the Gnostic systems.)

Epiphanius,\footnote{There are several prominent ecclesiastical writers of that name. Compare a list of them in Fabricius, l. c.} who achieved his great fame mainly by his learned and intolerant zeal for orthodoxy, was born near Eleutheropolis in Palestine, between 310 and 320, and died at sea, at a very advanced age, on his way back from Constantinople to Cyprus, in 403. According to an uncertain, though not improbable tradition, he was the son of poor Jewish parents, and was educated by a rich Jewish lawyer, until in his sixteenth year he embraced the Christian religion,\footnote{See the biography of his pupil John, ch. 2, in Migne’s ed. i. 25 sqq. Cave accepts this story, and it receives some support from the Palestine origin of Epiphanius, and from his knowledge of the Hebrew language, which was then so rare that Jeromewas the only father besides Epiphanius who possessed it.}—the first example, after St. Paul, of a learned Jewish convert and the only example among the ancient fathers; for all the other fathers were either born of Christian parents, or converted from heathenism.

He spent several years in severe ascetic exercises among the hermits of Egypt, and then became abbot of a convent near Eleutheropolis. In connection with his teacher and friend Hilarion he labored zealously for the spread of monasticism in Palestine.\footnote{He composed a eulogy on Hilarion, which, with some others of his works, is lost}

In the year 367 he was unanimously elected by the people and the monks bishop of Salamis (Constantia), the capital of the island of Cyprus. Here he wrote his works against the heretics, and took active part in the doctrinal controversies of his age. He made it his principal business to destroy the influence of the arch-heretic Origen, for whom he had contracted a thorough hatred from the anchorites of Egypt. On this mission he travelled in his old age to Palestine and Constantinople, and died in the same year in which Chrysostom was deposed and banished, an innocent sacrifice on the opposite side in the violent Origenistic controversies.\footnote{Comp. above, §§ 133 and 134.}

Epiphanius was revered even by his cotemporaries as a saint and as a patriarch of orthodoxy. Once as he passed through the streets of Jerusalem in company with bishop John, mothers brought their children to him that he might bless them, and the people crowded around him to kiss his feet and to touch the hem of his garment. After his death his name was surrounded by a halo of miraculous legends. He was a man of earnest, monastic piety, and of sincere but illiberal zeal for orthodoxy. His good nature easily allowed him to be used as an instrument for the passions of others, and his zeal was not according to knowledge. He is the patriarch of heresy-hunters. He identified Christianity with monastic piety and ecclesiastical orthodoxy and considered it the great mission of his life to pursue the thousand-headed hydra of heresy into all its hiding places. Occasionally, however, his fiery zeal consumed what was subsequently considered an essential part of piety and orthodoxy. Sharing the primitive Christian abhorrence of images, he destroyed a picture of Christ or some saint in a village church in Palestine; and at times he violated ecclesiastical order.

The learning of Epiphanius was extensive, but ill digested. He understood five languages: Hebrew, Syriac, Egyptian, Greek, and a little Latin. Jerome, who himself knew but three languages, though he knew these far better than Epiphanius, called him the Five-tongued,\footnote{\textgreek{Πεντάγλωττος}.} and Rufinus reproachfully says of him that he considered it his sacred duty as a wandering preacher to slander the great Origen in all languages and nations.\footnote{Hieron. Apol. adv. Rufinum, l. iii. c. 6 (Opera, tom. ii. 537, ed. Vall.) and l. ii. 21 and 22 (tom. ii. 515). Jeromesays that “papa” Epiphanius had read the six thousand [?] books of Origen, and in his apology against Rufinus and in his letters he speaks of him with great respect as a confederate in the war upon Origen. He acknowledges, however, that his statements need an accurate and careful verification. In his Liber de viris illustribus, cap. 114, he disposes of him very summarily with two sentences: “Epiphanius, Cypri Salaminae episcopus, scripsit adversus omnes haereses libros, et multa alia, quae ab eruditus propter res, a simplicioribus propter verba lectitantur. Superest usque hodie, et in extrema jam senectute varia cudit opera.”} He was lacking in knowledge of the world and of men, in sound judgment, and in critical discernment. He was possessed of a boundless credulity, now almost proverbial, causing innumerable errors and contradictions in his writings. His style is entirely destitute of beauty or elegance.

Still his works are of considerable value as a storehouse of the history of ancient heresies and of patristic polemics. They are the following:

1. The Anchor,\footnote{\textgreek{Αγκυρωτός}, Ancoratus, or Ancora fidei catholicae in tom, ii of Petavius; tom. iii. 11-236 of Migne.} a defence of Christian doctrine, especially of the doctrines of the Trinity, the incarnation, and the resurrection; in one hundred and twenty-one chapters. He composed this treatise a.d. 373, at the entreaty of clergymen and monks, as a stay for those who are tossed about upon the sea by heretics and devils. In it he gives two creeds, a shorter and a longer, which show that the addition made by the second ecumenical council to the Nicene symbol, in respect to the doctrine of the Holy Ghost and of the church, had already been several years in use in the church.\footnote{Anc. n. 119 and 120 (tom. iii. 23 sqq. ed. Migne).} For the shorter symbol, which, according to Epiphanius, had to be said at baptism by every orthodox catechumen in the East, from the council of Nicaea to the tenth year of Valentinian and Valens (a.d. 373), is precisely the same as the Constantinopolitan; and the longer is even more specific against Apollinarianism and Macedonianism, in the article concerning the Holy Ghost. Both contain the anathemas of the Nicene Creed; the longer giving them in an extended form.

2. The Panarium, or Medicine-chest,\footnote{\textgreek{Πανάριον}, Panarium (Panaria), sive Arcula, or Adversus lxxx. haereses (Petavius, tom. i. f. 1-1108; Migne, tom. i. 173-1200, and tom. ii. 10-832). Epiphanius himself names it \textgreek{πανάριον, εἴτ  οὖν κιβώτιον ἰατρικὸν καὶ θηριοδηκτικόν}, Panarium, sive Arculam Medicam ad eorum qui a serpentibus icti sunt remedium (Epist. ad Acacium et Paulum, in Oehler’s ed. i. p. 7).} which contains antidotes for the poison of all heresies. This is his chief work, composed between the years 374 and 377, in answer to solicitations from many quarters. And it is the chief hereseological work of the ancient church. It is more extensive than any of the similar works of Justin Martyr, Irenaeus, and Hippolytus before it, and of Philastrius (or Philastrus), Augustine, Theodoret, pseudo-Tertullian, pseudo-Jerome, and the author of Praedestinatus, after it.\footnote{Compare the convenient collection of the Latin writers De haeresibus, viz.: Philastrius, Augustine, the author of Praedestinatus (the first book), pseudo-Tertullian, pseudo-Jerome, Isdortis Hispalensis, and Gennadius (De ecclesiasticis dogmatibus), in the first volume of Franz Oehler’s Corpus haereseologicum, Berolini, 1856. This collection is intended to embrace eight volumes. Tom. ii. and iii. contain the anti-heretical works of Epiphanius; the remaining volumes are intended for Theodoret, pseudo-Origen, John of Damascus, Leontius, Timotheus, Irenaeus, and Nicetae Choniatae Thesaurus orthodoxae fidei.} Epiphanius brought together, with the diligence of an unwearied compiler, but without logical or chronological arrangement, everything he could learn from written or oral sources concerning heretics from the beginning of the world down to his time. But his main concern is the antidote to heresy, the doctrinal refutations, in which he believed himself to be doing God and the church great service, and which, with all their narrowness and passion, contain many good thoughts and solid arguments. He improperly extends the conception of heresy over the field of all religion; whereas heresy is simply a perversion or caricature of Christian truth, and lives only upon the Christian religion. He describes and refutes no less than eighty heresies,\footnote{Perhaps with a mystic reference to the eighty concubines in the Song of Songs, vi. 8: “Sexaginta sunt reginae et octoginta concubinae, et adolescentularum non est numerus. Una est columba mea, perfecta mea.” (Vulgate.)} twenty of them preceding the time of Christ.\footnote{Pseudo-Tertullian(in Libellus adversus omnes haereses), Philastrus, and pseudo-Hieronymus (Indiculus de haeresibus) likewise include the Jewish sects among the heresies; while Irenaeus, Augustine, Theodoret, and the unknown author of the Semi-Pelagian work Praedestinatus more correctly begin with the Christian sects. For further particulars, see the comparative tables of Lipsius, l.c. p. 4 ff.} The pre-Christian heresies are: Barbarism, from Adam to the flood; Scythism; Hellenism (idolatry proper, with various schools of philosophy); Samaritanism (including four different sects); and Judaism (subdivided into seven parties: Pharisees, Sadducees, Scribes, Hemerobaptists, Osseans, Nazarenes, and Herodians).\footnote{Epiphanius in his shorter work, the Anacephalaeosis, deviates somewhat from the order in the Panarion. His twenty heresies before Christ are as follows: \par Order in the Panarion: \par 1. Barbarismus, \par 2. Scythismus, \par 3. Hellenismus, \par 4. Judaismus, \par Hellenismi \par 5. Stoici, \par 6. Platonici, \par 7. Pythagorei, \par 8. Epicurei, \par Samaritismi \par 9. Samaritae, \par 10. Esseni, \par 11. Sebuaei, \par 12. Gortheni, \par 13. Dosithei, \par Judaismi \par 14. Saducaei, \par 15. Scribae, \par 16. Pharisaei, \par 17. Hemerobaptistae \par 18. Nazaraei, \par 19. Osseni or Ossaei, \par 20. Herodiani \par   \par Order in the Anacephalaeosis: \par 1. Barbarismus, \par 2. Scythismus, \par 3. Hellenismus, \par 4. Judaismus. \par 5. Samaritismus, \par Hellenismi \par 6. Pythagorei, \par 7. Platonici, \par 8. Stoici, \par 9. Epicurei, \par Sararitismi \par 10. Gortheni, \par 11. Sebuaei, \par 12. Esseni, \par 13. Dosithei, \par Judaismi \par 14. Scribae, \par 15. Pharisaei, \par 16. Sadducaei, \par 17. Hemerobaptistae, \par 18. Ossaei, \par 19. Nazaraei, \par 20. Herodiani \par} Among the Christian heresies, of which Simon Magus, according to ancient tradition, figures as patriarch, the different schools of Gnosticism (which may be easily reduced to about a dozen) occupy the principal space. With the sixty-fourth heresy Epiphanius begins the war upon the Origenists, Arians, Photinians, Marcellians, Semi-Arians, Pneumatomachians, Antidikomarianites, and other heretics of his age. In the earlier heresies he made large use, without proper acknowledgment, of the well-known works of Justin Martyr, Irenaeus, and Hippolytus, and other written sources and oral traditions. In the latter sections he could draw more on his own observation and experience.

3. The Anacephalaeosis is simply an abridgment of the Panarion, with a somewhat different order.\footnote{\textgreek{Ανακεφαλαίωσις}, or Epitome Panarii (tom. ii. 126, ed. Patav.; tom. ii. 834-886, ed. Migne).}

This is the proper place to add a few words upon similar works of the post-Nicene age.

About the same time, or shortly after Epiphanius (380), Philastrius or Philastrus, bishop of Brixia (Brescia), wrote his Liber de haeresibus (in 156 chapters).\footnote{Edited by \emph{J. A. Fabricius}, Hamburg, 1728; by \emph{Gallandi}, Bibliotheca, tom. vii. pp. 475-521; and by \emph{Oehler} in tom. i. of his Corpus haereseolog. pp. 5-185. The close affinity of Philastrus with Epiphanius is usually accounted for on the ground of the dependence of the former on the latter. This seems to have been the opinion of Augustine, Epistola 222 ad Quodvultdeum. But Lipsius (l.c. p. 29 ff.) derives both from a common older source, viz., the work of Hippolytus against thirty-two heresies and explains the offence of Epiphanius (who mentions Hippolytus only once) by the unscrupulousuess of the authorship of the age, which had no hesitation in decking itself with borrowed plumes.} He was still more liberal with the name of heresy, extending it to one hundred and fifty-six systems, twenty-eight before Christ, and a hundred and twenty-eight after. He includes peculiar opinions on all sorts of subjects: Haeresis de stellis coelo affixis, haeresis de peccato Cain, haeresis de Psalterii inequalitate, haeresis de animalibus quatuor in prophetis, haeresis de Septuaginta interpretibus, haeresis de Melchisedech sacerdote, haeresis de uxoribus, et concubinis Salomonis!

He was followed by St. Augustine, who in the last years of his life wrote a brief compend on eighty-eight heresies, commencing with the Simonians and ending with the Pelagians.\footnote{Liber de haeresibus, addressed to Quodvultdeus, a deacon who had requested him to write such a work. Augustine, in his letter of reply to Quodvultdeus (Ep. 222 in the Bened. edition) alludes to the work of Philastrus, whom he had seen with Ambrosein Milan, and to that of Epiphanius, and calls the latter “longe Philastrio doctiorem.” The work of Augustineis also embodied in Oehler’s Corpus haereseol. tom. i. pp. 189-225. The following is a complete list of the heresies of Augustineas given by him at the close of the preface: 1. Simoniani; 2. Menandriani; 3. Saturniniani; 4. Basilidiani; 5. Nicolaitae; 6. Gnostici; 7. Carpocratiani; 8. Cerinthiani, vel Merinthiani; 9. Nazaraei; 10. Hebio-naei; 11. Valentiniani; 12. Secundiani; 13. Ptolemaei; 14. Marcitae; 15. Colorbasii; 16. Heracleonitae; 17. Ophitae; 18. Caiani; 19. Sethiani; 20. Archontici; 21. Cerdoniani; 22. Marcionitae; 23. Apellitae; 24. Se-veriani; 25. Tatiani, vel Encratitae; 26. Cataphryges; 27. Pepuziani, alias Quintilliani; 28. Artotyritae; 29. Tessarescaedecatitae; 30. Alegi; 31. Adamiani; 32. Elcesaei et Sampsaei; 33. Theodotiani; 34. Mel-chisedechiani; 35. Bardesanistae; 36. Noëtiani; 37. Valesii; 38. Ca-thari, sive Novatiani; 39. Angelici; 40. Apostolici; 41. Sabelliani; 42. Ori-geniani; 43. Alii Origeniani; 44. Pauliani; 45. Photiniani; 46. Manichaei; 47. Hieracitae; 48. Meletiani; 49. Ariani; 50. Vadiani, sive Anthropo-morphitae; 51. Semiariani; 52. Macedoniani; 53. Aëriani; 54. Aëtiani, qui et Eunomiani; 55. Apollinaristae; 56. Antidicomarianitae; 57. Mas-saliani, sive Euchitae; 58. Metangismonitae; 59. Seleuciani, vel Her-miani; 60. Proclianitae; 61. Patriciani; 62. Ascitae; 63. Passaloryn-chitae; 64. Aquarii; 65. Coluthiani; 66. Floriniani; 67. De mundi statu dissentientes; 68. Nudis pedibus ambulantes; 69. Donatistae, sive Do-natiani; 70. Priscillianistae; 71. Cum hominibus non manducantes; 72. Rhetoriani; 73. Christi divinitatem passibilem dicentes; 74. Triformem deum putantes; 75. Aquam Deo coaeternam dicentes; 76. Imaginem Dei non esse animam dicentes; 77. Innumerabiles mundos opinantes; 78. Animas converti in daemones et in quaecunque animalia existi-mantes; 79. Liberationem omnium apud inferos factam Christi descen-sione credentes; 80. Christi de Patre nativitati initium temporis dantes; 81. Luciferiani; 82. Iovinianistae; 83. Arabici; 84. Helvidiani; 85. Pater-niani, sive Venustiani; 86. Tertullianistae; 87. Abeloitae; 88. Pelagiani, qui et Caelestiani.}

The unknown author of the book called Praedestinatus added two more heretical parties, the Nestorians and the Predestinarians, to Augustine’s list; but the Predestinarians are probably a mere invention of the writer for the purpose of caricaturing and exposing the heresy of an absolute predestination to good and to evil.\footnote{Corpus haereseol. i. 229-268. Comp. above, § 159.}

4. In addition to those anti-heretical works, we have from Epiphanius a biblical archeological treatise on the Measures and Weights of the Scriptures,\footnote{\textgreek{Περὶ μέτρων καὶ σταθμῶν}, De ponderibus et mensuris, written in 392. (Tom. ii. 158, ed. Petav.; tom. iii. 237, ed. Migne.)} and another on the Twelve Gems on the breastplate of Aaron, with an allegorical interpretation of their names.\footnote{\textgreek{Περὶ τῶν δώδεκα λίθων}, De xii. gemmis in veste Aaronis. (Tom. ii. 233, ed. Pet; iii. 293, ed. Migne.)}

A Commentary of Epiphanius on the Song of Songs was published in a Latin translation by Foggini in 1750 at Rome. Other works ascribed to him are lost, or of doubtful origin.

\xsection{John Chrysostom}

§ 170. John Chrysostom.

I. S. Joannis Chrysostomi. archiepiscopi Constantinopolitani, Opera omnia quae exstant vel quae ejus nomine circumferuntur, ad MSS. codices Gallic. etc. castigata, etc. (Gr. et Lat.). Opera et studio D. Bernardi de Montfaucon, monachi ordinis S. Benedicti e congregatione S. Mauri, opem ferentibus aliis ex eodem sodalitio monachis. Paris. 1718–’38, in 13 vols. fol. The same edition reprinted at Venice, 1734–’41, in 13 vols. fol. (after which I quote in this section); also at Paris by Sinner (Gaume), 1834–’39, in 13 vols. (an elegant edition, with some additions), and by J. P. Migne, Petit-Montrouge, 1859–’60, in 13 vols. Besides we have a number of separate editions of the Homilies, and of the work on the Priesthood, both in Greek, and in translations. A selection of his writings in Greek and Latin was edited by F. G. Lomler, Rudolphopoli, 1840, 1 volume. German translations of the Homilies (in part) by J. A. Cramer (Leipzig, 1748–’51), Feder (Augsburg, 1786), Ph. Mayer (Nürnberg, 1830), W. Arnoldi (Trier, 1831), Jos. Lutz (Tübingen, 1853); English translations of the Homilies on the New Testament in the Oxford Library of the Fathers, 1842–’53.

II. Palladius (a friend of Chrysostom and bishop of Helenopolis in Bithynia, author of the Historia Lausiaca; according to others a different person): Dialogus historicus de vita et conversatione beati Joannis Chrysostomi cum Theodoro ecclesiae Romanae diacono (in the Bened. ed. of the Opera, tom. xiii. pp. 1–89). Hieronymus: De viris illustribus, c. 129 (a very brief notice, mentioning only the work de sacerdotio). Socrates: H. E. vi. 3–21. Sozomen: H. E. viii. 2–23. Theodoret: H. E. v. 27–36. B. de Montfaucon: Vita Joannis Chrys. in his edition of the Opera, tom. xiii. 91–178. Testimonia Veterum de S. Joann. Chrys. scriptis, ibid. tom. xiii. 256–292. Tillemont: Mémoires, vol. xi. pp. 1–405. F. Stilting: Acta Sanctorum, Sept. 14 (the day of his death), tom. iv. pp. 401–709. A. Butler: Lives of Saints, sub Jan. 27. W. Cave: Lives of the Fathers, vol. iii. p. 237 ff. J. A. Fabricius: Biblioth. Gr. tom. viii. 454 sqq. Schröckh: Vol. x. p. 309 ff. A. Neander: Der heilige Chrysostomus (first 1821), 3d edition, Berlin, 1848, 2 vols. Abbé Rochet: Histoire de S. Jean Chrysostome. Par. 1866, 2 vols. Comp. also A. F. Villemain’s Tableau de l’éloquence chrétienne au IVe siècle. Paris, 1854.

John, to whom an admiring posterity since the seventh century has given the name Chrysostomus, the Golden-mouthed, is the greatest expositor and preacher of the Greek church, and still enjoys the highest honor in the whole Christian world. No one of the Oriental fathers has left a more spotless reputation; no one is so much read and so often quoted by modern commentators.

He was born at Antioch, a.d. 347.\footnote{Baur(Vorlesungen über die Dogmengeschichte, Bd. i. Abthlg. ii. p. 50) and others erroneously state the year 354 as that of his birth. Comp. Tillemont and Montfaucon (tom. xiii. 91).} His father was a distinguished military officer. His mother Anthusa, who from her twentieth year was a widow, shines with Nonna and Monica among the Christian women of antiquity. She was admired even by the heathen, and the famous rhetorician Libanius, on hearing of her consistency and devotion, felt constrained to exclaim: “Ah! what wonderful women there are among the Christians.”\footnote{\textgreek{Βαβαὶ, οἷαι παρὰ χριστιανοῖς γυναῖκες εἰσι} . Chrysostom himself relates this of his heathen teacher (by whom undoubtedly we are to understand Libanius), though, it Is true, with immediate reference only to the twenty years’ widowhood of his mother; Ad viduam juniorem, Opera, tom. i. p. 340. Comp. the remarks of Montfaucon in the Vita, tom. xiii. 92.} She gave her son an admirable education, and early planted in his soul the germs of piety, which afterwards bore the richest fruits for himself and for the church. By her admonitions and the teachings of the Bible he was secured against the seductions of heathenism.

He received his literary training from Libanius, who accounted him his best scholar, and who, when asked shortly before his death (395) whom he wished for his successor, replied: “John, if only the Christians had not carried him away.”

After the completion of his studies he became a rhetorician. He soon resolved, however, to devote himself to divine things, and after being instructed for three years by bishop Meletius in Antioch, he received baptism.

His first inclination after his conversion was to adopt the monastic life, agreeably to the ascetic tendencies of the times; and it was only by the entreaties of his mother, who adjured him with tears not to forsake her, that he was for a while restrained. Meletius made him reader, and so introduced him to a clerical career. He avoided an election to the bishopric (370) by putting forward his friend Basil, whom he accounted worthier, but who bitterly complained of the evasion. This was the occasion of his celebrated treatise On the Priesthood, in which, in the form of a dialogue with Basil, he vindicates his not strictly truthful conduct, and delineates the responsible duties of the spiritual office.\footnote{\textgreek{Περὶ ἱερωσύνης}. De sacerdotio libri vi. Separate editions are: That of \emph{Frobenius} at Basel, 1525, Greek, with a preface by \emph{Erasmus}; that of \emph{Hughes} at Cam. bridge, 1710, Greek and Latin, with the Life of Chrysostomby \emph{Cave}; that of J\emph{. A. Bengel}, Stuttgart, 1725, Greek and Latin, reprinted at Leipsic in 1825 and 1834; besides several translations into modern languages. Comp. above, § 51, p. 253.}

After the death of his mother he fled from the seductions and tumults of city life to the monastic solitude of the mountains near Antioch, and there spent six happy years in theological study and sacred meditation and prayer, under the guidance of the learned abbot Diodorus (afterwards bishop of Tarsus, † 394), and in communion with such like-minded young men as Theodore of Mopsuestia, the celebrated father of Antiochian (Nestorian) theology († 429). Monasticism was to him a most profitable school of experience and self-government; because he embraced this mode of life from the purest motives, and brought into it intellect and cultivation enough to make the seclusion available for moral and spiritual growth.

In this period he composed his earliest writings in praise of monasticism and celibacy, and his two long letters to the fallen Theodore (subsequently bishop of Mopsuestia), who had regretted his monastic vow and resolved to marry.\footnote{Compare Tillemont, Montfaucon, and Neander (l.c. i. p. 86 ff.).} Chrysostom regarded this small affair from the ascetic stand-point of his age as almost equal to an apostasy from Christianity, and plied all his oratorical arts of sad sympathy, tender entreaty, bitter reproach, and terrible warning, to reclaim his friend to what he thought the surest and safest way to heaven. To sin, he says, is human, but to persist in sin is devilish; to fall is not ruinous to the soul, but to remain on the ground is. The appeal had its desired effect, and cannot fail to make a salutary impression upon every reader, provided we substitute some really great offence for the change of a mode of life which can only be regarded as a temporary and abnormal form of Christian practice.

By excessive self-mortifications John undermined his health, and returned about 380 to Antioch. There he was immediately ordained deacon by Meletius in 386, and by Flavian was made presbyter. By his eloquence and his pure and earnest character he soon acquired great reputation and the love of the whole church.

During the sixteen or seventeen years of his labors in Antioch he wrote the greater part of his Homilies and Commentaries, his work on the Priesthood, a consolatory Epistle to the despondent Stagirius, and an admonition to a young widow on the glory of widowhood and the duty of continuing in it. He disapproved second marriage, not as sinful or illegal, but as inconsistent with an ideal conception of marriage and a high order of piety.

After the death of Nectarius (successor of Gregory Nazianzen), towards the end of the year 397, Chrysostom was chosen, entirely without his own agency, patriarch of Constantinople. At this post he labored several years with happy effect. But his unsparing sermons aroused the anger of the empress Eudoxia, and his fame excited the envy of the ambitious patriarch Theophilus of Alexandria. An act of Christian love towards the persecuted Origenistic monks of Egypt involved him in the Origenistic controversy, and at last the united influence of Theophilus and Eudoxia overthrew him. Even the sympathy of the people and of Innocent I., the bishop of Rome, was unavailing in his behalf. He died in banishment on the fourteenth of September, a.d. 407, thanking God for all.\footnote{Compare particulars above, § 134.} The Greeks celebrate his memorial day on the thirteenth of November, the Latins on the twenty-seventh of January, the day on which his remains in 438 were solemnly deposited in the Church of the Apostles in Constantinople with those of the emperors and patriarchs.

Persecution and undeserved sufferings tested the character of Chrysostom, and have heightened his fame. The Greek church honors him as the greatest teacher of the church, approached only by Athanasius and the three Cappadocians. His labors fall within the comparatively quiet period between the Trinitarian and the Christological controversies. He was not therefore involved in any doctrinal controversy except the Origenistic; and in that he had a very innocent part, as his unspeculative turn of mind kept him from all share in the Origenistic errors. Had he lived a few decades later he would perhaps have fallen under suspicion of Nestorianism; for he belonged to the same Antiochian school with his teacher Diodorus of Tarsus, his fellow-student Theodore of Mopsuestia, and his successor Nestorius. From this school, whose doctrinal development was not then complete, he derived a taste for the simple, sober, grammatico-historical interpretation, in opposition to the arbitrary allegorizing of the Alexandrians, while he remained entirely free from the rationalizing tendency which that school soon afterwards discovered. He is thus the soundest and worthiest representative of the Antiochian theology. In anthropology he is a decided synergist; and his pupil Cassian, the founder of Semi-Pelagianism, gives him for an authority.\footnote{Julianof Eclanum had already appealed several times to Chrysostomagainst Augustine, as Augustinenotes Contra Jul., and in the Opus imperfectum.} But his synergism is that of the whole Greek church; it had no direct conflict with Augustinianism, for Chrysostom died several years before the opening of the Pelagian controversy. He opposed the Arians and Novatians, and faithfully and constantly adhered to the church doctrine, so far as it was developed; but he avoided narrow dogmatism and angry controversy, and laid greater stress on practical piety than on unfruitful orthodoxy.\footnote{Niedner(Geschichte der christl. Kirche, 1846, p. 323, and in his posthumous Lehrbuch, 1866, p. 303) briefly characterizes him thus: “In him we find a most complete mutual interpenetration of theoretical and practical theology, as well as of the dogmatical and ethical elements, exhibited mainly in the fusion of the exegetical and homiletical. Hence his exegesis was guarded against barren philology and dogma; and his pulpit discourse was free from doctrinal abstraction and empty rhetoric. The introduction of the knowledge of Christianity from the sources into the practical life of the people left him little time for the development of special dogmas.}

Valuable as the contributions of Chrysostom to didactic theology may be, his chief importance and merit lie not in this department, but in homiletical exegesis, pulpit eloquence, and pastoral care. Here he is unsurpassed among the ancient fathers, whether Greek or Latin. By talent and culture he was peculiarly fitted to labor in a great metropolis. At that time a bishop, as he himself says, enjoyed greater honor at court, in the society of ladies, in the houses of the nobles, than the first dignitaries of the empire.\footnote{The \textgreek{τόπαρχοι}and \textgreek{ὕπαρχοι}, the praefect praetorio. Homil. iii. in Acta Apost.} Hence the great danger, of hierarchical pride and worldly conformity, to which so many of the prelates succumbed. This danger Chrysostom happily avoided. He continued his plain monastic mode of life in the midst of the splendor of the imperial residence, and applied all his superfluous income to the support of the sick and the stranger. Poor for himself, he was rich for the poor. He preached an earnest Christianity fruitful in good works, he insisted on strict discipline, and boldly attacked the vices of the age and the hollow, worldly, hypocritical religion of the court. He, no doubt, transcended at times the bounds of moderation and prudence, as when he denounced the empress Eudoxia as a new Herodias thirsting after the blood of John; but he erred “on virtue’s side,” and his example of fearless devotion to duty has at all times exerted a most salutary influence upon clergymen in high and influential stations. Neander not inaptly compares his work in the Greek church with that of Spener, the practical reformer in the Lutheran church of the seventeenth century, and calls him a martyr of Christian charity, who fell a victim in the conflict with the worldly spirit of his age.\footnote{In his monograph on Chrysostom, vol. i. p. 5.}

In the pulpit Chrysostom was a monarch of unlimited power over his hearers. His sermons were frequently interrupted by noisy theatrical demonstrations of applause, which he indignantly rebuked as unworthy of the house of God.\footnote{This Greek custom of applauding the preacher by clapping the hands and stamping the feet (called \textgreek{κρότος}, from \textgreek{κρούω}) was a sign of the secularization of the church after its union with the state. It is characteristic of his age that a powerful sermon of Chrysostomagainst this abuse was most enthusiastically applauded by his hearers!} He had trained his natural gift of eloquence, which was of the first order, in the school of Demosthenes and Libanius, and ennobled and sanctified it in the higher school of the Holy Spirit.\footnote{Karl Hase(Kirchengeschichte, § 104, seventh edition) truly says of Chrysostomthat “he complemented the sober clearness of the Antiochian exegesis and the rhetorical arts of Libanius with the depth of his warm Christian heart, and that he carried out in his own life, as far as mortal man can do it, the ideal of the priesthood which, in youthful enthusiasm, he once described.”} He was in the habit of making careful preparation for his sermons by the study of the Scriptures, prayer, and meditation; but he knew how to turn to good account unexpected occurrences, and some of his noblest efforts were extemporaneous effusions under the inspiration of the occasion. His ideas are taken from Christian experience and especially from the inexhaustible stores of the Bible, which he made his daily bread, and which he earnestly recommended even to the laity. He took up whole books and explained them in order, instead of confining himself to particular texts, as was the custom after the introduction of the pericopes. His language is noble, solemn, vigorous, fiery, and often overpowering. Yet he was by no means wholly free from the untruthful exaggerations and artificial antitheses, which were regarded at that time as the greatest ornament and highest triumph of eloquence, but which appear to a healthy and cultivated taste as defects and degeneracies. The most eminent French preachers, Bossuet, Massillon, and Bourdaloue, have taken Chrysostom for their model.

By far the most numerous and most valuable writings of this father are the Homilies, over six hundred in number, which he delivered while presbyter at Antioch and while bishop at Constantinople.\footnote{They are contained in vols. ii.-xii of the Benedictine edition.} They embody his exegesis; and of this they are a rich storehouse, from which the later Greek commentators, Theodoret, Theophylact, and Oecumenius, have drawn, sometimes content to epitomize his expositions. Commentaries, properly so called, he wrote only on the first eight chapters of Isaiah and on the Epistle to the Galatians. But nearly all his sermons on Scripture texts are more or less expository. He has left us homilies on Genesis, the Psalms, the Gospel of Matthew, the Gospel of John, the Acts, and all the Epistles of Paul, including the Epistle to the Hebrews. His homilies on the Pauline Epistles are especially esteemed.\footnote{A beautiful edition of the Homilies on the Pauline Epistles in Greek (but without the Latin version) has been recently published in connection with the Oxford Library of the Fathers under the title: S. Joannis Chrysostomi interpretatio omnium Epistolarum Paulinarum per homilias facta, Oxon. 1849-’52, 4 vols. The English translation has already been noticed.}

Besides these expository sermons on whole books of the Scriptures, Chrysostom delivered homilies on separate sections or verses of Scripture, festal discourses, orations in commemoration of apostles and martyrs, and discourses on special occasions. Among the last are eight homilies Against the Jews (against Judaizing tendencies in the church at Antioch), twelve homilies Against the Anomoeans (Arians), and especially the celebrated twenty and one homilies On the Statues, which called forth his highest oratorical powers.\footnote{The Homiliae xii contra Anomoeans de incomprehensibili Dei natura, and the Orationes viii adversus Judaeos are in the first, the Homilies xxi ad populum Antiochenum, de statuis, and the six Orationes de fato et providential in the second volume of the Bened. edition. The Homilies on the Statues are translated into English in the Oxford Library of the Fathers, 1842, 1 volume.} He delivered the homilies on the Statues at Antioch in 387 during a season of extraordinary public excitement, when the people, oppressed by excessive taxation, rose in rebellion, tore down the statues of the emperor Theodosius I., the deceased empress Flacilla, and the princes Arcadius and Honorius, dragged them through the streets, and so provoked the wrath of the emperor that he threatened to destroy the city—a calamity which was avoided by the intercession of bishop Flavian.

The other works of Chrysostom are his youthful treatise on the Priesthood already alluded to; a number of doctrinal and moral essays in defence of the Christian faith, and in commendation of celibacy and the nobler forms of monastic life;\footnote{Ad Theodorum lapsum; Adversus oppugnatores vitae monasticae; Comparatio regis et monachi; De compunctione cordis; De virginitate; Ad viduam juniorem, etc.,—all in the first volume of the Bened. edition together with the vi Libri de Sacerdotio; also in Lomler’s selection of Chrys. Opera praestantissima.} and two hundred and forty-two letters, nearly all written during his exile between 403 and 407. The most important of the letters are two addressed to the Roman bishop Innocent I., with his reply, and seventeen long letters to his friend Olympias, a pious widow and deaconess. They all breathe a noble Christian spirit, not desiring to be recalled from exile, convinced that there is but one misfortune,—departure from the path of piety and virtue, and filled with cordial friendship, faithful care for all the interests of the church, and a calm and cheerful looking forward to the glories of heaven.\footnote{The Epistles are in tom. iii. The Epistolae ad Olympiadem, and ad Innocentium are also included in Lomler’s selection (pp. 165-252). On Olympias, compare above, § 52, and especially Tillemont, tom. xi. pp. 416-440.}

The so-called Liturgy of Chrysostom, which is still in use in the Greek and Russian churches, has been already noticed in the proper place.\footnote{See above, § 99.}

Among the pupils and admirers of Chrysostom we mention as deserving of special notice two abbots of the first half of the fifth century: the elder Nilus of Sinai, who retired with his son from one of the highest civil stations of the empire to the contemplative solitude of Mount Sinai, while his wife and daughter entered a convent of Egypt;\footnote{Comp. S. P. N. Niliabbatis opera omnia, variorum curis, nempe Leonis Allatii, Petri Possini, etc., edita, nunc primum in unum collecta et ordinata, accurante \emph{J. P. Migne}, Par. 1860, 1 volume. (Patrol. Gr. tom. 79.)} and Isidore of Pelusium, or Pelusiota, a native of Alexandria, who presided over a convent not far from the mouth of the Nile, and sympathized with Cyril against Nestorius, but warned him against his violent passions.\footnote{Comp. S. Isidori PelusiotaeEpistolarum libri v, ed. \emph{Possinus} (Jesuit), republished by \emph{Migne}, Par. 1860. (Patrol. Gr. tom. 78, including the dissertation of H. Ag. Niemeyer: De Isid. Pel. vita, scriptis et doctrina, Hal. 1825.) It is not certain that Isidore was a pupil of Chrysostom, but he frequently mentions him with respect, and was evidently well acquainted with his writings. See the dissertation of Niemeyer, in Migne’s ed. p. 15 sq.} They are among the worthiest representatives of ancient monasticism, and, in a large number of letters and exegetical and ascetic treatises, they discuss, with learning, piety, judgment, and moderation, nearly all the theological and practical questions of their age.

\xsection{Cyril of Alexandria}

§ 171. Cyril of Alexandria.

I. S. Cyrillus, Alex. archiepisc.: Opera omnia, Gr. et Lat., cura et studio Joan. Auberti. Lutetiae, 1638, 6 vols. in 7 fol. The same edition with considerable additions by J. P. Migne, Petit-Montrouge, 1859, in 10 vols. (Patrol. Gr. tom. lxviii-lxxvii.). Comp. Angelo Mai’s Nova Bibliotheca Patrum, tom. ii. pp. 1–498 (Rom. 1844), and tom. iii. (Rom. 1845), where several writings of Cyril are printed for the first time, viz.: De incarnatione Domini; Explanatio in Lucam; Homiliae; Excerpta; Fragments of Commentaries on the Psalms, and the Pauline and Catholic Epistles. (These additional works are incorporated in Migne’s edition.) Cyrilli Commentarii in Lucca Evangelium quae supersunt, Syriace, e manuscriptis spud museum Britannicum edidit Rob. Payne Smith, Oxonii, 1858. The same also in an English version with valuable notes by R. P. Smith, Oxford, 1859, in 2 vols.

II. Scattered notices of Cyril in Socrates, Marius Mercator, and the Acts of the ecumenical Councils of Ephesus and Chalcedon. Tillemont: Tom. xiv. 267–676, and notes, pp. 747–795. Cellier: Tom. xiii. 241 sqq. Acta Sanctorum: Jan. 28, tom. ii. A. Butler: Jan. 28. Fabricius: Biblioth. Gr. ed. Harless, vol. ix. p. 446 sqq. (The Vita of the Bollandists and the Noticia literaria of Fabricius are also reprinted in Migne’s edition of Cyril, tom. i. pp. 1–90.) Schröckh Theil xviii. 313–354. Comp. also the Prefaces of Angelo Mai to tom. ii. of the Nova Bibl. Patrum, and of R. P. Smith to his translation of Cyril’s Commentary on Luke.

While the lives and labors of most of the fathers of the church continually inspire our admiration and devotion, Cyril of Alexandria makes an extremely unpleasant, or at least an extremely equivocal, impression. He exhibits to us a man making theology and orthodoxy the instruments of his passions.

Cyrillus became patriarch of Alexandria about the year 412. He trod in the footsteps of his predecessor and uncle, the notorious Theophilus, who had deposed the noble Chrysostom and procured his banishment; in fact, he exceeded Theophilus in arrogance and violence. He had hardly entered upon his office, when he closed all the churches of the Novatians in Alexandria, and seized their ecclesiastical property. In the year 415 he fell upon the synagogues of the very numerous Jews with armed force, because, under provocation of his bitter injustice, they had been guilty of a trifling tumult; he put some to death, and drove out the rest, and exposed their property to the excited multitude.

These invasions of the province of the secular power brought him into quarrel and continual contest with Orestes, the imperial governor of Alexandria. He summoned five hundred monks from the Nitrian mountains for his guard, who publicly insulted the governor. One of them, by the name of Ammon, wounded him with a stone, and was thereupon killed by Orestes. But Cyril caused the monk to be buried in state in a church as a holy martyr to religion, and surnamed him Thaumasios, the Admirable; yet he found himself compelled by the universal disgust of cultivated people to let this act be gradually forgotten.

Cyril is also frequently charged with the instigation of the murder of the renowned Hypatia, a friend of Orestes. But in this cruel tragedy he probably had only the indirect part of exciting the passions of the Christian populace which led to it, and of giving them the sanction of his high office.\footnote{Comp. above, § 6, p. 67, and Tillemont, tom. xiv. 274-’76. The learned, but superstitious and credulous Roman Catholic hagiographer, Alban Butler (Lives of the Saints, sub Jan. 28), considers Cyril innocent, and appeals to the silence of Orestes and Socrates. But Socrates, H. E. l. vii. c. 15, expressly says of this revolting murder: \textgreek{Τοῦτο οὐ μικρὸν μῶμον Κυρίλλῳ, καὶ τῇ τῶν Ἀλεξανδρέων ἐκκλησίᾳ εἰργάσατο}, and adds that nothing can be so contrary to the spirit of Christianity as the permission of murders and similar acts of violence. Walch, Schröckh, Gibbon, and Milman incline to hold Cyril responsible for the murder of Hypatia, which was perpetrated under the direction of a reader of his church, by the name of Peter. But the evidence is not sufficient. J. C. Robertson (History of the Christian Church, i. p. 401) more cautiously says: “That Cyril had any share in this atrocity appears to be an unsupported calumny; but the perpetrators were mostly officers of his church, and had unquestionably drawn encouragement from his earlier proceedings; and his character deservedly suffered in consequence.” Similarly W. Bright (A History of the Church from 313 to 451, p. 275): “Had there been no onslaught on the syna gogues, there would doubtless have been no murder of Hypatia.”}

From his uncle he had learned a strong aversion to Chrysostom, and at the notorious Synodus ad Quercum near Chalcedon, a.d. 403, he voted for his deposition. He therefore obstinately resisted the patriarchs of Constantinople and Antioch, when, shortly after the death of Chrysostom, they felt constrained to repeal his unjust condemnation; and he was not even ashamed to compare that holy man to the traitor Judas. Yet he afterwards yielded, at least in appearance, to the urgent remonstrances of Isidore of Pelusium and others, and admitted the name of Chrysostom into the diptychs\footnote{That is, the \textgreek{διπτυχα νεκρῶν}, or two-leaved tablets, with the list of names of distinguished martyrs and bishops, and other persons of merit, of whom mention was to be made in the prayers of the church. The Greek church has retained the use of diptychs to this day.} of his church (419), and so brought the Roman see again into communication with Alexandria.

From the year 428 to his death in 444 his life was interwoven with the Christological controversies. He was the most zealous and the most influential champion of the anti-Nestorian orthodoxy at the third ecumenical council, and scrupled at no measures to annihilate his antagonist. Besides the weapons of theological learning and acumen, he allowed himself also the use of wilful misrepresentation, artifice, violence, instigation of people and monks at Constantinople, and repeated bribery of imperial officers, even of the emperor’s sister Pulcheria. By his bribes he loaded the church property at Alexandria with debt, though he left considerable wealth even to his kindred, and adjured his successor, the worthless Dioscurus, with the most solemn religious ceremonies, not to disturb his heirs.\footnote{Dioscurus, however, did not keep his word, but extorted from the heirs of Cyril immense sums of money, and reduced them to extreme want. So one of Cyril’s relatives complained to the Conc, Chalc. Act. iii. in Hardouin, tom. ii. 406). A verification of the Proverb: Ill gotten, ill gone.}

His subsequent exertions for the restoration of peace cannot wipe these stains from his character; for he was forced to those exertions by the power of the opposition. His successor Dioscurus, however (after 444), made him somewhat respectable by inheriting all his passions without his theological ability, and by setting them in motion for the destruction of the peace.

Cyril furnishes a striking proof that orthodoxy and piety are two quite different things, and that zeal for pure doctrine may coëxist with an unchristian spirit. In personal character he unquestionably stands far below his unfortunate antagonist. The judgment of the Catholic historians is bound by the authority of their church, which, in strange blindness, has canonized him.\footnote{Even the monophysite Copts and Abyssinians celebrate his memory under the abbreviated name of Kerlos, and the title of Doctor of the World.} Yet Tillemont feels himself compelled to admit that Cyril did much that is unworthy of a saint.\footnote{Mémoires, xiv. 541: “S. Cyrille est Saint: mais on ne peut pas dire que toutes es actions soient saintes.”} The estimate of Protestant historians has been the more severe. The moderate and honest Chr. W. Franz Walch can hardly give him credit for anything good;\footnote{Comp. the description at the close of the fifth volume of his tedious but thorough Ketzerhistorie, where, after recounting the faults of Cyril, he exclaims, p. 932: “Can a man read such a character without a shudder? And yet nothing is fabricated here, nothing overdrawn; nothing is done but to collect what is scattered in history. And what is worst: I find nothing at all that can be said in his praise.” Schröckh(l.c. p. 352), in his prolix and loquacious way, gives an equally unfavorable opinion, and the more extols his antagonist Theodoret (p. 355 sqq.), who was a much more learned and pious man, but in his life-time was persecuted, and after his death condemned as a heretic, while Cyril was pronounced a saint.} and the English historian, H. H. Milman, says he would rather appear before the judgment-seat of Christ, loaded with all the heresies of Nestorius, than with the barbarities of Cyril.\footnote{History of Latin Christianity, vol. i. p. 210: “Cyril of Alexandria, to those who esteem the stern and uncompromising assertion of certain Christian tenets the one paramount Christian virtue, may be the hero, even the saint: but while ambition, intrigue, arrogance, rapacity, and violence, are proscribed as unchristian means—barbarity, persecution, bloodshed, as unholy and unevangelic wickednesses—posterity will condemn the orthodox Cyril as one of the worst heretics against the spirit of the Gospel. Who would not meet the judgment of the divine Redeemer loaded with the errors of Nestorius rather than the barbarities of Cyril?”}

But the faults of his personal character should not blind us to the merits of Cyril as a theologian. He was a man of vigorous and acute mind and extensive learning and is clearly to be reckoned among the most important dogmatic and polemic divines of the Greek church.\footnote{Baur(Vorlesungen über Dogmengeschichte, i. ii. p. 47) says of Cyril: “The current estimate of him is not altogether just. As a theologian he must be placed higher than he usually is. He remained true to the spirit of the Alexandrian theology, particularly in his predilection for the allegorical and the mystical, and he had a doctrine consistent with itself.”} Of his contemporaries Theodoret alone was his superior. He was the last considerable representative of the Alexandrian theology and the Alexandrian church, which, however, was already beginning to degenerate and stiffen; and thus be offsets Theodoret, who is the most learned representative of the Antiochian school. He aimed to be the same to the doctrine of the incarnation and the person of Christ, that his purer and greater predecessor in the see of Alexandria had been to the doctrine of the Trinity a century before. But he overstrained the supranaturalism and mysticism of the Alexandrian theology, and in his zeal for the reality of the incarnation and the unity of the person of Christ, he went to the brink of the monophysite error; even sustaining himself by the words of Athanasius, though not by his spirit, because the Nicene age had not yet fixed beyond all interchange the theological distinction between οὐσία and ὑπόστασις.\footnote{This is not considered by R. P. Smith, when, in the Preface to his English translation of Cyril’s Commentary on the Gospel of Luke from the Syriac (p. v.), he says, that Cyril never transcended Athanasius’ doctrine of \textgreek{μία φύσις τοῦ Θεοῦλόγου σεσαρκωμένη}, and that both are irreconcilable with the dogma of Chalcedon, which rests upon the Antiochian theology. Comp. §§ 137-140, above.}

And connected with this is his enthusiastic zeal for the honor of Mary as the virgin-mother of God. In a pathetic and turgid eulogy on Mary, which he delivered at Ephesus during the third ecumenical council, he piles upon her predicates which exceed all biblical limits, and border upon idolatry.\footnote{Encomium in sanctam Mariam Deiparam, in tom. v. Pars ii. p. 880 (in Migne’s ed. tom. x. 1029 sqq.).} “Blessed be thou,” says he, “O mother of God! Thou rich treasure of the world, inextinguishable lamp, crown of virginity, sceptre of true doctrine, imperishable temple, habitation of Him whom no space can contain, mother and virgin, through whom He is, who comes in the name of the Lord. Blessed be thou, O Mary, who didst hold in thy womb the Infinite One; thou through whom the blessed Trinity is glorified and worshipped, through whom the precious cross is adored throughout the world, through whom heaven rejoices and angels and archangels are glad, through whom the devil is disarmed and banished, through whom the fallen creature is restored to heaven, through whom every believing soul is saved.”\footnote{\textgreek{Δι  ἧς πᾶσα πνοὴ πιστεύουσα σώζεται}} These and other extravagant praises are interspersed with polemic thrusts against Nestorius.

Yet Cyril did not, like Augustine, exempt the Virgin from sin or infirmity, but, like Basil, he ascribed to her a serious doubt at the crucifixion concerning the true divinity of Christ, and a shrinking from the cross, similar to that of Peter, when he was scandalized at the bare mention of it, and exclaimed: “Be it far from thee, Lord!” (Matt. xvi. 22.) In commenting on John xix. 25, Cyril says: The female sex somehow is ever fond of tears,\footnote{\textgreek{Φιλόδακρυ}} and given to much lamentation .... It was the purpose of the holy evangelist to teach, that probably even the mother of the Lord Himself took offence\footnote{\textgreek{Εσκανδάλισε πάθος.}} at the unexpected passion; and the death upon the cross, being so very bitter, was near unsettling her from her fitting mind .... Doubt not that she admitted\footnote{\textgreek{Εἰσεδέξατο.}} some such thoughts as these: I bore Him who is laughed at on the wood; but when He said He was the true Son of the Omnipotent God, perhaps somehow He was mistaken.\footnote{\textgreek{Αλλ  υἱον ἑαυτὸν ἀληθινὸν εἶναι λέγων τοῦ πάντων κρατοῦντος Θεοῦ, τάχα που καὶ διεσφάλλετο.}} He said, ’I am the Life;’ how then has He been crucified? how has He been strangled by the cords of His murderers? how did He not prevail over the plot of His persecutors? why does He not descend from the cross, since He bade Lazarus to return to life, and filled all Judaea with amazement at His miracles? And it is very natural that woman,\footnote{Or woman’s nature, \textgreek{τὸ γύναιον}, which is sometimes used in a contemptuous sense, like the German \emph{Weibsbild}.} not knowing the mystery, should slide into some such trains of thought. For we should understand, that the gravity of the circumstances of the Passion was enough to overturn even a self-possessed mind; it is no wonder then if woman\footnote{\textgreek{Τὸ γύναιον.}} slipped into this reasoning.” Cyril thus understands the prophecy of Simeon (Luke ii. 35) concerning the sword, which, he says, “meant the most acute pain, cutting down the woman’s mind into extravagant thoughts. For temptations test the hearts of those who suffer them, and make bare the thoughts which are in them.”\footnote{Cyril, in Joann. lib. xii. (in Migne’s ed. of Cyril, vol. vii. col. 661 sq.). Dr. J. H. Newman(in his Letter to Dr. Pusey on his Eirenicon, Lond. 1866, p. 136) escapes the force of the argument of this and similar passages of Basil and Chrysostomagainst the Roman Mariolatry by the sophistical distinction, that they are not directed against the Virgin’s person, so much as against her nature (\textgreek{τὸ γύναιον}), of which the fathers had the low estimation then prevalent, looking upon womankind as the “varium et mutabile semper,” and knowing little of that true nobility which is exemplified in the females of the Germanic races, and in those of the old Jewish stock, Miriam, Deborah, Judith, Susanna. But it was to the human nature of Mary, and not to human nature in the abstract, that Cyril, whether right or wrong, attributed a doubt concerning the true divinity of her Son. I think there is no warrant for such a supposition in the accounts of the crucifixion, and the sword in the prophecy of Simeon means anguish rather than doubt. But this makes the antagonism of these Greek fathers with the present Roman Mariology only the more striking. Newman(l.c. p. 144) gratuitously assumes that the tradition of the sinlessness of the holy Virgin was obliterated and confused at Antioch and New Caesarea by the Arian troubles. But this would apply at best only to Chrysostomand Basil, and not to Cyril of Alexandria, who lived half a century after the defeat of Arianism at the second ecumenical council, and who was the leading champion of the \emph{theotokos} in the Nestorian controversy. Besides there is no clear trace of the doctrine of the sinlessness of Mary before St. Augustine, either among the Greek or Latin fathers; for the tradition of Mary as the second Eve does not necessarily imply that doctrine, and was associated in Irenaeus and Tertullianwith views similar to those expressed by Basil, Chrysostom, and Cyril. Comp. §§ 81 and 82, above.}

Aside from his partisan excesses, he powerfully and successfully represented the important truth of the unity of the person of Christ against the abstract dyophysitism of Nestorius.

For this reason his Christological writings against Nestorius and Theodoret are of the greatest importance to the history of doctrine.\footnote{Adversus Nestorii blasphemias contradictionum libri v (\textgreek{Κατὰ τῶν Νεστωρίου δυσφημιῶν πενταβίβλος ἀντίῤῥητος}); Explanatio xii capitum s. anathematismorum (\textgreek{Ἐπίλυσις τῶν δώδεκα κεφαλαίων}); Apologeticus pro xii capitibus adversus Orientales episcopos; Contra Theodoretum pro xii capitibus—all in the last volume of the edition of Aubert (in Migne, in tom. ix.).} Besides these he has left us a valuable apologetic work, composed in the year 433, and dedicated to the emperor Theodosius II., in refutation of the attack of Julian the Apostate upon Christianity;\footnote{Contra Julianum Apostatam libri x, tom. vi. in Aubert (tom. ix. in Migne); also in Spanheim’s Opera Juliani. Comp. §§ 4 and 9, above.} and a doctrinal work on the Trinity and the incarnation.\footnote{De S. Trinitate, et de incarnatione Unigeniti, etc., tom. v. Pars i. Not to be confounded with the spurious work Detrinitate, in tom. vi. 1-35, which combats the monothelite heresy, and is therefore of much later origin.} As an expositor he has the virtues and the faults of the arbitrary allegorizing and dogmatizing method of the Alexandrians, and with all his copiousness of thought he affords far less solid profit than Chrysostom or Theodoret. He has left extended commentaries, chiefly in the form of sermons, on the Pentateuch (or rather on the most important sections and the typical significance of the ceremonial law), on Isaiah, on the twelve Minor Prophets, and on the Gospel of John.\footnote{Tom. i.-iv.} To these must now be added fragments of expositions of the Psalms, and of some of the Epistles of Paul, first edited by Angelo Mai; and a homiletical commentary on the Gospel of Luke, which likewise has but recently become known, first by fragments in the Greek original, and since complete in a Syriac translation from the manuscripts of a Nitrian monastery.\footnote{By Angelo Mai and R. P. Smith. See the Literature above.} And, finally, the works of Cyril include thirty Easter Homilies (Homiliae paschales), in which, according to Alexandrian custom, he announced the time of Easter; several homilies delivered in Ephesus and elsewhere; and eighty-eight Letters, relating for the most part, to the Nestorian controversies.\footnote{The Homilies and Letters in tom. v. Pars ii. ed. Aubert (in Migne, with additions, in tom. x.).}

\xsection{Ephraem the Syrian}

§ 172. Ephraem the Syrian.

I. S. Ephraem Syrus: Opera omnia quae exstant Greece, Syriace, Latine, in sex tomos distributa, ad MSS. codices Vaticanos aliosque castigata, etc.: nunc primum, sub auspiciis S. P. Clementis XII. Pontificis Max. e Bibl. Vaticana prodeunt. Edited by the celebrated Oriental scholar J. S. Assemani (assisted by his nephew Stephen Evodius Assemani, 1732–’43, 6 vols. and the Maronite Jesuit Peter Benedict). Romae, fol. (vols. i.-iii. contain the Greek and Latin translations; vols. iv.-vi., which are also separately numbered i.-iii., the Syriac writings with a Latin version). Supplementary works edited by the Mechitarists, Venet. 1836, 4 vols. 8 vo. The hymns of Ephraem have also been edited by Aug. Hahn and Fr. L. Sieffert: Chrestomathia Syriaca sive S. Ephraemi carmina selecta, notis criticis, philologicis, historicis, et glossario locupletissimo illustr., Lips. 1825; and by Daniel: Thes. hymn. tom. iii. (Lips. 1855) pp. 139–268. German translation by Zingerle: Die heil. Muse der Syrer. Innsbruck, 1830. English translation by Henry Burgess: Select metrical Hymns and Homilies of Ephr. Syrus, transl. Lond. 1853, 2 vols. 12 mo. Comp. § 114, above.

II. Gregorius Nyss.: Vita et encomium S. Ephr. Syr. (in Opera Greg. ed. Paris. 1615, tom. ii. pp. 1027–1048; or in Migne’s ed. of Greg. tom. iii. 819–850, and in Ephr. Op. tom. i.). The Vita per Metaphrastem; several anonymous biographies; the Testimonia veterum and Judicia recentiorum; the Dissertation de rebus gestis, scriptis, editionibusque Ephr. Syr., etc., all in the first volume, and the Acta Ephraemi Syriaca auctore anonymo, in the sixth volume, of Assemani’s edition of the Opera Ephr. Jerome: Cat. vir. ill.c. 115. Sozomen: H. E. iii. c. 16; vi. 34. Theodoret: H. E. iv. 29. Acta Sanctorum for Fehr. i. (Antw. 1658), pp. 67–78. Butler: The Lives of the Saints, sub July 9. W. Cave: Lives of the Fathers, \&c. Vol. iii. 404–412 (Oxford ed. of 1840). Fabricius: Bibl. Gr. (reprinted in Assemani’s ed. of the Opera i. lxiii. sqq.). Lengerke: De Ephraemo Syro S. Scripturae interprete, Hal. 1828; De Ephr. arte hermeneutica, Regiom. 1831. Alsleben: Das Leben des h. Ephraem. Berlin, 1853. E. Rödiger: Art. Ephräm in Herzog’s Encykl. vol. iv. (1855), p. 85 ff.

Before we leave the Oriental fathers, we must give a sketch of Ephraem or Ephraim\footnote{The Greeks spell his name \textgreek{Ἐφραΐμ},,, the Latins Ephraem.} the most distinguished divine, orator, and poet, of the ancient Syrian church. He is called “the pillar of the church,” “the teacher,” “the prophet, of the Syrians,” and as a hymn-writer “the guitar of the Holy Ghost.” His life was at an early date interwoven with miraculous legends, and it is impossible to sift the truth from pious fiction.

He was born of heathen parents in Mesopotamia (either at Edessa or at Nisibis) in the beginning of the fourth century, and was expelled from home by his father, a priest of the god Abnil, for his leaning to Christianity.\footnote{This is the account of the Syriac Acta Ephraemi, in the sixth volume of the Opera p. xxiii sqq. But according to another account, which is followed by Butler and Cave, his parents were Christians, and dedicated him to God from the cradle.} He went to the venerated bishop and confessor Jacob of Nisibis, who instructed and probably also baptized him, took him to the council of Nicaea in 325, and employed him as teacher. He soon acquired great celebrity by his sacred learning, his zealous orthodoxy, and his ascetic piety. In 363, after the cession of Nisibis to the Persians, he withdrew to Roman territory, and settled in Edessa, which about that time became the chief seat of Christian learning in Syria.\footnote{On the early history of Christianity in Edessa, compare W. Cureton: Ancient Syriac Documents relative to the earliest Establishment of Christianity in Edessa and the neighboring Countries, from the Year after our Lord’s Ascension to the Beginning of the Fourth Century. Lond. 1866.} He lived a hermit in a cavern near the city, and spent his time in ascetic exercises, in reading, writing, and preaching to the monks and the people with great effect. He acquired complete mastery over his naturally violent temper, he denied himself all pleasures, and slept on the bare ground. He opposed the remnants of idolatry in the surrounding country, and defended the Nicene orthodoxy against all classes of heretics. He made a journey to Egypt, where he spent several years among the hermits. He also visited, by divine admonition, Basil the Great at Caesarea, who ordained him deacon. Basil held him in the highest esteem, and afterwards sent two of his pupils to Edessa to ordain him bishop; but Ephraem, in order to escape the responsible office, behaved like a fool, and the messengers returned with the report that he was out of his mind. Basil told them that the folly was on their side, and Ephraem was a man full of divine wisdom.

Shortly before his death, when the city of Edessa was visited by a severe famine, Ephraem quitted his solitary cell and preached a powerful sermon against the rich for permitting the poor to die around them, and told them that their wealth would ruin their soul, unless they made good use of it. The rich men felt the rebuke, and intrusted him with the distribution of their goods. Ephraem fitted up about three hundred beds, and himself attended to the sufferers, whether they were foreigners or natives, till the calamity was at an end. Then he returned to his cell, and a few days after, about the year 379, he expired, soon following his friend Basil.

Ephraem, says Sozomen, attained no higher clerical degree than that of deacon, but his attainments in virtue rendered him equal in reputation to those who rose to the highest sacerdotal dignity, while his holy life and erudition made him an object of universal admiration. He left many disciples who were zealously attached to his doctrines. The most celebrated of them were Abbas, Zenobius, Abraham, Maras, and Simeon, whom the Syrians regard as the glory of their country.\footnote{Sozomen, H. E. iii. 16. Cave(I. c. iii. 409) says of him: “He had all the virtues that can render a man great and excellent, and this that crowned all the rest, that he would not know it, nor cared to hear of it; being desirous, as Nyssen tells us, \textgreek{οὐ δοκεῖν, ἀλλ  εἶναι χρηστός}, not to seem, but to be really good.”}

Ephraem was an uncommonly prolific author. His fertility was prophetically revealed to him in his early years by the vision of a vine which grew from the root of his tongue, spreading in every direction to the ends of the earth, and was loaded with new and heavier clusters the more it was plucked. His writings consist of commentaries on the Scriptures, homilies, ascetic tracts, and sacred poetry. The commentaries and hymns, or metrical prose, are preserved in the Syriac original, and have an independent philological value for Oriental scholars. The other writings exist only in Greek, Latin, and Armenian translations. Excellent Greek translations were known and extensively read so early as the time of Chrysostom and Jerome. His works furnish no clear evidence of his knowledge of the Greek language; some writers assert his acquaintance with Greek, others deny it.\footnote{Sozomen and Theodoret expressly say that Ephraem was not acquainted with the Greek language, but used the Syriac “as a medium for reflecting the rays of divine grace.” According to the legend he was miraculously endowed with the knowledge of the Greek on his visit to Basil, who was in like manner inspired to greet him in Syriac.}

His commentaries extended over the whole Bible, “from the book of creation to the last book of grace,” as Gregory of Nyssa says. We have his commentaries on the historical and prophetical books of the Old Testament and the Book of Job in Syriac, and his commentaries on the Epistles of Paul in an Armenian translation.\footnote{Opera, tom. iv. and v., or vol. i. and ii. of the Opera Syr., and the supplements of the Mechitarists.} They have been but little used thus far by commentators. He does not interpret the text from the original Hebrew, but from the old Syriac translation, the Peshito.\footnote{He refers, however, occasionally to the original, as, for instance, ad Gen. i. 1: Interjecta particula \texthebrew{את }\par , quae in Hebraico textu hac loco legitur, idem valet, quod Syriacus articulus .” (Opera, vi. 116.) But such references prove no more than a superficial knowledge of Hebrew.}

His sermons and homilies, of which, according to Photius, he composed more than a thousand, are partly expository, partly polemical, against Jews, heathen, and heretics.\footnote{Opera, tom. i. ii. iii. and iv. Compare Photius, Bibl. Cod. 196.} They evince a considerable degree of popular eloquence; they are full of pathos, exclamations, apostrophes, antitheses, illustrations, severe rebuke, and sweet comfort, according to the subject; but also full of exaggerations, bombast, prolixity, and the superstitious of his age, such as the over-estimate of ascetic virtue, and excessive veneration of the Virgin Mary, the saints, and relics.\footnote{There is even a prayer to the holy Virgin (in Latin only) in his Works, tom. iii. p. 577; if it be genuine; for there are no other clear traces of such prayers before the fifth century. Mary is there addressed as “immaculata ... atque ab omni sorde ac labe peccati alienissima, virgo Dei sponsa, ac Domina nostra, ” etc.} Some of his sermons were publicly read after the Bible lesson in many Oriental and even Occidental churches.\footnote{Hieron, De script. eccl.c.115,}

His hymns were intended to counteract the influence of the heretical views of Bardesanes and his son Harmonius, which spread widely by means of popular Syrian songs. “When Ephraem perceived,” says Sozomen, “that the Syrians were charmed with the elegant diction and melodious versification of Harmonius, he became apprehensive, lest they should imbibe the same opinions; and therefore, although he was ignorant of Greek learning, he applied himself to the study of the metres of Harmonius, and composed similar poems in accordance with the doctrines of the church, and sacred hymns in praise of holy men. From that period the Syrians sang the odes of Ephraem, according to the method indicated by Harmonius.” Theodoret gives a similar account, and says, that the hymns of Ephraem combined harmony and melody with piety, and subserved all the purposes of valuable and efficacious medicine against the heretical hymns of Harmonius. It is reported that he wrote no less than three hundred thousand verses.\footnote{Sozomen, iii. 16: \textgreek{τριακοσια μυριάδα ἔπων,} —\textgreek{ἔπη}and \textgreek{στίχοι}is equivalent to verses or lines. Origen says of the Book of Job that it contains nearly 10,000 \textgreek{ἔπη}.} But, with the exception of his commentaries, all his Syriac works are written in verse, i.e., in lines of an equal number of syllables, and with occasional rhyme and assonance, though without regular metre.\footnote{Comp. Rödiger, in Herzog’s Encycl. vol. iv. p. 89, and the Observationes prosodicae of Hahnand Sieffertin their Chrestomathia Syriaca.}

II.—The Latin Fathers.

\xsection{Lactantius}

§ 173. Lactantius.

I. Lactantius, Lucius Caecilius Firmianus: Opera. First edition in venerabili monasterio Sublacensi, 1465. (Brunet: “Livre précieux, qui est en même temps la première édition de Lactance, et le premier ouvrage impr. en Italia avec date.”) Later editions by J. L. Brünemann, Lips. 1739; Le Brun and N. Lenglet Du Fresnoy, Par. 1748, 2 vols. 4to; F. E. a S. Xaverio, Rom. 1754–’9, and Migne, Par. 1844, in 2 vols. A convenient manual edition by O. Fridol. Fritzsche, in Gersdorf’s Bibliotheca Patrum ecclesiast. selecta, Lips. 1842, vol. x. and xi.

II. The introductory essays to the editions. Jerome: Cat. vir. illustr. c. 80. Notices in Dupin, Ceillier, Cave (Vol. iii. pp. 373–384), Schönemann (Biblioth. Patr. Lat. i. 177 sqq.), \&c. Möhler: Patrologie, i. pp. 917–933. On the Christology of Lactantius, comp. Dorner: Entwicklungsgeschichte der Lehre Von der Person Christi. Th. i. p. 761 ff.

Firmiamus Lactantius stands among the Latin fathers, like Eusebius among the Greek, on the border between the second period and the third, and unites in his reminiscences the personal experience of both the persecution and the victory of the church in the Roman empire; yet in his theological views he belongs rather to the ante-Nicene age.

According to his own confession he sprang from heathen parents. He was probably, as some have inferred from his name, a native of Firmum (Fermo) in Italy; he studied in the school of the rhetorician and apologist Arnobius of Sicca, and on this account has been taken by some for an African; he made himself known by a poetical work called Symposion, a collection of a hundred riddles in hexameters for table amusement; and he was called to Nicomedia by Dioclesian to teach Latin eloquence. But as this city was occupied mostly by Greeks, he had few hearers, and devoted himself to authorship.\footnote{He says of his heathen life, Inst. div. i. 1, that he trained youth by his “non ad virtutem, sed plane ad argutam malitiam.”} In his manhood, probably shortly before or during the persecution under Dioclesian, he embraced Christianity; he was witness of the cruel scenes of that persecution, though not himself a sufferer in it; and he wrote in defence of the hated and reviled religion.

Constantine subsequently (after 312) brought him to his court in Gaul, and committed to him the education of his son Crispus, whom the emperor caused to be executed in 326. At court he lived very simply, and withstood the temptations of luxury and avarice. He is said to have died in the imperial residence at Treves at a great age, about the year 330.

Jerome calls Lactantius the most learned man of his time.\footnote{Catal.c.80: “Lact. vir omnium suo tempore eruditissimus.” In Ep. 58 ad Paulinum (ed. Vall.), c. 10, he gives the following just view of him: “Lact. quasi quidam fluvius eloquentiae Tullianae, utinam tam nostra affirmare potuisset, quam facile aliena destruxit.” O. Friedol. Fritzsche, in the Praefatio of his edition of his Opera, thus estimates him: “Firm. Lactantius, qui Ciceronis felicissimus exstitit imitator, non solum sermonis castitate et elegantia orationisque flumine, sed, qua erat summa eruditione, rerum etiam copia et varietate inter reliquos ecclesiae latinae scriptores maxime eminuit, eoque factum est, ut, quamvis doctrinam ejus non satis esse sanam viros pios haud lateret, nunquam tamen prorsus negligeretur.”} His writings certainly give evidence of varied and thorough knowledge, of fine rhetorical culture, and particularly of eminent power of statement in clear, pure, and elegant style. In this last respect he surpasses almost all the Latin fathers, except Jerome, and has not unjustly been called the Christian Cicero.\footnote{Or, as Jerome, l. c., calls him: “Fluvius eloquentiae Tullianae.”} His is the famous derivation of the word religion from religare, defining it as the reunion of man with God, reconciliation; answering to the nature of Christianity, and including the three ideas of an original unity, a separation by sin, and a restoration of the unity again.\footnote{Instit. div. l. iv. cap. 28 (vol. i. p. 223, ed. Fritzsche): “Hoc vinculo pietatis obstricti Deo et religati sumus; unde ipsa religio nomen accepit, non ut Cicero interpretatus est, a relegendo” Cicero says, De natura deorum, ii. 28: “Qui omnia quae ad cultum deorum pertinerent, diligenter retmetarent et tamquam relegerent, \emph{religiosi} dicti sunt ex \emph{relegendo}, ut elegantes ex eligendo, itemque ex diligendo diligentes.” This derivation is not impossible, since we have \emph{legio} from \emph{legere}, and several nouns ending in io from verbs of the third conjugation, as \emph{regio, contagio, oblivio}. But the derivation of Lactantius gives a more correct and profound idea of religion, and etymologically it is equally admissible; for although \emph{religare} would rather yield the noun \emph{religatio}, yet we have \emph{optio} from \emph{optare}, \emph{rebellio} from \emph{rebellare} \emph{internecio}from \emph{internecare}, \&c. Augustine(Retract. i. 13), Jerome(Ad Amos, c. 9), and the majority of Christian divines have adopted the definition of Lactantius.}

But he is far more the rhetorician than the philosopher or theologian, and, as Jerome observes, has greater skill in the refutation of error than in the establishment of truth. The doctrinal matter of his writings, as in the case of his preceptor Arnobius, is very vague and unsatisfactory, and he does not belong to the narrower circle of the fathers, the authoritative teachers of the church. Pope Gelasius counted his works among the apocrypha, i.e., writings not ecclesiastically received.

Notwithstanding this, his Institutes, on account of their elegant style, have been favorite reading, and are said to have appeared in more than a hundred editions. His mistakes and errors in the exposition of points of Christian doctrine do not amount to heresies, but are mostly due to the crude and unsettled state of the church doctrine at the time. In the doctrine of sin he borders upon Manichaeism. In anthropology and soteriology he follows the synergism which, until Augustine, was almost universal. In the doctrine of the Trinity he was, like most of the ante-Nicene fathers, subordinationist. He taught a duplex nativitas of Christ, one at the creation, and one at the incarnation. Christ went forth from God at the creation, as a word from the mouth, yet hypostatically.\footnote{According to a statement of Jerome(Ep. 41 ad Pammach. et Ocean.) he denied the personality of the Holy Ghost.}

His most important work is his Divine Institutes, a comprehensive refutation of heathenism and defence of Christianity, designed to make Christianity better known among the cultivated classes, and to commend it by scholarship and attractive style.\footnote{Institutionum divinarum libri vii. The title was chosen with reference to the Institutiones juris civilis (i. 1). The several books then bear the following superscriptions: 1. De falsa religione; 2. De origine erroris; 3. De fan sapientia; 4. De vera sapientia; 5. De justitia; 6. De vero cultu; 7. De vita beata. Lactantius himself made an abstract of it under the title: Epitome ad Pentadium fratrem, in Fritzsche, Pars ii. pp. 114-171.} He seems to have begun the work during the Dioclesianic persecution, but afterwards to have enlarged and improved it about the year 321; for he dedicated it to the emperor, whom he celebrates as the first Christian prince.\footnote{L. i. c. 1: “Quod opus nunc nominis tui auspicio inchoamus, Constantineimperator maxims, qui primus Romanorum principum, repudiatis erroribus, majestatem Dei singularis ac veri cognovisti et honorasti, ” \&c. This passage, by the way, does not appear in all the codices. Comp. the note in the ed. of Fritzsche, Pars i. p. 3.}

To the same apologetic purpose was his work De morte, or mortibus, persecutorum, which is of some importance to the external history of the church.\footnote{In the ed. of Fritzsche, P. ii. pp. 248-286. This work is wanting in the earlier editions, and also in several manuscripts, and is therefore sometimes denied to Lactantius, e.g., by Dom de Nourry, in a learned dissertation on this question, reprinted in the Appendix to the second volume of Migne’s edition of Lactantius, p. 839 sqq. But its style, upon the whole, agrees with his; the work entirely suits his time and circumstances; and it is probably the same that Jeromecites under the name De persecutions. Jac. Burckhardt, in his monograph on Constantinethe Great, 1853, treats this book throughout as an untrustworthy romance, but without proof, and with an obvious aversion to all the fathers, similar to that of Gibbon.} It describes with minute knowledge, but in vehement tone, the cruel persecutions of the Christians from Nero to Dioclesian, Galerius, and Maximinus (314), and the divine judgments on the persecutors, who were compelled to become involuntary witnesses to the indestructible power of Christianity.

In his book De opificio Dei\footnote{In the ed. of Fritzsche, Pars ii. pp. 172-208.} he gives observations on the organization of the human nature, and on the divine wisdom displayed in it.

In the treatise De ira Dei\footnote{Ibid. ii. 208-247.} he shows that the punitive justice of God necessarily follows from his abhorrence of evil, and is perfectly compatible with his goodness; and he closes with an exhortation to live such a life that God may ever be gracious to us, and that we may never have to fear his wrath.

We have also from Lactantius various Fragmenta and Carmina de Phoenice, de Passione Domini, de resurrectione Domini, and one hundred Aenigmata, each of three hexameters.\footnote{Ibid. ii. p. 286 sqq. Other works of Lactantius, cited by Jerome, are lost.}

\xsection{Hilary of Poitiers}

§ 174. Hilary of Poitiers.

I. S. Hilarius Pictaviensis: Opera, studio et labore monach. S. Benedicti e congreg. S. Mauri. Paris, 1693, 1 vol. fol. The same ed. enlarged and improved by Scip. Maffei, Verona, 1730, 2 vols. fol. (reprinted in Venice, 1749). Am ed. by Fr. Overthür, Wirceburgi, 1785–’88, 4 vols.; and one by Migne, Petit-Montrouge, 1844–’45, in 2 vols. (Patrol. Lat. tom. ix. and x.).

II. The Praefatio et Vitae in the first vol. of the ed. of Maffei, and Migne (tom. i. 125 sqq.). Hieronymus: De viris illustr. c. 100. Tillemont (tom. vii.); Ceillier (tom. v.); and Butler, sub Jan. 14. Kling, in Herzog’s Encykl. vi. 84 ff. On the Christology of Hilary, comp. especially Dorner, Entwicklungsgeschichte, i. 1037 ff.

Hilary of Poitiers, or Pictaviensis, so named from his birth-place and subsequent bishopric in Southwestern France, and so distinguished from other men of the same name,\footnote{As Hilarius Arelatensis († 449), celebrated for his contest with pope Leo I.} was especially eminent in the Arian controversies for his steadfast confession and powerful defence of the orthodox faith, and has therefore been styled the “Athanasius of the West.”

He was born towards the end of the third century, and embraced Christianity in mature age, with his wife and his daughter Apra.\footnote{We have from him an Epistola ad Apram (or Abram in other manuscripts), filiam suam, written in 358, in tom. ii. 549 (ed. Migne). He sent to her his famous morning hymn: “Lucia largitor splendide.”} He found in the Holy Scriptures the solution of the riddle of life, which he had sought in vain in the writings of the philosophers. In the year 350 he became bishop of his native city, and immediately took a very decided stand against Arianism, which was at that time devastating the Gallic church. For this he was banished by Constantius to Phrygia in Asia Minor, where Arianism ruled. Here, between 356 and 361, he wrote his twelve books on the Trinity, the main work of his life.\footnote{De trinitate libri xii. (tom. i. 26-472, ed. Migne).} He was recalled to Gaul, then banished again, and spent the last years of his life in rural retirement till his death in 368.

We have from him, besides the theological work already mentioned several smaller polemic works against Arianism, viz., On Synods, or the Faith of the Orientals (358); fragments of a history of the Synod of Ariminum and Seleucia; a tract against the Arian emperor Constantius, and one against the Arian bishop Auxentius of Milan. He wrote also Commentaries on the Psalms (incomplete), and the Gospel of Matthew, which are partly a free translation of Origen,\footnote{Jerome(De viris illustr. c. 100) says of his Commentary on the Psalms: “In quo opere imitatus Origenem, nonnulla etiam de suo addidit,” and of the Commentary on Matthew and the tract on Job: “Quos de Graeco Origenis ad sensum transtulit.”} and some original hymns, which place him next to Ambrose among the lyric poets of the ancient church.

Hilary was a man of thorough biblical knowledge, theological depth and acuteness, and earnest, efficient piety. He had schooled himself in the works of Origen and Athanasius, but was at the same time an independent thinker and investigator. His language is often obscure and heavy, but earnest and strong, recalling Tertullian. He had to reproduce the profound thoughts of Athanasius and other Greek fathers in the Latin language, which is far less adapted to speculation than the copious, versatile, finely-shaded Greek. The incarnation of God was to him, as it was to Athanasius, the centre of theology and of the Christian life. He had an effective hand in the development of the dogma of the consubstantiality of the Son with the Father, and the dogma of the person of Christ. In this he was specially eminent for his fine use of the Gospel of John. But he could not get clear of subordinationism, nor call the Holy Ghost downright God. His Pneumatology, as well as his anthropology and soteriology, was, like that of all the fathers before Augustine, comparatively crude. In Christology he saw farther and deeper than many of his contemporaries. He made the distinction clear between the divine and the human in Christ, and yet held firmly to the unity of His person. He supposes a threefold birth of the Son of God: the eternal generation in the bosom of the Father, to whom the Son is equal in essence and glory; the incarnation, the humiliation of Himself to the form of a servant from the free impulse of love; and the birth of the Son of God out of the Son of Man in the resurrection, the transfiguration of the form of a servant into the form of God, at once showing forth again the full glory of God, and realizing the idea of humanity.\footnote{Klingsays, l.c. p. 94: ”Hilaryholds a most important place in the development of Christology, and his massive analysis contains fruitful germs which in the succeeding centuries have been only in part developed; profound and comprehensive thoughts, the stimulating and fertilizing power of which reaches down even into our own time; nor need our time be ashamed to learn from this ancient master, as well as from other teachers of that age.”}

\xsection{Ambrose}

§ 175. Ambrose.

I. S. Ambrosius Mediolanensis episcopus: Opera ad manuscriptos codices Vaticanos, Gallicanos, Belgicos, \&c., emendata, studio et labore monachorum ord. S. Benedicti e congreg. S. Mauri (Jac. du Fricke et Nic. de Nourry). Paris. 1686–’90, 2 vols. fol. This edition was reprinted at Venice, 1748–’51, in 4 vols. fol., and in 1781 in 8 vols. 4to, and by Abbé Migne in his Patrol., Petit-Montrouge, 1843, 2 tom. in 4 Parts with some additions. The Libri tres de officiis, and the Hexaëmeron of Ambrose have also been frequently published separately. A convenient edition of both is included in Gersdorf’s Bibliotheca Patrum Latinorum selecta, vols. viii. and ix. Lips. 1839. His hymns are found also in Daniel’s Thesaurus hymnolog tom. i. p. 12 sqq.

II. Paulinus (deacon of Milan and secretary of Ambrose): Vita S. Ambrosii (written by request of St. Augustine, derived from personal knowledge, from Marcella, sister of Ambrose, and several friends). The Vita of an anonymous writer, in Greek and Latin, in the Bened. ed. of the Opera. Both in the Appendix to tom. ii. ed. Benedictinae. Benedictini Editores: Vita Ambrosii ex ejus potissimum scriptis collecta et secundum chronologiae ordinem digesta, in the Bened. ed., in the Appendix to tom. ii., and in Migne’s reprint, tom. i. (very thorough and instructive). Comp. also the Selecta veterum testimonia de S. Ambr. in the same editions. The biographies of Hermant (1678), Tillemont (tom. x. pp. 78–306), Vagliano (Sommario degli archivescovi di Milano), Butler (sub Dec. 7), Schröckh, Böhringer, J. P. Silbert (Das Leben des heiligen Ambrosius, Wien, 1841).

Ambrose, son of the governor (praefectus) of Gaul, which was one of the three great dioceses of the Western empire, was born at Treves (Treviri) about 340, educated at Rome for the highest civil offices, and after greatly distinguishing himself as a rhetorician, was elected imperial president (praetor) of Upper Italy; whereupon Probus, prefect of Italy, gave him the remarkable advice, afterwards interpreted as an involuntary prophecy: “Go, and act not the judge, but the bishop.” He administered this office with justice and mildness, enjoying universal esteem.

The episcopal chair of Milan, the second capital of Italy, and frequently the residence of the emperors, was at that time occupied by the Cappadocian, Auxentius, the head of the Arian party in the West. Soon after the arrival of Ambrose, Auxentius died. A division then arose among the people in the choice of a successor, and a dangerous riot threatened. The governor considered it his duty to allay the storm. But while he was yet speaking to the people, the voice of a child suddenly rang out: “Let Ambrose be bishop!” It seemed a voice of God, and Arians and Catholics cried, Amen.

Ambrose was at that time a catechumen, and therefore not even baptized. He was terrified, and seized all possible, and even most eccentric, means to escape the responsible office. He was obliged to submit, was baptized, and eight days afterwards, in 374, was consecrated bishop of Milan. His friend, Basil the Great of Caesarea, was delighted that God had chosen such a man to so important a post, who counted noble birth, wealth, and eloquence loss, that he might win Christ.

From this time forward Ambrose lived wholly for the church, and became one of the greatest bishops of ancient Christendom, full of Roman dignity, energy, and administrative wisdom, and of the unction of the Holy Ghost. He began his work with the sale of his great estates and of his gold and silver for the benefit of the poor; reserving an allowance for his pious sister Marcella or Marcellina, who in early youth had taken the vow of virginity. With voluntary poverty he associated the strictest regimen of the ascetic spirit of his time; accepted no invitations to banquets; took dinner only on Sunday, Saturday, and the festivals of celebrated martyrs; devoted the greater part of the night to prayer, to the hitherto necessarily neglected study of the Scriptures and the Greek fathers, and to theological writing; preached every Sunday, and often in the week; was accessible to all, most accessible to the poor and needy; and administered his spiritual oversight, particularly his instruction of catechumens, with the greatest fidelity.

The Arians he vigorously opposed by word and act, and contributed to the victory of the Nicene faith in the West. In this work he behaved himself towards the Arian empress Justina with rare boldness, dignity, and consistency, in the heroic spirit of an Athanasius. The court demanded the cession of a catholic church for the use of the Arians, and claimed for them equal rights with the orthodox. But Ambrose asserted the entire independence of the church towards the state, and by perseverance came off victorious in the end. It was his maxim, that the emperor is in the church, but not over the church, and therefore has no right to the church buildings.

He did not meddle in secular matters, nor ask favor of the magistracy, except when he could put in a word of intercession for the unfortunate and for persons condemned to death in those despotic times. This enabled him to act the more independently in his spiritual office, as a real prince of the church, fearless even of the emperor himself. Thus he declared to the usurper Maximus, who desired church fellowship, that he would never admit him, unless he should do sincere penance for the murder of the emperor Gratian.

When the Roman prefect, Symmachus, the noblest and most eloquent advocate of the decaying heathenism of his time, implored the emperor Valentinian, in an apology for the altar of Victory which stood in the hall of the Roman senate, to tolerate the worship and the sanctuaries of the ancient gods, Ambrose met him with an admirable reply, and prevented the granting of his request.

The most imposing appearance of our bishop against the temporal power was in his dealing with Theodosius, when this truly great, but passionate and despotic, emperor, enraged at Thessalonica for a riot, had caused many thousand innocent persons to be put to death with the guilty, and Ambrose, interesting himself for the unfortunate, like a Nathan with David, demanded repentance of the emperor, and refused him the holy communion. “How wilt thou,” said he to him in the vestibule of the church, “how wilt thou lift up in prayer the hands still dripping with the blood of the murdered? How wilt thou receive with such hands the most holy body of the Lord? How wilt thou bring to thy mouth his precious blood? Get thee away, and dare not to heap crime upon crime.” When Theodosius appealed to David’s murder and adultery, the bishop answered: “Well, if thou hast imitated David in sin, imitate him also in repentance.”\footnote{“Qui sequutus es errantem, sequere corrigentem” Paulinus, Vita Ambr. c. 24.} The emperor actually submitted to ecclesiastical discipline, made public confession of his sin, and did not receive absolution until he had issued a law that the sentence of death should never be executed till thirty days after it was pronounced.\footnote{Paulinus, l. c. c. 24: “Quod ubi audivit clementissimus imperator, ita suscepit, ut publicam poenitentiam non abhorreret,” \&c. Ambrosehimself says in his funeral oration on Theodosius: “Stravit omne, quo utebatur insigne regium, deflevit in ecclesia publice peccatum suum, neque ullus postea dies fuit, quo non illum doleret errorem.” The main fact is beyond doubt; but the details are not all reliable, and may have been exaggerated for hierarchical ends.}

From this time the relation between Ambrose and Theodosius continued undisturbed, and the emperor is reported to have said afterwards with reference to the bishop, that he had recently found the first man who told him the truth, and that he knew only one man who was worthy to be bishop. He died in the arms of Ambrose at Milan in 395. The bishop delivered his funeral oration in which he tells, to his honor, that on his dying bed he was more concerned for the condition of the church than for himself, and says to the soldiers: “The faith of Theodosius was your victory; let your truth and faith be the strength of his sons. Where unbelief is, there is blindness, but where fidelity is, there is the host of angels.”

Two years after this, Ambrose himself was fatally sick. All Milan was in terror. When he was urged to pray God for a lengthening of his life, he answered: “I have so lived among you that I cannot be ashamed to live longer; but neither do I fear to die; for we have a good Lord.” During his sickness he had miraculous intimations and heard heavenly voices, and he himself related that Christ appeared to him smiling. His notary and biographer, the deacon Paulinus, who adorns his life throughout with miraculous incidents, tells us:\footnote{Vita Ambr. c. 42.} “Not long before his death, while he was dictating to me his exposition of the Forty-third Psalm, I saw upon his head a flame in the form of a small shield; hereupon his face became white as snow, and not till some time after did it return to its natural color.” In the night of Good Friday, on Saturday, the 4th of April, 397, he died, at the age of fifty-seven years, having first spent several hours, with his hands crossed, in uninterrupted prayer. Even Jews and pagans lamented his death. On the night of Easter following many were baptized in the church where his body was exposed Not a few of the newly baptized children saw him seated in the episcopal chair with a shining star upon his head. Even after his death he wrought miracles in many places, in proof of which Paulinus gives his own experience, credible persons, and documents.

Ambrose, like Cyprian before him, and Leo I. after him, was greatest in administration. As bishop he towered above the contemporary popes. As a theologian and author he is only a star of the second magnitude among the church fathers, yielding by far to Jerome and Augustine. We have from this distinguished prelate several exegetical, doctrinal, and ascetic works, besides homilies, orations, and letters. In exegesis he adopts the allegorical method entire, and yields little substantial information. The most important among his exegetical works are his homilies on the history of creation (Hexaëmeron, written 389), an Exposition of twenty-one Psalms (390–397), and a Commentary on the Gospel of Luke (386).\footnote{The exegetical works are in tom. i. of the Bened. ed., excepting Ambrosiaster, which is in the Appendix to tom. ii. Jeromehad a contemptuous opinion of his exegetical writings. In the preface to his translation of the thirty-nine Homilies of Origen on Luke, he compares the superficial and meagre Commentary of Ambroseon Luke to the croaking of a raven which makes sport of the colors of all other birds, and yet is itself dark all over (totus ipse tenebrosus). Against this attack Rufinus felt it his duty to defend Ambrose, “qui non solum Mediolanensis ecclesiae, verum etiam omnium ecclesiarum columna quaedam et turris inexpugnabilis fuit” (Invect. ii. adv. Hieron.). In his Catalogus vir. illustr. c. 124, Jeromedisposes of Ambrosewith the following frosty and equivocal notice: “Ambrosius Mediolanensis episcopus, usque in presentem diem scribit, de quo, quia superest, meum judicium subtraham, ne in alterutram partem aut adulatio in me reprehendatur, aut veritas.” In his Epistles, however, he occasionally makes favorable allusion to his ascetic writings which fell in with his own taste. Augustine, from a sense of gratitude to his spiritual father, always mentions his name with respect. The passages of Augustineon Ambroseare collected in the Selecta veterum testimonia at the beginning of the first tome of the Bened. edition. But the unfavorable notice of Jeromequoted above is omitted there.} The Commentary on the Pauline Epistles (Ambrosiaster so called or Pseudo-Ambrosius) which found its way among his works, is of uncertain authorship, perhaps the work of the Roman deacon Hilary under pope Damasus, and resembles in many respects the commentaries of Pelagius. Among his doctrinal writings his five books On Faith, three On the Holy Ghost, and six On the Sacraments (catechetical sermons on baptism, confirmation, and the eucharist) are worthy of mention. Among his ethical writings the work On Duties is the most important. It resembles in form the well-known work of Cicero on the same subject, and reproduces it in a Christian spirit. It is a collection of rules of living for the clergy, and is the first attempt at a Christian doctrine of morals, though without systematic method.\footnote{De officiis ministrorum in three books (in the Bened. ed. tom. ii. f. 1-142). Comp. F. Hassler: Ueber das Verhältniss der heidnischen und christlichen Ethik auf Grund einer Vergleichung des ciceronianischen Buches \emph{De officiis} mit dem gleichnamigen des heiligen Ambrosius, München, 1866.} Besides this he composed several ascetic essays: Three books on Virgins; On Virginity; On the Institution of the Virgin; On Exhortation to Virginity; On the Fall of a Consecrated Virgin, \&c., which contributed much to the spread of celibacy and monastic piety. Of his ninety-one Epistles several are of considerable historical interest.

In his exegesis and in his theology, especially in the doctrine of the incarnation and the Trinity, Ambrose is entirely dependent on the Greek fathers; most on Basil, whose Hexaëmeron he almost slavishly copied. In anthropology he forms the transition from the Oriental doctrine to the system of Augustine, whose teacher and forerunner he was. He is most peculiar in his ethics, which he has set forth in his three books De Officiis. As a pulpit orator he possessed great dignity, force, and unction, and made a deep impression on Augustine, to whose conversion he contributed a considerable share. Many mothers forbade their daughters to hear him lest he should induce them to lead a life of celibacy.

Ambrose has also a very important place in the history of worship, and did immortal service for the music and poetry of the church, as in a former section we have seen.\footnote{Paulinus, in Vita Ambr. c. 13, relates: “Hoc in tempore primum antiphonae hymni ac vigiliae in ecclesia Mediolanensi celebrari coeperunt. Cuius celebritatis devotio usque in hodiernum diem non solum in eadem ecclesia, verum per omnes pene occidentis provincias manet.”} Here again, as in theology and exegesis, he brought over the treasures of the Greek church into the Latin. The church of Milan uses to this day a peculiar liturgy which is called after him the ritus Ambrosianus.

\xsection{Jerome as a Divine and Scholar}

§ 176. Jerome as a Divine and Scholar.

Comp. the Literature at § 41; and especially the excellent monograph (which has since reached us) of Prof. Otto Zöckler: Hieronymus. Sein Leben und Wirken aus seinen Schriften dargestellt. Gotha, 1865.

Having already sketched the life and character of Jerome (born about 340, died in 419) in connection with the history of monasticism, we limit ourselves here to his theological and literary labors, in which he did his chief service to the church, and has gained the greatest credit to himself.

Jerome is the most learned, the most eloquent, and the most interesting author among the Latin fathers. He had by nature a burning thirst for knowledge,\footnote{As he himself says, Ep. 84, c. 3 (Opera, ed. Vallarsi, tom. i. 523): “Dum essem juvenis, miro discendi ferebar ardore, nee juxta quorundam praesumptionem ipse me docui.”} and continued unweariedly teaching, and learning, and writing, to the end of a very long life.\footnote{Sulpicius Severus, who describes from his own observation the learned seclusion of the aged Jeromeat Bethlehem, where, however, he was much interrupted and stimulated by the visits of Christians from all parts of the world, says of him, in Dial. i. 4: “Totus semper in lectione, totus in libris est; non die, non nocte requiescit; aut legit aliquid semper, aut scribit, ” \&c.} His was one of those intellectual natures, to which reading and study are as indispensable as daily bread. He could not live without books. He accordingly collected, by great sacrifices, a library for that time very considerable and costly, which accompanied him on his journeys.\footnote{He confesses that the purchase of the numerous works of Origen had exhausted his purse, Ep. 84, c. 3 (tom. i. 525): “Legi, inquam, legi Origenem, et, si in legendo crimen est, fateor; et nostrum marsupium Alexandrinae chartae evacuarunt.” When he saw, and was permitted to use, the library of Pamphiltus in Caesarea, with all the works of Origen, he thought he possessed more than the riches of Croesus (De viris illustr. c. 75).} He further availed himself of the oral instruction of great church teachers, like Apollinaris the Elder in Laodicea, Gregory Nazianzen in Constantinople, and Didymus of Alexandria, and was not ashamed to become an inquiring pupil in his mature age. His principle in studying was, in his own words: “To read the ancients, to test everything, to hold fast the good, and never to depart from the catholic faith.”\footnote{“Meum propositum est, antiquos legere, probare singula, retinere quae bona sunt, et a fide catholica numquam recedere.”}

Besides the passion for knowledge, which is the mother of learning, he possessed a remarkable memory, a keen understanding, quick and sound judgment, an ardent temperament, a lively imagination, sparkling wit, and brilliant power of expression. He was a master in all the arts and artifices of rhetoric, and dialectics. He, far more than Lactantius, deserves the name of the Christian Cicero, though he is inferior to Lactantius in classic purity, and was not free from the faulty taste, of his time. Tertullian had, indeed, long before applied the Roman language as the organ of Christian theology; Cyprian, Lactantius, Hilary, and Ambrose, had gone further on the same path; and Augustine has enriched the Christian literature with a greater number of pregnant sentences than all the other fathers together. Nevertheless Jerome is the chief former of the Latin church language, for which his Vulgate did a decisive and standard service similar to that of Luther’s translation of the Bible for German literature, and that of the authorized English Protestant version for English.\footnote{Ozanam(Histoire de la civilisationchrét. au 5. siècle, ii. 100) calls Jerome, ”Le maître de la prose chrétienne pour tous lea siècles suivants.” ZöcklerSays (l. a. p. 323): “As Cicero raised the language of his time to the classic grade, and cast it for all times in a model form, so, of the Western church fathers, Jeromewas the one to \emph{make the Latin language Christian, and Christian theology Latin}.” Erasmus placed him as an author in several respects even above Cicero.}

His scholarship embraced the Latin, Greek, and Hebrew languages and literature; while even Augustine had but imperfect knowledge of the Greek, and none at all of the Hebrew. Jerome was familiar with the Latin classics, especially with Cicero, Virgil, and Horace;\footnote{Virgil is quoted in the Letters of Jeromesome fifty times, in his other works much more frequently; Horace, in the Letters, some twenty times; of the prose writers Cicero more than all, next to him Varro, Sallust, Quintilian, Seneca, Suetonius, and Pliny. Virgil, however, is viewed by Jerome, and by Augustine, who likewise admired him greatly, simply as a great poet, and not, as he afterwards came to be considered in the Latin church, especially through the influence of Dante’s Divina Commedia, as a divine and prophet of heathenism.} and even after his famous anti-Ciceronian vision (which transformed him from a more or less secular scholar into a Christian ascetic and hermit) he could not entirely cease to read over the favorite authors of his youth, or at least to quote them from his faithful memory; thus subjecting himself to the charge of inconsistency, and even of perjury, from Rufinus.\footnote{Comp. § 41 above, and Zöcklerl.c. p. 45 ff., 156, and 325. It is certain that Jerome, after that dream of about 374, almost entirely suspended and even abhorred the study of the classics for fifteen years (comp. the Preface to his Commentary on the Galatians, written a. 388, Opera, tom. vii. 486, ed. Vallarsi), but that afterwards at Bethlehem he instructed the monks in grammaticis et humanioribus (Rufinus, Apol. ii. 8), and inserted quotations from the classics in his later writings, although mostly as reminiscences of his former reading (“quasi antiqui per nebulam somnii recordamur, ” as he says in the preface above referred to), and with the obvious intent of making profane literature subservient to the Bible (comp. his Epistola xxi. ad Damasum, cap. 13). Both Jeromeand Rufinus permitted themselves to be carried by passion to exaggerated assertions at the expense of truth.} Equally accurate was his knowledge of the literature of the church. Of the Latin fathers he particularly admired Tertullian for his powerful genius and vigorous style, though he could not forgive him his Montanism; after him Cyprian, Lactantius, Hilary, and Ambrose. In the Greek classics he was less at home; yet he shows acquaintance with Hesiod, Sophocles, Herodotus, Demosthenes, Aristotle, Theophrastus, and Galen. But in the Greek fathers he was well read, especially in Origen, Eusebius, Didymus, and Gregory Nazianzen; less in Irenaeus, Athanasius, Basil, and other doctrinal writers.

The Hebrew he learned with great labor in his mature years; first from a converted but anonymous Jew, during his five years’ ascetic seclusion in the Syrian desert of Chalcis (374–379); afterwards in Bethlehem (about 385) from the Palestinian Rabbi Bar-Anina, who, through fear of the Jews, visited him by night.\footnote{Ep. 84 ad Pammach. et Ocean. c. 3 (tom. i. 524, ed. Vallarsi): “Veni rursum Jerosolymam et Bethlehem. Quo labore, quo pretio Baraninam nocturnum babui praeceptorem! Timebat enim Judaeos, et mihi alterum exhibebat Nicodemum.”} This exposed him to the foolish rumor among bigoted opponents, that he preferred Judaism to Christianity, and betrayed Christ in preference to the new “Barabbas.”\footnote{So Rufinus wrested the name, with reference to Mark xv. 7. Comp. Rufinus, Apol. or Invect. ii. 12, and the answer of Jerometo these calumnies, in the Apol. adv. libros Ruf. l. i. c. 13 (tom. ii. 469).} He afterwards, in translating the Old Testament, brought other Jewish scholars to his aid, who cost him dear. He also inspired several of his admiring female pupils, like St. Paula and her daughter Eustochium, with enthusiasm for the study of the sacred language of the old covenant, and brought them on so far that they could sing with him the Hebrew Psalms in praise of the Lord. He lamented the injurious influence of these studies on his style, since “the rattling sound of the Hebrew soiled all the elegance and beauty of Latin speech.”\footnote{In the Preface to his Commentary on the Epistle to the Galatians: “Omnem sermonis elegantiam et Latini eloquii venustatem stridor Hebraicae lectionis sordidavit.” This, however, is to be understood cum grano salis.} Yet, on the other hand, he was by the same means preserved from flying off into hollow and turgid ornamentations, from which his earlier writings, such as his letters to Heliodorus and Innocentius, are not altogether free. Though his knowledge of Hebrew was defective, it was much greater than that of Origen, Epiphanius, and Ephraem Syrus, the only other fathers besides himself who understood Hebrew at all; and it is the more noticeable, when we consider the want of grammatical and lexicographical helps and of the Masoretic punctuation.\footnote{That there were at that time as yet no vowel-points or other diacritical signs in writing Hebrew words, has been proved against Buxtorf by L. Capellus, Morinus, and Clericus, and among modem Oriental scholars, especially by Hupfeld (Studien und Kritiken, 1830, p. 549 ff.). Comp. Zöckler, l.c. p. 345 f.}

Jerome, who unfortunately was not free from vanity, prided himself not a little upon his learning, and boasted against his opponent Rufinus, that he was “a philosopher, a rhetorician, a grammarian, a dialectician, a Hebrew, a Greek, a Latin, three-tongued,” that is, master of the three principal languages of the then civilized world.\footnote{Apol adv. Ruf. lib. iii. c. 6 (tom. ii. 537). His claim to be a philosopher may be questioned. In the same place he calls “papa” Epiphanius \textgreek{πεντάγλωττος}, a man of five tongues, because besides the three chief languages he also understood the Syriac and the Egyptian or Coptic. But his knowledge of the languages was far inferior to that of Jerome. Augustineregarded Jeromeas the most learned man among all mortals.“Quod Hieronymus nescivit,” he said, “nullus mortalium unquam scivit.” Comp. also the enthusiastic praise of Erasmus, quoted § 41, p. 206, who placed him far above all the fathers; while Luther acknowledged his learning indeed, but could not bear his monastic spirit, and judged him harshly and unjustly. Comp. M. Lutheri Colloquia, ed. H. Bindseil, 1863, tom. iii. 135, 149, 193; ii. 340, 349, 357.}

All these manifold and rare gifts and attainments made him an extremely influential and useful teacher of the church; for he brought them all into the service of an earnest and energetic, though monkishly eccentric piety. They gave him superior access to the sense of the Holy Scriptures, which continued to be his daily study to extreme old age, and stood far higher in his esteem than all the classics. His writings are imbued with Bible knowledge, and strewn with Bible quotations.

But with all this he was not free from faults as glaring as his virtues are shining, which disturb our due esteem and admiration. He lacked depth of mind and character, delicate sense of truth, and firm, strong convictions. He allowed himself inconsistencies of every kind, especially in his treatment of Origen, and, through solicitude for his own reputation for orthodoxy, he was unjust to that great teacher, to whom he owed so much. He was very impulsive in temperament, and too much followed momentary, changing impressions. Many of his works were thrown off with great haste and little consideration. He was by nature an extremely vain, ambitious, and passionate man, and he never succeeded in fully overcoming these evil forces. He could not bear censure. Even his later polemic writings are full of envy, hatred, and anger. In his correspondence with Augustine, with all assurances of respect, he everywhere gives that father to feel his own superiority as a comprehensive scholar, and in one place tells him that he never had taken the trouble to read his writings, excepting his Soliloquies and “some commentaries on the Psalms.” He indulged in rhetorical exaggerations and unjust inferences, which violated the laws of truth and honesty; and he supported himself in this, with a characteristic reference to the sophist Gorgias, by the equivocal distinction between the gymnastic or polemic style and the didactic.\footnote{Between \textgreek{γυμναστικῶς}scribere and \textgreek{δογματικῶς}scribere. Ep. 48 ad Pammachium pro libris contra Jovinianum, cap. 13.} From his master Cicero he had also learned the vicious rhetorical arts of bombast, declamatory fiction, and applause-seeking effects, which are unworthy of a Christian theologian, and which invite the reproach of the divine judge in that vision: “Thou liest! thou art a Ciceronian, not a Christian; for where thy treasure is, there thy heart is also.”

\xsection{The Works of Jerome}

§ 177. The Works of Jerome.

The writings of Jerome, which fill eleven folios in the edition of Vallarsi, may be divided into exegetical, historical, polemic doctrinal, and polemic ethical works, and epistles.\footnote{The Vallarsi edition, Verona, 1734-’42, and with improvements, Venet. 1766’72, is much more complete and accurate than the Benedictine or Maurine edition of \emph{Martianay} and \emph{Pouget}, in 5 vols. 1706, although this far surpassed the older editions of \emph{Erasmus}, and \emph{Marianus Victorius}. The edition of \emph{Migne}, Paris (Petit-Montrouge), 1845-’46, also in 11 volumes (tom. xxii.-xxx. of the Patrologia Lat.), notwithstanding the boastful title, is only an uncritical reprint of the edition of Vallarsi with unessential changes in the order of arrangement; the Vitae Hieronymi and the Testimonia de Hieronymo being transferred from the eleventh to the first volume, which is more convenient. Vallarsi, a presbyter of Verona, was assisted in his work by Scipio Maffei, and others. I have mostly used his edition, especially in the Epistles.}

I. The exegetical works stand at the head.

Among these the Vulgata,\footnote{The name \emph{Vulgata}, sc. \emph{editio}, \textgreek{κοινὴ ἔκδοσις}, i. e., the received text of the Bible, was a customary designation of the Septuagint, as also of the Latin Itala (frequently so used in Jeromeand Augustine), sometimes used in the bad sense of a vulgar, corrupt text as distinct from the original. The council of Trent sanctioned the use of the term in the honorable sense for Jerome’s version of the Bible. With the same right Luther’s version might be called the German, King James’ version the English Vulgate.} or Latin version of the whole Bible, Old Testament and New, is by far the most important and valuable, and constitutes alone an immortal service.\footnote{This is now pretty generally acknowledged. We add a few of the most weighty testimonies. Luther, who bore a real aversion to Jeromeon account of his fanatical devotion to monkery, still, in view of the invaluable assistance he received from the Vulgate in his own similar work, does him the justice to say: “St. Jeromehas personally done more and greater in translation than any one man will imitate.” Zöckler, l.c. p. 183, thinks: “The Vulgate is unquestionably the most important and most meritorious achievement of our author, the ripest fruit of his laborious studies, not only in the department of Hebrew, in which he leaves all other ecclesiastical authors of antiquity far behind, but also in that of Greek and of biblical criticism and exegesis in general, in which he excels at least all, even the greatest, of the Western fathers.” O. F. Fritzsche(in Herzog’s Encykl. vol. xvii. p. 435): “The severe judgment respecting the labor of Jeromesoftened with time, and, in fact, so swung to the opposite, that he was regarded as preserved from error by the guidance of the Holy Ghost. This certainly cannot be admitted, for the defects are palpably many and various. Yet criticism must acknowledge that Jeromeperformed a truly important service for his age; that he first gave the Old Testament to the West, and in a measure also the New, in a substantially pure form; put a stop, provisionally, to the confusion of the Bible text; and as a translator gave, on the whole, the true sense. He very properly aimed to be \emph{interpres}, not \emph{paraphrastes}, but in the great dissimilarity between the Hebrew and Latin idiom, he encountered the danger of slavish literalness. This he has \emph{in general} avoided, and has been able to keep a certain mean between too great strictness and too great freedom, so that the language, though everywhere showing the Hebrew tinge, would not at all offend, but rather favor, the reader \emph{of that day}. Yet it may be said that Jeromecould have done still better. It was not that reverence, caution, restrained him; to avoid offence, he adhered as closely as possible to the current version, especially in the New Testament. He sometimes let false translations stand, when they seemed harmless (” quod non nocebat, mutare noluimus ”), and probably followed popular usage in respect to phraseology; so that the style is not perfectly uniform. Finally, he did not always give himself due time, but worked rapidly. This is particularly true in the Apocrypha, of which, however, he had a very low estimate. Some parts he left entirely untouched, others he translated or revised very hastily.” Comp. also the opinion of the English scholar, B. F. Westcott, in W. Smith’s Dictionary of the Bible, vol. iii. pp. 1696 and 1714 f., who says among other things: “When every allowance has been made for the rudeness of the original Latin, and the haste of Jerome’s revision, it can scarcely be denied that the Vulgate is not only the most venerable but also the most precious monument of Latin Christianity. For ten centuries it preserved in Western Europe a text of Holy Scripture far purer than that which was current in the Byzantine church; and at the revival of Greek learning, guided the way towards a revision of the late Greek text, in which the best biblical critics have followed the steps of Bentley, with ever-deepening conviction of the supreme importance of the coincidence of the earliest Greek and Latin authorities.”}

Above all his contemporaries, and above all his successors down to the sixteenth century, Jerome, by his linguistic knowledge, his Oriental travel, and his entire culture, was best fitted, and, in fact, the only man, to undertake and successfully execute so gigantic a task, and a task which just then, with the approaching separation of East and West, and the decay of the knowledge of the original languages of the Bible in Latin Christendom, was of the highest necessity. Here, its so often in history, we plainly discern the hand of divine Providence. Jerome began the work during his second residence in Rome (382–385), at the suggestion of pope Damasus, who deserves much more credit for that suggestion than for his hymns. He at first intended only a revision of the Itala, the old Latin version of the Bible which came down from the second century, and the text of which had fallen into inextricable confusion through the negligence of transcribers and the caprice of correctors.\footnote{Jeromesays of the Itala: “Tot sunt exemplaria paene quot codices, ” and frequently complains of the “varietas” and “vitiositas” of the Codices Latini, which he charges partly upon the original translators, partly upon presumptuous revisers, partly upon negligent transcribers. Comp. especially his Praefat. in Evang. ad Damasum.} He finished the translation at Bethlehem, in the year 405, after twenty years of toil. He translated first the Gospels, then the rest of the New Testament, next the Psalter (which he wrought over twice, in Rome and in Bethlehem\footnote{Both versions continued in use, the former as the \emph{Psalterium Romanum}, the other as the \emph{Psalterium Gallicanum}, like the two English versions of the Psalms in the worship of the Anglican church.}), and then, in irregular succession, the historical, prophetic, and poetical books, and in part the Apocrypha, which, however, he placed decidedly below the canonical books. By this “labor pius, sed periculosa praesumtio,” as he called it, he subjected himself to all kinds of enmity from ignorance and blind aversion to change, and was abused as a disturber of the peace and falsifier of the Scripture;\footnote{Falsarius, sacrilegus, et corruptor Scripturae.} but from other sources he received much encouragement. The New Testament and the Psalter were circulated and used in the church long before the completion of the whole. Augustine, for example, was using the New Testament of Jerome, and urged him strongly to translate the Old Testament, but to translate it from the Septuagint.\footnote{Augustinefeared, from the displacement of the Septuagint, which he regarded as apostolically sanctioned, and as inspired, a division between the Greek and Latin church, but yielded afterwards, in part at least, to the correct view of Jerome, and rectified in his Retractations several false translations in his former works. Westcott, in his scholarly article on the Vulgate (in Smith’s Dictionary of the Bible, iii. 702), makes the remark: “There are few more touching instances of humility than that of the young Augustinebending himself in entire submission before the contemptuous and impatient reproof of the veteran scholar.”} Gradually the whole version made its way on its own merits, without authoritative enforcement, and was used in the West, at first together with the Itala, and after about the ninth century alone.

The Vulgate takes the first place among the Bible-versions of the ancient church. It exerted the same influence upon Latin Christendom as the Septuagint upon Greek, and it is directly or indirectly the mother of most of the earlier versions in the European vernaculars.\footnote{Excepting the Gothic version, which is older than Jerome, and the Slavonic, which comes down from Methodius and Cyril.} It is made immediately from the original languages, though with the use of all accessible helps, and is as much superior to the Itala as Luther’s Bible to the older German versions. From the present stage of biblical philology and exegesis the Vulgate can be charged, indeed, with innumerable faults, inaccuracies, inconsistencies, and arbitrary, dealing, in particulars;\footnote{It has been so censured long ago by Le Clere in his Quaestiones Hieronymianae,} but notwithstanding these, it deserves, as a whole, the highest praise for the boldness with which it went back from the half-deified Septuagint directly to the original Hebrew; for its union of fidelity and freedom; and for the dignity, clearness, and gracefulness of its style. Accordingly, after the extinction of the knowledge of Greek, it very naturally became the clerical Bible of Western Christendom, and so continued to be, till the genius of the Reformation in Germany, Switzerland, Holland, and England, returning to the original text, and still further penetrating the spirit of the Scriptures, though with the continual help of the Vulgate, produced a number of popular Bibles, which were the same to the evangelical laity that the Vulgate had been for many centuries to the Catholic clergy. This high place the Vulgate holds even to this day in the Roman church, where it is unwarrantably and perniciously placed on an equality with the original.\footnote{For particulars respecting the Vulgate, see H. Hody: De Bibliorum textibus originalibus, Oxon. 1705; Joh. Clericus: Quaestiones Hieronymianae, Amsterd. 1719 (who, provoked by the exaggerated praise of the Benedictine editor, Martianay, subjected the Vulgate to a sharp and penetrating though in part unjust criticism); Leander van Ess: Pragmatisch-kritische Geschichte der Vulgata, Tüb. 1824; the lengthy article \emph{Vulgata} by O. F. Fritzschein Herzog’s Theol. Encycl. vol. xvii. pp. 422-460; an article on the same subject by B. F. Westcottin W. Smith’s Dictionary of the Bible, 1863, vol. iii. pp. 688-718; and Zöckler: Hieronymus, pp. 99 ff.; 183 ff.; 343 ff. \par The text of the Vulgate, in the course of time, has become as corrupt as the text of the Itala was at the time of Jerome, and it is as much in need of a critical revision from manuscript sources, as the textus receptus of the Greek Testament. The authorized editions of SixtusV. and ClementXIII. have not accomplished this task. Martianay, in the Benedictine edition of Jerome’s work, did more valuable service towards an approximate restoration of the Vulgate in its original form from manuscript sources. Of late the learned Barnabite C. Vercellonehas commenced such a critical revision in Variae Lectiones Vulgatae Latin. Bibliorum editionis, tom. i. (Pentat.), Rome, 1860; tom. ii. Pars prior (to 1 Regg.), 1862. Westcott, in the article referred to, has made use of the chief results of this work, which may be said to create an epoch in the history of the Vulgate.}

The Commentaries of Jerome cover Genesis, the Major and Minor Prophets, Ecclesiastes, Job, some of the Psalms,\footnote{His seven treatises on Psalms x.-xvi. (probably translated from Origen), and his brief annotations to all the Psalms (commentarioli) are lost, but the pseudo-hieronymianum breviarium in Psalmos, a poor compilation of later times (Opera, vii. 1-588), contains perhaps fragments of these.} the Gospel of Matthew, and the Epistles to the Galatians, Ephesians, Titus, and Philemon.\footnote{Opera, tom. iii. iv. v. vi. and vii. Jeromededicated his commentaries and other writings mostly to those high-born ladies of Rome whom he induced to embrace the ascetic mode of life, as Paula, Eustochium, Marcella, \&c.h He received much encouragement from them in his labors;—such was the lively theological interest which prevailed in some female circles at the time. He was, however, censured on this account, and defended himself in the Preface to his Commentary on Zephaniah, tom. vi. 671, by referring to Deborah and Huldah, Judith and Esther, Anna, Elizabeth, and Mary, not forgetting the heathen Sappho, Aspasia, Themista, and the Cornelia Gracchorum, as examples of literary women.} Besides these he translated the Homilies of Origen on Jeremiah and Ezekiel, on the Gospel of Luke, and on the Song of Solomon. Of the last he says: “While Origen in his other writings has surpassed all others, on the Song of Solomon he has surpassed himself.”\footnote{Praef. in Homil. Orig. in Cantic. Cant. tom. iii. 500. Rufinus, during the Origenistic controversy, did not forget to remind him of this sentence.}

His best exegetical labors are those on the Prophets (Particularly his Isaiah, written a.d. 408–410; his Ezekiel, a.d. 410–415; and his Jeremiah to chap. xxxii., interrupted by his death), and those on the Epistles to the Galatians, Ephesians, and Titus, (written in 388), together with his critical Questions (or investigations) on Genesis. But they are not uniformly carried out; many parts are very indifferent, others thrown off with unconscionable carelessness in reliance on his genius and his reading, or dictated to an amanuensis as they came into his head.\footnote{He frequently excuses this “dictare quodcunque in buccam venerit,” by his want of time and the weakness of his eyes. Comp. Preface to the third book of his Comment. in Ep. ad Galat. (tom. vii. 486). At the close of the brief Preface to the second book of his Commentary on the Ep. to the Ephesians (tom. vii. 486), he says that he often managed to write as many as a thousand lines in one day (“interdum per singulos dies usque ad numerum mille versuum—i.e., here \textgreek{στίχοι —}-pervenire”).} He not seldom surprises by clear, natural, and conclusive expositions, while just on the difficult passages he wavers, or confines himself to adducing Jewish traditions and the exegetical opinions of the earlier fathers, especially of Origen, Eusebius, Apollinaris, and Didymus, leaving the reader to judge and to choose. His scholarly industry, taste, and skill, however, always afford a certain compensation for the defect of method and consistency, so that his Commentaries are, after all, the most instructive we have from the Latin church of that day, not excepting even those of Augustine, which otherwise greatly surpass them in theological depth and spiritual unction. He justly observes in the Preface to his Commentary on Isaiah: “He who does not know the Scriptures, does not know the power and wisdom of God; ignorance of the Bible is ignorance of Christ.”\footnote{“Qui nescit Scripturas, nescit Dei virtutem ejusque sapientiam; ignoratio Scripturarum ignoratio Christi est.”}

Jerome had the natural talent and the acquired knowledge, to make him the father of grammatico-historical interpretation, upon which all sound study of the Scriptures must proceed. He very rightly felt that the expositor must not put his own fancies into the word of God, but draw out the meaning of that word, and he sometimes finds fault with Origen and the allegorical method for roaming in the wide fields of imagination, and giving out the writer’s own thought and fancy for the hidden Wisdom of the Scriptures and the church.\footnote{Comp. particularly the Preface to the fifth book of his Commentary on Isaiah, and Ep. 53 ad Paulinum, c. 7.} In this healthful exegetical spirit he excelled all the fathers, except Chrysostom and Theodoret. In the Latin church no others, except the heretical Pelagius (whose short exposition of the Epistles of Paul is incorporated in the works of Jerome), and the unknown Ambrosiaster (whose commentary has found its way among the works of Ambrose), thought like him. But he was far from being consistent; he committed the very fault he censures in Eusebius, who in the superscription of his Commentary on Isaiah promised a historical exposition, but, forgetting the promise, fell into the fashion of Origen. Though he often makes very bold utterances, such as that on the original identity of presbyter and bishop,\footnote{In the Comm. on Tit. i. 5, and elsewhere, e.g., Epist. 69 ad Oceanum, c. 3, and Epist. 146 ad Evangelum, c. 1. Such assertions, which we find also in Ambrosiaster, Chrysostom, and Theodoret were not disputed at that time, but subsequently they gave rise to violent disputes between Episcopalians and Presbyterians. Comp. my History of the Apostolic Church, § 132.} and even shows traces of a loose view of inspiration,\footnote{He admits, for instance, chronological contradictions, or, at least inexplicable difficulties in the Gospel history (Ep. 57 ad Pammach. c. 7 and 8), and he even ventures unjustly to censure St. Paul for supposed solecisms, barbarisms, and weak arguments (Ep. 121 ad Alag.; Comment. in Gal. iii. 1; iv. 24; vi. 2; Comment. in Eph. iii. 3, 8, 13; Comment. in Tit. i. 3).} yet he had not the courage, and was too scrupulously concerned for his orthodoxy, to break with the traditional exegesis. He could not resist the impulse to indulge, after giving the historical sense, in fantastic allegorizing, or, as he expresses himself, “to spread the sails of the spiritual understanding.”\footnote{“Spiritualis intelligentiae vela panders,” or “spirituale aedificium super historiae fundamentum extruere,” or “quasi inter saxa et scopulos” (between Scylla and Charybdis), “sic inter historiam et allegoriam omtionis cursum flectere.”}

He distinguishes in most cases a double sense of the Scriptures: the literal and the spiritual, or the historical and the allegorical; sometimes, with Origen and the Alexandrians, a triple sense: the historical, the tropological (moral), and the pneumatical (mystical).

The word of God does unquestionably carry in its letter a living and life-giving spirit; and is capable of endless application to all times and circumstances; and here lies the truth in the allegorical method of the ancient church. But the spiritual sense must be derived with tender conscientiousness and self-command from the natural, literal meaning, not brought from without, as another sense beside, or above, or against the literal.

Jerome goes sometimes as far as Origen in the unscrupulous twisting of the letter and the history, and adopts his mischievous principle of entirely rejecting the literal sense whenever it may seem ludicrous or unworthy. For instance: By the Shunamite damsel, the concubine of the aged king David, he understands (imitating Origen’s allegorical obliteration of the double crime against Uriah and Bathsheba) the ever-virgin Wisdom of God, so extolled by Solomon;\footnote{Ep. 52 ad Nepotianum, c. 2-4. He objects against the historical construction, that it is absurd, inasmuch as the aged David, then seventy years old, might as well have warmed himself in the arms of Bathsheba, Abigail, and the other wives and concubines still living, considering that Abraham at a still more advanced age was content with his Sarah, Isaac with his Rebeccah. The Shunamite, therefore, must be “sapientia quae numquam senescit” (c. 4, tom. i. 258). Nevertheless, in another place, he understands the same passage literally, Contra Jovinian. l. i. c. 24 (tom. i. 274), where he mentions this and other sins of David, “non quod sanctis viris aliquid detrahere audeam, sed quod aliud sit in lege versari, aliud in evangelio.”} and the earnest controversy between Paul and Peter he alters into a sham fight for the instruction of the Antiochian Christians who were present; thus making out of it a deceitful accommodation, over which Augustine (who took just offence at such patrocinium mendacii) drew him into an epistolary controversy characteristic of the two men.”\footnote{Comp. Jerome’s Com. on Gal. ii. 11-14; Aug. Epp. 28, 40, and 82, or Epp. 56, 67, and 116 among the Epistles of Jerome(Opera, i. 300 sqq.; 404 sqq.; 761 sqq.) After defending for a long time his false interpretation, Jeromegave it up at last, a.d.415, in his Dial. contra Pelag. l. i. c. 22. Augustine, on the other hand, yielded his erroneous preference for a translation of the Old Testament from the Septuagint instead of the original Hebrew, although he continued to entertain an exaggerated estimate of the value of the Septuagint and the very imperfect Itala. Besides these two points of dispute the Origenistic errors were a subject of correspondence between these most distinguished fathers of the Latin Church.}

It is remarkable that Augustine and Jerome, in the two exegetical questions, on which they corresponded, interchanged sides, and each took the other’s point of view. In the dispute on the occurrence in Antioch (Gal. ii. 11–14), Augustine represented the principle of evangelical freedom and love of truth, Jerome the principle of traditional committal to dogma and an equivocal theory of accommodation; while in their dispute on the authority of the Septuagint Jerome held to true progress, Augustine to retrogression and false traditionalism. And each afterwards saw his error, and at least partially gave it up.

In the exposition of the Prophets, Jerome sees too many allusions to the heretics of his time (as Luther finds everywhere allusions to the Papists, fanatics, and sectarians); and, on the other hand, with the zeal he inherited from Origen against all chiliasm, he finds far too little reference to the end of, all things in the second coming of our Lord. He limits, for example, even the eschatological discourse of Christ in the twenty-fourth chapter of Matthew, and Paul’s prophecy of the man of sin in the second Epistle to the Thessalonians, to the destruction of Jerusalem.

Among the exegetical works in the wider sense belongs the book On the Interpretation of the Hebrew Names, an etymological lexicon of the proper names of the Old and New Testaments, useful for its time, but in many respects defective, and now worthless;\footnote{Liber de interpretatione nominum Hebraicorum, or De nominibus Hebr. (Opera, tom iii. 1-120). Clericus, in his Quaestiones Hieronymianae, severely criticised this book.} and a free translation of the Onomasticon of Eusebius, a sort of biblical topology in alphabetical order, still valuable to antiquarian scholarship.\footnote{Liber de situ et nominibus locorum Hebraicorum, usually cited under the title Eusebii Onomasticon (urbium et locorum S. Scripturae). Opera, tom. iii. 121-290. Comp. Clericus: Eusebii Onomasticon cum versione Hieronymi, Amstel. 1707, and a modern convenient edition in Greek and Latin by F. Larsowand G. Parthey, Berlin, 1862.}

II. The historical works, some of which we have already elsewhere touched, are important to the history of the fathers and the saints to Christian literature, and to the history of morals.

First among them is a free Latin reproduction and continuation of the Greek Chronicle of Eusebius; i.e., chronological tables of the most important events of the history of the world and the church to the year 379.\footnote{Opera, viii. 1-820, including the Greek fragments. There is added also the Chronicon of Prosper Aquitanus (pp. 821-856), and the Apparatus, Castigationes et Notae of Arn. Pontac. We must mention also the famous separate edition of Jerome’s Chronicle and its continuators by Joseph Scaliger: Thesaurus temporum Eusebii Pamphili, Hieronymi, Prosperi, etc., Lugd. Bat. 1606, ed. altera Amstel. 1658. Scaliger and Vallarsi have spent immense industry and acuteness in editing this work made very difficult by the many chronological and other blunders and the corruptions of the text caused by ignorant and careless transcribers. The Chronicle of Eusebius is now known also in an Armenian translation, edited by Angelo Mai, Rome, 1833. The Greek original is lost with the exception of a few fragments of Syncellus.} Jerome dictated this work quite fugitively during his residence with Gregory Nazianzen in Constantinople (a.d. 380). In spite of its many errors, it formed a very useful and meritorious contribution to Latin literature, and a principal source of the scanty historical information of Western Christendom throughout the middle age. Prosper Aquitanus, a friend of Augustine and defender of the doctrines of free grace against the Semi-Pelagians in Gaul, continued the Chronicle to the year 449; later authors brought it down to the middle of the sixth century.

More original is the Catalogue of Illustrious Authors,\footnote{Liber de illustribus viris, or De scriptoribus ecclesiasticis, frequently quoted by the title Catalogus. See Opera, ed. Vallarsi, tom. ii. 821-956, together with the Greek translation of Pseudo-Sophronius.} which Jerome composed in the tenth year of Theodosius (a.d. 392 and 393),\footnote{This date is given by himself, cap. 135, in which he speaks of his own writings.} at the request of his friend, an officer, Dexter. It is the pioneer in the history of theological literature, and gives, in one hundred and thirty-five chapters, short biographical notices of as many ecclesiastical writers, from the apostles to Jerome himself, with accounts of their most important works. It was partly designed to refute the charge of ignorance, which Celsus, Porphyry, Julian, and other pagans, made against the Christians. Jerome, at that time, was not yet so violent a heretic-hater, and was quite fair and liberal in his estimate of such men as Origen and Eusebius.\footnote{In the very first chapter he says of the Second Epistle of Peter that it was by most rejected as spurious “propter styli cum priore dissonantiam.” A thorough investigation, however, leads to a more favorable result as to the genuineness of this Epistle. He admits in his catalogue even heretics, as Tatian, Bardesanes, and Priscillian, also the Jews Philo and Josephus, and the heathen philosopher Seneca.} But many of his sketches are too short and meagre; even those, for example, of so important men as Cyprian, Athanasius, Basil the Great, Gregory of Nyssa, Epiphanius, Ambrose, and Chrysostom († 407).\footnote{Of Chrysostomhe merely says, cap. 129: “Joannes Antiochenae ecclesiae presbyter, Eusebii Emiseni Diodorique sectator, multa componere dicitar, de quibus \textgreek{περὶ ἱεροσύνης}tantum legi.” But afterwards, during the Origenistic controversies, he translated a passionate libel of Theophilus of Alexandria against Chrysostom, and praised it as a valuable book (Comp. Ep. 114 ad Theophilum, written 405). Fragments of this miserable Libellus Theophili contra Joannem Chrysost. are preserved in the Defensio trium capp. l. vi. by Facundus of Hermiane.} His junior cotemporary, Augustine, who had at that time already written several philosophical, exegetical, and polemic works, he entirely omits.

The Catalogue was afterwards continued in the same spirit by the Semi-Pelagian Gennadius of Marseilles, by Isidore of Seville, by Ildefonsus, and by others, into the middle age.

Jerome wrote also biographies of celebrated hermits, Paul of Thebes (a.d. 375), Hilarion, and the imprisoned Malchus (a.d. 390), in very graceful and entertaining style, but with many fabulous and superstitious accompaniments, and with extravagant veneration of the monastic life, which he aimed by these writings to promote.\footnote{Opera, tom. ii. 1 sqq. In most of the former editions these Vitae are wrongly placed among the Epistles. To the same class of writings belongs the translation of the Regula Pachomii. Characteristic is the judgment of Gibbon (ch. xxxvii. ad Ann. 370): “The stories of Paul, Hilarion, and Malchus by Jeromeare admirably told: and the only defect of these pleasing compositions is the want of truth and common sense.”} They were read at that time as eagerly as novels. These biographies, and several necrological letters in honor of deceased friends, such as Nepotian, Lucinius, Lea, Blasilla, Paulina, Paula, and Marcella are masterpieces of rhetorical ascetic hagiography. They introduce the legend ary literature of the middle age, with its indiscriminate mixture of history and fable, and its sacrifice of historical truth to popular edification.

III. Of the polemic doctrinal and ethical works\footnote{All in the second volume of the editions of Vallarsi (p. 171 sqq.) and Migne(p. 155 sqq.).} some relate to the Arian controversies, some to the Origenistic, some to the Pelagian. In the first class belongs the Dialogue against the schismatic Luciferians,\footnote{Altercatio Luciferiani et Orthodoxi, or Dialogus contra Luciferianos. The Luciferians had their name from Lucifer, bishop of Calaris in Sardinia (died 371), the head of the strict Athanasian party, who arbitrarily ordained Paulinus bishop of Antioch in opposition to the legitimate Meletius (362), because the latter had been elected by the Arian or Semi-Arian party, although immediately after his ordination he had given in his adhesion to the Nicene faith. Lucifer afterwards fell out with the orthodox and organized a new schismatic party, which adopted Novatian principles of discipline, but in the beginning of the fifth century gradually returned to the bosom of the Catholic church.} which Jerome wrote during his desert life in Syria (a.d. 379) on the occasion of the Meletian schism in Antioch; also his translation of the work of Didymus On the Holy Ghost, begun in Rome and finished in Bethlehem. His book Against Bishop John of Jerusalem (a.d. 399), and his Apology to his former friend Rufinus, in three books (a.d. 402–403), are directed against Origenism.\footnote{Besides these Jerometranslated several letters of Epiphanius and Theophilus of Alexandria against the Origenists, which have been incorporated by Vallarsi with the collection of Jerome’s Epistles.} In the third class belongs the Dialogue against the Pelagians, in three books (a.d. 415). Other polemic works, Against Helvidius (written in 383), Against Jovinian (a.d. 393), and Against Vigilantius (dictated rapidly in one night in 406), are partly doctrinal, partly ethical in their nature, and mainly devoted to the advocacy of the immaculate virginity of Mary, celibacy, vigils, relic-worship, and the monastic life.

These controversial writings, the contents of which we have already noted in the proper place, do the author, on the whole, little credit, and stand in striking contrast with his fame as one of the principal saints of the Roman church. They show an accurate acquaintance with all the arts of an advocate and all the pugilism of a dialectician, together with boundless vehemence and fanatical zealotism, which scruple over no weapons of wit, mockery, irony, suspicion, and calumny, to annihilate opponents, and which pursue them even after their death.\footnote{Of the dead Jovinianhe says (Adv. Vigil.c. 1): “Ille Romanae ecclesicae auctoritate damnatus, inter phasides aves et carnes suillas non tam emisit spiritum, quam cructavit.” He threatened his former friend Rufinus, whose language he had perverted into a threat to take his life, with a libel suit, and after his death in 410 he wrote in an ignoble sense of triumph (in the Prologue to his Commentary on Ezekiel): “Scorpius inter Enceladum et Porphyrionem Trinacriae humo premitur, et hydra multorum capitum contra nos aliquando sibilare cessavit.” From Jerome’s polemical writings one would form a most unfavorable opinion of Rufinus. Two divines of Aquileja, Fontanini and Maria de Rubeis, felt it their duty to vindicate his memory against unjust aspersions. Comp. Zöckler, l.c. p. 266 f. Augustine, in a letter to Jerome(Ep. Hieron. 110, c. 10), called it a “magnum et triste miraculum, ” that the friendship of Jeromeand Rufinus should have turned into such enmity, and urged him to reconciliation, but in vain. This change, however, is easily explained, since hatred is only inverted love. Rufinus, it must be remembered, had not spared Jerome, and charged him even with worse than heathen impiety for calling, in hyper-ascetic zeal, Paula, the mother of the nun Eustochium, the “mother-in-law of God” (socrus Dei). See his Ep. xxii. c. 20 ad Paulam.} And their contents afford no sufficient compensation for these faults. For Jerome was not an original, profound, systematic, or consistent thinker, and therefore very little fitted for a didactic theologian. In the Arian controversy he would not enter into any discussion of the distinction between οὐσίαand ὑπόστασις, and left this important question to the decision of the Roman bishop Damasus; in the Origenistic controversy he must, in his violent condemnation of all Origenists, contradict his own former view and veneration of Origen as the greatest teacher after the Apostles; and in the Pelagian controversy he was influenced chiefly by personal considerations, and drawn half way to Augustine’s side; for while he was always convinced of the universality of sin,\footnote{Comp. particularly the passage Dial adv. Pelag. l. ii. c. 4 (tom. ii. p. 744).} in reference to the freedom of the will and predestination he adopted synergistic or Semi-Pelagian views, and afterwards continued in the highest consideration among the Semi-Pelagians down to Erasmus.\footnote{Hence it is not accidental, that several writings of Pelagius, his Commentary on the Epistles of Paul (with some emendations), his Epistola ad Demetriadem de virginitate, his Libellus fidei addressed to pope Innocent, and the Epistola ad Celantiam matronam de rations pie vivendi (which was probably likewise written by him), found their way, by an irony of history, into the writings of Jerome, on a seeming resemblance in spirit and aim.}

He is equally unsatisfactory as a moralist and practical divine. He had no connected system of moral doctrine, and did not penetrate to the basis and kernel of the Christian life, but moved in the outer circle of asceticism and casuistry. Following the spirit of his time, he found the essence of religion in monastic flight from the world and contempt of the natural ordinances of God, especially of marriage; and, completely reversing sound principles, he advocated even ascetic filth as an external mark of inward purity.\footnote{“Difficile inter epulas servatur pudicitia. Nitens cutis sordidum ostendit animum.” So he wrote to two ladies, a mother and her daughter in Gaul, Ep. 117, c. 6 (tom. i. 786). St. Anthony, the patriarch of monks, and other saints of the desert were of the same opinion, who washed themselves but seldom and combed their hair but once in a year, on holy Easter (when they ought to have been eminently holy, that is, according to their notions, eminently slovenly). What a contrast this to our modern principle that cleanliness is next to godliness! We must, however, judge this catholic ascetic cynicism from the stand-point of antiquity. Even Socrates, starting from the principle that freedom from need was divine, despised undergarments and shoes, and contented himself with a miserable cloak. Yet he did not neglect cleanliness altogether, and censured his disciple Antisthenes, who ostentatiously wore a dirty and torn cloak, by reminding him: “Friend, vanity peeps out from the holes of thy cloak.” Man is by nature lazy and dirty. Industry and cleanliness are the fruit of discipline and civilization. In this respect Europe is in advance of Asia, the Teutonic races in advance of the Latin. The Italians call the English and Americans, soap-wasters. The use of soap and of the razor is a test of modern civilization.} Of marriage he had a very low conception, regarding it merely as a necessary evil for the increase of virgins. From the expression of Paul in 1 Cor. vii. 1: “It is good not to touch a woman,” he draws the utterly unwarranted inference: “It is therefore bad to touch one; for the only opposite of good is bad;” and he interprets the woe of the Lord upon those that are with child and those that give suck (Matt. xxiv. 19), as a condemnation of pregnancy in general, and of the crying of little children, and of all the trouble and fruit of the married life. The disagreeable fact of the marriage of Peter he endeavors to weaken by the groundless assumption that the apostle forsook his wife when he forsook his net, and, besides, that “he must have washed away the stain of his married life by the blood of his martyrdom.”\footnote{Compare the work Against Jovinian, l. i. c. 7, 10, 12, 13, 15, 16, 26, 33, etc., and several of his ascetic letters. Some of his utterances on the state of matrimony gave offence even to his monastic friends.}

In a letter, otherwise very beautiful and rich, to the young Nepotian,\footnote{Ep. 52 (i. 254 sqq.) de vita clericorum et monachorum, c. 5.} he gives this advice: “Let your lodgings be rarely or never visited by women. You must either ignore alike, or love alike, all the daughters and virgins of Christ. Nay, dwell not under the same roof with them, nor trust their former chastity; you cannot be holier than David, nor wiser than Solomon. Never forget that a woman drove the inhabitants of Paradise out of their possession. In sickness any brother, or your sister, or your mother, can minister to in the lack of such relatives, the church herself maintains many aged women, whom you can at the same time remunerate for their nursing with welcome alms. I know some who are well in the body indeed, but sick in mind. It is a dangerous service in any case, that is done to you by one whose face you often see. If in your official duty as a clergyman you must visit a widow or a maiden, never enter her house alone. Take with you only those whose company does you no shame; only some reader, or acolyth, or psalm-singer, whose ornament consists not in clothes, but in good morals, who does not crimp his hair with crisping pins, but shows chastity in his whole bearing. But privately or without witnesses, never put yourself in the presence of a woman.”

Such exhortations, however, were quite in the spirit of that age, and were in part founded in Jerome’s own bitter experience in his youth, and in the thoroughly corrupt condition of social life in the sinking empire of Rome.

While advocating these ascetic extravagancies Jerome does not neglect to chastise the clergy and the monks for their faults with the scourge of cutting satire. And his writings are everywhere strewn with the pearls of beautiful moral maxims and eloquent exhortations to contempt of the world and godly conduct.\footnote{Comp. a collection of the principal doctrinal and moral sentences of Jeromein Zöcklerp. 429 ff. and p. 458 ff.}

IV. The Epistles of Jerome, with all their defects are uncommonly instructive and interesting, and, in easy flow and elegance of diction, are not inferior to the letters of Cicero. Vallarsi has for the first time put them into chronological order in the first volume of his edition, and has made the former numbering of them (even that of the Benedictine edition) obsolete. He reckons in all a hundred and fifty, including several letters from cotemporaries, such as Epiphanius, Theophilus of Alexandria, Augustine, Damasus, Pammachius, and Rufinus; some of them written directly to Jerome, and some treating of matters in which he was interested. They are addressed to friends like the Roman bishop Damasus, the senator Pammachius, the bishop Paulinus of Nola, Theophilus of Alexandria, Evangelus, Rufinus, Heliodorus, Riparius, Nepotianus, Oceanus, Avitus, Rusticus, Gaudentius, and Augustine, and some to distinguished ascetic women and maidens like Paula, Eustochium, Marcella, Furia, Fabiola, and Demetrias. They treat of almost all questions of philosophy and practical religion, which then agitated the Christian world, and they faithfully reflect the virtues and the faults and the remarkable contrasts of Jerome and of his age.

Orthodox in theology and Christology, Semi-Pelagian in anthropology, Romanizing in the doctrine of the church and tradition, anti-chiliastic in eschatology, legalistic and ascetic in ethics, a violent fighter of all heresies, a fanatical apologist of all monkish extravagancies,—Jerome was revered throughout the catholic middle age as the patron saint of Christian and ecclesiastical learning, and, next to Augustine, as maximus doctor ecclesiae; but by his enthusiastic love for the Holy Scriptures, his recourse to the original languages, his classic translation of the Bible, and his manifold exegetical merits, he also played materially into the hands of the Reformation, and as a scholar and an author still takes the first rank, and as an influential theologian the second (after Augustine), among the Latin fathers; while, as a moral character, he decidedly falls behind many others, like Hilary, Ambrose, and Leo I., and, even according to the standard of Roman asceticism, can only in a very limited sense be regarded as a saint.\footnote{Comp. the various estimates of Jeromeat § 41 above; in Vallarsi, Opera Hier., tom. xi. 282-300, and in Zöckler, l.c. pp. 465-476. In the preface to his valuable monograph (p. v) Zöckler says: ”Jeromeis chiefly the orator and the scholar among the fathers. His life is essentially neither the life of a monk, nor a priest—for monk and priest he was only by the way—nor that of a saint—for he was no saint at all, at least not in the sense of the Roman church. It is from beginning to end the life of a \emph{scholar}, a life replete with literary studies and all sorts of scholarly enterprises.” This judgment we can subscribe only with two qualifications: he was as much a monk as a scholar, and exerted an extraordinary influence on the spread of monasticism in the West; and his reputation as a saint rests precisely on the \emph{Romish} overestimate of asceticism, as distinguished from the evangelical Protestant form of piety.}

\xsection{Augustine}

§ 178. Augustine.

I. S. Aurelii Augustini Hipponensis episcopi Opera … Post Lovaniensium theologorum recensionem [which appeared at Antwerp in 1577 in 11 vols.] castigatus [referring to tomus primus, etc.] denuo ad MSS. codd. Gallicanos, etc. Opera et studio monachorum ordinis S. Benedicti e congregatione S. Mauri [Fr. Delfau, Th. Blampin, P. Coustant, and Cl. Guesnié]. Paris, 1679–1700, xi tom. in 8 fol. vols. The same edition reprinted, with additions, at Antwerp, 1700–1703, 12 parts in 9 fol.; and at Venice, 1729–’34, in xi tom. in 8 fol. (this is the edition from which I have generally quoted; it is not to be confounded with another Venice edition of 1756–’69 in xviii vols. 4to, which is full of printing errors); also at Bassano, 1807, in 18 vols.; by Gaume fratres, Paris, 1836–’39, in xi tom. in 22 parts (a very elegant edition); and lastly by J. P. Migne, Petit-Montrouge, 1841–’49, in xii tom. (Patrol. Lat. tom. xxxii.-xlvii.). Migne’s edition (which I have also used occasionally) gives, in a supplementary volume (tom. xii.), the valuable Notitia literaria de vita, scriptis et editionibus Aug. from Schönemann’s Bibliotheca historico-literaria Patrum Lat. vol. ii. Lips. 1794, the Vindiciae Augustinianae of Norisius, and the writings of Augustine first published by Fontanini and Angelo Mai. But a thoroughly reliable critical edition of Augustine is still a desideratum. On the controversies relating to the merits of the Bened. edition, see the supplementary volume of Migne, xii. p. 40 sqq., and Thuillier: Histoire de la nouvelle ed. de S. Aug. par les PP. Bénédictins, Par. 1736. The first printed edition of Augustine appeared at Basle, 1489–’95; another, a. 1509, in 11 vols. (I have a copy of this edition in black letter, but without a title page); then the edition of Erasmus published by Frobenius, Bas. 1528–’29, in 10 vols. fol.: the Editio Lovaniensis, or of the divines of Louvain, Antw. 1577, in 11 vols., and often. Several works of Augustine have been often separately edited, especially the Confessions and the City of God. Compare a full list of the editions down to 1794 in Schönemann’s Bibliotheca, vol. ii. p. 73 sqq.

II. Possidius (Calamensis episcopus, a pupil and friend of Aug.): Vita Augustini (brief, but authentic, written 432, two years after his death, in tom. x. Append. 257–280, ed. Bened., and in nearly all other editions). Benedictini Editores: Vita Augustini ex ejus potissimum scriptis concinnata, in 8 books (very elaborate and extensive), in tom. xi. 1–492, ed. Bened. (in Migne’s reprint, tom. i. pp. 66–578). The biographies of Tillemont (Mém. tom. xiii.); Ellies Dupin (Nouvelle bibliothèque des auteurs ecclésiastiques, tom. ii. and iii.); P. Bayle (Dictionnaire historique et critique, art. Augustin); Remi Ceillier (Histoire générale des auteurs sacrés et ecclés., vol. xi. and xii.); Cave (Lives of the Fathers, vol. ii.); Kloth (Der heil. Aug., Aachen, 1840, 2 vols.); Böhringer (Kirchengeschichte in Biographien, vol. i. P. iii. p. 99 ff.); Poujoulat (Histoire de S. Aug. Par. 1843 and 1852, 2 vols.; the same in German by Fr. Hurter, Schaffh. 1847, 2 vols.); Eisenbarth (Stuttg. 1853); Ph. Schaff (St. Augustine, Berlin, 1854; English ed. New York and London, 1854); C. Bindemann (Der heil. Aug., vol. i. Berl. 1844; vol. ii. 1855, incomplete). Braune: Monica und Augustin. Grimma, 1846. Comp. also the literature at § 146, p. 783.

The Philosophy of Augustine is discussed in the larger Histories of Philosophy by Brucker, Tennemann, Rixner, H. Ritter (vol. vi. pp. 153–443), Huber (Philosophie der Kirchenväter), and in the following works: Theod. Gangauf: Metaphysische Psychologie des heil. Augustinus. 1ste Abtheilung, Augsburg, 1852. T. Théry: Le génie philosophique et littéraire de saint Augustin. Par. 1861. Abbé Flottes: Études sur saint Aug., son génie, son âme, sa philosophie. Par. 1861. Nourrisson: La philosophie de saint Augustin (ouvrage couronné par l’Institut de France), deuxième ed. Par. 1866, 2 vols.

It is a venturesome and delicate undertaking to write one’s own life, even though that life be a masterpiece of nature or of the grace of God, and therefore most worthy to be described. Of all autobiographies none has so happily avoided the reef of vanity and self-praise, and none has won so much esteem and love through its honesty and humility as that of St. Augustine.

The “Confessions,” which he wrote in the forty-sixth year of his life, still burning in the ardor of his first love, are full of the fire and unction of the Holy Ghost. They are a sublime effusion, in which Augustine, like David in the fifty-first Psalm, confesses to God, in view of his own and of succeeding generations, without reserve the sins of his youth; and they are at the same time a hymn of praise to the grace of God, which led him out of darkness into light, and called him to service in the kingdom of Christ.\footnote{Augustinehimself says of his Confessions: “Confessionum mearum libri tredecim et de malis et de bonis meis Deum laudant justum et bonum, atque in eum excitant humanum intellectum et affectum.” Retract. l. ii. c. 6.} Here we see the great church teacher of all times “prostrate in the dust, conversing with God, basking in his love; his readers hovering before him only as a shadow.” He puts away from himself all honor, all greatness, all beauty, and lays them gratefully at the feet of the All-merciful. The reader feels on every hand that Christianity is no dream nor illusion, but truth and life, and he is carried along in adoration of the wonderful grace of God.

Aurelius Augustinus, born on the 13th of November, 354,\footnote{He died, according to the Chronicle of his friend and pupil Prosper Aquitanus, the 28th of August, 430 (in the third month of the siege of Hippo by the Vandals); according to his biographer Possidius he lived seventy-six years. The day of his birth Augustinestates himself, De vita beata, § 6 (tom. i. 800): “Idibus Novembris mihi natalis dies erat.”} at Tagaste, an unimportant village of the fertile province Numidia in North Africa, not far from Hippo Regius, inherited from his heathen father, Patricius,\footnote{He received baptism shortly before his death.} a passionate sensibility, from his Christian mother, Monica (one of the noblest women in the history of Christianity, of a highly intellectual and spiritual cast, of fervent piety, most tender affection, and all-conquering love), the deep yearning towards God so grandly expressed in his sentence: “Thou hast made us for Thee, and our heart is restless till it rests in Thee.”\footnote{Conf. i. I: “Fecisti nos ad Te, et inquietum est cor nostrum, donee requiescat in Te.” In all his aberrations, which we would hardly know, if it were not from his own free confession, he never sunk to anything mean, but remained, like Paul in his Jewish fanaticism, a noble intellect and an honorable character, with burning love for the true and the good.} This yearning, and his reverence for the sweet and holy name of Jesus, though crowded into the background, attended him in his studies at the schools of Madaura and Carthage, on his journeys to Rome and Milan, and on his tedious wanderings through the labyrinth of carnal pleasures, Manichaean mock-wisdom, Academic skepticism, and Platonic idealism; till at last the prayers of his mother, the sermons of Ambrose, the biography of St. Anthony, and, above all, the Epistles of Paul, as so many instruments in the hand of the Holy Ghost, wrought in the man of three and thirty years that wonderful change which made him an incalculable blessing to the whole Christian world, and brought even the sins and errors of his youth into the service of the truth.\footnote{For particulars respecting the course of Augustine’s life, see my work above cited, and other monographs. Comp. also the fine remarks of Dr. Baurin his posthumous Lectures on Doctrine-History (1866), vol. i. Part ii, p. 26 ff. He compares the development of Augustinewith the course of Christianity from the beginning to his time, and draws a parallel between Augustineand Origen.}

A son of so many prayers and tears could not be lost, and the faithful mother who travailed with him in spirit with greater pain than her body had in bringing him into the world,\footnote{Conf. ix. c. 8: “Quae me parturivit et carne ut in hanc temporalem, et corde, ut in aeternam lucem nascerer.” L. v. 9: “Non enim satis eloquor, quid erga me habebat animi, et quanto majore sollicitudine me parturiebat spiritu, quam came pepererat.”} was permitted, for the encouragement of future mothers, to receive shortly before her death an answer to her prayers and expectations, and was able to leave this world with joy without revisiting her earthly home. For Monica died on a homeward journey, in Ostia at the mouth of the Tiber, in her fifty-sixth year, in the arms of her son, after enjoying with him a glorious conversation that soared above the confines of space and time, and was a foretaste of the eternal Sabbath-rest of the saints. She regretted not to die in a foreign land, because she was not far from God, who would raise her up at the last day. “Bury my body anywhere,” was her last request, “and trouble not yourselves for it; only this one thing I ask, that you remember me at the altar of my God, wherever you may be.”\footnote{Conf. l. ix. c. 11 “Tantum illud vos rogo, ut ad Domini altare memineritis mei, ubi fueritis.” This must be explained from the already prevailing custom of offering prayers for the dead, which, however, had rather the form of thanksgiving for the mercy of God shown to them, than the later form of intercession for them. Comp. above, § 84, p. 432 ff.} Augustine, in his Confessions, has erected to Monica the noblest monument that can never perish.

If ever there was a thorough and fruitful conversion, next to that of Paul on the way to Damascus, it was that of Augustine, when, in a garden of the Villa Cassiciacum, not far from Milan, in September of the year 386, amidst the most violent struggles of mind and heart—the birth-throes of the new life—he heard that divine voice of a child: “Take, read!” and he “put on the Lord Jesus Christ” (Rom. xiii. 14). It is a touching lamentation of his: “I have loved Thee late, Thou Beauty, so old and so new; I have loved Thee late! And lo! Thou wast within, but I was without, and was seeking Thee there. And into Thy fair creation I plunged myself in my ugliness; for Thou wast with me, and I was not with Thee! Those things kept me away from Thee, which had not been, except they had been in Thee! Thou didst call, and didst cry aloud, and break through my deafness. Thou didst glimmer, Thou didst shine, and didst drive away, my blindness. Thou didst breathe, and I drew breath, and breathed in Thee. I tasted Thee, and I hunger and thirst. Thou didst touch me, and I burn for Thy peace. If I, with all that is within me, may once live in Thee, then shall pain and trouble forsake me; entirely filled with Thee, all shall be life to me.”

He received baptism from Ambrose in Milan on Easter Sunday, 387, in company with his friend and fellow-convert Alypius, and his natural son Adeodatus (given by God). It impressed the divine seal upon the inward transformation. He broke radically with the world; abandoned the brilliant and lucrative vocation of a teacher of rhetoric, which he had followed in Rome and Milan; sold his goods for the benefit of the poor: and thenceforth devoted his rare gifts exclusively to the service of Christ, and to that service he continued faithful to his latest breath. After the death of his mother, whom he revered and loved with the most tender affection, he went a second time to Rome for several months, and wrote books in defence of true Christianity against false philosophy and the Manichaean heresy. Returning to Africa, he spent three years, with his friends Alypius and Evodius, on an estate in his native Tagaste, in contemplative and literary retirement.

Then, in 391, he was chosen presbyter against his will, by the voice of the people, which, as in the similar cases of Cyprian and Ambrose, proved to be the voice of God, in the Numidian maritime city of Hippo Regius (now Bona); and in 395 he was elected bishop in the same city. For eight and thirty years, until his death, he labored in this place, and made it the intellectual centre of Western Christendom.\footnote{He is still known among the inhabitants of the place as “the great Christian (Rumi Kebir). Gibbon(ch. xxxiii. ad Ann. 430) thus describes the place which became so famous through Augustine: ” The maritime colony of \emph{Hippo}, about two hundred miles westward of Carthage, had formerly acquired the distinguishing epithet of \emph{Regius}, from the residence of the Numidian kings; and some remains of trade and populousness still adhere to the modern city, which is known in Europe by the corrupted name of Bona.” See below, Fn126.}

His outward mode of life was extremely simple, and mildly ascetic. He lived with his clergy in one house in an apostolic community of goods, and made this house a seminary of theology, out of which ten bishops and many lower clergy went forth. Females, even his sister, were excluded from his house, and could see him only in the presence of others. But he founded religious societies of women; and over one of these his sister, a saintly widow, presided.\footnote{He mentions a sister, “soror mea, sancta proposita” [monasterii], without naming her Epist. 211, n. 4 (ed. Bened.), alias Ep. 109. He also had a brother by the name of Navigius.} He once said in a sermon, that he had nowhere found better men, and he had nowhere found worse, than in monasteries. Combining, as he did, the clerical life with the monastic, he became unwittingly the founder of the Augustinian order, which gave the reformer Luther to the world. He wore the black dress of the Eastern coenobites, with a cowl and a leathern girdle. He lived almost entirely on vegetables, and seasoned the common meal with reading or free conversation, in which it was a rule that the character of an absent person should never be touched. He had this couplet engraved on the table:

He often preached five days in succession, sometimes twice a day, and set it as the object of his preaching, that all might live with him, and he with all, in Christ. Wherever he went in Africa, he was begged to preach the word of salvation.\footnote{Possidius says, in his Vita Aug.: “Caeterum episcopatu suscepto multo instantius ac ferventius majore auctoritate, non in una tantum regione, sed ubicunque rogatus verisset verbum salutis alacriter ac suaviter, pullulante atque crescente Domini ecclesia, praedicavit.”} He faithfully administered the external affairs connected with his office, though he found his chief delight in contemplation. He was specially devoted to the poor, and, like Ambrose, upon exigency, caused the church vessels to be melted down to redeem prisoners. But he refused legacies by which injustice was done to natural heirs, and commended the bishop Aurelius of Carthage for giving back unasked some property which a man had bequeathed to the church, when his wife unexpectedly bore him children.

Augustine’s labors extended far beyond his little diocese. He was the intellectual head of the North African and the entire Western church of his time. He took active interest in all theological and ecclesiastical questions. He was the champion of the orthodox doctrine against Manichaean, Donatist, and Pelagian. In him was concentrated the whole polemic power of the catholicism of the time against heresy and schism; and in him it won the victory over them.

In his last years he took a critical review of his literary productions, and gave them a thorough sifting in his Retractations. His latest controversial works against the Semi-Pelagians, written in a gentle spirit, date from the same period. He bore the duties of his office alone till his seventy-second year, when his people unanimously elected his friend Heraclius to be his assistant and successor.

The evening of his life was troubled by increasing infirmities of body and by the unspeakable wretchedness which the barbarian Vandals spread over his country in their victorious invasion, destroying cities, villages, and churches, without mercy, and even besieging the fortified city of Hippo.\footnote{Possidius, c. 28, gives a vivid picture of the ravages of the Vandals, which have become proverbial. Comp. also Gibbon, ch. xxxiii.} Yet he faithfully persevered in his work. The last ten days of his life he spent in close retirement, in prayers and tears and repeated reading of the penitential Psalms, which he had caused to be written on the wall over his bed, that he might have them always before his eyes. Thus with an act of penance he closed his life. In the midst of the terrors of the siege and the despair of his people he could not suspect what abundant seed he had sown for the future.

In the third month of the siege of Hippo, on the 28th of August, 430, in the seventy-sixth year of his age, in full possession of his faculties, and in the presence of many friends and pupils, he passed gently and happily into that eternity to which he had so long aspired. “O how wonderful,” wrote he in his Meditations,\footnote{I freely combine several passages.} “how beautiful and lovely are the dwellings of Thy house, Almighty God! I burn with longing to behold Thy beauty in Thy bridal-chamber .... O Jerusalem, holy city of God, dear bride of Christ, my heart loves thee, my soul has already long sighed for thy beauty! .... The King of kings Himself is in the midst of thee, and His children are within thy walls. There are the hymning choirs of angels, the fellowship of heavenly citizens. There is the wedding-feast of all who from this sad earthly pilgrimage have reached thy joys. There is the far-seeing choir of the prophets; there the number of the twelve apostles; there the triumphant army of innumerable martyrs and holy confessors. Full and perfect love there reigns, for God is all in all. They love and praise, they praise and love Him evermore .... Blessed, perfectly and forever blessed, shall I too be, if, when my poor body shall be dissolved, ... I may stand before my King and God, and see Him in His glory, as He Himself hath deigned to promise: ’Father, I will that they also whom Thou hast given Me be with Me where I am; that they may behold My glory which I had with Thee before the world was.’ ” This aspiration after the heavenly Jerusalem found grand expression in the hymn De gloria et gaudiis Paradisi:

which is incorporated in the Meditations of Augustine, and the idea of which originated in part with him, though it was not brought into poetical form till long afterwards by Peter Damiani.\footnote{Comp. Daniel: Thesaurus hymnol. i. p. 116 sqq., and iv. p. 203 sq., and 116, above (p. 593, note 1).}

He left no will, for in his voluntary poverty he had no earthly property to dispose of, except his library; this he bequeathed to the church, and it was fortunately preserved from the depredations of the Arian barbarians.\footnote{Possidius says, Vita, c. 31: “Testamentum nullum fecit, quia unde faceret, pauper Dei non habuit. Ecclesiae bibliothecam omnesque codices diligenter posteris custodiendos semper jubebat.”}

Soon after his death Hippo was taken and destroyed by the Vandals.\footnote{The inhabitants escaped to the sea. There appears no bishop of Hippo after Augustine. In the seventh century the old city was utterly destroyed by the Arabians, but two miles from it Bona was built out of its ruins. Comp. Tillemont, xii i. 945, and Gibbon, ch. xxxiii. Gibbon says, that Bona, “in the sixteenth century, contained about three hundred families of industrious, but turbulent manufacturers. The adjacent territory is renowned for a pure air, a fertile soil, and plenty of exquisite fruits.” Since the French conquest of Algiers, Bona was rebuilt in 1832, and is gradually assuming a French aspect. It is now one of the finest towns in Algeria, the key to the province of Constantine, has a public garden, several schools, considerable commerce, and a population of over 10,000 of French, Moors, and Jews, the great majority of whom are foreigners. The relics of St. Augustinehave been recently transferred from Pavia to Bona. See the letters of abbé Sibour to Poujoulat sur la translation de la relique de saint Augustin de Pavie à Hippone, in Poujoulat’sHistoire de saint Augustin, tom. i. p. 413 sqq.} Africa was lost to the Romans. A few decades later the whole West-Roman empire fell in ruins. The culmination of the African church was the beginning of its decline. But the work of Augustine could not perish. His ideas fell like living seed into the soil of Europe, and produced abundant fruits in nations and countries of which he had never heard.\footnote{Even in Africa Augustine’s spirit reappeared from time to time, notwithstanding the barbarian confusion, as a light in darkness, first in Vigilius, bishop of Tapsus, who, at the close of the fifth century, ably defended the orthodox doctrine of the Trinity and the person of Christ, and to whom the authorship of the so-called Athanasian Creed has sometimes been ascribed; in Fulgentius, bishop of Ruspe, one of the chief opponents of Semi-Pelagianism, and the later Arianism, who with sixty catholic bishops of Africa was banished for several years by the Arian Vandals to the island of Sardinia, and who was called the Augustineof the sixth century died 533); and in Facundus of Hermiane (died 570), and Fulgentius Ferrandusand Liberatus, two deacons of Carthage, who took a prominent part in the Three Chapter controversy.}

Augustine, the man with upturned eye, with pen in the left hand, and a burning heart in the right (as he is usually represented), is a philosophical and theological genius of the first order, towering like a pyramid above his age, and looking down commandingly upon succeeding centuries. He had a mind uncommonly fertile and deep, bold and soaring; and with it, what is better, a heart full of Christian love and humility. He stands of right by the side of the greatest philosophers of antiquity and of modern times. We meet him alike on the broad highways and the narrow footpaths, on the giddy Alpine heights and in the awful depths of speculation, wherever philosophical thinkers before him or after him have trod. As a theologian he is facile princeps, at least surpassed by no church father, scholastic, or reformer. With royal munificence he scattered ideas in passing, which have set in mighty motion other lands and later times. He combined the creative power of Tertullian with the churchly spirit of Cyprian, the speculative intellect of the Greek church with the practical tact of the Latin. He was a Christian philosopher and a philosophical theologian to the full. It was his need and his delight to wrestle again and again with the hardest problems of thought, and to comprehend to the utmost the divinely revealed matter of the faith.\footnote{Or, as he wrote to a friend about the year 410, Epist. 120, c. 1, § 2 (tom. ii. p. 347, ed. Bened. Venet.; in older ed., Ep. 122): “Ut quod credis intelligas ... non ut fidem respuas, sed ea quae fidei firmitate jam tenes, etiam rationis luce conspicias.” He continues, ibid. c. 3: “Absit namque, ut hoe in nobis Deus oderit, in quo nos reliquis animalibus excellentiores creavit. Absit, inquam, ut ideo credamus, ne rationem accipiamus vel quaeramus; cum etiam credere non possemus, nisi rationales animas haberemus.” In one of his earliest works, Contra Academ. l. iii. c. 20, § 43, he says of himself: “Ita sum affectus, ut quid sit verum non credendo solum, sed etiam intelligendo apprehendere impatienter desiderem.”} He always asserted, indeed, the primacy of faith, according to his maxim: Fides praecedit intellectum; appealing, with theologians before him, to the well-known passage of Isaiah vii. 9 (in the LXX.): “Nisi credideritis, non intelligetis.” But to him faith itself was an acting of reason, and from faith to knowledge, therefore, there was a necessary transition.\footnote{Comp. De praed. Sanct. cap. 2, § 5 (tom. x. p. 792): “Ipsum credere nihil aliud est quam cum assensione cogitare. Non enim omnia qui cogitat, credit, cum ideo cogitant, plurique ne credant; sed cogitat omnia qui credit, et credendo cogitat et cogitando credit. Fides si non cogitetur, nulls est.” Ep. 120, cap. 1, § 3 (tom. ii. 347), and Ep. 137, c. 4, § 15 (tom. ii. 408): “Intellectui fides aditum aperit, infidelitas claudit.” Augustine’s view of faith and knowledge is discussed at large by Gangauf, Metaphysische Psychologie des heil Augustinus, i. pp. 31-76, and by Nourrisson, La philosophie de saint Augustin, tom. ii. 282-290.} He constantly looked below the surface to the hidden motives of actions and to the universal laws of diverse events. The metaphysician and the Christian believer coalesced in him. His meditatio passes with the utmost ease into oratio, and his oratio into meditatio. With profundity he combined an equal clearness and sharpness of thought. He was an extremely skilful and a successful dialectician, inexhaustible in arguments and in answers to the objections of his adversaries.

He has enriched Latin literature with a greater store of beautiful, original, and pregnant proverbial sayings, than any classic author, or any other teacher of the church.\footnote{Prosper Aquitanus collected from the works of Augustinea long list of sentences (see the Appendix to the tenth vol. of the Bened. ed. p. 223 sqq.), with reference to theological purport and the Pelagian controversies. We recall some of the best, which he has omitted: \par “Novum Testamentum in Vetere latet, Vetus in Novo patet.” \par “Distingue tempora, et concordabit Scriptura.” \par “Cor nostrum inquietum est, donec requiescat in Te.” \par “Da quod jubes, et jube quod vis.” \par “Non vincit nisi veritas, victoria veritatis est caritas.” \par “Ubi amor, ibi trinitas.” \par “Fides praecedit intellectum.” \par “Deo servire vera libertas est.” \par “Nulla infelicitas frangit, quem felicitas nulls corrumpit.” \par The famous maxim of ecclesiastical harmony: “In necessariis unitas, in dubiis (or non necessariis) libertas, in omnibus (in utrisque) caritas,”—which is often ascribed to Augustine, dates in this form not from him, but from a much later period. Dr. Lücke(in a special treatise on the antiquity of the author, the original form, etc., of this sentence, Göttingen, 1850) traces the authorship to Rupert Meldenius, an irenical German theologian of the seventeenth century.}

He had a creative and decisive hand in almost every dogma of the church, completing some, and advancing others. The centre of his system is the free redeeming grace of God in Christ, operating through the actual, historical church. He is evangelical or Pauline in his doctrine of sin and grace, but catholic (that is, old-catholic, not Roman Catholic) in his doctrine of the church. The Pauline element comes forward mainly in the Pelagian controversy, the catholic-churchly in the Donatist; but each is modified by the other.

Dr. Baur incorrectly makes freedom the fundamental idea of the Augustinian system (it much better suits the Pelagian), and founds on this view an ingenious, but only half true, comparison between Augustine and Origen. “There is no church teacher of the ancient period,” says he,\footnote{L.c.p. 30 sq.} “who, in intellect and in grandeur and consistency of view, can more justly be placed by the side of Origen than Augustine; none who, with all the difference in individuality and in mode of thought, so closely resembles him. How far both towered above their times, is most clearly manifest in the very fact that they alone, of all the theologians of the first six centuries, became the creators of distinct systems, each proceeding from its definite idea, and each completely carried out; and this fact proves also how much the one system has that is analogous to the other. The one system, like the other, is founded upon the idea of freedom; in both there is a specific act, by which the entire development of human life is determined; and in both this is an act which lies far outside of the temporal consciousness of the individual; with this difference alone, that in one system the act belongs to each separate individual himself, and only falls outside of his temporal life and consciousness; in the other, it lies within the sphere of the temporal history of man, but is only the act of one individual. If in the system of Origen nothing gives greater offence than the idea of the pre-existence and fall of souls, which seems to adopt heathen ideas into the Christian faith, there is in the system of Augustine the same overleaping of individual life and consciousness, in order to explain from an act in the past the present sinful condition of man; but the pagan Platonic point of view is exchanged for one taken from the Old Testament .... What therefore essentially distinguishes the system of Augustine from that of Origen, is only this: the fall of Adam is substituted for the pre-temporal fall of souls, and what in Origen still wears a heathen garb, puts on in Augustine a purely Old Testament form.”

The learning of Augustine was not equal to his genius, nor as extensive as that of Origen and Eusebius, but still considerable for his time, and superior to that of any of the Latin fathers, with the single exception of Jerome. He had received in the schools of Madaura and Carthage a good theoretical and rhetorical preparation for the forum, which stood him in good stead also in theology. He was familiar with Latin literature, and was by no means blind to the excellencies of the classics, though he placed them far below the higher beauty of the Holy Scriptures. The Hortensius of Cicero (a lost work) inspired him during his university course with enthusiasm for philosophy and for the knowledge of truth for its own sake; the study of Platonic and Neo-Platonic works (in the Latin version of the rhetorician Victorinus) kindled in him an incredible fire;\footnote{Adv. Academicos, l. ii. c. 2, § 5: “Etiam mihi ipsi de me incredibile incendium concitarunt.” And in several passages of the Civitas Dei (viii. 3-12; xxii. 27) he speaks very favorably of Plato, and also of Aristotle, and thus broke the way for the high authority of the Aristotelian philosophy with the scholastics of the middle age.} though in both he missed the holy name of Jesus and the cardinal virtues of love and humility, and found in them only beautiful ideals without power to conform him to them. His City of God, his book on heresies, and other writings, show an extensive knowledge of ancient philosophy, poetry, and history, sacred and secular. He refers to the most distinguished persons of Greece and Rome; he often alludes to Pythagoras, Plato, Aristotle, Plotin, Porphyry, Cicero, Seneca, Horace, Virgil, to the earlier Greek and Latin fathers, to Eastern and Western heretics. But his knowledge of Greek literature was mostly derived from Latin translations. With the Greek language, as he himself frankly and modestly confesses, he had, in comparison with Jerome, but a superficial acquaintance.\footnote{It is sometimes asserted that he had no knowledge at all of the Greek. So Gibbon, for example, says (ch. xxxiii.): “The superficial learning of Augustinewas confined to the Latin language.” But this is as much a mistake as the other assertion of Gibbon, that “the orthodoxy of St. Augustinewas derived from the Manichaean school.” In his youth he had a great aversion to the glorious language of Hellas (Conf. i. 14), and read the writings of Plato in a Latin translation (vii. 9). But after his baptism during his second residence in Rome, he took it up again with greater zest, for the sake of his biblical studies. In Hippo he had, while presbyter, good opportunity to advance in it, since his bishop, Aurelius, a native Greek, understood his mother tongue much better than the Latin. In his books he occasionally makes reference to the Greek. In his work Contra Jul. i. c. 6 § 21 (tom. x. 510), he corrects the Pelagian Julianin a translation from Chrysostom, quoting the original. “Ego ipsa verba Graeca quae a Joanne dicta sunt ponam: \textgreek{διὰ τοῦτο καὶ τὰ παιδία βαπτίζομεν, καίτοι ἁμαπτήματα οὐκ ἔχοντα}, quod est Latine: Ideo et infantes baptizamus, quamvis peccata non habentes.” Julianhad freely rendered this: “cum non sint coinquinati peccato,” and had drawn the inference: “Sanctus Joannes Constantinopolitanus negat esse in parvulis originale peccatum.” Augustinehelps himself out of the pinch by arbitrarily supplying \emph{propria} to \textgreek{ἁμαρτήματα}, so that the idea of sin inherited from another is not excluded. The Greek fathers, however, did not consider hereditary corruption to be proper sin or guilt at all, but only defect, weakness, or disease. In the City of God, lib. xix. c. 23, he quotes a passage from Porphyry’s \textgreek{ἐκ λογίων φιλοσοφία}. It is probable that he read Plotin, and the Panarion of Epiphanius or the summary of it, in Greek (while the Church History of Eusebius he knew only in the translation of Rufinus). But in his exegetical and other works he very rarely consults the Septuagint or Greek Testament, and was content with the very imperfect Itala or the improved version of Jerome. The Benedictine editors overestimate his knowledge of Greek. He himself frankly confesses that he knew very little of it, De Trinit. l. iii. Prooem. (“Graecae linguae non sit nobis tantus habitus, ut talium rerum libris legendis et intelligendis uno modo reperiamur idonei”), and Contra literas Petiliani (written in 400), l. ii. c. 38 (“Et ego quidem Graecae linguae perparum assecutus sum, et prope nihil”). On the philosophical learning of Augustinemay be compared Nourrisson, l.c. ii. p. 92 ff.} Hebrew he did not understand at all. Hence, with all his extraordinary familiarity with the Latin Bible, he made many mistakes in exposition. He was rather a thinker than a scholar, and depended mainly on his own resources, which were always abundant.\footnote{166 The following are some of the most intelligent and appreciative estimates of Augustine. Erasmus(Ep. dedicat. ad Alfons. archiep. Tolet. 1529) says, with an ingenious play upon the name Aurelius Augustinus: “Quid habet orbis christianus hoc scriptore magis \emph{aureum} vel \emph{augustius}? ut ipsa vocabula nequaquam fortuito, sed numinis providentia videantur indita viro. Auro sapientiae nihil pretiosius: fulgore eloquentiae cum sapientia conjunctae nihil mirabilius .... Non arbitror alium esse doctorem, in quem opulentus ille ac benignus Spiritus dotes suas omnes largius effuderit, quam in Augustinum.” The great philosopher Leibnitz(Praefat. ad Theodic. § 34) calls him “virum sane magnum et ingenii stupendi,” and “vastissimo ingenio praeditum.” Dr. Baur, without sympathy with his views, speaks enthusiastically of the man and his genius. Among other things be says (Vorlesungen über Dogmengeschichte, i. i. p. 61): “There is scarcely another theological author so fertile and withal so able as Augustine. His scholarship was certainly not equal to his mind; yet even that is sometimes set too low, when it is asserted that he had no acquaintance at all with the Greek language; for this is incorrect, though he had attained no great proficiency in Greek.” C. Bindemann(a Lutheran divine) begins his thorough monograph (vol. i. preface) with the well-deserved eulogium: “St. Augustineis one of the greatest personages in the church. He is second in importance to none of the teachers who have wrought most in the church since the apostolic time; and it can well be said that among the church fathers the first place is due to him, and in the time of the Reformation a Luther alone, for fulness and depth of thought and grandeur of character, may stand by his side. He is the summit of the development of the mediaeval Westem church; from him descended the mysticism, no less than the scholasticism, of the middle age; he was one of the strongest pillars of the Roman Catholicism, and from his works, next to the Holy Scriptures, especially the Epistles of Paul, the leaders of the Reformation drew most of that conviction by which a new age was introduced.” Staudenmaier, a Roman Catholic theologian, counts Augustineamong those minds in which an hundred others dwell (Scotus Erigena, i. p. 274). The Roman Catholic philosophers A. Güntherand Th. Gangauf, put him on an equality with the greatest philosophers, and discern in him a providential personage endowed by the Spirit of God for the instruction of all ages. A striking characterization is that of Dr. Johannes Huber(in his instructive work: Die Philosophie der Kirchenväter, Munich, 1859, p. 312 sq.): “Augustineis a unique phenomenon in Christian history. No one of the other fathers has left so luminous traces of his existence. Though we find among them many rich and powerful minds yet we find in none the forces of personal character, mind, heart, and will, so largely developed and so harmoniously working. No one surpasses him in wealth of perceptions and dialectical sharpness of thoughts, in depth and fervor of religious sensibility, in greatness of aims and energy of action. He therefore also marks the culmination of the patriotic age, and has been elevated by the acknowledgment of succeeding times as the first and the universal church father.—His whole character reminds us in many respects of Paul, with whom he has also in common the experience of being called from manifold errors to the service of the gospel, and like whom he could boast that he had labored in it more abundantly than all the others. And as Paul among the Apostles pre-eminently determined the development of Christianity, and became, more than all others, the expression of the Christian mind, to which men ever afterwards return, as often as in the life of the church that mind becomes turbid, to draw from him, as the purest fountain, a fresh understanding of the gospel doctrine,—so has Augustineturned the Christian nations since his time for the most part into his paths, and become pre-eminently their trainer and teacher, in the study of whom they always gain a renewal and deepening of their consciousness. Not the middle age alone, but the Reformation also, was ruled by him, and whatever to this day boasts of the Christian spirit, is connected at least in part with Augustine.” Nourrisson, the latest French writer on Augustine, whose work is clothed with the authority of the Institute of France, assigns to the bishop of Hippo the fast rank among the masters of human thought, alongside of Plato and Leibnitz, Thomas Aquinas and Bossuet. “Si une critique toujours respectueuse, mais d’une inviolable sincérité, est une des formes les plus hautes de l’admiration, j’estime, au contraire, n’avoir fait qu’exalter ce grand coeur, ce psychologue consolant et ému, ce métaphysicien subtil et sublime, en un mot, cet attachant et poétique génie, dont la place reste marquée, au premier rang, parmi le maîtres de la pensée humaine, à côté de Platon et de Descartes, d’Aristote et de saint Thomas, de Leibniz et de Bossuet.” (La philosophie de saint Augustin, Par. 1866, tom. i. p. vii.) Among English and American writers, Dr. Shedd, in the Introduction to his edition of an old translation of the Confessions (1860), has furnished a truthful and forcible description of the mind and heart of St. Augustine, as portrayed in this remarkable book.}

\xsection{The Works of Augustine}

§ 179. The Works of Augustine.

The numerous writings of Augustine, the composition of which extended through four and forty years, are a mine of Christian knowledge and experience. They abound in lofty ideas, noble sentiments, devout effusions, clear statements of truth, strong arguments against error, and passages of fervid eloquence and undying beauty, but also in innumerable repetitions, fanciful opinions, and playful conjectures of his uncommonly fertile brain.\footnote{Ellies Dupin(Bibliothèque ecclésiastique, tom. iii. lrepartie, p. 818) and Nourrisson(l.c. tom. ii. p. 449) apply to Augustinethe term \emph{magnus opinator}, which Cicero used of himself. There is, however, this important difference that Augustine, along with his many opinions on speculative questions in philosophy and theology, had very positive convictions in all essential doctrines, while Cicero was a mere ecclectic in philosophy.} His style is full of life and vigor and ingenious plays on words, but deficient in purity and elegance, and by no means free from wearisome prolixity and from that vagabunda loquacitas, with which his adroit opponent, Julian of Eclanum, charged him. He would rather, as he said, be blamed by grammarians, than not understood by the people; and he bestowed little care upon his style, though he many a time rises in lofty poetic flight. He made no point of literary renown, but, impelled by love to God and to the church, he wrote from the fulness of his mind and heart. The writings before his conversion, a treatise on the Beautiful (De Pulchro et Apto), the orations and eulogies which he delivered as rhetorician at Carthage, Rome, and Milan, are lost. The professor of eloquence, the heathen philosopher, the Manichaean heretic, the sceptic and freethinker, are known to us only, from his regrets and recantations in the Confessions and other works. His literary career for as commences in his pious retreat at Cassiciacum where he prepared himself for a public profession of his faith. He appears first, in the works composed at Cassiciacum, Rome, and near Tagaste, as a Christian philosopher, after his consecration to the priesthood as a theologian. Yet even in his theological works he everywhere manifests the metaphysical and speculative bent of his mind. He never abandoned or depreciated reason, he only subordinated it to faith and made it subservient to the defence of revealed truth. Faith is the pioneer of reason, and discovers the territory which reason explores.

The following is a classified view of his most important works, the contents of the most of which we have already noticed in former sections.\footnote{Possidiuscounts in all, including sermons and letters, one thousand and thirty writings of Augustine. On these see, above all, his Retractations, where he himself reviews ninety-three of his works (embracing two hundred and thirty-two books, see ii. 67), in chronological order; in the first book those which he wrote while a layman and presbyter, in the second those which he wrote when a bishop. Also the extended chronological index in Schönemann’sBiblioth. historico-literaria Patrum Latinorum, vol. ii. (Lips. 1794), p. 340 spq. (reprinted in the supplemental volume, xii., of Migne’s ed. of the Opera, p. 24 sqq.); and other systematic and alphabetical lists in the eleventh volume of the Bened. ed. (p. 494 sqq., ed. Venet.), and in Migne, tom. xi.}

I. Autobiographical works. To these belong the Confessions and the Retractations; the former acknowledging his sins, the latter his theoretical errors. In the one he subjects his life, in the other his writings, to close criticism; and these productions therefore furnish the best standard for judging of his entire labors.\footnote{For this reason the Benedictine editors have placed the Retractations and the Confessions at the head of his works.}

The Confessions are the most profitable, at least the most edifying, product of his pen; indeed, we may no doubt say, the most edifying book in all the patristic literature. They were accordingly, the most read even during his lifetime,\footnote{He himself says of them, Retract. l. ii. c. 6: “Multis fratribus eos [Confessionum libros tredecim] multum placuisse et placere scio.” Comp. De dono perseverantiae, c. 20: “Quid autem meorum opusculorum frequentius et delectabilius innotescere potuit quam libri Confessionum mearum?” Comp. Ep. 231 Dario comiti.} and they have been the most frequently published since.\footnote{Schönemann(in the supplemental volume of Migne’s ed. of Augustine, p. 134 sqq.) cites a multitude of separate editions of the Confessions in Latin, Italian, spanish, Portuguese, French, English, and German, from a.d.1475 to 1776. Since that time several new editions have been added. There are German translations by H. Kautz(R.C., Arnsberg, 1840), G. Rapp(Prot., 2d ed., Stuttg., 1847), and others. The best English edition is that of Dr. E. B. Pusey: The Confessions of S. Angustine, Oxford (first in 1838, as the first volume in the Oxf. Library of the Fathers, together with an edition of the Latin original). It is, however, as Dr. Pusey says, only a revision of the translation of Rev. W. Watts, D. D., London, 1650, accompanied with a long preface (pp. i-xxxv) and elucidations from Augustine’s works in notes and at the end (pp. 314-346). The edition of Dr. W. G. T. Shedd, Andover, 1860, is, as he says, “a reprint of an old translation by an author unknown to the editor, which was republished in Boston in 1843.” A cursory comparison shows, that this anonymous Boston reprint agrees almost word for word with Pusey’s revision of Watts, omitting his introduction and all his notes. Dr. Shedd has, however, added an excellent original introduction, in which he clearly and vigorously characterizes the Confessions and draws a comparison between them and the Confessions of Rousseau. He calls the former (p. xxvii) not inaptly the best commentary yet written upon the seventh and eighth chapters of Romans. “That quickening of the human spirit, which puts it again into vital and sensitive relations to the holy and eternal; that illumination of the mind, whereby it is enabled to perceive with clearness the real nature of truth and righteousness; that empowering of the will, to the conflict of victory—the entire process of restoring the Divine image in the soul of man—is delineated in this book, with a vividness and reality never exceeded by the uninspired mind.” ... “It is the life of God in the soul of a strong man, rushing and rippling with the freedom of the life of nature. He who watches can almost see the growth; he who listens can hear the perpetual motion; and he who is in sympathy will be swept along.”} A more sincere and more earnest book was never written. The historical part, to the tenth book, is one of the devotional classics of all creeds, and second in popularity only to the “Imitation of Christ,” by Thomas a Kempis, and Bunyan’s “Pilgrim’s Progress.” Certainly no autobiography is superior to it in true humility, spiritual depth, and universal interest. Augustine’s experience, as a heathen sensualist, a Manichaean heretic, an anxious inquirer, a sincere penitent, and a grateful convert, is reflected in every human soul that struggles through the temptations of nature and the labyrinth of error to the knowledge of truth and the beauty of holiness, and after many sighs and tears finds rest ad peace in the arms of a merciful Saviour. Rousseau’s “Confessions,” and Goethe’s “Truth and Poetry,” though written in a radically different spirit, may be compared with Augustine’s Confessions as works of rare genius and of absorbing interest, but, by attempting to exalt human nature in its unsanctified state, they tend as much to expose its vanity and weakness, as the work of the bishop of Hippo, being written with a single eye to the glory of God, raises man from the dust of repentance to a new and imperishable life of the Spirit.\footnote{Nourrisson(l.c. tom. i. p. 19) calls the Confessions “cet ouvrage unique, souvent imité, toujours parodié, où il s’accuse se condamne et s’humilie, prière ardente, récit entraînant, métaphysique incomparable, histoire de tout un monde qui se reflète dans l’histoire d’une âme.”}

Augustine composed the Confessions about the year 400. The first ten books contain, in the form of a continuous prayer and confession before God, a general sketch of his earlier life, of his conversion, and of his return to Africa in the thirty-fourth year of his age. The salient points in these books are the engaging history of his conversion in Milan, and the story of the last days of his noble mother in Ostia, spent as it were at the very gate of heaven and in full assurance of a blessed reunion at the throne of glory. The last three books (and a part of the tenth) are devoted to speculative philosophy; they treat, partly in tacit opposition to Manichaeism, of the metaphysical questions of the possibility of knowing God, and the nature of time and space; and they give an interpretation of the Mosaic cosmogony in the style of the typical allegorical exegesis usual with the fathers, but foreign to our age; they are therefore of little value to the general reader, except as showing that even abstract metaphysical subjects may be devotionally treated.

The Retractations were produced in the evening of his life (427), when, mindful of the proverb: “In the multitude of words there wanteth not sin,”\footnote{Prov. x. 19. This verse (ex multiloquio non effugies peccatum) the Semi-Pelagian Gennadius (De viris illustr. sub Aug.) applies against Augustinein excuse for his erroneous doctrines of freedom and predestination.} and remembering that we must give account for every idle word,\footnote{Matt. xii. 36.} he judged himself, that he might not be judged.\footnote{1 Cor. xi. 31. Comp. his Prologus to the two books of Retractationes.} He revised in chronological order the numerous works he had written before and during his episcopate, and retracted or corrected whatever in them seemed to his riper knowledge false or obscure. In all essential points, nevertheless, his theological system remained the same from his conversion to this time. The Retractations give beautiful evidence of his love of truth, his conscientiousness, and his humility.\footnote{J. Morell Mackenzie (in W. Smith’s Dictionary of Greek and Roman Biography and Mythology, vol. i. p. 422) happily calls the Retractations of Augustineone of the noblest sacrifices ever laid upon the altar of truth by a majestic intellect acting in obedience to the purest conscientiousness.”}

To this same class should be added the Letters of Augustine, of which the Benedictine editors, in their second volume, give two hundred and seventy (including letters to Augustine) in chronological order from a.d. 386 to a.d. 429. These letters treat, sometimes very minutely, of all the important questions of his time, and give us an insight of his cares, his official fidelity, his large heart, and his effort to become, like Paul, all things to all men.

When the questions of friends and pupils accumulated, he answered them in special works; and in this way he produced various collections of Quaestiones and Responsiones, dogmatical, exegetical, and miscellaneous (a.d. 390, 397, \&c.).

II. Philosophical treatises, in dialogue; almost all composed in his earlier life; either during his residence on the country-seat Cassiciacum in the vicinity of Milan, where he spent half a year before his baptism in instructive and stimulating conversation in a sort of academy or Christian Platonic banquet with Monica, his son Adeodatus, his brother Navigius, his friend Alypius, and some cousins and pupils; or during his second residence in Rome; or soon after his return to Africa.\footnote{In tom. i. of the ed. Bened., immediately after the Retractationes and Confessiones, and at the close of the volume. On these philosophical writings, see Brucker: Historia critics Philosophiae, Lips. 1766, tom. iii. pp. 485-507; H. Ritter: Geschichte der Philosophie, vol. vi. p. 153 ff.; Bindemann, l. c. p. 282 sqq.; Huber, l. c. p. 242 sqq.; Gangauf, l.c. p. 25 sqq., and Nourrison, l.c. ch. i. and ii. Nourrison makes the just remark (i. p. 53): “Si la philosophie est la recherche de la verité, jamais sans doute il ne s’est rencontré une âme plus philosophe que celle de saint Augustin. Car jamais âme n’a supporté avec plus d’impatience les anxiétés du doute et n’a fait plus d’efforts pour dissiper les fantômes de l’erreur.”}

To this class belong the works: Contra Academicos libri tres (386), in which he combats the skepticism and probabilism of the New Academy,—the doctrine that man can never reach the truth, but can at best attain only probability; De vita beata (386), in which he makes true blessedness to consist in the perfect knowledge of God; De ordine,—on the relation of evil to the divine order of the world\footnote{Or on the question: “Utrum omnia bona et mala divinae providentiae ordo contineat ? ” Comp. Retract. i. 3.} (386); Soliloquia (387), communings with his own soul concerning God, the highest good, the knowledge of truth, and immortality; De immortalitate animae (387), a continuation of the Soliloquies; De quantitate animae (387), discussing sundry questions of the size, the origin, the incorporeity of the soul; De musica libri vi (387389); De magistro (389), in which, in a dialogue with his son Adeodatus, a pious and promising, but precocious youth, who died soon after his return to Africa (389), he treats on the importance and virtue of the word of God, and on Christ as the infallible Master.\footnote{Augustine, in his Confessions (I. ix. c. 6), expresses himself in this touching way about this son of his illicit love: “We took with us [on returning from the country to Milan to receive the sacrament of baptism] also the boy Adeodatus, the son of my carnal sin. Thou hadst formed him well. He was but just fifteen years old, and he was superior in mind to many grave and learned men. I acknowledge Thy gifts, O Lord, my God, who createst all, and who canst reform our deformities; for I had no part in that boy but sin. And when we brought him up in Thy nurture, Thou, only Thou, didst prompt us to it; I acknowledge Thy gifts. There is my book entitled, De Magistro; he speaks with me there. Thou knowest that all things there put into his mouth were in his mind when he was sixteen years of age. That maturity of mind was a terror to me; and who but Thou is the artificer of such wonders? Soon Thou didst take his life from the earth; and I think more quietly of him now, fearing no more for his boyhood, nor his youth, nor his whole life. We took him to ourselves as one of the same age in Thy grace, to be trained in Thy nurture; and we were baptized together; and all trouble about the past fled from us.”} To these may be added the later work, De anima et ejus origine (419). Other philosophical works on grammar, dialectics (or ars bene disputandi), rhetoric, geometry, and arithmetic, are lost.\footnote{The books on grammar, dialectics, rhetoric, and the ten Categories of Aristotle, in the Appendix to the first volume of the Bened. ed., are spurious. For the genuine works of Augustineon these subjects were written in a different form (the dialogue) and for a higher purpose, and were lost in his own day. Comp. Retract. i. c. 6. In spite of this, Prantl(Geschichte der Logik im Abendlande, pp. 665-674, cited by Huber, l.c. p. 240) has advocated the genuineness of the Principia dialecticae, and Huberinclines to agree. Gangauf, l.c. p. 5, and Nourrisson, i. p. 37, consider them spurious.}

These works exhibit as yet little that is specifically Christian and churchly; but they show a Platonism seized and consecrated by the spirit of Christianity, full of high thoughts, ideal views, and discriminating argument. They were designed to present the different stages of human thought by which he himself had reached the knowledge of the truth, and to serve others as steps to the sanctuary. They form an elementary introduction to his theology. He afterwards, in his Retractations, withdrew many things contained in them, like the Platonic view of the pre-existence of the soul, and the Platonic idea that the acquisition of knowledge is a recollection or excavation of the knowledge hidden in the mind.\footnote{\textgreek{Ἡ μάθησις οὐκ ἄλλο τι ἢ ἀνάμνησις .} On his Plato, in the Phaedo, as is well known, rests his doctrine of pre-existence. Augustinewas at first in favor of the idea, Solil. ii. 20, n. 35; afterwards he rejected it, Retract. i. 4, § 4.} The philosopher in him afterwards yielded more and more to the theologian, and his views became more positive and empirical, though in some cases narrower also and more exclusive. Yet he could never cease to philosophize, and even his later works, especially De Trinitate and De Civitate Dei, are full of profound speculations. Before his conversion he, followed a particular system of philosophy, first the Manichaean, then the Platonic; after his conversion he embraced the Christian philosophy, which is based on the divine revelation of the Scriptures, and is the handmaid of theology and religion; but at the same time he prepared the way for the catholic ecclesiastical philosophy, which rests on the authority of the church, and became complete in the scholasticism of the middle age.

In the history of philosophy he deserves a place in the highest rank, and has done greater service to the science of sciences than any other father, Clement of Alexandria and Origen not excepted. He attacked and refuted the pagan philosophy as pantheistic or dualistic at heart; he shook the superstitions of astrology and magic; he expelled from philosophy the doctrine of emanation, and the idea that God is the soul of the world; he substantially advanced psychology; he solved the question of the origin and the nature of evil more nearly than any of his predecessors, and as nearly as most of his successors; he was the first to investigate thoroughly the relation of divine omnipotence and omniscience to human freedom, and to construct a theodicy; in short, he is properly the founder of a Christian philosophy, and not only divided with Aristotle the empire of the mediaeval scholasticism, but furnished also living germs for new systems of philosophy, and will always be consulted in the speculative establishment of Christian doctrines.

III. Apologetic works against Pagans and Jews. Among these the twenty-two books, De Civitate Dei, are still well worth reading. They form the deepest and richest apologetic work of antiquity; begun in 413, after the occupation of Rome by the Gothic king Alaric, finished in 426, and often separately published. They condense his entire theory of the world and of man, and are the first attempt at a comprehensive philosophy of universal history under the dualistic view of two antagonistic currents or organized forces, a kingdom of this world which is doomed to final destruction and a kingdom of God which will last forever.\footnote{In the Bened. ed. tom. vii. Comp. Retract. ii. 43, and above, § 12. The City of God and the Confessions are the only writings of Augustinewhich Gibbonthought good to read (chap. xxxiii.). Huber(l.c. p. 315) says: ”Augustine’s philosophy of history, as he presents it in his Civitas Dei, has remained to this hour the standard philosophy of history for the church orthodoxy, the bounds of which this orthodoxy, unable to perceive in the motions of the modern spirit the fresh morning air of a higher day of history, is scarcely able to transcend.” Nourrissondevotes a special chapter to the consideration of the two cities of Augustine, the City of the World and the City of God (tom. ii. 43-88). Compare also the Introduction to Saisset’sTraduction de la Cité de Dieu, Par. 1855.}

IV. Religious-Theological works of a general nature (in part anti-Manichaean): De utilitate credendi, against the Gnostic exaltation of knowledge (392); De fide et symbolo, a discourse which, though only presbyter, he delivered on the Apostles’ Creed before the council at Hippo at the request of the bishops in 393; De doctrina Christiana iv libri (397; the fourth book added in 426), a compend of exegetical theology for instruction in the interpretation of the Scriptures according to the analogy of the faith; De catechizandis rudibus, likewise for catechetical purposes (400); Enchiridion, or De fide, spe et caritate, a brief compend of the doctrine of faith and morals, which he wrote in 421, or later, at the request of Laurentius; hence also called Manuale ad Laurentium.

V. Polemic-Theological works. These are the most copious sources of the history of doctrine. The heresies collectively are reviewed in the book De haeresibus ad Quodvultdeum, written between 428 and 430 to a friend and deacon in Carthage, and giving a survey of eighty-eight heresies, from the Simonians to the Pelagians.\footnote{This work is also incorporated in the Corpus haereseologicum of Fr. Oehler, tom. i. pp. 192-225.} In the work De vera religione (390) Augustine proposed to show that the true religion is to be found not with the heretics and schismatics, but only in the catholic church of that time.

The other controversial works are directed against the particular heresies of Manichaeism, Donatism, Arianism, Pelagianism, and Semi-Pelagianism. Augustine, with all the firmness of his convictions, was free from personal antipathy, and used the pen of controversy in the genuine Christian spirit, fortiter in re, suaviter in modo. He understood Paul’s ἀληθεύειν ἐν ἀγάπῃ, and forms in this respect a pleasing contrast to Jerome, who probably had by nature no more fiery temperament than he, but was less able to control it. “Let those,” he very beautifully says to the Manichaeans, “burn with hatred against you, who do not know how much pains it costs to find the truth, how hard it is to guard against error;—but I, who after so great and long wavering came to know the truth, must bear myself towards you with the same patience which my fellow-believers showed towards me while I was wandering in blind madness in your opinions.”\footnote{Comp. Contra Epist. Manichaei quam vocant fundamenti, l. i. 2.}

1. The anti-Manichaean works date mostly from his earlier life, and in time and matter follow immediately upon his philosophical writings.\footnote{The earliest anti-Manichaean writings (De libero arbitrio; De moribus eccl. cath. et de moribus Manich.) are in tom. i. ed. Bened.; the latter in tom. viii.} In them he afterwards found most to retract, because he advocated the freedom of the will against the Manichaean fatalism. The most important are: De moribus ecclesiae catholicae, et de moribus Manichaeorum, two books (written during his second residence in Rome, 388); De vera religione (390); Unde malum, et de libero arbitrio, usually simply De libero arbitrio, in three books, against the Manichaean doctrine of evil as a substance, and as having its seat in matter instead of free will (begun in 388, finished in 395); De Genesi contra Manichaeos, a defence of the biblical doctrine of creation (389); De duabus animabus, against the psychological dualism of the Manichaeans (392); Disputatio contra Fortunatum (a triumphant refutation of this Manichaean priest in Hippo in August, 392); Contra Epistolam Manichaei quam vocant fundamenti (397); Contra Faustum Manichaeum, in thirty-three books (400–404); De natura boni (404), \&c.

These works treat of the origin of evil; of free will; of the harmony of the Old and New Testaments, and of revelation and nature; of creation out of nothing, in opposition to dualism and hylozoism; of the supremacy of faith over knowledge; of the, authority of the Scriptures and the church; of the true and the false asceticism, and other disputed points; and they are the chief source of our knowledge of the Manichaean Gnosticism and of the arguments against it. Having himself belonged for nine years to this sect, Augustine was the better fitted for the task of refuting it, as Paul was peculiarly prepared for the confutation of the Pharisaic Judaism. His doctrine of the nature of evil is particularly valuable, He has triumphantly demonstrated for all time, that evil is not a corporeal thing, nor in any way substantial, but a product of the free will of the creature, a perversion of substance in itself good, a corruption of the nature created by God.

2. Against the Priscillianists, a sect in Spain built on Manichaean principles, are directed the book Ad Paulum Orosium contra Priscillianistas et Origenistas (411);\footnote{Tom. viii. p. 611 sqq.} the book Contra mendacium, addressed to Consentius (420); and in part the 190th Epistle (alias Ep. 157), to the bishop Optatus, on the origin of the soul (418), and two other letters, in which he refutes erroneous views on the nature of the soul, the limitation of future punishments, and the lawfulness of fraud for supposed good purposes.

3. The anti-Donatistic works, composed between the years 393 and 420, argue against separatism, and contain Augustine’s doctrine of the church and church-discipline, and of the sacraments. To these belong: Psalmus contra partem Donati (a.d. 393), a polemic popular song without regular metre, intended to offset the songs of the Donatists; Contra epistolam Parmeniani, written in 400 against the Carthaginian bishop of the Donatists, the successor of Donatus; De baptismo contra Donatistas, in favor of the validity of heretical baptism (400); Contra literas Petiliani (about 400), against the view of Cyprian and the Donatists, that the efficacy of the sacraments depends on the personal worthiness and the ecclesiastical status of the officiating priest; Ad Catholicos Epistola contra Donatistas, vulgo De unitate ecclesiae (402); Contra Cresconium grammaticum Donatistam (406); Breviculus collationis cum Donatistis, a short account of the three-days’ religious conference with the Donatists (411); De correctione Donatistarum (417); Contra Gaudentium, Donat. Episcopum, the last anti-Donatistic work (420).\footnote{All these in tom. lx. Comp. above, §§ 69 and 70.}

4. The anti-Arian works have to do with the deity of Christ and of the Holy Ghost, and with the Holy Trinity. By far the most important of these are the fifteen books De Trinitate (400–416);—the most profound and discriminating production of the ancient church on the Trinity, in no respect inferior to the kindred works of Athanasius and the two Gregories, and for centuries final to the dogma.\footnote{Tom. viii. ed. Bened. p. 749 sqq. Comp. § 131, above. The work was stolen from him by some impatient friends before revision, and before the completion of the twelfth book, so that he became much discouraged, and could only be moved to finish it by urgent entreaties.} This may also be counted among the positive didactic works, for it is not directly controversial. The Collatio cum Maximino Ariano, an obscure babbler, belongs to the year 428.

5. The numerous anti-Pelagian works of Augustine are his most influential and most valuable. They were written between the years 412 and 429. In them Augustine, in his intellectual and spiritual prime, developes his system of anthropology and soteriology, and most nearly approaches the position of evangelical Protestantism: On the Guilt and the Remission of Sins, and Infant Baptism (412); On the Spirit and the Letter (413); On Nature and Grace (415); On the Acts of Pelagius (417); On the Grace of Christ, and Original Sin (418); On Marriage and Concupiscence (419); On Grace and Free Will (426); On Discipline and Grace (427); Against Julian of Eclanum (two large works, written between 421 and 429, the Second unfinished, and hence called Opus imperfectum); On the Predestination of the Saints (428); On the Gift of Perseverance (429); \&c.\footnote{Opera, tom. x., in two parts, with an Appendix. The same in Migne. Comp.\textbf{146-160}, above.}

VI. Exegetical works. The best of these are: De Genesi ad literam (The Genesis word for word), in twelve books, an extended exposition of the first three chapters of Genesis, particularly the history of the creation literally interpreted, though with many mystical and allegorical interpretations also (written between 401 and 415);\footnote{Tom. iii. 117-324. Not to be confounded with two other books on Genesis, in which he defends the biblical doctrine of creation against the Manichaeans. In this exegetical work he aimed, as he says, Retract. ii. c. 24, to interpret Genesis “non secundum allegoricas significationes, sed secundum rerum gestarum proprietatem.” The work is more original and spirited than the Hexaëmeron of Basil or of Ambrose.} Enarrationes in Psalmos (mostly sermons);\footnote{Tom. iv., the whole volume.} the hundred and twenty-four Homilies on the Gospel of John (416 and 417);\footnote{Tom. iii., 289-824.} the ten Homilies on the First Epistle of John (417); the Exposition of the Sermon on the Mount (393); the Harmony of the Gospels (De consensu evangelistarum, 400); the Epistle to the Galatians (394); and the unfinished commentary on the Epistle to the Romans.\footnote{All in tom. iii.}

Augustine deals more in lively, profound, and edifying thoughts on the Scriptures than in proper grammatical and historical exposition, for which neither he nor his readers had the necessary linguistic knowledge, disposition, or taste. He grounded his theology less upon exegesis than upon his Christian and churchly mind, saturated with Scriptural truths.

VII. Ethical or Practical and Ascetic works. Among these belong three hundred and ninety-six Sermones (mostly very short) de Scripturis (on texts of Scripture), de tempore (festival sermons), de sanctis (in memory of apostles, martyrs, and saints), and de diversis (on various occasions), some of them dictated by Augustine, some taken down by hearers.\footnote{Tom. v., which contains besides these a multitude (317) of doubtful and spurious sermons, likewise divided into four classes. To these must be added recently discovered sermons, edited from manuscripts in Florence, Monte Cassino, etc., by M. Denis(1792), O. F. Frangipane(1820), A. L. Caillau(Paris, 1836), and Angelo Mai(in the Nova Bibliotheca Patrum).} Also various moral treatises: De continentia (395); De mendacio (395), against deception (not to be confounded with the similar work already mentioned Contra mendacium, against the fraud-theory of the Priscillianists, written in 420); De agone Christiano (396); De opere monachorum, against monastic idleness (400); De bono conjugali adv. Jovinianum (400); De virginitate (401); De fide et operibus (413); De adulterinis conjugiis, on 1 Cor. vii. 10 sqq. (419); De bono viduitatis (418); De patientia (418); De cura pro mortuis gerenda, to Paulinus of Nola (421); De utilitate jejunii; De diligendo Deo; Meditationes; etc.\footnote{Most of them in tom. vi. ed. Bened. On the scripta deperdita, dubia et spuria of Augustine, see the index by Schönemann, l.c. p. 50 sqq., and in the supplemental volume of Migne’s edition, pp. 34-40. The so-called Meditations of Augustine(German translation by August Krohne, Stuttgart, 1854) are a later compilation by the abbot of Fescamp in France, at the close of the twelfth century, from the writings of Augustine, Gregory the Great, Anselm, and others.}

As we survey, this enormous literary labor, augmented by many other treatises and letters now lost, and as we consider his episcopal labors, his many journeys, and his adjudications of controversies among the faithful, which often robbed him of whole days, we must be really astounded at the fidelity, exuberance, energy, and perseverance of this father of the church. Surely, such a life was worth the living.

\xsection{The Influence of Augustine upon Posterity and his Relation to Catholicism and Protestantism}

§ 180. The Influence of Augustine upon Posterity and his Relation to Catholicism and Protestantism.

Before we take leave of this imposing character, and of the period of church history in which he shines as the brightest star, we must add some observations respecting the influence of Augustine on the world since his time, and his position with reference to the great antagonism of Catholicism and Protestantism. All the church fathers are, indeed, the common inheritance of both parties; but no other of them has produced so permanent effects on both, and no other stands in so high regard with both, as Augustine. Upon the Greek church alone has he exercised little or no influence; for this church stopped with the undeveloped synergistic anthropology of the previous age.\footnote{It betrays a very contracted, slavish, and mechanical view of history, when Roman Catholic divines claim the fathers as their exclusive property; forgetting that they taught a great many things which are as inconsistent with the papal as with the Protestant Creed, and knew nothing of certain dogmas (such as the infallibility of the pope, the seven sacraments, transubstantiation, purgatory, indulgences, auricular confession, the immaculate conception of the Virgin Mary, etc.), which are essential to Romanism. “I recollect well,” says Dr. Newman, the former intellectual leader of Oxford Tractarianism (in his Letter to Dr. Pusey on his Eirenicon, 1866, p. 5), “what an outcast I seemed to myself, when I took down from the shelves of my library the volumes of St. Athanasius or St. Basil, and set myself to study them; and how, on the contrary, when at length I was brought into Catholic communion, I kissed them with delight, with a feeling that in them I had more than all that I had lost, and, as though I were directly addressing the glorious saints, who bequeathed them to the Church, I said to the inanimate pages, ’You are now mine, and I am yours, beyond any mistake.’ ” With the same right the Jews might lay exclusive claim to the writings of Moses and the prophets. The fathers were living men, representing the onward progress and conflicts of Christianity in their time, unfolding and defending great truths, but not unmixed with many errors and imperfections which subsequent times have corrected. Those are the true children of the fathers who, standing on the foundation of Christ and the apostles, and, kissing the New Testament rather than any human writings, follow them only as far as they followed Christ, and who carry forward their work in the onward march of true evangelical catholic Christianity.}

1. Augustine, in the first place, contributed much to the development of the doctrinal basis which Catholicism and Protestantism hold in common against such radical heresies of antiquity, as Manichaeism, Arianism, and Pelagianism. In all these great intellectual conflicts he was in general the champion of the cause of Christian truth against dangerous errors. Through his influence the canon of Holy Scripture (including, indeed, the Old Testament Apocrypha) was fixed in its present form by the councils of Hippo (393) and Carthage (397). He conquered the Manichaean dualism, hylozoism, and fatalism, and saved the biblical idea of God and of creation and the biblical doctrine of the nature of sin and its origin in the free will of man. He developed the Nicene dogma of the Trinity, completed it by the doctrine of the double procession of the Holy Ghost, and gave it the form in which it has ever since prevailed in the West, and in which it received classical expression from his school in the Athanasian Creed. In Christology, on the contrary, he added nothing, and he died shortly before the great Christological conflicts opened, which reached their ecumenical settlement at the council of Chalcedon, twenty years after his death. Yet he anticipated Leo in giving currency in the West to the important formula: “Two natures in one person.”\footnote{He was summoned to the council of Ephesus, which condemned Nestorianism in 431, but died a year before it met. He prevailed upon the Gallic monk, Leporius, to retract Nestorianism. His Christology is in many points defective and obscure. Comp. Dorner’sHistory of Christology, ii. pp. 90-98. Jeromedid still less for this department of doctrine.}

2. Augustine is also the principal theological creator of the Latin-Catholic system as distinct from the Greek Catholicism on the one hand, and from evangelical Protestantism on the other. He ruled the entire theology of the middle age, and became the father of scholasticism in virtue of his dialectic mind, and the father of mysticism in virtue of his devout heart, without being responsible for the excesses of either system. For scholasticism thought to comprehend the divine with the understanding, and lost itself at last in empty dialectics; and mysticism endeavored to grasp the divine with feeling, and easily strayed into misty sentimentalism; Augustine sought to apprehend the divine with the united power of mind and heart, of bold thought and humble faith.\footnote{Wiggers(Pragmat. Darstellung des Augustinismus und Pelagianismus, i. p. 27) finds the most peculiar and remarkable point of Augustine’s character in his singular union of intellect and imagination, scholasticism and mysticism, in which neither can be said to predominate. So also Huber, l.c. p. 313.} Anselm, Bernard of Clairvaux, Thomas Aquinas, and Bonaventura, are his nearest of kin in this respect. Even now, since the Catholic church has become a Roman church, he enjoys greater consideration in it than Ambrose, Hilary, Jerome, or Gregory the Great. All this cannot possibly be explained without an interior affinity.\footnote{Nourrisson, the able expounder of the philosophy of Augustine, says (I. c. tom. i. p. iv): ”Je ne crois pas, qu’excepté saint Paul, aucun homme ait contribué davantage, par sa parole comme par ses écrits, à organiser, à interpréter, a répandre le christianisme; et, après saint Paul, nul apparemment, non pas même le glorieux, l’invincible Athanase, n’a travaillé d’une manière aussi puissante à fonder l’unité catholique.”}

His very conversion, in which, besides the Scriptures, the personal intercourse of the hierarchical Ambrose and the life of the ascetic Anthony had great influence, was a transition not from heathenism to Christianity (for he was already a Manichaean Christian), but from heresy to the historical, episcopally organized church, as, for the time, the sole authorized vehicle of the apostolic Christianity in conflict with those sects and parties which more or less assailed the foundations of the gospel.\footnote{On the catholic and ascetic character of his conversion and his religion, see the observations in my work on Augustine, ch. viii., in the \emph{German} edition.} It was, indeed, a full and unconditional surrender of his mind and heart to God, but it was at the same time a submission of his private judgment to the authority of the church which led him to the faith of the gospel.\footnote{We recall his famous anti-Manichaean dictum: “Ego evangelio non crederem, nisi me catholicae ecclesiae commoveret auctoritas.” The Protestant would reverse this maxim, and ground his faith in the church on his faith in Christ and in the gospel. So with the well-known maxim of Irenaeus: “Ubi ecclesia, ibi Spiritus Dei, et ubi Spiritus Dei, ibi ecclesia.” According to the spirit of Protestantism it would be said conversely: “Where the Spirit of God is, there is the church, and where the church is, there is the Spirit of God.”} In the same spirit he embraced the ascetic life, without which, according to the Catholic principle, no high religion is possible. He did not indeed enter a cloister, like Luther, whose conversion in Erfurt was likewise essentially catholic, but he lived in his house in the simplicity of a monk, and made and kept the vow of voluntary poverty and celibacy.\footnote{According to genuine Christian principles it would have been far more noble, if he had married the African woman with whom he had lived in illicit intercourse for thirteen years, who was always faithful to him, as he was to her, and had borne him his beloved and highly gifted Adeodatus; instead of casting her off, and, as he for a while intended, choosing another for the partner of his life, whose excellences were more numerous. The superiority of the evangelical Protestant morality over the Catholic asceticism is here palpable. But with the prevailing spirit of his age he would hardly have enjoyed so great regard, nor accomplished so much good, if he had been married. Celibacy was the bridge from the heathen degradation of marriage to the evangelical Christian exaltation and sanctification of the family life.}

He adopted Cyprian’s doctrine of the church, and completed it in the conflict with Donatism by transferring the predicates of unity, holiness, universality, exclusiveness, and maternity, directly to the actual church of the time, which, with a firm episcopal organization, an unbroken succession, and the Apostles’ Creed, triumphantly withstood the eighty or the hundred opposing sects in the heretical catalogue of the day, and had its visible centre in Rome. In this church he had found rescue from the shipwreck of his life, the home of true, Christianity, firm ground for his thinking, satisfaction for his heart, and a commensurate field for the wide range of his powers. The predicate of infallibility alone he does not plainly bring forward; he assumes a progressive correction of earlier councils by later; and in the Pelagian controversy he asserts the same independence towards pope Zosimus, which Cyprian before him had shown towards pope Stephen in the controversy on heretical baptism, with the advantage of having the right on his side, so that Zosimus found himself compelled to yield to the African church.\footnote{On Augustine’s doctrine of the church, see § 71, above, and especially the thorough account by R. Rothe: Anfänge der christl. Kirche und ihrer Verfassung, vol. i. (1837), pp. 679-711. ”Augustine,” says he, “decidedly adopted Cyprian’s conception [of the church] in all essential points. And once adopting it, he penetrated it in its whole depth with his wonderfully powerful and exuberant soul, and, by means of his own clear, logical mind, gave it the perfect and rigorous system which perhaps it still lacked” (p. 679 f).” Augustine’s conception of the doctrine of the church was about standard for succeeding times” (p. 685).}

He was the first to give a clear and fixed definition of the sacrament, as a visible sign of invisible grace, resting on divine appointment; but he knows nothing of the number seven; this was a much later enactment. In the doctrine of baptism he is entirely Catholic,\footnote{Respecting Augustine’s doctrine of baptism, see the thorough discussion in W. Wall’sHistory of Infant Baptism, vol. i. p. 173 ff. (Oxford ed. of 1862). His view of the slight condemnation of all unbaptized children contains the germ of the scholastic fancy of the \emph{limbus infantum} and the \emph{poena damni}, as distinct from the lower regions of hell and the \emph{poena sensus}.} though in logical contradiction with his dogma of predestination; but in the doctrine of the holy communion he stands, like his predecessors, Tertullian and Cyprian, nearer to the Calvinistic theory of a spiritual presence and fruition of Christ’s body and blood. He also contributed to promote, at least in his later writings, the Catholic faith of miracles,\footnote{In his former writings he expressed a truly philosophical view concerning miracles (De vera relig. c. 25, § 47; c. 50, § 98; De utilit. credendi, c. 16, § 34; De peccat. meritis et remiss. l. ii. c. 32, § 52, and De civit Dei, xxii. c. 8); but in his Retract. l. i. c. 14, § 5, he corrects or modifies a former remark in his book De utilit. credendi, stating that he did not mean to deny the continuance of miracles altogether, but only such great miracles as occurred at the time of Christ (“quia non tanta nec omnia, non quia nulla fiunt”). See above, §§ 87 and 88, and the instructive monograph of the younger Nitzsch(Lic. and Privatdocent in Berlin): Augustinus’ Lehre vom Wunder, Berlin, 1865 (97 pp.).} and the worship of Mary;\footnote{See above, §§ 81 and 82.} though he exempts the Virgin only from actual sin, not from original, and, with all his reverence for her, never calls her mother of God.\footnote{Comp. Tract. in Evang. Joannis, viii, c. 9, where he says: “Cur ergo ait matri filius: \emph{Quid} \emph{mihi et tibi est, mulier}? \emph{nondum venit hora mea} (John ii. 4). Dominus noster Jesus Christus et Deus erat et homo: \emph{secundum quod Deus erat, matrem non habebat}; secundum quod homo erat, habebat. Mater ergo [Maria] erat carnis, mater humanitatis, mater infirmitatis quam suscepit propter nos.” This strict separation of the Godhead from the manhood of Jesus in his birth from the Virgin would have exposed Augustinein the East to the suspicion of Nestorianism. But he died a year before the council of Ephesus, at which Nestorius was condemned.}

At first an advocate of religious liberty and of purely spiritual methods of opposing error, he afterwards asserted the fatal principle of the coge intrare, and lent the great weight of his authority to the system of civil persecution, at the bloody fruits of which in the middle age he himself would have shuddered; for he was always at heart a man of love and gentleness, and personally acted on the glorious principle: “Nothing conquers but truth, and the victory of truth is love.”\footnote{See above, § 27, p. 144 f. He changed his view partly from his experience that the Donatists, in his own diocese, were converted to the catholic unity “timore legum imperialium,” and were afterwards perfectly good Catholics. He adduces also a misinterpretation of Luke xiv. 23, and Prov. ix. 9: “Da sapienti occasionem et sapientior erit.” Ep. 93, ad Vincentium Rogatistam, § 17 (tom. ii. p. 237 sq. ed. Bened.). But he expressly discouraged the infliction of death on heretics, and adjured the proconsul Donatus, Ep. 100, by Jesus Christ, not to repay the Donatists in kind. “Corrigi eos cupimus, non necari.”}

Thus even truly great and good men have unintentionally, through mistaken zeal, become the authors of much mischief.

3. But, on the other hand, Augustine is, of all the fathers, nearest to evangelical Protestantism, and may be called, in respect of his doctrine of sin and grace, the first forerunner of the Reformation. The Lutheran and Reformed churches have ever conceded to him, without scruple, the cognomen of Saint, and claimed him as one of the most enlightened witnesses of the truth and most striking examples of the marvellous power of divine grace in the transformation of a sinner. It is worthy of mark, that his Pauline doctrines, which are most nearly akin to Protestantism, are the later and more mature parts of his system, and that just these found great acceptance with the laity. The Pelagian controversy, in which he developed his anthropology, marks the culmination of his theological and ecclesiastical career, and his latest writings were directed against the Pelagian Julian and the Semi-Pelagians in Gaul, who were brought to his notice by the two friendly laymen, Prosper and Hilary. These anti-Pelagian works have wrought mightily, it is most true, upon the Catholic church, and have held in check the Pelagianizing tendencies of the hierarchical and monastic system, but they have never passed into its blood and marrow. They waited for a favorable future, and nourished in silence an opposition to the prevailing system.

Even in the middle age the better sects, which attempted to simplify, purify, and spiritualize the reigning Christianity by return to the Holy Scriptures, and the reformers before the Reformation such as Wiclif, Russ, Wessel, resorted most, after the apostle Paul, to the bishop of Hippo as the representative of the doctrine of free grace.

The Reformers were led by his writings into a deeper understanding of Paul, and so prepared for their great vocation. No church teacher did so much to mould Luther and Calvin; none furnished them so powerful weapons against the dominant Pelagianism and formalism; none is so often quoted by them with esteem and love.\footnote{Lutherpronounced upon the church fathers (with whom, however, excepting Augustine, he was but slightly acquainted) very condemnatory judgments, even upon Basil, Chrysostom, and Jerome(for Jeromehe had a downright antipathy, on account of his advocacy of fasts, virginity, and monkery); he was at times dissatisfied even with Augustine, because he after all did not find in him his \emph{sola fide}, his \emph{articulus stantis et} \emph{cadentis ecclesiae}, and says of him: ”Augustineoften erred; he cannot be trusted. Though he was good and holy, yet he, as well as other fathers, was wanting in the true faith.” But this cursory utterance is overborne by numerous commendations; and all such judgments of Luther must be taken \emph{cum grano salis}. He calls Augustinethe most pious, grave, and sincere of the fathers, the patron of divines, who taught a pure doctrine and submitted it in Christian humility to the Holy Scriptures, etc., and he thinks, if he had lived in the sixteenth century, he would have been a Protestant (si hoc seculo viveret, nobiscum sentiret), while Jeromewould have gone with Rome. Compare his singular but striking judgments on the fathers in Lutheri colloquia, ed. H. E. Bindseil, 1863, tom. iii. 149, and many other places. Gangauf, a Roman Catholic (a pupil of the philosopher Günther), concedes (l.c. p. 28, note 13) that Luther and Calvin built their doctrinal system mainly on Augustine, but, as he correctly thinks, with only partial right. Nourrisson, likewise a Roman Catholic, derives Protestantism from a corrupted (!) Augustinianism, and very superficially makes Lutheranism and Calvinism essentially to consist in the denial of the freedom of the will, which was only one of the questions of the Reformation. ”On ne saurait le méconnaître, de l’Augustinianisme corrompu, mais enfin de l’Augustinianisme procède le Protestantisme. Car, sans parler de Wiclef et de Huss, qui, nourris de saint Augustin, soutiennent, avec le réalisme platonicien, la doctrine de la prédestination; Luther et Calvin ne font guére autre chose, dans leurs Principaux ouvages, que cultiver des semences d’Augustinianisme“ (l.c. ii. p. 176). But the Reformation is far more, of course, than a repristination of an old controversy; it is a new creation, and marks the epoch of modern Christianity which is different both from the mediaeval and from ancient or patristic Christianity.}

All the Reformers in the outset, Melancthon and Zwingle among them, adopted his denial of free will and his doctrine of predestination, and sometimes even went beyond him into the abyss of supralapsarianism, to cut out the last roots of human merit and boasting. In this point Augustine holds the same relation to the Catholic church, as Luther to the Lutheran; that is, he is a heretic of unimpeachable authority, who is more admired than censured even in his extravagances; yet his doctrine of predestination was indirectly condemned by the pope in Jansenism, as Luther’s view was rejected as Calvinism by the Form of Concord.\footnote{It is well known that Luther, as late as 1526, in his work, De servo arbitrio, against Erasmus, which he never retracted, proceeded upon the most rigorous notion of the divine omnipotence, wholly denied the freedom of the will, declared it a mere lie (merum mendacium), pronounced the calls of the Scriptures to repentance a divine irony, based eternal salvation and eternal perdition upon the secret will of God, and almost exceeded Calvin. See particulars in the books on doctrine-history; the inaugural dissertation of Jul. Müller: Lutheri de praedestinatione et libero arbitrio doctrina, Gott. 1832; and a historical treatise on predestination by Carl Beckin the Studien und Kritiken for 1847. We add, as a curiosity, the opinion of Gibbon(ch. xxxiii.), who, however, had a very limited and superficial knowledge of Augustine: “The rigid system of Christianity which he framed or restored, has been entertained, with public applause, and secret reluctance, by the Latin church. The church of Rome has canonized Augustine, and reprobated Calvin. Yet as the \emph{real} difference between them is invisible even to a theological microscope, the Molinists are oppressed by the authority of the saint, and the Jansenists are disgraced by their resemblance to the heretic. In the mean while the Protestant Arminians stand aloof, and deride the mutual perplexity of the disputants. Perhaps a reasoner, still more independent, may smile in \emph{his} tum when he peruses an Arminian commentary on the Epistle to the Romans.” Nourrisson(ii. 179), from his Romish standpoint, likewise makes Lutheranism to consist “essentiellement dans la question du libre arbitre.” But the principle of Lutheranism, and of Protestantism generally, is the supremacy of the Holy Scriptures as a rule of faith, and justification by free grace through faith in Christ.} For Jansenism was nothing but a revival of Augustinianism in the bosom of the Roman Catholic church.\footnote{On the mighty influence of Augustinein the seventeenth century in France, especially on the noble Jansenists, see the works on Jansenism, and also Nourrisson, l.c. tom. ii. pp. 186-276.}

The excess of Augustine and the Reformers in this direction is due to the earnestness and energy of their sense of sin and grace. The Pelagian looseness could never beget a reformer. It was only the unshaken conviction of man’s own inability, of unconditional dependence on God, and of the almighty power of his grace to give us strength for every good work, which could do this. He who would give others the conviction that he has a divine vocation for the church and for mankind, must himself be penetrated with the faith of an eternal, unalterable decree of God, and must cling to it in the darkest hours.

In great men, and only in great men, great opposites and apparently antagonistic truths live together. Small minds cannot hold them. The catholic, churchly, sacramental, and sacerdotal system stands in conflict with the evangelical Protestant Christianity of subjective, personal experience. The doctrine of universal baptismal regeneration, in particular, which presupposes a universal call (at least within the church), can on principles of logic hardly be united with the doctrine of an absolute predestination, which limits the decree of redemption to a portion of the baptized. Augustine supposes, on the one hand, that every baptized person, through the inward operation of the Holy Ghost, which accompanies the outward act of the sacrament, receives the forgiveness of sins, and is translated from the state of nature into the state of grace, and thus, qua baptizatus, is also a child of God and an heir of eternal life; and yet, on the other hand, he makes all these benefits dependent on the absolute will of God, who saves only a certain number out of the “mass of perdition,” and preserves these to the end. Regeneration and election, with him, do not, as with Calvin, coincide. The former may exist without the latter, but the latter cannot exist without the former. Augustine assumes that many are actually born into the kingdom of grace only to perish again; Calvin holds that in the case of the non-elect baptism is an unmeaning ceremony; the one putting the delusion in the inward effect, the other in the outward form. The sacramental, churchly system throws the main stress upon the baptismal regeneration to the injury of the eternal election; the Calvinistic and Puritan system sacrifices the virtue of the sacrament to the election; the Lutheran and Anglican system seeks a middle ground, without being able to give a satisfactory theological solution of the problem. The Anglican church allows the two opposite views, and sanctions the one in the baptismal service of the Book of Common Prayer, the other in her Thirty-nine Articles, which are moderately Calvinistic.

It was an evident ordering of God, that the Augustinian system, like the Latin Bible of Jerome, appeared just in that transitional period of history, in which the old civilization was passing away before the flood of barbarism, and a new order of things, under the guidance of the, Christian religion, was in preparation. The church, with her strong, imposing organization and her firm system of doctrine, must save Christianity amidst the chaotic turmoil of the great migration, and must become a training-school for the barbarian nations of the middle age.\footnote{Guizot, the Protestant historian and statesman, very correctly says in his Histoire générale de la civilization en Europe (Deuxième leçon, p. 45 sq. ed. Bruxelles, 1850): “S’il n’eût pas été une église, je ne sais ce qui en serait avenu au milieu de la chute de l’empire romain .... Si le christianisme n’eût été comme dans les premiers temps, qu’une croyance, un sentiment, une conviction individuelle, on peut croire qu’il aurait succombé au milieu de la dissolution de l’empire et de l’invasion des barbares. Il a succombé plus tard, en Asie et dans tous le nord de l’Afrique, sous une invasion de méme nature, sous l’invasion des barbares musulmans; il a succombé alors, quoiqu’il fût à l’état d’institution, d’église constituée. A bien plus forte raison le même fait aurait pu arriver au moment de la chute de l’empire romain. Il n’y avait alors aucun des moyens par lesquels aujourd’hui les influences morales s’établissent ou’ résistent indépendamment des institutions, aucun des moyens par lesquels une pure vérité, une pure idée acquiert un grand empire sur les esprits, gouverne les actions, détermine des événemens. Rien de semblable n’existait au IVesiècle, pour donner aux idées, aux sentiments personels, une pareille autorité. Il est clair qu’il fallait une société fortement organisée, fortement gouvernée, pour lutter contre un pareil désastre, pour sortir victorieuse d’un tel ouragan. Je no crois pas trop dire en affirmant qu’à la fin du IVeet au commencement du Vesiècle, c’est l’église chrétienne qui a sauvé le christianisme; c’est l’église avec ses institutions, ses magistrats, son pouvoir, qui s’est défendue vigoureusement contre la dissolution intérieure de l’empire, contre la barbarie, qui a conquis les barbares, qui est devenue le lien, le moyen, le principe de civilisation entre le monde romain et le monde barbare.“}

In this process of training, next to the Holy Scriptures, the scholarship of Jerome and the theology and fertile ideas of Augustine were the most important intellectual agent.

Augustine was held in so universal esteem that he could exert influence in all directions, and even in his excesses gave no offence. He was sufficiently catholic for the principle of church authority, and yet at the same time so free and evangelical that he modified its hierarchical and sacramental character, reacted against its tendencies to outward, mechanical ritualism, and kept alive a deep consciousness of sin and grace, and a spirit of fervent and truly Christian piety, until that spirit grew strong enough to break the shell of hierarchical tutelage, and enter a new stage of its development. No other father could have acted more beneficently on the Catholicism of the middle age, and more successfully provided for the evangelical Reformation than St. Augustine, the worthy successor of Paul, and the precursor of Luther and Calvin.

Had he lived at the time of the Reformation, he would in all probability have taken the lead of the evangelical movement against the prevailing Pelagianism of the Roman church. For we must not forget that, notwithstanding their strong affinity, there is an important difference between Catholicism and Romanism or popery. They sustain a similar relation to each other as the Judaism of the Old Testament dispensation, which looked to, and prepared the way for, Christianity, and the Judaism after the crucifixion and after the destruction of Jerusalem, which is antagonistic to Christianity. Catholicism covers the entire ancient and mediaeval history of the church, and includes the Pauline, Augustinian, or evangelical tendencies which increased with the corruptions of the papacy and the growing sense of the necessity of a “reformatio in capite et membris.” Romanism proper dates from the council of Trent, which gave it symbolical expression and anathematized the doctrines of the Reformation. Catholicism is the strength of Romanism, Romanism is the weakness of Catholicism. Catholicism produced Jansenism, popery condemned it. Popery never forgets and never learns anything, and can allow no change in doctrine (except by way of addition), without sacrificing its fundamental principle of infallibility, and thus committing suicide. But Catholicism may ultimately burst the chains of popery which have so long kept it confined, and may assume new life and vigor.

Such a personage as Augustine, still holding a mediating place between the two great divisions of Christendom, revered alike by both, and of equal influence with both, is furthermore a welcome pledge of the elevating prospect of a future reconciliation of Catholicism and Protestantism in a higher unity, conserving all the truths, losing all the errors, forgiving all the sins, forgetting all the enmities of both. After all, the contradiction between authority and freedom, the objective and the subjective, the churchly and the personal, the organic and the individual, the sacramental and the experimental in religion, is not absolute, but relative and temporary, and arises not so much from the nature of things, as from the deficiencies of man’s knowledge and piety in this world. These elements admit of an ultimate harmony in the perfect state of the church, corresponding to the union of the divine and human natures, which transcends the limits of finite thought and logical comprehension, and is yet completely realized in the person of Christ. They are in fact united in the theological system of St. Paul, who had the highest view of the church, as the mystical “body of Christ,” and “the pillar and ground of the truth,” and who was at the same time the great champion of evangelical freedom, individual responsibility, and personal union of the believer with his Saviour. We believe in and hope for one holy catholic apostolic church, one communion of saints, one fold, and one Shepherd. The more the different churches become truly Christian, or draw nearer to Christ, and the more they give real effect to His kingdom, the nearer will they come to one another. For Christ is the common head and vital centre of all believers, and the divine harmony of all discordant human sects and creeds. In Christ, says Pascal, one of the greatest and noblest disciples of Augustine, In Christ all contradictions are solved.

\xchapter{Addenda Et Corrigenda}

ADDENDA ET CORRIGENDA.

[In the additions to the literature I have followed the method of italicizing book-titles and words in foreign languages, as in the revised edition of vols. i. and ii. The same method will be carried out in all subsequent volumes.]

Page 11. Add To Literature On Constantine The Great:

Th. Zahn: Constantin der Grosse und die Kirche. Hannover, 1876. Demetriades: Die christl. Regierung und Orthodoxie Kaiser Constantin’s des Gr. Muenchen, 1878. Th. Brieger: Constantin der Gr. als Religionspolitiker. Gotha, 1880. E. L. Cutts: Constantine the Great. Lond. And N. Y., 1881. W. Gass: Konstantin der Gr. und seine Soehne, in Herzog,2 vii. (1881), 199–207. John Wordsworth: Const. the Gr. and his Sons, in Smith and Wace, i. 623–654. Edm. Stapfer: in Lichtenberger, iii. 388–393.—Comp. also vol. ii. p. 64–74, especially on the Edicts Of Toleration (only two, not three, as formerly assumed). Victor Schultze Geschichte des Untergangs des Griechisch-roemischen Heidenthums. Jena, 1887, Vol. i. 28–68.

Page 40. Add To Lit, on the heathen sources:

Juliani imperatoris Librorum contra Christianos quae supersunt. Collegit, recensuit, prolegomenis instruxit Car. Joa. Neumann. Insunt Cyrilli Alexandrini fragmenta syriaca ab Eberh. Nestle edita. Lips., 1880. Kaiser Julian’s Buecher gegen die Christen. Nach ihrer Wiederherstellung uebersetzt von Karl Joh. Neumann. Leipzig, 1880. 53 pages. This is Fasc. iii. of Scriptorum Graecorum qui Christianam impugnaverunt religionem quae supersunt, ed. by Neumann.

Page 40, bottom of the page. Add to works on Julian the Apostate:

Alb. De Broglie (R.C.), in the third and fourth vols. of his L’église et l’empire romain au quatrième siécle. Par., 4th ed., 1868. (Very full.) J. F. A. Muecke: Flavius Claudius Julianus. Nach den Quellen. Gotha, 1867 and 1869. 2 vols. (Full, painstaking, prolix, too much dependent on Ammianus, and partial to Julian.) Kellerbaum: Skizze der Vorgeschichte Julians, 1877. F. Rode: Gesch. der Reaction des Kaiser Julianus gegen die christl. Kirche. Jens, 1877. (Careful, partly against Teuffel and Muecke.) H. Adrien Naville: Julien l’apostate et sa philosophie du polythéisme. . Paris and Neuchatel, 1877. Comp. his art. in Lichtenbergers “Encyclop.,” vii. 519–525. Torquati: Studii storico-critici sulla vita … di Giuliano l’Apostata. Rom., 1878. G. H. Rendall: The Emperor Julian: Paganism and Christianity. Lond., 1879. J. G. E. Hoffmann: Jul. der Abtruennige, Syrische Erzaehlungen. Leiden, 1880. (Old romances reflecting the feelings of the Eastern Christians.) Comp. also art. on Jul. in the “Encycl. Brit.,” 9th ed., vol. xiii. 768–770 (by Kirkup); in Herzog2, vii. 285–296 (by Harnack); in Smith and Wace, iii. 484–524 (by Prebendary John Wordsworth, very full and fair).

Page 60. Add to literature:

Tillemont: Hist. des empereurs, tom. v. A. De Broglie, l.c. Victor Schultze: Gesch. d. Untergangs des gr. roem. Heidenthums, i. 209–400.

Page 81, last line, after Muenter, 1826, add:

; by C. Bursian, Lips., 1856; C. Halm, Vienna, 1867).

Page 93. Add as footnote 3:

3 Jerome, who was a shrewd observer of men and things, and witnessed the first effects of the union of church and state, says: “Ecclesia postquam ad Christianos principes venit, potentia quidem et divitiis major, sed virtutibus minor facta.’

Page 148. Add at the bottom of the page:

H. Weingarten: Der Ursprung des Moenchthums im nachconstantinischen Zeitalter. Gotha, 1877. See also his art. in Herzog2, x. 758 sqq. Ad. Harnack: Das Moenchthum, seine Ideale und seine Geschichte. Giessen, 1882.—Comp. vol. ii. ch. ix. p. 387 sqq.

Page 226. Add to footnote:

Ad. Franz, Marcus Aur. Cassiodorus Senator. Breslau, 1872.

Page 242, i 50, add:

See Lit. on clerical celibacy in vol. i. p. 403 sq., especially Theiner, Lea, and von Schulte.

Page 314. Add to Lit. on Leo the Great:

Friedrich (old Cath.): Zur aeltesten Geschichte des Primates in der Kirche. Bonn, 1879. Jos. Langen (old Cath.): Geschichte der roem. Kirche bis zum Pontificate Leo’s I. Bonn, 1881. Karl Mueller in Herzog2, viii. (1881), 551–563. C. Gore, in Smith and Wace, iii. (1882), 652–673. By the same: Leo the Great (Lond. Soc. for Promoting Christ Knowledge, 175 pages). On the literary merits of Leo, see Ebert: Geschichte der christl. Lat. Lit., vol. i. 447–449.

Page 329. .Add to § 64 the following:

LIST OF POPES AND EMPERORS

From Constantine the Great to Gregory the Great, a.d. 314–590.

Comp. the lists in vol. ii. 166 sqq., and vol. iv. 205 Sqq.

This list is based upon Jaffé’s Regesta, Potthast’s Biblioth. Hist. Medii Aevi, and Cardinal Hergenröther’s list, in his Kirchengesch., third ed. (1886), vol. iii. 1057 sqq.

Date

Pope

Emperor

Date

311–314

Melchiades

Constantine I, or The Great

306 (323)–337

314–335

Silvester I

336–337

Marcus

Constantine II (in Gaul)

337–340

337–352

Julius I

Constantius II (In the East)

337–350

Constans (In Italy)

352–66

Liberius

357

Filix II, Antipope

Constantius Alone

350–361

Julian

361–363

Jovian

363–364

366–843

Damasus

Valentinian I

364–375

Valens

364–378

366–367

Ursicinus, Antipope

Gratian

375–383

Valentinian II (in the West)

375–392

385–398

Siricius

Theodosius

379–395

398–402

Anastasius

Arcadius (in the East)

395–408

402–417

Innocent I

Honorius (in the West)

395–423

417–418

Zosimus

Theodosius II (E.)

408–450

418–422

Bonifacius

(418 Dec. 27)

(Eulalius, Antipope)

422–432

Coelestinus I

Valentinian III (W.)

423–455

432–440

Sixtus III

440–461

Leo I the Great

Marcian (E.)

450–457

Maximus Avitus (W.)

455–457

Majorian (W.)

457–461

Leo I. (E.)

457–474

461–468

Hilarus

Severus (W.)

461–465

Vacancy (W.)

465–467

468–483

Simplicius

Anthemius (W.)

467–472

Olybrius (W.)

472–473

Glycerius (W.)

473–474

Julius Nepos (W.)

474

Page 330, line 8 from below, read after Hefele (R.C.):

Conciliengeschichte, Freiburg i. B. 1855 sqq.; second revised ed. 1873 sqq., 7 vols., down to the Council of Florence (1447).

Page 353. Add to footnote:

The reign of Pope Pius IX. has added another Council to the Latin list of oecumenical Councils, that of the Vatican, 1870, which is counted as the twentieth (by Bishop Hefele, in the revised edition of his Conciliengesch., i. 60), and which decreed the infallibility of the Pope in all his official utterances, thereby superseding the necessity of future oecumencal Councils. It has given rise to the Old Catholic secession, headed by eminent scholars such as Döllinger, Reinkens, Reusch, Langen. See the author’s Creeds of Christendom, vol. i. 134 sqq.

Page 518. Add to Lit.

C. A. Hammond: Antient Liturgies (with introduction, notes, and liturgical glossary). Oxford, 1878. Ch. A. Swainson: Greek Liturgies, chiefly from Original Sources. Cambridge, 1884.

Page 541. § 103. Church Architecture:

On the history of Architecture in general, see the works of Kugler: Geschichte der Baukunst (1859, 3 vols.); Schnaase: Gesch. der Kunst (1843–66, 8 vols.); Lübke History of Art (Eng. transl. New York, 1877, 2 vols.); Viollet Le Duc: Lectures an Architecture (London, 1877), and his numerous works in French, including Dictionnaire De l’architecture Française (Paris, 1853–69, 10 vols.); James Fergusson: History of Architecture of all Countries from the earliest Times to the present (Lond., 1865; 2d ed., 1874, 4 vols.). On church architecture in particular: Richard Brown: Sacred Architecture; its Rise, Progress, and Present State (Lond., 1845); Kreuser: Der christl. Kirchenbau (Bonn, 1851); Hübsch: Altchristl. Kirchen (Karlsruhe, 1858–61); De Vogüé: Architecture civile et relig. du Ie au VIIe siècle (Paris, 1877, 2 vols.); Ch. E. Norton: Studies of Church Buildings in the Middle Ages (Now York, 1880). There are also special works on the basilicas in Rome, Constantinople, and Ravenna. See §§ 106 and 107.

Page 560. § 109. Crosses and Crucifixes.

Comp. the Lit. in vol. ii. §§ 75 and 77.

Page 563. Add to Lit.

Mrs. Jameson and Lady Eastlake: The History of Our Lord as exemplified in Works of Art (with illustrations). London, 1864; second ed. 1865. 2 vols. Also the works on Christian Art, and on the Catacombs quoted in vol. ii. §§ 75 and 82.

Page 622. Add to Lit., line 3 from below:

Eugene Revillout: Le Concile de Nicée d’après les textes coptes et les diverses collections canoniques. Paris, 1881. The works on Arianism and on Athanasius include accounts of the Council of Nicaea. On the Nicene Creed and its literature, see Schaff: Creeds of Christendom, vol. i. 12 sqq. and 24 sqq.; and the article of Ad. Harnack, in Herzog,2 vol. viii. (1881) 212–230, abridged in Schaff-Herzog (1886), ii. 1648 sqq.

Page 651. Add to Lit., line 13:

Theod. Zahn: Marcellus Von Ancyra. Gotha, 1867. (Zahn represents Marcellus as essentially orthodox and agreed with Irenaeus, but as seeking to gain a more simple and satisfactory conception of the truth from the Bible than the theology of the age presented. Neander, Dogmengesch., i. 275, had suggested a similar view.) W. Möller: Art. Marcellus in Herzog2 vol. ix. (1881), 279–282. (Partly in opposition to Zahn.) E. S. Ffoulkes, in Smith and Wace, iii. 808–813. (Ignores the works of Zahn and other German writers.)

Page 689. § 132. The Athanasian Creed. Add to Lit.:

A. P. Stanley: The Athanasian Creed. Lond., 1871. E. S. Ffoulkes: The Athanasian Creed. Lond., 1872. Ch. A. Heurtley: The Athanasian Creed. Oxf., 1872. (Against Ffoulkes.) J. R. Lumby: History of the Creeds. Cambridge, 1873; second ed. 1880. The Utrecht Psalter, a facsimile ed., published in London, 1875. This contains the oldest MS. of the Athan. Creed, which by Ussher and Waterland was assigned to the sixth century, but by recent scholars to the ninth century. C. A. Swainson: The Nicene and Apostles’ Creeds, together with an Account o f the Growth and Reception of the Creed of St. Athanasius. Lond., 1875. (Comp. his art. Creed in Smith and Wace, i. 711.) G. D. W. Ommaney: Early History of the Athan. Creed. An Examination of Recent Theories. Lond., 1875; 2d ed. 1880. Schaff: Creeds of Christendom, i. 34 sqq. and ii. 66–72, 555 sq. (With a facsimile of the oldest MS. from the Utrecht Psalter.)

Page 696.

The statements concerning the origin and age of the Athanasian Creed should be conformed to the authors views as expressed in his work on Creeds, i. 36. The latest investigations do not warrant us to trace it higher than the eighth or seventh century. The first commentary on it ascribed to Venantius Fortunatus, 570, is of doubtful genuineness, and denied to him by Gieseler, Ffoulkes, and others. The majority of recent Anglican writers, including Stanley, Swainson, and Lumby, assign the Creed to an unknown author in Gaul between a.d. 750 and 850, probably during the reign of Charlemagne (d. 814). Hardy and Ommaney plead for an earlier date. The question is not yet fully settled. The Creed consists of two parts, one on the Trinity and one on the Incarnation, which were afterward welded together by a third hand. The second part was found separately as a fragment of a sermon on the Incarnation, at Treves, in a MS. from the middle of the eighth century, and was first published by Prof. Swainson, 1871, and again in 1875.

Page 872. Add to Lit. on Eusebius:

Fr. Ad. Heinichen: Eusebii Pamphili Scripta Historica. New ed. Lips., 1868–70. 3 Tom. The third vol. (804 pages) contains Commentarii et Meletemata. The ample indexes and critical and explanatory notes make this the most useful edition of the Church History and other historical works of Eusebius. Dindorf’s ed., Lips., 1867 sqq., 4 vols., includes the two apologetic works. Best ed. of the Chronicle by Alfred Schöne: Eusebii Chronicorum libri II. Berol. 1866 and 1875. 2 Tom., 4°. Schöne was assisted by Petermann in the Armenian Version, and by Rödiger in the Syriac Epitome. He gives also the χρονογραφεῖον σύντομονof the year 853, the first part of which professes to be derived from the labors of Eusebius. Stein: Eusebius nach s. Leben, s. Schriften, und s. dogmatischen Charakter. Würzburg, 1859. Bishop Lightfoot: art. Eusebius of Caes. in Smith and Wace, vol. ii. (full and fair). Semisch: art. Eus. v. Caes. in Herzog,2 vol. iv. 390–398. A new translation of Eusebius, with commentary, by A. C. McGiffert, will appear, N. York, 1890.

Page 885. Add to Lit. on Athanasius:

G. R. Sievers: Athanasii Vita acephala (written before 412, first publ. by Maffei, 1738). Ein Beitrag zur Gesch. des Athan. In the “Zeitschr. für Hist. Theol.” (ed. by Kahnis). Gotha, 1868, pp. 89–162. Böhringer: Athanasius und Arius, in his Kirchengesch. in Biogr. Bd. vi., new ed. Leipz., 1874. Hergenröther (R.C.): Der heil. Athanas. der Gr. Cologne, 1877 (an essay, pages 24). L. Atzberger: Die Logoslehre des heil. Athanas. München, 1880. W. Möller: Art. Athan. in Herzog,2 i. 740–747. Lüdtke: in Wetzer and Welte, 2 i. (1882), 1534–1543. Gwatkin: Studies in Arianism. Cambr. 1882.

Page 890. Add to footnote at the bottom:

Villemain considers Athanasius the greatest man between the Apostles and Gregory VII., and says of him: “Sa vie, ses combats, son génie servirent plus à l’agrandissement du christianisme que toute la puissance de Constantin .... Athanase cherche le triomphe, et non le martyre. Tel qu’un chef de parti, tel qu’un général experimenté qui se sent nécessaire aux siens, Athan. ne s’expose que pour le succès, ne combat que pour vaincre, se retire quelque fois pour reparaître avec l’éclat d’un triomphe populaire.” (Tableau de l’éloquence chrétienne au IVe siècle, p. 92.)

Page 894 line 11. Add to Lit. on St. Basil:

Dörgens: Der heil. Basilius und die class. Studien. Leipz., 1857. Eug. Fialon. Étude historique et literaire sur S. Basile, suivie de l’hexaemeron. Paris, 1861. G. B. Sievers: Leben des Libanios. Berl., 1868 (p294 sqq.). Böhringer: Die drei Kappadozier oder die trinitarischen Epigonen (Basil, Gregory of Nyssa, and Gregory of Naz.), in Kirchengesch. in Biograph., new ed. Bd. vii. and viii. 1875. Weiss: Die drei grossen Kappadozier als Exegeten. Braunsberg, 1872. R. Travers Smith: St. Basil the Great. London, 1879. (Soc. for Promoting Christian Knowledge), 232 pages. Scholl: Des heil. Basil Lehre von der Gnade. Freib., 1881. W. Möller, in Herzog,2 ii. 116–121. E. Venables, in Smith and Wace, i. 282–297. Farrar: “Lives of the Fathers,” 1889. vol. ii. 1–55.

Page 904 line 7. Add to Lit. on Gregory of Nyssa:

Böhringer: Kirchengesch. in Biogr., new ed., vol. viii. 1876. G Herrmann: Greg. Nyss. Sententiae de salute adipiscenda. Halle, 1875. . T. Bergades: De universo et de anima hominis doctrina Gregor. Nyss. Leipz., 1876. W. Möller, in Herzog,2 v. 396–404. E. Venables, in Smith and Wace, ii. 761–768. A. Paumier, in Lichtenberger, 723–725. On his doctrine of the Trinity and the Person of Christ, see especially Baur and Dorner. On his doctrine of the apokatastasis and relation to Origen, see Möller, G. Herrmann, and Bergades. l.c. Farrar: “Lives of the Fathers,” (1889), ii. 56–83.

Page 909, line 4. Add to Lit. on Gregory of Nazianzus:

A. Grenier: La vie et les poésies de saint Grégoire de Nazianze. Paris, 1858. Böhringer: K. G. in Biogr., new ed., vol. viii. 1876. Abbé A. Benoît: Vie de saint Grégoire de Nazianze. Paris, 1877. J. R. Newman: Church of the Fathers, pp. 116–145, 551. Dabas: La femme au quatrième siècle dans les poésies de Grég. de Naz. Bordeaux, 1868. H. W. Watkins, in Smith and Wace, ii. 741–761. W. Gass, in Herzog,2 v. 392–396. A. Paumier, in Lichtenberger, v., 716–722. On his christology, see Neander, Baur and especially Dorner. His views on future punishment have been discussed by Farrar, and Pusey (see vol. ii. 612). Farrar:: “Lives of the Fathers,” i. 491–582.

Page 920, line 22. Add:

In one of his plaintive songs from his religious retreat, after lamenting the factions of the church, the loss of youth, health, strength, parents, and friends, and his gloomy and homeless condition, Gregory thus gives touching expression to his faith in Christ as the last and only comforter:

“Thy will be done, O Lord! That day shall spring, When at thy word, this clay shall reappear.

No death I dread, but that which sin will bring; No fire or flood without thy wrath I fear;

For Thou, O Christ, my King, art fatherland to me. My wealth, and might, and rest; my all I find in Thee.” 1

1 Πρὸς ἑαυτον, in Daniel’s Thesaurus Hymnol., iii., 11:

Χριστὲ ἄναξ, σὺ δέ μοι πάτρη, σθένος, ὄλβος, ἅπαντα,

Σοὶ δ  ἄρ  ἀναψύξαιμι βίον καὶ κήδε  ἀμείψας.

Page 924. After line 2, add to Lit. on Cyril of Jerusalem:

J. H. Newman: Preface to the Oxford transl. of Cyril in the “Library of the Fathers”(1839). E. Venables, in Smith and Wace, i. 760–763. C. Burk, in Herzog,2 iii. 416–418.

Page 933, line 4 from below. Add to Lit. on Chrysostom:

Villemain: L’éloquence chrétienne dans le quatrième siècle. Paris 1849; new ed. 1857. P. Albert: St. Jean Chrysostôme considéré comme orateur populaire. Paris, 1858. Abbé Rochet: Histoire de S. Jean Chrysostôme. Paris, 1866. 2 vols. Th. Förster: Chrysostomus in seinem Verhältniss zur antiochenischen Schule. Gotha, 1869. W. Maggilvray: John of the Golden Mouth. Lond., 1871. Am. Thierry: S. J. Chrysostôme et l’ imperatrice Eudoxie. 2d ed. Paris, 1874. Böhringer: Johann Chrysostomus und Olympias, in his K. G. in Biogr., vol. ix., new ed., 1876. W. R. W. Stephens: St. Chrysostom: his Life and Times. London, 1872; 3d ed., 1883. F. W. Farrar, in “Lives of the Fathers,” Lond., 1889, ii. 460–540.

Engl. translation of works of St. Chrys., edited by Schaff, N. York, 1889, 6 vols. (with biographical sketch and literature by Schaff).

Page 942, line 14. Add to Lit. on Cyril of Alex.:

A new ed. of Cyril’s works, including his Com. on the Minor Prophets, the Gospel of John, the Five Books against Nestorius, the Scholia on the Incarnation, etc., was prepared with great pains by Philip Pusey (son of Dr. Pusey). Oxf., 1868–81. In 5 vols Engl. trans. in the Oxford “Library of the Fathers.” 1874 sqq. See an interesting sketch of Ph. Pusey (d. 1880) and his ed. in the “Church Quarterly Review” (London), Jan., 1883, pp. 257–291.

Page 942, line 24. Add:

Hefele: Conciliengesch., vol. ii., revised ed. (1875), where Cyril figures very prominently, pp. 135, 157, 167 sqq., 247 sqq., 266 sqq., etc. C. Burk, in Herzog,2 iii. 418 sq. W. Bright: St. Cyrillus of Al., in Smith and Wace, i. 763–773.

Page 950. Add to Lit. on Ephraem:

Evangelii Concordantis Expositio facta a S. Ephraemo Doctore Syro. Venet., 1876. (A Commentary on Tatian’s Diatessaron, found in the Mechitarist Convent at Venice in an Armenian translation, translated into Latin, 1841, by Aucher, and published with an introduction by Prof. Mösinger of Salzburg.) Comp. also the art. Ephraem, in Herzog,2 iv. 255–261 (by Radiger, revised by Spiegel). In Smith and Wace, ii. 137–145 (by E. Venables).

Page 955. Add to Lit. on Lactantius:

English translation by W. Fletcher, in Clark’s “Ante-Nicene Library,” vols. xxi. and xxii. Edinb., 1871. For an estimate of his literary merits, see Ebert: Gesch. der christl. lat. Lit. Leipz., 1874 sqq., vol. i. 70–86. Ebert, in Herzog,2 viii. 364–366. Ffoulkes, in Smith and Wace, iii. 613–617.

Page 959, line 9. Add to Lit. on Hilary of Poitiers:

Reinkens: Hilarius von Poitiers. Schaffhausen, 1864. Semisch, in Herzog,2 vi. 416–427. Cazenove, in Smith and Wace, ii. 54–66, and his St. Hilary of Poitiers. Lond., 1883. (Soc. for Promot. Christian Knowledge.) Farrar: in “Lives of the Fathers” (1889), i. 426–467.

Page 961. Add to Lit on Ambrose,

Bannard: Histoire de S. Ambroise. Paris, 1871. Ebert: Gesch. der christl. lat. Lit., i. 135–176 (1874). Robinson Thornton: St. Ambrose: his Life, Times, and Teaching. Lond., 1879, 215 pages (Soc. for Promoting Christ. Knowledge). Plitt, in Herzog,2 i. 331–335. J. Ll. Davies, in Smith and Wace, i. 91–99. Cunitz, in Lichtenberger, i. 229–232. Farrar: “Lives of the Fathers ”(1889), ii. 84–149. On the hymns of Ambrose, Comp. especially Ebert, l. c.

Page 967. Add to Lit. on Jerome:

Amédée Thierry: St. Jérôme, la société chrétienne à Rome et l’emigration romaine en terre sainte. Par., 1867. 2 vols. (He says at the close: “There is no continuation of Jerome’s work; a few more letters of Augustine and Paulinus, and night falls on the West.”) Lübeck: Hieronymus quos noverit scriptores et ex quibus hauserit. Leipzig, 1872. Ebert: Gesch. der christl. lat. Lit. Leipz., 1874 sqq., i. 176–203 (especially on the Latinity of Jerome, in which he places him first among the fathers). Edward L. Cutts: St. Jerome. London, 1877 (Soc. for Promot. Chr. Knowledge), 230 pages. Zöckler, in Herzog,2 vi. 103–108. Cunitz, in Lichtenberger, vii. 243–250. Freemantle, in Smith and Wace, iii. 29–50. (”Jerome lived and reigned for a thousand years. His writings contain the whole spirit of the church of the middle ages, its monasticism, its contrast of sacred things with profane, its credulity and superstition, its subjection to hierarchical authority, its dread of heresy, its passion for pilgrimages. To the society which was thus in a great measure formed by him, his Bible was the greatest boon which could have been given. But he founded no school and had no inspiring power; there was no courage or width of view in his spiritual legacy which could break through the fatal circle of bondage to received authority which was closing round manki

On Jerome as a Bible translator, comp. F. Kaulen (R.C.): Geschichte der Vulgata. Mainz, 1869. Hermann Rönsch: Itala und Vulgata. Das Sprachidiom der urchristlichen Itala und der katholischen Vulgata. 2d ed., revised. Marburg, 1875. L. Ziegler: Die latein Bibelübersetzungen vor Hieronymus und die Itala des Augustinus. München, 1879. (He maintains the existence of several Latin versions or revisions before Jerome.) Westcott’s art. “Vulgate,” in Smith’s Dict. of the Bible. O. F. Fritzsche: Latein. Bibelübersetzungen, in the new ed. of Herzog, vol. viii. (1881), pp. 433–472. Westcott and Hort’s Greek Testament, vol. ii., lntrod., pp. 78–84.

Page 989, line 13. Add to Lit. on Augustine:

English translations of select works of Aug. by Dr. Pusey and others in the Oxford Library of the Fathers” : the Confessions, vol. i., 1839, 4th ed., 1853; Sermons, vol. xvi., 1844, and vol. xx., 1845; Short Treatises, vol. xxii., 1847; Expositions on the Psalms, vols. xxiv., xxv., xxx., xxxii., xxxvii., xxxix., 1847, 1849, 1850, 1853, 1854; Homilies on John, vols. xxvi. and xxix., 1848 and 1849. Another translation by Marcus Dods and others, Edinb. (T. and T. Clark), 1871–76, 15 vols., containing the City of God, the Anti-Donatist, the Anti-Pelagian, the Anti-Manichaean writings, Letters, On the Trinity, the Sermon on the Mount, and the Harmony of the Gospels, On Christian Doctrine, the Euchiridion, on Catechising, on Faith and the Creed, Lectures on John, and Confessions. The same revised with new translations and Prolegomena, edited by Philip Schaff, N. York, 1886–88, 8 vols. German translation of select writings of Aug. in the Kempten Bibliothek Der Kirchenväter, 1871–79, 8 vols.

On the same page, line 30. Substitute and add at the close of Lit.:

C. Bindemann: Der heil. Augustin. Berlin, 1844–55–69. 3 vols. Gangauf: Des heil. Aug. Lehre von Gott dem dreieinigen. Augsburg, 1866. Reinkens: Geschichtsphilosophie des heil. Augustin. Schaffhausen, 1866. Emil Feuerlein: Ueber die Stellung Augustin’s in der Kirchen- und Kulturgeschichte. 1869. (In v. Sybel’s “Hist. Zeitschrift” for 1869, vol. xi., 270–313. Ernst: Die Werke und Tugenden der Ungläubigen nach Augustin. Freib., 1872. Böhringer: Aurelius Augustinus, revised ed. Leipz., 1877–78. 2 parts. Aug. Dorner: Augustinus, sein Theol. System und seine religionsphilosophische Auschauung. Berlin, 1873. Ebert: Gesch. der christl. lat. Lit. Leipzig, 1874 sqq., vol. i. 203–243. Edward L. Cutts: St. Augustine. London (Soc. for Prom. Christian Knowledge), 1880. H. Reuter: Augustinische Studien, in Brieger’s “Zeitschrift für Kirchengesch.” for 1880–83 (four articles on Aug.’s doctrine of the church, predestination, the kingdom of God, etc.). Ch. H. Collett: St. Aug., a Sketch o f his Life and Writings as affecting the Controversy with Rome. Lond., 1883. W. Cunningham: S. Austin and his Place in Christian Thought (Hulsean Lectures for 1885), Cambridge, 1886 (283 pp.). James F. Spalding: The Teaching and Influence of Saint Augustine. N. York, 1886 (106 pp.). H. Reuter: Augustinische Studien, Gotha, 1887 (516 pp.; able, learned, and instructive). Ad. Harnack: Augustin’s Confessionen. Giessen, 1888 (31 pp., brief, but suggestive). F. W. Farrar, in his “Lives of the Fathers,” Lond. 1889, vol. ii. 298–460.

On the Philosophy of Aug., compare besides the works quoted on same page:

Erdmann: Grundriss der Gesch. der Philos., i. 231 sqq. Ueberweg: History of Philos. Engl. transl. by Morris, vol. i. 333–346. Ferraz: De la psychologie de S. Aug. 2d ed. Paris, 1869. Schütz: Augustinum non esse ontologum. Monast., 1867. G. Loesche: De Augustino Plotinizanto in doctrina de Deo disserenda. Jenae, 1880. (68 pages.)

\xchapter{Addenda to the Fifth Edition. 1893}

ADDENDA TO THE FIFTH EDITION. 1893

Page 2. Add to Literature:

The historical works of Eusebius, Socrates, Sozomenus, Theodoret, Rufinus, Jerome, and Gennadius, are translated with Introductions and Notes by various American and English scholars in the second series of the “Select Library of Nicene and Post-Nicene Fathers of the Christian Church,” edited by Schaff and Wace, New York and Oxford, vols. i.-iv., 1890–92.

Page 40.

Julian The Emperor, containing Gregory Nazianzen’s Two Invectives, and Libanius’ Monody, with Julian’s extant Theosophical Works. Translated by C. W. King, M. A. London, 1888. With notes and archaeological illustrations (pp. 288).

\xchapter{Addenda Et Corrigenda}

ADDENDA ET CORRIGENDA.

[In the additions to the literature I have followed the method of italicizing book-titles and words in foreign languages, as in the revised edition of vols. i. and ii. The same method will be carried out in all subsequent volumes.]

Page 11. Add To Literature On Constantine The Great:

Th. Zahn: Constantin der Grosse und die Kirche. Hannover, 1876. Demetriades: Die christl. Regierung und Orthodoxie Kaiser Constantin’s des Gr. Muenchen, 1878. Th. Brieger: Constantin der Gr. als Religionspolitiker. Gotha, 1880. E. L. Cutts: Constantine the Great. Lond. And N. Y., 1881. W. Gass: Konstantin der Gr. und seine Soehne, in Herzog,2 vii. (1881), 199–207. John Wordsworth: Const. the Gr. and his Sons, in Smith and Wace, i. 623–654. Edm. Stapfer: in Lichtenberger, iii. 388–393.—Comp. also vol. ii. p. 64–74, especially on the Edicts Of Toleration (only two, not three, as formerly assumed). Victor Schultze Geschichte des Untergangs des Griechisch-roemischen Heidenthums. Jena, 1887, Vol. i. 28–68.

Page 40. Add To Lit, on the heathen sources:

Juliani imperatoris Librorum contra Christianos quae supersunt. Collegit, recensuit, prolegomenis instruxit Car. Joa. Neumann. Insunt Cyrilli Alexandrini fragmenta syriaca ab Eberh. Nestle edita. Lips., 1880. Kaiser Julian’s Buecher gegen die Christen. Nach ihrer Wiederherstellung uebersetzt von Karl Joh. Neumann. Leipzig, 1880. 53 pages. This is Fasc. iii. of Scriptorum Graecorum qui Christianam impugnaverunt religionem quae supersunt, ed. by Neumann.

Page 40, bottom of the page. Add to works on Julian the Apostate:

Alb. De Broglie (R.C.), in the third and fourth vols. of his L’église et l’empire romain au quatrième siécle. Par., 4th ed., 1868. (Very full.) J. F. A. Muecke: Flavius Claudius Julianus. Nach den Quellen. Gotha, 1867 and 1869. 2 vols. (Full, painstaking, prolix, too much dependent on Ammianus, and partial to Julian.) Kellerbaum: Skizze der Vorgeschichte Julians, 1877. F. Rode: Gesch. der Reaction des Kaiser Julianus gegen die christl. Kirche. Jens, 1877. (Careful, partly against Teuffel and Muecke.) H. Adrien Naville: Julien l’apostate et sa philosophie du polythéisme. . Paris and Neuchatel, 1877. Comp. his art. in Lichtenbergers “Encyclop.,” vii. 519–525. Torquati: Studii storico-critici sulla vita … di Giuliano l’Apostata. Rom., 1878. G. H. Rendall: The Emperor Julian: Paganism and Christianity. Lond., 1879. J. G. E. Hoffmann: Jul. der Abtruennige, Syrische Erzaehlungen. Leiden, 1880. (Old romances reflecting the feelings of the Eastern Christians.) Comp. also art. on Jul. in the “Encycl. Brit.,” 9th ed., vol. xiii. 768–770 (by Kirkup); in Herzog2, vii. 285–296 (by Harnack); in Smith and Wace, iii. 484–524 (by Prebendary John Wordsworth, very full and fair).

Page 60. Add to literature:

Tillemont: Hist. des empereurs, tom. v. A. De Broglie, l.c. Victor Schultze: Gesch. d. Untergangs des gr. roem. Heidenthums, i. 209–400.

Page 81, last line, after Muenter, 1826, add:

; by C. Bursian, Lips., 1856; C. Halm, Vienna, 1867).

Page 93. Add as footnote 3:

3 Jerome, who was a shrewd observer of men and things, and witnessed the first effects of the union of church and state, says: “Ecclesia postquam ad Christianos principes venit, potentia quidem et divitiis major, sed virtutibus minor facta.’

Page 148. Add at the bottom of the page:

H. Weingarten: Der Ursprung des Moenchthums im nachconstantinischen Zeitalter. Gotha, 1877. See also his art. in Herzog2, x. 758 sqq. Ad. Harnack: Das Moenchthum, seine Ideale und seine Geschichte. Giessen, 1882.—Comp. vol. ii. ch. ix. p. 387 sqq.

Page 226. Add to footnote:

Ad. Franz, Marcus Aur. Cassiodorus Senator. Breslau, 1872.

Page 242, i 50, add:

See Lit. on clerical celibacy in vol. i. p. 403 sq., especially Theiner, Lea, and von Schulte.

Page 314. Add to Lit. on Leo the Great:

Friedrich (old Cath.): Zur aeltesten Geschichte des Primates in der Kirche. Bonn, 1879. Jos. Langen (old Cath.): Geschichte der roem. Kirche bis zum Pontificate Leo’s I. Bonn, 1881. Karl Mueller in Herzog2, viii. (1881), 551–563. C. Gore, in Smith and Wace, iii. (1882), 652–673. By the same: Leo the Great (Lond. Soc. for Promoting Christ Knowledge, 175 pages). On the literary merits of Leo, see Ebert: Geschichte der christl. Lat. Lit., vol. i. 447–449.

Page 329.

Add to § 64 the following:
